\documentclass[11pt,oneside]{article}
\usepackage[margin=1in]{geometry}
\usepackage{amsmath,amssymb,amsthm,mathtools,booktabs}
\usepackage[colorlinks=true,linkcolor=blue,urlcolor=blue]{hyperref}
\newtheorem{theorem}{Theorem}[section]
\newtheorem{proposition}[theorem]{Proposition}
\newtheorem{record}[theorem]{Record}
\newtheorem{firewall}[theorem]{Claim firewall}
\DeclareMathOperator{\SU}{SU}
\DeclareMathOperator{\Tr}{Tr}
\title{Renewal Yang--Mills Mass-Gap Research Notes\\
Post-Fifteenth-Archive Compact Lineage, Version V86 (branch f16.1)}
\author{Continuation with a scope-checked result ledger}
\date{Research continuation}
\begin{document}
\maketitle
\tableofcontents

\section{Context, restart, and the unchanged objective}
\label{f16v1:context}

The objective is a rigorous nontrivial interacting four-dimensional
continuum Yang--Mills theory with a positive physical Hamiltonian mass
gap. That objective remains unproved. We currently work with an exact
compact \(\SU(2)\) Wilson transfer, its gauge constraints, and particular
weak-field limits. Auxiliary Gaussian separation, a lattice transfer
ratio, and a calibrated physical continuum mass are different claims.

The hundred-version fifteenth lineage is preserved in
\begin{center}\small
\path{Renewal_Yang_Mills_Mass_Gap_Research_Notes_V100_f15.tex}.
\end{center}
This is a compact restart, not a replacement of that archive.
References such as f15 V99 specify sections of that file.
Prior hypotheses and corrections remain binding. The summary below
is a result map, not a fresh certification of every historical claim.

The required V100 checkpoint reread SM continuum V72, NS stability V26,
shell-mixing V16, parallel YM V5, SM source-selection V31,
regenerative stationarity V34 and exact-action Einstein bridge V24.
Their current versions and hashes were unchanged. No companion
provides the missing volume-uniform native YM spectral comparison.
The useful cross-program lesson is to control densities or local
quantities rather than demand a volume-independent tail for an
extensive count. The detailed checkpoint is in f15 V100.
The next companion review is this lineage's V5; its next normal
archive/restart is V100.

The recent upstream change is to retain the normalized magnetic
image of the cube carrier rather than freeze its internal variables
or project its image straight back to the old ground band.
The present note continues by computing the exact overlap of those
conditional vectors across a native bridge kinetic step.
Its weak-field response is not an invariant-subspace theorem:
even the resulting Gaussian compressed family fails the semigroup
law at finite time. We keep that discrepancy explicit.

\section{Major interfaces retained from f15}
\label{f16v1:ledger}

\begin{center}\small
\begin{tabular}{@{}p{0.15\textwidth}p{0.79\textwidth}@{}}
\toprule
Range & Established objects and scope limits\\
\midrule
V1--V11 &
Full compact source/entropy estimates, literal blocking, growing-depth
curvature comparisons and scale-correct pinned-fibre interfaces.
They do not give physical continuum coercivity.\\[1mm]
V12--V37 &
Local chart and source-response comparisons with normalization,
exceptional-set and entropy costs retained; obstructions to dropping
those costs. No unqualified extensive-source mixing theorem.\\[1mm]
V38--V56 &
Native compact conditional landscapes, coherent boundary minima,
reflected source laws and paid pressure comparisons.
Rare or low-action regions are not automatically harmless for spectra.\\[1mm]
V57--V65 &
Positive Wilson transfer and vacuum response estimates in their
selected regimes; soft microscopic spectral limits.
A collapsing one-step lattice gap is not a disproof of a positive
gap after independently calibrated physical time scaling.\\[1mm]
V66--V81 &
Exact decimation/conditional laws, local Dirichlet and bad-region
control, full Gaussian invariant sectors, and vacuum-independent
transfer/compound-operator certificates.
A local gap or a finite matrix check is not the full native gap.\\[1mm]
V82--V91 &
Exact Gauss-boundary charge decomposition, magnetic charge walks,
charged-sector domination, compact endpoint channels and paid
spectral tails. Literature and undecidability audits import
methods, not an unverified mass-gap conclusion.\\[1mm]
V92--V97 &
Complete native coupled finite-region kernels and their growing
boundary carriers; full internal oscillator spectra and
finite-volume complement approximations.\\[1mm]
V98--V100 &
Actual spatial assembly, a sharp auxiliary coarse electric edge,
exact magnetic Gram/dressing, a fixed-volume Gaussian response,
and extensive old-oscillator excitation budgets.
Volume-uniform native dressed dynamics remain open.\\
\bottomrule
\end{tabular}
\end{center}

\subsection{The exact spatial assembly and its unresolved dynamics}

Let the spatial torus have side \(2m\), \(m\ge4\), with
\(B=m^3\) disjoint vertex cubes. A link is internal when its base
coordinate in its own direction is even; otherwise it is a bridge.
There are \(12B\) of each. Of the \(24B\) plaquettes, \(6B\)
are internal, \(12B\) cross one block interface and \(6B\) cross
in both directions. The coarse graph has four parallel bridges
between each pair of positive-direction neighbouring blocks.

A rooted spanning tree in each cube gives exact Haar coordinates
\(g\in G^{7B}\), \(X\in G^{5B}\), \(Y\in G^{12B}\), \(G=\SU(2)\).
After nonroot Gauss reduction the whole physical Hilbert space is
\begin{equation}
 \mathcal H_{\rm phys}
       =L^2(G^{5B}\times G^{12B})^{G^B}.
 \label{f16v1:physical}
\end{equation}
The root group conjugates its internal chords and acts at all its
bridge endpoints. Mixed plaquette words are exactly the two-bridge,
two-internal-edge words or the four-bridge words from the original
lattice. No plaquette is discarded.

For \(s(U)=1-\frac12\Tr U\), define the normalized central convolution
kernel and bridge operator
\begin{equation}
 k_\gamma(U)=\frac{e^{-\gamma s(U)}}{z_\gamma},\quad
 z_\gamma=\int_Ge^{-\gamma s(U)}\,dH(U),\quad
 \mathcal C_\gamma=\bigotimes_{e\ {\rm bridge}} C_\gamma^{(e)}.
 \label{f16v1:kinetic}
\end{equation}
These are the Wilson kinetic kernels, not exactly heat kernels.
The full native transfer has the exact factorization
\begin{equation}
 T_{\rm full,\gamma}
 =M_\times\left[
       \left(\bigotimes_b T_{\Box,\gamma}^{(b)}\right)
                       \otimes\mathcal C_\gamma\right]M_\times,
 \quad M_\times=e^{-\gamma S_\times/2},
 \label{f16v1:full}
\end{equation}
where \(S_\times\) sums all \(18B\) mixed plaquette actions.
The regional transfer factors and the bridge factor commute before
the magnetic sandwich, but the magnetic factor cannot be omitted.

The ground-carrier compression of the cut transfer in f15 V98
has coarse electric operator
\[
 \mathcal D=\mathcal E_{\rm bridge}
       +\sum_b\sum_{v,w}(L_\Box^+)_{vw}J_{bv}\cdot J_{bw},
 \qquad \mathcal E_{\rm bridge}\le\mathcal D
                         \le4\mathcal E_{\rm bridge}.
\]
On the coarse physical space its first nonzero eigenvalue is \(19/2\).
This is an auxiliary operator. Its proved native growing-carrier
error \(O_R(B/\sqrt\gamma)\) is larger than the \(1/\gamma\)
scale of that auxiliary heat edge. The magnetic image contributes
an additional order-one effect. No full-transfer gap follows.

\subsection{The exact normalized carrier and the leading response}

Here are sufficient definitions to continue without reprinting the
finite-region proofs. In one cube let \(\mathsf B\) be the
root-deleted edge incidence matrix, with row \(z_t-z_s\);
let \(\mathsf J\) inject its five chord coordinates; and let
\(\mathsf Z\mathsf B=0,\ \mathsf Z\mathsf J=I\) be the cycle matrix.
Use the rooted axial tree of f15 V97. Put
\[
 \mathsf F=-(\mathsf B^{\mathsf T}\mathsf B)^{-1}
                     \mathsf B^{\mathsf T}\mathsf J,\quad
 \mathsf M=\mathsf Z\mathsf Z^{\mathsf T},\quad
 U=\mathsf J+\mathsf B\mathsf F=\mathsf Z^{\mathsf T}\mathsf M^{-1}.
\]
Let \(\mathsf H\) be the Gram of the six linearized cube faces and
\(T=\mathsf M^{1/2}\mathsf H\mathsf M^{1/2}\).
The spectrum of \(T\) is \(4,4,4,6,6\).
The covariance of one colour component of the squared internal
Gaussian ground vector is
\begin{equation}
 C=\mathsf M^{1/2}[2\sqrt{T+T^2/4}]^{-1}\mathsf M^{1/2},
 \qquad C_B=\bigoplus_b C.
 \label{f16v1:covariance}
\end{equation}
For \(X_i=\operatorname{Exp}(x_i/\sqrt\gamma)\) on the complete
chord charts, let \(\Phi_\gamma(X)\) be the positive normalized
product Gaussian vector with its Haar square-root Jacobian divided
out, and with its chart restriction normalized exactly. Its
squared \(x\)-density is the corresponding truncated Gaussian.

Set \(h_{bu}=\operatorname{Exp}((\mathsf F x_b)_u/\sqrt\gamma)\)
at nonroot ports and \(h_{bo}=I\). The moving bridge coordinates are
\[
                    (\Theta_XY)_e=h_s^{-1}Y_eh_t.
\]
At fixed \(X\) this is Haar preserving. Define
\begin{align}
 \widehat S_\gamma f(X,Y)
       &=\Phi_\gamma(X)f(\Theta_XY),\notag\\
 F_{\gamma,t}(Z)
       &=\int|\Phi_\gamma(X)|^2
              M_\times(X,\Theta_X^{-1}Z)^t\,dH^{5B}(X).
 \label{f16v1:gram}
\end{align}
The identities
\[
 \widehat S_\gamma^*M_\times^t\widehat S_\gamma=M_{F_{\gamma,t}},
 \qquad
 \mathcal I_\gamma=M_\times\widehat S_\gamma
                                    M_{F_{\gamma,2}^{-1/2}},
 \qquad \mathcal I_\gamma^*\mathcal I_\gamma=I
\]
are exact on the coarse physical carrier. At each finite regulator
\(F_{\gamma,2}>0\) is bounded below. It is \(F_{\gamma,2}\), not
\(F_{\gamma,1}^2\), that normalizes the magnetic image.

Let \(C_{\rm int}\) and \(D\) be the signed mixed-plaquette incidence
matrices on internal and bridge edges. In the moving frame the
internal linear embedding is \(\bigoplus_bU\), so put
\[
 A=C_{\rm int}\bigoplus_bU,\qquad
 K=AC_BA^{\mathsf T},\quad
 L=C_B^{1/2}A^{\mathsf T}AC_B^{1/2},\quad \rho=1/\sqrt2 .
\]
F15 V99 proves
\[
 0\le K,\quad\|K\|\le\rho<1,\quad
 \Tr K=B\left(\frac3{2\sqrt2}+\frac2{\sqrt{15}}\right)>B.
\]
With \(Z_e=\operatorname{Exp}(y_e/\sqrt\gamma)\),
\begin{equation}
 F_{\gamma,t}(Z)\longrightarrow
 F_t(y)=\det(I+(t/2)K)^{-3/2}
     e^{-(t/4)y^{\mathsf T}
       [D^{\mathsf T}(I+(t/2)K)^{-1}D\otimes I_3]y},
 \label{f16v1:limit}
\end{equation}
uniformly on bounded \(y\)-windows at fixed \(B\).
The proved absolute error is \(O_t(B(1+R^3)/\sqrt\gamma)
+O_t(Be^{-c\gamma})\). It is not a relative bound uniform in volume.

The normalized conditional limit of \(\mathcal I_\gamma\) is the
positive square root of a Gaussian with covariance and mean
\begin{equation}
 C_{\rm p}=(C_B^{-1}+A^{\mathsf T}A)^{-1},\qquad
 m_y=-(C_{\rm p}A^{\mathsf T}D\otimes I_3)y .
 \label{f16v1:posterior}
\end{equation}
The inverse has a block-distance convergent Neumann expansion:
for block sets \(E,F\) at distance \(d\),
\[
 \|P_EC_{\rm p}P_F\|\le\|C\|\rho^d/(1-\rho).
\]
The leading log determinant also has a uniformly summable
per-block closed-walk expansion. These are Gaussian response
statements, not a native non-Gaussian polymer theorem.

F15 V100 counts all internal oscillator levels in this dressed
vector. It proves exponentially small capture by every
\(o(B)\) total-degree cutoff at zero bridge field and an
exponentially accurate cutoff of order
\(B+\|Dy\|^2+\log\delta^{-1}\). Its compact projection handoff
takes \(\gamma\to\infty\) at each fixed \(B\); it does not control
the norm of the full excited-state transfer complement.

\section{New V1 result: exact kinetic compression by conditional affinity}
\label{f16v1:affinity}

For the derivation it is convenient to work first before root Gauss
restriction, on \(L^2(G^{12B})\) and
\(L^2(G^{5B}\times G^{12B})\) in \((X,Z)\) coordinates.
All maps below intertwine the root action, so their exact
operator identities restrict to \eqref{f16v1:physical}.
Define normalized conditional vectors
\begin{equation}
 a_{\gamma,Z}(X)
   =\frac{\Phi_\gamma(X)M_\times(X,\Theta_X^{-1}Z)}
                    {\sqrt{F_{\gamma,2}(Z)}},
 \qquad \int|a_{\gamma,Z}(X)|^2\,dH^{5B}(X)=1,
 \label{f16v1:conditional}
\end{equation}
and their affinity
\[
 \mathcal A_\gamma(Z,Z')
       =\int a_{\gamma,Z}(X)a_{\gamma,Z'}(X)\,dH^{5B}(X).
\]

\begin{theorem}[Exact native bridge kernel on the dressed carrier]
\label{f16v1:exact-kernel}
The compression of the bridge kinetic factor alone is
\begin{equation}
 \bigl(\mathcal I_\gamma^*\mathcal C_\gamma
                      \mathcal I_\gamma\bigr)(Z,Z')
   =\left[\prod_{e\ {\rm bridge}}
             k_\gamma(Z_eZ_e'^{-1})\right]\mathcal A_\gamma(Z,Z').
 \label{eq:f16v1:exact-kernel}
\end{equation}
The affinity is symmetric, positive semidefinite as a Gram kernel,
continuous, and satisfies \(0<\mathcal A_\gamma\le1\) and
\(\mathcal A_\gamma(Z,Z)=1\).
The compressed operator is a positive contraction.
The root gauge invariance is simultaneous in \(Z,Z'\), not
a claim that the affinity is separately invariant in each argument.
\end{theorem}

\begin{proof}
The bridge operator leaves \(X\) fixed. Under
\(Y_e=h_s Z_eh_t^{-1}\) and \(Y_e'=h_s Z_e'h_t^{-1}\),
\[
                  Y_eY_e'^{-1}
                       =h_s Z_eZ_e'^{-1}h_s^{-1}.
\]
Centrality of \(k_\gamma\) and Haar invariance remove both frame
factors. In these coordinates \(\mathcal I_\gamma f=a_{\gamma,Z}f(Z)\).
Integrating the internal coordinate in its compression gives
\eqref{eq:f16v1:exact-kernel}. The Gram assertion and the upper
bound are respectively the squared-norm identity and Cauchy--Schwarz.
Positivity follows also directly from compression of the positive
Wilson convolution. Its norm is at most one.

At finite \(\gamma,B\), the magnetic multiplier is strictly positive
and bounded, \(F_{\gamma,2}\) has a positive lower bound, and
\(\Phi_\gamma\) has norm one. Dominated convergence proves
continuity and strict positivity of the affinity. Equivariance
of the original carrier and kernel proves the root-action
statement and the restriction to invariant functions.
\end{proof}

This identity keeps the whole compact bridge configuration space.
It does not replace \eqref{f16v1:full} by its bridge factor.
In particular the cube kinetic dynamics and their internal
excited-state returns have not been integrated by this theorem.

\subsection{The leading affinity and its exact null directions}

In the next formulas the colour tensor \(I_3\) is suppressed:
when \(R\) acts on a \(36B\)-component bridge vector, it means
\(R\otimes I_3\).
Set
\begin{align}
 R&=D^{\mathsf T}A C_{\rm p}A^{\mathsf T}D
    =D^{\mathsf T}K(I+K)^{-1}D
    =D^{\mathsf T}D-H_{\rm eff},\notag\\
 H_{\rm eff}&=D^{\mathsf T}(I+K)^{-1}D .
 \label{f16v1:R}
\end{align}

\begin{theorem}[Conditional overlap and native local-kernel limit]
\label{f16v1:gaussian-affinity}
For the normalized Gaussian amplitudes \(\psi_y\) with squared
density \eqref{f16v1:posterior},
\begin{equation}
 \langle\psi_y,\psi_z\rangle
             =e^{-(y-z)^{\mathsf T}R(y-z)/8}.
 \label{eq:f16v1:gaussian-affinity}
\end{equation}
At fixed \(B\), \(\mathcal A_\gamma\) in rescaled bridge charts
converges to this expression uniformly on bounded pairs of
bridge windows. After including both square-root Haar chart
densities, the kernel \eqref{eq:f16v1:exact-kernel} converges there to
\begin{equation}
 (2\pi)^{-q/2}
   \exp\left[-\frac{|y-z|^2}{2}
                    -\frac{(y-z)^{\mathsf T}R(y-z)}8\right],
       \qquad q=36B.
 \label{f16v1:native-leading-kernel}
\end{equation}
Thus compression to any fixed bounded coordinate window
converges in Hilbert--Schmidt norm. No global norm or
volume-uniform spectral convergence is asserted.
\end{theorem}

\begin{proof}
Completing the common square of two Gaussian amplitudes with
covariance \(C_{\rm p}\) gives the overlap
\[
       \exp[-(m_y-m_z)^{\mathsf T}C_{\rm p}^{-1}(m_y-m_z)/8].
\]
Substitute \eqref{f16v1:posterior}; Woodbury inversion gives
\eqref{f16v1:R}. F15 V99's normalized-density convergence,
or f15 V100's \(L^2\) amplitude argument, gives the uniform
convergence of the affinities at fixed \(B\).

For completeness, in one bridge chart with Jacobian \(J_\gamma\),
ordinary Taylor expansion and Laplace normalization give
\[
 [J_\gamma(y)J_\gamma(z)]^{1/2}
 k_\gamma(\operatorname{Exp}(y/\sqrt\gamma)
                      \operatorname{Exp}(-z/\sqrt\gamma))
       \longrightarrow(2\pi)^{-3/2}e^{-|y-z|^2/2}
\]
uniformly on compact pairs. Here
\(s(\operatorname{Exp}v)=1-\cos|v|\);
\(\gamma s\) tends to \(|y-z|^2/2\).
The normalization follows by the same rescaling in
\(z_\gamma\); the inequality \(1-\cos r\ge2r^2/\pi^2\)
for \(0\le r\le\pi\) dominates the full chart integral.
The Haar Jacobian cancels to the displayed Euclidean prefactor.
Multiply over the fixed \(12B\) bridges and use
\eqref{eq:f16v1:gaussian-affinity}. Uniform convergence on
a bounded window implies square-integral kernel convergence.
This is a kinematic chart assertion before root averaging;
it does not identify a physical spectral subspace with that window.
\end{proof}

One has
\begin{equation}
 0\le R\le\frac{\rho}{1+\rho}D^{\mathsf T}D.
 \label{f16v1:Rbound}
\end{equation}
The covariance locality from \eqref{f16v1:posterior} and the
finite-range incidence factors give exponentially decaying
block entries of \(R\), uniformly in volume.

There is a stronger exact statement than simply
\(\ker D\subseteq\ker R\). Let \(\mathcal E\) consist of bridge
vectors constant on each set of four parallel bridges.
Then
\begin{equation}
 A^{\mathsf T}D y=0,\qquad R y=0,\qquad
 H_{\rm eff}y=D^{\mathsf T}Dy,\qquad \psi_y=\psi_0
                    \quad (y\in\mathcal E).
 \label{f16v1:common-bridge}
\end{equation}
Indeed a decorated digon takes the difference of two such
parallel values, so its curvature vanishes. A quadrilateral
has no internal edge, so the corresponding row of \(A\) is
zero. This proves the first identity; the others follow.
More generally \(\psi_{y+v}=\psi_y\) for \(v\in\mathcal E\).
The scalar magnetic normalizer still depends on the coarse
quadrilateral curvature; it has not disappeared.

In particular the transverse common-bridge cosine used in
f15 V99 has exactly
\[
       \frac{\langle y,H_{\rm eff}y\rangle}{\|y\|^2}
                           =2\sin^2(\pi/m),
\]
not merely the upper bound established there.
The leading fast conditional vector cannot generate a restoring
term in these coarse common-bridge directions.
This concerns a static linearized Hessian, not a family of
gapless native physical Hamiltonian states.

\section{Finite kinetic time is not fixed by the static normalization}
\label{f16v1:time}

The Gaussian limit allows a complete check of this issue.
Let \(G\) be any positive definite \(q\)-by-\(q\) constant matrix,
and let \(H_t^G\) act on the bridge variable, leaving the internal
variable fixed, with kernel
\[
 h_t^G(y-z)
 =(4\pi t)^{-q/2}\det(G)^{-1/2}
        e^{-(y-z)^{\mathsf T}G^{-1}(y-z)/(4t)}.
\]
The isometry \(\mathcal J f(u,y)=\psi_y(u)f(y)\) compresses it
to \(\mathcal B_t=\mathcal J^*H_t^G\mathcal J\).
The actual one-step local limit \eqref{f16v1:native-leading-kernel}
is the case \(G=I,t=1/2\).

\begin{theorem}[Exact finite-time diffusion change and its composition defect]
\label{f16v1:heat}
For every \(t>0\),
\begin{equation}
 \mathcal B_t=c_t H_t^{G_t},\qquad
 G_t=(G^{-1}+tR/2)^{-1},\qquad
 c_t=\det(I+(t/2)G^{1/2}RG^{1/2})^{-1/2}.
 \label{f16v1:heat-factor}
\end{equation}
Its operator norm on \(L^2(\mathbb R^q)\) is \(c_t\).
If \(R\ne0\), the norm-normalized family
\(c_t^{-1}\mathcal B_t\) is not a semigroup.
\end{theorem}

\begin{proof}
Multiply \(h_t^G(y-z)\) by the affinity
\eqref{eq:f16v1:gaussian-affinity}.
The exponent becomes
\(-(y-z)^{\mathsf T}(G^{-1}+tR/2)(y-z)/(4t)\).
Comparing Gaussian prefactors gives \(c_t\).
The Fourier multiplier is
\(c_t e^{-t p^{\mathsf T}G_t p}\), whose essential supremum is
\(c_t\).

For \(R\ne0\), positivity and inversion show
\(G_t-G_{2t}\) is nonzero positive semidefinite.
Thus for some \(p\),
\[
 e^{-2t p^{\mathsf T}G_t p}
                   \ne e^{-2t p^{\mathsf T}G_{2t}p}.
\]
These are the multipliers of
\((c_t^{-1}\mathcal B_t)^2\) and
\(c_{2t}^{-1}\mathcal B_{2t}\). Continuity in \(p\) makes this
a genuine operator inequality, proving the failure of the
semigroup law.
\end{proof}

The nonzero case actually occurs in the assembly.
For example take a bridge perturbation supported only on
the positive first-direction bridge from fine vertex \((1,0,0)\).
In its incident block, \(A^{\mathsf T}D\) applied to this vector
is nonzero. To see this without a numerical rank test, its
internal-edge source is, up to overall sign, the sum of the
second- and third-direction edges starting at port \((1,0,0)\).
The first/second-direction cube face at third coordinate zero
contains just the former edge, with nonzero signed coefficient.
Thus the source has nonzero circulation. Since
\(\ker U^{\mathsf T}\) is the cut space, its image under
\(U^{\mathsf T}\) cannot vanish. Positivity of \(C_{\rm p}\)
then implies \(R\ne0\).
The exact incidence checker supplies this literal relative-port
example on sides eight and twelve.

The failure is also transparent before taking limits:
\(\mathcal J^*H_s^G H_t^G\mathcal J\) differs from
\(\mathcal J^*H_s^G\mathcal J\mathcal J^*H_t^G\mathcal J\)
by the omitted intermediate internal complement.
A static isometry alone cannot remove that return.

\subsection{The infinitesimal check and its limited physical meaning}

There is a useful independent derivative check.
For bridge indices \(i,j\), real positivity and normalization give
\[
 \langle\psi_y,\partial_i\psi_y\rangle=0,\qquad
 \langle\partial_i\psi_y,\partial_j\psi_y\rangle=R_{ij}/4.
\]
Differentiate the Gaussian mean in
\eqref{f16v1:posterior} to verify the second identity.
For Schwartz \(f\), Fubini's theorem and this zero cross term yield
\begin{equation}
 \int \bigl|\nabla_y(\psi_y f)\bigr|_G^2\,du\,dy
  =\int|\nabla f|_G^2\,dy
                     +\frac14\Tr(GR)\int|f|^2\,dy .
 \label{f16v1:metric}
\end{equation}
Thus the compressed differential quadratic form has a constant
scalar addition, not a bridge-dependent mass potential.
It agrees with the derivative of \eqref{f16v1:heat-factor}:
\[
       \left.\frac d{dt}\log c_t\right|_{t=0}
                           =-\frac14\Tr(GR),\qquad G_0=G.
\]
Exponentiating that compressed differential form is nevertheless
not equal to the finite-time compression in
\eqref{f16v1:heat-factor}. The latter has \(G_t\), not \(G\).
For \(G=I\), the scalar addition is a vacuum normalization of
order at most \(B\), using \eqref{f16v1:Rbound} and the bounded
plaquette incidence. None of these chart statements fixes the
physical lattice time unit or the full interacting transfer gap.

\section{What must be controlled next}
\label{f16v1:next}

The direction is now more precise. The exact native affinity
in \eqref{eq:f16v1:exact-kernel} retains every bridge configuration
and its correct magnetic normalization. Its leading response
has a computable diffusion correction on relative-port directions,
while the fast conditional state is blind to all common-bridge
directions at this order. Repeated ground projection, even after
normalization, discards a definite intermediate return.

The next required calculation is the joint dressed compression
with the cube transfer factors in \eqref{f16v1:full}, retaining
their excited internal bands and the noncommuting moving frame
at both time ends. It must bound the discarded return relative
to the same full native top eigenvalue, uniformly in spatial
volume, or retain it through a controlled memory/Schur term.
A new scalar normalization or another fixed-volume Gaussian
determinant does not supply that estimate.

Beyond it remain nonlinear compact collective-mode coercivity,
bridge-field tail control, an independent physical-time calibration,
and nontrivial interacting continuum reconstruction. These are
requirements of the unchanged objective, not assumptions being
declared solved by the restart.

\section{Verification and preservation}
\label{f16v1:verification}

The diagnostic
\path{scripts/verify_ym_f16_v1_dressed_affinity.py}
checks independent noncommuting matrix completions,
the affinity and derivative metric, finite-time heat precisions,
the composition defect, a direct finite conditional-kernel
compression, and the complete mixed-plaquette incidence on
two actual spatial tori. It checks the common-bridge null
response and an explicit nonzero relative-port response.
It also reruns the inherited f15 V100 diagnostic chain.
All 10,918 new exact fixtures passed, together with the inherited
f15 V100 chain. These finite fixtures supplement proofs, not replace them.
The checker has 5,600 bytes and SHA--256
\begin{center}\scriptsize\ttfamily
CCBCBC50A687E88B74465051A0C7E9B283891911ADA928D910FA5C896B3EFC7B.
\end{center}

F15 V100 is preserved as an archive, with V99's body recoverable
byte for byte. Its source has 2,310,095 bytes and SHA--256
\begin{center}\scriptsize\ttfamily
D9349B9F87D70400E7A140A3BDC38D1BC60F8E3A12F830D8C37693CD6D3C30C5.
\end{center}
This V1 is the sole active main YM source.
Only \texttt{.tex} files are stored in \texttt{latex\_notes};
build products and visual checks are outside it.

\begin{firewall}
The new exact compact result is the conditional-affinity formula
for the dressed bridge kinetic factor. The subsequent Gaussian
overlap, finite-time diffusion change, null directions and
composition defect have proved fixed-volume local limits.
They are not a volume-uniform approximation of the full native
transfer or a positive physical continuum mass-gap theorem.
The full Yang--Mills goal remains unachieved.
\end{firewall}

\section{V2: joint cube dynamics in the dressed transfer}
\label{f16v2:context}

V1 computed the dressed bridge factor while holding the internal
variables fixed. We now include the cube motion at both time ends.
There are two separate corrections: nonroot gauge matching changes
the leading bridge covariance, and compression of the \emph{full}
magnetic sandwich uses the magnetic multiplier twice at each end.
Neither correction is contained in the static conditional affinity.

We first give an exact full-compact kernel. We then integrate its
entire leading internal Gaussian propagator, not just its ground
projection. The resulting two-end Schur matrices have
volume-independent locality and quadratic-form comparisons.
The native limit remains fixed-volume and local in the bridge
coordinates. This does not yet control the discarded native
transfer subspace or the interacting infrared modes.

\subsection{An exact two-ended compact kernel}

Write \(f_v(X)\) for V1's moving port frame, with \(f_o=I\) at
every block root, and put
\[
 \bar X_e=f_s(X)^{-1}X_ef_t(X),\qquad
 \bar X'_e=f_s(X')^{-1}X'_ef_t(X').
\]
This includes tree edges, whose original representatives are \(I\).
As in f15 V99, this is a pointwise gauge rewriting, not a change
of internal integration variables. Denote the exact mixed
multiplier in moving coordinates by
\[
 m_\gamma(X,Z)=M_\times(X,\Theta_X^{-1}Z),
\]
and the internal magnetic multiplier by
\(b_\gamma(X)=\exp[-\gamma S_{\rm int}(X)/2]\).
There are \(7B\) nonroot gauge variables.

\begin{proposition}[Full compact cut kernel with both moving frames]
\label{f16v2:exact-cut}
Before the remaining root Gauss restriction, the cut transfer has
kernel
\begin{align}
 \mathcal K_\gamma(X,Z;X',Z')
 &=b_\gamma(X)b_\gamma(X')\,\mathcal L_\gamma(X,Z;X',Z'),
 \notag\\
 \mathcal L_\gamma
 &=\int_{G^{7B}}
   \prod_{e\ {\rm internal}}
       k_\gamma(r_s^{-1}\bar X_e r_t\bar X_e'^{-1})
   \prod_{e\ {\rm bridge}}
       k_\gamma(r_s^{-1}Z_e r_t Z_e'^{-1})\,dH^{7B}(r).
 \label{f16v2:cut}
\end{align}
Here \(r_o=I\) at block roots. The exact kernel of the full
transfer is \(m_\gamma(X,Z)\mathcal K_\gamma
m_\gamma(X',Z')\). All formulas restrict to the root-invariant
physical space.
\end{proposition}

\begin{proof}
In the rooted tree variables at the two times write \(g'=gh\).
Nonroot Gauss reduction integrates \(h\) and leaves the kinetic
arguments \(h_s^{-1}X_eh_tX_e'^{-1}\) on internal edges and
\(h_s^{-1}Y_eh_tY_e'^{-1}\) on bridges.
This is the same exact Fourier-free Haar reduction used in
f15 V97; no boundary representation has been discarded.

Set \(h_v=f_v(X)r_vf_v(X')^{-1}\).
For fixed endpoints this is a product of Haar translations.
The internal argument becomes, up to conjugation by \(f_s(X')\),
\(r_s^{-1}\bar X_e r_t\bar X_e'^{-1}\).
On a bridge, substitute \(Y_e=f_s(X)Z_ef_t(X)^{-1}\) and its
primed counterpart. The same conjugation gives
\(r_s^{-1}Z_er_tZ_e'^{-1}\). Centrality of \(k_\gamma\) removes
the conjugations. The internal and mixed magnetic traces are
unchanged by the pointwise gauge rewriting.
This proves the kernel and the full sandwich.
\end{proof}

Let \(\mathcal I_\gamma\) be V1's exact normalized dressed
isometry. Distinguish
\[
 \mathcal B_{1,\gamma}=\mathcal I_\gamma^*
                  T_{\rm cut,\gamma}\mathcal I_\gamma,\qquad
 \mathcal B_{2,\gamma}=\mathcal I_\gamma^*
                  T_{\rm full,\gamma}\mathcal I_\gamma.
\]
Their exact kernels, for \(\ell=1,2\), are
\begin{equation}
 \begin{aligned}
 \mathcal B_{\ell,\gamma}(Z,Z')
 &=\frac{1}{\sqrt{F_{\gamma,2}(Z)F_{\gamma,2}(Z')}}\\
 &\quad\times\int \Phi_\gamma(X)\Phi_\gamma(X')\,
      m_\gamma(X,Z)^\ell m_\gamma(X',Z')^\ell
      \mathcal K_\gamma(X,Z;X',Z')\,dH^{5B}(X)dH^{5B}(X').
 \end{aligned}
 \label{f16v2:exact-compression}
\end{equation}
In particular, \(\ell=2\) is necessary for the full transfer:
\[
 \mathcal I_\gamma^*T_{\rm full,\gamma}\mathcal I_\gamma
 =M_{F_{\gamma,2}^{-1/2}}\widehat S_\gamma^*
      M_\times^2T_{\rm cut,\gamma}M_\times^2\widehat S_\gamma
                                      M_{F_{\gamma,2}^{-1/2}}.
\]
Using \(\ell=1\) instead would compress a different operator.

\subsection{The kinetic covariance produced by the moving nonroot vertices}

Use block-diagonal copies \(\mathsf B_B,\Sigma_B,\mathsf M_B,
\mathsf H_B,U_B\) of the single-cube matrices in V1 and f15 V97,
where \(\Sigma=(\mathsf B^{\mathsf T}\mathsf B)^{-1}\).
Let \(E\) be the \(12B\)-by-\(7B\) bridge endpoint incidence:
its row is \(z_t-z_s\), with each root coordinate set to zero.
Define
\begin{equation}
 G=I+E\Sigma_BE^{\mathsf T},\qquad
 L_h=\mathsf B_B^{\mathsf T}\mathsf B_B+E^{\mathsf T}E.
 \label{f16v2:G}
\end{equation}
The letter \(G\) in these matrix formulas denotes a covariance;
the group remains \(\SU(2)\).

\begin{theorem}[Exact leading kinetic square and a uniform covariance bound]
\label{f16v2:kinetic-square}
For one colour component, \(\delta x=x-x'\), \(\delta y=y-y'\),
and \(z\in\mathbb R^{7B}\),
\begin{align}
 &\|U_B\delta x+\mathsf B_Bz\|^2+\|\delta y+Ez\|^2
 \notag\\
 &\quad=\delta x^{\mathsf T}\mathsf M_B^{-1}\delta x
       +\delta y^{\mathsf T}G^{-1}\delta y
       +(z+L_h^{-1}E^{\mathsf T}\delta y)^{\mathsf T}
                    L_h(z+L_h^{-1}E^{\mathsf T}\delta y).
 \label{f16v2:completed-kinetic}
\end{align}
Moreover
\begin{equation}
  \det L_h=(\det\mathsf M)^B\det G,\qquad
                         I\le G\le30I.
 \label{f16v2:Gbounds}
\end{equation}
The matrix \(G\) couples only bridges incident to a common block.
\end{theorem}

\begin{proof}
The balanced internal embedding obeys
\(U_B^{\mathsf T}\mathsf B_B=0\) and
\(U_B^{\mathsf T}U_B=\mathsf M_B^{-1}\).
There is therefore no internal--port cross term.
Complete the \(z\)-square. Woodbury inversion gives
\[
 I-E L_h^{-1}E^{\mathsf T}
                    =(I+E\Sigma_BE^{\mathsf T})^{-1}=G^{-1}.
\]
The determinant lemma and the cube equality
\(\det(\mathsf B^{\mathsf T}\mathsf B)=\det\mathsf M\)
give the determinant statement.

Each row of \(E\) has absolute sum at most two. Each nonroot
port is incident to three bridges, so each column has absolute
sum three. Thus \(\|E\|^2\le6\).
The diagonal of the rooted inverse \(\Sigma\) consists of the
effective resistances from the root. For the cube these are
\(7/12\) at the three adjacent vertices, \(3/4\) at the three
face-opposite vertices and \(5/6\) at the opposite vertex.
They are the exact resistance values verified and derived
in f15 V98. Consequently
\[
 \|\Sigma_B\|=\|\Sigma\|\le\Tr\Sigma
                    =3(7/12)+3(3/4)+5/6=29/6.
\]
This proves \(G\le I+6(29/6)I=30I\).
Block diagonality of \(\Sigma_B\) gives the support assertion.
\end{proof}

The bound \(30\) is a convenient pre-root-quotient bound,
not a replacement for f15 V98's sharper physical electric
form comparison. The new point is that the bridge covariance
is \(G\), not \(I\). Fixing all temporal port variables to zero
would miss it.

For a positive \(d\)-by-\(d\) path covariance \(Q\), abbreviate
the three-colour Gaussian density by
\[
 \varphi_Q(v)
  =(2\pi)^{-3d/2}\det(Q)^{-3/2}
             e^{-v^{\mathsf T}(Q^{-1}\otimes I_3)v/2}.
\]
Write \(b_0(x)=e^{-x^{\mathsf T}(\mathsf H_B\otimes I_3)x/4}\).
The complete leading cut kernel is
\begin{equation}
 \mathcal K_0(x,y;x',y')
       =b_0(x)b_0(x')\,
           \varphi_{\mathsf M_B}(x-x')\varphi_G(y-y').
 \label{f16v2:leading-cut}
\end{equation}
The full leading kernel has in addition
\(e^{-|Ax+Dy|^2/4}\) at each end.

\begin{proposition}[Fixed-volume native convergence with all internal ends integrated]
\label{f16v2:native-limit}
At each fixed \(B\), after including all square-root Haar
chart densities, the kernel \eqref{f16v2:cut} converges uniformly
on bounded \(x,x',y,y'\) windows to \eqref{f16v2:leading-cut}.
The kernels \eqref{f16v2:exact-compression} converge uniformly
on bounded \(y,y'\) windows to the result of integrating the
complete Gaussian kernel \eqref{f16v2:leading-cut} with its
displayed endpoint weights and normalizers.
There is no internal excitation cutoff in this assertion.
It also gives Hilbert--Schmidt convergence after restricting
both bridge ends to a fixed bounded coordinate window.
\end{proposition}

\begin{proof}
Put \(r_v=\operatorname{Exp}(z_v/\sqrt\gamma)\).
The first-order kinetic words in \eqref{f16v2:cut} are exactly
\(U_B(x-x')+\mathsf B_Bz\) and \(y-y'+Ez\).
Their actions have leading exponent one half of the sum of
the two squared norms. Gaussian integration and
\eqref{f16v2:completed-kinetic} give the covariance in
\eqref{f16v2:leading-cut}; \eqref{f16v2:Gbounds} fixes its
normalizing determinant. The internal plaquette multiplier
converges to \(b_0\).

Here is the compact-integral justification.
When the endpoint charts tend to identity, the internal tree
kinetic words force every nonroot \(r_v\) to identity: each
tree has a fixed root. On the compact group the sum of their
squared geodesic distances controls the distances of all
\(r_v\), up to an endpoint error bounded by \(C/\gamma\)
on a fixed endpoint window. This follows successively along
the seven tree edges, by the bi-invariant triangle inequality
and Cauchy--Schwarz along each path. Since
\(1-\cos\theta\ge2\theta^2/\pi^2\) for \(0\le\theta\le\pi\),
the rescaled integrand is dominated by
\(C\exp(-c\sum_v|z_v|^2)\), with constants at fixed \(B\)
and window. The complement of any fixed small group
neighbourhood is exponentially small. Taylor expansion there,
followed by dominated convergence, proves the local kernel limit.

For the integration over all internal \(X,X'\), a separate
uniform bound avoids a hidden chart restriction.
Keep only the \(7B\) tree factors of the kinetic product and
bound the other \(17B\) factors by \(\|k_\gamma\|_\infty\).
Successive Haar integration of leaves gives one for the tree
product. The full integral is therefore at most
\(\|k_\gamma\|_\infty^{17B}\).
The normalization in \eqref{f16v1:kinetic} gives
\(\|k_\gamma\|_\infty\le C\gamma^{3/2}\) for \(\gamma\ge1\).
Each chart Jacobian is at most \(C\gamma^{-3/2}\).
The two half-densities on the \(17B\) reduced coordinates
cancel this power. Thus the cut chart kernel is bounded by
a constant depending on \(B\), uniformly throughout the
internal charts; magnetic multipliers are at most one.

The transformed \(\Phi_\gamma\) is exactly a normalized
truncation of the product Gaussian ground amplitude.
It therefore gives an integrable dominating product
\(C_B\Omega(x)\Omega(x')\).
On a fixed bridge window \(F_{\gamma,2}\) is bounded below
uniformly for sufficiently large \(\gamma\), by V1's
fixed-volume limit. Internal Gaussian tails are now uniform
in \(\gamma\). Apply the local limit on bounded internal
windows and then send those windows to infinity.
This proves uniform convergence of the compressed kernels.
The final Hilbert--Schmidt assertion follows on any bounded
bridge window. All constants may deteriorate with \(B\).
\end{proof}

\subsection{The two-end Schur formula retaining the full internal propagator}

The product cube Gaussian ground amplitude \(\Omega(x)\)
has squared covariance \(C_B\) from \eqref{f16v1:covariance}.
Write \(N=\mathsf M_B^{-1}\) and
\begin{equation}
 Q_0=\tfrac12(C_B^{-1}+\mathsf H_B),\qquad
 V_{\ell,+}=Q_0+\tfrac\ell2 A^{\mathsf T}A,\qquad
 V_{\ell,-}=V_{\ell,+}+2N,\qquad \ell=0,1,2.
 \label{f16v2:V}
\end{equation}
All are strictly positive. Define
\begin{align}
 P_{\ell,\pm}&=D^{\mathsf T}A V_{\ell,\pm}^{-1}A^{\mathsf T}D,
 \notag\\
 B_{\ell,\pm}
       &=\ell D^{\mathsf T}D-H_{\rm eff}
                              -\tfrac{\ell^2}{2}P_{\ell,\pm}.
 \label{f16v2:B}
\end{align}
Colour tensors are suppressed: each path matrix acts with \(I_3\)
on the actual bridge vectors.

Let \(r(q)=(1+q/2+\sqrt{q+q^2/4})^{-1}\), and put
\[
                   \lambda_\Box^0=r(4)^{9/2}r(6)^3.
\]
This is the top value of the entire single-cube Gaussian
transfer, not its physical mass. It is the
\href{https://dlmf.nist.gov/18.18.E28}{Mehler value}
already derived in f15 V97, with nine coordinates of parameter
four and six coordinates of parameter six.
Set
\begin{equation}
 \kappa_\ell
 =\frac{(\lambda_\Box^0)^B}{F_2(0)}
   \left[
    \frac{\det V_{0,+}\det V_{0,-}}
         {\det V_{\ell,+}\det V_{\ell,-}}
   \right]^{3/2},\qquad \ell=1,2.
 \label{f16v2:kappa}
\end{equation}

\begin{theorem}[Complete Gaussian dressed cut and full compressions]
\label{f16v2:schur}
The fixed-volume limits in Proposition~\ref{f16v2:native-limit} are
\begin{align}
 \mathcal B_{\ell,0}(y,y')
 &=\kappa_\ell\,\varphi_G(y-y') \notag\\
 &\quad\cdot
 \exp\left[-\tfrac18(y+y')^{\mathsf T}B_{\ell,+}(y+y')
           -\tfrac18(y-y')^{\mathsf T}B_{\ell,-}(y-y')\right],
       \qquad \ell=1,2.
 \label{f16v2:schur-kernel}
\end{align}
In particular \(\ell=2\) includes the entire internal Gaussian
propagator of the full magnetic transfer between the dressed ends.
No intermediate internal ground projection appears.
\end{theorem}

\begin{proof}
At one end the integrand of \eqref{f16v2:exact-compression}
tends to
\[
 \frac{\Omega(x)
           e^{-\ell|Ax+Dy|^2/4}b_0(x)}{\sqrt{F_2(y)}} .
\]
The two-end internal precision and source are
\[
 W_\ell=
 \begin{pmatrix}V_{\ell,+}+N&-N\\-N&V_{\ell,+}+N\end{pmatrix},
 \qquad
 \eta_\ell=-\frac\ell2
          \begin{pmatrix}A^{\mathsf T}Dy\\A^{\mathsf T}Dy'\end{pmatrix}.
\]
Changing to \((x+x')/\sqrt2,(x-x')/\sqrt2\) gives the two
precisions \(V_{\ell,+},V_{\ell,-}\). Gaussian integration
therefore contributes determinant
\((\det V_{\ell,+}\det V_{\ell,-})^{-3/2}\)
and exponent
\[
 \frac{\ell^2}{16}\left[
   (y+y')^{\mathsf T}P_{\ell,+}(y+y')
     +(y-y')^{\mathsf T}P_{\ell,-}(y-y')\right].
\]
The remaining endpoint exponent is
\[
 -\frac\ell4(y^{\mathsf T}D^{\mathsf T}Dy
                   +y'^{\mathsf T}D^{\mathsf T}Dy')
 +\frac14(y^{\mathsf T}H_{\rm eff}y
                   +y'^{\mathsf T}H_{\rm eff}y').
\]
The second term comes from \(F_2(y)^{-1/2}F_2(y')^{-1/2}\).
Combining these expressions gives exactly \eqref{f16v2:B}.

To fix the scalar without discarding any determinant, set the
endpoint magnetic power to zero in the internal integral.
It is \(\langle\Omega,K_{\Box,0}^{\otimes B}\Omega\rangle
=(\lambda_\Box^0)^B\). Dividing the general determinant by
this one and retaining \(F_2(0)^{-1}\) gives
\eqref{f16v2:kappa}. This proves the formula.
\end{proof}

This is a calculation of one-step compression, not a claim that
its spectrum equals the spectrum of \(T_{\rm full,\gamma}\).
V1's intermediate-complement warning remains in force.

\subsection{Uniform locality of the retained two-time fast response}

The inverses in \eqref{f16v2:V} are uniformly local even for the
full magnetic power \(\ell=2\).
For either sign put
\[
 Q_+=Q_0,\quad Q_-=Q_0+2N,\qquad
 Z_\pm=Q_\pm^{-1/2}A^{\mathsf T}A Q_\pm^{-1/2}.
\]
Since \(Q_\pm\ge C_B^{-1}/2\),
\[
                 0\le Z_\pm,\qquad \|Z_\pm\|\le2\rho .
\]
These matrices have only nearest-block off-diagonal entries.
Let \(c=1+2\rho\), \(\theta=1-c^{-1}<1\), and
\(T_\pm=I-(I+Z_\pm)/c\). Then \(0\le T_\pm\le\theta I\) and
\begin{equation}
 V_{2,\pm}^{-1}
   =c^{-1}Q_\pm^{-1/2}
                     \sum_{j=0}^{\infty}T_\pm^jQ_\pm^{-1/2}.
 \label{f16v2:local-inverse}
\end{equation}
For block sets \(E_0,F_0\) separated by distance \(d\), this gives
\[
 \|P_{E_0}V_{2,\pm}^{-1}P_{F_0}\|
                    \le\|Q_\pm^{-1}\|\theta^d.
\]
Indeed powers of degree less than \(d\) cannot connect the sets,
and \(c^{-1}/(1-\theta)=1\).
No volume or set-cardinality factor is needed.
Consequently \(P_{2,\pm}\) and \(B_{2,\pm}\) have exponentially
decaying block entries, with the fixed range shift of \(A,D\).
This controls the complete leading two-time fast response,
not the non-Gaussian native remainder.

\subsection{The full effective quadratic transfer and its surviving soft modes}

Write \(B_\pm=B_{2,\pm}\).
The lower bound here is relative to the actual coarse curl,
not relative to the identity.

\begin{theorem}[Volume-independent curvature and kinetic comparisons]
\label{f16v2:comparison}
One has
\begin{equation}
 \frac{D^{\mathsf T}D}{(1+\rho)(1+2\rho)}
       \le B_+\le B_-
       \le\frac{1+2\rho}{1+\rho}D^{\mathsf T}D .
 \label{f16v2:Bbounds}
\end{equation}
Define
\[
       C_{\rm eff}
       =\left[G^{-1}+\tfrac14(B_--B_+)\right]^{-1}.
\]
Then
\begin{equation}
 \frac{I}{1+16\rho/(1+2\rho)}
       \le C_{\rm eff}\le30I ,
 \label{f16v2:Ceffbounds}
\end{equation}
and, with
\(\widetilde\kappa=\kappa_2(\det C_{\rm eff}/\det G)^{3/2}\),
\begin{equation}
 \mathcal B_{2,0}(y,y')
   =\widetilde\kappa\,
      e^{-y^{\mathsf T}B_+y/4}
       \varphi_{C_{\rm eff}}(y-y')
      e^{-y'^{\mathsf T}B_+y'/4}.
 \label{f16v2:effective}
\end{equation}
\end{theorem}

\begin{proof}
Since \(V_{2,+}\ge C_B^{-1}/2+A^{\mathsf T}A\),
\[
 A V_{2,+}^{-1}A^{\mathsf T}
              \le2K(I+2K)^{-1}.
\]
The commuting functions of \(K\) obey the exact identity
\[
 2I-(I+K)^{-1}-4K(I+2K)^{-1}
                    =(I+K)^{-1}(I+2K)^{-1}.
\]
Substitute in \(B_+=2D^{\mathsf T}D-H_{\rm eff}-2P_{2,+}\)
and use \(0\le K\le\rho I\) for the lower bound.
The upper bound follows from \(P_{2,-}\ge0\) and
\(H_{\rm eff}\ge(1+\rho)^{-1}D^{\mathsf T}D\).
Inversion of \(V_{2,-}\ge V_{2,+}\) gives \(B_-\ge B_+\).

Furthermore
\[
 0\le B_--B_+
   =2D^{\mathsf T}A(V_{2,+}^{-1}-V_{2,-}^{-1})A^{\mathsf T}D
   \le\frac{4\rho}{1+2\rho}D^{\mathsf T}D.
\]
Every bridge belongs to four plaquettes; each mixed row has at
most four bridge entries. Thus \(\|D\|^2\le16\).
Together with \(I\le G\le30I\), this proves
\eqref{f16v2:Ceffbounds}. Finally split the exponent in
\eqref{f16v2:schur-kernel} into the two endpoint \(B_+\)
terms and the difference term \(B_--B_+\).
The Gaussian prefactor changes by the displayed determinant
ratio, proving \eqref{f16v2:effective}.
\end{proof}

These comparisons prove \(\ker B_+=\ker D\).
For the common-bridge subspace \(\mathcal E\) of V1,
\(A^{\mathsf T}Dy=0\) and \(H_{\rm eff}y=D^{\mathsf T}Dy\).
Thus
\[
 B_+y=B_-y=D^{\mathsf T}Dy,\qquad
            (B_--B_+)y=0\quad(y\in\mathcal E).
\]
The full leading cube propagation has not generated a constant
restoring term for those coarse collective modes.

For clarity, even the positive oscillator frequencies of
\eqref{f16v2:effective} have a soft finite-volume edge.
Let \(\nu_{\min}^+\) be the least positive eigenvalue of
\(C_{\rm eff}^{1/2}B_+C_{\rm eff}^{1/2}\).
Then
\begin{equation}
                     \nu_{\min}^+\le60\sin^2(\pi/m).
 \label{f16v2:soft}
\end{equation}
To prove it, use V1's common-bridge transverse cosine \(v\),
orthogonal in the Euclidean metric to \(\ker D\).
Its numerator is \(\langle v,D^{\mathsf T}Dv\rangle\).
Choose the representative \(v+w\), \(w\in\ker D\), orthogonal
to that kernel in the \(C_{\rm eff}^{-1}\) metric.
Its numerator is unchanged, and its denominator is at least
\(\|v+w\|^2/30\ge\|v\|^2/30\).
The generalized minimum principle and V1's exact quotient
\(2\sin^2(\pi/m)\) prove \eqref{f16v2:soft}.

The zero directions include residual linearized gauge directions
and torons. The full Euclidean Gaussian operator has free
diffusion in these directions, not a normalizable isolated vacuum.
If those directions are removed to discuss just its positive
oscillator factors, a parameter \(\nu>0\) has Mehler ratio
\((1+\nu/2+\sqrt{\nu+\nu^2/4})^{-1}\).
Equation~\eqref{f16v2:soft} concerns this specified Gaussian
factorization. It is neither a native physical gaplessness
theorem nor a choice of physical time scaling.

\subsection{The remaining full-operator return and verification}

The exact compact kernel and the full Gaussian Schur response
now include both time ends and all internal propagation within
one step. The missing spectral implication is still substantive:
a one-step compression does not bound the full complement.
With \(P_\gamma=\mathcal I_\gamma\mathcal I_\gamma^*\),
the actual two-step return is
\[
 \mathcal I_\gamma^*T_{\rm full,\gamma}
       (I-P_\gamma)T_{\rm full,\gamma}\mathcal I_\gamma .
\]
It is nonnegative, but no sufficiently small relative bound
uniform in spatial volume has been proved. The native bridge
tails and the nonlinear common-bridge dynamics must also be
kept under the same normalization. The fixed-volume domination
in Proposition~\ref{f16v2:native-limit} is not such a bound.

The next upstream target is to retain this return as a
two-time or Schur/memory operator and test a volume-uniform
relative comparison, rather than identify the compressed
Gaussian spectrum with the full transfer spectrum.
Independent physical-time calibration and interacting
four-dimensional continuum reconstruction remain open.

The diagnostic
\begin{center}\small
\path{scripts/verify_ym_f16_v2_joint_dressed_transfer.py}
\end{center}
checks noncommuting two-ended compact words, the actual
two-cube kinetic completion and determinant, full-torus port
incidences, the independent two-end Gaussian precision and
plus/minus Schur formulas, the matrix identity behind
\eqref{f16v2:Bbounds}, and the endpoint magnetic-square
distinction. Finite checks supplement the proofs; they are not
a continuum spectral certificate. All 4,582 new exact fixtures passed,
together with the inherited V1 chain. The diagnostic uses exact
elimination for the larger determinants, rather than the inherited
small-matrix cofactor routine. It has 8,154 bytes and SHA--256
\begin{center}\scriptsize\ttfamily
95087B3BF6B4AA9517C734E8A35B690B9074EF3D9355AC7125B2A22340C59DCF.
\end{center}

The predecessor V1 has 25,601 bytes and SHA--256
\begin{center}\scriptsize\ttfamily
A04BB24BC54769C31EF36F591282F6ABCF828DA25920C369DCB9DA9F1EF27097.
\end{center}
Its body is preserved byte for byte; removing this append and
restoring the title recovers that source. All 167 previously
protected files remain unchanged. The last companion
checkpoint is f15 V100 and the next is this lineage's V5.
Only \texttt{.tex} files belong in \texttt{latex\_notes};
build and visual-check products are kept outside that folder.

\begin{firewall}
V2 gives an exact full-compact joint kernel, a fixed-volume
native dressed-compression limit with all internal ends
integrated, and explicit volume-independent comparisons for
its complete leading two-time Gaussian response.
No small volume-uniform full-operator return, positive
physical continuum mass gap, or interacting continuum
Yang--Mills construction is proved. The goal remains unachieved.
\end{firewall}

\section{V3: extensive dressed-line leakage and a retained fast history}
\label{f16v3:context}

V2 computed one full step between the normalized dressed ends.
We now ask what happens before inserting another dressed projection.
The answer is not a small correction to the static conditional line:
even the complete leading Gaussian transfer changes its fast
covariance by an extensive amount. We prove a quantitative
pre-root-restriction obstruction to a small relative return.
We then integrate the fast variables along an arbitrary bridge
history without intermediate projection. Its covariance and
quadratic memory have volume- and time-length-independent decay.

This is a change in the next analytic target, not a mass-gap claim.
F15 V99--V100 concerned the magnetic image and its old oscillator
occupations; V1 concerned bridge-only affinity; V2 concerned the
joint one-step compression. The new calculation propagates the
dressed image by the full Gaussian transfer and retains arbitrarily
many fast time slices. The archive search also found spatial
Gaussian blocking memory in f11 V32--V33 and native observable
covariance estimates in f15 V61; neither is the conditional
fast-history calculation below.

\subsection{The propagated conditional vector is not the static dressed vector}

All statements in this subsection concern V2's leading Gaussian
model before the remaining root Gauss restriction.
Write \(\mathsf T_0\) for its full transfer on the internal and
bridge Euclidean coordinates, and
\[
 (\mathcal J_0 f)(x,y)=\psi_y(x)f(y),\qquad
 \psi_y(x)=\frac{\Omega(x)e^{-|Ax+Dy|^2/4}}{\sqrt{F_2(y)}} .
\]
The vector \(\psi_y\) has norm one in the internal coordinates.
Each displayed path matrix is tensored with \(I_3\) when acting
on the three colour components. Put
\begin{align}
 N&=\mathsf M_B^{-1},&
 S&=\mathsf H_B+A^{\mathsf T}A,&
 \Lambda&=A^{\mathsf T}D,\notag\\
 E_0&=N+\tfrac12(C_B^{-1}+\mathsf H_B),&
 E&=E_0+A^{\mathsf T}A,&
 P&=\tfrac12(C_B^{-1}+A^{\mathsf T}A).
 \label{f16v3:definitions}
\end{align}
Here \(E\) is V2's two-end diagonal block \(V_{2,+}+N\).
The precision in the amplitude \(\psi_z\), as opposed to its
squared probability density, is \(P\):
\[
 \psi_z(x)=c(z)\exp[-x^{\mathsf T}Px/2-x^{\mathsf T}\Lambda z/2],
                      \qquad c(z)>0.
\]
The sources and the positive scalar depend on the bridge field.

For fixed initial bridge coordinate \(y\) and final coordinate \(z\),
define the fast output feature
\[
 q_{z,y}(x)=\int \mathsf T_0(x,z;x',y)\psi_y(x')\,dx'.
\]
This definition has no intermediate fast projection.

\begin{proposition}[Full-step covariance mismatch]
\label{f16v3:propagation}
Set
\begin{equation}
 \Delta_f=N(E_0^{-1}-E^{-1})N\ge0.
 \label{f16v3:delta}
\end{equation}
There is a positive scalar \(c(z,y)\) such that
\begin{equation}
 q_{z,y}(x)=c(z,y)
 \exp\left[
   -\tfrac12x^{\mathsf T}(P+\Delta_f)x
   -x^{\mathsf T}\bigl(\tfrac12\Lambda z+NE^{-1}\Lambda y\bigr)
 \right].
 \label{f16v3:feature}
\end{equation}
Consequently the squared normalized output feature has covariance
\([2(P+\Delta_f)]^{-1}\otimes I_3\), independent of \(y,z\).
It is not the covariance \((2P)^{-1}\otimes I_3\) of the static
dressed line unless \(A=0\).
\end{proposition}

\begin{proof}
The full kernel is the product of its internal and bridge heat
densities and the two endpoint factors
\[
 \exp[-(x^{\mathsf T}\mathsf H_Bx+|Ax+Dz|^2)/4].
\]
After multiplication by \(\psi_y(x')\), the input precision is \(E\)
and its linear source is \(Nx-\Lambda y\).
Gaussian integration of \(x'\) leaves output precision
\[
 Q=N+\tfrac12S-NE^{-1}N
\]
and linear source \(-\Lambda z/2-NE^{-1}\Lambda y\).
The cube ground-state equation gives the exact identity
\begin{equation}
 N+\tfrac12(\mathsf H_B-C_B^{-1})=NE_0^{-1}N.
 \label{f16v3:riccati}
\end{equation}
For completeness, conjugate by \(\mathsf M_B^{1/2}\).
On a cube mode with parameter \(t=4\) or \(6\), the two sides
become
\[
 1+t/2-\sqrt{t+t^2/4}
       =\bigl(1+t/2+\sqrt{t+t^2/4}\bigr)^{-1};
\]
their product identity follows by squaring the radical.
This proves \eqref{f16v3:riccati}, hence \(Q=P+\Delta_f\).
Inversion reverses \(E\ge E_0\), so \(\Delta_f\ge0\).
If \(\Delta_f=0\), invertibility of \(N,E,E_0\) implies
\(A^{\mathsf T}A=0\). The covariance assertion follows by squaring
the amplitude and normalizing its positive Gaussian.
\end{proof}

\subsection{A volume-extensive obstruction to a small static return}

The covariance mismatch gives a uniform conditional overlap
bound, not just a nonzero example.

\begin{theorem}[Dressed-line retention after a full step]
\label{f16v3:retention}
For the actual parity-cube geometry with \(B=m^3\), define
\[
 \eta_B(z,y)=
       \frac{|\langle\psi_z,q_{z,y}\rangle|^2}
            {\|q_{z,y}\|_2^2}.
\]
Then, for every \(y,z\),
\begin{equation}
 0<\eta_B(z,y)\le
 \eta_B^{\rm cov}:=
 \frac{\det(I+Z)^{3/2}}{\det(I+Z/2)^3}
 \le e^{-c_*B}\le e^{-B/100000},
 \qquad
 c_*=\frac{50}{147}\frac{(25/1728)^2}{5}>10^{-5}.
 \label{f16v3:retention-bound}
\end{equation}
Here \(Z\) is the relative covariance-precision change defined
in the proof. The determinants retain all \(5B\) path modes,
with the displayed power accounting for all three colours.
\end{theorem}

\begin{proof}
Use bars to denote congruence by \(N^{-1/2}=\mathsf M_B^{1/2}\);
in particular
\[
 \overline E_0=N^{-1/2}E_0N^{-1/2},\quad
 W=N^{-1/2}A^{\mathsf T}AN^{-1/2},\quad
 \overline P=N^{-1/2}PN^{-1/2}.
\]
Then
\[
 \overline\Delta_f
   =\overline E_0^{-1}-(\overline E_0+W)^{-1},\qquad
 Z=\overline P^{-1/2}\overline\Delta_f\,\overline P^{-1/2}.
\]
The eigenvalues of \(Z\) are the generalized eigenvalues of
\((\Delta_f,P)\); no commuting-matrix assumption is being made.

The inherited cube spectrum and covariance give
\[
 5I\le\overline E_0\le8I,\qquad 2I\le\overline P\le6I.
\]
Indeed the exact bounds for \(\overline E_0\) are
\(3+2\sqrt2\) and \(4+\sqrt{15}\).
For the second assertion, the spectrum of the whitened
\(C_B^{-1}\) lies between \(4\sqrt2\) and \(2\sqrt{15}\), and
\(0\le W\le4I\).
This last inequality follows from
\(A=C_{\rm int}U_B\), \(\|C_{\rm int}\|\le2\), and
\(U_B^{\mathsf T}U_B=N\).
Moreover the exact diagonal Gram blocks
\((A^{\mathsf T}A)_{bb}=2N_{bb}\) from f15 V99 give
\[
                           \Tr W=10B.
\]
These are bounds for the actual geometry, not assumptions
on independently chosen toy matrices.

Woodbury inversion, or continuous extension if \(W\) is singular,
gives
\[
 \overline\Delta_f
 =\overline E_0^{-1}W^{1/2}
   (I+W^{1/2}\overline E_0^{-1}W^{1/2})^{-1}
                       W^{1/2}\overline E_0^{-1}
 \ge\frac59\,\overline E_0^{-1}W\overline E_0^{-1}.
\]
Taking traces and using \(\overline E_0\le8I\) yields
\[
 \Tr\overline\Delta_f\ge\frac{25}{288}B,\qquad
 \Tr Z\ge\frac{25}{1728}B,\qquad 0\le Z\le\frac1{10}I.
\]
For the last estimate use
\(\overline\Delta_f\le\overline E_0^{-1}\le I/5\) and
\(\overline P\ge2I\).
Cauchy--Schwarz on the \(5B\) eigenvalues gives
\[
                   \Tr Z^2\ge\frac{(25/1728)^2}{5}\,B.
\]

The squared overlap of two normalized Gaussian amplitudes with
precisions \(P\) and \(P+\Delta_f\) is at most the value when
their means coincide. Completing the square shows that unequal
means multiply this value by \(e^{-a}\) for a nonnegative
quadratic form \(a\). The coincident-mean determinant is precisely
\(\eta_B^{\rm cov}\) in \eqref{f16v3:retention-bound}; congruence changes
do not change it. For an eigenvalue \(0\le z\le1/10\),
\[
 \frac32\log\frac{(1+z/2)^2}{1+z}
 =\frac32\log\left(1+\frac{z^2}{4(1+z)}\right)
 \ge\frac{3z^2}{8(1+z/2)^2}
 \ge\frac{50}{147}z^2 .
\]
We used \(\log(1+u)\ge u/(1+u)\), which follows by integrating
its nonnegative derivative difference from zero.
Summation proves the stated exponential bound.
\end{proof}

Let
\[
 \mathsf B_0=\mathcal J_0^*\mathsf T_0\mathcal J_0,\quad
 \mathsf A_2=\mathcal J_0^*\mathsf T_0^2\mathcal J_0,\quad
 \mathsf R_0=\mathsf A_2-\mathsf B_0^2
 =\mathcal J_0^*\mathsf T_0(I-\mathcal J_0\mathcal J_0^*)
                              \mathsf T_0\mathcal J_0\ge0.
\]
Their diagonal kernels satisfy
\begin{align}
 \mathsf A_2(y,y)&=\int\|q_{z,y}\|_2^2\,dz,\notag\\
 \mathsf B_0^2(y,y)&=\int|\langle\psi_z,q_{z,y}\rangle|^2\,dz
                 \le e^{-c_*B}\mathsf A_2(y,y).
 \label{f16v3:diagonal}
\end{align}
All these diagonal integrals are finite. In fact the internal
cut transfer has norm \((\lambda_\Box^0)^B\); its magnetic
multipliers have norm at most one. Hence
\[
 \|q_{z,y}\|_2
       \le(\lambda_\Box^0)^B\varphi_G(z-y),
\]
whose square is integrable in \(z\), uniformly in \(y\).

For any bounded bridge window \(U\) of positive measure,
the resulting operators localized at their input are
Hilbert--Schmidt and
\begin{equation}
 \|(I-\mathcal J_0\mathcal J_0^*)\mathsf T_0
                      \mathcal J_0\,1_U\|_{\rm HS}^2
 \ge(1-e^{-c_*B})
              \|\mathsf T_0\mathcal J_0\,1_U\|_{\rm HS}^2.
 \label{f16v3:HS-loss}
\end{equation}
This follows by integrating \eqref{f16v3:diagonal}, or by
orthogonal decomposition of each output feature.
In particular a quadratic-form inequality
\(\mathsf R_0\le\varepsilon\mathsf A_2\) on the entire
pre-root-restriction coarse space requires
\(\varepsilon\ge1-e^{-c_*B}\). To see this, sandwich by \(1_U\)
and take the finite trace; the denominator is strictly positive.
Thus no fixed \(\varepsilon<1\) works for all volumes in this model.

This conclusion is stronger than a nonzero projection defect,
but its scope is important. It does not assert pointwise
positivity of every off-diagonal return kernel, nor does it
turn a diagonal comparison into a reverse operator inequality.
It does not apply automatically after physical Gauss
restriction, to a low-energy subspace, or to the native compact
transfer. In particular, V2's fixed-window one-step convergence
alone does not transfer this full two-step statement to that
native operator.

\subsection{Exact elimination along an arbitrary bridge history}

We now avoid intermediate projection altogether. Let \(n\ge1\),
fix the bridge history \(\mathbf y=(y_0,\ldots,y_n)\), and
integrate the internal coordinates \(x_0,\ldots,x_n\) in the
kernel of \(\mathcal J_0^*\mathsf T_0^n\mathcal J_0\).
Set \(E_{\rm i}=S+2N\), and define the \((n+1)\)-by-\((n+1)\)
block precision
\begin{equation}
 (W_n)_{00}=(W_n)_{nn}=E,\qquad
 (W_n)_{ii}=E_{\rm i}\ (0<i<n),\qquad
 (W_n)_{i,i+1}=(W_n)_{i+1,i}=-N,
 \label{f16v3:history-precision}
\end{equation}
with all other time blocks zero.
Its time blocks have \(5B\) path coordinates.
Let \(\mathcal L\mathbf y=(\Lambda y_0,\ldots,\Lambda y_n)\) and
\[
 a_n=\frac1{F_2(0)}
       \left[\det C_B\,(\det\mathsf M_B)^n\det W_n\right]^{-3/2}.
\]

\begin{theorem}[Complete fast-history kernel]
\label{f16v3:history}
After integration of every internal time slice, the exact
Gaussian history weight is
\begin{align}
 \mathcal W_n(\mathbf y)
 &=a_n\prod_{i=0}^{n-1}\varphi_G(y_i-y_{i+1})
       \exp\bigl(\mathcal E_n(\mathbf y)\bigr),\notag\\
 \mathcal E_n(\mathbf y)
 &=-\frac12\sum_{i=0}^n y_i^{\mathsf T}D^{\mathsf T}Dy_i
   +\frac14\bigl(y_0^{\mathsf T}H_{\rm eff}y_0
                         +y_n^{\mathsf T}H_{\rm eff}y_n\bigr)
   +\frac12(\mathcal L\mathbf y)^{\mathsf T}
                           W_n^{-1}(\mathcal L\mathbf y).
 \label{f16v3:history-weight}
\end{align}
The kernel of \(\mathcal J_0^*\mathsf T_0^n\mathcal J_0\) is
\(\int\mathcal W_n(\mathbf y)\,dy_1\cdots dy_{n-1}\).
No internal state or intervening bridge configuration has
been discarded. The determinant in \(a_n\) is part of this
identity, not a freely chosen spectral normalization.
\end{theorem}

\begin{proof}
Each interior time slice receives two magnetic endpoint
multipliers, whereas each external slice receives one
transfer multiplier and one multiplier from its dressed
end vector. Thus the mixed magnetic power is two at
\emph{every} slice. Each linear source is consequently
\(-\Lambda y_i\).
The internal magnetic power is two at interior slices;
at the two ends there is one such factor and one
\(\Omega\) amplitude. The kinetic differences contribute
\(\sum_{i=0}^{n-1}(x_i-x_{i+1})^{\mathsf T}N(x_i-x_{i+1})\).
These observations give exactly \eqref{f16v3:history-precision}.
Its quadratic form can also be written as
\begin{align*}
 \mathbf x^{\mathsf T}W_n\mathbf x
 ={}&x_0^{\mathsf T}(E-N)x_0+x_n^{\mathsf T}(E-N)x_n\\
 &+\sum_{i=1}^{n-1}x_i^{\mathsf T}Sx_i
   +\sum_{i=0}^{n-1}(x_i-x_{i+1})^{\mathsf T}N(x_i-x_{i+1}),
\end{align*}
which proves strict positivity.
Gaussian integration contributes the last term in
\(\mathcal E_n\) and \(\det(W_n)^{-3/2}\).
The scalar \((2\pi)^{15B(n+1)/2}\) from integration cancels
the \(n\) internal heat prefactors and the two ground-amplitude
prefactors. Their remaining factors are
\((\det\mathsf M_B)^{-3n/2}(\det C_B)^{-3/2}\).
Finally the two factors \(F_2^{-1/2}\) contribute \(F_2(0)^{-1}\)
and the displayed external \(H_{\rm eff}\) term.
This proves the history formula. Positivity permits the
iterated integrations; their interpretation as the stated
bounded-operator kernel follows from composition of the
positive Gaussian transfer kernels.
\end{proof}

For \(n=1\), the plus/minus transform has precisions
\(E-N=V_{2,+}\) and \(E+N=V_{2,-}\), so this is exactly V2,
including its scalar. For every \(n\) one may compute the
determinant without dropping its boundary terms by
\[
 R_0=E,\qquad R_i=(W_n)_{ii}-NR_{i-1}^{-1}N,\qquad
                   \det W_n=\prod_{i=0}^n\det R_i .
\]
All pivots are strictly positive by block Gaussian elimination.
The order of the matrix products matters.

\subsection{Uniform temporal and space--time locality of the fast memory}

Write \(\mathbb N=I_{n+1}\otimes N\) and
\(\widehat W_n=\mathbb N^{-1/2}W_n\mathbb N^{-1/2}\).
The graph for space--time locality has a vertex for each
(time slice, spatial block), with the nearest-time and
nearest-spatial-block edges.

\begin{theorem}[Length-independent conditional fast response]
\label{f16v3:locality}
For every \(B,n\),
\begin{equation}
                        4I\le\widehat W_n\le15I.
 \label{f16v3:uniform-spectrum}
\end{equation}
For sets \(\mathcal E,\mathcal F\) at graph distance \(d\),
\begin{equation}
 \|P_{\mathcal E}\widehat W_n^{-1}P_{\mathcal F}\|
                         \le\frac14(11/15)^d .
 \label{f16v3:spacetime}
\end{equation}
A sharper purely temporal bound is
\begin{equation}
 \|N^{1/2}(W_n^{-1})_{ij}N^{1/2}\|
       \le\frac{q^{|i-j|}}{p_*-2},
 \qquad p_*=3+2\sqrt2,\quad q=2/p_*<1.
 \label{f16v3:temporal}
\end{equation}
These are conditional fast-covariance bounds for arbitrary
bridge histories, not decay of the complete physical
Yang--Mills correlations.
\end{theorem}

\begin{proof}
Whitening the quadratic-form decomposition in the preceding
proof gives a time-path Laplacian, with norm at most four.
The interior on-site term is the whitened
\(\mathsf H_B+A^{\mathsf T}A\), between \(4I\) and \(10I\).
The external on-site term is the whitened \(E-N\), between
\((2+2\sqrt2)I\) and \((7+\sqrt{15})I<11I\).
This proves \eqref{f16v3:uniform-spectrum}.
All off-site entries of \(\widehat W_n\) join adjacent
space--time blocks: \(A^{\mathsf T}A\) has only nearest-block
entries, and the whitening is spatially block diagonal.
Thus \(T_n=I-\widehat W_n/15\) satisfies
\(0\le T_n\le(11/15)I\), with the same finite range.
The convergent expansion
\[
             \widehat W_n^{-1}=\frac1{15}\sum_{k\ge0}T_n^k
\]
has no connection between those sets before degree \(d\).
The geometric tail is \((11/15)^d/4\), proving
\eqref{f16v3:spacetime} without a set-cardinality factor.

For the temporal estimate, keep entire spatial slices as blocks.
The diagonal blocks of \(\widehat W_n\) are at least \(p_*I\):
the external ones are \(\overline E_0+W\), and the interior
ones are at least \(6I\). Its temporal off-diagonal blocks are
\(-I\). If \(\mathcal D_n\) is its time-block diagonal,
write \(\widehat W_n=\mathcal D_n-\mathcal A_n\), where
\(\mathcal A_n\) is the path adjacency with identity blocks.
Then
\(\|\mathcal D_n^{-1/2}\mathcal A_n\mathcal D_n^{-1/2}\|
\le2/p_*=q\).
This matrix need not be positive, but its norm-convergent
Neumann series is valid. A term of degree less than
\(|i-j|\) has zero \(ij\) block.
Multiplication by the two diagonal inverse square roots and
summation give
\(p_*^{-1}q^{|i-j|}/(1-q)=q^{|i-j|}/(p_*-2)\).
Undoing the whitening proves \eqref{f16v3:temporal}.
\end{proof}

There is also a curvature-relative truncation estimate.
Let \(\mathcal Q_n(\mathbf y)\) denote the last term in
\(\mathcal E_n\), and let \(\mathcal Q_n^{[r]}\) keep only
its ordered pairs with \(|i-j|\le r\), for an integer \(r\ge0\).
Since \(\|N^{-1/2}A^{\mathsf T}\|\le2\), the temporal bound yields
\begin{equation}
 |\mathcal Q_n(\mathbf y)-\mathcal Q_n^{[r]}(\mathbf y)|
 \le \frac{4}{p_*-2}\frac{q^{r+1}}{1-q}
                               \sum_{i=0}^n\|Dy_i\|^2 .
 \label{f16v3:memory-tail}
\end{equation}
Indeed each ordered matrix element is bounded by
\(4q^{|i-j|}\|Dy_i\|\|Dy_j\|/(p_*-2)\).
The omitted scalar matrix has row sum at most
\(2q^{r+1}/(1-q)\); its quadratic form is bounded by that
row sum times \(\sum_i\|Dy_i\|^2\). Include the factor
\(1/2\) in \(\mathcal Q_n\) to obtain the displayed constant.
No spatial volume or time-length factor multiplies this
relative curvature budget. Truncation is an estimate on the
quadratic action, not automatically an operator-norm or
positivity-preserving approximation.

For histories taking values in V1's common-bridge subspace,
\(\Lambda y_i=0\) at every time. The entire fast source
response \(\mathcal Q_n\), not merely its long-time tail,
then vanishes. The determinant \(a_n\) is independent of
the bridge history. This statement concerns the fixed-history
weight; integrating other bridge modes is a further operation.
It does not create a constant restoring term or settle the
native dynamics of those common modes.

\subsection{Direction of the next native estimate and verification}

The exact native counterpart of the return is still
\[
 \mathcal I_\gamma^*T_{\rm full,\gamma}
        (I-\mathcal I_\gamma\mathcal I_\gamma^*)
                         T_{\rm full,\gamma}\mathcal I_\gamma.
\]
V3 shows why asking for a small relative return on the entire
static dressed carrier is not a promising unrestricted
Gaussian target: \eqref{f16v3:HS-loss} goes the opposite way.
A physical low-energy restriction would require its own
estimate; it is not supplied by this obstruction.

The constructive replacement is the unprojected fast-history
integral. Its Gaussian precision remains uniformly coercive
even though the retained bridge dynamics has soft directions.
The next useful native theorem must control the non-Gaussian
conditional fast integral, its determinant or pressure
normalization, and its bridge-field tails in a norm compatible
with the physical Gauss restriction and the same full transfer
top value. The compact action is not globally convex in a
single chart, and V2's constants depend on volume. Consequently
\eqref{f16v3:memory-tail} is not yet a native uniform cluster or
spectral estimate. Physical-time calibration and an interacting
continuum construction remain additional open requirements.

The new diagnostic is
\begin{center}\small
\path{scripts/verify_ym_f16_v3_fast_history.py}.
\end{center}
Its 2,075 exact rational fixtures check the noncommuting
propagated precision and source, unequal-mean overlap
completion, full-history action and determinant recursion,
one-step agreement with V2, finite-range propagation, and
length-independent memory tails. They also check the rational
constants in the extensive leakage bound. The inherited V2
chain passed. Finite fixtures supplement the proofs and do
not certify a physical mass gap.
The checker has 7,716 bytes and SHA--256
\begin{center}\scriptsize\ttfamily
811A66019E7A8B4A78D6CF8C1F8F6B2882D318BBA5D2B57D6E4D5C274275EA02.
\end{center}

The predecessor V2 has 48,007 bytes and SHA--256
\begin{center}\scriptsize\ttfamily
D8BC72F2B7E92E94C2B8E3AF848D9580563E5F9302CA13E8843909C42AB80F6F.
\end{center}
Its body is preserved byte for byte: remove this append and
restore the title to recover it. The 168 previously protected
files are unchanged. The last companion checkpoint is f15 V100;
the next is f16 V5. Only \texttt{.tex} files belong in
\texttt{latex\_notes}; build products remain elsewhere.

\begin{firewall}
V3 proves extensive static dressed-line loss and a conditional,
unprojected fast-history formula with uniform Gaussian
space--time locality and curvature-relative temporal truncation.
The leakage obstruction is before physical Gauss restriction.
No native volume-uniform memory bound, nontrivial interacting
continuum construction, or positive physical Yang--Mills
mass gap is established. The full goal remains unachieved.
\end{firewall}

% BEGIN F16 V4 APPEND -- 2026-09-17
% Insert immediately before \end{document} in the supplied V3.
% This append uses only the theorem environments and packages of V3.

\clearpage
\section{V4: a native fast-history core and the nonlinear common-bridge return}
\label{f16v4:context}

V3 supplies an arbitrary-length Gaussian fast-history integral and rules
out a small relative return on the entire static dressed line before
root restriction. Repeating a one-step compression cannot repair that
obstruction. We therefore keep the nonroot temporal gauge variables as
\emph{additional fast variables}, rather than first integrating them into
a logarithmic kinetic kernel. This preserves the finite-range structure
of the exact nonlinear action.

The result below is a native, normalized, \emph{restricted} conditional
history theorem, uniform in spatial volume and history length. In
addition to exponential decay of its two-source temporal response, it
identifies an extra nonlinear small factor for distant common-bridge
sources. A separate pressure-density estimate pays the normalization of
the restricted integral. The passage from this core to the full compact
conditional remains an explicit, unestimated correction; it is not
absorbed into a choice of vacuum normalization.

\subsection{Dependency and non-duplication ledger}

The background audit used the four supplied sources: this lineage's V3,
f14 V100, f15 V100, and the Grand Tensor Monograph V076. It consisted of
section/result inventories and targeted reading of the relevant
interfaces, not a certification of every proof in those archives.
The following distinctions are essential.

\begin{center}\small
\begin{tabular}{@{}p{0.22\textwidth}p{0.72\textwidth}@{}}
\toprule
Inherited interface & Use here, without claiming it again as a new result\\
\midrule
f14 V64 & Nonlinear high-fibre concentration on a restricted convex
coordinate domain. This is not coverage of that domain by the full law.\\[1mm]
f15 V13 & Fixed product-ball conditionals, boundary-aware Witten
covariance, and weighted spatial response. We reuse its absolute
one-form realization.\\[1mm]
Monograph V076 & The direct-Hessian-minus-covariance identity and the
Witten representation, especially \texttt{thm:YM-Witten-router}. A return
of already integrated variables is charged once.\\[1mm]
f15 V97--V100 & Exact cube geometry, coupled compact kernels, moving
frames, and the magnetic dressed carrier.\\[1mm]
This lineage V2--V3 & The exact two-ended kernel, kinetic completion,
common-bridge cancellation, and Gaussian history precision.\\
\bottomrule
\end{tabular}
\end{center}

The new specialization is the complete nonlinear \emph{transfer-time}
conditional with the temporal port variables retained. Its two-source
common-bridge memory requires two small nonlinear couplings. Neither a
new abstract covariance identity nor a further Gaussian determinant is
being offered as the advance. No claim of general literature novelty is
made for the analytic tools.

\subsection{The complete native history density, including its scalar}

Work first before the remaining root Gauss restriction, with the exact
moving coordinates of Proposition~\ref{f16v2:exact-cut}. Fix
\(n\ge1\) and a bridge history \(\mathbf Z=(Z_0,\ldots,Z_n)\).
There is an internal chord configuration \(X_i\) at every slice and a
nonroot temporal gauge configuration \(r_t\in\SU(2)^{7B}\) on every
bond \(t\to t+1\). The root entries of \(r_t\) equal the identity.
Define the exact kinetic words
\begin{align}
 Q_{t,e}^{\rm int}
 &=r_{t,s}^{-1}\bar X_{t,e}r_{t,t(e)}\bar X_{t+1,e}^{-1},\notag\\
 Q_{t,e}^{\rm br}
 &=r_{t,s}^{-1}Z_{t,e}r_{t,t(e)}Z_{t+1,e}^{-1}.
 \label{f16v4:words}
\end{align}
Here \(s=s(e)\), \(t(e)\) denotes the spatial target vertex (not the
bond index), and \(\bar X\) has exactly V2's moving-frame meaning.
Let \(Q_{t,e}\) mean the appropriate word for any of the \(24B\)
fine spatial links.

Set
\[
 w_0=w_n=\tfrac12,\qquad w_i=1\quad(0<i<n).
\]
The factors in the native fixed-history integral are: \(m_\gamma^2\)
at every slice; \(b_\gamma\) at the two ends and
\(b_\gamma^2\) in the interior; one \(\Phi_\gamma\) at each end;
and all \(24Bn\) Wilson kinetic factors. Thus there is no internal
projection at an intermediate time.

We record the chart constants because subsequent logarithmic estimates
must refer to this same integral. In the convention of V1,
\[
 dH(\operatorname{Exp}(q/\sqrt\gamma))
  =c_{H,\gamma}j(q/\sqrt\gamma)\,dq,\qquad
 c_{H,\gamma}=\frac1{2\pi^2\gamma^{3/2}},\qquad
 j(a)=\left(\frac{\sin|a|}{|a|}\right)^2,
\]
on \(|q|<\pi\sqrt\gamma\), up to the Haar-null cut set. Write
\(X_i=\operatorname{Exp}(x_i/\sqrt\gamma)\),
\(r_t=\operatorname{Exp}(z_t/\sqrt\gamma)\), and
\(Z_i=\operatorname{Exp}(y_i/\sqrt\gamma)\), componentwise.
All path matrices below act with the colour factor \(I_3\), which is
suppressed.

The normalized full Euclidean cube amplitude and its exact chart mass are
\begin{align}
 \Omega(x)&=c_\Omega e^{-x^{\mathsf T}C_B^{-1}x/4},&
 c_\Omega&=(2\pi)^{-15B/4}\det(C_B)^{-3/4},\notag\\
 p_{\gamma,B}&=\int_{\mathcal P_x}\Omega(x)^2\,dx,&
 \Phi_\gamma(X(x))&=
       \frac{\Omega(x)}{\sqrt{p_{\gamma,B}J_{X,\gamma}(x)}}.
 \label{f16v4:amplitude}
\end{align}
Here \(\mathcal P_x\) is the complete product chord chart and
\(J_{X,\gamma}\) its Haar Jacobian. This is the carrier already chosen
in V1, not a newly substituted native vacuum.

On the complete internal and port chart product \(\mathcal P\), define
\begin{align}
 V_\gamma(x,z;\mathbf y)
 ={}&\tfrac14\bigl(x_0^{\mathsf T}C_B^{-1}x_0
                       +x_n^{\mathsf T}C_B^{-1}x_n\bigr)\notag\\
 &+\gamma\sum_{i=0}^n w_i S_{\rm int}(X_i)
   +\gamma\sum_{i=0}^n
             S_\times(X_i,\Theta_{X_i}^{-1}Z_i)\notag\\
 &+\gamma\sum_{t=0}^{n-1}\sum_{e\ {\rm fine}}s(Q_{t,e})\notag\\
 &-\sum_{i=0}^n w_i\sum_{a\ {\rm chord}}
                              \log j(x_{i,a}/\sqrt\gamma)
   -\sum_{t=0}^{n-1}\sum_{v\ {\rm nonroot}}
                              \log j(z_{t,v}/\sqrt\gamma).
 \label{f16v4:potential}
\end{align}
In particular, the nonlinear Wilson words and Haar factors are not
replaced by their Taylor polynomials in this definition. Put
\begin{align}
 Z_{\gamma,\mathrm{full}}(\mathbf y)
       &=\int_{\mathcal P}e^{-V_\gamma(x,z;\mathbf y)}\,dx\,dz,\notag\\
 J_{{\rm br},\gamma}(\mathbf y)
       &=\prod_{i=0}^n\prod_{e\ {\rm bridge}}
                         j(y_{i,e}/\sqrt\gamma)^{w_i},\notag\\
 c_{\gamma,B,n}
       &=p_{\gamma,B}^{-1}c_\Omega^2
                         (c_{H,\gamma}/z_\gamma)^{24Bn}.
 \label{f16v4:normalizers}
\end{align}

\begin{proposition}[Exact unprojected native history formula]
\label{f16v4:exact-history}
The full fixed-bridge-history weight in Euclidean bridge half-density
coordinates is
\begin{equation}
 \mathcal W_{\gamma,n}^{\mathrm{full}}(\mathbf y)
 =\frac{c_{\gamma,B,n}J_{{\rm br},\gamma}(\mathbf y)
                   Z_{\gamma,\mathrm{full}}(\mathbf y)}
        {\sqrt{F_{\gamma,2}(Z_0)F_{\gamma,2}(Z_n)}}.
 \label{f16v4:full-weight}
\end{equation}
Integrating \(y_1,\ldots,y_{n-1}\) over their complete principal charts
gives the chart kernel of
\(\mathcal I_\gamma^*T_{\rm full,\gamma}^{\,n}\mathcal I_\gamma\).
This is an identity at finite \(B,n,\gamma\), not a weak-field limit.
\end{proposition}

\begin{proof}
Compose the exact full kernel of Proposition~\ref{f16v2:exact-cut}
\(n\) times and attach the normalized amplitudes from V1 at the two
ends. Each interior magnetic multiplier occurs twice. At an external
slice its second occurrence comes from the dressed amplitude. This
produces precisely the internal and mixed action powers in
\eqref{f16v4:potential}.

At an external chord slice, the factor \(\Phi_\gamma\) cancels half
of the Haar Jacobian; interior chord slices retain the whole Jacobian.
The two external bridge half-densities and the interior bridge
integration measures have the same weights \(w_i\). Since
\(\sum_iw_i=n\), the constant Haar powers are \(5Bn\) from chords,
\(7Bn\) from temporal ports, and \(12Bn\) from bridges. Their total
is \(24Bn\), exactly the number of kinetic normalizers
\(z_\gamma^{-1}\). The two endpoint amplitudes contribute
\(p_{\gamma,B}^{-1}c_\Omega^2\). The nonconstant Haar factors are
those in \eqref{f16v4:potential} and \(J_{{\rm br},\gamma}\).
This proves the scalar as well as the action. Positivity justifies
composition and changes in integration order. The principal charts
cover the compact variables apart from null sets.
\end{proof}

\subsection{The augmented Gaussian reference and a fixed nonlinear core}

Avoid the notational collision with V3's endpoint matrix \(E\) by
writing \(E_{\rm br}\) for V2's bridge--nonroot incidence. Set
\[
 L_c=\mathsf B_B^{\mathsf T}\mathsf B_B,\qquad
 L_h=L_c+E_{\rm br}^{\mathsf T}E_{\rm br},\qquad
 \widehat L=L_c^{-1/2}L_hL_c^{-1/2}.
\]
All powers of \(L_c\) are spatially cube-block diagonal. Use the fast
coordinates
\[
       u_i=N^{1/2}x_i,\qquad v_t=L_c^{1/2}z_t,\qquad \xi=(u,v).
\]
Each \(u_{i,b}\) has dimension 15 and each \(v_{t,b}\) dimension 21.
Changing between \((x,z)\) and \(\xi\) has a constant Jacobian.
Partition functions below always use the original \(dx\,dz\) measure;
their normalized laws are pushed forward to \(\xi\). The constant
Jacobian cancels in those laws and in ratios with the same Gaussian
reference.

The quadratic reference potential, including its bridge-only part, is
\begin{align}
 V_0(x,z;\mathbf y)
 ={}&\tfrac14\bigl(x_0^{\mathsf T}C_B^{-1}x_0
                         +x_n^{\mathsf T}C_B^{-1}x_n\bigr)
       +\tfrac12\sum_{i=0}^n w_i x_i^{\mathsf T}\mathsf H_Bx_i
       +\tfrac12\sum_{i=0}^n|Ax_i+Dy_i|^2\notag\\
 &+\tfrac12\sum_{t=0}^{n-1}
       \left(\left|U_B(x_t-x_{t+1})+\mathsf B_Bz_t\right|^2
              +\left|y_t-y_{t+1}+E_{\rm br}z_t\right|^2\right).
 \label{f16v4:V0}
\end{align}
Because \(U_B^{\mathsf T}\mathsf B_B=0\), the fast Hessian in these
coordinates is the direct sum
\begin{equation}
 H_0:=D_\xi^2V_0
       =\widehat W_n\oplus(I_n\otimes\widehat L),\qquad
                         I\le H_0\le30I.
 \label{f16v4:H0}
\end{equation}
Indeed V3 gives \(4I\le\widehat W_n\le15I\); V2's bound
\(E_{\rm br}L_c^{-1}E_{\rm br}^{\mathsf T}\le29I\) gives
\(I\le\widehat L\le30I\). The \(u\)-diagonal time blocks are at
least \(p_*I\), where \(p_*=3+2\sqrt2\), and the only off-diagonal
time blocks of \(H_0\) are the adjacent \(u\)-blocks \(-I\).
There are no Gaussian \(u\)--\(v\) cross blocks.

Integrating all \(z_t\) in \eqref{f16v4:V0} by V2's kinetic completion
recovers V3's \(W_n\), bridge covariance \(G\), sources and determinant.
Equivalently, if \(Z_0=\int e^{-V_0}\,dx\,dz\) over the full Euclidean
fast space and \(c_{0,B,n}=c_\Omega^2(2\pi)^{-36Bn}\), then
\begin{equation}
       \mathcal W_n(\mathbf y)
        =\frac{c_{0,B,n}Z_0(\mathbf y)}
                       {\sqrt{F_2(y_0)F_2(y_n)}}.
 \label{f16v4:gaussian-consistency}
\end{equation}
This is a consistency check against Theorem~\ref{f16v3:history}, not a
new determinant theorem.

For \(R\ge1\) choose the \emph{separate-ball} product
\begin{equation}
 \mathcal O_R=
    \prod_{i=0}^n\prod_b\{|u_{i,b}|<R\}
       \ \times\!
    \prod_{t=0}^{n-1}\prod_b\{|v_{t,b}|<R\}.
 \label{f16v4:core}
\end{equation}
The two species must not be combined into a single ball: a combined
boundary would introduce a one-form boundary coupling absent from the
native action. The domain is fixed as the bridge history varies.
Let
\begin{equation}
 Z_{\gamma,R}(\mathbf y)=\int_{\mathcal O_R}e^{-V_\gamma}\,dx\,dz,
 \qquad
 \nu_{\gamma,R,\mathbf y}
       =\text{the normalized law on \(\mathcal O_R\) in \(\xi\)}.
 \label{f16v4:core-law}
\end{equation}
The notation \(\int_{\mathcal O_R}dx\,dz\) means integration over the
inverse linear image of \(\mathcal O_R\). Whenever this image is inside
\(\mathcal P\), this is \emph{exactly} the native augmented conditional
further conditioned on that event.

\begin{proposition}[Uniform native derivative certificate on the core]
\label{f16v4:derivatives}
There are finite constants \(C_*,C_1\ge1\), determined only by the
single-cube matrices and the finite list of native local words, such
that the following holds for every \(B,n\). Suppose
\begin{equation}
 \max_{i,e}|y_{i,e}|\le r,\qquad
 R\ge\max(1,r),\qquad
 \delta:=C_*R/\sqrt\gamma\le1/20.
 \label{f16v4:regime}
\end{equation}
The core is then inside the principal fast charts, and on its closure
\begin{align}
 |V_\gamma-V_0|&\le C_1B(n+1)R^3/\sqrt\gamma,\notag\\
 \|D_{(\xi,\mathbf y)}^2(V_\gamma-V_0)\|&\le\delta,\notag\\
       \tfrac{19}{20}I\le D_\xi^2V_\gamma&\le\tfrac{601}{20}I.
 \label{f16v4:derivative-bounds}
\end{align}
For a bridge direction \(h\) at a single time \(i\), put
\(J_{i,h}=\partial_{i,h}V_\gamma\). Then
\begin{equation}
       \|\nabla_\xi J_{i,h}\|_{L^\infty(\mathcal O_R)}
                                      \le12\|h\|.
 \label{f16v4:source-bound}
\end{equation}
Its derivative-index time support is contained in \(\{i-1,i\}\),
where unavailable indices are omitted and \(v_t\) is assigned time
\(t\). The exact Hessian has no \(v_t\)--\(v_s\) block for
\(t\ne s\); a \(u\)--\(v_t\) block can occur only at \(u_t,u_{t+1}\).
Finally, for V1's common-bridge space \(\mathcal E\),
\begin{equation}
 h\in\mathcal E\quad\Longrightarrow\quad
 \|\nabla_uJ_{i,h}\|_\infty\le\delta\|h\|,
              \qquad \|D_{uv}^2V_\gamma\|\le\delta.
 \label{f16v4:common-source}
\end{equation}
Here the second inequality does not require \(h\in\mathcal E\).
\end{proposition}

\begin{proof}
We give the local-to-global estimate so that a volume-dependent Taylor
constant is not being hidden in \(C_*\). After the fixed cube-linear
changes of variables, each action summand is a member of a finite list
of functions
\[
        a\longmapsto s\!\left(\prod_{\ell=1}^{k}
                       \operatorname{Exp}(T_\ell a)\right),
\]
where the product order, the signs and the matrices \(T_\ell\) are the
actual internal, mixed or kinetic words. Their lengths and dimensions
are bounded independently of \(B,n\). The moving frames are themselves
exponentials of fixed linear forms in the five chords of one cube.
Each summand has bounded third derivatives on a fixed neighbourhood of
zero. Its value and gradient vanish at zero, and its quadratic term is
one half of the squared sum of its signed linear arguments. Haar terms
\(-\log j\) have bounded derivatives on that neighbourhood, with value
and gradient zero at zero.

For completeness, constants can be fixed without estimating them
numerically. Choose a radius \(a_0>0\) such that all factors in this
finite word list, after the fixed linear changes, lie in their smooth
principal neighbourhood whenever every input block has norm at most
\(a_0\). Take the maxima of the third derivative norms of the action
words and the first three derivative norms of \(-\log j\) on the
closures of those local input sets. Multiply these finite maxima by
the maximal local input dimension and the maximal number of summands
incident to an input block. Choose \(C_*\) larger than the resulting
Hessian row-sum constant and \((20a_0)^{-1}\), and choose \(C_1\)
larger than the analogous value-sum constant. These constants use only
a finite local library, not a supremum over growing lattices.

Taylor's integral remainder for an action term, rescaled as
\(\gamma f(a/\sqrt\gamma)\), bounds its value remainder by
\(CR^3/\sqrt\gamma\) and its Hessian remainder by
\(CR/\sqrt\gamma\). The Haar contributions are bounded by
\(CR^2/\gamma\) and \(C/\gamma\), respectively, which fit those
bounds under \eqref{f16v4:regime}. The endpoint Gaussian terms are
already exactly quadratic. Every coordinate block meets only a fixed
number of terms: internal spatial and temporal words stay inside a
cube or an adjacent pair of times; a mixed two-bridge plaquette meets
two neighbouring cubes; a four-bridge plaquette has no fast variable;
a bridge kinetic word meets two neighbouring cubes on one time bond.
There are \(O(B(n+1))\) summands altogether. Summing the value bounds
and using the symmetric block row-sum bound for the Hessian proves
the first two estimates. Equation~\eqref{f16v4:H0} gives the third.
The chosen radius puts every individual \(x\) and \(z\) inside its
principal chart.

The Gaussian mixed source derivative has a \(u_i\) block
\(N^{-1/2}\Lambda h\), and \(v_{i-1},v_i\) blocks of the form
\(\pm L_c^{-1/2}E_{\rm br}^{\mathsf T}h\). The inherited bounds
\(\|N^{-1/2}A^{\mathsf T}\|\le2\), \(\|D\|\le4\), and
\(E_{\rm br}L_c^{-1}E_{\rm br}^{\mathsf T}\le29I\) therefore give
\[
 \|\nabla_\xi\partial_{i,h}V_0\|^2
              \le(64+2\cdot29)\|h\|^2=122\|h\|^2.
\]
The mixed Hessian remainder is at most \(\delta\), and
\(\sqrt{122}+1/20<12\), proving \eqref{f16v4:source-bound}.
The same argument and \(\Lambda h=0\) for \(h\in\mathcal E\)
prove the first part of \eqref{f16v4:common-source}; the Gaussian
\(u\)--\(v\) block vanishes, proving its second part.
All support assertions follow from the exact words, before Taylor
expansion. A temporal port variable occurs only on its own bond;
bridge differentiation meets \(u_i,v_{i-1},v_i\), and no other fast
time indices. This completes the certificate.
\end{proof}

The constants \(C_*,C_1\) are defined finite derivative constants, not
numerically optimized values. In particular, the small-field regime is
proved to exist; no practical numerical threshold for \(\gamma\) is
claimed. A bound on \(D_\xi^2V_\gamma\) outside this core has not been
used or established.

\subsection{Boundary-aware temporal response without a length loss}

We reuse f15 V13's absolute one-form covariance realization and the
monograph's Witten router. To fix the domain used here, let
\(\nu=\nu_{\gamma,R,\mathbf y}\), \(H(\xi)=D_\xi^2V_\gamma(\xi)\),
and let \(\mathfrak H\) have the closed quadratic form
\begin{equation}
 \mathfrak h[a]
  =\int|\nabla a|^2\,d\nu+\int a^{\mathsf T}H a\,d\nu
       +\sum_{\text{ball faces }\alpha}
             \int_{\partial_\alpha\mathcal O_R}
                    \mathrm{II}_\alpha(a_\alpha,a_\alpha)\,d\nu_\partial.
 \label{f16v4:witten-form}
\end{equation}
The component \(a_\alpha\) is tangential at its own ball boundary;
\(\mathrm{II}_\alpha\ge0\) is the outward second fundamental form.
This is the product absolute realization, including its corners. In
particular the derivative and boundary part, denoted \(\mathfrak D\),
is nonnegative and block diagonal in the individual coordinate-ball
labels. Multiplication by a constant projection onto any union of these
labels preserves its form domain. It need not be the same scalar
boundary realization for every component.

In this realization the inherited identity is
\begin{equation}
 \operatorname{Cov}_{\nu}(F,G)
       =\langle\nabla F,\mathfrak H^{-1}\nabla G\rangle_{L^2(\nu)}.
 \label{f16v4:covariance-identity}
\end{equation}
One obtains it by solving the mean-zero reflecting scalar equation,
using the scalar/one-form intertwining, and integrating by parts. The
strict lower bound in \eqref{f16v4:derivative-bounds} makes the one-form
inverse bounded. Smooth observables need not satisfy a boundary
condition. The product-ball construction and its nonnegative boundary
form are exactly the inherited boundary-aware mechanism; a primary
reference for the smooth-boundary Witten decomposition is
\href{https://arxiv.org/abs/1607.08714}{D. Le Peutrec,
\emph{On Witten Laplacians and Brascamp--Lieb's inequality on manifolds
with boundary}}. We now exploit the more specific temporal block
structure of the current native law.

Assign both \(u_i\) and \(v_i\), when present, the index-time \(i\).
For a set of such times \(I\), let \(\Pi_I\) project a vector field
onto those coordinate components. This is a projection in
\emph{derivative-index space}; the scalar coefficient functions of a
vector field may depend on every fast variable.

\begin{theorem}[Uniform nonlinear fast-history response]
\label{f16v4:temporal-inverse}
Under \eqref{f16v4:regime}, put \(\vartheta=9/19\). For any index-time
sets \(I,J\) at distance \(d\),
\begin{equation}
          \|\Pi_I\mathfrak H^{-1}\Pi_J\|
                                  \le2\vartheta^d.
 \label{f16v4:temporal-inverse-bound}
\end{equation}
There is also a space--time estimate. Give each \((i,b)\) the two
coordinate species when available, and use nearest spatial-cube and
nearest-time edges. For sets \(A_0,B_0\) of these labels,
\begin{equation}
 \|\Pi_{A_0}\mathfrak H^{-1}\Pi_{B_0}\|
       \le\frac{20}{19}
                 \left(\frac{601}{620}\right)^{d(A_0,B_0)}.
 \label{f16v4:spacetime-inverse-bound}
\end{equation}
Neither estimate has a cardinality, volume or duration prefactor.
Consequently \eqref{f16v4:temporal-inverse-bound} bounds a covariance
by \(2\vartheta^d\|\nabla F\|_2\|\nabla G\|_2\) whenever the two
gradients have the indicated index-time supports.
\end{theorem}

\begin{proof}
Let \(D_*\) be the constant coordinate operator equal to \(p_*I\)
on all \(u\)-components and to \(I\) on all \(v\)-components.
Write
\[
 \mathfrak H=\mathfrak A-\mathfrak C,\qquad
 \mathfrak A=\mathfrak D+\operatorname{diag}_{\rm time}H,\qquad
 \mathfrak C=-(H-\operatorname{diag}_{\rm time}H).
\]
The projection onto time-diagonal matrix blocks contracts operator
norm. Equation~\eqref{f16v4:H0} and the Hessian remainder give
\[
 \mathfrak A\ge(1-\delta)D_*,\qquad
 \|D_*^{-1/2}\mathfrak C D_*^{-1/2}\|
                   \le\frac2{p_*}+2\delta.
\]
For the second assertion, the Gaussian time-off-diagonal operator is
the path adjacency, with norm at most two, on the \(u\) species. The
remainder's off-diagonal part has norm at most \(2\delta\). The
nonnegative derivative and boundary terms preserve time blocks, so
\(\mathfrak A^{-1/2}\) preserves them too. Since \(2/p_*<7/20\),
\[
 \|\mathfrak A^{-1/2}\mathfrak C\mathfrak A^{-1/2}\|
       \le\frac{2/p_*+2\delta}{1-\delta}
       \le\frac9{19}=\vartheta<1.
\]
Bounded form perturbation and the norm-convergent Neumann series give
\[
 \mathfrak H^{-1}=\mathfrak A^{-1/2}
   \sum_{k\ge0}(\mathfrak A^{-1/2}\mathfrak C\mathfrak A^{-1/2})^k
                                                   \mathfrak A^{-1/2}.
\]
Every factor \(\mathfrak C\) changes the time label by at most one.
Terms with \(k<d\) therefore vanish between the stated projections.
Using \(\|\mathfrak A^{-1/2}\|^2\le20/19\), the tail is bounded by
\[
               \frac{20}{19}\frac{\vartheta^d}{1-\vartheta}
                           =2\vartheta^d.
\]
This proves \eqref{f16v4:temporal-inverse-bound} without assuming that
the off-diagonal operator is positive.

For the second estimate use instead
\(\mathfrak A_0=\mathfrak D+31I\) and
\(\mathfrak C_0=31I-H\). The derivative certificate implies
\(0\le\mathfrak C_0\le(601/20)I\), so
\[
 \|\mathfrak A_0^{-1/2}\mathfrak C_0\mathfrak A_0^{-1/2}\|
                                      \le601/620.
\]
The inverse square root is diagonal in all cube--time labels, and
\(\mathfrak C_0\) has range one on the displayed graph. The same
series argument has prefactor
\(31^{-1}/(1-601/620)=20/19\).
The locality assertion is about vector-component labels, not finite
propagation of the scalar diffusion in configuration space. Finally
apply \eqref{f16v4:covariance-identity} and Cauchy--Schwarz.
\end{proof}

Let
\[
 \mathcal A_{\gamma,R}(\mathbf y)=-\log Z_{\gamma,R}(\mathbf y),
 \qquad
 \mathcal W_{\gamma,n}^{R}
  =\frac{c_{\gamma,B,n}J_{{\rm br},\gamma}Z_{\gamma,R}}
            {\sqrt{F_{\gamma,2}(Z_0)F_{\gamma,2}(Z_n)}}.
\]
These definitions retain the full native endpoint normalizers.
The domain \(\mathcal O_R\) is independent of \(\mathbf y\), so
ordinary differentiation under its finite integral gives
\begin{equation}
 \partial_{i,h}\partial_{j,k}\mathcal A_{\gamma,R}
   =\mathbb E_\nu\partial_{i,h}\partial_{j,k}V_\gamma
         -\operatorname{Cov}_\nu(J_{i,h},J_{j,k}).
 \label{f16v4:pressure-hessian}
\end{equation}
There is no moving-boundary derivative in this formula. It is the
inherited exact response identity, not a second return to be subtracted
from an already integrated action.

\begin{proposition}[Two-source temporal memory of the native core]
\label{f16v4:memory}
Under \eqref{f16v4:regime}, for \(|i-j|\ge2\),
\begin{equation}
 |\partial_{i,h}\partial_{j,k}\mathcal A_{\gamma,R}|
       \le288\,\vartheta^{|i-j|-1}\|h\|\|k\|.
 \label{f16v4:generic-memory}
\end{equation}
The same statement holds with
\(\mathcal A_{\gamma,R}\) replaced by
\(-\log\mathcal W_{\gamma,n}^R\).
If \(K_{ij}\) denotes this bridge Hessian block, then for every
integer \(\ell\ge1\), and every sequence of bridge variations,
\begin{equation}
 \left|\sum_{|i-j|>\ell}\langle h_i,K_{ij}h_j\rangle\right|
       \le\frac{576}{1-\vartheta}\vartheta^\ell
                                         \sum_i\|h_i\|^2.
 \label{f16v4:generic-tail}
\end{equation}
This is a Hessian-tail estimate at the specified bridge history,
not an assertion about every higher connected response.
\end{proposition}

\begin{proof}
Every term of \(V_\gamma\) depends on at most two adjacent bridge
times. The direct mixed derivative in
\eqref{f16v4:pressure-hessian} is therefore zero when \(|i-j|\ge2\).
The source gradients lie in time sets \(\{i-1,i\}\) and
\(\{j-1,j\}\), at distance at least \(|i-j|-1\). Apply
Theorem~\ref{f16v4:temporal-inverse} and
\eqref{f16v4:source-bound}; \(2\cdot12^2=288\).
The logarithms of the bridge Jacobian and the two endpoint
normalizers are sums of single-time functions, and the remaining
scalar is history independent. None changes a distinct-time Hessian
block.

For the tail, the scalar majorant for the omitted ordered pairs has
row sum at most
\(2\cdot288\sum_{d>\ell}\vartheta^{d-1}
 =576\vartheta^\ell/(1-\vartheta)\).
The symmetric Schur bound, or \(2ab\le a^2+b^2\) term by term,
gives \eqref{f16v4:generic-tail}.
\end{proof}

\subsection{Distant common-bridge response needs two nonlinear couplings}

The generic estimate does not use \(\Lambda\mathcal E=0\).
That cancellation is not enough by itself to eliminate the temporal
ports: even at Gaussian order a bridge variation differentiates their
kinetic sources. The additional step is to retain their time-diagonal
one-form inverse before taking the \(u\)-Schur complement.

\begin{theorem}[A small common-bridge nonlocal return in the native core]
\label{f16v4:common-memory}
Assume \eqref{f16v4:regime} and set
\[
             K_{\rm c}=1+\frac{240}{19}=\frac{259}{19}.
\]
For \(h,k\in\mathcal E\) and \(|i-j|\ge2\),
\begin{equation}
 |\partial_{i,h}\partial_{j,k}\mathcal A_{\gamma,R}|
  \le2K_{\rm c}^{\,2}\delta^2
             \vartheta^{(|i-j|-2)_+}\|h\|\|k\|.
 \label{f16v4:common-memory-bound}
\end{equation}
The base bridge history need only satisfy its bound in
\eqref{f16v4:regime}; it need not lie in \(\mathcal E\).
The same estimate holds for \(-\log\mathcal W_{\gamma,n}^R\).
In particular the prefactor is \(O(R^2/\gamma)\), not merely order
one. For \(h_i\in\mathcal E\) and \(\ell\ge1\),
\begin{equation}
 \left|\sum_{|i-j|>\ell}\langle h_i,K_{ij}h_j\rangle\right|
  \le\frac{4K_{\rm c}^{\,2}}{1-\vartheta}\delta^2
                     \vartheta^{\ell-1}\sum_i\|h_i\|^2.
 \label{f16v4:common-tail}
\end{equation}
\end{theorem}

\begin{proof}
Decompose the one-form Hilbert space by species, not by scalar
conditional expectation:
\[
       \mathfrak H=
       \begin{pmatrix}\mathfrak H_{uu}&B_{uv}\\
                       B_{uv}^*&\mathfrak H_{vv}\end{pmatrix}.
\]
The derivative and boundary forms do not mix these species because
their balls are separate. Thus \(B_{uv}=D_{uv}^2V_\gamma\) is a
bounded multiplication operator with norm at most \(\delta\).
The exact native words imply that \(\mathfrak H_{vv}\) is block
diagonal in \emph{bond time}. Its inverse preserves those labels and
has norm at most \(20/19\).

Define the positive Schur operator in the form sense by
\[
 S_u=\mathfrak H_{uu}
             -B_{uv}\mathfrak H_{vv}^{-1}B_{uv}^*.
\]
For any vector fields \(g=(g_u,g_v)\), \(g'=(g'_u,g'_v)\), block
completion of the coercive form gives the exact identity
\begin{align}
 \langle g,\mathfrak H^{-1}g'\rangle
   ={}&\langle g_v,\mathfrak H_{vv}^{-1}g'_v\rangle\notag\\
     &+\langle g_u-B_{uv}\mathfrak H_{vv}^{-1}g_v,
            S_u^{-1}(g'_u-B_{uv}\mathfrak H_{vv}^{-1}g'_v)\rangle.
 \label{f16v4:species-schur}
\end{align}
All off-diagonal couplings are bounded, so this follows also by the
usual block triangular factorization on form domains. In particular
\(S_u^{-1}\) is precisely the \(uu\) corner of
\(\mathfrak H^{-1}\), not an independently normalized inverse.

Take \(g=\nabla J_{i,h}\). For common \(h\),
\eqref{f16v4:common-source} and \eqref{f16v4:source-bound} give
\[
 \|g_u-B_{uv}\mathfrak H_{vv}^{-1}g_v\|_2
       \le\left(1+\frac{20}{19}\,12\right)\delta\|h\|
       =K_{\rm c}\delta\|h\|.
\]
Its \(u\)-index support is contained in \(\{i-1,i,i+1\}\).
Indeed \(g_v\) is supported on bonds \(i-1,i\), the inverse
\(\mathfrak H_{vv}^{-1}\) preserves each bond label, and
\(B_{uv}\) only connects bond \(t\) to \(u_t,u_{t+1}\).
The same assertions hold for \(g'\) at \(j\).

For \(|i-j|\ge2\) the first term of
\eqref{f16v4:species-schur} vanishes: its two \(v\)-index supports
are disjoint. The corrected \(u\)-supports are separated by at least
\((|i-j|-2)_+\). Since \(S_u^{-1}\) is a corner of the full inverse,
Theorem~\ref{f16v4:temporal-inverse} bounds the remaining term by the
right side of \eqref{f16v4:common-memory-bound}.
The direct bridge mixed derivative is zero as before, so
\eqref{f16v4:pressure-hessian} proves the assertion. The normalizer
terms remain single-time. Summing the omitted scalar majorant gives
row sum
\[
  4K_{\rm c}^{\,2}\delta^2
       \sum_{d=\ell+1}^{\infty}\vartheta^{d-2}
    =\frac{4K_{\rm c}^{\,2}}{1-\vartheta}
                       \delta^2\vartheta^{\ell-1},
\]
which proves the tail estimate.
\end{proof}

The two small factors have different possible origins but the same
controlled size: a common source can acquire a small direct
\(u\)-component, or its order-one \(v\)-component can reach \(u\)
through the small \(uv\) coupling. To connect two nonadjacent source
times, both ends must pay one such factor. This conclusion uses the
exact nonlinear core; it is not obtained by declaring its covariance
Gaussian. It does not control the sign of a same-time restoring term,
and it is not a statement of physical gaplessness or of a mass gap.

\subsection{A paid pressure-density comparison for the restricted integral}

A covariance estimate does not fix the scalar part of a path integral.
We pay that scalar separately, without claiming a volume-independent
probability that every fast coordinate lies in the core.
Let \(\mathbb P_{0,\mathbf y}\) be the normalized Gaussian law with
potential \(V_0\), in the whitened coordinates \(\xi\), and let
\(\mu(\mathbf y)\) be its mean.

\begin{proposition}[A uniform local mean bound for the Gaussian reference]
\label{f16v4:gaussian-mean}
There is a finite geometric constant \(C_\mu\ge1\), independent of
\(B,n\), such that
\begin{equation}
       \max_\alpha|\mu_\alpha(\mathbf y)|\le C_\mu r
       \quad\hbox{when}\quad \max_{i,e}|y_{i,e}|\le r.
 \label{f16v4:mean-bound}
\end{equation}
Here \(\alpha\) runs over the separate \(u\)- and \(v\)-balls.
Every corresponding marginal Gaussian covariance is at most the
identity.
\end{proposition}

\begin{proof}
The covariance assertion follows from \(H_0\ge I\). The source
\(-D_\xi V_0(0;\mathbf y)\) is a finite-range linear function of the
bridge history and has norm at most \(C r\) on each ball, with a
fixed local constant. To obtain a block row-sum estimate for its
inverse, write
\[
       H_0^{-1}=\frac1{30}\sum_{k\ge0}(I-H_0/30)^k.
\]
The summand has norm at most \((29/30)^k\) and range at most \(k\)
on the cube--time graph. Thus the norm of an inverse block between
labels at distance \(d\) is at most \((29/30)^d\). The number of
labels at distance \(d\) from any fixed label is at most
\(C(d+1)^4\), uniformly on a three-dimensional periodic cube graph
times a finite time path, including the two species. Consequently the
sum of these block norms is at most the finite constant
\(C\sum_{d\ge0}(d+1)^4(29/30)^d\).
Multiplication by the local source bound proves
\eqref{f16v4:mean-bound}. This row-sum argument is needed: the
Euclidean norm bound \(\|H_0^{-1}\|\le1\) alone would give a
spurious volume factor for a uniformly bounded bridge field.
\end{proof}

\begin{theorem}[Core pressure per cube and time slice]
\label{f16v4:pressure}
Assume \eqref{f16v4:regime} and, in addition,
\begin{equation}
                  R\ge\max(12,2C_\mu r).
 \label{f16v4:mean-margin}
\end{equation}
Then
\begin{equation}
 \frac{1}{B(n+1)}
       \left|\log\frac{Z_{\gamma,R}(\mathbf y)}{Z_0(\mathbf y)}\right|
 \le C_1\frac{R^3}{\sqrt\gamma}
                  +4\,2^{21/2}e^{-R^2/16}.
 \label{f16v4:pressure-bound}
\end{equation}
This compares the exact native \emph{core} integral with the complete
Euclidean Gaussian fast integral. It does not compare the complete
compact native integral with that Gaussian integral.
\end{theorem}

\begin{proof}
Let \(a_R=2^{21/2}e^{-R^2/16}\). If \(g\) is a centered Gaussian
of dimension at most 21 and covariance at most \(I\), then
\[
 \mathbb P\{|g|>R/2\}
  \le e^{-R^2/16}\mathbb E e^{|g|^2/4}
  \le2^{21/2}e^{-R^2/16}=a_R.
\]
For \(R\ge12\), \(a_R<1/2\). Each event
\(\{|\xi_\alpha-\mu_\alpha|\le R/2\}\) is a symmetric convex
cylinder for the centered full Gaussian \(\xi-\mu\).
The Gaussian correlation inequality, iterated over these cylinders,
gives
\[
 \mathbb P_{0,\mathbf y}\!
     \left\{\forall\alpha:\ |\xi_\alpha-\mu_\alpha|\le R/2\right\}
                      \ge(1-a_R)^{(2n+1)B}.
\]
This use of the Gaussian correlation inequality is an external theorem,
not an independence assertion; see
\href{https://arxiv.org/abs/1408.1028}{T. Royen,
\emph{A simple proof of the Gaussian correlation conjecture extended to
multivariate gamma distributions}}. The centered event is contained in
\(\mathcal O_R\) by \eqref{f16v4:mean-bound} and
\eqref{f16v4:mean-margin}. Since \(\log(1-a)\ge-2a\) on
\([0,1/2]\),
\begin{equation}
 -4B(n+1)a_R\le
                \log\mathbb P_{0,\mathbf y}(\mathcal O_R)\le0.
 \label{f16v4:gaussian-core-cost}
\end{equation}
There is no lower bound here independent of the number of cells.

Set \(\epsilon_N=C_1B(n+1)R^3/\sqrt\gamma\).
The pointwise action remainder on the core yields
\[
 e^{-\epsilon_N}Z_0\mathbb P_{0,\mathbf y}(\mathcal O_R)
       \le Z_{\gamma,R}
       \le e^{\epsilon_N}Z_0\mathbb P_{0,\mathbf y}(\mathcal O_R).
\]
Take logarithms, use \eqref{f16v4:gaussian-core-cost}, and divide by
\(B(n+1)\). The argument controls the extensive scalar even when the
Gaussian probability of the all-good event is small.
\end{proof}

The remaining explicit scalar factors in the native weight can also be
separated without inventing a spectral normalization. Uniformly in
\(B,n\), for sufficiently large \(\gamma\),
\begin{align}
 \left|\log\frac{c_{\gamma,B,n}}{c_{0,B,n}}\right|
                &\le C Bn/\gamma+CB e^{-c\gamma},\notag\\
 |\log J_{{\rm br},\gamma}(\mathbf y)|
                &\le CBn r^2/\gamma.
 \label{f16v4:explicit-scalar-cost}
\end{align}
Here and below constants may depend on the fixed cube geometry but
not on \(B,n\). To verify the first estimate, rescale the defining
integral for \(z_\gamma\). Taylor expansion of \(1-\cos a\) and
\(j(a)\), with the bound
\(1-\cos a\ge2a^2/\pi^2\) on \([0,\pi]\), gives
\[
 \frac{z_\gamma}{c_{H,\gamma}}
                  =(2\pi)^{3/2}(1+O(\gamma^{-1})).
\]
For a direct domination argument, use Taylor expansion on
\(|q|\le\gamma^{1/8}\); the rescaled remainder is bounded by
\(C\gamma^{-1}(|q|^2+|q|^4)e^{-c|q|^2}\).
The complement is controlled by a Gaussian tail using the displayed
cosine bound. The cube Gaussian chart mass factorizes by cubes and
has \(-\log p_{\gamma,B}\le CB e^{-c\gamma}\), by fixed-dimensional
Gaussian tails for each of the five chord vectors. These facts prove
the first bound of \eqref{f16v4:explicit-scalar-cost};
\(|\log j(a)|\le C|a|^2\) on the small chart proves the second.

In particular the exact compensated comparison is
\begin{align}
 &\log\frac{\mathcal W_{\gamma,n}^{R}(\mathbf y)}
                    {\mathcal W_n(\mathbf y)}
   +\frac12\log\frac{F_{\gamma,2}(Z_0)F_{\gamma,2}(Z_n)}
                       {F_2(y_0)F_2(y_n)}\notag\\
 &\hspace{12mm}
   -\log\frac{c_{\gamma,B,n}}{c_{0,B,n}}
   -\log J_{{\rm br},\gamma}(\mathbf y)
            =\log\frac{Z_{\gamma,R}(\mathbf y)}{Z_0(\mathbf y)}.
 \label{f16v4:compensated-pressure}
\end{align}
The endpoint \(F_{\gamma,2}\) factors are kept literally. Theorem
\ref{f16v4:pressure} does not give a new uniform estimate for their
full compact-to-Gaussian ratios. They do not enter the separated-time
mixed response because their logarithms are single-time functions.

For fixed \(r\) and any fixed \(K>0\), the choice
\(R=K\sqrt{\log\gamma}\) satisfies all the stated conditions once
\(\gamma\) is sufficiently large. It gives a pressure-density error
\[
 O\!\left(\frac{(\log\gamma)^{3/2}}{\sqrt\gamma}
                             +\gamma^{-K^2/16}\right),
\]
while the common-bridge distant-time response has prefactor
\(O(\log\gamma/\gamma)\). These estimates are uniform in \(B,n\)
inside the indicated regime; no total-variation convergence of the
full growing-dimensional native law follows.

\subsection{The exact remaining correction and the next upstream target}

At every finite regulator and admissible bridge history set
\begin{equation}
 p_{\rm good}(\mathbf y)
       =\frac{Z_{\gamma,R}(\mathbf y)}
                        {Z_{\gamma,\mathrm{full}}(\mathbf y)}
       \in(0,1].
 \label{f16v4:good-probability}
\end{equation}
It is the actual probability of the specified joint \(x,z\) core under
the augmented native conditional. Positivity and the exact history
identity give
\begin{equation}
 \mathcal W_{\gamma,n}^{R}
         =p_{\rm good}\mathcal W_{\gamma,n}^{\mathrm{full}},\qquad
 -\log\mathcal W_{\gamma,n}^{\mathrm{full}}
         =-\log\mathcal W_{\gamma,n}^{R}+\log p_{\rm good}.
 \label{f16v4:full-core-split}
\end{equation}
The sign in this equation is important: restricting an integral makes
its negative logarithm larger. For separated bridge times, the full
mixed response is therefore exactly
\begin{equation}
 \partial_{i,h}\partial_{j,k}
            (-\log\mathcal W_{\gamma,n}^{\mathrm{full}})
 =\partial_{i,h}\partial_{j,k}
            (-\log\mathcal W_{\gamma,n}^{R})
          +\partial_{i,h}\partial_{j,k}\log p_{\rm good}.
 \label{f16v4:unpaid-response}
\end{equation}
Smoothness for these bridge derivatives follows at finite regulator
by differentiating the positive compact integral; the core integral
has a fixed domain. No estimate for the last term is being asserted.

This isolates a specific next interface. A useful V5 estimate would
control the \emph{local mixed derivatives} of
\(\log p_{\rm good}\), together with its pressure per cell, by a
native large-field decomposition that retains the temporal-port tails
and the actual bridge sources. It must control this correction rather
than count the same covariance return a second time. A uniform lower
bound on \(p_{\rm good}\) itself is not the appropriate demand:
even the Gaussian all-good probability can decrease exponentially in
\(B(n+1)\) while its pressure density and local response are controlled.
The next argument should therefore test local bad-region weights or
source-visible capacity, not repeat a Gaussian inverse estimate.
The archived pinned-cube compact bounds provide possible ingredients,
but do not by themselves bound this time-history correction.

There are two other boundaries on the claim. First, bounding every
bridge coordinate in an identity chart is not a gauge-invariant
coverage theorem for the physical root quotient. The full native
kernel identity is compatible with restriction to invariant
functions; the present local derivative estimate does not
\emph{automatically} survive that restriction or its boundary-field
integration. Second, \(\mathfrak H\) is an auxiliary conditional
one-form response operator. Its lower bound is not a spectral gap of
\(T_{\rm full,\gamma}\), and it has not been compared to that
transfer's top eigenvalue. Physical-time calibration and an
interacting continuum construction remain separate requirements.

\clearpage
\subsection{Verification, source preservation, and current status}

The newly supplied diagnostic is
\begin{center}\small
\path{verify_ym_f16_v4_native_history.py}.
\end{center}
Its 9,467 exact algebraic assertions passed. They check the cube
orthogonality and spectra; actual parity-cube plaquette and common-mode
incidences on fine tori of sides eight and ten; noncommuting finite
history elimination; endpoint and Haar power counts; rational temporal
contraction bounds; the two-species Schur source identity; and the universal quadratic
Wilson-word identity used in the local native expansion. The report is
\path{ym_v4_verification_report.json}.
The script has 8,442 bytes and SHA--256
\begin{center}\scriptsize\ttfamily
8C22B165F3A623A89F8E5EAB73A8B1F5291FF1575B4773344815847514EE0D06.
\end{center}
These finite fixtures check the indicated algebra and support structure.
They do not evaluate the analytic derivative suprema, independently
verify the Witten domain construction, or certify a physical mass gap.
The analytic arguments above, including their cited inherited tools,
are essential. No external proof-assistant verification was performed.

The supplied V3 source has 68,826 bytes and SHA--256
\begin{center}\scriptsize\ttfamily
05F39BD90C9796D694D109855BE56A508D2F2DF2857238F773B1F3E5B55F505E.
\end{center}
V4 appends this section before its final \verb|\end{document}| and changes
only the title version from V3 to V4. Removing this append and restoring
that title recovers the supplied V3 byte for byte. All four supplied
originals are retained unchanged. The older scripts and companion files
mentioned in the inherited notes were not supplied; their historical
checks are not represented here as newly rerun. Only the three supplied
background sources were available for the present companion audit.

\begin{firewall}
V4 proves duration- and volume-uniform nonlinear two-source memory for
a specified native fast-history \emph{core}, a smaller
\(O(R^2/\gamma)\) distant common-bridge return, and a paid core pressure
density relative to the full Gaussian fast integral. The exact
compact normalization and the remaining \(\log p_{\rm good}\)
correction are retained. No bound on that full-law correction,
gauge-invariant bridge-tail coverage, full-transfer spectral gap,
interacting continuum reconstruction, or positive physical
Yang--Mills mass gap is established by this appendix.
\end{firewall}
% END F16 V4 APPEND
% BEGIN F16 V5 APPEND -- 2026-09-17
% Insert immediately before \end{document} in the cumulative V4.
% All inherited text is preserved; only the cumulative title changes to V5.

\clearpage
\section{V5: complete compact fast pressure by rooted pins and sparse corridors}
\label{f16v5:context}

V4 leaves two different questions in its last term
\(\log p_{\rm good}\): its extensive value and its local source derivatives.
This appendix pays the first without assuming the second. The upstream
change is to undo the moving frame on the \emph{temporal integration
variables only}. The seven tree kinetic terms on a cube then become an
independent compact pin. Along with the internal magnetic pin, they remain
coercive after inter-cube and inter-time interactions are cut. This permits
an upper/lower comparison through finite boxes separated by sparse
corridors, instead of an unproved large-field cluster expansion.

The resulting estimate concerns the \emph{complete compact} fast-history
integral, at an arbitrary bridge history in a fixed rescaled bounded
window. It is uniform in the number of cubes and the history length. It
also holds for bounded one-cell quadratic fast sources. Convexity in
those auxiliary sources then gives an actual-law, spatially and temporally
averaged second-moment comparison. Neither conclusion is a local
bridge-source mixing theorem.

\subsection{The V5 checkpoint and the part that is genuinely new}

The available checkpoint inputs are the supplied f14 V100, f15 V100,
Grand Tensor Monograph V076, and this lineage through V4. Section
inventories and targeted passages were examined; this is not a claim to
have independently certified every historical proof. In particular:
\begin{center}\small
\begin{tabular}{@{}p{0.23\textwidth}p{0.71\textwidth}@{}}
\toprule
Inherited result & What is reused, and what is not inferred\\
\midrule
f15 V97 & The five-chord compact magnetic inequality and the rooted
seven-edge tree estimate are already proved. They are ingredients,
not new V5 theorems.\\[1mm]
f14 V95--V100 & Full Wilson pressure stability and thermal estimates
refer to their specified periodic measure, not to the present
fixed-bridge, dressed-end history conditional.\\[1mm]
f15 V59 & A pressure-value comparison must not be differentiated without
additional control. That limitation still applies to bridge derivatives.\\[1mm]
f15 V69--V71 & Native bad-region estimates there use different conditional
laws and an actual-vacuum exterior-energy budget. No such budget is
imported into this history law.\\[1mm]
Monograph V076 & The direct Hessian minus the returned covariance is
already available. We do not repackage that identity as a new result.\\[1mm]
This lineage V4 & The full history density, restricted-core response,
core pressure comparison and exact definition of \(p_{\rm good}\)
remain unchanged.\\
\bottomrule
\end{tabular}
\end{center}

The new conclusion is convergence of the \emph{full conditional fast
pressure} to its particular, complete Gaussian history pressure, with a
vanishing bound per cube and time slice. The proof does not assume a
uniform probability that every fast variable is small. The sparse-corridor
comparison is a proof device, not a replacement of the native transfer.
The other companion programs named in older checkpoints were not supplied
in this conversation; their historical checks have not been rerun.

For methodological context, partition-function stability via exact tree
coordinates is not new. See P.~A.~Faria da Veiga and M.~O'Carroll,
\href{https://arxiv.org/abs/1903.09829v2}{\emph{On Thermodynamic and
Ultraviolet Stability of Yang--Mills}}, and their
\href{https://arxiv.org/abs/2205.07376v2}{\emph{On Yang--Mills Stability
Bounds and Plaquette Field Generating Function}}. The former treats
\(\mathrm U(N)\) and \(\SU(N)\) with free boundaries; the latter's stated
setting is \(\mathrm U(N)\). We use neither as a black-box theorem for the
present \(\SU(2)\) conditional. The inequalities needed below are proved
for its actual words and normalization.

\subsection{Undoing temporal frames, without changing the compact integral}

Retain \(B=m^3\), \(m\ge4\), \(n\ge1\), and write
\[
 \mathcal N=B(n+1),\qquad
 \max_{i,e}|y_{i,e}|\le r,
 \qquad Z_i=\operatorname{Exp}(y_i/\sqrt\gamma),
\]
where \(r\ge0\) is fixed independently of \(B,n,\gamma\). Constants denoted
\(C_r\) may depend on this window and the fixed cube geometry, but not on
\(B,n\). All principal logarithms below are taken away from their
Haar-null cut sets.

Let \(f_{i,v}=f_v(X_i)\) be the moving port frame from V2--V4, and set
\begin{equation}
 h_{t,v}=f_{t,v}r_{t,v}f_{t+1,v}^{-1},\qquad
 h_{t,o}=I,\qquad
 Y_{i,e}=f_{i,s(e)}Z_{i,e}f_{i,t(e)}^{-1}.
 \label{f16v5:unmove}
\end{equation}
For fixed internal configurations, \(r\mapsto h\) is a product of Haar
translations. The internal kinetic word, up to conjugation, becomes
\[
 h_{t,s(e)}^{-1}X_{t,e}h_{t,t(e)}X_{t+1,e}^{-1};
\]
its bridge counterpart is
\[
 h_{t,s(e)}^{-1}Y_{t,e}h_{t,t(e)}Y_{t+1,e}^{-1}.
\]
For a tree edge, \(X_{t,e}=X_{t+1,e}=I\), so its word is exactly
\(h_{t,s(e)}^{-1}h_{t,t(e)}\). In particular it contains no internal chord
or bridge variable.

Use complete charts
\[
 X_i=\operatorname{Exp}(x_i/\sqrt\gamma),\qquad
 h_t=\operatorname{Exp}(a_t/\sqrt\gamma),\qquad \zeta=(x,a).
\]
A cell \(c=(i,b)\) contains \(x_{i,b}\in\mathbb R^{15}\) and also
\(a_{i,b}\in\mathbb R^{21}\) when \(i<n\). Thus its dimension \(d_c\)
is 36, except that the final-time cells have dimension 15. Set
\(d_{\rm f}=\sum_c d_c=15B(n+1)+21Bn\). The domain
\(\mathcal P'\) is the product of the principal three-dimensional group
charts, not a small-field restriction.

Let \(U_\gamma(\zeta;\mathbf y)\) be V4's classical exponent after
\eqref{f16v5:unmove}: it includes the two endpoint Gaussian quadratic
forms and every native Wilson term, but not the negative logarithms of
the Haar densities. Define
\begin{equation}
 J^f_\gamma(\zeta)=
  \prod_{i,b,j\ {\rm chord}}j(x_{i,b,j}/\sqrt\gamma)^{w_i}
  \prod_{t,b,v\ {\rm nonroot}}j(a_{t,b,v}/\sqrt\gamma),
 \qquad 0\le J^f_\gamma\le1.
 \label{f16v5:haar}
\end{equation}
Here the endpoint weights \(w_i=1/2\), interior weights \(w_i=1\), and
\(j\) have exactly their meanings in V4.

\begin{proposition}[Same full fast integral; different useful coordinates]
\label{f16v5:exact-unmoving}
At every finite regulator,
\begin{equation}
 Z_{\gamma,\mathrm{full}}(\mathbf y)
   =\int_{\mathcal P'}J^f_\gamma(\zeta)
                         e^{-U_\gamma(\zeta;\mathbf y)}\,d\zeta.
 \label{f16v5:full-integral}
\end{equation}
No history-dependent factor multiplies this identity. If \(U_0\) is its
complete quadratic limit, then
\begin{equation}
 Z_0(\mathbf y)=\int_{\mathbb R^{d_{\rm f}}}
                         e^{-U_0(\zeta;\mathbf y)}\,d\zeta
 \label{f16v5:Z0}
\end{equation}
is exactly V4's \(Z_0\), and hence the Gaussian integral in V3's history
formula, with its determinant intact.
\end{proposition}

\begin{proof}
Substitution of \(\bar X_{i,e}=f_{i,s}^{-1}X_{i,e}f_{i,t}\) in V4's
internal word gives its conjugate by \(f_{t+1,s}\) in the displayed
unmoving coordinates. The bridge calculation is identical. Class
functions remove these conjugations. Conditional Haar invariance gives
\(dH(r)=dH(h)\) before any chart is introduced. In complete charts this
replaces \(j(z/\sqrt\gamma)\,dz\) by
\(j(a/\sqrt\gamma)\,da\), with the same constant Haar powers on both
sides. The internal half-density weights and endpoint amplitudes have
not been changed. This proves \eqref{f16v5:full-integral}, including its
normalization. The dependence of \(h\) on the internal variables causes
no extra Jacobian: the conditional Haar changes are performed at each
fixed \(X\).

At quadratic order the coordinate change is
\begin{equation}
             a_t=z_t+\mathsf F_B(x_t-x_{t+1}).
 \label{f16v5:linear-unmoving}
\end{equation}
On the stacked fast space it is block triangular with identity diagonal,
so its determinant is one, also with all three colours. Substitute
\(z_t=a_t-\mathsf F_B(x_t-x_{t+1})\) in V4's \(V_0\).
This gives precisely \(U_0\), and the same Euclidean integral. This is a
change of variables in the full Gaussian integral, not a new choice of
its scalar prefactor.
\end{proof}

V4's event \(\mathcal O_R\) is \emph{not} a product of balls in these new
coordinates. We will not replace it by one. The new coordinates are used
to estimate the full integral; the already defined core is compared to it
only through their common \(Z_0\).

\subsection{Compact pins that survive every corridor cut}

Call the following part of the classical exponent its \emph{anchor}:
\begin{align}
 A_\gamma(\zeta)
 ={}&\frac14\left(x_0^{\mathsf T}C_B^{-1}x_0+
                            x_n^{\mathsf T}C_B^{-1}x_n\right)
      +\gamma\sum_{i,b}w_i S_\Box(X_{i,b})\notag\\
 &+\gamma\sum_{t,b}\sum_{e\in T_\Box}
                      s(h_{t,s(e)}^{-1}h_{t,t(e)}).
 \label{f16v5:anchor}
\end{align}
It is a sum of one-cell functions. Every remaining Wilson term is
nonnegative. Its fast carrier consists of cells in a box of side at most
two: the internal chord kinetic terms meet two adjacent times in one
cube; a mixed two-bridge plaquette meets neighbouring cubes; and a bridge
kinetic word meets at most two neighbouring cubes at two adjacent times.
Terms with no fast variable are retained as a separate scalar
\(b_\gamma(\mathbf y)\). There are at most \(C\mathcal N\) terms in all,
and at most a fixed number incident to any cell. The constants refer to
the literal word list, including its moving frames in \(Y\).

\begin{proposition}[A global independent pin for the augmented history]
\label{f16v5:global-pin}
On every complete chart,
\begin{equation}
 A_\gamma(\zeta)\ge a_*\|\zeta\|^2,
       \qquad a_*:=\frac1{6\pi^2},\qquad \gamma>0.
 \label{f16v5:pin}
\end{equation}
The same inequality holds for its quadratic limit \(A_0\). Deleting any
subset of the other positive terms, or setting any subset of cells to
zero in those other terms while retaining \emph{all} anchors, preserves
this pin. The zero of the anchor is the identity in every fast group.
\end{proposition}

\begin{proof}
The already proved f15 V97 cube inequality is
\(\sum_{j=1}^5s(X_j)\le3S_\Box(X)\).
For a principal angle \(\theta\in[0,\pi]\),
\(1-\cos\theta\ge2\theta^2/\pi^2\). Since \(w_i\ge1/2\),
\[
 \gamma\sum_{i,b}w_i S_\Box(X_{i,b})
                         \ge\frac1{3\pi^2}\sum_{i,b}|x_{i,b}|^2.
\]
The seven rooted tree paths have total depth 12. By the bi-invariant
quaternion chord triangle inequality and Cauchy--Schwarz along a path,
\[
 s(h_v)\le d_v\sum_{e\ {\rm on}\ o\to v}
                      s(h_{s(e)}^{-1}h_{t(e)}),\qquad
 \sum_v s(h_v)\le12\sum_{e\in T_\Box}
                      s(h_{s(e)}^{-1}h_{t(e)}).
\]
Consequently the tree terms in \eqref{f16v5:anchor} are at least
\(\sum_{t,b}|a_{t,b}|^2/(6\pi^2)\). Add the nonnegative endpoint
quadratics. This proves the pin and the zero assertion. Taking the
quadratic limit proves the Gaussian inequality on every finite
Euclidean vector. All the allowed changes leave these one-cell anchors
in place and leave the other terms nonnegative.
\end{proof}

This use of the inherited tree bound is the essential upstream step.
In the moving \(z\)-variables the tree kinetic terms appear to depend on
both internal time ends. In the full compact integral that apparent
coupling can be removed before estimating the normalization. No global
convexity of the nonlinear Wilson action is asserted.

We shall also allow an auxiliary one-cell quadratic source
\begin{equation}
 \mathcal Q_{\mathsf Q}(\zeta)=\sum_c\zeta_c^{\mathsf T}Q_c\zeta_c,
 \quad Q_c=Q_c^{\mathsf T},\quad\|Q_c\|\le1,
 \quad |t|\le t_*:=a_*/4.
 \label{f16v5:quadratic-source}
\end{equation}
It is included as \(-t\mathcal Q_{\mathsf Q}\) in the exponent. The
anchor then still exceeds \((a_*-|t|)\|\zeta\|^2\). Such a source is
kept in the one-cell anchor during every cut. It is an auxiliary moment
generating parameter; the native history law itself is always \(t=0\).
The source changes neither the integration measure nor the fast event.

\subsection{A complete finite-box comparison, with the dimension paid}

The next estimate applies to a block of \(M\ge1\) cells carrying their
full anchors, any retained subset of the native non-anchor terms, and
terms obtained by setting exterior fast cells to zero. Allow also
\eqref{f16v5:quadratic-source}. The domain of every integrated group is
its \emph{complete} principal chart. The block may include the first or
last history slice. Its number of real variables is at most \(36M\).
Fast-free terms can be split off as scalar factors; their number is
bounded by a fixed multiple of the cells or boundary cells to which they
are assigned.

Write its integral as \(\mathfrak Z_\gamma\), with Haar product as in
\eqref{f16v5:haar}, and let \(\mathfrak Z_0\) be the entire Euclidean
integral of its quadratic limit. All constants below are uniform in the
choice of retained terms, the source matrices, and \(|t|\le t_*\).

\begin{theorem}[Unrestricted finite-box Laplace comparison]
\label{f16v5:box-comparison}
There are constants \(C_r,K_r\ge1\) and \(\gamma_r\) such that, whenever
\(\gamma\ge\gamma_r\) and
\begin{equation}
                   K_r(M+\log\gamma)\le\gamma,
 \label{f16v5:box-size}
\end{equation}
one has
\begin{equation}
 \left|\log\frac{\mathfrak Z_\gamma}{\mathfrak Z_0}\right|
 \le \frac{C_r}{\sqrt\gamma}(M+\log\gamma)^{3/2}
                                      +C_r\gamma^{-4}.
 \label{f16v5:box-error}
\end{equation}
The integrated nonlinear variables are not conditioned to a small box.
The dimensional cost on the right is explicit.
\end{theorem}

\begin{proof}
Both classical exponents obey
\(U_\gamma,U_0\ge(a_*/2)|\zeta|^2\), with the quadratic source
included. We use this weaker pin for convenience. At \(|\zeta_c|\le1\) in each cell, the Gaussian exponent is
at most \(C_rM\). Indeed every word is from a fixed finite list, there
are \(O(M)\) terms, their local inputs are bounded, and their incidence
is uniformly bounded. The product of these unit balls has volume at
least \(e^{-CM}\). Thus
\begin{equation}
                  \mathfrak Z_0\ge e^{-C_rM}.
 \label{f16v5:box-lower}
\end{equation}
This product-ball lower bound is important; the volume of a unit ball
in the \emph{total} dimension would introduce an unnecessary
\(M\log M\) cost.

For a Wilson word \(f\), Taylor's integral formula for
\(\gamma f(w/\sqrt\gamma)\) gives a cubic error bounded by
\(C\gamma^{-1/2}|w|^3\). The constant can be taken uniform along the
whole real line segment of the input: differentiated exponentials have
bounds given by their fixed linear coefficients, and the remaining
matrix factors are unitary. Each factor here is an exponential of a
fixed linear form in \(x,a,\mathbf y\); no derivative of a newly chosen
logarithm of a product is used. The quadratic term is the squared
linearized word with its correct coefficient. Bounded local dimension
and bounded incidence give the global-on-the-block estimate
\begin{equation}
 |U_\gamma-U_0|
     \le\frac C{\sqrt\gamma}
               \left(\sum_c|\zeta_c|^3+Mr^3\right)
     \le\frac C{\sqrt\gamma}(|\zeta|^3+Mr^3).
 \label{f16v5:cubic-sum}
\end{equation}
Exactly quadratic endpoint and source terms cancel from this difference.
Frozen zero inputs and omitted terms do not worsen it. On
\(|\zeta|\le T\), \(T^2\le c\gamma\), the full Haar product obeys
\[
        0\le-\log J^f_\gamma(\zeta)\le C T^2/\gamma,
\]
by \(-\log j(v)\le C|v|^2\) in a fixed small chart.

Choose
\(T^2=K'_r(M+\log\gamma)\), with \(K'_r\) large enough. From the
pin, the estimate \(J^f_\gamma\le1\), and an elementary Gaussian
Chernoff bound, each full tail divided by \(\mathfrak Z_0\) is at most
\begin{align*}
 e^{C_rM}\int_{|\zeta|>T}e^{-a_*|\zeta|^2/2}\,d\zeta
 &\le\exp\{C'_rM-a_*T^2/4\}\\
 &\le e^{-M}\gamma^{-4}.
\end{align*}
This bounds the nonlinear tail on its principal domain and the Gaussian
tail on its entire Euclidean domain. Constants such as
\((4\pi/a_*)^{d/2}\), \(d\le36M\), were included in \(C'_rM\).
Enlarge \(K_r\) in \eqref{f16v5:box-size} so that the ball of radius
\(T\) lies inside every principal chart and inside the smaller chart
used for the Haar logarithm estimate.

On that total-radius ball the two integrands have logarithmic ratio at
most
\[
 \epsilon_M=\frac C{\sqrt\gamma}(T^3+Mr^3)+CT^2/\gamma
          \le\frac{C_r}{\sqrt\gamma}(M+\log\gamma)^{3/2}
\]
in absolute value. Writing \(\tau=e^{-M}\gamma^{-4}\), we obtain
\[
 e^{-\epsilon_M}(1-\tau)
       \le\frac{\mathfrak Z_\gamma}{\mathfrak Z_0}
       \le e^{\epsilon_M}+\tau.
\]
For \(\gamma\ge2\), taking logarithms costs at most
\(\epsilon_M+2\tau\). No assumption that \(\epsilon_M\) is small is
needed. This proves \eqref{f16v5:box-error}. The ball was used only to
compare integrands; the complement of the native integral was estimated,
not removed from its definition.
\end{proof}

Fast-free scalar terms satisfy the corresponding elementary bound
\(C_r\gamma^{-1/2}\) per term by the same Taylor argument. All finite-box
comparisons below retain these scalars.

\subsection{Uniform Gaussian boundary costs, not native mixing assumptions}

We need a Gaussian comparison only \emph{between two auxiliary cuts},
not between a cut spectrum and the full native spectrum. Let a corridor
be any set \(\mathcal C\) of cells. Define two modifications of the
classical exponent, always retaining every anchor and auxiliary source:
\[
 U_\gamma^{\mathrm N}=A_\gamma-t\mathcal Q_{\mathsf Q}
       +b_\gamma+\sum_{e:\,\operatorname{car}(e)\cap\mathcal C=\varnothing}
                                      V_{\gamma,e},
\]
and
\[
 U_\gamma^{\mathrm D}(\zeta)=A_\gamma(\zeta)
       -t\mathcal Q_{\mathsf Q}(\zeta)+b_\gamma
       +\sum_e V_{\gamma,e}(\zeta_{\mathcal C^c},0_{\mathcal C};\mathbf y).
\]
The symbols \(\mathrm N,\mathrm D\) denote, respectively, deletion of
terms touching the corridor and freezing its inputs \emph{in those
terms}. They are not claims about a physical boundary condition.
The corridor anchors are still integrated in the second expression.
Denote the corresponding Gaussian integrals by
\(Z_0^{\mathrm N}(t,\mathsf Q)\), \(Z_0^{\mathrm D}(t,\mathsf Q)\).

\begin{proposition}[Gaussian cost proportional to the changed cells]
\label{f16v5:gaussian-cut}
For \(|t|\le t_*\), \(\|Q_c\|\le1\), and \(\max|y_{i,e}|\le r\),
\begin{equation}
 \left|\log\frac{Z_0^{\mathrm N}(t,\mathsf Q)}{Z_0(t,\mathsf Q)}\right|
 +\left|\log\frac{Z_0^{\mathrm D}(t,\mathsf Q)}{Z_0(t,\mathsf Q)}\right|
                              \le C_r|\mathcal C|.
 \label{f16v5:gaussian-cut-bound}
\end{equation}
The same constant works for all corridor shapes, spatial volumes and
history lengths.
\end{proposition}

\begin{proof}
Every relevant interpolating Gaussian exponent retains the independent
anchor minus the source. All other terms are positive multiples of
squares of fixed finite-range linear forms plus bridge sources. The
interpolation for freezing is the convex combination of an original
square and the square with its corridor inputs zero. There are constants
\(h_-,h_+>0\), independent of all regulators and of the interpolation,
for which the full fast Hessian satisfies
\[
                          h_-I\le H\le h_+I.
\]
The lower bound follows from the anchor; the upper bound is the finite
word-incidence bound. Its fast linear source has norm at most \(C r\)
in each cell. All Hessians have range at most one in the cell graph whose
edges allow a unit change in each space--time coordinate.

Here is a local mean estimate that avoids a volume factor. With
\(h_+\) enlarged if necessary,
\[
 H^{-1}=h_+^{-1}\sum_{k\ge0}(I-H/h_+)^k,\qquad
 \|I-H/h_+\|\le q:=1-h_-/h_+<1.
\]
An inverse block between cells at graph distance \(d\) has norm at most
\(h_-^{-1}q^d\). A radius-\(d\) ball in the three-dimensional spatial
torus times the finite time path has at most \(C(d+1)^4\) cells,
uniformly in its side lengths; deleting cells or edges cannot create
shorter paths than in this ambient graph. Hence every inverse block row
has summable norm, bounded by
\(C h_-^{-1}\sum_{d\ge0}(d+1)^4q^d\). Multiplying by the local source
bound proves \(|\mu_c|\le C_r\) for every conditional Gaussian mean
cell. Every marginal covariance is at most \(h_-^{-1}I\).
It follows that the expectation of any one of the squared local words
is at most \(C_r\), including words with frozen zero inputs.

There are at most \(C|\mathcal C|\) changed terms. Differentiate the
logarithm of the Gaussian integral along the deletion interpolation:
its derivative is minus the expectation of the changed quadratic
potential. Integrating its absolute value costs at most
\(C_r|\mathcal C|\). Along the freezing interpolation, the changed
term is the difference of two nonnegative squares. Its absolute
expectation is bounded by the sum of their expectations, with the same
bound. Fast-free constants obey \(C_r\) per changed term as well.
All interpolated Gaussian integrals are finite and differentiable by the
uniform pin. This proves both estimates.
\end{proof}

These estimates use only the gapped \emph{fast Gaussian pin}. They do
not require a gap in the retained bridge quadratic form or a Witten
estimate for the full compact conditional.

\subsection{Sparse corridors and the complete native pressure}

Let \(\ell\ge8\) be an integer. On each spatial cycle and on the finite
time path, use no cut when its length is at most \(\ell\). Otherwise
partition it into consecutive intervals of length at most \(\ell\) and
at least \(\ell/2-1\), and mark the first two positions of each interval.
Let \(\mathcal C\) be the union of the resulting coordinate corridors.
Then
\begin{equation}
 |\mathcal C|\le\frac{16\mathcal N}{\ell},\qquad
     |\mathcal D_j|\le\ell^4
 \label{f16v5:corridors}
\end{equation}
for every connected interior component \(\mathcal D_j\). The first
constant is a convenient safe bound. Marking an end of the physical
time path is allowed and only strengthens the separation. Short
periodic directions remain intact and do not increase component size.
The width two ensures that a native local term cannot meet two different
interior components. Terms do not have to be pairwise interactions.

To see the count explicitly, a coordinate of length \(L>\ell\) can be
partitioned into \(k=\lceil L/\ell\rceil\) intervals of lengths
\(\lfloor L/k\rfloor\) or \(\lceil L/k\rceil\). Its marked fraction
is at most \(2k/L\le4/\ell\). Sum over the four coordinates. The
complement consists of products of the remaining coordinate intervals;
no range-one carrier crosses a width-two marked interval.

Denote the full and modified integrals, with the auxiliary source, by
\[
 Z_\gamma(t,\mathsf Q),\qquad
 Z_\gamma^{\mathrm N}(t,\mathsf Q),\qquad
 Z_\gamma^{\mathrm D}(t,\mathsf Q).
\]
Both modified integrals factor into complete compact integrals on
\(\mathcal D_j\), independent one-cell anchor integrals on
\(\mathcal C\), and their literal fast-free scalar factors.

\begin{theorem}[Full compact fast-pressure comparison, uniform in duration]
\label{f16v5:full-pressure}
Fix \(r<\infty\). There are \(C_r,K_r,\gamma_r\) such that, for all
\(B=m^3\), \(m\ge4\), \(n\ge1\), \(\gamma\ge\gamma_r\),
\(\max|y_{i,e}|\le r\), \(|t|\le t_*\), and symmetric
\(\|Q_c\|\le1\), the following holds. For every \(\ell\ge8\) with
\(K_r(\ell^4+\log\gamma)\le\gamma\),
\begin{equation}
 \frac1{\mathcal N}
    \left|\log\frac{Z_\gamma(t,\mathsf Q;\mathbf y)}
                       {Z_0(t,\mathsf Q;\mathbf y)}\right|
 \le C_r\left\{
       \frac{1+\log\gamma}{\ell}
       +\frac{\ell^2+(\log\gamma)^{3/2}}{\sqrt\gamma}
          \right\}.
 \label{f16v5:pressure-scale}
\end{equation}
In particular, for sufficiently large \(\gamma\),
\begin{equation}
 \frac1{\mathcal N}
    \left|\log\frac{Z_\gamma(t,\mathsf Q;\mathbf y)}
                       {Z_0(t,\mathsf Q;\mathbf y)}\right|
 \le C_r\varepsilon_\gamma,
 \qquad
 \varepsilon_\gamma=
       \gamma^{-1/6}(1+\log\gamma)^{2/3}.
 \label{f16v5:pressure-rate}
\end{equation}
At \(t=0\), \(Z_\gamma\) is exactly V4's
\(Z_{\gamma,\mathrm{full}}\), not its restricted core.
\end{theorem}

\begin{proof}
The upper native bound is immediate from positivity of every deleted
term:
\begin{equation}
                  Z_\gamma(t,\mathsf Q)
                            \le Z_\gamma^{\mathrm N}(t,\mathsf Q).
 \label{f16v5:native-upper}
\end{equation}
The source is retained in the anchors, so it does not affect this
ordering.

For the lower bound restrict only the corridor cells to
\(|\zeta_c|\le\eta\), \(\eta=\gamma^{-1/2}\). This small event is a
temporary lower integration bound, not the definition of the full law.
For a native Wilson term, differentiation with respect to a rescaled
fast coordinate has norm at most \(C\sqrt\gamma\) throughout the
complete principal charts. This follows by differentiating the product
of unitary exponentials; the fixed moving-frame matrices have bounded
linear coefficients. There is no need to bound any interior cell.
Changing all of a term's corridor inputs from zero to size at most
\(\eta\) therefore changes its action by at most \(C\). Only
\(C|\mathcal C|\) terms are affected. The one-cell anchor quadratics
and auxiliary source cost at most \(C\eta^2\) per corridor cell, and
its Haar density is bounded below by \(e^{-C\eta^2/\gamma}\).
Thus integration over these corridor balls gives
\begin{equation}
 Z_\gamma(t,\mathsf Q)
       \ge e^{-C(1+\log\gamma)|\mathcal C|}
                         Z_\gamma^{\mathrm D}(t,\mathsf Q).
 \label{f16v5:native-lower}
\end{equation}
Here are the normalization details. The ball in a corridor cell has
volume \(v_{d_c}\eta^{d_c}\), with \(d_c\le36\), whose negative
logarithm is at most \(C+18\log\gamma\). At frozen corridor inputs,
the remaining integral is the product of the interior Dirichlet
integrals and the fast-free factors. To turn it into
\(Z_\gamma^{\mathrm D}\), multiply by the complete independent
corridor anchor integrals. Each of these has logarithm bounded in
absolute value by a constant: the global pin gives its upper bound,
and a fixed unit ball with its positive small-chart Haar density gives
its lower bound. Its cell dimension is fixed. Dividing by these anchor
integrals only changes the constant in \eqref{f16v5:native-lower}.
All constants are uniform in \(t,\mathsf Q\). No coupling-dependent
factor proportional to \(\gamma|\mathcal C|\) is hidden here.

Apply Theorem~\ref{f16v5:box-comparison} to every factor of either
modified integral. The block sizes \(M_j\) sum to \(\mathcal N\),
with each at most \(\ell^4\); corridor anchors count as blocks of size
one. The number of nonempty blocks is at most \(\mathcal N\). Hence
\begin{align*}
 \sum_j(M_j+\log\gamma)^{3/2}
 &\le C\sum_j\{M_j^{3/2}+(\log\gamma)^{3/2}\}\\
 &\le C\mathcal N\{\ell^2+(\log\gamma)^{3/2}\}.
\end{align*}
Fast-free factors contribute at most \(C_r\mathcal N/\sqrt\gamma\)
to the comparison, also after freezing. The \(\gamma^{-4}\) errors
are absorbed by the same displayed budget. Consequently, for either
modification,
\begin{equation}
 |\log Z_\gamma^{\mathrm N,\mathrm D}
                   -\log Z_0^{\mathrm N,\mathrm D}|
 \le \frac{C_r\mathcal N}{\sqrt\gamma}
                         \{\ell^2+(\log\gamma)^{3/2}\}.
 \label{f16v5:sum-box-error}
\end{equation}
This compares \emph{complete} box integrals, not Gaussian integrals
with an unestimated nonlinear complement.

Combine \eqref{f16v5:native-upper}, \eqref{f16v5:native-lower},
\eqref{f16v5:sum-box-error} and
Proposition~\ref{f16v5:gaussian-cut}. Finally use
\eqref{f16v5:corridors} and divide by \(\mathcal N\).
This proves \eqref{f16v5:pressure-scale}. The argument also works when
there are no corridors because all four directions are shorter than
\(\ell\): there is then a single block of size at most \(\ell^4\).

Choose
\(\ell=\lceil\gamma^{1/6}(1+\log\gamma)^{1/3}\rceil\).
For large \(\gamma\), the finite-box condition holds since
\(\ell^4=O(\gamma^{2/3}(1+\log\gamma)^{4/3})=o(\gamma)\).
The first two dominant errors both have size
\(O(\gamma^{-1/6}(1+\log\gamma)^{2/3})\); the remaining logarithmic
term divided by \(\sqrt\gamma\) is smaller. This proves the rate.
The exponent is a conservative balance of the two proved costs, not a
claim of an optimal asymptotic expansion.
\end{proof}

The proof has not assumed translation invariance of the bridge history,
reflection positivity of its conditional law, convexity of the compact
potential, or an a priori small probability of bad cells. The corridors
are removed from the statement by the two-sided pressure comparison.
Their widths and the shrinking integration radius do not alter the
physical spatial or temporal regulator.

\subsection{The previously unpaid all-good pressure is now bounded}

Return to \(t=0\). Keep \(p_{\rm good}\) and \(Z_{\gamma,R}\) exactly
as defined in V4, in its moving, whitened fast variables. The following
uses the equality of the \emph{full} Gaussian integrals in
Proposition~\ref{f16v5:exact-unmoving}; it does not identify two different
small-field events.

\begin{theorem}[A vanishing full-law cost per cube and time slice]
\label{f16v5:good-pressure}
Under the hypotheses of V4's core pressure theorem and, in addition,
\(\gamma\ge\gamma_r\),
\begin{equation}
 0\le-\frac{\log p_{\rm good}(\mathbf y)}{\mathcal N}
 \le C_r\varepsilon_\gamma
       +C_1\frac{R^3}{\sqrt\gamma}
       +4\,2^{21/2}e^{-R^2/16}.
 \label{f16v5:good-pressure-bound}
\end{equation}
In particular, with \(R=4\sqrt{\log\gamma}\) and fixed \(r\), all the
required core conditions hold eventually and
\begin{equation}
 e^{-C_r\mathcal N\varepsilon_\gamma}
          \le p_{\rm good}(\mathbf y)\le1.
 \label{f16v5:good-exp}
\end{equation}
This is a full-native conditional probability, with the original
endpoint amplitudes and no changed fast normalizer.
\end{theorem}

\begin{proof}
By the exact definition,
\[
 -\log p_{\rm good}
     =\log\frac{Z_{\gamma,\mathrm{full}}}{Z_0}
                      -\log\frac{Z_{\gamma,R}}{Z_0}.
\]
Use \eqref{f16v5:pressure-rate} at \(t=0\) and
Theorem~\ref{f16v4:pressure}. Nonnegativity is exact because the core
is a restriction of the full integral. For the specified \(R\), the
additional errors are \(O((\log\gamma)^{3/2}/\sqrt\gamma)\) and
\(O(\gamma^{-1})\), both smaller than \(\varepsilon_\gamma\).
The fixed mean-margin and chart conditions in V4 hold for sufficiently
large \(\gamma\). This proves \eqref{f16v5:good-exp}.
\end{proof}

There is also an exact entropy interpretation. With
\(\nu_{\gamma,\mathrm{full},\mathbf y}\) denoting V4's augmented full
native conditional in the original moving variables,
\[
 \nu_{\gamma,R,\mathbf y}
    =\nu_{\gamma,\mathrm{full},\mathbf y}(\,\cdot\mid\mathcal O_R),
 \qquad
 D(\nu_{\gamma,R,\mathbf y}\Vert
                  \nu_{\gamma,\mathrm{full},\mathbf y})=-\log p_{\rm good}.
\]
This follows immediately from the conditional Radon--Nikodym derivative
\(1_{\mathcal O_R}/p_{\rm good}\). Thus its relative entropy per
cube and time slice vanishes at the proved rate. For each finite volume
their total-variation distance, in the convention \(\sup_A|\mu(A)-\nu(A)|\),
is exactly \(1-p_{\rm good}\). Consequently it tends to zero along
sequences with \(\mathcal N\varepsilon_\gamma\to0\). It is not
asserted to tend to zero uniformly in unrestricted growing volumes.

The scalar statement for the exact history weight retains all the
external normalizers as well:
\begin{align}
 &\log\frac{\mathcal W_{\gamma,n}^{\mathrm{full}}(\mathbf y)}
                   {\mathcal W_n(\mathbf y)}
  +\frac12\log\frac{F_{\gamma,2}(Z_0)F_{\gamma,2}(Z_n)}
                       {F_2(y_0)F_2(y_n)}\notag\\
 &\hspace{7mm}-\log\frac{c_{\gamma,B,n}}{c_{0,B,n}}
                 -\log J_{{\rm br},\gamma}(\mathbf y)
       =\log\frac{Z_{\gamma,\mathrm{full}}(\mathbf y)}{Z_0(\mathbf y)}.
 \label{f16v5:full-compensated}
\end{align}
The right side divided by \(\mathcal N\) is now bounded by
\(C_r\varepsilon_\gamma\). Equation~\eqref{f16v5:full-compensated}
is not permission to drop the exact \(F_{\gamma,2}\) endpoint factors.

\subsection{A justified consequence for actual full-law quadratic moments}

An arbitrary bridge derivative of a small pressure error is not small
by inference alone. The auxiliary sources
\eqref{f16v5:quadratic-source} provide a different, justified derivative:
the log partition function is convex in their \emph{linear coefficient}
\(t\), and the Gaussian comparison has a uniform second-derivative
bound. This yields the following averaged native statement.

Let \(\widetilde\nu_{\gamma,\mathbf y}\) be the probability law of
\(\zeta=(x,a)\) obtained by normalizing
\eqref{f16v5:full-integral}, and let \(\widetilde\nu_{0,\mathbf y}\)
be its full Gaussian reference. In a cell \(c\), define raw second-moment
matrices
\[
 M_{\gamma,c}=\mathbb E_{\widetilde\nu_\gamma}
                               \zeta_c\zeta_c^{\mathsf T},\qquad
 M_{0,c}=\mathbb E_{\widetilde\nu_0}
                               \zeta_c\zeta_c^{\mathsf T}.
\]
The tilde distinguishes these unmoving coordinates from V4's moving
whitened law. No mean is being set to zero at a nonzero bridge history.

\begin{theorem}[Full-native, averaged one-cell second moments]
\label{f16v5:second-moments}
For all sufficiently large \(\gamma\), uniformly in \(B,n\) and
\(\max|y_{i,e}|\le r\),
\begin{equation}
 \frac1{\mathcal N}\sum_c\|M_{\gamma,c}-M_{0,c}\|_1
           \le C_r\sqrt{\varepsilon_\gamma}
           =C_r\gamma^{-1/12}(1+\log\gamma)^{1/3}.
 \label{f16v5:moment-rate}
\end{equation}
Here \(\|\cdot\|_1\) is the finite-matrix trace norm. In particular
\begin{equation}
        \frac1{\mathcal N}
             \mathbb E_{\widetilde\nu_\gamma}\|\zeta\|^2\le C_r.
 \label{f16v5:actual-energy-density}
\end{equation}
This controls an average of complete one-cell moment matrices, not
every distant-time covariance block.
\end{theorem}

\begin{proof}
For fixed \(\mathsf Q=(Q_c)_c\), set
\[
 f_\gamma(t)=\mathcal N^{-1}\log Z_\gamma(t,\mathsf Q),\qquad
 f_0(t)=\mathcal N^{-1}\log Z_0(t,\mathsf Q).
\]
Theorem~\ref{f16v5:full-pressure} gives
\(\sup_{|t|\le t_*}|f_\gamma(t)-f_0(t)|
 \le\delta_\gamma:=C_r\varepsilon_\gamma\), with a constant independent
of \(\mathsf Q\). These are real, differentiable convex functions of
\(t\). Native differentiability holds at every finite regulator because
all principal fast coordinates have bounded norm there. Gaussian
integrability is uniform over this interval by the anchor.

If \(C_t,\mu_t\) are the Gaussian covariance and mean on the whole
fast space, and \(Q=\bigoplus_cQ_c\), direct Gaussian integration gives
\begin{align*}
 \operatorname{Var}_{0,t}(\mathcal Q_{\mathsf Q})
 &=2\Tr\bigl[(C_t^{1/2}QC_t^{1/2})^2\bigr]
                         +4\mu_t^{\mathsf T}QC_tQ\mu_t\\
 &\le C_r\mathcal N.
\end{align*}
For clarity, expand \(\zeta=\mu_t+g\): the centered quadratic and the
linear term in the centered Gaussian \(g\) have zero covariance by
oddness, their variances are the two displayed terms, and the constant
part does not contribute. The estimate uses
\(\|C_t\|\le C\), \(d_{\rm f}\le36\mathcal N\), and
\(\sum_c|\mu_{t,c}|^2\le C_r\mathcal N\), proved by the same uniform
Gaussian block-row argument as Proposition~\ref{f16v5:gaussian-cut}.
Thus \(0\le f_0''(t)\le C_r\) on the whole interval, including for
indefinite \(Q_c\).

For \(0<h\le t_*\), convexity traps \(f_\gamma'(0)\) between the
backward and forward secants. Comparing those secants with \(f_0\),
then using its second-derivative bound, yields
\[
 |f_\gamma'(0)-f_0'(0)|\le2\delta_\gamma/h+C_rh/2.
\]
Take \(h=\sqrt{\delta_\gamma}\) for large \(\gamma\); constants can
be absorbed into \(C_r\). Since differentiation of the partition
function gives
\[
 f_\gamma'(0)-f_0'(0)
       =\mathcal N^{-1}\sum_c\Tr\{Q_c(M_{\gamma,c}-M_{0,c})\},
\]
this expression has absolute value at most
\(C_r\sqrt{\varepsilon_\gamma}\), uniformly over all allowed \(Q_c\).
Choose in each cell the spectral sign of the symmetric matrix
\(M_{\gamma,c}-M_{0,c}\), with sign zero on its kernel. This is an
allowed norm-one symmetric matrix, and its trace pairing is the trace
norm. The uniformity in \(\mathsf Q\) permits this choice after the
moments are defined. This proves \eqref{f16v5:moment-rate}. The Gaussian
one-cell mean/covariance bounds and \(Q_c=I\) give
\eqref{f16v5:actual-energy-density}.
\end{proof}

There is a useful, weaker density consequence for V4's \emph{original}
fast labels. Let
\[
 N_R=\sum_{i,b}1_{\{|u_{i,b}|\ge R\}}
              +\sum_{t,b}1_{\{|v_{t,b}|\ge R\}},
 \qquad u=N^{1/2}x,\quad v=L_c^{1/2}z.
\]
For the full native conditional and any \(R\ge1\),
\begin{equation}
                   \frac{\mathbb E N_R}{\mathcal N}\le\frac{C_r}{R^2}.
 \label{f16v5:bad-density}
\end{equation}
Indeed, from \(r_t=f_t^{-1}h_tf_{t+1}\) and the bi-invariant geodesic
triangle inequality,
\[
 |z_{t,v}|\le|a_{t,v}|+|(\mathsf F_Bx_t)_v|
                              +|(\mathsf F_Bx_{t+1})_v|.
\]
This is a global inequality for principal angles, not a linearization
of \eqref{f16v5:linear-unmoving}. Each internal slice is used at most
twice; the cube matrices have fixed norms. Therefore
\[
 \sum_{i,b}|u_{i,b}|^2+\sum_{t,b}|v_{t,b}|^2
                                      \le C\|\zeta\|^2.
\]
Apply \(1_{\{|w|\ge R\}}\le|w|^2/R^2\) and
\eqref{f16v5:actual-energy-density}. With \(R=4\sqrt{\log\gamma}\),
the expected bad-label density is \(O_r((\log\gamma)^{-1})\).
There is no inferred exponential probability for every specified bad
set, nor a percolation or conditional-capacity conclusion.

\subsection{The remaining interface is now source-marked, not scalar}

The exact identity from V4 still reads, at nonadjacent bridge times,
\begin{align*}
 \partial_{i,h}\partial_{j,k}
       (-\log\mathcal W_{\gamma,n}^{\mathrm{full}})
  ={}&\partial_{i,h}\partial_{j,k}
       (-\log\mathcal W_{\gamma,n}^{R})\\
    &+\partial_{i,h}\partial_{j,k}\log p_{\rm good}.
\end{align*}
The value of \(-\log p_{\rm good}\) per cell is no longer unpaid.
Its \emph{local bridge mixed derivatives} are still unpaid. A small
pressure density does not permit their differentiation, and the
convex-source argument in Theorem~\ref{f16v5:second-moments} concerns
one-cell quadratic fast observables, not these nonlinear bridge writers.
Likewise, a small expected density of bad labels does not by itself
bound a source-visible connected bad region.

There is a specific constructive input for V6. All compact interior
boxes already have the full normalization estimate above; it is not
necessary to derive it again or introduce a new scalar adjustment.
A marked comparison must now preserve cancellations between the full
and core integrals near a chosen bridge-source carrier while controlling
how the sparse corridors or bad regions communicate with that carrier.
Applying the unmarked global bound to a single source would lose its
\(\mathcal N\) factor. The needed improvement is locality of the
\emph{normalizer's source response}, not another complete Gaussian
inverse or another global count.

The retained bridge window is still a coordinate window before the
remaining root Gauss restriction, not gauge-invariant coverage of the
physical bridge law. The full identity is the correct native one, but
the bound is conditional on these bridge histories. No comparison of a
full-transfer spectral gap with its top eigenvalue has been established.
Physical-time calibration and the interacting continuum construction
remain distinct requirements.

\subsection{Verification and exact source preservation}

The new diagnostic is
\begin{center}\small
\path{verify_ym_f16_v5_full_fast_pressure.py}.
\end{center}
It tests exact compact word identities, cube and tree pins, and the
determinant-one change of history coordinates. It also tests corridor
separation and counts, positive quadratic cuts, Gaussian source
derivatives, and the pressure and moment rates. All 89,339 assertions passed; most enumerate corridor support and separation
on finite grids. The separate available V4 rerun passed its 9,467
assertions. The V5 report is
\path{ym_v5_verification_report.json}. These are finite algebraic and
combinatorial diagnostics, not proof-assistant verification of the
analytic theorems. In particular they do not numerically optimize the
local derivative constants, replace the finite-box tail proof, establish
bridge-source mixing, or certify a mass gap.
The V5 checker has 14,365 bytes and SHA--256
\begin{center}\scriptsize\ttfamily
56942128B285E46AD12743BAEE41413BEA6E2B893BEEBEF4B7C74F2F9A786918.
\end{center}

The supplied V4 source has 111,673 bytes and SHA--256
\begin{center}\scriptsize\ttfamily
004C1B37DA71B3C6BF4BA25DFCDAFAF7B45B6CF744FEA8BA5A0C3A92BB5986D5.
\end{center}
This cumulative V5 changes only that source's title version and inserts
the present appendix immediately before its final
\verb|\end{document}|. Removing the appendix and restoring the title
recovers the supplied V4 byte for byte. The four original background
uploads and both V4 sources are retained unchanged. The available V4
diagnostic was extracted from the supplied package and rerun; older
scripts named in the archives were not supplied and have not been
represented as rerun. Build, diagnostic and hash results are provided
with the V5 package. This V5 checkpoint uses only the available companion
sources; the next scheduled checkpoint is V10.

\begin{firewall}
V5 proves a vanishing pressure-density error for the complete compact
fast-history conditional, with its Gaussian determinant and exact
native normalization retained. It therefore bounds the previously
unpaid value of \(-\log p_{\rm good}\) per cube and time slice.
Uniform auxiliary quadratic-source pressure also yields convergence of
averaged one-cell raw second-moment matrices under the full native law,
and a vanishing expected bad-label density at growing core radius.
No local bridge-source derivative bound for \(\log p_{\rm good}\),
full-law temporal mixing theorem, gauge-invariant bridge-tail coverage,
full-transfer spectral gap, interacting continuum construction, or
positive physical Yang--Mills mass gap is proved here.
\end{firewall}
% END F16 V5 APPEND
% BEGIN F16 V6 APPEND -- 2026-09-17
% Insert immediately before the final \end{document} in cumulative V5.
% No inherited proof is rewritten. The cumulative title alone changes to V6.

\clearpage
\section{V6: compensated likelihoods and source-marked full-law response}
\label{f16v6:context}

V5 pays the full fast pressure but correctly forbids differentiating its
error to obtain a local bridge susceptibility. We go upstream to a stronger
comparison of the \emph{normalized full laws}. A small positive power of the
native-to-Gaussian likelihood is integrable with a vanishing logarithmic
cost per cell. The negative Gaussian action in this likelihood tilt is
paid by an explicit redistribution of V5's compact anchors. It is not
silently treated as a positive Wilson interaction.

A Gaussian marginal-entropy inequality then distributes that cost over
bounded-overlap families of local source carriers. Together with global
bounds on the actual Wilson first derivatives, this yields a new
\emph{source-marked} conclusion: the mean absolute error of separated-time
bridge Hessian entries tends to zero, uniformly in spatial volume and
history length, for every fixed-overlap bank of localized source pairs.
For common-bridge sources the Gaussian entry vanishes exactly. Spatial
translation symmetry at the zero bridge history additionally gives an
individual-entry estimate uniform in spatial volume, with its explicit
history-length cost retained.

These conclusions do not assert exponentially decaying native memory.
The small error below is independent of separation, so it cannot be summed
over arbitrarily many lags. Nor does a bank average become a pointwise
bound at an arbitrary inhomogeneous bridge history.

\subsection{Audit, inherited tools, and the precise change of output}

The supplied V5 appendix was read in full. The available f14/f15 archives
and Grand Tensor Monograph V076 were inventoried and searched, with
relevant passages read rather than their entire proof collections
recertified. In particular, f15 V2 already converts controlled likelihood
moments to entropy; f15 V23 already has entropy budgets for other
conditional cores and an obstruction to inferring a functional inequality
from small entropy; f15 V37 separates baseline closeness from a new
extensive physical tilt; and f15 V59 prohibits differentiation of an
unmarked pressure error. None is claimed anew. The monograph's
first-source/Hessian identity is also inherited and charged only once.

The comparison here has a different pair of laws: the complete compact
\emph{fixed-bridge fast-history} law of V5 and its complete Euclidean
Gaussian reference, in the same unmoving fast coordinates. The new
analytic input is the compensated likelihood estimate for this pair.
The marginal-entropy step is a standard Gaussian entropy tool; a direct
proof of the exact version used below is included. Its application to the
actual nonlinear bridge sources, with their tails and centering retained,
is the marked advance. No general literature-novelty claim is made for
entropy subadditivity or the source-covariance identity.

\subsection{The same full native law and a compensated likelihood tilt}

Use V5's unmoving coordinates \(\zeta=(x,a)\), principal domain
\(\mathcal P'\), cell dimensions at most 36, and exact factors
\(J^f_\gamma,U_\gamma\). Abbreviate
\[
 \mathcal N=B(n+1),\qquad
 \epsilon_\gamma=\gamma^{-1/6}(1+\log\gamma)^{2/3},\qquad
 \omega_\gamma=\sqrt{\epsilon_\gamma},\qquad
 \max_{i,e}|y_{i,e}|\le r<\infty.
\]
All sufficiently-large-\(\gamma\) thresholds may depend on the fixed
window \(r\), but not on \(B=m^3\), \(m\ge4\), or \(n\ge1\).
Let the following be probability measures on the \emph{same} Euclidean
fast space:
\begin{equation}
 \begin{split}
  d\mu_\gamma(\zeta)
   &=Z_\gamma^{-1}\,1_{\mathcal P'}J^f_\gamma(\zeta)
                              e^{-U_\gamma(\zeta;\mathbf y)}\,d\zeta,\\
  d\mu_0(\zeta)&=Z_0^{-1}e^{-U_0(\zeta;\mathbf y)}\,d\zeta,
  \qquad Z_\gamma=Z_{\gamma,\mathrm{full}}(\mathbf y).
 \end{split}
 \label{f16v6:laws}
\end{equation}
The native law is extended by zero outside its complete charts. This does
not restrict it further. V5 proves that \(Z_\gamma\) and \(Z_0\) are
exactly V4's full integrals; in particular the Gaussian determinant is
unchanged. The external \(F_{\gamma,2}\) and bridge half-density factors
cancel in a normalized conditional on the fixed bridge history, not in
the full bridge history weight.

The Gaussian precision \(H=D_\zeta^2U_0\), covariance \(C=H^{-1}\), and
mean \(m\) obey
\begin{equation}
 h_- I\le H\le h_+I,\qquad
 \sup_c|m_c|\le C_r,
 \label{f16v6:gaussian-data}
\end{equation}
with fixed positive geometric constants \(h_-,h_+\). These are V5's
anchor, finite-incidence, and Gaussian inverse-row estimates in the
unmoving coordinates, not a new coercivity claim for the full native
law. Write \(\kappa=h_+/h_-\).

For \(\theta\ge0\), define the \emph{unnormalized} likelihood integral
\begin{equation}
 \mathcal Z_\gamma(\theta)
   =\int_{\mathcal P'}(J^f_\gamma)^{1+\theta}
       \exp\{-[(1+\theta)U_\gamma-\theta U_0]\}\,d\zeta.
 \label{f16v6:likelihood-integral}
\end{equation}
The sign of \(-\theta U_0\) is essential. V5's positive-term deletion
argument does not apply to this expression before the following
compensation.

\begin{proposition}[Positive local compensation for the likelihood integral]
\label{f16v6:compensation}
There is \(\theta_*>0\), depending only on the fixed local geometry, such
that uniformly for \(0\le\theta\le\theta_*\),
\begin{equation}
 \frac1{\mathcal N}
       |\log\mathcal Z_\gamma(\theta)-\log Z_0|
                           \le C_r\epsilon_\gamma.
 \label{f16v6:tilted-pressure}
\end{equation}
The complete native integration domain and the power
\((J^f_\gamma)^{1+\theta}\) are retained. The quadratic reference of this
integral is exactly \(U_0\), independent of \(\theta\).
\end{proposition}

\begin{proof}
Write V5's decomposition, with its fast-free scalar separated, as
\[
 U_\gamma=A_\gamma+\sum_{e\in\mathcal L}V_{\gamma,e}+b_\gamma,
 \qquad
 U_0=A_0+\sum_{e\in\mathcal L}V_{0,e}+b_0.
\]
Here \(A_\gamma\) is the sum of all one-cell anchors, every
\(V_{\gamma,e}\ge0\), and \(V_{0,e}\) is a nonnegative local square.
The index \(e\) in this proof labels a remaining interaction, not
necessarily one spatial link. Its nonempty fast carrier is \(S_e\).
There are fixed constants \(L,\Delta,a^+\) such that
\begin{equation}
 \begin{gathered}
  A_\gamma\ge a_*\|\zeta\|^2,\qquad
  a_*\|\zeta\|^2\le A_0\le a^+\|\zeta\|^2,\qquad
  a_*=(6\pi^2)^{-1},\\
  0\le V_{0,e}(\zeta;\mathbf y)
     \le Q_e(\zeta)+b_e(r),\qquad
  Q_e=L\sum_{c\in S_e}|\zeta_c|^2,\qquad b_e(r)\le Lr^2,\\
  \#\{e:c\in S_e\}\le\Delta.
 \end{gathered}
 \label{f16v6:compensation-data}
\end{equation}
Enlarge \(L\) to include the fixed number of bridge coordinates in a
local word. All these bounds are uniform over interior and endpoint
cell types. They follow respectively from V5's global pin, its
one-cell quadratic anchors, and the finite list of linearized words.

Set
\begin{align}
 A^\theta_\gamma
   &=(1+\theta)A_\gamma-\theta A_0-\theta\sum_e Q_e,\notag\\
 V^\theta_{\gamma,e}
   &=(1+\theta)V_{\gamma,e}
                     +\theta(Q_e+b_e-V_{0,e}),\notag\\
 b^\theta_\gamma
   &=(1+\theta)b_\gamma-\theta b_0-\theta\sum_e b_e.
 \label{f16v6:compensated-terms}
\end{align}
Then the following is an exact identity, including its scalar:
\begin{equation}
 (1+\theta)U_\gamma-\theta U_0
       =A^\theta_\gamma+\sum_eV^\theta_{\gamma,e}+b^\theta_\gamma.
 \label{f16v6:exact-compensation}
\end{equation}
Each new interaction is nonnegative by
\eqref{f16v6:compensation-data}. Its fast support is unchanged.
The compensated anchor is still a sum of one-cell functions. For example,
choose
\[
 \theta_*=
  \min\left(\frac12,\frac{a_*}{4(a^++L\Delta+a_*)}\right)>0.
\]
Since \(\sum_eQ_e\le L\Delta\|\zeta\|^2\), this choice gives
\begin{equation}
 A^\theta_\gamma\ge\tfrac34a_*\|\zeta\|^2,
 \qquad
 A^\theta_0=A_0-\theta\sum_e Q_e
                         \ge\tfrac34a_*\|\zeta\|^2.
 \label{f16v6:surviving-pin}
\end{equation}
The quadratic limits of the other compensated terms are
\[
 V^\theta_{0,e}=V_{0,e}+\theta(Q_e+b_e),\qquad
 b^\theta_0=b_0-\theta\sum_e b_e.
\]
They sum with \(A^\theta_0\) to \(U_0\). In particular the Gaussian
reference has not been altered to force agreement of partition functions.
The scalar \(b^\theta_\gamma\) need not be positive; it is split off,
retained exactly, and has a comparison error at most
\(C_r\mathcal N/\sqrt\gamma\).

We verify why V5's finite-box and corridor proof now applies, rather
than invoking positivity for a signed interaction. On any complete box
of \(M\) cells, after scalar factors are separated, the surviving
anchors give a fixed Gaussian lower pin in both exponents. Their
quadratic integrals have the product-unit-ball lower bound
\(e^{-C_rM}\). The termwise comparison remainders are
\[
 A^\theta_\gamma-A^\theta_0
      =(1+\theta)(A_\gamma-A_0),\qquad
 V^\theta_{\gamma,e}-V^\theta_{0,e}
      =(1+\theta)(V_{\gamma,e}-V_{0,e}).
\]
Thus V5's global cubic remainder and complete-box tail estimate have
uniform constants for \(\theta\in[0,\theta_*]\). The logarithm of the
small-chart Haar correction is merely multiplied by \(1+\theta\),
and globally \(0\le(J^f_\gamma)^{1+\theta}\le1\).
The resulting complete-box logarithmic error is still
\[
 C_r\gamma^{-1/2}(M+\log\gamma)^{3/2}+C_r\gamma^{-4}
\]
when the same constant-enlarged box-size condition holds. This includes
the entire compact complement, not just a small-field comparison.

For a corridor, deletion is now legal because all
\(V^\theta_{\gamma,e}\ge0\); all compensated anchors remain.
The lower integration bound freezes corridor inputs in these terms and
integrates those cells over radius \(\gamma^{-1/2}\), exactly as in V5.
The added and subtracted quadratics have derivatives bounded by
\(C\sqrt\gamma\) throughout the principal charts, as do the rescaled
Wilson terms. The change in a touched interaction therefore costs
\(C_r\) on that small corridor ball. The pin and a fixed unit-ball lower
bound give uniform upper and lower bounds on every independent
compensated-anchor integral. The resulting logarithmic corridor debit
is at most \(C_r(1+\log\gamma)|\mathcal C|\).

For the two Gaussian cuts, the compensated anchors retain their lower
Hessian floor and every other term is a nonnegative quadratic with
uniformly bounded finite-range Hessian. The Gaussian interpolation
argument of Proposition~\ref{f16v5:gaussian-cut} therefore costs
\(C_r|\mathcal C|\), with the same local mean bound. The extra constants
\(b_e\) are part of the literal changed terms and incur only this local
cost. No native mixing inequality has entered.

Consequently the existing corridor summation gives, for every admissible
\(\ell\),
\[
 \mathcal N^{-1}|\log\mathcal Z_\gamma(\theta)-\log Z_0|
 \le C_r\left\{\frac{1+\log\gamma}{\ell}
        +\frac{\ell^2+(\log\gamma)^{3/2}}{\sqrt\gamma}\right\}.
\]
Use V5's choice
\(\ell=\lceil\gamma^{1/6}(1+\log\gamma)^{1/3}\rceil\).
This proves \eqref{f16v6:tilted-pressure}. The work specific to this
extension is the exact compensated positive representation and the
checks above; the corridor construction itself is not a new result.
\end{proof}

\subsection{A native-to-Gaussian entropy density, in the useful direction}

For probabilities \(P\ll Q\), use
\[
 D(P\Vert Q)=\int\log\frac{dP}{dQ}\,dP,
 \qquad
 D_{1+\theta}(P\Vert Q)
   =\frac1\theta\log\int\left(\frac{dP}{dQ}\right)^{1+\theta}dQ
 \quad(\theta>0).
\]
These definitions refer to normalized probabilities.

\begin{theorem}[Complete native likelihood and relative entropy]
\label{f16v6:entropy}
For some fixed \(\theta_*>0\) and all sufficiently large \(\gamma\),
\begin{equation}
 0\le D(\mu_\gamma\Vert\mu_0)
      \le D_{1+\theta_*}(\mu_\gamma\Vert\mu_0)
      \le C_r\mathcal N\epsilon_\gamma.
 \label{f16v6:entropy-density}
\end{equation}
This is the full native conditional relative to the full Gaussian
conditional. It is neither a core entropy nor an assumption of a
uniformly small total entropy.
\end{theorem}

\begin{proof}
The density ratio is zero outside \(\mathcal P'\). Within the charts,
direct substitution of \eqref{f16v6:laws} gives the exact formula
\begin{equation}
 \int\left(\frac{d\mu_\gamma}{d\mu_0}\right)^{1+\theta}d\mu_0
       =\frac{Z_0^\theta\mathcal Z_\gamma(\theta)}{Z_\gamma^{1+\theta}}.
 \label{f16v6:renyi-identity}
\end{equation}
Its logarithm equals
\[
 [\log\mathcal Z_\gamma(\theta)-\log Z_0]
       -(1+\theta)[\log Z_\gamma-\log Z_0].
\]
Apply Proposition~\ref{f16v6:compensation} at the fixed
\(\theta=\theta_*\), and V5 at \(\theta=0\). Dividing by
\(\theta_*\) gives the last bound. Jensen under \(\mu_\gamma\),
applied to \(\exp(\theta_*\log(d\mu_\gamma/d\mu_0))\), gives the
relative-entropy bound; Jensen under \(\mu_0\) gives nonnegativity
of the likelihood logarithm. All these integrals are finite by the
compensated estimate. This is an elementary likelihood implication,
already used in f15 V2, applied to the newly controlled pair of laws.
\end{proof}

The orientation matters. At any finite \(\gamma\), the Gaussian gives
positive mass outside the principal native charts, whereas \(\mu_\gamma\)
gives that set mass zero. Hence
\(D(\mu_0\Vert\mu_\gamma)=+\infty\). A two-sided full entropy theorem
has not been asserted. V5's separate identity
\(D(\nu_{\gamma,R}\Vert\nu_{\gamma,\mathrm{full}})=-\log p_{\rm good}\)
compares different laws and has not been counted again.

\subsection{Distributing entropy over bounded-overlap marginal carriers}

For a set \(S\) of cells, let \(P_S\) be the coordinate projection onto
all real coordinates in those cells, and denote marginal probabilities
by a subscript \(S\). The diameter of \(S\) need not be bounded: a union
of two distant local source supports is allowed.

\begin{proposition}[Gaussian marginal-entropy localization; standard tool]
\label{f16v6:entropy-localization}
Let \(Q\) be any nondegenerate Gaussian with covariance \(C\), and let
\(P\ll Q\) have finite relative entropy. Suppose a finite family
\((S_\alpha)_{\alpha\in\mathcal A}\) of nonempty cell sets satisfies
\begin{equation}
        \sup_c\#\{\alpha:c\in S_\alpha\}\le q.
 \label{f16v6:overlap}
\end{equation}
Then, with \(\kappa_C=\lambda_{\max}(C)/\lambda_{\min}(C)\),
\begin{equation}
       \sum_\alpha D(P_{S_\alpha}\Vert Q_{S_\alpha})
                          \le q\kappa_C D(P\Vert Q).
 \label{f16v6:marginal-entropy}
\end{equation}
No independence among the Gaussian cell coordinates or native
coordinates is assumed.
\end{proposition}

\begin{proof}
Translate and whiten the full Gaussian. In standard Gaussian
coordinates, the standardized marginal map is
\[
 B_\alpha=C_{S_\alpha S_\alpha}^{-1/2}P_{S_\alpha}C^{1/2},
 \qquad B_\alpha B_\alpha^{\mathsf T}=I.
\]
Thus \(\Pi_\alpha=B_\alpha^{\mathsf T}B_\alpha\) is an orthogonal
projection. The overlap bound gives
\begin{align*}
 \sum_\alpha\Pi_\alpha
  &=C^{1/2}\left(\sum_\alpha P_{S_\alpha}^{\mathsf T}
              C_{S_\alpha S_\alpha}^{-1}P_{S_\alpha}\right)C^{1/2}\\
  &\le\lambda_{\min}(C)^{-1}C^{1/2}(qI)C^{1/2}
   \le q\kappa_C I.
\end{align*}
Let \(f\) be the density of the whitened \(P\) relative to standard
Gaussian measure \(g\), and put \(f_t=\mathsf P_t f\), where
\(\mathsf P_t\) is the Ornstein--Uhlenbeck semigroup with generator
\(\Delta-x\cdot\nabla\). Conditional expectation onto \(B_\alpha x\)
commutes with this semigroup. For a smooth positive density, the
marginal gradient is the conditional expectation of
\(B_\alpha\nabla f_t\). Conditional Cauchy--Schwarz therefore gives
\[
 I((f_t)_\alpha)
    \le\int\frac{|B_\alpha\nabla f_t|^2}{f_t}\,dg,
 \qquad
 \sum_\alpha I((f_t)_\alpha)\le q\kappa_C I(f_t),
\]
where \(I(f)=\int|\nabla f|^2/f\,dg\) is relative Gaussian Fisher
information. Integration by parts along the same semigroup yields
\[
 \operatorname{Ent}_g(f)=\int_0^\infty I(f_t)\,dt.
\]
Apply the identity to the full density and to every marginal, and
integrate the Fisher inequality. This proves
\eqref{f16v6:marginal-entropy}. For completeness, bounded densities suffice first. Smooth them for a
positive time, integrate the dissipation identity over \([s,T]\), and
let \(T\to\infty\): the Mehler formula gives convergence to the
constant density one, and boundedness gives entropy convergence by
dominated convergence. As \(s\downarrow0\), entropy lower
semicontinuity and semigroup Jensen imply convergence to the original
entropy. For a general finite-entropy density, use normalized
truncations \(\min(f,M)/\int\min(f,M)\,dg\). Their full entropies
converge to that of \(f\), and marginal entropy lower semicontinuity
passes the proved inequality to the limit. This also justifies the
stated dissipation use without any native Fisher-information assumption
at time zero. At each fixed regulator the full native likelihood is
in fact bounded on its compact support.
\end{proof}

This is the Gaussian relative-entropy version of the Fisher-information
projection method behind entropy Brascamp--Lieb inequalities. See
E.~A.~Carlen and D.~Cordero--Erausquin,
\href{https://arxiv.org/abs/0710.0870}{\emph{Subadditivity of the entropy
and its relation to Brascamp--Lieb type inequalities}}, especially the
Fisher-information and heat-flow argument in Section 3. The proof above
specifies the correlated Gaussian marginal constant used here. It is not
an approximate tensorization of \emph{native conditional} entropies and
not a native logarithmic Sobolev inequality.

For the laws \eqref{f16v6:laws}, write
\(d_\alpha=D((\mu_\gamma)_{S_\alpha}\Vert(\mu_0)_{S_\alpha})\).
The preceding results give
\begin{equation}
        \sum_\alpha d_\alpha\le C_r q\mathcal N\epsilon_\gamma.
 \label{f16v6:bank-entropy}
\end{equation}
If \(|S_\alpha|\le s\) for fixed \(s\), every reference marginal has
bounded dimension, uniformly bounded mean, and covariance bounds from
\eqref{f16v6:gaussian-data}, independently of the diameter of
\(S_\alpha\). This last fact is what permits distant source pairs.
The entropy budget is on marginals of the full native law, not on a
collection of separately normalized cavity laws.

\subsection{Unbounded local scores: the tail cost is included}

A total-variation comparison of small marginals would not suffice for
unbounded thermal scores. The following finite-dimensional estimate
makes the additional control explicit.

\begin{proposition}[Gaussian entropy controls affine scores and native remainders]
\label{f16v6:score-entropy}
Let \(Q\) be a Gaussian in dimension at most \(d_*\), with bounded mean
and fixed positive upper and lower covariance bounds. Let \(P\ll Q\)
have \(d=D(P\Vert Q)<\infty\). Suppose \(L_1,L_2\) are affine
functions with bounded coefficients, and \(R_1,R_2\) satisfy, for
\(j=1,2\),
\begin{equation}
 |R_j(z)|\le K(1+|z|),\qquad
 |R_j(z)|\le\frac K{\sqrt\gamma}(1+|z|^2),\qquad\gamma\ge1.
 \label{f16v6:two-remainder-bounds}
\end{equation}
Then, with constants depending only on the stated fixed bounds,
\begin{align}
 \mathbb E_P(1+|z|^2)&\le C(1+d),\notag\\
 \mathbb E_PR_j^2&\le C(d+\gamma^{-1}),\notag\\
 \big|\operatorname{Cov}_P(L_1+R_1,L_2+R_2)
                 -\operatorname{Cov}_Q(L_1,L_2)\big|
       &\le C\big(\sqrt d+d+\gamma^{-1/2}\big).
 \label{f16v6:score-entropy-bound}
\end{align}
In particular no native fourth-moment assumption is needed.
\end{proposition}

\begin{proof}
We use only the entropy variational inequality
\begin{equation}
 t\mathbb E_PF\le d+\log\mathbb E_Qe^{tF},\qquad t>0,
 \label{f16v6:entropy-variational}
\end{equation}
first for bounded truncations and then by monotone or dominated limits
when applicable. It follows from nonnegativity of relative entropy
against the normalized exponential tilt of \(Q\). Gaussian exponential
integrability of \(1+|z|^2\) at a fixed positive \(t\) proves the first
bound.

Choose \(t_0>0\) sufficiently small that the Gaussian exponential
moment below is uniformly finite. Both inequalities in
\eqref{f16v6:two-remainder-bounds} are used:
\begin{align*}
 \mathbb E_Q(e^{t_0R_j^2}-1)
  &\le t_0\mathbb E_Q[R_j^2e^{t_0R_j^2}]\\
  &\le\frac C\gamma\,
       \mathbb E_Q[(1+|z|^4)e^{C t_0(1+|z|^2)}]
   \le C\gamma^{-1}.
\end{align*}
Thus \(\log\mathbb E_Qe^{t_0R_j^2}\le C\gamma^{-1}\).
Equation~\eqref{f16v6:entropy-variational} proves the second bound,
including the contribution from arbitrarily large native coordinates.
Using only the quartic bound in \eqref{f16v6:two-remainder-bounds}
would not have justified this exponential integral.

For an affine \(L\), its centered Gaussian logarithmic moment generating
function is exactly a quadratic in the source. The entropy variational
inequality with both signs therefore gives
\[
     |\mathbb E_PL-\mathbb E_QL|\le C\sqrt d.
\]
For \(F=L_1L_2\), Gaussian exponential integrability of a quadratic
implies
\[
 \log\mathbb E_Q e^{t(F-\mathbb E_QF)}\le Ct^2
                          \quad (|t|\le t_1)
\]
for a fixed \(t_1>0\). One can verify this without a determinant formula:
use \(|F|\le C(1+|z|^2)\) and
\(e^u\le1+u+\tfrac12u^2e^{|u|}\); the linear expectation vanishes.
Optimizing \(d/|t|+C|t|\) on \((0,t_1]\) gives
\[
      |\mathbb E_PF-\mathbb E_QF|\le C(\sqrt d+d).
\]
Combining this product estimate with the two affine mean estimates,
including their product of size \(Cd\), proves
\[
  |\operatorname{Cov}_P(L_1,L_2)-\operatorname{Cov}_Q(L_1,L_2)|
                                      \le C(\sqrt d+d).
\]
Finally each cross covariance involving one \(R_j\) is bounded by
\(C\sqrt{(d+\gamma^{-1})(1+d)}\), using Cauchy--Schwarz and the
first two estimates. The covariance of the two remainders is bounded
by \(C(d+\gamma^{-1})\). Since
\[
 \sqrt{(d+\gamma^{-1})(1+d)}
                 \le C(\sqrt d+d+\gamma^{-1/2}),
\]
these bounds prove the last line of
\eqref{f16v6:score-entropy-bound}. All covariances are centered under
the indicated law; no conditional mean has been discarded.
\end{proof}

\subsection{The actual native bridge writers satisfy these bounds}

A localized bridge mark is \(\alpha=(i,h)\), where \(h\) is a real
bridge-coordinate direction at time \(i\), \(\|h\|\le1\), supported
on a fixed bounded number of bridges. It may carry all three colours.
Define, in V5's unmoving coordinates held fixed,
\[
   J_{\gamma,\alpha}(\zeta)=\partial_{i,h}U_\gamma(\zeta;\mathbf y),
   \qquad
   J_{0,\alpha}(\zeta)=\partial_{i,h}U_0(\zeta;\mathbf y).
\]
Neither \(J^f_\gamma\) nor the unmoving coordinate map depends on
\(\mathbf y\). These are therefore exactly the score insertions for
the fixed-domain fast integral, with no coordinate-derivative term
missing. A bridge-only part of a writer is a deterministic constant;
retaining or subtracting it does not affect its covariance.

\begin{proposition}[Complete-chart score certificate]
\label{f16v6:native-score}
For such a mark there is a finite fast cell carrier \(S_\alpha\), whose
cardinality is bounded by the marked bridge support and the local
geometry, on which both writers depend. The Gaussian writer is affine,
with uniformly bounded coefficients. On the complete principal charts,
\begin{equation}
 \begin{split}
  |J_{\gamma,\alpha}|+|J_{0,\alpha}|
                    &\le C_r(1+|\zeta_{S_\alpha}|),\\
  |J_{\gamma,\alpha}-J_{0,\alpha}|
                    &\le C_r\gamma^{-1/2}(1+|\zeta_{S_\alpha}|^2).
 \end{split}
 \label{f16v6:native-score-bounds}
\end{equation}
These same bounds hold for the analytic exponential-word expressions
extended to all Euclidean fast inputs. The carrier occupies at most the
times \(i-1,i,i+1\), omitting missing endpoint times. The constants are
independent of spatial volume, history length and location.
\end{proposition}

\begin{proof}
Every term involving the marked bridge variables is one of a fixed
finite list of rescaled functions
\[
    \gamma s\!\left(\prod_{j=1}^k
                       \operatorname{Exp}(T_j w/\sqrt\gamma)\right),
\]
where \(w\) contains only its fast variables and nearby bridge inputs.
The matrices \(T_j\), product order, and number of incident terms are
fixed. This is also true for the unmoving bridge words: the port frames
are exponentials of fixed linear forms in the adjoining chord inputs.
There is no principal logarithm of an intermediate product to
differentiate.

Duhamel's formula for derivatives of a matrix exponential shows that
its real skew-Hermitian linear arguments have globally bounded first,
second, and third derivatives with the corresponding fixed coefficient
norms. Indeed the derivative is an integral of products of unitary
matrices and the inserted coefficient matrices; at higher orders the
simplex volumes cancel the permutation count. Applying the product
rule to the finite word gives a uniform bound on its rescaled Hessian
and a third-derivative bound \(C/\sqrt\gamma\), throughout all real
inputs. Its gradient at zero vanishes. Integrating the Hessian from
zero bounds a first bridge derivative by \(C|w|\). Integrating the
third derivative in the Taylor expansion of that first derivative
bounds its difference from the quadratic-limit first derivative by
\(C|w|^2/\sqrt\gamma\). The endpoint Gaussian anchor is independent
of bridge variables and contributes nothing here.

Summing the fixed number of terms incident to the marked support and
using the bound \(r\) on nearby bridge inputs proves
\eqref{f16v6:native-score-bounds}. The bound on the difference by
\(C_r(1+|\zeta_{S_\alpha}|)\) also follows by adding the two first
bounds, so the remainder satisfies both hypotheses of
Proposition~\ref{f16v6:score-entropy}. Finally, a spatial mixed word
uses one bridge time, and a kinetic word uses only two adjacent bridge
times. In unmoving coordinates its frames involve the internal
coordinates at those same ends. This proves the stated fast carrier
and its duration-independent size.
\end{proof}

Let
\[
 \mathcal A_\gamma=-\log Z_{\gamma,\mathrm{full}},\qquad
 \mathcal A_0=-\log Z_0.
\]
For marks \(\alpha=(i,h),\beta=(j,k)\) with \(|i-j|\ge2\), every
native term has zero direct mixed bridge derivative. Exact
fixed-domain differentiation consequently gives
\begin{equation}
 \mathcal H_\gamma(\alpha,\beta)
       :=\partial_{i,h}\partial_{j,k}\mathcal A_\gamma
       =-\operatorname{Cov}_{\mu_\gamma}
                               (J_{\gamma,\alpha},J_{\gamma,\beta}),
 \label{f16v6:exact-marked-Hessian}
\end{equation}
and the same identity with subscript zero. This is the inherited
source/Hessian identity, not a new formal claim. Differentiation is
legitimate at each finite regulator: the principal fast domain is fixed,
the Haar weight is integrable, and the exponential-word derivatives
are bounded on its compact closure. There is an open bridge
neighbourhood of every bounded base point for large \(\gamma\).

The separated-time Hessian in
\eqref{f16v6:exact-marked-Hessian} is unchanged when \(\mathcal A_\gamma\)
is replaced by \(-\log\mathcal W_{\gamma,n}^{\mathrm{full}}\), or when
\(\mathcal A_0\) is replaced by \(-\log\mathcal W_n\). All external
logarithmic factors in V4 are sums of single-time functions or
history-independent scalars. In particular this assertion does not
require estimating or dropping either \(F_{\gamma,2}\) normalizer.

\subsection{An actual source-marked bank estimate, not a pressure derivative}

\begin{theorem}[Full-native separated-time Hessians on bounded-overlap banks]
\label{f16v6:marked-bank}
Let \((\alpha_a,\beta_a)_{a\in\mathcal I}\) be localized bridge-mark
pairs at nonadjacent times. Choose nonempty carriers
\(S_a=S_{\alpha_a}\cup S_{\beta_a}\) with \(|S_a|\le s\), and
suppose each fast cell lies in at most \(q\) such carriers, counting
multiplicity. For fixed \(s,q\), uniformly in \(B,n\), the positions
and separations of the marks, and every bounded bridge history,
\begin{equation}
 \frac1{\mathcal N}\sum_{a\in\mathcal I}
       |\mathcal H_\gamma(\alpha_a,\beta_a)
                         -\mathcal H_0(\alpha_a,\beta_a)|
                       \le C_{r,s}\,q\,\omega_\gamma.
 \label{f16v6:marked-bank-bound}
\end{equation}
The same estimate holds for the Hessians of the exact native and
Gaussian full history weights. No intermediate fast projection or
native small-field restriction occurs in the left side.
\end{theorem}

\begin{proof}
Set \(d_a=D((\mu_\gamma)_{S_a}\Vert(\mu_0)_{S_a})\). All
Gaussian marginals in this family have the uniform parameters required
in Proposition~\ref{f16v6:score-entropy}, by
\eqref{f16v6:gaussian-data} and \(|S_a|\le s\). Their diameter does
not affect these bounds. Apply that proposition on the union carrier,
with affine functions \(J_{0,\alpha_a},J_{0,\beta_a}\) and the two
native remainders supplied by
Proposition~\ref{f16v6:native-score}. The exact source identity gives
\begin{equation}
 |\mathcal H_\gamma(\alpha_a,\beta_a)
                    -\mathcal H_0(\alpha_a,\beta_a)|
                      \le C_{r,s}(\sqrt{d_a}+d_a+\gamma^{-1/2}).
 \label{f16v6:individual-entropy-bound}
\end{equation}
This inequality is also useful separately: it identifies the specific
\emph{marginal entropy} needed to control an individual pair.

Equation~\eqref{f16v6:bank-entropy} bounds
\(\sum_a d_a\) by \(C_r q\mathcal N\epsilon_\gamma\).
Nonemptiness and the overlap condition imply
\(|\mathcal I|\le q\mathcal N\). Cauchy--Schwarz therefore gives
\[
 \frac1{\mathcal N}\sum_a\sqrt{d_a}
     \le\sqrt{\frac{|\mathcal I|}{\mathcal N}
                    \frac{\sum_a d_a}{\mathcal N}}
     \le C_rq\sqrt{\epsilon_\gamma}.
\]
Sum \eqref{f16v6:individual-entropy-bound} and use
\(\epsilon_\gamma\le\omega_\gamma\) and
\(\gamma^{-1/2}\le\omega_\gamma\) for large \(\gamma\).
This proves \eqref{f16v6:marked-bank-bound}. The conclusion about
history weights follows from the exact single-time-normalizer
observation after \eqref{f16v6:exact-marked-Hessian}.

At no step was \(\log Z_\gamma-\log Z_0\) differentiated.
The pressure estimate supplied a controlled likelihood moment, which
supplied marginal entropy; the actual source insertions were then
estimated directly, including their unbounded parts.
\end{proof}

The normalization by \(\mathcal N\), rather than by
\(|\mathcal I|\), is part of the theorem. A bank with positive density
has a vanishing mean absolute error per mark. A lone inhomogeneous mark
would still inherit a possible \(\mathcal N\) loss through its
uncontrolled marginal entropy; the bank inequality does not conceal
that distinction. Repeating a fixed pair many times raises \(q\)
by the same factor and creates no pointwise improvement.

\subsection{Fixed-lag common-bridge memory and the full/core correction}

Let \(h^{b,\mu,v}\) be the common-bridge direction on the four parallel
bridges from cube \(b\) to \(b+e_\mu\), with value \(v/2\) on each,
where \(|v|=1\) is a colour vector. Thus \(\|h^{b,\mu,v}\|=1\) and
\(h^{b,\mu,v}\in\mathcal E\). Fix two such direction types, a spatial
cube displacement \(a\), and a lag \(d\) with \(2\le d\le n\).
For \(0\le i\le n-d\) and every cube \(b\), write
\[
 \alpha_{i,b}=(i,h^{b,\mu,v}),\qquad
 \beta_{i,b}=(i+d,h^{b+a,\nu,w}).
\]
The two types or colours need not agree. The displacement may depend on
the spatial torus and need not remain bounded.

\begin{theorem}[Full-law common-bridge response at any fixed lag, in density]
\label{f16v6:common-bank}
Uniformly in the above choices and in every bounded bridge history,
\begin{equation}
 \frac1{B(n+1)}\sum_{i=0}^{n-d}\sum_b
       |\partial_{\alpha_{i,b}}\partial_{\beta_{i,b}}
                (-\log\mathcal W_{\gamma,n}^{\mathrm{full}})|
                                      \le C_r\omega_\gamma.
 \label{f16v6:common-native}
\end{equation}
Moreover, for V4's \emph{original} core and its actual probability
\(p_{\rm good}\), with any admissible \(R\),
\begin{equation}
 \frac1{B(n+1)}\sum_{i=0}^{n-d}\sum_b
       |\partial_{\alpha_{i,b}}\partial_{\beta_{i,b}}
                                        \log p_{\rm good}|
 \le C_r\omega_\gamma
          +C\frac{R^2}{\gamma}
                         (9/19)^{(d-2)_+}.
 \label{f16v6:common-correction}
\end{equation}
In particular \(R=4\sqrt{\log\gamma}\) is admissible eventually and
makes the last right side at most \(C_r\omega_\gamma\).
\end{theorem}

\begin{proof}
The source carrier of a direction anchored at \((i,b)\) is contained
in the cells at times \(i-1,i,i+1\) and spatial cube
\(\ell^\infty\)-distance at most one from \(b\). This follows from
the literal mixed plaquette and kinetic words; it is a safe enlargement
of the smaller actual carrier. It has at most \(3\cdot3^3=81\)
cells. A union of the two carriers has at most 162 cells. For the
translated fixed-lag bank above, any given cell occurs in at most 81
first carriers and at most 81 second carriers, even with periodic
wrapping and time endpoints. Thus \(s=q=162\) are safe uniform choices;
large separation does not enlarge them.

The inherited Gaussian cancellation is exact. In V4's moving fast
coordinates, \(D_x\partial_{i,h}V_0=\Lambda h=0\) for a common
bridge direction. Its remaining fast source uses only the port bonds
\(i-1,i\). The Gaussian fast Hessian is block diagonal between the
internal species and the ports, and between different port bonds.
Consequently the Gaussian source covariance at times separated by at
least two is zero. This is precisely the common-bridge null response
of V1/V3 together with V4's retained-port kinetic completion, not a
new Gaussian mass argument. V5's determinant-one linear coordinate
change leaves this covariance unchanged. Hence
\[
         \mathcal H_0(\alpha_{i,b},\beta_{i,b})=0.
\]
Theorem~\ref{f16v6:marked-bank} now gives
\eqref{f16v6:common-native}.

Finally the exact full/core identity has the sign
\[
 \partial_\alpha\partial_\beta\log p_{\rm good}
   =\partial_\alpha\partial_\beta
                (-\log\mathcal W_{\gamma,n}^{\mathrm{full}})
     -\partial_\alpha\partial_\beta
                (-\log\mathcal W_{\gamma,n}^{R}).
\]
Apply the triangle inequality, the first part of this theorem, and
V4's common-bridge bound
\eqref{f16v4:common-memory-bound}. The latter has
\(\delta^2=C_*^2R^2/\gamma\), and the bank has at most
\(B(n+1)\) pairs. This proves
\eqref{f16v6:common-correction}. The stated growing radius satisfies
all V4 core and mean-margin conditions eventually at fixed \(r\).
No event in the unmoving coordinates has been substituted for that
original core.
\end{proof}

For other unit localized bridge directions the Gaussian entry need not
vanish. V3's inverse bound gives, at lag \(d\ge2\),
\[
 |\mathcal H_0(\alpha,\beta)|
       \le\frac{64}{p_*-2}\left(\frac2{p_*}\right)^d,
 \qquad p_*=3+2\sqrt2,
\]
using \(\|N^{-1/2}A^{\mathsf T}\|\le2\) and \(\|D\|\le4\).
For the corresponding fixed-overlap translated banks,
\eqref{f16v6:marked-bank-bound} thus gives a mean absolute native
bound \(C_r[(2/p_*)^d+\omega_\gamma]\). The second term is a
separation-independent error floor, not exponential mixing.

\subsection{A pointwise spatial-volume-uniform consequence at zero history}

The full native law, Gaussian reference, and original V4 core at
\(\mathbf y=0\) are invariant under spatial translations by whole
cubes. The endpoint amplitudes are identical products over cubes, and
the rooted tree and frame in each cube use the same relative port
convention. No temporal translation invariance is asserted for these
dressed-end histories.

\begin{theorem}[Individual common-source entries at zero bridge history]
\label{f16v6:homogeneous}
Fix common-bridge direction types as above, two times \(i,j\) with
\(|i-j|\ge2\), and a spatial displacement \(a\). At \(\mathbf y=0\),
for every cube \(b\),
\begin{equation}
 |\partial_{i,h^{b,\mu,v}}\partial_{j,h^{b+a,\nu,w}}
                    (-\log\mathcal W_{\gamma,n}^{\mathrm{full}})|
 \le C\left\{\sqrt{(n+1)\epsilon_\gamma}
                  +(n+1)\epsilon_\gamma+\gamma^{-1/2}\right\}.
 \label{f16v6:homogeneous-bound}
\end{equation}
The constant is independent of \(B,n,i,j,a\). In particular, if
\((n+1)\epsilon_\gamma\le1\) and \(R=4\sqrt{\log\gamma}\),
\begin{equation}
 \begin{split}
 &|\partial_{i,h^{b,\mu,v}}\partial_{j,h^{b+a,\nu,w}}
                    (-\log\mathcal W_{\gamma,n}^{\mathrm{full}})|\\
 &\quad+
 |\partial_{i,h^{b,\mu,v}}\partial_{j,h^{b+a,\nu,w}}
                                      \log p_{\rm good}|
                  \le C\sqrt{(n+1)\epsilon_\gamma}.
 \end{split}
 \label{f16v6:homogeneous-correction}
\end{equation}
Thus at each fixed history length these individual derivatives vanish
uniformly in spatial volume. The same is true along growing lengths
with \((n+1)\epsilon_\gamma\to0\).
\end{theorem}

\begin{proof}
Use the bank of the \(B\) translated pairs at these fixed two times,
with the carriers chosen as spatial translates of one another. Their
overlap is bounded by the same geometric constant. By exact spatial
translation symmetry, every marginal entropy \(d_b\) for these
carriers has the same value. Equation~\eqref{f16v6:bank-entropy}
therefore gives
\[
       d_b=B^{-1}\sum_{b'}d_{b'}
                              \le C(n+1)\epsilon_\gamma.
\]
Apply the individual estimate
\eqref{f16v6:individual-entropy-bound} and the exact common-bridge
Gaussian cancellation. This proves
\eqref{f16v6:homogeneous-bound}, with no spatial-volume factor.
It does not divide by the number of time slices: the dressed endpoints
break that symmetry.

The full/core identity and V4's common-source estimate add at most
\(C(R^2/\gamma)(9/19)^{(|i-j|-2)_+}\) to bound the second term in
\eqref{f16v6:homogeneous-correction}. If
\((n+1)\epsilon_\gamma\le1\), the linear entropy term in
\eqref{f16v6:homogeneous-bound} is bounded by its square root. At the
specified \(R\), both \(\gamma^{-1/2}\) and \(R^2/\gamma\) are
bounded by \(C\sqrt{(n+1)\epsilon_\gamma}\) for large \(\gamma\).
This proves the last assertion. Spatial symmetry is used only at the
stated zero bridge history; no symmetry of an arbitrary bounded
history has been assumed.
\end{proof}

\subsection{What remains, and the next noncircular target}

V6 has now estimated the actual mixed derivatives of
\(\log p_{\rm good}\) in two specified senses: a duration-uniform
\emph{density} bound on fixed-lag common-source banks at arbitrary
bounded histories, and an individual-entry bound uniform in spatial
volume at zero history, with its length dependence explicit. This is
stronger than the V5 scalar pressure and one-cell moment statements.
It is not the uniform pointwise, time-summable response originally
needed for the complete spectral program.

There are two precise places where the present proof loses that
stronger conclusion. For an individual pair at an inhomogeneous base
history, \eqref{f16v6:individual-entropy-bound} requires its own marginal
entropy \(d_a\), whereas the proved input is only the bank sum
\eqref{f16v6:bank-entropy}. Second, even a uniform small bound on every
such \(d_a\) would give a separation-independent covariance error via
the present argument. It would not control the sum over all remote
lags. Applying the fixed-lag bank theorem repeatedly costs the number
of lags; it is not a hidden summability theorem.

A useful V7 step is therefore a \emph{conditional, source-local}
entropy or covariance estimate for the connected part between two
separating collars, rather than another unmarked global pressure
comparison. The Gaussian marginal localization proved here provides
local normalization and unbounded-score budgets that such an argument
can reuse. A new proof must additionally show how native communication
through the intervening fast variables decays, or bound its connected
source-visible contribution. Replacing \(d_a\) by an assumed uniform
small number, using a native functional inequality not yet proved, or
differentiating the likelihood-pressure remainder would be circular.
The inherited f15 V23 counterexample already warns that small entropy
alone does not yield the missing native functional inequality; it is
not reproduced here.

All estimates remain for fixed rescaled bridge windows before the
remaining root Gauss restriction. The source directions used here are
coordinate directions, not newly identified physical Hamiltonian
states. Gauge-invariant bridge-tail coverage, a comparison with the
same full-transfer top eigenvalue, physical-time calibration and a
nontrivial interacting continuum construction remain additional
requirements.

\subsection{Verification and exact preservation}

The new diagnostic is
\begin{center}\small
\path{verify_ym_f16_v6_marked_entropy.py}.
\end{center}
Its 906 exact assertions passed. They check the compensated likelihood
identity and scalar signs, correlated Gaussian marginal-projection
matrices and overlap budgets, noncommuting history-coordinate source
cancellations, literal finite-torus source carriers, and the two source
Taylor orders. The finite carrier checks use fine spatial tori of sides
eight and ten, with endpoint and distant-pair support counts kept.
There are no floating-point fixtures in this checker. These tests
are diagnostics of the indicated algebra and support bookkeeping,
not formal verification of the analytic entropy, interpolation, or
source-tail arguments. The report is
\path{ym_v6_verification_report.json}.
The checker has 15,240 bytes and SHA--256
\begin{center}\scriptsize\ttfamily
598A2E0AAA1F37E32BFB5114B5E83B37C76F563888407158C4894E02B62ED9C5.
\end{center}

The V4 and V5 diagnostics available in the supplied V5 package were
rerun separately and passed 9,467 and 89,339 assertions, respectively.
Their exact input scripts, result reports and hashes are supplied.
No unavailable older checker is represented as rerun. The V5 predecessor
has 153,945 bytes and SHA--256
\begin{center}\scriptsize\ttfamily
7751C0A73C4344D29305D84C8180BFE25A97CF0F33696394113742B490CF4F82.
\end{center}

The cumulative V6 changes only the V5 title version and inserts this
appendix before its final \verb|\end{document}|. Removing this exact
appendix and restoring the title recovers V5 byte for byte. The input
hash manifest and the preservation/build reports accompany the
artifacts. The original uploads and earlier generated sources remain
unchanged. This is not a scheduled five-version companion checkpoint;
V5's next scheduled checkpoint remains V10.

\begin{firewall}
V6 proves a full-native likelihood and entropy-density comparison to
the complete Gaussian fast-history law, then a source-marked
separated-time Hessian comparison for bounded-overlap banks with
vanishing mean absolute error per cell. It estimates the actual
full/core common-source correction in that density sense and,
at zero bridge history, pointwise uniformly in spatial volume with
an explicit length cost. No arbitrary-history pointwise bound,
time-summable full-law memory estimate, native functional inequality,
gauge-invariant bridge-tail coverage, full-transfer spectral gap,
interacting continuum construction, or positive physical Yang--Mills
mass gap is established by this appendix.
\end{firewall}
% END F16 V6 APPEND

% BEGIN F16 V7 APPEND -- 2026-09-17
% Insert immediately before the final \end{document} in cumulative V6.
% The predecessor body is unchanged; the cumulative title becomes V7.

\clearpage
\section{V7: Haar-compensated common sources and whole-history trace-norm memory}
\label{f16v7:context}

V6 estimates the covariance of the original bridge scores by comparison
with their affine Gaussian limits. Its cross terms cost the square root
of the local entropy, and its separated-time error has no decay in the
separation. Repeating that argument at more lags is not an improvement.
Here we first change the \emph{exact source representation}, using Haar
integration by parts on the temporal ports. For common-bridge directions,
the new Gaussian score is deterministic. The complete nonlinear returned
covariance is consequently the Gram of two remainder scores, not a cross
term between a remainder and an order-one Gaussian score.

This yields a different output from V6. We bound the \emph{entire}
nonadjacent-time common-source Hessian in trace norm per cube and time
slice, simultaneously including all lags and spatial displacements.
The rate is $O_r(\epsilon_\gamma)$ rather than
$O_r(\sqrt{\epsilon_\gamma})$. The normalization by the number of cells
is essential: a trace-norm density is not a uniform operator norm, and
it is not an absolutely summable temporal row estimate. The same bound
is obtained with exactly finite spatial supports of logarithmically
growing radius. No native functional inequality is assumed.

\subsection{Dependencies and the non-duplication boundary}

The supplied V6 appendix was read in full. Section inventories and
focused searches of f14 V100, f15 V100 and the Grand Tensor Monograph
V076 were followed by reading the relevant interfaces. This is not an
independent certification of all archived proofs. Haar integration by
parts itself is already present in f14 V97, especially
\texttt{prop:v97-finite-Stein}, and in f15 V5--V6. The monograph's
\texttt{thm:YM-Wilsonian-signed-head} already keeps the full direct term
minus one covariance return. Those facts are inherited, not claimed anew.

The particular temporal-port lift below uses V2--V4's kinetic square and
V1's $A^{\mathsf T}D\mathcal E=0$. Its native energy estimate uses V6's
\emph{already proved} complete likelihood bound. There is no new corridor
pressure calculation. Gaussian logarithmic Sobolev and finite-matrix
trace-ideal arguments are standard tools; their precise versions are
proved or derived below. The new application is the exact compensated
common-source bank, its full-native quadratic energy budget, and the
whole-history trace-norm estimate for the actual full/core correction.

\subsection{A uniformly bounded temporal-port lift of every common source}

Keep the unmoving fast variables $\zeta=(x,a)$, full laws
$\mu_\gamma,\mu_0$, and potentials $U_\gamma,U_0$ of
\eqref{f16v6:laws}. Thus $a$ is the principal logarithm of V5's temporal
\emph{group} variable $h$, not V4's moving variable $r$. Set
\[
 \mathcal N=B(n+1),\qquad
 \epsilon_\gamma=\gamma^{-1/6}(1+\log\gamma)^{2/3},\qquad
 \max_{i,e}|y_{i,e}|\le r<\infty.
\]
Unless stated otherwise, all constants are independent of $B=m^3$,
$m\ge4$, $n\ge1$, and the admitted bridge history. Thresholds and
constants with subscript $r$ may depend on its fixed window.

The real common-source space is
\[
 \mathcal U=\bigoplus_{i=0}^n\mathcal E,
 \qquad \dim\mathcal U=9B(n+1)=9\mathcal N.
\]
Use the orthonormal basis from V6: at time $i$, each of the four parallel
bridges from $b$ to $b+e_\mu$ carries $e_k/2$, for
$\mu,k\in\{1,2,3\}$. Denote a basis index by $\alpha=(i,b,\mu,k)$.
The norm on $\mathcal U$ is the inherited Euclidean norm of the bridge
variations, not a physical Hamiltonian norm.

Write $E=E_{\rm br}$ in this appendix only, and retain
\[
 L_c=\mathsf B_B^{\mathsf T}\mathsf B_B,\qquad
 L_h=L_c+E^{\mathsf T}E,\qquad
 G^{-1}=I-EL_h^{-1}E^{\mathsf T}.
\]
Colour identity factors are suppressed. For a common history variation
$\mathbf v=(v_0,\ldots,v_n)$ put
$(d_t\mathbf v)_t=v_t-v_{t+1}$ and define the real linear map
\begin{equation}
 (\mathsf P\mathbf v)_t=-L_h^{-1}E^{\mathsf T}(v_t-v_{t+1}).
 \label{f16v7:lift}
\end{equation}
Its range is the $21Bn$-dimensional temporal-port Lie algebra. The matrix
$\mathsf P$ depends on the finite geometry, not on $\zeta$, $\mathbf y$
or $\gamma$.

\begin{proposition}[Common-source cancellation and finite-range lifts]
\label{f16v7:lift-prop}
Let $\partial_{\mathbf v}$ differentiate the retained bridge history.
In the complete Gaussian model,
\begin{equation}
 \left(\partial_{\mathbf v}
              +(\mathsf P\mathbf v)\cdot\partial_a\right)U_0
                  =g_{\mathbf v}(\mathbf y),
 \label{f16v7:Gaussian-cancellation}
\end{equation}
where the right side is independent of every fast variable and equals
\begin{equation}
 g_{\mathbf v}(\mathbf y)
 =\sum_{i=0}^n\langle v_i,D^{\mathsf T}Dy_i\rangle
   +\sum_{t=0}^{n-1}
       \langle v_t-v_{t+1},G^{-1}(y_t-y_{t+1})\rangle.
 \label{f16v7:g}
\end{equation}
Moreover $\|\mathsf P\|\le 2\|L_h^{-1}E^{\mathsf T}\|\le C$.
There are matrices $\mathsf P_\ell$, $\ell=0,1,\ldots$, supported within
spatial cube distance $\ell+1$ of each source and on its two adjacent
port bonds, with
\begin{equation}
 \|\mathsf P_\ell\|\le C,\qquad
 \|\mathsf P_\ell-\mathsf P\|\le C q_h^{\ell+1},\qquad 0<q_h<1.
 \label{f16v7:lift-approximation}
\end{equation}
All constants are independent of volume and duration.
\end{proposition}

\begin{proof}
Return only for this calculation to the Gaussian moving variables
$z_t=a_t-\mathsf F_B(x_t-x_{t+1})$ from V5. At fixed $x$,
$\partial_a=\partial_z$. The port-dependent part of the quadratic
potential is
\[
 \frac12\sum_t
 \{z_t^{\mathsf T}L_hz_t+
       2z_t^{\mathsf T}E^{\mathsf T}(y_t-y_{t+1})
                         +|y_t-y_{t+1}|^2\}.
\]
There is no $xz$ cross term by the balanced kinetic square. The magnetic
bridge derivative in a common direction has no $x$ term because
$A^{\mathsf T}Dv_i=0$. Its remaining term is
$\langle v_i,D^{\mathsf T}Dy_i\rangle$. Substitution of
\eqref{f16v7:lift} cancels every $z_t$ coefficient and leaves
\eqref{f16v7:g}, with the stated $G^{-1}$. This reuses the inherited
Gaussian cancellation; the next subsection implements it by a legitimate
\emph{compact} integration-by-parts field.

For the norm and support estimates take fixed geometric bounds
\[
 \lambda_-:=\lambda_{\min}(\mathsf B^{\mathsf T}\mathsf B)>0,
 \qquad \lambda_+:=\|\mathsf B^{\mathsf T}\mathsf B\|+6,
 \qquad q_h=1-\lambda_-/\lambda_+.
\]
V2 gives $\|E\|^2\le6$, so
$\lambda_-I\le L_h\le\lambda_+I$. Also $\|d_t\|\le2$.
Set
\[
 T_h=I-L_h/\lambda_+,\qquad
 Q_\ell=\lambda_+^{-1}\sum_{j=0}^{\ell}T_h^j,
 \qquad
 (\mathsf P_\ell\mathbf v)_t=-Q_\ell E^{\mathsf T}(v_t-v_{t+1}).
\]
Then $0\le T_h\le q_hI$, $\|Q_\ell\|\le\lambda_-^{-1}$ and
\[
 L_hQ_\ell=I-T_h^{\ell+1},\qquad
 \|Q_\ell-L_h^{-1}\|\le\lambda_-^{-1}q_h^{\ell+1}.
\]
The spatial blocks of $L_h$ have nearest-cube range; its powers have
the corresponding degree range. A localized common bridge source enters
$E^{\mathsf T}$ only at its endpoint cubes. These facts prove the stated
support and both norm bounds. No decay of a native covariance has been
used. Hereafter $\mathsf P_\infty=\mathsf P$ and
$q_h^{2(\infty+1)}=0$ are convenient conventions.
\end{proof}

\subsection{The exact compact source identity and its changed contact term}

For a temporal port $p=(t,b,v)$ and colour $k$, define the
\emph{left-translation derivative}
\[
 (X^\gamma_{p,k}F)(x,h;\mathbf y)
  =\left.\frac{d}{ds}
     F(x,\ldots,\operatorname{Exp}(s e_k/\sqrt\gamma)h_p,
                         \ldots;\mathbf y)\right|_{s=0}.
\]
This terminology specifies the flow, without assigning it an incorrect
Lie-bracket sign. These derivatives do not in general commute on a
single group. Their reference divergence is zero for product Haar.
They differentiate neither the chord variables nor their endpoint
amplitudes. In particular $X^\gamma_{p,k}$ must \emph{not} be replaced
by flat differentiation of $a_{p,k}$ at finite $\gamma$.

For $\ell\in\{0,1,\ldots,\infty\}$ and a basis mark $\alpha$, set
\begin{align}
 V_\alpha^{[\ell]}
   &=\sum_{p,k}(\mathsf P_\ell)_{p,k;\alpha}X^\gamma_{p,k},&
 \mathcal D_\alpha^{[\ell]}&=\partial_\alpha+V_\alpha^{[\ell]},
 \notag\\
 \widehat J_\alpha^{[\ell]}&=\mathcal D_\alpha^{[\ell]}U_\gamma,&
 r_\alpha^{[\ell]}&=\widehat J_\alpha^{[\ell]}-g_\alpha(\mathbf y),
 \label{f16v7:lifted-scores}
\end{align}
where $g_\alpha=g_{\mathbf v_\alpha}$. All the coefficients of $V$ are
fixed as the bridge history varies. Let $\mathbf r^{[\ell]}$ be the
complete vector of these $9\mathcal N$ residual scores.

\begin{theorem}[Full Haar-compensated Hessian, with all contact terms]
\label{f16v7:Haar-identity}
At each finite regulator, for the \emph{full} native conditional,
\begin{align}
 \partial_\alpha(-\log Z_\gamma)
     &=\mathbb E_\gamma\widehat J_\alpha^{[\ell]},\notag\\
 \partial_\alpha\partial_\beta(-\log Z_\gamma)
     &=\mathcal K^{[\ell]}_{\alpha\beta}
                              -\Gamma^{[\ell]}_{\alpha\beta},
 \label{f16v7:Hessian-decomposition}
\end{align}
where
\begin{align}
 \mathcal K^{[\ell]}_{\alpha\beta}
   &=\frac12\mathbb E_\gamma
      (\mathcal D_\alpha^{[\ell]}\mathcal D_\beta^{[\ell]}
        +\mathcal D_\beta^{[\ell]}\mathcal D_\alpha^{[\ell]})U_\gamma,
       \notag\\
 \Gamma^{[\ell]}_{\alpha\beta}
   &=\operatorname{Cov}_{\mu_\gamma}
                         (r_\alpha^{[\ell]},r_\beta^{[\ell]}),
       \qquad \Gamma^{[\ell]}\succeq0.
 \label{f16v7:contact-Gram}
\end{align}
If the source times obey $|i(\alpha)-i(\beta)|\ge2$, then
$\mathcal K^{[\ell]}_{\alpha\beta}=0$ \emph{exactly}. Thus every such
full-history Hessian entry is the negative of the indicated residual
covariance, also when $-\log Z_\gamma$ is replaced by
$-\log\mathcal W_{\gamma,n}^{\mathrm{full}}$.
\end{theorem}

\begin{proof}
Perform the port integrations in compact group variables before using
principal charts. At fixed $x$, their reference is product probability
Haar, up to a scalar independent of $x,h,\mathbf y$. The remaining chord
Jacobian factors depend on $x$ only. Haar invariance gives
\begin{equation}
 \mathbb E_\gamma[V_\alpha^{[\ell]}F]
    =\mathbb E_\gamma[F V_\alpha^{[\ell]}U_\gamma],
 \qquad \mathbb E_\gamma V_\alpha^{[\ell]}U_\gamma=0.
 \label{f16v7:Haar-parts}
\end{equation}
There is no artificial boundary term at the logarithm cut: it is not a
boundary of the compact port group. At finite $\gamma,B,n$, the relevant
functions and their port and bridge derivatives are bounded in the port
groups and in the closed principal chord ranges. The chord weights are
integrable. Fubini and dominated differentiation therefore justify all
identities, also if a port flow crosses its principal-chart cut.

Ordinary bridge differentiation gives the first line of
\eqref{f16v7:Hessian-decomposition}, by the second identity in
\eqref{f16v7:Haar-parts}. Differentiating that expectation in $\beta$
gives
\[
 \mathbb E_\gamma\partial_\beta\widehat J_\alpha^{[\ell]}
       -\operatorname{Cov}_\gamma
          (\widehat J_\alpha^{[\ell]},\partial_\beta U_\gamma).
\]
In the covariance substitute
$\partial_\beta U_\gamma=
 \widehat J_\beta^{[\ell]}-V_\beta^{[\ell]}U_\gamma$.
The added covariance is
$\mathbb E_\gamma V_\beta^{[\ell]}\widehat J_\alpha^{[\ell]}$
by \eqref{f16v7:Haar-parts}. This yields
\[
 \mathbb E_\gamma
       \mathcal D_\beta^{[\ell]}\mathcal D_\alpha^{[\ell]}U_\gamma
 -\operatorname{Cov}_\gamma
       (\widehat J_\alpha^{[\ell]},\widehat J_\beta^{[\ell]}).
\]
Interchange $\alpha,\beta$ and average. Subtraction of the deterministic
$g$'s does not change the covariance. This proves the symmetric contact
formula without assuming commuting group derivatives. Equivalently,
the expectation of the commutator contribution vanishes; it is another
Haar-divergence identity, not a pointwise zero.

Each native kinetic term uses ports on \emph{one} bond $t$ and bridge
variables only at $t,t+1$. A spatial magnetic term has no temporal port
and uses only one bridge time. The lift of a source at $i$ uses only
bonds $i-1,i$, even for $\ell=\infty$. For source times at least two
apart, these two bond sets are disjoint. Expanding the two $\mathcal D$'s
shows that every double derivative of every native term vanishes:
there is no common bond, no two nonadjacent bridge times in one term,
and no source-dependent lift coefficient. This proves the contact
support assertion. The endpoint Gaussian amplitude has no such
variation at all.

Finally V4's exact external factors have logarithms that are sums of
single-time functions or constants. Their Hessians vanish at these
separated times. They have not been estimated, omitted, or renormalized.
\end{proof}

The contact matrix has changed along with the score covariance. Merely
subtracting a Gaussian-looking random variable from the old scores,
while keeping the old direct Hessian, would be wrong. Equation
\eqref{f16v7:Hessian-decomposition} is a new representation of the
\emph{same} Hessian, not an extra negative return to add to V4 or to the
monograph's already integrated response. Positivity is asserted for
$\Gamma^{[\ell]}$, not for the contact or for its time-off-diagonal part.

\subsection{A whole-vector remainder certificate on complete charts}

Let $J_\gamma=(\partial_\alpha U_\gamma)_\alpha$ be the ordinary
common-source bank, and write
\[
 S_\gamma=(X^\gamma_{p,k}U_\gamma)_{p,k},\qquad
 J_0=(\partial_\alpha U_0)_\alpha,\qquad S_0=\partial_aU_0.
\]
The group-insertion formulas defining $S_\gamma$ make sense as analytic
exponential-word functions on all real $\zeta$; use those expressions
for Gaussian estimates outside the native charts. Their restriction to
the charts is precisely the compact derivative, with no change of law.
Put $e^{\rm br}=J_\gamma-J_0$ and $e^{\rm pt}=S_\gamma-S_0$.

\begin{proposition}[Small Gaussian energy and dimension-free Lipschitz norm]
\label{f16v7:remainder-certificate}
The complete residual vector obeys, uniformly in $\ell$,
\begin{equation}
 \sup_{\zeta}\|D_\zeta\mathbf r^{[\ell]}\|\le L_*,\qquad
 \mathbb E_0\|\mathbf r^{[\ell]}\|^2
       \le C_r\mathcal N
                   (\gamma^{-1}+q_h^{2(\ell+1)}),
 \label{f16v7:reference-energy}
\end{equation}
where $L_*\ge1$ is a fixed geometric constant. For $\ell=\infty$ the
second summand on the right is zero. These are statements on the
\emph{whole score vector}; there is no overlap constant growing with
$\ell$ or the number of marks.
\end{proposition}

\begin{proof}
V6 already proves the local bridge-score Taylor bound. For clarity we
check the additional port-insertion derivatives and the whole-vector
norm. Each scalar native term is a rescaled trace of a fixed product
of exponentials of real skew-Hermitian linear arguments. To evaluate a
port insertion, add one factor
$\operatorname{Exp}(s e_k/\sqrt\gamma)$ adjacent to the appropriate
$h_p$ or its inverse and differentiate in the added scalar $s$ at zero.
This avoids differentiating a logarithm or its singular chart Jacobian.
Duhamel differentiation and unitarity give global bounds on the
second derivatives of this rescaled word, and a bound
$C/\sqrt\gamma$ on its third derivatives. At zero all first
variation terms vanish. The quadratic inserted derivative is exactly
$\partial_{a_{p,k}}U_0$; both operations insert the same linear Lie
algebra variable into the linearized kinetic word.

Consequently every component $e_j$ of the two primitive remainder banks
has a fixed-size local fast carrier $C_j$ and satisfies
\begin{equation}
 |e_j(\zeta)|\le C_r\gamma^{-1/2}(1+|\zeta_{C_j}|^2),\qquad
 \|D_{\zeta_{C_j}}e_j\|\le C.
 \label{f16v7:primitive-remainders}
\end{equation}
The first estimate integrates the third-derivative bound twice; the
second uses the global second-derivative bounds for the native inserted
score and the Gaussian score separately. The same reasoning applies to
negative occurrences of a port: the adjacent inverse and its sign are
kept in the insertion. Neither the fast Haar Jacobian nor a derivative
of the endpoint amplitude is part of $S_\gamma$, since the integration
by parts was in Haar measure and varied ports only.

There are $9\mathcal N$ bridge components and $21Bn$ port components.
Each row of their Jacobian has a fixed number of entries, and each
fast coordinate appears in a fixed number of rows. The row and column
Schur bound therefore proves
\[
 \left\|D_\zeta\binom{e^{\rm br}}{e^{\rm pt}}\right\|\le C.
\]
This is an operator norm, not a sum of component Lipschitz constants.
The Gaussian law from V6 has uniformly bounded local means and
covariances. Its fourth moments on every fixed-size $C_j$ are bounded
by $C_r$. Squaring the first bound in
\eqref{f16v7:primitive-remainders} and summing the at most
$30\mathcal N$ components gives
\begin{equation}
 \mathbb E_0(\|e^{\rm br}\|^2+\|e^{\rm pt}\|^2)
                                  \le C_r\mathcal N/\gamma.
 \label{f16v7:primitive-energy}
\end{equation}

The exact Gaussian cancellation supplies the vector identity
\begin{equation}
 \mathbf r^{[\ell]}
   =e^{\rm br}+\mathsf P_\ell^{\mathsf T}e^{\rm pt}
                   +(\mathsf P_\ell-\mathsf P)^{\mathsf T}S_0.
 \label{f16v7:residual-synthesis}
\end{equation}
The last term is zero for the full lift. For a finite lift, $S_0$ is
an affine \emph{Gaussian score}: $\mathbb E_0 S_0=0$ and
\[
 \operatorname{Cov}_0(S_0)=H_{aa}=I_n\otimes L_h,
              \qquad H=D_\zeta^2U_0.
\]
Indeed $\nabla U_0=H(\zeta-m)$, so the covariance of its full gradient
is $H$, and its $a$ corner is the displayed matrix. Therefore
\[
 \mathbb E_0\| (\mathsf P_\ell-\mathsf P)^{\mathsf T}S_0\|^2
 \le \lambda_+\|\mathsf P_\ell-\mathsf P\|_{\rm HS}^2
 \le C\mathcal N q_h^{2(\ell+1)}.
\]
Use $\|\mathsf P_\ell\|\le C$ in
\eqref{f16v7:residual-synthesis} and
\eqref{f16v7:primitive-energy} to prove its squared-norm estimate.
For its Jacobian use additionally $\|D_\zeta S_0\|\le\|H\|\le h_+$.
The same operator-norm synthesis gives $L_*$ independently of $\ell$.
In particular a growing finite support has not been paid by a growing
entropy-overlap factor.
\end{proof}

\subsection{Transferring the residual energy, not differentiating pressure}

Here is the standard Gaussian estimate needed for the transfer. It is
stated separately to distinguish a Gaussian functional inequality from
an unproved native one.

\begin{proposition}[Gaussian exponential energy of a Lipschitz vector map]
\label{f16v7:Gaussian-energy}
Let $Q$ be any finite-dimensional Gaussian of covariance $C\preceq cI$.
Let $R$ be a smooth vector-valued map with $\|DR\|\le L$, and set
$F=\|R\|^2$, $a=2cL^2$. For $0<t<a^{-1}$,
\begin{equation}
 \log\mathbb E_Q e^{tF}\le\frac{t\mathbb E_QF}{1-at}.
 \label{f16v7:Gaussian-energy-mgf}
\end{equation}
The target dimension of $R$ does not enter the constant.
\end{proposition}

\begin{proof}
We use the Gaussian logarithmic Sobolev inequality in the convention
\[
 \operatorname{Ent}_Q(e^f)
       \le\frac12\mathbb E_Q[e^f\langle\nabla f,C\nabla f\rangle].
\]
For a direct derivation of this convention, whiten $Q$. V6's Gaussian
entropy dissipation identity is
$\operatorname{Ent}_g u=\int_0^\infty I(P_su)\,ds$.
The Mehler commutation $\nabla P_su=e^{-s}P_s\nabla u$ and
conditional Cauchy--Schwarz give
$I(P_su)\le e^{-2s}I(u)$. Integrating gives the factor $1/2$;
affine transport yields the displayed covariance form. This is the
classical Gaussian inequality, not a new native estimate. A primary
reference is L.~Gross,
\href{https://doi.org/10.2307/2373688}{\emph{Logarithmic Sobolev
Inequalities}}, American Journal of Mathematics 97 (1975), 1061--1083.

Now $|\nabla F|^2\le4L^2F$. Write
$\psi(t)=\log\mathbb E_Q e^{tF}$. The displayed inequality gives
\[
 t\psi'(t)-\psi(t)\le at^2\psi'(t).
\]
Hence $((1-at)\psi(t)/t)'\le0$. Its limit at zero is
$\mathbb E_QF$, proving \eqref{f16v7:Gaussian-energy-mgf}.
All differentiations are legitimate for $t<a^{-1}$: global
Lipschitz continuity gives $|R(z)|\le|R(0)|+L|z|$, and the strictly
subcritical Gaussian quadratic exponential has finite moments. Smooth
bounded truncations followed by limits justify the logarithmic Sobolev
step at the same parameters. The constant-map case is immediate.
\end{proof}

\begin{theorem}[A complete-native residual-energy and Gram budget]
\label{f16v7:native-energy}
Define
\[
 \eta_{\gamma,\ell}=\epsilon_\gamma+\gamma^{-1}
                                      +q_h^{2(\ell+1)}.
\]
For all sufficiently large $\gamma$, every admitted bridge history,
and every $\ell\in\{0,1,\ldots,\infty\}$, one has
\begin{align}
 \mathbb E_\gamma\|\mathbf r^{[\ell]}\|^2
                        &\le C_r\mathcal N\eta_{\gamma,\ell},
                              \notag\\
 0\le\operatorname{Tr}\Gamma^{[\ell]}
                        &\le C_r\mathcal N\eta_{\gamma,\ell}.
 \label{f16v7:native-energy-bound}
\end{align}
There is also a fixed $s_*>0$, independent of all these regulators and
of $\ell$, such that
\begin{equation}
 \log\mathbb E_\gamma
               e^{s_*\|\mathbf r^{[\ell]}\|^2}
                         \le C_r\mathcal N\eta_{\gamma,\ell}.
 \label{f16v7:native-mgf}
\end{equation}
In particular the probability that the residual energy exceeds
$(C_r\mathcal N\eta_{\gamma,\ell}+u)/s_*$ is at most $e^{-u}$,
$u\ge0$, with the constant from \eqref{f16v7:native-mgf}.
\end{theorem}

\begin{proof}
V6 provides, for the \emph{same full normalized laws},
\[
 D(\mu_\gamma\Vert\mu_0)
       \le D_p(\mu_\gamma\Vert\mu_0)
       \le C_r\mathcal N\epsilon_\gamma,
                 \qquad p=1+\theta_*>1.
\]
Their Gaussian covariance obeys $C\preceq h_-^{-1}I$.
Use Proposition~\ref{f16v7:Gaussian-energy} with
$R=\mathbf r^{[\ell]}$, $L=L_*$, $c=h_-^{-1}$, and fix
$t_0=(4cL_*^2)^{-1}$. Then
\[
 \log\mathbb E_0e^{t_0\|\mathbf r^{[\ell]}\|^2}
     \le 2t_0\mathbb E_0\|\mathbf r^{[\ell]}\|^2.
\]
The entropy variational inequality
$t_0\mathbb E_\gamma F\le
D(\mu_\gamma\Vert\mu_0)+\log\mathbb E_0e^{t_0F}$,
together with \eqref{f16v7:reference-energy}, proves the first bound.
Centering can only reduce the sum of squared scores, giving the trace
bound. Conditional means are not set to zero.

For the exponential bound write $p'=p/(p-1)$ and $s_*=t_0/p'$.
H\"older's inequality for the full density $f=d\mu_\gamma/d\mu_0$
gives
\begin{align*}
 \log\mathbb E_\gamma e^{s_*F}
  &\le \frac1p\log\mathbb E_0f^p
                   +\frac1{p'}\log\mathbb E_0e^{t_0F}\\
  &=\frac1{p'}D_p(\mu_\gamma\Vert\mu_0)
                   +\frac1{p'}\log\mathbb E_0e^{t_0F}\\
  &\le C_r\mathcal N\epsilon_\gamma
                   +2s_*\mathbb E_0F.
\end{align*}
Insert the same Gaussian energy bound. Exponential Markov proves the
last assertion. All constants are uniform in the polynomial lift degree.
The only logarithmic Sobolev inequality used was Gaussian; a native
logarithmic Sobolev or Poincar\'e inequality has not been assumed.
\end{proof}

For the full lift this is $O_r(\mathcal N\epsilon_\gamma)$, since
$\gamma^{-1}\le\epsilon_\gamma$ eventually. The same rate holds with
an exactly finite-range lift: choose any integer $\ell_\gamma\ge0$ with
\[
 \ell_\gamma+1\ge
           \frac{\log(\epsilon_\gamma^{-1})}{2\log(q_h^{-1})}.
\]
Its spatial radius is $O(\log\gamma)$, and its two-bond temporal
support never grows. This finite-range version concerns the actual
compact source identity, not a truncated native action or measure.

\subsection{The whole nonlocal history Hessian in trace norm}

Let $\Pi_i$ be projection onto the $9B$ common-source coordinates at
time $i$. For any matrix $M$ on $\mathcal U$, define its
nonadjacent-time part by
\[
 \operatorname{Off}_{>1}(M)
      =\sum_{|i-j|\ge2}\Pi_iM\Pi_j.
\]
We use the matrix trace norm $\|M\|_{S_1}$, the sum of singular values.
It is not the entrywise $\ell^1$ norm. Define the \emph{actual} memory
matrix and the original full/core correction matrix by
\begin{align}
 \mathsf M_\gamma
    &=\operatorname{Off}_{>1}
         \bigl(\nabla_{\mathcal U}^2
                   [-\log\mathcal W_{\gamma,n}^{\mathrm{full}}]\bigr),
                  \notag\\
 \mathsf C_{\gamma,R}
    &=\operatorname{Off}_{>1}
                  \bigl(\nabla_{\mathcal U}^2\log p_{\rm good}\bigr).
 \label{f16v7:memory-matrices}
\end{align}
The derivatives are evaluated at the chosen, possibly inhomogeneous,
bridge history. V4's $p_{\rm good}$ and its moving-coordinate core event
are unchanged.

\begin{theorem}[All-lag common-source memory and the full/core correction]
\label{f16v7:whole-memory}
Uniformly in spatial volume and history length, for all sufficiently
large $\gamma$,
\begin{equation}
 \boxed{\quad
   \frac{\|\mathsf M_\gamma\|_{S_1}}{B(n+1)}
                                  \le C_r\epsilon_\gamma.\quad}
 \label{f16v7:whole-memory-bound}
\end{equation}
This includes all time lags and all spatial displacements in one matrix.
For every core radius $R$ satisfying V4's derivative regime,
\begin{equation}
 \boxed{\quad
  \frac{\|\mathsf C_{\gamma,R}\|_{S_1}}{B(n+1)}
                  \le C_r\epsilon_\gamma+C\frac{R^2}{\gamma}.\quad}
 \label{f16v7:whole-correction-bound}
\end{equation}
In particular $R=4\sqrt{\log\gamma}$ is admissible eventually and the
second bound is also $C_r\epsilon_\gamma$. Either the full lift or the
finite lift $\ell_\gamma$ above proves these same statements about the
same full native Hessian.
\end{theorem}

\begin{proof}
The exact contact support gives, for every $\ell$,
\begin{equation}
 \mathsf M_\gamma=-\operatorname{Off}_{>1}(\Gamma^{[\ell]}).
 \label{f16v7:off-Gram}
\end{equation}
Although the full contact and Gram depend on the lift, their
nonadjacent-time combination in \eqref{f16v7:off-Gram} does not.
No sign is assigned to this off-diagonal matrix.

Here is a duration-independent trace-norm bound for this extraction.
Complexify the source space and put
$U_s=\sum_{i=0}^n e^{\mathrm i i s}\Pi_i$. The Fourier coefficients
\[
 \Phi_k(M)=\frac1{2\pi}\int_0^{2\pi}
                     e^{-\mathrm i k s}U_sMU_s^*\,ds
                 =\sum_{i-j=k}\Pi_iM\Pi_j
\]
have $\|\Phi_k(M)\|_{S_1}\le\|M\|_{S_1}$ by unitary invariance
and the triangle inequality. Thus
\begin{equation}
 \|\operatorname{Off}_{>1}(M)\|_{S_1}
   =\|M-\Phi_{-1}(M)-\Phi_0(M)-\Phi_1(M)\|_{S_1}
                           \le4\|M\|_{S_1}.
 \label{f16v7:mask}
\end{equation}
For a positive Gram, $\|\Gamma^{[\ell]}\|_{S_1}
=\operatorname{Tr}\Gamma^{[\ell]}$. Use the full lift in
Theorem~\ref{f16v7:native-energy}, or $q_h^{2(\ell_\gamma+1)}
\le\epsilon_\gamma$, to obtain \eqref{f16v7:whole-memory-bound}.
The factor four has not been multiplied by the number of lags.

For the core, V4's whole common-source tail estimate
\eqref{f16v4:common-tail}, with its cutoff index equal to one, gives
\[
 \left\|\operatorname{Off}_{>1}
     (\nabla_{\mathcal U}^2[-\log\mathcal W_{\gamma,n}^{R}])
                                     \right\|_{\rm op}
 \le\frac{4K_{\rm c}^2}{1-\vartheta}\delta^2,
 \qquad \vartheta=9/19,\quad \delta=C_*R/\sqrt\gamma.
\]
This estimate acts on arbitrary common-source history variations, not
only on localized basis marks. Since $\dim\mathcal U=9\mathcal N$,
its trace norm divided by $\mathcal N$ is at most
$36K_{\rm c}^2\delta^2/(1-\vartheta)$. Finally the exact V4 identity
has the sign
\[
 \nabla^2\log p_{\rm good}
     =\nabla^2[-\log\mathcal W_{\gamma,n}^{\mathrm{full}}]
                  -\nabla^2[-\log\mathcal W_{\gamma,n}^{R}].
\]
Apply $\operatorname{Off}_{>1}$ and the trace-norm triangle inequality.
This proves \eqref{f16v7:whole-correction-bound}. No endpoint normalizer
or chart coverage hypothesis was changed.
\end{proof}

\subsection{What the stronger matrix budget does and does not imply}

There are two useful, rigorous readings of the new bound. First it
improves V6's fixed-lag bank rate. For any family of pairs of canonical
common-source basis marks with each mark appearing in at most $q_0$
pair endpoints, Gram positivity gives
\[
 \sum_a|\Gamma^{[\infty]}_{\alpha_a\beta_a}|
 \le\frac12\sum_a
       (\Gamma^{[\infty]}_{\alpha_a\alpha_a}
                         +\Gamma^{[\infty]}_{\beta_a\beta_a})
 \le\frac{q_0}{2}\operatorname{Tr}\Gamma^{[\infty]}.
\]
Thus each translated fixed-lag, fixed-displacement common-source bank
from V6 now has mean absolute full-native response per cell at most
$C_r\epsilon_\gamma$, and its full/core correction at lag $d\ge2$
is at most
\[
 C_r\epsilon_\gamma
      +C(R^2/\gamma)(9/19)^{d-2}.
\]
The same conclusion for fixed unit colour combinations follows by their
norm-one synthesis from the three colour basis marks. At zero history,
spatial translation invariance makes
$\mathbb E_\gamma|r_{i,b,\mu,k}^{[\infty]}|^2$ independent of $b$.
The total energy bound consequently gives the individual-entry estimate
\begin{equation}
 |\partial_\alpha\partial_\beta
           [-\log\mathcal W_{\gamma,n}^{\mathrm{full}}]|\big|_{\mathbf y=0}
       \le C(n+1)(\epsilon_\gamma+\gamma^{-1}),
       \qquad |i(\alpha)-i(\beta)|\ge2.
 \label{f16v7:zero-entry}
\end{equation}
Here Cauchy--Schwarz was applied to the two \emph{residual} variances.
The length cost remains; no temporal stationarity is assumed.

Second, the whole-history bound controls the distribution of singular
values even when all lags are present. For any $\tau>0$,
\begin{equation}
 \frac{\#\{j:\sigma_j(\mathsf M_\gamma)>\tau\}}{9\mathcal N}
        \le\frac{C_r\epsilon_\gamma}{\tau}.
 \label{f16v7:singular-count}
\end{equation}
This follows by summing the singular values above $\tau$. As the matrix
is real symmetric, its spectral projection $P_\tau$ onto those absolute
eigenvalues obeys
\[
 \operatorname{rank}P_\tau\le C_r\mathcal N\epsilon_\gamma/\tau,
 \qquad \|\mathsf M_\gamma(I-P_\tau)\|_{\rm op}\le\tau.
\]
With $\tau=\sqrt{\epsilon_\gamma}$, the exceptional fraction and the
complement operator norm both vanish. The analogous statements hold for
$\mathsf C_{\gamma,R}$ with $\epsilon_\gamma$ replaced by
$\epsilon_\gamma+R^2/\gamma$.

These spectral projections live in the \emph{finite common-source
coordinate space}. They can depend on the complete base history. They
are not physical low-energy projections and cannot be discarded in a
Yang--Mills gap proof. A vanishing trace-norm density allows a coherent
exceptional direction with a large eigenvalue: a positive rank-one matrix
$\eta\mathbf1\mathbf1^{\mathsf T}$ in dimension $M$ has both trace and
operator norm $M\eta$. In particular neither
\eqref{f16v7:whole-memory-bound} nor the residual-energy exponential
moment supplies a bound on every normalized linear combination of
scores independent of volume.

\subsection{The remaining noncircular target}

The new estimate avoids V6's loss from summing its error floor over all
lags, but it does so in \emph{normalized trace norm}, not by proving
entrywise decay or temporal row summability. The small positive Gram
in \eqref{f16v7:contact-Gram} identifies the precise remaining object.
One possible stronger target is
\[
 \operatorname{Var}_{\mu_\gamma}
       \left(\sum_\alpha v_\alpha r_\alpha^{[\ell_\gamma]}\right)
                 \le a_\gamma\sum_\alpha|v_\alpha|^2,
 \qquad a_\gamma\longrightarrow0,
\]
uniformly in $v,B,n$ and the admitted histories. A time-weighted
connected-covariance estimate for these same residuals would additionally
address summability. Neither statement follows from the trace bound.

The important preparation for a collar or capacity argument is now
concrete: each residual can be built from exact compact Haar flows with
finite spatial radius $O(\log\gamma)$ and only two port times; its entire
native squared-energy density is small; and the Gaussian affine kinetic
score has already been removed with its contact cost retained. A future
argument must control \emph{coherent covariance} of these corrected
sources through the intervening native variables. Repeating the full
likelihood calculation or assuming a native functional inequality would
not supply that control. The source lift is not a physical gauge fixing
and does not establish gauge-invariant coverage of the retained bridge
window. The same-top full-transfer spectral comparison, physical-time
calibration, and interacting continuum construction remain separate. No infrared coercivity or positivity
for the changed contact matrix has been established either.

\clearpage
\subsection{Diagnostics, preservation, and claim scope}

The new diagnostic is
\begin{center}\small
\path{verify_ym_f16_v7_haar_memory.py}.
\end{center}
It checks the lift and its polynomial remainder, noncommuting Gaussian
source cancellation, exact compact Haar-derivative identities, the
contact-support bookkeeping, the whole-vector Gram synthesis, Fourier
block extraction, and the scalar logarithmic-moment inequalities used
above. 311 exact assertions passed; the family counts
and their scope are in \path{ym_v7_verification_report.json}.
Finite algebraic fixtures do not verify the analytic likelihood theorem,
Gaussian log-Sobolev theorem, or all native derivative suprema. They are
not proof-assistant verification and do not certify a physical mass gap.
The checker has 11,803 bytes and SHA--256
\begin{center}\scriptsize\ttfamily
EAA753817261B9830D904CC07586C2D4809BE46900B065AB91F063303032324F.
\end{center}

The cumulative V6 predecessor has 196,992 bytes and SHA--256
\begin{center}\scriptsize\ttfamily
CCB21705950AFCF7A33958811670B94E1E959459974DEF7E7D640F5ABA1B3D2D.
\end{center}
Only its title version changes and this exact appendix is inserted before
its final \verb|\end{document}|. Removing the appendix and restoring the
title recovers the supplied V6 byte for byte. Original and earlier
sources are checked against the input manifest. The supplied V4, V5 and V6 diagnostics were rerun unchanged and passed
9,467, 89,339 and 906 assertions respectively.
Unavailable historical scripts are not represented as rerun. This is not
the scheduled V10 companion checkpoint; only the supplied background
sources were available for the present targeted audit.

\begin{firewall}
V7 constructs exact Haar-compensated common-source derivatives and pays
their complete-native residual-energy and positive-Gram densities. It
bounds the entire nonadjacent-time common-source Hessian and the original
full/core correction in trace norm per cube and time slice, with all lags
included at once. A logarithmically growing finite spatial lift gives the
same rate. No volume-independent operator norm of that whole memory,
uniform arbitrary-history bound for every individual source, time-row
summability, native functional inequality, gauge-invariant bridge-tail
coverage, full-transfer spectral gap, interacting continuum theory, or
positive physical Yang--Mills mass gap is established here.
\end{firewall}
% END F16 V7 APPEND
% BEGIN F16 V8 APPEND -- 2026-09-17
% Insert immediately before the final \end{document} in cumulative V7.
% The predecessor body is unchanged; the cumulative title becomes V8.

\clearpage
\section{V8: a native fast-transfer estimate and time-summable coherent memory}
\label{f16v8:context}

V7's small trace-norm density permits a coherent exceptional direction.
Repeating its likelihood-to-energy argument cannot exclude that direction.
This appendix instead controls the \emph{actual propagation between two
corrected source insertions}. It obtains a small covariance operator and
geometric temporal decay, not merely a trace budget. There is a substantive
restriction: the bridge background is constant in time, and the spatial
cost is exponential in the number of cubes. That cost is displayed and
allows a logarithmically growing spatial volume at weak coupling. The
result is uniform in the entire history length, including the two original
dressed ends.

The argument is not a proof of the physical transfer gap. Its transfer
keeps the bridges fixed at every time and propagates only the complete
compact internal and temporal-port variables. No physical Gauss projection,
bridge integration, or physical time calibration is supplied by this
conditional construction.

\subsection{Audit and the restriction that makes a new estimate possible}

The supplied V7 appendix was read in full. The f14/f15 archives and Grand
Tensor Monograph V076 were searched through their result inventories and
relevant passages, rather than their whole proof collections recertified.
In particular, f15 V96 already proves fixed-region interacting spectral
convergence; f15 V97 already gives compact rooted pins and the full cube
oscillator spectrum. Finite-region Perron separation and Mehler
diagonalization are therefore \emph{not} being offered as new general
results. The monograph already distinguishes a conditional response from
the physical Perron law and forbids charging the same covariance twice.

The new object is the full assembled \emph{fixed-bridge fast transfer},
with all mixed plaquettes and all temporal port integrations present.
The new output is its quantitative, two-insertion estimate for V7's
Haar-compensated common sources under the actual dressed-end law. The
complete-chart error, the exact native top value, the dressed-end overlap,
and the source-insertion norms are all paid before taking arbitrarily
large time separation. None of V5--V7's extensive entropy estimates is
used to assert a native functional inequality.

Fix one bridge configuration in rescaled coordinates,
\begin{equation}
 \bar y\in\mathbb R^{36B},\qquad
 \max_e|\bar y_e|\le r<\infty,\qquad
 \bar Z_e=\operatorname{Exp}(\bar y_e/\sqrt\gamma),
 \qquad \mathbf y=(\bar y,\ldots,\bar y).
 \label{f16v8:static-history}
\end{equation}
There is no spatial homogeneity assumption on $\bar y$. All bridge
\emph{derivatives} are taken with respect to the independent time-slice
coordinates first, and only then evaluated at
\eqref{f16v8:static-history}. We are not differentiating only along the
diagonal submanifold of constant histories. Put
\[
 B=m^3,\quad m\ge4,\qquad n\ge1,\qquad \mathcal N=B(n+1).
\]
Constants may depend on $r$ and the fixed cube geometry. Dependence on
$B$ is never hidden in a volume-independent constant below. In particular
an estimate of the form $e^{c_rB}$ is not described as uniform in spatial
volume.

\subsection{The exact frozen-bridge transfer and its dressed-end law}

Use V5's unmoving compact temporal variables $h$, and write
$X=\operatorname{Exp}(x/\sqrt\gamma)$,
$h=\operatorname{Exp}(a/\sqrt\gamma)$ on their complete principal charts
$\mathcal P_x$, $\mathcal P_a$. The dimensions are $15B$ and $21B$.
At an internal tree edge set $X_e=I$; at a chord take its actual chord
variable. Define the original fine spatial representative
\[
 U_e(x;\bar y)=
 \begin{cases}
 X_e,&e\text{ internal},\\
 f_{s(e)}(X)\bar Z_e f_{t(e)}(X)^{-1},&e\text{ a bridge}.
 \end{cases}
\]
The moving frames $f_v$ are exactly those of V1--V7. Set
\begin{align}
 V_{\gamma,\bar y}(x)
   &=\gamma\{S_{\rm int}(X)+
                   S_\times(X,\Theta_X^{-1}\bar Z)\},\notag\\
 W_{\gamma,\bar y}(x,x',a)
   &=\gamma\sum_{e\ {\rm fine}}
       s\bigl(h_{s(e)}^{-1}U_e(x;\bar y)
                          h_{t(e)}U_e(x';\bar y)^{-1}\bigr),\notag\\
 \mathfrak E_{\gamma,\bar y}(x,x',a)
   &=\tfrac12 V_{\gamma,\bar y}(x)
         +W_{\gamma,\bar y}(x,x',a)
                          +\tfrac12 V_{\gamma,\bar y}(x').
 \label{f16v8:one-step-action}
\end{align}
Root entries of $h$ remain the identity. The sum has all $24B$ fine
spatial links; the two magnetic ends contain all $24B$ spatial plaquettes.
Let $j(v)=(\sin|v|/|v|)^2$ and abbreviate the products of these factors by
\[
 j_x(x)=\prod_{q\ {\rm chord}}j(x_q/\sqrt\gamma),\quad
 j_a(a)=\prod_{p\ {\rm nonroot}}j(a_p/\sqrt\gamma),\quad
 j_{\rm br}(\bar y)=\prod_{e\ {\rm bridge}}j(\bar y_e/\sqrt\gamma).
\]
The exact scalar for one step in internal and bridge half-density
coordinates is
\begin{equation}
 \rho_{\gamma,\bar y}
    =\left(\frac{c_{H,\gamma}}{z_\gamma}\right)^{24B}
                                    j_{\rm br}(\bar y),\qquad
 c_{H,\gamma}=(2\pi^2\gamma^{3/2})^{-1}.
 \label{f16v8:rho}
\end{equation}
Define an augmented step density and its integrated kernel by
\begin{align}
 d\mathfrak K_{\gamma,\bar y}(x,x';a)
   &=\rho_{\gamma,\bar y}
      1_{\mathcal P_x}(x)1_{\mathcal P_x}(x')1_{\mathcal P_a}(a)
      \sqrt{j_x(x)j_x(x')}j_a(a)
                   e^{-\mathfrak E_{\gamma,\bar y}(x,x',a)}\,da,\notag\\
 K_{\gamma,\bar y}(x,x')&=\int d\mathfrak K_{\gamma,\bar y}(x,x';a).
 \label{f16v8:kernel}
\end{align}
The operator $K_{\gamma,\bar y}$ acts on $L^2(\mathbb R^{15B},dx)$
with zero extension outside $\mathcal P_x$. The same letter for kernel
and operator will cause no ambiguity. Its endpoint vector is
\begin{equation}
 a_{\gamma,\bar y}(x)=
 \frac{1_{\mathcal P_x}(x)\Omega(x)
       e^{-\gamma S_\times(X,\Theta_X^{-1}\bar Z)/2}}
      {\sqrt{p_{\gamma,B}F_{\gamma,2}(\bar Z)}}.
 \label{f16v8:endpoint}
\end{equation}
All three quantities $\Omega,p_{\gamma,B},F_{\gamma,2}$ have their
inherited definitions. In particular the vector has norm one; it is not
replaced by a newly chosen native vacuum.

\begin{proposition}[An exact conditional transfer, not a compression]
\label{f16v8:exact-transfer}
The kernel \eqref{f16v8:kernel} is symmetric, positive semidefinite as
an operator, Hilbert--Schmidt, and strictly positive on the interiors of
its two principal internal charts. It has a simple positive top
eigenvalue $\lambda_{\gamma,\bar y}$ and a normalized positive top vector
$\phi_{\gamma,\bar y}$ on those charts. Moreover
\begin{equation}
 \mathcal W_{\gamma,n}^{\rm full}(\bar y,\ldots,\bar y)
   =\langle a_{\gamma,\bar y},
               K_{\gamma,\bar y}^{\,n}a_{\gamma,\bar y}\rangle.
 \label{f16v8:exact-power}
\end{equation}
Normalizing the augmented path integral on the right gives precisely
V4--V7's complete native fast-history conditional at
\eqref{f16v8:static-history}. There is no projection at an internal time.
\end{proposition}

\begin{proof}
Start with V2's exact full moving-frame kernel and perform V5's
conditional Haar change $h=f(X)r f(X')^{-1}$. It gives the kinetic
words in \eqref{f16v8:one-step-action}. At the two ends, internal
half-density factors contribute $c_{H,\gamma}^{5B}$ in total,
bridge half-densities contribute $c_{H,\gamma}^{12B}$, and the port
Haar integral contributes $c_{H,\gamma}^{7B}$. There are $24B$
kinetic normalizers $z_\gamma^{-1}$. This gives
\eqref{f16v8:rho} and all the nonconstant Haar factors in
\eqref{f16v8:kernel}. The two equal bridge half-densities give
$j_{\rm br}(\bar y)$, not its square.

Swapping $x,x'$ and replacing every $h_v$ by $h_v^{-1}$ leaves each
kinetic trace unchanged, by inversion and cyclicity; Haar measure is
invariant under this replacement. Thus the kernel is symmetric.
Its operator positivity is inherited from the full positive Wilson
transfer before the remaining root restriction, but requires a small
justification because evaluation at a bridge value is not an $L^2$
projection. In that full kernel test against
$u(x)\eta_\varepsilon(y-\bar y)$, with $u$ compactly supported inside
the internal charts and with a nonnegative bridge approximate identity
of integral one. The test has finite $L^2$ norm for each $\varepsilon$.
Positivity gives a nonnegative quadratic form. Local continuity of the
exact kernel lets $\varepsilon\downarrow0$ and yields
$\langle u,K_{\gamma,\bar y}u\rangle\ge0$. The kernel bound proved
below extends the assertion by density to all $L^2$.

For clarity the bound already follows from the inherited global pins.
The two internal magnetic ends supply the five-chord cube pins, and the
seven tree kinetic terms supply the independent rooted $h$ pin. Thus
\begin{equation}
 \mathfrak E_{\gamma,\bar y}(x,x',a)
       \ge a_*\bigl(|x|^2+|x'|^2+|a|^2\bigr),
             \qquad a_*=(6\pi^2)^{-1},
 \label{f16v8:step-pin}
\end{equation}
on the complete charts. All other terms are nonnegative. With every
Haar factor at most one and a finite prefactor, integrating $a$ gives
an integrable Gaussian square majorant in $x,x'$. This proves the
Hilbert--Schmidt assertion. Strict positivity follows directly from
the positive integrand on the open charts.

The usual compact positive-kernel argument can be given here directly.
The positive compact self-adjoint operator has a maximizing eigenvector.
Replacing a real eigenvector by its absolute value cannot decrease its
Rayleigh quotient, and increases it strictly if it has positive and
negative parts of positive measure. A nonnegative top eigenvector is
strictly positive by its integral equation. A second linearly independent
top eigenvector could be made orthogonal to it, and would change sign,
contradicting the strict inequality. This proves simplicity. The zero
extension adds zero spectrum, not another top eigenvector.

Finally expand the power in \eqref{f16v8:exact-power} using the augmented
kernel. An interior internal slice has a full Haar factor and two
magnetic half-weights. Each external slice has one internal half-weight
and the endpoint vector. The latter supplies the other mixed magnetic
half-weight and $\Omega/\sqrt{p_{\gamma,B}F_{\gamma,2}}$. The products
are exactly V4's action powers and normalizers, including its bridge
Jacobian power $j_{\rm br}(\bar y)^n$. This proves the identity and
the equality of normalized conditional laws.
\end{proof}

The positivity argument is the standard symmetric positive-kernel
Perron argument; see also R.~Jentzsch,
\href{https://doi.org/10.1515/crll.1912.141.235}
{\emph{\"Uber Integralgleichungen mit positivem Kern}} (1912).
The operator here is an exact bridge-diagonal kernel. It is not
$\mathcal I_\gamma^*T_{\rm full,\gamma}\mathcal I_\gamma$, and its
Perron value is not identified with the full physical transfer's top
value.

\subsection{The reference spectrum and the complete-chart error}

Suppress colour identities and put
\begin{align}
 S&=\mathsf H_B+A^{\mathsf T}A,&
 \Lambda&=A^{\mathsf T}D,&
 \xi_{\bar y}&=-S^{-1}\Lambda\bar y,\notag\\
 c(\bar y)&=\bar y^{\mathsf T}
            (D^{\mathsf T}D-\Lambda^{\mathsf T}S^{-1}\Lambda)\bar y,&
 T_S&=\mathsf M_B^{1/2}S\mathsf M_B^{1/2}.
 \label{f16v8:Gaussian-data}
\end{align}
The minimum of $x^{\mathsf T}\mathsf H_Bx+|Ax+D\bar y|^2$ is
$c(\bar y)\ge0$. V2's complete Gaussian kinetic integration gives
\begin{equation}
 K_{0,\bar y}(x,x')
  =\varphi_G(0)e^{-c(\bar y)/2}
    e^{-(x-\xi_{\bar y})^{\mathsf T}S(x-\xi_{\bar y})/4}
    \varphi_{\mathsf M_B}(x-x')
    e^{-(x'-\xi_{\bar y})^{\mathsf T}S(x'-\xi_{\bar y})/4}.
 \label{f16v8:Gaussian-kernel}
\end{equation}
The three-colour densities $\varphi$ are the ones already fixed in V2.
The reference endpoint is
\[
 a_{0,\bar y}(x)
       =F_2(\bar y)^{-1/2}\Omega(x)e^{-|Ax+D\bar y|^2/4}.
\]
Let $\phi_{0,\bar y}$ be the normalized positive ground vector of
\eqref{f16v8:Gaussian-kernel}.

\begin{proposition}[Quantitative frozen-background Gaussian comparison]
\label{f16v8:comparison}
There are finite constants $c_r,C_r\ge1$, independent of $B$, such that
for $\gamma\ge C_r$ and every admitted $\bar y$,
\begin{align}
 e^{-c_r B}&\le\lambda_{0,\bar y}\le e^{c_rB},&
 \frac{\lambda_1(K_{0,\bar y})}{\lambda_{0,\bar y}}
     &\le q_0:=3-2\sqrt2<\frac15,\notag\\
 \|K_{\gamma,\bar y}-K_{0,\bar y}\|_{\rm HS}
     &\le \frac{e^{c_rB}}{\sqrt\gamma},&
 \|a_{\gamma,\bar y}-a_{0,\bar y}\|_2
     &\le \frac{e^{c_rB}}{\sqrt\gamma},\notag\\
 \langle a_{0,\bar y},\phi_{0,\bar y}\rangle
     &\ge e^{-c_rB}.&&
 \label{f16v8:comparison-bounds}
\end{align}
The Gaussian top is, more precisely,
\begin{equation}
 \lambda_{0,\bar y}=\varphi_G(0)e^{-c(\bar y)/2}
             \prod_{j=1}^{5B}r(t_j)^{3/2},\qquad
 r(t)=\bigl(1+t/2+\sqrt{t+t^2/4}\bigr)^{-1},
 \label{f16v8:Gaussian-top}
\end{equation}
where $t_j$ are the eigenvalues of $T_S$ and $4\le t_j\le10$.
The errors include the complete compact complements, not a fixed-window
restriction of the internal variables.
\end{proposition}

\begin{proof}
The inherited whitened bounds for $\mathsf H_B$ and $A^{\mathsf T}A$
give $4I\le T_S\le10I$. Translate by $\xi_{\bar y}$, whiten by
$\mathsf M_B$, and diagonalize $T_S$. The scalar kernel of a mode of
parameter $t$ is
\[
 (2\pi)^{-1/2}
       \exp\{-\tfrac12(1+t/2)(u^2+v^2)+uv\}.
\]
Writing $a=1+t/2$ and $w=\sqrt{a^2-1}$, its ground precision is $w$,
its top eigenvalue is $(a+w)^{-1/2}=r(t)^{1/2}$, and its successive
Hermite eigenvalues have ratio $r(t)$. This is the inherited complete
Mehler calculation, with the general mixed magnetic matrix $S$ now
retained. For reference the identity is
\href{https://dlmf.nist.gov/18.18.E28}{NIST DLMF, equation 18.18.28}.
It proves \eqref{f16v8:Gaussian-top} and the ratio bound without deleting
non-singlet internal excitations.
Since $I\le G\le30I$, $c(\bar y)\le\|D\bar y\|^2\le CBr^2$,
and every $r(t_j)$ lies in the fixed interval $[r(10),r(4)]$, the
displayed exponential bounds on the top follow.

Here are details of the dimensional error. Write $w=(x,x',a)$, with
$51B$ real coordinates, and let $\mathfrak E_0$ be the quadratic limit of
\eqref{f16v8:one-step-action}. Both exponents obey
\eqref{f16v8:step-pin}, the second by taking the quadratic limit of the
compact anchors. Bounded-incidence differentiation of the actual
unitary exponential words gives, throughout the complete native charts,
\begin{equation}
 |\mathfrak E_{\gamma,\bar y}-\mathfrak E_{0,\bar y}|
     \le \frac{C}{\sqrt\gamma}(|w|^3+Br^3).
 \label{f16v8:global-error}
\end{equation}
Indeed each word has uniformly bounded third derivatives in its fixed
local inputs before rescaling. The sum of the resulting local cubic
norms is bounded by $C|w|^3+CBr^3$. No logarithm of a group product is
differentiated. The Haar products obey
\[
 0\le\sqrt{j_x(x)j_x(x')}j_a(a)\le1,\qquad
 1-\sqrt{j_x(x)j_x(x')}j_a(a)\le C|w|^2/\gamma.
\]
This uses $1-\sin u/u\le u^2/6$ and $0\le\sin u/u\le1$ on
$[0,\pi]$, followed by $1-\prod b_j\le\sum(1-b_j)$.
V4's one-link normalization also gives
\[
 \rho_{0}=(2\pi)^{-36B},\qquad
 \left|\log(\rho_{\gamma,\bar y}/\rho_0)\right|
                                      \le C_r B/\gamma.
\]
The bound on $j_{\rm br}(\bar y)$ uses its small fixed bridge chart;
the internal charts have not been shrunk.

Use $|e^{-u}-e^{-v}|\le|u-v|e^{-\min(u,v)}$ and the global pins.
On the common chart domain the difference of augmented integrands is
bounded by a fixed-degree polynomial in $|w|,B$, times
$e^{C_rB}\gamma^{-1/2}e^{-a_*|w|^2}$. The Gaussian complement of a
principal chart has some group coordinate of norm at least
$\pi\sqrt\gamma$. Its integral is bounded by
$e^{C_rB-c\gamma}$, and the same reasoning bounds its contribution to
the kernel's Hilbert--Schmidt norm. For example, split off half the
Gaussian exponent to pay that coordinate threshold and integrate the
other half over all $51B$ coordinates. This incurs an exponential in
$B$, not an unrecorded uniform constant.
Integrate $a$ and then take the $L^2(dx\,dx')$ norm. Fixed-degree
Gaussian moments in dimensions proportional to $B$ are a fixed
polynomial in $B$ times a constant to the power $B$. Such polynomials
are absorbed into $e^{c_rB}$. Also
$e^{-c\gamma}\le C\gamma^{-1/2}$. This proves the kernel error on the
whole zero-extended space.

For the endpoint, first compare the unnormalized vectors
\[
 q_\gamma=1_{\mathcal P_x}\Omega
           e^{-\gamma S_\times/2},\qquad
 q_0=\Omega e^{-|Ax+D\bar y|^2/4}.
\]
Their multipliers are at most one. The same cubic word remainder,
now integrated against $\Omega^2$, and its complete-chart Gaussian
tail give $\|q_\gamma-q_0\|_2\le e^{c_rB}/\sqrt\gamma$.
The exact identities are $\|q_\gamma\|_2^2=p_{\gamma,B}
F_{\gamma,2}(\bar Z)$ and $\|q_0\|_2^2=F_2(\bar y)$.
The latter is bounded below by $e^{-c_rB}$, either from its inherited
determinant formula or by integrating on a product of fixed one-cell
balls. Normalization uses the elementary inequality
\[
 \left\|\frac u{\|u\|}-\frac v{\|v\|}\right\|
                         \le\frac{2\|u-v\|}{\|v\|}
       \quad(u,v\ne0).
\]
This proves the endpoint bound with an enlarged exponent.

Finally both $a_{0,\bar y}$ and $\phi_{0,\bar y}$ are normalized
positive Gaussian amplitudes. Their squared-density covariance
spectra lie between two fixed positive constants, and the squared
norms of their means are at most $C_rB$. For $a_0$ use
$(C_B^{-1}+A^{\mathsf T}A)^{-1}$; for $\phi_0$ use
\[
 C_S=\mathsf M_B^{1/2}
           [2\sqrt{T_S+T_S^2/4}]^{-1}\mathsf M_B^{1/2},
       \qquad \text{mean }\xi_{\bar y}.
\]
Their product, integrated over $[-1,1]^{15B}$, is at least
$e^{-c_rB}$: the two determinant normalizers are exponential in $B$,
and both quadratic exponents are bounded there by $C_rB$. This proves
the overlap lower bound without asserting that the two amplitudes
coincide.
\end{proof}

\subsection{A native spectral ratio and an overlap valid for every duration}

\begin{theorem}[Frozen-background fast propagation with a paid spatial cost]
\label{f16v8:native-propagation}
There are $k_r\ge1$ and $\gamma_r$ such that, whenever
\begin{equation}
 \gamma\ge\gamma_r,\qquad
             \frac{e^{k_rB}}{\sqrt\gamma}\le\frac1{100},
 \label{f16v8:entry-regime}
\end{equation}
the exact operator and endpoint in \eqref{f16v8:kernel}--\eqref{f16v8:endpoint}
satisfy
\begin{equation}
 S_\gamma:=K_{\gamma,\bar y}/\lambda_{\gamma,\bar y}
       =P_\gamma+Q_\gamma,\qquad
 P_\gamma=|\phi_{\gamma,\bar y}\rangle
                 \langle\phi_{\gamma,\bar y}|,\qquad
 0\le Q_\gamma,\quad \|Q_\gamma\|\le\tau:=\frac13,
 \label{f16v8:spectral-split}
\end{equation}
with $P_\gamma Q_\gamma=0$. Moreover, for another fixed $b_r$,
\begin{equation}
 \beta_\gamma:=\langle a_{\gamma,\bar y},\phi_{\gamma,\bar y}\rangle
       \ge e^{-b_rB},\qquad
 \beta_\gamma^2\le
       Z_{\gamma,n}^{\rm norm}:=
          \langle a_{\gamma,\bar y},S_\gamma^n a_{\gamma,\bar y}\rangle
       \le1\quad(n\ge1).
 \label{f16v8:boundary-overlap}
\end{equation}
No factor depending on $n$ occurs in these bounds. Normalization always
uses the \emph{exact} $\lambda_{\gamma,\bar y}$.
\end{theorem}

\begin{proof}
The preceding proposition gives
\[
 \delta_{\gamma,B}:=
 \frac{\|K_{\gamma,\bar y}-K_{0,\bar y}\|}
                         {\lambda_{0,\bar y}}
              \le e^{c_rB}/\sqrt\gamma.
\]
Enlarge $k_r$ so that \eqref{f16v8:entry-regime} implies
$\delta_{\gamma,B}\le1/16$. The variational characterization of the
first two eigenvalues gives
\[
 (1-\delta_{\gamma,B})\lambda_{0,\bar y}
           \le\lambda_{\gamma,\bar y},\qquad
 \lambda_1(K_{\gamma,\bar y})
           \le(q_0+\delta_{\gamma,B})\lambda_{0,\bar y}.
\]
Thus their ratio is at most
$(1/5+1/16)/(1-1/16)=7/25<1/3$. This controls the complete compact
fast operator, including its excited states; it is not a band-compression
calculation. Operator positivity gives the nonnegative spectral split.

We must also pay the dressed ends. Let $P_0$ be the Gaussian top
projection. Projecting the native eigenvector equation onto $I-P_0$
gives
\[
 \|(I-P_0)\phi_{\gamma,\bar y}\|
       \le\frac{\delta_{\gamma,B}}{1-q_0-\delta_{\gamma,B}}.
\]
The two positive top eigenvectors have nonnegative inner product, so
$\|\phi_{\gamma,\bar y}-\phi_{0,\bar y}\|\le3\delta_{\gamma,B}$
in this regime. Combine this with the endpoint comparison and the
Gaussian overlap lower bound. Enlarging $k_r$ again makes the sum of
the two approximation errors at most half that lower bound, yielding
$\beta_\gamma\ge e^{-b_rB}$ for some $b_r$. Such a choice is possible
because all three costs are explicit exponentials in $B$ times
$\gamma^{-1/2}$.

Finally $S_\gamma^n=P_\gamma+Q_\gamma^n$ and $Q_\gamma^n\ge0$.
The lower bound in \eqref{f16v8:boundary-overlap} follows by retaining
the Perron term; the upper bound uses $\|S_\gamma\|=1$ and
$\|a_{\gamma,\bar y}\|=1$. No replacement of
$\lambda_{\gamma,\bar y}^n$ by $\lambda_{0,\bar y}^n$ has occurred.
Such a replacement would in general accumulate an error with duration.
\end{proof}

\subsection{Unbounded corrected scores as bounded transfer insertions}

Use V7's \emph{full} Haar lift $\mathsf P$ and its residuals
$r^{[\infty]}_\alpha$, with no new subtraction. A spatial common variation
$v\in\mathcal E$ at time $i$ gives
\[
 R_i(v)=\sum_{b,\mu,k}v_{b,\mu,k}
                     r^{[\infty]}_{i,b,\mu,k}.
\]
Coefficients refer to V7's orthonormal common-bridge basis, so their
Euclidean norm is $\|v\|$. Arbitrary spatially coherent $v$ are allowed.
For the current purpose the infinite lift means the exact finite-volume
inverse $L_h^{-1}$; it is bounded independently of $B$.

The variable $R_i(v)$ uses only the temporal window
\begin{equation}
 a_i=\max(0,i-1),\qquad b_i=\min(n,i+1),\qquad p_i=b_i-a_i\in\{1,2\}.
 \label{f16v8:source-window}
\end{equation}
It depends on the internal slices in that window and on its $p_i$ port
bonds. Its spatial dependence need not be confined to one cube.
Let $\mathbb A_i(v)$ be the kernel obtained by multiplying the product
of those $p_i$ augmented step densities by $R_i(v)$, then integrating all
ports and intermediate internal slices. Let $\mathbb B_i(v)$ use
$R_i(v)^2$ instead. These are signed and nonnegative-kernel insertion
operators, respectively; a nonnegative kernel is not asserted to be a
positive semidefinite operator. Define their exact normalized versions
\[
 A_i(v)=\lambda_{\gamma,\bar y}^{-p_i}\mathbb A_i(v),\qquad
 B_i(v)=\lambda_{\gamma,\bar y}^{-p_i}\mathbb B_i(v).
\]

\begin{proposition}[Two-insertion norms on the full compact domains]
\label{f16v8:insertions}
After increasing constants in \eqref{f16v8:entry-regime} if needed,
there are fixed $c_r$ such that, for every $n,i$ and $v\in\mathcal E$,
\begin{equation}
 \|A_i(v)\|\le\frac{e^{c_rB}}{\sqrt\gamma}\|v\|,\qquad
 \|B_i(v)\|\le\frac{e^{c_rB}}{\gamma}\|v\|^2.
 \label{f16v8:insertion-norms}
\end{equation}
The constants do not depend on the source time or history length. These
are full native insertion estimates, not estimates of observables under
a Gaussian probability substituted for the native law.
\end{proposition}

\begin{proof}
The additional point beyond V7's Gaussian energy budget is that a source
is now integrated against the complete \emph{step kernel pin}. We give
the pointwise and kernel estimates explicitly.

Separate one common-source derivative into its magnetic contribution at
$i$ and its kinetic contributions on bonds $i-1,i$ when present. At
Gaussian order the common magnetic fast coefficient vanishes by
$A^{\mathsf T}Dv=0$. On each port bond the exact lift cancels the
Gaussian fast coefficient separately, as in V7's kinetic square. At a
constant history the deterministic kinetic contribution is zero;
the deterministic magnetic contribution is precisely the corresponding
part of $g_\alpha(\mathbf y)$. Thus each of the three pieces of $R_i$
is a primitive nonlinear remainder or the bounded $L_h^{-1}E^{\mathsf T}$
synthesis of primitive port remainders.

For a local primitive component $e_j$, V7's complete-chart insertion
calculation gives
\[
 |e_j|\le C_r\gamma^{-1/2}(1+|w_{C_j}|^2).
\]
Here $w$ stacks the at most three internal slices and two port slices of
\eqref{f16v8:source-window}; its dimension is at most $87B$.
Each $C_j$ has fixed size and each primitive input occurs a bounded
number of times. Consequently
\[
 \sum_j|e_j|^2\le\frac{C_r}{\gamma}(B+|w|^4).
\]
For example the sum of squared local squared norms is bounded by their
sum times their maximum, and hence by $C|w|^4$. Applying the bounded
spatial lift and combining the at most three pieces proves the global
pointwise estimate
\begin{equation}
 |R_i(v)|\le\frac{C_r}{\sqrt\gamma}
                          (\sqrt B+|w|^2)\|v\|.
 \label{f16v8:window-remainder}
\end{equation}
No sum of absolute lift coefficients or support-volume factor is needed;
only its $\ell^2$ operator norm is used. The formula comes from group
insertions, so it remains valid across a principal port cut in its
compact interpretation. The analytic exponential-word expressions supply
the same polynomial majorant on all real inputs.

For $p_i=1$ or $2$, the product of augmented step densities is at most
\[
 e^{C_rB}\exp\{-a_*|w|^2\}
\]
with respect to its port and intermediate-slice Lebesgue measures, with
the endpoint internal coordinates left as kernel variables. This follows
by applying \eqref{f16v8:step-pin} to each step; interior slice pins can
be weakened to the displayed common lower bound. Every Haar factor is
at most one. The at most two factors $\rho_{\gamma,\bar y}$ cost only
$e^{C_rB}$.
Multiplication by \eqref{f16v8:window-remainder}, or by its square,
therefore leaves a fixed-degree polynomial times this Gaussian majorant.
Integrate the complete native port charts and intermediate internal
charts by enlarging them to Euclidean space in the majorant, and then
take the kernel Hilbert--Schmidt norm in its two endpoint variables.
The resulting bounds are
$e^{c_rB}\gamma^{-1/2}\|v\|$ and
$e^{c_rB}\gamma^{-1}\|v\|^2$ for the unnormalized insertions.
As in the one-step comparison, Gaussian moments of fixed degree and
dimension at most $87B$ cost an exponential in $B$. There is no time
length in this calculation.
The native top lower bound
$\lambda_{\gamma,\bar y}\ge\tfrac{15}{16}\lambda_{0,\bar y}
\ge e^{-C_rB}$ pays division by its first or second power. This proves
\eqref{f16v8:insertion-norms}. The argument integrates the entire native
domain; it has not thrown away a source-weighted large-field complement.
\end{proof}

\subsection{A two-window covariance identity with the boundary cost retained}

The following elementary transfer argument is recorded because using only
the stationary Perron law would omit the actual dressed-end cost. It
applies as well to observables of integrated auxiliary bond variables.

\begin{proposition}[Finite-history mixing from the exact normalized transfer]
\label{f16v8:two-window}
Let $S=P+Q$ be a positive self-adjoint contraction with positive kernel,
$P=|\phi\rangle\langle\phi|$, $PQ=0$, and
$0\le Q$, $\|Q\|\le\tau<1$. Let $a\ge0$, $\|a\|=1$, and
$\beta=\langle a,\phi\rangle>0$. Use the normalized length-$n$ path
law with these ends. Suppose real observables $F,G$ are inserted in
nonoverlapping bond windows, separated by $g\ge0$ unmarked steps,
and their normalized insertion operators are $A_F,A_G$. Then
\begin{equation}
 |\operatorname{Cov}(F,G)|
       \le4\beta^{-4}\|A_F\|\|A_G\|\tau^g.
 \label{f16v8:two-window-bound}
\end{equation}
For a single observable with square-insertion operator $B_F$,
\begin{equation}
             \mathbb E F^2\le\beta^{-2}\|B_F\|.
 \label{f16v8:one-window-square}
\end{equation}
No stationarity of the finite path law is required.
\end{proposition}

\begin{proof}
Write the left window as $[a_1,b_1]$, the right as $[a_2,b_2]$,
$p=b_1-a_1$, $q=b_2-a_2$, $g=a_2-b_1\ge0$. Set
\[
 u=S^{a_1}a,\quad v=S^{n-b_2}a,\quad
 u_0=S^pu,\quad u_F=A_F^*u,\quad
 v_0=S^qv,\quad v_G=A_Gv,\quad T=S^g.
\]
Then $\|u_0\|,\|v_0\|\le1$,
$\|u_F\|\le\|A_F\|$, and $\|v_G\|\le\|A_G\|$.
The common denominator is
$Z=\langle u_0,Tv_0\rangle=\langle a,S^na\rangle\ge\beta^2$.
The exact centered covariance is
\begin{equation}
 \operatorname{Cov}(F,G)=
 \frac{\langle u_F,Tv_G\rangle\langle u_0,Tv_0\rangle
       -\langle u_F,Tv_0\rangle\langle u_0,Tv_G\rangle}{Z^2}.
 \label{f16v8:covariance-minor}
\end{equation}
Replacing $T$ by the rank-one $P$ makes its numerator zero exactly.
For $g\ge1$, $\|T-P\|=\|Q^g\|\le\tau^g$; for $g=0$ use
$T=I$ and $\|I-P\|\le1=\tau^0$. Each of the two products in the
numerator changes by at most
$2\tau^g\|A_F\|\|A_G\|$, since $\|T\|,\|P\|\le1$.
Division by $Z^2\ge\beta^4$ proves
\eqref{f16v8:two-window-bound}. This cancellation includes the means
under the actual finite path law, not means borrowed from its stationary
limit.
For the square insertion, its expectation has numerator
$\langle a,S^{a_1}B_FS^{n-b_1}a\rangle$, bounded by $\|B_F\|$,
and denominator at least $\beta^2$. This proves
\eqref{f16v8:one-window-square}. Auxiliary bond integrations are part of
the inserted kernels and satisfy the same composition identities.
\end{proof}

\subsection{A coherent common-source bound and an absolutely summable time row}

Let $\Gamma_\gamma=\Gamma^{[\infty]}$ be exactly V7's full corrected
source Gram, evaluated at \eqref{f16v8:static-history}. Regard its time
blocks as operators on $\mathcal E$; there are $9B$ coordinates per block.
No spatial average or trace normalization is made in the next result.

\begin{theorem}[Complete-native coherent covariance on static backgrounds]
\label{f16v8:coherent-covariance}
There are constants $K_r\ge1$, $\gamma_r$ such that, for every $B=m^3$,
$m\ge4$, every $n\ge1$, every $\bar y$ in
\eqref{f16v8:static-history}, and
\begin{equation}
       \gamma\ge\gamma_r,\qquad
                  e^{K_rB}/\sqrt\gamma\le1/100,
 \label{f16v8:final-regime}
\end{equation}
one has, with $\tau=1/3$ and $A_{\gamma,B}=e^{K_rB}/\gamma$,
\begin{equation}
 \|\Gamma_{\gamma,ij}\|_{\mathcal E\to\mathcal E}
      \le A_{\gamma,B}\tau^{(|i-j|-2)_+},
              \qquad 0\le i,j\le n.
 \label{f16v8:covariance-blocks}
\end{equation}
In particular, for every common-source history vector
$\mathbf v=(v_0,\ldots,v_n)$,
\begin{equation}
 \boxed{
 \operatorname{Var}_{\mu_\gamma}
       \left(\sum_{i=0}^n R_i(v_i)\right)
           \le6A_{\gamma,B}\sum_{i=0}^n\|v_i\|^2.}
 \label{f16v8:coherent-bound}
\end{equation}
For every fixed $0\le\alpha<\log3$ there is also the weighted row bound
\begin{equation}
 \sup_i\sum_{j=0}^n e^{\alpha|i-j|}
                 \|\Gamma_{\gamma,ij}\|
  \le A_{\gamma,B}
       \left(1+2e^\alpha+
                    \frac{2e^{2\alpha}}{1-e^\alpha/3}\right).
 \label{f16v8:weighted-row}
\end{equation}
All duration dependence has been removed. Spatial dependence is the
explicit $e^{K_rB}$, not an unproved uniform constant.
\end{theorem}

\begin{proof}
Take unit spatial common variations $v,w$ at times $i,j$. If
$|i-j|\le1$, Cauchy--Schwarz, the square-insertion bound, and
\eqref{f16v8:boundary-overlap} give
\[
 |\operatorname{Cov}(R_i(v),R_j(w))|
                   \le\beta_\gamma^{-2}e^{c_rB}/\gamma.
\]
If $j-i\ge2$, their windows in \eqref{f16v8:source-window} have
disjoint bonds and gap $g=j-i-2$. They may share an endpoint internal
slice when $g=0$, which is allowed by
Proposition~\ref{f16v8:two-window}. Apply that proposition and the two
first-insertion norm bounds to obtain
\[
 |\operatorname{Cov}(R_i(v),R_j(w))|
       \le4\beta_\gamma^{-4}
               e^{2c_rB}\gamma^{-1}\tau^{j-i-2}.
\]
The same estimate holds in the reverse order by covariance symmetry.
Enlarge $K_r$ to absorb the finitely many constants and the powers of the
exponential overlap lower bound. Increasing it strengthens the entry
condition, so all the previous results remain valid. The estimates hold
for arbitrary unit spatial vectors, not just for one coordinate, and
therefore prove \eqref{f16v8:covariance-blocks} in block operator norm.

The scalar time majorant has row sum at most
\[
 1+4+2\sum_{d=3}^{\infty}\tau^{d-2}
                 =5+\frac{2\tau}{1-\tau}=6.
\]
Its symmetric Schur bound, or $2ab\le a^2+b^2$ applied to all time
pairs, proves \eqref{f16v8:coherent-bound}. Positivity of the full Gram
identifies its quadratic form with the displayed variance. Multiplying
the block estimate by $e^{\alpha|i-j|}$ and summing the corresponding
geometric series proves \eqref{f16v8:weighted-row}. There is no factor
$B(n+1)$ on the right of the coherent variance estimate.
\end{proof}

\begin{theorem}[Full history memory, original core correction, and time tails]
\label{f16v8:memory}
Under \eqref{f16v8:final-regime}, let $\mathsf M_\gamma$ and
$\mathsf C_{\gamma,R}$ be V7's actual nonadjacent-time Hessian matrices,
at the static background. Then
\begin{equation}
 \boxed{\quad \|\mathsf M_\gamma\|_{\rm op}
                         \le3A_{\gamma,B}.\quad}
 \label{f16v8:memory-op}
\end{equation}
More precisely, if $\mathsf M_\gamma^{>L}$ keeps only blocks with
$|i-j|>L$, for an integer $L\ge1$, then
\begin{equation}
 \|\mathsf M_\gamma^{>L}\|_{\rm op}
                         \le3A_{\gamma,B}(1/3)^{L-1}.
 \label{f16v8:memory-tail}
\end{equation}
For every radius $R$ obeying V4's derivative regime,
\begin{align}
 \|\mathsf C_{\gamma,R}\|_{\rm op}
        &\le3A_{\gamma,B}+C R^2/\gamma,\notag\\
 \|\mathsf C_{\gamma,R}^{>L}\|_{\rm op}
        &\le3A_{\gamma,B}(1/3)^{L-1}
                         +C(R^2/\gamma)(9/19)^{L-1}.
 \label{f16v8:correction-tail}
\end{align}
The constants are independent of $n$. The event, scalar factors, and
endpoint normalizers defining $p_{\rm good}$ are exactly V4's originals.
\end{theorem}

\begin{proof}
V7's full Haar-compensated identity, whose contact term has temporal range
one, gives
$\mathsf M_\gamma=-\operatorname{Off}_{>1}\Gamma_\gamma$ exactly.
At separation $d\ge2$ its block norm is at most
$A_{\gamma,B}\tau^{d-2}$. The scalar row sum of blocks with $d>L$ is
at most
\[
 2A_{\gamma,B}\sum_{d=L+1}^{\infty}\tau^{d-2}
          =\frac{2A_{\gamma,B}}{1-\tau}\tau^{L-1}
          =3A_{\gamma,B}\tau^{L-1}.
\]
This proves the first two claims by the same block Schur bound. Unlike
V7's trace-ideal extraction, this step uses actual separation decay.

V4's common-source core estimate acts on arbitrary spatial common
variations and gives a core block bound
$C(R^2/\gamma)(9/19)^{d-2}$ for $d\ge2$. Its time-tail row sum is
$C(R^2/\gamma)(9/19)^{L-1}$. Finally
\[
 \mathsf C_{\gamma,R}
   =\mathsf M_\gamma-
      \operatorname{Off}_{>1}\nabla_{\mathcal U}^2
                    [-\log\mathcal W_{\gamma,n}^{R}].
\]
Use the operator-norm triangle inequality, with the corresponding tail
mask when needed. The equality is the original full/core identity, not
a new covariance debit. External logarithmic factors depend on one time
only and have zero nonadjacent-time mixed derivatives; they have not
been dropped from either path integral.
\end{proof}

\subsection{What is gained, and the spatial limitation that remains}

For each fixed spatial $B$, the complete-native coherent covariance and
full memory bounds are $O_{r,B}(\gamma^{-1})$, with geometric time decay
and no history-length restriction. This improves the earlier fixed-lag
or zero-history bounds \emph{on the static-background subclass}; it does
not supersede V7's all-background, all-volume trace-density statement.
The present $\bar y$ may vary arbitrarily from one spatial bridge to
another inside its fixed rescaled window.

There is also a genuine growing-spatial-size consequence. With the same
$K_r$ as in the theorem, the regime
\begin{equation}
          B=m^3\le\frac{\log\gamma}{4K_r}
 \label{f16v8:mesoscopic}
\end{equation}
is admissible for sufficiently large $\gamma$, since the left side of
\eqref{f16v8:final-regime} is then at most $\gamma^{-1/4}$. It gives
$A_{\gamma,B}\le\gamma^{-3/4}$. Taking
$R=4\sqrt{\log\gamma}$ in V4's original core, all its fixed-$r$
conditions hold eventually and
\begin{equation}
 \|\Gamma_\gamma\|_{\rm op}
       +\|\mathsf M_\gamma\|_{\rm op}
       +\|\mathsf C_{\gamma,R}\|_{\rm op}
                   \le C_r\gamma^{-3/4},
 \label{f16v8:mesoscopic-conclusion}
\end{equation}
uniformly in \emph{all} $n\ge1$. The spatial side is therefore allowed
to grow like a constant times $(\log\gamma)^{1/3}$. The constants come
from conservative Gaussian integrals and derivative bounds; no useful
numerical entry coupling is claimed.

The source operator in \eqref{f16v8:coherent-bound} is the corrected
one from V7. The proof estimates every spatially coherent combination
within the specified volume, not a complement obtained by discarding
exceptional eigenvectors. All excited fast states and all compact port
tails remain in the transfer and its inserted kernels. In particular
there is no interchange of an uncontrolled $n\to\infty$ error with the
weak-field limit: the spectral split and boundary denominator hold at
each admitted finite $\gamma,B$ for every $n$.

Two restrictions are material. First, this proof uses one time-independent
kernel $K_{\gamma,\bar y}$ and its one Perron projection. An arbitrary
bounded history instead produces different consecutive kernels. Their
powers cannot be replaced by the spectral powers used here. This
appendix therefore gives no arbitrary-history temporal mixing theorem.
Second, the complete-chart and endpoint comparisons have paid costs
$e^{c_rB}$. They do not allow $B\to\infty$ at fixed $\gamma$, nor may
the small-volume covariance inequalities be tensorized across interacting
spatial regions. The fixed bridges are also not a gauge-invariant coverage
event for the integrated physical bridge law.

The next useful upstream target is a \emph{boundary-conditioned spatial
version} of the two-insertion estimate, preserving the exact normalization
while bounding how a spatial collar changes the source and endpoint
messages. The global internal pin supplies the present finite-volume
input, but arbitrary exterior fast fields can change the conditional
landscape, so the displayed torus estimate is not such a boundary-uniform
theorem. Separately, varying bridge histories require a normalized
inhomogeneous transfer estimate; a gap for each static operator alone
would not prove that estimate. These are concrete propagation problems,
not another request for scalar pressure or an assumption of native
functional inequalities. The full-transfer top comparison, physical time
calibration, infrared coercivity of the retained bridge dynamics, and an
interacting continuum construction remain outside the proved claims.

\clearpage
\subsection{Diagnostics, preservation, and scope}

The accompanying diagnostic is
\begin{center}\small
\path{verify_ym_f16_v8_static_transfer.py}.
\end{center}
It checks exact compact word reversal, noncommuting Gaussian square
completion, the scalar Mehler precision and ratio identities, finite
history insertion and centering formulas with non-Perron endpoints,
time-window and tail constants, and the logarithmic spatial entry regime.
All 609 exact assertions passed. Its finite positive-transfer examples
are tests of the transfer algebra, not numerical Wilson simulations. The
report separates the assertion families. The script has SHA--256
\begin{center}\scriptsize\ttfamily
92AA1F6C83ADE826AA35BA145CC6FDEACAB849C36626D2A0A7FB56B267009598.
\end{center} These tests do not evaluate the
suprema of native word derivatives, certify the infinite-dimensional
operator estimates with a proof assistant, or establish a physical mass
gap. The analytic proofs, with the restrictions above, are essential.

The supplied V4--V7 diagnostics were rerun unchanged and passed 9,467,
89,339, 906 and 311 assertions, respectively. Their reports are included.
Unavailable archive diagnostics are not represented as rerun. The V7
predecessor has 231,802 bytes and SHA--256
\begin{center}\scriptsize\ttfamily
36044181A0F2AB65DF5ABC97B95B0652891AAA923D08508DBF903277A3115BC3.
\end{center}
It is preserved byte for byte except for the cumulative
title version and insertion of this appendix immediately before its final
\verb|\end{document}|. Removing the exact appendix and restoring that title
recovers the supplied V7. The package includes source hashes, preservation
checks, and build/visual-review reports. This is not the scheduled V10
companion checkpoint.

\begin{firewall}
V8 proves geometric temporal decay and a small \emph{operator norm} for
V7's complete-native corrected-source covariance and nonadjacent-time
memory on time-independent bounded bridge backgrounds. The original
full/core correction has the same summable time-tail control. These
statements hold for arbitrary duration but with an explicit exponential
spatial cost, admitting $B=O_r(\log\gamma)$ in the stated weak-field
regime. No spatial-volume-uniform result at fixed coupling, arbitrary
inhomogeneous-history mixing, boundary-uniform spatial gluing,
gauge-invariant bridge coverage, full physical transfer gap, interacting
continuum Yang--Mills theory, or positive physical mass gap is established
by this appendix.
\end{firewall}
% END F16 V8 APPEND
% BEGIN F16 V9 APPEND -- 2026-09-17
% Insert immediately before the final \end{document} in cumulative V8.
% The predecessor body is unchanged; the cumulative title becomes V9.

\clearpage
\section{V9: exact history centering and inhomogeneous fast-transfer products}
\label{f16v9:context}

V8 controls coherent common-source memory for arbitrarily long histories,
but only when the bridge background is constant in time. A separate
spectral gap for each frozen operator does not control their product.
Here we remove that restriction without assuming that consecutive bridge
values are close. The upstream step is to center the \emph{whole Gaussian
history}, including its dressed ends, before comparing any native step.
Explicit linear multipliers then telescope between slices. Every resulting
Gaussian step is the same centered operator, while every resulting native
step is uniformly close to it. A direct nonautonomous product estimate
pays the exact product normalizer, rather than replacing a long native
product by a Gaussian power.

The spatial cost remains exponential in $B$. This appendix does not claim
a spatial-volume-uniform estimate at fixed coupling. What changes is the
background class: all bridge histories in the specified rescaled bounded
window are now admitted, with no bound on their duration or total temporal
variation. The covariance and time-tail estimates concern the same exact
Haar-corrected sources and the same original full/core probability as before.

\subsection{Dependency audit and the distinction from a global cone argument}

The supplied V8 appendix was read in full. The available archive inventories
were searched and the relevant interfaces inspected; this is not independent
certification of every earlier proof. We reuse V3's history precision and
Gaussian inverse locality, V5's exact unmoving coordinates and global pins,
V7's Haar source identity, and V8's complete-chart kernel and insertion
method. These are not new V9 results.

Two archived cautions matter here. F15 V57 explicitly restricts its repeated
transfer argument to a homogeneous baseline. F15 V91,
\texttt{f15v91:projective}, proves that endpoint multiplication does not
improve the global Hilbert-projective diameter of the full positive cone.
The present proof uses neither such a diameter estimate nor a claim that
multiplication alone improves it. It verifies an adapted Hilbert-space
slope ball around a specified Gaussian vector, with its exact perturbation
size and endpoint entry cost. The monograph's
\texttt{sec:YM-Perron-vacuum-metric} distinction between an actual slab law
and a stationary Perron law is retained.

For background on adapted cone methods, see H.~H.~Rugh,
\href{https://annals.math.princeton.edu/2010/171-3/p07}
{\emph{Cones and gauges in complex spaces: Spectral gaps and complex
Perron--Frobenius theory}}, Annals of Mathematics 171 (2010), 1707--1752.
The particular real-Hilbert-space product lemma needed here is proved
below, including non-self-adjoint factors and non-Perron endpoints. No
external theorem about the present native history is imported.

\subsection{The full history mean and its telescoping linear terms}

Throughout, fix
\[
 B=m^3,\quad m\ge4,\quad n\ge1,\qquad
 \mathbf y=(y_0,\ldots,y_n),\qquad \max_{i,e}|y_{i,e}|\le r<\infty.
\]
Constants with subscript $r$ may depend on this fixed window and the cube
geometry, never on $n$ or on the individual history. All path matrices act
with $I_3$ in colour. Let
\[
 N=\mathsf M_B^{-1},\quad S=\mathsf H_B+A^{\mathsf T}A,\quad
 \Lambda=A^{\mathsf T}D,\quad
 P_f=\tfrac12(C_B^{-1}+A^{\mathsf T}A).
\]
The symbol $P_f$ is an amplitude precision, not V7's port lift. Let $W_n$
be exactly V3's positive internal history precision and put
\begin{equation}
 \boldsymbol\mu=(\mu_0,\ldots,\mu_n)
       =-W_n^{-1}(\Lambda y_0,\ldots,\Lambda y_n).
 \label{f16v9:mean}
\end{equation}
Thus $\boldsymbol\mu$ is the internal mean of the complete Gaussian
conditional with its original dressed ends. It is not the list of minima
of the separate frozen potentials. In V5's unmoving temporal coordinates
$a_t$, its port mean is
\begin{equation}
 \nu_t=\mathsf F_B(\mu_t-\mu_{t+1})
               -L_h^{-1}E_{\rm br}^{\mathsf T}(y_t-y_{t+1}).
 \label{f16v9:port-mean}
\end{equation}
Use translated complete charts $u_i=x_i-\mu_i$, $b_t=a_t-\nu_t$.
These are translations in integration coordinates, not changes to a compact
configuration or to the definition of a source derivative.

Define the left and right linear terms of bond $t$ by
\begin{align}
 l_t&=N(\mu_t-\mu_{t+1})+\tfrac12(S\mu_t+\Lambda y_t),\notag\\
 r_t^{\rm lin}&=N(\mu_{t+1}-\mu_t)
                      +\tfrac12(S\mu_{t+1}+\Lambda y_{t+1}),
 \label{f16v9:bond-linear}
\end{align}
and define endpoint terms
\[
 d_L^{\rm lin}=P_f\mu_0+\tfrac12\Lambda y_0,\qquad
 d_R^{\rm lin}=P_f\mu_n+\tfrac12\Lambda y_n.
\]
The superscript distinguishes $r_t^{\rm lin}$ from a corrected score.

\begin{proposition}[Uniform mean bounds and exact telescoping]
\label{f16v9:centering}
There is $C_r<\infty$ independent of $B,n$ such that every cube block
of $\mu_i,\nu_t,l_t,r_t^{\rm lin},d_L^{\rm lin},d_R^{\rm lin}$ has norm
at most $C_r$. Moreover
\begin{equation}
 d_L^{\rm lin}+l_0=0,\qquad
 r_{i-1}^{\rm lin}+l_i=0\ (0<i<n),\qquad
 r_{n-1}^{\rm lin}+d_R^{\rm lin}=0.
 \label{f16v9:telescoping}
\end{equation}
No smallness of $y_t-y_{t+1}$ is assumed.
\end{proposition}

\begin{proof}
V3 gives a spatially block-local whitening of $W_n$ with spectrum in
$[4,15]$. Its inverse has block decay at most $C(11/15)^d$ on the
cube--time graph. A radius-$d$ ball has at most $C(d+1)^4$ labels,
uniformly on the spatial torus times the finite time interval. The inverse
therefore has a bounded sum of block norms in every row. The source
$\Lambda y_i$ has bounded norm per cube by the finite incidence and the
bound $r$. Undoing the fixed block whitening proves the assertion for
$\mu$. Likewise $L_h$ has a fixed positive lower spectral bound and
finite spatial range. Its inverse has a uniformly summable spatial block
row, by the Neumann expansion already used in V7. Equation
\eqref{f16v9:port-mean} gives the bound for $\nu$. The remaining matrices
in \eqref{f16v9:bond-linear} and the endpoint terms have fixed finite
range and bounded block rows, proving their bounds.

For an interior slice, the sum of the two displayed bond terms is
\[
 (S+2N)\mu_i-N\mu_{i-1}-N\mu_{i+1}+\Lambda y_i=0
\]
by \eqref{f16v9:mean}. At the first endpoint the coefficient of $\mu_0$
is $P_f+N+S/2=E$, exactly V3's endpoint diagonal; the source is
$\Lambda y_0$. The same argument applies at the last endpoint. This
proves all three identities, also for $n=1$ when there are no interior
slices. The port mean follows from V2's kinetic square: in moving
coordinates the minimizing port is
$-L_h^{-1}E_{\rm br}^{\mathsf T}(y_t-y_{t+1})$; V5's determinant-one
linear change adds $\mathsf F_B(\mu_t-\mu_{t+1})$.
\end{proof}

\subsection{Balanced kernels with one common Gaussian reference}

For two possibly different bridges $y,y'$, let $K_\gamma(y,y';x,x')$
be V8's augmented-step integral with these values at the two ends. More
explicitly, use its actual words with $U_e(x;y)$ and $U_e(x';y')$,
magnetic halves $V_{\gamma,y}(x)/2,V_{\gamma,y'}(x')/2$, and scalar
\begin{equation}
 \rho_\gamma(y,y')=
   (c_{H,\gamma}/z_\gamma)^{24B}
             [j_{\rm br}(y)j_{\rm br}(y')]^{1/2}.
 \label{f16v9:step-scalar}
\end{equation}
All internal and port charts are complete. This kernel is generally not
self-adjoint; its transpose is $K_\gamma(y',y)$. Its kernel is
nonnegative. No positive-semidefinite operator claim is needed for unequal
bridges. Exact composition gives
\begin{equation}
 \mathcal W_{\gamma,n}^{\rm full}(\mathbf y)
   =\langle a_{\gamma,y_0},
       K_\gamma(y_0,y_1)\cdots K_\gamma(y_{n-1},y_n)
                                      a_{\gamma,y_n}\rangle.
 \label{f16v9:inhomogeneous-path}
\end{equation}
The normalized endpoints are V8's literal vectors at their respective
bridge values. This follows by the same Haar powers as V8; it is not a
new replacement law.

The integrated Gaussian kernel is
\begin{equation}
 K_0(y,y';x,x')=
 \varphi_G(y-y')\,
 e^{-V_{0,y}(x)/4}\varphi_{\mathsf M_B}(x-x')e^{-V_{0,y'}(x')/4},
 \qquad
 V_{0,y}(x)=x^{\mathsf T}Sx+2x^{\mathsf T}\Lambda y+|Dy|^2.
 \label{f16v9:Gaussian-step}
\end{equation}
Set $K_*=K_0(0,0)$, $a_*=a_{0,0}$, and let
$\lambda_*,\phi_*$ be its exact top value and normalized positive vector.
In this appendix $a_*$ denotes this vector; the pin constant is denoted
$c_{\rm pin}=(6\pi^2)^{-1}$. V8's Mehler calculation gives
\begin{equation}
 S_*:=K_*/\lambda_*=P_*+Q_*,\quad
 P_*=|\phi_*\rangle\langle\phi_*|,\quad
 0\le Q_*,\quad P_*Q_*=0,\quad \|Q_*\|\le3-2\sqrt2<1/5.
 \label{f16v9:reference-split}
\end{equation}
No second oscillator diagonalization is being claimed.

Introduce positive scalars
\begin{equation}
 c_t=\frac{K_0(y_t,y_{t+1};\mu_t,\mu_{t+1})}{K_*(0,0)},\qquad
 d_L=\frac{a_{0,y_0}(\mu_0)}{a_*(0)},\qquad
 d_R=\frac{a_{0,y_n}(\mu_n)}{a_*(0)}.
 \label{f16v9:balancing-scalars}
\end{equation}
Define the balanced native kernels and endpoints on the same Euclidean
Hilbert space by
\begin{align}
 \widetilde K_t(u,v)
   &=c_t^{-1}e^{\langle l_t,u\rangle+\langle r_t^{\rm lin},v\rangle}
       K_\gamma(y_t,y_{t+1};u+\mu_t,v+\mu_{t+1}),\notag\\
 \widetilde a_L(u)
   &=d_L^{-1}e^{\langle d_L^{\rm lin},u\rangle}
                                      a_{\gamma,y_0}(u+\mu_0),\notag\\
 \widetilde a_R(u)
   &=d_R^{-1}e^{\langle d_R^{\rm lin},u\rangle}
                                      a_{\gamma,y_n}(u+\mu_n).
 \label{f16v9:balanced-native}
\end{align}
Their supports are the translated complete charts. They are not normalized
probability densities or unit vectors by definition.

\begin{theorem}[Exact balanced path and a uniform one-step comparison]
\label{f16v9:balanced-comparison}
The Gaussian versions of \eqref{f16v9:balanced-native} are exactly $K_*$
at every bond and $a_*$ at both ends. At finite regulator one has the
\emph{exact} identity
\begin{equation}
 \mathcal W_{\gamma,n}^{\rm full}(\mathbf y)
  =d_Ld_R\left(\prod_{t=0}^{n-1}c_t\right)\lambda_*^n
     \langle\widetilde a_L,T_0\cdots T_{n-1}\widetilde a_R\rangle,
 \qquad T_t=\widetilde K_t/\lambda_*.
 \label{f16v9:balanced-path}
\end{equation}
There are $C_r,c_r,\gamma_r$ independent of $B,n$ such that for
$\gamma\ge\gamma_r$,
\begin{align}
 \sup_t\|T_t-S_*\|_{\rm HS}
    +\|\widetilde a_L-a_*\|_2+\|\widetilde a_R-a_*\|_2
       &\le e^{C_rB}/\sqrt\gamma,\notag\\
 \langle a_*,\phi_*\rangle&\ge e^{-c_rB},\notag\\
 e^{-C_rB}\le c_t,d_L,d_R&\le e^{C_rB}.
 \label{f16v9:balanced-bounds}
\end{align}
The matrix element in \eqref{f16v9:balanced-path} is retained exactly; it
is not replaced by $1$ or by $\langle a_*,S_*^na_*\rangle$.
\end{theorem}

\begin{proof}
Expanding the quadratic kernel \eqref{f16v9:Gaussian-step} after the two
translations gives precisely the linear coefficients
\eqref{f16v9:bond-linear}. The remaining constant is fixed by evaluation
at $u=v=0$, which is the definition of $c_t$. Likewise the amplitude
precision of $a_{0,y}$ is $P_f$ and its fast linear coefficient is
$\Lambda y/2$. Its translated linear coefficient is $d_{L,R}^{\rm lin}$
and its constant is fixed by $d_{L,R}$. This proves the exact Gaussian
claims. Substitution into the native path integral cancels every linear
multiplier by \eqref{f16v9:telescoping}. The translations have Jacobian
one. This proves \eqref{f16v9:balanced-path} without a limit.

The uniform local mean bounds imply that the squared Euclidean norms of
all means and of all balancing coefficient vectors at one bond are at
most $C_rB$. The Gaussian exponent in $c_t$ is
\[
 -\tfrac14V_{0,y_t}(\mu_t)-\tfrac14V_{0,y_{t+1}}(\mu_{t+1})
 -\tfrac12|\mu_t-\mu_{t+1}|_N^2
 -\tfrac12|y_t-y_{t+1}|_{G^{-1}}^2.
\]
It lies in $[-C_rB,0]$. The Gaussian endpoint determinant, mean and
precision bounds from V8 similarly bound $\log d_L,\log d_R$ in
absolute value by $C_rB$. The overlap assertion is exactly V8's
zero-background overlap estimate.

We spell out why linear balancing is legitimate for the \emph{complete}
native integrals. Before integrating ports, translate them by $\nu_t$ as
in \eqref{f16v9:port-mean}. The port linear term in the Gaussian exponent
then vanishes, because that mean solves its bond stationarity equation.
The two remaining internal linear terms are $l_t,r_t^{\rm lin}$.
Thus the balanced Gaussian augmented exponent is exactly the zero-bridge
augmented exponent, plus the constant already removed in $c_t$.
Both original augmented actions have the global pin
\[
 c_{\rm pin}(|x|^2+|x'|^2+|a|^2),
\]
by the internal magnetic ends and independent rooted tree, as in V8.
After translation and the balancing linear multipliers, both integrands
are bounded by
\begin{equation}
 e^{C_rB}\exp[-c(|u|^2+|v|^2+|b|^2)]
 \label{f16v9:balanced-pin}
\end{equation}
for a fixed $c>0$. This follows twice from
$|z+m|^2\ge|z|^2/2-|m|^2$ and then from Young's inequality for the
balancing linear terms. All Haar factors are at most one and all scalar
powers cost at most $e^{C_rB}$.

The native action minus its Gaussian quadratic action obeys V8's global
word remainder; in the translated variables it is bounded by
$C_r\gamma^{-1/2}(B^{3/2}+|u|^3+|v|^3+|b|^3)$.
The nonconstant Haar-factor error is at most
$C_r\gamma^{-1}(B+|u|^2+|v|^2+|b|^2)$, and the scalar logarithmic
error is at most $C_rB/\gamma$. Using
$|e^{-a}-e^{-b}|\le|a-b|e^{-\min(a,b)}$ and
\eqref{f16v9:balanced-pin} gives a polynomial times a Gaussian majorant
for the difference. A translated principal-chart complement has a group
coordinate of norm at least $\pi\sqrt\gamma-C_r$, hence at least
$\pi\sqrt\gamma/2$ once $\gamma\ge\gamma_r$.
Its Gaussian tail costs $e^{C_rB-c\gamma}$, not an omitted remainder.
Integrate all ports and then take the kernel $L^2(du\,dv)$ norm.
Fixed-degree Gaussian moments in dimension $51B$ cost an exponential
in $B$. The resulting kernel error is $e^{C_rB}/\sqrt\gamma$;
$\lambda_*^{-1}\le e^{C_rB}$ pays the displayed normalization.

For the endpoints use the unnormalized native and Gaussian functions
\[
 q_\gamma=1_{\mathcal P_x}\Omega e^{-\gamma S_\times/2},\qquad
 q_0=\Omega e^{-|Ax+Dy|^2/4}.
\]
Their normalizing norms are at least
$e^{-C_rB}$: integrate their squares over a product of fixed small
one-cube balls, using bounded local inputs and $\gamma\ge\gamma_r$.
After the translation and the linear exponential multiplication,
the factor $\Omega$ still supplies a Gaussian majorant of the form
\eqref{f16v9:balanced-pin}. The same polynomial remainder proof bounds
the weighted unnormalized difference by $e^{C_rB}/\sqrt\gamma$.
The unweighted difference bounds the difference of the normalizing norms.
Divide by the just-proved lower bounds and by $d_{L,R}\ge e^{-C_rB}$.
This proves the two endpoint estimates. In particular we have not applied
an unbounded multiplication operator to an unweighted $L^2$ estimate.
\end{proof}

The centering depends on the complete finite bridge history and its two
ends. It is a comparison coordinate system, not a causal coarse state or
a new physical vacuum. All bridge derivatives defining the observables
are taken in the original coordinates \emph{before} this transformation.
Differentiating the balanced formula while omitting derivatives of its
means or multipliers would be incorrect; no such step is used below.

\subsection{A product lemma with an exact nonautonomous normalizer}

Here is the analytic step replacing the use of one static spectral power.
It is stated on a real Hilbert space, so it also applies to the
infinite-dimensional internal space. Write $P=|\phi\rangle\langle\phi|$.

\begin{theorem}[Uniform rank-one decomposition of near-reference products]
\label{f16v9:product}
Suppose $S=P+Q$, $\|\phi\|=1$, $Q=Q^*\ge0$, $PQ=0$, and
$\|Q\|\le1/5$. Let $(T_j)$ be any finite sequence of real bounded
operators satisfying $\|T_j-S\|\le\delta\le1/100$.
Neither self-adjointness nor commutation of the $T_j$ is assumed.
For an ordered interval product $T_I$ of length $k$, set
$p_I=\langle\phi,T_I\phi\rangle$. Then $p_I>0$ and
\begin{equation}
 \frac{T_I}{p_I}=|\phi+x_I\rangle\langle\phi+y_I|+E_I,
 \quad x_I,y_I\perp\phi,\quad
 \|x_I\|,\|y_I\|\le2\delta,\quad
 PE_I=E_IP=0,\quad \|E_I\|\le3^{-k}.
 \label{f16v9:product-decomposition}
\end{equation}
Consequently $\|T_I\|\le2p_I$. For adjoining intervals $I,J$ in that
product order,
\begin{equation}
 \frac{p_{IJ}}{p_Ip_J}=1+\langle y_I,x_J\rangle
            \in[1-4\delta^2,1+4\delta^2].
 \label{f16v9:quasi-multiplication}
\end{equation}
An empty interval has $T_I=I$, $p_I=1$, $x_I=y_I=0$, $E_I=I-P$;
the same assertions with $k=0$ hold. If $1\le k\le2$, then $p_I\ge1/2$.
\end{theorem}

\begin{proof}
Decompose a single factor relative to $\mathbb R\phi\oplus\phi^\perp$:
\[
 T_j=\begin{pmatrix}a_j&c_j^*\\ b_j&D_j\end{pmatrix},\qquad
 |a_j-1|,\|b_j\|,\|c_j\|\le\delta,\quad \|D_j\|\le1/5+\delta.
\]
Its projective slope map on the unit ball of $\phi^\perp$ is
\begin{equation}
 F_j(z)=\frac{b_j+D_jz}{a_j+\langle c_j,z\rangle}.
 \label{f16v9:slope}
\end{equation}
The denominator is at least $1-2\delta\ge49/50$. On this ball
\[
 \|F_j(z)\|\le\frac{1/5+2\delta}{1-2\delta}\le\frac{11}{49},
 \qquad
 \|DF_j(z)\|\le
 \frac{1/5+\delta+\delta(11/49)}{1-2\delta}<\frac13.
\]
Thus the unit ball is invariant and these maps contract its norm distance.
The smaller ball of radius $2\delta$ is also invariant, since
\[
 \frac{\delta+(1/5+\delta)2\delta}{1-\delta-2\delta^2}
                                                        \le2\delta.
\]
Exactly the same estimates hold for every adjoint factor.

Iterate from slope zero in the order in which the product acts on $\phi$.
All scalar coefficients are positive, proving $p_I>0$ and
$x_I=P^\perp T_I\phi/p_I$ with $\|x_I\|\le2\delta$.
The adjoint product gives $y_I=P^\perp T_I^*\phi/p_I$ with the same
bound. Differentiating the composed slope map at zero gives
\[
 DF_I(0)=P^\perp T_IP^\perp/p_I-x_Iy_I^*.
\]
The chain rule bounds its norm by $3^{-k}$. Extending this operator by
zero on $\phi$ gives $E_I$ and proves the exact block identity
\eqref{f16v9:product-decomposition}. For $k\ge1$ its norm is at most
$1+4\delta^2+1/3<2$; the empty case is immediate. Multiplying two such
block matrices in their top-left corner gives
\eqref{f16v9:quasi-multiplication} exactly. Finally at each slope step
the scalar factor is at least $1-\delta-2\delta^2\ge1-2\delta$;
for $k\le2$ their product is at least $(49/50)^2>1/2$.
\end{proof}

The scalar $p_I$ can grow or decay with $k$. It is \emph{not} replaced by
one. The theorem bounds the normalized remainder, rather than the generally
false estimate $\|T_I-S^k\|=O(\delta)$ independent of $k$. Nor is the
slope ball the full positive-function cone from f15 V91.

\begin{proposition}[Original-type endpoints and a denominator for all lengths]
\label{f16v9:endpoints}
In the preceding theorem suppose $\|a\|=1$,
$\langle a,\phi\rangle\ge\beta$ with $0<\beta\le1$, and
$\|a_L-a\|,\|a_R-a\|\le\eta$. Assume
\begin{equation}
 \delta\le10^{-6}\beta^4,\qquad \eta\le10^{-6}\beta^2.
 \label{f16v9:endpoint-entry}
\end{equation}
For every product of length $n\ge1$,
\begin{equation}
 \frac{\beta^2}{16}\,p_{[0,n)}
     \le\langle a_L,T_0\cdots T_{n-1}a_R\rangle
     \le8p_{[0,n)}.
 \label{f16v9:denominator}
\end{equation}
The factors need not have a common eigenvector. No stationarity of the
path ends is required.
\end{proposition}

\begin{proof}
Put $N_\beta=\lceil\log(128/\beta^2)/\log3\rceil$.
For $n\ge N_\beta$, the rank-one term in
\eqref{f16v9:product-decomposition}, paired with the endpoints, is at
least $(\beta-\eta-4\delta)^2\ge\beta^2/4$. Its remainder has absolute
value at most $4\,3^{-n}\le\beta^2/32$, since $\|a_L\|,\|a_R\|\le2$.
This proves the stated lower bound for these lengths.

For $1\le n<N_\beta$, use $N_\beta\le10/\beta^2$, which follows
from $\log x\le x-1$ for $x\ge1$. The telescoping product estimate gives
\[
 \|T_0\cdots T_{n-1}-S^n\|
   \le n\delta(1+\delta)^{n-1}\le2n\delta.
\]
Here $n\delta\le10^{-5}\beta^2$. The actual endpoint matrix element
therefore differs from $\langle a,S^na\rangle$ by at most
$6\eta+2n\delta<\beta^2/2$. The latter is at least $\beta^2$ because
$S^n=P+Q^n$ and $Q^n\ge0$. Also
$p_{[0,n)}\le(1+\delta)^n\le2$. This proves a lower bound
$\beta^2p_{[0,n)}/4$, stronger than required. Finally
$\|T_{[0,n)}\|\le2p_{[0,n)}$ and the endpoint norms give the upper
bound for all lengths. This proof treats short histories separately rather
than imposing an unproved Perron denominator on their endpoints.
\end{proof}

\subsection{Two source windows in the nonautonomous product}

The following consequence keeps the exact finite-path means. It applies
to positive integral kernels with auxiliary bond variables integrated into
insertions. It does not require their operators to be self-adjoint.

\begin{proposition}[Centered two-window estimate for the actual product law]
\label{f16v9:two-window}
Assume the preceding two results and that the kernels and endpoints define
a positive path law. Let $F,G$ be real observables on windows of one or
two bonds, with disjoint bonds and $g\ge0$ unmarked bonds between them.
Let $A_F,A_G$ be the kernels replacing their respective window products,
with the observables inserted; let $B_F$ insert $F^2$. There is an absolute
constant $C_0$ (for example $2^{28}$) such that
\begin{equation}
 |\operatorname{Cov}(F,G)|
       \le C_0\beta^{-4}\|A_F\|\|A_G\|3^{-g},\qquad
 \mathbb E F^2\le C_0\beta^{-2}\|B_F\|.
 \label{f16v9:nonautonomous-mixing}
\end{equation}
Both expectations use the \emph{actual} normalizer in
\eqref{f16v9:denominator}.
\end{proposition}

\begin{proof}
Split the unmarked full product into left segment $L$, first window $U$,
gap $G_0$, second window $V$, and right segment $R$, in that order.
Write $p_L,p_U,p_{G_0},p_V,p_R$ for their exact scalar coefficients.
At most four applications of \eqref{f16v9:quasi-multiplication}, and
$p_U,p_V\ge1/2$, give
\[
 Z':=\frac{\langle a_L,T_{[0,n)}a_R\rangle}{p_Lp_{G_0}p_R}
                                                      \ge\beta^2/128.
\]
The convention for empty products handles endpoints and $g=0$.
Set
\[
 \begin{gathered}
 u=T_L^*a_L/p_L,\quad v=T_Ra_R/p_R,\quad H=T_{G_0}/p_{G_0},\\
 u_0=T_U^*u,\quad u_F=A_F^*u,\quad
 v_0=T_Vv,\quad v_G=A_Gv.
 \end{gathered}
\]
Then $\|u\|,\|v\|\le4$, $\|u_0\|,\|v_0\|\le5$,
$\|u_F\|\le4\|A_F\|$, $\|v_G\|\le4\|A_G\|$, and
$Z'=\langle u_0,Hv_0\rangle$. Exact composition gives
\begin{equation}
 \operatorname{Cov}(F,G)=
 \frac{\langle u_F,Hv_G\rangle\langle u_0,Hv_0\rangle
       -\langle u_F,Hv_0\rangle\langle u_0,Hv_G\rangle}{(Z')^2}.
 \label{f16v9:minor}
\end{equation}
Replace $H$ by the rank-one part of
\eqref{f16v9:product-decomposition}. The numerator vanishes exactly.
The norms of both $H$ and that rank-one part are at most two, and their
difference has norm at most $3^{-g}$. Expanding the difference of the two
products bounds the numerator by
$3200\|A_F\|\|A_G\|3^{-g}$.
The displayed denominator lower bound makes $2^{28}$ a safe constant.
For a single window, the corresponding three-segment scalar comparison
gives $Z\ge\beta^2p_Lp_R/64$, while its square-insertion numerator has
absolute value at most $16p_Lp_R\|B_F\|$. This proves the second bound.
Auxiliary port integrations belong to the window kernels, so the same
identities apply to their observables. The centering in
\eqref{f16v9:minor} includes the actual two finite-history means.
\end{proof}

\subsection{Full-native remainder insertions after history centering}

Use exactly V7's full lift $\mathsf P$ and corrected common-source
residuals $r^{[\infty]}_\alpha$, evaluated at the original bridge history.
For $v\in\mathcal E$ at time $i$, set
\[
 R_i(v)=\sum_{b,\mu,k}v_{b,\mu,k}r^{[\infty]}_{i,b,\mu,k}.
\]
Its bond window is
$[\max(0,i-1),\min(n,i+1)]$, of length one or two. First form this
observable using original bridge derivatives and compact Haar insertions;
then express its arguments in the centered integration coordinates.
All centering means are held fixed during this evaluation.

Let $A_i(v),B_i(v)$ be its first and square insertion operators using the
balanced augmented kernels divided by $\lambda_*$ per bond. Thus they
replace the corresponding products of the $T_t$, with \emph{all} their
scalars $c_t$ and endpoint linear multipliers included.

\begin{proposition}[Uniform complete-chart insertion bounds at arbitrary histories]
\label{f16v9:insertions}
For sufficiently large $\gamma$ there is $C_r$ independent of $n,B,i$
and of the admitted history such that
\begin{equation}
 \|A_i(v)\|\le e^{C_rB}\gamma^{-1/2}\|v\|,\qquad
 \|B_i(v)\|\le e^{C_rB}\gamma^{-1}\|v\|^2.
 \label{f16v9:insertion-bounds}
\end{equation}
These are estimates for complete native kernels, not Gaussian expectations
used instead of native expectations.
\end{proposition}

\begin{proof}
V7's exact Gaussian cancellation holds for every bridge history. On each
adjacent bond, the lift cancels the fast affine kinetic score, leaving its
explicit deterministic contribution to $g_\alpha(\mathbf y)$; at a
common magnetic mark the fast affine coefficient is zero. Subtracting the
same $g_\alpha$ therefore leaves only primitive nonlinear remainders and
the bounded $L_h^{-1}E_{\rm br}^{\mathsf T}$ synthesis of port remainders,
also when $y_t-y_{t+1}\ne0$. The inserted exponential-word Taylor bound
from V7 is global on the complete charts. With $w$ the at most $87B$
original real fast coordinates in the window, it gives
\[
 |R_i(v)|\le C_r\gamma^{-1/2}(\sqrt B+|w|^2)\|v\|.
\]
Only the spatial lift's $\ell^2$ operator norm is used, not a sum of
absolute coefficients. After translation the window's mean vector has
squared norm at most $C_rB$. Thus the same expression is bounded by
$C_r\gamma^{-1/2}(B+|w_{\rm centered}|^2)\|v\|$.

The product of at most two balanced augmented kernels has the Gaussian
majorant \eqref{f16v9:balanced-pin}, with its constant enlarged. Multiplying
by this polynomial or its square, integrating the full translated port
and intermediate internal charts, and taking the Hilbert--Schmidt norm in
the two external variables proves \eqref{f16v9:insertion-bounds}.
The fixed-degree moments cost $e^{C_rB}$ as in V8. Division by one or two
powers of $\lambda_*\ge e^{-C_rB}$ is included. No whole-history
denominator or duration factor enters this local integration. The native
port derivatives are still Haar derivatives; translating their displayed
coordinate arguments has not replaced them by flat derivatives.
\end{proof}

\subsection{Arbitrary-history coherent memory and the unchanged core correction}

\begin{theorem}[Inhomogeneous common-source covariance, uniformly in duration]
\label{f16v9:coherent}
Fix $r<\infty$. There are finite $K_r\ge1,\gamma_r$ such that, whenever
\begin{equation}
 \gamma\ge\gamma_r,\qquad e^{K_rB}/\sqrt\gamma\le1/100,
 \qquad \max_{i,e}|y_{i,e}|\le r,
 \label{f16v9:entry}
\end{equation}
V7's exact full-native residual covariance, at \emph{this arbitrary
history}, satisfies
\begin{equation}
 \boxed{\quad
 \|\Gamma_{\gamma,ij}\|_{\mathcal E\to\mathcal E}
     \le A_{\gamma,B}(1/3)^{(|i-j|-2)_+},\qquad
 A_{\gamma,B}=e^{K_rB}/\gamma.\quad}
 \label{f16v9:covariance}
\end{equation}
Consequently, for every $n\ge1$ and every common-source sequence $v_i$,
\begin{equation}
 \boxed{\quad
 \operatorname{Var}_{\mu_{\gamma,\mathbf y}}
       \left(\sum_{i=0}^n R_i(v_i)\right)
      \le6A_{\gamma,B}\sum_{i=0}^n\|v_i\|^2.\quad}
 \label{f16v9:coherent-bound}
\end{equation}
For $0\le\alpha<\log3$, the weighted time-row bound is
\begin{equation}
 \sup_i\sum_j e^{\alpha|i-j|}\|\Gamma_{\gamma,ij}\|
 \le A_{\gamma,B}
       \left(1+2e^\alpha+\frac{2e^{2\alpha}}{1-e^\alpha/3}\right).
 \label{f16v9:weighted}
\end{equation}
No constancy, slow variation, or total-variation budget on $\mathbf y$
is assumed. The exponential spatial cost is part of the assertion.
\end{theorem}

\begin{proof}
Apply Theorem~\ref{f16v9:balanced-comparison} with the common reference
$S_*$ and $a_*$. Put $\beta=e^{-c_rB}$ using its overlap lower bound.
By enlarging $K_r$, the entry condition makes the uniform step error
at most $10^{-6}\beta^4$ and the endpoint errors at most
$10^{-6}\beta^2$. This is possible because every error is
$e^{C_rB}/\sqrt\gamma$, whereas $\beta$ is a fixed exponential in
$B$. The constants $10^{-6}$ are absorbed by enlarging a geometric
exponent, not by introducing a duration-dependent threshold.
The product and denominator results therefore apply for \emph{all} $n$.

The exact balancing identity changes no normalized path probability when
observables are translated along with their arguments. For unit spatial
$v,w$, square insertion and Cauchy--Schwarz give, when $|i-j|\le1$,
\[
 |\operatorname{Cov}(R_i(v),R_j(w))|
                         \le C\beta^{-2}e^{C_rB}/\gamma.
\]
For $j-i\ge2$ the two windows have gap $g=j-i-2$ and disjoint bonds.
Proposition~\ref{f16v9:two-window} and the two first insertion estimates
give
\[
 |\operatorname{Cov}(R_i(v),R_j(w))|
                  \le C\beta^{-4}e^{2C_rB}\gamma^{-1}3^{-(j-i-2)}.
\]
The proof allows windows sharing an internal endpoint when $g=0$.
Enlarge $K_r$ once more to absorb these finitely many spatial exponents.
The entry condition only becomes stronger. Taking the supremum over unit
$v,w$ proves the block operator bound. The scalar row sum of its majorant
is $5+2(1/3)/(1-1/3)=6$, proving the coherent variance estimate by the
block Schur bound. Summing with the indicated weights proves
\eqref{f16v9:weighted}. These last sums are the same elementary sums as
V8; the new input is that the full covariance estimate now holds without
one repeated native kernel.
\end{proof}

\begin{theorem}[Full history Hessian and original full/core tails]
\label{f16v9:memory}
Under \eqref{f16v9:entry}, let $\mathsf M_\gamma$ and
$\mathsf C_{\gamma,R}$ be V7's nonadjacent-time common-source Hessian
and its original full/core correction. For every integer $L\ge1$,
\begin{align}
 \|\mathsf M_\gamma\|_{\rm op}&\le3A_{\gamma,B},\notag\\
 \|\mathsf M_\gamma^{>L}\|_{\rm op}
                   &\le3A_{\gamma,B}(1/3)^{L-1},\notag\\
 \|\mathsf C_{\gamma,R}^{>L}\|_{\rm op}
   &\le3A_{\gamma,B}(1/3)^{L-1}
                   +C(R^2/\gamma)(9/19)^{L-1}.
 \label{f16v9:memory-tails}
\end{align}
Here $R$ satisfies V4's original derivative regime, and $>L$ keeps blocks
with $|i-j|>L$. In particular $L=1$ bounds the entire correction. These
bounds hold at every admitted history and every duration.
\end{theorem}

\begin{proof}
V7's exact Haar identity has contact range one for arbitrary bridge
histories. Therefore $\mathsf M_\gamma=-\operatorname{Off}_{>1}
\Gamma_\gamma$ before and after the present change of integration
coordinates. For $d>L\ge1$, summing the block bound gives
$2A_{\gamma,B}\sum_{d=L+1}^\infty3^{-(d-2)}
=3A_{\gamma,B}3^{-(L-1)}$. The block Schur bound proves the first two
claims. V4's common-source core bound already holds at arbitrary bounded
histories, with time tail $C(R^2/\gamma)(9/19)^{L-1}$. Subtract it using
the exact identity
\[
 \nabla^2\log p_{\rm good}
   =\nabla^2[-\log\mathcal W_{\gamma,n}^{\rm full}]
                         -\nabla^2[-\log\mathcal W_{\gamma,n}^{R}].
\]
This proves the last claim. The core event is not translated into a new
product event and is not replaced. All endpoint normalizers remain the
original ones; their separated-time Hessians vanish because they are
single-time factors, not because their values were discarded.
\end{proof}

\subsection{Scope gained and the next upstream spatial problem}

The entire static-background restriction of V8 has been removed for these
conditional common-source estimates. For each fixed $B$, the coherent
covariance and nonlocal memory are $O_{r,B}(\gamma^{-1})$ uniformly over
all bounded histories and all $n$. With
\begin{equation}
 B\le\frac{\log\gamma}{4K_r},\qquad R=4\sqrt{\log\gamma},
 \label{f16v9:mesoscopic}
\end{equation}
the entry and core conditions hold eventually, and
\begin{equation}
 \|\Gamma_\gamma\|_{\rm op}+\|\mathsf M_\gamma\|_{\rm op}
              +\|\mathsf C_{\gamma,R}\|_{\rm op}
                                    \le C_r\gamma^{-3/4}.
 \label{f16v9:mesoscopic-bound}
\end{equation}
This includes arbitrary temporal oscillation within the fixed rescaled
window. There is no $n\gamma^{-1/2}$ error in the conclusion. Individual
scalar normalizations can drift through a long product, but the exact
$p_I$ and the exact finite-path denominator were used at every step.

The remaining large spatial cost has three identifiable sources: the
complete-chart one-step Hilbert--Schmidt comparison, the weighted insertion
norms, and the overlap of the global dressed endpoint with the Gaussian
reference vector. They are all exponential in the number of spatial cubes.
A product estimate cannot remove a spatial cost already present in its
inputs. In particular the present inequality must not be tensorized across
interacting regions or extrapolated to $B\to\infty$ at fixed $\gamma$.

The next useful estimate is spatially local: a boundary-conditioned,
source-marked version of the complete-chart comparison and its endpoint
message bound, on a source region with a retained spatial collar. V9's
product lemma can then be reused without a static-background assumption;
another temporal product proof would not address the remaining cost.
Arbitrary exterior fast configurations can alter a conditional landscape,
so a torus pin by itself is not a boundary-uniform hypothesis. The original
bridge window is still not gauge-invariant coverage of the integrated
physical bridge law. Infrared contact coercivity, comparison with the same
full physical transfer top, physical-time calibration, and interacting
continuum reconstruction remain separate obligations.

\subsection{Diagnostics and preservation}

The accompanying diagnostic is
\begin{center}\small
\path{verify_ym_f16_v9_inhomogeneous_history.py}.
\end{center}
It tests noncommuting Gaussian history centering, both endpoint equations,
exact telescoping of linear multipliers and scalars, varying nonsymmetric
matrix products and their normalized Schur remainders, quasi-multiplicative
scalar identities, centered insertion formulas with unequal endpoints,
and the geometric time-tail constants. All 990 assertions passed.
Assertion families and counts are reported in \path{ym_v9_verification_report.json}. These are exact finite
algebraic checks, not Wilson simulations or proof-assistant verification
of the complete-chart estimates or the infinite-dimensional arguments.
The diagnostic source has SHA--256
\begin{center}\scriptsize\ttfamily
EF4E12653D522EB3F96547274A7AA981F64C4B4BD9398635DB39583BA5FC1657.
\end{center}

The cumulative V8 predecessor has 273,042 bytes and SHA--256
\begin{center}\scriptsize\ttfamily
D12A126849ACA4014F321B2D523BA2669C9C33503B839E86C9E13F010C54DC75.
\end{center}
Only its title version changes and this appendix is inserted before the
final \verb|\end{document}|. The package reports the byte-for-byte recovery
check, input hashes, build results, and reruns of the available V4--V8
scripts. Unavailable archive scripts are not represented as rerun. The
next scheduled companion checkpoint remains V10.

\begin{firewall}
V9 proves exact history centering, normalized nonautonomous fast-product
control, and small coherent common-source covariance with summable temporal
tails for every bounded rescaled bridge history. The original full/core
correction obeys the corresponding operator-norm bounds. The estimates are
uniform in history length but retain an exponential spatial cost, admitting
only the stated logarithmically growing spatial regime. No spatial-volume-
uniform bound at fixed coupling, arbitrary-boundary spatial gluing,
gauge-invariant bridge coverage, infrared contact coercivity, full physical
transfer gap, interacting continuum theory, or positive physical
Yang--Mills mass gap is established by this appendix.
\end{firewall}
% END F16 V9 APPEND
% BEGIN F16 V10 APPEND -- 2026-09-17
% Insert immediately before the final \end{document} in cumulative V9.
% The predecessor body is unchanged; only the cumulative title becomes V10.

\clearpage
\section{V10: finite-order local memory coefficients and a controlled native remainder}
\label{f16v10:context}

The supplied Navier--Stokes import memorandum starts from V3. The active
frontier is now V9: arbitrary bounded bridge histories have a native,
time-summable corrected-source covariance, but its comparison and endpoint
costs are exponential in the spatial volume. Starting again from the V3
Gaussian kernel would repeat work. We instead expand the \emph{already
identified exact corrected-source covariance}, retain its connected
coefficients, and improve only the remainder.

The principal new distinction is between two norms. Every fixed-order
coefficient has a volume- and duration-independent exponentially weighted
space--time row bound. Its \emph{native remainder} has a duration-independent
temporal block bound, still with a displayed exponential spatial cost.
Their combination gives a better quantitative native estimate in the
admitted growing-volume regime. It does not silently turn an asymptotic
coefficient bound into arbitrary-volume control of the full compact law.

\subsection{The V10 checkpoint and the external import actually used}

The primary source examined is OpenAI,
\href{https://cdn.openai.com/pdf/32d9f210-8b73-45\varepsilon0-91bc-82a30aef8a9a/navier-stokes.pdf}
{\emph{Finite Time Blowup for Navier--Stokes}}, released September 8, 2026.
Its Theorem 1.1 constructs smooth-forced finite-time blowup with bounded
kinetic energy. Proposition 9.6 increases the residual order by $1/10$;
Lemma 9.7 fixes the principal inverses and records finite derivative
dependencies. Lemma 8.7 separates two preserved moments from three defect
corrections. Section 5.4 and Lemma 5.4 use cutoff summation rather than
assuming convergence of the raw correction series. The accompanying
\href{https://github.com/openai/NavierStokesAndEuler}{Lean repository}
was inspected at its documentation level, not built or independently
verified here.

The useful import is a \emph{proof architecture}, not an NS estimate on a
YM measure. In particular, the native Wilson operator is fixed. Its genuine
response coefficients cannot be removed by selecting a different forcing,
adding arbitrary counterterms, or changing its probability law. Here a
``correction'' means adding the next coefficient to a specified
approximation of that fixed response. The exact residual, including its
normalizer, is then estimated anew.

The supplied V9 appendix and import memorandum were read in full. The
available f14/f15 archives and Grand Tensor Monograph V076 were inventoried
and their relevant passages checked. F14 V38 already treats connected
noise derivatives and regression for a different, uniformly convex shell
extension; it does not prove the full compact history theorem below.
The monograph's fixed-order Wilson moment bounds explicitly distinguish
finite constants at each order from an all-order summability estimate.
V6--V9 already supply the normalized laws, Haar sources, and history
balancing used here. None is presented again as a new construction.
Only the supplied companion sources were available at this V10 checkpoint;
other programs named in older ledgers were not supplied or recertified.

Several proposed imports need limits. Gauss constraints and the
$9B(n+1)$ common-source coordinates are not five fixed global moments and
are not projected out. Positive covariance does not make its nonadjacent-time
part positive. Auxiliary labels may organize overlapping terms but cannot
delete their native cross interactions. We establish arbitrary \emph{fixed}
order, not an order-uniform convergence theorem or a cutoff-summed replacement
for the native measure. The distinction between order gain and
order-dependent constants will remain explicit.

\subsection{One fixed law, one Gaussian inverse, and measured polynomial coefficients}

Use V6's full unmoving law, V7's full Haar lift, and V9's exact history
centering. To avoid collision with the earlier entropy rate, put
\[
 \varepsilon=\gamma^{-1/2},\qquad \mathcal N=B(n+1),\qquad
 \max_{i,a}|y_{i,a}|\le r<\infty.
\]
This $\varepsilon$ is the expansion parameter, distinct from V6's
entropy-rate sequence $\epsilon_\gamma$. The fast coordinates are $\zeta=(x,a)$, and
\begin{equation}
 d\mu_\varepsilon=Z_\varepsilon^{-1}1_{\mathcal P_\varepsilon}J_\varepsilon^f \exp\!\left(-U_\varepsilon\right)\,d\zeta,
 \qquad d\mu_0=Z_0^{-1}\exp\!\left(-U_0\right)\,d\zeta.
 \label{f16v10:law}
\end{equation}
These are the complete compact conditional and its complete Gaussian
reference, not a restricted core. The Gaussian precision $H$, covariance
$C=H^{-1}$, and mean $m$ satisfy V6's bounds
\[
 h_-I\le H\le h_+I,\qquad \sup_c|m_c|\le C_r.
\]
The covariance is independent of the bridge values; the mean generally is
not. Both depend on the finite geometry and its dressed time ends. No
inverse or Gaussian carrier is changed as the expansion order increases.

By a coefficient in this subsection we mean a finite Taylor coefficient at
$\varepsilon=0$ of the actual exponential-word expression. We do not yet assert a
convergent series. With $V_\varepsilon=U_\varepsilon-\log J_\varepsilon^f$ on its positive chart interior,
write formally
\begin{equation}
 V_\varepsilon\sim U_0+\sum_{q\ge1}\varepsilon^q V_q,\qquad
 r_{\varepsilon,\alpha}\sim\sum_{k\ge1}\varepsilon^k r_{\alpha,k}.
 \label{f16v10:jets}
\end{equation}
Here $r_{\varepsilon,\alpha}=\mathcal D^\varepsilon_\alpha U_\varepsilon-g_\alpha$ is precisely V7's
Haar-compensated residual. The measured coefficients $V_q$ include Haar;
the Haar insertion $\mathcal D^\varepsilon_\alpha U_\varepsilon$ does not acquire a second
Haar logarithmic derivative. The reference divergence in that identity
is Haar divergence.

\begin{proposition}[Actual local action and Haar-source coefficients]
\label{f16v10:local-jets}
For each fixed $q$, $V_q$ is a sum of fixed-range polynomials of total
fast/bridge degree at most $q+2$, with a bounded number of terms incident
to each fast cell. At fixed bounded bridge history their coefficients have
bounds depending on $q,r$, not on $B,n$.
For each $k$, $r_{\alpha,k}$ is the sum of a local bridge coefficient and
the same spatially exponentially local $\mathsf P^{\mathsf T}$ synthesis
of local port coefficients as in V7. These polynomials have degree at most
$k+1$ and occupy only the two port bonds adjacent to the source time,
and their incident internal slices.

In particular, if $U_q$ denotes the classical coefficient, then
\begin{align}
 V_1&=U_1,\notag\\
 V_2&=U_2+\frac13\sum_{i,b,j\ {\rm chord}}w_i|x_{i,b,j}|^2
                  +\frac13\sum_{t,b,v\ {\rm port}}|a_{t,b,v}|^2.
 \label{f16v10:haar-order-two}
\end{align}
For $\eta_\alpha=\mathsf P\mathsf e_\alpha$, with $\mathsf e_\alpha$ the canonical common-source basis vector, the first corrected score is
\begin{equation}
 r_{\alpha,1}
   =(\partial_\alpha+\eta_\alpha\cdot\partial_a)U_1
       +\frac12\sum_p[\eta_{\alpha,p},a_p]\cdot\partial_{a_p}U_0.
 \label{f16v10:first-score}
\end{equation}
The bracket is the Lie bracket in the fixed $\mathfrak{su}(2)$ convention.
In particular the second term in \eqref{f16v10:first-score} is not omitted.
\end{proposition}

\begin{proof}
Each classical summand is a fixed multiple of
$\varepsilon^{-2}s(\prod_j\operatorname{Exp}(\varepsilon T_jw))$, with fixed local linear maps
and the literal product order. Its constant and linear unscaled terms
vanish. The coefficient of $\varepsilon^q$ therefore has degree $q+2$. Endpoint
Gaussian terms are already quadratic. Every differentiated word has a
fixed carrier, with the same uniformly bounded incidence at every order.
For Haar, $-\log j(z)=|z|^2/3+O(|z|^4)$, and the powers $w_i$ are the
ones in V4--V9. This proves \eqref{f16v10:haar-order-two} and the asserted
local structure at higher fixed orders.

For a port, the flow is $h_p\mapsto\operatorname{Exp}(s\varepsilon\eta_p)h_p$.
At $h_p=\operatorname{Exp}(\varepsilon a_p)$, its Taylor jet in logarithmic
coordinates follows from
\[
 \varepsilon^{-1}\log\{\operatorname{Exp}(s\varepsilon\eta_p)\operatorname{Exp}(\varepsilon a_p)\}
   =a_p+s\eta_p+\frac{s\varepsilon}{2}[\eta_p,a_p]+O(\varepsilon^2s)+O(s^2).
\]
This is used only for coefficients at zero, not as a global coordinate
formula for the compact flow. Applying it to $U_\varepsilon$ proves the displayed
first coefficient; the order-zero part is exactly $g_\alpha$ by V7.
Higher coefficients can instead be obtained by inserting
$\operatorname{Exp}(s\varepsilon\eta_p)$ directly in each word and differentiating,
which avoids all logarithm singularities. A coefficient of order $k$ of
that inserted score has degree at most $k+1$. The fixed matrix $\mathsf P$
then synthesizes these primitive port coefficients. Its spatial decay and
two-bond temporal support were proved in V7 and do not deteriorate with $k$.
\end{proof}

There is a simple concrete recipe for $U_1$. Identify $\mathfrak{su}(2)$
with imaginary quaternions as in the inherited action convention. For one
ordered word $\prod_{j=1}^d\operatorname{Exp}(\varepsilon v_j)$,
\begin{equation}
 \varepsilon^{-2}s\left(\prod_j\operatorname{Exp}(\varepsilon v_j)\right)
 =\frac12\left|\sum_jv_j\right|^2
        +\varepsilon\sum_{i<j<k}\det(v_i,v_j,v_k)+O(\varepsilon^2).
 \label{f16v10:word-cubic}
\end{equation}
Indeed the scalar part of $v_iv_jv_k$ is
$-(v_i\mathbin\times v_j)\cdot v_k$; all other cubic scalar terms
vanish. Each negative occurrence is its signed $v_j$, in its actual
position. Summing with the original magnetic and kinetic weights gives
$U_1$. This recipe and \eqref{f16v10:first-score} make the leading memory
coefficient computable from the existing incidence and word data.

\subsection{Connected coefficients: the normalizer is part of every order}

Write $\kappa_0(F_1,\ldots,F_s)$ for the joint cumulant under the
\emph{same} Gaussian law $\mu_0$. It can be defined by the finite moment
partition formula; no exponential integrability of a cubic polynomial
at a nonzero source is assumed. For $j\ge2$ define
\begin{equation}
\begin{split}
 G_j(\alpha,\beta)
 =\sum_{\substack{k,l\ge1\\k+l\le j}}
   \sum_{m'=0}^{j-k-l}\frac{(-1)^{m'}}{m'!}
   \sum_{\substack{q_1,\ldots,q_{m'}\ge1\\
                   q_1+\cdots+q_{m'}=j-k-l}}
 \kappa_0(r_{\alpha,k},r_{\beta,l},V_{q_1},\ldots,V_{q_{m'}}).
\end{split}
 \label{f16v10:connected-coefficients}
\end{equation}
For $m'=0$ the inner sum has one term when $k+l=j$ and is otherwise
zero. The extensive $V_q$ in this formula is expanded as a sum of its
actual local terms when estimating it, not bounded by an extensive norm.

\begin{proposition}[Coefficient ownership and the first two corrections]
\label{f16v10:coefficient-identity}
At each fixed finite $B,n$, the asymptotic coefficients of the exact
normalized covariance $\Gamma_\varepsilon(\alpha,\beta)=
\operatorname{Cov}_{\mu_\varepsilon}(r_{\varepsilon,\alpha},r_{\varepsilon,\beta})$ are
\eqref{f16v10:connected-coefficients}. In particular
\begin{align}
 G_2(\alpha,\beta)&=\operatorname{Cov}_0(r_{\alpha,1},r_{\beta,1}),
 \notag\\
 G_3(\alpha,\beta)&=\operatorname{Cov}_0(r_{\alpha,1},r_{\beta,2})
    +\operatorname{Cov}_0(r_{\alpha,2},r_{\beta,1})
    -\kappa_0(r_{\alpha,1},r_{\beta,1},V_1).
 \label{f16v10:first-coefficients}
\end{align}
The $-\kappa_0$ term is a native measure correction, not an optional
observable correction. Adding any fast-independent scalar to any $V_q$
changes none of these covariance coefficients.
\end{proposition}

\begin{proof}
As an identity of formal series, differentiate twice in the two marking
variables $s,t$ the expression
\[
 \log\mathbb E_0\exp\left\{
 -\sum_{q\ge1}\varepsilon^qV_q+s\sum_{k\ge1}\varepsilon^kr_{\alpha,k}
                          +t\sum_{l\ge1}\varepsilon^lr_{\beta,l}\right\}.
\]
At each fixed coefficient only finitely many Gaussian polynomial moments
occur. The moment partition definition of the logarithm is exactly the
cumulant formula. Selecting the two marks and $m'$ ordered action
insertions gives the factor $(-1)^{m'}/m'!$ and the displayed degree
constraint. This is a formal identity, not a convergent integral of a
cubic exponential. Any cumulant of at least two entries with one constant
entry vanishes; thus fast-independent action scalars cancel, as they must
in a normalized covariance. Source constants cancel for the same reason.

To identify these with actual asymptotic coefficients at fixed $B,n$, use
the complete native pin, a smooth cutoff inside its principal charts, and
finite-order Taylor expansion under the resulting Gaussian majorant.
The detailed short-window version is proved below. The identical
finite-dimensional argument applies here at fixed $B,n$ with constants
allowed to depend on both. The omitted compact and Gaussian tails are
$O(\exp\!\left(-c/\varepsilon^2\right))$ times a finite constant. The partition function has positive
limit $Z_0$, so finite series division is valid. Uniqueness of finite
asymptotic coefficients proves the assertion. No duration-uniform error
is inferred in this identification step.
\end{proof}

For an explicit finite-matrix form of the leading answer, let
$\xi=\zeta-m$ and write the quadratic polynomial
\[
 r_{\alpha,1}=\xi^{\mathsf T}A_\alpha\xi
                         +b_\alpha^{\mathsf T}\xi+c_\alpha,
 \qquad A_\alpha=A_\alpha^{\mathsf T}.
\]
Wick contraction, with all means retained in $b_\alpha,c_\alpha$, gives
\begin{equation}
 G_2(\alpha,\beta)
       =2\Tr(A_\alpha C A_\beta C)+b_\alpha^{\mathsf T}C b_\beta.
 \label{f16v10:leading-matrix}
\end{equation}
The odd centered Gaussian cross terms are zero. No commutation of
$A_\alpha,A_\beta,C$ is required. The full matrix $G_2$ is positive
semidefinite. Set $M_j=-\operatorname{Off}_{>1}G_j$. These are the
coefficients of the actual nonadjacent-time common-source memory by V7's
exact contact support, not new transfers or counterterms.

\subsection{A fixed locality exponent for every finite coefficient}

Give fast cells their ambient space--time distance: nearest time and nearest
spatial cube steps, with a fixed enlargement of the step range when a
literal word meets a diagonal pair. Source coordinates inherit their
cube/time anchor; orientation and colour have fixed multiplicities. The
resulting graph has uniformly polynomial ball growth. Let $d(\alpha,\beta)$
denote its source distance. For a scalar source matrix define
\[
 \|G\|_{\mathrm{loc},b}
   =\max\left\{\sup_\alpha\sum_\beta \exp\!\left(b d(\alpha,\beta)\right)
                                      |G_{\alpha\beta}|,
               \sup_\beta\sum_\alpha \exp\!\left(b d(\alpha,\beta)\right)
                                      |G_{\alpha\beta}|\right\}.
\]
For time blocks acting on the full spatial common-source space set
\begin{equation}
 \|G\|_{\mathrm{time},b}
       =\sup_i\sum_j \exp\!\left(b|i-j|\right)\|G_{ij}\|_{\mathcal E\to\mathcal E}.
 \label{f16v10:time-norm}
\end{equation}
For symmetric block matrices this bounds their operator norm on the whole
history source space. We will use both norms, not identify them.

\begin{theorem}[Uniform linked-coefficient locality]
\label{f16v10:coefficient-locality}
There exists $a>0$, fixed independently of coefficient order, spatial volume
and duration, such that for every fixed $j\ge2$,
\begin{equation}
 \|G_j\|_{\mathrm{loc},4a}\le C_{j,r},\qquad
 \|G_j\|_{\mathrm{time},2a}+\|M_j\|_{\mathrm{time},2a}\le C_{j,r}.
 \label{f16v10:coefficient-norm}
\end{equation}
The constants may grow with $j$. We may fix this $a$ also to satisfy
$2a<\log(19/9)$ and $2a<\log3$.
\end{theorem}

\begin{proof}
The finite-range positive precision admits a norm-convergent Neumann
expansion. As in V5--V6 it gives covariance blocks
$\|C_{cd}\|\le C \exp\!\left(-\nu d(c,d)\right)$ for some fixed $\nu>0$.
The same argument for $L_h$ gives exponentially decaying spatial lift
entries; its temporal support is only the two adjacent bonds. By polynomial
ball growth, covariance and lift entries have bounded weighted absolute
row and column sums at some fixed exponent $\nu_1>0$. Take
$8a<\nu_1$, decreasing $a$ to satisfy the two stated temporal conditions.
Finite colour and coordinate multiplicities only change constants.

Expand each polynomial in \eqref{f16v10:connected-coefficients} in centered
Gaussian coordinates. Local means are bounded by $C_r$, so every
coefficient of these polynomials has a fixed-order bound. Wick's formula
represents their moments by pairings of these coordinates. The cumulant
partition formula cancels exactly every pairing whose graph on the
\emph{polynomial factors} is disconnected: its coefficient is the sum
of the partition coefficients over partitions coarser than its connected
components, which is zero unless there is one component. Self-pairings
within a factor do not connect factors. This argument proves the connected
pairing rule also when some polynomials have nonzero means.

Expand each $V_q$ into local terms. A local source term sits within a fixed
radius of its anchor; a synthesized port-source term has in addition its
lift coefficient, which is treated as an exponentially summable edge
from the source anchor to the port anchor. Every surviving pairing graph,
with these possible additional edges, connects the two marked anchors and
all its interaction anchors. Choose a spanning tree. Each tree edge carries
either one covariance with exponentially decaying block entries or one
lift coefficient. Other contractions have bounded absolute values and
can be absorbed into a constant depending on the finite total degree.
The finite stencil displacements change weights only by fixed-order
constants.

The triangle inequality on the tree path between the two marks gives
\[
 \exp\!\left(4a d(\alpha,\beta)\right)
       \le C_j\prod_{uv\ {\rm in\ the\ tree}}\exp\!\left(4a d(u,v)\right).
\]
Fix $\alpha$ and sum successively over leaf anchors, including $\beta$.
Every such sum is bounded by the same finite weighted edge row bound.
All unmarked interaction sites are summed in this way. At fixed $j$ there
are finitely many degree decompositions, monomials, contractions and trees,
because each action order is at least one and each marked order at least
one. Their number and all polynomial coefficients enter $C_{j,r}$, not
an exponent that worsens with $j$. This proves the first assertion; the
column assertion follows identically or by symmetry.

For fixed times $i,j'$, the block row and column sums are at most
$C_{j,r}\exp\!\left(-4a|i-j'|\right)$. Their Schur bound gives the same block operator
bound. Summing it against $\exp\!\left(2a|i-j'|\right)$ proves the second assertion with
a fixed geometric factor. Masking out contact times cannot increase any
of the absolute row bounds. This proves the assertion for $M_j$ as well.
The spatial dependence of $m(\mathbf y)$ does not create an uncontracted
edge: it is a bounded polynomial coefficient at the fixed base history,
not a variable differentiated in this cumulant estimate.
\end{proof}

This is linked Gaussian coefficient control for this \emph{native history
and its exact Haar-corrected writers}. It is not a native cluster expansion:
no assertion about convergence of the sum over $j$, or about a full native
large-field remainder, has entered its proof.

\subsection{Complete native short-window jets, with their spatial cost paid}

We next prove a finite-order error, rather than merely list formal
coefficients. Use V9's balanced step operators $T_t$, balanced endpoints,
and common Gaussian reference $S_*,a_*$. Their centering means and
balancing multipliers are fixed functions of the base bridge history,
independent of $\varepsilon$. Sources are formed in the original coordinates first,
as in V9. An insertion is the integral of that source observable against
one, two or three balanced augmented kernels, divided by $\lambda_*$ per
bond. All chart and normalization factors belong to these operators.

\begin{proposition}[Finite jets of complete native kernels and insertions]
\label{f16v10:short-jets}
For every fixed integer $J\ge2$ there are $C_{J,r},\varepsilon_{J,r}>0$ such that
the following estimates hold uniformly in $n$, the time location, and the
bounded bridge history when $0<\varepsilon\le \varepsilon_{J,r}$ and
$\varepsilon\,\exp\!\left(C_{J,r}B\right)$ is sufficiently small. Each balanced step and endpoint
has a Taylor polynomial through degree $J$, with zeroth terms $S_*,a_*$,
coefficient norms at most $\exp\!\left(C_{J,r}B\right)$, and remainders at most
\begin{equation}
                    \exp\!\left(C_{J,r}B\right)\varepsilon^{J+1}.
 \label{f16v10:short-error}
\end{equation}
The kernel estimates hold in Hilbert--Schmidt norm and endpoint estimates
in $L^2$. A first corrected-source insertion has such a jet starting at
order one, with an additional factor $\|v\|$. An insertion of the product
of two corrected sources on overlapping windows has a jet starting at
order two, with factor $\|v\|\|w\|$. Its union window has at most three
bonds. The estimates integrate the \emph{complete native domains}.
\end{proposition}

\begin{proof}
Here are the finite-differentiability and tail details. If $f(w)$ is an
unscaled Wilson word action, $f(0)=Df(0)=0$, and
\[
 \varepsilon^{-2}f(\varepsilon w)=\int_0^1(1-s)D^2f(s\varepsilon w)[w,w] ds.
\]
Unitary exponential differentiation bounds every fixed real derivative of
$f$ by a constant depending only on that order and the local word. Hence
its $k$th $\varepsilon$ derivative has absolute value at most $C_k|w|^{k+2}$,
also at $\varepsilon=0$ in this integral representation. Differentiating inserted
words gives the corresponding polynomial bounds for bridge and Haar-port
scores. These are direct insertions, not differentiated principal
logarithms. At most three bonds have dimension $O(B)$; bounded incidence
bounds their summed derivative polynomials by a fixed-order polynomial in
$|w|$ and $B$. The linear lift uses its bounded operator norm. Thus
coherent source directions cost $\|v\|$, or $\|v\|\|w\|$, rather than
an unrecorded number of source coordinates.

Choose a smooth product cutoff equal to one when each original fast group
coordinate has $|\varepsilon z|\le\pi/4$, and zero when any has
$|\varepsilon z|\ge\pi/2$. This cutoff is used only to prove the asymptotic estimate.
The original compact integral is not redefined. On its support all powers
of $j$ and their positive square roots are smooth functions of $ez$;
all their fixed derivatives have polynomial bounds. Differentiating a
cutoff contributes another polynomial in the fast coordinates. Its
nonconstant derivative support lies at $|z|\ge c/\varepsilon$.

V9's complete balanced pin bounds both the native and Gaussian integrands
by $\exp(C_rB-c|w_{\rm centered}|^2)$. It also applies at every real
parameter between zero and $\varepsilon$ wherever the cutoff is nonzero. The bounded
mean per group makes a cutoff or chart complement satisfy
$|w_{\rm centered}|\ge c/\varepsilon-C_r$. The integrals on those complements,
including any fixed polynomial source factor, are therefore bounded by
\[
                     \exp\!\left(C_{J,r}B-c'/\varepsilon^2\right).
\]
This bounds the difference between the cutoff and the \emph{full} native
integrals as well as the relevant Gaussian tails. Cutoff derivatives have
the same flat bound. Derivatives of the regularized whole-space integrand
are a fixed-order polynomial times a Gaussian majorant. Integrating its
Taylor remainder through degree $J$ thus costs
$\exp\!\left(C_{J,r}B\right)\varepsilon^{J+1}$. Fixed-degree Gaussian moments in dimensions
proportional to $B$ cost a fixed polynomial in $B$ times a constant to the
power $B$; both are included in the displayed exponential. Port and
intermediate-slice integrations followed by the endpoint kernel $L^2$
norm obey the same estimate by the Gaussian majorant. The flat term is
absorbed by increasing the fixed-order constant and decreasing $\varepsilon_{J,r}$.

The one-link scalar normalizer is treated by this same finite-dimensional
Laplace argument. Its normalized leading value is nonzero, so reciprocal
powers have finite jets; raising them to $24B$ costs at most an exponential
in $B$ at each fixed order. Endpoint norms are positive integrals with
lower bounds $\exp\!\left(-C_rB\right)$ from a product of one-cube balls, as in V9.
Finite series division for their reciprocal square roots is legitimate
under the stated entry condition. For balanced endpoints use the
\emph{weighted} Gaussian majorant before taking the norm, not an
unbounded exponential multiplier applied to an unweighted error.
The factors $F_{\gamma,2}$, the original endpoint chart mass, and the
bridge half densities are all included this way. The Gaussian balancing
scalars themselves are independent of $\varepsilon$ and are retained exactly.

V7's Gaussian cancellation makes a residual insertion zero at order zero.
A product of two such insertions has no term below order two. This remains
true on overlapping windows and at either time end. The same derivative
and tail estimates with these polynomial insertions prove all assertions.
No analyticity of the native integrals in complex $\varepsilon$ has been asserted.
\end{proof}

\subsection{Analytic dependence on bounded kernel arrays, not on a complex coupling}

Finite jets alone would give a factor proportional to duration if one
expanded an unnormalized long product term by term. The next lemma avoids
that loss by applying V9's exact product normalization before differentiating.
The analytic variable here is an array of \emph{bounded operators}. It is
not a complex-valued compact Wilson coupling.

\begin{proposition}[A uniform analytic two-window functional]
\label{f16v10:analytic-array}
Complexify the Hilbert space of V9. Keep $S_*=P_*+Q_*$ fixed, and let the
fixed real endpoint $a_*$ have $\|a_*\|=1$,
$\langle a_*,\phi_*\rangle\ge\beta>0$. There is an absolute $c>0$
such that on the complex array ball
\begin{equation}
 \sup_t\|T_t-S_*\|<\rho_B,
 \quad\|\ell_L-a_*^{\mathsf T}\|<\rho_B,
 \quad\|a_R-a_*\|<\rho_B,
 \qquad \rho_B=c\beta^4,
 \label{f16v10:array-ball}
\end{equation}
the exact denominator $\ell_LT_0\cdots T_{n-1}a_R$ never vanishes.
Here $\ell_L$ is a bounded complex-linear functional, not an
antiholomorphically varied adjoint vector.

The normalized centered two-window insertion functional is holomorphic
on this ball. For disjoint windows of at most two bonds separated by $g$
unmarked bonds it obeys
\begin{equation}
 |\mathcal F(T,\ell_L,a_R;A_F,A_G)|
       \le C\beta^{-4}\|A_F\|\|A_G\|3^{-g}.
 \label{f16v10:analytic-mixing}
\end{equation}
Normalized one-window expectations of length at most three satisfy the
corresponding bound without $3^{-g}$. Cauchy bounds for derivatives and
Taylor remainders on a smaller array ball cost powers of $\rho_B^{-1}$,
but no number of time slices.
\end{proposition}

\begin{proof}
We give the complex-domain check, which is not a consequence of a real
spectral gap alone. Split every factor relative to
$\mathbb C\phi_*\oplus\phi_*^\perp$. Its scalar block is within
$\delta$ of one, its row and column blocks have norm at most $\delta$,
and its lower block has norm at most $1/5+\delta$, where
$\delta=\sup\|T_t-S_*\|$. In the slope formula use the complex-linear
row functional instead of a conjugated variable. For $\|z\|\le1$,
\[
 |a+cz|\ge1-2\delta,\quad
 \|F(z)\|\le\frac{1/5+2\delta}{1-2\delta},\quad
 \|DF(z)\|<1/3\qquad(\delta\le1/100).
\]
The ball of radius $2\delta$ is invariant by the same scalar inequalities
as V9. Iterate columns and, separately, row functionals. No scalar factor
vanishes. For every interval,
\[
 T_I/p_I=(\phi_*+x_I)(\phi_*^{\mathsf T}+y_I)+E_I,
 \quad |p_I|>0,\quad \|x_I\|,\|y_I\|\le2\delta,
 \quad\|E_I\|\le3^{-|I|},
\]
where $y_I$ is a row functional annihilating $\phi_*$. The vanishing
corner properties are unchanged. The scalar multiplication identity is
$p_{IJ}/(p_Ip_J)=1+y_Ix_J$, whose modulus is between
$1-4\delta^2$ and $1+4\delta^2$. These are exact identities over the
complex field.

For long products, the rank-one endpoint pairing has modulus at least
$(\beta-\eta-4\delta)^2$, whereas the remainder has modulus at most
$4\,3^{-n}$. For short products compare with the positive real number
$\langle a_*,S_*^na_*\rangle\ge\beta^2$ by the telescoping norm
estimate. Exactly the short/long split of V9 gives
\[
 |\ell_LT_{[0,n)}a_R|\ge\frac{\beta^2}{16}|p_{[0,n)}|
\]
when $\delta\le10^{-6}\beta^4$ and
$\eta\le10^{-6}\beta^2$. Choose $c$ smaller than $10^{-6}$.
Positivity is used only for this fixed real reference in the short-product
comparison. The perturbed factors are not claimed to be positive operators
or kernels. The displayed inequality proves nonvanishing for every length.

The centered insertion is the same two-by-two numerator minor as V9,
divided by the square of this actual denominator. A rank-one middle
factor makes that numerator zero algebraically over $\mathbb C$ as well.
The error in the normalized middle product is at most $3^{-g}$.
The scalar quasi-multiplication inequalities, with absolute values, pay the
five segments and the two short windows exactly as in V9. This proves
\eqref{f16v10:analytic-mixing}. A single length-three window has
$|p_I|\ge(1-2\delta)^3>1/2$, so the one-window estimate follows by
the same three-segment calculation.

Products and insertions are continuous multilinear functions of their
bounded operator entries. Division by the nonzero denominator therefore
makes the functional holomorphic. Equip the finite array with its maximum
operator norm, not its sum norm. A complex disc in any unit array direction
of radius comparable to $\rho_B$ remains inside the ball. Cauchy's integral
formula bounds derivatives there by the uniform supremum above times the
corresponding inverse radius powers. The usual polarization, if mixed
derivatives are used, introduces only order-dependent constants. Neither
the radius nor the supremum contains the array length. In particular the
centered $3^{-g}$ factor remains in these derivative bounds.
\end{proof}

This step is deliberately different from assuming a complex strip for the
full compact partition function as a function of $\varepsilon$. The latter has not
been proved here. Real finite jets feed an analytic \emph{normalized
operator-product functional} on an independently established ball.

\subsection{A finite-order native remainder, with no duration loss}

\begin{theorem}[All fixed orders of the native coherent covariance]
\label{f16v10:native-expansion}
Fix $r<\infty$ and an integer $J\ge2$. There are finite
$K_{J,r},\gamma_{J,r}$, independent of $B,n$, such that if
\begin{equation}
 \gamma\ge\gamma_{J,r},\qquad
              \exp\!\left(K_{J,r}B\right)/\sqrt\gamma\le1/100,
 \qquad\max_{i,a}|y_{i,a}|\le r,
 \label{f16v10:entry}
\end{equation}
then for all time blocks
\begin{equation}
 \left\|\Gamma_{\gamma,ij}
       -\sum_{q=2}^{J}\gamma^{-q/2}(G_q)_{ij}\right\|
 \le \exp\!\left(K_{J,r}B\right)\gamma^{-(J+1)/2}
                         3^{-(|i-j|-2)_+}.
 \label{f16v10:block-remainder}
\end{equation}
After enlarging $K_{J,r}$, the same estimate summed in the fixed weighted
temporal norm is
\begin{align}
 \left\|\Gamma_\gamma-\sum_{q=2}^J\gamma^{-q/2}G_q
                             \right\|_{\mathrm{time},a}
       &\le \exp\!\left(K_{J,r}B\right)\gamma^{-(J+1)/2},\notag\\
 \left\|\mathsf M_\gamma-\sum_{q=2}^J\gamma^{-q/2}M_q
                             \right\|_{\mathrm{time},a}
       &\le \exp\!\left(K_{J,r}B\right)\gamma^{-(J+1)/2}.
 \label{f16v10:remainder-norm}
\end{align}
All orders use the same $C,\mathsf P$ and the same locality exponent $a$.
The constants and entry thresholds are allowed to depend on $J$.
\end{theorem}

\begin{proof}
Use V9's balanced representation. Its reference endpoint overlap is at
least $\beta=\exp\!\left(-c_rB\right)$. Proposition~\ref{f16v10:short-jets} supplies
polynomials $T_t^{[J]}(z)$, endpoint polynomials, and insertion polynomials
in an auxiliary \emph{complex} scalar $z$, using their real finite jets.
Their coefficients are bounded by $\exp\!\left(C_{J,r}B\right)$, uniformly over the
entire history array. Choose $R_B=\exp\!\left(-c_{J,r}B\right)$ with $c_{J,r}$ large
enough. On $|z|\le R_B$ all unmarked polynomials and endpoints lie in a
ball strictly smaller than \eqref{f16v10:array-ball}; this follows by
summing finitely many coefficient bounds. First insertion polynomials
have norm at most $\exp\!\left(C_{J,r}B\right)|z|\|v\|$, and overlapping-product
insertion polynomials at most $\exp\!\left(C_{J,r}B\right)|z|^2\|v\|\|w\|$.
The arbitrary-order constants only decrease this auxiliary radius by
another spatial exponential; no duration factor is present.

For separated source times use the analytic centered functional of the
preceding proposition. It has an analytic zero of order at least two at
$z=0$, and on this disc its modulus divided by $|z|^2$ is at most
$\exp\!\left(C_{J,r}B\right)3^{-(|i-j|-2)}\|v\|\|w\|$. Cauchy expansion of that
quotient shows that its remainder after covariance degree $J$ at
$0<\varepsilon\le R_B/2$ is at most
\[
 \exp\!\left(C'_{J,r}B\right)\varepsilon^{J+1}3^{-(|i-j|-2)}\|v\|\|w\|.
\]
This assertion concerns the polynomial kernel arrays evaluated at $\varepsilon$.
To replace them by the \emph{actual} arrays, use the short-jet errors
\eqref{f16v10:short-error} and the derivative bounds on the smaller
array ball. The latter have no duration factor and retain the same
separation factor. Insertion errors are multiplied by the other first
insertion, whose norm is $O(\varepsilon\,\exp\!\left(CB\right))$, and unmarked-array errors
by the two insertion norms. They are no larger than the displayed
$\varepsilon^{J+1}$ budget after increasing its spatial exponent. The real line
segments joining the actual arrays to the polynomial arrays remain inside
the larger ball by \eqref{f16v10:entry}. Thus no derivative of an
uncontrolled long-product error has been taken.

For $|i-j|\le1$, write the covariance as the normalized one-window
expectation of the product of the two sources, minus the product of their
two normalized expectations. The union of their windows has at most three
bonds. The same holomorphic and finite-jet argument gives the displayed
bound with separation factor one. The short-jet assertion covers this
product insertion, including its true lower-order zeros. It also covers
windows cut off by either dressed endpoint.

At each fixed $B,n$, this expansion computes the covariance of the exact
native residuals: balanced coordinates and scalar multipliers are only
changes in the evaluation of that same normalized path law. By
Proposition~\ref{f16v10:coefficient-identity}, uniqueness of fixed-finite
asymptotic coefficients identifies the coefficients with $G_q$, including
the action and normalizer terms. This identification does not transfer a
volume-dependent fixed-finite error; the uniform remainder was just proved
through the analytic array argument. All estimates hold for arbitrary
unit spatial $v,w$, so their supremum gives the block operator estimate.

Since $a<\log3$, the scalar sum
$\sum_j \exp\!\left(a|i-j|\right)3^{-(|i-j|-2)_+}$ is bounded independently of $i,n$.
Sum \eqref{f16v10:block-remainder} and absorb this fixed constant into
$K_{J,r}$. V7's exact identity
$\mathsf M_\gamma=-\operatorname{Off}_{>1}\Gamma_\gamma$ holds for
every coupling. Applying the same mask to each coefficient and to the
block bounds proves the second line of \eqref{f16v10:remainder-norm}.
All coupling-entry, overlap and jet costs are finitely many exponentials
in $B$ at fixed $J$, and are paid by enlarging $K_{J,r}$. None depends on
$n$ or on temporal variation of the admitted bridge history.
\end{proof}

The result gives a precise finite-order correction cycle. Define
$\mathsf M^{[J]}=\sum_{q=2}^J\gamma^{-q/2}M_q$ and
$\mathsf E^{[J]}=\mathsf M_\gamma-\mathsf M^{[J]}$. Then
\begin{equation}
 \mathsf E^{[J+1]}
       =\mathsf E^{[J]}-\gamma^{-(J+1)/2}M_{J+1},\qquad
 \|\mathsf E^{[J+1]}\|_{\mathrm{time},a}
       \le \exp\!\left(K_{J+1,r}B\right)\gamma^{-(J+2)/2}.
 \label{f16v10:cycle}
\end{equation}
The gain is one power of $\varepsilon=\gamma^{-1/2}$ at each fixed cycle. This is
an \emph{absolute finite-order hierarchy}, not the claim
$\|\mathsf E^{[J+1]}\|\le C\varepsilon\|\mathsf E^{[J]}\|$, which need not
hold if a lower-order residual happens to vanish. Every $M_q$ is retained
as a real response of the native theory; none is set to zero by definition.

\subsection{The first explicit correction and a local exported operator}

\begin{theorem}[Leading local native memory with a smaller coherent remainder]
\label{f16v10:first-cycle}
For fixed $r$ there are $K_r,\gamma_r$ such that in the entry regime
\eqref{f16v10:entry} for $J=2$,
\begin{equation}
 \mathsf M_\gamma=\gamma^{-1}M_2+\mathsf E_\gamma,
 \qquad \|M_2\|_{\mathrm{loc},4a}\le C_r,
 \qquad
 \|\mathsf E_\gamma\|_{\mathrm{time},a}
                           \le \exp\!\left(K_rB\right)\gamma^{-3/2}.
 \label{f16v10:first-native}
\end{equation}
The analogous expansion holds for $\Gamma_\gamma$ with $G_2$ in place
of $M_2$. In particular, throughout this regime,
\begin{equation}
 \|\Gamma_\gamma\|_{\mathrm{time},a}
              +\|\mathsf M_\gamma\|_{\mathrm{time},a}
                           \le C_r\gamma^{-1}.
 \label{f16v10:improved-coherent}
\end{equation}
The constants in this conclusion are independent of $B,n$ \emph{within
that coupling-dependent spatial regime}, not for unrestricted volumes.
If $M_2^{[L]}$ keeps only source entries at space--time distance at most
$L$, then
\begin{equation}
 \|\mathsf M_\gamma-\gamma^{-1}M_2^{[L]}\|_{\rm op}
       \le C_r\gamma^{-1}\exp\!\left(-4aL\right)
                              +\exp\!\left(K_rB\right)\gamma^{-3/2}.
 \label{f16v10:local-export}
\end{equation}
\end{theorem}

\begin{proof}
Apply the two preceding theorems with $J=2$. The leading coefficient is
explicitly \eqref{f16v10:leading-matrix} with the Haar term from
\eqref{f16v10:first-score}. Since
$\exp\!\left(K_rB\right)\gamma^{-1/2}\le1/100$, the remainder in
\eqref{f16v10:first-native} is at most $\gamma^{-1}/100$.
The uniform leading temporal row bound gives
\eqref{f16v10:improved-coherent}. Truncating its spatially and temporally
weighted entry rows has absolute row and column error at most
$C_r\exp\!\left(-4aL\right)$, by \eqref{f16v10:coefficient-norm}. Schur's bound and
the native temporal remainder prove \eqref{f16v10:local-export}.
This truncation changes the exported approximation, not the exact action,
measure, source definition or core event. Finite range here refers to pairs
of source indices. Coefficient values can still depend on the whole base
history through $m(\mathbf y)$; replacing exterior data requires its own
response estimate.
\end{proof}

For example, if $B\le(\log\gamma)/(4K_r)$, then the native remainder in
\eqref{f16v10:first-native} is at most $\gamma^{-5/4}$. The leading
memory and coherent covariance are $O_r(\gamma^{-1})$, for \emph{all}
history lengths and bounded histories. This improves the prior
$O_r(\gamma^{-3/4})$ mesoscopic bound without declaring its exponential
spatial error to be a uniform one.

The original full/core correction is also still controlled. With the same
fixed small $a$, V4's bound and V7's exact full/core identity give
\begin{equation}
 \|\mathsf C_{\gamma,R}\|_{\mathrm{time},a}
                   \le C_r\gamma^{-1}+C R^2/\gamma
 \label{f16v10:core-correction}
\end{equation}
for an admissible original core radius. With $R=4\sqrt{\log\gamma}$ this
is $O_r(\log\gamma/\gamma)$ in the stated spatial regime, uniformly in
$n$. The weighting is permissible because $a<\log(19/9)$. No expansion
or new event for $p_{\rm good}$ has been substituted. This last inequality
uses the existing core estimate rather than claiming a new all-order
expansion of a moving or growing cutoff boundary.

\subsection{Compatibility, the spatial obstruction, and the next useful step}

The finite-order hierarchy has three different safeguards. First, all
coefficients are those of the same fixed normalized compact conditional;
connected cumulants retain the actual source-dependent normalization.
For instance $-\log(Z_\varepsilon/Z_0)$ starts with
$\varepsilon\mathbb E_0V_1+\varepsilon^2(\mathbb E_0V_2-\operatorname{Var}_0(V_1)/2)$
at fixed finite volume. That scalar is not arbitrarily made zero.
Second, the changed contact matrix from V7 remains part of the exact
Hessian. Our expansion of nonadjacent times uses its \emph{exact} range,
not a missing contact correction. Third, all common-source coordinates
remain present; covariance under simultaneous colour rotations and spatial
translations, where the base history has those symmetries, is preserved by
Taylor coefficients of the exact expressions. No new physical Gauss
restriction is inferred from coordinate covariance.

The positive-cone suggestion has a precise limit. $G_2\succeq0$ is an
actual Gaussian covariance, but $M_2=-\operatorname{Off}_{>1}G_2$ has
zero time-diagonal and hence zero trace. A nonzero real symmetric
zero-trace matrix cannot be positive semidefinite. Therefore a nonzero
$M_2$ cannot be a positive combination of positive native sandwiches.
Higher $G_j$ may themselves be signed. A truncated series is not asserted
to define a positive transfer or probability law. Positivity stays with
the exact native law and its complete covariance.

What remains is not resolved by increasing $J$ formally. Although every
coefficient has a uniform local row bound, the proved full-native remainder
is only
\[
             \exp\!\left(K_{J,r}B\right)\gamma^{-(J+1)/2}
\]
in a temporal block norm, and its entry condition also depends on $J$.
For any fixed $J$ this does not allow $B\to\infty$ at fixed $\gamma$.
No common all-order entry domain, growth bound on $C_{j,r}$, convergent
series, or unique identification from a cutoff-summed asymptotic series
has been proved. A flat remainder invisible to every fixed Taylor
coefficient can still matter outside the admitted spatial regime. Unlike
a construction in which the forcing can be selected, the native Wilson
object cannot be replaced by a different realization having the same jets.

The new concrete input for the spatial problem is the local operator
$M_2$ and, at any fixed higher order, its linked successors. The next target
is a \emph{boundary-conditioned marked remainder estimate} with a retained
spatial collar, rather than another global Gaussian inverse or temporal
product lemma. It must pay the actual boundary and endpoint messages while
replacing the spatial exponential in the remainder by a source-local
bound. These coefficients allow that remainder to start beyond a chosen
finite order without forgetting what was subtracted. They do not provide
arbitrary-boundary compact convexity or a source-visible capacity bound.
The proposed $o(a_t)$ spectral handoff additionally requires an identified
physical-time scale and a comparison in the norm of the same physical
transfer; our source-coordinate covariance norm is not that norm.
Gauge-invariant bridge coverage, infrared contact coercivity and
interacting continuum reconstruction remain separate obligations.

\subsection{Diagnostics, preservation, and scope}

The accompanying diagnostic is
\begin{center}\small
\path{verify_ym_f16_v10_finite_order_memory.py}.
\end{center}
Its report separates exact quaternion/BCH and Haar coefficients, correlated
Gaussian polynomial moments, connected normalization coefficients, complex
near-reference products, finite path insertion jets, and the order/locality
bookkeeping. All 608 exact assertions passed; the assertion families are supplied in
the package report. The checker SHA--256 is
\begin{center}\scriptsize\ttfamily
F6F55A4534AFBA1A64E094B1D96D52CB6D5473762E3633C29E8998182BAA97B0.
\end{center} These checks are not formal verification of the analytic
kernel-jet, linked-coefficient, or Cauchy estimates. The native derivative
suprema and stage-dependent constants are proved finite but are not
numerically optimized. The NS Lean development has not been run here.

The V9 predecessor has 309,137 bytes and SHA--256
\begin{center}\scriptsize\ttfamily
26FF80027553957199D6DADEBBD7E5C6BE28E0F6B97FE5E2D070ECB57362BC50.
\end{center}
The cumulative V10 changes only V9's title version and inserts this exact
appendix before its final \verb|\end{document}|. The package records the
byte-for-byte predecessor recovery, unchanged input hashes, compilation,
render review and reruns of the supplied V4--V9 diagnostics. Those
unchanged scripts passed 9,467, 89,339, 906, 311, 609 and 990 assertions,
respectively. Unavailable
historical or companion scripts are not described as rerun. The next
scheduled five-version companion checkpoint is V15.

\begin{firewall}
V10 proves fixed-order, spatially local linked coefficients for the actual
Haar-corrected native memory, and a finite-order full-native remainder
uniform in duration with an explicit exponential spatial cost. The first
coefficient yields an $O_r(\gamma^{-1})$ coherent bound and a local exported
memory operator in the admitted logarithmic spatial regime. The original
full/core correction remains controlled there. The native measure and its
normalizers are unchanged. No unrestricted-volume remainder, arbitrary
spatial-boundary gluing theorem, convergent all-order realization,
gauge-invariant bridge coverage, physical transfer gap, or interacting
continuum Yang--Mills mass gap is established by this appendix.
\end{firewall}
% END F16 V10 APPEND

% BEGIN F16 V11 APPEND -- 2026-09-17
% Insert immediately before the final \end{document} in cumulative V10.
% No inherited theorem is edited. The cumulative title becomes V11.

\clearpage
\section{V11: complete conditional cylinders and spatial screening of the marked remainder}
\label{f16v11:context}

V10 localizes the coefficients, not the full native remainder, and explicitly
leaves exterior-data replacement unpaid. We now integrate a \emph{proper
spatial cylinder}, with its actual exterior fast history held fixed. All
crossing Wilson terms are retained. The interior still ranges over its
complete compact coordinates. Its comparison cost depends on the number of
\emph{interior cubes}, not on the size of the ambient torus. A boundary-marked
Gaussian comparison then identifies its leading memory with the already
computed full-history coefficient, up to an exponentially small collar error.
This gives a native conditional screening statement, not just a coefficient
with a finite matrix carrier.

Two limitations are declared at the outset. The exterior values that enter
the cylinder must lie in a fixed \emph{rescaled fast-field window}, at every
time. We do not prove that this event is typical under the actual exterior
posterior. Also a full covariance is not the average of conditional
covariances: the covariance of the conditional means must be retained. The
last subsection writes that additional return with its actual measure.

\subsection{Audit and the new interface}

The V10 appendix was read in full, and the supplied f14/f15 and monograph
inventories were searched with targeted passages read. In particular, f15
V13--V14 already prove boundary response and posterior-weighted transfer for
another source and another conditioned chart. The monograph's
\texttt{thm:YM-Wilson-blanket} and
\texttt{thm:YM-boundary-mediation} already record the Markov-blanket and
conditional-mean mechanisms. These identities are inherited tools, not the
new theorem. No earlier all-good probability or native functional inequality
is transferred to the present history law.

The new objects are the complete compact spatial-cylinder conditional of
V5's augmented history, its \emph{interior-port} Haar lift, and its
source-marked expansion with retained crossing interactions and dressed ends.
V9's normalized product lemma and V10's short-jet method are reused after
checking their hypotheses for this different conditional operator. Only the
first nonzero memory coefficient is needed. Repeating the all-order temporal
construction would not resolve the spatial issue.

The import memorandum's separation of order gain from locality cost is kept:
one power of the native expansion parameter is gained after extracting the
leading coefficient, while the cost is localized to a declared cylinder.
There is no new appeal to an NS theorem or to convergence of a formal series.
The inverse-decay argument used below is the elementary finite-range
positive-matrix argument already used in these notes; for classical context,
see S.~Demko, W.~F.~Moss and P.~W.~Smith,
\href{https://doi.org/10.1090/S0025-5718-1984-0758197-9}
{\emph{Decay rates for inverses of band matrices}}, Math. Comp. 43 (1984),
491--499. The boundary identities needed here are proved explicitly.

\subsection{The exact cylinder: condition the exterior, do not delete it}

Fix the ambient cube torus $\mathcal B=(\mathbb Z/m\mathbb Z)^3$, $B=m^3$,
$m\ge4$, and a nonempty proper set $D\subset\mathcal B$ of $M=|D|$ cubes.
The cylinder consists of all $x_{i,b}$, $0\le i\le n$, and all unmoving
$a_{t,b}$, $0\le t<n$, with $b\in D$. These are V5's five chord and seven
nonroot temporal variables per cube. Write $I$ for these free real
coordinate indices, and $O=I^c$ for the exterior indices. V5's unmoving
coordinates are important: $h_t=\operatorname{Exp}(a_t/\sqrt\gamma)$ has
independent rooted tree pins even after the exterior is fixed.

Let $\mathbf q=\zeta_O$ be the exterior fast data, held fixed in these
unmoving coordinates. Choose once and for all a geometric range $r_0\ge4$
large enough to contain the cube carriers of every primitive action word
and of every unlifted bridge mark. Let $\partial^+D$ be the exterior cubes
at spatial graph distance at most $r_0$ from $D$. It is a safe enlargement
of the actual collar. Assume
\begin{equation}
 \max_{i,e}|y_{i,e}|\le r,\qquad
 \sup_{c:\,b(c)\in\partial^+D}|q_c|\le Q,
 \qquad r,Q<\infty .
 \label{f16v11:windows}
\end{equation}
Here a cell vector has at most 36 real coordinates, with the final-time
cell having 15. The supremum includes every time and both species when
present. Exterior cubes beyond this collar are arbitrary. Constants
$C_{r,Q},K_{r,Q}$ will not depend on $B,M,n,D$, or on the individual histories.
Their dependence on the two fixed windows is not suppressed.

Split the exact classical exponent of V5 as
\begin{equation}
 U_\gamma(\zeta;\mathbf y)
  =U_{\gamma,D}(\zeta_I;\mathbf q,\mathbf y)
       +U_{\gamma,O}(\mathbf q;\mathbf y).
 \label{f16v11:action-split}
\end{equation}
Assign to the first term all anchors on free cells and all other primitive
terms whose fast carrier meets $I$; assign the remaining terms to the
second. Pure bridge terms can be assigned by their geometric carrier.
For definiteness assign every such term with carrier meeting $D$ to the
first term. Each original term is assigned once. Let $J^f_{\gamma,D}$ be
the product of V5's free chord Haar powers $w_i$ and free port Haar powers
one. Define
\begin{equation}
 Z_{\gamma,D}(\mathbf q;\mathbf y)
 =\int_{\mathcal P'_I}J^f_{\gamma,D}
                       e^{-U_{\gamma,D}}\,d\zeta_I,
 \qquad
 d\mu_{\gamma,D}^{\mathbf q}
 =\frac{1_{\mathcal P'_I}J^f_{\gamma,D}e^{-U_{\gamma,D}}}
        {Z_{\gamma,D}(\mathbf q;\mathbf y)}\,d\zeta_I.
 \label{f16v11:conditional-law}
\end{equation}
The domain is the product of the \emph{complete} principal charts for the
free groups. In particular this is not V4's core. Fast-independent factors
assigned in \eqref{f16v11:action-split} are part of $Z_{\gamma,D}$; they
cancel only in its normalized conditional probability.

Let $\mathcal A_D$ be the common-bridge basis marks whose spatial anchors
have distance greater than $r_0$ from $D^c$, at every available time. Their
entire unlifted action-source carrier is interior. The part
$U_{\gamma,O}$ has zero derivatives in these directions. We consider the
associated source space $\mathcal U_D$, with dimension at most $9M(n+1)$.
Empty source sets give vacuous estimates, not an exceptional case.

\begin{proposition}[An actual conditional with a surviving independent pin]
\label{f16v11:conditional-pin}
The measure \eqref{f16v11:conditional-law} is the regular conditional of
V6's full native fast-history law given $\zeta_O=\mathbf q$, with its
explicit positive-density version at every interior chart point. It depends
on the exterior only through the actual word collar, hence through
$\partial^+D$. All crossing interactions are retained, and
\begin{equation}
 U_{\gamma,D}\ge c_{\rm pin}\|\zeta_I\|^2,
 \qquad c_{\rm pin}=(6\pi^2)^{-1},
 \label{f16v11:pin}
\end{equation}
for its indicated nonnegative-term assignment, on the complete free charts.
The same pin holds for its quadratic limit. The number of primitive terms
per free cell and the real dimension per slice are bounded independently
of the ambient torus.
\end{proposition}
\begin{proof}
Product Haar measure factorizes between $I$ and $O$. Its chord half-density
powers at the two dressed ends do so as well. Substitution of
\eqref{f16v11:action-split} in the full integral gives
\eqref{f16v11:conditional-law} by ordinary normalization. Local word
support proves the collar assertion, including the moving frames, each
of which uses only the five chords of its own cube.

Every free cube retains all its internal magnetic plaquettes with weights
$w_i\ge1/2$, its two endpoint Gaussian quadratics when applicable, and
its seven tree kinetic terms on each free port bond. In unmoving variables
these last words are exactly $h_s^{-1}h_t$ within that cube, with root
$h_o=I$. V5's cube inequality and rooted depth-sum estimate therefore give
$|x|^2/(3\pi^2)$ and $|a|^2/(6\pi^2)$ on that cube. Exterior fixing changes
neither. All other assigned Wilson terms, including crossing terms with
fixed exterior entries, are nonnegative. Add these estimates. The quadratic
limit preserves the inequality. Each word touches a bounded number of
cubes, and each cube meets a bounded number of words, so all counts involve
$M$, not $B$.
\end{proof}

The pin is \emph{not} a claim that the full conditional potential is globally
convex. Its usefulness is complete-chart integration after the exterior has
been fixed. Also the kinetic normalizers of the original full transfer and
its external $F_{\gamma,2}$ factors have not changed: they cancel from the
fast conditional. They remain in the ambient path weight.

\subsection{Conditional Gaussian data and boundary-reaching inverse identities}

Use the same full Gaussian $N,S,\Lambda,H,C,m$ as V9--V10. Restrict the
spatial internal and port matrices by cube labels; subscripts $D,O$ in
spatial formulas denote these restrictions. In particular
\[
 S_D=S_{DD}=\mathsf H_D+A_D^{\mathsf T}A_D,\qquad
 L_D=(L_h)_{DD},\qquad E_D=E_{\rm br}|_{\text{port columns in }D}.
\]
The columns of $A_D$ are restricted, but its rows include \emph{all}
mixed plaquettes that meet those columns. Crossing plaquettes therefore
contribute to $S_D$. This is not the matrix obtained by deleting them.

The conditional Gaussian in the original unmoving coordinates has precision
$H_{II}$, covariance $C^D=(H_{II})^{-1}$, and mean
\begin{equation}
 m^D=m_I-C^DH_{IO}(\mathbf q-m_O).
 \label{f16v11:Gaussian-conditional}
\end{equation}
It is important that $C^D\ne C_{II}$ in general. Write, at exterior ports,
\[
 q^z_t=q^a_t-\mathsf F_O(q^x_t-q^x_{t+1}).
\]
This is the \emph{Gaussian linear} moving coordinate of the fixed exterior;
it is not asserted to be the exact finite-coupling logarithm of the moving
group. Its cell norms are bounded by $C Q$.

Inside $D$, change Gaussian variables by
$a_{t,D}=z_{t,D}+\mathsf F_D(x_{t,D}-x_{t+1,D})$. Conditional on the exterior,
the fast Gaussian precision separates as
\begin{equation}
 W_n^D\ \oplus\ (I_n\otimes L_D),
 \label{f16v11:separated-precision}
\end{equation}
where $W_n^D$ is the spatial principal restriction of V3's $W_n$. Its
internal and port linear sources, respectively, are
\begin{equation}
 b^x_i=\Lambda_Dy_i+(A^{\mathsf T}A)_{DO}q^x_i,
 \qquad
 b^z_t=E_D^{\mathsf T}(y_t-y_{t+1})+(L_h)_{DO}q^z_t.
 \label{f16v11:boundary-sources}
\end{equation}
The free internal Gaussian mean is $-(W_n^D)^{-1}\mathbf b^x$. All its
endpoint entries are retained; in particular $W_n^D$ is not made periodic
in time.

\begin{proposition}[Uniform conditional means, lifts, and screening identities]
\label{f16v11:Gaussian-screening}
There exist $c,C,\nu>0$, independent of $D,B,n$, such that all the indicated
conditional precisions and covariances have the inherited positive spectral
bounds, and
\begin{equation}
 \sup_{c\in I}|m^D_c|\le C_{r,Q},\qquad
 |m^D_c-m_c|\le C_{r,Q}e^{-\nu s_D(c)}.
 \label{f16v11:mean-screening}
\end{equation}
Here $s_D(c)$ is the spatial distance of its cube to $D^c$. The covariance
and lift differences are given by boundary-reaching identities
\begin{equation}
 \begin{aligned}
 C_{II}-C^D&=-C^D H_{IO}C_{OI},\\
 (L_h^{-1})_{DD}-L_D^{-1}
       &=-L_D^{-1}(L_h)_{DO}(L_h^{-1})_{OD}.
 \end{aligned}
 \label{f16v11:resolvent-boundary}
\end{equation}
Every factor inverse in these formulas has uniformly summable exponentially
weighted block rows and columns at a common smaller exponent. In particular
\begin{equation}
 \sum_{d\in I}e^{\nu d(c,d)/4}
       \|(C_{II}-C^D)_{cd}\|
       \le C e^{-\nu s_D(c)/4}.
 \label{f16v11:covariance-screening}
\end{equation}
The analogous spatial estimate holds for the port inverse, including its
zero extension outside $D$ when its row starts inside $D$.
\end{proposition}
\begin{proof}
Conditioning a positive Gaussian quadratic gives
\eqref{f16v11:Gaussian-conditional} by completing its $I$ square. The
block-diagonal, determinant-one change between unmoving and moving variables
is local to each spatial cube and uses only adjacent times. It therefore
commutes with the division into entire spatial cylinders. V4's separated
Gaussian square, with the exterior fixed, gives
\eqref{f16v11:separated-precision}--\eqref{f16v11:boundary-sources}.
Principal restriction preserves the bounds $h_-I\le H\le h_+I$ and the
spatially whitened bounds $4I\le\widehat W_n\le15I$. Similarly
$\lambda_-I\le L_D\le\lambda_+I$, using V7's constants. Also
$4I\le\mathsf M_D^{1/2}S_D\mathsf M_D^{1/2}\le10I$.

For any of these finite-range positive matrices $K$, with fixed spectral
bounds $k_-I\le K\le k_+I$, write
\[
 K^{-1}=k_+^{-1}\sum_{j\ge0}(I-K/k_+)^j.
\]
The summands have norm at most $(1-k_-/k_+)^j$ and range at most $j$ in
an appropriate fixed enlargement of the ambient cube--time graph. A ball
has at most $C(j+1)^4$ labels even with periodic spatial directions and
arbitrary time length. This proves exponential inverse block decay and
weighted row/column summability, also for every principal restriction.
This is a bound in the ambient distance; disconnected restricted regions
do not invalidate it. Local sources in
\eqref{f16v11:boundary-sources} are bounded by $C_{r,Q}$ per cell. The
row bound proves the first estimate in \eqref{f16v11:mean-screening}.

The $II$ block of $HC=I$ reads
$H_{II}C_{II}+H_{IO}C_{OI}=I$; multiply by $C^D$. This proves the first
identity in \eqref{f16v11:resolvent-boundary}; the second is identical.
The only nonzero entries of $H_{IO}$ join the fixed-width boundary.
Every product term in the first identity therefore travels from $c$ to
that boundary, and then to $d$. The sum of its two inverse-edge lengths,
up to a fixed range shift, is at least $s_D(c)$ and at least $d(c,d)$.
Reserve a quarter of the inverse exponent for each of these two distances
and sum the remaining factors by their weighted row bounds. This proves
\eqref{f16v11:covariance-screening} with constants enlarged for the range
shift. For the mean use \eqref{f16v11:Gaussian-conditional} and the same
boundary path, with $|q_o-m_o|\le C_{r,Q}$ on the actual collar. No value
of $q$ away from that collar is used. The port proof is spatially the same;
its exterior columns start beyond the boundary and have the same tail
bound. Choose a single smaller $\nu$ for the finitely many matrices.
\end{proof}

\subsection{A conditional Haar lift: do not differentiate a frozen port}

V7's unrestricted lift generally varies ports outside $D$. It cannot be
inserted into a conditional integration by parts after those ports have
been fixed. Use instead
\begin{equation}
 (\mathsf P_D\mathbf v)_t
     =-L_D^{-1}E_D^{\mathsf T}(v_t-v_{t+1}),
 \qquad \mathbf v\in\mathcal U_D.
 \label{f16v11:Dirichlet-lift}
\end{equation}
It acts only on the free \emph{unmoving} temporal groups $h_{t,b,v}$.
Its $\ell^2$ norm is uniformly bounded, its time support is the two adjacent
bonds, and its spatial coefficients have a uniform exponential row bound.
These follow from the preceding proposition and $\|d_t\|\le2$.
For a canonical source $\alpha$ let
\[
 \mathcal D^D_\alpha
  =\partial_\alpha+
     \sum_{p\in D,k}(\mathsf P_D)_{p,k;\alpha}X^\gamma_{p,k},
 \qquad
 g^D_\alpha=\partial_\alpha[-\log Z_{0,D}(\mathbf q;\mathbf y)],
 \qquad
 r^D_{\gamma,\alpha}=\mathcal D^D_\alpha U_{\gamma,D}-g^D_\alpha.
\]
Here $X^\gamma$ is V7's actual compact left-translation derivative. The
coefficients of $\mathsf P_D$ are independent of $\mathbf q,\mathbf y$
and all fast variables. All derivatives keep the exterior \emph{unmoving}
coordinates fixed.

\begin{theorem}[Exact conditional source representation]
\label{f16v11:Haar-conditional}
The Gaussian part of the lifted source is deterministic:
\begin{equation}
 (\partial_\alpha+(\mathsf P_D\mathsf e_\alpha)\cdot\partial_{a_D})
                          U_{0,D}=g^D_\alpha.
 \label{f16v11:Gaussian-cancellation}
\end{equation}
With $\Gamma^D_\gamma(\mathbf q)=
\operatorname{Cov}_{\mu_{\gamma,D}^{\mathbf q}}(\mathbf r^D_\gamma)$,
the exact conditional Hessian is
\begin{equation}
 \nabla^2_{\mathcal U_D}[-\log Z_{\gamma,D}]
       =\mathcal K^D_\gamma-\Gamma^D_\gamma,
 \qquad
 (\mathcal K^D_\gamma)_{ij}=0\quad (|i-j|\ge2).
 \label{f16v11:conditional-contact}
\end{equation}
The contact is the symmetrized expectation of
$\mathcal D^D_\alpha\mathcal D^D_\beta U_{\gamma,D}$, not the old direct
Hessian. In particular the conditional nonadjacent-time memory is exactly
$\mathsf M^D_\gamma(\mathbf q)=-\operatorname{Off}_{>1}\Gamma^D_\gamma(\mathbf q)$.
\end{theorem}
\begin{proof}
In Gaussian moving variables, common-source magnetic differentiation has
no free $x$ coefficient, because $\Lambda\mathcal E=0$ before restriction.
All action terms involving a free $x$ were retained, so exterior fixing
does not change that cancellation. The free port coefficient in a direction
$\mathbf v$ is $\sum_t(v_t-v_{t+1})^{\mathsf T}E_D z_{t,D}$.
The port Gaussian gradient is
$L_D z_{t,D}+b^z_t$. Substitution of
\eqref{f16v11:Dirichlet-lift} cancels every free port coefficient. The
remaining quantity is independent of the free variables. Its expectation
is the derivative of $-\log Z_{0,D}$, because the mean Gaussian score
$\partial_{a_D}U_{0,D}$ is zero. This proves
\eqref{f16v11:Gaussian-cancellation}, including its scalar value.

At finite coupling, condition the exterior first and perform only the free
port integrations in Haar variables. The field in
\eqref{f16v11:Dirichlet-lift} has zero Haar divergence. It does not vary
free chord weights or endpoint amplitudes. Thus
$\mathbb E_D[V_\alpha F]=\mathbb E_D[F V_\alpha U_{\gamma,D}]$ and
$\mathbb E_D V_\alpha U_{\gamma,D}=0$. Differentiating a normalized
expectation and using this identity is precisely V7's argument, with the
sum of ports now restricted to $D$. It gives
\eqref{f16v11:conditional-contact} with the symmetrized contact; it does
not assume that two group derivatives commute. Each kinetic word uses one
port bond and its two adjacent bridge times, even with fixed exterior
entries. Each magnetic word uses one bridge time. The two lift bond sets
are disjoint at nonadjacent source times, so the corresponding double
contact vanishes exactly. At every finite regulator the compact derivatives
and integrable chord weights justify differentiation, with no cut-boundary
term from the port logarithms.
\end{proof}

For an interior source the Gaussian scalar in
\eqref{f16v11:Gaussian-cancellation} can also be written explicitly. With
$E_O$ the exterior port columns of $E_{\rm br}$, it is the coefficient of
$\mathbf v$ in
\begin{align}
 g^D_{\mathbf v}
 ={}&\sum_i v_i^{\mathsf T}D^{\mathsf T}Dy_i\notag\\
 &+\sum_t(v_t-v_{t+1})^{\mathsf T}
 \left[(I-E_D L_D^{-1}E_D^{\mathsf T})(y_t-y_{t+1})
       +(E_O-E_D L_D^{-1}(L_h)_{DO})q^z_t\right].
 \label{f16v11:conditional-g}
\end{align}
All terms differentiated by these interior marks belong to $U_{\gamma,D}$.
There is no missing derivative of $U_{\gamma,O}$. The formula also shows
that the Gaussian exterior source message uses only the adjacent port
bonds. By the port inverse boundary identity,
\begin{equation}
 |g^D_\alpha-g_\alpha|\le C_{r,Q}e^{-\nu s_D(\alpha)/C},
 \label{f16v11:g-screen}
\end{equation}
where a fixed source-carrier range shift is absorbed into $C$. This is a
Gaussian screening estimate, not yet an estimate on the nonlinear mean.

\subsection{The leading coefficient with the exterior retained}

Put $\varepsilon=\gamma^{-1/2}$ as in V10. Taylor coefficients are taken
at fixed rescaled $\mathbf q,\mathbf y$. The exact conditional lifted score
has $r^D_{\gamma,\alpha}=\varepsilon r^D_{\alpha,1}+O(\varepsilon^2)$
at each fixed finite dimension. Its first coefficient is given by V10's
ordered-word and Haar-insertion formula, using $U_{1,D}$, the free port
variables, and $\mathsf P_D$. All frozen exterior coordinates remain in
the polynomial coefficients. For example the half-bracket term is
\[
 \frac12\sum_{p\in D}[\eta^D_{\alpha,p},a_p]
                                    \cdot\partial_{a_p}U_{0,D},
 \qquad \eta^D_\alpha=\mathsf P_D\mathsf e_\alpha.
\]
It is not discarded at a boundary port. Define
\[
 G^D_2(\mathbf q)_{\alpha\beta}
   =\operatorname{Cov}_{\mu^\mathbf q_{0,D}}
                     (r^D_{\alpha,1},r^D_{\beta,1}),\qquad
 M^D_2(\mathbf q)=-\operatorname{Off}_{>1}G^D_2(\mathbf q).
\]
If $\xi_D=\zeta_I-m^D$ and
$r^D_{\alpha,1}=\xi_D^{\mathsf T}A^D_\alpha\xi_D+
 b^{D\mathsf T}_\alpha\xi_D+c^D_\alpha$, the exact evaluation is
\begin{equation}
 G^D_2(\mathbf q)_{\alpha\beta}
 =2\Tr(A^D_\alpha C^D A^D_\beta C^D)
                  +b^{D\mathsf T}_\alpha C^D b^D_\beta.
 \label{f16v11:leading-conditional}
\end{equation}
The matrices can be noncommuting. This is the old Wick identity on the
\emph{new conditional} covariance and mean, not on $C_{II}$ and $m_I$.

Let $\mathcal A\subset\mathcal A_D$ have all its spatial anchors at least
$L+r_0$ from $D^c$, and let $\chi_{\mathcal A}$ project onto these sources
at every time. Use the source-distance and time-block norms from V10.

\begin{theorem}[Boundary-marked screening of the leading memory coefficient]
\label{f16v11:coefficient-screening}
There is $b>0$, fixed independently of $D,B,n,r,Q$, which can be chosen
smaller than V10's $a$, such that
\begin{equation}
 \|G^D_2(\mathbf q)\|_{\mathrm{loc},4b}
       +\|M^D_2(\mathbf q)\|_{\mathrm{loc},4b}\le C_{r,Q}.
 \label{f16v11:conditional-coefficient-locality}
\end{equation}
Moreover, with $G_2,M_2$ the \emph{full-history} coefficients of V10 at
this same bridge history,
\begin{align}
 &\left\|\chi_{\mathcal A}
   (G^D_2(\mathbf q)-G_2)\chi_{\mathcal A}\right\|_{\mathrm{time},b}
       \le C_{r,Q}e^{-bL},\notag\\
 &\left\|\chi_{\mathcal A}
   (M^D_2(\mathbf q)-M_2)\chi_{\mathcal A}\right\|_{\mathrm{time},b}
       \le C_{r,Q}e^{-bL}.
 \label{f16v11:coefficient-screening-bound}
\end{align}
The same estimate, with twice the constant, compares any two exterior
histories satisfying \eqref{f16v11:windows}. No bound on their difference
or temporal variation is required.
\end{theorem}
\begin{proof}
Use centered Gaussian coordinates in both polynomials. Their primitive
bridge and port pieces have bounded degree two, fixed local support and
bounded coefficients; the conditional means are bounded by the preceding
proposition. The only extended synthesis is by $\mathsf P_D^{\mathsf T}$,
whose entries have uniformly summable spatial exponential rows and two-bond
time support. Wick contraction gives one or two covariance edges between
the two nonconstant factors. It must connect them: internal contractions
and constant pieces cancel in a covariance. The same weighted-tree
summation as V10, now with $C^D$, proves
\eqref{f16v11:conditional-coefficient-locality} at a sufficiently small
fixed $b$. This argument has only finitely many quadratic contractions.

We describe the additional \emph{boundary mark} needed for comparison,
rather than inferring boundary independence from matrix-index locality.
Expand the full and conditional quadratic source polynomials as their
primitive local terms followed by their port-lift synthesis. An interior
primitive term with all inputs away from the boundary is the same ordered
word coefficient in both expressions. The possible differences are:

\begin{enumerate}
\item a synthesized port term outside $D$, or a primitive term with an
      exterior input replaced by its fixed value;
\item replacement of a full lift entry by its restricted inverse entry;
\item replacement of a mean $m_c$ by $m^D_c$;
\item replacement of a covariance $C_{cd}$ by $C^D_{cd}$.
\end{enumerate}

This list is exhaustive because the exponent and the port vector field
are obtained by literal restriction, not by changing word weights. In the
first case the source-to-port synthesis reaches the boundary. In the second
case use the port version of \eqref{f16v11:resolvent-boundary}. In the third
use \eqref{f16v11:Gaussian-conditional}, and in the fourth use the first
identity in \eqref{f16v11:resolvent-boundary}. These formulas represent the
difference by a path to a boundary index, not just by a norm estimate.
Frozen values and full/conditional means at that index are bounded by
$C_{r,Q}$. Entries $H_{IO}$ have fixed range and bounded row sums.

Telescope products in the Wick formula so that a difference term contains
one of the four marked differences, with other factors chosen from one of
the two sets of data. Its connected graph contains both source anchors and
a boundary visit. The total edge length is therefore at least the distance
between the sources and at least the distance from the first source to the
boundary, up to fixed word-range shifts. Reserve one part of the common
inverse/lift exponent for each of these two distances. Sum the remaining
edge factors using their exponential row bounds. With $b$ smaller again,
this gives the stronger row estimate
\begin{equation}
 \sum_{\beta\in\mathcal A_D}e^{2b d(\alpha,\beta)}
   |(G^D_2(\mathbf q)-G_2)_{\alpha\beta}|
       \le C_{r,Q}e^{-2b(s_D(\alpha)-r_0)_+}.
 \label{f16v11:boundary-row}
\end{equation}
There is an identical column bound when the second source is interior.
Polynomial growth of the ambient graph makes all boundary-time sums finite
with constants independent of $n$; no count of the entire boundary is used
instead of its decaying weights. Products of means are handled by telescoping
as above, so this also covers nonzero and inhomogeneous base histories.

For both sources in $\mathcal A$, the row and column estimates bound the
$(i,j)$ block operator norm by $C_{r,Q}e^{-2bL}e^{-2b|i-j|}$. Sum with
weight $e^{b|i-j|}$ to obtain \eqref{f16v11:coefficient-screening-bound}.
The time-off-diagonal mask preserves the absolute row bounds and gives the
memory assertion. Comparison through the same full coefficient proves the
last assertion. No exterior density or probability has been substituted.
\end{proof}

This screening statement compares coefficients of different \emph{conditional}
laws. It does not identify their native covariances with the full native
covariance. The role of the full coefficient is a common, already defined
local approximation, not a newly assigned vacuum.

\subsection{A complete native cylinder remainder with a local spatial cost}

The following is the marked native estimate. Its proof verifies the new
boundary hypotheses before reusing the temporal analytic argument.

\begin{theorem}[Conditional leading covariance and a uniform temporal remainder]
\label{f16v11:native-cylinder}
For fixed $r,Q$ there are $K_{r,Q},\gamma_{r,Q}<\infty$ such that, if
\begin{equation}
 \gamma\ge\gamma_{r,Q},\qquad
             e^{K_{r,Q}M}/\sqrt\gamma\le1/100,
 \label{f16v11:local-entry}
\end{equation}
then every exterior history obeying \eqref{f16v11:windows} satisfies
\begin{equation}
 \left\|\Gamma^D_{\gamma,ij}(\mathbf q)
              -\gamma^{-1}(G^D_2(\mathbf q))_{ij}\right\|
 \le e^{K_{r,Q}M}\gamma^{-3/2}
                   3^{-(|i-j|-2)_+}.
 \label{f16v11:conditional-block-remainder}
\end{equation}
After enlarging $K_{r,Q}$ the same bound summed in the fixed temporal norm
is
\begin{align}
 \|\Gamma^D_\gamma(\mathbf q)-\gamma^{-1}G^D_2(\mathbf q)
                                     \|_{\mathrm{time},b}
      &\le e^{K_{r,Q}M}\gamma^{-3/2},\notag\\
 \|\mathsf M^D_\gamma(\mathbf q)-\gamma^{-1}M^D_2(\mathbf q)
                                     \|_{\mathrm{time},b}
      &\le e^{K_{r,Q}M}\gamma^{-3/2}.
 \label{f16v11:conditional-native-remainder}
\end{align}
In particular these two exact conditional matrices have temporal norm at
most $C_{r,Q}\gamma^{-1}$ in this regime. The ambient $B$ is unrestricted,
and all estimates hold for every $n\ge1$. The cost in $M$ is not removed.
\end{theorem}
\begin{proof}
We check the ingredients of V9--V10 for this cylinder.

\emph{Gaussian reference and endpoints.} The free Gaussian after port
integration has the internal precision $W_n^D$ in
\eqref{f16v11:separated-precision}, with bounded sources
\eqref{f16v11:boundary-sources}. Its diagonal interior entries are
$S_D+2N_D$, its adjacent time entries are $-N_D$, and its two endpoint
entries are the restrictions of V3's $E$. Center its \emph{whole} internal
history at $\mu^D=-(W_n^D)^{-1}\mathbf b^x$, and its moving ports at
$-L_D^{-1}b^z_t$. In unmoving port coordinates the centering adds
$\mathsf F_D(\mu^D_t-\mu^D_{t+1})$. All centers have a bounded norm per
cube by Proposition~\ref{f16v11:Gaussian-screening}.

Completing the Gaussian squares yields linear terms at the two ends of
each bond. The internal stationarity equations cancel consecutive terms;
the two endpoint equations cancel the remaining terms against the dressed
endpoint amplitudes. This is V9's telescoping with $W_n^D$ and the actual
boundary sources, not with the full mean restricted to $D$. After this
balancing, every Gaussian step is one centered operator whose internal
magnetic matrix is $S_D$ and kinetic precision is $N_D$. The integrated port
constant contains the full $\det L_D$; it is retained as part of that step.
Its normalized transfer has the Mehler split $P_D^*+Q_D^*$ with
$\|Q_D^*\|\le3-2\sqrt2<1/5$, by
$4I\le\mathsf M_D^{1/2}S_D\mathsf M_D^{1/2}\le10I$.
Here $P_D^*$ denotes the rank-one reference projection, not a port lift.

The reference endpoint is the normalized centered amplitude with precision
$\tfrac12((C_B^{-1})_{DD}+A_D^{\mathsf T}A_D)$. Its overlap with the
centered reference ground vector is at least $e^{-cM}$, since both have
fixed positive covariance bounds in dimension $15M$. The native left and
right endpoint amplitudes contain the actual crossing mixed terms at times
zero and $n$ and the original product-cube Gaussian carrier restricted to
$D$. Their norms, and any scalar factors used in balancing, have upper and
lower bounds $e^{\pm C_{r,Q}M}$: integrate on product unit-cube balls for
the lower bounds and use the pin for upper bounds. They can be normalized
for application of the product lemma, but their exact norms multiply the
conditional partition function and are not set to one. Frozen exterior
weights cancel from a normalized conditional, not from its exterior
posterior.

\emph{Complete compact comparison.} A one-, two-, or three-bond conditional
window contains $O(M)$ free real coordinates and $O(M)$ fixed exterior
inputs from the collar. By \eqref{f16v11:pin} its free integrations have a
Gaussian majorant after the centering and linear balancing:
\[
             e^{C_{r,Q}M}e^{-c|w|^2}.
\]
The use of Young's inequality to absorb the linear terms is valid because
the squared norms of their coefficient vectors are $O_{r,Q}(M)$.
Every conditional Wilson term is a member of the original finite ordered
word list, with some arguments fixed at bounded $q$ and $y$. Its fixed
real derivatives are uniformly bounded before rescaling. The integral
Taylor formula for $\varepsilon^{-2}f(\varepsilon w)$ therefore gives
polynomial derivative majorants of each fixed order, exactly as in V10.
There are no variable principal logarithms of exterior products to
differentiate: the exterior unmoving groups are the given
$\operatorname{Exp}(\varepsilon q)$ and their fixed-frame products.

Use V10's smooth proof cutoff on the \emph{free} principal charts only.
On its complement at least one centered free group coordinate has norm
$c/\varepsilon-C_{r,Q}$; the pin pays this tail by
$e^{C_{r,Q}M-c'/\varepsilon^2}$. On the cutoff domain the Haar square roots,
port Haar densities, and their fixed derivatives have polynomial bounds.
Taylor expansion through degree two, followed by complete free integration,
gives balanced step and endpoint remainders at most
$e^{C_{r,Q}M}\varepsilon^3$ in Hilbert--Schmidt or $L^2$ norm. The
zeroth arrays are the common Gaussian reference. Fixed-degree Gaussian
moments cost a polynomial in $M$ times a constant to the power $M$.
No integration over the ambient exterior occurs in this estimate.

\emph{Marked insertions.} The source observables are formed in the original conditional coordinates
before balancing. No derivatives of the history-dependent centers or
multipliers are omitted by differentiating a transformed formula.
The local port lift has a dimension-free operator
norm. The inserted-word remainder certificate of V7, on a window with fixed
exterior entries, thus gives for arbitrary spatial common-source vectors
$v$ a first remainder insertion with norm
$e^{C_{r,Q}M}\varepsilon\|v\|$. A product of two such insertions starts at
order two and costs $e^{C_{r,Q}M}\varepsilon^2\|v\|\|w\|$.
The first and product insertions have jets through degree two with
$e^{C_{r,Q}M}\varepsilon^3$ remainders. These assertions follow by
multiplying the same Gaussian majorant by the bounded-degree insertion
polynomials. All free port tails are integrated. Constants from exterior
fields are bounded by the fixed $Q$, not a supremum over unrestricted
compact boundary data. Only free Haar derivatives are used.

\emph{Exact long-product normalization.} The preceding arrays, on a small
complex polynomial-jet disc of radius $e^{-c_{r,Q}M}$, lie in V10's
analytic array ball around the common Gaussian operator and endpoint.
That proposition uses only the reference spectral ratio, the endpoint
overlap, and maximum array norms; all have just been verified with $M$
in place of $B$. Its exact scalar normalizers are kept. Its centered
two-window numerator retains the factor $3^{-(|i-j|-2)}$ for separated
windows. Cauchy estimates on that disc and the real short-jet errors
therefore give \eqref{f16v11:conditional-block-remainder}, with no factor
of the number of bonds. For overlapping windows the product insertion
occupies at most three bonds and gives the same estimate with factor one.
This is an application of the proved analytic product functional, not an
assumption of analyticity of the native measure at complex coupling.

At fixed finite $D,n$, the leading covariance coefficient is necessarily
\eqref{f16v11:leading-conditional}, by ordinary finite-dimensional
normalized expansion and uniqueness of coefficients. In particular a new
boundary normalization is not being silently substituted for the conditional
one. The entry condition absorbs the finitely many overlap, jet and
Cauchy-radius costs, all exponential in $M$. Summation is legitimate for
$b<\log3$ and gives \eqref{f16v11:conditional-native-remainder}.
The exact contact support gives the memory bound. Finally
\eqref{f16v11:conditional-coefficient-locality} and
$e^{K_{r,Q}M}\gamma^{-1/2}\le1/100$ give the stated $O(\gamma^{-1})$
conditional norms.
\end{proof}

\subsection{A native collar-screening theorem and a computable local approximation}

\begin{theorem}[The first source-local marked native remainder]
\label{f16v11:collar-native}
Under \eqref{f16v11:windows} and \eqref{f16v11:local-entry}, for any source
set $\mathcal A$ of collar depth $L$ as above,
\begin{equation}
 \boxed{
 \left\|\chi_{\mathcal A}
    [\mathsf M^D_\gamma(\mathbf q)-\gamma^{-1}M_2(\mathbf y)]
                                      \chi_{\mathcal A}\right\|_{\mathrm{time},b}
 \le C_{r,Q}\gamma^{-1}e^{-bL}
                   +e^{K_{r,Q}M}\gamma^{-3/2}.}
 \label{f16v11:collar-bound}
\end{equation}
The analogous estimate holds with $\Gamma^D_\gamma$ and $G_2$.
Thus, for any two admitted exterior fast histories $\mathbf q,\mathbf q'$,
\begin{equation}
 \left\|\chi_{\mathcal A}
    [\mathsf M^D_\gamma(\mathbf q)-\mathsf M^D_\gamma(\mathbf q')]
                                      \chi_{\mathcal A}\right\|_{\mathrm{time},b}
 \le 2C_{r,Q}\gamma^{-1}e^{-bL}
                   +2e^{K_{r,Q}M}\gamma^{-3/2}.
 \label{f16v11:two-boundary-screening}
\end{equation}
All durations and ambient volumes are allowed. No native interior cutoff,
stationarity assumption, or probability bound for the exterior is used.
\end{theorem}
\begin{proof}
Insert the conditional leading coefficient $\gamma^{-1}M^D_2(\mathbf q)$.
The error to the native conditional is
\eqref{f16v11:conditional-native-remainder}; its comparison to the compressed
full coefficient is \eqref{f16v11:coefficient-screening-bound}. Compression
by the same spatial projection at every time cannot increase the time-block
norm. This proves \eqref{f16v11:collar-bound}. Compare both boundary histories
through the same $M_2(\mathbf y)$ to get
\eqref{f16v11:two-boundary-screening}. The proof for the complete residual
Gram is identical. The two errors have different meanings: the first is
boundary-reaching coefficient error, and the second is the complete native
remainder on the declared cylinder.
\end{proof}

For an explicit local recipe one can choose $\mathbf q=0$, compute
$C^D,m^D,\mathsf P_D$ and the literal cubic word tensors on $D$, then use
\eqref{f16v11:leading-conditional}. This gives $M^D_2(0)$ without knowing
fast fields beyond the collar. The same proof gives
\begin{equation}
 \left\|\chi_{\mathcal A}
 [\mathsf M^D_\gamma(\mathbf q)-\gamma^{-1}M^D_2(0)]
                              \chi_{\mathcal A}\right\|_{\mathrm{time},b}
 \le C_{r,Q}\gamma^{-1}e^{-bL}+e^{K_{r,Q}M}\gamma^{-3/2}.
 \label{f16v11:local-computation}
\end{equation}
The zero boundary is an \emph{approximation used for this coefficient};
no zero boundary is imposed on the native law on the left. The relevant
bridge history is still retained wherever it enters the cylinder words.

For example place a fixed source region inside a cube collar of spatial
radius $L+r_0$. There is a geometric $c_0$ such that $M\le c_0(L+r_0+1)^3$.
Choose $L=L_\gamma\to\infty$ slowly enough that
\begin{equation}
       K_{r,Q}c_0(L_\gamma+r_0+1)^3\le\tfrac14\log\gamma.
 \label{f16v11:collar-scale}
\end{equation}
For sufficiently large $\gamma$, the local entry condition holds and
\eqref{f16v11:collar-bound} is at most
\begin{equation}
 C_{r,Q}\gamma^{-1}e^{-bL_\gamma}+\gamma^{-5/4}.
 \label{f16v11:screening-scale}
\end{equation}
One may take $L_\gamma$ of order $(\log\gamma)^{1/3}$ with a sufficiently
small fixed coefficient. The ambient torus may be arbitrarily larger than
this collar. This is not an assertion for a collar at unrestricted $M$
with fixed coupling, and the first term is not a fixed algebraic improvement
in $\gamma$. On a torus too small to contain a proper such collar, the
previous full-torus results remain the applicable statements.

\subsection{Conditional source means and the return needed for actual gluing}

The covariance estimate alone is not enough to glue cylinders. There is
a useful additional mean estimate with the same boundary convention.
Let $a_{1,\alpha}(\mathbf y)=\mathbb E_0r_{\alpha,1}$ be the full Gaussian
first residual coefficient of V10, and set
\[
 j^D_{\gamma,\alpha}(\mathbf q)
       =\partial_\alpha[-\log Z_{\gamma,D}(\mathbf q;\mathbf y)],
 \qquad j^{\rm ref}_\alpha=g_\alpha+\gamma^{-1/2}a_{1,\alpha}.
\]
The reference value is independent of the exterior fast history.

\begin{proposition}[Native conditional mean message on the admitted collar]
\label{f16v11:mean-message}
After enlargement of the same constants and, if necessary, decrease of $b$,
each unit canonical source $\alpha$ of depth at least $L$ obeys
\begin{equation}
 |j^D_{\gamma,\alpha}(\mathbf q)-j^{\rm ref}_\alpha|
       \le C_{r,Q}e^{-bL}+e^{K_{r,Q}M}\gamma^{-1}.
 \label{f16v11:mean-message-bound}
\end{equation}
This bound is independent of duration and ambient volume, on the same
admitted boundary histories. It is not a bound on their probability.
\end{proposition}
\begin{proof}
The exact Haar identity gives
$j^D_{\gamma,\alpha}=g^D_\alpha+\mathbb E_Dr^D_{\gamma,\alpha}$.
The one-window version of the conditional short-jet and analytic-array
argument in Theorem~\ref{f16v11:native-cylinder}, now expanded only through
order one, gives
\[
 \left|\mathbb E_Dr^D_{\gamma,\alpha}
      -\gamma^{-1/2}\mathbb E_{0,D}^{\mathbf q}r^D_{\alpha,1}\right|
                           \le e^{K_{r,Q}M}\gamma^{-1}.
\]
There is no measure-correction term at this order because the order-zero
residual is identically zero. At higher orders such terms would be required.
The short insertion norm is uniform in its time location, including either
dressed end. The one-window normalized functional has the same
length-independent analytic domain used above.

The Gaussian scalar difference is \eqref{f16v11:g-screen}. To compare the
two first residual means, use the same quadratic polynomial and inverse
boundary identities as in the proof of
Theorem~\ref{f16v11:coefficient-screening}. There is now one source anchor
and a marked boundary visit; a quadratic mean consists only of its constant
term and one self-contraction. The mean, covariance and lift differences
all carry a boundary path, so
$|\mathbb E_{0,D}^{\mathbf q}r^D_{\alpha,1}-\mathbb E_0r_{\alpha,1}|
\le C_{r,Q}e^{-bL}$. This also covers a primitive term synthesized from a
boundary port. Combining the three bounds proves
\eqref{f16v11:mean-message-bound}. No derivative of a pressure-error bound
has been taken.
\end{proof}

For clarity about the remaining task, the exact exterior law at fixed
bridge history is
\begin{equation}
 d\bar\mu_\gamma(\mathbf q)
 =\frac{1_{\mathcal P'_O}J^f_{\gamma,O}(\mathbf q)
         e^{-U_{\gamma,O}(\mathbf q;\mathbf y)}
         Z_{\gamma,D}(\mathbf q;\mathbf y)}{Z_\gamma(\mathbf y)}\,d\mathbf q.
 \label{f16v11:posterior}
\end{equation}
In this formula \emph{all} exterior coordinates are integrated; its
projection to the word collar gives the same conditional messages. For two
interior marks at nonadjacent times, normalized differentiation or the
inherited total-covariance identity gives
\begin{equation}
 \boxed{
 \partial_\alpha\partial_\beta[-\log Z_\gamma]
 =-\mathbb E_{\bar\mu_\gamma}
          [\Gamma^D_{\gamma,\alpha\beta}(\mathbf q)]
   -\operatorname{Cov}_{\bar\mu_\gamma}
          (j^D_{\gamma,\alpha},j^D_{\gamma,\beta}).}
 \label{f16v11:actual-gluing}
\end{equation}
Indeed $U_{\gamma,O}$ has no derivative in either interior source direction;
therefore $\partial_\alpha[-\log Z_\gamma]=
\mathbb E_{\bar\mu_\gamma}j^D_{\gamma,\alpha}$. Differentiating again gives
the conditional Hessian minus the displayed covariance of conditional
means. The conditional Hessian at these times is exactly
$-\Gamma^D_{\gamma,\alpha\beta}$. At finite regulator all derivatives and
integrals are legitimate. The same identity holds for the full history
weight at nonadjacent times because its external logarithmic factors are
single-time. It is one decomposition of the original return, not an
additional debit to subtract from a previously integrated Hessian.

Let $\mathsf G_Q$ be the event in \eqref{f16v11:windows} for the entire
exterior collar history. The estimates in this appendix hold on
$\mathsf G_Q$. Neither $\bar\mu_\gamma(\mathsf G_Q^c)$ nor a source-weighted
moment on that event has been bounded here. Moreover even on the good
event, a small supremum for each scalar mean message does not itself give
a small operator norm or a summable temporal row for the covariance of
\emph{all} messages. A coherent boundary direction could remain.

For a quantitative ledger set
$d_\alpha(\mathbf q)=j^D_{\gamma,\alpha}(\mathbf q)-j^{\rm ref}_\alpha$ and
$\delta_L=C_{r,Q}e^{-bL}+e^{K_{r,Q}M}/\gamma$. The exact posterior obeys
\begin{equation}
 |\operatorname{Cov}_{\bar\mu_\gamma}(j^D_{\gamma,\alpha},j^D_{\gamma,\beta})|
 \le
 \sqrt{(\delta_L^2+\mathbb E[1_{\mathsf G_Q^c}|d_\alpha|^2])
       (\delta_L^2+\mathbb E[1_{\mathsf G_Q^c}|d_\beta|^2])}.
 \label{f16v11:bad-message-ledger}
\end{equation}
This is Cauchy--Schwarz after subtracting the deterministic reference and
splitting the second moments. It exposes, but does not estimate, the
exceptional source moments. The first term in
\eqref{f16v11:actual-gluing} has its own exceptional conditional-covariance
integral. Existing expected bad-label densities cannot be substituted for
these quantities, especially on an arbitrarily long cylinder.

The next noncircular target is consequently the \emph{actual exterior
message law} in \eqref{f16v11:posterior}: either a source-weighted collar
escape estimate plus a coherent boundary-message contraction, or a
multiscale decomposition that controls their combined connected return.
The compact conditional interior is no longer an unspecified input: its
marked remainder and leading coefficient now have the explicit local-size
cost in \eqref{f16v11:collar-bound}. Reproving a global temporal product
bound, or replacing the posterior by a bounded deterministic boundary,
would miss the remaining term. No physical Gauss restriction, bridge-window
coverage, infrared contact coercivity or same-top physical transfer
comparison is obtained from this conditioning argument.

\clearpage
\subsection{Diagnostics, preservation, and scope}

The accompanying diagnostic is
\begin{center}\small
\path{verify_ym_f16_v11_conditional_collar.py}.
\end{center}
It checks noncommuting Gaussian conditioning and covariance differences,
principal versus marginal inverse distinctions, restricted port-lift
cancellation with nonzero exterior data, endpoint weights and centering,
actual parity-cube crossing carriers, leading Gaussian polynomial
covariances, finite-range boundary-reaching inverse paths, and the exact
posterior covariance decomposition. All 463 assertions passed; their family counts and the script hash are in
\path{ym_v11_verification_report.json}. The unchanged V4--V10 scripts were
also rerun and passed 9,467, 89,339, 906, 311, 609, 990 and 608 assertions,
respectively. These finite checks do not certify analytic derivative
suprema, complete-chart Taylor remainders, or infinite-dimensional operator
estimates with a proof assistant. The theorems above rely on their analytic
proofs and on the explicitly identified inherited inputs.

The supplied V10 source is preserved byte for byte except for the cumulative
title version and insertion of this exact appendix before its final
\verb|\end{document}|. Removing this append and restoring the title recovers
V10 exactly. The package records original-source hashes, predecessor
recovery, compilation and render checks. Only diagnostics actually available
in the supplied packages are rerun and reported; no inaccessible historical
checker or NS Lean development is described as verified. The next scheduled
five-version companion checkpoint remains V15.

\begin{firewall}
V11 proves complete-native conditional cylinder estimates on a declared
bounded exterior fast-history window. The exponential comparison cost is
in the cylinder size $M$, not the ambient volume $B$, and no duration cost
is introduced. Boundary-marked coefficient screening yields a native
conditional memory approximation with error
$C_{r,Q}\gamma^{-1}e^{-bL}+e^{K_{r,Q}M}\gamma^{-3/2}$, and a corresponding
conditional mean-message estimate. All crossing interactions, free compact
tails and dressed endpoints are retained. No typicality of the admitted
collar, source-weighted control of its complement, coherent contraction of
the actual exterior-message covariance, unrestricted full-native spatial
gluing, physical transfer gap or interacting continuum Yang--Mills mass gap
is established by this appendix.
\end{firewall}
% END F16 V11 APPEND

% BEGIN F16 V12 APPEND -- 2026-09-17
% Insert immediately before the final \end{document} in cumulative V11.
% The predecessor body is unchanged; only its title version becomes V12.

\clearpage
\section{V12: posterior return without an all-good collar}
\label{f16v12:context}

V11 controls an interior cylinder on a bounded exterior-history window.
Requiring that window at every time is not a useful route to an estimate
uniform in duration: its complement can have large probability even when
local fluctuations are well controlled. We instead estimate the \emph{actual
posterior mean vector}, integrating all exterior fields. The two inputs are
V6's full-native relative entropy and the finite-range Markov specification
shared by the native law and its Gaussian reference. Conditional Gaussian
concentration is applied to a whole source vector, before summing over
cylinders. Applying it separately to every source would lose the duration.

The output has a deliberately different norm from V11's conditional theorem.
The posterior return, including its exceptional-field part, has a small
\emph{trace per space--time cell}; no uniform small operator norm is inferred.
We also match the conditional and full Haar lifts. This pays the exceptional
conditional-covariance integrals in the same density sense, without a
bounded-collar event. The final gluing identity keeps both terms under their
one actual exterior posterior. Nothing here promotes the conditional
$\gamma^{-1}$ expansion to an unrestricted-volume expansion.

\subsection{Dependencies and the upstream choice}

The supplied V11 appendix was read in full. The available f14/f15 archives
and Grand Tensor Monograph V076 were inventoried and relevant passages
examined; their entire proof collections have not been independently
recertified. F14 V67 already estimates entropy response for a different
restricted high conditional; f15 V1 already turns exponential-source control
into mean estimates, and f15 V13--V14 and V68 retain actual posterior weights
in other conditional comparisons. The monograph's Markov-blanket and
conditional-mean identities are inherited. None is presented as a new
abstract principle.

The model-specific input reused here is V6's entropy estimate for the
\emph{complete augmented fast-history law}, not entropy of a restricted core.
We also reuse V6's Gaussian marginal-entropy localization, V7's Gaussian
vector-energy inequality and exact Haar sources, and V11's literal spatial
conditioning and restricted port inverse. The new specialization combines
these inputs on entire-time spatial cylinders, with no exterior-field
cutoff, and keeps the source-vector and overlap costs explicit.
The import memorandum's separation of power gain from locality cost remains
binding: this is a change in how a remainder is measured, not another formal
order or an imported Navier--Stokes theorem.

For attribution, the Gaussian concentration used below follows from the
classical logarithmic Sobolev inequality of L.~Gross,
\href{https://doi.org/10.2307/2373688}{\emph{Logarithmic Sobolev Inequalities}},
Amer. J. Math. 97 (1975), 1061--1083. V6's marginal-entropy argument is an
instance of the entropy/Fisher-information method discussed by E.~A.~Carlen
and D.~Cordero--Erausquin,
\href{https://arxiv.org/abs/0710.0870}{\emph{Subadditivity of the entropy and
its relation to Brascamp--Lieb type inequalities}}. The precise Gaussian
consequences needed here are derived below or explicitly reused. No native
logarithmic Sobolev inequality is assumed.

\subsection{One full law, whole source vectors, and an exact local entropy cost}

Fix the same bounded bridge history as in V6--V11 and write
\[
 \mathcal N=B(n+1),\qquad
 \epsilon_\gamma=\gamma^{-1/6}(1+\log\gamma)^{2/3},\qquad
 \max_{i,e}|y_{i,e}|\le r<\infty.
\]
Here $B=m^3$, $m\ge4$, and $n\ge1$. On the same complete unmoving Euclidean
coordinate space $\zeta=(x,a)$ let
\begin{equation}
 P=\mu_{\gamma,\mathbf y},\qquad Q=\mu_{0,\mathbf y}=N(m,H^{-1}),
 \qquad h_-I\le H\le h_+I,\qquad
 D(P\Vert Q)\le C_r\mathcal N\epsilon_\gamma.
 \label{f16v12:inherited}
\end{equation}
The native density is extended by zero outside its complete principal charts,
as in V6. The threshold for \eqref{f16v12:inherited} depends on $r$, not on
$B,n$. No hypothesis $e^{KB}/\sqrt\gamma\ll1$ is used in this appendix.
The symbol $Q$ denotes the Gaussian probability, not a bound on exterior
coordinates. All exterior coordinates are integrated under $P$.

For a proper nonempty spatial cylinder $D$, let $I_D$ be all its free
space--time coordinate cells as in V11. Choose a fixed-width exterior
collar $C_D$ containing every exterior cell appearing in an action term
that meets $I_D$, and every nonzero block of $H_{I_D O_D}$. Write
$T_D=I_D\cup C_D$ and $O_D=I_D^c$. These are sets of \emph{cells}, with all
real components and all available times included. The width can be fixed
from the original finite word list, independently of $D,B,n,\gamma$.

Let $\mathcal A_D$ be a chosen bank of distinct canonical common-bridge
marks whose ordinary score derivative carriers are inside $I_D$. Define
\begin{equation}
 F_D(\zeta_{I_D})=(J_{\gamma,\alpha})_{\alpha\in\mathcal A_D},\qquad
 J_{\gamma,\alpha}=\partial_\alpha U_\gamma,\qquad
 p_D=|\mathcal A_D|.
 \label{f16v12:source-vector}
\end{equation}
These are the \emph{ordinary} source scores, not the compensated residuals.
Their Gaussian comparison below is the analytic exponential-word extension
of this same $F_D$, not a substitution of $J_0$ for $J_\gamma$.
Choose $L\ge1$ such that the fast derivative support of this bank has
spatial distance at least $L$, after a fixed word-range allowance, from the
interior cells coupled by $H$ to $C_D$. Source anchors of depth $L+2r_0$
with V11's sufficiently large $r_0$ meet this condition.

The exact native and reference messages, and a deterministic centering, are
\begin{equation}
 j_D(q)=\mathbb E_P[F_D\mid\zeta_{O_D}=q],\qquad
 f_D(q)=\mathbb E_Q[F_D\mid\zeta_{O_D}=q],\qquad
 a_D=\mathbb E_Q F_D.
 \label{f16v12:messages}
\end{equation}
In particular $j_D$ is V11's vector of derivatives of $-\log Z_{\gamma,D}$.
No interior source differentiates $U_{\gamma,O}$. The vectors $a_D$ merely
center estimates; their values are not chosen as a new vacuum or normalizer.

\begin{proposition}[Whole-bank Lipschitz bound and the correct conditional entropy]
\label{f16v12:local-input}
There is a fixed geometric $L_0<\infty$ such that
\begin{equation}
             \sup_\zeta\|D_\zeta F_D\|_{\rm op}\le L_0
 \label{f16v12:source-Lip}
\end{equation}
for every such bank, every $B,n,D$ and every $\gamma\ge1$. Both messages in
\eqref{f16v12:messages} depend on the exterior only through $q_{C_D}$.
Put
\begin{equation}
 d_T=D(P_{T_D}\Vert Q_{T_D}),\quad
 d_C=D(P_{C_D}\Vert Q_{C_D}),\quad k_D=d_T-d_C.
 \label{f16v12:local-entropies}
\end{equation}
Then
\begin{equation}
 k_D=\mathbb E_{P_{C_D}}
 D(P_{I_D\mid C_D}\Vert Q_{I_D\mid C_D})\ge0.
 \label{f16v12:chain}
\end{equation}
These are finite quantities at every regulator where
\eqref{f16v12:inherited} holds. The conditional divergence in this identity
is averaged against the \emph{native} collar marginal.
\end{proposition}
\begin{proof}
For a primitive action term $\gamma s(\prod_j\operatorname{Exp}
(T_jw/\sqrt\gamma))$, every second derivative in its real inputs is
bounded by a fixed constant on all real inputs. Duhamel differentiation
inserts fixed coefficient matrices into products of unitary matrices; the
two factors $\gamma^{-1/2}$ cancel the prefactor $\gamma$. Word lengths and
linear coefficients are fixed. Taking one bridge derivative gives a row of
the score Jacobian supported on a fixed number of fast coordinates. Each
fast coordinate appears in a fixed number of such rows, including the
four-bridge normalization of a canonical common mark. The block row and
column Schur bounds prove \eqref{f16v12:source-Lip}. Taking a subbank cannot
increase the operator norm. Endpoint Gaussian terms are independent of the
bridge marks; the fast Haar factors have no bridge dependence. This is the
whole-vector version of V6's local score certificate. It is not the sum of
$p_D$ scalar Lipschitz constants.

Native product-coordinate measure and finite-range action imply that the
conditional density on $I_D$ given $O_D$ depends only on $C_D$: exterior-only
terms cancel on normalization. The complete principal-chart indicator and
all Haar powers factor by cells, so they preserve this statement. The
Gaussian has the same property because $H_{I_D,O_D\setminus C_D}=0$.
Conditioning first on $C_D$ therefore gives the same interior kernels for
both laws. The log-density chain rule on $T_D=I_D\sqcup C_D$ gives
\eqref{f16v12:chain}; data processing from \eqref{f16v12:inherited} makes
$d_T,d_C$ finite. This argument does not invoke a sum of native conditional
entropies bounded by a full entropy without first identifying these patches.
\end{proof}

\subsection{Gaussian mean comparison without an output-dimension loss}

We record a standard Gaussian consequence in the needed vector form.

\begin{proposition}[A vector mean bound from one conditional entropy]
\label{f16v12:vector-mean}
Let $Q'$ be a finite-dimensional Gaussian of covariance at most $cI$, and
let $F$ take values in $\mathbb R^p$ with $\|DF\|\le L_F$. For any
$P'\ll Q'$ of finite relative entropy,
\begin{equation}
       \|\mathbb E_{P'}F-\mathbb E_{Q'}F\|^2
                     \le2cL_F^2D(P'\Vert Q').
 \label{f16v12:mean-transport}
\end{equation}
The dimension $p$ does not multiply the right side. For the messages above,
\begin{equation}
             \mathbb E_{P_{C_D}}\|j_D-f_D\|^2
                           \le2h_-^{-1}L_0^2 k_D.
 \label{f16v12:conditional-mean-error}
\end{equation}
\end{proposition}
\begin{proof}
For a unit vector $v$, the scalar $v^{\mathsf T}F$ is $L_F$-Lipschitz.
The Gaussian logarithmic Sobolev inequality in the convention used in V7
and its elementary Herbst integration give
\[
 \log\mathbb E_{Q'}\exp\{t(v^{\mathsf T}F-\mathbb E_{Q'}v^{\mathsf T}F)\}
                       \le ct^2L_F^2/2,\qquad t\in\mathbb R.
\]
For clarity, if $\psi(t)$ is this logarithmic moment generating function,
log-Sobolev gives $t\psi'(t)-\psi(t)\le ct^2L_F^2/2$; integrate
$(\psi(t)/t)'$ from zero, and use both signs. Gaussian integrability of
Lipschitz functions justifies this calculation by truncation.
The entropy variational inequality bounds a signed mean difference by
$D(P'\Vert Q')/t+ctL_F^2/2$. Optimize $t>0$. Choose $v$ in the direction
of the vector mean difference, rather than sum the scalar inequalities
for a basis. This proves \eqref{f16v12:mean-transport}.

The Gaussian interior conditional has covariance $(H_{I_DI_D})^{-1}
\preceq h_-^{-1}I$, independent of $q$. Its mean may be arbitrarily large;
the preceding inequality is translation independent. Apply it at each
native exterior value $q$, to the same analytic score bank $F_D$, and then
integrate the conditional divergence using \eqref{f16v12:chain}. The native
conditional is supported on its full compact chart; absolute continuity
with respect to the everywhere-positive Gaussian is sufficient. No exterior
amplitude bound is required.
\end{proof}

\begin{proposition}[The reference message is exponentially insensitive to its collar]
\label{f16v12:reference-message}
There are $C,\nu>0$, depending only on the fixed geometry and Gaussian
spectral bounds, such that
\begin{equation}
 \operatorname{Lip}(f_D)\le C e^{-\nu L},\qquad
 \mathbb E_{Q_{C_D}}\|f_D-a_D\|^2\le C p_D e^{-2\nu L}.
 \label{f16v12:Gaussian-message}
\end{equation}
Moreover under the \emph{native} collar marginal,
\begin{equation}
 \mathbb E_{P_{C_D}}\|f_D-a_D\|^2
                          \le C e^{-2\nu L}(p_D+d_C).
 \label{f16v12:native-reference-message}
\end{equation}
All times and arbitrary exterior amplitudes are included. The constants
contain neither the area of the boundary nor the duration.
\end{proposition}
\begin{proof}
Let $I=I_D$, $C_\partial=C_D$, and let $S_D^{\rm src}\subset I$ be the fast
derivative support of $F_D$. Gaussian conditioning realizes the free
coordinates as
\[
       m_I-H_{II}^{-1}H_{IC_\partial}(q-m_{C_\partial})+g,
             \qquad g\sim N(0,H_{II}^{-1}).
\]
Its covariance is independent of $q$. Differentiate this expectation, using
the global derivative bound of $F_D$. Only rows of $H_{II}^{-1}H_{IC_\partial}$
in $S_D^{\rm src}$ can contribute.

The nonzero rows of $H_{IC_\partial}$ lie in a fixed-width interior boundary
set $T$. A Neumann expansion of $H_{II}^{-1}$ has norm ratio
$\theta=1-h_-/h_+<1$ and finite spatial range per power. Before a number of
powers proportional to $L$, no term connects $S_D^{\rm src}$ to $T$.
Summing the remaining tail proves
\[
 \|P_{S_D^{\rm src}}H_{II}^{-1}H_{IC_\partial}\|
                                      \le C e^{-\nu L}.
\]
Here $\|H_{IC_\partial}\|\le h_+$; a set-to-set operator-norm estimate, not
a sum over all boundary cells, is being used. Powers can travel in time,
but cannot cross a spatial distance $L$ in fewer spatially ranged steps.
This proves the first bound in \eqref{f16v12:Gaussian-message} for every
$q$. Differentiation is justified by the Lipschitz majorant under the
Gaussian, without a bounded-mean hypothesis.

The full Gaussian marginal $Q_{C_D}$ has covariance at most $h_-^{-1}I$.
Gaussian Poincare applied componentwise, or its vector trace form, gives
\[
 \mathbb E_{Q_{C_D}}\|f_D-\mathbb E_{Q_{C_D}}f_D\|^2
       \le h_-^{-1}\mathbb E\|Df_D\|_{\rm HS}^2
       \le h_-^{-1}p_D\operatorname{Lip}(f_D)^2.
\]
The tower property identifies the mean as $a_D$.

It remains to change the \emph{collar marginal}, rather than quietly
replace it by the reference marginal. Set $R(q)=f_D(q)-a_D$ and
$\lambda_L=Ce^{-\nu L}$, chosen as an upper bound on its Lipschitz norm.
V7's Gaussian vector-energy inequality gives, for
$t_L=(4h_-^{-1}\lambda_L^2)^{-1}$,
\[
  \log\mathbb E_{Q_{C_D}}e^{t_L\|R\|^2}
                           \le2t_L\mathbb E_{Q_{C_D}}\|R\|^2.
\]
The entropy variational inequality therefore yields
\[
 \mathbb E_{P_{C_D}}\|R\|^2
 \le4h_-^{-1}\lambda_L^2d_C+2\mathbb E_{Q_{C_D}}\|R\|^2.
\]
This proves \eqref{f16v12:native-reference-message}. A constant message is
handled directly. The exponentially large allowed $t_L$ is legitimate:
it is paired with an exponentially \emph{small Lipschitz constant}, not
with an asserted Gaussian bound for the native measure.
\end{proof}

\subsection{An actual-posterior theorem, with no all-good event}

\begin{theorem}[Unrestricted posterior mean screening on a cylinder family]
\label{f16v12:posterior-screening}
For each cylinder and source bank just specified,
\begin{equation}
 \mathbb E_P\|j_D-a_D\|^2
       \le C\{k_D+e^{-2\nu L}(p_D+d_C)\}.
 \label{f16v12:single-message}
\end{equation}
The expectation uses the actual exterior posterior of the full native law.
Let $(D_s,\mathcal A_s)$ be any finite family whose enlarged cell sets
$T_s=I_{D_s}\cup C_{D_s}$ have overlap at most $q_0$, and whose derivative
supports all have the above depth $L$. Source occurrences are counted with
multiplicity. Put $p_{\rm tot}=\sum_s|\mathcal A_s|$ and stack
$\mathbf j=(j_{D_s})_s$, $\mathbf a=(a_{D_s})_s$, viewing every message as a
function on the \emph{same full probability space}. Then
\begin{equation}
 \boxed{
  \mathbb E_P\|\mathbf j-\mathbf a\|^2
   \le C_r q_0\mathcal N\epsilon_\gamma
                          +C p_{\rm tot}e^{-2\nu L}.}
 \label{f16v12:family-message}
\end{equation}
Consequently its actual posterior-return Gram obeys
\begin{equation}
 \boxed{
  \operatorname{Tr}\operatorname{Cov}_P(\mathbf j)
   \le C_r q_0\mathcal N\epsilon_\gamma
                          +C p_{\rm tot}e^{-2\nu L}.}
 \label{f16v12:return-trace}
\end{equation}
There is no fixed bound on $|q_c|$, no restriction on $|D_s|$ relative to
$\gamma$, and no duration bound. The result is not a supremum over exterior
histories, nor a covariance-operator bound independent of $\mathcal N$.
\end{theorem}
\begin{proof}
Split $j_D-a_D=(j_D-f_D)+(f_D-a_D)$ and use
$\|u+v\|^2\le2\|u\|^2+2\|v\|^2$, the conditional mean comparison and
\eqref{f16v12:native-reference-message}. This proves
\eqref{f16v12:single-message}. The precise first cost is $k_D=d_T-d_C$,
not $d_T$ counted a second time as a boundary divergence.

For the family, V6's Gaussian marginal-entropy localization gives
\[
 \sum_sD(P_{T_s}\Vert Q_{T_s})
      \le q_0\kappa D(P\Vert Q),\qquad \kappa=h_+/h_-.
\]
Its constant depends on overlap and the full Gaussian condition number,
not on marginal dimension. The sets $T_s$ can therefore include every time
of a long cylinder. Both $k_{D_s}$ and $d_{C_s}$ are at most $d_{T_s}$.
Sum \eqref{f16v12:single-message} and insert
\eqref{f16v12:inherited}. This proves \eqref{f16v12:family-message}.
Centering by the actual mean minimizes squared error, so
$\operatorname{Tr}\operatorname{Cov}_P(\mathbf j)
\le\mathbb E_P\|\mathbf j-\mathbf a\|^2$. Cross-covariances among
messages are retained in this Gram; independence is not assumed.
\end{proof}

The cancellation of the duration cost is substantive. A scalar application
of \eqref{f16v12:mean-transport} at every time would repeatedly charge the
same cylinder entropy, potentially multiplying it by $n+1$. Instead a
\emph{single} vector comparison pays each conditional entropy once, and the
bounded-overlap marginal theorem then distributes the already known full
entropy. This is not approximate tensorization of native conditional
entropies, and it is not a native functional inequality.

\begin{proposition}[The old exceptional mean term is included]
\label{f16v12:exceptional-means}
Use V11's deterministic reference
$j^{\rm ref}_\alpha=g_\alpha+\gamma^{-1/2}\mathbb E_0r_{\alpha,1}$ on every
chosen source bank. For any exterior events $E_s$, with no probability
assumption,
\begin{equation}
 \sum_s\mathbb E_P\!
       \left[1_{E_s}\|j_{D_s}-j^{\rm ref}_{\mathcal A_s}\|^2\right]
 \le C_rq_0\mathcal N\epsilon_\gamma
             +C_rp_{\rm tot}(e^{-2\nu L}+\gamma^{-1}).
 \label{f16v12:exceptional-mean}
\end{equation}
This applies in particular to failure of V11's entire-history bounded-collar
event. It bounds source-weighted escape in density, not escape probability.
\end{proposition}
\begin{proof}
For each ordinary score, V6's global Taylor certificate and bounded local
Gaussian moments give
$|\mathbb E_Q J_{\gamma,\alpha}-\mathbb E_QJ_{0,\alpha}|
\le C_r\gamma^{-1/2}$. V7's Gaussian Haar cancellation and Gaussian
integration by parts give $\mathbb E_QJ_{0,\alpha}=g_\alpha$.
The first residual coefficient has uniformly bounded Gaussian mean by its
quadratic local word representation and the summable port-lift row.
Thus $|a_\alpha-j^{\rm ref}_\alpha|\le C_r\gamma^{-1/2}$, uniformly
in the source location and duration. Sum its squared difference over the
bank. Dropping $1_{E_s}\le1$ in the squared-error upper bound and using
\eqref{f16v12:family-message} proves the claim. This argument never bounds
the probability of an all-time event by a union bound.
\end{proof}

\subsection{Geometry with bounded overlap at growing collar depth}

The overlap need not grow with the collar radius. Here is one explicit
construction when the torus is large enough to contain local cylinders.
Let $\rho$ be a fixed integer larger than all spatial word and derivative
ranges, and put $\ell=\lceil L+4\rho\rceil$. Suppose $m\ge16\ell$.
Partition each coordinate cycle into $k=\lfloor m/\ell\rfloor$ consecutive
intervals, distributing their lengths as equally as possible. Every length
is between $\ell$ and $2\ell$. Their products are disjoint core boxes.
For each core box enlarge by $L+2\rho$ in each spatial coordinate to obtain
$D_s$, and by a further $\rho$ to obtain a collar patch containing $T_s$.
Assign each canonical spatial common mark to its core box, at every time.

Every source derivative carrier then has the required depth, and
$p_{\rm tot}=9\mathcal N$. An enlarged coordinate interval has length at
most $2\ell+2L+6\rho\le4\ell$. Its starting core positions are separated
by at least $\ell$. Thus a point can belong to at most eight enlarged
intervals in one coordinate, a safe count also on the cycle. Hence
\begin{equation}
                         q_0\le8^3=512.
 \label{f16v12:cover-overlap}
\end{equation}
There is no extra temporal overlap factor: each spatial patch already
contains all available times. These patches overlap; their messages are
not conditionally independent by virtue of this cover. The next gluing
theorem imposes separation when independence is used.

With this cover, \eqref{f16v12:return-trace} yields
\begin{equation}
 \frac{\operatorname{Tr}\operatorname{Cov}_P(\mathbf j)}{\mathcal N}
               \le C_r(\epsilon_\gamma+e^{-2\nu L}).
 \label{f16v12:cover-rate}
\end{equation}
Taking $L\ge(2\nu)^{-1}\log(\epsilon_\gamma^{-1})$ makes the bound
$C_r\epsilon_\gamma$. This radius is $O(\log\gamma)$ and is \emph{not}
subject to V11's $e^{KM}/\sqrt\gamma$ entry condition, because no conditional
kernel expansion was used here. It must still fit geometrically in the
ambient torus. This does not extend V11's coefficient approximation to
cylinders of that size. On smaller tori the general family theorem remains
valid for whatever proper cylinders and depths actually fit; no nonexistent
collar is assumed.

\subsection{The conditional residuals also have an unrestricted posterior budget}

The exceptional conditional-covariance term in V11 can be paid in density
without applying its bounded-collar expansion. The point is to compare
\emph{compatible Haar fields}, rather than compare two differently centered
scores as unrelated observables.

For a common source $\alpha$ assigned to $D_s$, let $\eta_\alpha$ be its
full port coefficient column from V7 and let $\eta^D_\alpha$ be V11's
restricted coefficient, extended by zero to exterior ports. Write
$\Delta_\alpha=\eta_\alpha-\eta^D_\alpha$. At either adjacent bond put
$L=L_h$ only within the following block identities; here it denotes the
port matrix, not the collar depth. Since $E_O^{\mathsf T}v=0$ for an
interior source,
\begin{equation}
 \Delta_{\alpha,I}=-L_D^{-1}L_{DO}\eta_{\alpha,O},\qquad
 \Delta_{\alpha,O}=\eta_{\alpha,O},\qquad
                    (L\Delta_\alpha)_I=0.
 \label{f16v12:harmonic-lift}
\end{equation}
Let $e^{\rm pt}=S_\gamma-S_0$ be V7's primitive port-error vector, with
$S_\gamma=(X^\gamma_{p,k}U_\gamma)_{p,k}$ and $S_0=\partial_aU_0$.
Every evaluation uses the full configuration, with exterior values inserted
in a conditional expression when needed.

\begin{proposition}[Exact lift matching and an exponentially small mismatch]
\label{f16v12:lift-matching}
For every exterior chart value, not only bounded values,
\begin{equation}
           r^D_{\gamma,\alpha}
                =r^{[\infty]}_{\gamma,\alpha}
                                  -\Delta_\alpha^{\mathsf T}e^{\rm pt}.
 \label{f16v12:exact-residual-match}
\end{equation}
For a family with patch overlap at most $q_0$ and source depth at least $L$,
let $\Delta$ have all these columns, counting repeated marks as separate
columns. There are fixed $C,\nu'>0$ such that
\begin{equation}
               \|\Delta\|_{\rm op}\le C\sqrt{q_0}e^{-\nu' L}.
 \label{f16v12:Delta-op}
\end{equation}
The constants are uniform in domain shapes, duration and ambient volume.
After decreasing the earlier $\nu$, all exponential estimates may use a
single fixed exponent.
\end{proposition}
\begin{proof}
Restrict the equation $L\eta_\alpha=-E^{\mathsf T}d_t v_\alpha$ to free
port rows and compare it with the Dirichlet equation for $\eta^D_\alpha$.
This proves \eqref{f16v12:harmonic-lift}, including its sign. The full and
conditional Gaussian cancellation identities give, as an identity on the
whole fast space with the exterior inserted,
\[
                 g^D_\alpha-g_\alpha
                          =-\Delta_\alpha^{\mathsf T}S_0.
\]
The ordinary source is the same in both actions for these interior marks.
A free port differentiates no exterior-only term. Subtract the two exact
compact lifted scores and then their Gaussian constants to obtain
\eqref{f16v12:exact-residual-match}. There is no nonlinear comparison of
probability laws in this identity. In particular no flat derivative is used
in place of a compact Haar derivative; the only $S_0$ is the fixed Gaussian
coefficient.

We give the row/column estimate needed for a whole family. The full port
inverse and every principal inverse have uniformly exponentially decaying
entries in ambient spatial distance, with summable rows, as in V11. For an
exterior row $p$, $\Delta_{p\alpha}=\eta_{p\alpha}$; its path from the
source reaches $D_s^c$. For an interior row, the first identity in
\eqref{f16v12:harmonic-lift} expands it as an inverse path to a boundary
edge followed by a full inverse path to the source. Such a path has length
at least the source-to-row distance and at least the source depth, up to a
fixed stencil allowance. Splitting the available exponential weight between
these distances and summing the remaining convolution gives the uniform
entry bound
\[
  |\Delta_{p\alpha}|
       \le C e^{-\nu' L}e^{-\nu' d_{
m sp}(p,\alpha)}
              1_{\{p\text{ on a bond adjacent to }i(\alpha)\}}.
\]
Colour and port multiplicities are fixed. A source occurs at most $q_0$
times because its core lies in its patch. The absolute column sums are
therefore at most $Ce^{-\nu'L}$ and the absolute row sums at most
$Cq_0e^{-\nu'L}$. Schur's bound gives \eqref{f16v12:Delta-op}. This
argument does not multiply by the number of cylinders or by the duration.
\end{proof}

\begin{theorem}[Exceptional conditional Grams under their actual posteriors]
\label{f16v12:conditional-Gram-budget}
For every such cylinder family and every collection of exterior events
$E_s$,
\begin{align}
 \sum_s\mathbb E_P\|\mathbf r^D_{\gamma,\mathcal A_s}
                  -\mathbf r^{[\infty]}_{\gamma,\mathcal A_s}\|^2
        &\le C_rq_0\mathcal N\epsilon_\gamma e^{-2\nu L},\notag\\
 \sum_s\mathbb E_{\bar P_s}
        [1_{E_s}\operatorname{Tr}\Gamma^{D_s}_\gamma]
        &\le C_rq_0\mathcal N\epsilon_\gamma.
 \label{f16v12:conditional-Gram-estimates}
\end{align}
Here $\bar P_s$ is the actual exterior posterior for $D_s$, and each
conditional Gram is restricted to $\mathcal A_s$. Its complement outside
V11's bounded collar is included. No event probability is being estimated.
\end{theorem}
\begin{proof}
The primitive port vector has global Jacobian operator norm bounded by a
geometric constant and Gaussian squared energy at most
$C_r\mathcal N/\gamma$, by V7's primitive inserted-word certificate. Apply
V7's Gaussian vector-energy inequality and the entropy variational inequality
to this vector itself, using \eqref{f16v12:inherited}. At a fixed positive
parameter they give
\[
                   \mathbb E_P\|e^{\rm pt}\|^2
                                   \le C_r\mathcal N\epsilon_\gamma.
\]
This is the same proved energy-transfer argument, not a native functional
inequality. Combine it with \eqref{f16v12:exact-residual-match} and
\eqref{f16v12:Delta-op} to obtain the first bound. V7 also gives
$\mathbb E_P\|\mathbf r^{[\infty]}_\gamma\|^2\le
 C_r\mathcal N\epsilon_\gamma$. Restricting and repeating each source
at most $q_0$ times, followed by $\|u-v\|^2\le2\|u\|^2+2\|v\|^2$,
bounds the sum of uncentered conditional-residual energies by
$C_rq_0\mathcal N\epsilon_\gamma$. Conditional centering reduces that
energy. The tower property and $1_{E_s}\le1$ prove the second bound.
All conditional expectations and exterior integrals refer to the full native
law; no deterministic bounded boundary has replaced them.
\end{proof}

This estimate and \eqref{f16v12:exceptional-mean} pay both types of exceptional
term in V11's ledger, but only as sums of squared source energies or traces.
Their small parameter is $\epsilon_\gamma$, much larger than
$\gamma^{-1}$. They cannot be used to discard those terms in a leading
$\gamma^{-1}G_2$ coefficient theorem.

\subsection{Exact spatial gluing and its unconditioned trace-norm error}

For an exact gluing statement, now assume the interiors $D_s$ are separated
so that no native primitive term meets two of them. In particular a gap wider
than the fixed word range suffices. Let $I=\bigcup_s I_{D_s}$ and let $O=I^c$.
Every chosen source carrier and free Haar lift lies in its own cylinder.
Conditional on $O$, the interiors are independent, and the conditional
partition function factors into the literal $Z_{\gamma,D_s}$ times
exterior-only factors. The common exterior posterior is
\begin{equation}
 d\bar P(q)=Z_\gamma^{-1}
       1_{\mathcal P'_O}J^f_{\gamma,O}(q)e^{-U_{\gamma,O}(q;\mathbf y)}
               \prod_s Z_{\gamma,D_s}(q;\mathbf y)\,dq.
 \label{f16v12:joint-posterior}
\end{equation}
Each message $j_{D_s}$ is its earlier function of the word collar. Its
marginal integrals under this joint posterior agree with the corresponding
marginal integrals under the full native law.

Let $\chi$ select the union of the source banks, without repeats. Denote the
full nonadjacent-time memory by $\mathsf M_\gamma$ as in V7, and place each
conditional memory on its own source block. Put
\[
 \mathsf R_{\rm post}=\operatorname{Cov}_{\bar P}(\mathbf j),\qquad
 \mathsf L_D=\bigoplus_s\mathbb E_{\bar P}\mathsf M^{D_s}_\gamma(q).
\]
These are \emph{actual full-exterior averages}; they are not averages over
V11's good event, nor averages of its asymptotic coefficients.

\begin{theorem}[A full-native spatial gluing error in normalized trace norm]
\label{f16v12:gluing}
For the separated family,
\begin{equation}
       \chi\mathsf M_\gamma\chi-\mathsf L_D
                          =-\operatorname{Off}_{>1}\mathsf R_{\rm post}
 \label{f16v12:gluing-identity}
\end{equation}
is exact. If its collar patches have overlap at most $q_0$ and source depth
at least $L$, then
\begin{equation}
 \boxed{
  \frac{\|\chi\mathsf M_\gamma\chi-\mathsf L_D\|_{S_1}}{\mathcal N}
       \le C_rq_0\epsilon_\gamma
                     +C\frac{p_{\rm tot}}{\mathcal N}e^{-2\nu L}.}
 \label{f16v12:gluing-bound}
\end{equation}
Every time lag is included simultaneously. The trace norm is the sum of
singular values, not an entrywise absolute row norm. There is no claim that
the right side bounds the unnormalized operator norm.
\end{theorem}
\begin{proof}
At a fixed regulator differentiate the full normalized integral on the
separated interior source directions. Exterior-only action terms have no
such derivatives. The first derivative is the posterior mean of the
conditional first derivatives. Differentiating it once more gives
\[
 \nabla^2[-\log Z_\gamma]
    =\mathbb E_{\bar P}\bigoplus_s\nabla^2[-\log Z_{\gamma,D_s}]
                         -\operatorname{Cov}_{\bar P}(\mathbf j)
\]
on these sources. This is V11's exact total-covariance differentiation,
with the actual joint posterior \eqref{f16v12:joint-posterior}. Compact
integrability justifies it even at exterior values outside every bounded
rescaled window. Restrict to nonadjacent time blocks and use V11's exact
conditional contact range to get \eqref{f16v12:gluing-identity}. The full
history's external logarithmic factors are single-time functions and change
none of these blocks; they have not been dropped from the original weight.

The posterior Gram is positive. V7's Fourier block-extraction argument gives
$\|\operatorname{Off}_{>1}R\|_{S_1}\le4\|R\|_{S_1}$, independently
of duration. For positive $R$, its trace norm is its trace. Apply
\eqref{f16v12:return-trace} and divide by $\mathcal N$. This proves
\eqref{f16v12:gluing-bound} without summing a separation-independent error
once for every lag. The factorization of interior integrals did not
factorize the posterior, and no cross-message covariance was discarded.
\end{proof}

The bounded-overlap covering constructed above is not itself a separated
family. Its message-energy theorem holds without separation, but
\eqref{f16v12:gluing-identity} is asserted only for a separated collection.
One may take a separated subfamily, or finitely colour the bounded-degree
patch-intersection graph and apply the identity within a colour. This does
not authorize ignoring cross-colour response or tensorizing the native law.

\subsection{What is closed, and what is still needed}

The exterior-law terms in V11 are no longer wholly unestimated. For the
current full fast-history law, the actual posterior-mean Gram, the
source-weighted exceptional mean energy, and the exceptional conditional
Gram trace now have duration- and volume-uniform \emph{density} bounds.
Neither an all-good collar probability nor bounded conditional means on
arbitrary deterministic exterior data was needed. A collar of depth $L$
separates reference Gaussian communication, costing $e^{-2\nu L}$, from
native conditional-law mismatch, costing the actual entropy difference
$k_D$. Vector concentration and bounded overlap prevent either cost from
being multiplied by the number of time slices.

This is not a new complete-native $\gamma^{-1}$ expansion at arbitrary
volume. The floor $\epsilon_\gamma$ is much larger than $\gamma^{-1}$;
V11's marked coefficient remainder is still justified only on its stated
bounded exterior window and at its local-size entry condition. In
particular \eqref{f16v12:conditional-Gram-estimates} does not prove that the
exceptional conditional-covariance average is $o(\gamma^{-1})$.
No derivative of a scalar entropy or pressure error is being taken.

A trace controls total squared message energy but can allow that energy to
concentrate in a coherent source direction. It is therefore not enough to
bound the operator norm of the actual exterior return independent of
$B(n+1)$, and no exponentially weighted temporal row bound for that return
is asserted. V7's exceptional-direction warning remains in force. The
reference Gaussian inequalities in this appendix are not native
Poincare or logarithmic Sobolev inequalities.

The next useful target can be stated without the now-unnecessary all-time
collar gate. For the actual joint posterior and these exact messages, prove
an estimate on \emph{every} coefficient vector,
\[
 \operatorname{Var}_{\bar P}
       \left(\sum_{s,\alpha}v_{s,\alpha}j^D_{\gamma,\alpha}\right)
            \le a_\gamma\sum_{s,\alpha}|v_{s,\alpha}|^2,
                    \qquad a_\gamma\longrightarrow0,
\]
or a source-weighted connected return with summable space--time rows.
The present proof identifies its loss: $k_D$ is charged once to a whole
vector, yielding its squared-norm expectation, but it does not distribute
that cost according to each coherent direction's variance. A
source-resolved conditional contraction or an actual posterior functional
inequality would address that loss; another entropy-density estimate or
another formal coefficient would not. If the leading native coefficient is
also to survive spatial gluing, the exceptional budget must additionally be
improved below the $\gamma^{-1}$ scale.

All statements remain conditional on the prescribed bounded \emph{bridge}
history, before the physical root Gauss restriction. Removing a cutoff on
exterior \emph{fast} fields does not prove gauge-invariant coverage of the
integrated bridge law. Infrared contact coercivity, comparison with the same
full physical transfer top, physical-time calibration and an interacting
continuum construction remain separate obligations.

\subsection{Diagnostics and source preservation}

The new diagnostic is
\begin{center}\small
\path{verify_ym_f16_v12_posterior_screening.py}.
\end{center}
It checks Gaussian Markov conditioning and the conditional-entropy chain,
whole-vector mean comparison on translated correlated Gaussians,
boundary-reaching inverse factors, bounded-overlap spatial covers,
restricted/full port-lift matching, and exact posterior gluing on positive
finite laws. All 1,201 exact assertions passed, including 934 spatial-cover and duration-count
fixtures. Assertion families and counts are recorded in
\path{ym_v12_verification_report.json}. These are finite algebraic and
combinatorial diagnostics, not proof-assistant verification of the Gaussian
functional inequalities, native entropy input, or derivative suprema.
The analytic estimates above are the proofs; the finite checks do not
certify a physical mass gap.
The new checker has SHA--256
\begin{center}\scriptsize\ttfamily
A38B12309C4CA8CAB457AB53F820A6301EC24AA2F413D5BD4FE6B42CFA776F78.
\end{center}
The available V4--V11 diagnostics were rerun unchanged and passed 9,467,
89,339, 906, 311, 609, 990, 608 and 463 assertions, respectively.

The cumulative V12 changes only the V11 title version and appends this exact
section before its final \verb|\end{document}|. Removing the appendix and
restoring the title recovers V11 byte for byte. The package records the
input hashes, exact recovery, compilation, render checks, and actual reruns
of the available V4--V11 diagnostics. Unavailable archive diagnostics and
any external Lean development are not described as rerun. No predecessor
mathematical text or bibliography string is silently corrected. The next
scheduled five-version companion checkpoint remains V15.

\begin{firewall}
V12 proves an unrestricted-exterior posterior-mean screening estimate for
whole source banks, pays the old exceptional mean and conditional-Gram
terms in density, and gives an exact separated-cylinder spatial gluing
identity with a vanishing normalized trace-norm error. All fast exterior
fields and their actual posterior weights are retained, with no all-good
history event. The error floor is $\epsilon_\gamma$, not an
$o(\gamma^{-1})$ coefficient remainder. No unrestricted-volume small
operator norm or time-row summability of the posterior return, native
functional inequality, physical bridge coverage, same-top physical
transfer gap, or interacting continuum Yang--Mills mass gap is established
by this appendix.
\end{firewall}
% END F16 V12 APPEND

% BEGIN F16 V13 APPEND -- 2026-09-17
% Insert immediately before the final \end{document} in cumulative V12.
% The predecessor is not silently corrected; only its title becomes V13.

\clearpage
\section{V13: conditional resampling, a local-test obstruction, and a directional reduction}
\label{f16v13:context}

V12 bounds the trace density of the actual posterior return, but not its
largest eigenvalue. The new literature memorandum describes the earlier V3
frontier. We do not restart there. We instead introduce a checkable local
route to the \emph{directional} inequality left open in V12. The auxiliary
operator is the exact one-cell conditional resampling generator of the same
full compact history law. Its square-root-density realization is genuinely
local and frustration free, unlike an unverified claim about the physical
Wilson logarithmic transfer.

This appendix proves three new interfaces for the present law: a finite
local integral certificate for that auxiliary gap, a resampling-energy
estimate for the actual Haar residuals, and an operator-norm reduction for
the actual posterior return. The Gaussian auxiliary gap is determined on
\emph{all} polynomial chaoses. An exact cube witness shows that the simplest pair-row test fails even
for the Gaussian reference and, eventually, for the native law. A uniform
native gap by another criterion and uniform local fourth moments are
\emph{not} certified here. Consequently
this is a rigorous conditional route beyond V12, not a new unconditional
unrestricted-volume mass-gap or posterior-contraction theorem.

\subsection{What the literature changes, and what it does not}

The supplied V12 appendix and wider-literature memorandum were read in full.
Primary-source statements and selected technical sections were checked;
the accompanying literature audit records the depth of each review. The
following choices are based on applicability, not on the prestige of an
analogy.

\paragraph{Conditional dependence is the closest interface.}
Caputo--Menz--Tetali,
\href{https://arxiv.org/abs/1405.0608}{\emph{Approximate tensorization of
entropy at high temperature}}, derive functional inequalities from explicit
conditions on dependence. Anari--Liu--Oveis Gharan,
\href{https://arxiv.org/abs/2001.00303}{\emph{Spectral Independence in
High-Dimensional Expanders and Applications to the Hardcore Model}}, require
correlation control for the distribution and its conditionals. Neither
hypothesis is supplied by V12's small relative entropy density. Their discrete
statements are not quoted as compact-group theorems. Below we prove the
continuous projection inequality actually needed, with all conditioning and
clock conventions explicit.

\paragraph{Compatible local constraints.}
Gosset--Mozgunov's
\href{https://arxiv.org/abs/1512.00088}{\emph{Local gap threshold for
frustration-free spin systems}}, Lemm--Mozgunov's
\href{https://arxiv.org/abs/1801.08915}{\emph{Spectral gaps of
frustration-free spin systems with boundary}}, and Anshu--Arad--Vidick's
\href{https://arxiv.org/abs/1602.01210}{\emph{Simple proof of the detectability
lemma and spectral gap amplification}} concern compatible local constraints.
We construct such constraints for the \emph{auxiliary resampling law}; we do
not identify their gap with that of the physical transfer. The elementary
pair-angle criterion below works on the infinite-dimensional local group
spaces and is proved directly. It is not a new general Knabe theorem.

\paragraph{A more direct multiscale alternative than geometric analogy.}
Bauerschmidt--Bodineau's
\href{https://arxiv.org/abs/1907.12308}{\emph{Log-Sobolev inequality for the
continuum sine-Gordon model}} uses a multiscale Bakry--\'Emery criterion
controlled through the Polchinski equation, including non-log-concave
measures. This is a concrete alternative should a worst-boundary local-angle
test fail. It would require estimates on the actual renormalized potential
of the present law, not just its finite Taylor coefficients. It supplies
neither a Yang--Mills inequality nor permission to replace the compact law by
a convex extension. The more recent continuum spectral-independence work
\href{https://arxiv.org/abs/2601.18748}{\emph{Sampling Sphere Packings with
Continuum Glauber Dynamics}} concerns repulsive point processes, not this
fixed-coordinate gauge history.

\paragraph{Keep spectral elimination, but do not rename a covariance inverse.}
Bach--Ballesteros--Geisler's
\href{https://arxiv.org/abs/2508.07643}{\emph{The Spectral Renormalization Flow
Based on the Smooth Feshbach--Schur Map}} establishes a smooth spectral flow.
The monograph already retains vacuum-corrected Feshbach metrics. A physical
application still needs the complement inverse on its actual domain and the
same physical vacuum. Our resampling operator is not that physical
Feshbach operator.

\paragraph{Other proposals remain separate proof obligations.}
Kenig--Merle's method combines concentration compactness and rigidity. Its use
here needs a nonvanishing profile and a critical compactness argument; no such Yang--Mills profile
follows from the current local estimates. Perelman's Ricci-flow entropy is
not the relative entropy in V12. Tao's Collatz first-passage argument and
Bourgain--Gamburd expansion need their particular transport or convolution
and nonconcentration hypotheses. Nash--Moser smoothing needs a tame inverse,
not merely constants at each formal order. The audit records these primary
sources and the precise missing structures. None is used as an unproved
shortcut. We retain one active next test: local conditional resampling and
its actual source energy, rather than opening all these branches at once.

The archive audit found earlier conditional heat-bath and assumed-Dobrushin
interfaces, including f14's initial transition-family theorem and the
monograph's \texttt{thm:SMFS-Dobrushin-gap}. Those general constructions are
credited rather than claimed as absent. The new specialization below uses
V12's \emph{full augmented fast-history law}, its exact residual synthesis,
and its separated-cylinder posterior. All proofs keep the native measure
fixed.

\subsection{An exact local frustration-free operator for the full history law}

Keep V12's bounded bridge history, $\max_{i,e}|y_{i,e}|\le r$, and set
$\mathcal N=B(n+1)$. The fast cell $c=(i,b)$ contains V5's unmoving
coordinates $\zeta_c=(x_{i,b},a_{i,b})$ when $i<n$, and only $x_{n,b}$ at the
last slice. Dimensions are at most 36. Work on the complete principal-chart
product, with its null boundary sets ignored. Normalize each one-cell factor
of $J^f_\gamma\,d\zeta$ to a probability $m_{\gamma,c}$, and put
$m_\gamma=\bigotimes_c m_{\gamma,c}$. These normalizations multiply the
partition function by a scalar only. In particular the endpoint chord powers
$w_i=1/2$ and the interior powers one are retained. Then the same native law is
\[
 dP=Z_m^{-1}e^{-U_\gamma}\,dm_\gamma.
\]
The Gaussian endpoint amplitudes belong to $U_\gamma$, not to a newly chosen
vacuum. At fixed finite $B,n,\gamma$, $U_\gamma$ is bounded on this product
and $P$ is equivalent to $m_\gamma$ with bounded positive density ratio.

Define the bounded operators on $L^2(P)$
\begin{equation}
 E_c f=\mathbb E_P[f\mid\zeta_{c^c}],\qquad Q_c=I-E_c,
 \qquad \mathcal L_P=\sum_c Q_c,
 \qquad
 \mathcal E_P(f)=\sum_c\mathbb E_P\operatorname{Var}_P(f\mid\zeta_{c^c}).
 \label{f16v13:generator}
\end{equation}
Each cell is updated at rate one. There is no division by $\mathcal N$.
Let $\lambda_P$ be the infimum of $\mathcal E_P(f)/\operatorname{Var}_P f$
over nonconstant $f$.

\begin{proposition}[Local projections with the exact common state]
\label{f16v13:parent}
Each $Q_c$ is an orthogonal projection. The kernel of $\mathcal L_P$ consists
exactly of constants. Under the unitary
\[
 U_P:L^2(P)\longrightarrow L^2(m_\gamma),\qquad
 U_Pf=\Psi f,\qquad \Psi=(dP/dm_\gamma)^{1/2},
\]
the transformed projections have finite spatial and temporal interaction
range and annihilate the same state $\Psi$. If no primitive action term meets
both $c$ and $d$, then $E_cE_d=E_dE_c=E_{\{c,d\}}$, where the last operator
conditions on all variables except those two cells. Moreover, at every fixed
finite regulator,
\begin{equation}
 \lambda_P\ge e^{-\operatorname{osc}U_\gamma}>0.
 \label{f16v13:crude-gap}
\end{equation}
The bound in \eqref{f16v13:crude-gap} is generally exponentially poor in
$\gamma\mathcal N$ and is not the desired uniform result.
\end{proposition}
\begin{proof}
Conditional expectation is an orthogonal projection, and
$\|Q_cf\|_2^2=\mathbb E\operatorname{Var}(f\mid\zeta_{c^c})$.
For fixed exterior $z$, write the sum of terms meeting $c$ as $V_c(t;z)$.
The conditional amplitude is
\[
 \psi_c(t;z)=\frac{e^{-V_c(t;z)/2}}
 {\left(\int e^{-V_c(s;z)}dm_{\gamma,c}(s)\right)^{1/2}}.
\]
Direct multiplication gives
\begin{equation}
 (U_PE_cU_P^{-1}F)(t,z)
     =\psi_c(t;z)\int\psi_c(s;z)F(s,z)\,dm_{\gamma,c}(s).
 \label{f16v13:local-parent-kernel}
\end{equation}
Only cells in the fixed primitive word neighbourhood of $c$ enter $\psi_c$.
Thus this rank-one-in-$c$ operator acts trivially on every other exterior
coordinate. All transformed $Q_c$ kill $\Psi$ because all $Q_c$ kill one.
If no word meets $c,d$, their conditional pair density at fixed remaining
variables is a product. Integrating the two coordinates in either order
proves the commutation and joint-conditional identity, even when their
exterior neighbourhoods overlap.

For the kernel and gap assertions, let $p_-,p_+$ be the essential lower and
upper bounds of $dP/dm_\gamma$. Product variance tensorization gives
\[
 \operatorname{Var}_P f\le p_+\operatorname{Var}_{m_\gamma}f
 \le p_+\sum_c\mathbb E_{m_\gamma}\operatorname{Var}_{m_{\gamma,c}}f.
\]
Conditional variance averaged over the exterior is the infimum of
$\int|f-g(\zeta_{c^c})|^2dP$ over all exterior functions $g$. It is therefore
at least $p_-$ times the corresponding infimum under $m_\gamma$. Combining
these inequalities proves $\mathcal E_P(f)\ge(p_-/p_+)\operatorname{Var}_P f$.
Here $p_-/p_+\ge e^{-\operatorname{osc}U_\gamma}$. It also proves that the
kernel consists of constants. Approximation in $L^2$ extends the argument
from bounded functions. The complete principal angular domains have bounded
radius, and the number of positive action terms is $O(\mathcal N)$; hence
$\operatorname{osc}U_\gamma\le C_r\gamma\mathcal N$ for $\gamma\ge1$.
\end{proof}

This is a genuine local frustration-free \emph{auxiliary} operator. Its clock
is cell resampling, not lattice Euclidean time. Its constant state in $L^2(P)$
is not identified with the physical transfer vacuum. The map $U_P$ does not
assert a unitary equivalence to $-\log T_{\rm full,\gamma}$. This distinction
is essential even if a good gap for $\mathcal L_P$ is eventually proved.

\subsection{The reference gap is an all-chaos statement}

Let $Q=N(m,H^{-1})$ be V12's Gaussian law on the same Euclidean coordinate
partition, and let $\mathcal L_Q$ be \eqref{f16v13:generator} with Gaussian
conditionals. Write $D_H=\bigoplus_c H_{cc}$.

\begin{theorem}[Exact Gaussian resampling edge on the full function space]
\label{f16v13:Gaussian-gap}
The rate-one-per-cell Gaussian generator satisfies
\begin{equation}
 \lambda_Q=\lambda_{\min}(D_H^{-1/2}HD_H^{-1/2})
                         \ge h_-/h_+>0.
 \label{f16v13:Gaussian-edge}
\end{equation}
The equality is on all of $L^2(Q)$, not just on linear observables. It is
uniform in spatial volume, duration, and the Gaussian mean. It does not
compare $\mathcal L_Q$ with $\mathcal L_P$.
\end{theorem}
\begin{proof}
Whiten the Gaussian to $g=H^{1/2}(\zeta-m)$ and identify its first chaos with
the Euclidean coefficient space. Let $I_c$ inject the coordinates of cell
$c$. The innovation subspace orthogonal to all exterior linear functions
has orthogonal projection
\[
 R_c=H^{1/2}I_cH_{cc}^{-1}I_c^{\mathsf T}H^{1/2}.
\]
Indeed the columns of $H^{1/2}I_c$ are orthogonal to the exterior coordinate
columns of $H^{-1/2}$, and their Gram is $H_{cc}$. Therefore
\[
 \sum_cR_c=H^{1/2}D_H^{-1}H^{1/2}
\]
has the same eigenvalues as $D_H^{-1/2}HD_H^{-1/2}$. On the first chaos the
generator is exactly this sum, so its smallest eigenvalue is an upper bound
on the full gap.

On the $k$th homogeneous Gaussian chaos, conditional expectation onto
exterior coordinates is the restriction of $(I-R_c)^{\otimes k}$ to the
symmetric tensor space. One can see this by taking conditional expectations
of the Wick exponential
$\exp(\langle t,g\rangle-|t|^2/2)$ and comparing its homogeneous
coefficients. Projections acting on different tensor legs commute, and their
joint eigenvalues are zero or one. Hence
\[
 I-(I-R_c)^{\otimes k}\succeq\frac1k\sum_{j=1}^k(R_c)_j.
\]
Sum over $c$. The right side is at least $\lambda_Q I$ on each $k\ge1$,
where $\lambda_Q$ denotes the first-chaos value just computed. Hermite
polynomials are dense in Gaussian $L^2$, so the same lower bound holds on the
orthogonal complement of constants. Finally $H\succeq h_-I$ and
$D_H\preceq h_+I$ prove the stated comparison.
\end{proof}

Exact Gaussian sampler rates are classical; see Roberts--Sahu,
\href{https://doi.org/10.1111/1467-9868.00070}{\emph{Updating Schemes,
Correlation Structure, Blocking and Parameterization for the Gibbs Sampler}}
(1997). The direct proof above fixes the present continuous-time convention
and checks every chaos. V3's conditional covariance lower precision and
V8's fast transfer ratio concern different operators and are not substituted
for this calculation. Small $D(P\Vert Q)/\mathcal N$ still does not transfer
\eqref{f16v13:Gaussian-edge} to $P$.

\subsection{A local native conditional-angle criterion}

Fix two cells $c,d$. Conditional on all other cells $z$, use the full native
pair law $P_{cd}^z$. Define its maximal correlation
\begin{equation}
 \theta_{cd}(z)=\sup
 \left\{|\mathbb E[f(\zeta_c)g(\zeta_d)]|:
       \mathbb Ef=\mathbb Eg=0,\ \mathbb Ef^2=\mathbb Eg^2=1\right\}.
 \label{f16v13:theta}
\end{equation}
The expectations in this definition are under this \emph{conditional pair
law}. Define $\bar\theta_{cd}=\operatorname*{ess\,sup}_z\theta_{cd}(z)$,
with $\bar\theta_{cc}=0$. Only the finite word neighbourhood of the pair
enters. Noninteracting pairs have $\bar\theta_{cd}=0$ exactly. The essential
supremum allows every exterior fast value, not merely V11's bounded window.

\begin{theorem}[Weighted local-angle certificate for the native auxiliary gap]
\label{f16v13:angle-gap}
For positive weights $0<w_c\le1$, put
\begin{equation}
 \delta_w=\min_c\left(w_c-\sum_{d\ne c}w_d\bar\theta_{cd}\right).
 \label{f16v13:weighted-margin}
\end{equation}
If $\delta_w>0$, then
\begin{equation}
 \sum_c w_cQ_c\succeq\delta_w(I-\mathbb E_P),\qquad
                       \lambda_P\ge\delta_w.
 \label{f16v13:angle-conclusion}
\end{equation}
This assertion holds on the full compact-coordinate $L^2(P)$; no
finite-dimensional local-spin assumption or chart cutoff is required.
The criterion is sufficient, not necessary. A failed margin does not imply
a small actual gap.
\end{theorem}
\begin{proof}
First condition on the complement of one pair. On its mean-zero $L^2$
space let $\mathcal U,\mathcal V$ be the subspaces of centered functions of
$c$ and $d$, respectively. Their cross inner-product norm is
$\theta=\theta_{cd}(z)$. The synthesis map $(u,v)\mapsto u+v$ from
$\mathcal U\oplus\mathcal V$ has squared norm at most $1+\theta$, because
\[
 \|u+v\|^2\le\|u\|^2+\|v\|^2+2\theta\|u\|\|v\|
                  \le(1+\theta)(\|u\|^2+\|v\|^2).
\]
Its product with its adjoint is $E_c+E_d$ on this centered space (the
subspaces occur in the opposite cell order). Thus $Q_c+Q_d\succeq
(1-\theta)I$ there. Constants form its kernel. Spectral calculus yields
$(Q_c+Q_d)^2\succeq(1-\theta)(Q_c+Q_d)$, or
\begin{equation}
 Q_cQ_d+Q_dQ_c\succeq-\theta(Q_c+Q_d).
 \label{f16v13:anticommutator}
\end{equation}
This also holds on constants. Integrate the fiber inequality and replace
$\theta$ by $\bar\theta_{cd}$. Noninteracting pairs commute and have a
nonnegative anticommutator, consistently with $\bar\theta=0$.

For $\mathcal L_w=\sum_cw_cQ_c$, expand the square and use the last
inequality for each unordered pair:
\[
 \mathcal L_w^2\succeq
 \sum_c w_c\left(w_c-\sum_{d\ne c}w_d\bar\theta_{cd}\right)Q_c
                  \succeq\delta_w\mathcal L_w.
\]
The spectral measure of this bounded positive operator is supported on
$\{0\}\cup[\delta_w,\infty)$. Its kernel is constants by
Proposition~\ref{f16v13:parent} and positivity of every weight. Since
$\mathcal L_P\succeq\mathcal L_w$, this proves both claims. No covariance
bound on native test functions was assumed in this proof.
\end{proof}

There are finitely many pair and endpoint cell types at the fixed word range.
Thus a \emph{uniformly successful} family of local pair certificates and
bounded weights would give a uniform margin independent of $B,n$. Taking all
weights one gives $1-\sup_c\sum_d\bar\theta_{cd}$. Alternatively, lower
bounds for \eqref{f16v13:weighted-margin} are linear inequalities in the
weights; a proposed rational solution can be checked directly. This is
ordinary weighted projection bookkeeping, not an import of sieve weights.
The next proposition will show why this particular one-cell criterion
cannot provide the required weak-coupling margin.

\subsection{Finite local integral certificates with their unresolved continuum paid}

A small matrix sampled at selected exterior points is not by itself an
upper bound for $\bar\theta$. We now give a finite upper certificate that
also pays discretization and the complete exterior parameter range.

Use angular coordinates $u=\zeta_c/\sqrt\gamma$, $v=\zeta_d/\sqrt\gamma$,
and angular coordinates $z$ on their finite exterior word neighbourhood.
Each domain is a product of principal three-balls of radius $\pi$, with
null boundary sets ignored. The two-cell reference factors are normalized
products of the appropriate powers of $j$, independent of $z$. The exact
conditional potential consists of all terms meeting $c$ or $d$, with the
other inputs frozen; its density is
\[
 p_z(u,v)=\frac{e^{-V_z(u,v)}}{\int e^{-V_z}\,dm_cdm_d}.
\]
Here and in this subsection $m_c,m_d$ are the angular versions of
$m_{\gamma,c},m_{\gamma,d}$. The local exponential-word formulas extend
continuously to the chosen closed angular product. Fixed derivative and
incidence bounds give finite constants
\begin{equation}
 \operatorname{osc}_{u,v}V_z\le O_\gamma,\qquad
 \operatorname{Lip}_{u,v,z}(V)\le A_\gamma,
 \qquad O_\gamma+A_\gamma\le C_r(1+\gamma).
 \label{f16v13:finite-card-constants}
\end{equation}
The constants are uniform in ambient $B,n$ and all principal exterior
angles. Endpoint Gaussian quadratics are included in these bounds. They
are not claims of useful smallness. For a certificate uniform over bridge
values, include the finitely many bridge inputs in their prescribed compact
window among the parameter net coordinates, with its corresponding metric.

Partition each of the pair domains into finitely many Borel bins of diameter
at most $h$, omitting reference-null bins. At each parameter value define
\begin{equation}
 A(z)_{ab}=
  \frac{P_z(U_a\times V_b)}{\sqrt{P_z(U_a)P_z(V_b)}}.
 \label{f16v13:probability-matrix}
\end{equation}
Use an $h$-net $\mathcal Z_h$ of the entire closed local parameter domain.
It is permitted to be very large. Let $\epsilon_{\rm int}$ be a certified
operator-norm enclosure error for the matrices used in place of the exact
matrices \eqref{f16v13:probability-matrix}.

\begin{theorem}[An upper, not merely lower, local correlation certificate]
\label{f16v13:finite-card}
With the conventions above, the full pair correlation obeys
\begin{equation}
 \boxed{
 \bar\theta_{cd}\le\min\left\{1,
  \max_{z\in\mathcal Z_h}\sigma_2(\widehat A(z))
  +2e^{4O_\gamma}(e^{4A_\gamma h}-1)+\epsilon_{\rm int}\right\}.}
 \label{f16v13:finite-upper}
\end{equation}
The error is independent of the ambient volume and duration. It is not
claimed small on a practical mesh. At an exact fixed parameter,
$\sigma_1(A(z))=1$ and $\sigma_2(A(z))\le\theta_{cd}(z)$; without the
explicit remaining error the finite matrix gives the wrong inequality for
a gap certificate. Singular values are zero-padded when necessary.
\end{theorem}
\begin{proof}
Let $p_c,p_d$ be the marginal densities at fixed $z$. The normalized
conditional-expectation kernel on the fixed reference spaces is
\[
 K_z(u,v)=p_z(u,v)/\sqrt{p_c(u)p_d(v)}.
\]
It is compact since its positive kernel is bounded. Its top singular value
is one, with normalized singular functions $\sqrt{p_c},\sqrt{p_d}$;
conditional expectation is a contraction and maps these functions to each
other. Its second singular value is exactly \eqref{f16v13:theta} by restricting
to their orthogonal complements. Positivity of the pair density excludes
nonconstant perfectly correlated functions, though that strictness is not
needed for the bound below.

Equivalently, on the marginal probability spaces the kernel is
$k_z=p_z/(p_cp_d)$. Orthogonal projection onto functions constant on the
bins gives precisely \eqref{f16v13:probability-matrix}, in the normalized
indicator bases. Both constant functions remain in the projection, so the
top singular value remains one and the second cannot exceed the full one.
This projection includes its original marginals, not arbitrarily uniform
bin weights.

The potential oscillation gives
$e^{-O_\gamma}\le p_z,p_c,p_d\le e^{O_\gamma}$, so
$k_z\le e^{3O_\gamma}$. At fixed $z$, the logarithm of a marginal density
has Lipschitz constant at most $A_\gamma$ in its own coordinate: this follows
by integrating the pointwise ratio bound for $e^{-V}$. Therefore the
oscillation of $\log k_z$ on any product of two bins is at most
$4A_\gamma h$. The doubly projected kernel is the binwise conditional
average of $k_z$ under $P_c\otimes P_d$. Its difference from $k_z$ has
absolute value at most
$e^{3O_\gamma}(e^{4A_\gamma h}-1)$. Its Hilbert--Schmidt norm, in these
probability spaces, has the same upper bound. Singular-value stability
therefore bounds the full second singular value by that of $A(z)$ plus
this error.

For $z,z'$ at distance at most $h$, normalized joint density ratios have
logarithmic absolute value at most $2A_\gamma h$; integrating gives the
same bound for each marginal ratio. Hence
$|\log(K_z/K_{z'})|\le4A_\gamma h$. Both kernels on the fixed reference
spaces have supremum at most $e^{2O_\gamma}$. Their Hilbert--Schmidt
distance is at most $e^{2O_\gamma}(e^{4A_\gamma h}-1)$. This step uses a
fixed reference space, not a comparison of operators on two different
unidentified marginal $L^2$ spaces. Apply singular-value stability once
more to pass to a nearest net point. Add the two errors and the certified
matrix enclosure error, and bound their coefficient by $2e^{4O_\gamma}$.
This proves \eqref{f16v13:finite-upper}. All domains, including large angular
inputs, are covered by the bins and net. No chart-complement probability
has been dropped.
\end{proof}

The certificate uses a fixed number of local variables but can require a
mesh exponentially fine in $\gamma$. It is an existence and validation
interface, not a claim that this computation has been performed for the
native Wilson pairs. The diagnostic supplied here tests the matrix identities
on finite probability models; it supplies neither native integral enclosures
nor a successful value of \eqref{f16v13:weighted-margin}. A failure of this
worst-boundary pair test would motivate larger blocks or a multiscale
criterion, not a declaration that the physical or auxiliary gap vanishes.

\subsection{The actual cube rejects the naive pair row, already at Gaussian order}

Before spending effort on native integral enclosures, one should test whether
the proposed sufficient criterion can even recognize the reference model.
For the unmoving cells chosen above, the answer is negative. The obstruction
below is for the \emph{actual} parity-cube matrices, not an arbitrary Gaussian
countermodel.

\begin{proposition}[A certified Gaussian and native failure of the one-cell row test]
\label{f16v13:pair-obstruction}
For $m\ge4$ and $n\ge4$, consider the three interior cells at times $1,2,3$
in one fixed spatial cube. In the Gaussian reference, each of the two
consecutive pair conditionals has maximal correlation greater than $4/5$.
Consequently no positive weights can make
\eqref{f16v13:weighted-margin} positive for this Gaussian cell partition.

There is a finite $\gamma_0$, determined by the fixed local geometry, such
that at zero bridge history the same impossibility holds for the
\emph{native} worst-exterior pair coefficients for every
$\gamma\ge\gamma_0$, $m\ge4$, $n\ge4$. This is failure of this sufficient
row certificate, not a vanishing native spectral gap.
\end{proposition}
\begin{proof}
We provide a rational witness. For one colour, let $\mathsf B,\mathsf F,N$
and $\mathsf H$ be the inherited rooted cube matrices, with the five chords
ordered $(a,b,c,d,e)$ as in f15 V97 and nonroot vertices in lexicographic
order. Put
\[
 L=\mathsf B^{\mathsf T}\mathsf B+3I_7,
 \qquad K=N+\mathsf F^{\mathsf T}L\mathsf F.
\]
The port precision has diagonal spatial cube block $L$: each nonroot port
meets three bridges, and none of these bridges has both endpoints in the
same cube for $m\ge4$. The mixed magnetic diagonal Gram is $2N$, by f15
V99. Substituting $z_t=a_t-\mathsf F(x_t-x_{t+1})$ into the exact quadratic
kinetic form gives the diagonal interior cell block and next-time block
\begin{equation}
 A=\begin{pmatrix}
       \mathsf H+2N+2K&-\mathsf F^{\mathsf T}L\\
       -L\mathsf F&L
    \end{pmatrix},\qquad
 C=\begin{pmatrix}-K&0\\ L\mathsf F&0\end{pmatrix}.
 \label{f16v13:actual-pair-blocks}
\end{equation}
Thus conditioning all other cells gives pair precision
$\left(\begin{smallmatrix}A&C\\C^{\mathsf T}&A\end{smallmatrix}\right)$.
All its matrices are rational. Conditioning affects its mean, not this
precision. Endpoint amplitude precisions do not enter these interior blocks.

For two blocks with positive precision diagonal $A$, whitening each block
and then taking a singular-value decomposition of
$A^{-1/2}CA^{-1/2}$ decomposes the precision into scalar pairs
$\left(\begin{smallmatrix}1&s\\s&1\end{smallmatrix}\right)$.
Such a pair has linear correlation $-s$. In particular the maximal
correlation is at least
\[
 \frac{|u^{\mathsf T}Cv|}{\sqrt{(u^{\mathsf T}Au)(v^{\mathsf T}Av)}}
\]
for any two coefficient vectors $u,v$.

With coordinates ordered as five chords followed by seven ports, choose
\[
 \begin{split}
 u&=(0,0,0,0,0;\ 1,-1,2,0,1,-1,3)^{\mathsf T},\\
 v&=(1,-1,-3,-2,-4;\ 1,-1,2,0,1,-1,2)^{\mathsf T}.
 \end{split}
\]
Direct rational multiplication of the displayed inherited cube matrices gives
\begin{equation}
 u^{\mathsf T}Au=86,\quad v^{\mathsf T}Av=825/8,\quad
 u^{\mathsf T}Cv=303/4,\quad
 \frac{|u^{\mathsf T}Cv|^2}{(u^{\mathsf T}Au)(v^{\mathsf T}Av)}
    =\frac{30603}{47300}>\frac{16}{25}.
 \label{f16v13:rational-witness}
\end{equation}
The last strict inequality has difference $331/47300>0$.
The checker reconstructs the incidence matrices rather than importing
floating-point blocks. One colour suffices; the other colours are independent
copies in this Gaussian calculation. This proves both pair lower bounds.

If all row margins were positive, the three corresponding weights would
satisfy $w_1>(4/5)w_2$, $w_3>(4/5)w_2$, and
$w_2>(4/5)(w_1+w_3)$. Dropping all the other nonnegative row terms is
legitimate. Substitution gives $w_2>(32/25)w_2$, a contradiction. The exact
Gaussian gap in Theorem~\ref{f16v13:Gaussian-gap} remains positive: taking
absolute pair interactions has lost information needed to see it.

For the native assertion, fix the exterior fast coordinates and bridges to
zero in this pair conditional. All terms meeting the pair and all original
Haar factors are retained. Its fixed-dimensional complete-chart integral
has the independent magnetic and rooted port pins from V5. Its normalized
polynomial moments through degree two therefore converge to those of the
Gaussian pair in \eqref{f16v13:actual-pair-blocks}: Taylor expansion on fixed
rescaled balls gives local convergence, the complete Gaussian pin controls
their complements, and the limiting partition function is strictly positive.
This is an ordinary fixed-dimensional limit, not a growing-volume norm
comparison. There are only the two fixed local pair types here; their words
are independent of the ambient $m,n$ once these interior cells exist.

Choose Gaussian linear observables whose centered correlation is greater
than $4/5$, as guaranteed by the strict inequality above. Their variances
are positive. Moment convergence makes the corresponding native correlation
greater than $4/5$ for sufficiently large $\gamma$, with a threshold depending
only on these local words. At each such finite $\gamma$ the conditional
moments vary continuously with the local exterior principal angles near
zero. Thus the inequality holds on a nonempty open exterior neighbourhood,
which has positive native marginal probability. The essential supremum in
$\bar\theta$ is therefore greater than $4/5$, not merely its value on a
null conditioning event. Apply the same three-weight contradiction. No
claim of uniform convergence over all exterior fields was used.
\end{proof}

This check changes the next direction. The all-boundary one-cell test cannot
close the weak-coupling program on this partition; a finer mesh will not fix
it. The exact local integral certificates remain valid for a different
partition or update family. Useful alternatives are larger conditional
blocks, a coordinate/update organization that respects the balanced Gaussian
split, or multiscale functional inequalities retaining the signed Gaussian
metric rather than absolute pair rows. None of those native alternatives
is assumed successful here. This is precisely why a Gaussian or small-volume
preflight should precede a claimed literature-based global closure.


\subsection{Local resampling energy of the actual corrected sources}

Return to the unscaled notation $\zeta$ for V5's rescaled fast coordinates.
Define actual full-law moments
\begin{equation}
 m_2(P)=\sup_c\mathbb E_P(1+|\zeta_c|^2),\qquad
 m_4(P)=\sup_c\mathbb E_P(1+|\zeta_c|^4).
 \label{f16v13:moment-inputs}
\end{equation}
They are finite at each regulator. Uniformity in $B,n,\gamma$ is not asserted.
V5's averaged second moments and V12's entropy density do not supply the
supremum fourth-moment bound in \eqref{f16v13:moment-inputs}.

Let $e=(e^{\rm br},e^{\rm pt})$ be V7's two primitive remainder banks, where
$e^{\rm br}=J_\gamma-J_0$ for the common sources and
$e^{\rm pt}=S_\gamma-S_0$ for the Haar port insertions. Each primitive
component has a fixed-size local fast carrier $C_j$, bounded row and column
incidence, and the inherited \emph{complete-chart} estimate
\begin{equation}
 |e_j(\zeta)|\le C_r\gamma^{-1/2}(1+|\zeta_{C_j}|^2).
 \label{f16v13:primitive}
\end{equation}
For the full exact lift, V7 gives
$\mathbf r=e^{\rm br}+\mathsf P^{\mathsf T}e^{\rm pt}$ with
$\|\mathsf P\|\le C$. All these are actual native functions.

\begin{theorem}[Directional innovation energy, without a native gap assumption]
\label{f16v13:innovation}
For every primitive coefficient vector $b$ and every common-source vector
$v$, uniformly in volume and duration,
\begin{align}
 \mathcal E_P(b^{\mathsf T}e)&\le
                  C_r\frac{m_4(P)}{\gamma}\|b\|^2,\notag\\
 \mathcal E_P(v^{\mathsf T}\mathbf r)&\le
                  C_r\frac{m_4(P)}{\gamma}\|v\|^2.
 \label{f16v13:innovation-energy}
\end{align}
For the affine Gaussian port score $S_0=\partial_aU_0$, evaluated under the
native law, one also has
\begin{equation}
 \mathcal E_P(b^{\mathsf T}S_0)\le C m_2(P)\|b\|^2.
 \label{f16v13:linear-energy}
\end{equation}
No native covariance mixing estimate enters these inequalities.
\end{theorem}
\begin{proof}
At stationarity let $\zeta^{(c)}$ be obtained by replacing cell $c$ by an
independent draw from its \emph{actual} conditional given all other cells.
Both $\zeta$ and $\zeta^{(c)}$ have law $P$, and
\[
 \mathcal E_P(f)=\frac12\sum_c
                 \mathbb E_P|f(\zeta)-f(\zeta^{(c)})|^2.
\]
Let $J(c)$ be the primitive scores whose carrier contains $c$.
Its cardinality, and the number of cells in any one carrier, are bounded
by a fixed number $k_0$. Other score components are unchanged. Thus
\[
 |b^{\mathsf T}(e(\zeta)-e(\zeta^{(c)}))|^2
 \le k_0\sum_{j\in J(c)}|b_j|^2
                   |e_j(\zeta)-e_j(\zeta^{(c)})|^2.
\]
Stationarity bounds the last squared expectation by $4\mathbb E_P e_j^2$.
Using \eqref{f16v13:primitive} and the fixed number of coordinates in its
carrier gives $\mathbb E_Pe_j^2\le C_r m_4(P)/\gamma$.
Sum over $c$; every $j$ occurs at most $k_0$ times. This proves the first
inequality without summing all primitive variances for each direction.
The bounded synthesis $b=(v,\mathsf P v)$ proves the second.

For the affine score, write $b^{\mathsf T}S_0=l^{\mathsf T}\zeta+$constant.
The linear map $b\mapsto l$ is a corner of the Gaussian precision $H$, so
$\|l\|\le h_+\|b\|$. Replacing cell $c$ changes this expression by
$l_c^{\mathsf T}(\zeta_c-\zeta_c^{(c)})$. Its squared expectation is at most
$4|l_c|^2\mathbb E_P|\zeta_c|^2$. Sum the resampling identity and use the
bound on $l$. This proves \eqref{f16v13:linear-energy}.
\end{proof}

This is a local \emph{Dirichlet-energy} estimate on every coefficient vector,
not the covariance inequality V12 sought. It differs from V7's total residual
energy and from Gaussian energy transfer. Dividing it by a proved native
auxiliary gap would make it a directional covariance estimate; dividing by
a Gaussian gap without a comparison would not.

\subsection{A coherent posterior reduction on the same actual law}

Use a separated cylinder family and source selection $\chi$ exactly as in
V12's gluing theorem. There are no repeated source marks. All source carriers
are interior, source depths are at least $L$, and enlarged cylinder patches
have overlap at most $q_0$. Put $O$ for the common exterior of all these
interiors and $E_O=\mathbb E_P[\,\cdot\mid\zeta_O]$. The posterior is V12's
actual posterior with its conditional partition functions included.

Stack the restricted lifts and extend them by zero. Let $\Delta$ be their
difference from the corresponding columns of the full lift. The inherited
identities, with all exterior fields admitted, are
\begin{equation}
 \|\Delta\|\le C\sqrt{q_0}e^{-\nu L},\qquad
 \mathbf r^D=\chi\mathbf r-\Delta^{\mathsf T}e^{\rm pt},\qquad
 \mathbf g^D-\chi\mathbf g=-\Delta^{\mathsf T}S_0.
 \label{f16v13:matched}
\end{equation}
The last vector is $O$-measurable: the lift mismatch is harmonic on each
free cylinder. Equivalently it is the difference of its two Gaussian
conditional source constants. The actual conditional mean vector is
\begin{equation}
                  \mathbf j=\mathbf g^D+E_O\mathbf r^D.
 \label{f16v13:message-identity}
\end{equation}
Separation ensures that this one common conditional expectation gives each
cylinder's message. It must not be replaced by unrelated projections for
an overlapping cover.

\begin{theorem}[Gap-and-moment reduction of the coherent posterior return]
\label{f16v13:posterior-reduction}
For the exact finite-regulator native law and the separated family above,
\begin{align}
 \|\Gamma_\gamma\|_{\rm op}
   &\le\frac{C_r}{\lambda_P}\frac{m_4(P)}{\gamma},\notag\\
 \|\mathsf R_{\rm post}\|_{\rm op}
   &\le\frac{C_r q_0}{\lambda_P}
          \left(\frac{m_4(P)}{\gamma}+m_2(P)e^{-2\nu L}\right),\notag\\
 \|\chi\mathsf M_\gamma\chi-\mathsf L_D\|_{\rm op}
   &\le\frac{C_r q_0}{\lambda_P}
          \left(\frac{m_4(P)}{\gamma}+m_2(P)e^{-2\nu L}\right).
 \label{f16v13:directional-reduction}
\end{align}
Here $\Gamma_\gamma$ is V7's full residual covariance,
$\mathsf R_{\rm post}=\operatorname{Cov}_{\bar P}(\mathbf j)$ and
$\mathsf L_D=\bigoplus_s\mathbb E_{\bar P}\mathsf M^{D_s}_\gamma$ are V12's
literal matrices. The constant absorbs the factor four in the last bound.
There is no normalization by $\mathcal N$ or by the number of sources.

In particular, a \emph{separately verified} uniform auxiliary gap
$\lambda_P\ge\delta_0>0$ and local moment bound $m_4(P)\le M_4$ would give
an $O_r(\gamma^{-1}+e^{-2\nu L})$ coherent posterior and gluing bound.
This final implication does not assert that its two inputs have been proved.
\end{theorem}
\begin{proof}
The finite-regulator gap is positive by \eqref{f16v13:crude-gap}, so its
Poincar\'e inequality applies to every present square-integrable function.
Apply it to $v^{\mathsf T}\mathbf r$ and use
Theorem~\ref{f16v13:innovation}. Taking the supremum over unit $v$ proves
the first assertion. This step makes clear which \emph{native} gap is used.

For the posterior, \eqref{f16v13:matched} expresses any
$v^{\mathsf T}\mathbf r^D$ as $b^{\mathsf T}e$, with
$\|b\|^2\le Cq_0\|v\|^2$, using the bounded full lift and the mismatch
bound. Therefore
\[
 \operatorname{Var}_P(v^{\mathsf T}\mathbf r^D)
       \le C_r q_0 m_4(P)(\lambda_P\gamma)^{-1}\|v\|^2.
\]
Conditional expectation on the \emph{one common} exterior is a contraction
on centered $L^2(P)$. Consequently the same bound applies to
$\operatorname{Var}_P(E_O(v^{\mathsf T}\mathbf r^D))$. The random part of
$v^{\mathsf T}\mathbf g^D$ is $-(\Delta v)^{\mathsf T}S_0$; the deterministic
full $g$ changes no variance. The last inequality in
Theorem~\ref{f16v13:innovation} gives
\[
 \operatorname{Var}_P(v^{\mathsf T}\mathbf g^D)
       \le Cq_0m_2(P)\lambda_P^{-1}e^{-2\nu L}\|v\|^2.
\]
Use \eqref{f16v13:message-identity} and
$\operatorname{Var}(F+G)\le2\operatorname{Var}F+2\operatorname{Var}G$.
The law of any $O$-measurable function under $P$ is exactly its law under
$\bar P$, proving the second assertion. Native large fields and their
posterior weights have not been conditioned away.

Finally V12 proves the exact identity
$\chi\mathsf M_\gamma\chi-\mathsf L_D=
-\operatorname{Off}_{>1}\mathsf R_{\rm post}$. The Fourier phase extraction
used in V7 is contractive in operator norm as well as trace norm: with
$U_t=\sum_i e^{\mathrm iit}\Pi_i$, each diagonal-offset extraction is a
Fourier coefficient of $U_t A U_t^*$ and has norm at most $\|A\|$.
Subtracting offsets $-1,0,1$ gives
$\|\operatorname{Off}_{>1}A\|\le4\|A\|$. This proves the last bound for
all lags at once. It does not give an exponentially weighted temporal row.
Since $m_2(P)\le C m_4(P)$, the final conditional implication follows.
\end{proof}

A successful alternative local-gap certificate, combined with native local
moment estimates, would supply the two inputs of the last theorem without
assuming the desired posterior covariance itself. The one-cell absolute-row
criterion has instead been ruled out by Proposition~\ref{f16v13:pair-obstruction}. This is the point of the reduction. Merely
substituting the Gaussian edge, or using an empirical conditional-correlation
matrix without its upper error, would be circular.

\subsection{Concrete next tests and the boundary of the advance}

The Gaussian all-chaos gap, the native local projection construction, the
local-angle \emph{criterion}, its finite enclosure formula, the innovation
energy, and the posterior reduction are proved here. A native uniform gap by an applicable criterion and uniform local fourth
moments remain unverified; the one-cell pair-row margin is ruled out at
large coupling, rather than left as a hoped-for hypothesis. Inserting
only the crude finite-regulator gap and the compact bound
$m_4(P)\le C(1+\gamma^2)$ into \eqref{f16v13:directional-reduction} does not
improve V12. We therefore do not report a new unconditional
unrestricted-volume contraction or a preserved leading coefficient.

The immediate test has in fact rejected the simplest proposed route:
Proposition~\ref{f16v13:pair-obstruction} certifies two adjacent-time
correlations greater than $4/5$ using the actual cube matrices, and proves
that positive weights cannot rescue this cell partition. We therefore do
not propose spending the next iteration merely tightening that row bound.
A larger-block or balanced-update Gaussian test must first reproduce the
known all-chaos reference gap. Only then is a complete native local-angle
or multiscale comparison worth attempting. Failure of the pair test is not
failure of a native gap; it is a diagnosis of this sufficient criterion.

A separate local-moment estimate must use the \emph{same} native law. The
supremum over cells in \eqref{f16v13:moment-inputs} cannot be replaced by an
average or by Gaussian moments. Even successful gap-and-moment inputs would
not supply temporal row decay, an $o(\gamma^{-1})$ exceptional coefficient
budget, gauge-invariant coverage of the retained bridge law, or a comparison
to the same physical transfer top. Those obligations are unchanged.

The wider literature thus changes the active direction: rather than infer a
coherent bound from another small trace density, construct and test local
conditional constraints whose common measure is exact, and combine their
native gap with source-local innovation energy. The established Gaussian
edge is a reference check, not the missing native theorem. Physical
Yang--Mills reconstruction and mass generation remain separate.

\clearpage
\subsection{Verification and preservation}

The accompanying diagnostic is
\begin{center}\small
\path{verify_ym_f16_v13_conditional_angles.py}.
\end{center}
It checks finite exact conditional projections and their common state,
conditional pair correlations, anticommutator signs, weighted-gap algebra,
Gaussian block innovations and tensor-chaos inequalities, the exact actual-cube
pair obstruction, normalized
probability-matrix compression, and separated-posterior centering.
All 222 exact assertions passed, including 12 actual-cube
obstruction checks. The available unchanged V4--V12 diagnostics also passed
9,467, 89,339, 906, 311, 609, 990, 608, 463 and 1,201 assertions, respectively.
The reports distinguish these rational and symbolic fixtures from a computation
of native Wilson integral enclosures. No such native enclosure, uniform
local-moment certificate, or external formalization was run. Analytic proofs,
including the complete-domain discretization error, are not certified by
finite examples.

The cumulative V13 preserves V12 byte for byte except for the title version
and insertion of this exact appendix before its final \verb|\end{document}|.
The package records hashes, predecessor recovery, compilation, rendered-page
review, the new diagnostic, and actual reruns of the available V4--V12
scripts. Unavailable historical proofs and companion programs are not
represented as recertified. The next scheduled companion checkpoint remains
V15; this targeted literature audit is recorded separately.

\begin{firewall}
V13 supplies an exact local auxiliary parent for the native history law, an
all-chaos Gaussian reference gap, a finite local conditional-angle upper
certificate and its weighted gap implication, and a directional resampling
energy and posterior-operator reduction. It rules out the naive one-cell pair-row test at weak coupling but does
not certify a native uniform gap by a different route or the local
fourth-moment input. Accordingly no new
unconditional unrestricted-volume posterior contraction, time-row decay,
physical transfer-gap comparison, or interacting continuum Yang--Mills mass
gap follows from this appendix. Its auxiliary clock is not physical time.
\end{firewall}
% END F16 V13 APPEND

% BEGIN F16 V14 APPEND -- 2026-09-17
% Insert immediately before the final \end{document} in cumulative V13.
% No predecessor mathematics is edited; only its title version becomes V14.

\clearpage
\section{V14: balanced conditional updates, compact cancellation, and a native repair event}
\label{f16v14:context}

V13 rejects the absolute pair-row criterion for its unmoving cells, even
at Gaussian order. Its recommendation was to test a different update
organization before attempting a native comparison. We do that test here.
Separating whole chord cubes from individual \emph{moving} temporal ports
removes the artificial Gaussian time coupling. For this partition a rational
weighted pair certificate succeeds, with margin $2/39$ on the entire Gaussian
function space. It is not a new proof of V13's Gaussian gap for its old
partition: the updates, and hence the generator, are different.

The native test then distinguishes a second obstruction. Actual compact
exterior ports can cancel the fields pinning a neighbouring pair. The resulting
pair conditional has a correlation at least $7/8$ for $\gamma\ge4096$, even
though its balanced Gaussian pair correlation is $1/6$. Thus changing
coordinates alone does not validate a worst-exterior perturbation argument.
This is not just another negative test: an explicit neighbouring update repairs
this cancellation family with a uniform, same-clock support-energy floor.
We retain that repair rather than freezing the very boundary variable that
provides it. An integrated defect identity specifies what remains to extend
the local repair to the complete generator.

\subsection{Dependency audit and the choice not to repeat a failed criterion}

The supplied V13 appendix and literature audit were read in full. Two semantic
archive searches returned no results; the mounted source inventories and
relevant exact passages were then examined directly. The audit is targeted,
not independent recertification of the archives. F15 V99's balanced edge
embedding and diagonal mixed Gram are inherited. F15 V86 already exhibits
zero spatial staples and the resulting near-unit one-link transfer ratio;
\emph{the general zero-staple mechanism is not new}. The calculation here is
for two temporal integration ports of the current fixed-bridge history,
including its periodic fixed roots, the revised resampling partition, and an
explicit repair port. F14's early repair-event argument is also credited:
our simpler repair bound uses an exact one-port conditional and no diffusion
comparison or changed clock.

Roberts--Sahu's
\href{https://doi.org/10.1111/1467-9868.00070}{\emph{Updating Schemes,
Correlation Structure, Blocking and Parameterization for the Gibbs Sampler}}
is the relevant classical context for changing a Gaussian update
parameterization. Its primary bibliographic record and abstract were checked;
no theorem from its full paper is imported. Qin--Wang's
\href{https://arxiv.org/abs/2208.11299}{\emph{Spectral Telescope: Convergence
Rate Bounds for Random-Scan Gibbs Samplers Based on a Hierarchical Structure}}
provides a general-state-space conditional-gap framework, not just a discrete
analogy. Its primary abstract was screened here, not its full proof. We do
not invoke it without its conditional inputs. All inequalities used below
are proved directly or attributed to the exact inherited statement.

The useful part of the supplied wider-literature memorandum is its insistence
on compatible local constraints. The native measure remains fixed. We neither
assert that different Gibbs generators have the same gap nor turn their
resampling clock into physical Euclidean time. The fixed-order corrections
of V10 are not recomputed.

\subsection{A different local generator on exactly the same native law}

Use the full law of V12--V13 at a prescribed bridge history with
$\max_{i,e}|y_{i,e}|\le r$, and let $\mathcal N=B(n+1)$, $B=m^3$, $m\ge4$,
$n\ge1$. Keep the chord groups $X_i$ and change the unmoving temporal groups
$h_t$ back to the moving groups $r_t$ by the \emph{exact} relation
\begin{equation}
 h_{t,v}=f_v(X_t)r_{t,v}f_v(X_{t+1})^{-1},\qquad r_{t,o}=I.
 \label{f16v14:coordinates}
\end{equation}
The $f_v$ are the same cube frames as in V2--V12. Write $P^{\rm bal}$ for
the pushforward of $P$ to $(X,r)$. There are two update species:
\[
 \mathcal I_x=\{(i,b):0\le i\le n\},\qquad
 \mathcal I_r=\{(t,b,v):0\le t<n,\ v\ne o\}.
\]
An $x$ update resamples all five chords of one cube at one time, with all
other $(X,r)$ variables fixed. An $r$ update resamples one three-dimensional
group port, with all other $(X,r)$ variables fixed. Each update has rate one.
For $a\in\mathcal I_x\sqcup\mathcal I_r$ set
\begin{equation}
 \begin{gathered}
 \widehat E_a f=\mathbb E_{P^{\rm bal}}[f\mid\text{all coordinates except }a],
 \qquad \widehat Q_a=I-\widehat E_a,\\
 \widehat{\mathcal L}=\sum_a\widehat Q_a,\qquad
 \widehat{\mathcal E}(f)=\sum_a\|\widehat Q_a f\|_{L^2(P^{\rm bal})}^2.
 \end{gathered}
 \label{f16v14:generator}
\end{equation}
The notation $r_t$ for group ports is distinct from the scalar bridge bound
$r$ and from V7's residual scores.

\begin{proposition}[Exact balanced updates and their local realization]
\label{f16v14:same-law}
The transformation \eqref{f16v14:coordinates} preserves conditional product
Haar measure on the ports and leaves all chord endpoint factors unchanged.
The law $P^{\rm bal}$ is therefore the same normalized native probability in
new coordinates, not a comparison law. The projections in
\eqref{f16v14:generator} have a finite-range square-root-density realization,
a common null vector, and constants as their joint kernel. Their finite-regulator
gap $\widehat\lambda_P$ is positive.

Pulled back to the unmoving coordinates, an $x_{i,b}$ update changes only
$x_{i,b}$ and the ports $h_{i-1,b},h_{i,b}$ that exist. A port update changes
only its corresponding $h_{t,b,v}$. No global update is hidden in this
parameterization.
\end{proposition}
\begin{proof}
At fixed $X$, each map $r_{t,v}\mapsto h_{t,v}$ is a left and right Haar
translation. Its inverse is explicit. Fubini therefore gives an exact
measure change, including the constant Haar factors. The chord weights
$j(x/\sqrt\gamma)^{w_i}$ and the two original Gaussian endpoint amplitudes
are unaltered. Principal cuts have Haar measure zero and do not restrict the
group translations.

Use the normalized product of the inherited one-cube chord reference factors
and the individual port Haar probabilities as reference measure. The native
density relative to this reference is proportional to the exponential of the
original classical action, expressed in $(X,r)$. It is bounded above and
below by positive finite constants at a finite regulator. The action words
still have bounded spatial and temporal range, since every frame uses one
cube only. V13's conditional square-root calculation
\eqref{f16v13:local-parent-kernel} now applies to the new coordinate blocks:
each conjugated projection depends only on its finite word neighbourhood and
kills the same square-root density. Product variance comparison proves a
positive, possibly very poor finite-regulator gap and the constant kernel,
exactly as in Proposition~\ref{f16v13:parent}.

Finally inspect \eqref{f16v14:coordinates}. Changing $X_{i,b}$ can change the
frame at that slice only, and hence just the two indicated bonds. In V13's
unmoving cell partition this touches at most two cells, at times $i-1,i$,
with unavailable endpoints omitted. Each unmoving cell is touched by only a
fixed number of these updates. This also proves finite support after pulling
back. The coordinate change is unitary between the two probability $L^2$
spaces, but the old and new lists of conditional subspaces are different;
no equality of their generators or gaps follows.
\end{proof}

\subsection{The Gaussian pair calculation now has a positive margin}

In principal moving logarithms $z$, the Gaussian precision is the separated
one from V4,
\[
                 W_n\ \oplus\ (I_n\otimes L_h).
\]
Use whole $x_{i,b}$ blocks and individual three-colour $z_{t,b,v}$ blocks,
precisely the linear limit of the new native update partition. Mean shifts
change no Gaussian conditional correlations. Put $N=\mathsf M^{-1}$ for one
cube, and let $\mathsf H$ be its six-face Gram. The diagonal chord blocks are
\begin{equation}
 K_{\rm int}=\mathsf H+4N,\qquad
 K_{\rm end}=\tfrac12\mathsf C^{-1}+\tfrac12\mathsf H+3N.
 \label{f16v14:diagonals}
\end{equation}
These include the actual mixed diagonal $2N$ and the dressed endpoints.
The temporal off-diagonal chord block is $-N$; spatial off-diagonal chord
blocks at the same time come from one crossed face only.

\begin{proposition}[A rational face bound with the endpoint cost retained]
\label{f16v14:face}
Let $U=\mathsf J+\mathsf B\mathsf F$ be the inherited balanced edge embedding,
and let $V_F$ restrict its rows to the four edges of any cube face. Then
\begin{equation}
 V_F^{\mathsf T}V_F\preceq\frac9{80}K_{\rm int},\qquad
 K_{\rm end}\succeq\frac{39}{40}K_{\rm int},\qquad K_{\rm int}\succeq8N.
 \label{f16v14:face-bounds}
\end{equation}
Consequently the Gaussian pair maximal correlations for the new blocks are
bounded by $3/26$ for spatial chord neighbours and by $5/39$ for temporal
chord neighbours. Chord--port correlations are zero. Neighbouring individual
ports on the same fine spatial bond graph have pair correlation exactly $1/6$;
other port pairs have zero pair correlation.
\end{proposition}
\begin{proof}
Here is an explicit exact certificate for the face constant. In the inherited
chord order $(a,b,c,d,e)$, the rational matrices are
\[
 N=\frac1{24}\begin{pmatrix}
10&-4&1&-5&5\\-4&10&-1&5&-5\\1&-1&10&-2&-4\\
-5&5&-2&10&-4\\5&-5&-4&-4&10
\end{pmatrix},\qquad
 \mathsf H=\begin{pmatrix}
2&-1&0&-1&1\\-1&2&0&1&-1\\0&0&2&0&-1\\
-1&1&0&2&-1\\1&-1&-1&-1&2
\end{pmatrix}.
\]
For the face with first coordinate zero, in the positive edge order,
\[
 V_F=\frac1{24}\begin{pmatrix}
 -4&-2&5&-1&1\\-1&1&-4&-4&-2\\
 -2&-4&-5&1&-1\\1&-1&10&-2&-4
 \end{pmatrix}.
\]
Direct rational multiplication gives
\begin{equation}
 \det\!\left(tI-K_{\rm int}^{-1}V_F^{\mathsf T}V_F\right)
       =t(t-1/32)^2(t-1/60)(t-9/80).
 \label{f16v14:face-polynomial}
\end{equation}
The product is similar to a positive symmetric matrix. Thus its largest
eigenvalue is $9/80$. Cube automorphisms preserve the cycle-space inner
product and the face-action form and permute the six faces, up to oriented
edge signs. Transporting the chord coordinates gives a simultaneous
congruence, so the same generalized eigenvalues hold on every face. The
supplied checker also reconstructs and verifies all six faces separately.

Whitening by $N$ makes $\mathsf H$ have eigenvalues $t=4,6$, with their
inherited multiplicities. The whitened endpoint block is the function
$\sqrt{t+t^2/4}+t/2+3$. At $t=4$, its value is $\sqrt8+5\ge39/5$,
and at $t=6$ it is $\sqrt{15}+6\ge39/4$. These follow by squaring
$\sqrt8\ge14/5$ and $\sqrt{15}\ge15/4$. They prove the second bound;
the third follows from $t\ge4$.

For any two Gaussian blocks of positive precision diagonals $K_1,K_2$ and
cross block $C_{12}$, their conditional maximal correlation is
$\|K_1^{-1/2}C_{12}K_2^{-1/2}\|$. To verify that this is an upper as well
as a lower bound, whiten and take a singular-value decomposition of the cross
block. The law becomes a product of correlated scalar Gaussian pairs. The
conditional-expectation singular values on Hermite polynomials are products
of powers of the scalar correlations. On the centered space the largest is
the largest scalar correlation; Hermite density completes the argument.
This is also the two-block instance of V13's all-chaos analysis.

Across one spatial interface the cross block is $-V_+^{\mathsf T}OV_-$,
where $O$ matches the four edge orientations and is orthogonal. The face
bound and the endpoint comparison give a correlation at most
$(40/39)(9/80)=3/26$. For the temporal block $-N$, both diagonal blocks
are at least $(39/5)N$, giving $5/39$. There are no Gaussian cross-species
blocks. Finally, for individual nonroot ports the diagonal block of $L_h$
is $6I_3$, and each fine neighbouring nonroot port contributes $-I_3$.
A two-port conditional therefore has scalar correlation $1/6$ in each
colour. Fixed root neighbours contribute to the diagonal, not additional
random coordinates. The asserted zeros follow from the exact precision
support, not from marginal independence of all chord coordinates.
\end{proof}

\begin{theorem}[A successful balanced Gaussian local certificate]
\label{f16v14:Gaussian-margin}
Give every chord-cube update weight one. For a port at relative vertex
$v\in\{0,1\}^3\setminus\{0\}$, give its update weight
\begin{equation}
 w(v)=\begin{cases}7/10,&|v|_1=1,\\9/10,&|v|_1=2,\\1,&|v|_1=3.\end{cases}
 \label{f16v14:weights}
\end{equation}
The Gaussian weighted pair-row margin for this partition is at least $2/39$.
In particular its weighted and unweighted rate-one-per-block generators have
gap at least $2/39$ on all of their Gaussian $L^2$ spaces, uniformly in
$B,n$ and the prescribed bridge history.
\end{theorem}
\begin{proof}
A chord block has at most six spatial and two temporal chord neighbours.
Its row deficit is at least
\[
        1-6(3/26)-2(5/39)=2/39.
\]
An endpoint has fewer temporal neighbours, so the same safe bound applies.
There are no cross-species contributions at Gaussian order.

A fine nonroot port of Hamming weight one has four nonroot neighbours of
weight two and two fixed roots. A port of weight two has four neighbours of
weight one and two of weight three. A port of weight three has six neighbours
of weight two. Internal and bridge edges both count; every port has fine
degree six. With pair coefficient $1/6$, the three weighted deficits are
\[
 \frac7{10}-\frac46\frac9{10}=\frac1{10},\quad
 \frac9{10}-\frac16\left(4\frac7{10}+2\right)=\frac1{10},\quad
 1-\frac66\frac9{10}=\frac1{10}.
\]
Apply the projection anticommutator proof of
Theorem~\ref{f16v13:angle-gap} to these new Gaussian blocks. It works on the
complete $L^2$ space. Its constant kernel follows from Gaussian equivalence
to product Lebesgue measure, or the inherited all-chaos proof. The unweighted
form dominates the weighted one since all weights are at most one. This
proves the claim without a claim about the \emph{native} pair angles.
\end{proof}

This calculation separates a parameterization artifact from a native
obstruction. It repairs V13's failed Gaussian preflight; it does not justify
uniform convergence of full conditional-expectation operators over arbitrary
compact exterior data. The following exact test shows that such convergence
is false here.

\subsection{A native two-port cancellation including the fixed roots}

Identify $\SU(2)$ with the unit quaternions, write $\Re U$ for the scalar
part, and retain $s(U)=1-\Re U$. At one temporal bond fix all chord groups to
$I$ at its two ends and take the bridge history to be zero. The native
port-dependent action at that bond is exactly
\begin{equation}
      \gamma\sum_{\{u,v\}\text{ fine spatial edge}}s(r_u^{-1}r_v),
                  \qquad r_o=I\text{ at every cube root}.
 \label{f16v14:spin-action}
\end{equation}
Magnetic terms and chord endpoint amplitudes have no port dependence. This
is a particular actual conditional of the present law, not a new spin-model
replacement for its general background.

For neighbouring nonroot ports $c,d$, let $\mathcal H_c^{\setminus d}$ be
the sum of its five exterior neighbour quaternions, excluding $d$, and
define $\mathcal H_d^{\setminus c}$ similarly. Adjacent vertices in the cubic
graph have disjoint exterior-neighbour sets. A nonroot vertex of parity
weight one has two root neighbours; a vertex of weight two or three has none.
These facts hold without aliasing for $m\ge4$.

\begin{proposition}[Explicit compact field cancellation]
\label{f16v14:cancellation}
For every neighbouring nonroot pair there are exterior port values, respecting
all fixed roots, such that
$\mathcal H_c^{\setminus d}=\mathcal H_d^{\setminus c}=0$.
At those values the pair conditional is exactly
\begin{equation}
 d\pi_\gamma(R,S)=z_\gamma^{-1}e^{-\gamma s(R^{-1}S)}\,dH(R)dH(S).
 \label{f16v14:free-pair}
\end{equation}
Both marginals are probability Haar. Its maximal correlation is at least
\begin{equation}
 \rho_\gamma=\int\Re V\,k_\gamma(V)\,dH(V)
                      \ge1-512/\gamma\qquad(\gamma\ge1).
 \label{f16v14:pair-lower}
\end{equation}
The last constant is conservative and is not a numerical integral enclosure.
\end{proposition}
\begin{proof}
Let
\[
 I=(1,0,0,0),\qquad
 T_+=(-1/2,1/2,1/2,1/2),\qquad
 T_-=(-1/2,-1/2,-1/2,-1/2).
\]
All are unit quaternions and $T_++T_-=-I$. If there are no root neighbours,
assign the five neighbours the values $I,-I,I,T_+,T_-$. Their sum is zero.
If there are two fixed root neighbours, leave them $I$ and assign the three
free neighbours $-I,T_+,T_-$. The sum is again zero. The two neighbour
sets are disjoint, so these assignments are compatible. All other free
ports can be set to $I$.

For a fixed neighbour $A$, $s(R^{-1}A)=1-\Re(R^{-1}A)$ is affine in $A$.
Each five-neighbour contribution is therefore the constant five. The only
remaining variable action is $\gamma s(R^{-1}S)$, proving
\eqref{f16v14:free-pair}, including its normalizer. Equivalently take $R$
Haar and $S=RV$ with independent $V$ of law $k_\gamma dH$.
Conjugation symmetry gives $\mathbb EV=\rho_\gamma I$. Thus the functions
$2\Re R$ and $2\Re S$ have zero mean, variance one, and covariance
$\rho_\gamma$. The variance follows from rotational invariance of probability
Haar on the unit three-sphere: each of its four real coordinates has second
moment $1/4$.

Here are explicit bounds for $1-\rho_\gamma=\mathbb E_{k_\gamma}s(V)$.
The normalized angular Haar density is $(2/\pi)\sin^2\theta\,d\theta$ on
$[0,\pi]$. For $\gamma\ge1$, use $\theta\le\gamma^{-1/2}$,
$\sin\theta\ge2\theta/\pi$ and $s\le\theta^2/2$ to obtain
\[
 z_\gamma\ge\frac{8e^{-1/2}}{3\pi^3}\gamma^{-3/2}.
\]
Use $s\ge2\theta^2/\pi^2$, $s\le\theta^2/2$ and
$\sin\theta\le\theta$ for the numerator. Its upper bound is
\[
 \frac1\pi\int_0^\infty\theta^4e^{-2\gamma\theta^2/\pi^2}\,d\theta
       =\frac{3\pi^{9/2}}{32\sqrt2}\gamma^{-5/2}.
\]
Their ratio is $9e^{1/2}\pi^{15/2}/(256\sqrt2\,\gamma)$.
Using $e^{1/2}<2$, $\sqrt\pi<2$, $\sqrt2>1$, and $\pi<22/7$ bounds
this coefficient by $36(22/7)^7/256<512$. This proves
\eqref{f16v14:pair-lower}. No Gaussian approximation of the pair law was used.
\end{proof}

\begin{theorem}[The balanced native worst-exterior row still fails]
\label{f16v14:native-obstruction}
For zero bridge history, $m\ge4$, $n\ge1$ and $\gamma\ge4096$, the native
worst-exterior correlations for the two port pairs on the chain
\[
              (1,0,0),\quad(1,1,0),\quad(1,1,1)
\]
at any one temporal bond are at least $7/8$. No positive update weights can
make the worst-exterior pair-row margin positive for the balanced partition.
This is true despite Theorem~\ref{f16v14:Gaussian-margin}.
\end{theorem}
\begin{proof}
Apply the preceding construction separately to the two pairs. An essential
supremum for each pair is allowed to use its own exterior values; simultaneous
maximization is not required in a row bound. At finite coupling the conditional
expectations of the bounded continuous test functions $2\Re R,2\Re S$, and
of their products, depend continuously on the local compact exterior values
near these configurations and on chord inputs near the identity. Their
variances are nonzero there. Each such exterior neighbourhood has positive
native marginal probability. Values $-I$ may be represented by a principal
chart cut, but nearby group neighbourhoods minus the cut still have positive
Haar measure. Hence the essential supremum is at least the value in
\eqref{f16v14:pair-lower}, not just a value at a null conditioning point.

At $\gamma\ge4096$ this is at least $7/8$. Positive row margins would require
$w_1>(7/8)w_2$, $w_3>(7/8)w_2$, and
$w_2>(7/8)(w_1+w_3)$ after dropping all other nonnegative terms. They imply
$w_2>(49/32)w_2$, a contradiction. The calculation involves one existing
port bond only, so it also applies to a history of length one. It does not
bound the gap of the full resampling generator from above.
\end{proof}

The general slow zero-staple channel was already known in f15 V86. The new
conclusion is that it survives the present balanced \emph{two-port}
conditional even with the periodic rooted pins, and defeats this revised
worst-exterior certificate on the shortest history. A zero-field Laplace test
cannot upper-bound those native worst-exterior correlations.

\subsection{The actual cancellation witness has a uniformly available repair}

We now let one of the exterior variables move. This is an estimate for the
native generator, not for its Gaussian approximation. For general fixed
$X,\mathbf y$ and other ports, the conditional of one moving port $r_p$ has
the exact form
\begin{equation}
 d\pi_p(R\mid\text{rest})
       =\frac{e^{\gamma\Re(R^{-1}\mathcal H_p)}}{Z_p(\mathcal H_p)}dH(R),
                         \qquad |\mathcal H_p|\le6.
 \label{f16v14:one-port}
\end{equation}
Each of its six neighbouring terms supplies one unit-quaternion summand of
$\mathcal H_p$. For an oriented fine edge $e=(p,q)$ that summand is
$\bar U_e(i)r_q\bar U_e(i+1)^{-1}$; for $e=(q,p)$ it is
$\bar U_e(i)^{-1}r_q\bar U_e(i+1)$. Here $\bar U$ is the actual moving
fine representative, whose bridges are $Z$ and whose internal edges are
$\bar X$. These formulas follow by cyclicity and inversion of the real trace.
They retain all noncommuting factors. Neither magnetic weights nor chord
amplitudes depend on $r_p$.

Let
\[
 C_-=\{R:|R+I|\le1/24\},\qquad
 C_+=\{R:|R-I|\le1/24\},\qquad
 \mathcal R_p=\{\Re\mathcal H_p\ge1,\ r_p\in C_-\}.
\]
The field condition is measurable in the complement of $p$. These are
quaternion chord distances, not rescaled logarithmic distances.

\begin{theorem}[A same-law, rate-one repair floor on a compact witness family]
\label{f16v14:repair}
For every fixed bridge history and every $\gamma\ge1$,
\begin{equation}
 \pi_p(C_-\mid\text{rest})\le e^{-3\gamma/2}
                         \quad\text{when }\Re\mathcal H_p\ge1.
 \label{f16v14:escape}
\end{equation}
For any $f\in L^2(P^{\rm bal})$ supported on $\mathcal R_p$,
\begin{equation}
 \widehat{\mathcal E}(f)\ge
  \|\widehat Q_p f\|_2^2
   \ge(1-e^{-3\gamma/2})\|f\|_2^2\ge\tfrac12\|f\|_2^2.
 \label{f16v14:repair-floor}
\end{equation}
For the weights \eqref{f16v14:weights}, the weighted form has the analogous
lower bound $7\|f\|_2^2/20$.

This repair event contains a neighbourhood of an explicit two-port
zero-staple cancellation family, uniformly over all values of its two free
pair ports. No probability lower bound on that neighbourhood is required.
\end{theorem}
\begin{proof}
When $\Re\mathcal H_p\ge1$, the one-port exponent without $\gamma$ is at
most $-1+6/24=-3/4$ on $C_-$ and at least $1-6/24=3/4$ on $C_+$.
These two caps have the same positive Haar volume, by $R\mapsto-R$.
Lower-bound $Z_p$ by its integral on $C_+$ and upper-bound the numerator on
$C_-$. Their ratio proves \eqref{f16v14:escape}, without estimating the
normalizer by a different conditional law.

At a fixed exterior where $f$ can be nonzero, conditional Cauchy--Schwarz
and its support give
\[
  |\mathbb E_p f|^2\le\pi_p(C_-)\,\mathbb E_p|f|^2,
  \qquad \operatorname{Var}_p f\ge(1-e^{-3\gamma/2})\mathbb E_p|f|^2.
\]
Integrate in the actual exterior marginal. One term of the generator proves
\eqref{f16v14:repair-floor}. Since $e^{-3/2}<1/2$, the stated constants
follow, with $w_p\ge7/10$ for the weighted form.

To verify the final statement on the literal lattice, choose the pair
$c=(1,1,0)$, $d=(1,1,1)$. Assign each of their five exterior neighbours,
in increasing coordinate order, the values $I,-I,I,T_+,T_-$. On a torus of
side $2m$ this means, at the two relevant rows,
\[
\begin{array}{c|ccccc}
 & (0,1,0)&(1,0,0)&(1,1,2m-1)&(1,2,0)&(2,1,0)\\
\text{value}&I&-I&I&T_+&T_-\\ \hline
 &(0,1,1)&(1,0,1)&(1,1,2)&(1,2,1)&(2,1,1)\\
\text{value}&I&-I&I&T_+&T_-
\end{array}
\]
Leave all remaining ports at $I$ and all chord groups at $I$, with zero
bridges. The two pair exterior fields are zero as before. Take the repair
port $p=(1,0,0)$. Its five neighbours other than $c$ have sum $3I$:
$(0,0,0),(1,0,1),(1,0,2m-1),(1,2m-1,0),(2,0,0)$ carry
$I,-I,I,I,I$. Thus
\[
                 r_p=-I,\qquad\mathcal H_p=3I+r_c,
                 \qquad\Re\mathcal H_p\ge2
\]
for \emph{every} $r_c,r_d\in\SU(2)$. The roots in this list are already
fixed $I$, consistently with the prescription. Compactness of the two free
pair groups and continuity of the finite word list give a uniform exterior
and chord neighbourhood with $\Re\mathcal H_p>1$; restrict $r_p$ to $C_-$.
This proves that the entire pair-fibre witness has a repair neighbourhood.
The neighbourhood may be chosen in angular group coordinates. It need not
lie in a bounded thermal window, and no such restriction was used in the
one-port estimate.
\end{proof}

The support qualification in \eqref{f16v14:repair-floor} is essential. It is
not the false inequality $\|1_{\mathcal R_p}f\|_2^2\le C\widehat{\mathcal E}(f)$
for arbitrary $f$, which constants would contradict. A global argument must
pay localization commutators or use a compatible block comparison. What is
proved is that the explicit slow \emph{frozen pair} does not itself form a
low-energy supported state for the full resampling dynamics: the boundary
port it froze is a fast escape direction.

\subsection{Low exterior fields have a paid density cost, not a gap consequence}

For a neighbouring free pair $c,d$ at one port bond, define the two exterior
fields by omitting their mutual edge from \eqref{f16v14:one-port}. They are
sums of five transported unit quaternions even at general bounded bridge
history. Let
\[
 \mathcal B_{cd}=\{|\mathcal H_c^{\setminus d}|\le1\}
                     \cup\{|\mathcal H_d^{\setminus c}|\le1\}.
\]
Let $N_{\rm low}$ count these events over all neighbouring nonroot port
pairs and all temporal bonds, with each unoriented pair counted once.

\begin{proposition}[Actual-law low-field pair density]
\label{f16v14:density}
For sufficiently large $\gamma$ depending on the fixed bridge window,
\begin{equation}
                \frac{\mathbb E_P N_{\rm low}}{B(n+1)}
                                      \le C_r/\gamma.
 \label{f16v14:density-bound}
\end{equation}
This is an expectation under the complete native law, with every exterior
field integrated. It is not a bound on the probability of all-time good
fields, on a specified pair uniformly in location, or on source-weighted
Dirichlet defects.
\end{proposition}
\begin{proof}
If a sum of five unit quaternions has norm at most one, its distance from
$5I$ is at least four. By telescoping products in its five summands,
\[
 4\le\sum_{q\ne d}|I-A_{cq}|
      \le C\gamma^{-1/2}
           \left(r+\sum_{j\in S_{cd}}(|x_j|+|z_j|)\right),
\]
where $S_{cd}$ is a fixed-size neighbouring carrier in the two internal
slices and the one port bond. The internal moving edges are products of
exponentials of fixed cube-linear chord combinations, so
$|I-\bar X_e|\le C\gamma^{-1/2}\sum|x_j|$ globally on their principal
chord inputs. Also $|I-r_p|\le|z_p|/\sqrt\gamma$, and bridge distances are
at most $r/\sqrt\gamma$. Thus, after increasing the fixed-$r$ threshold,
$\mathcal B_{cd}$ implies
$\sum_{j\in S_{cd}}(|x_j|^2+|z_j|^2)\ge c\gamma$.

Every coordinate belongs to a bounded number of these carriers. The exact
relation $r_t=f_t^{-1}h_tf_{t+1}$ and the bi-invariant geodesic triangle
inequality give the inherited global bound
$\sum(|x|^2+|z|^2)\le C\sum(|x|^2+|a|^2)$, as in V5's bad-label proof.
V5's \emph{actual} second-moment density then bounds its expectation by
$C_r\mathcal N$. Sum the deterministic local implication and divide by
$\gamma\mathcal N$. This proves \eqref{f16v14:density-bound}. We reuse the
moment theorem; no new pointwise moment or independence assertion is made.
\end{proof}

Even this vanishing density does not suppress a worst-exterior correlation:
a small open event of positive probability already contributes to an
essential supremum. Nor may its probability be multiplied by a uniform bound
on arbitrary normalized test functions. The repair theorem controls a
supported function by its \emph{own} resampling energy instead.

\subsection{The directional source reduction survives the changed updates}

The old covariance matrices, posterior law, and Haar-compensated sources
remain those of V7 and V12. They are simply composed with
\eqref{f16v14:coordinates} when evaluated under $P^{\rm bal}$. Define the
native moments $m_2(P),m_4(P)$ in the original unmoving coordinates exactly
as in \eqref{f16v13:moment-inputs}. Uniform bounds on these moments are not
assumed established.

\begin{proposition}[Balanced-update innovation energy and posterior handoff]
\label{f16v14:energy}
For the primitive remainder bank $e=(e^{\rm br},e^{\rm pt})$ and the full
corrected source vector $\mathbf r$ of V7,
\begin{equation}
 \begin{aligned}
 \widehat{\mathcal E}(b^{\mathsf T}e)
       &\le C_r\frac{m_4(P)}\gamma\|b\|^2,\qquad
 \widehat{\mathcal E}(v^{\mathsf T}\mathbf r)
       \le C_r\frac{m_4(P)}\gamma\|v\|^2,\\
 \widehat{\mathcal E}(b^{\mathsf T}S_0)
       &\le C m_2(P)\|b\|^2.
 \end{aligned}
 \label{f16v14:energy-bound}
\end{equation}
Consequently V13's three directional bounds
\eqref{f16v13:directional-reduction} hold with $\widehat\lambda_P$ in place
of $\lambda_P$, and an enlarged geometric constant. They concern exactly
the same $\Gamma_\gamma,\mathsf R_{\rm post},\mathsf L_D$ and full memory.
No comparison between $\widehat\lambda_P$ and $\lambda_P$ is asserted.
\end{proposition}
\begin{proof}
Use stationary exact resampling in the balanced variables. Both the original
configuration and its updated configuration, pulled back to $(x,a)$, have
law $P$. By Proposition~\ref{f16v14:same-law}, each update changes only a fixed
number of unmoving cells, and each cell is changed by a fixed number of
update types. Each primitive remainder depends on a fixed finite carrier.
Thus every update changes a bounded number of primitive components, and
each component can change under only a bounded number of updates.

For the changed components use Cauchy--Schwarz in that fixed set and
stationarity to bound
$\mathbb E|e_j(\zeta)-e_j(\zeta')|^2\le4\mathbb E_P e_j^2$.
The complete-chart bound \eqref{f16v13:primitive} makes the latter at most
$C_rm_4(P)/\gamma$. The resampling identity
$\widehat{\mathcal E}(f)=\frac12\sum_a\mathbb E|f-f^{(a)}|^2$ now gives
the first estimate without a volume factor. The fixed bounded synthesis
$b=(v,\mathsf P v)$ gives the second. For an affine port score write
$b^{\mathsf T}S_0=l^{\mathsf T}\zeta+$constant with $\|l\|\le h_+\|b\|$.
Only finitely many cell terms change at one update. Apply the same argument
to those terms using $m_2(P)$; the bounded update incidence proves the third
estimate. Principal-logarithm jumps are allowed: stationarity, not a
small-displacement approximation, bounds them.

Apply the native Poincare inequality of the new finite-regulator generator
to these energies. V13's posterior proof then uses only the unchanged exact
lift-mismatch identities, one common exterior conditional expectation, and
its $L^2(P)$ contraction. Repeating those steps, not replacing that posterior,
gives the asserted handoff. Its use still needs a uniform native gap and
moments to obtain a useful uniform conclusion.
\end{proof}

\subsection{An integrated defect criterion that does not freeze the repair variables}

Let $w_a$ be \eqref{f16v14:weights} and set
$\widehat{\mathcal L}_w=\sum_a w_a\widehat Q_a$,
$\widehat{\mathcal E}_w(f)=\langle f,\widehat{\mathcal L}_w f\rangle$.
Define a deterministic symmetric table $\Theta_{ab}$ by the Gaussian
majorants $3/26,5/39,1/6$ above, and zero elsewhere. It has weighted row
margin $\delta_0=2/39$. The \emph{native} interaction graph can have additional
pairs, including cross-species pairs; retain them, with $\Theta_{ab}=0$.
For every pair in that finite-range graph let $\theta_{ab}(z)$ be its actual
conditional maximal correlation at exterior $z$, and define
\[
 d_{ab}(z)=(\theta_{ab}(z)-\Theta_{ab})_+.
\]
These nonnegative functions are measurable in the pair exterior. At a finite
regulator the normalized conditional kernels are bounded operators; they
provide a measurable version of the correlations. Put
\begin{equation}
 \mathfrak D(f)=\sum_{a<b}w_aw_b\,
 \mathbb E_P\left[d_{ab}(\omega_{(a,b)^c})
       \bigl(|\widehat Q_af|^2+|\widehat Q_bf|^2\bigr)\right].
 \label{f16v14:defect}
\end{equation}
Here $\omega=(X,r)$ denotes the balanced coordinates. The notation $P$ identifies
the same full law, not its Gaussian reference. Noninteracting pairs have no
defect and can be omitted.

\begin{theorem}[Same-law integrated defect identity and sufficient repair target]
\label{f16v14:integrated}
At each finite regulator, on the complete native $L^2$ space,
\begin{equation}
 \|\widehat{\mathcal L}_w f\|_2^2
       \ge\delta_0\widehat{\mathcal E}_w(f)-\mathfrak D(f).
 \label{f16v14:defect-identity}
\end{equation}
If, independently, constants $0\le u<\delta_0$ and $v\ge0$ are proved to obey
\begin{equation}
 \mathfrak D(f)\le u\widehat{\mathcal E}_w(f)
                         +v\|\widehat{\mathcal L}_w f\|_2^2
                              \quad\text{for every }f,
 \label{f16v14:repair-target}
\end{equation}
then
\begin{equation}
             \widehat\lambda_P\ge\frac{\delta_0-u}{1+v}.
 \label{f16v14:repair-gap}
\end{equation}
Neither \eqref{f16v14:escape} nor \eqref{f16v14:density-bound} alone proves
\eqref{f16v14:repair-target}; it is the still-unverified global input.
\end{theorem}
\begin{proof}
Condition on the exterior of a pair. V13's two-projection argument bounds
its anticommutator below by
$-\theta_{ab}(z)(\widehat Q_a+\widehat Q_b)$ on that fibre. Multiply by
$w_aw_b$ and integrate before taking any supremum. An exterior-measurable
coefficient commutes with both pair updates, so its quadratic form is its
expectation times the corresponding squared innovation, exactly as in
\eqref{f16v14:defect}. Use $\theta_{ab}\le\Theta_{ab}+d_{ab}$, expand
$\widehat{\mathcal L}_w^2$, and sum its deterministic part. Its coefficient
on $w_a\widehat Q_a$ is at least the verified margin $\delta_0$. This proves
\eqref{f16v14:defect-identity} on all $L^2$; all operators and coefficients
are bounded at finite regulator.

Combining with \eqref{f16v14:repair-target} gives
$(1+v)\widehat{\mathcal L}_w^2\succeq
(\delta_0-u)\widehat{\mathcal L}_w$.
The joint kernel is constants. Spectral calculus yields a gap at least
$(\delta_0-u)/(1+v)$, and the unweighted generator dominates the weighted
one. This proves \eqref{f16v14:repair-gap} without assuming the desired
posterior covariance. The deficit has not been multiplied by the probability
of a bad event independently of its test function.
\end{proof}

This criterion keeps the innovation where a conditional pair is bad and
allows the energy of the \emph{other} updates to pay for it. The explicit
repair port shows why that distinction is necessary. A complete proof would
need to classify or cover the relevant defect configurations, pay the
commutators introduced by such a localization, and control the resulting
constants in \eqref{f16v14:repair-target}. We have not shown that low staple
magnitude exhausts the defects: chord--port interactions and other native
conditional landscapes remain in \eqref{f16v14:defect}.

\subsection{Status and the next noncircular calculation}

There are three distinct conclusions. The old Gaussian row obstruction is
removed by exact local balanced updates, with an explicit successful margin.
The native worst-boundary row is nevertheless disproved by a compact
cancellation compatible with the rooted geometry. That cancellation family
has a native, rate-one repair event with a uniform support-energy floor.
These facts identify which variable a useful local block must retain:
freezing the cancellation's repair port produces a slow pair that the full
update leaves with the stated uniformly positive probability.

The next test should therefore include the boundary-repair variables in a
local block or in an integrated defect estimate, rather than improve the
same impossible worst-pair bound. The defect form and the adapted directional
energy estimate supply exact objects for that calculation. This is not a
claim that a finite block theorem has already been applied successfully.
Uniform native fourth moments, a global repair/auxiliary-gap bound, temporal
weighted-row decay of the posterior, and a leading-coefficient estimate at
unrestricted volume are still unproved. The existing physical bridge-window,
same-top transfer, infrared coercivity, and continuum obligations are
unchanged. No supported repair floor is identified with a physical mass.

\subsection{Diagnostics and exact predecessor preservation}

The accompanying script is
\begin{center}\small\path{verify_ym_f16_v14_balanced_repair.py}.\end{center}
It reconstructs the inherited cube matrices; checks the six rational face
characteristic polynomials, endpoint comparisons and weighted rows; checks
literal fine-torus roots and compatible quaternion cancellation assignments;
and verifies the repair-port sums, cap constants, and the integrated defect
algebra on finite positive laws. The report distinguishes exact diagnostics
from exploratory floating-point calculations. The latter are not proof
certificates and are not used for a theorem constant. Native continuum
integrals are bounded analytically above, not by an unreported quadrature.

The new diagnostic passed 622 exact assertions. Of these, 446 check
cancellation and compatibility with fixed roots on two finite tori. The available unchanged
V4--V13 diagnostics passed 9,467, 89,339, 906, 311, 609, 990, 608, 463,
1,201 and 222 assertions, respectively. The assertion families and script
hashes are recorded in the accompanying reports. These finite fixtures do
not formally verify the analytic proofs or establish a global native gap.

The cumulative source changes only V13's title version and inserts this
appendix before the final document-ending command. Removing the exact append
and restoring that title recovers the predecessor byte for byte. The package
records source hashes, preservation checks, compilation and visual review.
No inaccessible archive checker, external formalization or complete inherited
proof collection is described as recertified. The scheduled companion
checkpoint remains V15.

\begin{firewall}
V14 proves a successful balanced Gaussian pair certificate, an exact native
compact-boundary obstruction to that worst-exterior test, and a uniform
native support-energy repair floor on an explicitly identified witness
family. It also proves a low-field pair-density bound, adapted directional
source energies, and an integrated defect reduction with its global repair
input left explicit. The native law and all original normalizers are
unchanged. No uniform full-native auxiliary gap, unrestricted-volume coherent
posterior contraction, native local fourth-moment bound, physical transfer
gap, or interacting continuum Yang--Mills mass gap is established here.
\end{firewall}
% END F16 V14 APPEND
% BEGIN F16 V15 APPEND -- 2026-09-17
% Insert immediately before the final document-ending command in V14.
% No predecessor proof or bibliography is silently revised.

\clearpage
\section{V15: a complete three-port repair and a same-clock conditional gap}
\label{f16v15:context}

V14 identifies an escape port for a slow frozen pair, but its repair energy
floor applies only to functions supported on a specified event. Localizing an
arbitrary innovation to that event introduces an unpaid term. We avoid that
localization: integrate the escape port together with the pair and prove a
Poincar\'e inequality for \emph{every} function on this three-group conditional.
The final inequality uses the original rate-one single-port updates. A
collective update appears only in an explicitly paid intermediate comparison.

The exterior hypothesis is a fixed open angular neighbourhood of the actual
cancellation witness, not an all-exterior assertion. Within it, all three
ports range over their complete groups, including antipodal configurations.
The neighbourhood does not shrink with coupling, spatial volume, or duration.
Its radius and entry coupling are proved to exist but are not numerically
enclosed. This is a local conditional gap, not a global native or physical gap.

\subsection{The V15 checkpoint and the non-duplication boundary}

The V14 appendix and analytic audit were read in full. The supplied f14 V100,
f15 V100, Grand Tensor Monograph V076, and this lineage through V14 were
inventoried and relevant dependencies checked. Two semantic archive searches
were empty; this was not treated as evidence of absence. In particular, f14
V2 already proves same-clock \emph{supported} repair mechanisms. F15 V71 proves
a thermal conditional \emph{diffusion} gap on its own good-boundary law. Neither
is the three-moving-port resampling result below. The monograph's physical
Perron and vacuum-corrected Feshbach interfaces remain separate. This is a
targeted audit, not independent recertification of the entire archives.

The active obstruction is V14's integrated conditional-angle defect, not V3's
Gaussian inverse or V10's formal coefficients. The wider-literature memorandum's
compatible-local-constraint proposal is applied here to the exact auxiliary
operator. The other memorandum's distinction between order gain and locality
cost is retained: this iteration enlarges the class of test functions and
retains a previously frozen variable; it does not append a formal order to the
same uncontrolled error.

For context, Qin--Wang, \href{https://arxiv.org/abs/2208.11299}{\emph{Spectral
Telescope: Convergence Rate Bounds for Random-Scan Gibbs Samplers Based on a
Hierarchical Structure}}, Section 3.2 and Theorem 2, develops hierarchical
conditional-gap comparisons on general spaces. Its definitions and theorem
were checked, including the theorem page. The two-level comparison needed
below is proved directly with our clock, not quoted with a discrete random-scan
normalization. Liu--Wong--Kong,
\href{https://doi.org/10.1093/biomet/81.1.27}{\emph{Covariance structure of the
Gibbs sampler with applications to the comparisons of estimators and
augmentation schemes}} (1994), provides classical grouping/collapsing context;
only its primary summary was checked. No literature theorem identifies these
sampling gaps with a physical transfer gap.

Only supplied companions were available at this scheduled checkpoint. Other
programs named in older ledgers were not supplied or recertified. The next
five-version checkpoint is V20.

\subsection{The exact native three-port conditional}

Keep V14's balanced probability $P^{\rm bal}$. At one temporal bond use the
three nonroot fine vertices
\[
 c=(1,1,0),\qquad d=(1,1,1),\qquad p=(1,0,0).
\]
Their induced graph has exactly the edges $cd,cp$ for $m\ge4$. Condition every
other port and all chord groups. Each native port-edge interaction is bilinear
in its two quaternions, with fixed left/right unit factors. Since these two
edges form a tree, separate fixed left/right translations of the two leaves
make both couplings ordinary quaternion inner products. Keep the central port
unchanged and denote the resulting variables by $U,V,W$, respectively. These
changes preserve probability Haar and the three individual conditional update
operations; they are independent of the three free variables.

Write $U\cdot V=\Re(U^{-1}V)$ and $I=(1,0,0,0)$. The exact conditional is
\begin{equation}
 d\pi_\gamma^{A,B,D}(U,V,W)
 =\frac{e^{\gamma[U\cdot V+U\cdot W+A\cdot U+B\cdot W+D\cdot V]}}
 {\mathcal Z_\gamma(A,B,D)}\,dH(U)dH(V)dH(W),
 \label{f16v15:triplet}
\end{equation}
for exterior-dependent real quaternion fields with $|A|\le4$, $|B|,|D|\le5$.
The letter $D$ here is a four-vector, not V11's spatial cylinder. Each root
remains fixed. Exterior-only constants cancel on conditional normalization,
not in the exterior posterior. Magnetic and dressed-end factors have no
free-port dependence and are not altered. All kinetic terms meeting these
ports are included.

\begin{proposition}[Retaining the actual repair variable]
\label{f16v15:normal-form}
At V14's displayed witness, with only $c,d,p$ now left free,
\begin{equation}
                         A=I,\qquad B=3I,\qquad D=0.
 \label{f16v15:witness-fields}
\end{equation}
Freezing $W=-I$ recovers the unpinned two-port law proportional to
$e^{\gamma U\cdot V}dH(U)dH(V)$. The full three-port law instead has the unique
energy minimum $(I,I,I)$. The aligned fields depend continuously on the finite
exterior word data near the witness.
\end{proposition}
\begin{proof}
V14 writes every oriented edge contribution as the dot product of one port
with a left/right translate of its neighbour. Absorb the two such translates
on the leaves. A tree has no residual loop constraint. The central vertex has
four fixed neighbours, each leaf five, giving the stated field bounds.

At the witness all local chord frames and bridges are identities. V14's five
neighbours of $c$, excluding $d$, sum to zero and include $p=-I$. Removing $p$
leaves field $I$. The five exterior neighbours of $d$ sum to zero; the five
neighbours of $p$ other than $c$ sum to $3I$. This proves
\eqref{f16v15:witness-fields}. Setting $W=-I$ cancels $U\cdot W+I\cdot U$.
With all three variables free, the deficit from the maximum $6\gamma$ is
\[
 \gamma\{s(U^{-1}V)+s(U^{-1}W)+s(U)+3s(W)\}.
\]
It is nonnegative and vanishes exactly at $(I,I,I)$. Continuity follows from
the finite quaternion products. Group neighbourhoods at $-I$ remain legitimate
even though one principal logarithm has a Haar-null cut there.
\end{proof}

The following results hold on a fixed parameter set
\begin{equation}
 \mathcal T_\eta=\{(A,B,D):|A-I|\le\eta,\ |B-3I|\le\eta,\ |D|\le\eta\}.
 \label{f16v15:parameter-window}
\end{equation}
We will choose $0<\eta_0\le1/100$ independently of $\gamma,B,n$. The preimage of
its strict interior is a genuine open exterior neighbourhood of the witness.
No condition is imposed on the three free groups. Not every low-field or
large-correlation configuration is claimed to belong to this neighbourhood.

\subsection{The marginal denominators require their own compact estimate}

For a positive unnormalized pair density $f$ on $S^3\times S^3$, set
\begin{equation}
 m_x(x)=\int f(x,y)dH(y),\quad m_y(y)=\int f(x,y)dH(x),\quad
 \mathcal C(x,y)=\frac{f(x,y)}{\sqrt{m_x(x)m_y(y)}}.
 \label{f16v15:normalized-kernel}
\end{equation}
The total partition function cancels \emph{exactly}. Between the fixed Haar
$L^2$ spaces this is the normalized conditional-expectation kernel; its top
singular value is one with singular vectors the square roots of the normalized
marginals. Its second singular value is the pair maximal correlation. This
normalization is inherited from V13, not a new covariance convention.

\begin{proposition}[Uniform complete-chart normalized-kernel limit]
\label{f16v15:Laplace}
Let $f_{\gamma,\vartheta}=a_{\gamma,\vartheta}e^{\gamma F_\vartheta}$ on
$S^3\times S^3$, with $\vartheta$ in a compact parameter space. Assume the
family $F_\vartheta$ is continuous in the $C^3$ topology, with uniformly bounded
higher derivatives whenever they are used, and that the positive amplitudes
are uniformly bounded above and below and converge uniformly in $C^2$ to
positive $a_{0,\vartheta}$. Define
\[
 M_\vartheta(x)=\max_y F_\vartheta(x,y),\quad
 N_\vartheta(y)=\max_x F_\vartheta(x,y),\quad
 \mathcal D_\vartheta(x,y)=\tfrac12(M_\vartheta(x)+N_\vartheta(y))-F_\vartheta(x,y).
\]
Suppose $\mathcal D_\vartheta$ has exactly one zero
$(x_\vartheta,y_\vartheta)$, and the negative joint Hessian of $F_\vartheta$
there is uniformly positive. Then the kernels \eqref{f16v15:normalized-kernel},
in complete normal charts centered there and scaled by $\sqrt\gamma$, with
half-Haar Jacobians and zero extension, converge uniformly in
Hilbert--Schmidt norm to the normalized joint-Gaussian kernels with that
Hessian as precision. Their maximal correlations converge uniformly to the
corresponding Gaussian correlations.
\end{proposition}
\begin{proof}
The two terms $M-F,N-F$ are nonnegative. Their simultaneous zero means that
both coordinates are conditional maximizers. Every global maximizer is such
a zero, hence is the specified unique point. The points vary continuously:
compactness and continuity identify every subsequential limit of maximizers.
Use continuous tangent frames on finitely many local parameter neighbourhoods.

Write the negative joint Hessian in these frames as
$H=\left(\begin{smallmatrix}A_0&C_0\\C_0^{\mathsf T}&B_0\end{smallmatrix}\right)$.
At the joint point each individual conditional maximum is unique: any other
partial maximum with the other coordinate fixed would also be a global
maximum. Compactness excludes distant competing partial maxima in a small
neighbourhood. The implicit function theorem then describes the local partial
maxima using the positive diagonal Hessian blocks. Their quadratic deficits give
\begin{equation}
 \mathcal D(u,v)=\frac14\left(
 |u+A_0^{-1}C_0v|_{A_0}^2+
 |v+B_0^{-1}C_0^{\mathsf T}u|_{B_0}^2\right)+O(|(u,v)|^3).
 \label{f16v15:deficit-Hessian}
\end{equation}
Both squares zero imply $H(u,v)=0$, so this quadratic form is strictly
positive. Compactness makes its lower bound uniform. Away from a small
neighbourhood, continuity of the maxima and the unique-zero condition give
a uniform positive bound. Thus, in quaternion chord distances,
\begin{equation}
 \mathcal D_\vartheta(x,y)\ge c(|x-x_\vartheta|^2+|y-y_\vartheta|^2).
 \label{f16v15:deficit-pin}
\end{equation}
This argument does not require smoothness of $M,N$ away from their joint point.

A Hessian bound about \emph{any} partial maximizer, together with the positive
amplitude lower bound and integration on a ball of radius $\gamma^{-1/2}$,
gives the global marginal lower bounds
\[
 m_x(x)\ge c\gamma^{-3/2}e^{\gamma M_\vartheta(x)},\qquad
 m_y(y)\ge c\gamma^{-3/2}e^{\gamma N_\vartheta(y)}.
\]
They remain true for degenerate or multiple partial maxima away from the joint
point. Consequently
\[
 |\mathcal C_{\gamma,\vartheta}(x,y)|
       \le C\gamma^{3/2}e^{-\gamma\mathcal D_\vartheta(x,y)}.
\]
A complete normal chart on $S^3$ centered at a unit quaternion has
$dH=(2\pi^2)^{-1}\gamma^{-3/2}j(u/\sqrt\gamma)du$ on
$|u|<\pi\sqrt\gamma$. Both square-root factors in the kernel cancel
$\gamma^{3/2}$, and $j\le1$. Chord distance satisfies
$|x-x_\vartheta|\ge2|u|/(\pi\sqrt\gamma)$ on this entire chart. Thus the
zero-extended transformed kernel has the uniform majorant
\begin{equation}
                         C e^{-c'(|u|^2+|v|^2)}.
 \label{f16v15:normalized-majorant}
\end{equation}
In particular neither small marginal densities nor compact chart tails have
been discarded.

On every fixed rescaled ball, ordinary Taylor expansion about the moving
maximizer and Laplace expansion of the two marginals give the joint Gaussian
and its \emph{own} marginal normalizers. For these marginal expansions, the
partial Hessians are uniformly positive nearby; outside a fixed neighbourhood
the compact conditional phase has a uniform loss. These observations justify
uniform Laplace convergence there. The amplitude tends to its positive value
at the joint point and cancels from the normalized limit. The Gaussian
half-Jacobian normalization is also retained. Uniform convergence on balls
and the squared majorant \eqref{f16v15:normalized-majorant} prove the uniform
Hilbert--Schmidt limit, including the full compact complement. Singular-value
stability yields the last assertion. The Gaussian second singular value is
$\|A_0^{-1/2}C_0B_0^{-1/2}\|$: a singular-value decomposition of the whitened
precision and the Hermite basis give this value on all of centered $L^2$.
\end{proof}

The normalized deficit $\mathcal D$ is essential. A pin for $\max F-F$ alone
would not exclude another simultaneous conditional maximum carrying a large
normalized correlation in a small marginal region. We verify the stronger
condition for each of the following two reductions.

\subsection{The inner conditional remains controlled for every value of the third port}

Condition $V$ in \eqref{f16v15:triplet}, and put $C=A+V$. The free pair is
\begin{equation}
 d\pi_\gamma(U,W\mid V)\ \propto\
                e^{\gamma[U\cdot W+C\cdot U+B\cdot W]}dH(U)dH(W).
 \label{f16v15:inner-pair}
\end{equation}
The term $D\cdot V$ cancels here. In particular the worst old value $V=-I$
is not omitted.

\begin{proposition}[Uniform inner pair correlation]
\label{f16v15:inner}
For $|A-I|,|B-3I|\le1/100$ there is a finite $\gamma_1$, independent of
$A,B,V$, such that for all $V\in S^3$ and $\gamma\ge\gamma_1$ the maximal
correlation of \eqref{f16v15:inner-pair} is at most $3/4$. Thus its two
rate-one single-port resampling updates have gap at least $1/4$ on their
complete conditional $L^2$ space.
\end{proposition}
\begin{proof}
Here $\Re C\ge-1/100$, $|C|\le201/100$. The unique conditional maximizer in
$W$ is $W_B(u)=(B+u)/|B+u|$. For $|u|\le1$, $|B+u|\ge199/100$.
For $B=3I$ one has
\[
 \frac{3+u_0}{|3I+u|}\ge\frac{\sqrt8}{3}\ge\frac{14}{15}.
\]
Indeed $(3+u_0)^2-8|u_\perp|^2\ge(3u_0+1)^2\ge0$. Changing $B$ by
at most $1/100$ changes the normalized vector by at most $1/100$, using
$|3I+u|\ge2$ and
$|x/|x|-y/|y||\le2|x-y|/|y|$. Hence
\begin{equation}
           \Re W_B(u)\ge14/15-1/100>23/25.
 \label{f16v15:inner-cap}
\end{equation}
The same lower real part holds on line segments between two such vectors.
Thus $|C+w|\ge91/100$ on those segments.

The derivative of vector normalization has norm $1/|x|$. Applying it first
along the segment between two $u$'s in the unit ball, and then along the
segment between their $W_B$ values, proves that the best-response composition
\[
                u\longmapsto\frac{C+W_B(u)}{|C+W_B(u)|}
\]
is a contraction of $S^3$ in chord distance with constant at most
\begin{equation}
                          q_*:=\frac{10000}{18109}<1.
 \label{f16v15:inner-contraction}
\end{equation}
The sphere is complete. The contraction theorem gives a unique fixed point,
therefore a unique simultaneous conditional maximizer. Conversely every
global maximizer supplies such a point. No conditional-maximizer set is lost
by division, since \eqref{f16v15:inner-cap} makes $C+W_B(u)$ nonzero.

At this point the negative joint tangent Hessian is
\[
 \begin{pmatrix}aI_3&-T\\-T^{\mathsf T}&bI_3\end{pmatrix},\qquad
 a=|C+W|\ge91/100,\quad b=|B+U|\ge199/100,\quad\|T\|\le1.
\]
It has a uniform positive lower bound because $ab\ge18109/10000>1$, and all
upper bounds are fixed. Its Gaussian correlation is at most
$100/\sqrt{18109}<3/4$, the latter strict inequality being
$160000<9\cdot18109$. The compact parameter family consists of the two closed
field balls and every $V\in S^3$. Proposition~\ref{f16v15:Laplace}, with
amplitude one, applies uniformly. Its strict Gaussian margin proves the native
$3/4$ bound eventually.

For completeness, the standard two-projection implication does not change
the clock: for a pair with correlation $\theta$, the two centered
single-coordinate subspaces have synthesis norm at most $\sqrt{1+\theta}$.
Thus the sum of the two conditional-expectation projections is at most
$(1+\theta)I$ on the centered space, and the sum of the two rate-one
innovations is at least $(1-\theta)I$. This is V13's pair argument on the
present complete conditional.
\end{proof}

\subsection{Collapsing the repair port retains a nontrivial generated weight}

Integrate $W$ exactly. Set
\[
 Z(t)=\int_{S^3}e^{t\Re R}dH(R),\qquad t\ge0.
\]
The unnormalized marginal of $(U,V)$ is
\begin{equation}
 f_\gamma^{A,B,D}(U,V)
   =e^{\gamma[U\cdot V+A\cdot U+D\cdot V]}Z(\gamma|B+U|).
 \label{f16v15:collapsed}
\end{equation}
This generated factor must not be replaced by one. For $|B-3I|\le1/100$,
$199/100\le|B+U|\le401/100$. Define
\begin{equation}
 \begin{split}
 F_{A,B,D}(U,V)&=U\cdot V+A\cdot U+D\cdot V+|B+U|,\\
 a_{\gamma,B}(U)&=\gamma^{3/2}e^{-\gamma|B+U|}Z(\gamma|B+U|).
 \end{split}
 \label{f16v15:collapsed-phase}
\end{equation}
Then $f_\gamma=\gamma^{-3/2}a_{\gamma,B}e^{\gamma F}$. The common scalar
$\gamma^{-3/2}$ cancels in \eqref{f16v15:normalized-kernel}, but the amplitude
$a_{\gamma,B}$ does not disappear from the finite-coupling law.

\begin{proposition}[A complete outer-pair bound near the witness]
\label{f16v15:outer}
There exist $0<\eta_0\le1/100$ and $\gamma_2<\infty$ such that, throughout
$\mathcal T_{\eta_0}$ and for $\gamma\ge\gamma_2$, the actual collapsed pair
\eqref{f16v15:collapsed} has maximal correlation at most $2/3$. At the exact
witness its limiting correlation is $2/\sqrt{11}$.
\end{proposition}
\begin{proof}
First check the amplitude, including the uniformity used below. Angular Haar
measure is $(2/\pi)\sin^2\theta\,d\theta$ on $[0,\pi]$, so
\[
 e^{-t}Z(t)=\frac2\pi\int_0^\pi
               e^{-t(1-\cos\theta)}\sin^2\theta\,d\theta.
\]
Laplace expansion at zero gives
\begin{equation}
 a_{\gamma,B}(U)\longrightarrow
             \sqrt{2/\pi}\,|B+U|^{-3/2}
 \label{f16v15:amplitude-limit}
\end{equation}
uniformly with every fixed number of derivatives in $U,B$. To justify the
needed two derivatives directly, put $t=\gamma r$, with $r$ in the compact
positive interval above, and scale $\theta=v/\sqrt\gamma$ in the small-angle
integral. Differentiation in $r$ introduces powers of
$\gamma(1-\cos\theta)$; these have uniform polynomial Gaussian majorants
using $1-\cos\theta\ge2\theta^2/\pi^2$. The complement of a fixed small-angle
interval is exponentially small, also after these derivatives. Taylor
expansion and dominated convergence give the stated uniform derivative
limits; composition with the smooth $r=|B+U|$ gives the claim. The leading
constant is $(2/\pi)\int_0^\infty v^2e^{-rv^2/2}dv
=\sqrt{2/\pi}\,r^{-3/2}$. The amplitudes therefore have fixed positive upper
and lower bounds for sufficiently large $\gamma$.

We next verify the normalized deficit, not merely the total energy minimum.
At $(A,B,D)=(I,3I,0)$, write
\[
 a=1-U_0,\quad b=1-V_0,\quad d=1-U\cdot V,\qquad
 T(U)=U_0+\sqrt{10+6U_0}.
\]
Here $F(U,V)=U\cdot V+T(U)$, and $M(U)=1+T(U)$. With
$N(V)=\max_U F(U,V)$ and $E=\mathcal D(U,V)$, the nonnegative two
conditional gaps imply $2E\ge d$. Also $N(V)\ge V_0+5$ by testing $U=I$.
Concavity of the square root yields
$T(U)\le5-(7/4)a$. Consequently
\begin{equation}
                    2E\ge2d+\tfrac74 a-b.
 \label{f16v15:outer-pin-inequality}
\end{equation}
Chord-distance Young inequalities give
$b\le\tfrac32 a+3d$ and $b\le2a+2d$. From the first, the last displayed
bound implies $2E\ge a/4-d$; combine it with $2E\ge d$ to get $E\ge a/16$.
Since also $E\ge d/2$, one obtains the global bound
\begin{equation}
                  \mathcal D(U,V)\ge\frac{a+b}{52}.
 \label{f16v15:outer-global-pin}
\end{equation}
Thus the normalized deficit has exactly the zero $(I,I)$.

In ordinary orthonormal tangent coordinates at that point, the negative
Hessian of $F$ is
\begin{equation}
       H_{\rm out}=\begin{pmatrix}11/4&-1\\-1&1\end{pmatrix}\otimes I_3.
 \label{f16v15:outer-Hessian}
\end{equation}
Indeed $U\cdot V=1-|u-v|^2/2+O(|(u,v)|^3)$ and
$T(\operatorname{Exp}u)=5-7|u|^2/8+O(|u|^4)$. The Hessian is positive
and its Gaussian correlation is $2/\sqrt{11}<5/8$.

We now make this a fixed \emph{open-field} statement. The phases in
\eqref{f16v15:collapsed-phase} are smoothly dependent on $(A,B,D)$ near the
witness, since $|B+U|$ stays positive. Uniform convergence of phases implies
uniform convergence of $M,N$ and of their normalized deficits. The strict
pin \eqref{f16v15:outer-global-pin} places every zero of a nearby deficit in
a fixed small normal-coordinate neighbourhood of $(I,I)$. Choose this
neighbourhood convex in its Euclidean coordinates and small enough that the
joint Hessian of the baseline phase is strictly negative there. It remains
strictly negative for all sufficiently small parameter perturbations.
Every zero of a normalized deficit is a critical point of the joint phase.
Strict concavity permits at most one such point there; a global maximum
always provides one. Hence nearby deficits have exactly one zero and a
uniformly positive negative joint Hessian there. By continuity, shrink a
fixed $\eta_0>0$ so that every corresponding Gaussian pair correlation is at
most $5/8$.

The compact parameter set $\mathcal T_{\eta_0}$ and the amplitudes
\eqref{f16v15:amplitude-limit} now satisfy Proposition~\ref{f16v15:Laplace}.
Uniform Hilbert--Schmidt convergence makes the native second singular value
at most $5/8+1/24=2/3$ for $\gamma\ge\gamma_2$. At the exact witness its
limit is the specified $2/\sqrt{11}$. This proves an upper correlation
bound on the \emph{whole} pair $L^2$ space, not just convergence of a few
moments. All free compact tails and both marginal denominators were paid.
\end{proof}

\subsection{The full three-port gap in the original update clock}

For a probability $\pi$ on these three variables, write
\[
 \mathcal E_U(f)=\mathbb E_\pi\operatorname{Var}_\pi(f\mid V,W),\quad
 \mathcal E_V(f)=\mathbb E_\pi\operatorname{Var}_\pi(f\mid U,W),\quad
 \mathcal E_W(f)=\mathbb E_\pi\operatorname{Var}_\pi(f\mid U,V).
\]
The block quantity $\mathcal E_{UW}(f)=
\mathbb E_\pi\operatorname{Var}_\pi(f\mid V)$ is not one of the original
single-port updates. The next proof explicitly bounds it by those updates.

\begin{theorem}[Complete three-port conditional repair, for every test function]
\label{f16v15:triplet-gap}
There are $0<\eta_0\le1/100$ and a finite $\gamma_0$, depending only on this
fixed local conditional, such that for every $\gamma\ge\gamma_0$, every
$(A,B,D)\in\mathcal T_{\eta_0}$, and every $f\in L^2(\pi_\gamma^{A,B,D})$,
\begin{equation}
 \boxed{\quad
 \operatorname{Var}_{\pi_\gamma^{A,B,D}}f
       \le12\{\mathcal E_U(f)+\mathcal E_V(f)+\mathcal E_W(f)\}.
 \quad}
 \label{f16v15:three-port-gap}
\end{equation}
In particular the three original single-port resamplings, each at rate one,
have gap at least $1/12$. No support assumption on $f$, no free-port cutoff,
and no boundary probability estimate is used. The constants are independent
of the ambient spatial volume and history length.
\end{theorem}
\begin{proof}
The exact star density gives $W\perp V\mid U$. For a centered function
$h(U,W)$ and a centered function $g(V)$, this conditional independence gives
\[
 \mathbb E[h(U,W)g(V)]
   =\mathbb E[\mathbb E(h\mid U)\,g(V)].
\]
Conditional expectation decreases the $L^2$ norm. Hence maximal correlation
between $(U,W)$ and $V$ is no greater than that between $U$ and $V$ under the
collapsed law. The reverse inequality follows by taking $h$ to depend only
on $U$, so the two correlations are equal. Proposition~\ref{f16v15:outer}
and the two-projection argument therefore give
\begin{equation}
          \operatorname{Var}_\pi f\le
                  3\{\mathcal E_{UW}(f)+\mathcal E_V(f)\}.
 \label{f16v15:outer-block-gap}
\end{equation}
At every fixed $V$, Proposition~\ref{f16v15:inner} applies to the actual
conditional $(U,W)\mid V$, uniformly over all $V\in S^3$. Integrating its
variance inequality gives
\[
                 \mathcal E_{UW}(f)
                      \le4\{\mathcal E_U(f)+\mathcal E_W(f)\}.
\]
Insert this into \eqref{f16v15:outer-block-gap}. The coefficient on
$\mathcal E_V$ is actually three, so the common constant twelve is safe.
Take $\gamma_0=\max(\gamma_1,\gamma_2)$, increasing it if necessary for the
amplitude bounds. Conditional variance disintegration and continuity of
conditional projections extend the inequalities from bounded functions to
all $L^2$. Haar alignment of the two leaves is a separate coordinate bijection
on each free port and hence conjugates its single-port projection exactly.
Thus the result uses the original three port updates, not a changed clock.
\end{proof}

For comparison only, the complete leading precision of the exact witness is
\[
 \begin{pmatrix}3&-1&-1\\-1&1&0\\-1&0&4\end{pmatrix}\otimes I_3.
\]
Its determinant per colour is seven. Its rate-one single-port Gaussian gap is
$1-\sqrt{5/12}$ by V13's all-chaos formula. The Schur complement after
integrating $W$ is exactly \eqref{f16v15:outer-Hessian}. The proof of the
native theorem did not simply substitute this Gaussian gap: the two complete
normalized-kernel limits and the uniform conditional bound over $V$ were
needed before the comparison. V14's slow pair froze the variable that supplies
this additional pin.

\subsection{An exterior gate that acts on all native innovations}

Let $\mathcal R=\{c,d,p\}$ be the free triple, or a translated/permuted copy
whose induced graph and exterior field conditions are as above. At every
fixed exterior perform the specified leaf alignment and compute $(A,B,D)$.
Let $G_{\mathcal R}$ be the event that these fields lie in
$\mathcal T_{\eta_0}$. It is measurable in the complement of the \emph{entire}
triple. In particular it imposes no condition on $r_p$.
Let
\[
 E_{\mathcal R}=\mathbb E_{P^{\rm bal}}[\,\cdot\mid\omega_{\mathcal R^c}],
 \qquad Q_{\mathcal R}=I-E_{\mathcal R}.
\]
Use $\widehat Q_a$ for the original individual balanced updates from V14.

\begin{theorem}[Posterior-integrated repair with no support restriction]
\label{f16v15:gated}
For $\gamma\ge\gamma_0$ and every $f\in L^2(P^{\rm bal})$,
\begin{equation}
 \boxed{
 \mathbb E_P[1_{G_{\mathcal R}}|Q_{\mathcal R}f|^2]
       \le12\sum_{a\in\mathcal R}
                   \mathbb E_P[1_{G_{\mathcal R}}|\widehat Q_a f|^2].}
 \label{f16v15:gated-form}
\end{equation}
The expectations use the actual exterior posterior with all original
normalizers. At zero bridge history $G_{\mathcal R}$ contains a fixed open
neighbourhood of V14's witness in its exterior variables. The inequality
holds at general fixed bridge histories wherever the same aligned-field
gate is satisfied.
\end{theorem}
\begin{proof}
Condition the full native probability on $\mathcal R^c$. Its conditional on
the three free ports is precisely \eqref{f16v15:triplet}, because every term
meeting them was retained. On $G_{\mathcal R}$ apply
Theorem~\ref{f16v15:triplet-gap}, and integrate with respect to the actual
exterior marginal. The conditional variance on the left is the integral of
$|Q_{\mathcal R}f|^2$; each conditional single-coordinate variance on the
right is the corresponding squared original innovation. The field gate
commutes with these four projections, since it is exterior to all of them.
The open-neighbourhood statement follows from
Proposition~\ref{f16v15:normal-form} and the strict parameter window. No
constant reference exterior law has replaced this marginal.
\end{proof}

This improves V14's supported-function floor in a specific way. Both sides of
\eqref{f16v15:gated-form} vanish on constants, and the estimate applies to
arbitrary $f$, not only $f$ supported in the gate. It does \emph{not} assert
$\mathbb E[1_G|f|^2]\le C\widehat{\mathcal E}(f)$, which would fail on
constants. Nor is a gate depending on one of the three free variables inserted
and then commuted through its update. The formerly restricted repair variable
is now among the variables being resampled.

\subsection{Compatible local blocks can be added without an unrecorded clock}

Take a finite collection of such triples, with deterministic rates
$\alpha_{\mathcal R}\ge0$, and let $w_a$ be V14's weights. Suppose
\begin{equation}
  \kappa=\sup_a w_a^{-1}
              \sum_{\mathcal R\ni a}\alpha_{\mathcal R}<\infty.
 \label{f16v15:overlap}
\end{equation}
A translated finite set of shapes with bounded rates has such a constant
independent of $B,n$. Overlap is allowed; no independence is assumed. Define
\[
 \mathcal B_{\rm rep}=\sum_{\mathcal R}\alpha_{\mathcal R}
                         1_{G_{\mathcal R}}Q_{\mathcal R},\qquad
 \widehat{\mathcal L}_{\rm aug}=\widehat{\mathcal L}_w+\mathcal B_{\rm rep}.
\]
Each gated summand is a positive orthogonal projection because its indicator
commutes with $Q_{\mathcal R}$. It annihilates constants and has finite
interaction support in the native square-root-density representation.

\begin{proposition}[Same-clock comparison for a family of repaired blocks]
\label{f16v15:clock-comparison}
At every admitted coupling,
\begin{equation}
 0\preceq\mathcal B_{\rm rep}\preceq12\kappa\widehat{\mathcal L}_w,
 \qquad
 \widehat{\mathcal L}_w\preceq\widehat{\mathcal L}_{\rm aug}
             \preceq(1+12\kappa)\widehat{\mathcal L}_w.
 \label{f16v15:operator-comparison}
\end{equation}
Consequently an \emph{independently proved} gap $g>0$ for the augmented
operator would imply a gap at least $g/(1+12\kappa)$ for the original weighted
generator, and hence for the original unweighted balanced generator.
No such uniform augmented gap is asserted here.
\end{proposition}
\begin{proof}
The quadratic form of one gated projection is exactly the left side of
\eqref{f16v15:gated-form}. Apply that inequality, drop each gate only from
the nonnegative upper bound, sum the deterministic rates, and use
\eqref{f16v15:overlap}. This gives
\[
 \langle f,\mathcal B_{\rm rep}f\rangle
    \le12\sum_a\left(\sum_{\mathcal R\ni a}\alpha_{\mathcal R}\right)
                         \|\widehat Q_af\|_2^2
    \le12\kappa\langle f,\widehat{\mathcal L}_wf\rangle.
\]
Adding the original form proves \eqref{f16v15:operator-comparison}. Its common
kernel is constants since the augmented form contains the original one.
Comparison on their common centered space proves the stated implication.
No products of noncommuting gated operators or independence between different
triples were used. In particular block acceleration is not counted as free.
\end{proof}

\subsection{What has advanced, and the next global obligation}

The main new input is a full conditional $L^2$ theorem in place of a supported
repair floor. The slow two-port cancellation is not removed by a cutoff or
by changing its probability; its repairing boundary variable is retained.
A nontrivial collapsed normalizer, two complete normalized-kernel estimates,
and the original-clock comparison give the explicit constant $1/12$ on a
fixed open exterior-field neighbourhood. This is more than a Gaussian
preflight or a positive floor for a selected supported test function.

The conclusion does not cover every exterior and does not prove that these
gates exhaust V14's defects. Other port cancellation geometries, chord--port
correlations and chord-cube conditional landscapes remain. No probability
estimate for the gate is claimed or needed for the conditional theorem. A
small gate probability would not by itself pay the omitted spectral defect.
The exact global target is still V14's
\[
 \mathfrak D(f)\le u\widehat{\mathcal E}_w(f)
             +v\|\widehat{\mathcal L}_wf\|_2^2,
                       \qquad u<2/39.
\]
Neither \eqref{f16v15:gated-form} nor the comparison proposition establishes
that inequality for the uncovered sectors. An independent uniform gap for
$\widehat{\mathcal L}_{\rm aug}$ would be another route, but is not assumed.

The next noncircular calculation is therefore to organize compatible
repair-retaining blocks and bound the remaining defect sectors, keeping
actual gates and overlap costs. The present three-port result can enter that
calculation without a new supported-function localization error. Re-estimating
the old frozen pair cannot do so. V14's directional source-energy handoff is
unchanged; a uniform global native gap and a uniform native fourth moment
are still separate inputs. No posterior operator contraction or temporal
weighted-row theorem is newly inferred from the local conditional gap.
The prescribed bridge window, physical root Gauss restriction, same-top full
transfer comparison, physical-time calibration and interacting continuum
construction remain distinct obligations.

\clearpage
\subsection{Diagnostics, checkpoint preservation, and claim scope}

The accompanying exact diagnostic is
\begin{center}\small\path{verify_ym_f16_v15_three_port_repair.py}.\end{center}
It reconstructs the literal three-port fields and fixed-root constraints,
checks noncommuting tree-edge alignment, the full three-port and collapsed
Gaussian Hessians, the normalized-deficit algebra, and conditional-projection
and rate-one comparison identities on positive finite probability models.
All 434 exact assertions passed; the report records their families and source
hash.
These tests are not native integral enclosures or proof-assistant verification
of the normalized-kernel convergence argument. In particular $\eta_0$ and
$\gamma_0$ are not numerically certified, and no simulation is used to justify
a theorem constant.

The available V4--V14 diagnostics were rerun unchanged and passed 9,467,
89,339, 906, 311, 609, 990, 608, 463, 1,201, 222 and 622 assertions,
respectively. The package records those reports, input hashes, exact predecessor
recovery, compilation and rendered-page review.
The V15 cumulative source changes only V14's title version and inserts this
appendix before its final document-ending command. Removing that insertion
and restoring the title recovers V14 byte for byte. No earlier mathematical
claim or bibliography string is silently corrected. The V15 companion audit
covers only the supplied sources; no inaccessible historical checker or
external formalization is described as rerun. The next checkpoint is V20.

\begin{firewall}
V15 proves a uniform $1/12$ conditional gap for the three original rate-one
single-port updates on a fixed open exterior neighbourhood of an actual native
cancellation witness, with all free compact groups and all $L^2$ test functions
retained. It gives a posterior-integrated gated form inequality and an explicit
same-clock comparison for overlapping families of those blocks. It does not
prove gate coverage, a uniform full-native auxiliary gap, native local fourth
moments, unrestricted-volume coherent posterior contraction, a physical transfer
gap, or an interacting continuum Yang--Mills mass gap. The clock remains an
auxiliary resampling clock, not physical Euclidean time.
\end{firewall}
% END F16 V15 APPEND
% BEGIN F16 V16 APPEND -- 2026-09-17
% Insert immediately before the final document-ending command in V15.
% The predecessor body and bibliography strings are not edited.

\clearpage
\section{V16: local superstability and uniform native fast-field moments}
\label{f16v16:context}

V15 proves a complete three-port repair on one exterior neighbourhood, not
coverage of the integrated defect. Adding more neighbourhoods without a
coverage mechanism risks repeating that distinction. We take up the other
explicit input in V13--V15: uniform local fourth moments of the \emph{actual}
full fast-history law. The independent compact anchors already available in
V5 are stronger for this purpose than the global entropy comparison used in
V6--V12. They survive arbitrary finite sets of free space--time cells, not
only entire-time spatial cylinders.

We prove uniform local exponential moments, including a cube of inhomogeneous
quadratic sources. The proof first pays a conditional normalizer by a product
of fixed-radius cell balls. This gives a local energy inequality with an
\emph{actual random boundary-energy term}. Iterating that inequality through
successive balls, and using polynomial graph growth, removes the boundary
term after averaging in the same law. No small interaction coefficient,
translation invariance, conditional gap, or Gaussian replacement is used.

Consequences include the previously missing native fourth-moment bound,
exponential suppression of the specific low-field configurations of V14,
and $O(\gamma^{-1})$ residual and posterior-return \emph{trace densities} at
unrestricted volume. The global auxiliary gap remains unproved. A rare
configuration can still carry a large conditional correlation or a normalized
spectral test function; the probability estimates below do not erase that
obstruction.

\subsection{Dependencies, literature, and the non-duplication boundary}

The supplied V15 appendix and analytic audit were read in full. The current
lineage and the mounted f14/f15 archives and Grand Tensor Monograph V076 were
searched, with the relevant normalization, pin, and moment interfaces read.
The targeted archive search returned no semantic matches; direct text search
and identified source ranges were used instead. This is not independent
recertification of the entire archive. F15 V50 explicitly distinguishes its
total-action bound from a joint local bound on an arbitrary background. F15
V90 bounds local physical energy through a specified physical transfer-source
input. F14 V67 obtains joint tails for a different restricted conditional.
None of those conclusions is transplanted to the current fixed-bridge law.

For classical context, D.~Ruelle,
\href{https://doi.org/10.1007/BF01609400}{\emph{Probability estimates for
continuous spin systems}}, Commun. Math. Phys. 50 (1976), 189--194,
extends superstability probability estimates to continuous spins. The primary
publisher abstract and bibliographic record were checked here; its complete
proof is not invoked or described as reviewed. The present finite-graph
argument is proved below from its explicit hypotheses. ``Local
superstability'' here refers to those hypotheses and their moment consequence,
not a claim that every standard continuum-particle hypothesis has been verified.
The high-temperature or superquadratic-spin functional-inequality literature
is not used to supply a missing native gap.

This is the upstream alternative to another formal correction order. V5's
one-cell anchors, V7's exact Haar residual synthesis, and V12--V14's same-law
source and posterior identities are inherited. The new analytic input is the
local-energy argument under quadratic tilts and its uniform moment conclusion.
No NS assertion is needed. The memoranda's distinction between locality,
normalization, and order gain is retained; the next companion checkpoint is
still V20.

\subsection{Complete cells, local interactions, and a uniform source domain}

Return to the \emph{unmoving} complete coordinates of V5. Let
\[
 \mathcal V=\{(i,b):0\le i\le n,\ b\in(\mathbb Z/m\mathbb Z)^3\},
 \quad B=m^3,\quad m\ge4,\quad n\ge1,\quad
 \mathcal N=|\mathcal V|=B(n+1).
\]
A cell contains $\zeta_c=(x_{i,b},a_{i,b})$ for $i<n$ and just $x_{n,b}$
at the last time. Its dimension is 36 or 15. Set
\[
 e_c=|\zeta_c|^2,\qquad e_S=\sum_{c\in S}e_c,\qquad
 a_0=\frac1{6\pi^2},\qquad t_0=\frac{a_0}{4}=\frac1{24\pi^2}.
\]
The symbol $e_c$ in this appendix denotes cell energy, not a primitive
source-error component. Retain the bridge assumption
$\max_{i,e}|y_{i,e}|\le r<\infty$ and, for definiteness, take
$\gamma\ge\gamma_r:=1+r^2$. This keeps these bridge coordinates strictly
inside their principal charts. The moment argument has no further
volume- or duration-dependent entry threshold.

Let $d\nu_{\gamma,c}$ be the \emph{unnormalized} product of the original
complete principal-chart Lebesgue factors and their Haar powers in cell $c$.
These include chord powers $w_i=1/2$ at the dressed time ends, chord powers
one elsewhere, and port powers one. Thus
\[
 d\nu_\gamma(\zeta)=\prod_{c\in\mathcal V}d\nu_{\gamma,c}(\zeta_c)
                 =1_{\mathcal P'}J^f_\gamma(\zeta)\,d\zeta.
\]
No growing Haar-volume scalar is absorbed into a supposedly uniform constant.
The classical action has V5's exact decomposition
\begin{equation}
 U_\gamma=b_\gamma(\mathbf y)+\sum_c A_{\gamma,c}(\zeta_c)
                             +\sum_{\ell\in\mathcal L}V_{\gamma,\ell}(
                                  \zeta_{S_\ell};\mathbf y),
 \quad A_{\gamma,c}\ge a_0 e_c,\quad V_{\gamma,\ell}\ge0.
 \label{f16v16:decomposition}
\end{equation}
Here $S_\ell$ is a nonempty, uniformly bounded fast carrier. Word lengths,
carrier ranges, and the number of interactions incident to a cell are bounded
by fixed geometric constants. Fast-free terms remain in $b_\gamma$; it cancels
only when normalizing the fixed-bridge fast probability, not from the full
bridge-history weight.

For a real array $s=(s_c)$ with $|s_c|\le t_0$, define the exact auxiliary tilt
\begin{equation}
 dP_s=\frac{e^{\sum_c s_c e_c}}{\mathbb E_P e^{\sum_c s_c e_c}}\,dP,
 \qquad
 Z(s)=\int e^{-U_\gamma+\sum_c s_c e_c}\,d\nu_\gamma,
 \qquad P=P_0.
 \label{f16v16:tilts}
\end{equation}
All these are finite positive integrals at a fixed regulator. The sources
are moment-generating parameters, not new physical couplings or a replacement
for $P$. They leave the complete domains, endpoint amplitudes, and original
Haar factors unchanged. The tilted one-cell anchors remain at least
$3a_0 e_c/4$.

Join distinct cells if they occur together in a primitive carrier $S_\ell$.
Call this graph $\mathcal G$, adding no edges beyond a fixed geometric range.
Let $B_k(c)$ denote its ball of radius $k$, and let
$\partial S=\{d\notin S:\ d\sim c\text{ for some }c\in S\}$.
There is a fixed $C_{\rm gr}$ such that
\begin{equation}
 |B_k(c)|\le C_{\rm gr}(k+1)^4
       \quad(k\ge0),\qquad
 \partial B_k(c)\subset B_{k+1}(c)\setminus B_k(c).
 \label{f16v16:growth}
\end{equation}
Indeed each graph step changes spatial cube and time coordinates by a bounded
amount. A containing coordinate box gives this polynomial bound, including
wrapping and the two time endpoints. Saturation of a finite ball only improves
it. This argument does not assume a periodic time direction.

\begin{proposition}[Uniform local upper bounds and complete reference masses]
\label{f16v16:local-data}
There are $K_r<\infty$, $v_*>0$, and $G_*\ge1$, independent of $\gamma,B,n,s$,
such that
\begin{equation}
 A_{\gamma,c}\le K_r e_c,\qquad
 V_{\gamma,\ell}\le K_r\left(1+\sum_{c\in S_\ell}e_c\right),
 \label{f16v16:upper-quadratic}
\end{equation}
and, with $\mathbb B_c=\{|\zeta_c|\le1\}$,
\begin{equation}
 \nu_{\gamma,c}(\mathbb B_c)\ge v_*,\qquad
 \int e^{-a_0 e_c/2}\,d\nu_{\gamma,c}\le G_*.
 \label{f16v16:cell-mass}
\end{equation}
For every nonempty $S\subset\mathcal V$, the tilted conditional exponent is
exactly
\[
 U_{S,s}=\sum_{c\in S}(A_{\gamma,c}-s_ce_c)
           +\sum_{\ell:S_\ell\cap S\ne\varnothing}V_{\gamma,\ell}.
\]
It obeys $U_{S,s}\ge3a_0 e_S/4$, and on $\prod_{c\in S}\mathbb B_c$,
\begin{equation}
 U_{S,s}(\zeta_S;q)\le K'_r|S|+K'_r e_{\partial S}(q).
 \label{f16v16:trial-bound}
\end{equation}
Every crossing interaction is retained in this conditional exponent.
\end{proposition}
\begin{proof}
For unit quaternions $g_j$, $s(g)=|I-g|^2/2$ and telescoping give
\[
 \gamma s\left(\prod_{j=1}^k\operatorname{Exp}(w_j/\sqrt\gamma)\right)
 \le\frac12\left(\sum_{j=1}^k|w_j|\right)^2
 \le\frac{k}{2}\sum_{j=1}^k|w_j|^2.
\]
Every actual word in unmoving coordinates is a product of exponentials of
fixed local linear forms in fast inputs and the bounded bridge inputs.
Frames do not require taking a new logarithm of a product in this statement.
Bounded word length and local matrix norms prove
\eqref{f16v16:upper-quadratic}, with $r$ included in its constant. Endpoint
Gaussian quadratics are cube-local with a bounded coefficient matrix. The
anchor lower bound is V5's inherited cube and rooted-tree estimate, not a
new convexity assertion. Both bounds apply cell by cell.

On the unit cell ball every group input has angle at most one because
$\gamma\ge1$. Each factor $j(v)=(\sin|v|/|v|)^2$ is at least
$(\sin1)^2$, and a power $w\in[1/2,1]$ does not decrease this lower bound.
There are at most twelve group factors. The minimum of the positive
15- and 36-dimensional unit-ball volumes, times $(\sin1)^{24}$, is an
admissible $v_*$. Conversely $J^f_{\gamma,c}\le1$ on its entire chart, so
an ordinary Euclidean Gaussian integral gives an admissible $G_*$, for
example $\max\{1,(2\pi/a_0)^{18}\}$. No bound on the unweighted total
volume of a growing principal chart is used.

Conditioning $P_s$ on $S^c$ cancels exactly the exterior-only action terms,
exterior source factors, and exterior reference factors. The remaining
exponent is $U_{S,s}$. Every interaction meeting $S$ is included once.
Each such term has its exterior fast inputs in $\partial S$. On the
product of unit cell balls, bounded incidence bounds all free energies
and all constant bridge costs by $C_r|S|$, and bounds the exterior
contribution by $C_r e_{\partial S}$. The extra source term costs at most
$t_0|S|$ there. This proves \eqref{f16v16:trial-bound}. Positivity of all
non-anchor terms and $|s_c|\le t_0$ prove the lower bound for every exterior
value, not only for a bounded exterior window.
\end{proof}

\subsection{A local conditional normalizer with its boundary energy paid}

\begin{proposition}[Conditional exponential energy and an averaged boundary inequality]
\label{f16v16:conditional-energy}
There are constants $a_r,b_r\ge1$, independent of $B,n,\gamma,s$ and the
subset $S$, such that, at every admitted exterior value $q$,
\begin{equation}
 \log\mathbb E_{P_s}\left[e^{t_0 e_S}\mid\zeta_{S^c}=q\right]
 \le t_0 a_r|S|+t_0 b_r e_{\partial S}(q).
 \label{f16v16:conditional-mgf}
\end{equation}
Consequently, putting $u_c(s)=\mathbb E_{P_s}e_c$,
\begin{equation}
 \sum_{c\in S}u_c(s)\le a_r|S|+b_r\sum_{d\in\partial S}u_d(s).
 \label{f16v16:energy-boundary}
\end{equation}
Both expectations use the exact tilted native law. The conditional bound
itself is \emph{not} uniform in arbitrary exterior energy.
\end{proposition}
\begin{proof}
Write $Z_{S,s}(q)=\int e^{-U_{S,s}}\prod_{c\in S}d\nu_{\gamma,c}$.
Restrict its denominator only for a lower bound to the product of unit cell
balls, retaining every crossing term. Proposition~\ref{f16v16:local-data} gives
\[
 Z_{S,s}(q)\ge v_*^{|S|}
                  e^{-K'_r|S|-K'_r e_{\partial S}(q)}.
\]
The numerator with $e^{t_0 e_S}$ has exponent at most $-a_0 e_S/2$,
since $3a_0/4-t_0=a_0/2$. Integrating over the \emph{full} native charts,
and then enlarging only this positive upper bound to Euclidean space, gives
an upper bound $G_*^{|S|}$. Divide these two bounds to get
\eqref{f16v16:conditional-mgf}, for instance with
\[
 a_r\ge t_0^{-1}\bigl(\log G_* -\log v_*+K'_r\bigr),
 \qquad b_r\ge t_0^{-1}K'_r,
\]
where each constant is also enlarged to be at least one. Jensen's inequality
gives $t_0\mathbb E[e_S\mid q]\le\log\mathbb E[e^{t_0e_S}\mid q]$.
Integrate under the actual exterior marginal of $P_s$. This is
\eqref{f16v16:energy-boundary}. At a finite regulator every quantity is finite
by compactness, so neither truncation nor an unproved integrability estimate
has been used to start the argument.
\end{proof}

The denominator is the reason for the boundary term. Dropping it would give
a false all-exterior thermal estimate, contradicted by the unpinned pair
conditionals in V14. The next step averages and then removes this term; it
does not claim that the term vanishes pointwise.

\subsection{Polynomial growth closes the local moment inequality}

\begin{proposition}[Boundary iteration on a finite graph of polynomial growth]
\label{f16v16:ball-lemma}
Let a finite graph satisfy \eqref{f16v16:growth}. Suppose finite nonnegative
numbers $u_c$ satisfy, for every nonempty set $S$,
\[
 \sum_{c\in S}u_c\le a|S|+b\sum_{d\in\partial S}u_d,
 \qquad a,b>0.
\]
Set $\rho=b/(1+b)\in(0,1)$. Then every vertex obeys
\begin{equation}
 u_c\le M(a,b):=
 \frac{a C_{\rm gr}}{1+b}\sum_{j=0}^\infty\rho^j(j+1)^4
 =a C_{\rm gr}\frac{1+11\rho+11\rho^2+\rho^3}{(1-\rho)^4}.
 \label{f16v16:ball-bound}
\end{equation}
The interaction coefficient $b$ need not be smaller than one. No
translation symmetry or spectral-gap assumption is required.
\end{proposition}
\begin{proof}
Fix the vertex and write $F_k=\sum_{d\in B_k(c)}u_d$ and
$v_k=|B_k(c)|$. The boundary inclusion in \eqref{f16v16:growth} implies
\[
 F_k\le a v_k+b(F_{k+1}-F_k),\qquad
 F_k\le\rho F_{k+1}+\frac{a}{1+b}v_k.
\]
Iterating this scalar inequality gives, for every integer $R\ge1$,
\begin{equation}
 u_c=F_0\le\rho^R F_R+
                 \frac{a}{1+b}\sum_{j=0}^{R-1}\rho^j v_j.
 \label{f16v16:iteration}
\end{equation}
At this \emph{fixed finite graph}, $F_R\le\sum_d u_d<\infty$, so the first
term tends to zero as $R\to\infty$. This does not assume a uniform bound on
that total and does not interchange a volume limit with the iteration.
Apply the polynomial ball bound to the remaining sum. The identity
\[
 \sum_{j=0}^\infty\rho^j(j+1)^4
       =\frac{1+11\rho+11\rho^2+\rho^3}{(1-\rho)^5}
\]
follows by applying $(\rho\,d/d\rho+1)^4$ to $(1-\rho)^{-1}$.
Since $(1+b)^{-1}=1-\rho$, it proves \eqref{f16v16:ball-bound}.
A ball that has already saturated its connected component still satisfies
the same inequality with an empty exterior boundary. Thus small tori,
disconnected interaction graphs, and finite time endpoints cause no exception.
\end{proof}

The order of operations matters. A one-cell estimate alone would invite the
unsuccessful requirement $b\deg(\mathcal G)<1$. Instead the conditional
normalizer is estimated on every ball. The geometric decrease in
\eqref{f16v16:iteration} then beats polynomial ball growth. This proves a
moment bound even when a worst-boundary mixing row cannot be small.

\begin{theorem}[Uniform native local moments and joint quadratic-source pressure]
\label{f16v16:moments}
For every fixed bridge bound $r$, set $M_r=M(a_r,b_r)$ from
\eqref{f16v16:ball-bound}. For all $m\ge4$, $n\ge1$,
$\gamma\ge1+r^2$, all admitted bridge histories, and all arrays
$s\in[-t_0,t_0]^{\mathcal V}$,
\begin{equation}
                   \sup_c\mathbb E_{P_s}|\zeta_c|^2\le M_r.
 \label{f16v16:tilted-mean}
\end{equation}
For any two such arrays $s,t$,
\begin{equation}
 \left|\log\frac{Z(t)}{Z(s)}\right|
                           \le M_r\sum_c|t_c-s_c|.
 \label{f16v16:pressure-Lip}
\end{equation}
In particular, for arbitrary $0\le t_c\le t_0$ on any set of cells,
\begin{equation}
 \boxed{\quad
 \log\mathbb E_P e^{\sum_c t_c|\zeta_c|^2}
                              \le M_r\sum_c t_c.\quad}
 \label{f16v16:joint-mgf}
\end{equation}
Thus, for every fixed integer $k\ge1$,
\begin{equation}
 \sup_c\mathbb E_P|\zeta_c|^{2k}
       \le k!\,t_0^{-k}e^{t_0M_r},\qquad
 m_4(P)\le1+2t_0^{-2}e^{t_0M_r}<\infty.
 \label{f16v16:all-moments}
\end{equation}
These constants are independent of ambient volume, duration, and coupling in
the stated range. They concern the same $m_4(P)$ left open in V13--V15, not an
averaged moment or a moment of a reference law.
\end{theorem}
\begin{proof}
Apply the ball lemma to \eqref{f16v16:energy-boundary}, at each fixed $s$.
All constants in that inequality were uniform on the entire source cube,
which proves \eqref{f16v16:tilted-mean}. This step used only finite moments
from the compact regulator to start; the conclusion now makes them uniform.

At fixed regulator, compactness permits exact differentiation of the source
partition function along the straight segment $s+u(t-s)$, $0\le u\le1$:
\[
 \frac d{du}\log Z(s+u(t-s))
           =\sum_c(t_c-s_c)\mathbb E_{P_{s+u(t-s)}}e_c.
\]
The segment stays inside the same source cube. Bound the absolute value by
\eqref{f16v16:tilted-mean} and integrate. This proves
\eqref{f16v16:pressure-Lip} and, with $s=0$, \eqref{f16v16:joint-mgf}.
This is integration of controlled \emph{actual means}, not differentiation
of an unmarked asymptotic pressure error. Finally
$u^k\le k!t_0^{-k}e^{t_0u}$ for $u\ge0$ and
\eqref{f16v16:joint-mgf} on one cell prove
\eqref{f16v16:all-moments}. The definition of $m_4(P)$ in V13 is
$\sup_c\mathbb E_P(1+|\zeta_c|^4)$, so $k=2$ is the required input.
\end{proof}

This theorem does not give a native logarithmic Sobolev inequality, nor
covariance decay. It does not control moments uniformly under arbitrary
\emph{fixed} exterior configurations. The actual marginal, with its conditional
partition functions included, is essential. Its constants are conservative
expressions in fixed local bounds, not numerically optimized entry constants.

\subsection{Joint large-field tails, including a genuinely native cancellation family}

\begin{proposition}[Joint energy tails and exponentially rare low-field pairs]
\label{f16v16:tails}
For every cell set $S$, threshold $R\ge0$, and $u\ge0$,
\begin{align}
 P\{e_S\ge M_r|S|+u\}&\le e^{-t_0u},\notag\\
 P\{e_c\ge R^2\text{ for every }c\in S\}
             &\le\exp[-t_0(R^2-M_r)|S|].
 \label{f16v16:cell-tails}
\end{align}
Let $\mathcal B_{cd}$ be precisely V14's low-five-neighbour-field event on
one moving-port bond. There are fixed $c_1>0,C_r<\infty$ and a fixed-$r$
coupling threshold such that, for \emph{every} finite set $\mathcal A$ of
distinct neighbouring nonroot port pairs, at arbitrary times,
\begin{equation}
 P\left(\bigcap_{a\in\mathcal A}\mathcal B_a\right)
                     \le\exp[-(c_1\gamma-C_r)|\mathcal A|].
 \label{f16v16:low-field-joint}
\end{equation}
In particular $\mathbb E_P N_{\rm low}/\mathcal N
\le C\exp(C_r-c_1\gamma)$, and the same bound without the counting factor
holds for each specified pair. No independence of the events is assumed.
\end{proposition}
\begin{proof}
Apply exponential Markov with $t_c=t_0$ on $S$ in
\eqref{f16v16:joint-mgf}; both assertions in \eqref{f16v16:cell-tails}
follow immediately.

We retain the actual low-field definition, not all possible conditional-angle
defects. Each five-neighbour field is a sum of five transported unit
quaternions. If its norm is at most one, its distance from $5I$ is at least
four. The telescoping and geodesic comparison in V14 therefore gives, for
each event $a$, a fixed-size \emph{unmoving} cell set $S_a$ and $c_*>0$
such that
\begin{equation}
 \mathcal B_a\ \Longrightarrow\ e_{S_a}\ge c_*\gamma.
 \label{f16v16:low-field-energy}
\end{equation}
For completeness, $|I-r_p|$ is bounded by the sum of the chord distances of
$f_t^{-1},h_p,f_{t+1}$, and these are at most a fixed multiple of
$\gamma^{-1/2}(|a_p|+|x_{t,b}|+|x_{t+1,b}|)$. Moving internal edges have
the same bound in their own cube chords; bridge distances are at most
$r/\sqrt\gamma$. There are finitely many factors in the five transported
words. Cauchy--Schwarz on that fixed list, after absorbing $r/\sqrt\gamma$
for sufficiently large fixed-$r$ coupling, proves
\eqref{f16v16:low-field-energy}. This is the local version of V14's
deterministic implication, before summing the field count.

There are fixed integers $s_*,k_*$ with $|S_a|\le s_*$ and every cell in at
most $k_*$ such sets, counting all pair types and times. Put
$T=\bigcup_{a\in\mathcal A}S_a$. On the intersection event,
\[
 k_*e_T\ge\sum_{a\in\mathcal A}e_{S_a}\ge c_*\gamma|\mathcal A|,
 \qquad |T|\le s_*|\mathcal A|.
\]
Use \eqref{f16v16:joint-mgf} on $T$ and exponential Markov. The result is
\eqref{f16v16:low-field-joint} with $c_1=t_0c_*/k_*$ and
$C_r=t_0M_rs_*$. There are at most $C\mathcal N$ pair labels. Summing the
single-label probability proves the density statement. All coordinates,
including all time slices in any selected set, retain their actual native
law; a product exterior distribution was not inserted.
\end{proof}

A further precise geometric consequence is available. Give the low-field
labels any fixed finite-range adjacency of maximum degree $d_*\ge2$, and
put $p_\gamma=\exp(C_r-c_1\gamma)$. The probability that the bad connected
component of a specified label has size at least $k$ is at most
\begin{equation}
                    d_*^{2(k-1)}p_\gamma^k.
 \label{f16v16:cluster-tail}
\end{equation}
Indeed a connected set of size at least $k$ contains a connected subset of
exactly $k$ containing that label. Such subsets number at most
$d_*^{2(k-1)}$: choose a deterministic spanning tree and its depth-first walk
of length $2(k-1)$, which records the visited set injectively. Sum
\eqref{f16v16:low-field-joint}. This is a bad-\emph{field} cluster bound,
not an identification of all spectral defects with such clusters.

\subsection{Source-weighted tails and improved unrestricted-volume trace estimates}

Let $\mathbf e^{\rm prim}=(e^{\rm br},e^{\rm pt})$ be V7's primitive bridge
and compact Haar-port error banks, and let $\mathbf r$ be its full corrected
common-source residual vector. The exact synthesis is still
\[
 \mathbf r=e^{\rm br}+\mathsf P^{\mathsf T}e^{\rm pt}.
\]
The port lift has a volume-independent operator norm and exponentially
summable absolute columns. Each primitive component $e_j^{\rm prim}$ obeys
the inherited complete-chart bound
\[
 |e_j^{\rm prim}|\le C_r\gamma^{-1/2}
                       (1+e_{T_j}),\qquad |T_j|\le C,
\]
with bounded carrier incidence. These are statements about the same compact
Haar derivatives, not flat substitutes at finite coupling.

\begin{theorem}[Actual source moments, exceptional weights, and all-lag trace density]
\label{f16v16:source-energy}
For every fixed $p\ge1$, every primitive component and every canonical common
mark $\alpha$,
\begin{equation}
 \|e_j^{\rm prim}\|_{L^p(P)}+
 \|r_\alpha\|_{L^p(P)}\le C_{p,r}\gamma^{-1/2}.
 \label{f16v16:source-Lp}
\end{equation}
For any event $E$ in the full native probability space,
\begin{equation}
 \mathbb E_P[1_E|r_\alpha|^2]
       +\mathbb E_P[1_E|e_j^{\rm prim}|^2]
                          \le C_r\gamma^{-1}P(E)^{1/2}.
 \label{f16v16:marked-event}
\end{equation}
In particular \eqref{f16v16:low-field-joint} pays these fixed source weights
on a joint low-field event with an additional factor
$\exp[-(c_1\gamma-C_r)|\mathcal A|/2]$.
The full positive Gram $\Gamma_\gamma=\operatorname{Cov}_P(\mathbf r)$ and
V7's actual nonadjacent-time Hessian satisfy
\begin{equation}
 \boxed{\quad
 \operatorname{Tr}\Gamma_\gamma\le C_r\mathcal N/\gamma,
 \qquad \frac{\|\mathsf M_\gamma\|_{S_1}}{\mathcal N}\le C_r/\gamma.
 \quad}
 \label{f16v16:trace}
\end{equation}
Also every canonical entry has $|\Gamma_{\gamma,\alpha\beta}|\le C_r/\gamma$,
independently of both mark positions. Thus the same entry bound holds for the
full-history common-source mixed Hessian at nonadjacent times. It has no
decay in the separation. For V4's original core radius in its derivative
regime,
\begin{equation}
 \frac{\|\mathsf C_{\gamma,R}\|_{S_1}}{\mathcal N}
                         \le C_r/\gamma+C R^2/\gamma.
 \label{f16v16:core-trace}
\end{equation}
All these statements permit unrestricted $B,n$ in their fixed-$r$ coupling
regime. There is no $e^{KB}/\sqrt\gamma$ entry condition.
\end{theorem}
\begin{proof}
The uniform cell moments and the bounded size of $T_j$ give
$\|1+e_{T_j}\|_p\le C_{p,r}$. Apply the primitive estimate. For a single
full residual, Minkowski's inequality and the uniform absolute column sum of
$\mathsf P$ give the same $L^p$ bound. The latter summability follows from
the finite-range positive port matrix by the inverse decay already used in
V10--V12, and is unchanged at either dressed time endpoint. Cauchy--Schwarz
with $p=4$ proves \eqref{f16v16:marked-event}; this is a bound for the stated
observables, not arbitrary normalized $L^2$ functions.

There are $O(\mathcal N)$ primitive components and $9\mathcal N$ canonical
common sources. Sum their second moments, or use the bounded whole-vector
synthesis, to obtain $\mathbb E\|\mathbf r\|^2\le C_r\mathcal N/\gamma$.
Centering reduces this quantity and proves the trace bound. Individual
covariance entries are bounded by their two standard deviations. V7's
\emph{exact} contact identity is
$\mathsf M_\gamma=-\operatorname{Off}_{>1}\Gamma_\gamma$. Its Fourier
extraction bound is $\|\operatorname{Off}_{>1}A\|_{S_1}\le4\|A\|_{S_1}$,
independent of duration. Positivity of the full Gram gives
\eqref{f16v16:trace}. V4's original core operator-norm estimate, multiplied
by its $9\mathcal N$ source dimension, and the unchanged full/core identity
give \eqref{f16v16:core-trace}, exactly as in V7 but with the new energy rate.
No new event replaces the core, and no endpoint normalizer is deleted.
\end{proof}

These are sharper consequences than the $\epsilon_\gamma$ energy floor used
in V7 and V12, and than the time-length-dependent pointwise bounds at zero
history. They still do not identify the leading coefficient $G_2$ at arbitrary
volume. Small second moments of the residuals yield an $O(\gamma^{-1})$
scale; a uniform remainder $o(\gamma^{-1})$ and connected covariance decay
require additional arguments. The source-weighted exceptional estimate in
\eqref{f16v16:marked-event} likewise does not apply to every spectral
innovation without an appropriate norm or capacity estimate.

\subsection{The posterior input improves, while its coherent gap remains explicit}

Use V12's separated-cylinder family, without repeated sources, common
exterior projection $E_O$, selection $\chi$, depth $L$, and enlarged-patch
overlap $q_0$. All exterior values are integrated. The established exact
identities are
\begin{equation}
 \mathbf r^D=\chi\mathbf r-\Delta^{\mathsf T}e^{\rm pt},\qquad
 \mathbf g^D-\chi\mathbf g=-\Delta^{\mathsf T}S_0,\qquad
 \mathbf j=\mathbf g^D+E_O\mathbf r^D,\qquad
 \|\Delta\|\le C\sqrt{q_0}e^{-\nu L}.
 \label{f16v16:posterior-identities}
\end{equation}
Here $S_0=\partial_aU_0$ is the affine Gaussian port score \emph{evaluated
under the native law}. The conditional partition functions remain in the
actual exterior posterior. These are V12's harmonic lift identities, not a
new or extra covariance subtraction.

\begin{proposition}[Sharp-scale posterior trace and the remaining gap-only reduction]
\label{f16v16:posterior}
The actual posterior return and separated-cylinder gluing error obey
\begin{align}
 \frac{\operatorname{Tr}\mathsf R_{\rm post}}{\mathcal N}
       &\le C_rq_0(\gamma^{-1}+e^{-2\nu L}),\notag\\
 \frac{\|\chi\mathsf M_\gamma\chi-\mathsf L_D\|_{S_1}}{\mathcal N}
       &\le C_rq_0(\gamma^{-1}+e^{-2\nu L}).
 \label{f16v16:posterior-trace}
\end{align}
For any exterior events $E_s$, the exceptional conditional Gram integrals also
satisfy
\begin{equation}
 \sum_s\mathbb E_{\bar P_s}
          [1_{E_s}\operatorname{Tr}\Gamma_\gamma^{D_s}]
                                \le C_rq_0\mathcal N/\gamma.
 \label{f16v16:exceptional-Gram}
\end{equation}
Let $\widehat\lambda_P>0$ be the \emph{actual} balanced auxiliary gap of V14.
The following unnormalized operator bounds now have no unproved moment input:
\begin{align}
 \|\Gamma_\gamma\|_{\rm op}
                       &\le \frac{C_r}{\widehat\lambda_P\gamma},\notag\\
 \|\mathsf R_{\rm post}\|_{\rm op}
    +\|\chi\mathsf M_\gamma\chi-\mathsf L_D\|_{\rm op}
                       &\le\frac{C_rq_0}{\widehat\lambda_P}
                                        (\gamma^{-1}+e^{-2\nu L}).
 \label{f16v16:gap-only}
\end{align}
A uniform positive lower bound on $\widehat\lambda_P$ remains unproved;
\eqref{f16v16:gap-only} does not assert one.
\end{proposition}
\begin{proof}
The primitive and full residual energies just proved are
$O_r(\mathcal N/\gamma)$. Apply the bounded synthesis and lift-mismatch
estimate in \eqref{f16v16:posterior-identities}; this gives
$\mathbb E\|\mathbf r^D\|^2\le C_rq_0\mathcal N/\gamma$.
Because $S_0$ is a fixed finite-range affine map of the unmoving coordinates,
with bounded coefficients and fixed-$r$ constants per cell, the local
second moments also give $\mathbb E_P\|S_0\|^2\le C_r\mathcal N$.
Therefore
\[
 \mathbb E_P\|\mathbf j-\chi\mathbf g\|^2
 \le2\mathbb E\|\Delta^{\mathsf T}S_0\|^2
                +2\mathbb E\|E_O\mathbf r^D\|^2
 \le C_rq_0\mathcal N(\gamma^{-1}+e^{-2\nu L}).
\]
Use conditional $L^2$ contraction and centering to get the first trace bound.
V12's same-posterior identity
$\chi\mathsf M_\gamma\chi-\mathsf L_D=
-\operatorname{Off}_{>1}\mathsf R_{\rm post}$ and the factor-four trace
extraction give the second. For \eqref{f16v16:exceptional-Gram}, conditional
centering and the tower property give an upper bound by
$\sum_s\mathbb E\|\mathbf r^{D_s}\|^2$, which is bounded by the same
combined lift synthesis. The nonnegative gates can be dropped only from that
upper bound. This improves the rate of V12's exceptional trace budget, but
not to $o(\gamma^{-1})$.

Finally V14's directional innovation-energy inequalities hold for these
same balanced updates and for the same native moment definitions. Insert
\eqref{f16v16:all-moments}, in place of a previously unverified $m_2,m_4$,
to obtain
\[
 \widehat{\mathcal E}_P(v^{\mathsf T}\mathbf r)
      \le C_r\gamma^{-1}\|v\|^2,\qquad
 \widehat{\mathcal E}_P(b^{\mathsf T}S_0)\le C_r\|b\|^2.
\]
Their actual finite-regulator Poincar\'e inequality and V13--V14's exact
posterior reduction give \eqref{f16v16:gap-only}. The current theorem proves
the moment input, not the gap input. Replacing $\widehat\lambda_P$ by the
Gaussian margin $2/39$ would still be unjustified. Its crude finite-regulator
bound alone does not yield an unrestricted-volume contraction.
\end{proof}

\subsection{What changes in the program, and what is not inferred}

The separate full-native local moment obligation in V13--V15 is now paid,
uniformly over bounded inhomogeneous bridge histories, all durations, and all
spatial tori. It is paid without using V6's entropy-density theorem, V9's
product comparison, V15's repair neighbourhood, or an assumed native
functional inequality. Their respective conclusions remain useful but were
not hidden inputs to this moment proof. The argument uses the exact pinned
fast-coordinate specification; it is not a moment theorem for an unrestricted
physical gauge connection or for the integrated bridge law.

The low-field result is also stronger than a mean density: it controls any
prescribed joint family, and its connected components in a bounded-degree
label graph have an exponential size tail at sufficiently large coupling.
It is a concrete input for an eventual source-marked cluster or repair
argument. It does \emph{not} establish that V14's entire angle defect is
supported on those low-field events. Chord--port and other conditional
landscapes remain in that defect. Even a complete bad-event probability
estimate would not permit multiplying it by the squared innovation of an
arbitrary normalized test function independently of that function.

There is no contradiction with V14's compact cancellation. At its selected
fixed exterior, a port marginal can be Haar, with principal rescaled energy
of order $\gamma$. The conditional bound
\eqref{f16v16:conditional-mgf} then has a large exterior-energy debit. The
uniform conclusion is under the \emph{averaged full native law}. Every such
exterior neighbourhood can retain positive probability and still defeat a
worst-exterior correlation criterion. Neither V13's nor V14's impossibility
result is reversed.

The remaining active global task is therefore narrower: control the native
auxiliary gap, or directly the coherent source covariance, using compatible
repairs with a source-resolved treatment of the uncovered sectors. Uniform
local moments no longer need to be assumed in that reduction. Another
trace-density estimate would not exclude a coherent exceptional direction.
The all-lag matrices above have no new time-row decay at unrestricted volume,
and their $O(\gamma^{-1})$ bounds do not identify a leading coefficient with
an $o(\gamma^{-1})$ error. The existing infrared contact, gauge-invariant
bridge coverage, same-top physical transfer, calibrated physical time, and
interacting continuum obligations remain distinct.

\subsection{Diagnostics, preservation, and claim scope}

The accompanying diagnostic is
\begin{center}\small\path{verify_ym_f16_v16_local_moments.py}.\end{center}
It checks the ball-iteration algebra and its exact generating function,
finite-graph growth including saturated balls, conditional source and
normalizer identities on positive finite laws, and the conversion of local
energy witnesses into joint-event bounds with overlap. It also checks the
source/lift and posterior quadratic-form bookkeeping used by the corollaries.
All 1,077 exact assertions passed. The report separates those tests from
analytic hypotheses and inequalities.
The proof of the full-native moment theorem is the argument above, not an
inference from those finite examples.

The cumulative V15 is preserved byte for byte except for its title version
and insertion of this exact appendix before the final document-ending
command. The package supplies the exact recovery checker, the protected input
hashes, build and render reports, the new diagnostic, and actual reruns of the
available V4--V15 diagnostics, which passed 9,467, 89,339, 906, 311, 609,
990, 608, 463, 1,201, 222, 622, and 434 assertions, respectively.
No inaccessible archive program, external
formalization, or numerical enclosure of every native derivative bound is
represented as verified. The next companion checkpoint remains V20.

\begin{firewall}
V16 proves local exponential moment bounds for the complete native fast law
at a fixed bounded bridge history. These bounds are uniform in volume and
duration. They supply the missing fourth moment and joint low-field tails,
and improve residual and posterior trace densities to order $1/\gamma$.
The operator reduction still requires a uniform native auxiliary gap.

No such gap, exhaustive defect coverage, or spatially uniform time-row
mixing is established. Nor is a leading coefficient identified with an
error smaller than $1/\gamma$ at unrestricted volume. Physical bridge
coverage, the same-top physical transfer comparison, and an interacting
continuum Yang--Mills mass gap remain unproved.
\end{firewall}
% END F16 V16 APPEND

% BEGIN F16 V17 APPEND -- 2026-09-17
% Insert immediately before the final document-ending command in V16.
% No predecessor proof, claim, or bibliography string is silently revised.

\clearpage
\section{V17: local Stein recursion and coherent estimates on growing source banks}
\label{f16v17:context}

V16 supplies uniform moments of every fixed degree for the complete native
fast law. This changes what can be done without a native resampling gap.
Instead of further enlarging a repair neighbourhood, we apply a polynomial
Gaussian divergence identity directly under the native law. Local cutoffs
make integration by parts legitimate on the principal charts; V16 pays their
errors one coordinate group at a time. An all-chart, all-volume weak
expansion follows. Its coefficients are exactly V10's, not coefficients of a
new probability measure.

The new conclusion has a different norm from V10's. At every fixed order its
entrywise covariance remainder has a constant independent of both spatial
volume and history length. It has no separation decay. Nevertheless, combined
with the inherited locality of the coefficients, it gives operator-norm
control on every source bank of any prescribed polynomial size in coupling,
inside an otherwise unrestricted ambient history. No native auxiliary gap is
assumed. A completely delocalized coherent direction remains outside this
conclusion; we quantify that limitation instead of identifying the entrywise
norm with a full operator norm.

\subsection{Dependency audit and the change of method}

The supplied V16 appendix and analytic audit were read in full. Two targeted
semantic searches in the f14/f15/monograph sources returned no results. The
mounted inventories and exact text were then examined, not treated as empty.
F14 V98, \texttt{prop:v98-polynomial-Stein}, already gives a polynomial Haar
Stein error for its adapted curvature fields. F15 V64,
\texttt{f15v64:polynomial-Stein}, already gives location-uniform linear and
cubic Stein tests for its stationary, randomly framed curvature field. These
are explicitly credited. They do not give the present finite-order expansion
of the fully pinned augmented fast history: their laws, root-frame costs,
and harmonic-mode issues differ. This is a targeted non-duplication audit,
not recertification of the archives.

The relevant classical mechanism is approximation through a Stein identity,
not transferring a spectral gap between nearby measures. Barbour,
\href{https://doi.org/10.1007/BF01197887}{\emph{Stein's method for diffusion
approximations}} (1990), provides general context; only its primary summary
and bibliographic record were checked. Braverman--Dai--Fang,
\href{https://arxiv.org/abs/2012.02824}{\emph{High-order steady-state diffusion
approximations}}, derive recursive approximations through Stein/Poisson
equations. Their abstract and introductory stationary-identity setup were
reviewed, not their complete model-specific proofs. No theorem about their
Markov chains is transplanted here. The finite polynomial right inverse and
the native error estimates used below are proved directly.

The two import memoranda distinguish order gain from locality cost, and
compatible local constraints from an assumed physical gap. We keep those
distinctions. The present proof reuses V10's actual measured action and Haar
source coefficients, V11's inverse bounds, and V16's native moments. It does
not use V6's likelihood comparison, V9's product-entry condition, or V15's
repair gate. It does not reopen the failed worst-exterior pair-row test.
The next scheduled companion checkpoint remains V20.

\subsection{The exact measure and a polynomial norm with no hidden dimension}

Put $\varepsilon=\gamma^{-1/2}$ and fix an arbitrary admitted bridge history
$\max_{i,e}|y_{i,e}|\le r$. Let
\[
 P_\varepsilon=P_{\gamma,\mathbf y},\qquad
 Q=N(m,C),\qquad C=H^{-1},\qquad \xi=\zeta-m.
\]
These are the same complete unmoving native and Gaussian laws as V10 and V16.
The Gaussian mean and precision are independent of $\varepsilon$ at fixed
rescaled $\mathbf y$ and fixed geometry. There are fixed spectral bounds
$h_-I\preceq H\preceq h_+I$, bounded means per cell, and a fixed interaction
range for $H$. Its inverse has exponentially summable coordinate rows. In
particular, with scalar real indices $i,j$ and fixed coordinate multiplicity,
\begin{equation}
 \sup_i\sum_j|H_{ij}|\le K_H,\qquad
 \sup_i\sum_j|C_{ij}|\le K_C,
 \qquad \sup_i|m_i|\le M_r^{\rm G}.
 \label{f16v17:Gaussian-input}
\end{equation}
The constants do not depend on $B,n$. For clarity, the inverse row bound is
not inferred merely from $\|C\|_{2\to2}$. The finite-range Neumann expansion
of $H^{-1}$ gives exponential entry decay; polynomial growth of the cell
graph and its at most 36 coordinates per cell make those rows summable.
This is the inherited inverse-locality argument on the complete history,
including its actual dressed endpoints.

For a scalar polynomial $f(\xi)$ of degree at most $d$, collect its ordinary
monomials and set
\begin{equation}
 \|f\|_{\rm pol}=\sum_{\mathbf k}|f_{\mathbf k}|,
 \qquad f(\xi)=\sum_{|\mathbf k|\le d}f_{\mathbf k}\xi^{\mathbf k}.
 \label{f16v17:pol-norm}
\end{equation}
For a polynomial vector $g=(g_j)$ put
$\|g\|_{{\rm pol},1}=\sum_j\|g_j\|_{\rm pol}$.
The finite number of variables can be arbitrarily large. Degree, not that
number, is kept fixed. Repeated indices are included. The norm is
submultiplicative, and
\[
 \sum_j\|\partial_j f\|_{\rm pol}\le d\|f\|_{\rm pol}.
\]
A sum of many unrelated observables pays its actual coefficient sum; no
claim that this norm is dimension free for every such sum is made.

\begin{proposition}[Native polynomial moments and coordinate-local cut tails]
\label{f16v17:poly-moments}
For every fixed $d$ and $1\le p<\infty$,
\begin{equation}
 \|f\|_{L^p(P_\varepsilon)}+\|f\|_{L^p(Q)}
       \le C_{d,p,r}\|f\|_{\rm pol}.
 \label{f16v17:poly-Lp}
\end{equation}
If $g$ is one original three-coordinate group, and
$A_g=\{|\zeta_g|\ge(2\varepsilon)^{-1}\}$, then, for some fixed $c>0$,
\begin{equation}
 \mathbb E_{P_\varepsilon}[1_{A_g}|f|]
          \le C_{d,r}e^{-c/\varepsilon^2}\|f\|_{\rm pol}.
 \label{f16v17:cut-tail}
\end{equation}
Both estimates are uniform in the locations and number of variables appearing
in $f$ when its displayed coefficient norm is used. They hold for
$\gamma\ge1+r^2$.
\end{proposition}
\begin{proof}
V16 bounds every fixed moment of each $|\zeta_c|$, uniformly in its cell.
Subtracting the bounded Gaussian mean preserves these bounds for each
$\xi_i$. H\"older applied to a monomial of degree at most $d$, allowing
repeated coordinates, bounds its $L^p$ norm by a fixed coordinate moment of
order at most $dp$. Minkowski then gives the native part of
\eqref{f16v17:poly-Lp}. The Gaussian part follows from bounded coordinate
variances and the same argument. No independence of the native coordinates
is used.

The cell containing $g$ has energy at least $1/(4\varepsilon^2)$ on $A_g$.
V16's one-cell exponential moment gives
$P_\varepsilon(A_g)\le\exp(t_0M_r-t_0/(4\varepsilon^2))$.
Cauchy--Schwarz and the $L^2$ polynomial bound prove
\eqref{f16v17:cut-tail}, for example with $c=t_0/8$.
This estimates one group tail weighted by the actual polynomial. It is not
a probability estimate for a simultaneous all-good event over the history.
\end{proof}

\subsection{A finite Gaussian divergence solver, without a native inverse}

Define the Gaussian divergence on polynomial vector fields by
\[
 \delta_Q g=\sum_j\bigl((H\xi)_j g_j-\partial_jg_j\bigr).
\]
A useful right inverse can be constructed by finite algebra. For a monomial
$M=\xi_{i_1}\cdots\xi_{i_d}$ choose its first index by one fixed ordering,
put $h=\xi_{i_2}\cdots\xi_{i_d}$, and set
\begin{equation}
 \mathcal B M=\sum_{\ell=2}^d C_{i_1i_\ell}
                \prod_{k\ne1,\ell}\xi_{i_k},\qquad
 \mathcal S M=(C_{j i_1}h)_j+\mathcal S(\mathcal B M),
 \qquad \mathcal S1=0.
 \label{f16v17:solver-recursion}
\end{equation}
Empty sums are zero; extend linearly to collected polynomials. Recursion
reduces the degree by two and terminates. The letter $\mathcal B$ in this
subsection is this degree-lowering map, not a spatial cube set.

\begin{proposition}[A coefficient-summable Gaussian right inverse]
\label{f16v17:solver}
For every polynomial $f$ of degree at most $d$,
\begin{equation}
 \delta_Q\mathcal S f=f-\mathbb E_Qf,
 \qquad \deg\mathcal S f\le d-1,
 \qquad \|\mathcal S f\|_{{\rm pol},1}\le K_d\|f\|_{\rm pol}.
 \label{f16v17:solver-bound}
\end{equation}
For constant $f$ the vector is zero. The constants $K_d$ depend on $d,K_C$,
not on $B,n$. Thus this construction requires no inverse of a native
resampling generator or native differential operator.
\end{proposition}
\begin{proof}
Since $HC=I$, direct polynomial differentiation gives
\[
 \delta_Q(C_{\cdot i_1}h)=M-\mathcal B M.
\]
Ordinary Gaussian integration by parts gives
$\mathbb E_Q M=\mathbb E_Q\mathcal B M$. Applying the degree induction to
the second term of \eqref{f16v17:solver-recursion} therefore proves the exact
divergence identity. Equivalently this is Wick's finite recursion, not a
convergence assertion for an infinite series.

For one monomial the coefficient norm of the first vector is at most $K_C$.
The norm of $\mathcal B M$ is at most $(d-1)K_C$, and its degree is $d-2$.
Thus constants satisfying $k_0=0$, $k_1=K_C$ and
$k_d=K_C+(d-1)K_C k_{d-2}$ bound the construction on degree-$d$ monomials.
Taking $K_d=\max_{0\le l\le d}k_l$ and summing absolute coefficients proves
the bound. Coordinate ordering can change the right inverse, but not its
divergence. Later normalized coefficients will be independent of that choice.
\end{proof}

Only one covariance row is summed for each algebraic step. Replacing that
sum by a sum of unrelated coordinate error bounds would lose the point of
\eqref{f16v17:solver-bound}. No semigroup continuity on an infinite-dimensional
function space is needed for this finite polynomial construction.

\subsection{A native Stein identity with the principal-chart boundary paid}

On the positive chart interior retain V10's measured potential
\[
 V_\varepsilon=U_\varepsilon-\log J^f_\varepsilon,
 \qquad V_\varepsilon\sim U_0+\sum_{q\ge1}\varepsilon^qV_q.
\]
The $V_q$ here are \emph{exactly} \eqref{f16v10:jets}, including the
endpoint Haar powers. Each $\partial_jV_q$ has a fixed local carrier and
\begin{equation}
 \sup_j\|\partial_jV_q(m+\xi)\|_{\rm pol}\le L_{q,r},
 \qquad \deg(\partial_jV_q)\le q+1.
 \label{f16v17:local-drift-coefficients}
\end{equation}
Indeed differentiating an extensive local action retains only its finitely
many terms incident to $j$. Translation by the bounded $m$ preserves the
fixed-degree coefficient bound. No norm of the entire extensive $V_q$ is
used below.

\begin{proposition}[Full-native polynomial Stein expansion]
\label{f16v17:Stein}
For any fixed integer $N\ge0$ and any polynomial vector $g$ of degree at
most $d$, there are constants $C_{d,N,r},\gamma_{d,N,r}$ independent of
volume, duration, and the vector's support, such that
\begin{equation}
 \left|\mathbb E_{P_\varepsilon}\delta_Qg+
   \sum_{q=1}^N\varepsilon^q
           \mathbb E_{P_\varepsilon}(\nabla V_q\cdot g)\right|
 \le C_{d,N,r}\varepsilon^{N+1}\|g\|_{{\rm pol},1}
 \label{f16v17:Stein-estimate}
\end{equation}
for $\gamma\ge\gamma_{d,N,r}$. The sum is empty when $N=0$.
Every expectation uses the complete native law. No all-coordinate cutoff is
imposed on that law, and no derivative of an unmarked pressure error is taken.
\end{proposition}
\begin{proof}
We give the cutoff argument because flat integration by parts through the
logarithmic cut without its Haar term would be incorrect. Let $g(j)$ denote
the original group containing scalar coordinate $j$. Choose one smooth radial
function on $\mathbb R^3$, equal to one on radius $1/2$ and zero outside
radius one, and write
\[
 \chi_j(\zeta)=\chi(\varepsilon\zeta_{g(j)}),
 \quad 0\le\chi_j\le1,\qquad |\partial_j\chi_j|\le C\varepsilon.
\]
It cuts off only this coordinate group. For fixed values of all other groups,
$\chi_jg_j$ has support strictly inside its principal ball of radius
$\pi/\varepsilon$. The native Lebesgue density on that ball is smooth and
positive. Fubini and compact finite-regulator integrability therefore give
\[
 \mathbb E[\chi_j\partial_jg_j+(\partial_j\chi_j)g_j]
                    =\mathbb E[\chi_jg_j\partial_j V_\varepsilon].
\]
Other groups need not be in their small charts. Consequently the following
identity, before estimating anything, is exact:
\begin{equation}
 \begin{split}
 \mathbb E\delta_Q g={}&
 -\sum_j\mathbb E[\chi_j g_j
            (\partial_jV_\varepsilon-(H\xi)_j)]
 +\sum_j\mathbb E[(\partial_j\chi_j)g_j]\\
 &+\sum_j\mathbb E[(1-\chi_j)
                 ((H\xi)_jg_j-\partial_jg_j)].
 \end{split}
 \label{f16v17:exact-cut-identity}
\end{equation}
Here and for the remainder of this proof expectations are under
$P_\varepsilon$. This is a chart integration identity with its measured
Jacobian. It is not an assertion that a compact Haar derivative equals a
flat derivative. The exact Haar-defined source observables remain unchanged.

For the classical drift, every local word is
$\varepsilon^{-2}f(\varepsilon w)$, with $f(0)=Df(0)=0$ and uniformly
bounded real derivatives of every fixed order, by unitary exponential
insertion. Its coordinate derivative has a Taylor remainder of order
$\varepsilon^{N+1}$ bounded by a fixed local polynomial of degree at most
$N+2$. This bound holds for \emph{all} real inputs appearing along the
Taylor segment. Fixed incidence and the bounded bridge inputs therefore give,
for every fixed $p$,
\[
 \left\|\partial_jU_\varepsilon-\partial_jU_0
                 -\sum_{q=1}^N\varepsilon^q\partial_jU_q\right\|_{L^p(P_\varepsilon)}
                         \le C_{N,p,r}\varepsilon^{N+1}.
\]
The proof uses Proposition~\ref{f16v17:poly-moments}, not moments at an
intermediate changed coupling.

The Haar logarithm need not have bounded derivatives at its cut. On the
support of $\chi_j$, however, its only differentiated group obeys
$|\varepsilon\zeta_{g(j)}|\le1$. The smooth function
$-w\log j(\varepsilon\zeta_g)$, with the original $w\in\{1/2,1\}$, has
Taylor remainder bounded there by
$C_N\varepsilon^{N+1}(1+|\zeta_g|^{N+2})$ after one coordinate derivative.
It includes the order-two term in \eqref{f16v10:haar-order-two} and no
additional copy of it. Other Haar factors have zero $j$ derivative. Thus
\[
 \left\|\chi_j\left[
 \partial_jV_\varepsilon-(H\xi)_j
       -\sum_{q=1}^N\varepsilon^q\partial_jV_q\right]\right\|_{L^p(P_\varepsilon)}
                  \le C_{N,p,r}\varepsilon^{N+1}.
\]
For $N=0$ the classical difference begins at order one and the Haar
contribution at order two, so the same estimate applies.

Pair this drift error with $g_j$ by Cauchy--Schwarz. Removing $\chi_j$ from
the finite polynomial terms costs only an expectation on $A_{g(j)}$.
Every other term of \eqref{f16v17:exact-cut-identity} is likewise supported
on that event, including $\partial_j\chi_j$. The row bound for $H$, the
polynomial product rule, and \eqref{f16v17:cut-tail} bound the sum of these
errors by
\[
 C_{d,N,r}\bigl(\varepsilon^{N+1}+e^{-c/\varepsilon^2}\bigr)
                                    \sum_j\|g_j\|_{\rm pol}.
\]
At fixed order the flat term is absorbed into $\varepsilon^{N+1}$ by
increasing the constant and, if needed, the fixed-order coupling threshold.
There is no union bound over all coordinate groups and no multiplication by
$B(n+1)$. This proves \eqref{f16v17:Stein-estimate}.
\end{proof}

\subsection{A normalized finite-order expansion without a spatial entry condition}

For a polynomial $f$ define finite algebraic operators
\begin{equation}
 \mathcal T_q f=-\sum_j(\partial_jV_q)(\mathcal Sf)_j,
 \quad q\ge1;
 \qquad b_0(f)=\mathbb E_Q f,
 \quad b_k(f)=\sum_{q=1}^k b_{k-q}(\mathcal T_q f).
 \label{f16v17:coefficient-recursion}
\end{equation}
All expressions are evaluated in $\xi$, with $\zeta=m+\xi$ in $V_q$.
They use one fixed Gaussian covariance and fixed measured coefficients.
The maps obey
\begin{equation}
 \deg\mathcal T_q f\le d+q,\qquad
 \|\mathcal T_q f\|_{\rm pol}\le L_{q,r}K_d\|f\|_{\rm pol}.
 \label{f16v17:T-bound}
\end{equation}
This follows by multiplication and the \emph{sum} in
\eqref{f16v17:solver-bound}, not a sum of $\mathcal N$ local constants.

\begin{theorem}[All fixed orders of native polynomial expectations]
\label{f16v17:expectations}
For any fixed degree $d$ and order $N\ge0$, there are finite
$C_{d,N,r},\gamma_{d,N,r}$ such that for every finite spatial torus,
every duration, every admitted bridge history, and every degree-$d$
polynomial independent of $\varepsilon$ in the centered coordinates,
\begin{equation}
 \boxed{
 \left|\mathbb E_{P_\varepsilon}f-
                  \sum_{k=0}^N\varepsilon^k b_k(f)\right|
       \le C_{d,N,r}\varepsilon^{N+1}\|f\|_{\rm pol}.}
 \label{f16v17:expectation-expansion}
\end{equation}
The condition is only $\gamma\ge\gamma_{d,N,r}$; it contains no spatial
volume or duration. Also $|b_k(f)|\le C_{d,k,r}\|f\|_{\rm pol}$.
These are coefficients of the \emph{normalized native law}; action-only
constants and all normalization corrections are treated correctly.
\end{theorem}
\begin{proof}
Apply Proposition~\ref{f16v17:Stein} to $g=\mathcal S f$ and use its exact
divergence identity. For every fixed $N$ this gives
\begin{equation}
 \mathbb E_{P_\varepsilon}f
  =\mathbb E_Q f+\sum_{q=1}^N\varepsilon^q
                  \mathbb E_{P_\varepsilon}\mathcal T_qf
       +O_{d,N,r}(\varepsilon^{N+1}\|f\|_{\rm pol}).
 \label{f16v17:native-recursion}
\end{equation}
For $N=0$ this is already the theorem. Induct on $N$, applying the
order-$N-q$ statement to $\mathcal T_q f$. Its degree increases by at most
$q$, and its coefficient norm is bounded by \eqref{f16v17:T-bound}.
There are only $N$ terms. Collecting each total power gives exactly
\eqref{f16v17:coefficient-recursion}, and all errors are of order
$\varepsilon^{N+1}$. The coefficient bound follows by the same induction
from the Gaussian moment bound. Only finitely many degrees and thresholds
occur at each fixed $(d,N)$, so their maximum is independent of $B,n$.

Here is an algebraic check of coefficient ownership which does not assume
that an exponential of a cubic polynomial is integrable. At each fixed
finite regulator form the \emph{formal} normalized functional
\[
 \mathcal M_z(f)=
 \frac{\mathbb E_Q[f\exp(-\sum_{q\ge1}z^qV_q)]}
      {\mathbb E_Q[\exp(-\sum_{q\ge1}z^qV_q)]}.
\]
Every coefficient involves only finitely many Gaussian polynomial moments.
Gaussian integration by parts applied coefficient by coefficient gives
\[
 \mathcal M_z(\delta_Q g)
       =-\sum_{q\ge1}z^q\mathcal M_z(\nabla V_q\cdot g).
\]
With $g=\mathcal S f$, its coefficients satisfy the same triangular
recursion as $b_k$. The leading value is $\mathbb E_Q f$, so the coefficients
are uniquely those of \eqref{f16v17:coefficient-recursion}, independently of
the ordering used to choose $\mathcal S$. The already proved error estimate
identifies them with the actual native asymptotic coefficients. Thus the
formal calculation identifies ownership; it does not supply the error bound.
Fast-independent additions to $V_q$ disappear from both its gradient and this
normalized functional, as required.
\end{proof}

For example these normalized coefficients obey
\begin{equation}
 b_1(f)=-\operatorname{Cov}_Q(f,V_1),\qquad
 b_2(f)=-\operatorname{Cov}_Q(f,V_2)
                        +\tfrac12\kappa_Q(f,V_1,V_1).
 \label{f16v17:first-expectation-coefficients}
\end{equation}
The recursion uses local \emph{derivatives} of the extensive action, whereas
the equivalent cumulants keep its scalar normalization. Constants $f$ have
no corrections. A native sampler gap was not used to solve the Stein
equation: the right inverse is the finite Gaussian algebra above.

\subsection{Apply the expansion to the same exact Haar residuals}

Keep the literal residuals $r_{\varepsilon,\alpha}$ of V7 and their
coefficients $r_{\alpha,k}$ of V10. In particular the order-one coefficient
still includes
\[
 \frac12\sum_p[\eta_{\alpha,p},a_p]\cdot\partial_{a_p}U_0,
 \qquad \eta_\alpha=\mathsf P\mathsf e_\alpha.
\]
The new flat-chart integration identity does not remove this noncommutative
Haar-insertion term or differentiate a different source. All retained bridge
derivatives are formed in the original coordinates first.

\begin{proposition}[Uniform native source jets]
\label{f16v17:source-jets}
For every fixed $K\ge1$ and $1\le p<\infty$,
\begin{equation}
 \left\|r_{\varepsilon,\alpha}-
                 \sum_{k=1}^K\varepsilon^k r_{\alpha,k}\right\|_{L^p(P_\varepsilon)}
        \le C_{K,p,r}\varepsilon^{K+1},
 \qquad \|r_{\alpha,k}(m+\xi)\|_{\rm pol}\le C_{k,r}.
 \label{f16v17:source-jets-bound}
\end{equation}
These bounds are uniform in the source position, ambient volume, and duration.
The same statements hold for each primitive port or bridge remainder.
\end{proposition}
\begin{proof}
Insert the scaled Lie-algebra generator directly into each ordered compact
word, as in V7 and V10. Its fixed-order real Taylor remainder is bounded by
$C_{K,r}\varepsilon^{K+1}(1+|\zeta_T|^{K+2})$ on its fixed local carrier.
Uniform native moments from V16 bound its $L^p$ norm at the actual coupling.
The primitive polynomial coefficients have degree at most $k+1$ and bounded
coefficient norms after the bounded mean translation. The full residual is
the local bridge remainder plus the same $\mathsf P^{\mathsf T}$ synthesis
of port remainders. Its absolute column sum is uniformly finite. Minkowski
and the triangle inequality for polynomial coefficients therefore preserve
both estimates. There is no new native kernel comparison or exterior window.
\end{proof}

\begin{theorem}[Unrestricted-volume entrywise expansion of the native memory]
\label{f16v17:entry-expansion}
For every fixed $J\ge2$ there are $C_{J,r},\gamma_{J,r}$ independent of
$B,n$ such that for all admitted histories and $\gamma\ge\gamma_{J,r}$,
\begin{equation}
 \boxed{
 \sup_{\alpha,\beta}
 \left|\Gamma_{\gamma,\alpha\beta}
        -\sum_{q=2}^{J}\gamma^{-q/2}(G_q)_{\alpha\beta}\right|
       \le C_{J,r}\gamma^{-(J+1)/2}.}
 \label{f16v17:entry-remainder}
\end{equation}
The $G_q$ are exactly V10's connected coefficients
\eqref{f16v10:connected-coefficients}. The identical estimate holds for
$\mathsf M_\gamma$ and $M_q=-\operatorname{Off}_{>1}G_q$.
The norm in \eqref{f16v17:entry-remainder} is $\ell^1\to\ell^\infty$,
not an unrestricted operator norm or a time-weighted row norm.
\end{theorem}
\begin{proof}
Truncate both residuals through degree $J-1$. Their $L^2$ errors are
$O(\varepsilon^J)$ by Proposition~\ref{f16v17:source-jets}; their full and
truncated $L^2$ norms are $O(\varepsilon)$. Cauchy--Schwarz controls the
resulting covariance error by $C_{J,r}\varepsilon^{J+1}$.
For each product of source coefficients $r_{\alpha,k}r_{\beta,l}$ with
$k+l\le J$, apply Theorem~\ref{f16v17:expectations} through order
$J-k-l$. Apply it also to each factor when expanding the product of its
two means. Their polynomial coefficient norms are bounded uniformly, even
when their two source positions are arbitrarily distant. The omitted terms
with $k+l>J$ are bounded by fixed moments and cost
$O(\varepsilon^{J+1})$. There are finitely many terms at this fixed order.

The coefficient-ownership calculation in
Theorem~\ref{f16v17:expectations}, followed by subtracting the two native
mean expansions, is precisely the connected-cumulant formula of V10. Thus
its Gaussian term and all measure corrections are retained. In particular,
\[
 G_2=\operatorname{Cov}_Q(r_{\alpha,1},r_{\beta,1}),
\]
and $G_3$ contains the existing term
$-\kappa_Q(r_{\alpha,1},r_{\beta,1},V_1)$; it is not an expansion of the
observables alone. This proves \eqref{f16v17:entry-remainder} uniformly.
Finally V7's exact compact Haar identity, not the new chart identity, gives
$\mathsf M_\gamma=-\operatorname{Off}_{>1}\Gamma_\gamma$.
Applying that entry mask to the estimate proves the memory assertion.
No changed contact term or endpoint normalizer has been dropped.
\end{proof}

The new $J=2$ conclusion is an actual $o(\gamma^{-1})$ error for every
entry at unrestricted ambient volume:
\[
 \Gamma_{\gamma,\alpha\beta}
     =\gamma^{-1}(G_2)_{\alpha\beta}+O_r(\gamma^{-3/2}).
\]
This does not supersede V10's stronger temporal norm in its restricted
spatial regime. V10 pays $e^{K_{J,r}B}$ to retain geometric temporal decay of
the remainder; the present argument removes that spatial cost but does not
retain remainder decay. These are different theorems in different norms.

\subsection{A coherent conclusion on polynomially growing source banks}

V10 already proves $\|G_q\|_{\mathrm{loc},4a}\le C_{q,r}$ at one fixed
$a>0$, for each fixed $q$. We reuse that result, including its full-history
Gaussian mean and covariance, rather than repeat the linked-pairing proof.
It implies $\|G_q\|_{\rm op}+\|M_q\|_{\rm op}\le C_{q,r}$.
Let $\chi_{\mathcal A}$ select any finite set of distinct common-source
coordinates, with no restriction on their relative positions, and put
$k=|\mathcal A|$.

\begin{theorem}[Mixed-norm remainder and source-size control]
\label{f16v17:coherent}
Under the fixed-order hypotheses of Theorem~\ref{f16v17:entry-expansion},
write $R_{\gamma,J}=\Gamma_\gamma-\sum_{q=2}^J\gamma^{-q/2}G_q$.
For any real vectors $v,w$,
\begin{equation}
 |v^{\mathsf T}R_{\gamma,J}w|
       \le C_{J,r}\gamma^{-(J+1)/2}\|v\|_1\|w\|_1,
 \qquad
 \|\chi_{\mathcal A}R_{\gamma,J}\chi_{\mathcal A}\|_{\rm op}
       \le C_{J,r}k\gamma^{-(J+1)/2}.
 \label{f16v17:mixed-norm}
\end{equation}
The same estimates hold for the memory remainder. In particular
\begin{equation}
 \operatorname{Var}_{P_\varepsilon}(v^{\mathsf T}\mathbf r)
   \le C_{J,r}\left(\gamma^{-1}\|v\|_2^2
                +\gamma^{-(J+1)/2}\|v\|_1^2\right).
 \label{f16v17:variance}
\end{equation}
For every \emph{fixed} $p\ge0$ there are finite
$C_{p,r},\gamma_{p,r}$ such that all source banks with $k\le\gamma^p$
satisfy, at every ambient $B,n$ and every admitted history,
\begin{equation}
 \boxed{
 \begin{aligned}
 \|\chi_{\mathcal A}(\Gamma_\gamma-\gamma^{-1}G_2)
                         \chi_{\mathcal A}\|_{\rm op}
       &\le C_{p,r}\gamma^{-3/2},\\
 \|\chi_{\mathcal A}(\mathsf M_\gamma-\gamma^{-1}M_2)
                         \chi_{\mathcal A}\|_{\rm op}
       &\le C_{p,r}\gamma^{-3/2}.
 \end{aligned}}
 \label{f16v17:polynomial-banks}
\end{equation}
Consequently the restricted native covariance and memory have operator norms
at most $C_{p,r}/\gamma$. No average over source vectors or division by
$B(n+1)$ occurs.
\end{theorem}
\begin{proof}
Sum the entrywise remainder against $|v_\alpha w_\beta|$ to get the first
bound. On a bank of size $k$, $\|v\|_1\le\sqrt{k}\|v\|_2$ proves the
operator bound. The sum of the fixed-order coefficient operators has norm
at most $C_{J,r}/\gamma$ for $\gamma\ge1$, by the inherited absolute row
bounds. Its quadratic form plus the remainder gives \eqref{f16v17:variance}.
The memory mask preserves the entrywise estimate and the absolute coefficient
row bounds; positivity of the memory matrix is not asserted.

Given fixed $p$, choose a fixed integer $J\ge\max\{2,2p+2\}$. Then
$k\gamma^{-(J+1)/2}\le\gamma^{-3/2}$. The remaining finite sum of
coefficients with $q\ge3$ has operator norm at most
$C_{J,r}\gamma^{-3/2}$. Apply \eqref{f16v17:mixed-norm} and enlarge the
constant and threshold for this chosen $J$. This proves
\eqref{f16v17:polynomial-banks}. The argument is uniform over all choices
of the bank and all bounded bridge histories. The exponent $p$ is fixed
before taking the large-coupling limit; no bound on growth of the constants
as $p\to\infty$ is claimed.
\end{proof}

For $J=2$, the raw covariance remainder on a $k$-coordinate bank is
$C_rk\gamma^{-3/2}$, and the covariance itself is $O_r(\gamma^{-1})$
when $k\le\sqrt\gamma$. Higher fixed orders remove that particular source-size
restriction: \eqref{f16v17:polynomial-banks} identifies the same leading
operator on any prescribed polynomially growing bank. The rest of the
ambient torus and the full duration are unrestricted. This is not a theorem
about a growing cylinder on which all exterior fields have been fixed.

\begin{proposition}[What a remaining coherent obstruction must look like]
\label{f16v17:delocalization}
Define $k_{\rm eff}(v)=\|v\|_1^2/\|v\|_2^2$ for $v\ne0$.
For every fixed $J\ge2$ and unit vector $v$,
\begin{equation}
 \left|\gamma v^{\mathsf T}\Gamma_\gamma v-v^{\mathsf T}G_2v\right|
 \le C_{J,r}\gamma^{-1/2}
             +C_{J,r}\gamma^{-(J-1)/2}k_{\rm eff}(v).
 \label{f16v17:effective-support}
\end{equation}
Suppose along a sequence $\gamma\to\infty$ of admitted finite histories,
unit vectors $v_\gamma$ have the left side bounded below by a fixed
$\delta>0$. Then for every fixed $p\ge0$,
\begin{equation}
                  k_{\rm eff}(v_\gamma)/\gamma^p\longrightarrow\infty.
 \label{f16v17:superpolynomial}
\end{equation}
The same implication holds for $\mathsf M_\gamma,M_2$.
\end{proposition}
\begin{proof}
Multiply \eqref{f16v17:mixed-norm} by $\gamma$ and use the uniform operator
bounds for the higher coefficients. Their sum starts at
$\gamma^{-1/2}$, proving \eqref{f16v17:effective-support}.
For any fixed $p$, select a fixed $J$ with $(J-1)/2>p$.
Eventually the first term is at most $\delta/2$; hence
$k_{\rm eff}(v_\gamma)\ge c_{J,r,\delta}\gamma^{(J-1)/2}$.
Divide by $\gamma^p$. Each selected order has its own finite threshold,
which the sequence eventually crosses. No order is chosen as a function of
the instantaneous volume. The masked memory proof is identical.
\end{proof}

This is a necessary delocalization condition for an order-one failure of the
scaled leading-coefficient comparison, not a construction of such a failure
or a proof that it is impossible. There is no upper bound on ambient source
dimension. Very large volumes still admit vectors with this effective support.
The finite-source conclusion therefore does not give a global spectral gap.

\subsection{Separation, the posterior, and the next noncircular target}

The coefficients are exponentially local, whereas the new error is uniform
but unweighted. Combining the two gives, for every fixed $J\ge2$,
\begin{equation}
 |\Gamma_{\gamma,\alpha\beta}|
 \le C_{J,r}\gamma^{-1}e^{-4a d(\alpha,\beta)}
                          +C_{J,r}\gamma^{-(J+1)/2}.
 \label{f16v17:clustering-floor}
\end{equation}
The same estimate applies to the nonadjacent-time memory entries. Thus
separated response has an error floor smaller than any chosen algebraic
order, after choosing that fixed order. It is \emph{not} an absolutely
summable remainder at a fixed coupling. Summing the last term over all
source positions would again introduce the number of positions.

In particular this appendix does not replace V16's exact posterior
reduction by an unproved global one. V12's exterior messages are nonlinear
conditional expectations over entire cylinders; their covariance still
contains the actual exterior return. We do not identify that return with a
Gaussian marginal or infer a global sampler gap from the local Stein
recursion. V15's repair theorem remains a possible tool for a genuinely
source-resolved control of that return.

What has changed is the source-size frontier. A leading native coefficient
with an error $o(\gamma^{-1})$ is now identified at unrestricted ambient
volume in the entrywise norm, and in operator norm on every prescribed
polynomially growing source bank. Another repair neighbourhood or another
extensive trace budget was not needed. The unresolved upgrade is to an
\emph{absolutely summable connected remainder}, or an independent global
native inequality controlling the remaining delocalized directions. The
polynomial divergence recursion supplies coefficients with a full native
error, but its coefficient-$\ell^1$ estimate does not manufacture that
connected decay. It is important not to rename that loss as a closed
thermodynamic estimate.

All statements still concern a prescribed bounded rescaled bridge history,
before physical root Gauss restriction and integration of the bridges.
No new bridge-coverage theorem, physical-time calibration, infrared contact
coercivity, same-top physical transfer estimate, or interacting continuum
construction is obtained here. The truncations are approximations to an
observable covariance, not positive replacement transfer operators. No
convergence or unique infinite resummation of the formal series is asserted.

\clearpage
\subsection{Diagnostics, preservation, and claim scope}

The new diagnostic is
\begin{center}\small
\path{verify_ym_f16_v17_local_Stein.py}.
\end{center}
It tests the finite Gaussian divergence recursion with correlated rational
precisions, normalized perturbative moment and source-covariance coefficients,
measured Haar jets at both endpoint weights, the cutoff integration identity,
and the mixed source-norm and order bookkeeping. All 504 exact assertions
passed; the report records their families and the distinction between
fixtures and analytic proofs.
The fixtures are not native lattice simulations, enclosures of analytic
constants, or proof-assistant verification of the full-volume estimates.

The cumulative V16 is preserved byte for byte except for its title version
and insertion of this exact appendix before the final document-ending
command. The accompanying package records hashes, predecessor recovery,
actual diagnostic reruns, compilation and rendered-page inspection. No
previous bibliography string or inherited theorem is silently corrected.
The available V4--V16 diagnostics were rerun unchanged and passed 9,467,
89,339, 906, 311, 609, 990, 608, 463, 1,201, 222, 622, 434, and 1,077
assertions, respectively. Unavailable diagnostics are not represented as
rerun. The next scheduled five-version companion checkpoint is V20.

\begin{firewall}
V17 proves fixed-order polynomial expectation and corrected-source covariance
expansions under the full native fixed-bridge law, with entrywise errors
uniform in ambient volume and duration. It identifies the leading native
coefficient in operator norm on every prescribed polynomially growing source
bank and quantifies the delocalization required of a remaining coherent
failure. The proof uses V16's native moments and an explicit Gaussian
polynomial divergence solver, not a native sampler gap. No unrestricted
full-matrix operator bound, summable connected remainder, global auxiliary
gap, physical transfer gap or interacting continuum Yang--Mills mass gap
is established by this appendix.
\end{firewall}
% END F16 V17 APPEND

% BEGIN F16 V18 APPEND -- 2026-09-17
% Insert immediately before the final document-ending command in V17.
% The predecessor body is unchanged; only its cumulative title becomes V18.

\clearpage
\section{V18: exterior screening and full-history trace asymptotics}
\label{f16v18:context}

V17 controls polynomial tests but leaves its actual exterior messages unestimated
at arbitrary order. Repeating its entrywise expansion would not supply a
summable remainder. We instead allow an \emph{arbitrary square-integrable
exterior probe} in the native Stein identity. No derivative of that probe is
taken. Exponential localization of the Gaussian divergence solver pays the
coordinates on which integration by parts is not allowed. This gives a native
$L^2$ estimate for normalized conditional means, not just for finitely many
polynomial exterior tests.

The resulting posterior return is smaller than any prescribed algebraic order
in coupling \emph{in trace density}, with a logarithmic collar. A second step
uses separated, finite space--time blocks, rather than entire-time spatial
cylinders. The new message estimate pays their off-block covariance; V17 pays
their finite source banks; positivity pays the discarded corridors. Together
these prove a full-matrix normalized trace-norm expansion with the same
coefficients as V10, at unrestricted ambient volume and duration. This is not
a dimension-independent operator-norm expansion or exponential clustering at
fixed coupling. The distinction remains explicit throughout.

\subsection{Dependency audit and the norm selected by the new argument}

The supplied V17 appendix and its audit were read in full. The relevant V16
moment proof and the earlier exact source and posterior interfaces were checked
against the cumulative source. Two scoped semantic archive searches returned
no results. The mounted inventories were then searched, and identified
passages read: f14 V93's finite-volume Stein estimate, f15 V14's conditional-copy
posterior transfer, and the monograph's Wilson Markov-blanket and boundary
mediation statements. Those general mechanisms are credited, not presented
as missing from the archive. This is a targeted dependency review, not
independent recertification of the complete historical proofs.

V16's full-native local moments, V17's polynomial divergence solver and
finite-order expectation coefficients, V10's linked-coefficient locality,
and V7's exact Haar source synthesis are the analytic inputs. We do not assume
a native gap, a typical deterministic exterior window, or the desired
connected-correlation estimate. The new statements below are deductions from
these identified inputs. The next scheduled companion checkpoint remains V20.

The wider-literature memorandum correctly keeps summable interactions and the
subsequent global functional inequality as distinct requirements. The other
memorandum separates order gain from locality cost. The following proof does
not replace either requirement by an analogy. For context, M.~Lohmann,
\href{https://arxiv.org/abs/1411.1107}{\emph{Single Scale Cluster Expansions with
Applications to Many Boson and Unbounded Spin Systems}}, studies exponential
clustering with a massive Gaussian part and controlled remaining interactions.
Its abstract and introduction were reviewed, not its complete cluster theorem.
A.~Abdesselam, A.~Procacci and B.~Scoppola,
\href{https://arxiv.org/abs/0901.4756}{\emph{Clustering Bounds on N-Point
Correlations for Unbounded Spin Systems}}, gives related cluster-expansion
context; only the primary abstract was screened. No theorem from either paper
is imported into the compact Wilson law. In particular a weighted norm on
\emph{formal coefficients} is not a verified cluster norm for the full native
remainder. The conclusion here is a different norm, proved directly.

\subsection{Anchored polynomial norms and the inherited Gaussian solver}

Write $\varepsilon=\gamma^{-1/2}$, $\mathcal N=B(n+1)$ and retain
\[
 P=P_{\gamma,\mathbf y},\qquad Q=N(m,C),\qquad C=H^{-1},\qquad
 \xi=\zeta-m,\qquad \max_{i,e}|y_{i,e}|\le r.
\]
The native probability is the complete unmoving fast law of V5 and V16. Its
principal-chart domains, Haar powers, and two dressed endpoints are unchanged.
At fixed rescaled bridge data the Gaussian mean and precision do not depend
on $\varepsilon$. Expectations without a subscript in this appendix are under
$P$. Scalar-coordinate labels are denoted by $i,j$; $c(i)$ is their cell.

Use the ambient cell distance
\[
 d_*((t,b),(u,c))=|t-u|+d_{\mathbb T_m^3}(b,c),
\]
where spatial distance is nearest-neighbour torus distance. A fixed integer
$\rho\ge1$ bounds the diameter of each primitive carrier, the range of $H$,
and the carrier radius of each ordinary source or primitive error. Increasing
$\rho$ only changes fixed constants. Gaussian inverse locality and the port
inverse locality give one $\vartheta>0$ for which
\begin{equation}
 \sup_i\sum_j |C_{ij}|e^{\vartheta d_*(c(i),c(j))}<\infty.
 \label{f16v18:weighted-C}
\end{equation}
We choose $\vartheta$ smaller if necessary than the exponent in V10's
coefficient-locality theorem and the port-lift exponent. It is fixed
independently of all expansion orders, volumes and durations.

Fix an anchor cell $a$. For a monomial $\xi^{\mathbf k}$ put
\[
 R_a(\mathbf k)=\max_{i:k_i>0}d_*(a,c(i)),\qquad R_a(0)=0.
\]
For a polynomial $f=\sum f_{\mathbf k}\xi^{\mathbf k}$ and a polynomial vector
$g=(g_j)$ define
\begin{align}
 \|f\|_{a,\vartheta}
  &=\sum_{\mathbf k}|f_{\mathbf k}|e^{\vartheta R_a(\mathbf k)},\notag\\
 \|g\|_{a,\vartheta,1}
  &=\sum_{j,\mathbf k}|g_{j,\mathbf k}|
       e^{\vartheta\max\{R_a(\mathbf k),d_*(a,c(j))\}}.
 \label{f16v18:anchored-norms}
\end{align}
These are weighted coefficient norms, not native random norms. They charge
actual coefficient sums. No claim that arbitrary extensive polynomials have
bounded norms in \eqref{f16v18:anchored-norms} is made.

\begin{proposition}[Anchored propagation and the cost of unused coordinates]
\label{f16v18:anchored-solver}
For the \emph{same} polynomial map $\mathcal S$ as in
\eqref{f16v17:solver-recursion}, every polynomial of degree at most $d$ obeys
\begin{equation}
 \delta_Q\mathcal Sf=f-\mathbb E_Qf,\qquad
 \|\mathcal Sf\|_{a,\vartheta,1}\le K_{d,\vartheta}\|f\|_{a,\vartheta}.
 \label{f16v18:solver-weighted}
\end{equation}
The same operators $\mathcal T_q=-\nabla V_q\cdot\mathcal S$ as V17 obey
\begin{equation}
 \deg\mathcal T_qf\le d+q,\qquad
 \|\mathcal T_qf\|_{a,\vartheta}\le C_{d,q,r}\|f\|_{a,\vartheta}.
 \label{f16v18:T-weighted}
\end{equation}
Let $I$ be any union of complete coordinate cells and put
$L=d_*(a,I^c)$, with $L=\infty$ if $I^c$ is empty. If $g$ has bounded degree,
then, for every fixed finite $p$,
\begin{equation}
 \left\|\sum_{j:c(j)\notin I}
       [(H\xi)_jg_j-\partial_jg_j]\right\|_{L^p(P)}
       \le C_{d,p,r}e^{-\vartheta L}\|g\|_{a,\vartheta,1}.
 \label{f16v18:outside-divergence}
\end{equation}
All constants are independent of $I,B,n$.
\end{proposition}
\begin{proof}
Only the weighted estimates are new here; the exact divergence identity is
V17's finite algebra. For a degree-$d$ monomial with first selected index $i_1$
and remaining product $h$, the first vector of that recursion is
$C_{\cdot i_1}h$. Its radius after adding component $j$ is at most
$R_a(\mathbf k)+d_*(c(i_1),c(j))$. Equation \eqref{f16v18:weighted-C}
therefore bounds its weighted coefficient sum by a constant times the
original monomial weight. Each contraction removes two indices and cannot
increase the remaining radius. Its coefficient is bounded by $\sup|C_{ij}|$.
The terminating degree induction of V17 gives
\eqref{f16v18:solver-weighted}. Constants are sent to zero.

The polynomial $\partial_jV_q$ contains only the fixed number of original
interaction terms meeting $j$. Its variables lie within radius $\rho$ of
$c(j)$, its degree is at most $q+1$, and its coefficient norm after the bounded
mean translation is at most $C_{q,r}$. Multiplication with $g_j$ increases
the radius in the vector norm by at most $\rho$. Summing $j$ proves
\eqref{f16v18:T-weighted}. The measured coefficients $V_q$ include the original
Haar powers; no new local potential is substituted.

For outside components, the unweighted coefficient sum is at most
$e^{-\vartheta L}\|g\|_{a,\vartheta,1}$. Multiplication by $(H\xi)_j$ costs the
uniform absolute row bound for $H$; differentiation costs the fixed degree.
V17's polynomial $L^p$ estimate under the actual native law, followed by
Minkowski, gives \eqref{f16v18:outside-divergence}. The same proof bounds
$\sum_{j\notin I}(\partial_jV_q)g_j$ by
$C_{d,q,p,r}e^{-\vartheta L}\|g\|_{a,\vartheta,1}$. This last bound will pay
outside drift terms as well.
\end{proof}

The radius norm is chosen for this exterior test. It is not a norm forcing a
connected tree between two arbitrary source marks. No such conclusion is
built into its definition.

\subsection{Stein's identity tested against arbitrary native exterior functions}

Let $\mathcal F_O=\sigma(\zeta_{I^c})$. The test $h\in L^2(P,\mathcal F_O)$ may
depend on \emph{every} exterior coordinate and can be nonsmooth. Its norm is
under the actual exterior marginal of $P$, including all its conditional
partition functions. It need not be a polynomial, a cylinder function, or
bounded independently of the regulator.

\begin{theorem}[An exterior-tested native identity]
\label{f16v18:tested-Stein}
For every fixed degree $d$ and order $N\ge0$, every polynomial vector $g$ of
degree at most $d$, and every $h\in L^2(P,\mathcal F_O)$,
\begin{equation}
 \left|\mathbb E[h\,\delta_Qg]
       +\sum_{q=1}^{N}\varepsilon^q
                    \mathbb E[h\,\nabla V_q\cdot g]\right|
 \le C_{d,N,r}\bigl(\varepsilon^{N+1}+e^{-\vartheta L}\bigr)
                  \|h\|_2\|g\|_{a,\vartheta,1}.
 \label{f16v18:tested-identity}
\end{equation}
The coupling threshold depends only on $d,N,r$, not on the size, shape or
number of exterior coordinates. The empty sum convention applies for $N=0$.
\end{theorem}
\begin{proof}
Use V17's group-local cutoff
$\chi_j=\chi(\varepsilon\zeta_{g(j)})$, where $\chi$ is one on radius $1/2$
and zero beyond radius one. For $j\in I$ the whole group $g(j)$ is in $I$.
The factor $h$ is constant in that group, so ordinary coordinate integration
by parts, at fixed remaining coordinates, differentiates no part of $h$.
Initially take bounded measurable $h$. With the exact measured potential
$V_\varepsilon=U_\varepsilon-\log J_\varepsilon^f$, the identity is
\begin{align}
 \mathbb E[h\delta_Qg]
 ={}&-\sum_{j\in I}\mathbb E[
       h\chi_jg_j(\partial_jV_\varepsilon-(H\xi)_j)]\notag\\
 &+\sum_{j\in I}\mathbb E[h(\partial_j\chi_j)g_j]\notag\\
 &+\sum_{j\in I}\mathbb E[h(1-\chi_j)
                ((H\xi)_jg_j-\partial_jg_j)]\notag\\
 &+\sum_{j\notin I}\mathbb E[h((H\xi)_jg_j-\partial_jg_j)].
 \label{f16v18:exact-tested-cut}
\end{align}
Here $j\in I$ abbreviates $c(j)\in I$. The final line is an algebraic
remainder, not integration by parts on the forbidden exterior coordinates.
This distinction is the reason no regularity of $h$ is required.

The classical word-drift remainder through order $N$ has $L^4(P)$ norm at
most $C_{N,r}\varepsilon^{N+1}$ per coordinate, by the all-real Taylor
bounds and V16's local moments, exactly as in V17. On $\chi_j$'s support the
only differentiated Haar factor is a smooth function of its own group inside
radius one. Its original weight, including the endpoint weight $1/2$, is
retained. Its measured drift remainder has the same bound. Thus
\[
 \|\chi_j g_j[\partial_jV_\varepsilon-(H\xi)_j
                 -\sum_{q=1}^N\varepsilon^q\partial_jV_q]\|_2
       \le C_{d,N,r}\varepsilon^{N+1}\|g_j\|_{\rm pol}.
\]
Pair with $h$ by Cauchy--Schwarz. No moment bound under a new conditional or
intermediate coupling is being used.

For every fixed-degree polynomial $p$, V16's one-cell exponential moment and
H\"older imply
\[
 \|1_{\{|\zeta_{g(j)}|\ge(2\varepsilon)^{-1}\}}p\|_2
       \le C_{d,r}e^{-c/\varepsilon^2}\|p\|_{\rm pol}.
\]
This is the $L^2$ version of V17's local cut-tail estimate: use the fourth
moment of $p$ and the fourth root of the event probability. It pays the
second and third lines of \eqref{f16v18:exact-tested-cut}, and removal of
$\chi_j$ from the finite drift polynomials. Summation costs
$\sum_j\|g_j\|_{\rm pol}$, not the total number of coordinates.

Finally \eqref{f16v18:outside-divergence} and the outside drift bound from
its proof pay the last line and all $j\notin I$ polynomial drift terms by
$C_{d,N,r}e^{-\vartheta L}\|h\|_2\|g\|_{a,\vartheta,1}$. Absorb the
coordinate-local flat term into the fixed-order $\varepsilon^{N+1}$ budget.
This proves \eqref{f16v18:tested-identity} for bounded measurable $h$.
All other factors in its terms are in $L^2(P)$ by the estimates just given.
Truncating an arbitrary exterior $L^2$ function and passing in $L^2$ proves
the stated theorem. No all-group cutoff changes $P$ at any point.
\end{proof}

\subsection{Native conditional means, including their true normalizers}

Let $b_k(f)$ be precisely V17's coefficient recursion
\eqref{f16v17:coefficient-recursion}. In particular $b_0(f)=\mathbb E_Qf$;
no coefficient is fitted to the exterior. These numbers can depend on the
entire prescribed bridge history through $m$ and $C$, but that history is a
deterministic parameter, not the random fast exterior.

\begin{theorem}[All fixed orders of exterior conditional screening]
\label{f16v18:conditional-polynomials}
For every fixed polynomial degree $d$ and order $N\ge0$,
\begin{equation}
 \left\|\mathbb E[f\mid\mathcal F_O]
                  -\sum_{k=0}^N\varepsilon^k b_k(f)\right\|_2
 \le C_{d,N,r}\bigl(e^{-\vartheta L}+\varepsilon^{N+1}\bigr)
                         \|f\|_{a,\vartheta}.
 \label{f16v18:conditional-expansion}
\end{equation}
Consequently
\begin{equation}
 \|\mathbb E[f\mid\mathcal F_O]-\mathbb E f\|_2
 \le C_{d,N,r}\bigl(e^{-\vartheta L}+\varepsilon^{N+1}\bigr)
                         \|f\|_{a,\vartheta}.
 \label{f16v18:conditional-centered}
\end{equation}
These statements are uniform in ambient volume and duration, for sufficiently
large coupling depending only on the displayed fixed degree and order.
They are averaged in the native exterior law, not suprema over deterministic
exterior values.
\end{theorem}
\begin{proof}
Apply Theorem~\ref{f16v18:tested-Stein} to $g=\mathcal Sf$. Its divergence
identity gives, for every exterior $h$,
\begin{equation}
 \mathbb E[h f]=(\mathbb E h)\mathbb E_Q f
      +\sum_{q=1}^N\varepsilon^q\mathbb E[h\mathcal T_q f]
      +O\!\left((e^{-\vartheta L}+\varepsilon^{N+1})
                       \|h\|_2\|f\|_{a,\vartheta}\right).
 \label{f16v18:tested-recursion}
\end{equation}
For $N=0$ this is the assertion in dual form. Induct on $N$, applying the
order-$N-q$ formula to each $\mathcal T_q f$, with the \emph{same} exterior
$h$, the same set $I$, and the same anchor. Its weighted coefficient norm is
bounded by \eqref{f16v18:T-weighted}. Multiplying by $\varepsilon^q$ makes
every algebraic remainder order $N+1$; its spatial error is bounded by a
constant times $e^{-\vartheta L}$ because $\varepsilon\le1$. There are only
finitely many terms at each fixed order. Their coefficients are exactly
$b_k$, not a new conditional recursion. Thus
\[
 |\mathbb E[h(f-\sum_{k=0}^N\varepsilon^k b_k(f))]|
 \le C_{d,N,r}(e^{-\vartheta L}+\varepsilon^{N+1})
                         \|h\|_2\|f\|_{a,\vartheta}.
\]
Hilbert-space duality on $L^2(P,\mathcal F_O)$ proves
\eqref{f16v18:conditional-expansion}. Apply the orthogonal projection
$I-\mathbb E_P$ on that exterior space to obtain
\eqref{f16v18:conditional-centered}; all subtracted coefficients are constants.
This proof permits coefficients involving distant Gaussian variables inside
$\mathcal T_q f$: their anchored weights, not a false finite support claim,
pay the exterior tail at each step.
\end{proof}

This extends the class of \emph{tests}, not merely the polynomial order in
V17. The dual test can be the sign of a conditional message, or its centered
normalized value, and may inspect an arbitrarily large exterior. No regularity,
polynomial approximation, or reference probability for that test is required.
The factor $\varepsilon^{N+1}$ remains when $L$ grows at a fixed coupling;
this is not a proof of strong mixing as $L\to\infty$ with $\gamma$ fixed.

\begin{proposition}[Application to ordinary scores and Haar residuals]
\label{f16v18:source-screening}
Let $J_{\gamma,\alpha}=\partial_\alpha U_\gamma$ be an ordinary common-source
score, anchored at $a(\alpha)$, and let $r_{\gamma,\alpha}$ be the exact
Haar-corrected residual of V7. For every fixed $N\ge1$,
\begin{align}
 \|\mathbb E[J_{\gamma,\alpha}\mid\mathcal F_O]
                          -\mathbb E J_{\gamma,\alpha}\|_2
   &\le C_{N,r}(e^{-\vartheta L}+\varepsilon^{N+1}),\notag\\
 \|\mathbb E[r_{\gamma,\alpha}\mid\mathcal F_O]
                          -\mathbb E r_{\gamma,\alpha}\|_2
   &\le C_{N,r}(\varepsilon e^{-\vartheta L}+\varepsilon^{N+1}).
 \label{f16v18:source-centered}
\end{align}
The second bound also holds for every component of the primitive bank
$e^{\rm prim}=(e^{\rm br},e^{\rm pt})$.

If $J_{\gamma,\alpha}\sim\sum_{k\ge0}\varepsilon^kJ_{\alpha,k}$, define
$j_{\alpha,q}=\sum_{k=0}^q b_{q-k}(J_{\alpha,k})$. Then also
\begin{equation}
 \left\|\mathbb E[J_{\gamma,\alpha}\mid\mathcal F_O]
                         -\sum_{q=0}^N\varepsilon^qj_{\alpha,q}\right\|_2
       \le C_{N,r}(e^{-\vartheta L}+\varepsilon^{N+1}).
 \label{f16v18:source-expansion}
\end{equation}
The first coefficients are $j_{\alpha,0}=g_\alpha$ and
$j_{\alpha,1}=\mathbb E_Q r_{\alpha,1}$, as in V11's deterministic reference.
\end{proposition}
\begin{proof}
The ordinary score coefficients have fixed local carriers and bounded
fixed-degree coefficient norms. Their complete-chart Taylor remainder in
$L^2(P)$ is $O(\varepsilon^{N+1})$, by the same inserted-word bounds and
V16 moments used in V17. The primitive errors have the same properties but
start at order one. Their weighted coefficient norms are therefore uniformly
bounded at each fixed order. For the full Haar residual, the spatially
exponentially summable port-lift column preserves these weighted norms after
decreasing $\vartheta$ once, independently of order. Its exact native $L^2$
jet remainder remains that of Proposition~\ref{f16v17:source-jets}.

Apply Theorem~\ref{f16v18:conditional-polynomials} to each coefficient,
through the remaining order $N-k$, and use conditional $L^2$ contraction on
the native jet remainder. The spatial error is $O(e^{-\vartheta L})$ for
the ordinary score and $O(\varepsilon e^{-\vartheta L})$ for a bank starting
at order one. Centering is a contraction and proves
\eqref{f16v18:source-centered}. Collecting coefficients gives
\eqref{f16v18:source-expansion}.

Coefficient ownership is unchanged: V17's $b_k$ are normalized native
coefficients. In particular the direct formula for the first score correction
is $\mathbb E_QJ_{\alpha,1}-\operatorname{Cov}_Q(J_{\alpha,0},V_1)$.
The exact Haar identity $\mathbb E J_{\gamma,\alpha}
=g_\alpha+\mathbb E r_{\gamma,\alpha}$ identifies it with
$\mathbb E_Qr_{\alpha,1}$ by uniqueness of fixed-order coefficients.
The half-bracket term in $r_{\alpha,1}$ and the measure correction are not
removed. No derivative of a conditional approximation is taken.
\end{proof}

\subsection{The actual spatial-cylinder posterior has arbitrary-order trace accuracy}

Return first to the literal separated-cylinder setting of V12. Source marks
are interior to their cylinders, with ordinary score carriers strictly inside,
and their anchor distance to the exterior is at least $L$. There are no
repeated selected marks. The common exterior posterior retains
\[
 d\bar P(q)=Z_\gamma^{-1}1_{\mathcal P'_O}J^f_{\gamma,O}(q)
      e^{-U_{\gamma,O}(q;\mathbf y)}
                       \prod_sZ_{\gamma,D_s}(q;\mathbf y)\,dq.
\]
No conditional partition function in this expression is replaced. Let
$\mathbf j$ be the same vector of derivatives of $-\log Z_{\gamma,D_s}$,
$k$ its number of coordinates, $\mathsf R_{\rm post}=\operatorname{Cov}_{\bar P}
(\mathbf j)$, and $\mathsf L_D$ and $\chi$ be V12's objects.

\begin{theorem}[All-order posterior return and unchanged cylinder gluing]
\label{f16v18:posterior}
For every fixed $N\ge1$, at sufficiently large coupling depending only on
$N,r$, and for all ambient volumes and durations,
\begin{equation}
 \operatorname{Tr}\mathsf R_{\rm post}
     \le C_{N,r}k\bigl(e^{-2\vartheta L}+\gamma^{-(N+1)}\bigr),\qquad
 \|\chi\mathsf M_\gamma\chi-\mathsf L_D\|_{S_1}
     \le C_{N,r}k\bigl(e^{-2\vartheta L}+\gamma^{-(N+1)}\bigr).
 \label{f16v18:posterior-trace}
\end{equation}
Every time lag is included at once. In particular, for $k\le C\mathcal N$,
$L\ge(N+1)(2\vartheta)^{-1}\log\gamma$ gives both normalized trace bounds
$C_{N,r}\gamma^{-(N+1)}$. Such a collar must fit in the chosen cylinders;
no bound on their volume relative to coupling is otherwise imposed.
\end{theorem}
\begin{proof}
For an interior mark, exact normalized differentiation gives
$j^D_{\gamma,\alpha}=\mathbb E[J_{\gamma,\alpha}\mid\zeta_{D^c}]$.
No exterior-only action term is differentiated. Separation makes this the
same function as the conditional mean under the one common exterior $O$;
the other interiors have no interaction with this interior after $O$ is fixed.
Apply the first bound of \eqref{f16v18:source-centered} to each scalar
message. The law of that message under $P$ is precisely its law under the
actual $\bar P$. Sum the squared estimates and use
$(u+v)^2\le2u^2+2v^2$. The trace is the sum of those variances, proving
the first inequality. No covariance between different messages is claimed
zero.

The inherited exact identity is
$\chi\mathsf M_\gamma\chi-\mathsf L_D
=-\operatorname{Off}_{>1}\mathsf R_{\rm post}$.
The matrix on the right before extraction is positive, so its trace norm is
its trace. V7's duration-independent bound for that extraction is four.
This proves the second inequality. The single-time endpoint normalizers
remain in the original history weight; their nonadjacent mixed derivatives
vanish for the original reason, not because they were removed.
\end{proof}

These are actual exterior-message estimates, unlike a polynomial test of a
Gaussian exterior. They remove the $1/\gamma$ \emph{floor} from V16's posterior
trace bound at a chosen logarithmic collar. They do not identify V11's
conditional covariance for each deterministic exterior, or prove a uniform
native conditional gap. For a bank of size $k\le\gamma^p$, trace domination
also gives an unnormalized operator bound $C_{p,s,r}\gamma^{-s}$ for any
fixed $p,s>0$, by choosing $N+1\ge p+s$ and
$L\ge(p+s)(2\vartheta)^{-1}\log\gamma$. The constants and threshold depend
on the chosen orders. The entire unrestricted source bank can still be much
larger than $\gamma^p$.

\subsection{Finite space--time blocks and their exact connected remainder}

For the full-matrix conclusion, work first with the primitive errors, which
have finite carriers. A full lifted residual is spatially nonlocal and must
not be declared independent of another block just because its anchor lies
there. Write
\[
 e=e^{\rm prim},\qquad K_\gamma=\operatorname{Cov}_P(e),\qquad
 r=Te,\qquad T=(I,\mathsf P^{\mathsf T}),\qquad \|T\|\le C.
\]
There are at most $30\mathcal N$ primitive scalar components: $9\mathcal N$
common bridge errors and $21Bn$ port errors. The matrix $T$ is the original
Gaussian lift synthesis, independent of coupling at fixed geometry.

\begin{proposition}[Primitive coefficient inputs and a block return identity]
\label{f16v18:primitive-blocks}
The normalized coefficients $K_q$ of $K_\gamma$ satisfy, at every fixed
order $q\ge2$,
\begin{equation}
 \|K_q\|_{\mathrm{loc},\vartheta}\le C_{q,r},\qquad
 \sup_{u,v}\left|K_{\gamma,uv}-\sum_{q=2}^K\varepsilon^q(K_q)_{uv}\right|
                  \le C_{K,r}\varepsilon^{K+1}.
 \label{f16v18:primitive-inputs}
\end{equation}
Here the local norm is the maximum weighted absolute row and column sum on
primitive anchors. Moreover $G_q=TK_qT^{\mathsf T}$ exactly.

Let $(I_s)$ be separated free cell sets such that no native interaction meets
two sets, and select primitive banks $\mathcal A_s$ whose carriers lie inside
$I_s$ and whose anchor depth is at least $L$. Write $\Pi_s$ for their disjoint
coordinate projections, $\Pi=\sum_s\Pi_s$, and
$K_{\rm blk}=\sum_s\Pi_sK_\gamma\Pi_s$. Then
\begin{equation}
 \|\Pi K_\gamma\Pi-K_{\rm blk}\|_{S_1}
 \le C_{N,r}|\mathcal A|
          \bigl(\varepsilon^2e^{-2\vartheta L}
                        +\varepsilon^{2N+2}\bigr),\qquad N\ge1,
 \label{f16v18:block-return}
\end{equation}
where $\mathcal A=\bigcup_s\mathcal A_s$. The blocks may be finite in time;
they need not be V11's cylinders.
\end{proposition}
\begin{proof}
V17's native jet statement explicitly includes both primitive banks. Applying
its expectation expansion to their two products and means gives the entry
estimate in \eqref{f16v18:primitive-inputs}, by the same finite coefficient
calculation. Its coefficients are the connected Gaussian cumulants of the
primitive jets with the measured $V_q$. The linked-pairing proof of V10
applies to these coefficients without the port-lift edges: every primitive
carrier is local, covariance rows decay exponentially, and every surviving
Gaussian contraction connects the two marks. Summing a spanning tree bounds
the weighted row at one fixed exponent and each fixed order. This is a
specialization of that inherited argument, not a new native cluster
expansion. Since $r=Te$ exactly and $T$ is coupling independent, uniqueness
of the finite-dimensional coefficients gives $G_q=TK_qT^{\mathsf T}$.

Condition on $O=(\bigcup_s I_s)^c$. The product coordinate reference factors
and the local native action imply exact conditional independence of the
$I_s$. This includes the fractional Haar powers and the original dressed
endpoints, which are cell local. Put
$m_s=\mathbb E[e_{\mathcal A_s}\mid\mathcal F_O]$ and
$R=\operatorname{Cov}_P((m_s)_s)$. Total covariance gives
\[
 \Pi K_\gamma\Pi=\bigoplus_s\mathbb E\operatorname{Cov}
                      (e_{\mathcal A_s}\mid\mathcal F_O)+R.
\]
The $s$ diagonal block of this equality is the \emph{unconditional} covariance
$\Pi_sK_\gamma\Pi_s$, not a chosen boundary covariance. Hence exactly
\begin{equation}
 \Pi K_\gamma\Pi-K_{\rm blk}=R-\sum_s\Pi_sR\Pi_s.
 \label{f16v18:pinching-return}
\end{equation}
Each message $m_s$ agrees with conditioning on $I_s^c$ alone, by the same
finite-range specification. The primitive estimate in
\eqref{f16v18:source-centered} bounds the trace of $R$ by the right side
of \eqref{f16v18:block-return}, up to a fixed factor. Both $R$ and its block
pinching are positive and have the same trace. Their trace-norm difference
is at most twice that trace, proving the claim. No product approximation of
the actual exterior posterior has been used.
\end{proof}

The estimate is genuinely connected: it bounds the covariance \emph{between}
conditionally independent interiors through their actual posterior means.
It remains a trace estimate after summing a bank, not a bound on each
pair's remainder by its separation alone.

\begin{proposition}[A geometric partition with sparse discarded corridors]
\label{f16v18:partition}
For integers $L\ge\rho+1$ and $R\ge16(L+3\rho)$ there are separated free
cell sets $I_s$ and primitive banks $\mathcal A_s$ as above such that
\begin{equation}
 |I_s|\le(4R)^4,\qquad |\mathcal A_s|\le C R^4,\qquad
 |\mathcal A^c|\le C\mathcal N\frac{L+\rho}{R}.
 \label{f16v18:partition-bounds}
\end{equation}
The last complement is within the full primitive bank. Distinct selected
banks have anchor separation at least $L$. Constants are independent of
spatial periods and of the finite time length.
\end{proposition}
\begin{proof}
In each of the three coordinate cycles, and in the finite time path, leave a
direction uncut if its length is at most $4R$. Otherwise partition it into
$k=\lfloor\text{length}/R\rfloor$ consecutive intervals, distributing the
lengths as equally as possible. Each length is between $R$ and $2R$. In a
cycle the wrap-around cut is included. In the time direction only interior
cuts are used; the actual two history endpoints are not artificial cuts.
The products of these intervals, or complete uncut directions, are tiles of
at most $(4R)^4$ cells.

Remove a strip of width $\rho$ on each side of every genuine tile cut to
obtain $I_s$. Its gap from another free set exceeds the interaction range,
so no primitive action term meets two free sets. Select primitive anchors
in the same tile with distance at least $L+3\rho$ from every genuine cut.
Their complete carriers are inside $I_s$, and their distance to $I_s^c$ is
at least $L$. Where a direction is uncut and fills its entire cycle or time
path, it creates no exterior face and no depth condition in that direction.
This also covers the case of one tile filling the entire system; its exterior
is empty.

In a cut direction the number of removed coordinate positions is at most
$C(L+\rho)$ times the number of cuts, which is at most its length divided by
$R$. Taking the union of these strips in four dimensions bounds the discarded
cell count by $C\mathcal N(L+\rho)/R$. There are only a bounded number of
primitive types per anchor cell, including at time endpoints. Increasing the
strip allowance by the fixed carrier radius proves the last assertion of
\eqref{f16v18:partition-bounds}. Source counts per tile give the first two.
A path between selected anchors of different tiles must cross a genuine cut,
so their ambient distance is at least $L$. No translation symmetry of the
probability law is used.
\end{proof}

\subsection{All fixed orders in normalized trace norm on the entire history}

\begin{theorem}[Full-matrix native trace asymptotics at unrestricted volume]
\label{f16v18:full-trace}
For every fixed integer $J\ge2$ and fixed bridge window $r$, there are finite
$C_{J,r},\gamma_{J,r}$ independent of $B,n$ such that, for all admitted bridge
histories and $\gamma\ge\gamma_{J,r}$,
\begin{equation}
 \boxed{
 \frac1{\mathcal N}
 \left\|\Gamma_\gamma-\sum_{q=2}^J\gamma^{-q/2}G_q\right\|_{S_1}
 +\frac1{\mathcal N}
 \left\|\mathsf M_\gamma-\sum_{q=2}^J\gamma^{-q/2}M_q\right\|_{S_1}
       \le C_{J,r}\gamma^{-(J+1)/2}.}
 \label{f16v18:trace-expansion}
\end{equation}
Both matrices are the entire original common-source matrices; no bank is
omitted in the conclusion. There is no volume-dependent coupling-entry
condition. The norm is the sum of singular values divided by
$\mathcal N=B(n+1)$, not an unnormalized operator norm.
\end{theorem}
\begin{proof}
We prove the primitive assertion and then synthesize. All constants in this
proof can depend on fixed $J,r$. Decrease $\vartheta$ so that both anchored
screening and primitive coefficient locality use that exponent. Write
$K^{[J]}=\sum_{q=2}^J\varepsilon^qK_q$. Take
\begin{equation}
 L=\left\lceil\frac{J+1}{\vartheta}\log(\varepsilon^{-1})\right\rceil
                      +\rho+1,\qquad
 R=\lceil\varepsilon^{-2J}\rceil,\qquad K=9J.
 \label{f16v18:trace-scales}
\end{equation}
For sufficiently small $\varepsilon$, independently of $B,n$,
$R\ge16(L+3\rho)$. Apply Proposition~\ref{f16v18:partition}, and use its
coordinate projections $\Pi,\Pi_s$. There are four errors to pay.

\emph{Discarded native coordinates.} Positivity of $K_\gamma$ and V16's
component bound give $\Tr K_\gamma\le C\mathcal N\varepsilon^2$ and
$\Tr[(I-\Pi)K_\gamma]\le C|\mathcal A^c|\varepsilon^2$. The Schatten
Cauchy--Schwarz inequality, applied to the two factors of the positive square
root, gives
\begin{align}
 \|K_\gamma-\Pi K_\gamma\Pi\|_{S_1}
 &\le \Tr[(I-\Pi)K_\gamma]
      +2\sqrt{\Tr(\Pi K_\gamma)\Tr[(I-\Pi)K_\gamma]}\notag\\
 &\le C\mathcal N\varepsilon^2\sqrt{(L+\rho)/R}.
 \label{f16v18:corridor-cost}
\end{align}
If there are no discarded coordinates the error is zero. The last bound
uses $(L+\rho)/R\le1$. This is why a mere probability count is not multiplied
by an arbitrary spectral test function; positivity pays the actual covariance
cross block.

\emph{Exact off-block native return.} Apply
\eqref{f16v18:block-return} with screening order $N=J$. It gives
\begin{equation}
 \|\Pi K_\gamma\Pi-\sum_s\Pi_sK_\gamma\Pi_s\|_{S_1}
   \le C\mathcal N\bigl(\varepsilon^2e^{-2\vartheta L}
                                     +\varepsilon^{2J+2}\bigr).
 \label{f16v18:offblock-cost}
\end{equation}
The conditional covariance is not substituted for the unconditional block;
\eqref{f16v18:pinching-return} paid their distinction.

\emph{Finite-bank native expansion.} Put $k_s=|\mathcal A_s|\le CR^4$.
The order-$K$ entry error in \eqref{f16v18:primitive-inputs} has block operator
norm at most $C k_s\varepsilon^{K+1}$ and trace norm at most
$C k_s^2\varepsilon^{K+1}$. For the finite sum of coefficients with
$J<q\le K$, their uniform absolute row bounds give a trace error at most
$C k_s\varepsilon^{J+1}$. Since $\sum_s k_s\le C\mathcal N$,
\begin{equation}
 \left\|\sum_s\Pi_s(K_\gamma-K^{[J]})\Pi_s\right\|_{S_1}
     \le C\mathcal N
           \bigl(R^4\varepsilon^{K+1}+\varepsilon^{J+1}\bigr).
 \label{f16v18:bank-cost}
\end{equation}
This invokes V17 at a \emph{fixed} order $9J$, in the unrestricted ambient
law. It does not expand a conditional kernel on a fixed large exterior and
does not require a bound exponential in $R^4$.

\emph{Completion of the coefficient matrices.} Removing the discarded rows
and columns of $K_q$ changes a matrix of rank at most $2|\mathcal A^c|$;
its operator norm is at most twice the original norm. This costs at most
$C_q|\mathcal A^c|$ in trace norm. Between distinct selected banks all
anchors have distance at least $L$. The weighted row and column bounds for
$K_q$ therefore give operator norm at most $C_qe^{-\vartheta L}$ for the
remaining off-block part, and trace norm at most
$C_q\mathcal N e^{-\vartheta L}$. Sum the finitely many coefficients to obtain
\begin{equation}
 \left\|K^{[J]}-\sum_s\Pi_sK^{[J]}\Pi_s\right\|_{S_1}
      \le C\mathcal N\varepsilon^2
                     \bigl((L+\rho)/R+e^{-\vartheta L}\bigr).
 \label{f16v18:coefficient-completion}
\end{equation}
There is no positivity assumption on the higher $K_q$ in this argument.

Combining these four bounds gives the explicit finite-scale estimate
\begin{align}
 \frac{\|K_\gamma-K^{[J]}\|_{S_1}}{\mathcal N}
 \le C\bigl[&\varepsilon^2\sqrt{(L+\rho)/R}
       +\varepsilon^2 e^{-\vartheta L}
       +\varepsilon^{2J+2}\notag\\
       &+R^4\varepsilon^{K+1}+\varepsilon^{J+1}\bigr].
 \label{f16v18:four-costs}
\end{align}
For the choices \eqref{f16v18:trace-scales},
$e^{-\vartheta L}\le\varepsilon^{J+1}$,
$R^4\varepsilon^{9J+1}\le16\varepsilon^{J+1}$, and
\[
 \varepsilon^2\sqrt{(L+\rho)/R}
     \le C\varepsilon^{J+2}\sqrt{1+\log(\varepsilon^{-1})}
     \le C\varepsilon^{J+1}.
\]
This proves the primitive trace estimate. The geometric construction already
covers short directions and small tori; no nonexistent large block is assumed.
All thresholds are fixed-order thresholds, not functions of the ambient size.

Finally the exact synthesis and coefficient identity give
\[
 \Gamma_\gamma-\sum_{q=2}^J\varepsilon^qG_q
       =T(K_\gamma-K^{[J]})T^{\mathsf T}.
\]
Its trace norm is at most $\|T\|^2$ times the primitive trace norm. The
original Haar-contact identity gives
$\mathsf M_\gamma=-\operatorname{Off}_{>1}\Gamma_\gamma$ and
$M_q=-\operatorname{Off}_{>1}G_q$. That extraction has trace-norm bound four
for every duration. These two facts prove \eqref{f16v18:trace-expansion}.
Neither the measure nor the physical transfer has been changed by the proof
partition.
\end{proof}

\subsection{What the trace expansion establishes, and the remaining edge}

At its first order Theorem~\ref{f16v18:full-trace} gives
\begin{equation}
 \boxed{
 \frac{\|\gamma\Gamma_\gamma-G_2\|_{S_1}
              +\|\gamma\mathsf M_\gamma-M_2\|_{S_1}}{\mathcal N}
                           \le C_r\gamma^{-1/2}.}
 \label{f16v18:leading-trace}
\end{equation}
Thus the leading native coefficient is identified with an
$o(\gamma^{-1})$ normalized trace error on the \emph{entire} history, not
only entrywise or on a chosen small source bank. This replaces V16's
unidentified $O(\gamma^{-1})$ trace budget by a genuine fixed-order
asymptotic expansion in that norm. V17's stronger unnormalized bank estimates
remain separate and are used in the proof.

For every fixed $\delta>0$, the number of singular values of
$\gamma\Gamma_\gamma-G_2$ exceeding $\delta$ is at most
$C_r\mathcal N\gamma^{-1/2}/\delta$, and likewise for memory. This follows
by summing those singular values; it is a statement about the error matrix,
not an assertion that two noncommuting matrices share eigenvectors. A
vanishing proportion of exceptional directions can still include a single
large eigenvalue. Neither \eqref{f16v18:leading-trace} nor the analogous
higher-order formula bounds the unrestricted operator norm without its
volume factor.

The intermediate posterior result also has a more limited meaning than an
all-boundary conditional theorem. Large fast exteriors have been integrated
under the native marginal, not prohibited, and every message normalizer is
kept. But the bound contains an arbitrarily high \emph{fixed-order} coupling
floor at fixed $L$. No order-dependent constants have been controlled as
$N\to\infty$ with coupling fixed. Thus no assertion of strong mixing at all
distances, convergence of a series, or a unique infinite resummation follows.

The next unresolved question is an edge rather than an average: control the
remaining exceptional coherent directions, or prove a summable connected
native remainder at fixed sufficiently large coupling. The present block
argument cannot answer it by repeating the trace estimate. In
\eqref{f16v18:corridor-cost} it uses total covariance traces to pay discarded
rows, and in \eqref{f16v18:offblock-cost} it sums scalar message energies.
Those are the two locations where a directional, instead of trace, bound
would be needed. A source-resolved capacity/repair estimate or an independently
proved native functional inequality remains a distinct possible route.

V4's original full/core probability has not been redefined or given an
all-order boundary expansion. This appendix concerns the full native law
and its exact conditional messages. The prescribed bridge window also
remains: no integration over unrestricted physical bridges, root Gauss
restriction, physical-time calibration, infrared contact coercivity,
same-top full-transfer estimate, or interacting continuum construction is
supplied by a trace asymptotic. The auxiliary sampler gap is not used and
is not proved.

\subsection{Diagnostics, preservation, and claim scope}

The new diagnostic is
\begin{center}\small
\path{verify_ym_f16_v18_exterior_trace.py}.
\end{center}
It checks the weighted polynomial solver and its covariance path bounds.
Further tests cover the Stein identity on permitted coordinates, conditional
messages and posterior pinching on finite positive Markov laws, sparse
partitions including short directions, the four trace costs, and positive
matrix corridor and synthesis identities. All 819 exact assertions passed. Its report distinguishes exact algebraic
fixtures from analytic proofs. In particular a finite probability example is
not a native compact-integral enclosure, and no computation certifies all the
analytic constants or the inherited moment theorem.

The cumulative source changes only V17's title version and inserts this exact
appendix before the final document-ending command. Removing that insertion
and restoring the title recovers the predecessor byte for byte. The package
records the input hashes, preservation check, actual diagnostic runs,
compilation, and rendered-page inspection. The fourteen available V4--V17 diagnostics were rerun unchanged and passed.
Their individual reports and hashes are included. No inaccessible archive
checker or external formalization is described as verified. The new arguments explicitly depend on the identified V16--V17
analytic inputs, not on their finite diagnostic counts.

\begin{firewall}
V18 proves native polynomial identities with exterior tests. It gives
averaged screening estimates for conditional sources, with fixed-order
errors uniform in ambient volume and duration. These give arbitrary-order posterior trace accuracy
with a logarithmic collar. Separated finite space--time blocks then yield
fixed-order asymptotics of the full covariance and memory in trace norm per
cell, including an $o(1/\gamma)$ leading-coefficient error. No unrestricted
operator-norm remainder, fixed-coupling strong mixing or summable connected
remainder, uniform native sampler gap, physical transfer gap, or interacting
continuum Yang--Mills mass gap is established. All original native laws,
normalizers, and source definitions remain unchanged.
\end{firewall}
% END F16 V18 APPEND

% BEGIN F16 V19 APPEND -- 2026-09-17
% Insert immediately before the final document-ending command in V18.
% No inherited proof, normalization, or bibliography string is revised.

\clearpage
\section{V19: a positive shell return and one-sided coherent asymptotics}
\label{f16v19:context}

V18 identifies the full covariance in trace norm per cell. Its two losses at
an unrestricted spectral edge are explicit: discarded corridors are paid by
traces, and exterior messages are summed through their scalar squared norms.
We do not improve those same trace estimates. Instead, every primitive source
gets its own local ball, with no discarded sources. Subtracting the Gram of
its exact exterior conditional means leaves an \emph{exactly finite-range
matrix}. The subtracted Gram is positive; the finite-range matrix need not be.
This distinction gives a one-sided, unnormalized operator estimate at every
fixed order, and an explicit positive object carrying the remaining upper edge.

A second step resolves that object's change across scales through actual
annular innovations. Each scale has a positive finite-range Gram, although
cross-scale terms do not in general vanish. A dyadic synthesis estimate gives
a source-specific sufficient condition for the missing upper edge. The known
finite-order screening bounds pay transport to any prescribed polynomial
radius, but not an infinite sequence of radii at fixed coupling. We keep that
remaining hypothesis separate from the results proved here.

\subsection{Dependencies, non-duplication, and the literature distinction}

The supplied V18 appendix and audit were read in full, including the distinction
between primitive finite carriers and the full spatial Haar lift. The new
argument uses V16's native moments, V17's all-volume entry expansion, V18's
primitive exterior screening, and the exact finite-range specification of the
same law. It does not use V18's final corridor trace theorem to prove an
operator bound, nor assume a native functional inequality.

Two targeted semantic searches in the supplied f14/f15 and monograph sources
returned no matches; their mounted text inventories were then searched. The
monograph's shell--Doob covariance proposition already states the
martingale covariance formula for \emph{one common filtration}. F15 V56's
exact reference covariance decomposition already separates a retained-predictor
Gram from conditional variances for its reference law. These mechanisms are
credited, not claimed as new. Here each primitive has a \emph{different}
centered ball filtration. A common-filtration orthogonality formula cannot
be transferred to that collection without proof. An exact counterexample
below explains why. This is a targeted audit, not recertification of the
complete archive. The next scheduled companion checkpoint remains V20.

For context, Brydges--Guadagni--Mitter,
\href{https://arxiv.org/abs/math-ph/0303013}{\emph{Finite range Decomposition
of Gaussian Processes}}, construct positive finite-range decompositions for
specified Gaussian covariances. Bauerschmidt,
\href{https://arxiv.org/abs/1206.2212}{\emph{A simple method for finite range
decomposition of quadratic forms and Gaussian fields}}, gives a related
Green-form construction. Their primary abstracts and scope were checked;
neither complete theorem is imported. Our conditional means are native
non-Gaussian functions and their scales are not independent Gaussian fields.
The first page and introduction of Lohmann's
\href{https://arxiv.org/abs/1411.1107}{\emph{Single Scale Cluster Expansions}}
were also checked; no cluster-expansion hypotheses for this Wilson measure
are silently declared verified. The supplied memoranda's distinction between
power gain and locality cost, and between summable dependence and a physical
gap, remains binding.

\subsection{One native law and one finite shell for each primitive}

Keep the complete unmoving law $P=P_{\gamma,\mathbf y}$, the fixed bridge
window $\max_{i,e}|y_{i,e}|\le r$, and put
\[
 \varepsilon=\gamma^{-1/2},\qquad \mathcal N=B(n+1).
\]
Use V18's cell metric $d_*$ and a fixed integer $\rho\ge1$ larger than every
native interaction diameter and every primitive-source carrier radius. Let
$\mathcal J$ be the primitive scalar index set and $a(u)$ the anchor of $u$.
There are at most $30\mathcal N$ such indices and uniformly bounded index
multiplicity at each cell. Write
\begin{equation}
 F_u=e^{\rm prim}_u-\mathbb E_P e^{\rm prim}_u,\qquad
 K_\gamma=(\mathbb E_P F_uF_v)_{u,v\in\mathcal J},\qquad
 \Gamma_\gamma=T K_\gamma T^{\mathsf T},\quad T=(I,\mathsf P^{\mathsf T}).
 \label{f16v19:primitive}
\end{equation}
The exact synthesis $T$ has uniformly bounded operator norm and absolute row
sums. It is the original Haar-source synthesis, not a new random projection.
The original residuals and all their normalizers are unchanged.

For any integer $L\ge\rho$, define the cell ball and its actual word collar by
\[
 I_{u,L}=\{c:d_*(a(u),c)\le L\},\qquad
 C_{u,L}=\{c\notin I_{u,L}:c\text{ occurs in an interaction meeting }I_{u,L}\}.
\]
Then $C_{u,L}\subset\{L<d_*(a(u),c)\le L+\rho\}$. Balls are intersected
with the actual finite history; there is no artificial time continuation.
Let
\begin{equation}
 m_{u,L}=\mathbb E_P[F_u\mid\zeta_{I_{u,L}^c}],\qquad
 R_{\gamma,L}=(\mathbb E_P m_{u,L}m_{v,L})_{u,v},\qquad
 A_{\gamma,L}=K_\gamma-R_{\gamma,L}.
 \label{f16v19:shell-definitions}
\end{equation}
If the ball is the whole system, its conditional mean is zero. All matrices
are evaluated under the one full native probability. In particular the
conditional partition function is included in each $m_{u,L}$, and its actual
marginal is used in the Gram. Different balls can overlap arbitrarily.

\begin{theorem}[Exact finite-range subtraction with a positive return]
\label{f16v19:exact-local}
For every finite regulator and every $L\ge\rho$, $m_{u,L}$ is measurable
with respect to $\zeta_{C_{u,L}}$, and
\begin{equation}
 R_{\gamma,L}\succeq0,\qquad
 (A_{\gamma,L})_{uv}=0\quad
       \text{if }d_*(a(u),a(v))>2L+\rho.
 \label{f16v19:finite-range}
\end{equation}
Thus $K_\gamma=A_{\gamma,L}+R_{\gamma,L}$ exactly, with no discarded row or
column. This does not assert $A_{\gamma,L}\succeq0$.
For every fixed integer $N\ge1$, V18's analytic input gives
\begin{equation}
 s_\gamma(L):=\sup_u\|m_{u,L}\|_2
 \le C_{N,r}\bigl(\varepsilon e^{-\vartheta L}+\varepsilon^{N+1}\bigr),
 \label{f16v19:screening-input}
\end{equation}
for sufficiently large coupling independent of $B,n,L$. Consequently every
entry of $R_{\gamma,L}$ has absolute value at most $s_\gamma(L)^2$, and
$\Tr R_{\gamma,L}\le30\mathcal N s_\gamma(L)^2$.
\end{theorem}
\begin{proof}
The primitive carrier is inside $I_{u,L}$. Product coordinate reference
measure, including the fractional endpoint Haar powers, and the finite-range
native action imply that the conditional integral for $F_u$ given $I_{u,L}^c$
depends only on $C_{u,L}$. Every crossing term and its normalizer is retained.
Subtracting the deterministic full mean changes no measurability. This is the
same Markov-blanket identity used in V11--V18, now for overlapping balls.

Suppose the two anchors have distance greater than $2L+\rho$. The entire
carrier of $F_v$ is outside $I_{u,L}$, and $C_{u,L}$ is outside $I_{v,L}$.
Successive conditional expectations therefore give
\[
 \mathbb E_P F_uF_v=\mathbb E_P m_{u,L}F_v
                   =\mathbb E_P m_{u,L}m_{v,L}.
\]
No independence of overlapping balls or of their exterior marginals is used.
This proves the exact range statement. The return is a Gram, hence positive.
The primitive case of \eqref{f16v18:source-centered}, applied to this ball
and centered by the actual mean, proves \eqref{f16v19:screening-input}.
Cauchy--Schwarz and summation of the diagonal give the last assertions.
The proof is unchanged when a ball saturates the finite system.
\end{proof}

The matrix $A_{\gamma,L}$ is not defined by deleting interactions from $P$.
Its near entries still involve exact native integrals. Nor is it the Gram of
$F_u-m_{u,L}$ alone. If $d_u=F_u-m_{u,L}$, its correct expression is
\[
 (A_{\gamma,L})_{uv}=\mathbb E_P d_ud_v
              +\mathbb E_P d_um_{v,L}+\mathbb E_P m_{u,L}d_v.
\]
The cross terms are necessary because the conditional projections differ
with $u$. A finite-range matrix is not thereby a replacement transfer operator.

\subsection{The finite-range part has an unnormalized operator expansion}

Retain exactly the primitive coefficients $K_q$ from V18 and put
$K^{[J]}=\sum_{q=2}^J\varepsilon^qK_q$. Their exponentially weighted
absolute row and column bounds and V17's entry error are inherited as
\eqref{f16v18:primitive-inputs}. Decrease $\vartheta$ once so that these
coefficient bounds and \eqref{f16v19:screening-input} use the same exponent.

\begin{proposition}[A finite-scale operator estimate without corridors]
\label{f16v19:local-error}
For fixed integers $k\ge J\ge2$ and $N\ge1$, and $L\ge\rho$,
\begin{align}
 \|A_{\gamma,L}-K^{[J]}\|_{\rm op}
 \le C_{k,N,r}\bigl[&(1+L)^4
    (\varepsilon^{k+1}+\varepsilon^2e^{-2\vartheta L}
                                  +\varepsilon^{2N+2})\notag\\
    &+\varepsilon^{J+1}
       +\varepsilon^2 e^{-\vartheta(2L+\rho)}\bigr].
 \label{f16v19:local-error-bound}
\end{align}
The coefficient-sum term $\varepsilon^{J+1}$ may be omitted when $k=J$.
The coupling threshold and constants are independent of the ambient volume
and history length. There is no factor proportional to their source count.
\end{proposition}
\begin{proof}
Let $\mathsf B_L$ be the entry mask retaining pairs of anchor distance at
most $2L+\rho$. The exact range statement gives
\[
 A_{\gamma,L}-K^{[J]}
 =\mathsf B_L(K_\gamma-K^{[k]})
   +\mathsf B_L(K^{[k]}-K^{[J]})
   -\mathsf B_L R_{\gamma,L}-(I-\mathsf B_L)K^{[J]}.
\]
Here $I-\mathsf B_L$ denotes the complementary entry mask, not subtraction
of a source-space projection. We estimate every term by absolute rows and
columns, not by assuming that an arbitrary entry mask contracts operator norm.

Polynomial growth of the four-dimensional cell metric and bounded primitive
multiplicity imply at most $C(1+L)^4$ retained entries per row or column.
The entrywise remainder of order $k$ therefore costs
$C(1+L)^4\varepsilon^{k+1}$. The entry bound for $R_{\gamma,L}$ costs
$C(1+L)^4s_\gamma(L)^2$. The finitely many higher coefficients, when present,
have uniformly bounded absolute rows even after masking and cost
$C\varepsilon^{J+1}$. Finally the weighted coefficient rows give a tail cost
$C\varepsilon^2e^{-\vartheta(2L+\rho)}$. Schur's row-and-column bound and
\eqref{f16v19:screening-input} prove the assertion. Small systems have fewer
entries, so the same estimate applies without requiring a large ball to fit.
\end{proof}

\begin{theorem}[A positive native remainder and one-sided coherent asymptotics]
\label{f16v19:positive-remainder}
For every pair of fixed integers $J\ge2$ and $M\ge1$, there are a logarithmic radius
$L_\gamma=O_{J,M,r}(\log\gamma)$ and finite constants independent of $B,n$
such that, at sufficiently large coupling,
\begin{equation}
 \boxed{\quad
 \Gamma_\gamma=G^{[J]}+\mathsf R^{\rm sh}_{\gamma,L_\gamma}
                                      +\mathsf E_{\gamma,J,M},\qquad
 \mathsf R^{\rm sh}_{\gamma,L}:=T R_{\gamma,L}T^{\mathsf T}\succeq0,
 \quad}
 \label{f16v19:positive-split}
\end{equation}
where $G^{[J]}=\sum_{q=2}^J\varepsilon^qG_q$ uses the original coefficients,
and
\begin{equation}
 \|\mathsf E_{\gamma,J,M}\|_{\rm op}\le C_{J,M,r}\varepsilon^{J+1},\qquad
 \max_\alpha(\mathsf R^{\rm sh}_{\gamma,L_\gamma})_{\alpha\alpha}
       +\frac{\Tr\mathsf R^{\rm sh}_{\gamma,L_\gamma}}{\mathcal N}
       \le C_{J,M,r}\varepsilon^M.
 \label{f16v19:positive-budgets}
\end{equation}
In particular, on the \emph{entire} original source space,
\begin{equation}
 \boxed{\quad\Gamma_\gamma\succeq
           G^{[J]}-C_{J,r}\gamma^{-(J+1)/2}I.\quad}
 \label{f16v19:one-sided}
\end{equation}
Equivalently the negative spectral part of $\Gamma_\gamma-G^{[J]}$ has
operator norm at most $C_{J,r}\gamma^{-(J+1)/2}$. This is not an upper bound
on its positive part.
\end{theorem}
\begin{proof}
Choose $k=J+1$, $N=\max\{J+2,\lceil M/2\rceil\}$, and for
$0<\varepsilon<1/2$ take
\begin{equation}
 L_\gamma=\left\lceil\frac{N+2}{\vartheta}
                     \log(\varepsilon^{-1})\right\rceil+2\rho+1.
 \label{f16v19:log-radius}
\end{equation}
Then $e^{-\vartheta L_\gamma}\le\varepsilon^{N+2}$ and
$s_\gamma(L_\gamma)^2\le C_{N,r}\varepsilon^{2N+2}$.
In \eqref{f16v19:local-error-bound}, the first entry remainder is bounded by
$C(1+\log\varepsilon^{-1})^4\varepsilon^{J+2}$, which is at most
$C\varepsilon^{J+1}$. The other logarithm-weighted terms and the coefficient
tail are smaller after increasing fixed constants. Thus
$\|A_{\gamma,L_\gamma}-K^{[J]}\|_{\rm op}\le C\varepsilon^{J+1}$.

Apply the same deterministic synthesis $T$ to the exact primitive split.
V18's identity $G_q=TK_qT^{\mathsf T}$ identifies every coefficient.
Set $\mathsf E=T(A_{\gamma,L_\gamma}-K^{[J]})T^{\mathsf T}$; boundedness
of $T$ proves its norm estimate. Positivity and the trace ideal inequality
bound the returned trace by $C\mathcal N s_\gamma(L_\gamma)^2$.
For a canonical source, its returned variable is one row of $T$ applied to
$m_{u,L_\gamma}$. The uniform absolute row sum and Minkowski bound its
$L^2$ norm by $Cs_\gamma(L_\gamma)$, proving the diagonal estimate.
Since $2N+2\ge M$, these are \eqref{f16v19:positive-budgets}.

The error is symmetric. Its lower quadratic-form bound is
$-C\varepsilon^{J+1}I$, and the return is positive. This proves the
one-sided bound, taking one fixed $M$ to remove it from the notation there.
Spectral calculus for a finite symmetric matrix gives the equivalent
negative-part bound. There is no division by $\mathcal N$ in that operator
statement. No common eigenbasis for $\Gamma_\gamma,G^{[J]}$ is assumed.
\end{proof}

At the leading order this gives the new unrestricted-direction statement
\[
 \gamma\Gamma_\gamma\succeq G_2-C_r\gamma^{-1/2}I.
\]
It rules out a coherent \emph{downward} failure of the leading comparison.
A large upward failure can still live in the positive shell return. In fact
\begin{equation}
 \left|\|\Gamma_\gamma-G^{[J]}\|_{\rm op}
                 -\|\mathsf R^{\rm sh}_{\gamma,L_\gamma}\|_{\rm op}\right|
 \le C_{J,M,r}\varepsilon^{J+1}.
 \label{f16v19:edge-equivalence}
\end{equation}
This is the triangle inequality, not an assumption that the return is small.
The same decomposition for the original memory is
\begin{equation}
 \mathsf M_\gamma-M^{[J]}
   =-\operatorname{Off}_{>1}\mathsf R^{\rm sh}_{\gamma,L_\gamma}
                      +\mathsf E^{\rm mem},\qquad
 \|\mathsf E^{\rm mem}\|_{\rm op}\le4C_{J,M,r}\varepsilon^{J+1}.
 \label{f16v19:memory-split}
\end{equation}
The factor four is the inherited Fourier block-extraction estimate. The
extracted return need not be positive: the Loewner conclusion is for the
complete covariance, not for its time-off-diagonal mask. No contact term or
endpoint normalizer is changed.

\subsection{Two tempting positivity inferences are false}

\begin{proposition}[Different local filtrations do not supply a monotone covariance split]
\label{f16v19:counterexamples}
In a positive Gaussian Markov law, $K-R_L$ need not be positive even though
$R_L$ is a conditional-message Gram. Also $R_L-R_{L'}$ need not be positive
for $L<L'$, although every individual message variance decreases. These
statements concern general local conditional constructions, not counterexamples
to the native Yang--Mills bounds above.
\end{proposition}
\begin{proof}
For three centered Gaussian coordinates take precision
$H_{ii}=1$, $H_{ij}=3/4$ for $i\ne j$, and $C=H^{-1}$. For singleton
interiors, the vector of exterior means is $(I-H)\zeta$. Hence
\[
 C-(I-H)C(I-H)=2I-H.
\]
The precision is positive, with eigenvalues $5/2,1/4,1/4$. The displayed
local difference has eigenvalue $-1/2$ on $(1,1,1)$.

For failure of matrix monotonicity, instead use the path precision
\[
 H=\begin{pmatrix}1&-1/3&0\\-1/3&1&-1/3\\0&-1/3&1\end{pmatrix},\qquad
 C=\frac17\begin{pmatrix}8&3&1\\3&9&3\\1&3&8\end{pmatrix}.
\]
Singleton interiors give
\[
 R_0=\begin{pmatrix}1/7&2/21&1/7\\2/21&2/7&2/21\\1/7&2/21&1/7\end{pmatrix}.
\]
Radius-one interiors on this path give exterior means
$(\zeta_3/8,0,\zeta_1/8)$ and therefore
\[
 R_1=\begin{pmatrix}1/56&0&1/448\\0&0&0\\1/448&0&1/56\end{pmatrix}.
\]
All diagonal differences are nonnegative, but for $v=(1,0,-1)^{\mathsf T}$,
$v^{\mathsf T}(R_0-R_1)v=-1/32$. These are exact rational Gaussian
computations. The radius-zero convention is used only in these examples;
the native construction above continues to use $L\ge\rho$.
\end{proof}

Thus increasing $L$ is not an unproved Loewner-monotone renormalization flow.
Nor can one subtract successive $R_L$ as positive covariances and then sum
those differences. The appropriate positive objects are the Grams of actual
\emph{scalar} reverse-martingale increments, with cross-scale terms retained.

\subsection{Annular innovations: positivity at each scale, not between scales}

Introduce the actual synthesis operator into the one common probability space,
\[
 \mathcal M_L:\ell^2(\mathcal J)\longrightarrow L^2(P),\qquad
 \mathcal M_Lv=\sum_u v_u m_{u,L},\qquad R_{\gamma,L}=\mathcal M_L^*\mathcal M_L.
\]
Let $L_0\ge\rho$ be an integer and $L_j=2^jL_0$. Define
\[
 d_{u,j}=m_{u,L_j}-m_{u,L_{j+1}},\qquad
 b_j=\sup_u\|d_{u,j}\|_2,\qquad
 \mathcal D_j v=\sum_uv_ud_{u,j}.
\]
At every finite regulator these sequences are eventually zero, because all
balls eventually contain the whole system.

\begin{proposition}[Exact annular orthogonality and finite-range Grams]
\label{f16v19:annuli}
For every anchor and scale,
\begin{equation}
 \mathbb E_P[d_{u,j}\mid\zeta_{I_{u,L_{j+1}}^c}]=0,\qquad
 \|d_{u,j}\|_2^2=\|m_{u,L_j}\|_2^2-\|m_{u,L_{j+1}}\|_2^2.
 \label{f16v19:scalar-drop}
\end{equation}
The variable $d_{u,j}$ is measurable in the union of its two actual word
collars, contained in the ball of radius $L_{j+1}+\rho$. Its Gram is positive
and has range at most $2L_{j+1}+\rho$. In particular,
\begin{equation}
 \|\mathcal D_j\|_{\ell^2\to L^2(P)}
                  \le C(1+L_{j+1})^2 b_j.
 \label{f16v19:annular-norm}
\end{equation}
The exponent two is half the four-dimensional volume-growth exponent.
\end{proposition}
\begin{proof}
For a \emph{fixed} $u$, the sigma algebras of exteriors decrease with $L$.
The tower property gives
$\mathbb E[m_{u,L_j}\mid I_{u,L_{j+1}}^c]=m_{u,L_{j+1}}$.
Orthogonality of one projection difference gives both identities in
\eqref{f16v19:scalar-drop}. Markov-blanket measurability gives the collar
support. If the two anchors have distance greater than $2L_{j+1}+\rho$,
the entire support of $d_{v,j}$ is outside $I_{u,L_{j+1}}$. Conditioning
therefore makes their inner product zero. No independence assertion is needed.
The Gram is positive by definition, each entry has absolute value at most
$b_j^2$, and there are at most $C(1+L_{j+1})^4$ nonzero entries per row
and column. Schur's estimate bounds its operator norm by that count times
$b_j^2$. Taking a square root proves \eqref{f16v19:annular-norm}.
This calculation does not assert orthogonality for different source anchors
at different scales.
\end{proof}

\begin{theorem}[A source-specific dyadic criterion for the remaining coherent return]
\label{f16v19:dyadic}
At every finite regulator,
\begin{equation}
 \boxed{\quad
 \|R_{\gamma,L_0}\|_{\rm op}^{1/2}
 \le C\sum_{j\ge0}(1+L_{j+1})^2 b_j
 \le C\sum_{j\ge0}(1+L_{j+1})^2s_\gamma(L_j).
 \quad}
 \label{f16v19:dyadic-bound}
\end{equation}
For the full-source return the same right side is multiplied by $\|T\|$.
The sums are finite at each finite regulator. A uniform bound on the infinite
majorant would make the conclusion uniform in volume and duration.

In particular, the \emph{additional} source-specific estimate
\begin{equation}
 s_\gamma(L)\le A_\gamma(1+L)^{-2-\delta}
       \quad(L\ge L_0),\qquad \delta>0,
 \label{f16v19:screening-target}
\end{equation}
would imply
\begin{equation}
 \|\mathsf R^{\rm sh}_{\gamma,L_0}\|_{\rm op}
       \le C_\delta A_\gamma^2(1+L_0)^{-2\delta}.
 \label{f16v19:power-consequence}
\end{equation}
No assertion that \eqref{f16v19:screening-target} has been proved for the
native law is part of this theorem.
\end{theorem}
\begin{proof}
At a fixed finite system, the centered message at every sufficiently large
radius is zero. Hence $\mathcal M_{L_0}=\sum_j\mathcal D_j$ exactly as a
finite sum of operators into $L^2(P)$. Apply the triangle inequality in that
operator space and \eqref{f16v19:annular-norm}. The scalar variance-drop
identity gives $b_j\le s_\gamma(L_j)$, proving the second inequality.
No unproved infinite-volume tail triviality is used to remove the final term:
the telescoping is completed on each finite regulator before taking a uniform
upper bound. A use directly in infinite volume would have to retain its tail
projection unless separately shown zero.

For \eqref{f16v19:screening-target}, the summands are at most
$C A_\gamma(1+L_0)^{-\delta}2^{-j\delta}$. Sum the geometric series,
square, and use the bounded source synthesis. Cross-scale correlations are
paid by this triangle inequality, rather than set to zero.
\end{proof}

The target here is weaker in kind than a Poincar\'e inequality for every
native $L^2$ function: it involves only the two specified primitive error
banks and their own exterior conditional means. It is also an averaged
native statement, not a worst-exterior bound. If it held with
$A_\gamma\le C_r\varepsilon$ at all radii beyond a logarithmic $L_0$, the
positive split would give
\begin{equation}
 \|\Gamma_\gamma-\gamma^{-1}G_2\|_{\rm op}
 \le C_r\gamma^{-3/2}
       +C_{r,\delta}\gamma^{-1}(1+\log\gamma)^{-2\delta}
                  =o(\gamma^{-1}).
 \label{f16v19:conditional-edge}
\end{equation}
The memory operator would obey the same conclusion up to the factor four.
This implication does not supply exponential temporal row decay or a sampler
gap. A critical-power bound with
$L^{-2}(\log(2+L))^{-1-\delta}$ would also make the dyadic majorant summable,
with its corresponding slower tail. The ordinary volume-growth exponent
four has not been replaced by a claim of two-dimensional geometry.

\subsection{What is already proved across polynomially many radii}

Although \eqref{f16v19:screening-target} is unproved, the established
finite-order screening does give a directional statement across a long but
finite range of scales. This uses the annular operator bound, not a trace.

\begin{theorem}[Coherent transport from logarithmic to polynomial shells]
\label{f16v19:transport}
For fixed $p>0$ and $q>0$ there are a logarithmic integer
$L_\gamma=O_{p,q,r}(\log\gamma)$, a radius
$\gamma^p\le\mathcal R_\gamma<2\gamma^p$, and fixed constants such that
\begin{equation}
 \boxed{\quad
 \|\mathcal M_{L_\gamma}-\mathcal M_{\mathcal R_\gamma}\|_{\ell^2\to L^2(P)}
                      \le C_{p,q,r}\varepsilon^q.
 \quad}
 \label{f16v19:transport-bound}
\end{equation}
The larger radius can be chosen as a dyadic multiple of $L_\gamma$.
The estimate holds at every ambient volume and duration, for sufficiently
large coupling depending on the fixed exponents, not on the system size.
It remains valid after composition with $T^{\mathsf T}$.
\end{theorem}
\begin{proof}
Choose $N\ge\max\{1,\lceil4p+q\rceil\}$ in
\eqref{f16v19:screening-input}, and
\[
 L_\gamma=\left\lceil\frac{q+3}{\vartheta}
                        \log(\varepsilon^{-1})\right\rceil+2\rho+1.
\]
For sufficiently large coupling $L_\gamma<\gamma^p$. Choose the first dyadic
multiple at least $\gamma^p$, and telescope the annular innovations only up
to that radius. Their total operator norm is bounded by
\begin{align*}
 C_{N,r}\sum_{L_j<\mathcal R_\gamma}(1+2L_j)^2
                 (\varepsilon e^{-\vartheta L_j}+\varepsilon^{N+1})
 &\le C_{N,r}\bigl[
       \varepsilon(1+L_\gamma)^2e^{-\vartheta L_\gamma}
            +\mathcal R_\gamma^2\varepsilon^{N+1}\bigr]\\
 &\le C_{p,q,r}\varepsilon^q.
\end{align*}
The first sum is bounded by the corresponding integer exponential tail;
its logarithmic factor is absorbed by the extra power in the choice of
$L_\gamma$. The second is a finite geometric sum of squared dyadic radii.
Since $\mathcal R_\gamma^2\le4\varepsilon^{-4p}$, its power is at least
$N+1-4p\ge q$. Saturation of balls only makes later terms vanish. No factor
counting source coordinates or time slices occurs. Composing with the
bounded deterministic synthesis proves the last assertion.
\end{proof}

The theorem is an amplitude-level estimate. One cannot simply square it to
conclude $\|R_{\gamma,L_\gamma}-R_{\gamma,\mathcal R_\gamma}\|_{\rm op}
=O(\varepsilon^{2q})$: the cross term with a potentially large remote message
must be retained. The following relative enclosure pays that term.

\begin{theorem}[A relative positive-return enclosure at a remote shell]
\label{f16v19:remote-enclosure}
Fix $J\ge2$, $p>0$, and $h>0$. For suitable radii as in
Theorem~\ref{f16v19:transport}, write
$\tau=\varepsilon^h$ and
$\mathsf R^{\rm far}=T R_{\gamma,\mathcal R_\gamma}T^{\mathsf T}\succeq0$.
For all admitted geometries and sufficiently large fixed-order coupling,
\begin{equation}
 \boxed{
 \begin{aligned}
 G^{[J]}+(1-\tau)\mathsf R^{\rm far}-C\varepsilon^{J+1}I
   &\preceq\Gamma_\gamma\\
   &\preceq G^{[J]}+(1+\tau)\mathsf R^{\rm far}
                                      +C\varepsilon^{J+1}I .
 \end{aligned}}
 \label{f16v19:remote-order}
\end{equation}
The same underlying native functions define $\mathsf R^{\rm far}$; no
probability law is replaced by a reference shell law. The matrix is not
asserted small in operator norm.
\end{theorem}
\begin{proof}
Take $q\ge(J+1+h)/2$ in the transport theorem and enlarge its logarithmic
radius and fixed expansion order to also satisfy the local error estimate
of Theorem~\ref{f16v19:positive-remainder}. This is allowed: finitely many
fixed requirements are imposed, and the same estimates remain valid at a
larger logarithmic radius. Let
$\mathcal U=(\mathcal M_{L_\gamma}-\mathcal M_{\mathcal R_\gamma})T^{\mathsf T}$.
Its norm is at most $C\varepsilon^q$. For each source vector $v$, put
$X=\mathcal M_{\mathcal R_\gamma}T^{\mathsf T}v$ and $Y=\mathcal Uv$.
For $0<\tau<1$, the two elementary Hilbert-space inequalities are
\[
 (1-\tau)\|X\|^2-\tau^{-1}\|Y\|^2
 \le\|X+Y\|^2
 \le(1+\tau)\|X\|^2+(1+\tau^{-1})\|Y\|^2.
\]
Thus the two shell Grams are relatively enclosed, with additive error at most
$C\varepsilon^{2q-h}\|v\|^2$. Insert the positive split and its local
operator error. Since $2q-h\ge J+1$, these combine into
\eqref{f16v19:remote-order}. This explicitly pays the cross-scale terms
without bounding the norm of $X$ or silently changing the clock or measure.
\end{proof}

For fixed $J,p,h$, every unresolved coherent excess is therefore represented,
up to the displayed small relative and additive errors, by messages on a
shell at distance of order $\gamma^p$. If that shell already lies beyond the
finite system, the return is zero; this is consistent with V17's finite-bank
conclusion, not a new thermodynamic assumption. At unrestricted volume the
remote shell may still carry a nonzero, highly correlated vector of messages.

\subsection{What has advanced and the precise next target}

The new unconditional result is a \emph{one-sided coherent} asymptotic at
unrestricted volume, and an exact positive-return decomposition with a small
local operator error. It is not merely another trace budget. There are no
sparse discarded corridors and no omitted primitive sources. The returned
Gram has arbitrary fixed-order small diagonal and trace density, but its
largest eigenvalue is retained, not bounded by those densities.

The annular argument also gives the proved transport estimate through every
prescribed polynomial range of radii. To reach all scales at a fixed coupling,
one would need a summable bound on
\[
 \sum_{j\ge0}(1+2^{j+1}L_0)^2
 \sup_u\left\|m_{u,2^jL_0}-m_{u,2^{j+1}L_0}\right\|_2.
\]
This is a concrete source-resolved target for the actual compact law.
An averaged scalar decay faster than radius to the power minus two is a
sufficient condition, not a proved consequence of V18. Substituting its
fixed-order floor into this infinite sum would diverge. Nor can one choose
an order increasing with the scale while ignoring its constants and entry
thresholds. The counterexample to matrix monotonicity excludes another
possible shortcut.

The local repair results remain possible tools for this shell problem. The
new formulation does not require a native gap for every test function as an
input, but it does not prove that such a gap exists. A successful coherent
covariance bound would still be a statement about the fixed-bridge fast law,
not a physical Hamiltonian. Bridge-window coverage, root Gauss restriction,
infrared contact coercivity, the same-top physical transfer comparison,
physical-time calibration and an interacting continuum construction retain
their separate obligations. The positive shell return is not identified with
V12's one-common-exterior cylinder return; they are different exact matrices
under the same probability.

\subsection{Diagnostics, preservation, and scope}

The new diagnostic is
\begin{center}\small
\path{verify_ym_f16_v19_positive_shell.py}.
\end{center}
It checks exact shell measurability in an inhomogeneous finite Markov law,
finite-range covariance subtraction, the annular tower and variance-drop
identities, positive source synthesis, both rational Gaussian counterexamples,
finite-geometry support, and the relative-order and scale bookkeeping. Its
report distinguishes these finite fixtures from analytic proofs and from
native compact integral enclosures. The majority of assertions verify the
same collar value across individual configurations of the finite fixture;
the count is not a measure of proof completeness. The new diagnostic passed
3,948 exact assertions, including 3,530 collar-value comparisons. All fifteen
available V4--V18 diagnostics were rerun unchanged and passed; their actual
outputs and hashes are retained in the package.

The cumulative V19 changes only the V18 title version and appends this exact
section before the final document-ending command. The package includes a
byte-for-byte predecessor recovery check, input hashes, actual diagnostic
reruns, build reports and rendered-page review. No inherited proof is silently
corrected or independently recertified. The new analytic deductions explicitly
use the identified V16--V18 inputs, not their diagnostic assertion counts.
The next scheduled companion checkpoint remains V20.

\begin{firewall}
V19 proves an exact finite-range primitive covariance difference, a positive
native shell-return decomposition, and an unrestricted-volume lower Loewner
asymptotic with a genuine operator-norm error. It proves coherent transport
to polynomially remote shells and a relative positive-return enclosure.
The dyadic source-screening criterion for the remaining upper edge is proved
as an implication, not verified for the native law at all radii. No complete
two-sided operator expansion, summable native remainder, uniform sampler gap,
physical transfer gap or interacting continuum mass gap is established.
All original measures, normalizers, source definitions and physical caveats
remain unchanged.
\end{firewall}
% END F16 V19 APPEND


% BEGIN F16 V20 APPEND -- 2026-09-17
% Insert immediately before the final document-ending command in V19.
% No inherited mathematical statement or bibliography string is changed.

\clearpage
\section{V20: uniform boundary preparation and buffered contraction}
\label{f16v20:context}

V19 leaves a positive Gram of remote conditional messages at the unknown upper
edge. Increasing a fixed perturbative order again would not control its
infinite sequence of radii. We go upstream to a different conditional input:
\emph{prepare a deep interior under every compact exterior}, then compare its
entire marginal law, not just polynomial observables. The buffer is integrated,
not frozen. Its exterior may be at angular distance of order one from the
identity, including the cancellation configurations of V14.

The result is a volume-independent total-variation estimate for a core of
$k$ cells, with its \emph{linear} cost in $k$ explicit. It implies a conditional
maximal-correlation bound for every square-integrable core observable and every
incoming exterior law. This is stronger in test class and conditioning scope
than V18's averaged polynomial screening, but has only a first-order coupling
error. It does not give the missing all-radius bound for V19's whole source
bank. The separate roles of output size, buffer depth, and update clock are
kept explicit.

\subsection{The V20 checkpoint and the new analytic input}

The supplied V19 appendix and audit were read in full. At this scheduled
checkpoint the available f14 V100, f15 V100, Grand Tensor Monograph V076,
and current lineage were inventoried and targeted passages checked. Two scoped
semantic searches were empty; exact mounted-text searches and identified
source ranges were then used. This is not recertification of the entire
archive. In particular, f15 V22 already compares full compact core kernels in
total variation, with a confidence-event error, and constructs a separated
coupling. F14 V65 compares normalized nested chart laws. The monograph's
strong-electric Dobrushin theorem concerns a different regime and transfer.
None of these hypotheses is silently transferred to the present law.

The new preparation estimate keeps the terminal boundary term in V16's ball
recursion, rather than sending the radius to infinity after averaging the
whole native law. A second change replaces V17's polynomial right inverse
by a Gaussian divergence solver for \emph{bounded tests}. A common cutoff on
the observed core permits the proof to control the complete divergence,
without assuming a bounded Hessian of a Stein solution for arbitrary bounded
tests. Coordinate-local cutoffs are still used on the other differentiated
groups. All Haar powers and all cutoff terms are retained.

For attribution, the Gaussian Ornstein--Uhlenbeck/Stein representation is
classical; the Mehler formula and Lemma 3.3 in Nourdin--Peccati--R\'eveillac,
\href{https://arxiv.org/abs/0804.1889}{\emph{Multivariate normal approximation
using Stein's method and Malliavin calculus}}, were checked. That paper's
Wasserstein theorem is not quoted as our total-variation theorem. The
bounded-test gradient estimate actually used is proved below. Polyanskiy--Wu,
\href{https://arxiv.org/abs/1508.06025}{\emph{Strong data-processing inequalities
for channels and Bayesian networks}}, Section 2.1, Theorem 1 and (7)--(9),
records the domination of chi-square contraction by the Dobrushin coefficient
and its relation to maximal correlation on general spaces. We give a direct
Hilbert-space proof of this consequence with our precise kernels.

The two supplied import memoranda separate power gain from locality cost and
require a genuine dependence estimate before global coercivity. This appendix
provides such an estimate for a \emph{buffered marginal channel}, not for the
full collection of overlapping updates. Only supplied companions were
available at the checkpoint; other programs named in older ledgers were not
recertified. The next scheduled checkpoint is V25.

\subsection{The conditional law and finite-distance boundary preparation}

Use V16's unmoving cells, $e_c=|\zeta_c|^2$, graph $\mathcal G$, and constants
$a_0=(6\pi^2)^{-1}$, $t_0=a_0/4$. Fix
\[
 P=P_{\gamma,\mathbf y},\qquad \max_{i,e}|y_{i,e}|\le r,
 \qquad \gamma\ge1+r^2.
\]
At most twelve principal three-dimensional groups occur in a cell, so on the
complete compact coordinates, including every fixed exterior value,
\begin{equation}
                  0\le e_c\le E_*\gamma,\qquad E_*=12\pi^2.
 \label{f16v20:compact-cap}
\end{equation}
The graph joins cells sharing a non-anchor primitive interaction. It has the
polynomial ball bound of V16. Write $d_{\mathcal G}$ for its distance. Different
components do not interact. This metric is used in this appendix; every graph
step has bounded range in V18's ambient space--time metric.

Let $\Lambda$ be a nonempty set of free cells, and freeze all cells outside it
at $q$. Denote the explicit complete conditional law by
\[
 P_\Lambda^q(d\zeta_\Lambda)
  =Z_\Lambda(q)^{-1}e^{-U_\Lambda(\zeta_\Lambda;q)}
                                  \prod_{c\in\Lambda}d\nu_{\gamma,c}.
\]
Here $U_\Lambda$ contains every anchor on $\Lambda$ and every primitive
interaction meeting $\Lambda$, once each, with the original endpoints.
The factors $\nu_{\gamma,c}$ are V16's unnormalized complete chart factors.
There is no imposed fast-field window on $q$. We use this positive-density
version of the conditional integral at all principal exterior values; at
Haar-null chart boundaries this is an explicitly specified extension, not a
claim of uniqueness of a regular conditional version.

Define $d_c=d_{\mathcal G}(c,\Lambda^c)$. Set $d_c=\infty$ if the entire
component of $c$ is free, and $e^{-\kappa\infty}=0$. Tilts $s$ on the free cells
have $|s_c|\le t_0$, exactly as in V16; their law is denoted $P_{\Lambda,s}^q$.

\begin{theorem}[Conditional preparation with an exponentially attenuated boundary term]
\label{f16v20:preparation}
There are $B_r<\infty$ and $\kappa_r>0$, independent of $\Lambda,q,B,n,\gamma$
and the source array, such that
\begin{equation}
 \mathbb E_{P_{\Lambda,s}^q}e_c
             \le M_r+B_r\gamma e^{-\kappa_r d_c},\qquad c\in\Lambda.
 \label{f16v20:prepared-mean}
\end{equation}
Here $M_r$ may be V16's constant enlarged once. For every array
$0\le t_c\le t_0$ on the free cells,
\begin{equation}
 \log\mathbb E_{P_\Lambda^q}\exp\!\left(\sum_c t_c e_c\right)
       \le\sum_c t_c\bigl(M_r+B_r\gamma e^{-\kappa_r d_c}\bigr).
 \label{f16v20:prepared-mgf}
\end{equation}
In particular, at depths $d_c\ge\kappa_r^{-1}\log\gamma$, every fixed local
moment has a uniform bound, now for \emph{all} compact exterior values.
There also exist a fixed integer $D_r$ and $c_0>0$ such that every original
group $g$ in a cell of depth at least $D_r$ satisfies
\begin{equation}
 P_\Lambda^q\{|\zeta_g|\ge\sqrt\gamma/2\}
                         \le e^{t_0M_r-c_0\gamma}.
 \label{f16v20:prepared-cut-tail}
\end{equation}
The latter constant-depth bound concerns an angular cut event, not a uniform
thermal moment at constant depth.
\end{theorem}
\begin{proof}
Keep the proof of V16's conditional normalizer estimate, including all crossing
terms and its product unit-ball lower bound. It applies to any $S\subseteq
\Lambda$ conditional on all its other variables. Average only in
$P_{\Lambda,s}^q$. With
\[
 u_c=\mathbb E_{P_{\Lambda,s}^q}e_c\ (c\in\Lambda),\qquad
 u_c=e_c(q)\ (c\notin\Lambda),
\]
its conclusion is still
\begin{equation}
             \sum_{c\in S}u_c\le a_r|S|+b_r\sum_{d\in\partial S}u_d.
 \label{f16v20:conditional-set}
\end{equation}
This is not a new estimate on arbitrary fixed boundary values without a debit:
the exterior entries on the right remain their actual deterministic energies.
All $u_c$ also satisfy \eqref{f16v20:compact-cap}.

For $c\in\Lambda$, write $F_k=\sum_{d\in B_k(c)}u_d$ and
$\rho_r=b_r/(1+b_r)$. For $0\le k<d_c$, the ball $B_k(c)$ is entirely free,
so V16's same recurrence gives
\[
 F_k\le\rho_r F_{k+1}+\frac{a_r}{1+b_r}|B_k(c)|.
\]
If $d_c$ is finite, stop at $d_c$, not at infinity. The terminal term is
bounded by
\[
 \rho_r^{d_c}F_{d_c}
     \le E_*C_{\rm gr}\gamma(d_c+1)^4\rho_r^{d_c}.
\]
The sum of the preceding terms is at most V16's $M_r$. Choose
$\kappa_r=-\tfrac12\log\rho_r$ and
\[
 B_r\ge E_*C_{\rm gr}\sup_{j\ge0}(j+1)^4\rho_r^{j/2}<\infty.
\]
This proves \eqref{f16v20:prepared-mean}. If $d_c=\infty$, the recurrence
can be completed on its finite component and the terminal term tends to zero,
as in V16. No volume limit is exchanged with this iteration.

The proof is uniform over all signed free sources. Integrate the exact
finite-regulator derivatives of $\log Z_\Lambda(q;s)$ along the segment
from zero to $t$. This proves \eqref{f16v20:prepared-mgf}; no derivative of a
pressure approximation is taken. The local moment assertion follows by the
same exponential-moment inequality as V16. Finally choose $D_r$ such that
$B_r e^{-\kappa_rD_r}\le1/8$. A cut group has $e_c\ge\gamma/4$.
Exponential Markov at $t_0$ in that cell gives
\eqref{f16v20:prepared-cut-tail} with $c_0=t_0/8$.
\end{proof}

The preparation is uniform over the \emph{size} of the free set as well as
the exterior. A boundary cell may still have mean of order $\gamma$.
This is consistent with V14's unpinned pair, which had no free buffer between
the observed variables and its cancelling exterior. No global convexity of
the native compact potential has been asserted.

\subsection{A bounded-test Gaussian solver with a local core}

Let $\mathcal A\subseteq\Lambda$ be a nonempty set of $k$ complete cells, and
let $A$ denote all its real coordinates, with $d_A\le36k$. Put
\[
 L=d_{\mathcal G}(\mathcal A,\Lambda^c),\qquad
 Q=N(m,C),\quad C=H^{-1},\qquad
 Q_A=N(m_A,C_A),\quad C_A=C_{AA}.
\]
The Gaussian $Q$ is the original full-history reference. In particular
$C_A$ is a \emph{marginal} covariance, not $(H_{AA})^{-1}$ or the covariance
of a newly chosen conditional reference. It is independent of $q$.
All scalar indices below include their fixed colour multiplicities.

The inherited bounds on $H$ imply, at some fixed $\nu>0$,
\begin{equation}
 h_+^{-1}I\preceq C_A\preceq h_-^{-1}I,\qquad
 \sup_i\sum_j|C_{ij}|e^{\nu d_{\mathcal G}(c(i),c(j))}\le K_C,
 \qquad \sup_i|m_i|\le M_r^{\rm G}.
 \label{f16v20:Gaussian-local}
\end{equation}
For completeness, $H$ has range one in $\mathcal G$. The Neumann series
$H^{-1}=h_+^{-1}\sum_{l\ge0}(I-H/h_+)^l$ has geometric operator decay and
no entry before the corresponding graph distance. Polynomial ball growth
makes its exponentially weighted rows summable at a smaller exponent.
Across different components the entries vanish. Thus using the native graph
rather than the ambient metric introduces no unproved inverse bound.

For a smooth bounded $\varphi$ on $\mathbb R^{d_A}$ with bounded derivatives,
define
\begin{equation}
 u_\varphi(x)=\int_0^\infty
 \left[\mathbb E\varphi\bigl(m_A+e^{-t}(x-m_A)
                   +\sqrt{1-e^{-2t}}\,C_A^{1/2}G\bigr)-Q_A\varphi\right]dt,
 \quad G\sim N(0,I).
 \label{f16v20:OU-solver}
\end{equation}

\begin{proposition}[Bounded-test divergence without a Hessian test norm]
\label{f16v20:bounded-solver}
Set $b(x)=\nabla u_\varphi(x)$ and, on the entire coordinate space, set
$g_j(\zeta)=C_{jA}b(\zeta_A)$. Then
\begin{equation}
 \delta_Qg=\varphi(\zeta_A)-Q_A\varphi,\qquad
 \|b\|_\infty\le\sqrt{\pi/2}\,\|C_A^{-1/2}\|\,\|\varphi\|_\infty.
 \label{f16v20:bounded-solver-identity}
\end{equation}
Here $\delta_Q$ is the same full Gaussian divergence as in V17.
The components $g_j$ depend only on the core, even for $j\notin A$.
If $a_j=\|C_{jA}\|_2$, then
\begin{equation}
 \|g_j\|_\infty\le C a_j\|\varphi\|_\infty,\qquad
 \sum_j a_j\le Ck,\qquad \|g_A\|_\infty\le C\|\varphi\|_\infty.
 \label{f16v20:solver-size}
\end{equation}
None of these bounds uses a derivative norm of $\varphi$.
\end{proposition}
\begin{proof}
The generator of the Mehler semigroup in \eqref{f16v20:OU-solver} is
$\mathcal L_A=\operatorname{Tr}(C_A\nabla^2)-(x-m_A)\cdot\nabla$.
Integration of its time derivative gives
$-\mathcal L_Au_\varphi=\varphi-Q_A\varphi$. The integral converges at each
$x$; for large $t$, Gaussian translation and covariance comparison, or a
bounded derivative of the smooth test, gives an integrable exponential bound.
For any unit vector $h$, Gaussian integration by parts inside the Mehler
integral gives
\[
 |\partial_h P_t^A\varphi(x)|
 \le\sqrt{2/\pi}\,\frac{e^{-t}}{\sqrt{1-e^{-2t}}}
             |C_A^{-1/2}h|\,\|\varphi\|_\infty.
\]
The integral of the time factor is $\pi/2$, by the substitution $s=e^{-t}$.
This proves the gradient bound, uniformly in $x$ and in $d_A$. Initially the
smooth-test assumption justifies differentiations; its derivative constants
do not occur in this estimate.

Since $HC=I$, the drift term in $\delta_Qg$ is
$(\zeta_A-m_A)\cdot b$. Only indices $j\in A$ differentiate $b$, and their
sum is $\operatorname{Tr}(C_A\nabla^2u_\varphi)$. This proves the divergence
identity. The absolute column bound in \eqref{f16v20:Gaussian-local} gives
$\sum_j a_j\le\sum_{i\in A,j}|C_{ji}|\le K_Cd_A$. Finally
$\|g_A\|\le\|C_A\|\,\|b\|$ gives the last estimate.
\end{proof}

A generic multivariate bounded-test Stein solution need not have the uniform
Hessian bound needed for a Wasserstein-to-TV substitution. We do not use such
a bound. The next proof keeps all core divergence terms together; it never
estimates their individual second derivatives on a cutoff complement.

\subsection{The conditional drift error and the boundary reaching its core}

Write $V_\gamma=U_\gamma-\log J^f_\gamma$ on the open full chart, and
$\xi=\zeta-m$. Only derivatives in free coordinates are used; the exterior
is substituted after forming these derivatives. Use the same smooth group
cutoff $\chi_j=\chi(\zeta_{g(j)}/\sqrt\gamma)$ as V17: it equals one within
angular radius $1/2$ and vanishes beyond radius one. For $j\in\Lambda$ define
\[
                    D_j=\partial_jV_\gamma-(H\xi)_j.
\]
All expectations in this subsection are under $P_\Lambda^q$.

\begin{proposition}[An actual conditional drift estimate]
\label{f16v20:drift}
For every free coordinate $j$, any multiplier $0\le\eta_j\le1$ supported
where $|\zeta_{g(j)}|\le\sqrt\gamma$ obeys
\begin{equation}
 \mathbb E[\eta_j|D_j|]
                \le C_r\bigl(\gamma^{-1/2}
                    +\gamma^{1/2}e^{-\kappa_r d_{c(j)}}\bigr),
 \label{f16v20:drift-bound}
\end{equation}
after enlarging constants and, if necessary, decreasing $\kappa_r$.
There is $\mu_r>0$ such that the solver weights satisfy
\begin{align}
 \sum_{j\in\Lambda}a_j e^{-\kappa_r d_{c(j)}}
                              &\le C_r k e^{-\mu_r L},\notag\\
 \sum_{j\notin\Lambda}a_j
   +\sum_{j\in\Lambda:\ d_{c(j)}<D_r}a_j
                              &\le C_r k e^{-\mu_r L}.
 \label{f16v20:boundary-sums}
\end{align}
Membership of a scalar coordinate in $\Lambda$ abbreviates membership of its
cell. The constants do not count all exterior coordinates.
\end{proposition}
\begin{proof}
V17's global, first-order word-gradient remainder gives
\[
 |\partial_jU_\gamma-(H\xi)_j|
             \le C_r\gamma^{-1/2}(1+e_{N_j}),
\]
where $N_j$ has fixed size and lies within one graph step of $c(j)$.
This follows from differentiating the literal unitary exponential words;
quadratic endpoint terms cancel exactly. It is valid even when some inputs
are frozen exterior values. For a free neighbour apply
\eqref{f16v20:prepared-mean}; its distance to the boundary is at least
$d_{c(j)}-1$. An exterior neighbour occurs only at that boundary and has the
compact cap \eqref{f16v20:compact-cap}. This bounds the expected classical
remainder as claimed. On the multiplier's support the only differentiated
Haar logarithm is smooth inside its own angular unit ball, with its original
power $1$ or $1/2$. Its derivative is bounded by
$C\gamma^{-1}|\zeta_{g(j)}|\le C\gamma^{-1/2}$, proving
\eqref{f16v20:drift-bound}. No cut-singular logarithmic derivative is bounded
on the complete chart.

For $i\in A$ and $j\in\Lambda$, the path distance obeys
$d_{\mathcal G}(c(i),c(j))+d_{c(j)}\ge L$.
Choose $0<\mu_r\le\min(\nu,\kappa_r)/2$.
Multiply by the weighted covariance row in \eqref{f16v20:Gaussian-local}
and sum, first over $j$ and then over $i\in A$. This proves the first bound
using $a_j\le\sum_{i\in A}|C_{ji}|$. For an exterior $j$ the distance from
every $i\in A$ is at least $L$. For a free $j$ with $d_{c(j)}<D_r$ it is at
least $L-D_r$. The same row tail proves the second bound; $D_r$ is fixed.
If a core component has no boundary the corresponding terms vanish.
\end{proof}

\subsection{A full core law, uniform over all compact exteriors}

Use total variation in the convention
$d_{\rm TV}(\mu,\nu)=\sup_E|\mu(E)-\nu(E)|$.
Let $K_{\Lambda,A}(q,\cdot)$ be the marginal of $P_\Lambda^q$ on $\zeta_A$,
viewed on $\mathbb R^{d_A}$ by zero extension outside its original charts.

\begin{theorem}[Uniform native core total variation after an integrated buffer]
\label{f16v20:core-TV}
There are fixed $C_r<\infty$, $\mu_r>0$, and $L_r<\infty$ such that, for
$L\ge L_r$, $\gamma\ge1+r^2$, every finite ambient history and every
$\mathcal A\subseteq\Lambda$ as above,
\begin{equation}
 \boxed{
 \sup_q d_{\rm TV}\bigl(K_{\Lambda,A}(q,\cdot),Q_A\bigr)
 \le\min\left\{1,\ C_r k\left(\gamma^{-1/2}
                           +\gamma^{1/2}e^{-\mu_r L}\right)\right\}.}
 \label{f16v20:TV}
\end{equation}
The supremum includes every compact principal exterior value, not just values
with a bounded rescaled energy or a high-probability event. The buffer
$\Lambda\setminus\mathcal A$ is fully integrated. Both the native marginal
normalizer and the Gaussian marginal covariance are retained.
\end{theorem}
\begin{proof}
Take initially a smooth bounded test $\varphi$ with bounded derivatives and
use the solver of Proposition~\ref{f16v20:bounded-solver}. Scale so
$\|\varphi\|_\infty\le1$. Introduce \emph{one common core cutoff}
\[
 \kappa_A(\zeta_A)=\prod_{g\text{ in the core}}\chi(\zeta_g/\sqrt\gamma).
\]
It has $0\le\kappa_A\le1$, support strictly inside all core charts and
$\|\nabla\kappa_A\|\le C\gamma^{-1/2}\sqrt{k}$. Every core group has depth
at least $L$. For $L_r\ge D_r$, the union bound used \emph{only on this
observed core} and \eqref{f16v20:prepared-cut-tail} give
\begin{equation}
              \mathbb E(1-\kappa_A)\le C_r k e^{-c_0\gamma}.
 \label{f16v20:core-cut-mass}
\end{equation}
This is not an all-history event.

As $\delta_Qg=\varphi-Q_A\varphi$ is bounded by two, splitting its expectation
with $\kappa_A$ costs at most twice \eqref{f16v20:core-cut-mass}.
For $j\in A$, integrate by parts using $\kappa_Ag_j$.
For $j\in\Lambda\setminus A$, use $\kappa_A\chi_jg_j$; both $g_j$ and
$\kappa_A$ are independent of this coordinate. No exterior coordinate is
differentiated. These operations give the exact identity
\begin{align}
 \mathbb E[\kappa_A\delta_Qg]
 ={}&-\sum_{j\in A}\mathbb E[\kappa_A g_jD_j]
       -\sum_{j\in\Lambda\setminus A}
                          \mathbb E[\kappa_A\chi_jg_jD_j]\notag\\
 &+\mathbb E[\nabla\kappa_A\cdot g_A]
       +\sum_{j\in\Lambda\setminus A}
                          \mathbb E[\kappa_A(\partial_j\chi_j)g_j]\notag\\
 &+\sum_{j\in\Lambda\setminus A}
                         \mathbb E[\kappa_A(1-\chi_j)(H\xi)_jg_j]
       +\sum_{j\notin\Lambda}\mathbb E[\kappa_A(H\xi)_jg_j].
 \label{f16v20:bounded-cut-identity}
\end{align}
The last sum is an algebraic exterior term, not integration by parts at fixed
coordinates. There is no missing $\partial_jg_j$ there because $g_j$ depends
only on $\zeta_A$. The measured Haar derivatives are included in $D_j$.
Fubini on one interior group justifies every integration; the cutoff removes
its own logarithmic cut. Initially smooth tests suffice for this calculation.

The first line is at most
$C_r k(\gamma^{-1/2}+\gamma^{1/2}e^{-\mu_r L})$ by the drift proposition
and \eqref{f16v20:solver-size}. The first term on the second line is at most
$C\gamma^{-1/2}\sqrt{k}$ by the vector bound on $g_A$.
The remaining cutoff-derivative terms are at most
$C\gamma^{-1/2}\sum_j a_j\le Ck\gamma^{-1/2}$.

On the full compact coordinates, including the fixed exterior,
$|(H\xi)_j|\le C_r\sqrt\gamma$: use the bounded absolute row of $H$, the
bounded mean and \eqref{f16v20:compact-cap}.
In the third line, split free indices by $d_{c(j)}\ge D_r$ and its complement.
For the former, $1-\chi_j$ requires the group cut event; its probability is
at most $C_r e^{-c_0\gamma}$. Their total is at most
$C_r k\sqrt\gamma e^{-c_0\gamma}\le C_r k\gamma^{-1/2}$.
For the latter, and for the last exterior sum, use
\eqref{f16v20:boundary-sums} and the compact bound instead. They cost at most
$C_r k\sqrt\gamma e^{-\mu_r L}$. The discarded expectation
$\mathbb E[(1-\kappa_A)\delta_Qg]$ has the same required bound.
We have proved
\[
 |P_\Lambda^q\varphi(\zeta_A)-Q_A\varphi|
          \le C_rk(\gamma^{-1/2}+\gamma^{1/2}e^{-\mu_r L})
\]
with a constant depending on no derivative norm of the test.

Both core measures are finite regular Borel measures on a Euclidean space
and have Lebesgue densities (one extended by zero). Bounded indicators can
therefore be approximated in their joint $L^1$ norm by smooth bounded tests.
Apply the preceding uniform test bound and pass to that limit. This proves
the total-variation assertion, absorbing its factor-of-two convention in
$C_r$. Although the approximating tests can depend on $q$, the bound does not,
so taking the supremum is legitimate. At no point was a small conditional
partition function replaced by a Gaussian denominator.
\end{proof}

The common core cutoff is important. Cutting each core component separately
and estimating the individual Hessian terms would reintroduce a test-regularity
loss. Here the omitted part multiplies the \emph{whole bounded divergence}.
The proof does not assert finite relative Fisher information of a marginal
with a fractional Haar power at its compact cut.

For example,
\begin{equation}
 L\ge L_r+\mu_r^{-1}\log\gamma
 \quad\Longrightarrow\quad
 \sup_q d_{\rm TV}(K_{\Lambda,A}(q),Q_A)\le C_rk/\sqrt\gamma.
 \label{f16v20:log-buffer}
\end{equation}
A core of any fixed number of cells is thus prepared uniformly, independent
of the size or energy of the frozen exterior. A growing core pays its stated
$k$ cost; this is not a bound for an arbitrary macroscopic output bank.

\subsection{An incoming-law-independent conditional contraction}

For a Markov kernel $K:Z\longrightarrow X$ on standard Borel spaces define
its diameter
\[
                 \delta(K)=\sup_{z,z'}d_{\rm TV}(K(z),K(z')).
\]
A law $\pi$ on $Z$ gives the joint probability
$\mathbb J_\pi(dz,dx)=\pi(dz)K(z,dx)$ and the output marginal $\mu=\pi K$.

\begin{proposition}[Dobrushin diameter bounds all-observable correlation]
\label{f16v20:channel-lemma}
For every incoming law $\pi$ and every $F\in L^2(\mu)$,
\begin{equation}
 \operatorname{Var}_\pi(KF)\le\delta(K)\operatorname{Var}_\mu F,
 \qquad \operatorname{corr}_{\max}(Z,X)^2\le\delta(K).
 \label{f16v20:channel-variance}
\end{equation}
The adjoint conditional expectation has the same squared norm on centered
spaces. In particular, if $\pi'\ll\pi$ has finite chi-square divergence, then
\begin{equation}
      \chi^2(\pi'K\Vert\pi K)\le\delta(K)\chi^2(\pi'\Vert\pi).
 \label{f16v20:chi-square}
\end{equation}
This is a standard data-processing consequence; it requires no lower bound
on either incoming probability and no bound on the test function's sup norm.
\end{proposition}
\begin{proof}
Let $\mathsf K:L^2(\mu)\to L^2(\pi)$ be conditional expectation and let
$\mathsf K^*$ be its adjoint, the reverse conditional kernel.
The operator $B=\mathsf K^*\mathsf K$ on $L^2(\mu)$ is a positive
self-adjoint Markov contraction with invariant law $\mu$. Its transition
rows are mixtures of the rows of $K$, hence their total-variation diameter
is at most $\delta(K)$. This remains true for regular conditional versions
on a full-measure set; extend null rows by any fixed mixture.
For bounded $F$, the oscillation contraction gives
\[
 \|B^l(F-\mu F)\|_2\le\operatorname{osc}(B^l F)
             \le\delta(K)^l\operatorname{osc}(F).
\]
If the centered positive operator $B$ had spectral mass above $\delta(K)$,
a bounded centered $F$, by density, would have a nonzero projection onto an
interval $[a,1]$ with $a>\delta(K)$. The spectral theorem would then give
$\|B^lF\|_2\ge a^l\|1_{[a,1]}(B)F\|_2$, contradicting the last estimate.
Thus $B\preceq\delta(K)I$ on centered space. Since
$\|\mathsf KF\|_2^2=\langle F,BF\rangle$, this proves the first inequality
and the maximal-correlation statement, including unbounded $L^2$ tests.
The adjoint norm is equal. Finally the centered density of $\pi'K$ relative
to $\pi K$ is $\mathsf K^*(d\pi'/d\pi-1)$. Its squared norm is the output
chi-square divergence. Apply the same adjoint norm bound.
\end{proof}

\begin{theorem}[A complete native buffered channel, for every exterior law]
\label{f16v20:buffered-contraction}
For the exact conditional kernel $K_{\Lambda,A}$, set
\begin{equation}
 \Delta_{\gamma,k,L}
   =\min\{1,\ C_r k(\gamma^{-1/2}+\gamma^{1/2}e^{-\mu_rL})\},
 \label{f16v20:Delta}
\end{equation}
with $C_r$ enlarged from Theorem~\ref{f16v20:core-TV}. Then
$\delta(K_{\Lambda,A})\le\Delta_{\gamma,k,L}$, and for \emph{every} law
$\pi$ on all compact exterior parameters,
\begin{equation}
 \boxed{
 \operatorname{Var}_\pi\!\bigl(\mathbb E[F(\zeta_A)\mid q]\bigr)
       \le\Delta_{\gamma,k,L}
                         \operatorname{Var}_{\pi K_{\Lambda,A}}F.}
 \label{f16v20:native-channel}
\end{equation}
This holds for all output $L^2$ functions, including arbitrary rare-event
weights with finite variance. It also holds for every vector of such functions
as a Loewner inequality between its two covariance matrices, with the same
factor. In particular, in the actual native joint law of core and exterior,
\begin{equation}
 \operatorname{corr}_{\max}(\zeta_A,\zeta_{\Lambda^c})
                           \le\sqrt{\Delta_{\gamma,k,L}}.
 \label{f16v20:native-rho}
\end{equation}
The conclusion survives arbitrary reweighting of the \emph{exterior law},
or further fixing outside $\Lambda$. It does not authorize extra conditioning
inside the core or its integrated buffer.
\end{theorem}
\begin{proof}
The triangle inequality through the \emph{same} Gaussian marginal $Q_A$
bounds the distance between any two native rows by twice the right side of
\eqref{f16v20:TV}. Enlarge the constant, cap at one, and apply
Proposition~\ref{f16v20:channel-lemma}. Test that scalar inequality on every
linear combination of a finite vector to obtain the covariance order.
For the actual law, $\pi$ is precisely its exterior marginal, including
$Z_\Lambda(q)$ in its density. No replacement exterior probability is used.
For a different incoming $\pi$, the conditional kernel remains exactly the
same; the lemma was uniform in $\pi$. Further conditioning outside $\Lambda$
only restricts incoming parameters and changes their law. Fixing a buffer
variable instead changes $K_{\Lambda,A}$ itself, which the theorem has not
allowed.
\end{proof}

On the collapsed two-block probability $\mathbb J_\pi$ of $(\zeta_A,q)$,
the two rate-one conditional resamplings consequently have gap at least
$1-\sqrt{\Delta_{\gamma,k,L}}$. Indeed, the two centered coordinate subspaces
have angle norm at most $\sqrt\Delta$; V13's two-projection argument gives
this bound on their full joint $L^2$ space. For a logarithmic buffer and fixed
$k$, sufficiently large coupling makes $\Delta\le1/4$ and gives gap at least
$1/2$ for \emph{this collapsed two-block generator}.

This is not a gap for V14's original local sampler. The buffer has been
integrated out; updating the core conditional on $q$ resamples its buffer
implicitly, and the other collapsed update can resample the entire far
exterior. No comparison returning these two operations to the original
single-port/chord clock is asserted. This limitation is mathematical, not a
choice of terminology.

\subsection{Separated cores: a simultaneous same-law coupling}

\begin{proposition}[Uniform buffered-core coupling with no confidence exception]
\label{f16v20:coupling}
Let free cell sets $\Lambda_s$ be separated so that no primitive interaction
meets two of them. Let $\mathcal A_s\subset\Lambda_s$ have $k_s$ cells and
depth $L_s\ge L_r$. Set
\[
 \eta_s=\min\{1,C_r k_s(\gamma^{-1/2}
                              +\gamma^{1/2}e^{-\mu_rL_s})\}.
\]
On a probability extension of the \emph{unchanged} native law, reference core
vectors $\widehat\zeta_{A_s}$ can be added such that their joint law is
$\bigotimes_s Q_{A_s}$, independently of the common exterior, and for every
subfamily $S$,
\begin{equation}
 \mathbb P\{\zeta_{A_s}\ne\widehat\zeta_{A_s}\text{ for all }s\in S\}
                      \le\prod_{s\in S}\eta_s.
 \label{f16v20:joint-coupling}
\end{equation}
\end{proposition}
\begin{proof}
Condition on the actual common exterior of the separated free sets. Their
interior laws factor exactly by the finite-range specification. For each
core couple its conditional marginal to $Q_{A_s}$ by the common-part
(maximal) coupling, independently between blocks at this fixed exterior.
Its mismatch probability is the corresponding total variation, at most
$\eta_s$ for every exterior by Theorem~\ref{f16v20:core-TV}. The reference
marginals do not depend on the exterior, so their conditional product law
remains the stated independent product after exterior averaging. Fill every
buffer using its original conditional law given its core and the exterior;
this retains the original native full-field law. Conditional multiplication
of the mismatch probabilities, followed by exterior averaging, proves the
claim. Measurability follows from the explicit densities and the measurable
common-part construction on Euclidean spaces. The same construction was used
in f15 V22; the new input here is a uniform bound without a confidence-bad
boundary term. The reference cores are auxiliary random variables, not a
replacement for the integrated native probability.
\end{proof}

\subsection{Why a local channel certificate is not yet the whole positive return}

Two finite examples make the remaining operations precise. They are examples
of probability channels, not counterexamples to the native theorems.

\begin{proposition}[Input reweighting, output conditioning, and coherent aggregation differ]
\label{f16v20:limits}
A channel can have arbitrarily small Dobrushin diameter while conditioning on
an event in its \emph{output} produces diameter arbitrarily close to one.
Also many individually contracting output channels, conditionally independent
given one input, can have a conditional-message Gram whose largest eigenvalue
grows with their number.
\end{proposition}
\begin{proof}
For binary input $Z=\pm1$ and output $X\in\{*,+,-\}$ take the positive rows
\[
 K(+)=\bigl(1-\eta,\eta(1-\sigma),\eta\sigma\bigr),\qquad
 K(-)=\bigl(1-\eta,\eta\sigma,\eta(1-\sigma)\bigr),
 \quad0<\eta<1,\quad0<\sigma<1/2.
\]
Their diameter is $\eta(1-2\sigma)$. Conditional on the output event
$X\ne *$, the diameter is $1-2\sigma$. For $\eta=1/100$ and
$\sigma=1/1000$ these are respectively $499/50000$ and $499/500$.
Input-independent probability of the conditioning event does not prevent
this change of the channel. Thus arbitrary exterior reweighting in
Theorem~\ref{f16v20:buffered-contraction} is not arbitrary output or buffer
conditioning.

For the second assertion let $Z=\pm1$ be uniform and, conditional on $Z$,
let $X_1,\ldots,X_N\in\{\pm1\}$ be independent with
$\mathbb E[X_i\mid Z]=aZ$, $0<a<1$. Each individual channel has diameter $a$
and squared maximal correlation $a^2$. Nevertheless
\[
 \operatorname{Cov}\bigl((\mathbb E[X_i\mid Z])_{i=1}^N\bigr)
                       =a^2\mathbf1\mathbf1^{\mathsf T}
\]
has operator norm $Na^2$, while every diagonal entry is $a^2$.
Conditional independence of the outputs did not make their conditional means
independent under the actual incoming law.
\end{proof}

The new native result provides a true all-observable channel contraction after
a logarithmic buffer, even with an adversarial compact exterior. This is a
usable input for a future block or information-percolation argument: it does
not first assume the required native $L^2$ gap. But the output-size factor in
\eqref{f16v20:Delta} cannot be suppressed. A shell surrounding a large ball
has growing dimension. Applying the theorem to that whole shell eventually
makes the bound trivial at fixed coupling. Applying it separately to small
pieces does not permit a free tensorization through their correlated incoming
law, as the example shows.

In particular the estimate does not verify V19's dyadic all-radius condition.
For a fixed local primitive it gives a relative $L^2$ screening inequality
under any incoming exterior law, but the power gain for that primitive under
the unmodified native law is weaker than V18's fixed-order estimates. Its
advance is the \emph{class of tests and the boundary uniformity}, not a new
arbitrary-order expansion. No previous coefficient or shell Gram is redefined.

The next noncircular target is a \emph{local composition theorem} for these
buffered channels with the actual shared boundary messages and output counts
retained, or a comparison of a compatible buffered block dynamics to the
original update clock. Conditioning inside the buffer, multiplying unrelated
contraction factors, or treating a sparse mismatch coupling as independence
of the native error fields would each omit a required step. The original
positive shell-return norm, and the physical same-top transfer problem,
therefore remain distinct tasks.

\subsection{Diagnostics, checkpoint preservation, and claim scope}

The accompanying diagnostic is
\begin{center}\small
\path{verify_ym_f16_v20_buffered_channels.py}.
\end{center}
It checks the finite stopped boundary recurrence, finite-graph distance and
boundary paths, the core-supported Gaussian divergence, the exact combined
cutoff identity with fixed exterior parameters, channel reversals and
chi-square quadratic forms, conditional couplings, and both limiting examples.
The new checker passed 401 exact assertions. The report records their families
and distinguishes these fixtures from analytic proofs. They do not enclose every native integral or verify the
analytic estimates with a proof assistant.

The cumulative V20 changes only V19's title version and appends this exact
section before its final document-ending command. The package records input
hashes, byte-for-byte predecessor recovery, new diagnostic output, the actual
reruns of the available V4--V19 diagnostics, and compilation/render checks.
No inaccessible companion program or external formalization is represented as
run. The current proof depends on V5/V16's exact independent anchors and local
normalizer estimates, and the identified Gaussian and word-derivative bounds;
it does not use V18--V19's final asymptotic theorems to supply a native gap.
The scheduled V20 review is limited to the supplied sources; the next
checkpoint is V25.

\begin{firewall}
V20 proves deep conditional moment bounds for every compact exterior, a
complete-core total-variation comparison with the original Gaussian marginal,
and a buffered all-observable maximal-correlation/chi-square contraction
uniform in the incoming exterior law. It also supplies a separated-core
coupling without a confidence exception. The output-size cost and integrated
buffer are explicit. No all-radius contraction of an unrestricted source bank,
uniform gap for the original local native sampler, same-top physical transfer
gap, bridge-window removal, or interacting continuum mass gap is inferred.
All native conditional normalizers, free compact domains, Haar factors and
original source definitions are retained.
\end{firewall}
% END F16 V20 APPEND


% BEGIN F16 V21 APPEND -- 2026-09-17
% Insert immediately before the final document-ending command in V20.
% No predecessor proof, bibliography, source definition, or normalizer is edited.

\clearpage
\section{V21: buffered temporal composition and polynomial-width operator asymptotics}
\label{f16v21:context}

V20 proves a buffered channel estimate but warns that unrelated channels cannot
be multiplied through a shared incoming law. We take a composition for which
the required conditional factorization is exact: \emph{complete time slices}
of the augmented native history. They are Markov separators. The continuation
message at each step is retained, including its endpoint and normalization;
no stationary endpoint or Gaussian transition replaces it.

There is also an upstream improvement before composition. V20 bounded its
Stein drift by summing absolute scalar errors. Its conditional preparation
actually supplies second moments of that drift sufficiently far inside the
buffer. Keeping the drift vector intact reduces the core-size cost from
$k/\sqrt\gamma$ to $\sqrt{k/\gamma}$. This permits spatial cross-sections
$B\le c_r\gamma$, rather than the exponential-in-$B$ entry restriction of
V8--V10. Time remains unrestricted. At this maximal width the proved decay
length is $O_r(\log\gamma)$. In every strictly sublinear power-width regime,
a fixed positive temporal decay exponent is available.

The resulting covariance expansion is \emph{two-sided in a weighted temporal
operator-row norm} on the entire source history. It is not a new unrestricted
spatial-volume theorem: the cross-section restriction is essential in the
proof and is stated with every main conclusion. No gap for the original local
sampler or the physical transfer is inferred from the conditional channels.

\subsection{Dependency check and the deliberate restriction of the composition}

The supplied V20 appendix and analytic audit were read in full. Its prepared
conditional moments, bounded-test Gaussian solver, exact combined cutoff
identity, and general-state-space channel lemma are the principal inputs.
V16 supplies full-native local moments; V17 supplies entrywise jets and
polynomial-bank estimates; V10 supplies locality of the normalized coefficients;
V7 supplies the unchanged Haar sources and contact identity. The present
proof does not use V19's one-sided order bound to assume the missing upper
bound. It does not use an unproved native functional inequality.

Two scoped semantic archive searches were empty. Mounted exact-text searches
were then used, and the identified f14 V52 Dobrushin passage and the
monograph's strong-electric slab statement were checked. Those are different
coupling regimes and physical operators. Their temporal mixing principles
are credited, not claimed absent; their hypotheses do not supply the current
weak-coupling conditional input. The current V8--V10 temporal results are also
not forgotten: their exponential spatial cost is precisely what the present
buffered comparison avoids. This is a targeted dependency audit, not an
independent certification of the historical analytic collection. The next
scheduled companion checkpoint remains V25.

For literature context, Polyanskiy--Wu,
\href{https://arxiv.org/abs/1508.06025}{\emph{Strong data-processing inequalities
for channels and Bayesian networks}}, Section 2.1 and Section 4, treats
single-channel contraction and its propagation through an actual network.
The definition of serial composition, the single-channel theorem, and the
network theorem and proof were examined, including their displayed PDF pages.
We only use the elementary serial-chain implication, proved below on the
present standard Borel spaces. No finite-alphabet network or percolation
result is imported without its hypotheses. The supplied memoranda's
separation of power gain from locality cost is retained: the new locality
input is a normalized conditional channel, not another formal coefficient.

\subsection{A square-root core-size estimate from the same bounded-test identity}

Use exactly V20's notation $P_\Lambda^q$, $\mathcal A$, $A$, $k=|\mathcal A|$,
$L=d_{\mathcal G}(\mathcal A,\Lambda^c)$, $Q_A=N(m_A,C_{AA})$, and its
explicit conditional-density version on complete principal charts. In
particular all fast exterior values $q$ are admitted. Let
$\varepsilon=\gamma^{-1/2}$. V20's preparation constants are denoted by
$\kappa_r,B_r,M_r,D_r$, and its Gaussian inverse row exponent by $\nu>0$.
After enlarging $D_r$ once, define
\begin{equation}
 d_\gamma=D_r+2+\left\lceil\kappa_r^{-1}\log\gamma\right\rceil,
 \qquad
 b_\gamma=d_\gamma+L_r+2+
                   \left\lceil2\nu^{-1}\log\gamma\right\rceil.
 \label{f16v21:buffer-length}
\end{equation}
All logarithms below are natural. We may increase a fixed-$r$ lower coupling
threshold so that $\gamma\ge e$, $b_\gamma\ge1$, and
\begin{equation}
                         b_\gamma\le C_{b,r}\log\gamma.
 \label{f16v21:buffer-growth}
\end{equation}
These quantities do not depend on $B,n,\Lambda,q,k$.

\begin{proposition}[Prepared second moments of the measured drift]
\label{f16v21:drift-L2}
Let $D_j=\partial_jV_\gamma-(H(\zeta-m))_j$, with the exact measured potential
and the cutoffs of V20. For a multiplier $0\le\eta_j\le1$ supported where the
original group containing $j$ has angular radius at most one,
\begin{equation}
 \|\eta_jD_j\|_{L^2(P_\Lambda^q)}\le
 \begin{cases}
 C_r\varepsilon,&d_{c(j)}\ge d_\gamma,\\
 C_r\sqrt\gamma,&j\in\Lambda.
 \end{cases}
 \label{f16v21:drift-two-regions}
\end{equation}
Consequently, for any choice of such multipliers and $L\ge d_\gamma$, the
vector indexed by the real core coordinates,
$W_i=\sum_{j\in\Lambda}C_{ij}\eta_jD_j$, satisfies
\begin{equation}
 \mathbb E_{P_\Lambda^q}|W_A|
 \le C_r\sqrt{k}\left(\varepsilon+
                   \sqrt\gamma e^{-\nu(L-d_\gamma)}\right).
 \label{f16v21:vector-drift}
\end{equation}
The multipliers need not be independent, deterministic, or local to $j$.
\end{proposition}
\begin{proof}
Every cell in the fixed drift carrier of a coordinate with depth at least
$d_\gamma$ is free and has depth at least
$\kappa_r^{-1}\log\gamma+D_r$. V20's conditional exponential moment
\eqref{f16v20:prepared-mgf}, on each such cell, bounds its fourth moment
uniformly in every compact $q$. The inherited complete-chart classical bound
is $|\partial_jU_\gamma-\partial_jU_0|\le
C_r\varepsilon(1+e_{N_j})$. Its $L^2$ norm is therefore $C_r\varepsilon$.
On the stated support the only differentiated Haar logarithm is smooth and
its derivative has magnitude at most $C\varepsilon$; the endpoint power
$1/2$ is retained. This proves the first case of
\eqref{f16v21:drift-two-regions}.

For the second case, a scaled Lie or chart derivative of each original
classical word has magnitude at most $C\sqrt\gamma$. The bounded local
incidence, bounded bridge window, exactly quadratic endpoint terms, and
compact cap $|\zeta_c|^2\le E_*\gamma$ preserve this bound. The Gaussian
linear drift has the same bound by its absolute row sum. The cut Haar term
again costs at most $C\varepsilon$. This proves the second case without
claiming thermal moments at arbitrary near-boundary conditionals.

For each $i\in A$, Minkowski in $L^2$, followed by the inverse row bound,
gives
\[
 \|W_i\|_2\le C_r\varepsilon\sum_j|C_{ij}|
       +C_r\sqrt\gamma\sum_{j:\ d_{c(j)}<d_\gamma}|C_{ij}|
 \le C_r\left(\varepsilon+
                 \sqrt\gamma e^{-\nu(L-d_\gamma)}\right).
\]
The last step uses the distance from every core cell to every shallow cell.
Summing the \emph{squared} coordinate bounds and using
$\mathbb E|W_A|\le(\sum_{i\in A}\mathbb E W_i^2)^{1/2}$ proves
\eqref{f16v21:vector-drift}, since $|A|\le36k$. No orthogonality or independence
of drift coordinates has been used.
\end{proof}

\begin{theorem}[Square-root output cost for a complete buffered marginal]
\label{f16v21:sqrt-core}
There are constants $C_{0,r}\ge1$ and $\gamma_{0,r}<\infty$ such that, at all
$\gamma\ge\gamma_{0,r}$ and for every core at depth $L\ge b_\gamma$,
\begin{equation}
 \boxed{\quad
 \sup_q d_{\rm TV}(K_{\Lambda,A}(q),Q_A)
                  \le\min\{1,C_{0,r}\sqrt{k/\gamma}\}.
 \quad}
 \label{f16v21:sqrt-TV}
\end{equation}
The Gaussian is the original full-history \emph{marginal}; the native law,
conditional normalizer, complete free domains, and both endpoint factors are
unchanged. No restriction on ambient volume or duration occurs in this theorem.
\end{theorem}
\begin{proof}
Keep V20's bounded-test solver $b=\nabla u_\varphi$ and
$g_j=C_{jA}b(\zeta_A)$, with $\|b\|_\infty\le C\|\varphi\|_\infty$.
Take $\|\varphi\|_\infty\le1$. Use precisely the common core cutoff
$\kappa_A$ and the exact identity \eqref{f16v20:bounded-cut-identity}.
The first line of that identity can be written as
\[
       -\mathbb E\, b(\zeta_A)\cdot W_A,
\]
with multipliers $\eta_j=\kappa_A$ on core indices and
$\eta_j=\kappa_A\chi_j$ on free noncore indices. They satisfy the support
conditions just proved. This whole line therefore costs
$C_r\sqrt{k}(\varepsilon+\sqrt\gamma e^{-\nu(L-d_\gamma)})$,
not the sum of $k$ unrelated error bounds.

For the common-core derivative use the unchanged bounds
$|\nabla\kappa_A|\le C\varepsilon\sqrt{k}$ and $\|g_A\|_\infty\le C$.
For the remaining cutoff derivatives, write their sum as
\[
 \mathbb E\,\kappa_A b(\zeta_A)\cdot
       \left(\sum_{j\in\Lambda\setminus A}C_{ij}\partial_j\chi_j\right)_{i\in A}.
\]
Each vector component is at most $C\varepsilon\sum_j|C_{ij}|$, so this
term also costs $C\varepsilon\sqrt{k}$.

The noncore cutoff-complement and fixed-exterior terms have the same vector
form. For a free coordinate of depth at least the fixed $D_r$,
\[
 \|(1-\chi_j)(H(\zeta-m))_j\|_2
               \le C_r\sqrt\gamma e^{-c_0\gamma/2},
\]
by the compact drift bound and V20's prepared angular-tail estimate.
For shallower free coordinates and for fixed exterior coordinates use the
compact drift bound and the exponentially small covariance row tail. For
each core component this costs at most
$C_r\sqrt\gamma(e^{-c_0\gamma/2}+e^{-\nu(L-D_r)})$.
Taking the vector second moment bounds these two lines by $\sqrt{k}$ times
that expression. All estimates remain valid after multiplication by
$\kappa_A$.

Only the complement of the \emph{whole bounded divergence} is estimated by
its core union bound, exactly as in V20. Its cost is $C_r k e^{-c_0\gamma}$.
Combining the terms gives, for $L\ge d_\gamma$ and a core deep enough for the
fixed-depth tail bound,
\begin{equation}
 |P_\Lambda^q\varphi-Q_A\varphi|
 \le C_r\sqrt{k}\left(\varepsilon+
                   \sqrt\gamma e^{-\nu(L-d_\gamma)}\right)
                           +C_r k e^{-c\gamma}.
 \label{f16v21:sqrt-pre-bound}
\end{equation}
The exponential noncore term has been absorbed into the first term by
increasing a fixed constant. At $L\ge b_\gamma$, the boundary term is at most
$\gamma^{-3/2}$, by \eqref{f16v21:buffer-length}. If $k\le\gamma$, then
$k e^{-c\gamma}\le C\sqrt{k/\gamma}$; if $k>\gamma$, the proposed capped
bound is already the trivial bound one. Approximation of bounded indicators
by smooth bounded tests in the two measures' joint $L^1$ norm gives total
variation as in V20. No derivative norm of $\varphi$, Fisher-information
assumption at a Haar cut, or changed reference probability is needed.
\end{proof}

The square-root estimate still has an output cost. We now pay it for a
complete slice, rather than claim that infinitely many independently tested
cores form a contracting channel for free.

\subsection{Complete time slices are exact Markov separators}

Let $Z_i=(\zeta_{i,b})_b$ denote all fast cells at time $i$. For $i<n$ its
real dimension is $36B$, and for $i=n$ it is $15B$. These are entire spatial
slices, not independent spatial cubes. The actual word list in V5 has no
interaction meeting times with separation greater than one. The full
probability can therefore be written exactly as
\begin{equation}
 dP=Z^{-1}\prod_{i=0}^n\psi_i(z_i)\,d\nu_i(z_i)
                         \prod_{i=0}^{n-1}W_i(z_i,z_{i+1}),
 \label{f16v21:chain-density}
\end{equation}
where $\nu_i$ is the product of the original cell factors at time $i$,
$\psi_i>0$ contains the within-slice action, and $W_i>0$ contains all
interactions joining the two indicated slices. Source-free scalar factors
cancel only in this normalized fast probability. The original bridge-history
weight is not redefined. No spatial interaction is deleted.

\begin{proposition}[Exact nonstationary transitions with their continuation messages]
\label{f16v21:Markov}
Set $h_n=1$ and recursively define
\begin{equation}
 h_i(x)=\int W_i(x,y)\psi_{i+1}(y)h_{i+1}(y)\,d\nu_{i+1}(y),\qquad
 K_i(x,dy)=\frac{W_i(x,y)\psi_{i+1}(y)h_{i+1}(y)}{h_i(x)}\,d\nu_{i+1}(y).
 \label{f16v21:continuation}
\end{equation}
Then $P$ is the inhomogeneous Markov-chain law with initial measure
$Z^{-1}\psi_0h_0\nu_0$ and transitions $K_i$. Every $h_i$ is finite and
strictly positive at the finite regulator. The same statement, with its
actual terminal factor, holds on an interval whose right exterior is fixed.
There is an analogous backward construction.
\end{proposition}
\begin{proof}
Assign each local action term to its single time or to its unique adjacent
time pair to obtain \eqref{f16v21:chain-density}. Every original complete
coordinate factor is finite; the classical exponent is bounded at a fixed
compact regulator. Thus the recursive integrals are finite and positive.
Their displayed quotients integrate to one. Multiplication of the initial
measure and all transitions cancels each $h_i$ against its literal neighbor,
leaving exactly \eqref{f16v21:chain-density}. In particular the right dressed
endpoint belongs to $\psi_n$, not an imposed invariant measure. A fixed right
boundary contributes its actual edge factor to the terminal integrand and
the same calculation applies. Reading the product from the other end gives
the backward statement, with the original left endpoint. This argument does
not assume equality of different $K_i$, self-adjointness of an individual
transition on a fixed space, or a lower bound on $h_i$ uniform in duration.
\end{proof}

Choose once a fixed width constant
\begin{equation}
 c_r=\min\{1/36,(8C_{0,r})^{-2}\}>0.
 \qquad B\le c_r\gamma.
 \label{f16v21:width}
\end{equation}
Every theorem below that uses contraction assumes \eqref{f16v21:width} and
$\gamma\ge\gamma_{0,r}$. This permits polynomially growing spatial tori but
not arbitrary $B$ at fixed coupling. Define
\begin{equation}
 \delta_{\gamma,B}=2C_{0,r}\sqrt{B/\gamma}\le\tfrac14,
 \qquad q_0=\tfrac14.
 \label{f16v21:delta}
\end{equation}
The constant $q_0$ here is a channel bound, not a cylinder-overlap count.

\begin{theorem}[Serial contraction under the actual endpoint law]
\label{f16v21:temporal-channel}
For $0\le i<j\le n$, let $K_{i,j}$ be the actual conditional transition from
$Z_i$ to $Z_j$ in \eqref{f16v21:chain-density}. Its Dobrushin diameter obeys
\begin{equation}
 \delta(K_{i,j})\le
             \delta_{\gamma,B}^{\lfloor(j-i)/b_\gamma\rfloor}.
 \label{f16v21:diameter-products}
\end{equation}
For all square-integrable $F$ measurable at times at most $i$, and $G$
measurable at times at least $j$, under the actual full native law,
\begin{equation}
 \boxed{
 |\operatorname{Cov}_P(F,G)|\le
 \delta_{\gamma,B}^{\lfloor(j-i)/b_\gamma\rfloor/2}
                 \sqrt{\operatorname{Var}_PF\operatorname{Var}_PG}.}
 \label{f16v21:past-future}
\end{equation}
The statement is uniform in all bounded, arbitrarily time-varying bridge
histories and every duration. No Gaussian transition or stationary endpoint
law is substituted.
\end{theorem}
\begin{proof}
Given the complete past through time $t$, the conditional future is the
native conditional on the free set consisting of all cells at times $t+1$
through $n$. Its marginal at time $t+b_\gamma$ has $B$ cells. Each graph step
changes time by at most one, so this core has depth at least $b_\gamma$;
the physical end at $n$ is free, not an artificial frozen boundary.
Theorem~\ref{f16v21:sqrt-core} compares every such row with the same original
Gaussian slice marginal. Two rows therefore have distance at most
$\delta_{\gamma,B}$. By the exact Markov property this is precisely the
$b_\gamma$-step transition, including the future continuation message in
\eqref{f16v21:continuation}.

For any two probability measures $\alpha,\beta$ and any Markov kernel $K$,
\[
 d_{\rm TV}(\alpha K,\beta K)
                     \le\delta(K)d_{\rm TV}(\alpha,\beta).
\]
To verify it on general spaces, write the Jordan decomposition of
$\alpha-\beta$ as $a(\alpha_+-\alpha_-)$, where
$a=d_{\rm TV}(\alpha,\beta)$ and both positive parts are probabilities.
Mixtures of kernel rows have pairwise distance at most $\delta(K)$, by
integrating the row comparison. Multiplication by $a$ proves the assertion.
Consequently diameters multiply under serial composition. Partition
$[i,j]$ into intervals of length $b_\gamma$ and one shorter remainder;
any remaining kernel has diameter at most one. Exact chain disintegration
proves \eqref{f16v21:diameter-products}.

V20's channel lemma bounds the centered conditional-expectation norm by the
square root of this diameter, for the \emph{actual} law of $Z_i$. The Markov
separator property replaces $F$ and $G$ in their covariance by their
conditional expectations on $Z_i$ and $Z_j$, respectively. These projections
decrease their variances. This proves \eqref{f16v21:past-future}. Equivalently
one can multiply the centered $L^2$ norms between the actual, possibly
unequal, successive marginal spaces. No repeated comparison of unnormalized
kernels, accumulating scalar approximation, or independence of spatial
coordinates occurs.
\end{proof}

\subsection{Two boundary messages can also be kept, rather than conditioned away}

Let $l<a\le b<r$ be time indices. Freeze the two exterior parts up to $l$
and from $r$ onward, and integrate all cells at times $l+1,\ldots,r-1$.
A missing left or right boundary means the original natural endpoint is
retained. Let $\mathcal K_{[a,b]}^{l,r}$ be the conditional marginal of the
whole central time window $[a,b]$. Set $d_-=a-l$, $d_+=r-b$, and use
$\delta^{\lfloor\infty/b_\gamma\rfloor}=0$ for a missing boundary.

\begin{theorem}[Two-sided temporal screening without an algebraic distance floor]
\label{f16v21:two-sided}
When every existing boundary is at distance at least $b_\gamma$ from the
central window,
\begin{equation}
 \delta(\mathcal K_{[a,b]}^{l,r})\le
 \min\{1,\delta_{\gamma,B}^{\lfloor d_-/b_\gamma\rfloor}
                  +\delta_{\gamma,B}^{\lfloor d_+/b_\gamma\rfloor}\}.
 \label{f16v21:two-boundary-TV}
\end{equation}
This is the diameter over the joint pair of boundary values. Thus for every
incoming law of those boundary values and every central-window $L^2$
observable $F$, the variance of its exterior conditional mean is at most
the right side of \eqref{f16v21:two-boundary-TV} times its variance in that
incoming-law joint distribution. The entire central window may have
arbitrary length; no factor counts its cells in this conclusion.
\end{theorem}
\begin{proof}
Vary the left boundary while keeping the right one fixed. Use the exact
forward continuation kernels with that fixed right factor. Each full
$b_\gamma$ step ending no later than $a$ has its target at least
$b_\gamma$ from both existing frozen boundaries: the left boundary is its
current starting slice, and the fixed right boundary is at least $r-a\ge
b_\gamma$ away. The free future interval with its actual right factor is
therefore a conditional to which Theorem~\ref{f16v21:sqrt-core} applies.
Serial contraction bounds the change in the law of $Z_a$ by
$\delta_{\gamma,B}^{\lfloor d_-/b_\gamma\rfloor}$. Given $Z_a$ and the fixed
right boundary, the continuation through $b$ is a common Markov kernel,
independent of the changed left boundary. Adding the rest of the central
window cannot enlarge total variation. This proves the left-boundary bound
for the whole window. Reverse time, with the left boundary now fixed, for
the right-boundary bound. Change two boundary pairs through one intermediate
pair and use the triangle inequality.

The pair of exterior parts enters only through the two boundary slices by
\eqref{f16v21:chain-density}. Apply V20's general incoming-law channel lemma
to the resulting joint-boundary kernel. This proves the variance assertion.
The fixed opposite boundary is retained in each transition before its
contraction is estimated; it is not treated as a harmless later conditioning
of a previously estimated output. The missing-boundary cases use only the
remaining comparison, and both missing boundaries give a constant kernel.
\end{proof}

This is a true fixed-coupling, all-distance temporal statement in the width
regime \eqref{f16v21:width}, not a fixed-order power floor. It does not assert
an unrestricted four-dimensional mixing theorem.

\begin{proposition}[V19's ball messages now decay at every sufficiently large radius]
\label{f16v21:radial-screening}
Let $s_\gamma(L)$ denote precisely V19's centered primitive-message norm and
let $D_m=3\lfloor m/2\rfloor$ be the spatial torus diameter in its metric
$d_*$. Let $\rho$ be a fixed temporal/carrier allowance. If
$h=L-D_m-\rho\ge b_\gamma+2$, then
\begin{equation}
 s_\gamma(L)\le C_r\varepsilon\,
 q_0^{\lfloor(h-1)/b_\gamma\rfloor/2}.
 \label{f16v21:radial-bound}
\end{equation}
The bound is zero when the relevant ball fills the entire history.
In particular V19's dyadic majorant is finite for every admitted coupling,
uniformly in duration at a fixed permitted spatial width. No infinite-volume
tail-triviality premise is used.
\end{proposition}
\begin{proof}
The ball of radius $L$ around a source contains every spatial cell at times
within $L-D_m$ of that source time, intersected with $[0,n]$. Its actual
primitive carrier lies in a central time window of fixed radius $\rho$.
The exterior sigma algebra of the ball is a sub-sigma algebra of the exterior
of this full temporal slab. Conditional expectation onto the smaller algebra
cannot increase the centered $L^2$ norm. Apply
Theorem~\ref{f16v21:two-sided} to the slab. Each existing exterior face is at
least $h-1$ from the central window. V16 bounds the primitive standard
deviation by $C_r\varepsilon$, and the square root of the sum of two
boundary terms costs at most $\sqrt2$ times the indicated power. Missing
physical time faces add no term. If both are missing the conditional mean of
a centered source is zero. The fixed finite sum of dyadic scales before
$D_m+\rho+b_\gamma+2$ has finite terms; afterwards exponential decay beats
the squared-radius factor in V19's criterion. This proves the last assertion
without exchanging infinite-volume and martingale limits.
\end{proof}

\subsection{Whole-spatial-source temporal rows, with the original observables}

Group the original common sources by time, so that each time block has $9B$
coordinates. Their exact Haar residual vector at time $i$, denoted $r_i$, is
measurable at times in $[i-\sigma,i+\sigma]\cap[0,n]$ for a fixed integer
$\sigma$. This follows from the ordinary score's two incident kinetic bonds
and the same temporally local port lift; its spatial inverse is not declared
finite-range. Every component and coefficient remains the one from V7/V10.

For a symmetric time-block matrix define
\begin{equation}
 \|A\|_{{\rm tr},a}=
             \sup_i\sum_{j=0}^n e^{a|i-j|}\|A_{ij}\|_{\rm op},\qquad
                  a_\gamma=\frac{\log2}{4b_\gamma}.
 \label{f16v21:time-row}
\end{equation}
The blocks act on the \emph{entire} spatial common-source space.
Symmetry makes the analogous column bound equal, so this norm controls the
full operator norm. Its subscript means temporal row, not matrix trace.

\begin{proposition}[Coherent covariance decay between arbitrary source time blocks]
\label{f16v21:coherent-decay}
In \eqref{f16v21:width}, for sufficiently large fixed-$r$ coupling,
\begin{equation}
 \|\Gamma_{\gamma,ij}\|_{\rm op}\le
 C_r\varepsilon^2
 \delta_{\gamma,B}^{\lfloor(|i-j|-2\sigma)_+/b_\gamma\rfloor/2}
 \le C_r\varepsilon^2
             2^{-\lfloor(|i-j|-2\sigma)_+/b_\gamma\rfloor}.
 \label{f16v21:coherent-decay-bound}
\end{equation}
The constants are independent of $B,n,i,j$ in this regime.
\end{proposition}
\begin{proof}
Each single-time source bank has $9B\le\gamma/4$ coordinates by
\eqref{f16v21:width}. V17's fixed polynomial-bank theorem, for example with
its fixed exponent one, gives
$\operatorname{Var}_P(v^{\mathsf T}r_i)\le C_r\varepsilon^2|v|^2$ for every
spatial source vector $v$. This is an actual native covariance bound; it is
not inferred from summing $9B$ separate variances. The inherited coefficient
operator bound and V17's bank error supply it also at either dressed time
endpoint. For two separated time carriers apply
\eqref{f16v21:past-future} with their actual intervening time gap. For
intersecting carriers ordinary covariance Cauchy--Schwarz suffices. Take
the supremum over unit spatial vectors and use $\delta_{\gamma,B}\le1/4$.
The full spatially nonlocal lift causes no difficulty because it changes no
temporal carrier outside the stated fixed allowance.
\end{proof}

By itself this proposition bounds the temporal row norm by
$C_r b_\gamma/\gamma$. The next theorem removes that logarithmic amplitude
loss by using the actual local coefficients on the short temporal band.

\begin{theorem}[Two-sided full-history operator-row asymptotics at polynomial width]
\label{f16v21:row-asymptotics}
For every fixed integer $J\ge2$, there are $C_{J,r},\gamma_{J,r}<\infty$
such that, for all $B\le c_r\gamma$, all durations $n\ge1$, all prescribed
bridge histories in the fixed window, and $\gamma\ge\gamma_{J,r}$,
\begin{equation}
 \boxed{
 \left\|\Gamma_\gamma-\sum_{q=2}^J\gamma^{-q/2}G_q\right\|_{{\rm tr},a_\gamma}
 +\left\|\mathsf M_\gamma-\sum_{q=2}^J\gamma^{-q/2}M_q\right\|_{{\rm tr},a_\gamma}
                \le C_{J,r}\gamma^{-(J+1)/2}.}
 \label{f16v21:row-expansion}
\end{equation}
In particular both complete source matrices have operator norms at most
$C_r/\gamma$, and their leading-coefficient errors are at most
$C_r\gamma^{-3/2}$ in the displayed weighted norm. No normalization by
$B(n+1)$ and no restriction on the total number of time-source coordinates
occurs. The exponent $a_\gamma$ is positive but is of order
$1/\log\gamma$, not a uniform positive exponent at the full linear width.
\end{theorem}
\begin{proof}
Write $\varepsilon=\gamma^{-1/2}$ and $G^{[J]}=\sum_{q=2}^J\varepsilon^qG_q$.
The inherited coefficient locality gives constants $a_*>0$ and $C_{q,r}$ such
that
\[
 \|(G_q)_{ij}\|_{\rm op}\le C_{q,r}e^{-a_*|i-j|},
 \qquad \|G_q\|_{{\rm tr},a_*/2}\le C_{q,r}.
\]
Indeed for each time block the scalar row and column sums are bounded by the
weighted full source-row estimate times $e^{-a_*|i-j|}$; Schur's inequality
bounds that block, and the remaining time geometric sum is finite. This is
not the invalid exchange of a source supremum with an unbounded time sum.
Increase the fixed coupling threshold so that $a_\gamma\le a_*/2$.

Take the fixed entry expansion order $K=2J+8$ in V17 and split time lags at
\begin{equation}
 D=2\sigma+
 \left\lceil\frac{4(J+2)}{\log2}\,b_\gamma\log(\varepsilon^{-1})\right\rceil.
 \label{f16v21:lag-cut}
\end{equation}
The entry error on one $9B$ by $9B$ block is at most
$C_{K,r}B\varepsilon^{K+1}\le C_{K,r}\varepsilon^{K-1}$ in operator norm,
by scalar rows and columns. Summing this error on $|i-j|\le D$ with its
weight gives at most
\[
 C_{J,r}(1+D)e^{a_\gamma D}\varepsilon^{K-1}
 \le C_{J,r}(1+\log\varepsilon^{-1})^2
                                \varepsilon^{K-J-3}
 \le C_{J,r}\varepsilon^{J+1}.
\]
Here $b_\gamma=O_r(\log\varepsilon^{-1})$ and
$e^{a_\gamma D}\le C_r\varepsilon^{-(J+2)}$ were used. The finitely many
coefficients with $J<q\le K$ cost $C_{J,r}\varepsilon^{J+1}$ in the entire
weighted row norm, rather than once for every lag.

For $|i-j|>D$, \eqref{f16v21:coherent-decay-bound} and
$2^{-\lfloor x\rfloor}\le2e^{-(\log2)x}$ bound the weighted native row tail by
\[
 C_r\varepsilon^2 b_\gamma
 \exp\!\left[-\frac{3\log2}{4b_\gamma}(D-2\sigma)\right]
 \le C_{J,r}b_\gamma\varepsilon^{3J+8}
 \le C_{J,r}\varepsilon^{J+1}.
\]
The desired coefficient tail is at most
$C_{J,r}\varepsilon^2e^{-a_*D/2}\le C_{J,r}\varepsilon^{J+1}$,
after a fixed threshold increase. Combine short and long lags. These bounds
are valid if the actual duration is shorter than $D$: the tail is then empty.

Finally the same exact Haar-contact identity gives
$\mathsf M_\gamma=-\operatorname{Off}_{>1}\Gamma_\gamma$ and
$M_q=-\operatorname{Off}_{>1}G_q$. This mask only removes entire time blocks,
so it cannot increase the temporal absolute row norm. It proves the second
term of \eqref{f16v21:row-expansion}. Higher coefficients retain their actual
native measure corrections; the core-size argument did not change them.
\end{proof}

\subsection{A fixed temporal exponent below the maximal width}

\begin{theorem}[Uniform positive temporal weight on every sublinear power cross-section]
\label{f16v21:sublinear}
Fix $0\le p<1$ and $A_0<\infty$. There is $a_{p,r}>0$, independent of
$J,B,n,\gamma$, such that, for each fixed $J\ge2$, throughout
\begin{equation}
                   B\le A_0\gamma^p,
                   \qquad\gamma\ge\gamma_{J,p,A_0,r},
 \label{f16v21:sublinear-width}
\end{equation}
the conclusion \eqref{f16v21:row-expansion} holds with $a_{p,r}$ in place of
$a_\gamma$ and constant $C_{J,p,A_0,r}$. Thus the leading-coefficient error is
$O_{p,A_0,r}(\gamma^{-3/2})$ in a \emph{fixed} exponentially weighted temporal
row norm on the entire history. The case $p=0$ means a fixed bounded spatial
cross-section; it is not an assertion for $B=1$ only.
\end{theorem}
\begin{proof}
At a fixed-$p,A_0,r$ coupling threshold,
\[
 B\le c_r\gamma,\qquad
 \delta_{\gamma,B}\le2C_{0,r}\sqrt{A_0}\,
                \gamma^{-(1-p)/2}\le\gamma^{-(1-p)/4}.
\]
For a temporal gap $g\ge2b_\gamma$, the floor bound
$\lfloor g/b_\gamma\rfloor\ge g/(2b_\gamma)$ and
\eqref{f16v21:buffer-growth} give
\[
 \delta_{\gamma,B}^{\lfloor g/b_\gamma\rfloor/2}
                  \le e^{-\lambda_{p,r}g},
             \qquad\lambda_{p,r}=\frac{1-p}{16C_{b,r}}>0.
\]
Choose once
\[
 0<a_{p,r}\le
 \min\{\lambda_{p,r}/4,\ a_*/4,\ (4C_{b,r})^{-1}\}.
\]
This choice does not depend on $J$. Repeat the short/long lag proof with
$K=2J+8$ and
\[
 D=2b_\gamma+2\sigma+
       \left\lceil\frac{J+3}{a_{p,r}}\log(\varepsilon^{-1})\right\rceil.
\]
Then $D=O_{J,p,r}(\log\varepsilon^{-1})$ and
$e^{a_{p,r}D}\le C_r\varepsilon^{-(J+4)}$ because
$b_\gamma\le2C_{b,r}\log\varepsilon^{-1}$.
The near error is at most
$C(1+D)\varepsilon^{K-J-5}=O(\varepsilon^{J+1})$.
The higher coefficients retain their uniform weighted row bounds. On the
far band the new genuine native exponential has rate at least $4a_{p,r}$,
so its weighted tail is at most
$C\varepsilon^2e^{-3a_{p,r}(D-2\sigma)}=O(\varepsilon^{J+1})$;
the coefficient tail has the same bound. There is now no factor $b_\gamma$
from that geometric sum. The time mask is unchanged. This proves the
statement with no order-dependent loss of temporal exponent.
\end{proof}

\subsection{What has closed for the shell return, and what remains spatial}

\begin{proposition}[The positive shell edge is small in the admitted width regime]
\label{f16v21:return-closure}
Fix $J\ge2$ and use a logarithmic radius satisfying V19's local operator-error
conditions at that order. In \eqref{f16v21:width}, with a sufficiently large
fixed-order coupling threshold,
\begin{equation}
       \|\mathsf R^{\rm sh}_{\gamma,L_\gamma}\|_{\rm op}
                          \le C_{J,r}\gamma^{-(J+1)/2}.
 \label{f16v21:return-small}
\end{equation}
For any desired fixed power the radius and fixed orders can therefore be
chosen so that this \emph{actual} positive return is smaller than that power.
The choice of radius is not asserted independent of the requested power.
\end{proposition}
\begin{proof}
V19's exact positive split is
$\Gamma_\gamma-G^{[J]}=\mathsf R^{\rm sh}_{\gamma,L_\gamma}
+\mathsf E_{\gamma,J,M}$, with
$\|\mathsf E_{\gamma,J,M}\|_{\rm op}\le C\gamma^{-(J+1)/2}$ at its specified
logarithmic radius. The new two-sided row estimate bounds
$\|\Gamma_\gamma-G^{[J]}\|_{\rm op}$ by the same power. The triangle
inequality proves \eqref{f16v21:return-small}, taking one fixed $M$.
This argument bounds the previous positive matrix rather than redefining it.
It is consistent with the direct all-radius screening in
Proposition~\ref{f16v21:radial-screening}; neither proof uses an assumed native
sampler gap.
\end{proof}

The spatial improvement is substantial but explicit. V9--V10's native
operator estimates require an entry bound of the form
$e^{KB}/\sqrt\gamma\le1/100$. Here the width can be as large as $c_r\gamma$
for a deteriorating temporal exponent, or $A_0\gamma^p$, $p<1$, for a fixed
positive temporal exponent. The duration and total source count are
unrestricted in either case. At fixed $\gamma$, however, arbitrary spatial
volume is still not allowed. The constants are proved finite from the
inherited local estimates, not optimized numerical thresholds.

This is not the false tensorization warned against in V20. All outputs on a
time slice are included in one vector and their size is paid once. The
intermediate sigma algebras are actual complete-slice separators, and every
future endpoint message stays inside each conditional channel. Arbitrary
conditioning inside a step's buffer is not used. Two-sided screening
re-estimates the channels with the opposite boundary retained, rather than
conditioning a previously contracted output after the fact.

The next spatial question is sharper: replace the cost of a \emph{complete}
separator of $B$ cells by a local spatial composition with its shared
boundary data retained. The square-root bound alone cannot do that once
$B/\gamma$ is unbounded. Repeating the same temporal composition would not
remove this remaining cross-section cost. V19's unrestricted-volume positive
return, V12's separate spatial-cylinder posterior, and a gap for the original
single-port/chord sampler are therefore not declared closed by the current
restricted-width result.

Nothing here integrates the retained physical bridges outside their prescribed
window, proves root Gauss coverage, controls the infrared contact term, or
identifies a conditional fast-channel contraction with the same-top physical
Wilson transfer. A decay exponent in lattice time, especially one that can
deteriorate as $1/\log\gamma$, is not a calibrated physical mass. The
interacting continuum construction and its nontriviality remain distinct
obligations. The native measure and every original source definition are
unchanged throughout.

\subsection{Diagnostics, predecessor preservation, and proof status}

The accompanying diagnostic is
\begin{center}\small\path{verify_ym_f16_v21_temporal_composition.py}.\end{center}
It checks the vector regrouping behind the square-root size bound, exact
nonstationary path disintegrations with unequal endpoint factors, conditioned
continuation kernels, serial diameter contraction, two-boundary screening,
source-window support, and the lag/order arithmetic. The new diagnostic passed 351 exact assertions. The report separates
exact rational and symbolic fixtures from analytic estimates. No fixture is
presented as native Wilson quadrature, a numerical enclosure of the entry
constants, or a proof-assistant certification of V20's premises.

The supplied V20 cumulative source is preserved byte for byte except for one
title-version change and insertion of this exact appendix before its final
document-ending command. The package records exact recovery, input hashes,
actual reruns of all seventeen available V4--V20 diagnostics, the new diagnostic, compilation,
and rendered-page inspection. Old sources are not silently corrected. These
new analytic deductions depend on the identified inherited pin, preparation,
Stein, source, and coefficient results; the archive as a whole is not
independently recertified. The next scheduled companion checkpoint is V25.

\begin{firewall}
V21 proves a square-root core-size total-variation bound under every compact
exterior, a valid serial composition through complete time slices, and
all-distance temporal contraction in the explicit regime $B\le c_r\gamma$.
It proves full-history two-sided covariance and memory asymptotics in a
weighted temporal operator-row norm, with a fixed positive temporal exponent
on every strictly sublinear power cross-section. The old positive shell
return is correspondingly bounded in that regime. It does not prove
unrestricted spatial-volume contraction at fixed coupling, a uniform gap for
the original local sampler, a physical transfer gap, or an interacting
continuum Yang--Mills mass gap. All endpoint normalizers and all compact fast
fields are retained; the bridge-history window and physical clock obligations
are unchanged.
\end{firewall}
% END F16 V21 APPEND


% BEGIN F16 V22 APPEND -- 2026-09-17
% Insert immediately before the final document-ending command in V21.
% No predecessor theorem, normalization, source, or bibliography is altered.

\clearpage
\section{V22: native boundary comparison beyond the Gaussian channel floor}
\label{f16v22:context}

V21 pays the discrepancy between a native core and a Gaussian marginal before
composing its temporal channels. That discrepancy includes a local nonlinear
correction which is present even when the exterior has no influence. It is
therefore not the right floor for comparing \emph{two native boundary rows}.
We retain the nonlinear core integral instead. A small-angular comparison is
used only inside a finite integrated buffer; the actual conditional law pays
its complement, uniformly over every compact exterior. The resulting channel
bound has an exponentially small large-field term, not an intrinsic
$\gamma^{-1/2}$ floor.

This is an upstream change in the reference used for boundary comparison,
not another order of Gaussian perturbation theory. It permits an
arbitrary prescribed polynomial spatial cross-section, with a fixed positive
temporal exponent, for the full native covariance expansion. A separate
all-observable temporal-channel statement allows an exponentially growing
cross-section. Neither assertion is an unrestricted-spatial-volume theorem
at fixed coupling. The local spatial composition problem remains explicit.

\subsection{Dependency review and the native comparison being changed}

The supplied V21 appendix and audit were read in full. Its remaining
complete-separator cost, V20's conditional preparation, and V4's product-ball
one-form covariance identity were checked against the cumulative source.
Two scoped semantic archive searches returned no results. Exact mounted-text
searches and the identified f15 V22 passages were then used. F15 V22 already
proves a restricted core density comparison through a pinned-complement
covariance identity, and a transfer to the full law with a confidence
exception. We credit both mechanisms. The new ingredients are their
application to the current unmoving specification with \emph{uniform}
compact-boundary preparation, a square-root core-size response bound, and an
exponentially small finite-buffer exception. No archive-wide novelty or
independent recertification claim is made.

For boundary-aware analytic context, D.~Le Peutrec,
\href{https://arxiv.org/abs/1607.08714}{\emph{On Witten Laplacians and
Brascamp--Lieb's inequality on manifolds with boundary}}, derives the
Neumann/Dirichlet form identities for smooth boundaries. Its introductory
identity and boundary hypotheses were checked. The product-ball realization
with corners used below is the one already specified in V4 and f15 V13;
we do not pretend that a smooth-boundary theorem alone states that product
result. The elementary resolvent estimate needed here is supplied below.

A spatial-composition literature check also examined the definitions and
initial theorems of Spinka's
\href{https://arxiv.org/abs/1803.10578}{\emph{Finitary codings for spatial
mixing Markov random fields}}. The distinction between weak spatial mixing,
strong spatial mixing, and a grand coupling is relevant. Its finite-valued,
translation-invariant coding theorem is not applied to this inhomogeneous
compact law. The abstract of Chazottes--Redig--V\"ollering,
\href{https://arxiv.org/abs/1012.2489}{\emph{Poincar\'e inequality for Markov
random fields via disagreement percolation}}, was screened only. Its required
subcritical disagreement coupling has not been constructed here.

The import memoranda's separate ledgers for power gain and locality cost
remain useful; neither a physical gap nor an all-scale cluster theorem is
an input. The current proof uses V5/V16's exact specification and bounds,
V20's all-boundary moment preparation, and the inherited product-convex
covariance realization. V17 and V10 are used only for the subsequent source
asymptotics. The next scheduled companion checkpoint remains V25.

\subsection{A comparison domain, not a replacement for the native law}

Use V20's interaction graph $\mathcal G$, distance $d_{\mathcal G}$, and
complete unmoving cells. Write
\[
 P=P_{\gamma,\mathbf y},\quad \varepsilon=\gamma^{-1/2},\quad
 \mathcal N=B(n+1),\quad \max_{i,e}|y_{i,e}|\le r.
\]
A cell has 36 or 15 real coordinates. The full measured potential on its
open principal charts is $V_\gamma=U_\gamma-\log J^f_\gamma$. Its Gaussian
fast precision is the same $H$, with $h_-I\preceq H\preceq h_+I$.
Every non-anchor interaction has graph diameter one. No new logarithm of a
matrix product is used in differentiating the local exponential words.

Choose a fixed angular radius $0<\eta_*<1/4$, small enough as specified
below, and put
\[
 \mathcal D_{\gamma,c}=\{\zeta_c:|\zeta_c|\le\eta_*\sqrt\gamma\},
 \qquad \mathcal D_{\gamma,W}=\prod_{c\in W}\mathcal D_{\gamma,c}.
\]
Each is a ball in a \emph{whole unmoving cell}. All its original group
coordinates have angular radius at most $\eta_*$, so their Haar factors,
including the endpoint powers $1/2$, are smooth and positive nearby.
This is an auxiliary comparison domain. It is neither V4's event
$\mathcal O_R$ nor a redefinition of $p_{\rm good}$.

For a set $W$, let $C=\partial_{\mathcal G}W$ be its actual exterior word
collar, with redundant adjacent cells harmless. For $z_C$ in the corresponding
product of small balls define the \emph{restricted native} conditional
\begin{equation}
 \pi_W^z(d\zeta_W)=
 \frac{1_{\mathcal D_{\gamma,W}}e^{-U_W(\zeta_W;z)}
                         \prod_{c\in W}d\nu_{\gamma,c}}
      {Z_W^{\rm box}(z)}.
 \label{f16v22:restricted-law}
\end{equation}
Every term meeting $W$ is retained, including crossing terms. Exterior-only
constants cancel in this conditional probability, not from the full exterior
posterior. The free domain is independent of $z$. Its measured potential
includes the complete original Haar factors on $W$.

\begin{proposition}[Uniform restricted convexity and a spatial inverse]
\label{f16v22:convexity}
There are $\eta_*>0$, $\gamma_r<\infty$, $\lambda>0$, and $M>\lambda$, with
constants independent of $W,B,n$, such that for $\gamma\ge\gamma_r$ and every
admitted $z$, throughout
$\mathcal D_{\gamma,W}$,
\begin{equation}
       \lambda I\preceq\nabla_W^2V_\gamma^z\preceq MI.
 \label{f16v22:H-bounds}
\end{equation}
Scalar mixed derivatives involving a free and a collar coordinate have
uniform absolute row and column bounds and vanish beyond graph distance one.
The same Hessian statements hold when any subset of the free cells is fixed
anywhere in its small balls.

The reflecting scalar law satisfies
$\operatorname{Var}_{\pi_W^z}f\le\lambda^{-1}\mathbb E|\nabla f|^2$.
For its absolute one-form covariance realization $\mathfrak H$, and cell
component projections $P_c,P_d$ on any remaining free set,
\begin{equation}
 \|P_c\mathfrak H^{-1}P_d\|
       \le\lambda^{-1}(1-\lambda/M)^{d_{\mathcal G}(c,d)}.
 \label{f16v22:inverse}
\end{equation}
Distances in the full graph give a valid, possibly weaker, bound after
pinning. The covariance identity is the inherited one from V4.
\end{proposition}
\begin{proof}
For each original classical term, two coordinate derivatives of
$\gamma f(w/\sqrt\gamma)$ have a uniformly bounded third-derivative remainder
relative to its quadratic limit. On the displayed product, every local fast
argument has angular magnitude at most a fixed multiple of $\eta_*$; each
bridge argument is at most $r/\sqrt\gamma$. Bounded word length, local
linear-map norms, and incidence give
\[
 \|\nabla_W^2 U_W-H_{WW}\|\le C(\eta_*+r/\sqrt\gamma).
\]
The estimate is by bounded symmetric rows, not by the number of cells.
The exactly quadratic endpoint terms are retained. The Hessian of each
$-w\log j(\zeta_g/\sqrt\gamma)$ is bounded by $C/\gamma$ on this domain,
with the original $w=1$ or $1/2$. Choose $\eta_*$ first and then a fixed-$r$
threshold so that the total difference is at most $h_-/2$. One may take
$\lambda=h_-/2$ and enlarge $M$. Mixed rows with collar coordinates follow
from the same local derivative bounds. Principal submatrices preserve the
lower and upper bounds after more cells are fixed. No convexity outside
these small angular domains is asserted.

Use exactly the product absolute form described in
\eqref{f16v4:witten-form}: its derivative and ball-boundary part
$\mathfrak D$ is nonnegative and block diagonal in cell component labels.
The outward second fundamental forms are nonnegative; the component at its
own ball boundary is tangential. Cell projections preserve the form domain.
The scalar/one-form intertwining gives
\[
 \operatorname{Cov}(f,g)=
       \langle\nabla f,(\mathfrak D+\nabla^2V)^{-1}\nabla g\rangle.
\]
The inverse lower bound gives the Poincar\'e inequality. These statements use
the inherited product-ball boundary realization, including its corners,
not boundary-free integration by parts.

Put $\mathfrak A=\mathfrak D+MI$ and $\mathfrak B=MI-\nabla^2V$.
Then $0\preceq\mathfrak B\preceq(M-\lambda)I$,
$\mathfrak A^{-1/2}$ preserves cell component labels, and
\[
 \mathfrak H^{-1}=\mathfrak A^{-1/2}
 \sum_{l\ge0}(\mathfrak A^{-1/2}\mathfrak B\mathfrak A^{-1/2})^l
 \mathfrak A^{-1/2}.
\]
The series converges in operator norm, with ratio at most
$1-\lambda/M$. Its $l$th term cannot connect component labels at distance
greater than $l$. Summing the geometric tail proves \eqref{f16v22:inverse}.
The differential coefficients may depend on all free variables; the range
statement is about \emph{derivative indices}, not finite propagation in
configuration space. Product balls of either cell dimension and their
principal-subset products satisfy the same argument.
\end{proof}

\subsection{Compare native rows without paying their common nonlinear correction}

Let $A$ denote all coordinates of a nonempty cell set $\mathcal A\subset W$,
$k=|\mathcal A|$, and assume $d_{\mathcal G}(\mathcal A,C)\ge\ell+1$,
$\ell\ge3$. Write $\pi_{W,A}^z$ for the core marginal of
\eqref{f16v22:restricted-law}. When $C$ is empty this marginal is independent
of exterior data and every boundary-comparison error below is zero.

\begin{theorem}[Restricted native boundary response with square-root output cost]
\label{f16v22:box-TV}
There are $C_r<\infty$ and $a>0$, independent of $W,A,B,n,\gamma$, such that
\begin{equation}
 \sup_{z,z'\in\mathcal D_{\gamma,C}}
 d_{\rm TV}(\pi_{W,A}^z,\pi_{W,A}^{z'})
            \le\min\{1,C_r\sqrt{\gamma k}\,e^{-a\ell}\}.
 \label{f16v22:box-boundary}
\end{equation}
The distance convention is $\sup_E|\mu(E)-\nu(E)|$. The reference in this
comparison is native, not Gaussian. No expansion of the restricted
partition function is used.
\end{theorem}
\begin{proof}
Join $z$ to $z'$ by its straight segment $z_s$. It stays in the product of
collar balls and has scalar displacements at most $2\eta_*\sqrt\gamma$.
Let $T=W\setminus\mathcal A$. At fixed core $x$, disintegrate the restricted
native density over the \emph{fixed} product domain on $T$. Write $p_s(x)$
for the normalized core density and
\[
 S_s=\partial_s V_W^{z_s},\qquad
 q_s(x)=\mathbb E_{\pi_T^{x,z_s}}S_s.
\]
All denominators are positive and all differentiated integrands smooth on
these compact sets. Normalized differentiation gives, exactly,
\begin{equation}
 \partial_s\log p_s(x)=-q_s(x)+\mathbb E_{\pi_{W,A}^{z_s}}q_s,
 \qquad
 \partial_{x_i}q_s=\mathbb E\partial_{x_i}S_s
                       -\operatorname{Cov}(S_s,\partial_{x_i}V_W^{z_s}).
 \label{f16v22:score-identity}
\end{equation}
The covariance is in the pinned-complement conditional on $T$. In particular
both the inner and outer normalizers have been retained. This is the same
mixed-normalization mechanism as f15 V22, on the present conditional.

The direct derivative in the second identity vanishes by separation.
For a collar scalar coordinate $j$, the free gradients of
$\partial_{z_j}V$ and $\partial_{x_i}V$ have fixed carriers within one graph
step of their respective cells and bounded norms. The inherited covariance
identity and \eqref{f16v22:inverse} consequently give
\[
 |\operatorname{Cov}_{\pi_T^{x,z_s}}
                    (\partial_{z_j}V,\partial_{x_i}V)|
            \le C e^{-a_0d_{\mathcal G}(c(i),c(j))}
\]
for some $a_0>0$, after absorbing the fixed two-step allowance. The case of
empty $T$ has zero covariance. Fix $0<a<a_0/2$. Polynomial graph growth and
bounded coordinate multiplicity imply
\[
 \sum_{j\in C}e^{-a_0d_{\mathcal G}(c(i),c(j))}
                  \le C e^{-a\ell}.
\]
Summing the actual collar displacement therefore proves
\begin{equation}
 \sup_x|\partial_{x_i}q_s(x)|\le C_r\sqrt\gamma e^{-a\ell},
 \qquad \sup_x|\nabla_Aq_s(x)|^2\le C_r\gamma k e^{-2a\ell}.
 \label{f16v22:marginal-score-gradient}
\end{equation}
This uses a distance-summed row bound, not the collar's total cardinality.

The core marginal inherits the full restricted Poincar\'e inequality by
lifting a function of $x_A$ to $W$. Hence
$\operatorname{Var}_{\pi_{W,A}^{z_s}}q_s\le C_r\gamma k e^{-2a\ell}$.
For every bounded measurable $f$ with $|f|\le1$, differentiation of its
expectation, not of $f$, yields
\[
 \left|\partial_s\mathbb E_{\pi_{W,A}^{z_s}}f\right|
       =|\operatorname{Cov}(f,q_s)|
       \le C_r\sqrt{\gamma k}e^{-a\ell}.
\]
Finite-dimensional compact differentiation justifies this also for
measurable bounded $f$ by dominated convergence. Integrate in $s$ and take
indicator tests. This proves \eqref{f16v22:box-boundary}. No derivative norm
of the observed function and no native \emph{full-law} Poincar\'e inequality
has been assumed.
\end{proof}

\subsection{Pay the entire comparison buffer under every actual compact boundary}

Use V20's complete $P_\Lambda^q$, with an arbitrary compact exterior $q$.
For $\mathcal A\subseteq\Lambda$ and integer $\ell\ge3$, set
\[
 W=\{c:d_{\mathcal G}(c,\mathcal A)\le\ell\},\qquad
 C=\partial_{\mathcal G}W,\qquad
 G_{W,C}=\{\zeta_c\in\mathcal D_{\gamma,c}\text{ for all }c\in W\cup C\}.
\]
Balls and collars refer to the actual finite graph. Assume
\begin{equation}
       L=d_{\mathcal G}(\mathcal A,\Lambda^c)\ge\ell+D_*+2,
 \label{f16v22:buffer-geometry}
\end{equation}
where $D_*$ is fixed below. In particular $W\cup C\subseteq\Lambda$;
no artificially frozen time endpoint is introduced.

\begin{theorem}[Complete buffered rows with an exponentially small native exception]
\label{f16v22:native-TV}
There are $D_*<\infty$, $c_*>0$, and $C_r<\infty$, independent of ambient
volume, duration, core size, and frozen exterior, such that for
$\gamma\ge\max\{1+r^2,\gamma_r\}$ and geometry \eqref{f16v22:buffer-geometry},
\begin{equation}
 \boxed{
 \sup_q d_{\rm TV}(K_{\Lambda,A}(q),\lambda_{W,A})
  \le\min\{1,C_r[\sqrt{\gamma k}e^{-a\ell}
                       +k(1+\ell)^4e^{-c_*\gamma}]\},
 \quad \lambda_{W,A}:=\pi_{W,A}^{0}.}
 \label{f16v22:native-buffer-bound}
\end{equation}
The reference is the explicitly normalized restricted \emph{native} integral
at zero collar. It can depend on $\gamma,W,\mathbf y$ but not on $q$.
The target is the complete native core marginal. All compact fast fields
are retained there. Consequently its Dobrushin diameter satisfies the same
bound with $C_r$ enlarged, as does its squared maximal correlation for
\emph{every} incoming exterior law.
\end{theorem}
\begin{proof}
Let $M_r,B_r,\kappa_r$ be the conditional preparation constants from
\eqref{f16v20:prepared-mgf}. Choose a fixed $D_*$ so that
$B_r e^{-\kappa_rD_*}\le\eta_*^2/2$, and put
$c_*=t_0\eta_*^2/2>0$. At any cell of depth at least $D_*$,
exponential Markov in that actual conditional gives
\[
 P_\Lambda^q\{e_c>\eta_*^2\gamma\}
     \le\exp(t_0M_r-c_*\gamma).
\]
This is the fixed-depth angular tail, not a claim that arbitrary boundary
cells have uniform thermal moments. Every cell of $W\cup C$ has this depth
by \eqref{f16v22:buffer-geometry}. Polynomial ball growth gives
$|W\cup C|\le Ck(1+\ell)^4$. The union bound therefore yields
\begin{equation}
 \sup_q P_\Lambda^q(G_{W,C}^{\,c})
                 \le C_r k(1+\ell)^4e^{-c_*\gamma}.
 \label{f16v22:buffer-tail}
\end{equation}
This event involves only the chosen finite comparison buffer, not all cells
of the history.

Condition the \emph{actual} law further on $G_{W,C}$; it has positive
probability. Given all variables outside $W$, its conditional inside $W$ is
exactly \eqref{f16v22:restricted-law} at the observed small collar $z_C$.
Thus its core marginal is a mixture of $\pi_{W,A}^z$ over a posterior that
retains the conditional normalizers $Z_W^{\rm box}(z)$ and the actual gate
weights. No property of that posterior except its support is needed.
Every row in this mixture is within \eqref{f16v22:box-boundary} of
$\lambda_{W,A}$. The complete law and its restriction differ in total
variation by $P_\Lambda^q(G_{W,C}^{\,c})$, and taking a marginal cannot
increase this distance. The triangle inequality proves
\eqref{f16v22:native-buffer-bound}. There is no division by a small
estimated good-event probability.

Compare two native rows through this \emph{same} reference to bound the
diameter. V20's proved general-space channel lemma then gives its
all-incoming-law, all-$L^2$ maximal-correlation consequence. The good event
is not commuted through any original update or differentiated in a retained
bridge source. It is eliminated by the exact mixture argument before the
claim about the full native kernel is made.
\end{proof}

This removes a comparison floor, not all large fields from the target law.
The nonlinear core correction is present on both restricted native boundary
rows and cancels in their comparison. It was not zero when either row was
compared with $Q_A$. The estimates \eqref{f16v22:box-boundary} and
\eqref{f16v22:buffer-tail} keep these two different errors separate.

\begin{proposition}[Any fixed polynomial output and any fixed power of contraction]
\label{f16v22:polynomial-core}
Fix $p\ge0$, $s>0$ and $A_0<\infty$. There is
$b_{\gamma,p,s}=O_{p,s,r}(\log\gamma)$ such that every core with
$k\le A_0\gamma^p$ and depth at least $b_{\gamma,p,s}$ has
\begin{equation}
          \delta(K_{\Lambda,A})\le\gamma^{-s}
 \label{f16v22:arbitrary-power-channel}
\end{equation}
for $\gamma\ge\gamma_{p,s,A_0,r}$. There is no restriction on the ambient
volume or duration in this local-core statement.
\end{proposition}
\begin{proof}
Take
\[
 \ell=3+\left\lceil a^{-1}\bigl(s+(p+1)/2+1\bigr)\log\gamma\right\rceil,
 \qquad b_{\gamma,p,s}=\ell+D_*+2.
\]
The first term of \eqref{f16v22:native-buffer-bound} is at most a fixed
constant times $\gamma^{-s-1}$. The second is a fixed polynomial times
$e^{-c_*\gamma}$, hence also at most that power at a fixed threshold.
Absorb the constants and the two-row factor by increasing the threshold.
All exponents were fixed before the coupling limit. No order-dependent
Taylor expansion was used in this proof.
\end{proof}

\subsection{Valid temporal composition on much larger separators}

We now reuse, rather than reprove, V21's exact complete-slice Markov
factorization, including every continuation message and endpoint factor.
Its serial and two-boundary proofs require a uniform diameter for a
buffered slice under every compact exterior, including the fixed opposite
endpoint. Theorem~\ref{f16v22:native-TV} supplies exactly that input with a
better boundary-comparison scale.

\begin{theorem}[All polynomial spatial widths have a fixed temporal decay rate]
\label{f16v22:polynomial-time}
Fix $p\ge0$ and $A_0<\infty$, and assume
\begin{equation}
                B\le A_0\gamma^p.
 \label{f16v22:polynomial-width}
\end{equation}
At a fixed-$p,A_0,r$ coupling threshold, set $b_\gamma=b_{\gamma,p,2}$.
There is $C_{b,p,r}$ with $b_\gamma\le C_{b,p,r}\log\gamma$, and the
actual complete-slice transitions satisfy
\begin{equation}
 \delta(K_{i,j})\le\gamma^{-2\lfloor(j-i)/b_\gamma\rfloor},
 \qquad
 |\operatorname{Cov}_P(F,G)|
 \le\gamma^{-\lfloor(j-i)/b_\gamma\rfloor}
                         \sqrt{\operatorname{Var}_PF\operatorname{Var}_PG}.
 \label{f16v22:polynomial-time-bound}
\end{equation}
Here $F$ and $G$ are arbitrary past/future $L^2$ observables separated by
$i,j$, as in V21. For $j-i\ge2b_\gamma$, their normalized correlation is
at most $e^{-\lambda_{p,r}(j-i)}$ for a fixed $\lambda_{p,r}>0$.
Two-sided central-window screening holds with the sum of the two respective
powers in the diameter bound and no central-window-size factor.
\end{theorem}
\begin{proof}
Observe one complete output slice, so that its core size is $k=B$.
Given the complete past, the future interval is the original conditional
with its natural endpoint, not a transition with that endpoint discarded.
A step of length $b_\gamma$ has graph depth at least $b_\gamma$. Apply
Proposition~\ref{f16v22:polynomial-core} with $s=2$.
V21's exact serial proof multiplies these diameters and gives the first
inequality. Its all-observable channel argument gives the square-root
power in the covariance bound. For a gap $g\ge2b_\gamma$,
\[
 \gamma^{-\lfloor g/b_\gamma\rfloor}
       \le\exp[-g\log\gamma/(2b_\gamma)]
       \le\exp[-g/(2C_{b,p,r})].
\]
A smaller fixed rate can be selected for later estimates. The two-boundary
argument of Theorem~\ref{f16v21:two-sided} applies without modification:
retain the right boundary in every forward conditional, then reverse time
and retain the left boundary. No previously contracted output is
subsequently conditioned inside its buffer. All constants are independent
of the duration, which is unrestricted.
\end{proof}

\begin{theorem}[An exponential-width channel result, distinct from source asymptotics]
\label{f16v22:exponential-time}
There are $\beta_r>0$, $c_r'>0$, and a buffer $b_\gamma^{\rm exp}=O_r(\gamma)$
such that, for
\begin{equation}
                B\le e^{\beta_r\gamma},\qquad \gamma\ge\gamma_r',
 \label{f16v22:exponential-width}
\end{equation}
the actual $b_\gamma^{\rm exp}$-step slice diameter is at most
$e^{-c_r'\gamma}$. Hence all past/future $L^2$ correlations have a fixed
positive exponential decay rate at gaps at least $2b_\gamma^{\rm exp}$.
The same two-boundary conclusion holds. This is not a covariance expansion
in this exponential-width regime.
\end{theorem}
\begin{proof}
In \eqref{f16v22:native-buffer-bound}, choose
\[
 \beta_r=c_*/8,\qquad
 \ell=3+\lceil c_*\gamma/(4a)\rceil,\qquad
 b_\gamma^{\rm exp}=\ell+D_*+2.
\]
For $k=B\le e^{c_*\gamma/8}$ the two error terms are bounded respectively
by $C_r\sqrt\gamma e^{-3c_*\gamma/16}$ and
$C_r(1+\gamma)^4e^{-7c_*\gamma/8}$. After a fixed threshold increase the
two-row diameter is at most $e^{-c_*\gamma/16}$. Take $c_r'=c_*/16$.
The exact serial argument gives a correlation factor
$\exp[-c_r'\gamma\lfloor g/b_\gamma^{\rm exp}\rfloor/2]$.
Since $b_\gamma^{\rm exp}\le C_r'\gamma$, it is at most
$\exp[-c_r'g/(4C_r')]$ when $g\ge2b_\gamma^{\rm exp}$.
The opposite-boundary proof is the same as above.

For gaps shorter than this buffer the general bound is only one. An
exponentially wide source slice is not a polynomial bank in V17. Accordingly
neither a coupling-independent short-gap prefactor nor an all-source
operator expansion at this width is inferred from this theorem.
\end{proof}

\subsection{Fixed-weight operator asymptotics at every polynomial cross-section}

Retain V21's temporal-row norm, original common-source blocks, and fixed
source time-carrier allowance $\sigma$. For a symmetric time-block matrix,
\[
 \|A\|_{{\rm tr},b}=\sup_i\sum_j e^{b|i-j|}\|A_{ij}\|_{\rm op}.
\]
This is not a trace norm. Each block has $9B$ coordinates and includes the
whole spatially nonlocal Haar lift, whose time support remains bounded.

\begin{theorem}[Two-sided native expansion for every fixed polynomial width]
\label{f16v22:row-expansion}
For every fixed $p\ge0$ there is $a_{p,r}>0$, independent of expansion order,
coupling, volume, and duration, such that for every fixed $J\ge2$ and
$A_0<\infty$, throughout \eqref{f16v22:polynomial-width} at a sufficiently
large $\gamma_{J,p,A_0,r}$,
\begin{equation}
 \boxed{
 \begin{aligned}
 &\left\|\Gamma_\gamma-\sum_{q=2}^J\gamma^{-q/2}G_q
                          \right\|_{{\rm tr},a_{p,r}}\\
 &\quad+\left\|\mathsf M_\gamma-\sum_{q=2}^J\gamma^{-q/2}M_q
                          \right\|_{{\rm tr},a_{p,r}}
              \le C_{J,p,A_0,r}\gamma^{-(J+1)/2}.
 \end{aligned}}
 \label{f16v22:row-asymptotics}
\end{equation}
The coefficients, sources, complete native law and endpoints are unchanged.
Their full operator norms are $O_{p,A_0,r}(\gamma^{-1})$. The
leading-coefficient errors are $O_{p,A_0,r}(\gamma^{-3/2})$, for every
duration. There is no division by $\mathcal N$ and no source-vector average.
\end{theorem}
\begin{proof}
V17's coherent bank theorem applies to a single time block: for large enough
fixed-$p,A_0$ coupling, $9B\le\gamma^{p+1}$. Thus every spatial vector $v$
satisfies
\[
 \operatorname{Var}(v^{\mathsf T}r_i)\le C_{p,A_0,r}\varepsilon^2|v|^2.
\]
This is the native bank theorem, not the sum of separate component variances.
Together with Theorem~\ref{f16v22:polynomial-time}, it gives
\[
 \|\Gamma_{\gamma,ij}\|\le C_{p,A_0,r}\varepsilon^2,
 \qquad
 \|\Gamma_{\gamma,ij}\|\le
 C_{p,A_0,r}\varepsilon^2 e^{-\lambda_{p,r}(|i-j|-2\sigma)}
 \quad\text{if }|i-j|-2\sigma\ge2b_\gamma.
\]
The Gaussian coefficient locality of V10 gives a fixed $a_*>0$ and
$\|G_q\|_{{\rm tr},a_*/2}\le C_{q,r}$ at every fixed order, just as in V21.
Choose once
\[
 0<a_{p,r}\le\min\{\lambda_{p,r}/4,\ a_*/4,\ (4C_{b,p,r})^{-1}\}.
\]
This choice is independent of $J$ and $A_0$; the entry threshold may depend
on both. Set $b=a_{p,r}$ and
\[
 D=2b_\gamma+2\sigma+
       \left\lceil (J+4)b^{-1}\log(\varepsilon^{-1})\right\rceil,
 \qquad K=2J+2\lceil p\rceil+12.
\]
Here $D=O_{J,p,r}(\log\varepsilon^{-1})$ and
$e^{bD}\le C_{p,r}\varepsilon^{-(J+5)}$, since
$b_\gamma\le2C_{b,p,r}\log\varepsilon^{-1}$.

On $|i-j|\le D$, V17's order-$K$ entry error has block operator norm at
most $C_{K,p,A_0,r}B\varepsilon^{K+1}\le
C_{K,p,A_0,r}\varepsilon^{K+1-2p}$ by scalar rows and columns. Summing over
the near band with its weight costs at most
\[
 C(1+D)e^{bD}\varepsilon^{K+1-2p}
      \le C(1+\log\varepsilon^{-1})\varepsilon^{K-J-2p-4}
      \le C\varepsilon^{J+1}.
\]
The coefficients of orders $J+1$ through $K$ cost
$C\varepsilon^{J+1}$ in their entire weighted row norm, not once per lag.

On the far band the genuine native decay has rate at least $4b$, so its
weighted sum is at most
$C\varepsilon^2 e^{-3b(D-2\sigma)}=O(\varepsilon^{J+1})$.
The desired coefficient tail is bounded the same way, using its rate at
least $4b$. Short actual histories simply have an empty far band.
This proves the first term of \eqref{f16v22:row-asymptotics}.
The exact identity $\mathsf M_\gamma=-\operatorname{Off}_{>1}\Gamma_\gamma$
and $M_q=-\operatorname{Off}_{>1}G_q$ removes entire time blocks and cannot
increase this row norm, proving the second. Symmetry and Schur's bound give
the full operator consequences. No higher coefficient is required to be
positive and no hard comparison gate is differentiated in a bridge source.
\end{proof}

V19's same positive shell return is therefore small, at its stipulated
logarithmic radius, throughout \emph{every fixed polynomial} width regime:
subtract its unchanged local-error decomposition from
\eqref{f16v22:row-asymptotics}. The resulting norm is
$C_{J,p,A_0,r}\gamma^{-(J+1)/2}$. This is an enlargement of V21's
$B\le c_r\gamma$ closure, not a redefinition of that return or of V12's
different one-common-exterior cylinder Gram.

\subsection{What remains spatial and what the comparison has not changed}

The new conditional comparison removes an intrinsic native--Gaussian error
from a native--native boundary question. Its two actual costs are
\[
           \sqrt{\gamma k}e^{-a\ell}
                  \quad\text{and}\quad k(1+\ell)^4e^{-c_*\gamma}.
\]
The first is a boundary-response estimate within one normalized nonlinear
comparison law. The second is the probability of leaving its finite buffer,
under the complete conditional at the actual exterior. An artificial
convex extension, a new physical action, and an assumed full-native
functional inequality are not used.

The local-core theorem is uniform in ambient volume. Its complete-slice
applications are not: at fixed $\gamma$, sufficiently large $B$ makes the
exception term ineffective. The exponential-width theorem has a buffer of
order $\gamma$ and only a relative all-observable channel conclusion at that
width. Finite-order source-bank information cannot be used on exponentially
large slices without a new norm estimate. The polynomial-width expansion
keeps the exponent $p$ fixed before its coupling limit; it is not permissible
to choose $p$ from the instantaneous unrestricted volume and ignore the
resulting constants and thresholds.

Likewise, for a fixed local core, sending $\ell$ to infinity at fixed
coupling does not make our bound tend to zero: the counted finite-buffer
exception eventually dominates. The next spatial task is to replace that
whole-buffer union by a \emph{connected propagation} estimate or a compatible
local dynamics with its actual update cost. The uniform native kernel and
its exponentially small exception are inputs to such an argument, not its
conclusion. Neither pairwise maximal couplings nor separate shell tests
supply an unproved grand coupling or tensorization through a shared exterior.

The original local resampling gap remains unproved. All statements concern
the prescribed bounded-bridge fast law. They do not remove the bridge window,
perform the physical Gauss reduction, prove an infrared contact floor, identify
a conditional transition with the physical transfer at its same top, or
calibrate physical time. The interacting continuum construction and a
physical Yang--Mills mass gap retain their separate obligations.

\subsection{Diagnostics, exact preservation, and analytic status}

The accompanying script is
\begin{center}\small\path{verify_ym_f16_v22_native_boundary.py}.\end{center}
It checks the normalized marginal mixed-derivative signs, covariance
localization algebra on exact Gaussian fixtures, conditional restriction and
posterior mixture identities on finite positive nonlinear laws, graph-buffer
depths and growth, and the polynomial/exponential-width and time-row power
bookkeeping. The new diagnostic passed 258 exact rational, integer, and symbolic assertions.
Its output distinguishes finite algebra from the analytic estimates; no finite
fixture is a native compact-integral enclosure. The available V4--V21
diagnostics were rerun without source changes and all eighteen passed; their
actual outputs and hashes are included in the package.

The cumulative source changes only the single V21 title version and inserts
this exact appendix before its final document-ending command. Removing that
insertion and restoring the title recovers V21 byte for byte. The package
records the initial input hashes, recovery, compilation, and rendered-page
inspection. These checks do not certify the inherited analytic archive. The
new deductions explicitly rely on the identified pinned specification,
V20 preparation, and product-ball covariance realization. In particular the
small-angular radius, comparison constants, and coupling thresholds are
proved to exist by the estimates above, not numerically enclosed or formally
verified. The next scheduled companion checkpoint remains V25.

\begin{firewall}
V22 proves a uniform complete-native buffered boundary comparison with an
exponentially small finite-buffer exception and no native--Gaussian power
floor. It gives arbitrarily strong fixed-power channel contraction for every
polynomial-size core. Exact temporal composition then proves fixed-weight
source covariance and memory asymptotics for every fixed polynomial spatial
width, and a separate all-observable temporal channel bound for a specified
exponential width. It does not establish unrestricted spatial volume at fixed
coupling, a gap in the original local sampler clock, a physical transfer gap,
or an interacting continuum mass gap. Every native normalizer and original
source remains unchanged; the small-angular restriction is paid and removed
from the target law, not silently imposed on it.
\end{firewall}
% END F16 V22 APPEND


% BEGIN F16 V23 APPEND -- 2026-09-17
% Insert immediately before the final document-ending command in V22.
% No predecessor theorem, probability law, source, or bibliography is edited.

\clearpage
\section{V23: common regeneration and exponentially wide coherent histories}
\label{f16v23:context}

V22 supplies a native comparison with an exponentially small finite-buffer
exception. Two consequences have not yet been extracted. First, a comparison
of densities, rather than only their total variation, supplies \emph{one common
subprobability measure for every compact boundary}. This constructs a genuine
simultaneous coupling, including uncountably many incoming boundary values.
It does not by itself control the dependence of the filled buffer. Second,
the same exponential error, combined with the exact two-shell identity of V19,
gives a native spatial covariance envelope with a \emph{nonperturbative} floor.
Only that floor, not the finite-order local error, pays the size of a coherent
source bank. Exponentially large banks can therefore be treated without an
expansion order growing with the volume.

These inputs close the source-asymptotic question that V22 explicitly left
open at exponential spatial width. There are fixed $\beta_r,a_r>0$ for which
all fixed orders of the complete covariance and memory hold in an exponentially
weighted temporal operator-row norm when $B\le e^{\beta_r\gamma}$, for every
duration. The width and temporal exponents are fixed independently of expansion
order; error constants and entry thresholds may depend on that order. This is
still not arbitrary spatial volume at fixed coupling. The all-boundary
regeneration is a new input for that remaining spatial problem, not a proof
that its overlapping buffers have already been composed.

\subsection{Source review and the distinction between a pair coupling and a common part}

The complete V22 appendix and analytic audit were read. The inputs used below
are its restricted mixed-normalizer identity, the componentwise core-score
gradient estimate, and its actual full-conditional buffer tail. V16 supplies
native source moments, V17 fixed-order entry coefficients, V19 the exact local
shell identity, and V21--V22 the actual temporal Markov composition. No global
native sampler gap is assumed. All conclusions are deductions from these
identified analytic inputs, not independent certification of their full
historical proofs.

Two targeted semantic searches of the f14/f15/Grand Tensor sources returned
no results. Mounted text searches and identified passages were then examined.
The monograph already gives a finite-card positive shell minorization and
warns that its constants may collapse with the carrier. F15 V22 already uses
pairwise common-part couplings on separated cores. Neither general mechanism
is claimed new. Here the common component is \emph{quantitatively near one},
uniform over the entire compact boundary family, and is constructed before any
boundary is sampled. The current calculation does not equate the monograph's
physical shell channels with the present fixed-bridge fast kernels.

For classical attribution, Spinka,
\href{https://arxiv.org/abs/1803.10578}{\emph{Finitary codings for spatial mixing
Markov random fields}}, Section 2.2, distinguishes a grand coupling from
separate optimal pair couplings. Its common-overlap formula and the associated
warning, as well as the scope of Theorems 1.2 and 4.1, were examined. Its
finite-valued, translation-invariant finitary-coding theorem is \emph{not}
imported for this inhomogeneous continuous law. The common randomization used
here is given directly by finite-dimensional measurable kernels. The
memoranda's separate ledgers for power gain and locality cost remain binding;
no Navier--Stokes theorem is an analytic premise. The next companion checkpoint
remains V25.

\subsection{A common native density floor, not an inference from total variation}

Keep V22's notation: $\mathcal A\subset W\subset\Lambda$ is a core of $k$
complete unmoving cells, $W$ is its $\ell$-neighbourhood in $\mathcal G$,
$C=\partial_{\mathcal G}W$, and
\[
 d_{\mathcal G}(\mathcal A,\Lambda^c)\ge\ell+D_*+2,\qquad \ell\ge3.
\]
The actual full conditional is $P_\Lambda^q$, its core kernel is
$K_{\Lambda,A}(q,\cdot)$, and its prescribed bridge history satisfies
$\max_{i,e}|y_{i,e}|\le r$. All compact exterior values $q$ are admitted in
V20's explicit principal-chart version. The restricted native row
$\pi_{W,A}^z$ and reference $\lambda_{W,A}=\pi_{W,A}^{0}$ use the same
whole-cell comparison balls of radius $\eta_*\sqrt\gamma$. They are not
Gaussian laws, and this event is not V4's original $\mathcal O_R$.

\begin{proposition}[Restricted density ratios with their output cost]
\label{f16v23:density-ratio}
At V22's fixed-$r$ entry threshold, there is $C_r<\infty$ such that, for
all small collar values $z,z'$ and all core points in their common comparison
domain,
\begin{equation}
 e^{-D_{\gamma,k,\ell}}\le
 \frac{d\pi_{W,A}^z}{d\pi_{W,A}^{z'}}(x)
 \le e^{D_{\gamma,k,\ell}},
 \qquad D_{\gamma,k,\ell}=C_r\gamma k e^{-a\ell}.
 \label{f16v23:ratio}
\end{equation}
Here $a>0$ may be decreased once from V22 and is independent of size and
coupling. If the collar is empty, take $D_{\gamma,k,\ell}=0$.
\end{proposition}
\begin{proof}
Use the \emph{same} normalized density $p_s$ along the collar segment and
$q_s(x)=\mathbb E_{\pi_{W\setminus A}^{x,z_s}}\partial_sV_W^{z_s}$ as in
\eqref{f16v22:score-identity}. That identity and
\eqref{f16v22:marginal-score-gradient} give
\[
 \left|\partial_{x_i}\log\frac{p_1(x)}{p_0(x)}\right|
 =\left|\int_0^1-\partial_{x_i}q_s(x)\,ds\right|
 \le C_r\sqrt\gamma e^{-a\ell}.
\]
There are at most $36k$ real core coordinates. Join any two core points by
their segment inside the product of whole-cell balls. Each scalar displacement
is at most $2\eta_*\sqrt\gamma$, so the oscillation of the log ratio is at
most $C_r\gamma k e^{-a\ell}$. Both densities are normalized and strictly
positive on the common compact comparison domain. Hence the minimum of the
log ratio is at most zero and its maximum at least zero: otherwise its
exponential could not have integral one against $p_0$. Its absolute value is
therefore at most its oscillation. This proves \eqref{f16v23:ratio}.
Both marginal normalizers, the core Haar factors, and the pinned-complement
covariance used in the gradient estimate remain those of V22. This proof
uses no functional inequality for the complete native law.
\end{proof}

Define
\begin{equation}
 u_{\gamma,k,\ell}=C_r k(1+\ell)^4e^{-c_*\gamma},\qquad
 \alpha_{\gamma,k,\ell}=(1-u_{\gamma,k,\ell})_+
                                      e^{-D_{\gamma,k,\ell}}.
 \label{f16v23:alpha}
\end{equation}
The constants can be enlarged to include the actual V22 buffer-tail bound.
Always $0\le\alpha\le1$ and
\begin{equation}
 1-\alpha\le\min\{1,u_{\gamma,k,\ell}+D_{\gamma,k,\ell}\}.
 \label{f16v23:failure}
\end{equation}

\begin{theorem}[One common submeasure for every compact exterior]
\label{f16v23:minorization}
For every Borel core set $E$ and every compact exterior $q$,
\begin{equation}
 \boxed{\qquad
 K_{\Lambda,A}(q,E)\ge\alpha_{\gamma,k,\ell}\lambda_{W,A}(E).
 \qquad}
 \label{f16v23:common-part}
\end{equation}
The reference and weight are independent of $q$. This is an inequality of
measures, not just a total-variation statement. The target remains the
complete native marginal. No good-event probability or conditional partition
function is replaced by a Gaussian denominator.
\end{theorem}
\begin{proof}
Let $G$ be V22's event that every cell in $W\cup C$ lies in its comparison
ball. The full conditional satisfies $P_\Lambda^q(G)\ge1-u$ for every $q$.
Conditional on $G$ and the actual exterior of $W$, its core law is exactly a
restricted row $\pi_{W,A}^z$. The mixing posterior retains all the original
conditional partition functions, as in V22. Every such row dominates
$e^{-D}\lambda_{W,A}$ by \eqref{f16v23:ratio}. Therefore
\[
 P_\Lambda^q(\zeta_A\in E,G)
 \ge P_\Lambda^q(G)e^{-D}\lambda_{W,A}(E)
 \ge(1-u)_+e^{-D}\lambda_{W,A}(E).
\]
Discard only the nonnegative contribution of $G^c$. This proves the measure
inequality on every Borel set. The empty-collar case has a constant restricted
row and is covered with $D=0$. The inequality is trivial but valid when $u\ge1$.
\end{proof}

The price $\gamma k$ in a log-density oscillation is larger than the
$\sqrt{\gamma k}$ in V22's TV comparison. They prove different statements.
For fixed $k\le A_0\gamma^p$ and any fixed $s>0$, choosing
\[
 \ell=3+\left\lceil a^{-1}(s+p+2)\log\gamma\right\rceil
\]
makes $1-\alpha\le\gamma^{-s}$ at a fixed threshold. There is still no
unrestricted-output assertion at fixed coupling.

\subsection{A simultaneous native regeneration map, with its buffer retained}

\begin{theorem}[All-boundary coalescence and exact full-block filling]
\label{f16v23:grand}
There is a measurable random map with independent uniform seeds such that:
for each $q$ its output on $\Lambda$ has exactly the law $P_\Lambda^q$;
and on a measurable seed event of probability $\alpha$, its core output is
\emph{the same for every compact boundary $q$ simultaneously}, with law
$\lambda_{W,A}$. The probability is uniform also when the incoming boundary
is any measurable function of previously exposed randomness, provided the
new seeds are independent of it.

Applying this map as a full $\Lambda$ conditional resampling preserves the
original native probability. It does not assert that the filled buffer is
independent of $q$ on the coalescence event, or that replacing only the core
and keeping the old buffer would preserve that probability.
\end{theorem}
\begin{proof}
For $0\le\alpha<1$ define the exact residual kernel
\[
 K^{\rm res}(q,dx)
       =\frac{K_{\Lambda,A}(q,dx)-\alpha\lambda_{W,A}(dx)}{1-\alpha}.
\]
It is a positive probability kernel by \eqref{f16v23:common-part}. It is
measurable in $q$: the full and restricted finite-dimensional integrals have
measurable densities. Let $H(q,x,dz)$ be the actual conditional law of the
remaining free cells given the core $x$ and the exterior $q$. Choose arbitrary
measurable values at null core parameters. The exact disintegration is
\[
 P_\Lambda^q(dx,dz)
 =\{\alpha\lambda_{W,A}(dx)+(1-\alpha)K^{\rm res}(q,dx)\}H(q,x,dz).
\]
Every kernel on these finite-dimensional standard Borel spaces can be sampled
by a measurable uniform-seed map. One direct construction orders its real
coordinates, uses measurable conditional distribution functions, and takes
their generalized inverses; values on null conditional sets can be fixed
arbitrarily. Apply this construction to $\lambda$, $K^{\rm res}$, and $H$.

With a fresh Bernoulli decision of success probability $\alpha$, draw the
core from $\lambda$ on success and from $K^{\rm res}(q)$ on failure, then
fill the rest using $H(q,x)$. Use the \emph{same} success decision and the
same $\lambda$ seed for all $q$. On success the equality of core outputs is
pathwise for all $q$, so no uncountable intersection of separately almost-sure
statements is needed. The final buffer seed is allowed to depend on $q$ through
its sampling map. Integrating the seeds gives exactly the displayed
conditional disintegration for every $q$. If $\alpha=1$, omit the residual
branch; if $\alpha=0$, the claimed common event is empty.

Conditional resampling of the entire $\Lambda$ from its actual law preserves
$P$ by the tower property. A deterministic sequence of such updates, or an
update schedule chosen independently of the current field, therefore also
preserves it. A field-dependent schedule would need a separate invariance
argument. Fresh decisions in a sequence have conditional failure probabilities
$1-\alpha_t$ given the past; for deterministic weights, any prescribed family
of failures has probability $\prod_t(1-\alpha_t)$. This is a statement about
algorithmic labels, not independence of native error fields. The buffer
conditional $H$ can still transmit incoming information even on success.
\end{proof}

For separated blocks the maps can be used independently after fixing their
one common exterior. Their success labels are then independent of that
exterior, unlike a posterior estimate that merely says an event is usually
good. For overlapping updates, use fresh seeds in their actual order; do not
replace the resulting joint law by independent simultaneous buffer outputs.
These are genuine common random maps, but not yet a subcritical exploration
of all their dependencies.

There is a simple reason not to infer this theorem from V22's TV bound alone.
A positive regularization of the standard grand-coupling example in Spinka's
Section 2.2 has rows on $\{1,\ldots,N\}$
\[
 k_i(j)=\frac{\theta}{N}
              +\frac{1-\theta}{N-1}1_{\{j\ne i\}},\qquad0<\theta<1.
\]
Each pair has TV distance $(1-\theta)/(N-1)$, tending to zero with $N$.
Nevertheless the mass of \emph{any} submeasure common to all rows is at most
$\sum_j\min_i k_i(j)=\theta$. The density floor above, rather than pairwise
closeness alone, proves the near-unit simultaneous coalescence. This is a
classical coupling warning, not a counterexample to the native specification.

\subsection{A nonperturbative spatial envelope for the actual primitive sources}

Return to V19's ambient cell distance
\[
 d_*((t,b),(s,c))=|t-s|+d_{\mathbb T_m^3}(b,c),\qquad
 \varepsilon=\gamma^{-1/2},\quad\mathcal N=B(n+1).
\]
Every edge of $\mathcal G$ has $d_*$ length at most a fixed $r_0$. Every
primitive error has a carrier of fixed $d_*$ radius $r_s$ and a uniformly
bounded number of cells. Let $F_u=e_u^{\rm prim}-\mathbb E_Pe_u^{\rm prim}$,
$K_\gamma=\operatorname{Cov}_P(e^{\rm prim})$, and
\[
 m_{u,R}=\mathbb E_P[F_u\mid\zeta_{I_{u,R}^c}],
 \qquad I_{u,R}=\{c:d_*(a(u),c)\le R\}.
\]
These are precisely the centered shell messages of V19, with their actual
conditional normalizers. V16 gives $\operatorname{Var}_P F_u\le C_r\varepsilon^2$.

\begin{proposition}[Source screening with an exponentially small coupling floor]
\label{f16v23:nonpert-screen}
There are fixed $R_0,b,c,C_r>0$ such that, for all sufficiently large
fixed-$r$ coupling and every $R\ge R_0$,
\begin{equation}
 \sup_u\|m_{u,R}\|_2^2
 \le C_r\varepsilon^2
       \min\{1,\gamma e^{-bR}+(1+R)^4e^{-c\gamma}\}.
 \label{f16v23:screen}
\end{equation}
A ball covering the whole system has centered message zero. There are also
fixed $b_1,\kappa>0$ such that, for every pair of primitive marks,
\begin{equation}
 |(K_\gamma)_{uv}|
 \le C_r\varepsilon^2
       \left[\gamma e^{-b_1d_*(a(u),a(v))}+e^{-\kappa\gamma}\right].
 \label{f16v23:primitive-envelope}
\end{equation}
The same bound, after one fixed decrease of $b_1$, holds for each entry of
$\Gamma_\gamma=TK_\gamma T^{\mathsf T}$ in its original common-source anchors.
Constants do not depend on volume, duration, or the marks.
\end{proposition}
\begin{proof}
Use the complete primitive carrier as the observed core; its cell count is
uniformly bounded. Its distance in $\mathcal G$ to $I_{u,R}^c$ is at least
$(R-r_s)/r_0$. Choose
\[
 \ell=\left\lfloor(R-r_s)/r_0\right\rfloor-D_*-3.
\]
Choose $R_0$ also larger than every primitive carrier allowance. For
$R\ge R_0$ this is at least three and satisfies the exact buffer geometry
of \eqref{f16v23:common-part}. Its common component has failure probability
at most $C_r[\gamma e^{-bR}+(1+R)^4e^{-c\gamma}]$, by
\eqref{f16v23:failure}. A common component of mass $\alpha$ bounds the
kernel diameter by $1-\alpha$. V20's proved all-incoming-law channel lemma,
applied with the \emph{actual} exterior marginal and the fixed observable
$F_u$, bounds the squared conditional norm by $(1-\alpha)\operatorname{Var}F_u$.
This proves \eqref{f16v23:screen}. No native fourth moment at an adversarial
boundary or inverse of a native generator is used.

For two anchors at distance $d>2R+\rho$, V19's exact two-shell identity is
\[
 \mathbb E_P F_uF_v=\mathbb E_P m_{u,R}m_{v,R}.
\]
Here $\rho$ is one fixed ambient interaction/carrier allowance. The first
message is a function of its finite exterior word collar, which lies outside
the second ball; two conditional expectations justify the identity. It does
not assert that the collars are independent. Choose, whenever $d$ is large
enough,
\[
 R=\min\left\{\left\lfloor\frac{d-\rho-1}{3}\right\rfloor,
                                      \lceil\gamma\rceil+R_0\right\}.
\]
This obeys $d>2R+\rho$. Before the radius reaches its cap, the first term of
\eqref{f16v23:screen} is bounded by $C_r\gamma e^{-b_1d}$ for a fixed
$b_1>0$. The second is at most $C_r(1+\gamma)^4e^{-c\gamma}$, which is
bounded by a fixed exponential $C_re^{-c\gamma/2}$. After saturation, both
terms are at most $C_re^{-\kappa\gamma}$ for some fixed $\kappa>0$.
Cauchy--Schwarz in the exact two-shell identity proves
\eqref{f16v23:primitive-envelope}. The finitely many shorter distances are
covered by the original variance bound, enlarging $C_r$. Empty exteriors
cause no exception.

The original deterministic synthesis $T=(I,\mathsf P^{\mathsf T})$ has
uniform exponentially weighted absolute rows in these anchor labels, as
used in V18--V19. Decrease $b_1$ below its weight. The triangle inequality
for $d_*$ bounds
\[
 \sum_{u,v}|T_{\alpha u}T_{\beta v}|
                       e^{-b_1d_*(a(u),a(v))}
                    \le C e^{-b_1d_*(a(\alpha),a(\beta))}.
\]
The constant floor is preserved by the two absolute row sums. This proves the
full-source assertion. The spatial lift is not falsely assigned finite support.
\end{proof}

In particular, at the \emph{single} linear radius
$R_\gamma=\lceil\gamma\rceil+R_0$, decrease $\kappa$ once so that
\begin{equation}
 \sup_u\|m_{u,R_\gamma}\|_2^2\le C_r\varepsilon^2e^{-\kappa\gamma}.
 \label{f16v23:linear-screen}
\end{equation}
Unlike an arbitrarily high fixed-order floor, this estimate has specified
exponential dependence on coupling. It still does not tend to zero as
$R\to\infty$ at fixed coupling by the displayed bound alone.

\subsection{Only the nonperturbative floor pays coherent source size}

The next estimate uses the same normalized coefficients $G_q$ and
$M_q=-\operatorname{Off}_{>1}G_q$ as V10. All fixed coefficient rows decay
exponentially in $d_*$, and V17's entry expansion is uniform in ambient volume.
Let $\|v\|_1$ be the counting coefficient norm on the original canonical
common-source coordinates, without volume normalization.

\begin{theorem}[An operator error plus an exponentially small coherent remainder]
\label{f16v23:mixed}
For every fixed $J\ge2$ there are $C_{J,r},\gamma_{J,r}<\infty$ such that,
for every admitted finite history with $\gamma\ge\gamma_{J,r}$ and all real
or complex source vectors $v,w$,
\begin{equation}
 \boxed{
 \left|v^*\left(\Gamma_\gamma-\sum_{q=2}^J\gamma^{-q/2}G_q\right)w\right|
 \le C_{J,r}\gamma^{-(J+1)/2}\|v\|_2\|w\|_2
       +C_r\gamma^{-1}e^{-\kappa\gamma}\|v\|_1\|w\|_1.}
 \label{f16v23:mixed-bound}
\end{equation}
The same estimate holds for $\mathsf M_\gamma$ and $M_q$, with enlarged fixed
constants. The exponent $\kappa>0$ is independent of $J$. Consequently, for
every fixed $0<\beta<\kappa$ and any source bank $\mathcal A$ with
$|\mathcal A|\le e^{\beta\gamma}$,
\begin{equation}
 \left\|\chi_{\mathcal A}
   \left(\Gamma_\gamma-\sum_{q=2}^J\gamma^{-q/2}G_q\right)
       \chi_{\mathcal A}\right\|_{\rm op}
                  \le C_{J,\beta,r}\gamma^{-(J+1)/2},
 \label{f16v23:exponential-bank}
\end{equation}
and likewise for memory. The ambient system may be arbitrarily larger and
all source positions are allowed. The fixed order does not grow with the bank.
\end{theorem}
\begin{proof}
Choose a fixed $b_2>0$ below the envelope and coefficient decay exponents, so
polynomial graph growth gives an exponential absolute row tail at exponent
$b_2$ for either exponential envelope. Put
\[
 D_J=\left\lceil (J+3)b_2^{-1}\log(\varepsilon^{-1})\right\rceil+R_0,
 \qquad K=J+1.
\]
On pairs with $d_*\le D_J$, apply V17 at order $K$. The number of entries
per row and column is at most $C(1+D_J)^4$. The entry remainder therefore
has absolute row and column sums at most
$C_{J,r}(1+D_J)^4\varepsilon^{J+2}\le C_{J,r}\varepsilon^{J+1}$.
The extra coefficient $\varepsilon^{J+1}G_{J+1}$ has that same global row
bound without a distance count.

On pairs with $d_*>D_J$, bound the native covariance by
\eqref{f16v23:primitive-envelope} in its full-source form and the finitely
many desired coefficients by their inherited exponential rows. Their
exponential parts have absolute row and column sums at most
\[
 C_{J,r}(\varepsilon^2\gamma+\varepsilon^2)e^{-b_2D_J}
                         \le C_{J,r}\varepsilon^{J+1}.
\]
The remaining entry bound is precisely
$C_r\varepsilon^2e^{-\kappa\gamma}$. The nonnegative envelopes with small
row and column sums bound a bilinear form by their Schur operator norm times
$\|v\|_2\|w\|_2$. The constant entry floor costs instead
$C_r\varepsilon^2e^{-\kappa\gamma}\|v\|_1\|w\|_1$. This proves
\eqref{f16v23:mixed-bound} without an unrestricted sum of that floor.
All estimates used absolute coefficients, so complex vectors cause no change.

For memory the exact identity masks entries of this same error. It preserves
both nonnegative entry envelopes and their absolute row bounds; hence the
identical argument applies. It is not assumed to preserve positivity of a
matrix. On a bank of size $k$, $\|v\|_1\le\sqrt{k}\|v\|_2$. For
$k\le e^{\beta\gamma}$ the extra term is at most
$C_r\gamma^{-1}e^{-(\kappa-\beta)\gamma}\|v\|_2\|w\|_2$, smaller than
any specified fixed algebraic power at a fixed-$J,\beta,r$ threshold. This
proves \eqref{f16v23:exponential-bank}.
\end{proof}

At leading order, arbitrary banks of that exponential size obey
\[
 \|\chi_{\mathcal A}(\Gamma_\gamma-\gamma^{-1}G_2)\chi_{\mathcal A}\|
 \le C_{\beta,r}\gamma^{-3/2}.
\]
The same holds for memory. Their actual covariances are
$O_{\beta,r}(\gamma^{-1})$ in operator norm. This is not derived by
multiplying an algebraic entry error by an exponentially large bank.

A sharper necessary condition on a remaining coherent obstruction follows.
For $v\ne0$ put $k_{\rm eff}(v)=\|v\|_1^2/\|v\|_2^2$. For unit vectors,
\begin{equation}
 \left|\gamma v^*\Gamma_\gamma v-v^*G_2v\right|
 \le C_r\gamma^{-1/2}+C_re^{-\kappa\gamma}k_{\rm eff}(v).
 \label{f16v23:effective-support}
\end{equation}
An order-one failure along $\gamma\to\infty$ must therefore have
$k_{\rm eff}(v_\gamma)\ge c\,e^{\kappa\gamma}$ eventually. This improves
V17's superpolynomial necessary delocalization, but does not exclude such
vectors in unrestricted volume. The same conclusion applies to memory.

\subsection{The exponential-width temporal channel now has a coherent input}

V22's exponential-width channel theorem supplies fixed constants
$\beta_{\rm ch},d_{\rm ch},C_{\rm ch}>0$ and actual buffers $b_\gamma^{\rm exp}$
with
\begin{equation}
 B\le e^{\beta_{\rm ch}\gamma},\qquad
 c_{\rm ch}\gamma\le b_\gamma^{\rm exp}\le C_{\rm ch}\gamma,
 \qquad \delta(K_{i,i+b_\gamma^{\rm exp}})\le e^{-d_{\rm ch}\gamma}.
 \label{f16v23:channel-input}
\end{equation}
The positive lower constant $c_{\rm ch}$ follows from V22's explicit linear
choice of its comparison radius. Every continuation message and both endpoints
remain in these native transitions. Its past/future correlation bound yields
a fixed $\lambda_{\rm ch}>0$ with relative correlation at most
$e^{-\lambda_{\rm ch}g}$ whenever the gap $g\ge2b_\gamma^{\rm exp}$.
These are inherited serial-chain facts, not a new composition across spatial
pieces.

Fix once
\begin{equation}
 0<\beta_r\le\min\{\beta_{\rm ch},\kappa/8\},\qquad
                         B\le e^{\beta_r\gamma}.
 \label{f16v23:width}
\end{equation}
At sufficiently large coupling a full single-time source bank, of size $9B$,
is within the exponential-bank theorem with some exponent strictly below
$\kappa$. Thus its coherent variance is at most
$C_r\varepsilon^2\|v\|_2^2$. The full Haar source has its original finite
temporal carrier allowance $\sigma$, although its spatial lift is nonlocal.
Combining this coherent input with \eqref{f16v23:channel-input} gives
\begin{equation}
 \|\Gamma_{\gamma,ij}\|_{\rm op}
 \le C_r\varepsilon^2e^{-\lambda_{\rm ch}(|i-j|-2\sigma)}
 \quad\text{when }|i-j|-2\sigma\ge2b_\gamma^{\rm exp}.
 \label{f16v23:far-time}
\end{equation}
At every lag, spatial row summation of the new native envelope also gives
\begin{equation}
 \|\Gamma_{\gamma,ij}\|_{\rm op}
 \le C_r\varepsilon^2
              \bigl[\gamma e^{-b_3|i-j|}+B e^{-\kappa\gamma}\bigr]
 \label{f16v23:middle-time}
\end{equation}
for a fixed $b_3>0$. The spatial exponential row sum is bounded independently
of torus size; the constant floor pays exactly $B$. These two estimates have
different roles. Neither replaces the coherent short-time variance by a sum
of $B$ independent-source variances.

\begin{theorem}[All fixed orders on an exponentially wide complete history]
\label{f16v23:exponential-history}
There are fixed $a_r>0$ and $\beta_r>0$, independent of expansion order, such
that for every fixed $J\ge2$ there are finite $C_{J,r},\gamma_{J,r}$ such
that, for all $\gamma\ge\gamma_{J,r}$ in \eqref{f16v23:width} and every duration,
\begin{equation}
 \boxed{
 \begin{aligned}
 &\left\|\Gamma_\gamma-\sum_{q=2}^J\gamma^{-q/2}G_q
                                          \right\|_{{\rm tr},a_r}\\
 &\quad+\left\|\mathsf M_\gamma-\sum_{q=2}^J\gamma^{-q/2}M_q
                                          \right\|_{{\rm tr},a_r}
                  \le C_{J,r}\gamma^{-(J+1)/2}.
 \end{aligned}}
 \label{f16v23:row-asymptotics}
\end{equation}
Here $\|A\|_{{\rm tr},a_r}=\sup_i\sum_j e^{a_r|i-j|}\|A_{ij}\|_{\rm op}$
is the \emph{temporal operator-row norm}, not trace norm. Each block includes
the entire spatial common-source space. There is no normalization by
$B(n+1)$ or restriction on the total time-source count. In particular the
full covariance and memory are $O_r(\gamma^{-1})$ in this norm and their
leading-coefficient errors are $O_r(\gamma^{-3/2})$.
\end{theorem}
\begin{proof}
Let $a_*>0$ be a fixed temporal coefficient exponent, so that each fixed
$G_q$ has finite row norm at exponent $a_*$. This follows from V10's full
space--time absolute rows, as already explained in V21. Put
$D_1=2b_\gamma^{\rm exp}+2\sigma$ and choose a fixed $C_1$ with
$D_1\le C_1\gamma$. Choose once
\begin{equation}
 0<a_r\le\min\{\lambda_{\rm ch}/4,b_3/4,a_*/4,\kappa/(8C_1)\}.
 \label{f16v23:weight}
\end{equation}
For the given $J$, set
\[
 D_0=2\sigma+\left\lceil(J+5)a_r^{-1}\log(\varepsilon^{-1})\right\rceil,
 \qquad K=2J+8.
\]
At a fixed-$J,r$ threshold $D_0<D_1$. Split the temporal sum into three bands.

For $|i-j|\le D_0$, apply \eqref{f16v23:mixed-bound} at order $K$ to unit
vectors supported in time blocks $i,j$. Its error block is bounded by
\[
 C_{K,r}\varepsilon^{K+1}
                      +C_r B\varepsilon^2e^{-\kappa\gamma}.
\]
There is \emph{no} factor $B$ on the algebraic term. Summing that term with
its temporal weight costs
\[
 C_{J,r}(1+D_0)e^{a_rD_0}\varepsilon^{K+1}
 \le C_{J,r}(1+\log\varepsilon^{-1})\varepsilon^{K-J-4}
 \le C_{J,r}\varepsilon^{J+1}.
\]
The coefficients of orders $J+1,\ldots,K$ cost
$C_{J,r}\varepsilon^{J+1}$ in their \emph{entire} weighted row norm.
They are not charged separately for every near lag.

On $D_0<|i-j|\le D_1$, use \eqref{f16v23:middle-time} for the native matrix
and the coefficient row tails. The exponential part has weighted sum at most
\[
 C_{J,r}\varepsilon^2\gamma e^{-(b_3-a_r)D_0}
               +C_{J,r}\varepsilon^2e^{-(a_*-a_r)D_0}
                         \le C_{J,r}\varepsilon^{J+1}.
\]
The constant floor on \emph{both} bands up to $D_1$ has total cost at most
\begin{equation}
 C_r\varepsilon^2(1+D_1)e^{a_rD_1}B e^{-\kappa\gamma}
 \le C_r\varepsilon^2(1+C_1\gamma)e^{-3\kappa\gamma/4}.
 \label{f16v23:middle-floor}
\end{equation}
Here $\beta_r\le\kappa/8$ and $a_rC_1\le\kappa/8$ were used. This is
smaller than every prescribed fixed power of $\varepsilon$; its entry
threshold may depend on $J$.

For $|i-j|>D_1$, the genuine temporal estimate \eqref{f16v23:far-time}
has weighted tail at most
\[
 C_r\varepsilon^2e^{-(\lambda_{\rm ch}-a_r)(D_1-2\sigma)}.
\]
Since $b_\gamma^{\rm exp}\ge c_{\rm ch}\gamma$, this is exponentially small
in $\gamma$. The coefficient tail is likewise exponentially small. If the
actual duration is shorter, the missing bands are simply empty. The three
bounds prove the covariance assertion. The unchanged contact identity
$\mathsf M_\gamma=-\operatorname{Off}_{>1}\Gamma_\gamma$ removes entire
time blocks and contracts this absolute row norm, giving memory. The exponent
\eqref{f16v23:weight} was fixed before $J$ was chosen. All source and
measure-correction coefficients remain those of the original native law.
\end{proof}

\subsection{One linear shell has an exponentially small coherent return}

The following identifies the old positive object more strongly than an
order-dependent logarithmic radius. It uses exactly
\[
 \mathsf R^{\rm sh}_{\gamma,R}
 =T\operatorname{Cov}_P((m_{u,R})_u)T^{\mathsf T}
\]
from V19, not V12's different one-common-exterior cylinder return.

\begin{proposition}[Exponential return bound at one order-independent radius]
\label{f16v23:linear-return}
After a possible fixed decrease of $a_r,\beta_r$, there is $\kappa_r'>0$ such
that, with $R_\gamma=\lceil\gamma\rceil+R_0$,
\begin{equation}
 \|\mathsf R^{\rm sh}_{\gamma,R_\gamma}\|_{{\rm tr},a_r}
                      \le C_r\gamma^{-1}e^{-\kappa_r'\gamma}
 \label{f16v23:linear-return-bound}
\end{equation}
throughout \eqref{f16v23:width}, uniformly in duration. At unrestricted
ambient volume, its diagonal and trace density still obey
$C_r\gamma^{-1}e^{-\kappa\gamma}$, but that alone is not an operator bound.
\end{proposition}
\begin{proof}
By \eqref{f16v23:linear-screen}, every primitive message has variance at most
$C_r\varepsilon^2e^{-\kappa\gamma}$. Its actual collar is contained in the
space--time ball of radius $R_\gamma+\rho$. A time bank of such messages has
at most $30B$ entries, so its covariance operator norm is at most
$C_rB\varepsilon^2e^{-\kappa\gamma}$ by its trace. This use of a trace is
explicit and is now paid by an exponential coupling margin.

Two such time banks have separated temporal carriers once their time gap
exceeds $2(R_\gamma+\rho)$. Apply the actual exponential-width past/future
channel after a further gap $2b_\gamma^{\rm exp}$. Covariance
Cauchy--Schwarz handles the shorter lags. Write
$G_\gamma=2(R_\gamma+\rho)+2b_\gamma^{\rm exp}\le C_G\gamma$.
The resulting weighted row sum is at most
\[
 C_rB\varepsilon^2e^{-\kappa\gamma}
                         (1+G_\gamma)e^{a_rG_\gamma},
\]
provided $a_r\le\lambda_{\rm ch}/4$. Decrease $a_r$ so that
$a_rC_G\le\kappa/8$ and retain $\beta_r\le\kappa/8$.
The last expression is at most $C_r\varepsilon^2e^{-\kappa\gamma/2}$
after a fixed threshold increase. The original synthesis $T$ has a fixed
finite temporal range and bounded spatial block norms, hence a uniformly
bounded temporal row norm at this weight. Row-norm multiplication transfers
the bound through $T$ and $T^{\mathsf T}$. This proves the proposition.
At unrestricted volume, sum the scalar diagonal bounds and use bounded
absolute synthesis rows, or its operator norm and the trace ideal inequality,
to get the stated diagonal and trace-density bounds.
\end{proof}

\subsection{What remains spatial, and why common regeneration is not yet closure}

The new coherent remainder is
\[
 C_{J,r}\gamma^{-(J+1)/2}\|v\|_2\|w\|_2
       +C_r\gamma^{-1}e^{-\kappa\gamma}\|v\|_1\|w\|_1,
\]
not an algebraic error multiplied by the number of coordinates. It permits
exponentially large source banks and, using V22's actual temporal channels,
all-duration operator-row asymptotics at one fixed exponential spatial width.
The constants $\beta_r,a_r,\kappa$ are proved positive from the inherited
analytic bounds but have not been numerically enclosed or optimized. They
cannot be identified with a physical running-coupling coefficient.

At fixed $\gamma$, an unrestricted source vector can still have arbitrarily
large $\ell^1/\ell^2$ ratio, and $B$ can exceed $e^{\beta_r\gamma}$. The
constant floor cannot be summed over that class. The new grand coupling also
does not solve this by declaration. On a successful core refresh, the exact
buffer law $H(q,x)$ can retain dependence on the incoming boundary. Subsequent
overlapping updates can query precisely that buffer. Independent success
labels do not make those filled fields independent, and changing only a core
while retaining its old buffer is not the original conditional resampling.
A proof of subcritical \emph{dependency propagation}, or a paid comparison of
compatible block dynamics with the local clock, remains necessary.

Thus the simultaneous all-boundary map supplies an input explicitly missing
from V22, while the nonperturbative source estimate closes its separate
exponential-width covariance question. Neither is presented as unrestricted
spatial mixing. The original local sampler gap, V4's unchanged full/core
boundary derivatives, and V12's distinct spatial-cylinder posterior are not
silently redefined. Bridge-window removal, physical root Gauss restriction,
infrared contact coercivity, same-top physical transfer comparison, independent
physical-time calibration, and interacting continuum reconstruction retain
their separate obligations.

\subsection{Diagnostics, source preservation, and claim scope}

The accompanying diagnostic is
\begin{center}\small\path{verify_ym_f16_v23_regeneration_coherence.py}.\end{center}
It checks normalized density-ratio algebra, common submeasures and residual
kernels, simultaneous finite randomizations with actual conditional buffer
filling, exact preservation of a positive finite Gibbs law, and a failure of
core-only replacement. Further checks cover the two-shell covariance identity,
spatial-to-coherent envelopes, exponential support and three-band exponent
bookkeeping, and the distinction between pairwise and common overlap. The
diagnostic passes 291 exact assertions; its report records the families.
These are finite rational, integer, and symbolic diagnostics, not native compact-integral enclosures or
formal verification of the analytic estimates.

The cumulative V22 is preserved byte for byte except for one literal title
version and this exact terminal insertion. The package records predecessor
recovery, protected input hashes, build and rendering results, and the actual
reruns of available earlier diagnostics. No inaccessible historical checker
or external formalization is represented as run. The new deductions depend
on the identified V16--V22 analytic inputs, especially the preparation and
restricted product-ball covariance theorem. The entire archive has not been
independently recertified. The next scheduled companion checkpoint remains V25.

\begin{firewall}
V23 proves a uniform common native submeasure and a simultaneous all-boundary
core regeneration map, with the full buffer filled from its actual conditional.
It proves a native spatial covariance envelope with an exponentially small
coupling floor, mixed coherent-error estimates and fixed-order expansions on
exponentially large source banks. Exact temporal composition then gives
all-duration covariance and memory asymptotics in a fixed weighted temporal
operator-row norm for a specified exponential spatial width, and an
exponentially small positive return at one linear radius. It does not prove
unrestricted spatial volume at fixed coupling, subcritical composition of
all overlapping buffers, the original local sampler gap, a physical transfer
gap, or an interacting continuum Yang--Mills mass gap. All native normalizers,
compact fields and original sources remain unchanged.
\end{firewall}
% END F16 V23 APPEND

% BEGIN F16 V24 APPEND -- 2026-09-17
% Insert immediately before the final document-ending command in V23.
% No predecessor assertion, source, probability, or bibliography is edited.

\clearpage
\section{V24: a paid local-clock comparison and boundary-sensitive block contraction}
\label{f16v24:context}

V23 constructs a simultaneous core refresh but correctly retains dependence
of its filled buffer on the incoming boundary. Increasing its common-core
probability does not, by itself, control that dependence. We address the
\emph{update operator} before making another covariance expansion. For complete
time intervals the required boundary-sensitive coupling can be proved: its
expected number of affected slices is bounded by the channel buffer length,
not by the resampled interval length. Overlapping exact interval resamplings
then have a uniform spectral gap in a stated per-slice clock.

There is a second, separate calculation. Returning this operator to V14's
original rate-one balanced chord/port updates enlarges its port support by one
temporal bond. We keep that enlargement and prove the form comparison using
the actual complete conditional density. The resulting original-clock gap is
positive uniformly in duration in the admitted width regimes, but has an
explicit, pessimistic coupling/width cost. It is \emph{not} a positive gap
uniform along the weak-coupling limit. The last part identifies the local
boundary-response quantity needed to run the same construction spatially,
and shows why the coarse common-core certificate alone does not verify it.

\subsection{Dependencies, non-duplication, and which gap is under discussion}

The supplied V23 appendix and audit were read in full. Two targeted semantic
searches of the current lineage and supplied f14/f15/Grand Tensor files were
empty; mounted text searches and identified passages were then checked.
V15 already compares its gated three-port updates to the original clock; that
form-comparison principle is credited. Its gates neither cover the present
intervals nor prove their gap. The monograph's random-scan and blocking
countertheorems distinguish a per-cell clock from a single global update and
warn that taking a power of an unestimated transfer creates no gap. F14 V2's
same-clock local repair likewise concerns a different conditional and support.
These statements remain binding. This is a targeted audit, not independent
recertification of the archives; the next scheduled checkpoint remains V25.

Path coupling is a classical method; see Bubley--Dyer,
\href{https://doi.org/10.1109/SFCS.1997.646111}{\emph{Path coupling: A technique
for proving rapid mixing in Markov chains}}. Its primary abstract and metadata
were checked; no inaccessible full-paper theorem is invoked. Spinka,
\href{https://arxiv.org/abs/1803.10578}{\emph{Finitary codings for spatial mixing
Markov random fields}}, Section 4 and Theorems 4.1--4.3, makes a stronger
all-boundary-subset sensitivity requirement explicit for grand couplings.
The relevant definitions and contraction proof were reviewed. Its finite-valued,
translation-invariant coding theorem is not applied to this inhomogeneous
continuous law. We prove below the pair-coupling/spectral statement actually
used, directly on finite products of standard Borel spaces.

The analytic native input is V22's buffered \emph{temporal channel} under the
actual continuation law, including an arbitrarily fixed opposite endpoint.
V21 supplies that Markov factorization. The new spectral proof does not assume
a gap of any native generator, use V23's covariance expansion to infer one,
or substitute the Gaussian gap. The full-domain comparison to the local
clock uses only the exact coordinate transformation, product reference
factors and bounded local classical action. No small-field gate is imposed.

\subsection{A boundary-input ledger for exact conditional updates}

Let $\mu$ be a probability with strictly positive density relative to a product
of $N$ probability measures on standard Borel coordinate spaces. Null parameter
values can be assigned the explicit versions of the kernels when available.
For a coordinate block $D$, put
\[
 E_Df=\mathbb E_\mu[f\mid\omega_{D^c}],\qquad Q_D=I-E_D,
 \qquad \mathcal E_D(f)=\langle f,Q_Df\rangle_\mu.
\]
The positive generator for a finite indexed family of blocks with deterministic
rates $r_D>0$ is $\mathcal L_{\rm bl}=\sum_D r_DQ_D$. Repeated blocks may carry
separate indices and rates. These are \emph{full} conditional resamplings,
not replacement of a core while retaining its previous buffer.

Choose fixed positive coordinate weights $w_j$ and define the bounded measurable
weighted Hamming metric
\[
 d_w(x,y)=\sum_j w_j1_{\{x_j\ne y_j\}}.
\]
Use fixed measurable conditional-kernel versions on the full product carrier,
and require the following comparison for every such pair of parameters,
including the intermediate configurations in a sequence of coordinate changes.
For each $j\notin D$ suppose that when two exteriors differ only at $j$, a
coupling of the two \emph{entire} conditional outputs on $D$ has expected
inside distance at most $I_{D,j}$. Set $I_{D,j}=0$ when that exterior coordinate
cannot affect this kernel. This is a bound on a specified boundary change,
not the diameter of an arbitrary collection of core marginals. Define
\begin{equation}
 g=\inf_j\left\{\sum_{D\ni j}r_D
              -\frac1{w_j}\sum_{D\not\ni j}r_DI_{D,j}\right\}.
 \label{f16v24:influence-margin}
\end{equation}

\begin{theorem}[Boundary-sensitive block contraction in its stated clock]
\label{f16v24:general-gap}
If $g>0$, then on the complete $L^2(\mu)$ space
\begin{equation}
 \mathcal L_{\rm bl}\succeq g(I-\mathbb E_\mu),\qquad
 \operatorname{Var}_\mu f\le g^{-1}\sum_Dr_D\mathcal E_D(f).
 \label{f16v24:general-gap-bound}
\end{equation}
This holds on finite continuous product spaces as well as finite alphabets.
No global measurable grand coupling or prior spectral-gap premise is needed.
\end{theorem}
\begin{proof}
Let $R=\sum_Dr_D$ and $K=R^{-1}\sum_Dr_DE_D$, a self-adjoint Markov operator
preserving $\mu$. Consider configurations differing only at $j$. If $j\in D$,
the two exterior rows are identical, and the entire output can be coupled
identically: the distance drops from $w_j$ to zero. If $j\notin D$, the old
$w_j$ remains, and the expected additional distance is at most $I_{D,j}$.
Average the specified couplings over the chosen block. By
\eqref{f16v24:influence-margin}, the expected resulting distance is at most
$(1-g/R)w_j$. Here $0<g\le R$.

For a bounded measurable $f$, write
\[
 \operatorname{Lip}_w(f)=
 \max_j\sup_{x_{j^c}=y_{j^c},\ x_j\ne y_j}
                    \frac{|f(x)-f(y)|}{w_j}.
\]
It is bounded by $2\|f\|_\infty/\min_jw_j$. Changing one coordinate at a
time bounds $|f(x)-f(y)|$ by $\operatorname{Lip}_w(f)d_w(x,y)$. The pair
couplings therefore give
\[
 \operatorname{Lip}_w(Kf)\le(1-g/R)\operatorname{Lip}_w(f).
\]
Iteration and Poissonization, using $\mathcal L_{\rm bl}=R(I-K)$, give
\begin{equation}
 \operatorname{osc}(e^{-t\mathcal L_{\rm bl}}f)
       \le \left(\sum_jw_j\right)e^{-gt}\operatorname{Lip}_w(f).
 \label{f16v24:osc-decay}
\end{equation}
This argument needs only the existence of the pair couplings: no measurable
selection over all pairs is used to define the Markov operator.

For centered bounded $f$, invariance of $\mu$ bounds its semigroup $L^2$ norm
by the right side of \eqref{f16v24:osc-decay}. If the nonnegative self-adjoint
operator $\mathcal L_{\rm bl}$ had nonzero centered spectral subspace in
$[0,a]$, where $a<g$, density of bounded centered functions would give a test
$f$ with nonzero projection there. The spectral theorem would give a lower
bound $e^{-at}$ times its projected norm, contradicting
\eqref{f16v24:osc-decay} as $t\to\infty$. The finite constant $\sum_jw_j$ is
held fixed before this time limit; no volume-independent total-variation
prefactor is asserted. Constants are fixed, so this proves the operator
inequality and then the variance inequality for every $L^2$ function.
\end{proof}

The death term in \eqref{f16v24:influence-margin} records how often an already
different coordinate is actually refreshed. The birth term counts the
\emph{whole output}, including its buffer, for one changed incoming coordinate.
A high probability of refreshing only its central part does not supply that
birth term at a useful scale.

\subsection{A complete temporal interval has bounded boundary susceptibility}

Return to the unchanged native fixed-bridge probability $P=P_{\gamma,\mathbf y}$,
with $\max_{i,e}|y_{i,e}|\le r$. Its complete unmoving slice is $Z_i$,
$0\le i\le n$, as in \eqref{f16v21:chain-density}. A slice includes all $B$
cubes, and its size is \emph{not} counted as independent spins in the next
Hamming metric. Write $b\ge1$ for an integer channel buffer. We use the
following already proved input:
\begin{equation}
 \delta(K^{v}_{t,t+b})\le q\le\tfrac14
 \quad\hbox{whenever the output is at least $b$ from a fixed right boundary.}
 \label{f16v24:channel-input}
\end{equation}
The analogous backward statement holds with the left boundary fixed. Natural
free endpoints require no artificial boundary allowance. The superscript $v$
means the actual fixed right factor is retained inside every continuation
kernel; it is not a conditioning imposed after an estimate.

V22 proves \eqref{f16v24:channel-input} in two usable regimes. For every fixed
$p\ge0,A_0<\infty$, $B\le A_0\gamma^p$ permits $b=b_{\gamma,p,2}=O_{p,r}(\log\gamma)$
and $q\le\gamma^{-2}$ at its stated threshold. Separately,
$B\le\exp(\beta_{\rm ch}\gamma)$ permits $b=b_\gamma^{\rm exp}=O_r(\gamma)$ and
$q\le\exp(-d_{\rm ch}\gamma)$. Increase only those fixed thresholds so that
$q\le1/4$. No source-bank expansion is needed for this input.

\begin{proposition}[An exact interval coupling, including the near-boundary remainder]
\label{f16v24:interval-coupling}
Let $I=[a,c]\cap\mathbb Z\subseteq[0,n]$ be nonempty. If two full conditional
laws on $I$ differ only in the left boundary slice, with the right boundary
fixed or natural, there is a coupling of their \emph{entire} slice paths for
which
\begin{equation}
 \mathbb E\sum_{i\in I}1_{\{Z_i\ne Z_i'\}}\le4b.
 \label{f16v24:boundary-susceptibility}
\end{equation}
The same is true for a right-boundary change. The bound is independent of
$|I|,B,n$, and the actual boundary values, in the stated width regime.
Both endpoint changes can be paid with the sum of their separate budgets.
\end{proposition}
\begin{proof}
Use the actual forward Markov disintegration of the interval, with its fixed
right endpoint included in its continuation messages. Starting at its left
boundary, put $T=|I|$ and take
$M=\max\{0,\lfloor T/b\rfloor-1\}$ complete steps of length $b$.
The residual segment has fewer than $2b$ slices, and for $T\ge2b$ has
at least $b$ slices. If $T<2b$, couple the interval arbitrarily; its
Hamming cost is already less than $2b$.

For every complete step before the last segment, the target is at least $b$
from the fixed right endpoint. Given its two starting values, couple its
endpoint marginals by the usual common-part coupling; when the starting values
are equal use identical samples. Equation \eqref{f16v24:channel-input} bounds
the conditional probability of unequal step endpoints, given unequal starting
values, by $q$. Given both endpoints of a segment, fill its internal slices
from their \emph{exact} conditional bridge laws. These are the marginals of
the original interval disintegration. Their normalizers are not discarded.
If the starting values are equal, the entire segment may be drawn identically.
If they are unequal, its inside cost is at most $b$, even when its last slices
coalesce. Once the sampled step endpoints agree, the full remaining right
conditional can be coupled identically; the common fixed right factor has
not changed.

If $M$ complete steps were used, their expected total cost is at most
$b\sum_{j=0}^{M-1}q^j$. The residual segment costs at most $2bq^M$.
Thus the total is at most $b/(1-q)+2b\le(10/3)b<4b$. The construction retains
the native law in each component by Markov disintegration and conditional
bridge filling. These are measurable finite-dimensional probability kernels;
their common-part couplings use their actual densities. No simultaneous
coupling of all interval boundaries is required. The case $M=0$ is included.
Reverse time for a right change. For two changed endpoints, concatenate the
two comparisons through a path law with one endpoint from each, and use
conditional gluing of couplings and the metric triangle inequality. Natural
endpoints create no additional term.
\end{proof}

The terminal $2b$ term is essential. A forward step too close to the fixed
right boundary was \emph{not} proved contracting by V22. We stop before that
step and pay its whole output, instead of treating the opposite conditioning
as harmless.

\subsection{Overlapping interval resampling with unit coverage at every time}

Set $\ell=16b$. Use the indexed clipped intervals
\begin{equation}
 I_s=[s,s+\ell-1]\cap[0,n],\qquad 1-\ell\le s\le n,
 \qquad
 \mathcal L_{\rm slab}=\frac1\ell\sum_{s=1-\ell}^{n}(I-E_{I_s}).
 \label{f16v24:slab-generator}
\end{equation}
Here $E_{I_s}$ conditions on all unmoving cells outside the \emph{entire}
space--time slab $I_s\times\mathbb T_m^3$. Each indexed summand has rate
$1/\ell$. Some clipped intervals can coincide when $n+1<\ell$; they are
\emph{not} deduplicated. Every time slice belongs to exactly $\ell$ indexed
intervals, so its coverage rate is exactly one, including either dressed end.
There is no factor $1/(n+1)$ in this clock.

\begin{theorem}[A full-$L^2$ native slab gap with no duration loss]
\label{f16v24:slab-gap}
Under \eqref{f16v24:channel-input}, the exact reversible generator
\eqref{f16v24:slab-generator} satisfies
\begin{equation}
 \boxed{\quad \mathcal L_{\rm slab}\succeq\tfrac12(I-\mathbb E_P),\qquad
 \operatorname{Var}_P f\le\frac2\ell
           \sum_{s=1-\ell}^{n}\mathbb E_P\operatorname{Var}_P(f\mid Z_{I_s^c}).
 \quad}
 \label{f16v24:slab-gap-bound}
\end{equation}
It holds for every $f\in L^2(P)$, not only local sources or polynomial tests.
No small-field condition is imposed on a conditional exterior.
\end{theorem}
\begin{proof}
Treat complete slices as coordinates and give each weight one in
Theorem~\ref{f16v24:general-gap}. A changed slice $j$ is contained in exactly
$\ell$ indexed intervals: $j-\ell+1\le s\le j$. Its death rate is one.
For a slab not containing $j$, its conditional law can depend on this change
only if $j$ is an existing immediate left or right boundary. There are at
most two such indexed intervals, namely $s=j+1$ and $s=j-\ell$ when they
exist. All other exterior changes cancel from the conditional by the exact
temporal Markov property. Proposition~\ref{f16v24:interval-coupling} bounds
each affected output cost by $4b$. Hence the margin is at least
\[
                    1-\frac{8b}{\ell}=\frac12.
\]
Apply Theorem~\ref{f16v24:general-gap}. Short histories and clipped endpoints
were included in both counts. Strict positivity of the native density and
all its original endpoint factors ensures these are its own orthogonal
conditional projections. No independence between simultaneous overlapping
slabs and no product of their grand couplings is asserted.
\end{proof}

The gap belongs to a \emph{newly defined auxiliary sampler}, not to a power of
the physical transfer. Its time normalization has been specified explicitly.
An individual move may change $B\ell$ cells. A comparison to original local
updates is therefore still required, and is now paid rather than assumed.

\subsection{The exact cost of returning to the original balanced updates}

Let $\omega=(X,r)$ be V14's balanced compact variables. One elementary update
resamples all five chords of one cube at one time, or one nonroot moving
port. Denote their rate-one projections by $\widehat Q_\alpha$ and set
\[
 \widehat{\mathcal L}_P=\sum_\alpha\widehat Q_\alpha,
 \qquad \widehat{\mathcal E}_P(f)=\sum_\alpha
             \mathbb E_P\operatorname{Var}_P(f\mid\omega_{\alpha^c}).
\]
This is the original unweighted balanced generator, not the whole-slice
sampler. Its weighted version is a separate clock and is not substituted.
The exact relation to unmoving ports is
\begin{equation}
             r_t=f_t(X_t)^{-1}h_t f_{t+1}(X_{t+1}).
 \label{f16v24:frame-change}
\end{equation}
All frame operations are within their original spatial cubes.

\begin{proposition}[Bounded classical oscillation gives a conditional clock comparison]
\label{f16v24:finite-comparison}
For any finite set $D$ of elementary balanced update groups, its complete
conditional satisfies
\begin{equation}
 \mathbb E_P\operatorname{Var}_P(f\mid\omega_{D^c})
 \le \exp(C_r\gamma|D|)
             \sum_{\alpha\in D}\mathbb E_P
                       \operatorname{Var}_P(f\mid\omega_{\alpha^c}).
 \label{f16v24:conditional-clock}
\end{equation}
The constant is independent of $D,B,n$ and every fixed exterior. Its potentially
large exponential is retained; it is not a coupling-uniform local gap.
\end{proposition}
\begin{proof}
In balanced variables the reference measure is a product over chord-cube
blocks and nonroot ports. The chord factors retain exactly the original
principal-chart Lebesgue measure with Haar exponent $1$ or $1/2$; moving ports
retain Haar measure by conditional Haar invariance. Normalize these individual
finite reference measures only for this comparison and call their product on
$D$ $\nu_D$. Their normalization scalars cancel in the conditional density.
No ratio of a fractional endpoint Haar factor to full Haar is bounded here.

The remaining conditional density is proportional to $e^{-U_D}\nu_D$,
where $U_D$ is the \emph{classical} action of all terms meeting $D$.
Every Wilson term lies between $0$ and $2\gamma$, has fixed carrier and bounded
incidence in these coordinates. V5/V14's cube-local frame dependence preserves
these bounds. Each endpoint quadratic is bounded by $C\gamma$ per incident
chord cube on its complete principal domain. Consequently
\[
            \operatorname{osc}U_D\le C_r\gamma|D|.
\]
Exterior-only terms cancel, but every crossing term remains. This bound uses
neither convexity nor a typical-exterior event.

For completeness, if $\mu_D=p\nu_D$ and $m\le p\le M$, then
$M/m\le e^{\operatorname{osc}U_D}$. The variational characterization of
variance and the product variance inequality give
\[
 \operatorname{Var}_{\mu_D}f
 \le M\operatorname{Var}_{\nu_D}f
 \le M\sum_{\alpha\in D}\inf_{g\text{ measurable off }\alpha}
                             \int|f-g|^2d\nu_D
 \le\frac M m\sum_{\alpha\in D}
                  \mathbb E_{\mu_D}\operatorname{Var}_{\mu_D}(f\mid\omega_{D\setminus\alpha}).
\]
The product inequality follows, for example, by revealing independent
coordinates successively and applying conditional Jensen to each martingale
difference. The infimum formula proves the last inequality even for unbounded
$L^2$ tests, by truncation. Conditioning the other variables of $D$ restores
the original elementary single-group conditional. Integrate over the actual
exterior marginal to obtain \eqref{f16v24:conditional-clock}. All measures have
the same full support up to reference-null sets, so the bounded density
comparison is legitimate.
\end{proof}

\begin{proposition}[A slab changes its preceding balanced port as well]
\label{f16v24:slab-enlargement}
For an unmoving interval $I=[a,c]$, let $D(I)$ contain its balanced chord cubes
at times $a,\ldots,c$ and moving ports at times $a-1,\ldots,c$, omitting
nonexistent ports at $-1$ or $n$. Then
\begin{equation}
 \mathbb E_P\operatorname{Var}_P(f\mid Z_{I^c})
 \le\mathbb E_P\operatorname{Var}_P(f\mid\omega_{D(I)^c}),\qquad
 |D(I)|\le8B(|I|+1).
 \label{f16v24:enlargement}
\end{equation}
For the clipped indexed slabs, each chord-cube update belongs to exactly
$\ell$ sets $D(I_s)$ and each moving-port update to exactly $\ell+1$.
\end{proposition}
\begin{proof}
The unmoving exterior fixes all $X_i,h_i$ outside $I$. Formula
\eqref{f16v24:frame-change} then fixes every balanced variable outside the
stated $D(I)$: if $t\notin[a-1,c]$, both $X_t,X_{t+1}$ and $h_t$ are fixed.
Thus $\sigma(\omega_{D(I)^c})\subseteq\sigma(Z_{I^c})$. Conditional expectation
is an orthogonal projection, so conditioning on the latter, larger sigma
algebra decreases the averaged conditional variance. This proves the inequality;
it does \emph{not} identify the two conditional kernels.

There are $B|I|$ chord-cube groups and at most $7B(|I|+1)$ port groups.
For coverage, a chord time $i$ occurs in $\ell$ intervals as already counted.
A port time $t\in[0,n-1]$ occurs when $I_s$ contains $t$ or $t+1$; the union
of start indices is $t-\ell+1\le s\le t+1$, of size $\ell+1$. Both time
indices exist, so clipping changes neither count. In general omitting the
preceding port would be false: changing $X_a$ can change $r_{a-1}$ with
$h_{a-1}$ held fixed. The elementary updates are therefore not the same in
moving and unmoving coordinates.
\end{proof}

\begin{theorem}[A duration-uniform bound in the original rate-one clock]
\label{f16v24:original-gap}
Under the channel and width hypotheses above, with $\ell=16b$,
\begin{equation}
 \mathcal L_{\rm slab}\preceq
 \left(1+\ell^{-1}\right)e^{C_r\gamma B(\ell+1)}\widehat{\mathcal L}_P,
 \qquad
 \boxed{\quad\widehat\lambda_P\ge
               \frac14 e^{-C_r\gamma B(\ell+1)}.\quad}
 \label{f16v24:original-gap-bound}
\end{equation}
All inequalities hold on the complete native $L^2$ space, with original endpoint
factors. In particular, for every fixed $p,A_0$,
\begin{equation}
 B\le A_0\gamma^p\quad\Longrightarrow\quad
       \widehat\lambda_P\ge\tfrac14
                      e^{-C_{p,r}\gamma B\log\gamma},
 \label{f16v24:polynomial-gap}
\end{equation}
and in the specified exponential-width channel regime,
\begin{equation}
 B\le e^{\beta_{\rm ch}\gamma}\quad\Longrightarrow\quad
       \widehat\lambda_P\ge\tfrac14 e^{-C_r\gamma^2B}.
 \label{f16v24:exponential-gap}
\end{equation}
The entry thresholds are those of the corresponding channels, enlarged by
fixed requirements only. These are positive bounds uniform in duration, not
positive bounds uniform in $\gamma$ or unrestricted spatial volume.
\end{theorem}
\begin{proof}
Use \eqref{f16v24:enlargement} for each slab and then
\eqref{f16v24:conditional-clock} on its enlarged balanced block. Its group
count is at most $8B(\ell+1)$, absorbed into $C_r$. Sum the rates $1/\ell$.
The exact coverage counts bound the rate assigned to any elementary update
by $1+1/\ell$. This proves the operator-form comparison. Combine it with
\eqref{f16v24:slab-gap-bound} on centered space. The sharper lower constant
is $[2(1+1/\ell)]^{-1}$; $1/4$ is a uniform simplification. Insert the inherited
bounds $b=O_{p,r}(\log\gamma)$ or $b=O_r(\gamma)$ to obtain the two displayed
regimes. No source-covariance estimate was used in the proof of this gap.
\end{proof}

This pays a genuinely missing clock step for a concrete dynamics, but not at
a useful weak-coupling constant. Inserting \eqref{f16v24:original-gap-bound}
into V16's gap-only source estimate is generally \emph{worse} than V23's direct
source theorem. It is not reported as an improved covariance bound. Its new
content is all-$L^2$ coercivity in the original rate-one clock without a duration
factor, and an exact identification of its expensive finite-block comparison.
A sharp estimate for those finite conditional blocks could improve this cost;
the present proof has not established such an estimate.

\subsection{Why common-core success alone does not verify the spatial ledger}

For the spatial diagnostic below, the Hamming coordinates are complete
\emph{unmoving cells}, not the elementary balanced chord/port groups.
Returning a spatial cell-block dynamics to those elementary updates would
require the analogous frame-support enlargement as well.
The same general criterion can be tested on local space--time cell blocks. Its
missing quantity is $I_{D,j}$, the expected \emph{entire} resampled output
disagreement due to a \emph{single} changed exterior input. V23's common map
already supplies a valid but generally inadequate upper bound. If the block
has core $A$, fringe $D\setminus A$, and common core mass $\alpha$, couple the
two full outputs with its same success decision and seeds. On success the
core agrees, but every fringe coordinate may still disagree. On failure every
coordinate of $D$ may disagree. For unit weights this proves
\begin{equation}
 I_{D,j}\le U_D:=|D\setminus A|+(1-\alpha)|A|,
                  \qquad j\in\partial D.
 \label{f16v24:crude-fringe}
\end{equation}
It uses the actual conditional filling of the fringe. It is not an assertion
that different native error fields are independent.

\begin{proposition}[The whole-fringe estimate has the wrong spatial scaling]
\label{f16v24:fringe-obstruction}
For translated nearest-neighbour boxes $V_R=\{1,\ldots,R\}^4$ and core
$A_R=\{2,\ldots,R-1\}^4$, with $R\ge3$, normalize each box rate by $R^{-4}$.
The uniform influence table supplied solely by \eqref{f16v24:crude-fringe}
would give the certified margin
\begin{equation}
 1-\frac{8}{R}\left[R^4-(R-2)^4+(1-\alpha)(R-2)^4\right]<0,
 \label{f16v24:fringe-margin}
\end{equation}
even at $\alpha=1$. This is failure of this \emph{sufficient bound}, not proof
that the actual dynamics has a nonpositive gap or that better couplings cannot
work. The geometry is an explicit four-dimensional comparison template, not
an assertion that the native interaction graph has exactly these eight faces.
\end{proposition}
\begin{proof}
A site belongs to $R^4$ translated boxes and lies on the exterior nearest-neighbour
collar of $8R^3$ translated boxes, away from finite-size identifications.
Thus coverage is one, whereas the uniform birth estimate is $8U_D/R$.
Even at $\alpha=1$, the two opposite layers in one coordinate give
$R^4-(R-2)^4\ge2R^3$. The displayed margin is therefore at most
$1-16R^2<0$. Larger fringes only worsen this particular table. No lower bound
on the \emph{true} $I_{D,j}$ has been inferred from its upper bound.
\end{proof}

For comparison, if a translated finite-range four-dimensional box family
could be given \emph{full native} pair couplings satisfying
\begin{equation}
 I_{D,j}\le\Xi_\gamma\quad\hbox{for each changed collar input,}\qquad
 \frac{|\partial D|}{|D|}\le\frac{C_\rho}{R},
 \label{f16v24:spatial-target}
\end{equation}
with $\Xi_\gamma$ independent of the box side $R$ and the ambient size, then
$R\ge2C_\rho\Xi_\gamma$ would give $g\ge1/2$ in unit-coverage block time.
This follows directly from Theorem~\ref{f16v24:general-gap}. An actual cover on
a finite path or small torus must use its true coverage and boundary counts;
no translation-invariance of $P$ is needed, but those counts cannot be omitted.
A product-density comparison on correctly enlarged balanced supports would
then pay the elementary-clock cost; a cell update is not identified with an
elementary balanced update. This is a \emph{conditional criterion};
\eqref{f16v24:spatial-target} is not proved for the unrestricted native law.

The temporal argument verifies precisely such a boundary-sensitive quantity
with $\Xi=4b$ and only two incoming slices. Its success did not come from
multiplying unrelated core refresh probabilities. Spatially, the near-boundary
fringe can carry a change tangentially, and the common-core event alone places
no localized bound on that propagation. Spinka's boundary-subset condition
records the analogous issue for grand couplings; our spectral criterion needs
only pairwise single-input couplings, but it still needs their full-output cost.

\subsection{What changes in the program and what remains open}

There are two distinct new outputs. First, the complete conditional interval
has an actual coupling with affected-slice budget $4b$, uniformly in its
length and in both compact boundary values. This constructs a reversible
full-block sampler with gap $1/2$ and \emph{unit coverage per slice}. Second,
the exact moving/unmoving support change and complete density comparison return
it to the original balanced sampler with its explicit exponential price.
The latter bound is uniform in duration within the stated polynomial or
exponential spatial-width regime. It applies to every native $L^2$ function,
not just the chosen response bank.

Neither result gives an original-sampler gap uniform in coupling. Nor does it
improve V23's source asymptotics. The primitive/common-source covariances, V19's
positive shell return and V12's different cylinder return have not been
redefined. In particular the exponential cost in
\eqref{f16v24:original-gap-bound} cannot be suppressed by saying that each
single update has rate one. Overlap, coordinate conversion, and conditional
mixing are separate parts of the clock ledger, and all have been paid here.

The next spatial calculation is sharper than another common-core estimate:
bound the total resampled disagreement caused by one changed collar input,
with the genuine buffer law retained, or obtain an equivalent directional
operator estimate. The whole-fringe bound fails that test already on the
simplest four-dimensional box. Improving only its common success probability
cannot fix its surface-volume scaling. Alternatively, sharpening the native
finite-block comparison would improve the original-clock result in the existing
width regime, but would not itself remove the complete-separator cost.

All clocks in this appendix are auxiliary resampling clocks. No statement
identifies them with Euclidean transfer time. The fixed rescaled bridge window,
root Gauss restriction, infrared contact coercivity, same-top physical Wilson
transfer comparison, independently calibrated physical time, and nontrivial
interacting continuum construction remain separate obligations. No physical
Yang--Mills mass gap is claimed.

\subsection{Diagnostics, preservation, and proof status}

The new diagnostic is
\begin{center}\small\path{verify_ym_f16_v24_clock_sensitivity.py}.\end{center}
It checks exact finite conditional projections, the weighted influence ledger,
normalized inhomogeneous path kernels with a retained opposite endpoint,
clipped-interval coverage and boundary counts, the preceding-port enlargement,
product-density comparisons and the constants in the fringe and clock estimates.
The new checker passed 1,128 exact assertions. The finite path and spin
fixtures are not native Wilson integral quadratures.
Native compact integral suprema and entry thresholds are not enclosed by those
fixtures. The proof of the theorems is the argument above with its identified
analytic premises, not a consequence of the finite assertion count.

All twenty supplied V4--V23 diagnostics were rerun unchanged and passed.
The package records their actual outputs, source hashes, the new diagnostic,
predecessor checks, compilation and rendered-page inspection. The cumulative source changes
only the title version and inserts this exact appendix before the final
document-ending command; its reverse operation recovers V23 byte for byte.
No inherited claim, normalization or bibliography has been silently changed.
The scope of the inherited analytic archive has not been independently
recertified. The next scheduled companion checkpoint remains V25.

\begin{firewall}
V24 proves an all-observable gap for exact overlapping temporal-slab resampling
in a specified unit-coverage clock. It pays the conversion to the original
rate-one balanced chord/port sampler and obtains a positive duration-uniform
bound with an explicit coupling/width-dependent exponential cost. It proves
a boundary-input sensitivity criterion and diagnoses the insufficient
whole-fringe bound. It does not prove unrestricted spatial mixing, a
noncollapsing original-sampler gap as coupling increases, the missing spatial
sensitivity estimate, a physical transfer gap, or an interacting continuum
mass gap. Every native conditional, compact buffer and endpoint factor is
retained.
\end{firewall}
% END F16 V24 APPEND

% BEGIN F16 V25 APPEND -- 2026-09-17
% Insert immediately before the final document-ending command in V24.
% The predecessor body and all its probability/source conventions are unchanged.

\clearpage
\section{V25: interface-cost comparison and an original-clock finite-width baseline}
\label{f16v25:context}

V24 pays for converting a whole native slab update to elementary balanced
updates by the oscillation of the \emph{entire} conditional action. That costs
an exponential in the slab volume. It is not necessary to pay every interaction
in that comparison. Exchange the coordinates of two configurations in a fixed
order. Terms wholly on one side of the moving cut cancel \emph{exactly} between
the two configurations and their complementary recombinations. Only terms
crossing the cut remain. We give this comparison for continuous product spaces,
with the conversion from an independent marginal draw to the actual one-site
conditional included explicitly.

For the native slab this replaces a volume cost by the smaller of a temporal
and a spatial interface cost. It also makes a different, elementary temporal
input useful: compact bounded interactions give a very small but positive
one-step common component under every continuation message. Combining these
two finite arguments gives a positive lower bound for the original rate-one
balanced sampler, uniform in duration, for \emph{every} finite spatial width.
The lower bound decays exponentially with that width and with coupling. It is
not a thermodynamic spectral floor, and it does not replace V23's sharper
source estimates in their stated regimes.

\subsection{V25 checkpoint: source ownership and the changed comparison}

The supplied V24 appendix and analytic audit were read in full. The V5 exact
unmoving action and its local carriers, V14's balanced coordinates, and V24's
conditional-support enlargement were checked. The mounted f14 V100, f15 V100,
and Grand Tensor V076 sources were searched for cutwidth, coordinate-splicing,
and canonical-path comparisons after two scoped semantic searches returned no
results. The f15 hit for ``canonical path indicators'' concerns geometric
repair paths, not this Gibbs variance comparison. The monograph's random-scan
and blocking warnings remain binding. A failed search is not evidence of
archive-wide novelty; this is a targeted checkpoint, not recertification of
all historical analytic claims or of unavailable companion programs.

The method is classical. Berger--Kenyon--Mossel--Peres,
\href{https://arxiv.org/abs/math/0308284}{\emph{Glauber Dynamics on Trees and
Hyperbolic Graphs}}, Proposition 1.1 and its Section 2.1 proof, relates an Ising
relaxation bound to the number of interactions crossing an ordering cut. The
statement and the canonical-path proof were examined, including their rate
convention. We do not quote its finite-spin theorem for a continuous Wilson
law. The finite-product proof below treats arbitrary probability reference
factors and many-coordinate interactions directly. No lower bound on a point
mass of a continuous spin is required.

The two supplied import memoranda propose separating order gain from locality
cost, and retaining a genuine dependence estimate before a global functional
inequality. Here there is no new perturbative order and no new global spatial
coupling. The changed object is the cost of the \emph{same-law clock comparison}.
The new unrestricted-finite-width baseline below uses bounded classical words,
product references, exact temporal disintegration, and V24's proved block-gap
criterion. It does not use the V16 moment theorem, the small-angular one-form
realization, or the V17--V23 asymptotic conclusions. The sharper versions may
reuse V22's better temporal buffer as a separate input. The next scheduled
companion checkpoint is V30.

\subsection{A two-copy identity that charges only an interface}

Consider a finite product of standard Borel spaces, with probability reference
\(\nu=\bigotimes_{i=1}^{N}\nu_i\), and the exact probability
\[
 d\mu=Z^{-1}e^{-U}\,d\nu,\qquad U=\sum_X\Phi_X.
\]
The finitely many interaction functions are bounded and measurable. One-site
terms may first be incorporated into their own reference factor; this changes
neither \(\mu\) nor its conditional resampling. All remaining \(X\)'s can
therefore be taken to have at least two elements. A fixed exterior parameter,
if present, is substituted in each interaction before its free support is
counted. There is no restriction on those exterior values.

Fix an ordering, relabelled \(1,\ldots,N\). Write
\[
 S_i=\{1,\ldots,i-1\},\quad
 c_i=\sum_{X:\,X\cap S_i\ne\varnothing,\ X\cap S_i^c\ne\varnothing}
                  \operatorname{osc}\Phi_X,\quad
 s_i=\sum_{X\ni i}\operatorname{osc}\Phi_X,
\]
where the complement is within this free product. Define
\begin{equation}
 \mathfrak c_\prec=\max_i c_i,\qquad
 \mathfrak s=\max_i s_i.
 \label{f16v25:interface-costs}
\end{equation}
The second quantity is a one-coordinate interaction cost. It is not the
oscillation of all terms in \(U\).

For configurations \(x,y\), let
\[
 z^i=(y_1,\ldots,y_i,x_{i+1},\ldots,x_N),\quad z^0=x,
 \qquad w^{i-1}=(x_{S_i},y_{S_i^c}).
\]
Thus \(z^{i-1}=(y_{S_i},x_{S_i^c})\). These are coordinate exchanges, not
interpolations on the group or approximations of a conditional density.

\begin{proposition}[Exact two-copy cancellation]
\label{f16v25:splice}
Let \(p=d\mu/d\nu\). The coordinate exchange is a measurable
involution preserving \(\nu\otimes\nu\). Its density ratio is
\begin{equation}
 \frac{p(x)p(y)}{p(z^{i-1})p(w^{i-1})}
 =\exp\{-U(x)-U(y)+U(z^{i-1})+U(w^{i-1})\}
 \le e^{2c_i}.
 \label{f16v25:splice-ratio}
\end{equation}
Every term with support on only one side of the cut cancels, even if it is
large, nonconvex, or depends on a fixed exterior. The partition functions
cancel exactly in this two-copy ratio.
\end{proposition}
\begin{proof}
At each coordinate either exchange the two entries or leave them alone. Its
reference factor is \(\nu_i\otimes\nu_i\), so product Fubini proves invariance
and the map is its own inverse. A term supported in \(S_i\) has its two values
exchanged; one supported in \(S_i^c\) has them retained. Both contributions to
the exponent are zero. For a crossing term its two new values minus its two
old values are at most twice its oscillation. Sum only those bounds. This
proves the identity and inequality. The argument does not commute the factors
inside a noncommutative Wilson word: a whole scalar word potential is evaluated
at each of the four configurations before it is bounded.
\end{proof}

\begin{theorem}[Interface comparison in rate-one time]
\label{f16v25:path-comparison}
Let \(\mathcal E_i^\mu(f)=\mathbb E_\mu
\operatorname{Var}_\mu(f\mid\omega_{i^c})\). For every \(f\in L^2(\mu)\),
\begin{equation}
 \boxed{\quad
 \operatorname{Var}_\mu f
 \le N\sum_{i=1}^N e^{2c_i+2s_i}\mathcal E_i^\mu(f)
 \le N e^{2\mathfrak c_\prec+2\mathfrak s}
                         \sum_i\mathcal E_i^\mu(f).
 \quad}
 \label{f16v25:path-bound}
\end{equation}
This bounds the inverse gap of the sum of rate-one conditional projections.
No extra division by \(N\) is part of the generator. The factor \(N\) in this
bound is the cost of the chosen telescoping path, not a change of update clock.
\end{theorem}
\begin{proof}
Initially let \(f\) be bounded. The independent-copy variance identity and
Cauchy--Schwarz along the path give
\[
 \operatorname{Var}_\mu f
 =\tfrac12\mathbb E_{\mu\otimes\mu}|f(y)-f(x)|^2
 \le\tfrac N2\sum_i\mathbb E_{\mu\otimes\mu}
                       |f(z^i)-f(z^{i-1})|^2.
\]
Use Proposition~\ref{f16v25:splice} in the \(i\)th term. Writing \(z=z^{i-1}\)
and \(w=w^{i-1}\), that term is at most
\[
 e^{2c_i}\int\mu(dz)\mu(dw)
                          |f(z^{i\leftarrow w_i})-f(z)|^2.
\]
The incoming \(w_i\) in this expression has the \emph{unconditional marginal}
\(\mu_i^{\rm m}\). It is not the native heat-bath conditional at \(z_{i^c}\).
We pay that distinction before interpreting a Dirichlet form.

For two exteriors \(a,a'\), the difference of the one-coordinate interaction
potentials as a function of its free value has oscillation at most \(2s_i\).
The exact normalized conditional density ratio is consequently between
\(e^{-2s_i}\) and \(e^{2s_i}\). This follows by dividing the two exponential
weights and their two integrals; the normalized log ratio lies between the
minimum and maximum of the unnormalized log ratio minus one scalar in that
same interval. Averaging one of the conditional measures over its actual
exterior marginal gives
\begin{equation}
             \mu_i^{\rm m}(db)\le e^{2s_i}\mu_i(db\mid z_{i^c}).
 \label{f16v25:marginal-conditional}
\end{equation}
No lower bound on the mass of an individual \(b\) is used.

Substitute \eqref{f16v25:marginal-conditional}. The two free coordinates are
now independent draws from the \emph{same} conditional on \(z_{i^c}\); the
inner squared difference has integral twice its conditional variance.
This cancels the initial factor \(1/2\) and proves the first inequality of
\eqref{f16v25:path-bound}. The second uses the maxima. Truncating a real or
complex \(L^2\) function and using contraction of conditional expectation in
\(L^2\) extends the result. Bounded interaction weights give equivalent
finite reference and native \(L^2\) spaces, so all intermediate integrals are
legitimate. The proof covers continuous, atomic, or mixed reference factors.
\end{proof}

For an independent product this particular path estimate loses \(N\), although
product tensorization has constant one. The theorem is deliberately a robust
comparison, not a claim of sharpness. Its advantage over a whole-density bound
is the cancellation of all noncrossing interactions from the exponential.

\subsection{Actual balanced supports and two native ordering cuts}

Keep \(P=P_{\gamma,\mathbf y}\), \(B=m^3\), \(m\ge4\), \(n\ge1\), and the
prescribed bridge window. We use \(\gamma\ge1+r^2\) throughout for consistency
with the existing lineage. The original elementary coordinates are the five
chords of one cube at one time, regarded as \emph{one} update group, and each
of its seven nonroot moving ports. These are precisely V14's rate-one updates.
Their product reference factors retain all original Haar powers; the two
endpoint chord factors have power \(1/2\), not power one.

Assign each elementary group to its cube and its time. There are at most eight
such groups per cube--time cell. The exact balanced frame relation is local to
a cube. Consequently the remaining scalar classical word potentials have a
fixed bounded carrier diameter, a fixed incidence per elementary group, and
oscillation at most \(C\gamma\) per word. The endpoint quadratic is a sum of
one-cube functions and can be included in that cube's reference. After a
conditional exterior is fixed, any word with only one free elementary group
can likewise be included in its reference; all words with two or more free
groups are retained as interactions. This is a factorization of the same
conditional probability, not a cutoff or a reference law substituted for it.

\begin{proposition}[Two interface bounds for the exact enlarged slab]
\label{f16v25:native-cuts}
For an unmoving interval \(I=[a,c]\), let \(D(I)\) be the exact enlarged balanced
support of V24: chords at \(a,\ldots,c\), and moving ports at
\(a-1,\ldots,c\), with unavailable endpoints omitted. Put \(T=|I|\). Then
\begin{equation}
 |D(I)|\le8B(T+1),\qquad
 \inf_\prec\mathfrak c_\prec(D(I))+\mathfrak s(D(I))
 \le C_r\gamma\min\{B,m^2(T+1)\}.
 \label{f16v25:cuts}
\end{equation}
The bound is uniform in every fixed balanced exterior, all compact fast
coordinates, and the original endpoint factors. Periodic spatial seams are
included in the cut count.
\end{proposition}
\begin{proof}
The group count is V24's exact support count, not a newly shrunken support.
Let \(R_*\) bound the carrier range of the balanced words in every coordinate
direction, and let the fixed incidence bound be \(M_*\). Neither depends on
the regulator or the prescribed bridge values in their window. Assign every
word crossing an ordering cut to one incident group in its active layers;
there are at most \(M_*\) assigned words per such group.

First order by time, keeping all groups of a cube--time cell consecutive.
At any prefix, words crossing the cut must meet either the partially processed
time layer or at most \(R_*\) neighboring time layers on either side. There
are at most \(C B\) groups there, regardless of \(T\). A word with all free
groups earlier or all later cancels even if it also contains frozen variables.
There is no time-periodic seam. This gives \(\mathfrak c_\prec\le C\gamma B\).

Alternatively order by the first spatial cube coordinate. An ordering prefix
has fully processed spatial layers, at most one partially processed layer,
and unprocessed layers. A crossing word must meet a fixed-width band at the
moving cut or a fixed-width band at the periodic seam. Each band contains at
most \(C m^2(T+1)\) groups. Terms inside the partial layer are included by
counting that entire layer. If the period is smaller than the band width,
count the actual period once; the bound only improves after enlarging the
fixed constant. Missing chord or port groups at either temporal end can be
removed from this ordering without adding a crossing word. Thus this second
ordering gives \(C\gamma m^2(T+1)\).

Every elementary group meets at most a fixed number of remaining interactions,
so \(\mathfrak s\le C\gamma\). Since the minimum in
\eqref{f16v25:cuts} is at least one, absorb this last term and select the better
of the two orderings. No spatial translation invariance of \(P\) has been used.
\end{proof}

\begin{theorem}[Finite native block comparison with an interface, not volume, exponential]
\label{f16v25:native-comparison}
For every \(f\in L^2(P)\), the actual enlarged-block conditional obeys
\begin{align}
 &\mathbb E_P\operatorname{Var}_P(f\mid\omega_{D(I)^c})\notag\\
 &\qquad\le
 8B(T+1)\exp\!\left(C_r\gamma\min\{B,m^2(T+1)\}\right)
       \sum_{\alpha\in D(I)}
                 \mathbb E_P\operatorname{Var}_P(f\mid\omega_{\alpha^c}).
 \label{f16v25:native-block}
\end{align}
There is no restriction on the exterior values or on the ratio between \(B\)
and \(\gamma\). The target is still the complete native conditional.
\end{theorem}
\begin{proof}
Condition on the balanced exterior of \(D(I)\). Normalize the original product
reference factors, with the one-free-group terms included as just described.
Theorem~\ref{f16v25:path-comparison} applies to this finite product. Its
interaction bounds are uniform in the exterior by
Proposition~\ref{f16v25:native-cuts}. Fixing the other coordinates of \(D(I)\)
restores exactly the original elementary conditional on the right. Integrate
over the actual exterior marginal. Its conditional partition functions are
not replaced or estimated separately. This proves the statement.
\end{proof}

In particular the exponential cost is \(C\gamma B\), not
\(C\gamma B(T+1)\), when the temporal cut is chosen. When \(T+1\ll m\), the
spatial cut improves it further to \(C\gamma m^2(T+1)\). The remaining linear
factor counts elementary path steps. This does not prove a spatial coupling
of different boundary rows; it compares forms at one and the same boundary.

\subsection{A sharper paid return to the original update clock}

Let \(b\) be any valid V24 temporal channel buffer with step diameter at most
\(1/4\). Its complete-slab generator with \(\ell=16b\) has gap at least
\(1/2\), as already proved there. Its indexed clipped intervals, including
multiplicities, remain unchanged. Define
\[
       W_m(\ell)=\min\{B,m^2(\ell+1)\}.
\]

\begin{theorem}[Interface-priced original-clock comparison]
\label{f16v25:clock}
Under that channel hypothesis, on the entire native \(L^2(P)\) space,
\begin{equation}
 \mathcal L_{\rm slab}
 \preceq 8B(\ell+1)(1+\ell^{-1})
                 e^{C_r\gamma W_m(\ell)}\widehat{\mathcal L}_P.
 \label{f16v25:form-clock}
\end{equation}
\begin{equation}
 \boxed{\quad
 \widehat\lambda_P
 \ge \frac{e^{-C_r\gamma W_m(\ell)}}{32B(\ell+1)}.
 \quad} \label{f16v25:clock-gap}
\end{equation}
The unweighted balanced generator \(\widehat{\mathcal L}_P\) is still the sum
of the original rate-one five-chord and single-port innovations.
\end{theorem}
\begin{proof}
V24 proves the exact sigma-algebra inclusion
\(\sigma(\omega_{D(I)^c})\subseteq\sigma(Z_{I^c})\), including the preceding
port. Thus the slab conditional variance is bounded by the left side of
\eqref{f16v25:native-block}. Apply that estimate to each clipped interval and
use \(T\le\ell\). In the indexed family a chord group is covered \(\ell\)
times, and a moving port \(\ell+1\) times. Rates are \(1/\ell\), so the
maximum aggregate elementary rate is \(1+1/\ell\), not the number of slabs.
This proves \eqref{f16v25:form-clock}. The slab gap gives the lower constant
\([16B(\ell+1)(1+1/\ell)]^{-1}\); use \(1+1/\ell\le2\) for the displayed
simplification. No variance estimate for a selected source was used.
\end{proof}

For example, insert the already proved V22 buffers. In their respective
large-coupling regimes this gives
\begin{align}
 B\le A_0\gamma^p\quad&\Longrightarrow\quad
 \widehat\lambda_P\ge
 \frac{\exp[-C_{p,r}\gamma\min\{B,m^2\log\gamma\}]}{C_{p,r}B\log\gamma},
 \label{f16v25:polynomial-clock}\\
 B\le e^{\beta_{\rm ch}\gamma}\quad&\Longrightarrow\quad
 \widehat\lambda_P\ge
 \frac{\exp[-C_r\gamma\min\{B,m^2\gamma\}]}{C_rB\gamma}.
 \label{f16v25:exponential-clock}
\end{align}
In the first line \(p,A_0\) are fixed before taking the coupling limit; its
entry threshold may depend on both. Constants in the buffer length can be
absorbed outside and inside the exponential: for \(C\ge1\),
\(\min\{B,Cx\}\le C\min\{B,x\}\). These are better comparison costs than
V24's volume exponent, but they still vanish as \(\gamma\) increases. They
are not substituted into V23's source theorem as a better weak-coupling
covariance estimate.

\subsection{An elementary channel available at every finite spatial width}

The next argument is separate from all small-field preparation and expansion
results. Retain the exact native slice factorization
\[
 dP=Z^{-1}\prod_{i=0}^n\psi_i(z_i)\,d\nu_i(z_i)
                              \prod_{i=0}^{n-1}W_i(z_i,z_{i+1}).
\]
Each factor \(W_i=e^{-U_{i,i+1}}\) contains only the original classical terms
joining the two adjacent slices. There are at most \(C B\) such terms, each
with bounded oscillation \(C\gamma\). Fix
\begin{equation}
 A_{\gamma,B}=C_r\gamma B\ge
                     \sup_i\operatorname{osc}U_{i,i+1},\qquad
 \alpha_0=e^{-A_{\gamma,B}}.
 \label{f16v25:raw-cost}
\end{equation}
All within-slice weights, including the actual endpoint reference factors,
remain in \(\psi_i\nu_i\). Scalar shifts of a pair potential do not change
its normalized transition. Shift it so that \(e^{-A_{\gamma,B}}\le W_i\le1\).

\begin{proposition}[Bounded pair words contract under every continuation message]
\label{f16v25:raw-channel}
For any actual finite continuation message, including an arbitrarily fixed
opposite endpoint, every one-step slice transition has a common submeasure
of mass at least \(\alpha_0\). Its diameter is at most \(1-\alpha_0\).
Consequently a valid V24 buffer at \emph{every} finite width is
\begin{equation}
 b_0=\left\lceil(\log4)e^{A_{\gamma,B}}\right\rceil,
 \qquad \delta(K_{i,i+b_0})\le\tfrac14,
 \qquad b_0\le3e^{A_{\gamma,B}}.
 \label{f16v25:raw-buffer}
\end{equation}
The estimate is uniform in duration. It is usually far weaker than V22's
channel where that stronger theorem applies.
\end{proposition}
\begin{proof}
Write the transition, with its actual right continuation, as
\[
 K_i(x,dy)=\frac{W_i(x,y)\psi_{i+1}(y)h_{i+1}(y)\,d\nu_{i+1}(y)}{h_i(x)}.
\]
The measure \(\psi_{i+1}h_{i+1}\nu_{i+1}\) has a positive finite integral at
each finite regulator. Normalize it and call the resulting probability
\(\lambda_i\). This is a native probability retaining the continuation;
it is not a Gaussian or stationary replacement. Then
\[
 K_i(x,dy)=\frac{W_i(x,y)}{\int W_i(x,z)\,\lambda_i(dz)}\lambda_i(dy)
                                \ge\alpha_0\lambda_i(dy).
\]
The denominator is at most one. A fixed opposite boundary changes \(h\) and
hence \(\lambda_i\), but not the inequality for any row \(x\). Reverse time
for the other orientation. Exact serial composition gives diameter at most
\((1-\alpha_0)^b\le e^{-b\alpha_0}\), proving the middle assertion of
\eqref{f16v25:raw-buffer}. The ceiling and \(\log4<2\) give its final bound.
If an interval is too short to contain a full buffer, V24's clipped-interval
and terminal-segment argument already pays its entire length. No long interval
is assumed to fit into a shorter history.
\end{proof}

\begin{theorem}[Duration-uniform original-clock baseline at every finite cross-section]
\label{f16v25:all-width}
There are finite constants \(C_r,c_r>0\), independent of \(B,n,\gamma\), such
that the original balanced sampler satisfies
\begin{equation}
 \boxed{\qquad
 \widehat\lambda_{\gamma,\mathbf y,B,n}
                 \ge \frac{c_r}{B}e^{-C_r\gamma B}>0
 \qquad}
 \label{f16v25:all-width-gap}
\end{equation}
for every \(m\ge4\), \(B=m^3\), \(n\ge1\),
\(\max_{i,e}|y_{i,e}|\le r\), and \(\gamma\ge1+r^2\). There is no polynomial
or exponential upper restriction on \(B\) relative to coupling. The bound is
uniform in duration \emph{at a fixed cross-section}, not uniformly positive
as \(B\to\infty\) or \(\gamma\to\infty\).
\end{theorem}
\begin{proof}
Use Proposition~\ref{f16v25:raw-channel}, V24's interval coupling and its exact
clipped-slab gap with \(b=b_0\), \(\ell=16b_0\). None of their conditional
arguments needs the stronger channel's width hypothesis once the required
diameter is supplied independently. This is a new application of those
proved finite arguments, not deletion of a hypothesis from their earlier
width-specialized statements.

Apply Theorem~\ref{f16v25:clock} and choose its temporal cut, so that
\(W_m(\ell)\le B\). Since
\(\ell+1\le49e^{A_{\gamma,B}}\), we obtain
\[
 \widehat\lambda_P\ge
 \frac{1}{1568B}\exp[-C_r\gamma B-A_{\gamma,B}].
\]
Absorb the two fixed interface constants into one \(C_r\) and take, for
example, \(c_r=1/1568\). All finite reference factors, endpoint messages and
conditional normalizers were kept. In particular no exponentially small
one-step common mass has been renamed a coupling-independent quantity. The
large buffer contributes its exponential only through the remaining linear
path prefactor, not a second exponential in its volume.
\end{proof}

This theorem does not require V16--V23's analytic comparison inputs. It uses
the finite bounded-word specification and V24's separately proved Markov,
path-coupling, and support identities. Thus it supplies a more elementary
baseline against which later sharp claims can be checked. Its decay
\(e^{-C\gamma B}\) is still far too costly for the unrestricted spatial
spectral edge. The phrase ``every finite width'' must not be replaced by
``uniform in width.''

\subsection{The exact congestion left behind by the interface supremum}

There is also a useful same-law object underneath the exponential majorant.
For the general finite product above, fix \(i\), use the cut \(S_i\), and set
\[
 x=(w_{S_i},z_{S_i^c}),\qquad y=(z_{S_i},w_{S_i^c}),\qquad
 Q_i(z,w)=\frac{p(x)p(y)}{p(z)p(w)}.
\]
Let \(\mu_i^{\rm m}\) be the actual coordinate marginal. With positive
conditional-density versions, define
\begin{equation}
 a_i(z,b)=
 \frac{d\mu_i^{\rm m}}{d\mu_i(\cdot\mid z_{i^c})}(b)
       \mathbb E_\mu[Q_i(z,W)\mid W_i=b].
 \label{f16v25:congestion}
\end{equation}
It is an average over the complementary configuration under its
\emph{actual conditional law}, with no Gaussian normalization. Values on
reference-null pairs may be fixed arbitrarily. Write
\[
 d\mathbb B_i(z,b)=\mu(dz)\mu_i(db\mid z_{i^c}).
\]
This is the exact edge measure of a rate-one native conditional update.

\begin{proposition}[Normalized edge congestion, and the weight it must multiply]
\label{f16v25:edge-identity}
The nonnegative density \(a_i\) obeys
\begin{equation}
 \int a_i\,d\mathbb B_i=1,\qquad 0\le a_i\le e^{2c_i+2s_i},
 \label{f16v25:edge-normalization}
\end{equation}
and, for every bounded \(f\),
\begin{align}
 &\mathbb E_{\mu\otimes\mu}|f(z^i)-f(z^{i-1})|^2\notag\\
 &\qquad=
 \int a_i(z,b)|f(z^{i\leftarrow b})-f(z)|^2\,d\mathbb B_i(z,b).
 \label{f16v25:edge-exact}
\end{align}
Consequently the first line of the path proof can retain these exact weighted
innovations instead of replacing \(a_i\) by its supremum. Its mean one does
\emph{not} authorize that replacement by one inside a squared innovation.
\end{proposition}
\begin{proof}
Apply the reference-preserving exchange without estimating its ratio. The
resulting integral is against \(\mu(dz)\mu(dw)Q_i(z,w)\). Disintegrate the
second copy at \(w_i=b\), then convert its marginal to the exact conditional
\(\mu_i(db\mid z_{i^c})\). This is exactly
\eqref{f16v25:congestion} and proves \eqref{f16v25:edge-exact}. Apply the same
change of variables with integrand one to obtain
\(\int a_i\,d\mathbb B_i=\int d\mu\,d\mu=1\).
Proposition~\ref{f16v25:splice} and
\eqref{f16v25:marginal-conditional} give its upper bound. All these statements
refer to the same normalized law. The last warning follows because a
nonnegative factor with mean one can correlate with the displayed squared
increment; the identity includes that correlation.
\end{proof}

This identifies a sharper potential use of native geometry: control the
\emph{weighted innovations} in \eqref{f16v25:edge-exact}, or choose paths
whose interaction cuts have a better source-dependent cost. An unweighted
mean of \(a_i\), a rare bad-event probability, or an estimate for a different
core conditional does not supply that comparison. The present proof uses the
uniform supremum and retains its price. Normalization of \(a_i\) is global
under \(\mathbb B_i\), not row normalization at each \(z\); it is not
a replacement Markov transition.

For a simple check on the scope of boundedness alone, take two binary spins
with law \(\mu_J(x,y)\propto e^{Jxy}\), \(J>0\), and both updates at rate one.
The centered function \(x+y\) has generator eigenvalue \(1-\tanh J\); the
other centered eigenvalues are \(1+\tanh J\) and \(2\). Hence its gap is
exactly \(1-\tanh J\to0\). This finite-spin example is not a native
Yang--Mills counterexample. It shows why a uniform positive low-temperature
gap cannot follow from compact support, positivity and bounded finite-range
interactions alone. Native-specific information is needed for that stronger
conclusion.

\subsection{Checkpoint conclusions and the next noncircular task}

The new cost ledger distinguishes three pieces: a same-law interface ratio,
a marginal-to-conditional conversion, and a linear path length. For native
slabs the first costs the smaller of a spatial and temporal section, not the
whole slab volume. The previous port enlargement and exact coverage remain
paid. This gives sharper original-clock estimates in the already controlled
width regimes, and a baseline positive duration-uniform gap at \emph{every}
finite width without using the perturbative hierarchy.

The spatial single-input susceptibility asked for in V24 has not been proved.
The new path comparison does not couple two different boundary laws and
cannot be relabelled as that result. Nor do the elementary one-step common
components solve spatial composition: their mass decays with the full slice
width. Repeating the crude whole-fringe estimate or improving only a core
success probability would still leave the same obstruction.

The next useful alternatives are now explicit: pay the actual weighted
interface congestion using native structure, or establish a spatial
single-input full-output coupling whose affected volume stays controlled.
Either would be new information beyond boundedness. An all-boundary
conditional gap uniform in coupling is not assumed. The expensive baseline
is not inserted into a gap-only covariance bound to claim improvement over
V23. The stronger covariance and memory conclusions already proved there
remain limited to their stated width regimes, with all their original
coefficients and normalizers unchanged.

At this checkpoint the main native finite-product and update-clock inputs
are isolated from the more demanding V16--V23 analytic chain. This does not
certify that chain independently or invalidate it. The supplied memoranda
remain proposals rather than evidence for their suggested physical arrows.
No assertion about an external Navier--Stokes theorem is a premise here.
Only the supplied f14/f15/Grand Tensor and current-lineage companions were
reviewed; no unavailable program is described as checked.

All resampling clocks in this section are auxiliary. No physical Gauss
projection, unrestricted retained-bridge integration, infrared contact floor,
same-top physical-transfer comparison, physical-time calibration or
nontrivial interacting continuum limit follows from
\eqref{f16v25:all-width-gap}. In particular its width dependence has not been
exchanged for a uniform continuum mass.

\clearpage
\subsection{Diagnostics, preservation, and precise verification status}

The diagnostic is
\begin{center}\small\path{verify_ym_f16_v25_interface_clock.py}.\end{center}
The new diagnostic passed 1,292 exact assertions. It checks two-copy
cancellation for finite many-body potentials, nonuniform product references,
marginal-to-conditional conversion, and normalized edge congestion. An exact
continuous polynomial-density example is included. The native geometry checks
use a stated dominating support envelope of the literal fine-lattice words;
they include preceding ports and periodic seams, not numerical Wilson
integrals. Counts and fixture scope are recorded in the report. These finite
checks do not enclose analytic constants or constitute proof-assistant
verification. All twenty-one available V4--V24 diagnostics were rerun with
their sources unchanged and passed. Their actual outputs and hashes are
included in the package.

The cumulative V24 is preserved byte for byte except for its one title-version
change and insertion of this exact appendix before the final document-ending
command. A reconstruction script checks the reverse operation. The package
contains the input hashes, checkpoint and analytic audit, executable finite
diagnostics, their actual outputs, and compilation/render checks. No inherited
proof or bibliography is silently corrected. The next companion checkpoint
is V30.

\begin{firewall}
V25 proves a continuous-product interface-cost variance comparison and applies
it to the original balanced native supports, including their preceding ports
and periodic seams. It sharpens the paid slab-to-local clock conversion and
obtains a positive duration-uniform original-sampler bound at every finite
spatial width, with explicit exponential decay in coupling times width. It
also identifies an exactly normalized, innovation-weighted interface
congestion. It does not prove a gap uniform in width or coupling, the missing
spatial boundary-susceptibility theorem, unrestricted coherent covariance
control, a physical transfer gap, or an interacting continuum mass gap.
\end{firewall}
% END F16 V25 APPEND

% BEGIN F16 V26 APPEND -- 2026-09-17
% Insert immediately before the final document-ending command in V25.
% No inherited proof, probability, source, or normalization is edited.

\clearpage
\section{V26: single-input spatial coupling on a controlled collar}
\label{f16v26:context}

V25 leaves two distinct targets: control its actual innovation-weighted
interface density, or construct a spatial coupling that pays the whole output
for one changed collar input. We first test the density route. Its exact
factorization shows that its second moment can grow exponentially even for a
product of uniformly gapped interacting pairs. This does not disprove a
native weighted-form estimate, but it rules out interpreting an unweighted
congestion moment as a proxy for the gap. We then pursue the second target.

The new estimate is a \emph{single-input, full-output coupling} of the same
native conditionals. It is uniform in ambient volume and duration, but it
requires a declared bound $Q$ on the rescaled fields of the actual collar.
Within that window, the expected number of different output cells is at most
a constant times the changed input displacement, plus a counted exponentially
small exception. All free cells, including their large fields, are retained.
A finite colour sweep is used only to smooth test functions inside an
auxiliary convex conditional; it is not substituted for the original sampler.
A positive influence ledger on these admitted rows is not asserted to hold on
the excluded boundary rows.

\subsection{Dependencies and the direction selected by the audit}

The complete V25 appendix, analytic audit, and checkpoint were read. Two scoped
semantic searches of the current source, f14 V100, f15 V100, and Grand Tensor
V076 returned no matches. Exact mounted-text searches followed. The monograph's
pressure-flow congestion concerns graph currents, not V25's two-copy density.
F14 V99's action-density susceptibility concerns a particular gauge-invariant
observable, not a coupling of entire conditional outputs. These distinctions,
and the existing V24 single-input criterion, are retained. This is a targeted
non-duplication review, not an independent certification of the archives.

The analytic inputs to the positive result are precisely V16/V20's
conditional set-energy inequality and V22's restricted native Hessian and
product-ball covariance realization. In particular, the absolute one-form
boundary realization with corners is inherited, not newly certified. V17's
all-order expansion and V23's exponential-width covariance theorem are not
used to prove this coupling. V25's better finite-width gap remains unchanged.
The next scheduled companion checkpoint is V30.

For context, Dorlas--Savoie,
\href{https://arxiv.org/abs/2305.19371}{\emph{On the Wasserstein distance and the
Dobrushin uniqueness theorem}}, formulates block dependence using transport
costs; its Section 5 makes the all-boundary hypothesis explicit. That theorem
is not invoked for the present boundary-restricted, inhomogeneous law. Its
transport definitions, block-criterion statement, and duality statement were
examined, not its complete proof. The coloured-sweep derivative argument and
the finite-product coupling conversion actually needed here are proved below.
The two supplied import memoranda remain methodological proposals. A local
power estimate is not silently promoted to global dependence or physical time.

\subsection{The splice density can be large in a uniformly gapped product}

Use the finite-product probability $\mu$ and coordinate order of V25. For
coordinate $i$ let $S=\{1,\ldots,i-1\}$, $T=S^c$, and let
$\mu_S,\mu_T$ denote its actual marginals. At fixed finite regulator the
positive bounded density hypotheses of V25 give all Radon--Nikodym derivatives
below. For its exact edge probability
\[
 d\mathbb B_i(z,b)=\mu(dz)\mu_i(db\mid z_{i^c}),
\]
keep \emph{exactly} V25's density $a_i$.

\begin{proposition}[A marginal factorization of the actual splice congestion]
\label{f16v26:factorization}
Up to $\mathbb B_i$-null sets,
\begin{equation}
 a_i(z,b)=
 \frac{d(\mu_S\otimes\mu_T)}{d\mu}(z)
 \frac{d\mu_i(\cdot\mid z_S)}{d\mu_i(\cdot\mid z_{i^c})}(b).
 \label{f16v26:congestion-factor}
\end{equation}
The first factor destroys the correlations across the ordering cut. The
second restores the correct conditional distribution of the incoming value.
Their product is globally normalized on the update-edge space, as in V25;
neither factor is dropped from a squared innovation.
\end{proposition}
\begin{proof}
Before the $i$th replacement, $z=(y_S,x_T)$ and the incoming value is $b=y_i$,
where $x,y$ are independent with law $\mu$. Their joint law is exactly
\[
 \mu_T(dz_T)\mu_S(dz_S)\mu_i(db\mid z_S).
\]
Its density relative to $\mathbb B_i$ is the right side of
\eqref{f16v26:congestion-factor}. V25's unestimated coordinate-exchange
identity identifies the same joint law as $a_i\,d\mathbb B_i$.
Uniqueness of the density proves the assertion. This is an identity between
normalized marginals; it does not replace any conditional by a product law.
\end{proof}

\begin{proposition}[Exponential congestion moments do not diagnose a small gap]
\label{f16v26:congestion-test}
Fix $0<\theta<1$. Take $M$ independent pairs $(X_j,Y_j)\in\{-1,1\}^2$ with
probability $(1+\theta X_jY_j)/4$, together with an independent uniform spin
$V$. Order the coordinates as $X_1,\ldots,X_M,V,Y_1,\ldots,Y_M$.
For the $V$ update, V25's density is
\begin{equation}
 a_V=\prod_{j=1}^M(1+\theta X_jY_j)^{-1},\qquad
 \mathbb E_{\mathbb B_V}a_V=1,\qquad
 \mathbb E_{\mathbb B_V}a_V^2=(1-\theta^2)^{-M}.
 \label{f16v26:product-moments}
\end{equation}
Nevertheless the sum of the original rate-one single-spin updates has gap
$1-\theta$, independently of $M$. A termwise inequality replacing $a_V$ by
a uniformly bounded constant against its own $V$-innovation is impossible.
This is a finite-spin counterfixture for that proposed inference, not a
counterexample to native Yang--Mills mixing.
\end{proposition}
\begin{proof}
At the cut just before $V$, the two coordinate marginals are products of
uniform spins. The incoming $V$ law is uniform under both conditionings.
Equation \eqref{f16v26:congestion-factor} gives the product displayed above.
For each pair, $XY=\pm1$ with probabilities $(1\pm\theta)/2$. Its inverse
factor has mean one and second moment $(1-\theta^2)^{-1}$. Independence proves
\eqref{f16v26:product-moments}.

One pair has generator eigenvalues $0,1-\theta,1+\theta,2$: its linear
modes are $X+Y,X-Y$, and the remaining centered mode is $XY-\theta$.
The independent $V$ generator has nonzero eigenvalue one. Tensor products
and sums of these commuting generators therefore give gap $1-\theta$.
For the last assertion set $H=\{X_jY_j=-1\text{ for all }j\}$ and
$f=V1_H$. Its $V$ squared increment is supported on $H$, where
$a_V=(1-\theta)^{-M}$. The ratio of its weighted and unweighted $V$-edge
integrals is exactly that number. The other spin innovations are not zero;
this example does not rule out a comparison using their energies too.
\end{proof}

Thus a bound on a raw congestion moment is not selected as the next sufficient
criterion. The remaining argument follows the spatial single-input target
itself, with its boundary window stated rather than hidden in an average.

\subsection{Preparation for a bounded rescaled collar, including boundary cells}

Let $D$ be a nonempty finite set of complete \emph{unmoving} cells, let
$C=\partial_{\mathcal G}D$ be its actual interaction collar, and put $N_D=|D|$.
Let $P_D^q$ denote V20's complete native conditional, with all terms meeting
$D$ and all original endpoint factors retained. Only $q_C$ affects this law.
Assume
\begin{equation}
 \max_{c\in C}|q_c|\le Q,\qquad Q\ge0,\qquad
 \max_{i,e}|y_{i,e}|\le r.
 \label{f16v26:collar-window}
\end{equation}
No value farther away than $C$ is restricted. Write $e_c=|\zeta_c|^2$ and
retain $t_0=1/(24\pi^2)$. The parameter $Q$ is a bound on \emph{rescaled}
whole-cell Euclidean coordinates, not an angular bound independent of coupling.

\begin{proposition}[Uniform full conditional moments at a controlled collar]
\label{f16v26:prepared-Q}
For the signed source cube $|s_c|\le t_0$ on $D$, the same native conditional
satisfies
\begin{equation}
 \sup_{c\in D}\mathbb E_{P_{D,s}^q}e_c\le A_r(1+Q^2),\qquad
 \log\mathbb E_{P_D^q}\exp\!\left(\sum_{c\in D}t_ce_c\right)
       \le A_r(1+Q^2)\sum_{c\in D}t_c\quad(0\le t_c\le t_0).
 \label{f16v26:conditional-moments}
\end{equation}
Here and below constants do not depend on $D,B,n,Q,\gamma$, unless explicitly
indicated, and $\gamma\ge1+r^2$. In particular, for V22's fixed comparison
radius $\eta_*$, there are $c_*>0,A_r<\infty$ such that
\begin{equation}
 P_D^q\{\exists c\in D:\ e_c>\eta_*^2\gamma\}
 \le N_D\exp\{A_r(1+Q^2)-c_*\gamma\}.
 \label{f16v26:conditional-tail}
\end{equation}
Unlike arbitrary-compact-boundary preparation, this estimate needs no depth
allowance: it includes the first free layer next to the bounded collar.
\end{proposition}
\begin{proof}
The inherited local normalizer inequality, averaged under $P_{D,s}^q$, gives
for every $T\subseteq D$
\[
 \sum_{c\in T}u_c\le a_r|T|+b_r\sum_{d\in\partial T}\widetilde u_d,
 \qquad u_c=\mathbb E_{P_{D,s}^q}e_c.
\]
On $D$ put $\widetilde u=u$, and outside $D$ set $\widetilde u=Q^2$.
This dominates each frozen collar energy that can occur in the inequality;
values assigned beyond the collar are an auxiliary scalar extension, not a
change of the probability or its boundary. For an arbitrary ambient cell set
$S$, apply the inequality to $T=S\cap D$. Boundary terms lying in $S\setminus
D$ cost at most $b_rQ^2|S\setminus D|$; the others lie in $\partial S$.
Consequently
\[
 \sum_{c\in S}\widetilde u_c
 \le [a_r+(1+b_r)Q^2]|S|+b_r\sum_{d\in\partial S}\widetilde u_d.
\]
Iterate on ambient graph balls as in V16, with
$\rho_r=b_r/(1+b_r)<1$. Polynomial ball growth bounds the resulting series
$\sum_{k\ge0}\rho_r^k(k+1)^4$. At the fixed finite regulator the terminal
term tends to zero, since $\widetilde u\le\max(E_*\gamma,Q^2)$ and the ambient
vertex set is finite. This proves the first bound, uniformly throughout the
signed source cube. Integrating the \emph{actual} derivatives of the conditional
$\log Z_D(q;s)$ along $s=ut$, $0\le u\le1$, proves the second. Exponential
Markov in each free cell and a union bound prove
\eqref{f16v26:conditional-tail}, after absorbing $t_0$ in $A_r$ and taking
$c_*=t_0\eta_*^2$. No conditional partition function is approximated.
\end{proof}

\subsection{A finite sweep turns Hamming oscillations into local derivatives}

For the comparison only, use exactly V22's restricted native law $\pi_D^q$
on the fixed product $\mathcal D_{\gamma,D}$ of whole-cell balls of radius
$\eta_*\sqrt\gamma$. At its original coupling threshold and for collars in
these balls, retain its constants $\lambda,M>0$, bounded graph degree $\Delta$,
bounded mixed Hessian blocks, and absolute one-form covariance identity.
In particular all one-cell conditionals have Poincar\'e constant at most
$\lambda^{-1}$, and the inverse one-form blocks are at most
$\lambda^{-1}e^{-a_0d_{\mathcal G}(c,d)}$ for a fixed $a_0>0$.
These are the explicitly inherited analytic inputs, not full-law convexity.

For a real bounded function $f$ define its cell oscillation by
\[
 o_c(f)=\sup_{x_{c^c}=x'_{c^c}}|f(x)-f(x')|.
\]
Choose a proper colouring of the free interaction graph with $s\le\Delta+1$
colours $C_1,\ldots,C_s$. Each primitive carrier is a clique in that graph,
so, given their exterior, cells in one colour are conditionally independent.
Let $E_{C_l}^q$ be their exact conditional expectation and define
$f_0=f$, $f_l=E_{C_l}^qf_{l-1}$, $g_q=f_s$.

\begin{proposition}[A local derivative estimate with no derivative norm of the test]
\label{f16v26:sweep}
There are fixed $h,C_r<\infty$ such that, whenever $f$ is smooth and independent
of $q$, and $o_c(f)\le w_c$ with $w_c\ge0$,
\begin{align}
 \pi_D^qg_q&=\pi_D^qf,\notag\\
 \|\nabla_cg_q\|_\infty
       &\le C_r\sum_{d\in D:\ d_{\mathcal G}(c,d)\le h}w_d,\notag\\
 \|\nabla_{q_j}g_q\|_\infty
       &\le C_r\sum_{d\in D:\ d_{\mathcal G}(j,d)\le h}w_d,
                    \qquad j\in C.
 \label{f16v26:sweep-gradients}
\end{align}
The constants depend on the local analytic bounds, not the number of colours
actually occupied, the size of $D$, the radius $\eta_*\sqrt\gamma$, or any
initial derivative of $f$. This sweep is used for smoothing and invariance
only; neither its spectral gap nor its contraction is assumed.
\end{proposition}
\begin{proof}
For a colour $C_l$ and $j\notin C_l$, changing cell $j$ changes only the
conditional factors in $C_l\cap N(j)$. Product coupling of those factors,
with no better bound than probability one of a mismatch, gives
\begin{equation}
 o_j(E_{C_l}F)\le o_j(F)+\sum_{i\in C_l\cap N(j)}o_i(F),
       \qquad o_j(E_{C_l}F)=0\quad(j\in C_l).
 \label{f16v26:osc-propagation}
\end{equation}
A telescoping change of the output coordinates proves the same inequality
without selecting couplings. Thus after at most $s$ layers every oscillation
is bounded by a fixed local sum of the original $w$'s: the propagation range
increases by at most one per layer, with row and column sums at most
$(1+\Delta)^s$.

Differentiation of a \emph{normalized} one-cell conditional gives
\[
 \nabla_jE_iF=E_i\nabla_jF-\operatorname{Cov}_i(F,\nabla_jV),\qquad j\ne i.
\]
For a unit vector $v$ in cell $j$, conditional covariance Cauchy--Schwarz and
the one-cell Poincar\'e inequality give
\[
 |\operatorname{Cov}_i(F,\partial_{j,v}V)|
 \le \tfrac12 o_i(F)\lambda^{-1/2}
                  \|\nabla_i\partial_{j,v}V\|_\infty
 \le C_r o_i(F).
\]
This is zero for nonneighbours. For a whole colour, differentiate its product
conditional; the same terms are summed only over $C_l\cap N(j)$, with the
other conditional factors integrated. Hence
\begin{equation}
 \|\nabla_jE_{C_l}F\|_\infty
 \le\|\nabla_jF\|_\infty+C_r\sum_{i\in C_l\cap N(j)}o_i(F)
                       \quad(j\notin C_l).
 \label{f16v26:gradient-propagation}
\end{equation}
When $j\in C_l$ the output is independent of its old value, so its derivative
is zero. Each free cell has such a layer. All dependence on the initial
$\nabla_jf$ is therefore erased before the finite sweep ends. Subsequent terms
in \eqref{f16v26:gradient-propagation} and
\eqref{f16v26:osc-propagation} prove the free derivative bound with, for example,
$h=2s+2$. A collar derivative is never erased, but starts at zero because $f$
is independent of $q$. The same recursion, with only the finitely many free
neighbours of $j$ contributing, proves the collar bound. Each colour
expectation preserves $\pi_D^q$, proving invariance. All derivatives keep the
original measured Haar terms and the conditional normalizers. No uniform
Hessian bound for an arbitrary nonsmooth test has been used.
\end{proof}

\subsection{One changed collar input has a localized full-output response}

For weights $w=(w_c)_{c\in D}\ge0$, put
\[
 d_w(x,x')=\sum_{c\in D}w_c1_{\{x_c\ne x'_c\}},\qquad
 \mathsf W_w(\mu,\nu)=\inf_{\Pi\in\mathrm{Cpl}(\mu,\nu)}\int d_w\,d\Pi.
\]
This is a cost on whole-cell outputs, not Euclidean displacement cost. In
particular $\mathsf W_{\boldsymbol1}$ is the minimum expected number of unequal
cells. Separate optimal couplings of the scalar marginals would not bound it.

\begin{theorem}[Weighted boundary response in the restricted native law]
\label{f16v26:restricted-response}
For two collar values $q,q'$ in the same small-angular product, differing only
at cell $j\in C$, every bounded measurable $f$ with $o_c(f)\le w_c$ satisfies
\begin{equation}
 |\pi_D^{q'}f-\pi_D^qf|
       \le C_r|q'_j-q_j|\sum_{c\in D}w_c e^{-a d_{\mathcal G}(c,j)}
 \label{f16v26:weighted-response}
\end{equation}
for one fixed $a>0$. Its constants are independent of the block and the
collar size. No global gap of a complete native sampler is an input.
\end{theorem}
\begin{proof}
First take a smooth test. Along the collar segment $q(t)$ let $v=q'_j-q_j$,
$S_t=\partial_tV_D^{q(t)}$, and use the \emph{parameter-dependent} invariant
sweep $g_t$ above. Keeping both its derivative and the density derivative gives
\begin{equation}
 \frac{d}{dt}\pi_D^{q(t)}f
     =\pi_D^{q(t)}(\partial_tg_t)
                        -\operatorname{Cov}_{\pi_D^{q(t)}}(g_t,S_t).
 \label{f16v26:swept-response}
\end{equation}
The first term must not be dropped: the smoothing kernel depends on the
boundary. Its absolute value is at most
$C_r|v|\sum_{d_{\mathcal G}(c,j)\le h}w_c$ by
\eqref{f16v26:sweep-gradients}.

The free gradient of $S_t$ has support within one graph step of $j$ and each
block has norm at most $C_r|v|$. Insert the inherited one-form covariance
identity, its componentwise exponential inverse bound, and
\eqref{f16v26:sweep-gradients}. For a fixed $0<a<a_0$ this bounds the second
term by
\[
 C_r|v|\sum_{x\in D}\sum_{c:\ d_{\mathcal G}(c,x)\le h}
                             w_c e^{-a_0d_{\mathcal G}(x,j)}
 \le C_r|v|\sum_{c\in D}w_c e^{-a d_{\mathcal G}(c,j)}.
\]
The interchange of sums uses only a fixed-radius ball count. The finite-range
first term obeys the same bound after increasing $C_r$. Integrate in $t$.
Every intermediate collar remains in its declared convex product domain.

To extend to bounded measurable tests, extend $f$ outside each closed ball by
coordinatewise nearest-point projection and convolve with a product of smooth
approximate identities. The extension and convolution do not increase any
cell oscillation. At almost every interior point the approximation converges
to $f$; both endpoint probabilities have Lebesgue densities. Bounded dominated
convergence therefore applies in their joint $L^1$ space. The estimate has no
derivative norm of the test and survives the limit. The finite regulator is
held fixed during this approximation; its constants in
\eqref{f16v26:weighted-response} are already size independent.
\end{proof}

\begin{proposition}[From all weighted tests to one spatially localized coupling]
\label{f16v26:coupling-conversion}
There is a coupling $(X,X')$ of $\pi_D^q,\pi_D^{q'}$ such that simultaneously
for every $c\in D$,
\begin{equation}
 \mathbb P\{X_c\ne X'_c\}
       \le C_r|q'_j-q_j|e^{-a d_{\mathcal G}(c,j)}.
 \label{f16v26:restricted-coupling}
\end{equation}
Consequently
\begin{equation}
 \mathsf W_{\boldsymbol1}(\pi_D^q,\pi_D^{q'})\le C_r|q'_j-q_j|.
 \label{f16v26:restricted-total}
\end{equation}
This is one coupling for a \emph{fixed pair} of complete boundary rows, not a
grand coupling over all boundary values.
\end{proposition}
\begin{proof}
Here is a finite-product version of the transport step, to avoid changing the
measurable spaces to the nonseparable discrete topology. Partition each compact
cell ball by nested finite grids of mesh tending to zero, choosing grid faces
of zero mass for both one-cell marginals. Weighted Hamming cost on the finite
product of labels has finite Kantorovich dual equal to the supremum over
functions of coordinate oscillation at most $w_c$. Lifting such a function to
the continuous product preserves these oscillations, so
\eqref{f16v26:weighted-response} bounds the optimal finite cost by
$\sum_cw_c b_c$, where $b_c=C_r|q'_j-q_j|e^{-a d(c,j)}$.
Lift an optimal label coupling using the actual conditional distributions in
each product atom. On the compact product, extract a weak limit of these
couplings. For any fixed grid, the cost is continuous except on marginal-null
faces, so its integral in the limit is at most $\sum_cw_cb_c$. The grid costs
increase to the full weighted Hamming cost; monotone convergence proves
$\inf_\Pi\sum_cw_c\Pi\{X_c\ne X'_c\}\le\sum_cw_cb_c$.

For completeness this implies one coupling satisfying all component bounds.
The set of vectors dominating the component disagreement probabilities of some
coupling is convex and upward closed. It is closed: compactness of the
couplings and lower semicontinuity of each open-set indicator
$1_{\{X_c\ne X'_c\}}$ preserve each limiting upper constraint. If the vector
$b$ were outside this set, finite-dimensional separation would give a
nonnegative weight $w$ for which every coupling has cost strictly greater than
$\sum_cw_cb_c$, contradicting the preceding transport bound. This proves
\eqref{f16v26:restricted-coupling}. Polynomial graph growth makes
$\sup_j\sum_c e^{-ad(c,j)}$ finite, uniformly in $D$ and the finite ambient
history. Summing proves \eqref{f16v26:restricted-total}. All conditionals in
this argument retain their complete restricted joint law; scalar marginal
couplings were not independently combined.
\end{proof}

\subsection{The full native block: the restriction is paid in the coupling}

We now return to $P_D^q$ on its entire original compact domain. Assume
\eqref{f16v26:collar-window} for $q$ and $q'$, differing only at $j$, and require
\begin{equation}
 Q\le\eta_*\sqrt\gamma,
       \qquad \gamma\ge\max\{1+r^2,\gamma_r^{\rm box}\}.
 \label{f16v26:entry}
\end{equation}
This is a genuine restriction on the deterministic collar. It is not inherited
from V20's all-compact-boundary comparison and must not be omitted. Set
\[
 \varepsilon_D(Q)=
    \min\{1,N_D\exp[A_r(1+Q^2)-c_*\gamma]\}.
\]

\begin{theorem}[Single-input full-output susceptibility on a controlled collar]
\label{f16v26:full-coupling}
There is a coupling $(Y,Y')$ of the \emph{complete} native laws $P_D^q,P_D^{q'}$
for which all cells satisfy
\begin{equation}
 \mathbb P\{Y_c\ne Y'_c\}
 \le C_r|q'_j-q_j|e^{-ad_{\mathcal G}(c,j)}+2\varepsilon_D(Q).
 \label{f16v26:full-component}
\end{equation}
In particular the full expected resampled disagreement has the bound
\begin{equation}
 \boxed{
 I_{D,j}(q,q'):=\mathsf W_{\boldsymbol1}(P_D^q,P_D^{q'})
 \le \min\{N_D,\ C_r|q'_j-q_j|
                    +2N_D^2 e^{A_r(1+Q^2)-c_*\gamma}\}.}
 \label{f16v26:full-output}
\end{equation}
The constants are independent of ambient volume, duration, block shape, and
its number of cells. That number occurs only in the stated exceptional term.
No fast variable inside $D$, including its fringe, has been fixed or omitted.
\end{theorem}
\begin{proof}
Let $G_D$ be the event that every free cell belongs to its comparison ball.
It has positive probability. The exact identity
\[
       P_D^q(\,\cdot\mid G_D)=\pi_D^q
\]
holds with the same original action, crossing words, and product references.
Equation \eqref{f16v26:conditional-tail} gives
$d_{\rm TV}(P_D^q,\pi_D^q)=P_D^q(G_D^c)\le\varepsilon_D(Q)$, and the same
holds at $q'$. Couple each full row to its restricted row by a whole-vector
maximal coupling. Between the two restricted rows use the one coupling in
Proposition~\ref{f16v26:coupling-conversion}. Disintegration glues these three
couplings while preserving every full and restricted marginal. For each cell,
a mismatch of the full outputs requires either a restricted-pair mismatch or
one of the two whole-vector coupling failures. This gives
\eqref{f16v26:full-component}. Summation, the exponential graph-row bound, and
the trivial cost $N_D$ prove \eqref{f16v26:full-output}. On either exceptional
branch the full field remains drawn from its actual conditional distribution.
No independence of failures, changed exterior law, or division by an estimated
good-event probability is used.
\end{proof}

For fixed $Q$ and $N_D\le e^{\beta\gamma}$ with any fixed $0<2\beta<c_*$,
the exceptional total cost is exponentially small, at a threshold depending
on $Q,\beta,r$. More generally if $Q_\gamma=\gamma^\alpha$ with
$0<\alpha<1/2$ and $N_D$ is any prescribed polynomial in $\gamma$, then
\begin{equation}
 I_{D,j}(q,q')\le C_r|q'_j-q_j|+e^{-c_*\gamma/2}
 \label{f16v26:moderate-output}
\end{equation}
eventually, uniformly throughout that collar window. Constants multiplying
the polynomial and its exponent enter the threshold. For a single changed
input $|q'_j-q_j|\le2Q_\gamma$, this is a side-length-independent leading
susceptibility $2C_rQ_\gamma$, not the whole-fringe cost of V24.

\subsection{The admitted-row influence ledger is positive; the excluded rows remain}

For clarity we give the actual finite-geometry count used to compare the new
budget to V24. Choose a fixed integer $\rho\ge1$ bounding every native cell
interaction in the ambient coordinatewise torus/path distance. On a spatial
cycle of length $m>R$ use every translated cyclic interval of $R$ cells; if
$m\le R$ use the whole cycle once. On the time path use V24's indexed clipped
intervals of length $R$. Form product blocks $D$. Give each product index rate
$R^{-(k+1)}$, where $k$ is the number of cut spatial directions. In this
isotropic torus $k$ is zero or three; the notation also makes the counting
transparent. Repeated clipped blocks keep their indices and rates.

\begin{proposition}[A finite cover and a positive moderate-collar table]
\label{f16v26:moderate-ledger}
For $R\ge2\rho$, every cell has coverage rate one, $|D|\le R^4$, and the total
rate of blocks not containing a given cell $j$ whose conditional can depend on
$j$ is at most $C_\rho/R$. Fix $0<\alpha<1/2$. There is a fixed $K_r$ such that
with $Q_\gamma=\gamma^\alpha$ and $R_\gamma=\lceil K_r(1+Q_\gamma)\rceil$,
the single-input table supplied by \eqref{f16v26:full-output}, \emph{only on
pairs of collars within this window}, satisfies
\begin{equation}
 1-\sum_{D:\,j\notin D}r_DI_{D,j}(q,q')\ge\tfrac12
 \label{f16v26:good-ledger}
\end{equation}
at a fixed-$\alpha,r$ large-coupling threshold. Equation
\eqref{f16v26:good-ledger} is not a global native block-generator inequality:
its table is not proved on arbitrary compact exterior rows.
\end{proposition}
\begin{proof}
A site belongs to $R$ indexed intervals in each cut direction, including at
both time endpoints, and to one interval in an uncut cycle. Product rates
therefore give unit coverage. If $j$ is outside a product block but affects it,
some cut coordinate puts it outside its interval at distance at most $\rho$.
There are at most $2\rho$ such interval indices in that coordinate. In each
other cut coordinate there are at most $R+2\rho$ indices whose interval comes
within $\rho$ of $j$; an uncut cycle contributes one. Summing over at most four
coordinates and dividing by the coverage count gives
\[
 \sum_{D:\,j\notin D,\ j\in\partial D}r_D
       \le\frac{8\rho}{R}(1+2\rho/R)^3\le C_\rho/R.
\]
This counts periodic seams and clipped time intervals. It is an upper bound
for the literal word collar, not an assertion that every nearby box interacts.
The size estimate follows since each interval has length at most $R$.

Take $K_r\ge8C_\rho C_r+2\rho+1$, enlarging it if necessary for fixed
constants. Then the leading birth cost from
$|q'_j-q_j|\le2Q_\gamma$ is at most $1/4$. The exceptional cost per block is
at most $2R_\gamma^8\exp[A_r(1+\gamma^{2\alpha})-c_*\gamma]$, which tends to
zero exponentially since $2\alpha<1$. Its aggregate is at most $1/4$
eventually. The entry condition \eqref{f16v26:entry} also holds eventually.
This proves the restricted-row ledger and no more. The all-row hypothesis of
Theorem~\ref{f16v24:general-gap} is still missing, so it is not applied.
\end{proof}

Under the \emph{actual full native law} V16 does give, for each block's collar,
\begin{equation}
 P\{\max_{c\in\partial D}|\zeta_c|>Q_\gamma\}
       \le |\partial D|\exp[t_0M_r-t_0\gamma^{2\alpha}].
 \label{f16v26:bad-collar-prob}
\end{equation}
For any family with disjoint collars, its simultaneous bad-collar probability
is at most the product of the corresponding right sides: expand choices of
one witness per collar and apply V16's joint exponential moment to those
distinct cells. This needs neither independence nor a product posterior.
The statement is not made after arbitrary exterior conditioning.

Neither \eqref{f16v26:bad-collar-prob} nor the positive admitted-row table
bounds the omitted squared innovations for all spectral test functions.
Arbitrary normalized functions may concentrate on the excluded collars, and
field-dependent gating would require an invariance and commutator calculation.
The new coupling closes the tangential whole-fringe loss \emph{inside the
stated boundary class}; it does not erase that class from the hypotheses.

\subsection{Remaining target and exact scope of this continuation}

There are two genuine advances. The congestion test identifies a failure of
one tempting norm in an exactly gapped model. More importantly, the new spatial
coupling controls every output cell at once after one collar cell changes,
first with exponential localization in the restricted native comparison and
then under the full compact conditional with its exception paid. It does not
merely sum unrelated scalar total variations. Ambient volume and history
length are unrestricted throughout; the deterministic collar window and the
finite-block exception are explicit.

The next noncircular target is the remaining \emph{rough-collar contribution}
in a same-law, all-test-function estimate, or a boundary repair/exploration
that extends the single-input coupling to those rows. Increasing an expansion
order does not supply it. Replacing the true conditional rows by their
small-field restrictions, inserting their rare-event probability independently
of a squared innovation, or applying V24's all-row theorem to only the good
rows would all skip a required step. The coupled outputs do not define a new
physical transfer or a new field law.

V25's all-finite-width baseline and V23's direct covariance/memory results
retain their original scopes. No new unrestricted covariance operator bound,
uniform original-clock sampler gap, or improvement of physical spectral ratios
is asserted. The prescribed bridge window, physical Gauss reduction, infrared
contact coercivity, same-top transfer comparison, physical-time calibration,
and interacting continuum construction remain separate obligations.

\subsection{Diagnostics, preservation, and verification boundary}

The accompanying diagnostic is
\begin{center}\small\path{verify_ym_f16_v26_spatial_coupling.py}.\end{center}
It tests the exact marginal formula for congestion; the gapped product-pair
obstruction; normalized conditional derivative and sweep bookkeeping; finite
transport duality and simultaneous disagreement constraints; the full/restricted
coupling mixture; and the clipped finite-cover counts. Analytic inequalities
in the preceding proofs are not certified by the finite fixtures. In
particular no native compact-integral enclosure, numerical value of the
one-form constants, or proof-assistant verification is claimed. The precise
executed test families and available predecessor reruns are recorded in the
machine-readable reports, rather than used as a measure of proof completeness.

The cumulative source changes only V25's title version and inserts this exact
appendix before its final document-ending command. Removing that insertion and
restoring the title recovers V25 byte for byte. The package records preservation,
actual diagnostic outputs, compilation, and rendered-page inspection. These
new deductions are conditional on their identified inherited analytic inputs;
the entire historical lineage has not been independently recertified. The
next scheduled companion checkpoint remains V30.

\begin{firewall}
V26 proves a full-output single-input spatial coupling on a declared
rescaled-collar window, with a localized restricted-native part and an explicitly
paid full-native exception. It proves a positive influence ledger only for
those admitted boundary pairs, not a global native gap. It also shows that
unweighted splice congestion can grow exponentially in a uniformly gapped
product model. No arbitrary-compact-boundary spatial coupling, control of the
excluded innovation-weighted defect, unrestricted coherent covariance theorem,
physical transfer gap, or interacting continuum mass gap is established.
All original laws, normalizers, source definitions and clocks are retained.
\end{firewall}
% END F16 V26 APPEND

% BEGIN F16 V27 APPEND -- 2026-09-17
% Insert immediately before the final document-ending command in V26.
% No inherited statement, normalizer, source, or update clock is changed.

\clearpage
\section{V27: buffered rough-collar spectral localization}
\label{f16v27:context}

V26 controls a whole conditional output after one collar input changes, but
only on a declared collar window. The probability of its excluded collars is
small under the native law; an arbitrary squared spectral test can nevertheless
concentrate there. We do not multiply that probability by an unrelated
innovation. Instead, we use a buffered conditional expectation to separate an
arbitrary test into an exterior-measurable part and a part paid by an actual
resampling energy. Uniform conditional preparation controls the first part.
A local, interface-priced comparison pays the second in the original clock.

This proves an all-test-function rough-collar estimate and an operator bound
on its intersection with every low-energy spectral subspace. A family of
separated repair regions also gives a simultaneous rough-count estimate with
no factor counting the regions in its energy term. The target remains the
complete native law. These are localization and killed-form statements, not
a global gap theorem. The original-clock comparison is expensive, and an
extensive term still prevents the final rough-innovation bound from closing
V26's unrestricted spatial argument.

\subsection{Dependencies and the distinction from an unweighted tail}

The complete V26 appendix and audit were read. The new proof uses V20's
all-compact-boundary preparation, V24's exact moving/unmoving support inclusion,
and V25's finite-product interface comparison. It does not use V26's restricted
one-form response theorem to assume that rough collars are already repaired.
It needs neither a global native gap nor the arbitrary-order source expansions.
Every conclusion is a deduction from the expressly identified inherited inputs,
not an independent recertification of their full analytic lineage.

Two semantic archive searches were empty. Exact mounted-text searches and
identified source ranges were then checked. F14 V1 already discusses
original-clock killed repair on specified visible events; its negative-curvature
geometry is not transplanted here. The Grand Tensor monograph already separates
continuous killing from endpoint sampling. Those principles are credited. The
new event family consists of arbitrary rough cells and their collars in the
present fixed-bridge history, with preparation uniform over every exterior of
the repair region. The energy price and the number of repair regions are
controlled separately. The next scheduled checkpoint remains V30.

For context, Goel--Montenegro--Tetali,
\href{https://arxiv.org/abs/math/0505690}{\emph{Mixing Time Bounds via the Spectral
Profile}}, Definition 1.4 and (1.4), distinguish a restricted Dirichlet eigenvalue
from a global gap. The definitions and adjacent argument were examined; their
finite-state mixing-time theorem is not invoked on this continuous law.
Spinka's \href{https://arxiv.org/abs/1803.10578}{\emph{Finitary codings for
spatial mixing Markov random fields}}, Theorem 1.2, requires common overlap
along unbounded spatial scales. Its finite-valued coding result is not supplied
by V23's fixed-coupling finite-buffer bound. The present step estimates spectral
weights instead of asserting that this missing coupling composition holds.

\subsection{One conditional projection and a genuinely weighted rare event}

On an arbitrary probability space let $E$ be conditional expectation onto a
sigma algebra $\mathcal F$, $Q=I-E$, and let $M_H$ multiply by the indicator
of an event $H$. Suppose
\begin{equation}
                 P(H\mid\mathcal F)\le p\quad\hbox{a.s.},\qquad 0\le p\le1.
 \label{f16v27:conditional-smallness}
\end{equation}
The event need not be independent of $\mathcal F$.

\begin{proposition}[Conditional smallness with its energy debit]
\label{f16v27:projection}
For every real or complex $f\in L^2(P)$,
\begin{equation}
 \|M_HE\|\le\sqrt p,\qquad
 \|M_Hf\|_2\le\|Qf\|_2+\sqrt p\,\|Ef\|_2
                  \le\|Qf\|_2+\sqrt p\,\|f\|_2.
 \label{f16v27:projection-bound}
\end{equation}
On functions supported on $H$ there is the sharper form estimate
\begin{equation}
       \|Qf\|_2^2\ge(1-p)\|f\|_2^2,\qquad f=M_Hf.
 \label{f16v27:killed-projection}
\end{equation}
These statements neither replace a probability by an energy nor assume an
inequality for every mean-zero function.
\end{proposition}
\begin{proof}
For $g=Eg$, conditional expectation gives
$\|M_Hg\|_2^2=\mathbb E[P(H\mid\mathcal F)|g|^2]\le p\|g\|_2^2$.
Decompose $f=Ef+Qf$, use this estimate on the first component and the
contraction $\|M_HQf\|_2\le\|Qf\|_2$ on the second. This proves
\eqref{f16v27:projection-bound}. Taking adjoints gives $\|EM_H\|\le\sqrt p$.
For $f=M_Hf$, orthogonality of $E,Q$ consequently gives
$\|Qf\|_2^2=\|f\|_2^2-\|Ef\|_2^2\ge(1-p)\|f\|_2^2$.
No differentiability, finite alphabet, or lower bound on a point mass is used.
\end{proof}

For example, if $P(H)=p>0$ and $f=p^{-1/2}1_H$, then
$\|f\|_2=\|1_Hf\|_2=1$. Its unweighted event probability alone says nothing
small about that weighted mass. When \eqref{f16v27:conditional-smallness}
holds, \eqref{f16v27:killed-projection} instead forces an actual conditional
innovation of size at least $1-p$. Constants are also consistent with
\eqref{f16v27:projection-bound}: their energy is zero and their event mass is
paid by the probability term.

\subsection{Prepare each rough-cell event inside a local repair region}

Retain exactly $P=P_{\gamma,\mathbf y}$ on the complete unmoving charts, with
$B=m^3$, $m\ge4$, $n\ge1$, $\max_{i,e}|y_{i,e}|\le r$, and
$\gamma\ge1+r^2$. All norms below refer to $L^2(P)$. The original balanced
generator, transported unitarily to this space, is
\[
 \widehat{\mathcal L}=\sum_\alpha\widehat Q_\alpha,\qquad
 \widehat{\mathcal E}(f)=\langle f,\widehat{\mathcal L}f\rangle.
\]
Every original five-chord cube or single nonroot port still has rate one.
For a set $S$ of unmoving cells let $E_S$ condition on its complete exterior,
$Q_S=I-E_S$, and $\mathcal E_S(f)=\|Q_Sf\|_2^2$.

Let $e_c=|\zeta_c|^2$ and define the same type of rough-cell event as V26 by
\[
       H_c(Q)=\{e_c>Q^2\},\qquad Q\ge0.
\]
Fix a positive integer $r_0$ bounding the coordinatewise space--time displacement
of every edge of the native interaction graph $\mathcal G$. Denote V20's
preparation constants by $\kappa_r,B_r,M_r$. Choose
\begin{equation}
 b_\gamma=2+\left\lceil\kappa_r^{-1}\log\gamma\right\rceil,\qquad
 w_\gamma=r_0(b_\gamma+2),\qquad h_\gamma=2w_\gamma+5.
 \label{f16v27:repair-size}
\end{equation}
Let $S_c$ be the product of the coordinatewise spatial-cycle balls and the
clipped time interval of radius $w_\gamma$ around $c$. A short spatial cycle
is taken in full, once. Natural time endpoints are retained, not frozen.
Then $c$ is at graph distance at least $b_\gamma$ from $S_c^c$, unless its
component is entirely free, in which case the distance is infinite.

\begin{proposition}[A roughness bound uniform in the actual repair exterior]
\label{f16v27:uniform-tail}
With $\overline M_r=M_r+B_r$ and $t_0=1/(24\pi^2)$, put
\begin{equation}
       p_\gamma(Q)=\min\{1,\exp[t_0(\overline M_r-Q^2)]\}.
 \label{f16v27:p}
\end{equation}
The exact conditional-density versions satisfy, for every compact exterior $q$,
\begin{equation}
 P_{S_c}^{q}(H_c(Q))\le p_\gamma(Q),\qquad
 \mathbb E_{P_{S_c}^{q}}e_c^k
       \le M_{k,r}:=k!t_0^{-k}e^{t_0\overline M_r}\quad(k\ge1).
 \label{f16v27:uniform-preparation}
\end{equation}
There is no small-field restriction on $q$ or on the free repair output.
The radius $h_\gamma=O_r(1+\log\gamma)$ is independent of $Q,B,n$.
\end{proposition}
\begin{proof}
V20's exact conditional moment-generating estimate, with a source only at
$c$, bounds $\mathbb E_{P_{S_c}^{q}}e^{t_0 e_c}$ by
\[
       \exp\{t_0(M_r+B_r\gamma e^{-\kappa_r b_\gamma})\}
                         \le e^{t_0\overline M_r}.
\]
A graph step changes each ambient coordinate by at most $r_0$, so the repair
box has the claimed depth, including periodic identifications and clipped
natural ends. Exponential Markov and $x^k\le k!t_0^{-k}e^{t_0x}$ give the
two conclusions. The result concerns the complete conditional and includes
its exact normalizer, Haar powers, and every crossing term. Its uniformity in
$q$ is inherited from preparation, not from an averaged bad-event estimate.
\end{proof}

For $Q_\gamma=\gamma^\alpha$, $0<\alpha<1/2$ fixed, one has
$p_\gamma(Q_\gamma)\le e^{-(t_0/2)\gamma^{2\alpha}}$ eventually.
This is V26's growing rescaled threshold, but the exterior of $S_c$ is now
arbitrary: the observed rough cell sits inside an integrated repair buffer.
No all-row susceptibility estimate for the original block $D$ has thereby
been inferred.

\subsection{The repair has a local original-clock cost}

Let $\widetilde S_c$ denote the elementary balanced support needed for the
unmoving update $S_c$. It contains each chord cube in $S_c$ and the moving
ports at times $t$ for which that cube's cell at $t$ or $t+1$ is in $S_c$.
Thus its preceding ports are included exactly as in V24. All these statements
are about support inclusion, not equality of balanced and unmoving kernels.

\begin{proposition}[Local interface price and overlap, independent of ambient size]
\label{f16v27:clock}
There are constants independent of $B,n,\gamma,c$ such that, with
$h=h_\gamma$,
\begin{equation}
 \begin{split}
 |\widetilde S_c|&\le C h^4,\\
 \mathcal E_{S_c}(f)&\le K_0(\gamma)
            \sum_{\alpha\in\widetilde S_c}\|\widehat Q_\alpha f\|_2^2,
       \qquad K_0(\gamma)=C_r h^4e^{C_r\gamma h^3},\\
 \sup_\alpha\#\{c:\alpha\in\widetilde S_c\}&\le\varkappa_\gamma:=C h^4.
 \end{split}
 \label{f16v27:local-clock}
\end{equation}
Define $K_1(\gamma)=\varkappa_\gamma K_0(\gamma)$, enlarging fixed constants
once. For any nonnegative deterministic cell weights $a_c$,
\begin{equation}
 \sum_c a_c\mathcal E_{S_c}(f)
 \le K_0(\gamma)\sum_\alpha
          \left(\sum_{c:\alpha\in\widetilde S_c}a_c\right)
                          \|\widehat Q_\alpha f\|_2^2
 \le K_1(\gamma)\|a\|_\infty\widehat{\mathcal E}(f).
 \label{f16v27:weighted-clock}
\end{equation}
In particular $K_1(\gamma)\le C_rh^8e^{C_r\gamma h^3}$. This large
coefficient is retained; the original update rates are not accelerated.
\end{proposition}
\begin{proof}
The exact frame relation gives
$\sigma(\omega_{\widetilde S_c^c})\subseteq\sigma(\zeta_{S_c^c})$.
Projection monotonicity first bounds $\mathcal E_{S_c}$ by the conditional
variance on $\widetilde S_c$. Apply V25's continuous-product interface theorem
to that balanced conditional, for every value of its exterior. Its original
product reference retains the fractional endpoint Haar powers and the one-cube
endpoint factors.

Order free elementary groups by time. The spatial cross-section of this local
box has at most $C h^3$ cells. Only a fixed number of active time layers, including
a partially processed layer, can meet an ordering cut. Bounded native word
incidence and oscillation give cut cost $C_r\gamma h^3$ and incident cost
$C_r\gamma$. The path prefactor is $|\widetilde S_c|\le C h^4$.
A spatial cycle fully contained in a short box has no more than $h$ distinct
coordinates in that direction; counting a complete layer includes its periodic
seam. Missing time groups and preceding-port enlargement respect the same
bounds. V25 therefore gives the second inequality of
\eqref{f16v27:local-clock} after integrating the actual balanced exterior.

For a fixed chord group, centers whose $S_c$ contain it are in a coordinatewise
box of size at most $C h^4$. For a fixed moving port, either of its two adjacent
chord times may cause membership; the union has the same bound with a fixed
factor. Periodic saturation reduces the number of distinct centers. This proves
the overlap bound. Sum the nonnegative form estimates with weights $a_c$ to
obtain \eqref{f16v27:weighted-clock}. No factor counting all repair centers
appears in the energy coefficient.
\end{proof}

\subsection{An all-test-function rough-collar inequality and a low-energy projector}

Set $p=p_\gamma(Q)$. For a finite cell set $C$, including any actual collar,
write $H_C(Q)=\bigcup_{c\in C}H_c(Q)$. No separation or diameter restriction is
imposed on $C$. Let
\[
              \Pi_{\le E}=1_{[0,E]}(\widehat{\mathcal L}),\qquad E\ge0.
\]
This is a spectral projection of the original auxiliary sampler, and includes
its constant mode. It is not a physical energy projection.

\begin{theorem}[Rough-collar spectral weights with a paid energy term]
\label{f16v27:weighted-rough}
For every $f\in L^2(P)$ and every nonnegative deterministic weight array $a$,
\begin{equation}
 \boxed{
 \sum_c a_c\|1_{H_c(Q)}f\|_2^2
 \le\left[
       \sqrt{K_1(\gamma)\|a\|_\infty\widehat{\mathcal E}(f)}
       +\sqrt{p\sum_ca_c}\,\|f\|_2
       \right]^2.}
 \label{f16v27:weighted-rough-bound}
\end{equation}
Consequently, at unrestricted spatial volume and duration,
\begin{equation}
 \boxed{
 \|1_{H_C(Q)}\Pi_{\le E}\|_{2\to2}
 \le\min\{1,\sqrt{K_1(\gamma)E}+\sqrt{|C|p_\gamma(Q)}\}.}
 \label{f16v27:spectral-collar}
\end{equation}
The energy coefficient does not depend on $|C|$ or its diameter. Its explicit
coupling dependence is part of the assertion. No native gap is assumed in
defining or estimating $\Pi_{\le E}$.
\end{theorem}
\begin{proof}
For each $c$, Proposition~\ref{f16v27:uniform-tail} gives
$P(H_c\mid\zeta_{S_c^c})\le p$. Decompose
$1_{H_c}f=1_{H_c}Q_{S_c}f+1_{H_c}E_{S_c}f$.
Minkowski in the finite direct sum with weights $a_c$, followed by
Proposition~\ref{f16v27:projection}, gives
\[
 \left(\sum_ca_c\|1_{H_c}f\|_2^2\right)^{1/2}
 \le\left(\sum_ca_c\mathcal E_{S_c}(f)\right)^{1/2}
       +\left(p\sum_ca_c\|E_{S_c}f\|_2^2\right)^{1/2}.
\]
Use \eqref{f16v27:weighted-clock} and contraction of each conditional
expectation. This proves the first assertion for arbitrary complex tests.
For $a=1_C$, $1_{H_C}\le\sum_{c\in C}1_{H_c}$, so the same bound controls
$\|1_{H_C}f\|_2$. On the range of $\Pi_{\le E}$, spectral calculus gives
$\widehat{\mathcal E}(f)\le E\|f\|_2^2$. Taking the operator norm proves
\eqref{f16v27:spectral-collar}. All inequalities precede any spectral
restriction; rare-event weights are not excluded from the test class.
\end{proof}

For fixed $0<\alpha<1/2$ and $|C|\le A_0\gamma^s$, with $s,A_0$ fixed,
\begin{equation}
 \|1_{H_C(\gamma^\alpha)}\Pi_{\le E}\|
       \le \sqrt{K_1(\gamma)E}+C_{A_0,s,r}e^{-c_r\gamma^{2\alpha}}
 \label{f16v27:polynomial-collar}
\end{equation}
eventually. In particular an arbitrary normalized spectral vector with
$K_1(\gamma)E\to0$ cannot put order-one mass on that rough collar. This is
an operator statement on the complete low-energy band, not a statement about
a chosen finite list of moments. The small-energy scale here can be much
smaller than an algebraic power of $\gamma^{-1}$.

\begin{proposition}[A localized killed floor, not a global gap]
\label{f16v27:killed}
For $|C|p<1$ and every $f$ supported on $H_C(Q)$,
\begin{equation}
 \widehat{\mathcal E}(f)\ge
      \frac{(1-\sqrt{|C|p})^2}{K_1(\gamma)}\|f\|_2^2.
 \label{f16v27:collar-killed}
\end{equation}
This is a lower bound for the compressed operator on $L^2(H_C,P)$.
For the polynomial collars above, it is eventually at least
$1/(4K_1(\gamma))$.
\end{proposition}
\begin{proof}
Set $f=1_{H_C}f$ in \eqref{f16v27:weighted-rough-bound} with $a=1_C$,
take square roots, and rearrange. At a finite regulator the original
generator is bounded and positive, so its compression to the indicated closed
subspace has exactly the displayed restricted quadratic form. The final
assertion uses $|C|p\le1/4$. This compression describes continuous killing
on leaving the event, not observations of the unmodified process at separated
times with intervening returns allowed.
\end{proof}

The explicit original-clock lower bound supplied by this comparison has scale
$C_r^{-1}h_\gamma^{-8}e^{-C_r\gamma h_\gamma^3}$ in the present estimate.
It is positive independently of ambient size at fixed coupling, but does not
stay positive as coupling grows. An unrestricted growing union may have
$|C|p\ge1$; the theorem does not then assert a killed floor.

\subsection{Separated repair regions: a joint rough-count estimate without a count debit}

Choose a finite set $\mathcal J$ of centers such that no native interaction
meets two of their $S_c$ and their enlarged balanced supports $\widetilde S_c$
are disjoint. A sufficient condition is a fixed multiple of $h_\gamma$
separation in an actual finite-geometry packing. Set $U=\bigcup_{c\in\mathcal J}S_c$
and $M=|\mathcal J|$. This is one common exterior. Different centers need not
have the same conditional distribution.

For an increasing set $\mathcal A\subseteq\{0,1\}^{\mathcal J}$ define
\[
 H_{\mathcal A}=\{(1_{H_c(Q)})_{c\in\mathcal J}\in\mathcal A\},\qquad
 b_{\mathcal A}=\operatorname{Ber}(p)^{\otimes M}(\mathcal A).
\]
Here increasing means coordinatewise increasing. The Bernoulli product is a
comparison of indicators, not a substitute for the native field law.

\begin{theorem}[Simultaneous rough-event localization in the original clock]
\label{f16v27:joint}
For all $f\in L^2(P)$,
\begin{equation}
 \boxed{
 \|1_{H_{\mathcal A}}f\|_2
 \le\sqrt{K_0(\gamma)\widehat{\mathcal E}(f)}
                         +\sqrt{b_{\mathcal A}}\,\|f\|_2.}
 \label{f16v27:joint-weight}
\end{equation}
The coefficient $K_0(\gamma)$ is independent of $M$. Moreover,
\begin{equation}
 \|1_{H_{\mathcal A}}\Pi_{\le E}\|
       \le\min\{1,\sqrt{K_0(\gamma)E}+\sqrt{b_{\mathcal A}}\},\qquad
 \widehat{\mathcal E}(f)\ge
       \frac{1-b_{\mathcal A}}{K_0(\gamma)}\|f\|_2^2
       \quad(f=1_{H_{\mathcal A}}f).
 \label{f16v27:joint-spectral}
\end{equation}
In particular all $M$ selected cells being rough costs $b_{\mathcal A}=p^M$.
For $0<p<\theta<1$, at least $\theta M$ of them being rough obeys
\begin{equation}
 b_{\mathcal A}\le e^{-M D_{\rm B}(\theta\Vert p)},\qquad
 D_{\rm B}(\theta\Vert p)=
   \theta\log\frac\theta p+(1-\theta)\log\frac{1-\theta}{1-p}.
 \label{f16v27:chernoff}
\end{equation}
No independence of the actual exterior posterior is asserted.
\end{theorem}
\begin{proof}
At fixed $\zeta_{U^c}$, the exact product reference and the absence of a word
meeting two free components factor the conditional into the original
$P_{S_c}^{q_c}$. Its normalizer is the product of their normalizers only
\emph{at this fixed exterior}. The indicators are consequently conditionally
independent Bernoulli variables with success parameters at most $p$ by
Proposition~\ref{f16v27:uniform-tail}. Coupling each with a Bernoulli($p$)
using independent uniforms proves
$P(H_{\mathcal A}\mid\zeta_{U^c})\le b_{\mathcal A}$.

Product conditional variance tensorization at the same exterior gives
\[
 \|Q_Uf\|_2^2\le\sum_{c\in\mathcal J}\mathcal E_{S_c}(f)
       \le K_0(\gamma)\widehat{\mathcal E}(f).
\]
The last inequality uses disjoint enlarged balanced supports; there is no
factor $M$ and no triangle inequality summing $M$ norms. Apply
Proposition~\ref{f16v27:projection} with $E=E_U$. Its supported-function
version gives the stronger second inequality of
\eqref{f16v27:joint-spectral}. Spectral calculus gives the first. For the
count bound, conditional multiplication yields
$\mathbb E e^{t\sum1_{H_c}}\le(1-p+pe^t)^M$ under each exterior-tilted
comparison. Exponential Markov and minimizing over $t>0$ gives
\eqref{f16v27:chernoff}. Only this auxiliary product bound is used; averaging
over a correlated exterior does not make the actual fields independent.
\end{proof}

\subsection{What a low-energy change of density can and cannot do}

The last theorem has an exact change-of-law interpretation. Let $\|f\|_2=1$,
put $g=E_Uf$, and let $e=\|Q_Uf\|_2^2$. If $e<1$, define
\[
         d\nu_f=|f|^2dP,\qquad d\widetilde\nu_f=\frac{|g|^2}{1-e}\,dP.
\]
The second density depends only on the one common exterior.

\begin{proposition}[Exterior reweighting approximation for spectral tests]
\label{f16v27:tilt}
The measures just defined obey
\begin{equation}
        d_{\rm TV}(\nu_f,\widetilde\nu_f)\le\sqrt e
                 \le\sqrt{K_0(\gamma)\widehat{\mathcal E}(f)}.
 \label{f16v27:tilt-TV}
\end{equation}
Conditionals inside $U$ under $\widetilde\nu_f$ are exactly the original
native conditionals. Thus their rough indicators retain the comparison in
Theorem~\ref{f16v27:joint}. The corresponding claim for $\nu_f$ pays the
explicit error above; it does not hold by invariance alone.
\end{proposition}
\begin{proof}
Let $v=g/\sqrt{1-e}$. Orthogonality gives
$\|v\|_2=1$ and $\langle f,v\rangle=\sqrt{1-e}$, real and nonnegative.
Cauchy--Schwarz applied to
$||f|^2-|v|^2|\le|f-v||f+v|$ yields
\[
 \tfrac12\int\bigl||f|^2-|v|^2\bigr|\,dP
 \le\tfrac12\|f-v\|_2\|f+v\|_2=\sqrt e.
\]
Exterior measurability of $g$ leaves the conditional laws inside $U$ unchanged
on their nonzero exterior weights. The energy bound was proved above.
\end{proof}

A local moment consequence is also available. On the complete charts
$e_c\le E_*\gamma$, $E_*=12\pi^2$. For any normalized $f$ with
$\widehat{\mathcal E}(f)\le E$, decomposition with its single repair projection
gives
\begin{equation}
 \left(\mathbb E_P[e_c^k|f|^2]\right)^{1/2}
 \le \sqrt{M_{k,r}}+(E_*\gamma)^{k/2}\sqrt{K_0(\gamma)E},\qquad k\ge1.
 \label{f16v27:spectral-moments}
\end{equation}
Indeed the exterior-measurable part uses the uniform conditional moment in
\eqref{f16v27:uniform-preparation}, whereas the orthogonal part uses the
compact cap and its paid energy. These are moments under the test-weighted
law, not V16's unweighted moments substituted for them.

\subsection{The rough innovation form is bounded, but not yet absorbed}

To display the remaining extensive loss, let $D$ run over any deterministic
finite family of unmoving update blocks with rates $r_D\ge0$. Put
$C_D=\partial_{\mathcal G}D$ and
\[
 \mathfrak D_Q(f)=\sum_Dr_D\|1_{H_{C_D}(Q)}Q_Df\|_2^2,
       \qquad a_c=\sum_{D:\,c\in C_D}r_D.
\]
This is a specified rough-collar innovation form. It is not asserted equal
to every error in a future path-coupling or noncommuting projection argument.

\begin{proposition}[A same-law all-test-function bound on the rough innovations]
\label{f16v27:rough-form}
For every $f\in L^2(P)$,
\begin{equation}
 \boxed{
 \mathfrak D_Q(f)\le
 \left[
 \sqrt{K_1(\gamma)\|a\|_\infty\widehat{\mathcal E}(f)}
     +\sqrt{p\sum_ca_c\,\operatorname{Var}_P f}
 \right]^2.}
 \label{f16v27:rough-form-bound}
\end{equation}
For V26's unit-coverage product boxes of side $R$, its actual collar count
therefore gives
\begin{equation}
 \mathfrak D_Q(f)\le\frac{C_\rho}{R}
     \left[\sqrt{K_1(\gamma)\widehat{\mathcal E}(f)}
          +\sqrt{p\mathcal N\operatorname{Var}_P f}\right]^2,
          \qquad\mathcal N=B(n+1).
 \label{f16v27:extensive-loss}
\end{equation}
No bad-event probability is multiplied independently of an innovation.
Nevertheless this bound does not supply an unrestricted-volume absorption.
\end{proposition}
\begin{proof}
The event $H_{C_D}$ is exterior to $D$. Its multiplication operator commutes
with $E_D$ and $Q_D$. Writing $f_0=f-Pf$, contraction gives
\[
 \|1_{H_{C_D}}Q_Df\|_2
   =\|Q_D(1_{H_{C_D}}f_0)\|_2\le\|1_{H_{C_D}}f_0\|_2.
\]
Use $1_{H_{C_D}}\le\sum_{c\in C_D}1_{H_c}$, sum the deterministic rates,
and apply \eqref{f16v27:weighted-rough-bound} to $f_0$. The energy is unchanged
and $\|f_0\|_2^2=\operatorname{Var}_Pf$. This proves
\eqref{f16v27:rough-form-bound}, including vanishing on constants. V26's
finite-cover proof bounds $a_c\le C_\rho/R$, hence
$\sum_ca_c\le C_\rho\mathcal N/R$, with periodic seams and clipped ends
already counted. Substitution proves \eqref{f16v27:extensive-loss}.
\end{proof}

The remaining two costs are now explicit. First, $K_1(\gamma)$ is a pessimistic
interface comparison, not a coupling-uniform constant. Second, at fixed
coupling the term $p\mathcal N$ cannot be made small uniformly in the ambient
history. The separated-region theorem avoids a count factor for one specified
joint event; it does not turn every large union into a rare event. These facts
prevent application of V24's all-row contraction theorem to V26's good rows
alone. The proof supplies a real weighted estimate where there was previously
only a probability estimate, but not the missing global absorption.

\subsection{Scope, next target, and verification}

V27 excludes concentration on declared rough collars from sufficiently low
original-sampler spectral bands, with an explicit local energy price. It also
controls arbitrary increasing rough events on separated repair regions without
charging their number to that price. The repair exteriors are completely
unrestricted, and no field-dependent update schedule is introduced. All
conditional probabilities, product references, compact outputs, and original
source definitions remain unchanged. The supplied import notes' distinction
between a probability/power gain and a locality or functional-inequality cost
is implemented in the two different terms of
\eqref{f16v27:rough-form-bound}.

The next useful step must reduce or absorb those terms in the actual global
operator, or extend spatial boundary propagation through the remaining rough
rows. Another unweighted rough-collar tail would not do that. The killed
operator $1_H\widehat{\mathcal L}1_H$ excludes paths as soon as they leave $H$;
$1_He^{-t\widehat{\mathcal L}}1_H$ permits returns and is a different object.
No improvement of the global gap, V23's covariance norms, or the physical
transfer follows simply from the killed floor. In particular a low-energy
projection here is not a Yang--Mills physical Hilbert-space projection.

The accompanying diagnostic is
\begin{center}\small\path{verify_ym_f16_v27_rough_spectral.py}.\end{center}
It checks finite conditional projections and weighted inequalities, original
clock and overlap counts, conditional product domination with correlated
exteriors, spectral-band and density-change formulas, and a distinction between
killed and returning dynamics. Its counts record exact finite fixtures, not
native integral enclosures or formal certification of analytic constants.
The actual outputs, predecessor diagnostic reruns, hashes, reconstruction,
and build/render records are packaged separately.

The cumulative V26 source is preserved except for its one title version and
this terminal insertion. Removing the insertion and reversing that title change
recovers V26 byte for byte. The analytic deductions rely on the identified
preparation and bounded-word comparison inputs; the full archive has not been
independently recertified. The next scheduled companion checkpoint is V30.

\begin{firewall}
V27 proves an all-test-function rough-collar estimate, localized original-clock
killed bounds, and low-energy spectral-projection estimates. It proves a
simultaneous rough-count bound on separated repair regions and an exact
exterior-reweighting approximation with its energy error. The original-clock
coefficient is explicit and can grow rapidly with coupling; the global rough
innovation bound retains an extensive remainder. No all-boundary spatial
susceptibility, unrestricted global gap, new unrestricted covariance theorem,
physical Gauss coverage, bridge-window removal, infrared contact coercivity,
same-top physical transfer comparison, physical-time calibration, or interacting
continuum mass gap is inferred.
\end{firewall}
% END F16 V27 APPEND
% BEGIN F16 V28 APPEND -- 2026-09-18
% Insert immediately before the final document-ending command in V27.
% No predecessor proof, native probability, source, or elementary clock is edited.

\clearpage
\section{V28: nonextensive rough-innovation absorption by local repair dynamics}
\label{f16v28:context}

V27's final rough-innovation estimate contains $p\mathcal N\operatorname{Var}f$.
That loss occurs after replacing each block innovation by the whole centered
function. We retain the innovation instead. Before applying conditional
smallness, subtract the expectation outside the \emph{union} of the updated
block and its repair region. That component is annihilated by the block
innovation. The probability term then multiplies a local excess conditional
energy, not the full variance. No commutation of overlapping conditional
expectations is needed.

The resulting estimate has no ambient-volume remainder. With two explicitly
normalized auxiliary update families it gives a vanishing relative form bound
for the \emph{same} rough-collar innovation form defined in V27. In particular,
rough-collar updates may be suppressed in that augmented dynamics at a small
relative form cost, while the repair and enlarged-block updates remain active.
The probability law is unchanged, and all output fields remain integrated.
This is not a proof of a gap for either augmented dynamics. Returning their
energies to the original elementary clock still has a large local cost, which
is stated separately.

\subsection{Dependencies and the upstream cancellation}

The supplied V27 appendix and audit were read completely. The new main
inequality uses its conditional preparation theorem
\ref{f16v27:uniform-tail}, and the exact finite-cover geometry of V26. It does
not use V26's restricted-collar coupling to assume that rough rows contract.
V25's interface comparison is used only when returning the new auxiliary
forms to the original balanced clock. V17--V23's covariance expansions and a
native sampler gap are not inputs to the absorption argument.

Two scoped semantic searches returned no matches. Mounted-text searches and
identified passages were then checked. The Grand Tensor monograph already
uses nested conditional-expectation Pythagoras in its score-cutoff loss formula
(NL.12)--(NL.13), and its nested-source identity (NS.6). That general principle
is credited, not claimed new. Its use here subtracts a particular common
exterior inside an \emph{event-weighted block innovation}. V27's estimate
\eqref{f16v27:extensive-loss} remains valid; the present estimate is a different,
stronger interface, not a silent correction of the predecessor.

Reversible Dirichlet-form comparison and monotone-system censoring are distinct
methods. Peres--Winkler,
\href{https://arxiv.org/abs/1112.0603}{\emph{Can extra updates delay mixing?}},
addresses monotone spin systems and specified initial laws. Its primary abstract
and scope were checked; no theorem from its full proof is used here. We prove
reversibility and the form inequalities for the present exterior-gated updates
directly. No order on compact group fields, stochastic domination of trajectories,
or mixing-time censoring inequality is assumed. The review is targeted, not
independent certification of the analytic archive. The next checkpoint remains
V30. The supplied import memoranda's distinction between probability gain,
locality cost, and a subsequent global spectral theorem is retained.

\subsection{A union projection removes the irrelevant exterior component}

On the unchanged finite native probability space, for a set $A$ of unmoving
cells write
\[
 E_Af=\mathbb E_P[f\mid\zeta_{A^c}],\qquad Q_A=I-E_A,
 \qquad \mathcal E_A(f)=\|Q_Af\|_2^2.
\]
All conditional expectations are orthogonal projections on the same $L^2(P)$.
The next proposition holds on any probability space with these sigma-algebra
inclusions; it does not require product independence or a positive-density
comparison.

\begin{proposition}[Nested localization of an actual innovation]
\label{f16v28:nested}
Let $D,S$ be cell sets, let $U=D\cup S$, and let $H$ be measurable outside
$D$. Suppose $P(H\mid\zeta_{S^c})\le p$ almost surely, with $0\le p\le1$.
Then for every real or complex $f\in L^2(P)$,
\begin{equation}
 \boxed{
 \|1_HQ_Df\|_2
 \le \sqrt{\mathcal E_S(f)}+
       \sqrt{p\{\mathcal E_U(f)-\mathcal E_S(f)\}}.}
 \label{f16v28:nested-bound}
\end{equation}
The quantity in braces is nonnegative and equals
$\|(E_S-E_U)f\|_2^2$. In particular, for every $t>0$ the following is a
quadratic-form inequality:
\begin{equation}
 1_HQ_D\preceq (1+t)Q_S+(1+t^{-1})p(E_S-E_U).
 \label{f16v28:nested-Young}
\end{equation}
No identity $E_SE_D=E_DE_S$ is required.
\end{proposition}
\begin{proof}
Since $U$ contains $S$ and $D$, $E_SE_U=E_U=E_DE_U$ and their reversed
products have the same values. Multiplication $M=1_H$ commutes with $E_D$,
because $H$ is exterior to $D$. Decompose
\[
 f=Q_Sf+(E_S-E_U)f+E_Uf.
\]
The last component disappears from $MQ_Df$. Using the commutation only with
$M$ gives the exact two-term identity
\begin{equation}
 MQ_Df=Q_DM Q_Sf+Q_DM(E_S-E_U)f.
 \label{f16v28:two-term}
\end{equation}
The first term has norm at most $\|Q_Sf\|_2$. The second intermediate vector
$(E_S-E_U)f$ is exterior-measurable for $S$. Conditional smallness therefore
bounds its image under $M$ by $\sqrt p\|(E_S-E_U)f\|_2$. The contraction
$Q_D$ does not increase that norm. This proves \eqref{f16v28:nested-bound}.
Nested orthogonal projections give
\[
 \|(E_S-E_U)f\|_2^2
 =\|E_Sf\|_2^2-\|E_Uf\|_2^2
 =\mathcal E_U(f)-\mathcal E_S(f).
\]
Finally $MQ_D$ is itself a positive orthogonal projection, since its two
factors commute. Its quadratic form is $\|MQ_Df\|_2^2$. Apply
$(a+b)^2\le(1+t)a^2+(1+t^{-1})b^2$ to obtain
\eqref{f16v28:nested-Young}. These arguments are bounded-operator identities
on the whole $L^2$ space, including constants and rare-event-weighted tests.
\end{proof}

The cancellation of $E_Uf$ is essential. Its norm can be large, but it has no
innovation in $D$. Charging that norm independently for every collar witness
would recreate V27's extensive term. The subtraction changes neither the
observable $Q_Df$ nor the conditional probability used to estimate $H$.

\subsection{The same rough form now has a local-energy remainder}

Use V27's actual repair regions $S_c$, rough events $H_c(Q)$, and probability
\[
 p=p_\gamma(Q)=\min\{1,e^{t_0(\overline M_r-Q^2)}\},\qquad
 t_0=(24\pi^2)^{-1}.
\]
The preparation theorem gives $P(H_c(Q)\mid\zeta_{S_c^c})\le p$ with all
compact exterior fields retained. Let $D$ run over a deterministic indexed
family, with rates $r_D\ge0$ and actual collars $C_D=\partial_{\mathcal G}D$.
Keep exactly V27's form
\[
 \mathfrak D_Q(f)=\sum_Dr_D\|1_{H_{C_D}(Q)}Q_Df\|_2^2,
 \qquad a_c=\sum_{D:c\in C_D}r_D.
\]
Repeated blocks keep their indices. Define two positive forms
\begin{align}
 \mathcal A(f)&=\sum_c a_c\mathcal E_{S_c}(f),\notag\\
 \mathcal B(f)&=\sum_Dr_D\sum_{c\in C_D}
       \{\mathcal E_{D\cup S_c}(f)-\mathcal E_{S_c}(f)\}.
 \label{f16v28:AB}
\end{align}
Each difference in the second line is a positive nested-projection form.

\begin{theorem}[Nonextensive localization of rough innovations]
\label{f16v28:local-form}
At every admitted finite regulator and for all $f\in L^2(P)$,
\begin{equation}
 \boxed{\quad
 \mathfrak D_Q(f)^{1/2}\le\mathcal A(f)^{1/2}
                                      +p^{1/2}\mathcal B(f)^{1/2}.
 \quad}
 \label{f16v28:local-form-bound}
\end{equation}
There is no $\operatorname{Var}_Pf$ or total-cell count on the right. If
$D^+$ is any declared set containing $D\cup S_c$ for every $c\in C_D$, then
\begin{equation}
 \mathcal B(f)\le\sum_Dr_D|C_D|\mathcal E_{D^+}(f).
 \label{f16v28:B-upper}
\end{equation}
The target rough form, its gates, and its update rates have not been changed.
\end{theorem}
\begin{proof}
The event $H_c$ is exterior to $D$ for $c\in C_D$. Apply
Proposition~\ref{f16v28:nested} to each pair $(D,S_c)$. The indicator of a
union is at most the sum of its member indicators, so
\[
 \mathfrak D_Q(f)\le\sum_Dr_D\sum_{c\in C_D}
                                      \|1_{H_c}Q_Df\|_2^2.
\]
Minkowski in this finite weighted direct sum gives
\eqref{f16v28:local-form-bound}; the two sums of squared terms are exactly
\eqref{f16v28:AB}. If $D\cup S_c\subseteq D^+$, monotonicity of averaged
conditional variance gives $\mathcal E_{D\cup S_c}\le\mathcal E_{D^+}$.
Dropping the nonnegative subtracted energy proves \eqref{f16v28:B-upper}.
There is no independence assertion and no interchange of a growing-volume
limit with a conditional expectation.
\end{proof}

In particular, the probability multiplies an energy of a \emph{local union},
not the variance of the entire history. Converting that energy to the original
clock is still an additional estimate. The next construction keeps its local
update rate explicit before paying that comparison.

\subsection{Two repair families with bounded coverage}

Retain V27's box parameters $w=w_\gamma$, $h=h_\gamma=2w+5$, and choose an
integer $\rho\ge1$ dominating both the native interaction range and the frame
support allowance in ambient coordinatewise distance. Put $v=w+\rho+2$.
Choose an integer $R\ge2(v+\rho+1)$.

Use V26's indexed product cover $\mathcal B_R$. In a spatial cycle with
$m>R$ take every cyclic interval of length $R$; when $m\le R$ take the whole
cycle once. In time take all nonempty clipped intervals
$[s,s+R-1]\cap[0,n]$, $1-R\le s\le n$, without deduplication. Product rates
are $r_D=R^{-(k+1)}$, where $k$ is the number of cut spatial directions.
Let $D^+$ expand each coordinate interval of $D$ by $v$, saturating a short
cycle and clipping at natural time endpoints. Set
\[
 \kappa_{\rm rep}=h^4,\qquad \kappa_+=(1+2v/R)^4,
\]
\begin{equation}
 \begin{split}
 \mathcal L_R&=\sum_{D\in\mathcal B_R}r_DQ_D,\\
 \mathcal L_{\rm rep}&=\kappa_{\rm rep}^{-1}\sum_cQ_{S_c},\qquad
 \mathcal L_+=\kappa_+^{-1}\sum_{D\in\mathcal B_R}r_DQ_{D^+}.
 \end{split}
 \label{f16v28:generators}
\end{equation}
Their forms have the matching symbols $\mathcal E_R,\mathcal E_{\rm rep},
\mathcal E_+$. These are exact complete conditional updates in the same
unmoving coordinates. They are not elementary balanced updates.

\begin{proposition}[Finite geometry and actual rate bounds]
\label{f16v28:geometry}
Every cell has coverage exactly one under $\mathcal L_R$ and coverage at most
one under each of $\mathcal L_{\rm rep}$ and $\mathcal L_+$. Moreover,
\begin{equation}
 \begin{split}
 D\cup S_c&\subseteq D^+\quad(c\in C_D),\qquad |D^+|\le(R+2v)^4,\\
 a_*:=\sup_c a_c&\le \frac{8\rho}{R}(1+2\rho/R)^3,\\
 b_*:=\max_D|C_D|&\le8\rho(R+2\rho)^3.
 \end{split}
 \label{f16v28:geometry-bounds}
\end{equation}
All bounds include periodic seams, short cycles, clipped temporal ends and
repeated indices.
\end{proposition}
\begin{proof}
V26's product count gives unit coverage for $D$, and the stated upper bound
on $a_*$. A collar cell lies at coordinatewise distance at most $\rho$ from
$D$. Its repair box has coordinatewise radius $w$, so the declared enlargement
contains $S_c$, including the temporal and periodic saturation cases.

For a fixed cell, at most $(2w+1)^4\le h^4$ repair centers have a box containing
it. In a cut spatial direction, at most $R+2v$ starting indices have an enlarged
interval containing that cell. If the enlargement saturates the cycle, the
number is $m\le R+2v$. In an uncut direction there is only one interval. The
same $R+2v$ bound holds for indexed clipped time intervals by restricting the
possible start indices. Product rates therefore give enlarged coverage at most
$(1+2v/R)^{k+1}\le\kappa_+$. The displayed normalizations prove the rate claims.

The actual collar is contained in the coordinatewise $\rho$-enlargement minus
$D$. Assign a point of that difference to one coordinate in which it lies
outside the corresponding interval. There are at most $2\rho$ positions in
that coordinate and at most $R+2\rho$ in each other coordinate, for at most
four choices. Saturated directions and missing physical time faces only remove
positions. This proves the bound on $b_*$. No invariance of the probability
under spatial translations was used in these counts.
\end{proof}

\begin{theorem}[Rough-update absorption in explicitly normalized local forms]
\label{f16v28:absorption}
Define
\begin{equation}
 u_{\gamma,R}=\kappa_{\rm rep}a_*,\qquad
 z_{\gamma,R,Q}=p_\gamma(Q)b_*\kappa_+,\qquad
 \theta_{\gamma,R,Q}=u_{\gamma,R}+z_{\gamma,R,Q}.
 \label{f16v28:theta}
\end{equation}
Then, on the entire native $L^2(P)$ space,
\begin{equation}
 \boxed{
 \begin{aligned}
 \mathfrak D_Q(f)^{1/2}
 &\le\sqrt{u_{\gamma,R}\mathcal E_{\rm rep}(f)}
                  +\sqrt{z_{\gamma,R,Q}\mathcal E_+(f)},\\
 \mathfrak D_Q(f)
 &\le\theta_{\gamma,R,Q}
                  \{\mathcal E_{\rm rep}(f)+\mathcal E_+(f)\}.
 \end{aligned}}
 \label{f16v28:absorption-bound}
\end{equation}
In particular,
\begin{equation}
 \theta_{\gamma,R,Q}\le
 C_\rho\left(\frac{h_\gamma^4}{R}
           +p_\gamma(Q)R^3(1+2v/R)^4\right).
 \label{f16v28:theta-scale}
\end{equation}
There is no $B$, $n$, $\mathcal N$, or number of update indices in this
coefficient. All compact field values, including rough collars, remain covered.
\end{theorem}
\begin{proof}
By definition of $a_*$,
$\mathcal A(f)\le a_*\sum_c\mathcal E_{S_c}(f)
=u_{\gamma,R}\mathcal E_{\rm rep}(f)$.
Equation \eqref{f16v28:B-upper} gives
$p\mathcal B(f)\le p b_*\sum_Dr_D\mathcal E_{D^+}(f)
=z_{\gamma,R,Q}\mathcal E_+(f)$.
Insert these two bounds in Theorem~\ref{f16v28:local-form}. Cauchy--Schwarz
in the two-dimensional vector of square roots proves the second line.
Use \eqref{f16v28:geometry-bounds} and $R\ge2\rho$ for
\eqref{f16v28:theta-scale}. The price for conditioning an event on its repair
exterior has been paid as an actual energy before this comparison; an
unweighted event probability has never been multiplied by a separate test norm.
\end{proof}

\subsection{A polylogarithmic choice compatible with the controlled-row coupling}

\begin{proposition}[One small coefficient with two finite local update scales]
\label{f16v28:scales}
At sufficiently large fixed-$r$ coupling, take
\begin{equation}
 R_\gamma=\lceil h_\gamma^5\rceil,\qquad
 Q_\gamma^2=\overline M_r+\frac{20}{t_0}\log h_\gamma.
 \label{f16v28:scale-choice}
\end{equation}
Then $p_\gamma(Q_\gamma)=h_\gamma^{-20}$, all the required geometry holds, and
\begin{equation}
 \theta_\gamma\le C_r(h_\gamma^{-1}+h_\gamma^{-5})
             =O_r((\log\gamma)^{-1})\longrightarrow0.
 \label{f16v28:theta-small}
\end{equation}
The repair side is $O_r(\log\gamma)$; the enlarged blocks have side
$O_r((\log\gamma)^5)$ and at most $O_r((\log\gamma)^{20})$ cells.
In the separate V26 controlled-collar theorem, these same $R_\gamma,Q_\gamma$
give a positive admitted-row margin eventually. That statement is still only
about its admitted pairs of deterministic collars.
\end{proposition}
\begin{proof}
The probability identity follows directly from its definition. Here $h_\gamma$
is bounded above and below by positive fixed-$r$ multiples of $\log\gamma$
eventually; $v=O(h_\gamma)$. Thus $R_\gamma\ge2(v+\rho+1)$ eventually and
$\kappa_+$ is bounded. In \eqref{f16v28:theta-scale} the two terms are at most
$C h_\gamma^{-1}$ and $C h_\gamma^{15-20}$, respectively.
All update sizes follow from the actual side lengths, not a volume limit.

For the last statement $Q_\gamma=O_r(\sqrt{\log h_\gamma})$, so
$Q_\gamma\le\eta_*\sqrt\gamma$ eventually. V26's full-output estimate on
admitted rows has budget
\[
 2C_rQ_\gamma+2R_\gamma^8
                  e^{A_r(1+Q_\gamma^2)-c_*\gamma}.
\]
Multiply it by the actual affected-block rate $C_\rho/R_\gamma$.
The first part tends to zero. The second is a fixed power of $h_\gamma$
times $e^{-c_*\gamma}$ and also tends to zero. Thus the admitted-row margin
is eventually at least $1/2$. This calculation does not extend the coupling
to excluded collars or identify its error with an arbitrary spectral defect.
\end{proof}

Alternatively, retain V26's $Q_\gamma=\gamma^\alpha$ and
$R_\gamma=\lceil K_r(1+\gamma^\alpha)\rceil$, $0<\alpha<1/2$ fixed.
Equation \eqref{f16v28:theta-scale} then gives
$\theta_\gamma\le C_r\gamma^{-\alpha}(1+\log\gamma)^4+
C_r\gamma^{3\alpha}e^{-c_r\gamma^{2\alpha}}$ eventually. The polylogarithmic
choice above uses smaller blocks at the expense of a slower vanishing form
error. Neither choice changes the native measure.

\subsection{Exterior-gated updates are reversible, with a relative form bound}

Let $M_D=1_{H_{C_D}(Q)}$ and define
\begin{equation}
 \begin{split}
 \mathcal L_{\rm bad}&=\sum_Dr_DM_DQ_D,\qquad
 \mathcal L_{\rm good}=\sum_Dr_D(1-M_D)Q_D,\\
 \mathcal L_{\rm aug}&=\mathcal L_R+\mathcal L_{\rm rep}+\mathcal L_+,\\
 \mathcal L_{\rm keep}&=\mathcal L_{\rm good}+\mathcal L_{\rm rep}+\mathcal L_+
                     =\mathcal L_{\rm aug}-\mathcal L_{\rm bad}.
 \end{split}
 \label{f16v28:gated-generators}
\end{equation}
The label ``keep'' means that repair and enlarged-block moves are always
retained. It does not mean that only good configurations occur in the process.

\begin{theorem}[Small relative deletion of rough block updates]
\label{f16v28:relative}
These are bounded nonnegative self-adjoint generators at each finite regulator,
with invariant probability $P$. The gate changes an update rate, not its
conditional output law. Each cell has total coverage rate at most three under
$\mathcal L_{\rm aug}$, and no larger under $\mathcal L_{\rm keep}$.
When $\theta=\theta_{\gamma,R,Q}<1$,
\begin{equation}
 \boxed{
 \begin{aligned}
 0\preceq\mathcal L_{\rm bad}
      &\preceq\theta(\mathcal L_{\rm rep}+\mathcal L_+),\\
 (1-\theta)\mathcal L_{\rm aug}
      &\preceq\mathcal L_{\rm keep}\preceq\mathcal L_{\rm aug}.
 \end{aligned}}
 \label{f16v28:relative-order}
\end{equation}
In fact
$\mathcal L_{\rm keep}\succeq
\mathcal L_R+(1-\theta)(\mathcal L_{\rm rep}+\mathcal L_+)$.
Both full augmented generators have constants as their common null space.
These comparisons do not give a new size-uniform lower bound on their gaps.
\end{theorem}
\begin{proof}
The indicator $M_D$ is measurable outside $D$, so it commutes with $E_D,Q_D$.
Thus $M_DQ_D$ and $(1-M_D)Q_D$ are positive orthogonal projections, and
\[
 \langle f,M_DQ_Df\rangle=\|M_DQ_Df\|_2^2.
\]
Their jump forms are obtained by multiplying the ordinary reversible
conditional-update edge measure by the exterior indicator. That indicator
has the same value before and after the update. Detailed balance, preservation
of constants and the Markov property therefore hold without a field-independent
rate assumption. Equivalently their semigroups make the declared conditional
refresh at rate $r_D$ only while the exterior gate allows it. The remaining
two families are ordinary deterministic-rate conditional refreshes.

The first operator bound is Theorem~\ref{f16v28:absorption}, since the form
of $\mathcal L_{\rm bad}$ is exactly $\mathfrak D_Q$. Subtract it from
$\mathcal L_{\rm aug}$ to get the stronger lower bound and then the second
line of \eqref{f16v28:relative-order}. Coverage was proved above. Both
generators contain every repair $Q_{S_c}$ at positive rate. A function in
their common kernel is measurable outside each $S_c$, hence outside each
individual cell. The strictly positive finite-product density of the native
law makes such a function constant, up to null sets. This can also be seen
by successive single-cell conditional expectations under the equivalent
product reference. None of these statements estimates a thermodynamic gap.
\end{proof}

\begin{proposition}[Gap and shifted-resolvent comparison]
\label{f16v28:resolvent}
Let $\lambda_{\rm aug},\lambda_{\rm keep}$ be the variational gaps on the
common centered space. Then
\begin{equation}
 (1-\theta)\lambda_{\rm aug}\le\lambda_{\rm keep}\le\lambda_{\rm aug}.
 \label{f16v28:gap-relative}
\end{equation}
For every $s>0$, put $C_s=\mathcal L_{\rm aug}+sI$. On the full space,
\begin{equation}
 \begin{split}
 0\preceq C_s^{1/2}\bigl[(\mathcal L_{\rm keep}+sI)^{-1}
            -(\mathcal L_{\rm aug}+sI)^{-1}\bigr]C_s^{1/2}
 \preceq \frac{\theta}{1-\theta}I.
 \end{split}
 \label{f16v28:resolvent-bound}
\end{equation}
No operator ordering of their finite-time semigroups is inferred.
\end{proposition}
\begin{proof}
Take infima of the form comparison over centered unit vectors to get
\eqref{f16v28:gap-relative}. Also
$(1-\theta)C_s\preceq\mathcal L_{\rm keep}+sI\preceq C_s$.
Conjugate by $C_s^{-1/2}$ and invert these strictly positive operators. The
result is between $I$ and $(1-\theta)^{-1}I$; subtract $I$ to obtain
\eqref{f16v28:resolvent-bound}. No commutation of different local projections is
used. A small relative error alone does not supply a positive lower gap for
either side of the comparison.
\end{proof}

\subsection{The original clock and the remaining global test}

\begin{proposition}[The auxiliary rates have a paid original-clock comparison]
\label{f16v28:clock}
With $L=R+2v+3$ there is an ambient-size-independent bound
\begin{equation}
 \mathcal L_{\rm keep}\preceq\mathcal L_{\rm aug}
       \preceq\mathcal C_r(\gamma,R,h)\widehat{\mathcal L},\qquad
 \mathcal C_r(\gamma,R,h)=C_r L^4e^{C_r\gamma L^3}.
 \label{f16v28:clock-bound}
\end{equation}
Consequently a separately proved gap $g>0$ for either augmented generator
would imply an original-clock gap at least $g/\mathcal C_r$.
No such new uniform augmented gap is proved here.
\end{proposition}
\begin{proof}
For each unmoving update set, include its chord cubes and every moving port
whose time or following time at that cube lies in the set. This is the exact
preceding-port enlargement from V24. Its balanced exterior sigma algebra is
contained in the unmoving exterior sigma algebra, so its averaged conditional
variance dominates that of the unmoving update. Each resulting support has
at most $CL^4$ elementary groups and a time-major ordering cut of at most
$CL^3$ groups. The bounded classical word oscillations and V25's interface
comparison therefore give the local factor $C_rL^4e^{C_r\gamma L^3}$.
All product references retain the endpoint Haar exponents and endpoint factors.

Each of the three normalized unmoving families has coverage at most one.
An elementary chord belongs only through its own cell. An elementary moving
port can belong through either of two adjacent time cells, so the corresponding
aggregate balanced rate is at most two per family. Sum the three nonnegative
comparisons and absorb the factor six in $C_r$. This proves
\eqref{f16v28:clock-bound}. If one of the augmented forms has a lower bound
$g\operatorname{Var}f$, its upper comparison gives the stated implication.
The comparison factor is generally large; it has not been declared free.
\end{proof}

For the polylogarithmic scales, this bound is at worst
$C_rh_\gamma^{20}\exp(C_r\gamma h_\gamma^{15})$ after fixed constants are
enlarged. It need not improve V25's finite-width bound. The new conclusion is
\emph{relative local-form absorption without an extensive remainder}, not a
better direct source covariance obtained from this crude clock comparison.

Two finite checks identify what the new conclusion does not allow. For an
independent uniform sign $X$ and a Bernoulli $Y$ with $P(Y=1)=p$, take
$D=\{X\}$, $H=\{Y=1\}$, and $f=X1_H$. Then
\[
 \|1_HQ_Df\|_2^2=\mathcal E_D(f)=p.
\]
Thus a bound $\mathfrak D_Q\le o(1)\mathcal E_R$ cannot follow from event
rarity alone. The repair energy in \eqref{f16v28:absorption-bound} is necessary.

Also, even when an independent $X$ has identical conditional laws at all
values of $Y$, an $X$ update gated by $1_{\{Y=0\}}$ has different rates at
$Y=0$ and $Y=1$. Comparing configurations with the same $X=+1$ and those two
values of $Y$, an allowed refresh creates an $X$ mismatch with probability
$1/2$ while the other configuration does not refresh. The ungated boundary
influence was zero, but this gate-switching contribution is not. Hence V26's
admitted-row table is not by itself a path-coupling proof for
$\mathcal L_{\rm keep}$. Detailed balance proved above does not erase this
term. The ungated repair and enlarged-block families also need to be included
in any global spectral argument.

The next concrete target is therefore an independently proved coercive or
boundary-response estimate for the \emph{displayed retained generator}, with
its gate switching and its two ungated families included. Equivalently one may
seek a sharper all-boundary argument for the original blocks. The rough form
has now been absorbed relative to explicit local auxiliary energies at all
volumes, but it has not been identified with every defect in such a future
argument. No infinite-order expansion, unweighted rough probability, or
unverified commutator cancellation is used to bridge that gap.

\subsection{Diagnostics, preservation, and claim boundary}

The accompanying exact diagnostic is
\begin{center}\small\path{verify_ym_f16_v28_nested_absorption.py}.\end{center}
It checks noncommuting conditional projections, their nested union and energy
difference, event-weighted Young forms, positive rough/good generators,
finite-cover and enlargement rates, relative form and resolvent inequalities,
and the two limiting examples above. A continuous positive polynomial-density
fixture checks the normalized identities without assigning point masses to a
continuous coordinate. These are finite algebraic checks, not native compact
integral enclosures, formal verification, or numerical certification of the
inherited preparation constants.

The package records actual diagnostic outputs and any executed predecessor
reruns. The cumulative source changes only V27's title version and adds this
exact terminal section. Removing it and restoring that title recovers V27
byte for byte. No previous theorem or bibliography entry is silently edited.
The new estimates rely on the specified V27 preparation and finite conditional
identities; the optional good-row calculation retains V26's analytic hypotheses.
The full inherited archive has not been independently recertified.

\begin{firewall}
V28 removes the extensive full-variance remainder from the specified V27
rough-innovation estimate. It proves small relative form absorption into two
explicitly normalized local repair families, a reversible exterior-gated
retained dynamics, and gap/resolvent comparisons with the augmented dynamics.
Its update radii can be chosen polylogarithmically in coupling, with an
ambient-size-independent vanishing relative error. It pays a separate,
pessimistic conversion to the original elementary clock. It does not prove a
new uniform augmented or original sampler gap, all-boundary spatial
susceptibility, unrestricted coherent covariance control, or a physical mass
gap. The native measure, compact outputs, endpoint factors, original sources,
and original elementary rates are unchanged. Bridge-window removal, physical
Gauss restriction, infrared contact coercivity, same-top transfer comparison,
physical-time calibration and interacting continuum construction remain separate.
\end{firewall}
% END F16 V28 APPEND

% BEGIN F16 V29 APPEND -- 2026-09-18
% Insert immediately before the final document-ending command in V28.
% No inherited theorem, probability, source, or elementary rate is altered.

\clearpage
\section{V29: fixed-depth angular repairs and a sharper original-clock bound}
\label{f16v29:context}

V28 controls the removal of rough updates relative to two helper families, but
relative closeness supplies no spectral edge. We return to the cost of mixing
one complete native conditional. For this purpose its exceptional event can be
chosen at a \emph{fixed angular radius}, rather than at the shrinking angular
radius used to control every local boundary response. Fixed-depth buffers then
suffice. Their original-coordinate comparison costs $e^{C_r\gamma}$, not an
exponential in the volume of a growing slab.

There are two separate steps. A restricted convex native conditional has a
uniform whole-cell heat-bath variance bound. A weighted repair estimate transfers
that bound to the complete conditional, including its large fields, when the
\emph{actual} collar has a smaller fixed angular radius. We do not remove this
collar condition by probability alone. For complete temporal slabs, V28's nested
innovation identity and an independent two-overlapping-slab comparison absorb
the rough-collar contribution back into the same slab form. This yields
\[
 \widehat\lambda_P\ge c_r e^{-C_r\gamma},\qquad
 B\le e^{\beta_r\gamma},\quad n\ge1,
\]
with $\beta_r>0$ fixed and no $B$ or duration in the bound. Its cross-section
restriction is essential; its lower bound still vanishes with coupling. A
consequence controls V28's displayed retained generator in that same regime.
No unrestricted spatial gap or physical mass is inferred.

\subsection{Dependencies and the change of exceptional scale}

The supplied V28 appendix and audit were read completely. The decisive new
inputs are proved below; their inherited premises are the V16/V20 conditional
set-energy inequality on complete charts, V22's restricted measured Hessian
bounds on products of whole-cell balls, the bounded-word product reference, and
V22/V24's all-boundary temporal channels and exact slab clock. V28's nested
identity is reused, not reasserted as a new result. The higher-order source
expansions are not needed in this proof of a sampler gap.

Two semantic searches returned no results. Targeted mounted-text searches and
identified passages were then examined. F15 V70 already uses a Neumann Bochner
calculation on a convex spherical cap and a separate full-law argument. The
monograph's Bochner--Schur construction concerns a different reduced-fibre
response operator. Neither result is a theorem about the current all-test
whole-cell conditional update, and neither is substituted for it. V25's
all-finite-width baseline remains valid outside the new, stronger regime.
This is a dependency review, not independent certification of the archives.
The next scheduled companion checkpoint remains V30.

For contemporary attribution, Ascolani--Lavenant--Zanella,
\href{https://arxiv.org/abs/2410.00858v2}{\emph{Entropy contraction of the Gibbs
sampler under log-concavity}}, Corollary 3.7, proves an approximate variance
tensorization with a condition-number cost. The paper's assumptions, rate
convention, and that corollary were checked. Its Euclidean theorem is not
applied to a hard product boundary without argument. We give the weaker
product-ball estimate needed here directly, using the explicitly inherited
reflecting domain and its nonnegative curvature terms. No entropy contraction
for the complete native law is claimed. The supplied import memoranda remain
method proposals: their separation of locality cost from power gain, and of an
auxiliary spectral edge from the physical transfer, is preserved.

\subsection{Fixed-depth repair at two declared angular radii}

Retain $P=P_{\gamma,\mathbf y}$, $B=m^3$, $m\ge4$, $n\ge1$,
$\max_{i,e}|y_{i,e}|\le r$, and $\gamma\ge1+r^2$. Let
$e_c=|\zeta_c|^2\le E_*\gamma$, $E_*=12\pi^2$, and
$t_0=(24\pi^2)^{-1}$. Write
\[
 \mathcal L_{\rm cell}=\sum_c Q_{\{c\}},\qquad
 \mathcal E_{\rm cell}(f)=\sum_c\|Q_{\{c\}}f\|_2^2.
\]
A coordinate here is one complete \emph{unmoving} cell, of dimension 36 or 15.
This is not yet V14's elementary balanced chord/port clock.

Fix V22's small angular comparison radius $\eta_0:=\eta_*>0$ and its constants
$0<\lambda\le M$. We will choose a smaller fixed radius $\eta_1>0$. For any
finite free cell set $\Omega$, let $\mu=P_\Omega^q$ denote its complete native
conditional, and set
\[
 G_\Omega=\{e_c\le\eta_0^2\gamma\text{ for every }c\in\Omega\}.
\]
This event is new notation for this comparison only. It does not redefine V4's
$\mathcal O_R$, $p_{\rm good}$, or the rough threshold in V28.

\begin{proposition}[Fixed repair boxes with a controlled actual boundary]
\label{f16v29:angular-preparation}
There exist fixed $\eta_1\in(0,\eta_0)$, a fixed integer $v_0$, and positive
constants $A_r,c_0,C_r$, all independent of coupling and sizes, with the
following property. Suppose
\begin{equation}
       |q_d|\le\eta_1\sqrt\gamma\quad(d\in\partial_{\mathcal G}\Omega).
 \label{f16v29:angular-collar}
\end{equation}
For each $c\in\Omega$ let $S_c^\Omega$ be the intersection of $\Omega$ with a
coordinatewise box of radius $v_0$ about $c$, including short cycles in full
and clipping only at the natural ends. Then
\begin{equation}
 \mu(e_c>\eta_0^2\gamma\mid\zeta_{\Omega\setminus S_c^\Omega})
       \le p_0(\gamma):=\min\{1,A_r e^{-c_0\gamma}\}
 \label{f16v29:angular-tail}
\end{equation}
for every compact value of the other free variables, in the explicit
conditional-density versions. The original frozen $q$ remains fixed. Moreover
\begin{equation}
 \sum_{c\in\Omega}\mathcal E_{S_c^\Omega}^{\mu}(f)
       \le K_0(\gamma)\mathcal E_{\rm cell}^{\mu}(f),
       \qquad K_0(\gamma)\le C_r e^{C_r\gamma}.
 \label{f16v29:constant-repair-energy}
\end{equation}
The free set can have holes or disconnected components. No fast coordinate
inside it is restricted in these statements.
\end{proposition}
\begin{proof}
We spell out the boundary preparation, because a cell near the actual boundary
of $\Omega$ has no full free buffer in all directions. Denote the constants in
\eqref{f16v20:conditional-set} by $a_r,b_r$, and put
$\rho=b_r/(1+b_r)<1$. Fix $c$, condition on the exterior of $S_c^\Omega$,
and allow arbitrary signed sources of size at most $t_0$ on this free set.
Within it let $u_d$ be the actual mean energy. At frozen cells in $\Omega$
use their actual deterministic energy, bounded by $E_*\gamma$; outside
$\Omega$ extend the scalar array by $Q^2=\eta_1^2\gamma$. Only collar values
of this extension enter any native interaction, and these dominate the
actual prescribed $q$.

For an ambient set $A$ lying inside the coordinate box, apply the native set
inequality to $A\cap\Omega$. As in V26's scalar extension argument, adding the
entries outside $\Omega$ gives
\[
 \sum_{d\in A}\widetilde u_d
 \le [a_r+(1+b_r)Q^2]|A|+
                           b_r\sum_{d\in\partial A}\widetilde u_d.
\]
For graph balls of radii $k<K$ contained in that coordinate box, iteration gives
\begin{equation}
 u_c\le C_r^{(0)}+C_r^{(1)}Q^2
       +E_*C_{\rm gr}\gamma(K+1)^4\rho^K.
 \label{f16v29:stopped-mixed-boundary}
\end{equation}
Here $C_r^{(0)},C_r^{(1)}$ come from the convergent series
$\sum_{k\ge0}\rho^k(k+1)^4$ and do not depend on $K$ or the regulator.
The terminal compact cap is legitimate even when the ball crosses the
artificial repair boundary. If the finite component is already filled, one
may continue the same recurrence or take its zero terminal limit.

First choose $\eta_1$ so small that
$C_r^{(1)}\eta_1^2\le\eta_0^2/8$. Then choose a \emph{fixed} $K$ so large
that $E_*C_{\rm gr}(K+1)^4\rho^K\le\eta_0^2/8$. The native interaction has
fixed coordinatewise range, so a fixed $v_0$ contains all these graph balls.
Thus the right side of \eqref{f16v29:stopped-mixed-boundary} is at most
$C_r^{(0)}+\eta_0^2\gamma/4$. It holds throughout the signed source cube.
Integrating the actual conditional $\log Z$ derivative for a source only at
$c$ yields
\[
 \mathbb E e^{t_0e_c}\le
       \exp[t_0C_r^{(0)}+t_0\eta_0^2\gamma/4].
\]
Exponential Markov proves \eqref{f16v29:angular-tail}, for example with
$c_0=t_0\eta_0^2/2$ and a fixed prefactor. This argument keeps both the
moderate \emph{actual} exterior and the arbitrary artificial repair exterior.

For \eqref{f16v29:constant-repair-energy}, work in unmoving whole cells. Each
$S_c^\Omega$ has a fixed maximum number of cells and every cell occurs in only
a fixed number of these sets. Relative to the original product cell reference,
with one-cell classical terms included in their own factors, all remaining
words meeting a repair have total oscillation at most $C_r\gamma$. The bounded-density variance comparison of V24, or V25's
interface version, therefore bounds one repair energy by $C_re^{C_r\gamma}$
times the sum of its complete single-cell energies. This uses original cell
factors, including fractional endpoint Haar powers; no ratio to pure Haar at
its cut is taken. Sum the bounded overlap. Conditional on the remaining cells
these single-cell laws are exactly those of $\mu$. The radius and overlap are
fixed before coupling varies, so all fixed geometric factors can be included
in $C_r$. Arbitrary holes only reduce each repair's group count.
\end{proof}

The radius in this proposition is constant because the event being excluded
has energy of order $\gamma$. It does not supply uniform thermal moments at
constant depth, and it does not replace V27's logarithmic repair for a much
smaller rough threshold.

\subsection{A restricted heat-bath estimate and an all-test transfer to the full law}

\begin{proposition}[A nonsharp variance tensorization on the restricted product]
\label{f16v29:box-heatbath}
Let $\pi$ have smooth measured potential $V$ on a finite product of the
whole-cell balls, with $\lambda I\preceq\nabla^2V\preceq MI$. For its
rate-one whole-cell conditional updates,
\begin{equation}
 \operatorname{Var}_\pi f\le K_{\rm box}
                 \sum_c\mathbb E_\pi\operatorname{Var}_\pi(f\mid\zeta_{c^c}),
       \qquad K_{\rm box}=1+M/\lambda,
 \label{f16v29:box-tensorization}
\end{equation}
for every $f\in L^2(\pi)$. The constant has no number-of-cells factor.
\end{proposition}
\begin{proof}
Use the reflecting scalar realization on this product, with positive generator
$\mathcal A=-\Delta+\nabla V\cdot\nabla=\sum_c\mathcal A_c$. Its gap is at
least $\lambda$ by the inherited product-ball Brascamp--Lieb identity. For a
smooth function $u$ satisfying the Neumann condition on each ball face, the
full and blockwise integrated Bochner identities have the same nonnegative
boundary-curvature sum. The full squared Hessian contains every block-diagonal
squared Hessian as part of a sum of squares. Since $(\nabla^2V)_{cc}\preceq MI$
and $\nabla^2V\succeq0$, subtracting these identities gives
\begin{equation}
       \sum_c\|\mathcal A_cu\|_2^2
       \le\|\mathcal Au\|_2^2+M\langle u,\mathcal Au\rangle.
 \label{f16v29:partial-bochner}
\end{equation}
In particular no row-wise absolute diagonal-dominance condition is imposed.
The boundary curvature on a cell face acts on that cell's tangential
components only; cross-cell tangents have zero curvature. This is why the
boundary terms agree rather than introduce a number-of-cells loss.

The domain passage uses the same product reflecting realization as V4/V22.
At each finite regulator $V$ and its derivatives are bounded on the compact
product. Finite tensor sums of smooth Neumann eigenfunctions of the individual
balls form a graph core for the product Neumann Laplacian; the bounded smooth
first-order drift is relatively bounded with arbitrarily small relative bound.
Thus the same core is a graph core for $\mathcal A$ on weighted $L^2$.
Inequality \eqref{f16v29:partial-bochner} makes the partial generators Cauchy
under graph-core approximation. It therefore extends to $u$ in the full
scalar generator domain. Conditional integration by parts also extends
$E_c\mathcal A_cu=0$. This argument neither imposes a boundary condition on
$f$ nor substitutes the smooth-boundary theorem for an unmentioned corner
realization.

For centered $f\in L^2(\pi)$ solve $u=\mathcal A^{-1}f$. Then
$\langle u,\mathcal Au\rangle\le\lambda^{-1}\|f\|_2^2$ and
\[
 \|f\|_2^2
   =\sum_c\langle (I-E_c)f,\mathcal A_cu\rangle
   \le\left(\sum_c\|(I-E_c)f\|_2^2\right)^{1/2}
        \left(\sum_c\|\mathcal A_cu\|_2^2\right)^{1/2}.
\]
Insert \eqref{f16v29:partial-bochner} and divide by $\|f\|_2$, unless it is
zero. This proves \eqref{f16v29:box-tensorization}; complex tests follow by
real and imaginary parts. The constant is deliberately not optimized.
\end{proof}

\begin{proposition}[Variance transfer with the exceptional test mass paid]
\label{f16v29:variance-gluing}
Let $\mu$ be a finite-product probability, $G=\bigcap_cG_c$ a positive-probability
product event, and $\pi=\mu(\cdot\mid G)$. Suppose $\pi$ has whole-cell variance
tensorization constant $K_{\rm box}$ and, for every $g\in L^2(\mu)$,
\begin{equation}
 \|1_{G^c}g\|_2\le
       \sqrt{K\mathcal E_{\rm cell}^\mu(g)}+\sqrt\epsilon\|g\|_2,
             \qquad \epsilon\le\tfrac14.
 \label{f16v29:gluing-input}
\end{equation}
Then
\begin{equation}
 \operatorname{Var}_\mu f\le
       (2K_{\rm box}+4K)\mathcal E_{\rm cell}^\mu(f).
 \label{f16v29:gluing-output}
\end{equation}
No density lower bound of $\pi$ on $G^c$, and no invariance of $G$ under the
full conditional updates, is assumed.
\end{proposition}
\begin{proof}
For every coordinate $c$, conditional variance decomposition gives the exact
one-sided comparison
\begin{equation}
 \mu(G)\mathbb E_\pi\operatorname{Var}_\pi(f\mid\zeta_{c^c})
          \le\mathbb E_\mu\operatorname{Var}_\mu(f\mid\zeta_{c^c}).
 \label{f16v29:restriction-energy}
\end{equation}
Indeed condition on $z=\zeta_{c^c}$, restrict to $G_{c^c}$, and retain the
factor $\mu(G_c\mid z)$ multiplying the variance under the normalized
restricted conditional. It is one term of the full conditional variance.
Averaging proves \eqref{f16v29:restriction-energy}; no conditional normalizer
has been suppressed.

Put $a=\pi f$, $g=f-a$, and $T=\|g\|_{L^2(\mu)}$. Tensorization and
\eqref{f16v29:restriction-energy} bound
$\|1_Gg\|_2^2\le K_{\rm box}\mathcal E_{\rm cell}^\mu(f)$. The other part
is bounded by \eqref{f16v29:gluing-input}; all energies are unchanged by $a$.
Thus, with $\mathcal E=\mathcal E_{\rm cell}^\mu(f)$,
\[
 T^2\le K_{\rm box}\mathcal E+
                  (\sqrt{K\mathcal E}+\sqrt\epsilon T)^2
       \le (K_{\rm box}+2K)\mathcal E+2\epsilon T^2.
\]
Rearrange and use $\operatorname{Var}_\mu f\le T^2$. For an arbitrary $L^2$
test its restriction belongs to $L^2(\pi)$ because $\mu(G)>0$. The same
calculation, or truncation, justifies every conditional variance. Constants
and tests supported entirely on $G^c$ are both included.
\end{proof}

\begin{theorem}[Complete native conditional gap without an exponential size cost]
\label{f16v29:conditional-gap}
There are $A_r,c_0,C_r>0$ and a fixed-$r$ entry threshold such that, under
\eqref{f16v29:angular-collar}, for every finite $\Omega$ satisfying
\begin{equation}
                 |\Omega|A_re^{-c_0\gamma}\le\tfrac14,
 \label{f16v29:conditional-size}
\end{equation}
the \emph{complete} native conditional $\mu=P_\Omega^q$ obeys
\begin{equation}
 \boxed{\quad
 \operatorname{Var}_\mu f\le K_{\rm nat}(\gamma)
                     \mathcal E_{\rm cell}^{\mu}(f),
          \qquad K_{\rm nat}(\gamma)\le C_re^{C_r\gamma}.
 \quad}
 \label{f16v29:conditional-native-gap}
\end{equation}
The constants have no dependence on the shape or size of $\Omega$, ambient
volume, history length, or the admitted collar values. The collar condition
and \eqref{f16v29:conditional-size} are genuine hypotheses.
\end{theorem}
\begin{proof}
Use $G_\Omega$ defined above. The restricted law is exactly V22's native
product-ball conditional, and its Hessian has the same $\lambda,M$ when its
actual collar obeys \eqref{f16v29:angular-collar}. Hence
Proposition~\ref{f16v29:box-heatbath} applies with a fixed $K_{\rm box}$.
For every $c\in\Omega$, \eqref{f16v29:angular-tail} and V27's conditional
projection estimate give
\[
 \|1_{\{e_c>\eta_0^2\gamma\}}g\|_2
       \le\|Q_{S_c^\Omega}g\|_2+\sqrt{p_0}\|E_{S_c^\Omega}g\|_2.
\]
Direct-sum Minkowski and the union bound therefore give
\[
 \|1_{G_\Omega^c}g\|_2
       \le\sqrt{K_0(\gamma)\mathcal E_{\rm cell}^{\mu}(g)}
                                 +\sqrt{|\Omega|p_0}\|g\|_2.
\]
Apply Proposition~\ref{f16v29:variance-gluing} and
\eqref{f16v29:conditional-size}. This proves the theorem with
$K_{\rm nat}=2K_{\rm box}+4K_0$, bounded by the displayed exponential.
The probability term is used on the \emph{same test} before it is absorbed.
The interior is not replaced by its restriction. No all-field convexity or
pre-existing native gap is a premise.
\end{proof}

For instance every prescribed polynomial $|\Omega|$ is allowed eventually,
and an exponential $|\Omega|\le e^{\beta_0\gamma}$ is allowed for fixed
$\beta_0<c_0$. The cost $e^{C_r\gamma}$ remains. This theorem is not the
unproved all-compact-boundary conditional gap: surface configurations outside
\eqref{f16v29:angular-collar} have not been included.

\subsection{Fixed-depth repair of the actual slab collar}

We now work under the full native law $P$, with no conditioning or restriction
on its fast fields. Put
\[
 H_c^{(1)}=\{e_c>\eta_1^2\gamma\}.
\]
This is a fixed angular threshold, not V28's $Q_\gamma$. V20's complete
conditional preparation gives fixed-radius boxes $T_c$ and constants
$A_r',c_1>0$ such that
\begin{equation}
 P(H_c^{(1)}\mid\zeta_{T_c^c})\le
       p_1(\gamma):=\min\{1,A_r'e^{-c_1\gamma}\}.
 \label{f16v29:collar-smallness}
\end{equation}
To see this explicitly, choose a fixed graph depth $d_1$ with
$B_re^{-\kappa_rd_1}\le\eta_1^2/4$ in
\eqref{f16v20:prepared-mgf}. Exponential Markov at $t_0$ gives
\eqref{f16v29:collar-smallness}, for example with
$c_1=t_0\eta_1^2/2$. A fixed coordinate radius contains that depth, including
saturated cycles and natural ends. All compact $T_c$ exteriors are admitted.
As in \eqref{f16v29:constant-repair-energy}, bounded size and overlap yield
\begin{equation}
       \sum_c\mathcal E_{T_c}(f)
          \le K_{\rm rep}(\gamma)\mathcal E_{\rm cell}(f),
          \qquad K_{\rm rep}(\gamma)\le C_re^{C_r\gamma}.
 \label{f16v29:angular-repair-clock}
\end{equation}
Fix an integer $v$ larger than this coordinate radius plus two. It is fixed
independently of $\gamma,B,n$.

Use V22's exponential-width temporal channel with buffer $b=b_\gamma^{\rm exp}$,
$\ell=16b$, and $b\ge2v+2$. This is arranged by increasing the fixed entry
threshold. In the admissible channel regime $b\le C_r\gamma$. Let
\[
 I_s=[s,s+\ell-1]\cap[0,n],\quad
 J_s=[s-v,s+\ell-1+v]\cap[0,n],\quad 1-\ell\le s\le n.
\]
The corresponding updates contain all $B$ spatial cubes. Use the indexed
slab generator of V24, denoted here by $\mathcal L_{\rm sl}$, with rates
$1/\ell$. Let $C_s$ be its actual cell collar. Since every native interaction
meets at most two adjacent times,
\begin{equation}
 |C_s|\le2B,\qquad
       \sup_c\frac1\ell\#\{s:c\in C_s\}\le\frac2\ell,
       \qquad I_s\cup T_c\subseteq J_s\quad(c\in C_s).
 \label{f16v29:slab-collar-counts}
\end{equation}
Spatial components are never declared independent in these statements.

\begin{proposition}[A fixed enlargement is dominated by the original slab family]
\label{f16v29:enlargement-return}
In the above temporal channel regime, as quadratic forms on the full native
space,
\begin{equation}
       \frac1\ell\sum_s Q_{J_s}\preceq4\mathcal L_{\rm sl}.
 \label{f16v29:enlargement-domination}
\end{equation}
All compact exteriors and both endpoint factors are retained.
\end{proposition}
\begin{proof}
For each $s$ set $A_s=I_{s-v}$ and $B_s=I_{s+v}$, assigning the empty set
and zero innovation to an index outside the nonempty indexed family. Their
union is $J_s$ because $\ell>2v$. If one exclusive part is empty, one block
already equals $J_s$ and the needed inequality is immediate.

Otherwise condition on the actual exterior of $J_s$. Inside it the conditional
is the same inhomogeneous native time chain with the two fixed outside factors,
when present. The exclusive left and right parts are separated by an overlap
of length at least $\ell-2v$. Choose a $b$-step target inside this overlap,
at least $b$ from the right frozen endpoint. V22's actual forward channel,
with that right factor retained, has diameter at most $1/4$. Its
all-incoming-law $L^2$ consequence and Markov conditioning bound the maximal
correlation of the two exclusive parts by $1/2$. Appending the remaining
right path is a common conditional kernel and cannot increase this bound.
No estimate is applied to an output too close to the right frozen boundary.

On this fixed-exterior fibre, $E_{A_s}$ and $E_{B_s}$ are the projections onto
functions of the right and left exclusive parts. On centered space their
ranges have angle at most $1/2$. For vectors $x,y$ in those two centered
ranges, $|\langle x,y\rangle|\le\tfrac12\|x\|\|y\|$. Thus the
synthesis $(x,y)\mapsto x+y$ has squared norm at most $3/2$, and its
product with its adjoint gives $E_{A_s}+E_{B_s}\preceq3I/2$ on centered
space. Consequently $Q_{A_s}+Q_{B_s}\succeq\tfrac12 I$ there; constants
have zero on both sides.
Equivalently, before averaging, $Q_{J_s}\preceq2(Q_{A_s}+Q_{B_s})$ as a
conditional form. Integrate the actual exterior. Finally each shifted index
$s-v$ or $s+v$ occurs at most once in its own shifted sum of nonempty intervals.
Summing rates therefore gives \eqref{f16v29:enlargement-domination}. Repeated
clipped intervals retain their distinct indices, so very short histories are
covered by the same calculation.
\end{proof}

The proposition is independent of the gap comparison we are about to prove.
Its input is the already established conditional channel, not a presumed
original-clock lower bound. It is precisely the comparison that would be
missing for analogous overlapping spatial blocks.

\subsection{Rough-slab absorption closes without an elementary volume exponential}

Define exterior gates $M_s=1_{H_{C_s}^{(1)}}$, and write the actual slab energy
as $\mathcal E_{\rm sl}=\mathcal E_{\rm sl}^{\rm good}
+\mathcal E_{\rm sl}^{\rm bad}$, with
\[
 \mathcal E_{\rm sl}^{\rm bad}(f)
          =\frac1\ell\sum_s\|M_sQ_{I_s}f\|_2^2.
\]
These are exterior gates for \emph{these slabs}; none of V28's existing gates
is redefined.

\begin{proposition}[Local repair of all rough temporal boundary rows]
\label{f16v29:slab-absorption}
For every native $L^2$ test,
\begin{equation}
 \bigl(\mathcal E_{\rm sl}^{\rm bad}(f)\bigr)^{1/2}
 \le
 \sqrt{\frac{2K_{\rm rep}(\gamma)}{\ell}\mathcal E_{\rm cell}(f)}
       +\sqrt{8Bp_1(\gamma)\mathcal E_{\rm sl}(f)}.
 \label{f16v29:slab-rough}
\end{equation}
If also $B\ell A_re^{-c_0\gamma}\le1/4$, then
\begin{equation}
           \mathcal E_{\rm sl}^{\rm good}(f)
                    \le K_{\rm nat}(\gamma)\mathcal E_{\rm cell}(f).
 \label{f16v29:slab-good}
\end{equation}
The first inequality includes arbitrary rough collars. The second is integrated
only on its declared good collars, using their actual posterior weights.
\end{proposition}
\begin{proof}
Apply V28's nested inequality to $D=I_s$, $S=T_c$ and $H=H_c^{(1)}$ for each
$c\in C_s$, with conditional probability \eqref{f16v29:collar-smallness}.
The first sum of local energies is at most
$2K_{\rm rep}\mathcal E_{\rm cell}/\ell$ by
\eqref{f16v29:slab-collar-counts}. For the second, keep the local union until
using $I_s\cup T_c\subseteq J_s$ and drop only its nonnegative subtracted
energy. Its upper bound is
\[
 p_1\frac1\ell\sum_s |C_s|\mathcal E_{J_s}(f)
       \le 2Bp_1\frac1\ell\sum_s\mathcal E_{J_s}(f)
       \le8Bp_1\mathcal E_{\rm sl}(f).
\]
Direct-sum Minkowski proves \eqref{f16v29:slab-rough}. No full variance or
number of time slices enters this estimate.

On a good collar, the conditional inside $I_s$ satisfies
\eqref{f16v29:angular-collar} and has at most $B\ell$ cells. Apply
Theorem~\ref{f16v29:conditional-gap} at that exterior. Multiply its inequality
by the exterior indicator and integrate. Conditioning further on other cells
in the slab restores the original full single-cell conditional on the right.
Dropping its nonnegative gate gives an upper bound by those actual cell
energies. Each cell is in exactly $\ell$ indexed slabs, so summing their
rates gives \eqref{f16v29:slab-good}. The restriction to the good boundary was
not removed from a conditional theorem without payment: the preceding rough
form is exactly the complementary contribution.
\end{proof}

\begin{theorem}[An original-clock gap with no cross-section in its exponential cost]
\label{f16v29:main}
There exist fixed $\beta_r,c_r,C_r>0$ and $\gamma_r<\infty$ such that for every
\begin{equation}
 \gamma\ge\gamma_r,\qquad B=m^3\le e^{\beta_r\gamma},\quad m\ge4,\quad n\ge1,
                  \quad\max_{i,e}|y_{i,e}|\le r,
 \label{f16v29:width}
\end{equation}
the full native whole-cell sampler and the original unweighted rate-one
balanced sampler satisfy
\begin{equation}
 \boxed{\qquad
   \lambda_{\rm cell}\ge c_re^{-C_r\gamma},\qquad
   \widehat\lambda_P\ge c_re^{-C_r\gamma}.
 \qquad}
 \label{f16v29:main-gap}
\end{equation}
The constants do not depend on $B$ or $n$ within \eqref{f16v29:width}.
All native $L^2$ observables, complete compact coordinates, and original
endpoint factors are included. The rate of every elementary balanced update
remains one.
\end{theorem}
\begin{proof}
Let $\beta_{\rm ch}>0$ be the width exponent of V22's actual exponential-width
channel. Choose once
\[
        0<\beta_r<\min\{\beta_{\rm ch},c_0/4,c_1/4\}.
\]
Its buffer obeys $b\le C_r\gamma$. Increase the fixed threshold so that
$b\ge2v+2$, $B\ell A_re^{-c_0\gamma}\le1/4$, and
$16Bp_1(\gamma)\le1/2$ throughout \eqref{f16v29:width}. These follow from
strictly positive exponential margins; the factors $B$ and $B\ell$ have not
been ignored.

Square \eqref{f16v29:slab-rough} with $(x+y)^2\le2x^2+2y^2$, add
\eqref{f16v29:slab-good}, and rearrange:
\begin{equation}
 (1-16Bp_1)\mathcal E_{\rm sl}(f)
 \le\left(K_{\rm nat}+\frac{4K_{\rm rep}}\ell\right)
                                    \mathcal E_{\rm cell}(f).
 \label{f16v29:closed-slab-comparison}
\end{equation}
Both energy coefficients are at most $C_re^{C_r\gamma}$. V24's independently
proved slab gap is at least $1/2$ in its unit-coverage clock, hence
\[
 \operatorname{Var}_P f\le2\mathcal E_{\rm sl}(f)
                  \le C_re^{C_r\gamma}\mathcal E_{\rm cell}(f).
\]
This proves the first inequality of \eqref{f16v29:main-gap} without
conditioning away rough boundary rows or using a presumed native gap.

To return to the original balanced coordinates, one free unmoving cell
$(i,b)$ changes its five-chord group and, when present, the seven moving
ports at times $i-1$ and $i$ in that cube. Thus at most fifteen elementary
balanced groups are needed. The exact frame identity gives the same
exterior-sigma-algebra inclusion as V24. On this fixed support the original
product references and bounded classical words give a conditional variance
comparison of cost $C_re^{C_r\gamma}$. Each elementary chord occurs through
one cell and each moving port through at most two. Summing therefore gives
\begin{equation}
                    \mathcal L_{\rm cell}
                       \preceq C_re^{C_r\gamma}\widehat{\mathcal L}.
 \label{f16v29:cell-balanced}
\end{equation}
No fractional Haar factor is replaced by full Haar in this comparison.
Combining with the first inequality, and enlarging $C_r$, proves the second.
There is no hidden factor $1/[B(n+1)]$ in either generator. In particular,
\eqref{f16v29:main-gap} is not the old whole-slab density estimate with its
size deleted from the exponent; that estimate has been replaced by the
conditional gap, local rough repair, and overlap comparison proved above.
\end{proof}

The last theorem improves the \emph{scaling of the proved clock cost} inside
an explicit exponential-width regime. It does not assert that its unevaluated
constants are numerically optimal. At fixed coupling it allows only a bounded
set of spatial cross-sections. As $\gamma\to\infty$ its lower bound still
tends to zero; it is not a coupling-uniform gap or a calibrated physical mass.
V25's weaker all-finite-width baseline remains available for widths beyond
\eqref{f16v29:width}.

\subsection{A consequence for V28's actual retained generator}

\begin{proposition}[A paid spectral edge for the displayed retained dynamics]
\label{f16v29:retained}
Keep \emph{exactly} V28's $R_\gamma=\lceil h_\gamma^5\rceil$, gates, normalized
helper families, and generators $\mathcal L_{\rm aug},\mathcal L_{\rm keep}$.
For sufficiently large coupling in \eqref{f16v29:width},
\begin{equation}
 \boxed{\quad
    \lambda_{\rm keep}\ge R_\gamma^{-4}c_re^{-C_r\gamma},\qquad
    \lambda_{\rm aug}\ge\lambda_{\rm keep}.
 \quad}
 \label{f16v29:retained-gap}
\end{equation}
In particular both lower bounds can be written $c_r'e^{-C_r'\gamma}$ after
changing fixed constants. This is a statement in the width regime, not an
unrestricted spatial gap for either generator.
\end{proposition}
\begin{proof}
V28 proves, once its $\theta_\gamma<1$, the stronger form inequality
\[
 \mathcal L_{\rm keep}\succeq
           \mathcal L_R+(1-\theta_\gamma)
                         (\mathcal L_{\rm rep}+\mathcal L_+)
                 \succeq\mathcal L_R.
\]
For a cell $c\in D$, nested projections give $Q_c\preceq Q_D$. Thus
$\sum_{c\in D}Q_c\preceq |D|Q_D\preceq R_\gamma^4Q_D$. Sum the original
indexed unit-coverage rates to get
$\mathcal L_{\rm cell}\preceq R_\gamma^4\mathcal L_R$.
Theorem~\ref{f16v29:main} and variational gap comparison prove
\eqref{f16v29:retained-gap}. The polynomial logarithmic factor can be absorbed
into $e^{C_r'\gamma}$ at a fixed threshold.

This proof uses the full retained generator and V28's already established
rough-form absorption. It does not analyze only its good output rows, omit
gate switching from a path coupling, or switch off either helper family.
Detailed balance and the coverage clocks remain V28's. The new lower edge
comes from an independent ungated temporal-slab argument followed by form
comparison; it is not deduced circularly from relative closeness alone.
\end{proof}

\subsection{What has closed and what remains spatial}

The three comparisons have different jobs: fixed-depth angular repair bounds
the full conditional mixing cost by $e^{C_r\gamma}$; the overlapping-slab
argument pays the enlarged union energy; and nested rough-innovation
localization absorbs it into the same slab form. The final conversion to
balanced updates has fixed support. No duration factor or exponential cost
proportional to $B$ occurs in the resulting bound.

The argument still counts $B$ collar cells and $B\ell$ cells inside one slab.
Their exponentially small angular exceptions are paid only while
$Bp_1$ and $B\ell p_0$ are small. Arbitrary spatial volume at fixed coupling
would destroy this certificate. No common scale is taken to infinity with
those factors silently held fixed. Spatially, the analogue of
\eqref{f16v29:enlargement-domination} with useful, arbitrary-boundary constants
is not supplied here. Nor is a gap for a full conditional at every compact
surface field inferred from the controlled angular-collar theorem.

The remaining productive directions are therefore an unrestricted spatial
union-energy comparison or a direct spatial dependence estimate, and,
separately, a sharper fixed-support repair cost if a nonvanishing
weak-coupling sampler floor is desired. Another temporal power iteration
cannot remove the $B$-dependent admissibility conditions. Substitution of the
present $e^{-C_r\gamma}$ bound into a gap-only source inequality is not
advertised as an improvement over V23's direct $1/\gamma$ covariance bound.
All source definitions, $p_{\rm good}$, and prior conditional return matrices
are unchanged.

Every clock here is auxiliary. The fixed bridge window is retained. Root
Gauss coverage, unrestricted bridge integration, infrared contact coercivity,
comparison to the physical transfer at its own top, independent physical-time
calibration, and nontrivial interacting continuum reconstruction remain
separate obligations. No Yang--Mills physical mass gap follows from an
uncalibrated original-sampler lower bound.

\subsection{Diagnostics, preservation, and verification scope}

The accompanying diagnostic is
\begin{center}\small\path{verify_ym_f16_v29_angular_clock.py}.\end{center}
It checks finite conditional projections and restriction energies, the
variance-gluing algebra, Gaussian block Bochner identities, stopped-boundary
bookkeeping, exact interval unions and multiplicities, retained-endpoint
conditional angles, nonextensive slab absorption, the fifteen-group frame
support, and the original-rate conversions. The diagnostics are rational,
integer, or symbolic fixtures, not native compact-integral enclosures or
proof-assistant verification. They do not certify the inherited product-ball
domain or evaluate the constants $\eta_1,v_0,c_0,c_1,\beta_r,C_r$.

The package records the actual tests and any executed predecessor reruns,
source hashes, exact reconstruction, final compilation, and rendered-page
inspection. V28 is preserved except for the one title version and this exact
terminal insertion. The next companion checkpoint remains V30. The new
analytic deductions rely on the explicitly identified set-energy and
reflecting-domain inputs; the historical archive has not been independently
recertified.

\begin{firewall}
V29 proves full-conditional mixing on a declared angular collar and an
original-clock lower bound $c_re^{-C_r\gamma}$ for
$B\le e^{\beta_r\gamma}$, uniformly in duration. It also controls V28's
retained generator in that regime. These deductions retain every native
normalizer, compact output, endpoint factor, source, and elementary rate.
They do not give unrestricted spatial volume, a coupling-uniform sampler gap,
or a physical transfer or interacting continuum mass gap.
\end{firewall}
% END F16 V29 APPEND

% BEGIN F16 V30 APPEND -- 2026-09-18
% Insert immediately before the final document-ending command in V29.
% The predecessor, its native laws, sources, and elementary clock are preserved.

\clearpage
\section{V30: fixed-range block coercivity and spatial screening}
\label{f16v30:context}

V29's exponential elementary-clock cost enters when a fixed angular repair is
resolved into individual cell updates. It is not necessary to pay that cost
when the repair is itself an allowed, exactly normalized block update. We keep
such blocks at one fixed radius, chosen before coupling varies. A buffered-cover
variance inequality then has a coupling-independent constant. Combining it with
V29's independent temporal overlap argument gives a genuine positive gap for a
\emph{fixed-range} native block generator. No block radius or rate grows with
coupling, cross-section, or duration.

The spatial-width condition remains: the result holds for
$B\le e^{\beta_r\gamma}$, with a fixed $\beta_r>0$, and every duration. Within
that regime the new local gap also gives a finite-layer contraction and
exponential \emph{space--time} screening at every radius, without an algebraic
or nonperturbative distance floor. This yields the original coefficient
expansions in a weighted absolute space--time row norm. It is not an
unrestricted-spatial-volume theorem and it does not turn the block clock into
the original elementary or physical clock.

\subsection{V30 checkpoint and the precise dependency change}

The complete V29 appendix and audit were read. Its restricted product-ball
variance theorem, two angular repair scales, and retained-boundary slab
comparison are the analytic premises used here. V29's final original-clock gap
is \emph{not} an input to the new block gap. The V25 checkpoint and the supplied
f14 V100, f15 V100, and Grand Tensor V076 were included in targeted section and
text inventories. Empty semantic searches were not interpreted as absence of
prior work. Neither the whole archive nor unavailable companion programs are
represented as independently certified; the separate checkpoint records the
scope. The next scheduled checkpoint is V35.

V13 already identifies conditional projections as a local frustration-free
auxiliary parent and cites detectability. We reuse that construction. The
projection-product estimate below is the classical argument of
Anshu--Arad--Vidick,
\href{https://arxiv.org/abs/1602.01210}{\emph{A simple proof of the detectability
lemma and spectral gap amplification}}, Lemma 2 and Corollary 3. Their
statement and proof were examined. The proof uses only bounded projections,
so we give it on the present Hilbert space rather than assume a finite local
spin dimension. No new general detectability or Knabe theorem is claimed.
The new input that makes this principle quantitative here is a native gap for
blocks of fixed radius in an explicitly normalized clock.

The import memoranda's distinction between locality cost and perturbative
order remains important. The block-gap proof uses no higher-order expansion.
V16's moments and V17's entry expansion are used only afterward, to convert
spatial screening into the same source coefficients. The memoranda's final
physical-transfer arrow is not an analytic premise.

\subsection{A buffered cover of one conditional, without a microscopic comparison}

Keep $P=P_{\gamma,\mathbf y}$, $B=m^3$, $m\ge4$, $n\ge1$, and
$\max_{i,e}|y_{i,e}|\le r$. For a free set $\Omega$ let $\mu=P_\Omega^q$.
Retain V29's fixed radii $\eta_1<\eta_0$, the event
$G_\Omega=\bigcap_{c\in\Omega}\{e_c\le\eta_0^2\gamma\}$, its repairs
$S_c^\Omega$, and
\[
 p_0(\gamma)=\min\{1,A_re^{-c_0\gamma}\},\qquad
 K_{\rm box}=1+M/\lambda,\qquad K_{\rm buf}=2K_{\rm box}+4.
\]
For a free subset $A\subseteq\Omega$, conditional expectation and its
innovation under $\mu$ are denoted $E_A^\mu,Q_A^\mu$. Their form is
$\mathcal E_A^\mu(f)=\|Q_A^\mu f\|_\mu^2$. The complete actual collar is
required to obey
\begin{equation}
 |q_d|\le\eta_1\sqrt\gamma\quad(d\in\partial_{\mathcal G}\Omega),
 \qquad |\Omega|p_0(\gamma)\le\tfrac14.
 \label{f16v30:conditional-regime}
\end{equation}
There is no restriction on any free coordinate in $\mu$.

\begin{proposition}[Regrouping the restricted coordinates]
\label{f16v30:grouping}
Let $\pi=\mu(\cdot\mid G_\Omega)$ and let
$\Omega=V_1\sqcup\cdots\sqcup V_t$ be any partition into unions of complete
cells. If $V_j\subseteq A_j\subseteq\Omega$, then
\begin{equation}
 \operatorname{Var}_\pi f\le K_{\rm box}
                    \sum_{j=1}^t\mathcal E_{A_j}^\pi(f).
 \label{f16v30:restricted-cover}
\end{equation}
The constant is independent of the dimensions of the groups and their number.
\end{proposition}
\begin{proof}
Apply the proof of Proposition~\ref{f16v29:box-heatbath} with the reflecting
partial generators grouped as $\mathcal A_{V_j}=\sum_{c\in V_j}\mathcal A_c$.
The full Hessian-square still contains the sum of all within-group
Hessian-squares. Each principal measured-Hessian block is bounded above by
$MI$, while the full measured Hessian is nonnegative. Every cell face appears
in exactly one group, so the same nonnegative boundary-curvature sum cancels
in the full and grouped Bochner identities. Thus
\[
 \sum_j\|\mathcal A_{V_j}u\|_\pi^2
 \le\|\mathcal Au\|_\pi^2+M\langle u,\mathcal Au\rangle_\pi.
\]
The graph-core argument is exactly V29's inherited product reflecting
realization. With $u=\mathcal A^{-1}(f-\pi f)$,
$E_{V_j}^\pi\mathcal A_{V_j}u=0$ and Cauchy--Schwarz give the variance bound
with $V_j$ in place of $A_j$. Conditional-projection nesting gives
$\mathcal E_{V_j}^\pi\le\mathcal E_{A_j}^\pi$. Empty groups can be omitted.
No connectedness of a group, absolute diagonal dominance, or smooth-boundary
replacement for the product domain is required.
\end{proof}

\begin{theorem}[Complete native tensorization over a buffered cover]
\label{f16v30:buffered-cover}
Assume \eqref{f16v30:conditional-regime} and the inherited fixed-$r$ threshold.
Let $A_1,\ldots,A_t\subseteq\Omega$ admit a partition
$\Omega=V_1\sqcup\cdots\sqcup V_t$ such that
\begin{equation}
                 S_c^\Omega\subseteq A_j\quad(c\in V_j).
 \label{f16v30:buffer-cover-condition}
\end{equation}
Then, on the complete conditional space,
\begin{equation}
 \boxed{\quad
 \operatorname{Var}_\mu f\le K_{\rm buf}
                         \sum_{j=1}^t\mathcal E_{A_j}^\mu(f).
 \quad}
 \label{f16v30:buffered-cover-bound}
\end{equation}
No factor depending on coupling, the number of covering blocks, or the size
of $\Omega$ occurs in $K_{\rm buf}$. The two hypotheses in
\eqref{f16v30:conditional-regime} have not been removed.
\end{theorem}
\begin{proof}
Put $H_j=\bigcup_{c\in V_j}\{e_c>\eta_0^2\gamma\}$.
V29's \emph{uniform conditional} angular repair bound and
$S_c^\Omega\subseteq A_j$ imply, by the tower property,
\[
            \mu(H_j\mid\zeta_{\Omega\setminus A_j})
                          \le |V_j|p_0.
\]
This does not require independence of the repairs or their events. Decompose
a test $g$ into $Q_{A_j}^\mu g+E_{A_j}^\mu g$. The conditional projection
estimate from V27, the union bound, and direct-sum Minkowski give
\begin{equation}
 \|1_{G_\Omega^c}g\|_\mu
 \le\left(\sum_j\mathcal E_{A_j}^\mu(g)\right)^{1/2}
                       +\sqrt{|\Omega|p_0}\,\|g\|_\mu.
 \label{f16v30:block-exception}
\end{equation}
Crucially, each repair is already contained in an allowed update. We have not
estimated that update by the sum of elementary cell energies.

The exact restriction-energy comparison of V29 holds for any block $A_j$:
\[
 \mu(G_\Omega)\mathcal E_{A_j}^\pi(f)
                                    \le\mathcal E_{A_j}^\mu(f).
\]
To check its normalizer, condition on the exterior of $A_j$, restrict that
exterior to $G_{\Omega\setminus A_j}$, and retain
$\mu(G_{A_j}\mid\zeta_{\Omega\setminus A_j})$ multiplying the restricted
conditional variance. It is one nonnegative part of the full conditional
variance. The comparison uses the same product event, also when the $A_j$
overlap.

Set $g=f-\pi f$, $T=\|g\|_\mu$, and
$\mathcal F=\sum_j\mathcal E_{A_j}^\mu(f)$. Proposition~\ref{f16v30:grouping}
and this restriction comparison give $\|1_{G_\Omega}g\|_\mu^2\le
K_{\rm box}\mathcal F$. Equation~\eqref{f16v30:block-exception} then yields
\[
 T^2\le K_{\rm box}\mathcal F+
       (\sqrt{\mathcal F}+\sqrt{|\Omega|p_0}T)^2
 \le(K_{\rm box}+2)\mathcal F+2|\Omega|p_0T^2.
\]
Absorb the last term and use $\operatorname{Var}_\mu f\le T^2$.
Restriction of an $L^2(\mu)$ function lies in $L^2(\pi)$ because
$\mu(G_\Omega)>0$. Thus the proof includes constants, complex functions,
and tests supported on exceptional configurations.
\end{proof}

\begin{proposition}[A spatial union comparison, with the rough exterior retained]
\label{f16v30:spatial-union}
Suppose $U=A\cup B$ has a two-block cover satisfying
\eqref{f16v30:buffer-cover-condition}. On each admitted angular exterior of $U$,
provided $|U|p_0\le1/4$, its complete conditional obeys
\begin{equation}
                 Q_U\preceq K_{\rm buf}(Q_A+Q_B).
 \label{f16v30:good-union}
\end{equation}
For example, split an ordinary box along one coordinate into two overlapping
boxes with overlap of at least $2v_0+2$ cells. Assign each center to a side
containing its $v_0$-repair. This includes spatial splits, not just temporal
slabs.

Under the actual, unrestricted native law $P$, let $C=\partial_{\mathcal G}U$,
use V29's fixed-radius $T_c$ and $p_1=\min(1,A'_re^{-c_1\gamma})$, and let
$U^+\supseteq U\cup T_c$ for every $c\in C$. For any $t>0$,
\begin{equation}
 \begin{split}
 Q_U\preceq {}&K_{\rm buf}(Q_A+Q_B)
             +(1+t)\sum_{c\in C}Q_{T_c}\\
             &+(1+t^{-1})p_1|C|Q_{U^+}.
 \end{split}
 \label{f16v30:rough-union}
\end{equation}
This is an all-test-function form inequality, not an all-boundary version of
\eqref{f16v30:good-union}. The condition $|U|p_0\le1/4$ remains.
\end{proposition}
\begin{proof}
The first claim is Theorem~\ref{f16v30:buffered-cover} on the conditional fibre;
conditional variance on $U$ gives its operator interpretation. For the box
example choose an assignment plane in the middle of the overlap. Its distance
from the opposite side's end is at least $v_0$; clipping at the actual boundary
of $U$ removes rather than adds repair cells.

Let $M=1_{\cup_{c\in C}\{e_c>\eta_1^2\gamma\}}$. It is exterior to $U,A,B$.
Integrating \eqref{f16v30:good-union} on $M=0$ gives
$(1-M)Q_U\preceq K_{\rm buf}(Q_A+Q_B)$. For $MQ_U$, use the union bound and
V28's nested Young inequality on $(U,T_c)$, with the actual conditional
smallness $P(e_c>\eta_1^2\gamma\mid T_c^c)\le p_1$. It gives
\[
 MQ_U\preceq(1+t)\sum_{c\in C}Q_{T_c}
 +(1+t^{-1})p_1\sum_{c\in C}(E_{T_c}-E_{U\cup T_c}).
\]
Each excess projection is nonnegative and bounded above by $Q_{U^+}$.
This proves \eqref{f16v30:rough-union} without dropping the rough exterior
or assuming that overlapping conditional expectations commute.
\end{proof}

\subsection{One fixed-radius native generator and its fixed coverage clock}

Choose a fixed integer $v_*$ at least the radii of both V29 repair families
$S_c^\Omega$ and $T_c$. Enlarge it once for the fixed coordinate-range allowances.
For every ambient cell $c$, let $U_c$ be its coordinatewise box of radius $v_*$,
saturating short spatial cycles and clipping at natural time endpoints. Put
\begin{equation}
 \kappa_*=(2v_*+1)^4,\qquad
 \mathcal L_* =\kappa_*^{-1}\sum_c Q_{U_c},\qquad
 \mathcal E_*(f)=\kappa_*^{-1}\sum_c\mathcal E_{U_c}(f).
 \label{f16v30:fixed-generator}
\end{equation}
Each indexed update has rate $1/\kappa_*$. The index is the center, even if two
clipped boxes coincide. Each cell has coverage at most one and at least
$1/\kappa_*$. Both the radius and normalization are independent of coupling,
volume, and duration. A move resamples the \emph{whole} native $U_c$ conditional,
including its buffer and every compact field.

\begin{theorem}[A noncollapsing fixed-range auxiliary gap in the width regime]
\label{f16v30:fixed-gap}
There exist fixed $\beta_r>0$ and $\gamma_r<\infty$ such that, for
\begin{equation}
 \gamma\ge\gamma_r,\qquad B=m^3\le e^{\beta_r\gamma},\quad
 m\ge4,\quad n\ge1,\quad\max_{i,e}|y_{i,e}|\le r,
 \label{f16v30:width}
\end{equation}
the complete native generator \eqref{f16v30:fixed-generator} satisfies
\begin{equation}
 \boxed{\quad \mathcal L_*\succeq g_*(I-\mathbb E_P),\qquad
 g_*:=\frac{1}{4\kappa_*(K_{\rm buf}+4)}>0.\quad}
 \label{f16v30:fixed-gap-bound}
\end{equation}
The lower bound does not deteriorate with coupling, spatial width within
\eqref{f16v30:width}, or duration. It is not a gap claim at unrestricted width
or in the elementary balanced clock.
\end{theorem}
\begin{proof}
Use V29's actual temporal buffer $b=O_r(\gamma)$, $\ell=16b$, its unchanged
clipped slabs $I_s$, and its independently proved gap
$\operatorname{Var}_P f\le2\mathcal E_{\rm sl}(f)$. Retain V29's fixed
slab enlargement and its form bound
$\ell^{-1}\sum_s Q_{J_s}\preceq4\mathcal L_{\rm sl}$.
No elementary or fixed-block gap was used to establish either input.

On a good angular collar of $I_s$, apply
Theorem~\ref{f16v30:buffered-cover} with blocks $A_c=I_s\cap U_c$ and
assigned sets $V_c=\{c\}$, $c\in I_s$. Each contains $S_c^{I_s}$.
At this fixed exterior it gives conditional variance at most
$K_{\rm buf}\sum_{c\in I_s}\mathcal E_{I_s\cap U_c}^{\mu}(f)$,
provided $B\ell p_0\le1/4$. Integrating the good gate gives the corresponding
\emph{global} subset innovations; then $Q_{I_s\cap U_c}\preceq Q_{U_c}$.
Each center occurs in exactly $\ell$ indexed slabs, so
\begin{equation}
           \mathcal E_{\rm sl}^{\rm good}(f)
                        \le K_{\rm buf}\kappa_*\mathcal E_*(f).
 \label{f16v30:slab-good}
\end{equation}
The gate was dropped only from a nonnegative upper bound after conditioning.

Repeat the exact nested calculation in Proposition~\ref{f16v29:slab-absorption},
but keep $Q_{T_c}$ as a block energy instead of comparing it to cell energies.
Since $T_c\subseteq U_c$,
$\sum_c\mathcal E_{T_c}\le\kappa_*\mathcal E_*$. The unchanged collar counts
and the unchanged enlargement return give
\begin{equation}
 \sqrt{\mathcal E_{\rm sl}^{\rm bad}(f)}
 \le\sqrt{\frac{2\kappa_*}{\ell}\mathcal E_*(f)}
                          +\sqrt{8Bp_1\mathcal E_{\rm sl}(f)}.
 \label{f16v30:slab-bad}
\end{equation}
Choose $\beta_r<\min(\beta_{\rm ch},c_0/4,c_1/4)$ once. At a fixed threshold,
$B\ell p_0\le1/4$ and $16Bp_1\le1/2$ hold throughout
\eqref{f16v30:width}. Squaring \eqref{f16v30:slab-bad}, adding
\eqref{f16v30:slab-good}, and absorbing yields
\[
 (1-16Bp_1)\mathcal E_{\rm sl}(f)
 \le\kappa_*(K_{\rm buf}+4/\ell)\mathcal E_*(f).
\]
Thus $\operatorname{Var}_P f\le4\kappa_*(K_{\rm buf}+4)\mathcal E_*(f)$.
All constants used to choose $U_c$ preceded the coupling limit. No hidden
factor $1/[B(n+1)]$, exponentially large repair rate, or comparison of a full
slab with elementary updates occurs in this proof.
\end{proof}

\begin{proposition}[Locality and the elementary-clock boundary]
\label{f16v30:local-parent}
The generator $\mathcal L_*$ preserves $P$, is reversible, and has constants
as its kernel. In the exact square-root-density realization on the original
product reference it is a sum of bounded fixed-support projections with a
common null vector. Each summand fails to commute with at most $\mathfrak d_*$ others,
for a fixed finite $\mathfrak d_*$ independent of coupling and regulator.

In the original balanced clock one still has only the paid comparison
\begin{equation}
       \mathcal L_*\preceq C_re^{C_r\gamma}\widehat{\mathcal L}.
 \label{f16v30:elementary-clock}
\end{equation}
Consequently the new constant $g_*$ does not assert a coupling-independent
lower bound for $\widehat{\mathcal L}$.
\end{proposition}
\begin{proof}
Conditional projections preserve the exact measure. Every cell belongs to its
own $U_c$, so a function in the joint kernel is exterior-measurable for every
single cell. Equivalence with the finite product reference makes it constant.
For clarity, let $p_U(x\mid z)$ be the actual conditional density relative to
that product reference on $U$. Conjugating $E_U$ by multiplication by
$\sqrt{dP/d\nu}$ gives the fibre kernel
\[
 F(x,z)\longmapsto
 \sqrt{p_U(x\mid z)}\int\sqrt{p_U(y\mid z)}F(y,z)\,d\nu_U(y).
\]
The exterior marginal cancels. This kernel depends only on $U$ and its actual
interaction collar. The rank-one fibre projector fixes the global square-root state; its
complement annihilates that state and is a local positive projection. This is V13's auxiliary-parent
construction, now with fixed-radius blocks and a quantitative gap.

A fixed coordinate enlargement of $U_c$ contains its kernel support. The graph
of overlapping enlarged supports has uniformly bounded degree and admits a
fixed number of colors with disjoint supports in each color. Disjoint supports
commute, giving the asserted bound $\mathfrak d_*$ (take $\mathfrak d_*\ge1$).

For \eqref{f16v30:elementary-clock}, include all chord groups of $U_c$ and all
moving ports whose cell at either adjacent chord time lies in $U_c$. This
includes the preceding ports. The support has a fixed number of elementary
groups. Its complete classical oscillation is $C_r\gamma$ relative to the
original product factors, including fractional endpoint Haar powers. The fixed
conditional comparison and bounded overlap give \eqref{f16v30:elementary-clock}.
This factor, unlike the radius and block rates, can diverge with coupling.
\end{proof}

\subsection{A finite-layer contraction and all-radius spatial screening}

Color the enlarged operator supports just described using $s_*<\infty$ colors.
Put
\[
 \mathcal S=\left(\prod_{c\in C_{s_*}}E_{U_c}\right)\cdots
             \left(\prod_{c\in C_1}E_{U_c}\right),\qquad
 \Pi_P f=Pf.
\]
A color can be updated in parallel: its disjoint blocks have no interaction
between them after their one common exterior is fixed. All conditional laws
and the subsequent layers remain native. One sweep has a fixed number of
layers, not a single globally rate-one elementary update.

\begin{theorem}[Native sweep contraction and relative space--time clustering]
\label{f16v30:clustering}
Throughout \eqref{f16v30:width}, set
\begin{equation}
           q_*^2=\frac{\mathfrak d_*^2}{\mathfrak d_*^2+\kappa_*g_*}<1.
 \label{f16v30:sweep-rate}
\end{equation}
Then $\|\mathcal S-\Pi_P\|\le q_*$ and
$\|\mathcal S^k-\Pi_P\|\le q_*^k$. There are fixed $a_r,C_r>0$ such that
for arbitrary $F,G\in L^2(P)$ measurable in cell sets $A,B$,
\begin{equation}
 \boxed{\quad
 |\operatorname{Cov}_P(F,G)|
 \le C_re^{-a_r d_*(A,B)}
               \sqrt{\operatorname{Var}_PF\operatorname{Var}_PG}.
 \quad}
 \label{f16v30:relative-clustering}
\end{equation}
Here $d_*((i,b),(j,c))=|i-j|+d_{\mathbb T_m^3}(b,c)$ as in V19.
There is no cardinality prefactor, no distance-independent coupling floor,
and no boundary conditioning beyond the stated native law.
\end{theorem}
\begin{proof}
We give the standard projection-product argument to specify its Hilbert-space
scope. Order the blocks by the color sweep, write $Q_i=Q_{U_i}$, and put
$f_i=(I-Q_i)f_{i-1}$. Orthogonality gives
\[
 \sum_i\|Q_if_{i-1}\|^2=\|f\|^2-\|\mathcal Sf\|^2.
\]
At the end of update $i$, $Q_if_i=0$. A later commuting projection cannot
increase $\|Q_if_j\|$; a noncommuting update can increase it by at most
$\|Q_jf_{j-1}\|$. Cauchy--Schwarz over at most $\mathfrak d_*$ such updates and then
summing over $i$ therefore give
\[
 \sum_i\|Q_i\mathcal Sf\|^2
       \le \mathfrak d_*^2\{\|f\|^2-\|\mathcal Sf\|^2\}.
\]
For centered $f$, invariance keeps $\mathcal Sf$ centered. The gap of
$\sum_iQ_i=\kappa_*\mathcal L_*$ bounds the left side below by
$\kappa_*g_*\|\mathcal Sf\|^2$. Rearranging proves
\eqref{f16v30:sweep-rate}; constants are fixed by every factor, so powers follow.
This is the detectability estimate, not an additional gap assumption.

For clustering work in the exact product-reference representation, and write
$\Psi=\sqrt{dP/d\nu}$. Multiplication by a local bounded observable has its
ordinary cell support and commutes with the conjugation. Let $\upsilon_*$ denote a
fixed upper bound on the propagation distance of one full color sweep in
$d_*$. Such a bound exists since every layer has disjoint fixed-diameter
operator supports. In the expression
\[
              \langle\Psi,F^*\mathcal S^k G\Psi\rangle
\]
(remove the conjugation from the notation only), delete each projector outside
the expanding support of $F$ by commuting it to the left vacuum. After $k$
sweeps the remaining projectors lie within distance $k\upsilon_*$ of $A$. If
$k\upsilon_*<d_*(A,B)$, they commute with $G$ and can then be moved to the right
vacuum. The result is exactly $\langle\Psi,F^*G\Psi\rangle$.

Apply this identity to centered $F,G$. The sweep norm bounds their covariance
by $q_*^k\|F\Psi\|\|G\Psi\|$. Take, for example,
$k=\lfloor d_*(A,B)/(2\max(1,\upsilon_*))\rfloor$, with $k=0$ when the supports
meet. This gives \eqref{f16v30:relative-clustering} with
$a_r=-\log(q_*)/[2\max(1,\upsilon_*)]$ and an enlarged fixed prefactor. All
projector deletions use bounded local operators fixing \emph{the same} state.
No independence of disjoint native supports is asserted. Truncate local
$L^2$ functions and pass in $L^2(P)$ to remove boundedness. This retains their
supports and proves the all-$L^2$ assertion.
\end{proof}

The sweep and the resulting clustering do not refer to the physical Wilson
Hamiltonian. A sweep entails a number of local operations proportional to
the system size, organized in a fixed number of parallel layers. Nor does
\eqref{f16v30:relative-clustering} assert a bound after arbitrary new pinning:
that operation changes the native probability and its generator.

\subsection{Original source coefficients in a weighted space--time row norm}

Keep V7's primitive errors, deterministic synthesis $T=(I,\mathsf P^{\mathsf T})$,
and covariance $\Gamma_\gamma=TK_\gamma T^{\mathsf T}$. The primitive carriers
have fixed radius. V16 bounds their variances by $C_r\gamma^{-1}$. Applying
\eqref{f16v30:relative-clustering} and summing the exponentially localized
synthesis rows gives, at a fixed exponent $a_0>0$,
\begin{equation}
 |(K_\gamma)_{uv}|\le C_r\gamma^{-1}e^{-a_0d_*(u,v)},\qquad
 |(\Gamma_\gamma)_{\alpha\beta}|
                    \le C_r\gamma^{-1}e^{-a_0d_*(\alpha,\beta)}.
 \label{f16v30:native-spatial-envelope}
\end{equation}
Anchors are used in this notation; their fixed carrier allowances are included
in $C_r$. Decrease $a_0$ below the original synthesis exponent. The common
source is not falsely assigned a finite spatial carrier.

For a scalar matrix in canonical source coordinates define
\[
 \|A\|_{\mathrm{st},a}
       =\sup_\alpha\sum_\beta
          e^{a d_*(\alpha,\beta)}|A_{\alpha\beta}|.
\]
This is an absolute \emph{space--time} row norm, not a trace norm or the earlier
temporal norm of spatial operator blocks.

\begin{theorem}[Full space--time row asymptotics in the same width regime]
\label{f16v30:source-expansion}
There is a fixed $a>0$, independent of coupling and expansion order, such that
for every fixed $J\ge2$, at a sufficiently large $\gamma_{J,r}$ in
\eqref{f16v30:width},
\begin{equation}
 \boxed{
 \begin{aligned}
 &\left\|\Gamma_\gamma-\sum_{j=2}^J\gamma^{-j/2}G_j\right\|_{\mathrm{st},a}
 \\
 &\quad+
 \left\|\mathsf M_\gamma-\sum_{j=2}^J\gamma^{-j/2}M_j\right\|_{\mathrm{st},a}
                   \le C_{J,r}\gamma^{-(J+1)/2}.
 \end{aligned}}
 \label{f16v30:source-row}
\end{equation}
The coefficients are precisely V10's normalized coefficients. In particular
both native norms are $O_r(\gamma^{-1})$, and their leading-coefficient errors
are $O_r(\gamma^{-3/2})$. All rows and spatial displacements are included,
for every duration, with no source-count normalization.
\end{theorem}
\begin{proof}
Choose $a>0$ so small that \eqref{f16v30:native-spatial-envelope}, polynomial
four-dimensional ball growth, and the coefficient locality of V10 give
\[
 \|\Gamma_\gamma\|_{\mathrm{st},4a}\le C_r\varepsilon^2,
 \qquad \|G_j\|_{\mathrm{st},4a}\le C_{j,r},\qquad
 \varepsilon=\gamma^{-1/2}.
\]
For the chosen $J$ put
$D=\lceil(J+2)(3a)^{-1}\log(\varepsilon^{-1})\rceil$ and $K=2J+6$.
On pairs with distance at most $D$, use V17's order-$K$ entry expansion.
At most $C(1+D)^4$ source labels occur in each row there, so its weighted
remainder is at most
\[
 C_{K,r}(1+D)^4 e^{aD}\varepsilon^{K+1}
 \le C_{J,r}(1+\log\varepsilon^{-1})^4
                       \varepsilon^{K+1-(J+2)/3}
 =O_{J,r}(\varepsilon^{J+1}).
\]
The extra coefficients of orders $J+1$ through $K$ are bounded using their
whole row norms. On the complement use genuine native and coefficient
exponential rows: the weighted tail is at most
$C_{J,r}\varepsilon^2e^{-3aD}=O_{J,r}(\varepsilon^{J+1})$.
There is no flat remainder to sum and no full spatial slice count.
The exact identity $\mathsf M_\gamma=-\operatorname{Off}_{>1}\Gamma_\gamma$
and the same identity for $M_j$ remove entries and cannot increase this
absolute norm. They are not claimed to preserve positivity.
\end{proof}

\begin{proposition}[All-radius screening of the same shell return]
\label{f16v30:shell}
For V19's unchanged centered primitive messages
$m_{u,R}=\mathbb E_P[F_u\mid\zeta_{I_{u,R}^c}]$, there are fixed $a',C_r>0$
such that, throughout \eqref{f16v30:width} and all $R\ge R_0$,
\begin{equation}
 \sup_u\|m_{u,R}\|_2\le C_r\gamma^{-1/2}e^{-a'R},\qquad
 \|\mathsf R^{\rm sh}_{\gamma,R}\|_{\rm op}
                                  \le C_r\gamma^{-1}e^{-a'R}.
 \label{f16v30:shell-bound}
\end{equation}
A ball covering the whole system has zero centered message. No spatial-torus
diameter must first be crossed and no coupling-dependent floor remains.
\end{proposition}
\begin{proof}
Dualize \eqref{f16v30:relative-clustering} over centered exterior $L^2$ tests.
Their supports are separated from the primitive carrier by $R$ minus a fixed
allowance. This proves the first bound, using V16's variance estimate.
The message is measurable on its finite interaction collar, contained in the
ball of radius $R+\rho$. For two anchors farther apart than $2R+2\rho$, apply
\eqref{f16v30:relative-clustering} again to their messages; for closer anchors
use their small variances. Polynomial ball growth then bounds the absolute
row sum of their Gram by
$C_r\gamma^{-1}(1+R)^4e^{-2a'R}$ after initially choosing a common smaller
exponent. Absorb the polynomial by a final decrease of $a'$. The original
bounded synthesis $T$ transfers its operator bound to
$\mathsf R^{\rm sh}_{\gamma,R}=T\operatorname{Cov}(m_R)T^{\mathsf T}$.
This is V19's object, not V12's different common-exterior posterior.
\end{proof}

\subsection{Checkpoint conclusions and the unresolved extrapolation}

The coupling loss and the width restriction are different issues. The first
is removed for one fixed-range auxiliary block generator by retaining its
repairs as actual moves. A microscopic comparison still costs
$C_re^{C_r\gamma}$. The second survives: the gap proof used both
$B\ell p_0\le1/4$ and $16Bp_1\le1/2$, as well as V22's complete-slice
channel. None holds for arbitrary spatial width at fixed coupling merely by
choosing a larger constant in the final estimate.

The native spatial union estimate \eqref{f16v30:good-union} advances the local
overlap question on a declared outer collar. Its all-test extension
\eqref{f16v30:rough-union} retains a repair form and the larger-union return.
A uniform spatial iteration must control those terms, or establish an
all-boundary local finite-size criterion. The new periodic/cylindrical gap
cannot be inserted as the gap of such an arbitrarily pinned patch: the
conditional parent changes at its boundary. This is the next locality task,
not a new perturbative coefficient or another temporal blocking step.

A simple clock check prevents a different overstatement. For independent
pairs with law $(1+\theta XY)/4$, the original two single-spin rate-one updates
have gap $1-\theta$, while refreshing the entire pair at rate $1/2$ has gap
$1/2$ and unit-or-smaller coverage. The latter does not keep the former away
from zero as $\theta\uparrow1$. This generic fixture does not establish an
obstruction for the native model; it shows why
\eqref{f16v30:elementary-clock} remains a separate obligation even with fixed
block size. Likewise the colored sweep has a declared parallel-layer clock,
not physical Euclidean time.

The supplied memoranda suggested compatible local constraints and a downstream
physical transfer comparison. We now have a quantitative fixed-range native
parent in the stated regime, and a corresponding spatial source norm; neither
is the physical Wilson logarithmic transfer. The retained bridge window,
physical Gauss reduction, infrared contact coercivity, same-top physical
normalization, independent physical-time calibration, and nontrivial interacting
continuum reconstruction remain separate. The V30 checkpoint records the
current dependency map and the limits of the archive review.

\subsection{Diagnostics and predecessor preservation}

The new diagnostic is
\begin{center}\small\path{verify_ym_f16_v30_fixed_block_screening.py}.\end{center}
It tests buffered-cover conditioning, exact restriction energies, noncommuting
conditional projections, fixed-radius coverage, temporal absorption constants,
projection-product loss, exact local square-root kernels, and the separation
between a block and an elementary clock. Finite rational and symbolic fixtures
are not native compact-integral enclosures or proof-assistant verification.
The new diagnostic passed 206 exact assertions. All 26 available V4--V29
diagnostics were rerun unchanged and passed. Their actual reports and hashes
are included separately from the analytic arguments.

Only one literal title version and this terminal insertion differ from V29.
Deleting the insertion and reversing the title change recovers V29 byte for
byte. The compilation and render records refer to the final files. No earlier
proof, source, probability, or reference is silently amended. The present
proofs explicitly depend on the inherited angular preparation and reflecting
product-domain statements; the entire lineage has not been independently
recertified. The next companion checkpoint is V35.

\begin{firewall}
V30 proves a coupling-independent gap for an explicitly normalized fixed-range
native block sampler in a specified exponential spatial-width regime, uniformly
in duration. It proves a controlled-collar buffered-cover and spatial-union
inequality, with an explicit all-test rough-exterior remainder. A standard
projection-product argument then gives a finite-layer contraction, relative
space--time clustering at all radii, and the original source expansions in a
weighted absolute space--time row norm in that regime. It does not prove a
coupling-independent elementary gap, unrestricted spatial width at fixed
coupling, an arbitrary-pinning finite-size criterion, a physical transfer gap,
or an interacting continuum mass gap. All block operations use the same
complete native conditionals and retain their actual rates and normalizers.
\end{firewall}
% END F16 V30 APPEND
% BEGIN F16 V31 APPEND -- 2026-09-18
% Insert immediately before the final document-ending command in V30.
% No predecessor proof, native measure, source, or elementary rate is changed.

\clearpage
\section{V31: the complete retained-bridge action and an in-window conditional gap}
\label{f16v31:context}

V30 proves spatial screening under the complete fixed-bridge fast law, in its
specified width regime. It does not prove that arbitrarily pinned fast patches
have the same gap. Rather than use that unavailable implication, we extract a
different interface that the new screening does support: the \emph{complete}
retained-bridge action, not only the common-source memory submatrix. Its Gaussian
part is retained exactly. The nonlinear correction has a small, exponentially
summable Hessian and an exact finite-volume interaction representation. At a
sufficiently small fixed bridge window this yields a conditional-resampling gap
for the actual interior bridge-path weight, uniformly under further bridge
pinning \emph{within that window}.

Two qualifications are essential. The source space below has all
$36B(n+1)$ real bridge coordinates, rather than only the common-source bank.
Its leading Gaussian action is generally not small. Also, a bridge-window path
probability is not the unrestricted physical path law, and pinning a retained
bridge is not pinning a fast coordinate. Neither the spatial-width restriction
nor the physical-transfer obligations are removed.

\subsection{Dependencies, normalization, and non-duplication}

The full V30 appendix, analytic audit, and checkpoint were read. Its
Theorem~\ref{f16v30:clustering} is used as an inherited analytic premise, not
independently recertified here. V16 supplies native local exponential moments;
V6 supplies the complete-chart ordinary score certificate and the fixed-domain
normalization. The Gaussian precision and source matrices are those of V4--V5.
The new leading-order comparison below does not require an infinite series,
V17's higher-order remainder, or a new native gap assumption.

The supplied f14 V40 already constructs interaction germs by an ordered,
two-radius telescope. Its activation changes which fast variables are active,
and its conclusion concerns a different restricted radial remainder and its
gauge-orbit Hessian. That general telescoping principle is credited. Here all
fast variables stay integrated; only retained bridge inputs outside a declared
carrier are set to zero when defining a coefficient function. The representation
is for the exact complete fast integral on the whole specified bridge window,
not an assertion of new gauge-invariant germs or an infinite-volume potential.
The source and archive searches were targeted, not exhaustive novelty checks.

For classical context, the paper by Barbieri, G\'omez, Marcus, Meyerovitch
and Taati, \href{https://arxiv.org/abs/2001.03880}{\emph{Gibbsian representations
of continuous specifications}}, distinguishes
norm-summable interaction representations from stronger symmetry requirements.
Its abstract and model scope were checked. Its symbolic-configuration
theorem is not invoked for the present continuous variables. All required
finite telescopes and estimates are proved below. The Dobrushin conclusion uses
the already proved finite-product criterion of V24, with its rates explicit.
The next scheduled companion checkpoint remains V35.

Let $\mathcal V=\{(i,e):0\le i\le n,\ e\text{ a bridge}\}$; each $y_u$ has
three real components. There are $12B$ bridges per slice. Assign a bridge to its
source cube; the time plus torus distance of these anchors is denoted $d(u,v)$.
Distinct bridge types can have the same anchor. This bounded multiplicity only
changes fixed constants, and every anchor ball has at most $C(1+L)^4$ labels.
The same convention, with its bounded component multiplicities, is used for
fast scalar indices. All original word carriers have bounded diameter.

For a fixed $r>0$, write
\[
 \mathcal Y_r=\prod_{u\in\mathcal V}\{y_u\in\mathbb R^3:|y_u|\le r\},
 \qquad \varepsilon=\gamma^{-1/2}.
\]
Unless otherwise stated, every result is in V30's regime
\begin{equation}
 \gamma\ge\gamma_r,\qquad B=m^3\le e^{\beta_r\gamma},\quad m\ge4,
 \quad n\ge1,\quad \mathbf y\in\mathcal Y_r.
 \label{f16v31:regime}
\end{equation}
Thresholds may be increased by fixed-$r$ requirements. No estimate below permits
$B\to\infty$ at one fixed coupling outside this regime.

Use the unchanged classical unmoving potential $U_\gamma$ and fast Haar factor
$J^f_\gamma$. Define
\begin{equation}
 \begin{split}
 \mathcal F_\gamma(\mathbf y)
   &=-\log\frac{Z_{\gamma,\mathrm{full}}(\mathbf y)}
                       {Z_{\gamma,\mathrm{full}}(0)},\\
 \mathcal F_0(\mathbf y)&=-\log\frac{Z_0(\mathbf y)}{Z_0(0)},\\
 \mathcal R_\gamma&=\mathcal F_\gamma-\mathcal F_0.
 \end{split}
 \label{f16v31:actions}
\end{equation}
The normalization subtracts one exact value; it does not replace a native
normalizer by a Gaussian one. In the original fast Euclidean coordinates,
write the already defined quadratic form as
\begin{equation}
 \begin{gathered}
 U_0(\zeta;\mathbf y)=\tfrac12\zeta^{\mathsf T}H\zeta
       +\zeta^{\mathsf T}L\mathbf y+\tfrac12\mathbf y^{\mathsf T}B_0\mathbf y,\\
 C=H^{-1},\qquad m=-CL\mathbf y,\qquad K=B_0-L^{\mathsf T}CL.
 \end{gathered}
 \label{f16v31:Gaussian-blocks}
\end{equation}
These definitions are by differentiation of V4's exact quadratic reference,
transported by V5's triangular change of variables. They introduce no fitted
matrix. Gaussian integration gives $\mathcal F_0=\tfrac12\mathbf y^{\mathsf T}K
\mathbf y$; its determinant cancels only in this same-Gaussian value ratio.
$H,L,B_0$ have fixed finite range and bounded absolute rows and columns.
$H$ has a uniform positive lower bound; $C$ has exponentially summable rows.
No positive lower bound on $K$ is assumed.

\subsection{An exact covariance identity with a locally paid chart defect}

Put $V_\gamma=U_\gamma-\log J^f_\gamma$. Its derivative near a principal cut
must not be used in an uncut integration by parts. Choose a smooth radial
function $\chi$ on $\mathbb R^3$ equal to one for $|x|\le1/4$, zero for
$|x|\ge1/2$, and between zero and one elsewhere. For a fast scalar index $j$,
let $g(j)$ be its original three-coordinate group and set
\[
 \chi_j(\zeta)=\chi(\zeta_{g(j)}/\sqrt\gamma),\qquad
 \xi=\zeta-m.
\]
The same cutoff is used for the three components of one group. Define
\begin{equation}
 d_{\gamma,j}
   =\chi_j\partial_jV_\gamma-\partial_j\chi_j-(H\xi)_j,
 \qquad
 \Theta_\gamma=\operatorname{diag}_j\mathbb E_P(1-\chi_j),
 \quad \Sigma_\gamma=\operatorname{Cov}_P(\zeta).
 \label{f16v31:drift-defect}
\end{equation}
The product $\chi_j\partial_jV_\gamma$ is set to zero off its interior support.
Here $P=P_{\gamma,\mathbf y}$ is the same complete native fast law. These
$d_{\gamma,j}$ are auxiliary integration-by-parts defects, not a redefinition
of V7's Haar-corrected bridge scores.

\begin{proposition}[Complete-law Gaussian identities with the chart term retained]
\label{f16v31:Stein-identity}
At every finite regulator in the stated regime,
\begin{equation}
 \begin{split}
 H(\mathbb E_P\zeta-m)&=-\mathbb E_P d_\gamma,\\
 H\Sigma_\gamma&=I-\Theta_\gamma
                     -\operatorname{Cov}_P(d_\gamma,\zeta).
 \end{split}
 \label{f16v31:Stein-exact}
\end{equation}
Each $d_{\gamma,j}$ has a fixed local fast carrier and
\begin{equation}
 \|d_{\gamma,j}\|_{L^2(P)}\le C_r\varepsilon,
 \qquad 0\le(\Theta_\gamma)_{jj}\le C_re^{-c_r\gamma}.
 \label{f16v31:drift-bounds}
\end{equation}
The bounds are uniform in source history, position, volume, and duration in
\eqref{f16v31:regime}.
\end{proposition}
\begin{proof}
Integrate the derivative of $\chi_j f e^{-V_\gamma}$ in its own group.
The cutoff is compactly supported inside that group's principal ball, so its
boundary term is zero. All other groups remain on their full charts. For
$f=1$ or $f=\zeta_k-\mathbb E_P\zeta_k$ this gives
\[
 \mathbb E_P[(\chi_j\partial_jV_\gamma-\partial_j\chi_j)f]
                       =\mathbb E_P[\chi_j\partial_j f].
\]
Substitute \eqref{f16v31:drift-defect}. Centering in the second choice gives
$H\Sigma_\gamma+\operatorname{Cov}_P(d_\gamma,\zeta)
=\operatorname{diag}(\mathbb E_P\chi_j)$. This proves both identities with
their signs. No flat derivative crosses a logarithm cut.

For the estimates, split
\[
 d_{\gamma,j}=\chi_j(\partial_jV_\gamma-\partial_jU_0)
              +(\chi_j-1)\partial_jU_0-\partial_j\chi_j.
\]
The classical part of the first term is bounded by
$C_r\varepsilon(1+|\zeta_T|^2)$ on a fixed local carrier $T$.
Indeed every word is a fixed ordered product of exponentials of fixed linear
forms, with its original $\gamma$ prefactor. The real third-derivative bound
from V6 applies to a fast as well as a bridge derivative. The endpoint
quadratic is exact in $U_0$ and cancels. On the cutoff support the measured
Haar derivative is bounded by $C\gamma^{-1}|\zeta_{g(j)}|$; the endpoint
power $1/2$ only decreases this bound. There is no intermediate principal
logarithm to differentiate in a word. V16's native fourth moments therefore
bound the first term in $L^2$ by $C_r\varepsilon$.

The other terms are supported where $|\zeta_{g(j)}|\ge\sqrt\gamma/4$.
V16's one-cell exponential moment gives exponentially small probability.
The linear polynomial $\partial_jU_0=(H\xi)_j$ has a bounded local $L^4$ norm,
and $|\partial_j\chi_j|\le C\varepsilon$. Cauchy--Schwarz pays their $L^2$
tails without assuming that a rare event is independent of its polynomial.
The same probability bound controls $\Theta_\gamma$. All differentiated
quantities have fixed carriers; the Gaussian mean is a deterministic bounded
function of the prescribed history. Haar factors and their powers have been
retained in the integration identity.
\end{proof}

For a possibly rectangular scalar matrix define a weighted absolute-row norm
using its two anchor families. All row and column bounds asserted below hold
with the same fixed exponential weight, after a decrease if needed. Supremum
over $\mathbf y$ may be placed \emph{inside} each absolute entry: the local
estimates give a common pointwise exponential envelope first.

\begin{theorem}[A summable native mean and covariance comparison]
\label{f16v31:fast-moments}
There is $a>0$ such that
\begin{equation}
 \sup_{\mathbf y\in\mathcal Y_r}\|\mathbb E_P\zeta-m\|_\infty\le C_r\varepsilon,
 \quad
 \sup_j\sum_k e^{a d(j,k)}
       \sup_{\mathbf y\in\mathcal Y_r}|(\Sigma_\gamma-C)_{jk}|
                                      \le C_r\varepsilon.
 \label{f16v31:fast-covariance}
\end{equation}
The same bound holds for weighted columns. This is a consequence of V30's
complete-law screening, not a new assumption that the native law is Gaussian.
\end{theorem}
\begin{proof}
V16 bounds the variance of each $\zeta_k$ uniformly. V30's relative clustering
applies to the fixed local carrier of $d_{\gamma,j}$ and to $\zeta_k$.
Using \eqref{f16v31:drift-bounds}, it gives
\[
 |\operatorname{Cov}_P(d_{\gamma,j},\zeta_k)|
                       \le C_r\varepsilon e^{-a_0d(j,k)}.
\]
Carrier allowances are absorbed into $C_r$. Polynomial graph growth makes
this rectangular covariance summable in both directions at a smaller weight.
From \eqref{f16v31:Stein-exact},
\begin{equation}
 \Sigma_\gamma-C=-C\Theta_\gamma
                    -C\operatorname{Cov}_P(d_\gamma,\zeta).
 \label{f16v31:covariance-resolvent}
\end{equation}
The original $C$ has summable weighted rows and columns. Matrix multiplication
and the triangle inequality for anchor distance prove the covariance bound.
The first line of \eqref{f16v31:Stein-exact} and the absolute rows of $C$
prove the mean bound. Products of the exponential envelopes justify moving
the supremum inside the sums. The inverse used here is only $H^{-1}$; no
inverse of an unproved native generator has been introduced.
\end{proof}

\subsection{The complete bridge Hessian, with its direct term}

For bridge scalar marks $\alpha,\beta$ put
\[
 J_{\gamma,\alpha}=\partial_\alpha U_\gamma,
 \qquad B_{\gamma,\alpha\beta}=\partial_\alpha\partial_\beta U_\gamma.
\]
These are ordinary coordinate derivatives at fixed unmoving fast variables.
$J^f_\gamma$ is independent of the bridge history. Exact differentiation of
the positive compact integral therefore gives
\begin{equation}
 \partial_\alpha\partial_\beta\mathcal F_\gamma
 =\mathbb E_P B_{\gamma,\alpha\beta}
        -\operatorname{Cov}_P(J_{\gamma,\alpha},J_{\gamma,\beta}).
 \label{f16v31:Hessian-exact}
\end{equation}
This is the existing normalized Hessian identity, now used on \emph{all}
bridge directions and at all separations. The direct term cannot be dropped
on equal-time or neighboring source carriers.

For a $3$-by-$3$ bridge-block matrix field $M(\mathbf y)$ define
\begin{equation}
 \mathfrak N_{a,r}(M)=\sup_u\sum_v e^{a d(u,v)}
                    \sup_{\mathbf y\in\mathcal Y_r}\|M_{uv}(\mathbf y)\|_{\rm op}.
 \label{f16v31:Hessian-norm}
\end{equation}

\begin{theorem}[The full native bridge action is a small Hessian perturbation]
\label{f16v31:action-Hessian}
For a fixed $a>0$ and $C_r<\infty$ in \eqref{f16v31:regime},
\begin{equation}
 \boxed{\qquad
 \mathfrak N_{a,r}(D_y^2\mathcal F_\gamma-K)
              =\mathfrak N_{a,r}(D_y^2\mathcal R_\gamma)
              \le C_r\gamma^{-1/2}.
 \qquad}
 \label{f16v31:action-smallness}
\end{equation}
Moreover $\mathcal F_\gamma(0)=\mathcal R_\gamma(0)=0$ and
$D_y\mathcal F_\gamma(0)=D_y\mathcal R_\gamma(0)=0$. In particular,
\begin{equation}
 |\mathcal R_\gamma(\mathbf y)|
       \le C_r\gamma^{-1/2}\sum_u|y_u|^2,
 \qquad
 \sup_u|\partial_{y_u}\mathcal R_\gamma(\mathbf y)|
       \le C_r r\gamma^{-1/2}.
 \label{f16v31:action-size}
\end{equation}
No derivative of an approximate pressure is used. The matrix $K$ includes
all Gaussian direct and returned terms and is not replaced by the common-mode
memory coefficient $G_2$.
\end{theorem}
\begin{proof}
The unchanged V6 certificate gives
$J_\gamma=L^{\mathsf T}\zeta+B_0\mathbf y+\delta J_\gamma$, where each
$\delta J_{\gamma,\alpha}$ has fixed carrier and $L^2(P)$ norm at most
$C_r\varepsilon$. Its variance obeys the same bound. All rows and columns of
$L$ have bounded finite support. Expand the covariance in
\eqref{f16v31:Hessian-exact} using this identity. The linear--linear part
differs from $L^{\mathsf T}CL$ by
$L^{\mathsf T}(\Sigma_\gamma-C)L$, bounded by
Theorem~\ref{f16v31:fast-moments}. Each linear--remainder covariance has a
pointwise bound $C_r\varepsilon e^{-a_0d(\alpha,\beta)}$ by V30's clustering.
The remainder--remainder term has the stronger $C_r\varepsilon^2$ prefactor.
Their weighted rows and columns are summable uniformly over $\mathcal Y_r$.

If no classical word contains both marks, $B_{\gamma,\alpha\beta}$ and the
corresponding entry of $B_0$ vanish. Otherwise the same global real
third-derivative certificate yields
$|B_{\gamma,\alpha\beta}-(B_0)_{\alpha\beta}|
\le C_r\varepsilon(1+|\zeta_T|)$ on a fixed carrier. Native moments bound
its expectation; only a fixed number of such entries occur per row.
Consequently the entire Hessian is
$B_0-L^{\mathsf T}CL+O(\varepsilon)$ in the stated norm. Conversion from
scalar colors to $3$-by-$3$ operator blocks costs only a fixed factor.

Global simultaneous conjugation of every group preserves the exact fast
integral, its complete charts, Haar powers, Wilson words and color-isotropic
endpoint amplitudes. It sends every $y_u$ through the same adjoint rotation.
At $\mathbf y=0$, each of the separate three-component gradient vectors must
therefore be fixed by every rotation in $\mathrm{SO}(3)$, and hence is zero.
The Gaussian action has the same property. This argument asserts only global
conjugation, not gauge invariance of the bounded bridge window under independent
root transformations. Integrating the actual Hessian of $\mathcal R_\gamma$
along $t\mathbf y$, $0\le t\le1$, gives \eqref{f16v31:action-size}; symmetry
and the absolute-row bound control its quadratic form. The interpolation stays
inside the convex product $\mathcal Y_r$.
\end{proof}

\subsection{An exact, exponentially summable retained-input interaction}

Order the bridge sites once and for all. Let $P_u$ be the predecessors of $u$,
$B_l(u)=\{v:d(u,v)\le l\}$, and $A_{u,l}=P_u\cap B_l(u)$ for integers
$l\ge0$. For $S\subseteq\mathcal V$, let $\mathbf y^S$ keep the inputs in $S$
and set all other bridge inputs to zero. Define
\begin{equation}
 \begin{split}
 D_{u,l}(\mathbf y)&=
 \mathcal R_\gamma(\mathbf y^{A_{u,l}\cup\{u\}})
                -\mathcal R_\gamma(\mathbf y^{A_{u,l}}),\\
 \Phi_{u,0}&=D_{u,0},\qquad
 \Phi_{u,l}=D_{u,l}-D_{u,l-1}\quad(l\ge1).
 \end{split}
 \label{f16v31:interaction-definition}
\end{equation}
These functions use complete native fast integrals at the displayed histories.
Fast variables outside $B_l(u)$ have \emph{not} been frozen or removed.

\begin{theorem}[A normalized action, not only an entrywise response expansion]
\label{f16v31:interaction}
The finite-radius sum terminates and gives the exact identity
\begin{equation}
 \boxed{\quad
 \mathcal F_\gamma(\mathbf y)
       =\tfrac12\mathbf y^{\mathsf T}K\mathbf y
                               +\sum_{u,l\ge0}\Phi_{u,l}(\mathbf y).
 \quad}
 \label{f16v31:interaction-exact}
\end{equation}
Each term depends only on bridge inputs in $B_l(u)$, vanishes at zero, and for
some fixed $b>0$ obeys
\begin{equation}
 \boxed{\quad
 \sup_v\sum_{u,l:\,v\in B_l(u)}
       |B_l(u)|e^{b\operatorname{diam}B_l(u)}
                  \|\Phi_{u,l}\|_{\infty,\mathcal Y_r}
                         \le C_r\gamma^{-1/2}.
 \quad}
 \label{f16v31:interaction-norm}
\end{equation}
A tail with $l>L$ has the stronger bound $C_r\gamma^{-1/2}e^{-bL}$ after a
possible fixed decrease of $b$. Constants are uniform in the regime, but the
coefficient functions can depend on the ambient regulator and the chosen
ordering. No infinite-volume consistency or independent-root gauge symmetry
of each term is claimed.
\end{theorem}
\begin{proof}
Once $l$ reaches the diameter of the finite anchor graph, every predecessor
has entered. Telescoping over $l$, and then over the site ordering, yields
$\sum_{u,l}\Phi_{u,l}=\mathcal R_\gamma(\mathbf y)-\mathcal R_\gamma(0)$.
The value at zero is exactly zero. Dependence on the retained inputs follows
from their definition; no spatial factorization of a native integral is used.

For $l\ge1$, add the sites in $A_{u,l}\setminus A_{u,l-1}$ one at a time.
The change in the $u$ increment when a site $v$ is added is a mixed finite
difference. Its exact double-integral formula is bounded by
\[
 |y_u|\,|y_v|\sup_{\mathbf z\in\mathcal Y_r}
              \|\partial_{y_u}\partial_{y_v}\mathcal R_\gamma(\mathbf z)\|
                 \le C_r r^2\varepsilon e^{-a d(u,v)}.
\]
Intermediate backgrounds differ with $v$; the supremum is already inside the
Hessian envelope in \eqref{f16v31:Hessian-norm}. The shell contains at most
$C(1+l)^4$ sites. Decreasing the rate gives
$\|\Phi_{u,l}\|_\infty\le C_r\varepsilon e^{-a_1l}$.
At $l=0$ there are only boundedly many co-anchored bridge types. The gradient
at zero vanishes, so the Hessian bound on this finite carrier gives
$\|\Phi_{u,0}\|_\infty\le C_r\varepsilon$ as well.

For fixed $v$ and $l$, there are at most $C(1+l)^4$ possible anchors $u$ with
$v\in B_l(u)$; also $|B_l(u)|\le C(1+l)^4$ and its diameter is at most $2l$.
Choosing $2b<a_1$ proves the norm bound by a polynomial times an exponential
series. Leave a strict exponential margin to obtain the tail assertion.
These counts remain valid on saturated tori, short histories, and at both
endpoints. They do not require a translation-invariant choice of ordering.
\end{proof}

The Gaussian quadratic term itself has exponentially summable two-site
coefficients, since $K=B_0-L^{\mathsf T}CL$. Thus the theorem supplies an
interaction representation of the \emph{whole} normalized fast action, with a
small correction to a generally order-one Gaussian interaction. The individual
$\Phi$ are smooth on the admitted window, but \eqref{f16v31:interaction-norm}
is a supremum-norm bound, not an unproved derivative norm for each term.
At fixed coupling the width restriction still prevents an unrestricted spatial
limit. A family of finite-volume local carrier functions is not automatically
one consistent infinite-volume Gibbs interaction.

\subsection{The endpoint factors and the exact interior bridge-window law}

V4's complete history weight has not changed. Its exact logarithmic ratio is
\begin{equation}
 \begin{split}
 -\log\frac{\mathcal W_{\gamma,n}^{\mathrm{full}}(\mathbf y)}
                  {\mathcal W_{\gamma,n}^{\mathrm{full}}(0)}
   ={}&\mathcal F_\gamma(\mathbf y)-\log J_{{\rm br},\gamma}(\mathbf y)\\
     &+\tfrac12\log\frac{F_{\gamma,2}(Z_0)}{F_{\gamma,2}(I)}
      +\tfrac12\log\frac{F_{\gamma,2}(Z_n)}{F_{\gamma,2}(I)}.
 \end{split}
 \label{f16v31:complete-weight}
\end{equation}
The last two functions are retained exactly. This appendix does not prove an
exponentially summable spatial representation for those endpoint normalizers.
No assertion about the full endpoint-varying Hessian is obtained by dropping
them. $J_{{\rm br},\gamma}(0)=1$ in the inherited convention.

Now fix both bridge endpoints, in the stated window, and let $n\ge2$.
Normalize the original positive path integrand over
$|y_{i,e}|\le r$, $1\le i\le n-1$. The resulting probability is exactly
\begin{equation}
 d\nu_{\gamma,r}^{\mathrm{win}}(\mathbf y_{\rm int})
  =\frac1{\mathcal Z_{\gamma,r}^{\mathrm{win}}}
    e^{-\mathcal F_\gamma(\mathbf y)}
       \prod_{u\,\mathrm{interior}}
       1_{\{|y_u|\le r\}}j(y_u/\sqrt\gamma)\,dy_u.
 \label{f16v31:bridge-window-law}
\end{equation}
The two dressed-end factors, their bridge half-Haar factors and the scalar
$c_{\gamma,B,n}$ cancel here because \emph{both endpoints are fixed}. The
remaining bridge Haar factors stay in their original product reference.
This is a window-conditioned, endpoint-fixed version of the original path
weight, not the unrestricted physical path probability. A further conditioning
of any subset of interior bridge sites also remains an exact kernel of this
same window law.

\begin{proposition}[Finite-range correction with normalized conditional error]
\label{f16v31:truncation}
Replace only $\mathcal R_\gamma$ in \eqref{f16v31:bridge-window-law} by
$\sum_{u,l\le L}\Phi_{u,l}$ and keep $K$, the endpoints and the Haar reference
unchanged. For every one-site conditional, at any admissible exterior,
\begin{equation}
 d_{\rm TV}(\nu_{\gamma,r}^{\rm win}(dy_v\mid\mathbf y_{v^c}),
                    \nu_{\gamma,r}^{[L]}(dy_v\mid\mathbf y_{v^c}))
                         \le C_r\gamma^{-1/2}e^{-bL}.
 \label{f16v31:conditional-truncation}
\end{equation}
There is no corresponding ambient-size-independent assertion for total variation
of the two \emph{joint} history laws.
\end{proposition}
\begin{proof}
Only omitted terms whose declared carrier contains $v$ can vary with $y_v$.
Their total oscillation is at most twice their summed supremum norms, bounded
by Theorem~\ref{f16v31:interaction}. For two normalized densities with the
same reference, a potential-difference oscillation $\delta$ implies TV at most
$\tanh(\delta/4)\le\delta/4$, as proved immediately below. This includes both
one-site normalizers and proves the bound. Summing a small per-site error over
an arbitrarily large joint carrier would be a different estimate.
\end{proof}

\subsection{A checkable influence margin, and a small fixed-window consequence}

\begin{proposition}[Oscillation, normalization, and a one-bridge influence bound]
\label{f16v31:TV-influence}
Two positive densities proportional to $e^{-V}$ and $e^{-W}$ with a common
reference have TV distance at most $\tanh(\operatorname{osc}(V-W)/4)$.
For the exact window law, set
\[
 \mathcal A_{uv}=\sup_{\mathbf y\in\mathcal Y_r}
                  \|\partial_{y_u}\partial_{y_v}\mathcal F_\gamma(\mathbf y)\|,
 \qquad u\ne v.
\]
Its influence when only exterior bridge $v$ changes obeys
\begin{equation}
 c_{uv}\le\tanh(r^2\mathcal A_{uv})\le r^2\mathcal A_{uv}.
 \label{f16v31:influence}
\end{equation}
The bound holds for all in-window exterior values and after arbitrary further
pinning of retained bridge sites within that window.
\end{proposition}
\begin{proof}
Let $R=d\mu/d\nu$, with $\nu$ the second normalized density. Its extrema
$m\le R\le M$ satisfy $m\le1\le M$ and $M/m\le e^\delta$, where
$\delta=\operatorname{osc}(V-W)$. The chord bound on $(R-1)_+$ gives
$d_{\rm TV}(\mu,\nu)\le(M-1)(1-m)/(M-m)$. Maximizing this expression subject
to the ratio bound gives $(e^{\delta/2}-1)/(e^{\delta/2}+1)$; the degenerate
case is zero. This proves the claim without ignoring either normalizer.

Changing $y_v$ between two values in its radius-$r$ ball changes the one-site
potential at $u$ by a function whose oscillation in $y_u$ is at most
$4r^2\mathcal A_{uv}$. Integrate the actual mixed derivative along the two
straight segments; both stay in $\mathcal Y_r$. All terms independent of
$y_u$, including the fixed endpoint factors, cancel from that conditional.
The native one-site Haar reference is the same in the two rows. Apply the
first part. Additional bridge pinning only removes free coordinates; it does
not change the uniform derivative estimate. No fast coordinate is pinned by
this argument.
\end{proof}

Write
\[
 L_0=\sup_u\sum_{v\ne u}\|K_{uv}\|<\infty,\qquad
 q_0(r)=\sup_u\sum_{v\ne u}\tanh(r^2\|K_{uv}\|).
\]
The constants refer to the actual complete Gaussian Schur matrix. Suprema may
be taken over all finite histories; its local Gaussian bounds make $L_0$ finite.
By the symmetric Hessian envelope, row and column influences are both bounded
by $q_0(r)+C_r r^2\gamma^{-1/2}$. Small correction alone does not say that
$q_0(r)<1$ at an arbitrary $r$.

\begin{theorem}[An exact retained-path conditional gap inside a certified window]
\label{f16v31:bridge-gap}
Whenever
\begin{equation}
 q_0(r)+C_r r^2\gamma^{-1/2}\le q<1,
 \label{f16v31:margin}
\end{equation}
the sum of rate-one single-bridge conditional innovations for
$\nu_{\gamma,r}^{\mathrm{win}}$, and for any of its further in-window bridge
conditionals, has gap at least $1-q$.

In particular, fix the V30 constants for the ambient bound $r=1$, and put
\begin{equation}
 r_* =\min\left\{1,\,[4(1+L_0)]^{-1/2}\right\}>0.
 \label{f16v31:small-window}
\end{equation}
At a fixed sufficiently large coupling in that same width regime, every
$0<r\le r_*$, every pair of fixed endpoints in $\mathcal Y_r$, and every
further retained-bridge pinning in the window satisfy
\begin{equation}
 \boxed{\quad
 \operatorname{Var}_{\nu} f
       \le2\sum_{u\,\mathrm{free}}
              \mathbb E_{\nu}\operatorname{Var}_{\nu}(f\mid y_{u^c}).
 \quad}
 \label{f16v31:window-gap}
\end{equation}
The constant has no number-of-bridges or duration factor. This is a new auxiliary
clock on retained bridge paths, not the fast block clock, the balanced elementary
clock, or physical Euclidean time.
\end{theorem}
\begin{proof}
The matrix in \eqref{f16v31:influence} has a symmetric nonnegative dominating
table, since the mixed Hessian blocks are transposes. The function $\tanh$ is
one-Lipschitz on the nonnegative axis. The Hessian remainder bound therefore
gives the stated row and column margin. Apply V24's finite-product
boundary-sensitive contraction theorem with one bridge site per update, unit
rates and unit weights. The death rate is one and the birth column sum is at
most $q$. Every intermediate exterior in a sequence of one-site changes is
still in the product window. The same conclusion holds on any conditional
subproduct because its dominating table is a principal submatrix. Empty free
sets give the vacuous variance identity.

For \eqref{f16v31:small-window}, use $q_0(r)\le r^2L_0\le1/4$. The constants
for $r=1$ dominate all smaller windows. Increase only a fixed threshold so that
$C_1r_*^2\gamma^{-1/2}\le1/4$. Then $q\le1/2$ uniformly as claimed. No small
Gaussian eigenvalue is inverted and no diagonal direct term is discarded.
The positivity of the native Haar reference on the small bridge balls supplies
explicit conditional-density versions at every in-window parameter. The
physical probability outside this retained window is not estimated here.
\end{proof}

\subsection{What remains beyond the retained window}

The exact interaction keeps every fast field integrated and the complete
Gaussian bridge action. The all-pinning theorem concerns retained bridges
inside a small fixed window, not the arbitrary compact fast patches left
open by V30. Nothing proves that the unrestricted retained path typically
stays in this window. A small nonlinear correction alone also does not make
the leading Gaussian influence sum less than one on a larger window.

The next distinct tasks are to prove or bypass that influence margin on a
larger admissible domain, obtain gauge-invariant coverage, and control the
endpoint functions retained in \eqref{f16v31:complete-weight}. Compatibility
of the carrier functions across ambient regulators is not established.
Neither an unrestricted spatial limit nor a physical transfer or continuum
mass statement follows from this finite-window auxiliary clock.

\subsection{Verification and preservation}

The script
\begin{center}\small\path{verify_ym_f16_v31_effective_bridge_action.py}\end{center}
passes 1,040 exact assertions: compact integration by parts with a nonzero
cutoff defect, normalized Schur/Hessian identities, many-body retained-input
telescopes, conditional density ratios, all-pinning finite probability fixtures,
and weighted-norm bookkeeping. These are algebraic tests, not native Wilson
integral enclosures or formal verification. Execution reports state the scope.

Only V30's literal title version and this terminal insertion are changed;
reversing them recovers V30 byte for byte. No inherited proof, endpoint factor,
source, or reference is silently edited. The analytic deductions use V16's
moments and V30's screening together with the identified Gaussian and ordered-word
certificates. The inherited lineage is not independently recertified. The next
checkpoint remains V35.

\begin{firewall}
V31 proves the complete bridge-Hessian comparison, an exact summable interaction
correction, and normalized one-site truncation bounds. Its retained-path gap
requires a checked influence margin; a sufficiently small fixed window satisfies
it, including further bridge pinning within that window. Width and window
restrictions remain. Endpoint normalizers cancel only in the endpoint-fixed
conditional. No arbitrary-fast-pinning theorem, physical-time identification,
or Yang--Mills mass gap is inferred.
\end{firewall}
% END F16 V31 APPEND

% BEGIN F16 V32 APPEND -- 2026-09-18
% Insert immediately before the final document-ending command in V31.
% No predecessor assertion, native probability, endpoint, source, or clock is edited.

\clearpage
\section{V32: boundary coercivity beyond the small bridge window}
\label{f16v32:context}

V31 bounds the complete retained-action Hessian, but its conditional gap uses
an absolute influence sum that is small only on a sufficiently small window.
We do not try to improve that sum term by term. The leading Schur action is
nonnegative, although it need not be strictly positive. We retain this signed
matrix information and use the boundary of the already specified retained
window to control its soft directions. The negative curvature of the native
remainder is then a perturbation of a \emph{boundary} coercivity estimate,
not of a nonexistent positive lower eigenvalue of the Gaussian Schur matrix.

For every prescribed finite bridge radius, this proves a positive rate-one
retained heat-bath gap, uniform in duration and in the inherited spatial-width
regime. It survives every further retained-bridge pinning in that window.
The radius is fixed before coupling is taken large; the constants deteriorate
with it. All fast coordinates remain completely integrated. No estimate of
unrestricted bridge coverage, arbitrary fast pinning, or physical mass follows.

\subsection{Dependencies, ownership, and the boundary being used}

The complete V31 appendix and audit were read. We use its full Hessian
estimate and endpoint-fixed law. V29's partial Bochner/heat-bath conversion
is reused, but the needed coercivity is proved here from the boundary rather
than assumed from a positive Hessian. Two unsuccessful semantic searches were
followed by mounted-text checks and identified f14 V64 and f15 V70--V71
passages. Their different marginal and observable-oscillation results are not
substituted for the present theorem. This is a targeted review, not an
exhaustive novelty certificate. The next checkpoint remains V35.

For context, Esposito--Nitsch--Trombetti,
\href{https://arxiv.org/abs/1110.2960}{\emph{Best constants in Poincar\'e
inequalities for convex domains}}, records the sharp Neumann diameter bound;
we reviewed its statement and prove a weaker elementary ball estimate below.
Ascolani--Lavenant--Zanella,
\href{https://arxiv.org/abs/2410.00858v2}{\emph{Entropy contraction of the Gibbs
sampler under log-concavity}}, Corollary 3.7 and Section 6, distinguishes its
strongly convex and merely convex settings. Those statements and assumptions
were checked, not the full proof. No whole-space entropy theorem is imported
across our hard boundary, and no entropy tensorization is claimed.

Retain the V31 law with both bridge endpoints fixed. Let $I$ be any subset of
its interior bridge sites, and pin the other retained sites at values $q$ in
the same radius-$r$ balls. Write $y=(x_I,q)$ and
\begin{equation}
 d\nu_{\gamma,r}^{I,q}(x)
 =\frac{e^{-V_{\gamma}^{I,q}(x)}}{Z_{\gamma,r}^{I,q}}
                  \prod_{u\in I}1_{\{|x_u|\le r\}}\,dx_u,\qquad
 V_{\gamma}^{I,q}(x)=\mathcal F_\gamma(x,q)
                   -\sum_{u\in I}\log j(x_u/\sqrt\gamma).
 \label{f16v32:law}
\end{equation}
Each free site has three real coordinates. Constants depending only on $q$
cancel from this conditional, including the two dressed endpoint functions
because both endpoints remain fixed. The interior Haar density in
\eqref{f16v32:law} is not deleted: it is included in the measured potential.
This is exactly V31's probability, not a convex extension or a replacement
Gaussian law. An empty $I$ has only the vacuous variance identity.

\subsection{The Schur action has no negative directions; it need not have a mass}

Let $K,H,L,B_0$ remain the matrices in \eqref{f16v31:Gaussian-blocks}. Put
\[
 L_K=\sup_u\sum_v\|K_{uv}\|_{\rm op}<\infty,
 \qquad
 \delta_{\gamma,r}=C_r\gamma^{-1/2},
\]
where $C_r$ is a valid constant in \eqref{f16v31:action-smallness}, enlarged
for its unweighted block norm if necessary. $L_K$ is a uniform bound over
the finite Gaussian histories, including diagonal blocks. It is not the
small-window off-diagonal bound $L_0$.

\begin{proposition}[Signed curvature and one-site oscillation budgets]
\label{f16v32:budgets}
The original Schur matrix satisfies $K\succeq0$, without an asserted positive
lower edge. There is a fixed finite Haar constant $C_H$ such that, whenever
\[
 \gamma\ge\max\{4r^2,C_H,\gamma_r\},\qquad
 B\le e^{\beta_r\gamma},\qquad \delta_{\gamma,r}\le1,
\]
every conditional \eqref{f16v32:law} satisfies
\begin{equation}
 \begin{gathered}
 D^2V_{\gamma}^{I,q}\succeq-\delta_{\gamma,r}I,
 \qquad (D^2V_{\gamma}^{I,q})_{uu}\preceq M_*I_3,
 \qquad M_*=L_K+2,\\
 \sup_{x_{u^c}}\operatorname{osc}_{|x_u|\le r}
                  V_{\gamma}^{I,q}(x)
       \le\Omega_r,\qquad \Omega_r=(2L_K+3)r^2.
 \label{f16v32:budgets-bound}
 \end{gathered}
\end{equation}
All bounds include arbitrary values of the retained pins within the radius-$r$
window. They contain no number of free sites.
\end{proposition}
\begin{proof}
The original quadratic $U_0$ is a sum of nonnegative quadratic limits of
Wilson actions and the nonnegative endpoint quadratics. Completing its fast
square gives the exact identity
\[
 U_0(\zeta;y)=\tfrac12(\zeta+H^{-1}Ly)^{\mathsf T}
                  H(\zeta+H^{-1}Ly)+\tfrac12y^{\mathsf T}Ky.
\]
Since $H\succ0$, minimize in $\zeta$ to obtain $K\succeq0$. No determinant
or scalar normalization supplies this conclusion; it follows from the
nonnegative quadratic action. Principal restrictions preserve it.

For $t=|x|/\sqrt\gamma<\pi$, the inherited Haar factor is
$j(x/\sqrt\gamma)=(\sin t/t)^2$. Set $\phi(t)=2\log(t/\sin t)$. Its radial
Hessian eigenvalues before scaling are
\[
 \phi''(t)=2(\csc^2t-t^{-2})>0,\qquad
 \frac{\phi'(t)}t=2(t^{-2}-t^{-1}\cot t)>0.
\]
Both extend continuously at zero, with value $2/3$. The first sign follows
from $\sin t<t$; the second follows from
$\sin t-t\cos t>0$ on $(0,\pi)$. Their maximum on $[0,1/2]$ defines a
finite $C_H$. Hence the measured one-site Haar Hessian is between zero
and $C_H/\gamma$, and its gradient vanishes at zero.

V31 gives $\|D^2\mathcal R_\gamma\|_{\rm op}\le\delta_{\gamma,r}$ from its
symmetric block-row bound. Add the nonnegative Haar Hessian and restrict to
$I$ for the first two estimates. Its proof also gives
$|\partial_{y_u}\mathcal F_\gamma(y)|\le r(L_K+\delta_{\gamma,r})$
throughout the full product window: integrate the Hessian along $ty$ from its
zero gradient at zero. A change of one bridge has length at most $2r$,
so the action oscillation at that site is at most $2r^2(L_K+1)$. The Haar
potential ranges from zero to at most $C_Hr^2/(2\gamma)\le r^2/2$.
The stated $\Omega_r$ is a conservative sum. Fixed retained pins do not
invalidate the estimate along the original full-window segment. No fast
variable is pinned in this step.
\end{proof}

\subsection{A one-ball normal-trace estimate, with a local density comparison}

The next calculation supplies the lower edge missing from a semidefinite $K$.
It concerns vector fields, not just a scalar Poincar\'e inequality.

\begin{proposition}[Boundary exclusion of constant vector fields]
\label{f16v32:ball}
Let $D=B_r\subset\mathbb R^d$, $1\le d\le3$, and let $v\in H^1(D;\mathbb R^d)$
have normal trace $v\cdot n=0$ on $\partial D$. Then
\begin{equation}
 \int_D|v|^2\,dx\le19r^2\int_D|Dv|_{\rm HS}^2\,dx.
 \label{f16v32:ball-vector}
\end{equation}
For a positive continuous weight $w$ on $\overline D$ with
$\sup w/\inf w\le e^\Omega$, the same normal-trace class obeys
\begin{equation}
 \int_D|Dv|^2w\,dx
 \ge \frac{e^{-\Omega}}{19r^2}\int_D|v|^2w\,dx.
 \label{f16v32:weighted-ball}
\end{equation}
The ratio is local to this one ball. It is not a density ratio on a growing
product of balls.
\end{proposition}
\begin{proof}
For a scalar $a\in H^1(D)$, the independent uniform-copy variance formula,
the line-segment fundamental theorem, and $|x-y|\le2r$ give
\[
 \int_D|a-\overline a|^2\,dx
          \le2^{d+1}r^2\int_D|\nabla a|^2\,dx.
\]
Here is the dimension constant explicitly. In the segment parameter $t\le1/2$
change $x$ to $(1-t)x+ty$ at fixed $y$; its Jacobian cost is at most $2^d$.
For $t\ge1/2$ change $y$ at fixed $x$. Convexity keeps both image domains in
$D$. The squared diameter and the factor $1/2$ in the variance formula give
$2^{d+1}r^2$. Approximation gives the $H^1$ statement.

Apply this to each component of $v$. Normal trace zero and the divergence
theorem give $\int_Dv=-\int_Dx\operatorname{div}v$. Therefore
\[
 |D|\,|\overline v|^2\le r^2\int_D|\operatorname{div}v|^2
                         \le dr^2\int_D|Dv|^2.
\]
Adding mean and centered parts yields the constant $2^{d+1}+d\le19$.
The trace/divergence identity holds by the $H^1$ trace theorem and density.
Finally bound the weighted derivative energy below by $\inf w$ times the
unweighted energy and the weighted norm above by $\sup w$ times its
unweighted value. Normalization of $w$ cancels. This proves
\eqref{f16v32:weighted-ball}.
\end{proof}

\subsection{A reflecting estimate for nearly convex product-window laws}

Consider $\mu=Z^{-1}e^{-V}dx$ on $\mathcal D=\prod_{u=1}^N B_r^{d_u}$,
$1\le d_u\le3$, with $V$ smooth on a neighborhood of its closure. Suppose
\begin{equation}
 D^2V\succeq-\delta I,\qquad (D^2V)_{uu}\preceq MI,
 \qquad \sup_{x_{u^c}}\operatorname{osc}_{x_u}V\le\Omega,
 \quad M\ge0.
 \label{f16v32:abstract-hypotheses}
\end{equation}
Set $\kappa=e^{-\Omega}/(19r^2)$ and assume $0\le\delta<\kappa$.
The scalar reflecting generator is $\mathcal A=-\Delta+\nabla V\cdot\nabla$.
For a one-form $v=(v_u)$ its absolute form domain consists of $H^1$ vector
fields with $v_u\cdot n_u=0$ on each face belonging to ball $u$. The form is
\begin{equation}
 \begin{split}
 \mathfrak a_1(v,v)={}&\int_{\mathcal D}\bigl(|Dv|^2+
                                      v^{\mathsf T}D^2Vv\bigr)d\mu\\
 &+\sum_u\int_{\partial_u\mathcal D}
            \mathrm{II}_u(v_{u,\mathrm{tan}},v_{u,\mathrm{tan}})\,d\sigma_\mu .
 \end{split}
 \label{f16v32:oneform}
\end{equation}
Here $d\sigma_\mu=Z^{-1}e^{-V}d\sigma$ and the ball second fundamental form is
$\mathrm{II}_u=r^{-1}I$ on its own tangential directions (zero in dimension
one). Cross-cell tangent directions have zero face curvature. All boundary
terms are retained.

\begin{theorem}[Boundary coercivity and whole-$L^2$ variance tensorization]
\label{f16v32:abstract}
Under \eqref{f16v32:abstract-hypotheses}, the positive one-form realization
$\mathcal A_1$ and scalar reflecting law satisfy, with the variance inequality for $f\in H^1(\mu)$,
\begin{equation}
 \mathcal A_1\succeq(\kappa-\delta)I,\qquad
 \operatorname{Var}_\mu f\le\frac1{\kappa-\delta}
                                  \int|\nabla f|^2d\mu .
 \label{f16v32:diffusion}
\end{equation}
For the exact rate-one block conditional expectations $E_u$, with one ball
per update, every $f\in L^2(\mu)$ satisfies
\begin{equation}
 \boxed{\quad
 \operatorname{Var}_\mu f
 \le\frac{\kappa+M}{\kappa-\delta}
                   \sum_u\|(I-E_u)f\|_2^2.
 \quad}
 \label{f16v32:heatbath}
\end{equation}
No strictly positive Hessian lower bound, diagonal dominance, or division of
the update rates by $N$ is required.
\end{theorem}
\begin{proof}
Fix all variables other than $x_u$. The normal trace of $v_u$ on this ball
is zero, and its actual conditional density has ratio at most $e^\Omega$.
Equation \eqref{f16v32:weighted-ball}, followed by disintegration, gives
\[
 \int|D_uv_u|^2d\mu\ge\kappa\int|v_u|^2d\mu.
\]
Sum in $u$. The full $|Dv|^2$ contains these terms, the face terms in
\eqref{f16v32:oneform} are nonnegative, and the curvature term is bounded
below by $-\delta\|v\|^2$. This proves the one-form floor with no global
density comparison.

We specify the domain passage. At each fixed $N$, the positive smooth weight
is bounded above and below on the compact product; these particular bounds
may depend on $N$ and are used only to identify Sobolev domains. The scalar
Neumann realization is the closed form $\int|\nabla u|^2d\mu$. Finite tensor
sums of smooth Neumann functions on the individual balls form a graph core
for the unweighted product Laplacian. The smooth bounded drift is relatively
bounded with arbitrarily small relative bound, by interpolation between its
gradient and Laplacian norms, so the same graph-domain approximation applies
to $\mathcal A$. On these functions, componentwise integration by parts gives
\begin{equation}
 \|\mathcal Au\|_2^2=\mathfrak a_1(du,du).
 \label{f16v32:Bochner}
\end{equation}
There is no codimension-two corner integral: the product identity is a sum
over its ball faces. The face terms follow by tangential differentiation of
$\partial_nu=0$. In particular $du$ has exactly the normal trace required
above. The closed forms extend the identity and its inequality by graph
approximation. Equivalently, the associated absolute Witten realization
satisfies $d\mathcal A=\mathcal A_1d$ on its natural domain; the integration
identity against smooth normal-zero vector fields proves this before closure.

Thus $\|\mathcal Au\|^2\ge(\kappa-\delta)\langle u,\mathcal Au\rangle$.
The compact reflecting scalar resolvent has a complete discrete spectrum.
Its kernel is constants because the product is connected and its Dirichlet
form vanishes only on constants. Applying the last inequality to nonconstant
eigenfunctions gives its gap at least $\kappa-\delta$. This proves
\eqref{f16v32:diffusion}. The dimension-dependent compact-domain facts used
in this passage introduce no dimension-dependent constant into the inequality.

For the heat bath, write $\mathcal A=\sum_i\mathcal A_i$. Full and
blockwise Bochner identities have the same nonnegative boundary sum. The
sum of block-diagonal squared Hessians is no greater than the full squared
Hessian. The difference of their potential contributions is at most
$(M+\delta)\int|\nabla u|^2d\mu$. Consequently
\begin{equation}
 \sum_i\|\mathcal A_i u\|_2^2
 \le\|\mathcal Au\|_2^2+(M+\delta)\langle u,\mathcal Au\rangle.
 \label{f16v32:partial}
\end{equation}
The same graph approximation extends this estimate and the exact conditional
identity $E_i\mathcal A_i u=0$. For centered $f\in L^2(\mu)$ solve
$u=\mathcal A^{-1}f$. Then
\[
 \|f\|_2^2=\sum_i\langle(I-E_i)f,\mathcal A_i u\rangle,
 \qquad \langle u,\mathcal Au\rangle
                      \le(\kappa-\delta)^{-1}\|f\|_2^2.
\]
Cauchy--Schwarz and \eqref{f16v32:partial} give the factor
$1+(M+\delta)/(\kappa-\delta)=(\kappa+M)/(\kappa-\delta)$.
The calculation is in the complete $L^2$ space; the test $f$ is not required
to satisfy any reflecting boundary condition. Real and imaginary parts cover
complex tests. This is the V29 conversion with a newly established boundary
floor, not a reversal of a one-site Poincar\'e inequality.
\end{proof}

\subsection{Every fixed retained radius, including all in-window pinning}

Define the explicit, conservative boundary constant
\begin{equation}
 \kappa_r=\frac{\exp[-(2L_K+3)r^2]}{19r^2},\qquad
 g_r=\frac{\kappa_r}{2(\kappa_r+L_K+2)}>0.
 \label{f16v32:constants}
\end{equation}
They do not contain a lower eigenvalue of $K$.

\begin{theorem}[Retained heat-bath gap at every prescribed finite window]
\label{f16v32:main}
For every fixed $0<r<\infty$ there are $\beta_r>0$ and a finite
$\Gamma_r$ such that, for
\begin{equation}
 \gamma\ge\Gamma_r,\quad B=m^3\le e^{\beta_r\gamma},\quad m\ge4,\quad n\ge2,
 \label{f16v32:regime}
\end{equation}
every endpoint-fixed law \eqref{f16v32:law}, and every further retained pinning
within its radius-$r$ window, obeys
\begin{equation}
 \boxed{\quad
 \operatorname{Var}_{\nu}f\le g_r^{-1}
             \sum_{u\in I}\mathbb E_\nu\operatorname{Var}_\nu(f\mid x_{u^c}),
 \qquad f\in L^2(\nu).
 \quad}
 \label{f16v32:main-gap}
\end{equation}
The reflecting bridge diffusion has gap at least $\kappa_r/2$. Every
single-bridge conditional update has rate one. There is no source-count,
spatial-volume, or duration multiplier. In contrast to V31, no condition
$q_0(r)<1$ or $r\le r_*$ is imposed.
\end{theorem}
\begin{proof}
Use the inherited V31 constants for the actual chosen radius $r$, not its
constants for radius one at a larger domain. Enlarge the threshold so that
\[
 \gamma\ge\max\{4r^2,C_H,\gamma_r\},\qquad
 \delta_{\gamma,r}\le\min\{1,\kappa_r/2\}.
\]
For example, the further requirement
$\gamma\ge[2C_r/\min(1,\kappa_r)]^2$ suffices. No numerical size for this
threshold is asserted. Proposition~\ref{f16v32:budgets} and
Theorem~\ref{f16v32:abstract}, with $M=M_*$, prove the conclusion.
All subsequent pinning takes a principal Hessian restriction and keeps the
same one-site oscillation bound. Its domain is still a product of the same
balls, including when the set of free sites has holes. Normalization constants
depending on pinned values disappear from derivatives but are never removed
from the probability. Empty free sets are vacuous. The proof does not change
or condition any of the integrated fast coordinates.
\end{proof}

This bypasses the absolute influence test, not the retained cutoff itself.
The radius $r$ is rescaled: its angular size is $r/\sqrt\gamma$, not a
fixed positive angular fraction of the compact group. The radius is chosen
first; $g_r$ can be very small, and $\Gamma_r$ can
be large. The result does not give uniform constants as $r\to\infty$, a
quantitative growing-radius law $r=r(\gamma)$, or a larger spatial-width
regime than that inherited from V31. It also does not prove a gap for the
endpoint-varying path law by dropping the dressed endpoint functions.

\subsection{Local response on those conditionals without a Dobrushin margin}

The same argument retains spatial information. We record a derivative-class
statement rather than equating it to a total-variation or Hamming coupling.
Fix $r$. The V31 Hessian envelope, the Gaussian rows, and the diagonal Haar
term give constants $a_0,M_{a_0}<\infty$, uniform at sufficiently large
coupling, such that the off-diagonal measured Hessian blocks obey
\[
 \sup_u\sum_{v\ne u}e^{a_0d(u,v)}
                   \sup_{x,q}\|(D^2V)_{uv}(x)\|\le M_{a_0}.
\]
These are rows of the complete retained action, including its generally
nonzero Gaussian interaction. Its range need not be finite.

\begin{proposition}[Localized one-form inverse and pinned response]
\label{f16v32:response}
There are $a_r,C_r'>0$, independent of sizes and of the in-window pinning,
such that all laws in Theorem~\ref{f16v32:main} satisfy
\begin{equation}
 \|P_u\mathcal A_1^{-1}P_v\|\le
                   C_r'e^{-a_rd(u,v)}.
 \label{f16v32:inverse}
\end{equation}
For $H^1$ observables the exact covariance identity therefore gives
\begin{equation}
 |\operatorname{Cov}_\nu(F,G)|
 \le C_r'\sum_{u,v\in I}e^{-a_rd(u,v)}
                  \|\nabla_uF\|_{L^2(\nu)}\|\nabla_vG\|_{L^2(\nu)}.
 \label{f16v32:gradient-covariance}
\end{equation}
If a retained pin $q_j$ changes to $q_j'$ within its ball, while all other
pins are held fixed, a free-coordinate Lipschitz observable $F$ independent
of that pin satisfies
\begin{equation}
 |\nu^{I,q'}F-\nu^{I,q}F|
 \le C_r'|q_j'-q_j|\sum_{u\in I}e^{-a_rd(u,j)}
                                \operatorname{Lip}_u(F).
 \label{f16v32:pinned-response}
\end{equation}
Here $\operatorname{Lip}_u$ is the Euclidean Lipschitz seminorm in one
three-component bridge. This is not an all-bounded-test influence coefficient.
\end{proposition}
\begin{proof}
The diagonal gradient form, the normal-trace constraints, and the face curvature
in \eqref{f16v32:oneform} are block diagonal in the one-form component label.
They commute with multiplication of component $u$ by $e^{t d(u,v)}$.
Only the Hessian multiplication has off-diagonal component blocks. Conjugating
it by these weights changes its operator by a bounded matrix whose row and
column norms are at most
\[
 M_{a_0}\sup_{s\ge0}(e^{ts}-1)e^{-a_0s}
                    \le 2tM_{a_0}/a_0\quad(0<t\le a_0/2).
\]
Triangle inequality for anchor distance and the block Schur bound justify both
row and column estimates. At each finite regulator the conjugating weights
are bounded and invertible; their size is not used in the perturbation bound.
Choose a fixed $t>0$ with this bound at most $\kappa_r/4$. Since
$\mathcal A_1\succeq\kappa_r/2$, the inverse of the conjugated operator has
norm at most $4/\kappa_r$ by its convergent resolvent perturbation. Unweighting
one block yields \eqref{f16v32:inverse}. This uses a proved one-form floor,
not an assumed inverse for an arbitrarily pinned fast patch.

The absolute reflecting identity gives
\[
 \operatorname{Cov}_\nu(F,G)
       =\langle dF,\mathcal A_1^{-1}dG\rangle.
\]
For smooth tests it follows from the scalar Neumann Poisson equation and
$d\mathcal A=\mathcal A_1d$ proved by the form integration above. Form
approximation extends it to $H^1$ tests. Apply the block inverse bound for
\eqref{f16v32:gradient-covariance}.

For the last claim, interpolate the actual pin linearly and put
$S_t=\partial_tV^{I,q(t)}$. Differentiation of the normalized conditional
integral, on its unchanged compact product, gives
\[
       \frac d{dt}\nu^{I,q(t)}F=-\operatorname{Cov}_{\nu^{I,q(t)}}(F,S_t).
\]
The derivative of the normalizer is exactly the centering on the right.
The gradient of $S_t$ in component $v$ is the free/pinned mixed Hessian
block applied to $q_j'-q_j$. It has the inherited exponential envelope.
Convolve it with \eqref{f16v32:inverse}; polynomial anchor-ball growth bounds
the convolution after a fixed decrease of its exponent. Integrate $t$.
Approximation of Lipschitz tests inside the convex product preserves the
coordinate Lipschitz bounds and gives the stated class. All intermediate
pins remain in the same radius-$r$ balls.
\end{proof}

In particular the retained-coordinate covariance has an exponentially summable
block-row bound at each fixed $r$, uniformly in all its admitted pinning.
These coordinates have order-one scale: no $1/\gamma$ covariance bound is
claimed for them. The small $1/\gamma$ covariance in earlier sections belongs
to the fast Haar-corrected residual sources, a different observable bank.

\subsection{A local comparison with the normalized truncated Gaussian action}

Let $\nu_{0,r}^{I,q}$ denote the normalized probability on the \emph{same}
retained product with potential $\tfrac12 y^{\mathsf T}Ky$ and Lebesgue
reference. It exists even if $K$ is singular, because the domain is compact.
It is not the untruncated Gaussian measure obtained by inverting $K$.

\begin{proposition}[Normalized local bridge-law comparison]
\label{f16v32:Gaussian-local}
For each fixed $r$, in the regime above, every free-coordinate Lipschitz $F$
satisfies
\begin{equation}
 |\nu_{\gamma,r}^{I,q}F-\nu_{0,r}^{I,q}F|
       \le C_r'\gamma^{-1/2}\sum_{u\in I}\operatorname{Lip}_u(F).
 \label{f16v32:Gaussian-local-bound}
\end{equation}
Consequently the Wasserstein distance between the marginals on a selected
set $A\subseteq I$, with cost $\sum_{u\in A}|x_u-x_u'|$, is at most
$C_r'|A|\gamma^{-1/2}$. The bound is uniform in the surrounding native
volume and duration. It is not total variation of the complete history laws.
\end{proposition}
\begin{proof}
Write $h_\gamma(x)=-\log j(x/\sqrt\gamma)$ and use the normalized
interpolation
\[
 V_t=\tfrac12 y^{\mathsf T}Ky+tW,\qquad
 W=\mathcal R_\gamma(y)+\sum_{u\in I}h_\gamma(x_u),\qquad 0\le t\le1.
\]
Its two endpoints are exactly the two displayed probabilities. Intermediate
$t$ is a comparison device, not a replacement native Wilson model. All
Hessian, conditional-oscillation and one-form estimates above hold uniformly
along this segment: $K\succeq0$, the Haar Hessian is nonnegative, and the
negative remainder is at most $\delta_{\gamma,r}$. Moreover
\[
 \sup|\nabla_u W|\le C_r r\gamma^{-1/2}+C_Hr/\gamma.
\]
Differentiation keeps the actual partition function at each $t$ and yields
$d\nu_tF/dt=-\operatorname{Cov}_{\nu_t}(F,W)$. Apply
\eqref{f16v32:gradient-covariance} and sum its exponentially decaying row
against the last uniform component bound. This proves
\eqref{f16v32:Gaussian-local-bound}. Compact-metric Kantorovich duality,
or finite-grid duality followed by compactness as in V26, gives the stated
Wasserstein consequence. The number of selected coordinates is kept, rather
than being mistaken for a size-free bound on a joint law.
\end{proof}

\subsection{Why the radius cannot be removed by the new gap theorem}

For a precise non-native test, consider
\[
 d\mu_r(x,y)=Z_r^{-1}e^{-(x-y)^2/2}
             1_{[-r,r]^2}\,dx\,dy,\qquad r\ge2.
\]
Its quadratic Hessian is positive semidefinite and has an exact constant-mode
null direction. The product-window theorem gives a positive heat-bath gap at
each fixed $r$. Nevertheless that gap is at most $C/r^2$. To see this, use
$f=x+y$. Each one-coordinate conditional is a unit-variance Gaussian truncated
to an interval. The one-dimensional reflecting Bochner identity with curvature
one gives its variance at most one, so $\mathcal E_{\rm HB}(f)\le2$.

Put $s=(x+y)/2$, $t=x-y$. The Jacobian is one and
$|s|\le r-|t|/2$, $|t|\le2r$. The normalizer is at most
$2r\sqrt{2\pi}$. The subregion $|t|\le1$,
$r/2\le|s|\le3r/4$ is admissible for $r\ge2$; it gives
\[
 \operatorname{Var}_{\mu_r}(x+y)=4\mathbb E_{\mu_r}s^2
       \ge \frac{e^{-1/2}}{2\sqrt{2\pi}}r^2.
\]
Thus the Rayleigh quotient is at most $4e^{1/2}\sqrt{2\pi}/r^2$.
This example does not claim a null mode or slow dynamics in the native model.
It proves that fixed-window gaps, even with a convex quadratic action, do not
justify a uniform radius-removal theorem. In this example the unbounded
quadratic density is not even normalizable.

The hard boundary is essential: it excludes a constant vector field from
\eqref{f16v32:ball-vector}. It supplies no positive mass in $K$. The local
comparison pays $e^{\Omega_r}$, not $e^{N\Omega_r}$, but deteriorates with $r$.

V32 bypasses V31's small-window influence condition at every fixed rescaled
radius. Cutoff removal or gauge-invariant coverage is now the next retained
variable problem, with the Gaussian soft structure and the endpoint functions
still present. A quantitative growing $r=r(\gamma)$ requires growth bounds
for the inherited $C_r$ and thresholds. The arbitrary-fast-pinning problem
remains separate; none of these auxiliary clocks is physical Euclidean time.

\subsection{Validation and predecessor preservation}

The diagnostic
\begin{center}\small\path{verify_ym_f16_v32_boundary_coercivity.py}\end{center}
passes 534 exact rational and symbolic assertions. It includes normal-trace
polynomial fields, nonzero weighted Bochner face terms, semidefinite Schur
completions, normalized interpolation, and the radius-removal example.
All 28 supplied V4--V31 diagnostics were rerun unchanged and passed. These
are finite fixtures, not native Wilson integral enclosures or verification
of the Sobolev-domain passage. Actual outputs and hashes are packaged.

Only the title version and this terminal appendix differ from V31; reversing
them recovers it byte for byte. No inherited proof, endpoint, Haar power,
source or probability is silently amended. The package records the dependency
audit and final presentation checks. The historical analytic lineage has not
been independently recertified. The next checkpoint remains V35.

\begin{firewall}
V32 proves a rate-one retained heat-bath gap at every fixed finite rescaled
radius, with all further in-window retained pinning, plus derivative-class
spatial response and a local truncated-Gaussian comparison. The hard boundary,
radius-dependent constants, fixed endpoints, and inherited width restriction
remain. No positive mass in $K$, arbitrary fast-pinning theorem, quantitative
cutoff removal, physical transfer gap, or interacting continuum mass gap is
inferred. All native normalizers and fully integrated fast fields are retained.
\end{firewall}
% END F16 V32 APPEND

% BEGIN F16 V33 APPEND -- 2026-09-18
% Insert immediately before the final document-ending command in V32.
% All predecessor statements, measures, endpoints, and clocks are preserved.

\clearpage
\section{V33: exact period dynamics and the retained-cutoff obstruction}
\label{f16v33:context}

V32 obtains a gap at every fixed retained radius by using the hard boundary.
Before trying to remove that boundary, one must identify which directions it
confines in the actual Gaussian reference. We return to the complete quadratic
history, not to another absolute influence estimate. The static magnetic null
space is already known from f15 V99. The new calculation follows its spatially
constant period directions through the \emph{actual} kinetic covariance and all
history times. They form nine exactly decoupled Gaussian bridges with increment
variance two. Their normalization and diffusion cannot be replaced by a mass.

This calculation has two consequences of different types. First, the existing
native compact-window comparison gives an upper bound on the probability of a
fixed all-time retained window under the \emph{unrestricted} endpoint-fixed native
path law. This is a one-sided consequence of nested conditioning; it needs no
unproved native tail bound. Second, the complete Gaussian history controls the
combined temporal-difference and magnetic-curvature observables in a norm that
has no radius or history-length cost. These statements select a more precise
cutoff target: distinguish raw connection amplitudes from their derivatives and
retain the nonlinear period variables rather than confining them by definition.

\subsection{Sources, ownership, and separated analytic dependencies}

The V32 appendix and analytic audit were read completely. Its statement for
every fixed radius remains unchanged. The f15 V99 proposition on
$\ker D$ already identifies coarse gradients and three constant directional
periods, of dimension $B+2$ per colour. That static classification is used and
credited below, not claimed as a new theorem. The f14 V54 discussion separately
treats nonlinear constant modes at fixed four-torus volume; neither its compact
matrix integral nor its normalization is the marginal of the present history.
V2's kinetic completion and V4's quadratic action supply the matrices.

Empty semantic searches were followed by identified passage reads and mounted
text searches. This is targeted review, not certification of the archives.
The next checkpoint remains V35.
For external context, Fodor--Holland--Kuti--Nogradi--Wong,
\href{https://arxiv.org/abs/1208.1051}{\emph{The Yang--Mills gradient flow in finite
volume}}, Section 2, separates constant fields from perturbatively integrated
nonzero modes. Its introduction and mode decomposition were examined, not imported as a lattice
theorem. The external PDF screenshot was unavailable. The import memoranda's
power/locality and physical-time distinctions remain binding.

The quadratic results below use only the identified finite matrices and exact
Gaussian integration. The later \emph{native} probability bound additionally
uses V31's action-size estimate on each fixed bridge window. That estimate
inherits V16 and V30; those analytic dependencies are not recertified here.
No untruncated native Gaussian limit is assumed.

\subsection{The full quadratic history controls differences and curvature}

Suppress colour tensors $I_3$ when writing spatial matrices. Keep the original
$N=\mathsf M_B^{-1}$, $A,D$, $W_n$, $G$, and $\Lambda=A^{\mathsf T}D$.
The matrix called $K$ throughout this section is V31's \emph{complete retained
history} Schur matrix, not V1's differently named static auxiliary matrix.
Let $\mathsf T_n\mathbf y=(y_t-y_{t+1})_{t=0}^{n-1}$ and
$\mathsf C_n\mathbf y=(Dy_i)_{i=0}^{n}$. Define
\[
 \mathsf M_n=
 \mathsf T_n^{\mathsf T}(I_n\otimes G^{-1})\mathsf T_n
                   +\tfrac13\mathsf C_n^{\mathsf T}\mathsf C_n.
\]
All coordinates in these maps are the original bridge coordinates. Both external
bridge slices are included until they are explicitly fixed.

\begin{theorem}[Two-sided kinetic--curvature comparison for the complete Schur action]
\label{f16v33:quadratic}
At every finite $B=m^3$, $m\ge4$, and $n\ge1$, the exact Gaussian action obeys
\begin{equation}
 \boxed{\quad
 \mathsf M_n\preceq K\preceq
 \mathsf T_n^{\mathsf T}(I_n\otimes G^{-1})\mathsf T_n
                          +\mathsf C_n^{\mathsf T}\mathsf C_n.
 \quad}
 \label{f16v33:quadratic-comparison}
\end{equation}
In particular its kernel consists exactly of time-independent histories with
$Dy_i=0$. Its dimension is $3(B+2)$ by the inherited static classification.
The comparison has no coupling parameter and no hard retained boundary.
\end{theorem}
\begin{proof}
Eliminate the temporal ports in \eqref{f16v4:V0} using the unchanged completion
\eqref{f16v2:completed-kinetic}. This leaves the bridge kinetic term
$\frac12\sum_t(y_t-y_{t+1})^{\mathsf T}G^{-1}(y_t-y_{t+1})$ and the internal
functional
\[
 \frac12\sum_i\left\{x_i^{\mathsf T}H_i x_i+|Ax_i+Dy_i|^2\right\}
       +\frac12\sum_t(x_t-x_{t+1})^{\mathsf T}N(x_t-x_{t+1}),
\]
where $H_i=w_i\mathsf H_B+\frac12 1_{\{0,n\}}(i)C_B^{-1}$.
The original endpoint factors give $w_0=w_n=1/2$; interior weights are one.
The inherited cube bounds give
\[
 H_i\succeq2N,\qquad A N^{-1}A^{\mathsf T}\preceq4I.
\]
Indeed $N^{-1/2}\mathsf H_BN^{-1/2}\succeq4I$ and
$A=C_{\rm int}U_B$, $U_B^{\mathsf T}U_B=N$, $\|C_{\rm int}\|\le2$.

The determinant after integrating the $x$ variables is independent of the
bridge history. It cancels in $Z_0(\mathbf y)/Z_0(0)$, so
$\mathcal F_0(\mathbf y)=\frac12\mathbf y^{\mathsf T}K\mathbf y$ is the minimum
of this complete reduced quadratic functional. Drop only the nonnegative
internal kinetic term for a lower bound. At a single time its remaining minimum
is
\[
 \tfrac12(Dy_i)^{\mathsf T}(I+AH_i^{-1}A^{\mathsf T})^{-1}Dy_i
                            \ge\tfrac16|Dy_i|^2,
\]
since $AH_i^{-1}A^{\mathsf T}\preceq2I$. Choosing every $x_i=0$ instead gives
the upper bound. This proves \eqref{f16v33:quadratic-comparison}; no matrix
commutativity is assumed. Since $G\succ0$, its two lower terms vanish exactly
when $y_i$ is time independent and belongs to $\ker D$. The upper bound proves
the converse. The dimension and the gradient/period description are the already
proved f15 V99 result, transported constantly through time.
\end{proof}

A useful exact expression, consistent with this variational proof, is
\begin{equation}
 \mathcal F_0(\mathbf y)=
 \tfrac12\sum_t|y_t-y_{t+1}|_{G^{-1}}^2+
 \tfrac12\sum_i|Dy_i|^2-
 \tfrac12(\Lambda y_i)_i^{\mathsf T}W_n^{-1}(\Lambda y_i)_i.
 \label{f16v33:action-formula}
\end{equation}
It follows from the same integration, or from V3 and V4's Gaussian consistency
identity. There is no dressed-end subtraction in $\mathcal F_0$: the latter is
the fast partition-value ratio, before dividing the physical history weight by
its two dressed-end normalizers. Fixing endpoints makes those factors constant.

\subsection{The actual port metric on constant directional periods}

There are four bridge sites on each positive coarse edge $(b,\mu)$.
Let $h_{\mu,a}$ take the value $(2\sqrt B)^{-1}e_a$ on \emph{every} bridge of
direction $\mu$, and zero on the other directions, where $1\le\mu,a\le3$ and
$e_a$ is a colour unit vector. These nine vectors are orthonormal in the original
$36B$-dimensional bridge Euclidean space. Write $\mathcal H$ for their span.
They are the spatially constant periods; they do not include the $B-1$ gradient
modes in each colour.

\begin{proposition}[Exact harmonic kinetic coefficient]
\label{f16v33:metric}
For every $h\in\mathcal H$,
\begin{equation}
             Dh=0,\qquad Gh=2h,\qquad G^{-1}h=\tfrac12h.
 \label{f16v33:harmonic-metric}
\end{equation}
Consequently $\mathcal H^{n+1}$ is a reducing subspace of $K$, and, if
$y_i=\sum_{\mu,a}a_{i,\mu,a}h_{\mu,a}$,
\begin{equation}
             \mathcal F_0(\mathbf y)=\frac14
                    \sum_{t=0}^{n-1}|a_{t+1}-a_t|^2.
 \label{f16v33:period-action}
\end{equation}
The factor two in \eqref{f16v33:harmonic-metric} uses the actual nonroot ports.
Freezing those ports would incorrectly replace it by one.
\end{proposition}
\begin{proof}
The first identity follows from the inherited common-bridge curl description.
For the kinetic calculation use $G=I+E\Sigma_BE^{\mathsf T}$ from V2, where
$\Sigma_B=(\mathsf B_B^{\mathsf T}\mathsf B_B)^{-1}$ is block diagonal.
It suffices to work with one unnormalized period $h$, equal to one on its
chosen directional bridges and zero on others. In a cube write
$u=(u_1,u_2,u_3)\in\{0,1\}^3$, with root $u=0$. Set $z_{b,u}=-u_\mu$.

At a nonroot vertex, $E^{\mathsf T}h$ is $+1$ if $u_\mu=0$ and $-1$ if
$u_\mu=1$: the exterior edge is respectively incoming or outgoing. The internal
cube Laplacian applied to $z$ has exactly these values; the other two directions
contribute zero. The root value of $z$ is zero, so
$\mathsf B_B^{\mathsf T}\mathsf B_B z=E^{\mathsf T}h$ on its actual rooted
space. Across a bridge, $Ez$ is one in direction $\mu$ and zero otherwise.
Thus $E\Sigma_BE^{\mathsf T}h=h$. Periodic seams satisfy the same identity.
Restore normalization and colour to obtain \eqref{f16v33:harmonic-metric}.

In \eqref{f16v33:action-formula} the two magnetic terms annihilate the period
space at each time. The kinetic matrix preserves that space and its Euclidean
orthogonal complement because it is symmetric and has the displayed eigenvalue.
Therefore all cross terms vanish, proving both reduction and
\eqref{f16v33:period-action}. This is not merely evaluation of the quadratic
form on one selected path.
\end{proof}

\subsection{The untruncated endpoint-fixed Gaussian and its nine exact bridges}

For $n\ge2$ fix $y_0=y_n=0$ and let $K^{\rm int}$ be the principal retained
precision on the $36B(n-1)$ free scalar coordinates. Denote by $L_n$ the
$(n-1)$-dimensional Dirichlet time Laplacian, with diagonal two and adjacent
entries minus one. Since $I\preceq G\preceq30I$, Theorem~\ref{f16v33:quadratic}
gives
\begin{equation}
        K^{\rm int}\succeq\tfrac1{30}(L_n\otimes I_{36B})\succ0.
 \label{f16v33:dirichlet-lower}
\end{equation}
Thus the normalized Gaussian probability $Q_{B,n}^{00}$ on the entire interior
Euclidean bridge space exists. This does not assert the existence of a Gaussian
with free, unanchored endpoints, whose static kernel was identified above.

\begin{theorem}[Exact normalized period factor and temporal soft scale]
\label{f16v33:bridges}
Under $Q_{B,n}^{00}$ the nine period amplitudes
$a_{i,\mu,a}=\langle h_{\mu,a},y_i\rangle$ are independent of the complementary
bridge history and have joint density
\begin{equation}
 \frac{\prod_{t=0}^{n-1}p_1^{(9)}(a_{t+1}-a_t)}{p_n^{(9)}(0)}
                    \prod_{i=1}^{n-1}da_i,
 \qquad p_t^{(d)}(z)=(4\pi t)^{-d/2}e^{-|z|^2/(4t)}.
 \label{f16v33:bridge-density}
\end{equation}
Their means are zero and their covariance is
\begin{equation}
 \mathbb E[a_{i,\mu,a}a_{j,\nu,b}]
 =2\left(\min(i,j)-\frac{ij}{n}\right)\delta_{\mu\nu}\delta_{ab}.
 \label{f16v33:bridge-covariance}
\end{equation}
With nonzero fixed endpoints, their mean is the linear endpoint interpolation
and this covariance is unchanged. In particular $K^{\rm int}$ has the exact
period eigenvalues $2\sin^2(k\pi/(2n))$, $1\le k<n$, each with at least ninefold
multiplicity, and
\begin{equation}
 \frac{2}{15}\sin^2\frac\pi{2n}
 \le\lambda_{\min}(K^{\rm int})
 \le2\sin^2\frac\pi{2n}.
 \label{f16v33:soft-scale}
\end{equation}
No positive lower edge uniform in duration can hold for this untruncated
retained Gaussian precision, even at fixed spatial width.
\end{theorem}
\begin{proof}
Apply the orthogonal period/complement change of coordinates at every time.
Proposition~\ref{f16v33:metric} eliminates all cross terms, including terms
involving fixed endpoints. The density consequently factors after normalization.
The period action is \eqref{f16v33:period-action}; convolution of the displayed
heat densities gives its exact denominator $p_n^{(9)}(a_n-a_0)$. At zero endpoints
this is \eqref{f16v33:bridge-density}. It is not an adjustable determinant.

The inverse of $L_n$ has entries $\min(i,j)-ij/n$. Direct second differences,
with values zero at times zero and $n$, verify $L_nL_n^{-1}=I$.
The period precision is $L_n/2$, proving the covariance. Completing the square
about the linear interpolation proves the endpoint assertion. The sine vectors
diagonalize $L_n$, with eigenvalues $4\sin^2(k\pi/(2n))$.
Equation~\eqref{f16v33:dirichlet-lower} proves the lower bound in
\eqref{f16v33:soft-scale}; the reducing period space proves the upper bound.
\end{proof}

For completeness, the sum of rate-one three-component bridge heat baths of
$Q_{B,n}^{00}$ also has gap at most $30\sin^2(\pi/(2n))$. Indeed choose a unit
first temporal period eigenvector $v$ and test $f=v^{\mathsf T}y$.
Its variance is $[2\sin^2(\pi/(2n))]^{-1}$. Every diagonal bridge block of
$K^{\rm int}$ is at least $I_3/15$ by \eqref{f16v33:dirichlet-lower}, so the
sum of its exact Gaussian conditional variances is at most $15\|v\|^2=15$.
The Rayleigh quotient gives the stated upper bound. This sampler and the
precision are not physical Hamiltonians. The result does not contradict V32:
its hard window is absent here.

\subsection{A native upper bound on fixed all-time bridge-window coverage}

Let $\boldsymbol\nu_{\gamma,B,n}^{00}$ be the normalized \emph{original}
endpoint-fixed retained path measure with $y_0=y_n=0$, integrated over every
interior principal bridge chart. It is the positive unprojected history
integrand of V4, with all fast fields integrated; it is not a newly Gauss-projected
physical probability. In its chart it is proportional to
\[
 e^{-\mathcal F_\gamma(\mathbf y)}
           \prod_{u\ {
m interior}}j(y_u/\sqrt\gamma)\,dy_u.
\]
The scalar and dressed-end factors cancel because the endpoints are fixed.
For a finite rescaled radius $r$ put
\[
 G_r=\{|y_u|\le r\text{ for every interior bridge }u\},\qquad
 M_{B,n}=12B(n-1).
\]
When $r<\pi\sqrt\gamma$, this is an event in the original full compact path.
The event concerns the \emph{entire} retained history, not one local marginal.

\begin{proposition}[Nested-window comparison with the full native normalizer retained]
\label{f16v33:native-comparison}
Fix $B,n$ and $0<r\le R<\infty$. At the V31 entry threshold for the actual
radius $R$, and $\gamma\ge\max(4R^2,C_H)$, define
\[
 d_{\gamma,R}=M_{B,n}R^2
                  \left(C_R\gamma^{-1/2}+\frac{C_H}{2\gamma}\right).
\]
Then
\begin{equation}
 \boxed{\quad
 \boldsymbol\nu_{\gamma,B,n}^{00}(G_r)
 \le\min\left\{1,\ e^{2d_{\gamma,R}}
               \frac{Q_{B,n}^{00}(G_r)}{Q_{B,n}^{00}(G_R)}\right\}.
 \quad}
 \label{f16v33:nested-probability}
\end{equation}
In particular,
\begin{equation}
 \limsup_{\gamma\to\infty}
       \boldsymbol\nu_{\gamma,B,n}^{00}(G_r)\le Q_{B,n}^{00}(G_r).
 \label{f16v33:native-limsup}
\end{equation}
The limit holds with $B,n,r$ fixed. It is an upper bound, not convergence of
the unrestricted native path to an untruncated Gaussian.
\end{proposition}
\begin{proof}
On $G_R$, V31 gives
$|\mathcal F_\gamma-\mathcal F_0|\le C_R\gamma^{-1/2}\sum_u|y_u|^2$.
V32's elementary Haar estimate gives
$0\le-\log j(y_u/\sqrt\gamma)\le C_H|y_u|^2/(2\gamma)$.
Thus the difference of the two measured potentials is bounded in absolute
value by $d_{\gamma,R}$. The \emph{normalized} likelihood ratio between
$\boldsymbol\nu_\gamma(\cdot\mid G_R)$ and $Q_{B,n}^{00}(\cdot\mid G_R)$
is consequently between $e^{-2d_{\gamma,R}}$ and $e^{2d_{\gamma,R}}$.
Both restricted normalizers are included.

Since $G_r\subseteq G_R$,
$\boldsymbol\nu_\gamma(G_r)=\boldsymbol\nu_\gamma(G_R)
\boldsymbol\nu_\gamma(G_r\mid G_R)\le\boldsymbol\nu_\gamma(G_r\mid G_R)$.
This proves \eqref{f16v33:nested-probability}, without estimating the full
normalizer or any native tail outside $G_R$. For fixed $R,B,n$, the inherited
width condition is eventually satisfied and $d_{\gamma,R}\to0$. Take this
limsup first. Only afterward send $R\to\infty$; the proper Gaussian of
\eqref{f16v33:dirichlet-lower} has $Q_{B,n}^{00}(G_R)\to1$. This proves
\eqref{f16v33:native-limsup} in the stated order of limits.
\end{proof}

A fully explicit Gaussian denominator bound is available:
\begin{equation}
 Q_{B,n}^{00}(G_R)\ge1-6M_{B,n}e^{-R^2/(45n)}.
 \label{f16v33:outer-window}
\end{equation}
Indeed \eqref{f16v33:dirichlet-lower} bounds every scalar variance by
$30i(n-i)/n\le15n/2$. If a three-vector exceeds $R$, one coordinate exceeds
$R/\sqrt3$. The scalar Gaussian exponential bound and a union bound prove
\eqref{f16v33:outer-window}, without independence of the bridge coordinates.
For example $R^2\ge45n\log(12M_{B,n})$ makes this lower bound at least $1/2$.
The native constant $C_R$ and its threshold still depend on this chosen radius.

\begin{theorem}[An actual native obstruction to uniform all-time window coverage]
\label{f16v33:cutoff}
For fixed $B=m^3$, $m\ge4$, and $r>0$, every even $n\ge2$ satisfies
\begin{equation}
 \boxed{\quad
 \limsup_{\gamma\to\infty}
 \boldsymbol\nu_{\gamma,B,n}^{00}(G_r)
 \le\left[\min\left\{1,4r\sqrt{\frac{B}{\pi n}}\right\}\right]^9.
 \quad}
 \label{f16v33:midpoint-cutoff}
\end{equation}
There is also a temporal skeleton bound. Put
$\tau=\max\{1,\lceil16Br^2\rceil\}$ and take $n=k\tau$, $k\ge2$. Then
\begin{equation}
 \limsup_{\gamma\to\infty}
 \boldsymbol\nu_{\gamma,B,n}^{00}(G_r)
       \le\min\{1,k^{9/2}2^{-9(k-1)}\}.
 \label{f16v33:skeleton-cutoff}
\end{equation}
Consequently
\begin{equation}
 \inf_{n\ge2}\limsup_{\gamma\to\infty}
              \boldsymbol\nu_{\gamma,B,n}^{00}(G_r)=0.
 \label{f16v33:no-uniform-coverage}
\end{equation}
In particular there exist $n_j\to\infty$ and $\gamma_j\to\infty$, at fixed $B$,
for which the native probability of this fixed all-time window tends to zero.
No quantitative joint growth rate for $n_j,\gamma_j$ is asserted.
\end{theorem}
\begin{proof}
On $G_r$, for each period and each interior time,
\[
 |a_{i,\mu,a}|=\left|\frac1{2\sqrt B}
             \sum_{e:\,\mathrm{dir}(e)=\mu}y_{i,e}^a\right|
                                          \le2\sqrt B\,r=:A_r^{\rm per}.
\]
At $i=n/2$, Theorem~\ref{f16v33:bridges} gives nine independent scalar
Gaussians of variance $n/2$. Each has density at most $(\pi n)^{-1/2}$.
Their simultaneous interval probability is at most the right side of
\eqref{f16v33:midpoint-cutoff}. The full event is a subset of that event.
Apply \eqref{f16v33:native-limsup}.

For the sharper skeleton calculation, one scalar period sampled at times
$\tau,2\tau,\ldots,(k-1)\tau$ has transition density $p_\tau^{(1)}$ and
exact bridge denominator $p_{k\tau}^{(1)}(0)$. Uniformly in its incoming value,
\[
 \int_{-A_r^{\rm per}}^{A_r^{\rm per}}p_\tau^{(1)}(x-y)\,dy
    \le\frac{A_r^{\rm per}}{\sqrt{\pi\tau}}\le\frac12.
\]
Bound the final transition density by $(4\pi\tau)^{-1/2}$, integrate the
preceding $k-1$ transitions with this row bound, and divide by
$p_{k\tau}^{(1)}(0)$. The resulting probability is at most
$\sqrt k\,2^{-(k-1)}$. The nine scalar period bridges are independent, so
raise this bound to the ninth power. Again $G_r$ implies this skeleton event,
and \eqref{f16v33:native-limsup} transfers its upper bound. Finally choose
$k_j\to\infty$ and, for each resulting finite $n_j$, choose a sufficiently
large $\gamma_j$ making the limsup bound hold within $1/j$. Increase these
couplings successively. This is a diagonal choice, not a uniform-radius
estimate on $C_R$ or an exchange of the two limits.
\end{proof}

This native conclusion uses only the nested-window upper comparison; it does
not prove that a \emph{single-time} native marginal is non-tight, because that
event is not contained in every $G_R$. It does not contradict V32's positive
gap for the law already conditioned on $G_r$. An all-time coordinate window
and local gauge-invariant coverage are different targets. Linear period
coordinates are not nonlinear gauge-invariant Wilson holonomies, and no
physical mass conclusion follows from \eqref{f16v33:no-uniform-coverage}.

\subsection{A positive radius-free Gaussian interface for the next native estimate}

The same quadratic calculation supplies a norm that does not pay the period
amplitude. Let $\Delta y_t=y_t-y_{t+1}$, keep zero endpoints, and form
\[
 \mathcal B\mathbf y=
 \left((G^{-1/2}\Delta y_t)_{t=0}^{n-1},
                         (3^{-1/2}Dy_i)_{i=1}^{n-1}\right).
\]
Its components are assembled together, including cross covariances.

\begin{theorem}[Uniform difference--curvature covariance, without a retained radius]
\label{f16v33:derivatives}
Under the proper untruncated Gaussian $Q_{B,n}^{00}$,
\begin{equation}
                  \operatorname{Cov}(\mathcal B\mathbf y)\preceq I.
 \label{f16v33:derivative-Gram}
\end{equation}
In particular, for arbitrary coefficient banks $a_t,b_i$,
\begin{equation}
 \boxed{\quad
 \operatorname{Var}\left(\sum_t a_t^{\mathsf T}\Delta y_t+
                                \sum_i b_i^{\mathsf T}Dy_i\right)
 \le\sum_t a_t^{\mathsf T}Ga_t+3\sum_i|b_i|^2
 \le30\sum_t|a_t|^2+3\sum_i|b_i|^2.
 \quad}
 \label{f16v33:derivative-bound}
\end{equation}
The constants are independent of $B,n$ and there is no hard window. The same
covariance bounds hold after arbitrary further \emph{Gaussian} retained
pinning, and for nonzero fixed endpoints, with the appropriate means subtracted.
For the exact period factors,
\begin{equation}
 \operatorname{Cov}(a_{t+1}-a_t,a_{s+1}-a_s)
                  =2(\delta_{ts}-n^{-1})I_9.
 \label{f16v33:increment-covariance}
\end{equation}
Thus the period amplitudes grow through the bridge, but their increments have
uniform covariance with the endpoint constraint retained.
\end{theorem}
\begin{proof}
On the endpoint-fixed space, \eqref{f16v33:quadratic-comparison} says
$\mathcal B^{\mathsf T}\mathcal B\preceq K^{\rm int}$. Hence
$\|\mathcal B(K^{\rm int})^{-1/2}\|\le1$, and multiplication by its adjoint
proves \eqref{f16v33:derivative-Gram}. Apply it to the output coefficient vector
$((G^{1/2}a_t)_t,(\sqrt3b_i)_i)$; this proves the first inequality in
\eqref{f16v33:derivative-bound}, including all mixed terms. The second is
$G\preceq30I$. Further Gaussian pinning takes the same principal precision and
columns of $\mathcal B$; their matrix inequality and the proof persist.
Nonzero endpoint values change linear terms, not the precision.
Finally apply two temporal differences to \eqref{f16v33:bridge-covariance},
or condition $n$ independent variance-two increments on their zero sum. This
gives \eqref{f16v33:increment-covariance}, including its negative off-diagonal
entries and its exact zero mode in the sum of the increments.
\end{proof}

\paragraph{What would be sufficient, and what is not supplied.}
On any fixed product window with endpoints fixed, an independently proved
native lower bound $D^2V\succeq(1-\eta)K^{\rm int}$, $0\le\eta<1$, would
preserve \eqref{f16v33:derivative-Gram} with factor $(1-\eta)^{-1}$ by the
reflecting covariance identity and its nonnegative face terms. This is a
\emph{relative} curvature requirement. V31 supplies instead the absolute
bound $D^2V\succeq K^{\rm int}-\delta_{\gamma,r}I$. Using only that bound
requires $\delta_{\gamma,r}$ smaller than a constant times $n^{-2}$ to obtain
the relative implication. It therefore does not close a duration-uniform
native derivative theorem. No such additional relative bound is asserted here.

The next noncircular task is to formulate the native comparison in temporal
differences and curvature while retaining the compact nonlinear periods and
the exact gauge/endpoint data. The period/complement orthogonal change of
coordinates used for a Gaussian does not turn the native product domain into
a product in those variables. Its moving fibres, density, and conditional
normalizers would have to be kept. The f14 fixed-four-torus constant-mode
analysis is a useful separate input, not a ready-made history marginal.
Nothing here removes the spatial-width restriction from the native estimates,
identifies an auxiliary clock with physical time, or proves a physical gap.

\subsection{Verification and preservation}

The accompanying diagnostic is
\begin{center}\small\path{verify_ym_f16_v33_period_bridge.py}.\end{center}
It checks the literal rooted-cube and periodic bridge incidence identity,
Gaussian Schur comparisons, temporal inverse and normalization identities,
period reduction, derivative covariance and the finite-window comparison
algebra. The checker passes 226 exact assertions. All 29 available V4--V32 checkers were
rerun unchanged and passed. These are finite algebraic fixtures, not native
integral enclosures or formal verification. Actual reports are packaged.

Only V32's one title version and this terminal insertion are changed. Their
reversal recovers V32 byte for byte. The cutoff theorem depends on V31; the quadratic theorems do not depend on
V30's native gap. The archive has not been independently recertified.

\begin{firewall}
V33 identifies nine exact period Gaussian bridges in the complete retained
history, proves a two-sided kinetic--curvature comparison, and obtains
radius-free Gaussian derivative covariance bounds. A normalized nested-window
argument gives a one-sided native obstruction to uniform fixed all-time window
coverage in a stated order of limits. It does not claim an untruncated native
Gaussian limit, a native all-radius derivative inequality, local native
non-tightness, a gauge-invariant cutoff-removal theorem, a thermodynamic gap,
or a physical Yang--Mills mass gap. Every native normalizer used in the comparison
is retained, and all physical clocks and Gauss/continuum obligations remain
separate.
\end{firewall}
% END F16 V33 APPEND

% BEGIN F16 V34 APPEND -- 2026-09-18
% Insert before the final document-ending command of the cumulative V33.
% The preceding text, all original measures, sources and update clocks are preserved.

\clearpage
\section{V34: endpoint anchoring and quantitative native cutoff removal}
\label{f16v34:context}

V33 obtains only an upper comparison for the probability of a retained window:
its proof does not estimate the full native normalizer. We go upstream to that
missing estimate. With both retained endpoints fixed at the identity, V5's
compact fast pins and the original bridge kinetic words also confine the
retained coordinates. The confinement constant deteriorates as the inverse
square of the duration. This is sufficient for a complete normalized Laplace
comparison, but not for a duration-uniform mass statement.

The resulting theorem removes the retained cutoff at every fixed regulator,
and supplies an explicit, conservative simultaneous regime
$Bn\log(n+1)=o(\log\gamma)$. It upgrades V33's one-sided comparison to total
variation and polynomial-moment convergence of the full native joint law.
Consequences include a native difference--curvature covariance comparison, a
comparison for actual nonlinear Wilson-word observables, and a quantitative
period-bridge limit. The latter also shows that a single midpoint bridge cannot
remain in a fixed rescaled window along a specified slowly growing history.
Every assertion concerns the endpoint-fixed, unprojected history, not the
physical Gauss-reduced vacuum.

\subsection{Dependencies and the exact joint probability}

The supplied V33 appendix and audit were read completely. Its period metric,
Gaussian history and derivative bank are reused, not recomputed as new results.
The analytic starting points here are substantially earlier: V4's exact history
integrand, V5's complete-chart compact pin and frame relation, and V6's bounded
ordered-word derivatives. The new proof uses neither V16's conditional-moment
iteration, V30's screening, nor V31's fixed-window action comparison. In
particular no growing-radius estimate for the constants $C_r$ is assumed.
Those earlier results keep their own scopes and are not silently amended.

Scoped semantic searches returned no results; mounted exact-text searches and
identified passages were then checked. The inherited tree-pin and finite-window
Laplace arguments are credited. No exhaustive novelty or archive certification
is asserted. For context, Faria da Veiga--O'Carroll,
\href{https://arxiv.org/abs/1903.09829}{\emph{On Thermodynamic and Ultraviolet
Stability of Yang--Mills}}, uses exact tree coordinates for partition bounds.
Spokoiny, \href{https://arxiv.org/abs/2204.11038}{\emph{Dimension free
non-asymptotic bounds on the accuracy of high dimensional Laplace approximation}},
keeps dimension costs in a normalized approximation. Their primary abstracts
and scope were checked, not their complete proofs. Neither theorem is imported
for this history. The finite-dimensional domination and normalization estimates
needed below are proved directly. The next companion checkpoint remains V35.

Fix $B=m^3$, $m\ge4$, $n\ge2$, and $Z_0=Z_n=I$. Use V5's unmoving fast
coordinates $\zeta=(x,a)$ and all interior retained coordinates $y$ on their
\emph{complete} principal charts. Put $w=(\zeta,y)$ and
\begin{equation}
 d=d_{B,n}=15B(n+1)+21Bn+36B(n-1)=(72n-21)B.
 \label{f16v34:dimension}
\end{equation}
Let $\mathcal P_\gamma\subset\mathbb R^d$ be their product of
three-dimensional radius-$\pi\sqrt\gamma$ balls. Its boundary has measure zero.
With exactly V5's classical action, define
\begin{equation}
 \begin{split}
 \mathcal J_\gamma(w)&=J^f_\gamma(\zeta)
                 \prod_{i=1}^{n-1}\prod_{e\ {\rm bridge}}
                           j(y_{i,e}/\sqrt\gamma),\\
 q_\gamma(w)&=1_{\mathcal P_\gamma}(w)
                 \mathcal J_\gamma(w)e^{-U_\gamma(\zeta;y)},\\
 \mathcal Z_\gamma&=\int_{\mathbb R^d}q_\gamma(w)\,dw,
 \qquad d\mathbb P_\gamma=\mathcal Z_\gamma^{-1}q_\gamma(w)\,dw.
 \end{split}
 \label{f16v34:joint-law}
\end{equation}
Chord endpoint Haar powers are $1/2$, all other free-group powers are one.
The two endpoint Gaussian amplitudes remain in $U_\gamma$. The factors
$c_{\gamma,B,n}$ and $F_{\gamma,2}(I)^{-1}$ cancel only when this fixed-endpoint
probability is normalized. Thus its retained marginal is exactly V33's
$\boldsymbol\nu_{\gamma,B,n}^{00}$, not a new probability or a replacement of
its physical transfer normalization.

Write $U_0(w)=\tfrac12w^{\mathsf T}\mathcal H_{B,n}w$ for the original full
quadratic limit with these fixed endpoints, and set
\[
 q_0=e^{-U_0},\qquad
 \mathcal Z_0=(2\pi)^{d/2}(\det\mathcal H_{B,n})^{-1/2},\qquad
 d\mathbb Q=\mathcal Z_0^{-1}q_0\,dw.
\]
Positive definiteness and this normalizer are verified below. Its retained
marginal is V33's $Q_{B,n}^{00}$, since the fast Gaussian integration and its
determinant are unchanged.

\subsection{The compact kinetic terms transmit the endpoint pin}

Let $\operatorname{dist}(U,V)\in[0,\pi]$ be the bi-invariant geodesic distance
on $\SU(2)$ in the convention $s(U)=1-\Re U$. Thus
$\operatorname{dist}(I,\operatorname{Exp}(z/\sqrt\gamma))=|z|/\sqrt\gamma$
on a principal chart and
$\operatorname{dist}(I,U)^2\le(\pi^2/2)s(U)$.

\begin{proposition}[A complete joint pin with its duration cost]
\label{f16v34:compact-pin}
There is a fixed $c>0$, independent of $B,n,\gamma$, such that on every
complete principal chart with $Z_0=Z_n=I$,
\begin{equation}
 \boxed{\qquad U_\gamma(\zeta;y)\ge\frac{c}{n^2}
                    (\|\zeta\|^2+\|y\|^2),\qquad \gamma>0.\qquad}
 \label{f16v34:pin}
\end{equation}
The same inequality holds for $U_0$. Moreover
$\mathcal H_{B,n}\preceq H_*I$ for a fixed finite $H_*$.
The identity is the unique zero of the classical endpoint-fixed action in
these coordinates. No retained window is imposed.
\end{proposition}
\begin{proof}
V5's anchor is independent of $y$ and every other classical term is
nonnegative. Its literal proof therefore gives, for all compact retained values,
\[
 U_\gamma\ge a_*\|\zeta\|^2,\qquad a_*=(6\pi^2)^{-1}.
\]
We are not extending V5's radius-dependent pressure estimates. Only its
explicit, retained-input-independent anchor is used here.

Let $z_{t,v}$ be the principal coordinate of the original moving port $r_{t,v}$.
V5's exact relation $r_t=f_t^{-1}h_tf_{t+1}$ and its global angle estimate give
\[
 \sum_{t,v}|z_{t,v}|^2\le C_f\|\zeta\|^2,
 \qquad C_f=3+6\|\mathsf F_\Box\|^2.
\]
Indeed $|z_{t,v}|\le|a_{t,v}|+|(\mathsf F_Bx_t)_v|
+|(\mathsf F_Bx_{t+1})_v|$; each internal slice occurs at most twice.
This is a group-distance bound, not the linearized frame change.

For an actual bridge kinetic word
$Q_{t,e}=r_{t,s}^{-1}Z_{t,e}r_{t,t(e)}Z_{t+1,e}^{-1}$, the exact rearrangement
\[
 Z_{t,e}Z_{t+1,e}^{-1}
 =r_{t,s}Q_{t,e}Z_{t+1,e}r_{t,t(e)}^{-1}Z_{t+1,e}^{-1}
\]
and the bi-invariant triangle inequality imply
\[
 \gamma\operatorname{dist}(Z_{t,e},Z_{t+1,e})^2
 \le3|z_{t,s}|^2+3|z_{t,t(e)}|^2+\tfrac32\pi^2\gamma s(Q_{t,e}).
\]
Each nonroot fine vertex has at most three incident bridges. Sum over all
bridges and time bonds; roots have zero port value. Since the bridge kinetic
terms are among the nonnegative terms of $U_\gamma$,
\begin{equation}
 \gamma\sum_{t,e}\operatorname{dist}(Z_{t,e},Z_{t+1,e})^2
 \le C_{\nabla} U_\gamma,
 \qquad C_{\nabla}=9C_f/a_*+3\pi^2/2.
 \label{f16v34:metric-increments}
\end{equation}
No two noncommutative factors have been interchanged.

For each bridge apply the Dirichlet time-path inequality to the scalar sequence
$u_i=\operatorname{dist}(I,Z_{i,e})$, with $u_0=u_n=0$.
The reverse triangle inequality gives
$|u_{i+1}-u_i|\le\operatorname{dist}(Z_{i+1,e},Z_{i,e})$.
The least Dirichlet eigenvalue is $4\sin^2(\pi/(2n))\ge4/n^2$, whence
\[
 \|y\|^2=\gamma\sum_{i,e}u_i^2
 \le\frac{n^2}{4}\gamma\sum_{t,e}
                  \operatorname{dist}(Z_{t,e},Z_{t+1,e})^2.
\]
Combine with the fast anchor. One permissible constant is
$c=(a_*^{-1}+C_{\nabla}/4)^{-1}$. All root, seam and endpoint incidences are
included in the counts. Taking the quadratic limit on a fixed finite vector
proves the same lower bound for $U_0$, hence positive definiteness of
$\mathcal H_{B,n}$. The quadratic form is a sum of fixed-length local squares
with bounded incidence and the original bounded endpoint quadratics. Its
absolute block rows are uniformly bounded, giving $H_*<\infty$.
\end{proof}

The $n^{-2}$ loss is not an avoidable notation issue: V33's exact period
bridges exhibit that temporal soft scale in the quadratic limit. The estimate
pins the connection through its actual endpoints; it is not a curvature-relative
mass or a pin on a history with its endpoints left free.

\subsection{A global Taylor and Haar estimate on the original charts}

\begin{proposition}[Unrestricted chart remainder with bounded word incidence]
\label{f16v34:Taylor}
There is a fixed finite $C_W$, independent of $B,n,\gamma$, such that, on
$\mathcal P_\gamma$,
\begin{equation}
 |U_\gamma(w)-U_0(w)|\le C_W\gamma^{-1/2}\|w\|^3,
 \qquad
 0\le1-\mathcal J_\gamma(w)\le\frac{\|w\|^2}{3\gamma}.
 \label{f16v34:Taylor-bounds}
\end{equation}
The latter bound retains all half-Haar endpoint powers.
\end{proposition}
\begin{proof}
Every classical Wilson term is $\gamma s$ of a fixed ordered product of
exponentials of fixed local linear forms in $w/\sqrt\gamma$. The frames have
this same form in the specified unmoving charts. Duhamel differentiation of
unitary factors gives a globally bounded third derivative with coefficient
$C\gamma^{-1/2}$. Its value and first derivative at zero vanish, and its
quadratic term is the corresponding term of $U_0$. Taylor's integral formula
therefore bounds one remainder by $C\gamma^{-1/2}\|w_T\|^3$ on its fixed
carrier $T$. Bounded carrier size and incidence yield
$\sum_T\|w_T\|^3\le C\sum_g|w_g|^3\le C\|w\|^3$.
The endpoint quadratic is exact and has no remainder. This is the ordered-word
certificate already used in V6; there is no principal logarithm of an
intermediate product to differentiate.

For $0\le t\le\pi$, $0\le\sin(t)/t\le1$ and
$1-\sin(t)/t\le t^2/6$. A Haar factor of power $1/2$ is $\sin(t)/t$;
a full factor is its square and has defect at most $t^2/3$.
For numbers in $[0,1]$, $1-\prod a_j\le\sum(1-a_j)$.
Apply this to the complete joint product, with $t=|w_g|/\sqrt\gamma$.
No logarithm of the vanishing Haar density is estimated at the cut.
\end{proof}

Put $a=c/n^2$. For any real $s\ge0$ define the explicit Gaussian radial integral
\begin{equation}
 M_s(d,a)=\int_{\mathbb R^d}|w|^s e^{-a|w|^2}\,dw
 =\pi^{d/2}a^{-(d+s)/2}
                \frac{\Gamma((d+s)/2)}{\Gamma(d/2)}.
 \label{f16v34:radial-moment}
\end{equation}
For an integer $p\ge0$ let
\[
 A_p=C_W\{M_3+M_{p+3}\}
       +\tfrac13\{M_2+M_{p+2}\}
       +\pi^{-1}\{M_1+M_{p+1}\},
\]
where these moments have arguments $(d,a)$. All quantities are finite and the
dimension is explicit.

\begin{proposition}[Unnormalized weighted comparison and the actual partition mass]
\label{f16v34:unnormalized}
For $\gamma\ge1$,
\begin{equation}
 \int(1+|w|^p)|q_\gamma-q_0|\,dw\le A_p\gamma^{-1/2}.
 \label{f16v34:unnormalized-bound}
\end{equation}
Consequently
$|\mathcal Z_\gamma-\mathcal Z_0|\le A_0\gamma^{-1/2}$ and
$\mathcal Z_\gamma/\mathcal Z_0\to1$ at fixed $B,n$.
\end{proposition}
\begin{proof}
On the complete chart, both exponents are at least $a|w|^2$. Hence
$|e^{-U_\gamma}-e^{-U_0}|\le|U_\gamma-U_0|e^{-a|w|^2}$.
The Haar defect in Proposition~\ref{f16v34:Taylor} is paid separately. This gives
\[
 |q_\gamma-q_0|\le
 [C_W\gamma^{-1/2}|w|^3+(3\gamma)^{-1}|w|^2]e^{-a|w|^2}
 \quad\text{on }\mathcal P_\gamma.
\]
Outside that product, $|w|\ge\pi\sqrt\gamma$ and $q_\gamma=0$.
Use $1_{\mathcal P_\gamma^c}\le |w|/(\pi\sqrt\gamma)$ and
$q_0\le e^{-a|w|^2}$ there. Integrating proves the stated formula. In
particular the actual full chart mass, not a restricted mass, is compared.
\end{proof}

\subsection{Complete native convergence, including a simultaneous growing regime}

Write $p_\gamma=q_\gamma/\mathcal Z_\gamma$ and $p_0=q_0/\mathcal Z_0$.
The subscript zero here labels the Gaussian probability, not zero coupling.

\begin{theorem}[Quantitative full-native Laplace comparison]
\label{f16v34:full-limit}
For each fixed $B,n$ and every integer $p\ge0$,
\begin{equation}
 \boxed{\quad
 \int_{\mathbb R^d}(1+|w|^p)|p_\gamma-p_0|\,dw
            \le C_{p,B,n}\gamma^{-1/2}
 \quad}
 \label{f16v34:weighted-limit}
\end{equation}
at sufficiently large coupling. Thus the complete native joint probability
converges in total variation, with all fixed polynomial moments, to the exact
untruncated joint Gaussian. Its retained marginal converges to $Q_{B,n}^{00}$
in the same sense. No fast or retained cutoff remains in these probabilities.

There are constants $C_p<\infty$, independent of $B,n,\gamma$, for which the
right side may be bounded by
\begin{equation}
 \gamma^{-1/2}\exp\{C_p d_{B,n}\log(n+1)\},
 \label{f16v34:dimension-rate}
\end{equation}
provided the analogous expression with $C_0$ is at most a fixed sufficiently
small constant. In particular, along any sequence with
\begin{equation}
            Bn\log(n+1)=o(\log\gamma),
 \label{f16v34:joint-regime}
\end{equation}
the weighted difference in \eqref{f16v34:weighted-limit} tends to zero for
every fixed $p$. This is a conservative slowly growing regime, not V30's much
larger fixed-window width regime or a uniform-in-duration theorem.
\end{theorem}
\begin{proof}
If $A_0/\sqrt\gamma\le\mathcal Z_0/2$, the actual normalizer satisfies
$\mathcal Z_\gamma\ge\mathcal Z_0/2$. Let
$N_p=\int(1+|w|^p)q_0(w)\,dw$. Direct subtraction of the two normalized
densities gives
\begin{equation}
 \int(1+|w|^p)|p_\gamma-p_0|\,dw
 \le\frac2{\mathcal Z_0\sqrt\gamma}
                    \left(A_p+\frac{A_0N_p}{\mathcal Z_0}\right).
 \label{f16v34:normalized-bound}
\end{equation}
This includes both partition functions. Marginalization contracts the weighted
bound for weights depending only on retained variables.

We give the dimension estimate to specify a genuine simultaneous regime.
The upper Hessian bound gives
$\mathcal Z_0\ge(2\pi/H_*)^{d/2}$. Each ratio $M_s/\mathcal Z_0$ is at most
\[
 (H_*/(2a))^{d/2}a^{-s/2}
                         \frac{\Gamma((d+s)/2)}{\Gamma(d/2)}.
\]
For fixed integer $s$, the last ratio is at most
$C_s(d+s)^{s/2}$. For even $s$ this is the finite gamma-function product;
for odd $s$, Cauchy--Schwarz between the adjacent even radial moments proves it.
Also $\mathcal H_{B,n}\succeq2aI$, so the Gaussian moment
$N_p/\mathcal Z_0$ is at most $1+C_p a^{-p/2}(d+p)^{p/2}$.
Since $a=c/n^2$ and $d\ge1$, all these factors are bounded by
$\exp\{C_p d\log(n+1)\}$ after fixed constants are enlarged. Inserting in
\eqref{f16v34:normalized-bound} proves \eqref{f16v34:dimension-rate} and the
stated normalizer threshold. Relation \eqref{f16v34:joint-regime} makes each
fixed exponent $C_p d\log(n+1)=o(\log\gamma)$, so the error is
$\gamma^{-1/2+o(1)}$. No radius-dependent inherited threshold is used.
\end{proof}

At fixed $B,n$, the compact domination also gives a uniform native tail, once
$\mathcal Z_\gamma\ge\mathcal Z_0/2$:
\begin{equation}
 \mathbb P_\gamma\{|w|>R\}
 \le \frac2{\mathcal Z_0}(2\pi/a)^{d/2}e^{-aR^2/2}.
 \label{f16v34:tail}
\end{equation}
The dimension and $a=c/n^2$ are deliberately visible. This proves tightness
as $\gamma\to\infty$ at each fixed regulator; it is not a bound uniform as
$n\to\infty$. The estimate also gives uniform exponential moments with any
fixed exponent smaller than $a$, at that fixed regulator.

\subsection{Period diffusion now follows under the full native law}

\begin{theorem}[Two-sided window comparison and a local midpoint consequence]
\label{f16v34:periods}
For each fixed $B,n,r$, with exactly the probabilities and event of V33,
\begin{equation}
 \lim_{\gamma\to\infty}
       \boldsymbol\nu_{\gamma,B,n}^{00}(G_r)=Q_{B,n}^{00}(G_r).
 \label{f16v34:window-equality}
\end{equation}
The full vector of V33's nine period amplitudes converges in total variation
and moments to its exactly normalized Gaussian bridge law. For even $n$ and
any specified bridge $e$ at time $n/2$,
\begin{equation}
 \lim_{\gamma\to\infty}\operatorname{Var}_{\boldsymbol\nu_\gamma}
                 (y_{n/2,e}^{a})\ge\frac{n}{8B},
 \label{f16v34:local-variance}
\end{equation}
for each colour component. The corresponding one-bridge probability obeys
\begin{equation}
 \lim_{\gamma\to\infty}\boldsymbol\nu_\gamma
        \{|y_{n/2,e}|\le r\}
 \le\left[\min\left\{1,4r\sqrt{\frac B{\pi n}}\right\}\right]^3.
 \label{f16v34:local-cutoff}
\end{equation}
These are rescaled coordinate statements, not claims about nonlinear
period holonomies or the physical mass.
\end{theorem}
\begin{proof}
Total variation proves \eqref{f16v34:window-equality} for the original event,
with no exceptional-boundary qualification. Orthogonal linear projection onto
the period coordinates contracts total variation, and the weighted convergence
with $p=2$ proves covariance convergence. The native means vanish by simultaneous
global conjugation at identity endpoints; alternatively the same convergence
centers them correctly.

V33's reducing Gaussian sector has midpoint amplitude covariance $(n/2)I_9$.
A bridge of direction $\mu$ has coefficient $1/(2\sqrt B)$ in each of its
three period directions. The complementary Gaussian contributes a nonnegative
covariance and is independent of those amplitudes. Thus its three-component
marginal covariance is $\sigma^2I_3$ with $\sigma^2\ge n/(8B)$. Colour
isotropy follows from the actual quadratic colour tensor $I_3$. This proves
the variance assertion. Each independent scalar Gaussian component has
interval probability at most $2r/\sqrt{2\pi\sigma^2}\le
4r\sqrt{B/(\pi n)}$. The ball event is a subset of the three interval events.
Apply total variation to obtain \eqref{f16v34:local-cutoff}.
\end{proof}

\begin{proposition}[An explicit slowly growing native period process]
\label{f16v34:process}
Fix $B$. For sufficiently large $\gamma$, choose the even integer
\[
 n_\gamma=2\left\lfloor\frac{\sqrt{\log\gamma}}2\right\rfloor.
\]
Linearly interpolate V33's native amplitudes $a_i/\sqrt{n_\gamma}$ at times
$i/n_\gamma$. Under the \emph{full} endpoint-fixed native law this continuous
nine-component process converges weakly to the centered continuous Gaussian
bridge with covariance
\begin{equation}
              2(\min(s,t)-st)I_9,\qquad 0\le s,t\le1.
 \label{f16v34:process-limit}
\end{equation}
For every fixed $r>0$, the native probability
$\boldsymbol\nu_{\gamma,B,n_\gamma}^{00}\{|y_{n_\gamma/2,e}|\le r\}$
tends to zero, as does the all-time window probability. The native midpoint
variance is at least $n_\gamma/(8B)-o(1)$.
\end{proposition}
\begin{proof}
For this choice $Bn_\gamma\log(n_\gamma+1)=o(\log\gamma)$.
Theorem~\ref{f16v34:full-limit} therefore compares the entire increasing-dimensional
native law and its Gaussian with error $\gamma^{-1/2+o(1)}$, including second
moments. Under that Gaussian, the interpolated amplitudes have exactly the law
of the polygonal interpolation of one continuous variance-two Gaussian bridge
at the uniform mesh. On a common realization these polygonal paths converge
uniformly by continuity. Pushforward contracts total variation, so the native
path laws have the same weak limit. This is a probabilistic time parametrization
of a specified limit, not calibrated physical time.

The proof of \eqref{f16v34:local-cutoff} applies to the Gaussian at each
$n_\gamma$, and its right side tends to zero. The total-variation error tends
to zero as well. The all-time window is a subset of this single-time event.
The weighted $p=2$ bound proves the stated absolute $o(1)$ covariance error.
Angular coordinates still satisfy $y/\sqrt\gamma\to0$ in this regime; their
rescaled spreading does not mean escape from a fixed angular neighbourhood.
\end{proof}

\subsection{Positive native derivative and Wilson-word covariance statements}

Retain V33's assembled map
$\mathcal B y=((G^{-1/2}\Delta y_t)_t,(Dy_i/\sqrt3)_i)$.
It obeys $\mathcal B^{\mathsf T}\mathcal B\preceq K^{\rm int}$ and has bounded
operator norm, independent of $B,n$. Define
\[
 \varepsilon_{p,B,n,\gamma}
       =\gamma^{-1/2}\exp\{C_p d_{B,n}\log(n+1)\}.
\]
Constants can be enlarged once when passing to a fixed finite-degree observable.

\begin{theorem}[Native covariance in the derivative bank and the actual word bank]
\label{f16v34:fields}
At the normalizer threshold above,
\begin{equation}
 \operatorname{Cov}_{\boldsymbol\nu_\gamma}(\mathcal B y)
              \preceq(1+\varepsilon_{2,B,n,\gamma})I.
 \label{f16v34:native-derivative}
\end{equation}
In particular the complete mixed bank of temporal differences and magnetic
curvature has the V33 bound with factor $1+\varepsilon_{2,B,n,\gamma}$.
There is no retained radius. The error tends to zero at fixed $B,n$, and in
\eqref{f16v34:joint-regime}; no error uniform in all durations is asserted.

There is also a genuinely nonlinear observable statement on the joint law.
List each original classical Wilson word once, including its original magnetic
weight $\omega_\ell\in\{1/2,1\}$ and all kinetic words. Define its quaternion
vector observable by
\begin{equation}
 \mathcal X_{\gamma,\ell}(w)
   =\sqrt{\gamma\omega_\ell}\,\boldsymbol{\operatorname{Im}}W_\ell(w/\sqrt\gamma).
 \label{f16v34:word-fields}
\end{equation}
Then, with constants of the same explicit dimension-growth type,
\begin{equation}
 \operatorname{Cov}_{\mathbb P_\gamma}(\mathcal X_\gamma)
                         \preceq(1+\varepsilon_{4,B,n,\gamma})I.
 \label{f16v34:word-covariance}
\end{equation}
Every covariance is centered under its indicated native law and includes the
cross covariances of the assembled bank. These word vectors are not claimed
to be gauge-invariant scalar observables; their frames and order are unchanged.
\end{theorem}
\begin{proof}
The weighted second-moment comparison bounds the operator-norm difference of
retained covariances: for a unit vector $v$, $(v^{\mathsf T}y)^2\le|w|^2$.
The Gaussian derivative covariance is at most $I$ by V33. The uniform norm of
$\mathcal B$ and the weighted comparison prove \eqref{f16v34:native-derivative}.
Equivalently, testing its output against $((G^{1/2}a_t)_t,(\sqrt3b_i)_i)$
bounds the native variance by
$(1+\varepsilon_{2,B,n,\gamma})(30\sum_t|a_t|^2+3\sum_i|b_i|^2)$.

Let $D_\ell w$ be the first quaternion-vector differential of word $\ell$
at the identity, and let $\mathcal D$ stack $\sqrt{\omega_\ell}D_\ell$.
The exact quadratic action has the form
\[
 U_0(w)=\tfrac12|\mathcal D w|^2+E_{\rm end}(w),\qquad E_{\rm end}\ge0,
\]
so $\mathcal D^{\mathsf T}\mathcal D\preceq\mathcal H_{B,n}$. Gaussian
integration therefore gives
$\operatorname{Cov}_{\mathbb Q}(\mathcal D w)
=\mathcal D\mathcal H_{B,n}^{-1}\mathcal D^{\mathsf T}\preceq I$.
No word covariance is counted independently of the others in this inequality.

The ordered unitary second-derivative estimate gives the complete-chart bound
$\|\mathcal X_\gamma-\mathcal D w\|\le C\gamma^{-1/2}|w|^2$ by bounded
incidence, and $\|\mathcal X_\gamma\|+\|\mathcal D w\|\le C|w|$.
The full-law moment estimates control its covariance error using moments
through degree four. The normalizer threshold and Gaussian moment bounds
above bound all these moments by the same
$\exp\{C_p d\log(n+1)\}$ type of constant. Combine with the weighted-law
comparison for $\mathcal D w$. This proves \eqref{f16v34:word-covariance}.
Neither a relative lower Hessian for the nonlinear retained action nor a
native diffusion inverse has been assumed.
\end{proof}

\begin{proposition}[A native retained-sampler upper test, not a physical spectrum]
\label{f16v34:gap-test}
Let $\lambda_{\gamma,B,n}^{\rm br}$ be the variational gap of the sum of
rate-one \emph{full} retained-bridge heat baths under
$\boldsymbol\nu_{\gamma,B,n}^{00}$. At fixed $B,n$,
\begin{equation}
       \limsup_{\gamma\to\infty}\lambda_{\gamma,B,n}^{\rm br}
                         \le30\sin^2\frac\pi{2n}.
 \label{f16v34:native-gap-upper}
\end{equation}
On the sequence of Proposition~\ref{f16v34:process}, the gap is $O(n_\gamma^{-2})$.
This is not an assertion of update-operator convergence.
\end{proposition}
\begin{proof}
Use V33's unit first temporal period eigenvector $v$ and the linear test
$f=v^{\mathsf T}y$. Its Gaussian variance is
$[2\sin^2(\pi/(2n))]^{-1}$. At a retained site $u$, use as an admissible
\emph{native} exterior predictor the exact Gaussian conditional mean
$g_u(y_{u^c})=\mathbb E_Q[f\mid y_{u^c}]$. The minimizing property of native
conditional expectation gives
\[
 \mathbb E_{\boldsymbol\nu_\gamma}\operatorname{Var}(f\mid y_{u^c})
 \le \mathbb E_{\boldsymbol\nu_\gamma}|f-g_u|^2.
\]
These are quadratic polynomials. Their Gaussian sum is
$\sum_uv_u^{\mathsf T}(K^{\rm int}_{uu})^{-1}v_u\le15$, using V33's
$K^{\rm int}_{uu}\succeq I_3/15$. The weighted convergence gives the fixed
regulator limsup after division by the convergent positive variance.
For the simultaneous assertion, the coefficient operator of
$\sum_u|f-g_u|^2$ is uniformly bounded: the diagonal inverses are at most
15 and the rows of $K^{\rm int}$ are uniformly bounded. Hence the weighted
error still tends to zero in \eqref{f16v34:joint-regime}. The denominator is
$[2\sin^2(\pi/(2n_\gamma))]^{-1}+o(1)$, proving the stated order.
This is a variational upper estimate using fixed Gaussian predictors, not
conditional-kernel or spectral convergence inferred from total variation.
\end{proof}

\subsection{What has been removed and what still requires new information}

V34 pays the complete compact tail and the actual full normalizer at fixed
endpoint data. The Gaussian and native laws now agree in a two-sided,
untruncated comparison at fixed regulator and along the explicitly slow
simultaneous regime. V33's merely one-sided cutoff argument is thereby
strengthened, not withdrawn. The new local midpoint statement was not a
consequence of that earlier nested-window upper bound alone.

The rate has a large dimension price and a duration-dependent coercivity scale.
Neither disappears from the theorem. In particular the positive native field
covariance statement is not the sought unrestricted-duration relative
curvature estimate. The next useful target is to reduce the endpoint/dimension
loss specifically in the derivative bank, or keep the nonlinear compact period
sector nonperturbatively while controlling its conditional complement. Another
fixed-window boundary gap would not address either task.

The two supplied memoranda distinguish perturbative gain from locality and
normalization costs, and an auxiliary functional inequality from its physical
transfer implication. Here the Taylor power is $\gamma^{-1/2}$, while the
explicit cost is $\exp\{C_p Bn\log(n+1)\}$. Those ledgers are not conflated.
All root gauge and physical endpoint issues remain: the law is the original
\emph{unprojected} identity-endpoint history, not a stationary physical vacuum.
A long-history bridge-sampler upper bound does not prove a gapless physical
Hamiltonian; even massive physical transfer bridges can have algorithmic slow
modes. No physical time is set by the interpolation in
\eqref{f16v34:process-limit}.

\subsection{Verification and preservation}

The accompanying diagnostic is
\begin{center}\small\path{verify_ym_f16_v34_endpoint_cutoff.py}.\end{center}
It checks exact group-word rearrangements, the literal bridge-degree and frame
energy counts, scalar metric-path bounds, Haar endpoint exponents, Gaussian
normalizers, weighted-density algebra, period covariances and variational
predictors. The new script passes 259 exact assertions. All 30 supplied V4--V33 checkers
were rerun unchanged and passed. These are finite rational and symbolic
fixtures, not native Wilson integral enclosures or formal verification. Actual
outputs and hashes are recorded separately.

Only the single title version and this terminal insertion differ from V33.
Deleting the insertion and restoring that title recovers the supplied V33 byte
for byte. No old formula, source, proof or reference is silently edited. The
compact pin uses the inherited finite cube and frame estimates; this is not an
independent recertification of the complete archive. The next checkpoint is V35.

\begin{firewall}
V34 proves full endpoint-anchored compact coercivity, a quantitative normalized
joint native Gaussian comparison without a retained cutoff, and polynomial
moment convergence at fixed regulator and in the stated slow simultaneous
regime. It obtains native derivative and Wilson-word covariance comparisons,
a quantitative period process, local rescaled-window escape and a native
retained-sampler upper test. It does not prove uniformity at arbitrary spatial
volume or duration, a stationary gauge-invariant vacuum limit, unrestricted
relative native curvature, a physical transfer gap, or an interacting continuum
mass gap. All complete integration domains, original Haar powers, endpoint
normalizers and sampler rates used here are retained.
\end{firewall}
% END F16 V34 APPEND

% BEGIN F16 V35 APPEND -- 2026-09-18
% Insert immediately before the final document-ending command of V34.
% No predecessor assertion, native measure, endpoint factor, or clock is changed.

\clearpage
\section{V35: polynomial-regime native comparison and relative field covariance}
\label{f16v35:context}

V34 proves a complete normalized comparison, but charges its coarse radial
majorant throughout the integration domain. The resulting exponential
size factor limits its simultaneous application to logarithmically growing
histories. We separate two tasks which need different estimates. The coarse
compact pin pays only the tails. On an explicitly chosen interior ball, the
actual measured Hessian is compared with the \emph{full endpoint-fixed}
Gaussian precision. A normalized square-root-density estimate then measures
the correction by its score, rather than by the absolute oscillation of an
extensive action. This gives polynomially growing native histories and a
relative covariance estimate, including the temporal soft directions.

The comparison ball is a proof device, not a new retained window or a change
to V4's core event. Its complement is paid under both complete laws. The
endpoint and size conditions remain substantive. In particular, the new
covariance comparison does not give a native pointwise relative Hessian bound,
a physical transfer comparison, or a duration-uniform mass.

\subsection{V35 checkpoint and unchanged probability}

The complete supplied V34 appendix and audit were read. The new argument uses
its compact endpoint pin, V5's independent fast anchor, the actual ordered-word
derivative bounds, and V33's finite Gaussian matrices. It does not use the
V16 moment iteration, V30 spatial-screening theorem, or V31's radius-dependent
constants. Those results remain separate, with their original hypotheses.
The companion checkpoint records targeted inventories of the current lineage,
f14 V100, f15 V100, and Grand Tensor V076; it is not independent certification
of their proofs or a claim that unavailable companions were examined. The
next scheduled checkpoint is V40.

The interior/tail split is already present in V5's finite-box Laplace proof.
The monograph also uses normalized Hellinger embeddings. Neither generic
principle is claimed new. The additional step here is a score estimate in the
endpoint Gaussian metric, together with an explicitly removed comparison ball.
For the boundary-aware variance principle, see Kolesnikov--Milman,
\href{https://arxiv.org/abs/1310.2526}{\emph{Brascamp--Lieb type inequalities on
weighted Riemannian manifolds with boundary}}, Theorem 1.2(1) at $N=\infty$.
Its boundary statement and Reilly formula were checked; the special estimate
used below is also proved directly. Spokoiny's
\href{https://arxiv.org/abs/2204.11038}{\emph{Dimension free non-asymptotic bounds
on the accuracy of high dimensional Laplace approximation}} was screened for
context only, not used as a black-box native theorem.

Fix identity retained endpoints, $B=m^3$, $m\ge4$, $n\ge2$. Keep exactly V34's
joint probability $\mathbb P_\gamma$, its Gaussian $\mathbb Q$, and their full
normalizers $\mathcal Z_\gamma,\mathcal Z_0$. Write
\[
 w=(\zeta,y)\in\mathbb R^d,\qquad d=(72n-21)B,\qquad
 U_0(w)=\tfrac12 w^{\mathsf T}\mathcal H w,\qquad \mathcal C=\mathcal H^{-1}.
\]
Thus $\mathcal H$ is the full \emph{joint} precision, not the fast conditional
precision $H$, nor the retained Schur precision $K^{\rm int}$. All chord
endpoint Haar powers are $1/2$; all other free Haar powers are one. All fast
and interior bridge groups remain integrated over their complete principal
charts. Scalar dressed-end factors cancel only in the same endpoint-fixed
probability, not from a proposed physical transfer normalization.

Shrink V34's fixed constant $c$ and enlarge $H_*$ once so that
$0<c\le1/4$, $H_*\ge1$. Put $a=c/n^2$. Its global estimates are
\begin{equation}
 U_\gamma(w),U_0(w)\ge a|w|^2,\qquad
 2aI\preceq\mathcal H\preceq H_*I,\qquad 0\le\mathcal J_\gamma\le1.
 \label{f16v35:global-data}
\end{equation}
The first bound for $U_\gamma$ is only on its complete principal product, as
in V34; the native density is zero off that product.

\subsection{Local Gaussian scales and a relative interior Hessian}

\begin{proposition}[Gaussian coordinate scales and local measured derivatives]
\label{f16v35:local-data}
There are fixed $C_2,C_H,C_G<\infty$, independent of $B,n,\gamma$, such that
\begin{equation}
 \|\mathcal C\|\le(2a)^{-1},\qquad
 \mathcal C_{\zeta\zeta}\preceq(2a_*)^{-1}I,\qquad
 \mathcal C_{yy}\preceq30(L_n^{-1}\otimes I),\qquad
 \sup_g\|\mathcal C_{gg}\|\le C_G n.
 \label{f16v35:Gaussian-scales}
\end{equation}
Here $g$ labels an original three-component group and $a_*=(6\pi^2)^{-1}$.
On $|w|\le R\le\sqrt\gamma/2$, set
\[
 V_\gamma=U_\gamma-\log\mathcal J_\gamma,\qquad W_\gamma=V_\gamma-U_0.
\]
Then
\begin{equation}
 \|D^2W_\gamma(w)\|\le C_2R/\sqrt\gamma+C_H/\gamma,
 \qquad
 |\nabla W_\gamma(w)|^2\le
 \frac{C_2}{\gamma}\sum_g|w_g|^4+
 \frac{C_H}{\gamma^2}|w|^2.
 \label{f16v35:local-derivatives}
\end{equation}
The constants account for the fixed local carriers and their actual incidence,
not for independent words or independent coordinates.
\end{proposition}
\begin{proof}
The first bound is \eqref{f16v35:global-data}. To prove the fast marginal bound,
minimize the quadratic $U_0$ in $y$ at fixed $\zeta$. V5's anchor gives
$\inf_y U_0\ge a_*|\zeta|^2$, so its fast marginal precision is at least
$2a_*I$. This is a marginal statement, not a substitution of the conditional
fast covariance. The retained marginal is exactly V33's Gaussian, whose
precision is at least $(L_n\otimes I)/30$. The diagonal of $L_n^{-1}$ is
$i(n-i)/n\le n/4$. These two facts give the group bound, including the
unmoving temporal ports and the endpoint chord groups.

For a scalar derivative in one group, only a fixed number of ordered words
contribute. V34's unitary third-derivative argument gives a Hessian remainder
at most $C\gamma^{-1/2}|w_T|$ and a gradient remainder at most
$C\gamma^{-1/2}|w_T|^2$ on each such fixed carrier $T$. The zero and linear
Taylor terms vanish; the endpoint quadratic is exact. The block-row Schur
bound, fixed carrier size, and bounded row and column incidence prove the first
classical estimate and
$|\nabla(U_\gamma-U_0)|^2\le C\gamma^{-1}\sum_g|w_g|^4$.
This retains the local sum instead of replacing it by $|w|^4$.

On angles at most $1/2$, the measured Haar Hessian is positive and bounded
above by $C_H/\gamma$ on each three-group, with gradient bounded by
$C_H|w_g|/\gamma$. The endpoint exponent $1/2$ is included. Increasing the
fixed constants proves \eqref{f16v35:local-derivatives}. No logarithm of the
Haar density is used near its zero at the original principal cut.
\end{proof}

Fix $s\ge1$, $\gamma\ge2$, and define the explicit tail budget and radius
\begin{equation}
 \begin{split}
 \Lambda_s={}&1+\frac d2\log\frac{2H_*}{a}
       +\log\left[1+\left(\frac{d+6}{a}\right)^4\right]
       +(s+1)\log\gamma,\qquad R_s^2=\frac{2\Lambda_s}{a},\qquad
 \mathcal B_s=\{|w|\le R_s\}.
 \end{split}
 \label{f16v35:radius}
\end{equation}
The finite-parameter entry conditions used below are
\begin{equation}
 R_s\le\sqrt\gamma/2,\qquad
 C_2R_s/\sqrt\gamma+C_H/\gamma\le a.
 \label{f16v35:entry}
\end{equation}
They are checkable in the inherited fixed constants. They imply
\begin{equation}
       \tfrac12\mathcal H\preceq D^2V_\gamma\preceq2\mathcal H
            \quad\hbox{on }\mathcal B_s.
 \label{f16v35:interior-curvature}
\end{equation}
Indeed a symmetric error of norm at most $a$ is at most $\mathcal H/2$ in
absolute quadratic form. Both potentials and their gradients vanish at zero,
so $V_\gamma\le2U_0$ on this ball. This is an \emph{interior} comparison;
no such native Hessian inequality is asserted on the full principal product.

\begin{proposition}[Complete normalized tails, paid before the interior comparison]
\label{f16v35:tails}
Under \eqref{f16v35:entry},
\begin{equation}
 \mathcal Z_\gamma\ge2^{-d/2-1}\mathcal Z_0,
 \qquad
 \mathbb E_{\mathbb P_\gamma}[(1+|w|^8)1_{\mathcal B_s^c}]
 +\mathbb E_{\mathbb Q}[(1+|w|^8)1_{\mathcal B_s^c}]
 \le2\gamma^{-s}.
 \label{f16v35:tails-bound}
\end{equation}
The normalizer inequality is only a lower bound. In particular it is not a
claim that $\mathcal Z_\gamma/\mathcal Z_0$ is close to one uniformly in the
new growing regime.
\end{proposition}
\begin{proof}
The Gaussian of precision $2\mathcal H$ has
$\mathbb E|w|^2\le d/(4a)$. The radius \eqref{f16v35:radius} has
$R_s^2\ge d/(2a)$, so at least half this Gaussian lies in $\mathcal B_s$.
The ball lies strictly inside every original principal group chart.
Since $V_\gamma\le2U_0$ there,
\[
 \mathcal Z_\gamma\ge\int_{\mathcal B_s}e^{-2U_0}
           \ge2^{-d/2-1}\mathcal Z_0.
\]
This uses the actual full native mass and its original Haar powers.

For either tail, use the global classical pin and $\mathcal J_\gamma\le1$;
off the native chart the native density remains zero. For $t=|w|>R_s$,
$e^{-at^2}\le e^{-aR_s^2/2}e^{-at^2/2}$. The exact radial eighth moment gives
\[
 \int(1+|w|^8)e^{-a|w|^2/2}\,dw
 \le(2\pi/a)^{d/2}\left[1+\left((d+6)/a\right)^4\right].
\]
Also $\mathcal Z_0\ge(2\pi/H_*)^{d/2}$. Dividing by the proved native mass
lower bound bounds its weighted tail by
\[
 2\exp\left\{\frac d2\log(2H_*/a)
 +\log[1+((d+6)/a)^4]-\frac{aR_s^2}{2}\right\}
 =2e^{-1}\gamma^{-(s+1)}\le\gamma^{-s}.
\]
The Gaussian needs no factor $2^{d/2+1}$ and obeys the same weaker bound.
Their sum proves the result. All eighth-degree tail weights are explicit;
the size loss has been paid only in choosing the radius, not multiplied into
the later interior error.
\end{proof}

\subsection{Normalized square roots instead of a global likelihood-ratio bound}

We record the needed covariance and normalization calculation on a smooth
Euclidean ball. It does not use a product-ball one-form locality claim.

\begin{proposition}[A convex-domain variance bound and a normalized Hellinger estimate]
\label{f16v35:Hellinger}
Let $D$ be a bounded smooth convex domain, and let $\mu\propto e^{-V}$ be smooth
and positive there, with $D^2V\succeq A\succ0$. Then
\begin{equation}
       \operatorname{Var}_\mu f
          \le\mathbb E_\mu[\nabla f^{\mathsf T}A^{-1}\nabla f].
 \label{f16v35:convex-Poincare}
\end{equation}
If $Q_D\propto e^{-w^{\mathsf T}\mathcal H w/2}$ and
$P_D\propto e^{-w^{\mathsf T}\mathcal H w/2-W(w)}$ are two normalized laws on
that same domain, put
$h^2(P_D,Q_D)=\int(\sqrt{p_D}-\sqrt{q_D})^2$. Then
\begin{equation}
 \boxed{\quad h^2(P_D,Q_D)
       \le\tfrac12\mathbb E_{P_D}
                   [\nabla W^{\mathsf T}\mathcal H^{-1}\nabla W],
       \qquad d_{\rm TV}(P_D,Q_D)\le h(P_D,Q_D).\quad}
 \label{f16v35:Hellinger-bound}
\end{equation}
This statement includes both normalizers. It does not require an estimate
that their ratio is near one, or a small supremum norm of $W$.
\end{proposition}
\begin{proof}
For centered smooth $f$, solve the weighted Neumann Poisson equation
$\mathcal Au=f$. The integrated Bochner identity has the nonnegative outward
second-fundamental-form term on $\partial D$ and implies
$\|f\|_2^2\ge\mathbb E_\mu[\nabla u^{\mathsf T}A\nabla u]$.
Integration by parts and Cauchy--Schwarz in the constant metric $A$ give
\[
 \|f\|_2^2=\mathbb E_\mu[\nabla f\cdot\nabla u]
 \le\bigl(\mathbb E_\mu[\nabla f^{\mathsf T}A^{-1}\nabla f]\bigr)^{1/2}
                         \|f\|_2.
\]
This proves \eqref{f16v35:convex-Poincare}. At a fixed finite dimension the
smooth positive weight on the compact smooth ball gives the usual Neumann
realization and Poisson inverse on constants' complement; equivalence of
Sobolev norms is used only for this domain statement. Smooth approximation
extends the inequality to $H^1$. The boundary term was retained before its
nonnegative sign was used.

Apply this inequality under $Q_D$ to $b=\sqrt{p_D/q_D}$.
One has $\mathbb E_{Q_D}b^2=1$ and
$\nabla b=-b\nabla W/2$. If $m=\mathbb E_{Q_D}b$, then $0<m\le1$ and
\[
 h^2=2(1-m)\le2(1-m^2)=2\operatorname{Var}_{Q_D}b
       \le\tfrac12\mathbb E_{P_D}[\nabla W^{\mathsf T}\mathcal C\nabla W].
\]
The density ratio in $b$ is \emph{normalized}; its scalar derivative is zero
in $w$, not omitted from its definition. Cauchy--Schwarz applied to
$(\sqrt p-\sqrt q)(\sqrt p+\sqrt q)$ gives $d_{\rm TV}\le h$ in this convention.
\end{proof}

Let $P_s=\mathbb P_\gamma(\cdot\mid\mathcal B_s)$ and
$Q_s=\mathbb Q(\cdot\mid\mathcal B_s)$. The subscript here labels the tail
budget, not an interpolation of the native couplings. Simultaneous global
conjugation acts by a common orthogonal rotation on all three-groups and
preserves both balls and both densities. Thus $\mathbb E_{P_s}w=
\mathbb E_{Q_s}w=0$, as also holds for the two complete laws.

\begin{proposition}[Directional moments and the local-score estimate]
\label{f16v35:score}
Under \eqref{f16v35:entry}, for every real coefficient vector $u$,
\begin{equation}
 \mathbb E_{P_s}e^{t u^{\mathsf T}w}
       \le e^{t^2u^{\mathsf T}\mathcal C u},\qquad
 \mathbb E_{P_s}|u^{\mathsf T}w|^4
       \le C(u^{\mathsf T}\mathcal C u)^2.
 \label{f16v35:directional-moments}
\end{equation}
The same weaker bounds hold under $Q_s$. In particular
\begin{equation}
 \mathbb E_{P_s}[\nabla W_\gamma^{\mathsf T}\mathcal C\nabla W_\gamma]
                   \le C\frac{d n^4}{\gamma}.
 \label{f16v35:score-cost}
\end{equation}
\end{proposition}
\begin{proof}
Every linear tilt of $P_s$ has the same Hessian lower bound
$\mathcal H/2$ and the same convex ball. Apply
\eqref{f16v35:convex-Poincare} to $u^{\mathsf T}w$ under that tilted law.
The second derivative of its log moment-generating function is at most
$2u^{\mathsf T}\mathcal C u$. Its value and first derivative at zero are zero,
which proves the first estimate by two integrations. Exponential Markov and
integration of its two-sided Gaussian tail give the fourth-moment estimate.
This does not assume that the coordinates are independent. The argument for
$Q_s$ has the stronger Hessian $\mathcal H$.

The local Gaussian variance in \eqref{f16v35:Gaussian-scales} therefore gives
$\mathbb E_{P_s}|w_g|^4\le C n^2$ for each fixed three-group, and
$\mathbb E_{P_s}|w|^2\le Cdn$. Use the \emph{local} sum in
\eqref{f16v35:local-derivatives} to get
$\mathbb E_{P_s}|\nabla W_\gamma|^2\le Cdn^2/\gamma$.
Finally $\|\mathcal C\|\le Cn^2$ proves \eqref{f16v35:score-cost}.
The $n^2$ inverse cost is real; it has not been replaced by a positive mass.
\end{proof}

\subsection{Polynomially growing complete native histories}

Set
\[
            e_{d,n,\gamma,s}=\sqrt{\frac{dn^4}{\gamma}}+\gamma^{-s}.
\]

\begin{theorem}[Full native total variation with a polynomial size cost]
\label{f16v35:TV}
Under the explicit entry conditions \eqref{f16v35:entry},
\begin{equation}
 \boxed{\qquad d_{\rm TV}(\mathbb P_\gamma,\mathbb Q)
               \le C e_{d,n,\gamma,s}.\qquad}
 \label{f16v35:TV-bound}
\end{equation}
The same bound holds for their complete retained marginals and every measurable
pushforward. No proof cutoff remains in these laws.
There are fixed positive constants such that, for each fixed $s$, a sufficient
large-coupling entry condition is
\begin{equation}
 n^6\{d\log(n+1)+(s+1)\log\gamma\}\le c_0\gamma,
 \label{f16v35:polynomial-entry}
\end{equation}
with $c_0$ sufficiently small in the named fixed geometric constants.
In particular, the estimate tends to zero along any sequence with
\begin{equation}
 \frac{Bn^7\log(n+1)+n^6\log\gamma}{\gamma}\longrightarrow0.
 \label{f16v35:polynomial-regime}
\end{equation}
If $B=O(\gamma^b)$ and $n=O(\gamma^\alpha)$ with $b,\alpha\ge0$,
$b+7\alpha<1$ is a sufficient strict exponent condition.
\end{theorem}
\begin{proof}
Propositions~\ref{f16v35:Hellinger} and \ref{f16v35:score} bound
$d_{\rm TV}(P_s,Q_s)$ by $C\sqrt{dn^4/\gamma}$. For any probability $\mu$,
\[
 d_{\rm TV}(\mu,\mu(\cdot\mid\mathcal B_s))=\mu(\mathcal B_s^c).
\]
Use this identity twice and \eqref{f16v35:tails-bound} to obtain
\eqref{f16v35:TV-bound}. The principal-chart complement is included in those
tails. This is not a comparison of conditioned laws relabelled as full laws.

For the entry condition note that $a=c/n^2$ and
\[
 \Lambda_s\le C_s\{d\log(n+1)+(s+1)\log\gamma\}.
\]
Here $\log(d+6)\le C d$; there is no hidden exponential in $d$. Hence
\[
 R_s\le C_s n\sqrt{d\log(n+1)+(s+1)\log\gamma}.
\]
A sufficiently small
constant in \eqref{f16v35:polynomial-entry} ensures
$C_2R_s/\sqrt\gamma\le a/2$, $C_H/\gamma\le a/2$, and
$R_s\le\sqrt\gamma/2$. Finite fixed thresholds cover the other elementary
requirements. With $d=(72n-21)B$, \eqref{f16v35:polynomial-regime} implies this
entry and $dn^4/\gamma\to0$. The strict exponent example follows directly,
including rounded admissible spatial sides and bounded time lengths.
\end{proof}

The gain is not a claim of an optimal high-dimensional Laplace rate. One still
pays $n^2$ in the inverse and a coarse $n^{-2}$ pin in preparing the comparison
ball. The improvement over V34 is that its exponential radial ratio now
occurs inside the tail logarithm, whereas the normalized interior score pays
$\sqrt d$. No higher perturbative order or convergence of a series is used.

\begin{theorem}[Relative covariance in the original Gaussian metric]
\label{f16v35:covariance}
Let $\Sigma_\gamma=\operatorname{Cov}_{\mathbb P_\gamma}(w)$. Under the same
entry conditions,
\begin{equation}
 \boxed{\quad
 \|\mathcal H^{1/2}\Sigma_\gamma\mathcal H^{1/2}-I\|_{
 \mathrm{op}}\le C e_{d,n,\gamma,s}.\quad}
 \label{f16v35:relative-covariance}
\end{equation}
In particular, with $\eta=C e_{d,n,\gamma,s}$,
\begin{equation}
 (1-\eta)\mathcal C\preceq\Sigma_\gamma\preceq(1+\eta)\mathcal C.
 \label{f16v35:Loewner}
\end{equation}
The retained marginal obeys the same comparison with
$(K^{\rm int})^{-1}$. This is a \emph{covariance} conclusion; it is not a
pointwise lower bound on the full native Hessian outside $\mathcal B_s$.
\end{theorem}
\begin{proof}
Fix $u$ with $u^{\mathsf T}\mathcal C u=1$, and set $X=u^{\mathsf T}w$.
The fourth moments of $X$ under $P_s,Q_s$ are bounded by a numerical constant
by \eqref{f16v35:directional-moments}. Factor the density difference into
square roots and use Cauchy--Schwarz:
\[
 |\mathbb E_{P_s}X^2-\mathbb E_{Q_s}X^2|
 \le h(P_s,Q_s)
       \{2\mathbb E_{P_s}X^4+2\mathbb E_{Q_s}X^4\}^{1/2}
 \le C\sqrt{dn^4/\gamma}.
\]
Moreover $|u|^2\le H_*$, since $\mathcal C\succeq H_*^{-1}I$.
Equation~\eqref{f16v35:tails-bound} pays the missing $X^2$ integrals by
$C\gamma^{-s}$; the factors $P(\mathcal B_s)$ and $Q(\mathcal B_s)$ in the
interior expectations differ from one by the same bound. All four means are
zero by simultaneous global conjugation (for the Gaussian, also by reflection).
Thus $|\operatorname{Var}_{\mathbb P_\gamma}X-1|\le C e_{d,n,\gamma,s}$.
The estimate is uniform over every such $u$, proving both matrix statements.
Taking principal covariance blocks gives the retained assertion. No factor
counting the number of tested directions is introduced.
\end{proof}

\begin{theorem}[The same untruncated native difference and Wilson-word banks]
\label{f16v35:fields}
Retain V33's assembled map $\mathcal B$ and V34's actual nonlinear bank
$\mathcal X_{\gamma,\ell}=\sqrt{\gamma\omega_\ell}\,
\boldsymbol{\operatorname{Im}}W_\ell$, with every original word weight.
If $e_{d,n,\gamma,s}\le1$, then
\begin{equation}
 \begin{split}
 \operatorname{Cov}_{\boldsymbol\nu_\gamma}(\mathcal B y)
                     &\preceq(1+C e_{d,n,\gamma,s})I,\\
 \big\|\operatorname{Cov}_{\mathbb P_\gamma}(\mathcal X_\gamma)
       -\mathcal D\mathcal C\mathcal D^{\mathsf T}\big\|_{\rm op}
                     &\le C e_{d,n,\gamma,s}.
 \end{split}
 \label{f16v35:field-bounds}
\end{equation}
Here $\mathcal D$ is the original first-differential word bank from V34;
$\mathcal D\mathcal C\mathcal D^{\mathsf T}\preceq I$. Both estimates include
all cross covariances and refer to the complete native laws, without a retained
window. No automatic gauge invariance of the quaternion-vector words is claimed.
\end{theorem}
\begin{proof}
Use \eqref{f16v35:Loewner} and the unchanged inequalities
$\mathcal B^{\mathsf T}\mathcal B\preceq K^{\rm int}$ and
$\mathcal D^{\mathsf T}\mathcal D\preceq\mathcal H$.
The first gives the first assertion, including the bound
\[
 \operatorname{Var}_{\boldsymbol\nu_\gamma}
 (\textstyle\sum_t a_t^{\mathsf T}\Delta y_t+\sum_i b_i^{\mathsf T}Dy_i)
 \le(1+C e_{d,n,\gamma,s})
                (30\sum_t|a_t|^2+3\sum_i|b_i|^2).
\]
The second compares the \emph{linear} word-bank covariance in operator norm.
For the nonlinear bank, retain the local-carrier version of V34's Taylor bound:
\[
 \|\mathcal X_\gamma-\mathcal D w\|^2
                    \le\frac C\gamma\sum_g|w_g|^4.
\]
It holds on the full principal charts by ordered unitary differentiation and
bounded incidence. On $\mathcal B_s$, the expected sum is at most $Cdn^2$.
Off that ball, its expectation is at most
$\mathbb E_{\mathbb P_\gamma}|w|^4 1_{\mathcal B_s^c}\le\gamma^{-s}$.
Thus the vector $L^2$ error is at most $C\sqrt{dn^2/\gamma}+C\gamma^{-s/2}/
\sqrt\gamma$, which is bounded by $C e_{d,n,\gamma,s}$: the second term is already
dominated by $\sqrt{dn^4/\gamma}$ for $s\ge1$, $d\ge1$.
Centering contracts $L^2$. For each unit output vector the linear bank has
variance at most $1+C e_{d,n,\gamma,s}$; Cauchy--Schwarz pays both mixed
terms and the squared error. This proves the second assertion. No observable
is replaced before the nonlinear error is paid.
\end{proof}

\subsection{Algebraically long native period bridges and the sampler distinction}

\begin{proposition}[Period limits at polynomial durations]
\label{f16v35:periods}
Fix $B$ and $0<\alpha<1/7$, and let
$n_\gamma=2\lfloor\gamma^\alpha/2\rfloor$. Under the complete native law,
linearly interpolate V33's nine period amplitudes $a_i/\sqrt{n_\gamma}$ at
$i/n_\gamma$. Their laws converge weakly in $C([0,1];\mathbb R^9)$ to the
centered Gaussian bridge with covariance $2(\min(s,t)-st)I_9$.
For each fixed $r>0$ and specified midpoint bridge,
\begin{equation}
 \boldsymbol\nu_{\gamma,B,n_\gamma}^{00}(|y_{n_\gamma/2,e}|\le r)\to0,
 \qquad
 \operatorname{Var}_{\boldsymbol\nu_\gamma}(y_{n_\gamma/2,e}^a)
               \ge(1-o(1))\frac{n_\gamma}{8B}.
 \label{f16v35:midpoint}
\end{equation}
In the general entry regime, whenever $\eta=C e_{d,n,\gamma,s}<1$, the native
rate-one full retained-bridge sampler satisfies
\begin{equation}
 \lambda_{\gamma,B,n}^{\rm br}
      \le\frac{1+\eta}{1-\eta}\,30\sin^2\frac\pi{2n}.
 \label{f16v35:sampler-test}
\end{equation}
These are neither a physical-time calibration nor spectral convergence inferred
from total variation.
\end{proposition}
\begin{proof}
For this sequence \eqref{f16v35:polynomial-regime} holds, and
$e_{d,n_\gamma,\gamma,s}=O(\gamma^{-(1-5\alpha)/2})+\gamma^{-s}$.
The complete Gaussian period factor and its normalization are exactly V33's;
they are not recomputed or replaced here. Its polygonal rescaling is distributed
as uniform-mesh sampling of a continuous variance-two Gaussian bridge.
On a common realization these interpolations converge uniformly. Total
variation contracts under the measurable projection and interpolation at each
$\gamma$, so Theorem~\ref{f16v35:TV} transfers the weak process limit.
The limiting path law is not asserted to be a total-variation limit of
finite-dimensional polygonal path laws.

The Gaussian midpoint ball bound from V34 is
$[\min(1,4r\sqrt{B/(\pi n)})]^3$. Add the new total-variation error; both
terms vanish. The reducing period covariance contributes $n/(8B)$ to every
scalar midpoint bridge. The retained part of \eqref{f16v35:Loewner} therefore
gives the variance inequality. The all-time fixed window is a subset of this
midpoint event and also has probability tending to zero.

For \eqref{f16v35:sampler-test}, reuse V34's first temporal period test $f$ and
its exact Gaussian exterior predictors $g_u$. The native conditional-variance
minimum is at most $\mathbb E_{\boldsymbol\nu_\gamma}|f-g_u|^2$.
Each difference is linear, so the \emph{relative} retained covariance bound
gives a summed energy at most $(1+\eta)15$, with no accumulated coefficient
error. The native test variance is at least
$(1-\eta)[2\sin^2(\pi/(2n))]^{-1}$. Division proves the displayed bound.
No native update kernel is identified with its Gaussian counterpart.
\end{proof}

\subsection{Checkpoint: which costs have changed and which have not}

V34's full-density error $\gamma^{-1/2}\exp\{C_p d\log(n+1)\}$ remains a
valid predecessor estimate. The present result replaces it, for total variation
and the stated Gaussian-normalized covariance tests, by a polynomial error
under a new explicit entry condition. For example $n\asymp\gamma^{1/8}$ at
fixed $B$ is now allowed, with error $O(\gamma^{-3/16})$ in these norms.
This is not a claim that all \emph{unscaled} polynomial-weighted $L^1$ errors
vanish on every such sequence. Higher raw moments carry their actual Gaussian
scales; only the tested normalized moments and field banks were estimated here.

There is also a normalization boundary. The square-root argument compares two
already normalized interior densities. Its gradient does not see an additive
constant in the action. Consequently normalized-law closeness does not by itself
supply an $o(1)$ absolute pressure correction or a physical Perron-normalizer
comparison in the new growing regime. V34's fixed-regulator partition conclusion
is unchanged. The lower mass bound in \eqref{f16v35:tails-bound} is used only to
pay complete tails before normalization; it is not promoted to a sharp pressure
asymptotic.

The new rate still pays the temporal soft eigenvalue in its score inverse and
in its global comparison radius. A sharper derivative-relative score estimate,
or a treatment of nonlinear compact periods that does not expand their entire
history at once, is the next plausible improvement. The exponent $1/7$ is a
sufficient threshold of this proof, not a physical threshold or an optimality
claim. Neither an additional fixed-window boundary gap nor another temporal
power removes these costs.

All conclusions concern the original unprojected identity-endpoint history.
The prior fixed-bridge width theorems and arbitrary-fast-pinning problem retain
their scopes. Root Gauss reduction, physical endpoint/vacuum normalization,
independent physical time, unrestricted regulators, and interacting continuum
reconstruction remain separate. The supplied import memoranda distinguish
power gain from locality/normalization cost and keep the final physical-transfer
arrow conditional; this checkpoint retains both distinctions.

\subsection{Verification and preservation}

The diagnostic is
\begin{center}\small\path{verify_ym_f16_v35_polynomial_comparison.py}.\end{center}
It tests normalized square-root identities, a nonzero Neumann boundary term,
local-incidence derivative sums, Gaussian marginal scales and noncommuting
matrix comparisons, exact tail arithmetic, the relative covariance handoff,
and the retained Gaussian-predictor upper test. These are finite algebraic
checks, not native compact-integral enclosures or a formal proof of the analytic
lemmas. The new diagnostic passes 254 exact assertions. All 31 supplied
V4--V34 checkers were rerun unchanged and passed; their actual outputs and
hashes are reported separately from inherited reports.

Only the literal title version and this terminal insertion change V34.
Removing them recovers the supplied predecessor byte for byte. The original
sources and companion inventories remain unchanged. The analytic dependencies
are explicitly the compact-pin, finite-word, and Gaussian-matrix inputs above;
no independent certification of the full archive is claimed. The V35 companion
checkpoint records review depth and the next checkpoint V40.

\begin{firewall}
V35 gives a polynomial-regime total-variation comparison of the complete
endpoint-fixed native law, a relative native covariance estimate in its actual
Gaussian metric, and complete nonlinear field-bank bounds. It extends the
native period process and sampler upper test to specified algebraic durations.
The comparison ball is removed with quantified tails; its interior convexity
is not asserted globally. No unrestricted-duration estimate, sharp growing-regime
pressure asymptotic, physical Gauss or vacuum identification, calibrated
physical transfer gap, or interacting continuum mass gap is proved.
\end{firewall}
% END F16 V35 APPEND

% BEGIN F16 V36 APPEND -- 2026-09-18
% Insert immediately before the final document-ending command of V35.
% No predecessor proof, probability, scalar normalization, or clock is edited.

\clearpage
\section{V36: absolute normalization, cubic cancellation, and even-observable accuracy}
\label{f16v36:context}

V35 deliberately compares normalized probabilities. Its square-root score cannot
see an additive action constant, so its improved growing-regime estimate does
not determine the absolute partition ratio. We address that separate scalar
problem before trying to extrapolate another sampling gap. The action is already
anchored at the identity, and its leading cubic correction is odd under global
Euclidean reflection. Its Gaussian mean vanishes, but its variance does not.
Keeping that variance and the actual Haar correction gives the first absolute
normalization coefficient, with a remainder controlled in V35's polynomial
regime. The same parity improves the relative covariance estimate for linear
coordinates without claiming a better total-variation rate for arbitrary tests.

All statements concern the original identity-endpoint, unprojected joint history.
The comparison ball will again be removed. No physical Perron normalization,
root Gauss reduction, unrestricted regulator limit, or mass gap is inferred.

\subsection{Dependencies, ownership, and the fixed scalar convention}

The V35 appendix, analytic audit, and checkpoint were read completely. We use
its explicit comparison ball, complete tail estimates, full joint Gaussian
metric, and convex-domain variance inequality. The next companion checkpoint
remains V40. The new estimates below do not use V30's screening or any assumed
native global spectral gap. Two empty semantic searches were followed by
mounted-text inventories and an identified f14 V82 read; they are not evidence
of archive-wide novelty.

Thermodynamic interpolation already occurs in f14 V82, for a different
chart law weighted by a determinant. The cubic--quartic log-integral
coefficient is also a standard cumulant calculation. For primary-source context,
Greeff, \href{https://arxiv.org/abs/2009.14321}{\emph{Free Energy Perturbation Theory
at Low Temperature}}, Sections I--II, records the cumulant expansion and its
sensitivity to constant energy shifts. Katsevich,
\href{https://arxiv.org/abs/2306.07262v3}{\emph{The Laplace approximation accuracy
in high dimensions: a refined analysis and new skew adjustment}},
Theorem 2.17 and Corollary 2.25, explains the cancellation of the odd leading
correction for even observables. Those statements and surrounding assumptions
were inspected in primary HTML. We prove the needed finite interpolation
identities and estimates directly; neither paper supplies the native constants,
complete-chart tails, or regulator regime used here. No convergent cumulant
series is assumed.

Keep exactly V34--V35's notation
\[
 w=(\zeta,y),\quad d=(72n-21)B,\quad
 U_0=\tfrac12w^{\mathsf T}\mathcal H w,\quad
 \mathcal C=\mathcal H^{-1},\quad \epsilon=\gamma^{-1/2},
\]
with $B=m^3$, $m\ge4$, $n\ge2$, and both retained endpoints the identity.
Let $\mathbb P_\gamma$ and $\mathbb Q$ be their unchanged complete joint laws,
and $\mathcal Z_\gamma,\mathcal Z_0$ their unchanged chart masses. Thus
$U_\gamma(0)=U_0(0)=0$ and $\mathcal J_\gamma(0)=1$ are exact conventions, not
post hoc adjustments. Let $\lambda_g\in\{1/2,1\}$ be each original three-group's
Haar power; $1/2$ occurs on the dressed chord endpoints only.

Fix $s\ge1$, apply V35's radius construction with tail exponent $t=s+4$, and
write $\mathbb B=\{|w|\le R_t\}$. Assume its exact entry conditions
\eqref{f16v35:entry} at this exponent, and, for convenience,
\begin{equation}
 \rho:=\frac{dn^4}{\gamma}\le1.
 \label{f16v36:rho}
\end{equation}
No different full law is defined by this ball. On it put
\[
 W=U_\gamma-\log\mathcal J_\gamma-U_0,\qquad
 d\mu_\theta=Z_{\theta,\mathbb B}^{-1}
       1_{\mathbb B}e^{-U_0-\theta W}\,dw,\quad 0\le\theta\le1.
\]
Here $\mu_0=\mathbb Q(\cdot\mid\mathbb B)$ and
$\mu_1=\mathbb P_\gamma(\cdot\mid\mathbb B)$, including their own normalizers.
Intermediate $\theta$ is a comparison parameter, not physical coupling or time.
All constants denoted $C$ below depend only on fixed word geometry; tail-entry
thresholds may depend on the fixed exponent $s$.

\subsection{The odd cubic and the measured quartic coefficient}

For an original word let $f_\ell(w)=s(W_\ell(w))$, with its actual ordered
exponentials and action weight $\omega_\ell$. Define the homogeneous polynomials
\[
 P_3(w)=\frac1{3!}\sum_\ell\omega_\ell
                    D^3f_\ell(0)[w,w,w],\qquad
 U_4(w)=\frac1{4!}\sum_\ell\omega_\ell D^4f_\ell(0)[w,w,w,w],
\]
and the measured coefficient
\begin{equation}
 P_4(w)=U_4(w)+\frac13\sum_g\lambda_g|w_g|^2.
 \label{f16v36:P4}
\end{equation}
Exactly quadratic endpoint amplitudes contribute to $U_0$, not to these jets.
These polynomials are coefficients for the \emph{joint} endpoint-fixed integral;
they are not V10's common-source matrices $G_j$.

\begin{proposition}[Local jets with their Haar and parity terms]
\label{f16v36:jets}
On $\mathbb B$, with a fixed geometry constant,
\begin{align}
 |W-\epsilon P_3-\epsilon^2P_4|
   &\le C\epsilon^3\sum_g|w_g|^5+C\epsilon^4\sum_g|w_g|^4,
       \label{f16v36:jet-value}\\
 |\nabla(W-\epsilon P_3)|^2
   &\le C\epsilon^4\sum_g(|w_g|^6+|w_g|^2),
       \label{f16v36:jet-gradient}\\
 |\nabla P_3|^2&\le C\sum_g|w_g|^4.\notag
\end{align}
The polynomial $P_3$ is odd and $P_4$ is even. In particular
$\mathbb E_{\mu_0}P_3=\mathbb E_{\mathbb Q}P_3=0$.
\end{proposition}
\begin{proof}
The ordered-unitary derivative argument of V6/V34 applies at orders four and
five: a differentiated exponential is an integral of unitary products with the
specified inserted matrices. At every fixed order its norm is bounded by the
product of those fixed matrix norms. Taylor's integral remainder on each word
therefore has the asserted local degree, with its original power of $\epsilon$.
There is no principal logarithm of an intermediate product. Summing the finite
carriers and their bounded row and column incidences gives the displayed
bounds, rather than a factor counting all words. The same argument applies to
one derivative of the fourth-order remainder.

For $|w_g|/\sqrt\gamma\le1/2$,
\[
 -\lambda_g\log j(w_g/\sqrt\gamma)
 =\frac{\lambda_g\epsilon^2}{3}|w_g|^2
                         +O(\epsilon^4|w_g|^4).
\]
The derivative of this term is $O(\epsilon^2|w_g|)$. This proves
\eqref{f16v36:P4} and its contribution to both estimates. The endpoint half
powers are not replaced by full Haar. Homogeneity gives the two parities, and
both Gaussian measures in the last claim are reflection invariant. The full
native measure need not be reflection invariant: oddness is used only where
that symmetry is actually available.
\end{proof}

\subsection{Uniform interpolation moments and centered cumulants}

Every $\mu_\theta$ has Hessian at least $\mathcal H/2$, since this holds at
both interpolation endpoints. The same is true after any real linear tilt.
Global simultaneous conjugation preserves $\mathbb B$, $U_0$, and $W$, and
hence makes the \emph{untilted} coordinate means zero. V35's variance argument
therefore gives, for every fixed $k$,
\begin{equation}
 \mathbb E_{\mu_\theta}|u^{\mathsf T}w|^k
       \le C_k(u^{\mathsf T}\mathcal C u)^{k/2},\qquad
 \mathbb E_{\mu_\theta}|w_g|^k\le C_kn^{k/2}.
 \label{f16v36:directional}
\end{equation}
The first estimate is used only for integer $k\ge1$. It follows by integrating
the sub-Gaussian tail obtained from the linear-tilt moment-generating function.
It is uniform in $\theta$, not an identification of $\mu_\theta$ with a Gaussian.

\begin{proposition}[Local-score cumulant bounds]
\label{f16v36:cumulants}
Uniformly for $0\le\theta\le1$,
\begin{align}
 \operatorname{Var}_{\mu_\theta}W&\le C\rho,&
 \mathbb E_{\mu_\theta}|W-\mu_\theta W|^4&\le C\rho^2,&
 |\kappa_{3,\mu_\theta}(W)|&\le C\rho^{3/2},
 \label{f16v36:W-moments}\\
 \operatorname{Var}_{\mu_\theta}(W-\epsilon P_3)
      &\le C\epsilon^4dn^5,&
 \operatorname{Var}_{\mu_0}P_3&\le Cdn^4,&
 |\mu_0P_4|&\le Cdn^2.
 \label{f16v36:jet-moments}
\end{align}
Here $\kappa_3(W)=\mathbb E[(W-\mathbb EW)^3]$. No independence of local
scores or summands is used.
\end{proposition}
\begin{proof}
Write $G_A=\nabla A^{\mathsf T}\mathcal C\nabla A$. The variance inequality
under $\mu_\theta$ reads $\operatorname{Var}A\le2\mathbb EG_A$.
By V35's local score estimate, \eqref{f16v36:directional}, and
$\|\mathcal C\|\le Cn^2$,
\[
 \mathbb EG_W\le C\epsilon^2dn^4=C\rho,\qquad
 \mathbb EG_W^2\le C\epsilon^4d^2n^8=C\rho^2.
\]
For the second inequality, use Cauchy--Schwarz on pairs of local fourth powers
and their eighth moments, retaining $d^2$ explicitly. It is not a decorrelation
argument.

For any smooth $A$, put $A_c=A-\mathbb EA$, $M_4=\mathbb EA_c^4$ and
$G=G_A$. Poincar\'e applied also to $A_c^2$ gives
\[
 M_4\le4(\mathbb EG)^2+8\mathbb E[A_c^2G]
       \le4\mathbb EG^2+8\sqrt{M_4\mathbb EG^2}.
\]
Solving the scalar quadratic yields $M_4\le81\mathbb EG^2$.
H\"older gives $|\kappa_3(A)|\le M_4^{3/4}$, proving
\eqref{f16v36:W-moments}. All functions are smooth on the fixed compact
comparison ball, so these applications are within the variance form domain.

The gradient estimates in Proposition~\ref{f16v36:jets}, the local sixth and
fourth moments, and $\|\mathcal C\|\le Cn^2$ give the first two assertions of
\eqref{f16v36:jet-moments}. The local quartic and quadratic moments give its
last assertion. The identical polynomial estimates hold under the full
Gaussian $\mathbb Q$ when no $\theta$ occurs.
\end{proof}

\subsection{An absolute partition ratio with the full chart mass retained}

\begin{theorem}[Polynomial-regime absolute normalization]
\label{f16v36:pressure}
Under the entry conditions above,
\begin{equation}
 \boxed{\quad
  \left|\log\frac{\mathcal Z_\gamma}{\mathcal Z_0}\right|
       \le C(\rho+\gamma^{-s}).
 \quad}
 \label{f16v36:pressure-bound}
\end{equation}
In particular the actual complete chart masses have ratio tending to one in
V35's polynomial regime. This is a new scalar estimate, not an inference from
total variation of normalized laws.
\end{theorem}
\begin{proof}
Let $\psi(\theta)=\log(Z_{\theta,\mathbb B}/Z_{0,\mathbb B})$.
Differentiation of these finite smooth integrals keeps every normalizer and gives
\begin{equation}
 \psi'=-\mu_\theta W,\qquad
 \psi''=\operatorname{Var}_{\mu_\theta}W,\qquad
 \psi'''=-\kappa_{3,\mu_\theta}(W).
 \label{f16v36:thermodynamic}
\end{equation}
The exact second-order integral identity is
\[
 \psi(1)=-\mu_0W+\int_0^1(1-\theta)
                              \operatorname{Var}_{\mu_\theta}W\,d\theta.
\]
Oddness eliminates $\mu_0P_3$. Equations \eqref{f16v36:jet-value} and
\eqref{f16v36:directional} give
\[
 |\mu_0W|\le C\{\epsilon^2dn^2+\epsilon^3dn^{5/2}
                                      +\epsilon^4dn^2\}\le C\rho.
\]
The variance term is at most $C\rho$. Notice that smallness of the mean is
proved from the anchored action and parity; it does not follow from its gradient
alone.

For the complete masses the exact relation is
\begin{equation}
 \log\frac{\mathcal Z_\gamma}{\mathcal Z_0}
  =\psi(1)+\log\mathbb Q(\mathbb B)-\log\mathbb P_\gamma(\mathbb B).
 \label{f16v36:full-mass}
\end{equation}
V35's actual complete tails at exponent $t=s+4$ bound the last two terms by
$C\gamma^{-t}$. Thus neither the native chart complement nor either full
normalizer has been omitted. This proves the theorem.
\end{proof}

Define two coefficients by expectations in the \emph{untruncated} original
joint Gaussian, not the ball-conditioned measure:
\begin{equation}
 \mathfrak p_{B,n}=\tfrac12\mathbb E_{\mathbb Q}P_3^2
                                  -\mathbb E_{\mathbb Q}P_4,
 \qquad
 \mathfrak e_{B,n}=\tfrac12\mathbb E_{\mathbb Q}P_3^2\ge0.
 \label{f16v36:coefficients}
\end{equation}
These coefficients can depend on $B,n$. They contain all original word
incidences, kinetic terms, and endpoint Haar powers.

\begin{theorem}[The first absolute coefficient with a quantitative remainder]
\label{f16v36:coefficient}
The coefficients in \eqref{f16v36:coefficients} satisfy
$|\mathfrak p_{B,n}|+\mathfrak e_{B,n}\le Cdn^4$, and
\begin{equation}
 \boxed{\quad
 \left|\log\frac{\mathcal Z_\gamma}{\mathcal Z_0}
                  -\frac{\mathfrak p_{B,n}}\gamma\right|
       \le C(\rho^{3/2}+\gamma^{-s}).
 \quad}
 \label{f16v36:pressure-expansion}
\end{equation}
At fixed $B,n$, this gives the coefficient of $\gamma^{-1}$ with an
$O_{B,n}(\gamma^{-3/2})$ remainder by choosing $s\ge3/2$.
The bound is not asserted sharp in fixed dimension.
\end{theorem}
\begin{proof}
The coefficient bound follows from Gaussian Poincar\'e for $P_3$ and the
local fourth moments for $P_4$. The same calculation on $\mathbb B$ was proved
in Proposition~\ref{f16v36:cumulants}.
Taylor's formula for \eqref{f16v36:thermodynamic} gives
\[
 \psi(1)=-\mu_0W+\tfrac12\operatorname{Var}_{\mu_0}W
       -\tfrac12\int_0^1(1-\theta)^2\kappa_{3,\mu_\theta}(W)\,d\theta.
\]
The last integral is $O(\rho^{3/2})$. The jet mean satisfies
\[
 |\mu_0W-\epsilon^2\mu_0P_4|
       \le C(\epsilon^3dn^{5/2}+\epsilon^4dn^2).
\]
For $R=W-\epsilon P_3$, Cauchy--Schwarz and
\eqref{f16v36:jet-moments} give
\[
 |\operatorname{Var}_{\mu_0}W-\epsilon^2\mu_0P_3^2|
 \le C\epsilon^3dn^{9/2}+C\epsilon^4dn^5.
\]
Every displayed error is bounded by $C\epsilon^3d^{3/2}n^6=C\rho^{3/2}$
for $d\ge1$, $n\ge2$, and $\epsilon\le1$.

It remains to replace the Gaussian ball coefficients by the full ones.
Bounded word incidence gives $|P_3|^2\le C|w|^6$ and
$|P_4|\le C(|w|^4+|w|^2)$. The actual Gaussian eighth-degree tail from V35
therefore pays their omitted integrals by $C\gamma^{-t}$. The conditioning
normalizer adds at most $Cdn^4\gamma^{-t}$. After multiplication by
$\epsilon^2$, this is $C(\epsilon^2+\rho)\gamma^{-t}\le C\gamma^{-t}$.
Equation~\eqref{f16v36:full-mass} completes the proof. Only a finite Taylor
formula on a compact ball has been used; an unbounded cubic exponential has
not been declared integrable, and no infinite series is summed.
\end{proof}

\paragraph{Explicit Gaussian contractions.}
Use the symmetric coefficient tensors defined by
\[
 P_3=T_{ijk}w_iw_jw_k/6,\qquad U_4=A_{ijkl}w_iw_jw_kw_l/24,
\]
with repeated scalar indices summed. Set
$v_k=T_{ijk}\mathcal C_{ij}$. Gaussian pairings give
\begin{equation}
 \begin{split}
 \mathfrak p_{B,n}={}&
 \tfrac1{12}T_{ijk}T_{abc}\mathcal C_{ia}\mathcal C_{jb}\mathcal C_{kc}
 +\tfrac18v^{\mathsf T}\mathcal C v
 -\tfrac18 A_{ijkl}\mathcal C_{ij}\mathcal C_{kl}\\
 &-\tfrac13\sum_g\lambda_g\Tr\mathcal C_{gg}.
 \end{split}
 \label{f16v36:wick}
\end{equation}
There are six fully cross pairings and nine one-cross pairings in the cubic
square, and three pairings in the quartic mean. In this native identity-endpoint
setting $v=0$: $\mathbb E_{\mathbb Q}\nabla P_3=v/2$ is a vector in each colour
group fixed by every simultaneous adjoint rotation, and hence is zero. The
formula is displayed before this simplification to retain its general pairing
convention. No numerical enclosure or evaluation of the entire native tensor
contraction is claimed.

As an exact normalization check, one isolated Wilson factor with full Haar has
$P_3=0$, $U_4=-|w|^4/24$, and Gaussian covariance $I_3$. Its coefficient is
$15/24-3/3=-3/8$. Replacing that \emph{reference factor} by a half-Haar factor
instead gives $15/24-3/6=1/8$. These one-group checks are not the full-history
coefficient; they show why its actual endpoint powers cannot be discarded.

\subsection{Relative entropy of the complete law, with its support direction}

\begin{theorem}[Forward native relative entropy and its first coefficient]
\label{f16v36:entropy}
In the same regime,
\begin{equation}
 \boxed{\quad
 0\le D(\mathbb P_\gamma\Vert\mathbb Q)\le C(\rho+\gamma^{-s}),\qquad
 \left|D(\mathbb P_\gamma\Vert\mathbb Q)
                    -\frac{\mathfrak e_{B,n}}\gamma\right|
       \le C(\rho^{3/2}+\gamma^{-s}).
 \quad}
 \label{f16v36:entropy-bound}
\end{equation}
For every finite coupling the reverse entropy
$D(\mathbb Q\Vert\mathbb P_\gamma)$ is infinite, because the untruncated
Gaussian has positive mass outside the native principal product.
\end{theorem}
\begin{proof}
On the common comparison ball, exact normalization gives
\[
 D(\mu_1\Vert\mu_0)=\psi'(1)-\psi(1)
       =\int_0^1\theta\operatorname{Var}_{\mu_\theta}W\,d\theta.
\]
This proves the $O(\rho)$ bound there. Taylor's formula for $\psi'$ and $\psi$
and Proposition~\ref{f16v36:cumulants} give
\[
 D(\mu_1\Vert\mu_0)=\tfrac12\operatorname{Var}_{\mu_0}W
                                      +O(\rho^{3/2}).
\]
The coefficient substitutions in the preceding proof give
$\mathfrak e_{B,n}/\gamma$ with the same error. The quartic mean and Haar mean
cancel from this leading entropy coefficient, not from the pressure coefficient.

The full entropy needs an additional tail check: $-\log\mathcal J_\gamma$
is unbounded near a principal cut. The complete ordered-word second-derivative
bound gives $U_\gamma+U_0\le C|w|^2$ on the native chart. Its tail expectation
is paid by V35. For the logarithmic Haar term, use the pointwise identity and
bound
\[
 0\le-\mathcal J_\gamma\log\mathcal J_\gamma\le e^{-1},
 \qquad(0\log0:=0).
\]
The same radial tail integral and lower bound for $\mathcal Z_\gamma$ used in
V35 then give
\begin{equation}
 \mathbb E_{\mathbb P_\gamma}[|W|1_{\mathbb B^c}]
                                      \le C\gamma^{-t}.
 \label{f16v36:entropy-tail}
\end{equation}
No Gaussian expectation of the singular logarithm is substituted here.

Let $p=\mathbb P_\gamma(\mathbb B)$, $q=\mathbb Q(\mathbb B)$. On the ball,
$\log(d\mathbb P_\gamma/d\mathbb Q)
=\log(d\mu_1/d\mu_0)+\log(p/q)$. Off it, on the native support, the same
logarithm is $-W-\log(\mathcal Z_\gamma/\mathcal Z_0)$.
Integrating this exact split, using \eqref{f16v36:entropy-tail},
$1-p,1-q\le\gamma^{-t}$ and Theorem~\ref{f16v36:pressure}, changes the ball
entropy by at most $C\gamma^{-s}$. This proves both claims for the full law.
The reverse assertion is immediate from the support mismatch and is not
contradicted by either total variation convergence or the finite forward entropy.
\end{proof}

The chain rule, with the same actual marginals and conditionals, now yields
\begin{equation}
 \begin{split}
 D(\mathbb P_\gamma\Vert\mathbb Q)
 ={}&D(\boldsymbol\nu_\gamma\Vert Q_{B,n}^{00})\\
 &+\mathbb E_{\boldsymbol\nu_\gamma}
 D(\mathbb P_\gamma(d\zeta\mid y)\Vert\mathbb Q(d\zeta\mid y)).
 \end{split}
 \label{f16v36:entropy-chain}
\end{equation}
Both nonnegative terms are therefore at most $C(\rho+\gamma^{-s})$. This is
an average under the \emph{actual retained marginal}, not a supremum over all
retained boundary histories and not a new arbitrary-fast-pinning theorem.

\subsection{Even observables have a second-order covariance error}

\begin{theorem}[A sharper relative covariance bound in the same Gaussian metric]
\label{f16v36:covariance}
For $\Sigma_\gamma=\operatorname{Cov}_{\mathbb P_\gamma}(w)$,
\begin{equation}
 \boxed{\quad
 \|\mathcal H^{1/2}\Sigma_\gamma\mathcal H^{1/2}-I\|_{\rm op}
                  \le C(\rho+\gamma^{-s}).
 \quad}
 \label{f16v36:relative-covariance}
\end{equation}
This improves V35's $O(\sqrt\rho+\gamma^{-s})$ bound for this particular even
observable class; it does not claim an improved total-variation rate for
arbitrary events. The retained covariance satisfies the corresponding
relative bound with $(K^{\rm int})^{-1}$.
\end{theorem}
\begin{proof}
Fix $u$ with $u^{\mathsf T}\mathcal C u=1$, and set $F=(u^{\mathsf T}w)^2$.
The directional moment estimates give $\operatorname{Var}_{\mu_\theta}F\le C$
uniformly in $u,\theta$. For $f(\theta)=\mu_\theta F$ the exact derivatives are
\[
 f'(\theta)=-\operatorname{Cov}_{\mu_\theta}(F,W),\qquad
 f''(\theta)=\mathbb E_{\mu_\theta}
                 [(F-\mu_\theta F)(W-\mu_\theta W)^2].
\]
The second equality retains both derivative-of-normalizer terms through
centering. Reflection of $\mu_0$ makes $\operatorname{Cov}_{\mu_0}(F,P_3)=0$.
Consequently
\[
 |f'(0)|\le C\sqrt{\operatorname{Var}_{\mu_0}(W-\epsilon P_3)}
       \le C\epsilon^2\sqrt{dn^5}\le C\rho.
\]
Cauchy--Schwarz and \eqref{f16v36:W-moments} give
$|f''(\theta)|\le C\rho$. Taylor's integral formula therefore bounds
$|\mu_1F-\mu_0F|$ by $C\rho$.

Since $|u|^2\le H_*$, V35's complete weighted tails pay all omitted second
moments and the two conditioning probabilities uniformly in $u$ by
$C\gamma^{-t}$. The complete and conditioned coordinate means vanish by global
conjugation, so second moments are exactly the needed variances. The result is
uniform in every Gaussian-normalized direction; taking the supremum proves
\eqref{f16v36:relative-covariance}. Principal covariance blocks give the retained
statement. No spectral convergence of a sampler is inferred from it.
\end{proof}

\begin{proposition}[Derivative banks, marked Gaussian predictors, and growth rates]
\label{f16v36:handoff}
Put $\eta=C(\rho+\gamma^{-s})$. The unchanged V33 assembled difference--curvature
map and the unchanged linear word bank obey
\begin{equation}
 \operatorname{Cov}_{\boldsymbol\nu_\gamma}(\mathcal B y)\preceq(1+\eta)I,
 \qquad
 \|\operatorname{Cov}_{\mathbb P_\gamma}(\mathcal D w)
                     -\mathcal D\mathcal C\mathcal D^{\mathsf T}\|
                                                    \le\eta.
 \label{f16v36:linear-fields}
\end{equation}
The retained-sampler upper test of V35 remains valid with this smaller $\eta$
when $\eta<1$. No second-order rate for the \emph{nonlinear} word bank is
asserted by these linear statements.

A sufficient simultaneous regime remains
\begin{equation}
 \frac{Bn^7\log(n+1)+n^6\log\gamma}{\gamma}\longrightarrow0.
 \label{f16v36:regime}
\end{equation}
It has not been enlarged. Within it, the complete chart partition ratio tends
to one, the forward entropy tends to zero, and the relative covariance error
is $O(Bn^5/\gamma+\gamma^{-s})$. At fixed $B$ and $n\asymp\gamma^{1/8}$,
these bounds are $O(\gamma^{-3/8})$ for $s\ge1$, and the remainder after the
explicit pressure coefficient is $O(\gamma^{-9/16})$.
\end{proposition}
\begin{proof}
The matrix inequalities $\mathcal B^{\mathsf T}\mathcal B\preceq K^{\rm int}$
and $\mathcal D^{\mathsf T}\mathcal D\preceq\mathcal H$ were proved in V33--V34.
Apply the new relative covariance theorem to their original input spaces.
V35's upper sampler test uses only such a relative covariance comparison and
Gaussian conditional predictors as admissible native predictors, so the same
proof applies. The nonlinear word itself has a separate even second-order jet;
its products have not been estimated here at the improved rate.

Use V35's explicit polynomial entry with $t=s+4$ instead of $s$.
Condition \eqref{f16v36:regime} supplies that entry and $\rho\to0$.
Since $d=(72n-21)B$, the displayed powers follow by substitution. The radius and
its entry cost are unchanged except for this fixed tail exponent. In particular
this is not a new duration-uniform mass estimate or a removal of the soft
$n^2$ inverse from the proof.
\end{proof}

\subsection{The scalar handoff and remaining physical normalization}

Let $\mathcal A_{\gamma,B,n}^{00}$ denote the exact identity-endpoint compressed
kernel written in V4's bridge half-density chart. Its original factorization is
\[
 \mathcal A_{\gamma,B,n}^{00}
    =c_{\gamma,B,n}F_{\gamma,2}(I)^{-1}\mathcal Z_\gamma.
\]
Writing the analogous Gaussian coefficient as
$c_{0,B,n}F_2(0)^{-1}\mathcal Z_0$ gives the exact scalar separation
\begin{equation}
 \begin{split}
 \log\frac{\mathcal A_{\gamma,B,n}^{00}}{\mathcal A_{0,B,n}^{00}}
 ={}&\log\frac{\mathcal Z_\gamma}{\mathcal Z_0}
   +\log\frac{c_{\gamma,B,n}}{c_{0,B,n}}\\
 &+\log\frac{F_2(0)}{F_{\gamma,2}(I)}.
 \end{split}
 \label{f16v36:physical-scalars}
\end{equation}
The first term is now quantitatively estimated in the polynomial regime; the
last two are not silently set to zero or replaced by a chosen vacuum constant.
A diagonal endpoint kernel is also not a Perron eigenvalue. The normalized
probability, its chart mass, and the physical transfer at its own spectral top
remain distinct objects.

V35's warning about additive constants remains mathematically correct: replacing
$U_\gamma$ by $U_\gamma+a_\gamma$ preserves every normalized law and covariance
but changes $\log\mathcal Z_\gamma$ by $-a_\gamma$. V36 controls the scalar
because the original action convention is fixed and its mean correction is
computed, not because that warning has been bypassed by terminology.

The next structural task is still a source-relative treatment of the temporal
inverse or the compact nonlinear periods, together with the untouched physical
scalar factors where a transfer comparison is intended. The present second-order
cancellations do not enlarge the valid space--time regime, give a stationary
root-Gauss-reduced vacuum, calibrate physical time, or construct an interacting
continuum theory. They close a specific missing normalization estimate and
improve its accompanying even-observable comparison.

\subsection{Verification and preservation}

The diagnostic is
\begin{center}\small\path{verify_ym_f16_v36_pressure_parity.py}.\end{center}
It checks the signs and weights in exact normalized interpolation, Gaussian
cubic--quartic contractions, a genuine ordered quaternion cubic, full versus
half-Haar coefficients, forward entropy and its direction, and the even-observable
cancellation. Finite algebraic fixtures are not enclosures of the native
high-dimensional Wilson integral or formal verification of the analytic proof.
Actual test outputs, predecessor reruns, hashes, and rendering records are
reported in the package, not inferred from their historical descriptions.

Only one literal title version and this terminal insertion change V35. Removing
them recovers the predecessor byte for byte. The deductions inherit V35's
interior curvature and complete tails, with their compact-pin and finite-word
premises. The entire analytic archive has not been independently recertified.
The next companion checkpoint remains V40.

\begin{firewall}
V36 estimates the absolute endpoint-fixed chart partition ratio and its first
Gaussian coefficient in the existing polynomial regime. It proves finite
forward relative entropy, its leading coefficient, and a second-order relative
covariance estimate for the linear-coordinate bank. The reverse full entropy
is infinite because of the support mismatch. All complete probabilities,
endpoint Haar powers, Gaussian soft directions, and original scalar factors
are retained. No stronger total-variation rate, unrestricted space--time limit,
physical Perron normalization, calibrated transfer gap, or interacting
Yang--Mills continuum mass gap follows.
\end{firewall}
% END F16 V36 APPEND

% BEGIN F16 V37 APPEND -- 2026-09-18
% Insert immediately before the final document-ending command of V36.
% Every predecessor assertion, endpoint factor, measure and clock is preserved.

\clearpage
\section{V37: complete scalar normalization of the dressed endpoint kernel}
\label{f16v37:context}

V36 estimates the full joint chart mass but leaves two scalar factors in
\eqref{f16v36:physical-scalars}. Before proving a new spectral implication, we
return to their definitions. The finite-chart normalization of the chosen cube
carrier cancels \emph{exactly} against its magnetic endpoint normalization.
What remains is one static native integral and the explicitly normalized
one-group Wilson factor. Estimating these gives the complete identity-endpoint
compressed-kernel ratio, not merely the ratio of normalized path probabilities.
The carrier and the transfer themselves are unchanged.

This is a scalar closure of the displayed identity, not a spectral closure of
Yang--Mills. The result is at fixed identity endpoints, before the remaining
root Gauss restriction, and in the polynomial regulator regime inherited from
V35--V36. A diagonal kernel of a power does not determine its Perron eigenvalue
or its first excited spectral ratio. That limitation is retained quantitatively
below instead of renaming the computed scalar a vacuum energy.

\subsection{Exact endpoint cancellation before any asymptotic estimate}

The complete V36 appendix and analytic audit were read. Its joint pressure
coefficient and error estimate are used as inherited inputs. We also checked
the literal definitions \eqref{f16v1:gram}, \eqref{f16v4:amplitude}, and
\eqref{f16v4:normalizers}. The general cancellation is elementary; its
quantitative application to the two outstanding factors is the purpose here.
The static integral is not one cube in isolation: all mixed plaquettes of its
spatial assembly remain coupled. Targeted searches of the active source and the
supplied f14/f15/Grand Tensor archives are recorded in the audit, not presented
as an exhaustive novelty or proof certificate. The next checkpoint remains V40.

The one-variable Bessel formulas used for the kinetic factor are classical.
The NIST Digital Library of Mathematical Functions, equations
\href{https://dlmf.nist.gov/10.32.E2}{10.32.2} and
\href{https://dlmf.nist.gov/10.40.E1}{10.40.1}, with coefficients from
\href{https://dlmf.nist.gov/10.17.E1}{10.17.1}, were checked in primary HTML.
We give a direct real-integral remainder bound as well; no lattice theorem is
imported from those special-function identities. The supplied memoranda's
separation of scalar normalization, locality cost and physical spectral
comparison remains in force.

Let $m\ge4$, $B=m^3$, and write $d_{\rm e}=15B$. In scalar coordinates
$C_B$ and the spatial matrices carry their original colour tensor $I_3$.
For a complete compact bridge configuration $Z$, define
\begin{equation}
 \begin{split}
  E_{\gamma,B}(x;Z)&=\tfrac12x^{\mathsf T}C_B^{-1}x
       +\gamma S_\times(X(x),\Theta_{X(x)}^{-1}Z),\\
  \mathcal S_{\gamma,B}(Z)&=
       \int_{\mathcal P_x}e^{-E_{\gamma,B}(x;Z)}\,dx .
 \end{split}
 \label{f16v37:static-mass}
\end{equation}
The domain $\mathcal P_x$ is the original product of principal chord charts.
There is \emph{no Haar density} in this displayed integral. This is not an
approximation: the reason is the exact square-root Jacobian already built into
$\Phi_\gamma$. Haar factors in the full joint history are not removed.

\begin{proposition}[The carrier chart mass cancels from the full dressed weight]
\label{f16v37:cancellation}
At every finite $B,n,\gamma$ and every admitted compact retained history,
\begin{equation}
 F_{\gamma,2}(Z)
     =\frac{c_\Omega^2}{p_{\gamma,B}}\mathcal S_{\gamma,B}(Z),
 \qquad
 \mathcal W^{\rm full}_{\gamma,n}(\mathbf y)
 =\left(\frac{c_{H,\gamma}}{z_\gamma}\right)^{24Bn}
 \frac{J_{{\rm br},\gamma}(\mathbf y)Z_{\gamma,\rm full}(\mathbf y)}
 {\sqrt{\mathcal S_{\gamma,B}(Z_0)\mathcal S_{\gamma,B}(Z_n)}}.
 \label{f16v37:cancelled-weight}
\end{equation}
No chart probability or exponentially small carrier mass has been estimated in
this identity. Every mixed plaquette, endpoint Gaussian, and original kinetic
normalizer is retained.
\end{proposition}
\begin{proof}
The exact change to chord charts gives
$|\Phi_\gamma(X)|^2dH^{5B}(X)=p_{\gamma,B}^{-1}\Omega(x)^2dx$.
Also $M_\times^2=e^{-\gamma S_\times}$, with the original moving bridge input.
Substitute these identities into \eqref{f16v1:gram}; the first equality follows.
Each of the two square-root endpoint factors contains
$c_\Omega/\sqrt{p_{\gamma,B}}$. Their product cancels
$p_{\gamma,B}^{-1}c_\Omega^2$ in $c_{\gamma,B,n}$. Substitution into V4's exact
weight proves the second equality. The cancellation is of identical scalar
factors, not a replacement of $F_{\gamma,2}$ by its limiting determinant.
\end{proof}

At the identity define
\[
 E_{\gamma,B}(x)=E_{\gamma,B}(x;I),\qquad
 E_{0,B}(x)=\tfrac12 x^{\mathsf T}H_{\rm e}x,\qquad
 H_{\rm e}=C_B^{-1}+A^{\mathsf T}A,\qquad C_{\rm e}=H_{\rm e}^{-1}.
\]
Thus $H_{\rm e}^{-1}$ is the inherited static posterior covariance; it is not
the joint endpoint-history covariance $\mathcal C$ of V35. Put
\[
 \mathcal S_{0,B}=\int_{\mathbb R^{d_{\rm e}}}e^{-E_{0,B}}dx,
 \qquad Q_{\rm e}=N(0,C_{\rm e}).
\]
The original Gaussian normalization is
$F_2(0)=c_\Omega^2\mathcal S_{0,B}$. In particular its determinant is exactly
that in V1, with all colours. With both retained endpoints fixed at $I$, let
$\mathcal A_{\gamma,B,n}^{00}$ be V36's unchanged compressed diagonal kernel.
Integrating the interior bridges in \eqref{f16v37:cancelled-weight} gives
\begin{equation}
 \mathcal A_{\gamma,B,n}^{00}
    =\left(\frac{c_{H,\gamma}}{z_\gamma}\right)^{24Bn}
                   \frac{\mathcal Z_\gamma}{\mathcal S_{\gamma,B}(I)},
 \qquad
 \mathcal A_{0,B,n}^{00}
    =(2\pi)^{-36Bn}\frac{\mathcal Z_0}{\mathcal S_{0,B}}.
 \label{f16v37:kernel-factors}
\end{equation}
The joint masses $\mathcal Z_\gamma,\mathcal Z_0$ are exactly V34--V36's masses,
not the fixed-bridge fast integrals $Z_{\gamma,\rm full}(\mathbf y)$.

\subsection{The static endpoint integral has its own uniform Gaussian scale}

Fix a requested tail power $s\ge1$ and put $t=s+4$. There are fixed constants
$0<a_{\rm e}\le1/4$, $H_{\rm e,+}\ge1$, and $C_{\rm e,2}$ such that
\begin{equation}
 E_{\gamma,B},E_{0,B}\ge a_{\rm e}|x|^2,\qquad
 2a_{\rm e}I\preceq H_{\rm e}\preceq H_{\rm e,+}I,
 \qquad
 \|D^2(E_{\gamma,B}-E_{0,B})\|
             \le C_{\rm e,2}|x|/\sqrt\gamma.
 \label{f16v37:static-data}
\end{equation}
The native inequality is on the whole principal product. The first term of
$E_{\gamma,B}$ supplies the pin independently of the retained input; the mixed
action is nonnegative. The fixed cube covariance and bounded mixed-word
incidence give the two Gaussian bounds. The last bound follows from the same
ordered-unitary third derivatives used in V6. These constants do not involve
$n$ or $B$.

For explicit entry bookkeeping set
\begin{equation}
 \begin{split}
 L_{{\rm e},t}={}&1+\tfrac{d_{\rm e}}2\log(2H_{\rm e,+}/a_{\rm e})
  +\log[1+((d_{\rm e}+6)/a_{\rm e})^4]+(t+1)\log\gamma,\\
 R_{{\rm e},t}^2={}&2L_{{\rm e},t}/a_{\rm e},\qquad
 \mathbb B_{\rm e}=\{|x|\le R_{{\rm e},t}\}.
 \end{split}
 \label{f16v37:static-radius}
\end{equation}
Assume $\gamma\ge2$ and
\begin{equation}
 R_{{\rm e},t}\le\sqrt\gamma/2,\qquad
 C_{\rm e,2}R_{{\rm e},t}/\sqrt\gamma\le a_{\rm e},\qquad
 \sigma:=d_{\rm e}/\gamma\le1.
 \label{f16v37:static-entry}
\end{equation}
A sufficiently small fixed constant in
$d_{\rm e}+(t+1)\log\gamma\le c_{\rm e}\gamma$ is sufficient, after a fixed
threshold. This condition is separately stated, not inferred from a
fixed-volume limit.

\begin{proposition}[Complete static tails and interpolation moments]
\label{f16v37:static-preparation}
Let $P_{\rm e,\gamma}$ be the normalized density proportional to
$1_{\mathcal P_x}e^{-E_{\gamma,B}}$. Under \eqref{f16v37:static-entry},
\begin{equation}
 \mathcal S_{\gamma,B}(I)\ge2^{-d_{\rm e}/2-1}\mathcal S_{0,B},\qquad
 \mathbb E_{P_{\rm e,\gamma}}[(1+|x|^8)1_{\mathbb B_{\rm e}^c}]
 +\mathbb E_{Q_{\rm e}}[(1+|x|^8)1_{\mathbb B_{\rm e}^c}]
 \le2\gamma^{-t}.
 \label{f16v37:static-tails}
\end{equation}
On the ball, every normalized interpolation
$\mu_\theta^{\rm e}\propto\exp[-E_{0,B}-\theta(E_{\gamma,B}-E_{0,B})]$,
$0\le\theta\le1$, has curvature at least $H_{\rm e}/2$ and local moments
$\mathbb E|x_g|^k\le C_k$ at each original chord group, for every fixed $k$.
\end{proposition}
\begin{proof}
The Hessian error is at most $a_{\rm e}I\preceq H_{\rm e}/2$. Both exponents
and gradients vanish at zero. Hence $E_{\gamma,B}\le2E_{0,B}$ on the ball.
The Gaussian of precision $2H_{\rm e}$ has mean square norm at most
$d_{\rm e}/(4a_{\rm e})$, and at least half its mass lies in that ball.
Integrating proves the lower bound on the actual static mass.

For either weighted tail use the complete pin
$e^{-a_{\rm e}|x|^2}$, splitting it into two equal exponents outside the ball.
The radial eighth moment is at most
$(2\pi/a_{\rm e})^{d_{\rm e}/2}[1+((d_{\rm e}+6)/a_{\rm e})^4]$.
Also $\mathcal S_{0,B}\ge(2\pi/H_{\rm e,+})^{d_{\rm e}/2}$.
Division by the proved native mass gives the tail bound
$2e^{-1}\gamma^{-(t+1)}\le\gamma^{-t}$ for each law. The native density is
zero off its principal product; that complement is included. No product Haar
factor is inserted into this static integral.

Each linear tilt of an interpolation has the same Hessian floor on the same
convex ball. Apply the constant-metric variance inequality proved in V35 to a
linear observable under that tilt, and integrate the second log-moment
derivative. Untilting gives zero means by simultaneous colour conjugation.
Thus $\mathbb E\exp(u^{\mathsf T}x)\le\exp(u^{\mathsf T}C_{\rm e}u)$.
The fixed bound on $C_{\rm e}$ gives all stated local moments. This argument
requires neither spatial independence nor a global convexity claim outside
the comparison ball.
\end{proof}

For the mixed action only, define its classical cubic and quartic jets
$P_3^{\rm e},P_4^{\rm e}$ by
\[
 E_{\gamma,B}-E_{0,B}
     =\gamma^{-1/2}P_3^{\rm e}+\gamma^{-1}P_4^{\rm e}
          +\text{Taylor remainder}.
\]
They are the third and fourth homogeneous differentials of the literal mixed
Wilson words at $Z=I$, with their original weights. The confining Gaussian is
exactly quadratic. In contrast to V36's joint $P_4$, there is no additional
$\frac13\sum_g\lambda_g|x_g|^2$ here: the static Haar factors already cancelled.
Put
\begin{equation}
 \mathfrak s_B
   =\tfrac12\mathbb E_{Q_{\rm e}}(P_3^{\rm e})^2
                      -\mathbb E_{Q_{\rm e}}P_4^{\rm e}.
 \label{f16v37:static-coefficient}
\end{equation}

\begin{theorem}[Static endpoint pressure with all mixed couplings retained]
\label{f16v37:static-pressure}
Under \eqref{f16v37:static-entry}, $|\mathfrak s_B|\le C d_{\rm e}$ and
\begin{equation}
 \begin{split}
 \left|\log\frac{\mathcal S_{\gamma,B}(I)}{\mathcal S_{0,B}}\right|
     &\le C(\sigma+\gamma^{-s}),\\
 \left|\log\frac{\mathcal S_{\gamma,B}(I)}{\mathcal S_{0,B}}
                     -\frac{\mathfrak s_B}{\gamma}\right|
     &\le C(\sigma^{3/2}+\gamma^{-s}).
 \end{split}
 \label{f16v37:static-pressure-bound}
\end{equation}
In the original flat chord half-density representation, the normalized static
endpoint vector also obeys
\begin{equation}
 \left\|\frac{1_{\mathcal P_x}e^{-E_{\gamma,B}/2}}
                   {\sqrt{\mathcal S_{\gamma,B}(I)}}
             -\frac{e^{-E_{0,B}/2}}{\sqrt{\mathcal S_{0,B}}}\right\|_2^2
           \le C(\sigma+\gamma^{-s}).
 \label{f16v37:endpoint-vector}
\end{equation}
This is the original dressed endpoint vector at $Z=I$, not an isolated-cube
vacuum. The coefficient can depend on $B$, although its displayed bound is
linear in $B$.
\end{theorem}
\begin{proof}
Write $\epsilon=\gamma^{-1/2}$ and $W_{\rm e}=E_{\gamma,B}-E_{0,B}$.
Fixed word carriers and local incidence give, on the ball,
\[
 |\nabla W_{\rm e}|^2\le C\epsilon^2\sum_g|x_g|^4,\quad
 |\nabla(W_{\rm e}-\epsilon P_3^{\rm e})|^2
                       \le C\epsilon^4\sum_g|x_g|^6,
\]
\[
 |W_{\rm e}-\epsilon P_3^{\rm e}-\epsilon^2P_4^{\rm e}|
                       \le C\epsilon^3\sum_g|x_g|^5.
\]
The original unitary derivative formulas justify these estimates without a
logarithm of a product. Under every $\mu_\theta^{\rm e}$, the preceding local
moments and bounded $C_{\rm e}$ imply
\[
 \operatorname{Var}W_{\rm e}\le C\epsilon^2d_{\rm e},\quad
 \mathbb E|W_{\rm e}-\mathbb EW_{\rm e}|^4
                                      \le C\epsilon^4d_{\rm e}^2,\quad
 \operatorname{Var}(W_{\rm e}-\epsilon P_3^{\rm e})
                                      \le C\epsilon^4d_{\rm e}.
\]
For clarity, the centered fourth-moment estimate is obtained as in V36 by applying
Poincar\'e to the centered observable and its square; it is not a
factorization of local correlations. It bounds the third centered moment
by $C\sigma^{3/2}$. Gaussian Poincar\'e also gives
$\mathbb E_{Q_{\rm e}}(P_3^{\rm e})^2\le Cd_{\rm e}$, since the cubic is odd.
The quartic mean is bounded by $Cd_{\rm e}$.

Let $\psi_{\rm e}(\theta)$ be the log partition ratio of these ball integrals.
The exact identities $\psi_{\rm e}'=-\mathbb EW_{\rm e}$,
$\psi_{\rm e}''=\operatorname{Var}W_{\rm e}$ and
$\psi_{\rm e}'''=-\mathbb E(W_{\rm e}-\mathbb EW_{\rm e})^3$ are V36's
finite thermodynamic identities, now under the static laws. Oddness at
$\theta=0$ removes the cubic mean, not its variance. Taylor's formula at zero
therefore gives
\[
 \psi_{\rm e}(1)=\epsilon^2\left\{
     \tfrac12\mathbb E_{\mu_0^{\rm e}}(P_3^{\rm e})^2
                         -\mathbb E_{\mu_0^{\rm e}}P_4^{\rm e}\right\}
                         +O(\sigma^{3/2}).
\]
In detail, the fifth-order mean and the cubic/higher-jet covariance each cost
at most $C\epsilon^3d_{\rm e}$; the squared higher jet costs
$C\epsilon^4d_{\rm e}$. They are bounded by $C\sigma^{3/2}$ when
$d_{\rm e}\ge1$, $\sigma\le1$. No infinite cumulant expansion is invoked.

The eighth-degree Gaussian tails replace the two ball-conditioned coefficients
by their full $Q_{\rm e}$ expectations. Their normalizer corrections, multiplied
by $\epsilon^2$, are at most $C(\epsilon^2+\sigma)\gamma^{-t}$.
Finally the exact mass identity adds
$\log Q_{\rm e}(\mathbb B_{\rm e})-
 \log P_{\rm e,\gamma}(\mathbb B_{\rm e})$.
Proposition~\ref{f16v37:static-preparation} pays these terms and proves
\eqref{f16v37:static-pressure-bound}.

For \eqref{f16v37:endpoint-vector}, apply V35's normalized square-root
Poincar\'e calculation on the ball. Its score expectation is at most
$C\epsilon^2d_{\rm e}$. The square-root distance from a probability to its
ball conditioning is squared distance
$2(1-\sqrt{P(\mathbb B_{\rm e})})\le2P(\mathbb B_{\rm e}^c)$.
The triangle inequality and the two paid tails prove the full-space bound.
The exact flat representative follows by multiplying $a_{\gamma,I}(X)$ by
$\sqrt{J_{X,\gamma}}$ and using Proposition~\ref{f16v37:cancellation}.
\end{proof}

\subsection{The Wilson kinetic normalization is an explicit scalar}

Define
\begin{equation}
 \mathfrak b_\gamma=\frac{z_\gamma}{c_{H,\gamma}(2\pi)^{3/2}}.
 \label{f16v37:kinetic-ratio}
\end{equation}
The symbol $\mathfrak b_\gamma$ is a scalar, distinct from the internal magnetic
multiplier $b_\gamma(X)$ of V2.

\begin{proposition}[Exact one-group factor and a uniform logarithmic remainder]
\label{f16v37:kinetic}
For $\gamma>0$,
\begin{equation}
 z_\gamma=\frac{2e^{-\gamma}I_1(\gamma)}{\gamma},\qquad
 \mathfrak b_\gamma=\sqrt{2\pi\gamma}\,e^{-\gamma}I_1(\gamma).
 \label{f16v37:Bessel}
\end{equation}
There are numerical $C,\gamma_*<\infty$ such that, for $\gamma\ge\gamma_*$,
\begin{equation}
 \left|\log \mathfrak b_\gamma+\frac3{8\gamma}+\frac3{16\gamma^2}\right|
                                      \le C\gamma^{-3}.
 \label{f16v37:kinetic-expansion}
\end{equation}
In particular $-24Bn\log \mathfrak b_\gamma=9Bn/\gamma+O(Bn/\gamma^2)$,
with a constant independent of $B,n$.
\end{proposition}
\begin{proof}
The normalized class angle density is $(2/\pi)\sin^2\theta\,d\theta$.
The integral representation for $I_1$ gives the first equality; inserting the
original $c_{H,\gamma}=(2\pi^2\gamma^{3/2})^{-1}$ gives the second.
One may also set $u=\gamma(1-\cos\theta)$ directly, obtaining
\begin{equation}
 \mathfrak b_\gamma=\frac2{\sqrt\pi}\int_0^{2\gamma}
       e^{-u}u^{1/2}\left(1-\frac{u}{2\gamma}\right)^{1/2}du.
 \label{f16v37:kinetic-real}
\end{equation}
On $0\le u\le\gamma$, Taylor's formula gives the square root
$1-u/(4\gamma)-u^2/(32\gamma^2)$ with remainder at most
$C u^3/\gamma^3$. The remaining interval and the omitted tails of the three
polynomial integrals are exponentially small. The gamma moments with shape
$3/2$ are $\mathbb Eu=3/2$ and $\mathbb Eu^2=15/4$. Thus
\[
 \mathfrak b_\gamma=1-\frac3{8\gamma}-\frac{15}{128\gamma^2}+O(\gamma^{-3}).
\]
It is bounded away from zero eventually. Taylor's formula for the logarithm
gives the second coefficient $-15/128-9/128=-3/16$ and a uniform remainder.
This proves the required estimate without assuming convergence of a Bessel
asymptotic series. Raising the normalized factor to its actual $24Bn$ power
multiplies its logarithm; it does not produce independent uncontrolled errors.
\end{proof}

\subsection{The complete identity-endpoint kernel ratio}

Retain V36's joint coefficient $\mathfrak p_{B,n}$ and
$\rho=dn^4/\gamma$, $d=(72n-21)B$. Define
\begin{equation}
 \mathfrak a_{B,n}=\mathfrak p_{B,n}-\mathfrak s_B+9Bn.
 \label{f16v37:complete-coefficient}
\end{equation}
The joint $\mathfrak p_{B,n}$ includes its original half-Haar endpoint terms.
The static $\mathfrak s_B$ has no Haar term for the exact reason proved above.
The final summand is the one-group kinetic normalization, not an adjustable
vacuum subtraction.

\begin{theorem}[All scalar factors in the dressed identity-endpoint comparison]
\label{f16v37:main}
Assume V36's entry conditions at the requested $s\ge1$,
\eqref{f16v37:static-entry}, $\gamma\ge\gamma_*$, and $\rho\le1$.
Then the original kernels satisfy the exact identity
\begin{equation}
 \log\frac{\mathcal A_{\gamma,B,n}^{00}}{\mathcal A_{0,B,n}^{00}}
 =\log\frac{\mathcal Z_\gamma}{\mathcal Z_0}
  -\log\frac{\mathcal S_{\gamma,B}(I)}{\mathcal S_{0,B}}
                                      -24Bn\log \mathfrak b_\gamma.
 \label{f16v37:exact-scalar}
\end{equation}
Moreover $|\mathfrak a_{B,n}|\le Cdn^4$ and
\begin{equation}
 \boxed{\quad
 \left|\log\frac{\mathcal A_{\gamma,B,n}^{00}}{\mathcal A_{0,B,n}^{00}}
                    -\frac{\mathfrak a_{B,n}}\gamma\right|
     \le C\left\{\rho^{3/2}+\sigma^{3/2}+\frac{Bn}{\gamma^2}
                        +\gamma^{-s}\right\}
     \le C(\rho^{3/2}+\gamma^{-s}).\quad}
 \label{f16v37:kernel-expansion}
\end{equation}
In particular $|\log(\mathcal A_{\gamma,B,n}^{00}/\mathcal A_{0,B,n}^{00})|
\le C(\rho+\gamma^{-s})$. The entire ratio tends to one in V35's stated
polynomial regime. No scalar factor in V36's decomposition is left unestimated
for this identity-endpoint ratio.
\end{theorem}
\begin{proof}
Divide the two exact formulas \eqref{f16v37:kernel-factors}; this proves
\eqref{f16v37:exact-scalar} before any approximation. V36 supplies the joint
mass expansion with error $C(\rho^{3/2}+\gamma^{-s})$.
Theorem~\ref{f16v37:static-pressure} supplies the static term; the kinetic
proposition supplies $9Bn/\gamma$ with error $CBn/\gamma^2$. Their signs
in \eqref{f16v37:complete-coefficient} follow from the exact quotient.

Since $d_{\rm e}=15B\le d$ and $n\ge2$, $\sigma\le\rho$. Also
$Bn/\gamma^2\le C\rho^{3/2}$ for $\gamma\ge2$, $B\ge1$, $n\ge2$.
The coefficient bound follows from V36 and $|\mathfrak s_B|\le15CB$.
As $\rho\le1$, the last claim follows. V35's regime implies
$B/\gamma\to0$ and $\log\gamma/\gamma\to0$, so it eventually satisfies the
separate static entry condition as well. The constants and finite thresholds
in these estimates are not numerically optimized or enclosed.
\end{proof}

It is also legitimate to leave the exactly known kinetic factor unexpanded:
\[
 \left|\log\frac{\mathcal A_{\gamma,B,n}^{00}}{\mathcal A_{0,B,n}^{00}}
       +24Bn\log \mathfrak b_\gamma
       -\frac{\mathfrak p_{B,n}-\mathfrak s_B}{\gamma}\right|
                 \le C(\rho^{3/2}+\sigma^{3/2}+\gamma^{-s}).
\]
This is an exact resummation of one explicit scalar only. It is not a new
resummation of the coupled field theory. Neither $\mathfrak p_{B,n}$ nor
$\mathfrak s_B$ is claimed numerically evaluated by its Gaussian definition.

\begin{proposition}[Length ratios cancel the static endpoint integral exactly]
\label{f16v37:length-ratio}
At one common $B,\gamma$ and any two integer durations $n,k\ge2$, set
$\mathcal R_{n,k}=\mathcal A_{\gamma,B,n}^{00}\mathcal A_{0,B,k}^{00}/
(\mathcal A_{0,B,n}^{00}\mathcal A_{\gamma,B,k}^{00})$.
Then
\begin{equation}
 \log\mathcal R_{n,k}
 =\log\frac{\mathcal Z_{\gamma,B,n}}{\mathcal Z_{0,B,n}}
  -\log\frac{\mathcal Z_{\gamma,B,k}}{\mathcal Z_{0,B,k}}
                                     -24B(n-k)\log \mathfrak b_\gamma.
 \label{f16v37:length-exact}
\end{equation}
If both durations satisfy V36's entry conditions, this has coefficient
$[\mathfrak p_{B,n}-\mathfrak p_{B,k}+9B(n-k)]/\gamma$ and error at most
\[
 C\{\rho_n^{3/2}+\rho_k^{3/2}+B|n-k|/\gamma^2+\gamma^{-s}\},
 \qquad \rho_j=(72j-21)Bj^4/\gamma.
\]
No static approximation or division by a possibly tiny approximate endpoint
factor is needed for this ratio. The errors for the two durations are added;
they are not asserted to cancel or to gain a factor $|n-k|$.
\end{proposition}
\begin{proof}
The same static factor occurs in both exact diagonal kernels, since the
carrier and endpoints are unchanged. Subtract \eqref{f16v37:exact-scalar} at
the two durations and apply V36 and \eqref{f16v37:kinetic-expansion}.
\end{proof}

\subsection{What this closes, and the next spectral obligation}

At fixed $B$ and $n\asymp\gamma^{1/8}$, the full log-ratio bound is
$O(\gamma^{-3/8})$ and the error after its explicit first coefficient is
$O(\gamma^{-9/16})$, with $s\ge1$. These are the same joint scales as V36,
now for the entire original dressed endpoint quantity rather than its chart
mass alone. The result does not enlarge the admissible growth regime.

There is no remaining carrier-chart-mass problem in this scalar: its
cancellation is exact. The next normalization question is instead a spectral
one. A diagonal compressed kernel and finite-duration ratios are not the
spectral top of the full transfer. An elementary example makes the distinction
precise. For $0<a<b<1$, take
$T_0=\operatorname{diag}(a,0)$, $T_1=\operatorname{diag}(a,b)$ and
$f_\epsilon=(\sqrt{1-\epsilon},\sqrt\epsilon)$. Their moments differ by
$\epsilon b^j$ at every positive integer $j$, yet their top eigenvalues are
$a$ and $b$. At any finite list of durations the relative moment differences
can be made arbitrarily small by choosing $\epsilon>0$ sufficiently small.
This counterfixture is not a native model; it excludes an algebraic inference
from a controlled finite family of scalar moments to a spectral-top estimate.

Thus a useful next transfer result must control the relevant physical
spectral measure, its vacuum overlap and its complementary channels, or provide
an appropriate same-top operator comparison. The existing source-relative
temporal-inverse problem is also unaffected. The memoranda already separate
an auxiliary functional inequality from this physical-transfer implication;
closing the literal scalar identity does not remove that separation.
Root Gauss averaging, endpoint-varying matrix elements, independent physical
time calibration, unrestricted regulators and an interacting continuum
construction remain unpaid. The covariance and process theorems of V35--V36
are not replaced by a claim about the physical spectrum.

\subsection{Verification and predecessor preservation}

The accompanying exact diagnostic is
\begin{center}\small\path{verify_ym_f16_v37_scalar_closure.py}.\end{center}
It checks the carrier/Jacobian cancellation, non-factorized static Gaussian
normalizers, normalized interpolation and Wick coefficients without a spurious
static Haar term, the single-group real integral coefficients, signs and powers
in the full scalar quotient, length-ratio cancellation, and the finite-moment
spectral counterfixture. The actual report records assertion counts and scope;
these finite tests are not native high-dimensional integral enclosures or a
formal proof of the analytic estimates.

Only the literal version in the title and this terminal insertion change V36.
Removing those edits recovers its bytes exactly. All available predecessor
checkers are copied without modification; their rerun outputs and hashes are
reported separately from inherited reports. The static proof uses fixed-cube
Gaussian coercivity and the original mixed-word derivatives, not a new spatial
mixing hypothesis. The final joint scalar theorem explicitly uses V36's pressure
estimate. The full analytic lineage has not been independently recertified.
The next scheduled companion checkpoint remains V40.

\begin{firewall}
V37 cancels the original carrier chart mass exactly, estimates the complete
static magnetic endpoint mass, and includes the Wilson kinetic normalization
with its correct coefficient. It thereby controls every scalar factor in the
identity-endpoint dressed diagonal-kernel comparison in the inherited polynomial
regime. It also gives a static endpoint-vector comparison and exact endpoint
cancellation in length ratios. It does not prove a spectral-top estimate,
root-Gauss-reduced vacuum convergence, an unrestricted duration or volume theorem,
a physical time identification, or a Yang--Mills continuum mass gap. No original
measure, carrier, endpoint factor or update clock is changed.
\end{firewall}
% END F16 V37 APPEND

% BEGIN F16 V38 APPEND -- 2026-09-18
% Insert immediately before the final document-ending command of V37.
% The predecessor body, endpoint factors, transfer and clocks are unchanged.

\clearpage
\section{V38: endpoint-weighted spectral edges and vacuum visibility}
\label{f16v38:context}

V37 closes the scalar normalization of the identity-endpoint kernel. It does
not identify that scalar with a spectral top. We now load the entire sequence
into the spectral measure of the \emph{actual} full transfer, using one exact
transfer to regularize the endpoint distribution. This produces a new native
conclusion: at fixed spatial width and the admitted polynomial durations, the
endpoint-weighted spectrum has a Gamma edge profile relative to the Gaussian
threshold. It also proves that this particular filtered endpoint vector has
vanishing overlap with the native Perron vector along that limit. Thus scalar
accuracy is not the missing vacuum-preparation estimate.

The conclusion concerns the unprojected root-coordinate transfer used by the
preceding history integrals. Its Perron top agrees with its root-invariant top,
but its endpoint spectral measure need not be supported in that physical
subspace. We give the exact projection split. Neither a native upper spectral
edge, a first-excited ratio, nor a physical time scale is obtained by deleting
that distinction.

\subsection{Dependencies and the endpoint vector that is actually in the Hilbert space}

The supplied V37 appendix and audit were read completely. Its scalar theorem
is the only native asymptotic input below. V2 supplies the full Gaussian
transfer, with its actual kinetic metric; the static null-space classification
is the inherited f15 V99 result used in V33. The Grand Tensor's transfer-moment
and source-cyclic reconstruction, and f15 V80/V89's overlap and all-duration
warnings, already contain the relevant general spectral principles. We use
those principles, not present them as new general theorems. The new objects
are this endpoint sequence's exact Hilbert-space realization, quantitative
Perron visibility, full Gaussian edge exponent, and the native edge limit.
Targeted archive searches and passage reads are recorded separately; they do
not certify the entire archive. The next companion checkpoint remains V40.

For classical input, the Mehler identity is
\href{https://dlmf.nist.gov/18.18.E28}{DLMF 18.18.28}; its normalization is checked
below. Blossier--Della Morte--von Hippel--Mendes--Sommer,
\href{https://arxiv.org/abs/0902.1265}{\emph{On the generalized eigenvalue method
for energies and matrix elements in lattice field theory}}, provides context
for the distinction between correlation data, overlap and spectral extraction.
The abstract and introductory spectral representation were inspected; no
physical overlap or uniform gap hypothesis from that paper is imported here.
The moment-convergence argument needed below is proved directly.

Fix $B=m^3$, $m\ge4$, and $\gamma>0$. Work on the unchanged pre-root-restriction
Hilbert space
\[
 \mathcal H_{\mathrm{rt},B}=L^2(\SU(2)^{5B}\times\SU(2)^{12B},dH),\qquad
 T_\gamma=T_{\mathrm{full},\gamma}.
\]
The inherited operator is a positive self-adjoint contraction with a strictly
positive bounded integral kernel at each finite regulator. It is compact.
Let $\Lambda_\gamma=\|T_\gamma\|$ and let $\Omega_\gamma>0$ be its unit Perron
vector. Positivity here includes operator positivity, not just positive kernel
entries. Let $\Pi_{\rm G}$ be the orthogonal average over the remaining root
gauge action. The transfer commutes with that action. Uniqueness and positivity
of $\Omega_\gamma$ imply
\begin{equation}
 \Pi_{\rm G}\Omega_\gamma=\Omega_\gamma,\qquad
 \|T_\gamma|_{\operatorname{Ran}\Pi_{\rm G}}\|=\Lambda_\gamma.
 \label{f16v38:same-top}
\end{equation}
This is the inherited finite-regulator Perron argument, not a new continuum
Gauss theorem.

Let $a_{\gamma,I}(X)$ be the original unit conditional endpoint vector of V1.
A delta at the retained identity is \emph{not} an $L^2$ vector. Define instead
\begin{equation}
 b_\gamma^{\rm end}(Q)=j_{\rm br,0}^{1/2}
    \int T_\gamma(Q;X,I)a_{\gamma,I}(X)\,dH^{5B}(X),\qquad
 j_{\rm br,0}=c_{H,\gamma}^{12B},\qquad
 f_\gamma=b_\gamma^{\rm end}/\|b_\gamma^{\rm end}\|.
 \label{f16v38:loaded-vector}
\end{equation}
Here $Q=(X',Z')$ is a full root-coordinate slice. The scalar $j_{\rm br,0}$
is the coarse bridge-chart Jacobian at zero. It accounts for the two endpoint
half-densities in the diagonal chart kernel, rather than being absorbed into
an unspecified source normalization.

\begin{proposition}[One common spectral measure for the complete endpoint sequence]
\label{f16v38:moments}
The vector $b_\gamma^{\rm end}$ is nonzero, positive and in $\mathcal H_{\mathrm{rt},B}$. For every
integer $j\ge0$,
\begin{equation}
 \|b_\gamma^{\rm end}\|^2=\mathcal A_{\gamma,B,2}^{00},\qquad
 m_{\gamma,j}:=\langle f_\gamma,T_\gamma^j f_\gamma\rangle
       =\frac{\mathcal A_{\gamma,B,j+2}^{00}}
                    {\mathcal A_{\gamma,B,2}^{00}}
       =\int_{[0,1]}\lambda^j\,d\mu_\gamma(\lambda),
 \label{f16v38:moment-identity}
\end{equation}
where $\mu_\gamma$ is a probability spectral measure. No intermediate dressed
projection is inserted into $T_\gamma^j$.
\end{proposition}
\begin{proof}
The kernel is bounded at fixed regulator and $\|a_{\gamma,I}\|_1\le1$ on the
probability Haar space. Thus the displayed integral is bounded and belongs to
$L^2$; strict kernel positivity makes it nonzero. Symmetry, composition and
Tonelli give
\[
 \langle b_\gamma^{\rm end},T_\gamma^j b_\gamma^{\rm end}\rangle
 =j_{\rm br,0}\int a_{\gamma,I}(X)
 T_\gamma^{j+2}(X,I;X',I)a_{\gamma,I}(X')\,dH(X)dH(X').
\]
This is exactly the coarse diagonal kernel in Haar coordinates times its
zero-bridge chart Jacobian. V4's half-density formula identifies it with
$\mathcal A_{\gamma,B,j+2}^{00}$. Taking $j=0$ and then normalizing proves
the first two identities. The spectral theorem for the same positive
contraction gives the last one. This regularization is performed before any
use of moments; the delta itself is never treated as a unit vector.
\end{proof}

For $n\ge3$ define its endpoint-weighted probability
\begin{equation}
 d\mu_{\gamma,n}(\lambda)
   =\frac{\lambda^{n-2}\,d\mu_\gamma(\lambda)}{m_{\gamma,n-2}}.
 \label{f16v38:tilt}
\end{equation}
It has no mass at zero, and its $j$th moment is
$\mathcal A_{\gamma,B,n+j}^{00}/\mathcal A_{\gamma,B,n}^{00}$.
Equivalently it is the spectral measure of the unit vector
\[
 v_{\gamma,n}=T_\gamma^{(n-2)/2}f_\gamma/
                      \|T_\gamma^{(n-2)/2}f_\gamma\|.
\]
Fractional powers here are spectral calculus of the original positive transfer,
not a newly chosen physical interpolation. Even $n$ uses only integer powers.

\subsection{A quantitative native top weight, with its actual size cost}

The full kernel factors as
$T_\gamma(Q,Q')=b_{\rm mag}(Q)L_\gamma(Q,Q')b_{\rm mag}(Q')$, where
$b_{\rm mag}=b_\gamma(X)m_\gamma(X,Z)>0$ denotes only the original magnetic
multiplier. This notation is distinct from the vector $b_\gamma^{\rm end}$ above.
The root-fixed integral $L_\gamma$ contains exactly $E=24B$ Wilson kinetic
factors. Therefore
\[
 \ell_\gamma\le L_\gamma(Q,Q')\le u_\gamma,\qquad
 \ell_\gamma=z_\gamma^{-E}e^{-2\gamma E},\quad
 u_\gamma=z_\gamma^{-E},\quad u_\gamma/\ell_\gamma=e^{48\gamma B}.
\]
These inequalities hold on the full compact configuration space; they require
neither a small-field event nor a comparison with a Gaussian.

\begin{theorem}[Perron visibility and finite-duration top brackets]
\label{f16v38:visibility}
At every finite $B,\gamma$,
\begin{equation}
 q_\gamma:=|\langle\Omega_\gamma,f_\gamma\rangle|^2
       \ge \operatorname{sech}^2(48\gamma B)
       \ge e^{-96\gamma B}.
 \label{f16v38:overlap}
\end{equation}
The same lower bound holds after normalizing $\Pi_{\rm G}f_\gamma$. For every
integer $j\ge1$,
\begin{equation}
 m_{\gamma,j}^{1/j}\le\Lambda_\gamma
 \le\min\{1,e^{96\gamma B/j}m_{\gamma,j}^{1/j}\}.
 \label{f16v38:top-bracket}
\end{equation}
Thus these moments determine the same physical top when all durations are
available at one fixed regulator. Their finite-duration upper certificate has
logarithmic width at most $96\gamma B/j$, not a size-free width.
\end{theorem}
\begin{proof}
The quotient of the positive vector $b_\gamma^{\rm end}$ by $b_{\rm mag}$ lies between
$\ell_\gamma A$ and $u_\gamma A$ for a positive scalar $A$. The Perron equation
places $\Omega_\gamma/b_{\rm mag}$ between
$\ell_\gamma A'$ and $u_\gamma A'$ for another positive scalar. Consequently
$h=b_\gamma^{\rm end}/\Omega_\gamma$ has essential range $[m,M]$ with
$M/m\le(u_\gamma/\ell_\gamma)^2$. No bound on the magnetic multiplier itself
has been multiplied into this ratio.

For a positive random variable $m\le h\le M$,
$(h-m)(M-h)\ge0$ gives
$\mathbb Eh^2\le(m+M)\mathbb Eh-mM$.
Minimizing $(\mathbb Eh)^2/\mathbb Eh^2$ over $\mathbb Eh\in[m,M]$ yields
$4mM/(m+M)^2$. Apply this under $\Omega_\gamma^2dH$ to obtain
\[
 \frac{|\langle\Omega_\gamma,b_\gamma^{\rm end}\rangle|^2}{\|b_\gamma^{\rm end}\|^2}
 \ge\frac{4R^2}{(1+R^2)^2}=\operatorname{sech}^2(\log R),\qquad
 R=e^{48\gamma B}.
\]
The final elementary bound in \eqref{f16v38:overlap} follows from
$\cosh t\le e^t$ for $t\ge0$. The gauge projection preserves the numerator
and cannot increase the denominator; it is nonzero by the positive numerator.
Finally the top atom and the upper spectral support give
$q_\gamma\Lambda_\gamma^j\le m_{\gamma,j}\le\Lambda_\gamma^j$.
Take roots and use the overlap bound. No excited spectral ratio is estimated.
\end{proof}

For later use the projection split is exact. Put
$\theta_{\gamma,n}=\|\Pi_{\rm G}v_{\gamma,n}\|^2$. Commutation gives
\begin{equation}
 \mu_{\gamma,n}=\theta_{\gamma,n}\mu^{\rm G}_{\gamma,n}
                  +(1-\theta_{\gamma,n})\mu^{\perp}_{\gamma,n}.
 \label{f16v38:Gauss-mixture}
\end{equation}
The components are the spectral measures of the two normalized projections;
a zero-weight component is omitted. The second has no atom at $\Lambda_\gamma$.
There is no estimate here making $\theta_{\gamma,n}$ close to one. In particular,
root projection cannot be inserted into V37's scalar asymptotic without new
matrix-element information.

\subsection{The full Gaussian threshold and its polynomial endpoint factor}

We next compute a \emph{fixed-$B$} Gaussian long-duration asymptotic. In this
subsection every spatial matrix includes its three-colour tensor. On
$q=(x,y)\in\mathbb R^{51B}$ put
\begin{equation}
 \mathsf A_{\rm kin}=\operatorname{diag}(\mathsf M_B,G),\qquad
 \mathsf S_B=\begin{pmatrix}
  \mathsf H_B+A^{\mathsf T}A&A^{\mathsf T}D\\
  D^{\mathsf T}A&D^{\mathsf T}D
 \end{pmatrix},\qquad
 \mathsf V_B=\mathsf A_{\rm kin}^{1/2}\mathsf S_B\mathsf A_{\rm kin}^{1/2}.
 \label{f16v38:Gaussian-transfer-data}
\end{equation}
These are the complete \emph{uncompressed} Gaussian transfer data of V2.
Let $\nu_1,\ldots,\nu_{51B}$ be the eigenvalues of $\mathsf V_B$ and set
\begin{equation}
 \Lambda_{{\rm G},B}
 =\prod_{\nu_j>0}\left(1+\frac{\nu_j}{2}
                    +\sqrt{\nu_j+\nu_j^2/4}\right)^{-1/2},\qquad
 h_B=3(B+2),\qquad a_B=h_B/2.
 \label{f16v38:Gaussian-top}
\end{equation}
The letter $a_B$ here is an edge exponent, not a carrier or V37's coefficient
$\mathfrak a_{B,n}$.

\begin{theorem}[Gaussian edge exponent for the actual dressed endpoint]
\label{f16v38:Gaussian-edge}
For every fixed $B$, the norm of the full Gaussian transfer $T_0$ is
$\Lambda_{{\rm G},B}\in(0,1)$. There are $c_B>0$ and real $d_{B,1},d_{B,2}$
such that, as integer $n\to\infty$,
\begin{equation}
 \mathcal A_{0,B,n}^{00}
 =c_B\Lambda_{{\rm G},B}^{\,n}n^{-a_B}
       \{1+d_{B,1}/n+d_{B,2}/n^2+O_B(n^{-3})\}.
 \label{f16v38:Gaussian-asymptotic}
\end{equation}
The null count $h_B$ includes the $3(B-1)$ gradient directions and nine period
directions before root Gauss reduction. In particular the exponent is not
merely $9/2$. No uniformity of the constants as $B\to\infty$ is asserted.
\end{theorem}
\begin{proof}
By the inherited positive internal stiffness,
$q^{\mathsf T}\mathsf S_Bq=x^{\mathsf T}\mathsf H_Bx+|Ax+Dy|^2$ vanishes
exactly when $x=0$ and $Dy=0$. The f15 V99 kernel classification therefore gives
nullity $h_B$. Congruence by the positive kinetic square root preserves that
nullity. Whitening the kinetic covariance and orthogonally diagonalizing
$\mathsf V_B$ are unitary coordinate changes in $L^2$ with their determinants
included. The transfer becomes a tensor product of one-dimensional kernels
\[
 t_\nu(u,v)=(2\pi)^{-1/2}
  \exp[-(u-v)^2/2-\nu(u^2+v^2)/4].
\]
For $\nu>0$, let $\xi_\nu=\operatorname{arcosh}(1+\nu/2)$ and
$\omega_\nu=\sqrt{\nu+\nu^2/4}$. Mehler's identity gives eigenvalues
$e^{-(k+1/2)\xi_\nu}$ on the oscillator basis of frequency $\omega_\nu$.
One may check the ground factor directly by integrating
$t_\nu(u,v)e^{-\omega_\nu v^2/2}$; its factor is
$(1+\nu/2+\omega_\nu)^{-1/2}$. For $\nu=0$, unitary Fourier transformation
gives multiplier $e^{-p^2/2}$ with norm one, not a normalizable ground vector.
This proves \eqref{f16v38:Gaussian-top} and the norm assertion.

The Gaussian counterpart $b_0=T_0(\psi_0\delta_0)$ of
\eqref{f16v38:loaded-vector} is a strictly positive centered nondegenerate
Gaussian function in all $51B$ coordinates: integrate the original static
Gaussian in its internal endpoint, keeping the positive kinetic square for
the retained output. Its norm square is $\mathcal A_{0,B,2}^{00}$.
After the preceding unitary changes it is still a positive Gaussian. Denote
its coordinates by $(z,s)$, with $z\in\mathbb R^{h_B}$ free and $s$ massive,
and let $\varphi_0(s)>0$ be the unit massive ground function. Then
\[
 F_B(z)=\int\varphi_0(s)\widetilde b_0(z,s)\,ds
\]
is a positive nondegenerate Gaussian. With unitary Fourier convention,
$\widehat F_B(0)>0$ and
$|\widehat F_B(p)|^2=C_B^{\rm F}e^{-p^{\mathsf T}D_Bp}$ for some $D_B\succ0$.
These constants are determined by the original endpoint, not freely chosen.

Put $N=n-2$. The massive ground contribution to
$\langle b_0,T_0^Nb_0\rangle$ is exactly
\[
 \Lambda_{{\rm G},B}^{N}C_B^{\rm F}(2\pi)^{h_B/2}
                        \det(NI+2D_B)^{-1/2}.
\]
Every other massive component has its transfer norm at most
$\Lambda_{{\rm G},B}e^{-\xi_{\min}}$, where the minimum is over the finitely
many positive $\nu_j$. Its total contribution is bounded by
$\Lambda_{{\rm G},B}^{N}e^{-N\xi_{\min}}\|b_0\|^2$, using the orthogonal
spectral decomposition, not an independent-mode assumption for $b_0$.
Expand $\det(I+(2D_B-2I)/n)^{-1/2}$ through order two with its finite third-order
remainder. The exponentially small massive contribution can be included in
$O_B(n^{-3})$. This proves \eqref{f16v38:Gaussian-asymptotic}. It also proves
$c_B=\Lambda_{{\rm G},B}^{-2}C_B^{\rm F}(2\pi)^{h_B/2}>0$.
\end{proof}

The determinant expression, rather than an unspecified $O(1/n)$ error, also
justifies the adjacent-ratio expansions
\begin{equation}
 \begin{split}
 \frac{\mathcal A_{0,B,n+1}^{00}}{\Lambda_{{\rm G},B}\mathcal A_{0,B,n}^{00}}
        &=1-a_B/n+O_B(n^{-2}),\\
 \frac{\mathcal A_{0,B,n}^{00}\mathcal A_{0,B,n+2}^{00}}
          {(\mathcal A_{0,B,n+1}^{00})^2}-1
        &=a_B/n^2+O_B(n^{-3}).
 \end{split}
 \label{f16v38:Gaussian-ratios}
\end{equation}
Their errors are not obtained by differentiating an uncontrolled asymptotic.

\subsection{A native endpoint-weighted spectral law}

Fix $B$ henceforth. Let $n=n_\gamma\to\infty$ be integers satisfying
\begin{equation}
 \frac{Bn_\gamma^7\log(n_\gamma+1)+n_\gamma^6\log\gamma}{\gamma}\longrightarrow0.
 \label{f16v38:regime}
\end{equation}
Every fixed integer multiple of $n_\gamma$, and the two adjacent durations,
then satisfy V37's entry conditions eventually. One concrete choice is
$n_\gamma=2\lfloor\gamma^\alpha/2\rfloor$, $0<\alpha<1/7$.

\begin{theorem}[Native Gamma edge profile without a claimed top identification]
\label{f16v38:native-edge}
Under $\mu_{\gamma,n_\gamma}$, define
\begin{equation}
 X_\gamma=-n_\gamma\log(\lambda/\Lambda_{{\rm G},B}),\qquad
 S_\gamma=n_\gamma(1-\lambda/\Lambda_{{\rm G},B}).
 \label{f16v38:edge-variables}
\end{equation}
Both variables converge in distribution to the Gamma law with shape $a_B$ and
rate one, whose density is
\begin{equation}
 1_{\{x>0\}}\,\frac{x^{a_B-1}e^{-x}}{\Gamma(a_B)}\,dx.
 \label{f16v38:Gamma}
\end{equation}
In addition,
\begin{equation}
 \mathbb E_{\mu_{\gamma,n_\gamma}}S_\gamma\longrightarrow a_B,
 \qquad
 \operatorname{Var}_{\mu_{\gamma,n_\gamma}}S_\gamma\longrightarrow a_B.
 \label{f16v38:edge-moments}
\end{equation}
At finite coupling these edge variables can be negative; the native top is not
replaced by $\Lambda_{{\rm G},B}$. The limit is a theorem about the spectral
measure of the specified endpoint vector of the \emph{native} transfer.
\end{theorem}
\begin{proof}
Write $A_{\gamma,n}=\mathcal A_{\gamma,B,n}^{00}$ for this proof only.
V37 gives
$|\log(A_{\gamma,n}/A_{0,n})|\le C(\rho_n+\gamma^{-s})$,
where $\rho_n=(72n-21)Bn^4/\gamma$. Choose $s=2$.
For each fixed integer $k\ge0$, put
$Y_\gamma=(\lambda/\Lambda_{{\rm G},B})^{n_\gamma}$. The exact tilted moments are
\begin{equation}
 \mathbb E Y_\gamma^k
  =\Lambda_{{\rm G},B}^{-kn_\gamma}
                     \frac{A_{\gamma,(k+1)n_\gamma}}{A_{\gamma,n_\gamma}}
  \longrightarrow(k+1)^{-a_B}.
 \label{f16v38:edge-power-moments}
\end{equation}
Use \eqref{f16v38:Gaussian-asymptotic} and the errors at \emph{both} durations.
No cancellation of those errors is assumed, and $k$ is fixed before the limit.
All original scalar factors are included through V37; a comparison of normalized
path probabilities alone would not imply this formula.

For completeness these moments determine an actual probability limit.
The bounded first moments give tightness of $Y_\gamma\ge0$. The bounded
$(k+1)$st moments give uniform integrability of its $k$th moments. Hence every
subsequential limit has moments $(k+1)^{-a_B}$. It is supported on $[0,1]$:
positive mass above $1+\delta$ would force its high moments to grow exponentially,
contrary to that formula. On $[0,1]$, polynomials determine a probability measure
by uniform approximation of continuous functions. The candidate $e^{-G}$,
$G\sim\operatorname{Gamma}(a_B,1)$, has exactly these moments. Thus
$Y_\gamma\Rightarrow e^{-G}$. The limit has no mass at zero. The continuous
mapping theorem applied to $-\log Y_\gamma$ proves the first limit, and uniform
convergence of $n(1-e^{-x/n})$ to $x$ on bounded $x$ proves the second.

For the two moment assertions let $r_n=A_{\gamma,n+1}/A_{\gamma,n}$ and
$v_n=A_{\gamma,n}A_{\gamma,n+2}/A_{\gamma,n+1}^2-1\ge0$. They are exactly
the mean spectral value and its relative variance under $\mu_{\gamma,n}$.
Equation \eqref{f16v38:Gaussian-ratios}, with the V37 errors added, gives
\[
 r_n/\Lambda_{{\rm G},B}=1-a_B/n+O_B(n^{-2})+O(\rho_{n+2}+\gamma^{-2}),
 \qquad
 v_n=a_B/n^2+O_B(n^{-3})+O(\rho_{n+2}+\gamma^{-2}).
\]
The exact formulas are
$\mathbb E S_\gamma=n(1-r_n/\Lambda_{{\rm G},B})$ and
$\operatorname{Var}S_\gamma=n^2(r_n/\Lambda_{{\rm G},B})^2v_n$.
The regime implies $n^2(\rho_{n+2}+\gamma^{-2})\to0$, proving
\eqref{f16v38:edge-moments}. This is a finite-difference use of quantified
scalar errors, not differentiation of those errors.
\end{proof}

For each fixed $x>0$, the exact spectral probability satisfies
\[
 \mu_{\gamma,n_\gamma}
 \{\lambda\ge\Lambda_{{\rm G},B}e^{-x/n_\gamma}\}
 \longrightarrow \frac1{\Gamma(a_B)}\int_0^x t^{a_B-1}e^{-t}\,dt>0.
\]
Weight above $\Lambda_{{\rm G},B}e^{x/n_\gamma}$ tends to zero. These are
weights, not assertions that no eigenvalue exists outside the stated band.

\subsection{What the edge law proves about the native top and the filtered vacuum weight}

\begin{proposition}[A one-sided top estimate and failure of endpoint Perron preparation]
\label{f16v38:consequences}
Along \eqref{f16v38:regime}, at fixed $B$,
\begin{equation}
 n_\gamma\bigl[\log(\Lambda_{{\rm G},B}/\Lambda_\gamma)\bigr]_+
      \longrightarrow0,
 \qquad
 |\langle\Omega_\gamma,v_{\gamma,n_\gamma}\rangle|^2\longrightarrow0.
 \label{f16v38:top-and-weight}
\end{equation}
The first statement also concerns the physical root-invariant top by
\eqref{f16v38:same-top}. It is not an upper estimate for that top. The second
statement concerns the actual unprojected endpoint-filtered vector; it is not
a claim that the normalized root projection has the same limit.
For even durations, let $\pi_{\gamma,n}^{\rm mid}$ be the full slice marginal
at time $n/2$ of the original identity-endpoint return law, in root coordinates.
Then
\begin{equation}
 d\pi_{\gamma,n}^{\rm mid}=|v_{\gamma,n}|^2dH,\qquad
 d_{\rm TV}(\pi_{\gamma,n_\gamma}^{\rm mid},\Omega_\gamma^2dH)
                       \longrightarrow1.
 \label{f16v38:midpoint-nonequilibrium}
\end{equation}
This is a gauge-dependent conditioned midpoint law, not a statement about its
restriction to gauge-invariant observables.
The effective scalar energy also satisfies
\begin{equation}
 -\log\frac{\mathcal A_{\gamma,B,n_\gamma+1}^{00}}
                 {\mathcal A_{\gamma,B,n_\gamma}^{00}}
   =-\log\Lambda_{{\rm G},B}+\frac{a_B}{n_\gamma}+o(n_\gamma^{-1}).
 \label{f16v38:effective-energy}
\end{equation}
The right side is measured against the Gaussian threshold, not an identified
native vacuum energy or a calibrated physical energy.
\end{proposition}
\begin{proof}
For each $x>0$ the positive limiting weight just proved ensures
$\Lambda_\gamma\ge\Lambda_{{\rm G},B}e^{-x/n_\gamma}$ eventually. Taking an
arbitrarily small fixed $x$ gives the first assertion.

A tight sequence of real probability measures converging to a non-atomic law
has maximum atomic mass tending to zero. To verify this directly, an atom of
mass at least $\epsilon$ cannot escape every compact set by tightness. Choose a
convergent subsequence of its locations inside a compact set; the portmanteau
inequality on arbitrarily small closed intervals about the limiting location
would then give an atom of at least $\epsilon$ in the limit. This is impossible
for the Gamma density. Apply this to the law of $X_\gamma$. Its atom at
$-n_\gamma\log(\Lambda_\gamma/\Lambda_{{\rm G},B})$ has weight exactly
$|\langle\Omega_\gamma,v_{\gamma,n_\gamma}\rangle|^2$, since the Perron
projection is one-dimensional. The conclusion also applies to the total weight
of any fixed number of eigenvalues. No rate for this vanishing is claimed.
For even $n$, splitting the exact return integral at its midpoint gives two
identical positive half-path amplitudes. Their squared norm is the full return
normalizer, so the normalized midpoint density is precisely $|v_{\gamma,n}|^2$.
The Hellinger affinity of this density with $\Omega_\gamma^2$ equals
$\langle\Omega_\gamma,v_{\gamma,n}\rangle$, since both vectors are positive.
For probability densities $p,q$, the pointwise bound $\min(p,q)\le\sqrt{pq}$
gives $d_{\rm TV}(p,q)\ge1-\int\sqrt{pq}$. The vanishing squared overlap
therefore proves \eqref{f16v38:midpoint-nonequilibrium}.
Finally the adjacent-ratio estimate in the preceding proof and
$n(\rho_{n+2}+\gamma^{-2})\to0$ give \eqref{f16v38:effective-energy}.
\end{proof}

There is no contradiction with \eqref{f16v38:overlap}. At each fixed $B,\gamma$,
all-duration filtering eventually selects the positive Perron vector. Here
$\gamma$ changes while $n_\gamma=o(\gamma)$, and the proved overlap lower bound
is exponentially small. Its finite-duration top bracket has logarithmic width
$96\gamma B/(n_\gamma-2)$, which does not vanish in this regime. This diagnoses
the weakness of that \emph{certificate}, not a necessary mixing-time lower bound
or an upper bound on the true overlap.

Equation \eqref{f16v38:Gauss-mixture} is especially important now. If
$\theta_{\gamma,n}$ is small, root projection followed by renormalization can
radically change the spectral law and its vacuum weight. The fixed-$B$ Gaussian
exponent above counts linearized gauge-gradient directions before a physical quotient.
It cannot be relabelled a physical particle multiplicity, a toron-only exponent,
or a root-Gauss-reduced Gamma law. No such projected law is proved here.

\subsection{Direction after the spectral calculation and verification boundary}

The immediate unresolved task is now more specific than obtaining another
scalar coefficient: control the root-projected endpoint spectral measure and
its overlap, or analyze the compact soft sector on durations beyond the present
Gaussian entrance. V37's scalar sequence does have a genuine native spectral
meaning, but its polynomial-duration limit is not Perron dominated. Another
unprojected finite-time coefficient will not by itself reverse
\eqref{f16v38:top-and-weight}. The older complementary-channel and same-top
operator obligations remain, as do physical time and continuum reconstruction.

The accompanying diagnostic is
\begin{center}\small\path{verify_ym_f16_v38_endpoint_spectrum.py}.\end{center}
It checks the endpoint moment offset and Jacobian cancellation, normalized
spectral tilting, root projection, the sharp positive-ratio overlap inequality,
Mehler normalization, Gaussian heat powers, moment identification and finite
error exponents. These are finite algebraic fixtures, not numerical enclosures
of the native transfer spectrum or a formal proof of the analytic limit.
Execution and preservation reports state the actual scope. Only the title
version and this terminal insertion change V37. No earlier claim is silently
corrected, and no archive-wide proof certification is asserted.

\begin{firewall}
V38 realizes the exact dressed endpoint sequence as native transfer moments,
quantifies its Perron visibility with a paid coupling/volume cost, and derives
a Gamma endpoint-weighted edge law in the inherited polynomial regime at fixed
spatial width. It proves a one-sided native top estimate and vanishing Perron
weight of that unprojected filtered vector, not a physical excited-state gap.
Root-projected spectral weights, a native upper edge, unrestricted regulators,
physical vacuum identification, time calibration and an interacting continuum
mass gap remain unproved. No intermediate dressed projection or changed scalar
normalization is inserted into the actual transfer powers.
\end{firewall}
% END F16 V38 APPEND

% BEGIN F16 V39 APPEND -- 2026-09-18
% Insert before the final document-ending command in the cumulative V38.
% All predecessor proofs, measures, endpoint factors and clocks remain unchanged.

\clearpage
\section{V39: root projection, orbit weights, and physical endpoint preparation}
\label{f16v39:context}

V38 leaves the weight of the root-invariant component uncontrolled. Its
unprojected endpoint edge law cannot therefore be assigned a physical exponent
by subtracting the number of gauge directions. We return to the compact root
action itself. An exact forest disintegration separates amplitude projection
from averaging a probability density. It also allows a quantitative estimate
of the actual projected norm. At fixed spatial width that norm has a polynomial,
not merely exponentially small, coupling scale.

The resulting estimates do not identify a projected Gamma law. They instead
give two-sided power bounds on the projected endpoint moments. Comparing three
admitted durations proves that, on a smaller but explicit polynomial range,
the normalized \emph{root-projected} filtered endpoint still has vanishing
Perron overlap. Unlike the corresponding V38 conclusion, its midpoint and
stationary densities are both gauge invariant. This is a failure of the
specified source to prepare the finite-regulator vacuum on those durations,
not a conclusion about the physical mass gap.

\subsection{Dependencies and the unchanged compact root action}

The V38 appendix and audit were read completely. The native inputs below are
V35's complete-law tails and local linear-tilt moments, V37's scalar comparison,
and V38's endpoint loading and fixed-width Gaussian asymptotic. We do not assume
an asymptotic for a root-twisted endpoint matrix element. The forest/Haar
mechanism already appears in f15 V10 and is credited. Its application here is
to the remaining \emph{coarse root} group and to a generally non-invariant
endpoint density, not to a Wilson measure asserted to factor along these
orbits. Targeted searches of the active source and supplied archives were
followed by identified passage reads; this is not an exhaustive novelty audit.
The next scheduled checkpoint remains V40.

For context, Sections II--III.1 of Burbano--Bauer's
\href{https://arxiv.org/html/2409.13812}{maximal-tree formulation} distinguish
based gauge fixing from the remaining global constraint. Bahr--Thiemann's
\href{https://arxiv.org/html/0709.4636}{non-Abelian projected-state construction},
equations (2.4)--(2.5), uses normalized compact group averaging. These classical
constructions are context, not native gap hypotheses. We prove the identities
and estimates used here directly, in the inherited Wilson orientation.

Let $G=\SU(2)$, $B=m^3$, $m\ge4$, and retain
\[
 \mathcal H_{\rm rt}=L^2(G^{5B}\times G^{12B},dH),\quad
 T_\gamma=T_{{\rm full},\gamma},\quad
 \Pi=\Pi_{\rm G},\quad \Lambda_\gamma=\|T_\gamma\|.
\]
Write $\mathsf U_g$ for the root action, with
\[
 X_{b,a}\mapsto g_bX_{b,a}g_b^{-1},\qquad
 Z_e\mapsto g_{s(e)}Z_eg_{t(e)}^{-1}.
\]
The moving frames are equivariant under these conjugations. The full transfer
commutes with this action and has the same Perron top on $\operatorname{Ran}\Pi$,
as already established in V38. Let $f_\gamma$ be its \emph{same} once-loaded
unit endpoint vector. For even $n\ge2$ put
\begin{equation}
 v_{\gamma,n}=
 \frac{T_\gamma^{(n-2)/2}f_\gamma}
 {\|T_\gamma^{(n-2)/2}f_\gamma\|},\qquad
 \theta_{\gamma,n}=\|\Pi v_{\gamma,n}\|^2,
 \qquad v^{\rm G}_{\gamma,n}=\frac{\Pi v_{\gamma,n}}{\sqrt{\theta_{\gamma,n}}}.
 \label{f16v39:vectors}
\end{equation}
All are positive; the first is invariant under simultaneous diagonal
conjugation, because the identity endpoint and its original carrier are.
We use even durations so these vectors are obtained by positive integer
transfer powers. No new continuum-time interpolation is required.

\subsection{An exact forest disintegration of projected amplitudes}

Choose a spanning tree $\mathcal F$ on the connected coarse multigraph, using
one actual bridge for each tree edge, and choose one root $o$. Set
\[
 r_B=B-1,\qquad D_B=3(B-1),\qquad \ell_B=D_B/2.
\]
Here $D_B$ is a dimension, not the Fourier covariance matrix used in V38.
Let $t_b$ be the ordered product from $o$ to $b$ along $\mathcal F$, with
$t_o=I$ and inverses on reverse traversals. Define
\[
 \widetilde X_{b,a}=t_bX_{b,a}t_b^{-1},\qquad
 \widetilde Z_e=t_{s(e)}Z_et_{t(e)}^{-1}.
\]
Every tree value is $I$. Let $\eta$ collect all $5B$ transformed chords and
$12B-r_B$ non-tree bridges, with product Haar measure $d\nu(\eta)$.
There is still a global simultaneous conjugation on $\eta$.

\begin{proposition}[Haar fibres and the actual orbit weight]
\label{f16v39:fibres}
This is a global bijection and an exact probability-Haar identity:
\begin{equation}
 dH(X,Z)=\prod_{b\ne o}dH(t_b)\,d\nu(\eta).
 \label{f16v39:Haar}
\end{equation}
For any positive, diagonal-conjugation-invariant unit vector $v$, set
\[
 r(\eta)=\int v(t,\eta)^2\,dt,\qquad
 a(\eta)=\int v(t,\eta)\,dt,\qquad
 k(\eta)=\frac{a(\eta)^2}{r(\eta)},\qquad
 \theta=\int a(\eta)^2\,d\nu(\eta).
\]
Then $0<k\le1$ almost everywhere and $\theta=\|\Pi v\|^2=\int kr\,d\nu$.
The unprojected density restricted to invariant observables and the normalized
projected density have respective quotient representatives
\begin{equation}
 r(\eta)\,d\nu(\eta),\qquad
 \frac{k(\eta)}{\theta}r(\eta)\,d\nu(\eta).
 \label{f16v39:reweight}
\end{equation}
In particular their total variation on invariant events is exactly
\begin{equation}
 \frac12\int r(\eta)\left|\frac{k(\eta)}{\theta}-1\right|d\nu(\eta).
 \label{f16v39:quotient-TV}
\end{equation}
The normalized projected density uniquely maximizes Hellinger affinity with
$v^2dH$ among invariant probability densities, and that maximum is
$\sqrt\theta$.
\end{proposition}
\begin{proof}
The inverse is $X_{b,a}=t_b^{-1}\widetilde X_{b,a}t_b$ and
$Z_e=t_{s(e)}^{-1}\widetilde Z_et_{t(e)}$. In particular a tree bridge is
$t_s^{-1}t_t$. Successive tree Haar translations and then the left/right
translations of all non-tree variables prove \eqref{f16v39:Haar}.
No word is reordered and no determinant is introduced.

The based root group $g_o=I$ changes $t_b$ to $t_bg_b^{-1}$ and leaves $\eta$
fixed. Thus its average sends $v$ to $a(\eta)$. Write a general root element
as $g_b=h_bq$, with $h_o=I$. Diagonal invariance of $v$ shows that full root
averaging equals this based average; $a,r,k$ are themselves invariant under
the remaining simultaneous conjugation of $\eta$. Cauchy--Schwarz on the
probability Haar fibre gives $a^2\le r$. Integration yields the formula for
$\theta$ and both normalized quotient densities. The sign set of $k/\theta-1$
is invariant, so the ordinary quotient total variation is also the supremum
over invariant events. Fibrewise the coefficient is
\[
 k(\eta)=\left(\int\sqrt{v(t,\eta)^2/r(\eta)}\,dt\right)^2.
\]
It is the squared affinity of the actual conditional orbit density with Haar;
that density is not asserted to be Haar.

If $q$ is any invariant probability density, its square root is an invariant
unit vector. Hence $\int v\sqrt q\,dH=\langle\Pi v,\sqrt q\rangle\le\sqrt\theta$.
Equality in Cauchy--Schwarz requires $\sqrt q=\Pi v/\sqrt\theta$ almost
everywhere, proving the final assertion. The transfer and its kinetic energy
have not been changed by the coordinate identity.
\end{proof}

Small $\theta$ alone does not bound \eqref{f16v39:quotient-TV}. For example, if
$v(t,\eta)=u(t)b(\eta)$ with both factors normalized, $k$ is constant and the
two quotient probabilities coincide, even if $\theta$ is small. With an
$\eta$-dependent orbit profile the reweighting is genuine. Probability averaging
and amplitude projection must therefore remain distinct.

\subsection{A polynomial lower bound from the literal kinetic factors}

\begin{proposition}[Orbit Fisher bound and root-projection visibility]
\label{f16v39:lower}
At every finite $B$, $\gamma\ge1$ and even $n\ge2$, there is a constant
$c_B>0$, independent of $n,\gamma$, such that
\begin{equation}
 \theta_{\gamma,n}\ge c_B\gamma^{-\ell_B}.
 \label{f16v39:lower-theta}
\end{equation}
More precisely, for a unit Lie algebra generator acting at a single coarse
root, with skew derivative $\mathcal D_b$,
\begin{equation}
 \|\mathcal D_b v_{\gamma,n}\|_2^2\le48\gamma B.
 \label{f16v39:Fisher}
\end{equation}
These estimates use no small-field event or Gaussian comparison.
\end{proposition}
\begin{proof}
Consider any positive output $v=T_\gamma h/\|T_\gamma h\|$, allowing the
initial endpoint to be a positive finite measure. The vector in
\eqref{f16v39:vectors} has this form. Along a one-root path $e^{t\tau}$,
$|\tau|=1$, the output magnetic factor is exactly invariant. In the kinetic
integrand each of its $24B$ words contains at most two factors $e^{\pm t\tau}$.
The second derivative of its scalar quaternion part has absolute value at
most four: the two second derivatives and two cross derivatives each have
operator norm at most one. Hence the logarithm of the complete kinetic
integrand has second derivative at least $-96\gamma B$.

Integrating the nonroot Haar variables and the positive input preserves this
lower bound on the logarithm, since
\[
 \frac{d^2}{dt^2}\log\int e^{l_t}\,d\sigma
     =\mathbb E_t l_t''+\operatorname{Var}_t(l_t')\ge-96\gamma B.
\]
Here $d\sigma$ includes the unchanged positive input and all other factors.
At a fixed regulator derivatives in these compact group directions are bounded
relative to the integrand, so differentiation is justified. Thus
$\mathcal D_b^2\log v\ge-96\gamma B$.
The action is Haar preserving. Integrating $\mathcal D_b(v^2\mathcal D_b\log v)$
gives zero and proves
\[
 \|\mathcal D_bv\|_2^2
       =-\tfrac12\int v^2\mathcal D_b^2\log v\,dH\le48\gamma B.
\]
This is integration on the compact group, not through a principal-chart cut.
Along an orbit the original equivariant frames have no new cut derivative.

For $h_o=I$, write $h_b=\operatorname{Exp}(x_b)$. Unitarity, the fundamental
theorem along each one-parameter subgroup, and a product telescoping estimate give
\[
 \|\mathsf U_hv-v\|_2\le\sqrt{48\gamma B}\sum_{b\ne o}|x_b|.
\]
On $|x_b|\le\delta=(16r_B\sqrt{B\gamma})^{-1}$ the squared distance is at most
$3/16$, so the real inner product is at least $29/32$. It is nonnegative for
all $h$, by positivity of $v$. The normalized $\SU(2)$ cap satisfies, for
$0<\delta\le1$,
\[
 H\{|x|\le\delta\}
 =\frac2\pi\int_0^\delta\sin^2t\,dt\ge\frac8{3\pi^3}\delta^3.
\]
The $r_B$ based group factors are independent Haar variables. Since
$\theta=\int\langle v,\mathsf U_hv\rangle dh$, integrating over these caps
proves \eqref{f16v39:lower-theta}; for example one may take
\[
 c_B=\frac12\left[\frac8{3\pi^3}(16r_B\sqrt B)^{-3}\right]^{r_B}.
\]
The constant is conservative, but is explicit and contains no duration.
\end{proof}

\subsection{The complementary upper bound from the actual midpoint concentration}

Fix $B$ for the remainder of the asymptotic statements. Choose an integer
$p_B>\ell_B$, for instance $p_B=\lceil\ell_B\rceil+1$, and use V35's complete
comparison-ball entry at tail exponent $s_B=p_B+2$. This is a fixed exponent
once $B$ is fixed. Its asymptotic regime is still V38's fixed-$B$ polynomial
regime. Constants below may depend on $B$; no uniform-in-$B$ estimate is claimed.

\begin{theorem}[The actual root-component scale]
\label{f16v39:upper}
For all sufficiently large couplings satisfying that explicit V35 entry, and
all admitted even durations,
\begin{equation}
 \boxed{\qquad
 c_B\gamma^{-\ell_B}\le\theta_{\gamma,n}
       \le C_B(n/\gamma)^{\ell_B},\qquad \ell_B=\tfrac32(B-1).
 \qquad}
 \label{f16v39:theta-bounds}
\end{equation}
In particular, along the fixed-$B$ regime of V38, $\theta_{\gamma,n_\gamma}\to0$.
At the fixed loading duration $n=2$ it has matching upper and lower powers:
$c_B\gamma^{-\ell_B}\le\theta_{\gamma,2}\le C_B\gamma^{-\ell_B}$.
No leading coefficient or exact $n$-dependence of $\theta_{\gamma,n}$ is asserted.
\end{theorem}
\begin{proof}
For a full slice define the positive tree weight
\[
 A(Z)=1+\frac\gamma n\sum_{e\in\mathcal F}\operatorname{dist}(I,Z_e)^2.
\]
In the forest coordinates this depends only on $t$, not on $\eta$. Its negative
power has an exact Haar upper bound. Tree edge values are independent Haar
under the product fibre reference; their change from $t$ is triangular. Using
complete principal charts only for this reference integral gives
\begin{equation}
 \begin{split}
 J_p:=\int A(t)^{-p}\,dt
 &\le (2\pi^2)^{-r_B}(n/\gamma)^{D_B/2}
           \int_{\mathbb R^{D_B}}(1+|z|^2)^{-p}dz\\
 &= (2\pi^2)^{-r_B}\pi^{D_B/2}
       \frac{\Gamma(p-D_B/2)}{\Gamma(p)}(n/\gamma)^{\ell_B},
       \qquad p>\ell_B.
 \end{split}
 \label{f16v39:negative-moment}
\end{equation}
The Haar Jacobian is bounded above by one; the full native density is not
replaced by this reference.

Weighted Cauchy--Schwarz separately on every orbit gives
\[
 \left(\int v(t,\eta)dt\right)^2
       \le J_p\int v(t,\eta)^2A(t)^pdt.
\]
After integrating $\eta$ this yields the general inequality
$\theta\le J_p\,\mathbb E_{v^2dH}A^p$.
By V38 the native midpoint density for even $n$ is exactly
$v_{\gamma,n}^2dH$. In the complete joint path its tree weight is
\[
 A=1+n^{-1}\sum_{e\in\mathcal F}|y_{n/2,e}|^2.
\]
On V35's joint comparison ball, its proved linear-tilt bound implies
$\mathbb E|y_{n/2,e}|^{2p_B}\le C_{p_B}n^{p_B}$. Summing the fixed $r_B$ sites
therefore bounds the conditional expectation of $A^{p_B}$ by $C_B$.
On the complement, compactness gives $A^{p_B}\le C_B\gamma^{p_B}$ for $n\ge2$,
whereas the complete native tail probability is at most $\gamma^{-s_B}$.
Consequently
\[
 \mathbb E_{v_{\gamma,n}^2dH} A^{p_B}\le C_B+C_B\gamma^{p_B-s_B}\le C_B'.
\]
This uses high \emph{interior} moments derived from the already proved
sub-Gaussian bound; it does not claim an unproved high-degree complete-tail
estimate. The finite cap pays the complement. Equations
\eqref{f16v39:negative-moment} and \eqref{f16v39:lower-theta} prove the theorem.
\end{proof}

The upper bound also implies that the full unprojected midpoint density is far
in total variation from \emph{every} root-invariant probability density: the
bound is at least $1-\sqrt{\theta_{\gamma,n}}$ by Proposition~\ref{f16v39:fibres}.
This observation is gauge dependent and does not alone say that invariant
observables fail to be stationary. The following spectral argument is needed
to obtain a statement for the projected source itself.

\subsection{Root-projected native moments with both polynomial factors retained}

Define the exact normalized physical source and its moments by
\begin{equation}
 f_\gamma^{\rm G}=\Pi f_\gamma/\sqrt{\theta_{\gamma,2}},\qquad
 M_{\gamma,n}^{\rm G}
    =\langle f_\gamma^{\rm G},T_\gamma^{n-2}f_\gamma^{\rm G}\rangle,
 \qquad n\ge2.
 \label{f16v39:projected-moments}
\end{equation}
There is no insertion of a new carrier between transfer powers. Commutation
and \eqref{f16v38:moment-identity} give, for even $n$,
\begin{equation}
 M_{\gamma,n}^{\rm G}
   =\frac{\mathcal A_{\gamma,B,n}^{00}}{\mathcal A_{\gamma,B,2}^{00}}
                         \frac{\theta_{\gamma,n}}{\theta_{\gamma,2}}.
 \label{f16v39:exact-projected-moment}
\end{equation}
The normalization $\theta_{\gamma,2}$ is essential.

\begin{theorem}[Two-sided projected power bounds, not a projected edge law]
\label{f16v39:moment-envelope}
Let $a_B=3(B+2)/2$ and $\Lambda_{{\rm G},B}$ retain exactly their V38 meanings.
At fixed $B$, there exist $0<c_B'\le C_B'<\infty$ and $n_0(B)$ such that,
for even $n\ge n_0(B)$ satisfying the above entry and V37's scalar error bound
with that error at most one,
\begin{equation}
 \boxed{\quad
 c_B'\Lambda_{{\rm G},B}^{n-2}n^{-a_B}
 \le M_{\gamma,n}^{\rm G}
 \le C_B'\Lambda_{{\rm G},B}^{n-2}n^{-9/2}.
 \quad}
 \label{f16v39:projected-envelope}
\end{equation}
The constants are independent of $n,\gamma$ in this admitted set.
\end{theorem}
\begin{proof}
The new bounds at $n$ and two give
$c_B''\le\theta_{\gamma,n}/\theta_{\gamma,2}\le C_B''n^{\ell_B}$.
V37's complete scalar comparison, including all kinetic and endpoint factors,
and V38's fixed-$B$ Gaussian asymptotic give
\[
 c_B'''\Lambda_{{\rm G},B}^{n-2}n^{-a_B}
 \le\frac{\mathcal A_{\gamma,B,n}^{00}}{\mathcal A_{\gamma,B,2}^{00}}
 \le C_B'''\Lambda_{{\rm G},B}^{n-2}n^{-a_B}.
\]
Errors at both durations are paid by their absolute bounds. Insert them into
\eqref{f16v39:exact-projected-moment} and use $a_B-\ell_B=9/2$.
This identity of exponents gives the upper envelope, \emph{not} an asymptotic
formula with exponent $9/2$ or a root-projected Gamma distribution.
\end{proof}

\subsection{A three-duration test excludes projected vacuum preparation}

The next step uses only positive spectral measures, the same native top, and
the envelope just proved. It does not transfer V38's unprojected edge law
through a projection of vanishing norm.

\begin{proposition}[Finite three-duration Perron-weight certificate]
\label{f16v39:three-times}
Suppose even $n_0(B)\le n\le M\le L$ all satisfy
\eqref{f16v39:projected-envelope}. Enlarge a fixed $C_B^0$ if necessary. Then
\begin{equation}
 \boxed{\quad
 |\langle\Omega_\gamma,v^{\rm G}_{\gamma,n}\rangle|^2
 \le C_B^0 n^{a_B}M^{-9/2}
 \exp\left[\frac{(M-n)(a_B\log L+C_B^0)}{L-2}\right].
 \quad}
 \label{f16v39:atom-bound}
\end{equation}
The bound may be replaced by its minimum with one. Every moment and top in
this inequality belongs to the root-invariant restriction of the original
$T_\gamma$; that top is $\Lambda_\gamma$.
\end{proposition}
\begin{proof}
The variational moment bound at $L$ yields
\[
 \Lambda_\gamma\ge (M_{\gamma,L}^{\rm G})^{1/(L-2)}
 \ge\Lambda_{{\rm G},B}
        \exp[-(a_B\log L+C_B^0)/(L-2)].
\]
If $q_\gamma^{\rm G}=|\langle\Omega_\gamma,f_\gamma^{\rm G}\rangle|^2$, then
positivity of its spectral measure gives
$M_{\gamma,M}^{\rm G}\ge q_\gamma^{\rm G}\Lambda_\gamma^{M-2}$.
The actual filtered top weight is
$q_\gamma^{\rm G}\Lambda_\gamma^{n-2}/M_{\gamma,n}^{\rm G}$.
Therefore it is at most
\[
 \frac{M_{\gamma,M}^{\rm G}}
 {M_{\gamma,n}^{\rm G}\Lambda_\gamma^{M-n}}.
\]
Apply the upper envelope at $M$, the lower envelope at $n$, and the just-proved
lower bound on the \emph{native} top. The powers of $\Lambda_{{\rm G},B}$ cancel
exactly, giving \eqref{f16v39:atom-bound}. No upper estimate for $\Lambda_\gamma$
or assumed rate of projected filtering has been used.
\end{proof}

\begin{theorem}[Projected endpoint nonequilibrium in an explicit smaller entrance]
\label{f16v39:physical-nonequilibrium}
Fix $B=m^3$, $m\ge4$, and
\begin{equation}
 0<\alpha<\frac{3}{7(B+2)},\qquad
 n_\gamma=2\lfloor\gamma^\alpha/2\rfloor.
 \label{f16v39:short-regime}
\end{equation}
Then for the normalized root-projected filtered source in
\eqref{f16v39:vectors},
\begin{equation}
 |\langle\Omega_\gamma,v^{\rm G}_{\gamma,n_\gamma}\rangle|^2\longrightarrow0,
 \qquad
 d_{\rm TV}\left(|v^{\rm G}_{\gamma,n_\gamma}|^2dH,
                                      \Omega_\gamma^2dH\right)\longrightarrow1.
 \label{f16v39:physical-limit}
\end{equation}
Both densities in the second expression are root invariant, so the total
variation is attained by invariant events. It is thus a genuine distinction
on the invariant configuration algebra for this projected return source.
\end{theorem}
\begin{proof}
Choose fixed exponents
\[
 \frac{B+2}{3}\alpha<\beta<\chi<\frac17,
 \quad M_\gamma=2\lfloor\gamma^\beta/2\rfloor,
 \quad L_\gamma=2\lfloor\gamma^\chi/2\rfloor.
\]
They exist precisely under the strict bound \eqref{f16v39:short-regime}.
At fixed $B$ all three durations satisfy the inherited V35--V37 entries
at every fixed tail exponent eventually. Moreover
\[
 \frac{M_\gamma\log L_\gamma}{L_\gamma}\to0,\qquad
 n_\gamma^{a_B}M_\gamma^{-9/2}
           =O_B(\gamma^{a_B\alpha-(9/2)\beta})\to0.
\]
Proposition~\ref{f16v39:three-times} proves the vanishing projected overlap.
The projected vector is positive, so its probability affinity with
$\Omega_\gamma^2dH$ is that overlap before squaring. The inequality
$\min(p,q)\le\sqrt{pq}$ gives the total-variation conclusion. Since both
densities are invariant, the set where one exceeds the other is invariant.

For interpretation, the same vector is the exact midpoint amplitude of the
return prepared by projecting the original once-loaded endpoint on \emph{both}
sides. Indeed $\Pi T_\gamma=T_\gamma\Pi$ and the normalizing power is
\eqref{f16v39:projected-moments}. This is not obtained by averaging the density
of the unprojected bridge: Proposition~\ref{f16v39:fibres} describes that
additional distinction explicitly. The definition changes the source by the
required Gauss projection, not the transfer or its spectral top.
\end{proof}

One explicit exponent choice is
$\alpha=1/[7(B+2)]$, $\beta=1/14$, $\chi=3/28$. It gives the proved bound
$C_B\gamma^{-3/28}(1+o(1))$ on the squared projected Perron overlap.
The duration exponent is small and not asserted optimal or uniform in $B$.
At fixed coupling and regulator, sufficiently long filtering still selects the
Perron vector; this theorem changes the coupling and duration together.

\subsection{What is resolved and the next distinct estimate}

V39 determines that the root-invariant component in V38 really vanishes, with
explicit polynomial upper and lower bounds. The projected moment sequence is
therefore not approximated by the unprojected one at a relative error inferred
from ordinary total variation. The orbit affinity $k(\eta)$ is the exact missing
weight on invariant observables. Controlling its dependence on the remaining
variables would give information beyond the present power envelope.

The shorter-duration theorem also shows that root projection alone does not
prepare the vacuum for this source. A new projected edge law, a useful longer-time
vacuum-overlap estimate, or a compact nonlinear-period analysis is still needed.
The number $9/2$ above is an envelope exponent, not a proved projected density
of states, particle multiplicity, or physical spectral gap. Likewise none of
the clocks used in earlier sampling estimates is substituted for Euclidean
transfer time. This maintains the supplied memoranda's separate normalization,
locality, and physical-transfer requirements. No excited spectral ratio,
unrestricted-regulator theorem, time calibration, or interacting continuum
construction follows from the present endpoint result.

\subsection{Verification and preservation}

The new diagnostic is
\begin{center}\small\path{verify_ym_f16_v39_root_orbits.py}.\end{center}
It passes 144 exact assertions in 48 named families, including noncommutative
forest identities, orbit weights, the kinetic derivative bound, projected
moment normalization, and three-duration arithmetic. All 35 supplied V4--V38
checkers were rerun unchanged and passed. These are finite fixtures, not native
spectral enclosures or formal proof verification; actual outputs are packaged.

Removing this insertion and reversing the one title version recovers V38 byte
for byte. The new spectral conclusions inherit V35's midpoint moments and
V37--V38's scalar/Gaussian comparisons. The orbit and kinetic arguments have
separate compact proofs. The analytic lineage has not been independently
recertified. The next checkpoint remains V40.

\begin{firewall}
V39 proves polynomial bounds for the actual root-projected norm, an exact
orbit-affinity reweighting, native projected moment envelopes, and failure of
projected endpoint Perron preparation on a stated smaller polynomial duration
range at fixed spatial width. It does not derive a projected Gamma law, remove
the remaining invariant orbit weights, estimate a physical excited-state gap,
or identify a continuum Hamiltonian or its time scale. The source projection
is explicit; the native transfer, its same physical top, all scalar factors,
and the remaining nonlinear compact variables are retained.
\end{firewall}
% END F16 V39 APPEND

% BEGIN F16 V40 APPEND -- 2026-09-18
% Insert before the final document-ending command of the supplied cumulative V39.
% All predecessor bodies, native measures, endpoint factors, and clocks are preserved.

\clearpage
\section{V40: a direct root-orbit integral and the projected spectral edge}
\label{f16v40:context}

V39 estimates the norm of the root projection between two different powers of
its duration. Dividing an ordinary density approximation by that small norm
would not improve the estimate. We instead include the compact root average in
the original positive endpoint integral before making any approximation. The
based root variables become additional, explicitly normalized integration
variables. One anchored endpoint still confines them. At Gaussian order their
integral can be evaluated exactly by a temporal affine shift in the inherited
gradient subspace. This supplies the missing relative normalization rather than
an assumed gauge-volume factor.

At fixed spatial width we obtain the leading projected norm, a normalized
projected moment asymptotic, and a native root-invariant Gamma edge law of shape
$9/2$ throughout the existing polynomial-duration entrance. This strengthens
V39's power envelope and its smaller-duration non-preparation theorem. The
transfer and its physical top are unchanged. The result is not a density of
states, an estimate of the excited-state gap, or a physical time calibration.

\subsection{V40 checkpoint, sources, and the exact additional integral}

The supplied V39 appendix and analytic audit were read completely. The current
cumulative source, f14 V100, f15 V100, and Grand Tensor V076 were inventoried and
searched; identified forest, scalar, and Gaussian-history passages were checked.
This is a targeted checkpoint, not independent certification of those archives.
The next scheduled checkpoint is V45. The forest quotient and the distinction
between probability averaging and amplitude projection remain owned by V39 and
its cited f15 V10 construction. We do not repeat them as new results.

For primary-source context, Burbano--Bauer,
\href{https://arxiv.org/html/2409.13812}{\emph{Gauge Loop-String-Hadron Formulation
on General Graphs}}, Sections II--III.1, separates based gauge fixing from the
remaining global conjugation. Bahr--Thiemann,
\href{https://arxiv.org/html/0709.4636}{\emph{Gauge-invariant coherent states II}},
Section 2, defines compact group projection; its discussion of exceptional orbit
geometry is a reason not to assume a uniform stationary-phase formula. Those
primary HTML statements were checked. Neither a coherent-state asymptotic nor a
Hamiltonian truncation is imported. The compact domination and Gaussian
integral needed here are proved below. The supplied import memoranda remain
methodological proposals, not premises supplying a physical spectral theorem.

Fix $B=m^3$, $m\ge4$, and one coarse root $o$. Put
\[
 r_B=B-1,\qquad k_B=3r_B,\qquad \ell_B=k_B/2.
\]
For $h=(h_b)_{b\ne o}$, set $h_o=I$ and
\begin{equation}
 Z(h)_e=h_{s(e)}h_{t(e)}^{-1}.
 \label{f16v40:pure-endpoint}
\end{equation}
Use the original orientation, all four parallel bridges, and every periodic
seam. Let $a_{\gamma,Z}(X)$ be V1's normalized dressed endpoint vector. Its exact
equivariance gives
\[
 F_{\gamma,2}(Z(h))=F_{\gamma,2}(I),\qquad
 a_{\gamma,Z(h)}(hXh^{-1})=a_{\gamma,I}(X).
\]
Indeed the carrier's five chord vectors in a cube undergo one common adjoint
rotation, which preserves their Gaussian quadratic and their Haar Jacobians;
the mixed multiplier and its actual integral transform covariantly.

Let $\mathcal A^{\rm G}_{\gamma,B,n}$ denote the original identity-endpoint
return with a root projection on both endpoints. More precisely it is
$j_{\rm br,0}\langle \Pi s_\gamma,T_\gamma^n\Pi s_\gamma\rangle$, where
$s_\gamma=a_{\gamma,I}\delta_I$ is interpreted through V38's one-step loading.
For $n\ge2$ this is a finite positive scalar. Equivalently it is
$\langle b_\gamma^{\rm end},T_\gamma^{n-2}\Pi b_\gamma^{\rm end}\rangle$.
This uses no inner product of an unsmoothed delta. Since $s_\gamma$ is invariant
under simultaneous conjugation, full root averaging equals based averaging.
Commutation and $\Pi^2=\Pi$ reduce the two averages to one:
\begin{equation}
 \mathcal A^{\rm G}_{\gamma,B,n}
 =j_{\rm br,0}\int_{G^{r_B}}
    \langle s_\gamma,T_\gamma^n\mathsf U_hs_\gamma\rangle\,dH^{r_B}(h).
 \label{f16v40:one-average}
\end{equation}
All identities involving endpoint distributions are justified by the bounded
kernel after one loading, then Tonelli. No residual global orientation is fixed.

Write $h_b=\operatorname{Exp}(\xi_b/\sqrt\gamma)$ on its complete principal
chart, and let $w=(\zeta,y_1,\ldots,y_{n-1})$ have exactly its V34 meaning.
Substitute $Z_0=I$, $Z_n=Z(h)$ into the \emph{original} classical action in
unmoving coordinates, leaving the endpoint group products ordered. Denote the
result by $U^+_\gamma(w,\xi)$. Define
\begin{equation}
 \begin{split}
 d_+&=(72n-21)B+3(B-1)=(72n-18)B-3,\\
 \mathcal J^+_\gamma(w,\xi)
   &=J^f_\gamma(\zeta)
     \prod_{i=1}^{n-1}\prod_{e\ {\rm bridge}}j(y_{i,e}/\sqrt\gamma)
     \prod_{b\ne o}j(\xi_b/\sqrt\gamma),\\
 \mathcal Z^+_{\gamma,B,n}
   &=\int_{\mathcal P^+_\gamma}
              \mathcal J^+_\gamma e^{-U^+_\gamma}\,dw\,d\xi .
 \end{split}
 \label{f16v40:root-mass}
\end{equation}
Here $\mathcal P^+_\gamma$ is the product of complete three-group charts.
There is no retained-bridge Haar factor at the final pure-gauge endpoint: that
endpoint is not integrated against bridge Haar. Its variables are the $r_B$
root groups, whose full Haar factors and constant density are retained instead.
Chord endpoint powers remain $1/2$.

\begin{proposition}[Exact projected scalar before the weak-field limit]
\label{f16v40:exact}
At every finite regulator with $n\ge2$,
\begin{equation}
 \boxed{\quad
 \mathcal A^{\rm G}_{\gamma,B,n}
 =c_{H,\gamma}^{r_B}
   \left(\frac{c_{H,\gamma}}{z_\gamma}\right)^{24Bn}
   \frac{\mathcal Z^+_{\gamma,B,n}}{\mathcal S_{\gamma,B}(I)}.
 \quad}
 \label{f16v40:exact-scalar}
\end{equation}
For even $n$, this equals
$\mathcal A^{00}_{\gamma,B,n}\theta_{\gamma,n}$. Consequently
\begin{equation}
 \theta_{\gamma,n}
   =c_{H,\gamma}^{r_B}\frac{\mathcal Z^+_{\gamma,B,n}}{\mathcal Z_{\gamma,B,n}}.
 \label{f16v40:theta-exact}
\end{equation}
The same scalar ratio defines the corresponding spectral projection weight at
other integer durations. This is an identity for the original transfer, not a
replacement endpoint ensemble.
\end{proposition}
\begin{proof}
In the root average, the right endpoint distribution is exactly the dressed
vector at $Z(h)$ by equivariance. V37's cancellation of $p_{\gamma,B}$ and
$c_\Omega^2$ therefore applies, and both static endpoint denominators equal
$\mathcal S_{\gamma,B}(I)$. The kernel in \eqref{f16v40:one-average} uses Haar
endpoint coordinates. Multiplication by $j_{\rm br,0}$ gives the two prescribed
zero-chart half-densities. If the nonzero final endpoint is temporarily written
in its bridge chart, the conversion back divides by its square-root Jacobian.
It cancels precisely the final product
$\prod_e j(y_{n,e}/\sqrt\gamma)^{1/2}$ in V4's chart formula. At cut values one can instead use the Haar kernel directly; the same
identity holds without a final bridge logarithm.

The initial endpoint has unit nonconstant Jacobian. All interior bridge Haar
factors and the fast endpoint half powers remain those in
\eqref{f16v40:root-mass}. The $r_B$ based root integrations contribute
$c_{H,\gamma}^{r_B}$ and the displayed full $j(\xi_b/\sqrt\gamma)$ factors.
This proves \eqref{f16v40:exact-scalar}, including the unchanged $24Bn$ kinetic
normalizers. Comparing with V37's unprojected formula and V38's loaded-vector
identity proves \eqref{f16v40:theta-exact}. No intermediate carrier projection
is introduced. Positivity permits every rearrangement of these integrals.
\end{proof}

\subsection{The Gaussian root integral is an exact temporal affine shift}

Let $\mathsf d_{\rm rt}:\mathbb R^{k_B}\to\mathbb R^{36B}$ be the based
root-gradient injection
\[
 (\mathsf d_{\rm rt}\xi)_e=\xi_{s(e)}-\xi_{t(e)},\qquad \xi_o=0.
\]
All matrices include colour. The inherited static classification gives
$D\mathsf d_{\rm rt}=0$, and connectedness makes $\mathsf d_{\rm rt}$ injective.
Define the positive matrix
\begin{equation}
 \mathsf J_{\rm rt}=\mathsf d_{\rm rt}^{\mathsf T}G^{-1}\mathsf d_{\rm rt}
 \succ0,
 \qquad
 \kappa_B=\frac{(2\pi)^{k_B/2}}
                  {(2\pi^2)^{r_B}\sqrt{\det\mathsf J_{\rm rt}}}.
 \label{f16v40:root-metric}
\end{equation}
This is the original kinetic $G$, with the nonroot ports integrated, not an
identity metric chosen on the root graph. The determinant includes all colours.
Let $U_0^+$ be the quadratic limit of $U^+_\gamma$ and
$\mathcal Z^+_{0,B,n}=\int e^{-U_0^+}\,dw\,d\xi$.

\begin{theorem}[Exact Gaussian orbit mass]
\label{f16v40:Gaussian}
At every fixed finite $B,n$,
\begin{equation}
 \boxed{\quad
 \frac{\mathcal Z^+_{0,B,n}}{\mathcal Z_{0,B,n}}
       =\frac{(2\pi n)^{k_B/2}}{\sqrt{\det\mathsf J_{\rm rt}}}.
 \quad}
 \label{f16v40:Gaussian-ratio}
\end{equation}
After fast Gaussian integration, the variables $\xi$ and
$\widetilde y_i=y_i-(i/n)\mathsf d_{\rm rt}\xi$, $1\le i<n$, are independent;
the latter have the zero-endpoint retained Gaussian and
$\xi\sim N(0,n\mathsf J_{\rm rt}^{-1})$. This is a statement about the exact
quadratic integral. It is not a factorization of nonlinear native orbit fibres.
\end{theorem}
\begin{proof}
At quadratic order the final endpoint is
$y_n=\mathsf d_{\rm rt}\xi$. The conjugation of an endpoint chord by $h_b$ has
no first-order change at the identity and does not change its original endpoint
quadratic. V33's complete action formula applies with these free final data.
The magnetic terms and their integrated fast return depend on $Dy_i$ only;
they are unchanged by the displayed shift because $D\mathsf d_{\rm rt}=0$.
The kinetic cross term vanishes exactly:
\[
 \sum_{t=0}^{n-1}
 (\widetilde y_t-\widetilde y_{t+1})^{\mathsf T}
                          G^{-1}\mathsf d_{\rm rt}\xi=0,
 \qquad \widetilde y_0=\widetilde y_n=0.
\]
Thus the retained action splits as
\begin{equation}
 \mathcal F_0(0,y_1,\ldots,y_{n-1},\mathsf d_{\rm rt}\xi)
 =\mathcal F_0(0,\widetilde y_1,\ldots,\widetilde y_{n-1},0)
                  +\frac1{2n}\xi^{\mathsf T}\mathsf J_{\rm rt}\xi.
 \label{f16v40:affine-split}
\end{equation}
The fast Gaussian determinant is independent of the retained history, including
its endpoint, and is the same determinant in the two integrals. The triangular
change $(y,\xi)\mapsto(\widetilde y,\xi)$ has determinant one. Integrating the
last Gaussian factor proves \eqref{f16v40:Gaussian-ratio} with its exact
normalization and the independence statement. No simultaneous diagonalization
of $G$ with the magnetic matrices is needed.
\end{proof}

In particular, set
\begin{equation}
 M_{0,n}^{\rm G}:=
 \left(\frac n2\right)^{\ell_B}
              \frac{\mathcal A^{00}_{0,B,n}}{\mathcal A^{00}_{0,B,2}}.
 \label{f16v40:Gaussian-moment}
\end{equation}
V38's full Gaussian endpoint determinant, not a gauge-dimension guess, now gives
\begin{equation}
 M_{0,n}^{\rm G}
 =\widetilde c_B\Lambda_{{\rm G},B}^{n-2}n^{-9/2}
  \{1+\widetilde d_{B,1}/n+\widetilde d_{B,2}/n^2+O_B(n^{-3})\},
 \qquad \widetilde c_B>0.
 \label{f16v40:Gaussian-edge}
\end{equation}
The residual simultaneous conjugation has not been divided out as an additional
three-dimensional translation. The remaining nine free Gaussian directions
transform under it, while the endpoint source is already invariant.

\subsection{Compact domination for the root-averaged integral}

The new integral must be estimated in its own right. A small absolute error
for an unprojected law cannot be divided by $\gamma^{-\ell_B}$.

\begin{proposition}[One anchored endpoint confines the based root variables]
\label{f16v40:pin}
For each fixed $B$ there are $c_B^+>0$, $H_B^+<\infty$, independent of $n,\gamma$,
such that on the complete root-extended principal product, with $u=(w,\xi)$,
\begin{equation}
 U^+_\gamma(u),\ U^+_0(u)\ge\frac{c_B^+}{n^2}|u|^2,
 \qquad 2c_B^+n^{-2}I\preceq\mathcal H^+\preceq H_B^+I,
 \label{f16v40:pin-bound}
\end{equation}
where $\mathcal H^+=D^2U_0^+$. The identity is the unique zero in these based
coordinates. For $\mathcal C^+=(\mathcal H^+)^{-1}$,
\begin{equation}
 \|\mathcal C^+\|\le C_Bn^2,
 \qquad \sup_g\|\mathcal C^+_{gg}\|\le C_Bn,
 \qquad \mathcal C^+_{\zeta\zeta}\preceq(2a_*)^{-1}I.
 \label{f16v40:Gaussian-scales}
\end{equation}
The subscript $g$ in the middle bound denotes an original three-group, including
one of the new based root coordinates, not a random gauge transformation.
\end{proposition}
\begin{proof}
V5's exact independent anchor gives $U^+_\gamma\ge a_*|\zeta|^2$ at every
retained input. V34's metric-increment estimate used neither final anchoring
nor a retained window, so it still gives
\[
 \gamma\sum_{t,e}\operatorname{dist}(Z_{t,e},Z_{t+1,e})^2
                                      \le C_\nabla U^+_\gamma.
\]
Starting only from $Z_0=I$, telescoping and Cauchy--Schwarz bound the sum of all
interior squared distances by $n^2$ times the sum of squared increments, and
bound the final squared distances by $n$ times that sum. Along a coarse spanning
tree, the ordered product of the final values \eqref{f16v40:pure-endpoint} from
$o$ to $b$ is $h_b^{-1}$. The triangle inequality consequently gives
\[
 \sum_{b\ne o}\operatorname{dist}(I,h_b)^2
 \le(B-1)^2\sum_{e\in\mathcal F}\operatorname{dist}(I,Z_{n,e})^2.
\]
Principal coordinates satisfy $|\xi_b|^2=\gamma\operatorname{dist}(I,h_b)^2$.
Combining these inequalities proves the native lower bound, for example with a
constant no larger than
$[a_*^{-1}+C_\nabla(1+(B-1)^2)]^{-1}$. This constant is independent of duration
but not spatial width. Zero action forces all fast anchors and increments to
vanish; then $Z_i=I$ and every $h_b=I$ by the based tree. No additional saddle
or global-conjugation volume has been deleted.

The fixed-vector quadratic limit gives the same lower bound for $U_0^+$.
Original finite word carriers and the extra endpoint products give the upper
Hessian bound. Minimizing over $(y,\xi)$ retains the fast anchor and proves the
last bound in \eqref{f16v40:Gaussian-scales}. For the retained blocks use
\eqref{f16v40:affine-split}: their covariance is the original zero-endpoint
covariance, with local size $O(n)$, plus
$(i/n)^2n\mathsf d_{\rm rt}\mathsf J_{\rm rt}^{-1}\mathsf d_{\rm rt}^{\mathsf T}$.
The root covariance is $n\mathsf J_{\rm rt}^{-1}$. At fixed $B$ these give the
middle bound. This is a marginal calculation, not a conditional covariance
substituted for the full one.
\end{proof}

\begin{theorem}[Normalized root-extended comparison and absolute mass]
\label{f16v40:mass}
Fix $B$ and a tail power $s\ge2$. There are constants depending on $B,s$ but
not on $n,\gamma$ such that the following statements hold under the entry
conditions specified in the proof and $\rho_+=d_+n^4/\gamma\le1$:
\begin{equation}
 \left|\log\frac{\mathcal Z^+_{\gamma,B,n}}{\mathcal Z^+_{0,B,n}}\right|
       \le C_B(\rho_++\gamma^{-s}),\qquad
 d_{\rm TV}(\mathbb P^+_\gamma,\mathbb Q^+)
       \le C_B(\sqrt{\rho_+}+\gamma^{-s}).
 \label{f16v40:mass-bounds}
\end{equation}
Here the laws normalize their complete respective integrals; all root groups
remain integrated. A sufficient asymptotic entry, at fixed $B$, is
\begin{equation}
 \frac{Bn^7\log(n+1)+n^6\log\gamma}{\gamma}\longrightarrow0.
 \label{f16v40:regime}
\end{equation}
This is not a theorem uniform over growing spatial widths.
\end{theorem}
\begin{proof}
We verify the new integral's hypotheses rather than applying V36 to a different
law by name. Shrink $c_B^+\le1/4$, enlarge $H_B^+\ge1$, and put
$a=c_B^+/n^2$, $t=s+4$. In the formula for V35's radius replace $(d,c,H_*)$ by
$(d_+,c_B^+,H_B^+)$:
\[
 R^2=\frac2a\left\{1+\frac{d_+}{2}\log\frac{2H_B^+}{a}
  +\log\left[1+\left(\frac{d_++6}{a}\right)^4\right]+(t+1)\log\gamma\right\}.
\]
The precise entry is $R\le\sqrt\gamma/2$ and
$C_{2,B}R/\sqrt\gamma+C_{H,B}/\gamma\le a$. The constants are the finite
ordered-word and Haar derivative bounds for this explicit endpoint
substitution. All new factors are $\operatorname{Exp}(\pm\xi_b/\sqrt\gamma)$;
there is no logarithm of their product to differentiate. Bounded incidence
gives on this ball, for $W^+=U^+_\gamma-\log\mathcal J^+_\gamma-U_0^+$,
\[
 \|D^2W^+\|\le C_{2,B}R/\sqrt\gamma+C_{H,B}/\gamma,\qquad
 |\nabla W^+|^2\le\frac{C_B}{\gamma}\sum_g|u_g|^4
                       +\frac{C_B}{\gamma^2}|u|^2.
\]
Thus $D^2(U_0^++\vartheta W^+)\succeq\mathcal H^+/2$ for $0\le\vartheta\le1$.
All these laws and all linear tilts are on one unchanged convex ball. Their
untilted means vanish by simultaneous colour conjugation, which preserves the
based condition $h_o=I$. The convex-boundary variance identity proved in V35
yields directional moments in metric $\mathcal C^+$ and therefore local fourth
moments $O_B(n^2)$, using \eqref{f16v40:Gaussian-scales}. It follows that
\[
 \mathbb E_{\vartheta}
 [\nabla W^+{}^{\mathsf T}\mathcal C^+\nabla W^+]
             \le C_B\rho_+,\qquad
 \operatorname{Var}_{\vartheta}W^+\le C_B\rho_+.
\]

The compact lower bound \eqref{f16v40:pin-bound}, $0\le\mathcal J^+_\gamma\le1$,
and the Gaussian of precision $2\mathcal H^+$ give
$\mathcal Z^+_\gamma\ge2^{-d_+/2-1}\mathcal Z^+_0$ by integrating over the
ball. The same radial eighth-moment calculation as V35, with the displayed
parameters, bounds each complete weighted tail by $\gamma^{-t}$. This includes
the entire root-chart complement. No exterior good event is assumed.

For the scalar, the classical cubic $P_3^+$ is odd under simultaneous Euclidean
reflection. Ordered fourth derivatives and the Haar quadratic give
\[
 |W^+-\gamma^{-1/2}P_3^+|
       \le\frac{C_B}{\gamma}\sum_g(|u_g|^4+|u_g|^2).
\]
The centered Gaussian on the ball has zero cubic mean. Hence
$|\mathbb E_0W^+|\le C_Bd_+n^2/\gamma\le C_B\rho_+$. Differentiating the
\emph{normalized} interpolating integrals gives the exact identity
\[
 \log(Z_{1,\mathbb B}^+/Z_{0,\mathbb B}^+)
 =-\mathbb E_0W^+
   +\int_0^1(1-\vartheta)\operatorname{Var}_{\vartheta}W^+\,d\vartheta.
\]
The two complete tail probabilities restore the full scalar by adding
$\log\mathbb Q^+(\mathbb B)-\log\mathbb P^+_\gamma(\mathbb B)$. This proves
the first bound, including its actual normalizers. The normalized
square-root-density argument of V35 bounds total variation on the ball by
$C_B\sqrt{\rho_+}$; adding the two paid tails proves the other bound.

Finally $R\le C_{B,s}n\sqrt{d_+\log(n+1)+\log\gamma}$. A sufficiently small
constant in
$n^6[d_+\log(n+1)+(s+5)\log\gamma]\le c_{B,s}\gamma$
ensures the stated finite entry. At fixed $B$, \eqref{f16v40:regime} implies
this and $\rho_+\to0$. This verifies the parameter range without importing
an unproved estimate for a root-twisted boundary row.
\end{proof}

\subsection{Relative projected normalization, without dividing an absolute error}

\begin{theorem}[The leading native root weight and projected moments]
\label{f16v40:projection}
At fixed $B$, at the preceding entry and the original joint entry, for even
$n\ge2$,
\begin{equation}
 \boxed{\quad
 \left|\log\frac{\theta_{\gamma,n}}
                  {\kappa_B(n/\gamma)^{\ell_B}}\right|
       \le C_B(\rho_++\gamma^{-s}).
 \quad}
 \label{f16v40:theta-asymptotic}
\end{equation}
For V39's \emph{same} normalized projected loading $f_\gamma^{\rm G}$, its
moments at all integer $n\ge2$ satisfy
\begin{equation}
 \boxed{\quad
 \left|\log\frac{M_{\gamma,n}^{\rm G}}{M_{0,n}^{\rm G}}\right|
        \le C_B(\rho_++\gamma^{-s}).
 \quad}
 \label{f16v40:moment-asymptotic}
\end{equation}
In particular the leading duration dependence of $\theta$ is now identified,
and $M_{\gamma,n}^{\rm G}$ has the $n^{-9/2}$ leading factor in
\eqref{f16v40:Gaussian-edge} when \eqref{f16v40:regime} holds and $n\to\infty$.
Neither assertion is obtained by dividing a total-variation error by $\theta$.
\end{theorem}
\begin{proof}
Insert \eqref{f16v40:Gaussian-ratio} into \eqref{f16v40:theta-exact} and use
$c_{H,\gamma}=(2\pi^2\gamma^{3/2})^{-1}$. The remaining logarithm is exactly
\[
 \log(\mathcal Z^+_\gamma/\mathcal Z^+_0)
                       -\log(\mathcal Z_\gamma/\mathcal Z_0).
\]
The new theorem pays the first term and V36 pays the second; its original
$\rho\le\rho_+$. This proves \eqref{f16v40:theta-asymptotic}.

More economically, the projected moment can be compared without either static
endpoint approximation or the unprojected joint-mass expansion. The exact
formula \eqref{f16v40:exact-scalar} at $n$ and two gives
\[
 M_{\gamma,n}^{\rm G}
  =\left(\frac{c_{H,\gamma}}{z_\gamma}\right)^{24B(n-2)}
                \frac{\mathcal Z^+_{\gamma,B,n}}{\mathcal Z^+_{\gamma,B,2}}.
\]
Both $c_{H,\gamma}^{r_B}$ and $\mathcal S_{\gamma,B}(I)$ cancel exactly. Its
Gaussian counterpart is \eqref{f16v40:Gaussian-moment}. Their logarithmic ratio
is the difference of the two new mass errors minus
$24B(n-2)\log\mathfrak b_\gamma$, with precisely V37's scalar
$\mathfrak b_\gamma$. Since $|\log\mathfrak b_\gamma|\le C/\gamma$ eventually,
this last cost is $O(Bn/\gamma)$ and is bounded by $C_B\rho_+$. All three
errors are paid. This proves \eqref{f16v40:moment-asymptotic}, including the
loading normalization, with no assumed cancellation of approximation errors.
\end{proof}

The result identifies an orbit-integrated weight, not the pointwise function
$k(\eta)$ of V39. The latter need not become constant. It also leaves the
original midpoint probability and the square of its projected amplitude distinct.
For an invariant bounded midpoint multiplication $F$ and even $n$, however,
commutation gives the exact useful identity
\[
 \frac{\langle\Pi s_\gamma,T_\gamma^{n/2}F T_\gamma^{n/2}\Pi s_\gamma\rangle}
      {\langle\Pi s_\gamma,T_\gamma^n\Pi s_\gamma\rangle}
 =\frac{\int dh\,\langle s_\gamma,T_\gamma^{n/2}F
                       T_\gamma^{n/2}\mathsf U_hs_\gamma\rangle}
        {\int dh\,\langle s_\gamma,T_\gamma^n\mathsf U_hs_\gamma\rangle}.
\]
Thus the added positive integral retains the correct invariant-observable
weight; it is not probability averaging substituted for amplitude projection.

\subsection{The native root-invariant Gamma edge and the longer non-preparation range}

Let $\mu_\gamma^{\rm G}$ be the probability spectral measure of
$f_\gamma^{\rm G}$ for the unchanged root-invariant restriction of $T_\gamma$.
Define
\[
 d\mu_{\gamma,n}^{\rm G}(\lambda)
  =\frac{\lambda^{n-2}\,d\mu_\gamma^{\rm G}(\lambda)}{M_{\gamma,n}^{\rm G}}.
\]
Its top atom is at the same native $\Lambda_\gamma$ as in V38. The reference
$\Lambda_{{\rm G},B}$ is not substituted for it.

\begin{theorem}[Projected endpoint edge law at the original polynomial durations]
\label{f16v40:edge}
Fix $B$. Let even $n=n_\gamma\to\infty$ satisfy \eqref{f16v40:regime}, using
fixed tail powers large enough for the adjacent-ratio tests. Under the
\emph{native root-invariant} probability $\mu_{\gamma,n_\gamma}^{\rm G}$,
\begin{equation}
 -n_\gamma\log(\lambda/\Lambda_{{\rm G},B})
     \ \Rightarrow\ \operatorname{Gamma}(9/2,1),\qquad
 n_\gamma(1-\lambda/\Lambda_{{\rm G},B})
     \ \Rightarrow\ \operatorname{Gamma}(9/2,1).
 \label{f16v40:Gamma}
\end{equation}
The second variable has mean and variance both tending to $9/2$. The Gamma
parameter convention is shape, then rate. For every $0<\alpha<1/7$ the
admissible example $n_\gamma=2\lfloor\gamma^\alpha/2\rfloor$ gives
\begin{equation}
 \begin{split}
 |\langle\Omega_\gamma,v^{\rm G}_{\gamma,n_\gamma}\rangle|^2&\longrightarrow0,\\
 d_{\rm TV}(|v^{\rm G}_{\gamma,n_\gamma}|^2dH,\Omega_\gamma^2dH)
                                                   &\longrightarrow1.
 \end{split}
 \label{f16v40:nonpreparation}
\end{equation}
Both densities in the second line are invariant. This extends V39's smaller
range $\alpha<3/[7(B+2)]$ at each fixed $B$, not uniformly in growing $B$.
\end{theorem}
\begin{proof}
The spectral identity and \eqref{f16v40:moment-asymptotic} imply, for every
fixed integer $j\ge0$,
\[
 \mathbb E_{\mu_{\gamma,n}^{\rm G}}
       (\lambda/\Lambda_{{\rm G},B})^{jn}
 =\Lambda_{{\rm G},B}^{-jn}
         \frac{M_{\gamma,(j+1)n}^{\rm G}}{M_{\gamma,n}^{\rm G}}
                  \longrightarrow(j+1)^{-9/2}.
\]
Every fixed multiple of the duration satisfies the entry eventually. The
native scalar errors at both times tend to zero; they are not differentiated
or assumed to cancel. As in V38's proved moment argument, the first moment
gives tightness of $Y=(\lambda/\Lambda_{{\rm G},B})^n$, and the next moment
provides uniform integrability of each preceding one. The limiting moments
force support in $[0,1]$ and determine the law there by polynomial density.
They are the moments of $e^{-G_*}$ for $G_*\sim\operatorname{Gamma}(9/2,1)$.
There is no atom at zero, so taking $-\log Y$ and then the locally uniform map
$n(1-e^{-x/n})$ proves \eqref{f16v40:Gamma}. Negative edge deficits at finite
coupling are permitted; this argument estimates spectral weights, not absence
of eigenvalues above the Gaussian threshold.

For the claimed moments, the determinant expansion
\eqref{f16v40:Gaussian-edge} gives adjacent-ratio errors $O_B(n^{-2})$ and
second-difference errors $O_B(n^{-3})$, with leading constants $9/2$. Adding
the new native logarithmic errors costs $O_B(d_+n^4/\gamma+\gamma^{-s})$.
The regime gives $n^2(d_+n^4/\gamma+\gamma^{-s})\to0$ when, for example,
$s\ge2$. The exact mean and variance formulas used in V38 therefore prove
mean and variance convergence for the linear deficit. This step is separate
from weak convergence.

A tight sequence converging to a nonatomic law has maximal atomic weight
tending to zero, by the compact-subsequence and closed-interval argument in
V38. The Perron eigenvalue of the actual restricted transfer is simple, so its
atom is exactly the squared overlap in \eqref{f16v40:nonpreparation}. At even
durations the filtered projected vector is positive. Its probability affinity
with $\Omega_\gamma^2$ is that overlap before squaring, and
$d_{\rm TV}(p,q)\ge1-\int\sqrt{pq}$ proves the second line. The maximizing
sign event is invariant because both densities are invariant. No spectral
convergence of operators or a native upper-top estimate is inferred.
\end{proof}

For completeness the corresponding projected scalar energy is
\[
 -\log\frac{M_{\gamma,n+1}^{\rm G}}{M_{\gamma,n}^{\rm G}}
   =-\log\Lambda_{{\rm G},B}+\frac{9}{2n}+o_B(n^{-1})
\]
along the same regime. This is measured against the Gaussian threshold, not an
identified physical ground energy. The function $k(\eta)$ remains relevant to
other quotient observables, but a direct positive root integration has paid
its total contribution to these moments without assuming its constancy.

\subsection{Checkpoint: the compact soft sector is now the outstanding scale}

V40 obtains a projected Gamma law by integrating the based group before the
limit. It does not subtract gauge dimensions from V38's exponent without a
calculation: \eqref{f16v40:Gaussian-ratio} supplies the exact determinant and
the native mass theorem supplies its relative error, including the root Haar
volume. The root variables are not assigned an independent Gaussian law in the
native model. The approximate Gaussian separation is used only after complete
compact domination has been proved for the new integral.

The remaining nine-dimensional entrance sector is still nonlinear and compact
at longer scales. The theorem leaves arbitrary durations, growing spatial
width, quantitative vacuum-preparation times, the first excited spectral ratio,
and independently calibrated physical time unpaid. A useful next target is
an exact treatment of that compact period sector together with its integrated
complement, or a same-top spectral comparison valid beyond this entrance.
Another unprojected or projected finite-order scalar coefficient alone would
not supply it. The supplied memoranda's distinction between power gain and
locality/normalization cost, and between auxiliary estimates and physical
transfer, remains binding. No Yang--Mills continuum mass-gap claim is made.

\subsection{Verification and preservation}

The diagnostic
\begin{center}\small\path{verify_ym_f16_v40_root_integral.py}\end{center}
checks finite noncommutative endpoint identities, all root/Haar powers,
noncommuting Gaussian affine shifts and determinants, normalization before
projection, complete-tail and parameter arithmetic, and the spectral moments.
It distinguishes literal periodic incidence checks from finite reference
fixtures. These checks are not native integral enclosures or formal proof
verification. The audit records actual executions and review depth, separately
from inherited reports.

Removing this insertion and reversing the title version recovers V39 byte for
byte. The analytic inputs are the compact anchor, ordered-word bounds, finite
comparison mechanism, and V38's Gaussian asymptotic, with the new hypotheses
checked above. The full lineage is not independently recertified. The next
checkpoint is V45.

\begin{firewall}
V40 proves the exact positive root-averaged endpoint representation, an explicit
Gaussian root determinant, and a relative compact-to-Gaussian mass estimate at
fixed spatial width. It identifies the leading native projected norm and the
root-invariant endpoint Gamma edge of shape $9/2$, with vanishing projected
Perron weight throughout the stated polynomial entrance. It does not prove a
full physical density of states, a gap or gaplessness of the calibrated
Hamiltonian, stationary vacuum convergence, unrestricted regulators, or an
interacting continuum construction. All original transfer powers, endpoint
carriers, root Haar factors, and physical normalizations are retained.
\end{firewall}
% END F16 V40 APPEND

% BEGIN F16 V41 APPEND -- 2026-09-18
% Insert before the final document-ending command of the supplied cumulative V40.
% Predecessor assertions, endpoint measures, normalizations, and clocks are preserved.

\clearpage
\section{V41: the native spectral top and central-flat vacuum localization}
\label{f16v41:context}

V40 controls a projected endpoint spectral measure, but does not supply the
matching upper bound on the native spectral top. Extending that moment
calculation to longer durations would require new estimates. We instead return
to the original one-step Wilson operator on the full fine-link space. Its
positive kernel and its sum-of-squares magnetic action permit a local operator
comparison near every flat connection, including singular flat strata. The
comparison never inverts a small magnetic eigenvalue uniformly over the flat
set. A positive winding expansion then identifies which flat backgrounds have
the largest harmonic transfer norm.

The result is a fixed-spatial-regulator theorem: the actual physical Perron
value converges to V38's Gaussian threshold, and its stationary probability
concentrates on the eight central-flat gauge orbits. This is not an estimate of
the first excited ratio. No rate sufficient to center V40's spectral variable
at the native top is asserted. The localization lengths and constants below
are allowed to depend on the spatial regulator.

\subsection{Sources, exact operator, and the change of viewpoint}

The supplied V40 appendix, analytic audit, and checkpoint were read completely.
The f15 V91 fine-link transfer formula and the f14 discussion of flat backgrounds
were checked at identified passages. The f14 flat-holonomy cancellation concerns
a different quotient determinant and must not be identified with the present
transfer threshold. Searches of the current lineage and supplied archives are
recorded in the audit; these are targeted checks, not an exhaustive novelty or
proof certificate. The next scheduled companion checkpoint remains V45.

Harmonic localization is a classical method. Klein--Rosenberger,
\href{https://arxiv.org/abs/1706.06357}{\emph{Harmonic Approximation of Difference
Operators}}, supplies related context, but its lattice Schr\"odinger hypotheses
are not substituted for a positive compact integral operator with a singular
zero set. We prove the required norm statement directly. For the Bessel
identities in the winding calculation we use the primary formulas
\href{https://dlmf.nist.gov/10.32.E1}{DLMF 10.32.1} and
\href{https://dlmf.nist.gov/10.35.E2}{DLMF 10.35.2}; their relevant integral and
series manipulations are also justified below. No external mass-gap theorem is
an input.

Put $L=2m$, $m\ge4$, $B=m^3$, $V=L^3=8B$ and $E=3V=24B$.
On $\mathcal X_L=\SU(2)^E$, with probability product Haar $dH$, retain
\begin{equation}
 \begin{split}
 S(U)&=\sum_{p\ {\rm spatial}}\left(1-\Re U_p\right),\qquad
 b_\gamma(U)=e^{-\gamma S(U)/2},\\
 \overline T_\gamma&=M_{b_\gamma}C_\gamma^{\otimes E}M_{b_\gamma},\qquad
 C_\gamma(U,U')=z_\gamma^{-1}e^{-\gamma(1-\Re UU'^{-1})}.
 \end{split}
 \label{f16v41:native}
\end{equation}
Here $\Re$ is the scalar quaternion part, so it is the original normalized
fundamental trace. Every original spatial plaquette and kinetic scalar is
included. At finite $L,\gamma$, the operator is compact, positive self-adjoint,
and positivity improving. Its unit Perron vector $\Omega_\gamma>0$ is gauge
invariant by uniqueness. Consequently its norm is the physical top. The exact
rooted-tree reduction in f15 V98/V1 identifies that physical operator with the
root-invariant restriction used in V38--V40. Hence
\[
 \Lambda_\gamma=\|\overline T_\gamma\|
   =\|T_\gamma|_{\operatorname{Ran}\Pi_{\rm G}}\|.
\]
We use the full-link space to estimate this same top, not to replace it by a
row-normalized Markov kernel or a carrier compression. The stationary vacuum
probability in these coordinates is $\Omega_\gamma^2dH$.

\subsection{A narrow-kernel localization identity with its error paid}

Let $\varepsilon=\gamma^{-1/2}$, and use the round product metric for which a
unit imaginary quaternion has norm one. Write $D_L=3E$ for the dimension of
$\mathcal X_L$. Near a point $p$, choose local coordinates whose derivative at
$p$ is an orthonormal identification with $\mathbb R^{D_L}$. In the flat
half-density representation, the normalized kinetic kernel has the following
local properties. Given $\delta>0$, after shrinking the coordinate neighborhood,
for all $x,y$ in it and sufficiently small $\varepsilon$,
\begin{equation}
 c_\gamma^{\rm flat}(x,y)
 \le(1+\delta)(2\pi\varepsilon^2)^{-D_L/2}
       \exp\left[-\frac{(1-\delta)|x-y|^2}{2\varepsilon^2}\right].
 \label{f16v41:local-kernel}
\end{equation}
The corresponding rescaled kernel at the diagonal converges to the normalized
Euclidean Gaussian. These assertions include the local measure density and the
actual $z_\gamma^E$: V37's one-group formula gives
$z_\gamma=(2\pi^2)^{-1}(2\pi\varepsilon^2)^{3/2}(1+O(\varepsilon^2))$.
Also $\sum_e(1-\Re U_eU_e'^{-1})=\frac12\sum_e|U_e-U_e'|^2$ in quaternion
coordinates. Its pullback, and the smooth product volume, give
\eqref{f16v41:local-kernel} by a local $C^1$ near-isometry. The chart can be
chosen with a convex coordinate image. There is no undetermined leading
normalization in this estimate.

\begin{proposition}[Finite partition localization for the actual transfer]
\label{f16v41:IMS}
Let $\chi_0,\ldots,\chi_N$ be real smooth functions on $\mathcal X_L$ with
$\sum_j\chi_j^2=1$. Then
\begin{equation}
 \left\|\overline T_\gamma-
       \sum_jM_{\chi_j}\overline T_\gamma M_{\chi_j}\right\|
       \le C_{L,\chi}\gamma^{-1}.
 \label{f16v41:IMS-bound}
\end{equation}
The partition and $L$ are fixed before $\gamma\to\infty$.
\end{proposition}
\begin{proof}
The difference kernel is exactly
\[
 \frac12\sum_j(\chi_j(U)-\chi_j(U'))^2
 b_\gamma(U)c_\gamma(U,U')b_\gamma(U').
\]
It has nonnegative entries; operator positivity of this difference is not
asserted. The squared differences are bounded by $C_\chi\operatorname{dist}(U,U')^2$.
Since $b_\gamma\le1$, the Schur bound is at most $C_\chi$ times the second
moment of the product Wilson displacement. For one group,
$1-\cos\theta\ge2\theta^2/\pi^2$ on $[0,\pi]$, its normalization is at least
$c\gamma^{-3/2}$, and its weighted radial numerator is at most
$C\gamma^{-5/2}$. Thus that moment is $O(\gamma^{-1})$. Summing the $E$ group
moments proves \eqref{f16v41:IMS-bound}, with its fixed-regulator cost.
\end{proof}

\subsection{Normal coordinates without a regular-flat-stratum assumption}

Let $\mathcal F_L=\{U:S(U)=0\}$. It is a compact set, not presumed smooth.
In the ambient quaternion space define the smooth constraint map
\[
 \mathcal W(U)=(U_p-I)_{p\ {\rm spatial}},\qquad
 S(U)=\tfrac12|\mathcal W(U)|^2.
\]
For $p\in\mathcal F_L$, let $H_p=(D\mathcal W_p)^*D\mathcal W_p$ on its
orthonormal tangent space. Define the continuous harmonic threshold
\begin{equation}
 \mathfrak h(p)=
 \prod_{\nu\in\operatorname{spec}(H_p),\ \nu>0}
 \left(1+\frac\nu2+\sqrt{\nu+\nu^2/4}\right)^{-1/2}.
 \label{f16v41:harmonic-threshold}
\end{equation}
A zero eigenvalue contributes one. Thus the formula is continuous also where
$\operatorname{rank}H_p$ changes. Its positive factors are V38's normalized
Mehler factors, not an exponentiated Schr\"odinger frequency substituted for the
discrete transfer.

\begin{proposition}[Exact magnetic lower bound in a pointwise normal chart]
\label{f16v41:normal-chart}
For every $p\in\mathcal F_L$ there is a smooth local coordinate system
$x=(x_+,x_0)$ centered at $p$, orthonormal to first order, such that
\begin{equation}
 S(U(x))\ge\tfrac12\sum_{i=1}^{r}\nu_i x_{+,i}^2,
 \qquad r=\operatorname{rank}H_p,
 \label{f16v41:normal-lower}
\end{equation}
where $\nu_1,\ldots,\nu_r$ are the positive eigenvalues of $H_p$.
This is an exact inequality on the neighborhood, with no negative remainder.
\end{proposition}
\begin{proof}
Write the nonzero singular decomposition of $D\mathcal W_p$ as
$D\mathcal W_p v_i=\sqrt{\nu_i}e_i$, with both systems orthonormal. Set
$x_{+,i}=\langle e_i,\mathcal W(U)\rangle/\sqrt{\nu_i}$, and take as $x_0$
the components of an initial logarithmic chart along $\ker D\mathcal W_p$.
The combined derivative at $p$ is an orthogonal isomorphism. The finite
inverse function theorem gives a local coordinate system. Orthogonal
projection of $\mathcal W(U)$ onto the span of the $e_i$ proves
\eqref{f16v41:normal-lower}. If $r=0$, use the original orthonormal chart and
$S\ge0$. No inverse of a positive eigenvalue uniform over $p$ is required;
only a chart at the individual point is constructed.
\end{proof}

\begin{theorem}[Fixed-regulator harmonic principle for the native top]
\label{f16v41:harmonic-principle}
For the original operator \eqref{f16v41:native},
\begin{equation}
 \boxed{\qquad
 \lim_{\gamma\to\infty}\Lambda_\gamma
       =\max_{p\in\mathcal F_L}\mathfrak h(p).
 \qquad}
 \label{f16v41:top-principle}
\end{equation}
Let $\mathcal F_L^*$ be the maximizing set. For every open neighborhood
$O\supset\mathcal F_L^*$,
\begin{equation}
 \int_{\mathcal X_L\setminus O}\Omega_\gamma(U)^2\,dH(U)\longrightarrow0.
 \label{f16v41:abstract-localization}
\end{equation}
No positive Hessian floor in the null directions, Morse--Bott hypothesis, or
uniform rank on $\mathcal F_L$ is assumed.
\end{theorem}
\begin{proof}
For the upper bound fix a small $\delta>0$. At each flat $p$, use the preceding
chart, shrunk so \eqref{f16v41:local-kernel} holds. Kernelwise domination and
\eqref{f16v41:normal-lower} bound the localized operator norm by
\begin{equation}
 (1+\delta)(1-\delta)^{-D_L/2}
 \prod_{i=1}^{r}
 \left(1+\frac{\nu_i}{2(1-\delta)}+
 \sqrt{\frac{\nu_i}{1-\delta}+\frac{\nu_i^2}{4(1-\delta)^2}}
 \right)^{-1/2}.
 \label{f16v41:local-norm}
\end{equation}
Indeed this is the norm of the dominating whole-space Gaussian after scaling
$x$ by $\sqrt{1-\delta}/\varepsilon$. The zero directions give normalized heat
convolution of norm one; the raw kinetic prefactor gives
$(1-\delta)^{-D_L/2}$. For complex tests use their absolute values before this
positive-kernel comparison. No Loewner ordering is inferred from pointwise
ordering alone.

A finite subcover of the compact flat set suffices. Add a smooth partition
piece supported away from it; its restricted norm is at most $e^{-\gamma s_0}$
for some $s_0>0$, because both magnetic factors are bounded by
$e^{-\gamma s_0/2}$ and $C_\gamma^{\otimes E}$ has norm one. Take a smooth
quadratic partition subordinate to this cover. Proposition~\ref{f16v41:IMS}
then bounds the global Rayleigh quotient by the largest local norm plus
$O_{L,\chi}(\gamma^{-1})$. At fixed $L$ the positive eigenvalues are bounded,
and the expression \eqref{f16v41:local-norm} converges uniformly in those
values, including zero, to \eqref{f16v41:harmonic-threshold} as $\delta\downarrow0$.
First send $\gamma\to\infty$ for the chosen finite cover, then
$\delta\downarrow0$. This proves the upper bound without localizing every
configuration into a coupling-dependent tube around a singular stratum.

For the lower bound fix a flat $p$ and split an orthonormal logarithmic tangent
chart into the range and kernel of $H_p$, of dimensions $r$ and $k$.
Choose real smooth compactly supported unit functions $f$ on $\mathbb R^r$
and $g$ on $\mathbb R^k$. When $k=0$ omit the second factor. Put
$t_\varepsilon=\varepsilon^{3/4}$ and use the half-density trial function
\[
 \phi_\varepsilon(x_+,x_0)=
 \varepsilon^{-r/2}t_\varepsilon^{-k/2}
 f(x_+/\varepsilon)g(x_0/t_\varepsilon).
\]
It has norm one for small $\varepsilon$ and support inside the chart. Taylor's
formula for the constraint map, not a uniform quadratic lower bound, gives
\[
 \varepsilon^{-1}\mathcal W(U(\varepsilon u,t_\varepsilon z))
    =D\mathcal W_p(u,0)+O(t_\varepsilon^2/\varepsilon+\varepsilon)
    =D\mathcal W_p(u,0)+o(1)
\]
uniformly on its rescaled compact support. Thus both magnetic factors tend
to the required normal Gaussian multipliers. In the kinetic integral put
$x'=x+\varepsilon v$. The soft argument changes by
$(\varepsilon/t_\varepsilon)v_0\to0$, while the hard argument changes by
$v_+$. The actual normalized kernel and volume converge to the Euclidean
Gaussian displacement. On every bounded set of $v$ the convergence is
uniform; \eqref{f16v41:local-kernel} and the Schur estimate for the tail
$|v|>R$ bound the omitted bilinear form by $C e^{-cR^2}$, uniformly for small
$\varepsilon$. Integrate first at bounded $v$, then send $R\to\infty$.
The Rayleigh quotient tends to the normal Gaussian transfer quotient of $f$,
times $\|g\|_2^2=1$. Smooth compactly supported $f$ approximate its massive
Gaussian ground function. Its supremum is $\mathfrak h(p)$. This proves the
lower bound. The soft trial scale is a variational device, not an asserted
native fluctuation scale. The trial need not be gauge invariant: equality of
the full and physical tops at each coupling is used only for the largest
eigenvalue, not for excited spectral values.

For localization, choose $O_0$ with
$\mathcal F_L^*\subset O_0\subset\overline O_0\subset O$. Compactness and
continuity give a strict threshold loss $\eta>0$ on
$\mathcal F_L\setminus O_0$. In the preceding upper construction choose charts
centered in $\mathcal F_L\cap O_0$ to lie in $O$, and use charts with loss at
least $\eta$ for the remaining flat set. The nonflat partition piece is also
in the lower-norm family. The sum of squared lower-norm pieces is at least one
on $\mathcal X_L\setminus O$. With chart errors made less than an arbitrary
$\delta'>0$ and less than $\eta/2$, the Perron Rayleigh identity gives
\[
 \Lambda_\gamma\le h_*+\delta'
       -\tfrac\eta2\int_{O^c}\Omega_\gamma^2\,dH+o_\gamma(1),
 \qquad h_*:=\max\mathfrak h.
\]
Use the already proved top limit, then let $\delta'\downarrow0$. This proves
\eqref{f16v41:abstract-localization}. All limits in this proof fix $L$ first.
\end{proof}

\subsection{All compact flat backgrounds and their exact transverse thresholds}

The periodic lattice has no twisted boundary conditions. Flatness makes
parallel transport path independent on the universal cover. Its three
fundamental holonomies commute, and conversely a commuting triple specifies a
flat gauge orbit. For $\SU(2)$ write it, up to simultaneous conjugation, as
\[
 H_\mu=\operatorname{Exp}(\vartheta_\mu e_3),\qquad
 \varphi_\mu=2\vartheta_\mu\pmod {2\pi},\qquad \mu=1,2,3.
\]
The factor two is the adjoint charged character. Choosing commuting $L$th
roots gives a constant-link representative, or one may put all twists on
three seams. Gauge changes are isometries of the fine-link tangent metric.
Define
\begin{equation}
 \begin{split}
 \lambda_L(k,\varphi)&=4\sum_{\mu=1}^3
      \sin^2\!\left(\frac{2\pi k_\mu+\varphi_\mu}{2L}\right),
       \qquad k\in\{0,\ldots,L-1\}^3,\\
 \xi(t)&=\operatorname{arcosh}(1+t/2),\qquad
 F_L(\varphi)=\sum_k\xi(\lambda_L(k,\varphi)).
 \end{split}
 \label{f16v41:dispersion}
\end{equation}

\begin{proposition}[The flat-background transfer threshold with all colors]
\label{f16v41:flat-spectrum}
For a flat connection with adjoint twists $\varphi$,
\begin{equation}
 \mathfrak h(\varphi)=e^{-\mathcal E_L(\varphi)},\qquad
 \mathcal E_L(\varphi)=F_L(0)+2F_L(\varphi).
 \label{f16v41:flat-energy}
\end{equation}
At the trivial flat connection,
\begin{equation}
 e^{-3F_L(0)}=\Lambda_{{\rm G},B},
 \label{f16v41:reference-match}
\end{equation}
with exactly V38's Gaussian threshold and its original kinetic normalization.
\end{proposition}
\begin{proof}
At a flat point, differentiating the literal plaquette gives the covariant
lattice curl $d_{1,p}$; hence $H_p=d_{1,p}^*d_{1,p}$. Gauge the background
away on the universal cover. The neutral color has periodic boundary condition
and the two complex charged color lines have twists $\varphi$ and $-\varphi$.
In a Fourier mode the covariant derivative is
$d_\mu=e^{i(2\pi k_\mu+\varphi_\mu)/L}-1$. The curl symbol obeys
\[
 d_1^*d_1=|d|^2 I_3-dd^*,
\]
with eigenvalues $\lambda_L,\lambda_L,0$ if $|d|>0$, and three zeros
otherwise. Complexification preserves the trace of a real spectral function;
the two charged lines must both be counted. One half of the positive-mode sum
of $\xi$ is therefore $F_L(0)+F_L(\varphi)+F_L(-\varphi)$.
The change $k\mapsto-k$ shows $F_L(-\varphi)=F_L(\varphi)$, proving
\eqref{f16v41:flat-energy}. Gauge zero directions contribute a factor one to
the Gaussian norm, not a fictitious oscillator ground state.

To match the inherited rooted coordinates, let $R$ be the linear quotient map
from all fine-link tangent coordinates to V2's balanced chord/bridge slice
coordinates. Its kernel is the nonroot gauge-gradient space. The exact
Gaussian integration of those directions says its kinetic covariance is
$RR^*=\operatorname{diag}(\mathsf M_B,G)$; this is the original kinetic
completion, not identity on the quotient. The fine curl factors through $R$,
so $H_I=R^*\mathsf S_BR$ with the same magnetic matrix as V38. The nonzero
spectra of $R^*\mathsf S_BR$ and
$(RR^*)^{1/2}\mathsf S_B(RR^*)^{1/2}$ coincide, by the nonzero singular-value
identity. The deleted eigenvalues are zero and have norm factor one. Thus
V38's product of Mehler factors is precisely the fine-link product just
computed. This proves \eqref{f16v41:reference-match}, including normalization.
\end{proof}

\subsection{A positive winding expansion selects the central holonomies}

\begin{theorem}[Strict comparison of every flat threshold]
\label{f16v41:winding}
There are strictly positive coefficients $a_L(z)$, $z\in\mathbb Z^3\setminus\{0\}$,
with $\sum_{z\ne0}a_L(z)<\infty$, such that
\begin{equation}
 \boxed{\quad
 F_L(\varphi)-F_L(0)
   =\sum_{z\ne0}a_L(z)\{1-\cos(z\cdot\varphi)\},\qquad
 a_L(z)=L^3\int_0^\infty
   \frac{e^{-8t}I_0(2t)}{t}
       \prod_{\mu=1}^3 I_{Lz_\mu}(2t)\,dt.
 \quad}
 \label{f16v41:winding-formula}
\end{equation}
Here $I_{-j}=I_j$. Consequently $F_L(\varphi)\ge F_L(0)$, with equality
exactly at $\varphi\in(2\pi\mathbb Z)^3$. In particular, putting
$A_L=a_L(1,0,0)>0$,
\begin{equation}
 \mathcal E_L(\varphi)-\mathcal E_L(0)
       \ge4A_L\sum_{\mu=1}^3(1-\cos\varphi_\mu).
 \label{f16v41:strict-loss}
\end{equation}
The coefficients and this contrast are at fixed $L$, not uniform ultraviolet
or thermodynamic constants.
\end{theorem}
\begin{proof}
The integral for $I_0$ gives, for $s>0$,
\[
 \int_0^\infty e^{-(s+2)t}I_0(2t)\,dt
 =\frac1\pi\int_0^\pi\frac{d\theta}{s+2-2\cos\theta}
 =\frac1{\sqrt{s(s+4)}}=\xi'(s).
\]
The middle integral is elementary after the tangent half-angle substitution.
Integrating in $s$, using positivity and $\xi(0)=0$, gives
\begin{equation}
 \xi(s)=\int_0^\infty(1-e^{-st})\frac{e^{-2t}I_0(2t)}{t}\,dt.
 \label{f16v41:subordination}
\end{equation}
The integral is finite: at zero its integrand is bounded, and the elementary
bound $e^{-2t}I_0(2t)\le C(1+t)^{-1/2}$ follows from the same angle integral
and $1-\cos\theta\ge2\theta^2/\pi^2$.

The Bessel Fourier series, whose integer coefficients are positive for $t>0$,
gives the exact twisted heat trace
\[
 H_L(t,\varphi):=\sum_k e^{-t\lambda_L(k,\varphi)}
 =L^3e^{-6t}\sum_{z\in\mathbb Z^3}
       \left(\prod_\mu I_{Lz_\mu}(2t)\right)e^{iz\cdot\varphi}.
\]
The series is absolutely convergent at every positive $t$. Pairing $z$ and
$-z$ shows $H_L(t,0)-H_L(t,\varphi)$ is a sum of nonnegative terms
$L^3e^{-6t}\prod I_{Lz_\mu}(2t)(1-\cos(z\cdot\varphi))$.
Insert this difference into \eqref{f16v41:subordination} and use Tonelli to
obtain \eqref{f16v41:winding-formula}.

Absolute summability requires a separate check. The sum of the nonzero winding
coefficients before integration equals
$H_L(t,0)-L^3e^{-6t}I_0(2t)^3$. For $t\le1$ it is $O_L(t^L)$:
expand the Bessel generating series, and note that a nonzero displacement
multiple of $L$ needs at least $L$ steps. The sum of the remaining positive
series is bounded by its exponential majorant. For $t\ge1$ the difference
is between zero and $H_L(t,0)\le L^3$. Multiplication by
$e^{-2t}I_0(2t)/t$ is integrable in both ranges. This proves
$\sum a_L<\infty$, hence uniform convergence in $\varphi$. Each coefficient
is strictly positive by its positive integral. Keeping the six windings
$z=\pm e_\mu$ proves \eqref{f16v41:strict-loss}. Their equality condition
forces all $\varphi_\mu$ to be multiples of $2\pi$, and the converse is immediate.
\end{proof}

The threshold has a useful singularity check. Near $\varphi=0$,
\begin{equation}
 \mathcal E_L(\varphi)-\mathcal E_L(0)
             =\frac{2}{L}|\varphi|+O_L(|\varphi|^2).
 \label{f16v41:cusp}
\end{equation}
Indeed the $k=0$ charged term is
$\xi(\lambda_L(0,\varphi))=|\varphi|/L+O_L(|\varphi|^3)$.
The remaining finite sum is analytic and even, so its change is
$O_L(|\varphi|^2)$. This cusp comes from quantizing charged modes whose
harmonic frequency vanishes at the central point. It is not a new independent
classical potential to add to a model that still retains those same modes.
In particular a smooth Gaussian elimination with one positive gap on the
entire compact flat family would miss this rank change.

\subsection{The physical top and its stationary central-flat limit}

Let $\mathcal C_\sigma$ be the flat gauge orbit with fundamental holonomies
$\sigma_\mu I$, where $\sigma\in\{+1,-1\}^3$, and put
$\mathcal C_*:=\bigcup_\sigma\mathcal C_\sigma$. These are eight disjoint
compact orbits, not eight isolated points of the fine-link space.

\begin{theorem}[Two-sided native top identification and fixed-volume vacuum localization]
\label{f16v41:main}
At every fixed $B=m^3$, $m\ge4$, for the original normalized Wilson transfer,
\begin{equation}
 \boxed{\qquad
 \lim_{\gamma\to\infty}\Lambda_\gamma
     =\Lambda_{{\rm G},B}
     =\exp\left[-3\sum_{k\in\{0,\ldots,2m-1\}^3}
       \operatorname{arcosh}\left(1+2\sum_{\mu=1}^3
                      \sin^2\frac{\pi k_\mu}{2m}\right)\right].
 \qquad}
 \label{f16v41:top-limit}
\end{equation}
For every open neighborhood $O\supset\mathcal C_*$,
\begin{equation}
 \int_{O^c}\Omega_\gamma^2\,dH\longrightarrow0.
 \label{f16v41:vacuum-concentration}
\end{equation}
More precisely the complete stationary probabilities have the weak limit
\begin{equation}
 \boxed{\qquad
 \Omega_\gamma^2dH\ \Rightarrow\
             \frac18\sum_{\sigma\in\{\pm1\}^3}\mathfrak m_\sigma,
 \qquad}
 \label{f16v41:vacuum-limit}
\end{equation}
where $\mathfrak m_\sigma$ is normalized gauge-orbit Haar measure on
$\mathcal C_\sigma$. Equivalently every continuous gauge-invariant observable
$F$ converges in expectation to $\frac18\sum_\sigma F(U_\sigma)$, for any
representatives $U_\sigma\in\mathcal C_\sigma$.
\end{theorem}
\begin{proof}
Every compact periodic $\SU(2)$ flat connection has the commuting-holonomy
form used above. The winding theorem shows that its harmonic norm is at most
the trivial value, with equality exactly when every fundamental holonomy is
central. Thus the maximizing set in Theorem~\ref{f16v41:harmonic-principle}
is precisely $\mathcal C_*$. Apply that theorem and
\eqref{f16v41:reference-match}; these prove \eqref{f16v41:top-limit} and
\eqref{f16v41:vacuum-concentration}.

For the weights, multiply every link crossing one chosen dual periodic plane
by $-I$. Every spatial plaquette contains zero or two such factors. Both its
action and the normalized kinetic kernel are unchanged when the operation is
applied to both endpoints. These three commuting transformations generate the
usual $\mathbb Z_2^3$ center symmetry and permute the eight $\mathcal C_\sigma$
transitively. They preserve product Haar and positivity. Uniqueness of the
positive normalized Perron vector implies its invariance under them, as well
as under the local gauge group.

Every subsequential weak limit exists by compactness of $\mathcal X_L$, is
supported on $\mathcal C_*$ by the concentration assertion, and inherits these
continuous symmetries. Each $\mathcal C_\sigma$ is a single compact gauge
orbit, whose unique invariant probability is the Haar pushforward
$\mathfrak m_\sigma$. The center symmetry forces all eight weights equal.
This identifies every subsequential limit and proves \eqref{f16v41:vacuum-limit}.
No singular gauge section or Faddeev--Popov determinant is introduced. In
particular the limiting probability is singular relative to full link Haar;
no $L^2$ limit of normalized vacuum vectors is asserted.
\end{proof}

\subsection{What the new top estimate does and does not close}

The equality in \eqref{f16v41:top-limit} supplies the missing fixed-$B$ upper
comparison of V38--V40. It is a theorem about the same native physical top,
not about a surrogate local sampler. The proof uses no endpoint moment
approximation and no assumed vacuum overlap. It also identifies the weak
stationary probability, whereas V40 studies the nonstationary return prepared
by one particular endpoint source. Those laws can have total variation tending
to one while both concentrate near the same central-flat set on different
shrinking scales. Weak concentration does not determine that scale.

The theorem is deliberately qualitative in $\gamma$. In particular it does
\emph{not} imply $n_\gamma\log(\Lambda_\gamma/\Lambda_{{\rm G},B})\to0$.
Therefore V40's Gamma variables must still be centered at their declared
Gaussian threshold unless a rate is proved separately. Neither a rate for the
first Perron correction nor an estimate of the first excited ratio follows
from the normal-mode maximum principle. A positive finite limiting top is not
a positive mass above that top.

The next useful calculation is a uniform effective analysis in neighborhoods
of the central-flat orbits, keeping their nonquadratic soft directions and
center sectors. The cusp \eqref{f16v41:cusp} identifies why a Gaussian normal
elimination valid only off those orbits is insufficient there. Obtaining a
subleading spectral scale, its low-level operator, and a comparison to the
actual same-top transfer is a different task from evaluating another endpoint
coefficient. No physical time is defined by $\gamma$ or by an auxiliary update
clock. Fixed-lattice weak coupling, spatial removal, physical time calibration,
and an interacting continuum construction remain separate limits.

\subsection{Verification and preservation}

The accompanying diagnostic is
\begin{center}\small\path{verify_ym_f16_v41_flat_top.py}.\end{center}
Its exact fixtures check the localization identity, normal-coordinate
sum-of-squares bound, Gaussian prefactors, covariant-curl symbols, quotient
nonzero-spectrum identity, winding-series coefficients, color/center counts,
and the limiting statement's scope. Floating-point exploratory checks, when
used, are reported separately and are not integral or spectral enclosures.
The new checker passes 98 exact assertions. All 37 available V4--V40 checkers
were rerun unchanged and passed. No finite checker is described as formal
verification of the analytic localization proof.

Removing this exact insertion and reversing the one title-version change
recovers the supplied V40 byte for byte. The earliest exact fine-link transfer,
normalization, and rooted Gaussian identities are inherited inputs; the
singular-stratum operator comparison and winding argument are proved in this
section. The entire prior analytic archive is not independently recertified.
The next companion checkpoint remains V45.

\begin{firewall}
V41 proves the fixed-spatial-regulator limit of the original physical Perron
value, identifies its maximizing flat backgrounds, and proves weak convergence
of its stationary probability to the equal mixture of eight central-flat gauge
orbits. It does not estimate the first excitation, the subleading Perron scale,
a uniform-in-volume gap, a stationary fluctuation rate, a physical time, or an
interacting continuum limit. Its Gaussian threshold is exactly the inherited
one. All original kinetic scalars, magnetic words, and gauge constraints used
in the top comparison are retained; no endpoint source or sampling clock is
substituted for the transfer.
\end{firewall}
% END F16 V41 APPEND

% BEGIN F16 V42 APPEND -- 2026-09-18
% Insert before the final document-ending command of the cumulative V41.
% No predecessor proof, measure, scalar normalization, or clock is edited.

\clearpage
\section{V42: forbidden-region decay and exponentially close center sectors}
\label{f16v42:context}

V41 identifies the native top and eight symmetry-related regions supporting its
stationary probability, but does not estimate an excitation. Before constructing
a subleading oscillator model, we examine an implication of those eight regions
for the \emph{same} transfer. A fixed neighborhood of all maximizing flat orbits
has a complement whose compressed transfer norm stays strictly below the top.
Exponential weighting of the actual narrow Wilson kernel then gives quantitative
forbidden-region decay. The resulting center-character quasimodes lie in the
periodic-gauge-invariant Hilbert space and produce exponentially close sector
edges. No sampling generator or new kinetic normalization is introduced.

The conclusion is at fixed spatial regulator. It is an \emph{upper} bound on
seven center-charged excitation energies, not a positive lower mass bound, an
exact tunneling asymptotic, or a result for the center-neutral excitation gap.
An exact Schur reduction is also available on the union of the eight regions;
its entries remain operators on their complete local spaces, not an assumed
eight-dimensional effective Hamiltonian.

\subsection{Dependencies, exact symmetries, and the spectral question}

The V41 appendix and audit were read completely. The inherited input is its
fixed-regulator harmonic principle, including the strict maximizing set and
the original normalized fine-link kernel. The extension of its localization
argument to a compressed forbidden set is made explicit below. We do not use
V35--V40's endpoint asymptotics or any earlier auxiliary sampling gap. Targeted
searches of the supplied archives were followed by mounted-text inspection;
they are not an exhaustive novelty or proof audit. The reversible ground-state
transform already occurs in f15 V78 and its cited predecessors. The monograph's
Schur/Feshbach calculus is likewise credited, rather than renamed a new general
construction. The next scheduled companion checkpoint remains V45.

For context, Klein--L\'eonard--Rosenberger,
\href{https://arxiv.org/abs/1309.3918}{\emph{Agmon-type estimates for a class of
jump processes}}, Section 1, relates exponentially weighted estimates for
nonlocal operators to multiwell spectral splitting. Its hypotheses include
nondegenerate point wells in a different space. They are not assumed here:
our wells are compact gauge orbits in a singular flat set. The bounded-kernel
and resolvent argument needed below is proved directly. The center/electric-flux
sector terminology is consistent with de Forcrand--von Smekal,
\href{https://arxiv.org/abs/hep-lat/0107018}{\emph{'t Hooft Loops, Electric Flux
Sectors and Confinement in SU(2) Yang--Mills Theory}}; its numerical or thermal
conclusions are not inputs to the fixed-lattice theorem below.

Fix $L=2m$, $m\ge4$, $B=m^3$, $E=24B$, and
$\mathcal X_L=\SU(2)^E$ with probability Haar measure. Keep
\[
 T_\gamma=\overline T_\gamma
   =M_{b_\gamma}C_\gamma^{\otimes E}M_{b_\gamma},\qquad
 b_\gamma=e^{-\gamma S/2},\qquad
 \Lambda_\gamma=\|T_\gamma\|,\qquad T_\gamma\Omega_\gamma
       =\Lambda_\gamma\Omega_\gamma,\quad \|\Omega_\gamma\|_2=1.
\]
Here $T_\gamma$ is V41's full fine-link operator, not its root-coordinate
matrix before a further restriction. Its unique positive Perron vector is
periodic-gauge invariant, so its top is the same physical top. Let
\[
 \lambda_*:=\Lambda_{{\rm G},B}>0,\qquad
 \mathcal C_*:=\bigcup_{\sigma\in\{\pm1\}^3}\mathcal C_\sigma
\]
retain V41's meanings. Its theorem gives $\Lambda_\gamma\to\lambda_*$ and
identifies exactly these maximizing flat orbits.

For $\eta\in\mathsf Z:=\{\pm1\}^3$, let $z_\eta$ multiply every positive
$\mu$-link crossing the chosen periodic $\mu$-seam by $\eta_\mu I$.
These are commuting Haar-preserving isometries. They commute with periodic
gauge transformations, preserve each original plaquette and the two-end
kinetic kernel, and send $\mathcal C_\sigma$ to $\mathcal C_{\eta\sigma}$.
They are not identified with periodic gauge transformations: fundamental
Wilson-loop traces distinguish their actions. Write
\[
 (\mathsf Z_\eta f)(U)=f(z_\eta^{-1}U),\qquad
 \chi_a(\eta)=\prod_{\mu=1}^3\eta_\mu^{a_\mu},\quad a\in\{0,1\}^3,
\]
and define the physical character spaces
\begin{equation}
 \mathcal H_a=\{f\in L^2(\mathcal X_L)^{\rm gauge}:
            \mathsf Z_\eta f=\chi_a(\eta)f\ \hbox{for every }\eta\}.
 \label{f16v42:sectors}
\end{equation}
They are mutually orthogonal reducing subspaces. The Perron vector belongs
to $\mathcal H_0$. All constants below may depend on this fixed $L$ and on
fixed neighborhoods; none is claimed uniform in spatial volume.

\subsection{Exponential weighting of the literal normalized Wilson kernel}

Let $d_L$ be the round product geodesic distance. The raw kinetic density,
before the two magnetic factors, is denoted $c_\gamma(U,V)$; each of its rows
integrates to one with the \emph{original} $z_\gamma^E$.

\begin{proposition}[Uniform displacement weights and a bounded conjugation]
\label{f16v42:weighted-kernel}
There are finite $K_L,A_L$ and $b_L>0$ such that, for all $\gamma\ge1$,
\begin{align}
 \sup_U\int \sqrt\gamma\,d_L(U,V)e^{\sqrt\gamma d_L(U,V)}
                      c_\gamma(U,V)\,dH(V)&\le K_L,
                      \label{f16v42:exp-moment}\\
 \sup_U\int_{d_L(U,V)\ge r}c_\gamma(U,V)\,dH(V)
                      &\le A_Le^{-b_L\gamma r^2},\qquad r\ge0.
                      \label{f16v42:jump-tail}
\end{align}
If $\phi$ is a real $1$-Lipschitz function on $\mathcal X_L$, and
$0\le\alpha\le1$, multiplication by $W_\alpha=e^{\alpha\sqrt\gamma\phi}$
is bounded and invertible at each fixed coupling, and
\begin{equation}
 \|W_\alpha T_\gamma W_\alpha^{-1}-T_\gamma\|
             \le K_L\alpha.
 \label{f16v42:conjugation}
\end{equation}
For measurable sets $A,D$ separated by distance at least $r$,
$\|1_A T_\gamma1_D\|\le A_Le^{-b_L\gamma r^2}$.
\end{proposition}
\begin{proof}
For one displacement group, write its angle as $\theta\in[0,\pi]$.
The exact class density is proportional to
$e^{-\gamma(1-\cos\theta)}\sin^2\theta\,d\theta$.
The inequalities $1-\cos\theta\ge2\theta^2/\pi^2$ and
$z_\gamma\ge c\gamma^{-3/2}$, the latter obtained by integrating
$0\le\theta\le\gamma^{-1/2}$, show that the rescaled radial density is
bounded by $C x^2e^{-2x^2/\pi^2}\,dx$, $x=\sqrt\gamma\theta$.
The product displacement distance is $(\sum_e\theta_e^2)^{1/2}$.
Its exponential moment and its Gaussian tail follow by integrating this
fixed-dimensional product bound and splitting a Gaussian exponent in half.
Haar invariance makes the estimates independent of $U$. Symmetry supplies the
same column bounds. No large dimension is hidden as a volume-independent
constant.

The absolute difference kernel in \eqref{f16v42:conjugation} is bounded by
\[
 c_\gamma(U,V)\bigl(e^{\alpha\sqrt\gamma d_L(U,V)}-1\bigr)
 \le\alpha\sqrt\gamma d_L(U,V)e^{\sqrt\gamma d_L(U,V)}c_\gamma(U,V),
\]
because $0<b_\gamma\le1$. The row and column Schur estimates prove the norm
bound. This is not a Loewner inequality for an asymmetric weighted kernel.
The separated-support estimate follows from the same Schur argument using
\eqref{f16v42:jump-tail}. Only a Lipschitz weight is required; no derivative
through a principal logarithm cut is taken.
\end{proof}

\subsection{The complement of fixed central neighborhoods has a genuine barrier}

\begin{proposition}[A forbidden-region norm gap at the same native top]
\label{f16v42:barrier}
Let $\mathcal O$ be a fixed open neighborhood of $\mathcal C_*$, put
$P=1_{\mathcal O}$ and $Q=I-P$. There exist
$0<\delta<\lambda_*/8$ and $\gamma_0<\infty$ such that, for
$\gamma\ge\gamma_0$,
\begin{equation}
 \|QT_\gamma Q\|\le\lambda_*-4\delta,
       \qquad \Lambda_\gamma\ge\lambda_*-\delta.
 \label{f16v42:barrier-gap}
\end{equation}
The first norm is on the complete fine-link space; it consequently bounds its
physical restriction too.
\end{proposition}
\begin{proof}
The compact set $\mathcal F_L\setminus\mathcal O$ contains no maximizing
flat point. V41's continuous harmonic threshold therefore has a strict maximum
below $\lambda_*$ there, if that set is nonempty. Repeat the upper-bound
partition in Theorem~\ref{f16v41:harmonic-principle} for functions supported
in $\mathcal O^c$: use finitely many normal charts centered on these remaining
flat points, a piece inside $\mathcal O$ which contributes zero, and pieces
whose relevant supports are separated from the full flat set. The last pieces
have a positive action floor. The local norm errors can be fixed smaller than
one quarter of the strict threshold difference; the partition commutator then
tends to zero by Proposition~\ref{f16v41:IMS}. It follows that
$\limsup\|QT_\gamma Q\|<\lambda_*$. If no flat point remains, compactness
gives a positive action floor directly. The already-proved top limit and a
smaller positive $\delta$ give \eqref{f16v42:barrier-gap}.

This is not obtained from the small probability of $\mathcal O^c$. It is an
operator norm bound on all functions supported there. Nor is $S$ asserted
positive on that whole complement: noncentral \emph{flat} connections can
remain there, and their lower harmonic thresholds are essential.
\end{proof}

\begin{theorem}[Exponential localization of a whole native spectral band]
\label{f16v42:Agmon}
With $\mathcal O,\delta$ as above, put
\[
 \lambda_c=\lambda_*-\delta,\qquad
 \mathsf E_\gamma=1_{[\lambda_c,1]}(T_\gamma),\qquad
 \alpha=\min\{1,\delta/(1+K_L)\}>0.
\]
There is $C_{L,\mathcal O}<\infty$, independent of coupling and spectral-band
rank, such that
\begin{equation}
 \boxed{\quad
 \big\|e^{\alpha\sqrt\gamma d_L(\cdot,\mathcal O)}\mathsf E_\gamma\big\|
       \le C_{L,\mathcal O}.
 \quad}
 \label{f16v42:band-bound}
\end{equation}
In particular,
\begin{equation}
 \|1_{\{d_L(\cdot,\mathcal O)\ge r\}}\mathsf E_\gamma\|
       \le C_{L,\mathcal O}e^{-\alpha r\sqrt\gamma},\qquad
 \int e^{2\alpha\sqrt\gamma d_L(U,\mathcal O)}\Omega_\gamma(U)^2dH(U)
       \le C_{L,\mathcal O}^2.
 \label{f16v42:vacuum-tail}
\end{equation}
All these statements concern the actual full transfer and its own top.
\end{theorem}
\begin{proof}
Write $W=e^{\alpha\sqrt\gamma d_L(\cdot,\mathcal O)}$ and
$T_\alpha=WT_\gamma W^{-1}$. On $Q L^2$ set
$A=QT_\gamma Q$ and $A_\alpha=QT_\alpha Q$.
The preceding propositions give
$\|A_\alpha\|\le\lambda_*-3\delta$ and $\|T_\alpha\|\le1+\delta$.
Let $S_\gamma=T_\gamma|_{\operatorname{Ran}\mathsf E_\gamma}$, so
$S_\gamma\succeq\lambda_c I$. Regard
$X=WQ\mathsf E_\gamma$ as an operator from that spectral space to $QL^2$.
Since $WP=P$, the exact intertwining equation is
\[
 X S_\gamma-A_\alpha X
       =QT_\alpha P\mathsf E_\gamma=:B_\alpha.
\]
It has the norm-convergent solution
\begin{equation}
 X=\int_0^\infty e^{tA_\alpha}B_\alpha e^{-tS_\gamma}\,dt,
       \qquad \|X\|\le\frac{1+\delta}{2\delta}.
 \label{f16v42:Sylvester}
\end{equation}
Indeed differentiate $e^{tA_\alpha}Xe^{-tS_\gamma}$. Its limit at infinity
is zero, since $\lambda_c-\|A_\alpha\|\ge2\delta$; at each fixed coupling
$X$ is bounded before this uniform estimate is taken. This proves the integral
identity and bound without summing separate eigenvector estimates. Adding
$P\mathsf E_\gamma$ proves \eqref{f16v42:band-bound}. The Perron vector is
in the band by \eqref{f16v42:barrier-gap}. The tail conclusions follow by
multiplication. The integration parameter in \eqref{f16v42:Sylvester} is a
resolvent proof parameter, not a new physical time.
\end{proof}

\subsection{An exact eight-region Schur map, with all returning paths retained}

Choose $r_0>0$ so small that the closed $4r_0$-neighborhoods of the eight
$\mathcal C_\sigma$ are disjoint. This is possible because they are disjoint
compact subsets of the fixed compact metric space. Let
$\mathcal N_{j,\sigma}=\{d_L(U,\mathcal C_\sigma)<jr_0\}$ and
$\mathcal N_j=\bigcup_\sigma\mathcal N_{j,\sigma}$.
Apply the preceding results to $\mathcal O=\mathcal N_1$ and keep their
$\delta,\lambda_c,\alpha$. Put $P_\sigma=1_{\mathcal N_{1,\sigma}}$,
$P=\sum_\sigma P_\sigma$ and $Q=I-P$.

\begin{theorem}[Uniform forbidden complement and exponentially small inter-region returns]
\label{f16v42:Feshbach}
For $\gamma\ge\gamma_0$ and real $z\ge\lambda_c$, the operator
$z-QT_\gamma Q$ is positive and invertible on $QL^2$. On $PL^2$ define
\begin{equation}
 \mathfrak F_\gamma(z)=PT_\gamma P+
    PT_\gamma Q(z-QT_\gamma Q)^{-1}QT_\gamma P.
 \label{f16v42:Feshbach-map}
\end{equation}
Then $z$ is an eigenvalue of $T_\gamma$ if and only if it is an eigenvalue
of $\mathfrak F_\gamma(z)$, with multiplicities preserved. More generally,
$z-T_\gamma$ is invertible if and only if
$zP-\mathfrak F_\gamma(z)$ is invertible. Eigenvectors satisfy
\begin{equation}
 \psi=u+(z-QT_\gamma Q)^{-1}QT_\gamma P u,\quad u=P\psi,
 \qquad
 \|\psi\|^2=\langle u,(P-\mathfrak F_\gamma'(z))u\rangle.
 \label{f16v42:reconstruction}
\end{equation}
For distinct regions $\sigma,\tau$, uniformly in such $z$,
\begin{equation}
 \boxed{\quad
 \|P_\sigma\mathfrak F_\gamma(z)P_\tau\|
 \le C_{L,r_0}
 e^{-\alpha\sqrt\gamma\operatorname{dist}(\mathcal N_{1,\sigma},
                                           \mathcal N_{1,\tau})}.
 \quad}
 \label{f16v42:interwell}
\end{equation}
All statements restrict to periodic-gauge-invariant functions. Center symmetry
permutes the diagonal blocks by exact unitary equivalence.
\end{theorem}
\begin{proof}
The unweighted complement inverse has norm at most $(3\delta)^{-1}$ by
\eqref{f16v42:barrier-gap}. Block Gaussian elimination of $z-T_\gamma$
gives \eqref{f16v42:Feshbach-map} and the equivalences. The complement component
of an eigenvector is forced by its second block equation. Differentiating the
inverse gives
\[
 \mathfrak F_\gamma'(z)
   =-PT_\gamma Q(z-QT_\gamma Q)^{-2}QT_\gamma P\preceq0,
\]
which proves the norm identity. In particular the returning term is positive,
and its energy derivative is not omitted from reconstruction.

For the off-diagonal estimate take the different weight
$W_\tau=e^{\alpha\sqrt\gamma d_L(\cdot,\mathcal N_{1,\tau})}$.
The same conjugation bound gives
$\|QW_\tau T_\gamma W_\tau^{-1}Q\|\le\lambda_*-3\delta$.
Its resolvent at $z\ge\lambda_c$ has norm at most $(2\delta)^{-1}$.
Conjugating \eqref{f16v42:Feshbach-map}, using that all multiplication
projections commute with the weight, gives norm at most
\[
 (1+\delta)+(1+\delta)^2/(2\delta).
\]
On $P_\tau$ the weight is one; on $P_\sigma$ its inverse is at most the
exponential in \eqref{f16v42:interwell}. This proves that estimate, including
arbitrarily many returns through the forbidden region. Gauge and center
commutation follows because the regions and the original transfer have those
symmetries. No replacement probability or conditional row normalization is used.
\end{proof}

The Schur identity itself is standard. The new input is a \emph{proved} native
complement interval and its exponential cross-region control. Each diagonal
block still acts on an infinite-dimensional space and retains every soft and
hard coordinate inside that region. This is not an eight-state spectral model;
no inverse on the complement of just eight chosen vectors is supplied.

\subsection{Physical center-character quasimodes and sector edges}

There are smooth periodic-gauge-invariant functions $g_\sigma$ with
$0\le g_\sigma\le1$, equal to one on $\mathcal N_{3,\sigma}$, supported in
$\mathcal N_{4,\sigma}$, and satisfying
$g_\sigma(z_\eta U)=g_{\eta\sigma}(U)$.
To construct them, start with a smooth bump between the closures of the two
neighborhoods at $\sigma=(1,1,1)$, average over the compact gauge group, and
translate by the center action. Isometry preserves its plateau and support.
The eight supports are disjoint. For nonzero $a\in\{0,1\}^3$ put
\begin{equation}
 F_a(U)=\sum_\sigma\chi_a(\sigma)g_\sigma(U),\qquad
 u_{\gamma,a}=\frac{F_a\Omega_\gamma}{\|F_a\Omega_\gamma\|_2},
 \quad u_{\gamma,0}=\Omega_\gamma.
 \label{f16v42:quasimodes}
\end{equation}
The denominator is nonzero at all sufficiently large couplings.

\begin{theorem}[At least eight same-top eigenvalues with an exponential splitting bound]
\label{f16v42:sector-edges}
At fixed $L$ there are $C_L,\kappa_L>0$ and a finite threshold such that the
eight vectors in \eqref{f16v42:quasimodes} are orthonormal physical vectors,
$u_{\gamma,a}\in\mathcal H_a$, and
\begin{equation}
 \|(T_\gamma-\Lambda_\gamma)u_{\gamma,a}\|_2
                   \le C_Le^{-\kappa_L\sqrt\gamma}.
 \label{f16v42:quasimode-error}
\end{equation}
Let $\Lambda_{\gamma,a}=\|T_\gamma|_{\mathcal H_a}\|$. Then for every
nontrivial character $a$,
\begin{equation}
 \boxed{\quad
 0<1-\frac{\Lambda_{\gamma,a}}{\Lambda_\gamma}
       \le C_Le^{-\kappa_L\sqrt\gamma},\qquad
 0<-\log\frac{\Lambda_{\gamma,a}}{\Lambda_\gamma}
       \le C_Le^{-\kappa_L\sqrt\gamma}.
 \quad}
 \label{f16v42:sector-splitting}
\end{equation}
In particular, if $\lambda_{0,\gamma}=\Lambda_\gamma\ge\lambda_{1,\gamma}
\ge\cdots$ are the physical eigenvalues counted with multiplicity, then
\begin{equation}
 0<-\log(\lambda_{j,\gamma}/\Lambda_\gamma)
       \le C_Le^{-\kappa_L\sqrt\gamma},\qquad 1\le j\le7.
 \label{f16v42:eight-eigenvalues}
\end{equation}
The estimate asserts at least one near-top eigenvalue in each center sector,
not that there are exactly eight eigenvalues in this interval.
\end{theorem}
\begin{proof}
Center invariance of $\Omega_\gamma$ and covariance of the $g_\sigma$ give
$\mathsf Z_\eta(F_a\Omega_\gamma)=\chi_a(\eta)F_a\Omega_\gamma$.
Characters with different labels are orthogonal, including their orthogonality
to the Perron vector for $a\ne0$. Since $F_a^2=1$ on $\mathcal N_3$ and
$|F_a|\le1$ everywhere, \eqref{f16v42:vacuum-tail} gives
\[
 1-C_Le^{-4\alpha r_0\sqrt\gamma}
              \le\|F_a\Omega_\gamma\|_2^2\le1.
\]
In particular no small approximate vector norm is divided out.

The eigenvector equation gives the exact commutator identity
\[
 (\Lambda_\gamma-T_\gamma)(F_a\Omega_\gamma)
                         =[M_{F_a},T_\gamma]\Omega_\gamma.
\]
Split the input into $1_{\mathcal N_2}\Omega_\gamma$ and its complement.
The second part has norm at most $C_Le^{-\alpha r_0\sqrt\gamma}$, and the
commutator has norm at most two. If the first input lies in
$\mathcal N_{2,\sigma}$, the difference $F_a(U)-F_a(V)$ is zero when the
output lies in $\mathcal N_{3,\sigma}$. Otherwise $d_L(U,V)\ge r_0$.
The separated-jump Schur estimate therefore bounds this restricted commutator
by $2A_Le^{-b_L\gamma r_0^2}$. Combining the two parts and the norm bound proves
\eqref{f16v42:quasimode-error}, for a smaller fixed $\kappa_L>0$.

The sector restriction is compact positive self-adjoint. Its Rayleigh quotient
on $u_{\gamma,a}$ is at least
$\Lambda_\gamma-C_Le^{-\kappa_L\sqrt\gamma}$ and at most
$\Lambda_\gamma$. For sufficiently large coupling this is positive, and hence
the sector norm is an attained eigenvalue. It is strictly smaller than
$\Lambda_\gamma$ for $a\ne0$, by simplicity and center invariance of the
Perron vector. Since $\Lambda_\gamma\ge\lambda_*/2$ eventually, division and
$-\log(1-x)\le2x$ for $0\le x\le1/2$ prove
\eqref{f16v42:sector-splitting}. The eight mutually orthogonal eigenvectors,
one from each sector, give \eqref{f16v42:eight-eigenvalues} by the min--max
principle. They need not be the only near-top vectors.
\end{proof}

The proof uses the existing vacuum as a variational factor; it does not provide
a numerical representation of that vector. Its trial multipliers are genuinely
gauge invariant. This is why the result concerns physical sector eigenvalues,
unlike a non-invariant trial justified only by equality of the two Perron tops.
The positive constants and entry threshold are not optimized or numerically
enclosed. No matching tunneling lower bound or leading tunneling action is
claimed.

\subsection{Stationary center memory and the neutral-sector distinction}

Use the \emph{existing} ground-state transform, now on the literal fine-link
space, with $d\mu_\gamma=\Omega_\gamma^2dH$:
\begin{equation}
 (\mathcal P_\gamma f)(U)=\frac{1}{\Lambda_\gamma\Omega_\gamma(U)}
          \int T_\gamma(U,V)\Omega_\gamma(V)f(V)\,dH(V).
 \label{f16v42:Doob}
\end{equation}
This is a reversible positive contraction with stationary law $\mu_\gamma$;
multiplication by $\Omega_\gamma$ is unitary to $T_\gamma/\Lambda_\gamma$.
Its transition is not obtained by dividing rows of $T_\gamma$ by their sums.
It is one original transfer step at its \emph{actual} Perron top.

\begin{proposition}[Exponentially long stationary center correlations]
\label{f16v42:memory}
For each nontrivial $a$, $\mu_\gamma F_a=0$. For every integer $N\ge0$,
\begin{equation}
 0\le 1-
 \frac{\mathbb E_{\mu_\gamma}[F_a(U_0)F_a(U_N)]}
      {\mathbb E_{\mu_\gamma}F_a^2}
 \le\min\{1,C_LN e^{-\kappa_L\sqrt\gamma}\},
 \label{f16v42:correlation}
\end{equation}
where $(U_j)$ is the stationary chain of \eqref{f16v42:Doob}. In particular,
the normalized correlation tends to one for
$N_\gamma=o(e^{\kappa_L\sqrt\gamma})$.
\end{proposition}
\begin{proof}
The normalized correlation is
$\langle u_{\gamma,a},(T_\gamma/\Lambda_\gamma)^N u_{\gamma,a}\rangle$.
The operator is a positive contraction. Its spectral theorem and
$0\le1-x^N\le N(1-x)$, $0\le x\le1$, bound the deficit by $N$ times the
one-step Rayleigh deficit. Apply \eqref{f16v42:quasimode-error} and
$\Lambda_\gamma\ge\lambda_*/2$. This is a stationary correlation estimate,
not a claim that a finite-time endpoint source has prepared that stationary
law or that an entire path stays in a well without a separate event estimate.
\end{proof}

\begin{proposition}[The center-charged band is invisible to center-neutral insertions]
\label{f16v42:selection}
If a bounded gauge-invariant multiplication $F$ is invariant under every
$z_\eta$, then $F\Omega_\gamma\in\mathcal H_0$ and
\[
 \langle\psi_a,F\Omega_\gamma\rangle=0
      \quad(a\ne0,\ \psi_a\in\mathcal H_a).
\]
In particular this holds for every bounded function of a fixed collection of
contractible Wilson-loop traces. Conversely a multiplication with character
$\chi_a$ creates a vector in $\mathcal H_a$ from the vacuum.
\end{proposition}
\begin{proof}
Apply $\mathsf Z_\eta$ inside the inner product. The Perron vector is invariant,
and a nontrivial character has a group element with value $-1$. This forces the
claimed inner product to vanish. The center seam is a closed mod-two cocycle;
its pairing with a contractible loop is zero. Equivalently every elementary
plaquette crosses each seam an even number of times, and a contractible loop
bounds a lattice surface. Centrality allows the signs to cancel without
reordering noncommutative link factors. Its trace is therefore center invariant.
The last assertion is the same character multiplication identity.
\end{proof}

There are thus seven nontrivial center-charge channels with small transfer
energies, but no estimate here for the first excited value in $\mathcal H_0$.
This distinction cannot be erased by calling all eight sectors gauge equivalent.
The center-neutral spectral problem may contain additional near-top states;
no isolated eight-state cluster or complement gap after removing just these
quasimodes is proved. The full operator-valued map
\eqref{f16v42:Feshbach-map} deliberately retains those states.

\subsection{What changes at the spectral frontier}

V42 supplies an actual same-top excited-sector comparison at fixed spatial
regulator, together with exponential forbidden-region localization and an exact
regional elimination whose cross-region entries are exponentially small. It
does not determine the common shift $\Lambda_\gamma-\lambda_*$. In particular
V40's entrance law still cannot be recentered at the native top on its moving
$1/n_\gamma$ scale without a separate rate for that common shift.

If an independently specified physical time step is $a_t(\gamma)>0$, then the
center-sector energies obey only the converted upper bound
\[
 0<E_a(\gamma):=-a_t(\gamma)^{-1}
       \log(\Lambda_{\gamma,a}/\Lambda_\gamma)
       \le\frac{C_Le^{-\kappa_L\sqrt\gamma}}{a_t(\gamma)}.
\]
Whether this tends to zero cannot be decided without that time scale; a fixed
number of spatial lattice sites is not a fixed physical-volume continuum
limit either. No continuum gaplessness or positive continuum mass is inferred.
The charged characters are often described as finite-volume electric-flux
sectors. They must not be substituted for the center-neutral local excitation
channel.

The next useful problem is now inside the diagonal blocks of
\eqref{f16v42:Feshbach-map}: identify their nonquadratic soft operator and a
quantitative remainder at the same spectral top, or obtain a lower splitting
estimate with the actual returning paths. The growing-regulator and
physical-time costs remain separate from this fixed-regulator exponential
bound. This respects the supplied memoranda's separate locality, normalization,
and physical-transfer obligations; no external Navier--Stokes theorem is a
premise of this iteration.

\subsection{Verification and preservation}

The diagnostic is
\begin{center}\small\path{verify_ym_f16_v42_center_splitting.py}.\end{center}
It checks exact center/plaquette and winding-loop parities on the actual finite
geometry, character projectors, a noncommuting weighted Sylvester fixture,
Schur reconstruction with its derivative metric, and a strictly positive
finite multiwell transfer with a known Perron vector. It also tests stationary
correlation and sector-selection identities, and a counterfixture showing why
an eighth near-top eigenvalue does not isolate an eight-state cluster. These
are finite algebraic tests, not native Wilson integral or spectral enclosures,
and not formal verification of the analytic localization argument.

The package records actual diagnostic executions, source hashes, proof-review
scope, and rendering checks. Only the literal title version and this terminal
insertion differ from V41. Removing them recovers its exact bytes. The new
native conclusions depend on V41's harmonic principle and strict maximizing
set; those inputs are stated rather than independently recertified for the
whole preceding archive. The next scheduled checkpoint remains V45.

\begin{firewall}
V42 proves exponential localization outside fixed central-flat neighborhoods,
a same-native-top Schur reduction with exponentially small inter-region blocks,
and at least one exponentially close eigenvalue in each of eight physical
center-character sectors at fixed lattice size. It gives upper bounds on the
seven charged excitation ratios and stationary center correlations, not a lower
mass bound, a center-neutral gap, an exact tunneling rate, an isolated
eight-dimensional low-energy model, or a continuum result. All magnetic words,
kinetic scalars, gauge constraints, returning paths, and the actual Perron
normalization used in these deductions are retained.
\end{firewall}
% END F16 V42 APPEND

% BEGIN F16 V43 APPEND -- 2026-09-18
% Insert before the final document-ending command of cumulative V42.
% Predecessor bodies, native kernels, probabilities, and clocks are unchanged.

\clearpage
\section{V43: a single-region spectral pencil and the neutral return metric}
\label{f16v43:context}

V42 proves exponentially close charged sector tops, but does not separate a
neutral excitation from an additional local soft state. Its regional Schur map
retains exactly the space needed for that question. We now use the center
symmetry on this \emph{operator-valued} map, without selecting one vector per
region. All near-top levels, counted with multiplicity, can be compared with
one common single-region nonlinear eigenvalue problem. Its error is exponential
and has no level-count factor. A second calculation identifies the correct
same-top neutral gap on that region: its norm is the returning-state metric,
not the flat norm of a bare compression.

These are deductions about the original finite-lattice spectrum, conditional
on V42's stated forbidden-complement bounds. They do not evaluate the remaining
single-region soft spectrum. In particular no algebraic neutral-gap scale,
finite-dimensional eight-state model, or continuum mass is asserted.

\subsection{Inputs, ownership, and the retained physical spaces}

The complete supplied V42 appendix and analytic audit were read. Its complement
interval and exponential inter-region estimate are the native inputs. The
Grand Tensor's signed Feshbach inertia and two-scale gap compiler, and f15
V78's warning about an unpaid vertical norm, were checked at identified
passages. Those general mechanisms are credited, not claimed new. The present
step folds the actual eight regions into their physical character spaces,
counts every level in a fixed spectral interval, and pays the returning norm
in a neutral-gap and source-resolvent comparison. Two empty semantic searches
were followed by mounted exact-text inventories; they do not establish absence
of related results. The next companion checkpoint remains V45.

For primary context, Dusson--Sigal--Stamm,
\href{https://arxiv.org/html/2105.02058}{\emph{The Feshbach--Schur map and
perturbation theory}}, Theorem 1.2, Corollary 1.3 and Proposition 1.4,
describe isospectral reconstruction and spectral fixed points. Their setting
uses a finite-rank retained projection for perturbation estimates. Our retained
space is infinite dimensional. We prove the compact-operator counting and
comparison statements below directly, and do not assume differentiability of
ordered eigenvalue branches at crossings. The supplied wider-literature
memorandum recommends keeping the full return; it does not supply the local
neutral gap that remains to be computed here.

Fix $L=2m$, $m\ge4$. Use V42's unchanged physical transfer on
\[
 T=T_\gamma,\qquad
 \mathcal H^{\rm phys}=L^2(\mathcal X_L,dH)^{\rm gauge}.
\]
Let $\Lambda=\Lambda_\gamma$ and $\Omega=\Omega_\gamma$ be its actual top and
unit positive vacuum. Keep V42's disjoint regions
$\mathcal O_\sigma=\mathcal N_{1,\sigma}$, their projections $P_\sigma$,
$P=\sum_\sigma P_\sigma$, and $Q=I-P$. All these projections preserve the
physical space. Write
\[
 c=\lambda_*-\delta>0,\qquad A_Q=QTQ,\qquad
 R(z)=(z-A_Q)^{-1}\quad\hbox{on }Q\mathcal H^{\rm phys}.
\]
For sufficiently large coupling the inherited bounds are
\begin{equation}
 \|A_Q\|\le c-3\delta,\qquad \Lambda\ge c,
 \qquad z\ge c\ \Longrightarrow\ \|R(z)\|\le(3\delta)^{-1}.
 \label{f16v43:input}
\end{equation}
Let $\mathsf Z=\{\pm1\}^3$, with identity $e=(1,1,1)$ and character
$\chi_a(\eta)=\prod_\mu\eta_\mu^{a_\mu}$. Its exact unitary action is
$\mathsf Z_\eta$, and $\mathcal H_a$ has V42's meaning. No center operation
is quotiented as a periodic gauge transformation.

Define the one-region physical space and its extension by zero by
\[
 \mathfrak h=L^2(\mathcal O_e,dH)^{\rm gauge},\qquad
 \iota:\mathfrak h\longrightarrow P\mathcal H^{\rm phys}.
\]
It retains every soft and hard variable of that open region. The isometries
\begin{equation}
 W_a u=8^{-1/2}\sum_{\eta\in\mathsf Z}
                    \chi_a(\eta)\mathsf Z_\eta\iota u
 \label{f16v43:Wa}
\end{equation}
have ranges $P\mathcal H_a$, are mutually orthogonal, and together exhaust
$P\mathcal H^{\rm phys}$. This follows from the disjoint translated supports
and the eight-character orthogonality identity. No local state is truncated.

\subsection{An exact character decomposition with differentiated return control}

Retain V42's map
\[
 \mathfrak F(z)=PTP+PTQ R(z)QTP.
\]
On the common space $\mathfrak h$ put
\begin{equation}
 A(z)=\iota^*\mathfrak F(z)\iota,\qquad
 B_\eta(z)=\iota^*\mathfrak F(z)\mathsf Z_\eta\iota,
 \qquad F_a(z)=W_a^*\mathfrak F(z)W_a.
 \label{f16v43:blocks}
\end{equation}
Here $B_e=A$. The letter $A(z)$ is a regional operator, not an endpoint scalar.

\begin{proposition}[Exact sector pencils and their exponentially small differences]
\label{f16v43:Fourier}
For every real $z\ge c$,
\begin{equation}
 \boxed{\quad F_a(z)=A(z)+\sum_{\eta\ne e}\chi_a(\eta)B_\eta(z).
 \quad}
 \label{f16v43:Fourier-formula}
\end{equation}
The operators $F_a(z),A(z)$ are positive, compact and self-adjoint, and
\begin{equation}
 0\preceq-F_a'(z)\preceq K I,\qquad
 0\preceq-A'(z)\preceq K I,\qquad K=(3\delta)^{-2}.
 \label{f16v43:monotonicity}
\end{equation}
There are fixed $C,\kappa>0$, depending on $L$ and the regions but not on
coupling or level count, such that, with
$\epsilon_\gamma=C e^{-\kappa\sqrt\gamma}$,
\begin{equation}
 \sup_{z\ge c}\bigl(\|F_a(z)-A(z)\|+
                         \|F_a'(z)-A'(z)\|\bigr)\le\epsilon_\gamma.
 \label{f16v43:exponential}
\end{equation}
No positivity of an individual $B_\eta$ in operator order is required.
\end{proposition}
\begin{proof}
Insert \eqref{f16v43:Wa} twice. Center commutation turns each term into
$\iota^*\mathfrak F\mathsf Z_{\eta^{-1}\tau}\iota$; each relative label occurs
eight times. This proves \eqref{f16v43:Fourier-formula}, including its
normalization. Since the group has exponent two, $\mathsf Z_\eta^*=\mathsf Z_\eta$.
Commutation makes $B_\eta$ self-adjoint, although it need not be positive.
Positivity of $F_a,A$ follows by compression of $\mathfrak F\succeq0$.
Compactness follows from compactness of $T$ and boundedness of $R(z)$.

Differentiate the actual inverse:
\[
 \mathfrak F'(z)=-PTQ R(z)^2QTP.
\]
The bound $\|T\|\le1$ proves \eqref{f16v43:monotonicity}. V42 already bounds
each off-diagonal block of $\mathfrak F$. Its same exponential conjugation
also bounds the derivative. In detail, with a weight of distance from the
source region, the conjugated complement inverse has norm at most
$(2\delta)^{-1}$ and the conjugated transfer has norm at most $1+\delta$.
Consequently the conjugate of $\mathfrak F'$ has norm at most
$(1+\delta)^2/(2\delta)^2$. Restriction to a distinct target region restores
the same exponential separation factor as for $\mathfrak F$. The minimum
separation of the fixed regions is positive. Sum the seven bounds and enlarge
$C$ to obtain \eqref{f16v43:exponential}. This calculation retains the inverse
square and all its complementary returns; it does not differentiate an
unspecified big-O remainder.
\end{proof}

At a sector eigenvalue $z\ge c$, the exact reconstruction from $u\in\mathfrak h$
is
\begin{equation}
 J_a(z)u=W_a u+R(z)Q T W_a u,
 \qquad J_a(z)^*J_a(z)=I-F_a'(z).
 \label{f16v43:reconstruction}
\end{equation}
Its residual is $(T-z)J_a(z)u=W_a(F_a(z)-z)u$. More generally, for two distinct
real parameters $z,w\ge c$,
\begin{equation}
 J_a(z)^*J_a(w)
     =I-\frac{F_a(z)-F_a(w)}{z-w}.
 \label{f16v43:cross-metric}
\end{equation}
These identities follow by multiplication and the resolvent identity. They
keep the norm and mutual inner products of the returning full vectors. An
orthonormal list of regional vectors in the flat norm is not automatically
an orthonormal list of reconstructed physical states.

\subsection{Counting every near-top level, without an isolated-cluster premise}

For a compact positive operator $C$, let $n_+(C-z)$ count its eigenvalues
strictly larger than $z$, with multiplicity. For $z\ge c$ define
\[
 N_a(z)=\dim 1_{(z,1]}(T|_{\mathcal H_a}),\qquad
 N_{\rm reg}(z)=n_+(A(z)-zI).
\]
These numbers are finite. Their threshold convention excludes eigenvalues
exactly equal to $z$.

\begin{theorem}[Exact inertia and uniform sector counting brackets]
\label{f16v43:counting}
For $z\ge c$,
\begin{equation}
 \boxed{\quad N_a(z)=n_+(F_a(z)-zI).\quad}
 \label{f16v43:inertia}
\end{equation}
For $z-\epsilon_\gamma\ge c$,
\begin{equation}
 \boxed{\quad
 N_{\rm reg}(z+\epsilon_\gamma)\le N_a(z)
                       \le N_{\rm reg}(z-\epsilon_\gamma).
 \quad}
 \label{f16v43:count-bracket}
\end{equation}
In particular, if the two regional counts equal $k$, every physical center
sector contains exactly $k$ eigenvalues above $z$, and the entire physical
space contains exactly $8k$ there. No estimate of $k$ is assumed or supplied
by symmetry alone.
\end{theorem}
\begin{proof}
Restrict $z-T$ to $\mathcal H_a=(P\mathcal H_a)\oplus(Q\mathcal H_a)$.
Its second diagonal block is positive and invertible by \eqref{f16v43:input}.
The bounded triangular congruence implementing Schur elimination transforms
its form into the direct sum of $z-F_a(z)$ and that positive block. An
invertible change of variables preserves the maximal dimension of a negative
definite subspace: apply the change and its inverse to such subspaces. Thus
the finite negative index of $z-T$ is exactly that of $z-F_a(z)$, proving
\eqref{f16v43:inertia} without a finite-dimensional retained projection.
This is the inherited inertia principle in its present compact setting.

The norm estimate and monotonicity give
\[
 A(z+\epsilon_\gamma)-(z+\epsilon_\gamma)I
 \preceq F_a(z)-zI
 \preceq A(z-\epsilon_\gamma)-(z-\epsilon_\gamma)I.
\]
Monotonicity of the finite positive index and \eqref{f16v43:inertia} prove the
bracket. Summing the eight orthogonal reducing sector counts proves its last
claim. Strict inequalities in the spectral intervals are retained at both
ends; no generic absence of a threshold eigenvalue is presumed.
\end{proof}

Let $\nu_j(z)$ be the nonincreasingly ordered eigenvalues of $A(z)$, repeated
with multiplicity. For each $j$ with $\nu_j(c)>c$, define $r_j$ to be the unique
solution
\begin{equation}
                      \nu_j(r_j)=r_j,\qquad r_j>c.
 \label{f16v43:roots}
\end{equation}
These roots have an energy-independent interpretation as well. Define the
native one-region compression
\begin{equation}
 T^{\rm one}=(P_e+Q)T(P_e+Q)
       \quad\hbox{on }(P_e+Q)\mathcal H^{\rm phys}.
 \label{f16v43:one-region-transfer}
\end{equation}
It retains the chosen region and the entire forbidden complement, but deletes
the other seven retained regions. Its Schur map is exactly $A(z)$, so its
eigenvalues above $c$ are precisely the $r_j$. Thus the reference below is an
actual compressed native transfer, not an assumed oscillator Hamiltonian.
Its top is not substituted for the full physical top without the stated
error comparison. Powers of this compression restrict slice-time paths; no
continuous-time killing or new physical clock is implied.

These roots are in descending order. Existence follows because $A(z)$ is
bounded as $z\to\infty$, while $z$ increases without bound. Uniqueness follows
from monotonicity. In fact
\[
 t-s\le[t-\nu_j(t)]-[s-\nu_j(s)]\le(1+K)(t-s),\qquad t>s\ge c.
\]
This uses the min--max principle and \eqref{f16v43:monotonicity}; it remains
valid at crossings and repeated eigenvalues. At $z=\Lambda$, positivity of
$\Lambda-T$ implies $\Lambda P-\mathfrak F(\Lambda)\succeq0$, so every root
is at most $\Lambda\le1$.

\begin{theorem}[One regional ladder for all eight native sectors]
\label{f16v43:ladders}
If $r_j>c+\epsilon_\gamma$, the $j$th eigenvalue $\lambda_{a,j}$ of
$T|_{\mathcal H_a}$ exists above $c$ and satisfies
\begin{equation}
 \boxed{\qquad |\lambda_{a,j}-r_j|\le\epsilon_\gamma
                         \quad\hbox{for every }a.\qquad}
 \label{f16v43:ladder-bound}
\end{equation}
The indexing here starts at $j=1$ in \emph{each} sector; in the neutral sector
$\lambda_{0,1}=\Lambda$. It counts multiplicity and introduces no factor $j$.
In particular corresponding sector levels differ by at most $2\epsilon_\gamma$.
For two consecutive roots above this threshold,
\begin{equation}
 |(\lambda_{a,j}-\lambda_{a,j+1})-(r_j-r_{j+1})|
                                      \le2\epsilon_\gamma.
 \label{f16v43:level-spacing}
\end{equation}
At sufficiently large coupling $r_1>c+\epsilon_\gamma$, and
\begin{equation}
 \left|\log\frac{\Lambda}{\lambda_{a,j}}-
                       \log\frac{r_1}{r_j}\right|
                       \le\frac{2\epsilon_\gamma}{c}.
 \label{f16v43:log-ladder}
\end{equation}
\end{theorem}
\begin{proof}
Let $\nu_{a,j}(z)$ be the ordered eigenvalues of $F_a(z)$. They satisfy
$|\nu_{a,j}(z)-\nu_j(z)|\le\epsilon_\gamma$. Each function
$z-\nu_{a,j}(z)$ is continuous and increases by at least the increase in $z$.
At $r_j-\epsilon_\gamma$ it is nonpositive, and at
$r_j+\epsilon_\gamma$ it is nonnegative. Thus its unique root lies in that
interval. Exact inertia identifies that root with the $j$th physical sector
eigenvalue, including possible multiplicities. One can make the interval
strictly larger and then decrease it to include an equality at an endpoint.
This avoids assuming simple or differentiable eigenbranches.

The two triangle inequalities give \eqref{f16v43:level-spacing} and the
cross-sector bound. For the top-root existence, $F_0(c)\succeq F_0(\Lambda)$ and
$\nu_{0,1}(\Lambda)=\Lambda$. Thus
$\nu_1(c)\ge\Lambda-\epsilon_\gamma>c$ eventually, since
$\Lambda\to c+\delta$. The already proved root comparison then puts
$r_1$ a fixed distance above $c$. All arguments of logarithms in
\eqref{f16v43:log-ladder} are at least $c$. The mean value bound for $\log$
and two applications of \eqref{f16v43:ladder-bound} prove it.
\end{proof}

The theorem applies to the second \emph{neutral} eigenvalue when that level is
in the declared interval. It is not limited to the seven charged tops already
controlled in V42. For example, if $r_2>c+\epsilon_\gamma$, the additive
neutral deficit satisfies
\[
 |(\Lambda-\lambda_{0,2})-(r_1-r_2)|\le2\epsilon_\gamma.
\]
If a regional counting calculation instead gives
$N_{\rm reg}(z-\epsilon_\gamma)\le1$ with $c+\epsilon_\gamma\le z<\Lambda$,
then $\lambda_{0,2}\le z$, so $\Lambda-\lambda_{0,2}\ge\Lambda-z$.
These are proved implication rules, not an evaluation of either regional
count or the difference $r_1-r_2$ for the Wilson operator.

\subsection{The neutral excitation uses the returning-state norm}

A frozen spectral pencil still has an exact same-top meaning if its metric is
retained. Work now only in $\mathcal H_0$, the periodic-gauge-invariant,
center-neutral sector. All complementary inverses below are restricted to
$Q\mathcal H_0$. Write
\[
 D=\Lambda-QTQ,\quad B_0=QTW_0,\quad
 S=\Lambda I-F_0(\Lambda),\quad
 M=I-F_0'(\Lambda)=I+B_0^*D^{-2}B_0.
\]
The label $0$ on $F_0,W_0$ means the neutral character, not a Gaussian model.
Put
\begin{equation}
 J=W_0+D^{-1}B_0,\qquad
 u_0=W_0^*P\Omega,\qquad e_0=M^{1/2}u_0,
 \qquad \mathcal K_{\rm reg}=M^{-1/2}SM^{-1/2}.
 \label{f16v43:neutral-data}
\end{equation}
Then $J^*J=M$, $Ju_0=\Omega$, and $\|e_0\|=1$. Define
\[
 g_\gamma=\inf_{v\perp e_0,\,\|v\|=1}
                    \langle v,\mathcal K_{\rm reg}v\rangle,
 \qquad d_\gamma=\Lambda-\|QTQ|_{Q\mathcal H_0}\|\ge3\delta,
\]
with $d_\gamma=\Lambda$ when the complementary restriction is zero.
At each fixed coupling the compactness and the simple native Perron kernel
imply $g_\gamma>0$. This fixed-regulator positivity is not a useful uniform
lower bound. Let
$\Delta_\gamma^{\rm neu}=\Lambda-\lambda_{0,2}$ be the actual neutral
additive excitation gap.

\begin{theorem}[A paid same-top comparison for the complete neutral gap]
\label{f16v43:neutral-gap}
The quantities above obey
\begin{equation}
 \boxed{\qquad
 \frac{d_\gamma g_\gamma}{d_\gamma+g_\gamma}
       \le\Delta_\gamma^{\rm neu}\le g_\gamma,
 \qquad I\preceq M\preceq(1+K)I.\qquad}
 \label{f16v43:neutral-gap-bound}
\end{equation}
In particular, along any fixed-$L$ sequence on which $g_\gamma\to0$,
\begin{equation}
 \frac{\Delta_\gamma^{\rm neu}}{g_\gamma}\longrightarrow1,\qquad
 -\log(\lambda_{0,2}/\Lambda)
                      =\frac{g_\gamma}{\Lambda}(1+o(1)).
 \label{f16v43:neutral-relative}
\end{equation}
No hypothesis that $g_\gamma\to0$, or a formula for its rate, is proved here.
The theorem holds on all neutral form vectors, not just a selected source bank.
\end{theorem}
\begin{proof}
Block completion at the \emph{native} top gives, for every $u\in\mathfrak h$
and $w\in Q\mathcal H_0$,
\begin{equation}
 \langle Ju+w,(\Lambda-T)(Ju+w)\rangle
                       =\langle u,Su\rangle+\langle w,Dw\rangle.
 \label{f16v43:completion}
\end{equation}
The graph range of $J$ need not be orthogonal to $Q\mathcal H_0$.
Nevertheless its distance from the actual vacuum satisfies
\[
 \operatorname{dist}(Ju,\mathbb C\Omega)^2
   =\operatorname{dist}(M^{1/2}u,\mathbb C e_0)^2
   \le g_\gamma^{-1}\langle u,Su\rangle.
\]
Every neutral $\psi$ has a unique decomposition $\psi=Ju+w$ by taking
$u=W_0^*P\psi$. The triangle inequality for the distance to the vacuum,
followed by scalar Cauchy--Schwarz, yields
\[
 \operatorname{dist}(\psi,\mathbb C\Omega)^2
 \le\left(\sqrt{\langle u,Su\rangle/g_\gamma}
          +\sqrt{\langle w,Dw\rangle/d_\gamma}\right)^2
 \le(g_\gamma^{-1}+d_\gamma^{-1})
               \langle\psi,(\Lambda-T)\psi\rangle.
\]
This proves the lower bound. It does not assert the false full-vector bound
$\Lambda-T\succeq d_\gamma Q$.

For the upper bound take a unit $v\perp e_0$ approaching the infimum defining
$g_\gamma$, and put $\psi=JM^{-1/2}v$. It is a unit neutral vector orthogonal
to $\Omega$, and its Rayleigh energy is exactly
$\langle v,\mathcal K_{\rm reg}v\rangle$. The variational principle gives
$\Delta_\gamma^{\rm neu}\le g_\gamma$. The metric bounds follow from
\eqref{f16v43:monotonicity}. Finally $M-I$ and $F_0(\Lambda)$ are compact,
so $\mathcal K_{\rm reg}$ is a compact perturbation of $\Lambda I$.
Its kernel corresponds exactly to the one-dimensional native kernel by
\eqref{f16v43:completion}; zero is isolated. This also justifies the claimed
fixed-coupling positivity of $g_\gamma$. Since $d_\gamma\ge3\delta$ and
$\Lambda\to\lambda_*>0$, the two limits in
\eqref{f16v43:neutral-relative} follow from the bracket and the scalar logarithm.
\end{proof}

\paragraph{Why the metric is not optional.}
Consider the finite reference form, unrelated to native matrix elements,
\[
 L_\tau=\begin{pmatrix}
 0&0&0\\ 0&\tau+b^2/d&-b\\ 0&-b&d
 \end{pmatrix},\qquad \tau,d>0.
\]
Retain the first two coordinates. The Schur form is
$S=\operatorname{diag}(0,\tau)$, but its returning norm is
$M=\operatorname{diag}(1,1+b^2/d^2)$. Its first positive eigenvalue satisfies
\[
 \lambda_1(L_\tau)=\frac{\tau}{1+b^2/d^2}+O_{b,d}(\tau^2).
\]
The metric gap has exactly that leading coefficient. Using the bare gap
$\tau$ would miss an arbitrary fixed factor. This is consistent with the
f15 V78 counterfixture; the present proof pays the graph norm and uses the
full completed form rather than assuming an unpaid vertical inequality.

\subsection{Neutral observables retain their complementary source and vacuum subtraction}

Let $f\in\mathcal H_0$ satisfy $\langle\Omega,f\rangle=0$. For example
$f=(F-\mu_\gamma F)\Omega$ for a bounded gauge- and center-invariant
multiplication $F$, with $d\mu_\gamma=\Omega^2dH$. For real $z>\Lambda$ put
\[
 b_f(z)=W_0^*Pf+W_0^*TQ R(z)Qf.
\]
The symbol $b_f$ denotes an effective source, not the magnetic multiplier.

\begin{proposition}[Exact source resolvent on the single-region pencil]
\label{f16v43:source}
For every real $z>\Lambda$,
\begin{equation}
 \boxed{\quad
 \langle f,(z-T)^{-1}f\rangle
 =\langle Qf,R(z)Qf\rangle+
 \langle b_f(z),(zI-F_0(z))^{-1}b_f(z)\rangle.
 \quad}
 \label{f16v43:source-resolvent}
\end{equation}
At the native top $b_f(\Lambda)=J^*f$ is orthogonal to $u_0$. With the
Moore--Penrose inverse on the corresponding vacuum complement,
\begin{equation}
 \langle f,(\Lambda-T)^\dagger f\rangle
 =\langle Qf,D^{-1}Qf\rangle+
 \langle M^{-1/2}b_f(\Lambda),
       \mathcal K_{\rm reg}^{\dagger}M^{-1/2}b_f(\Lambda)\rangle.
 \label{f16v43:source-top}
\end{equation}
All inverses in this identity exist on their declared complements at each
fixed finite coupling. No bound uniform in coupling on
$\mathcal K_{\rm reg}^{\dagger}$ is inferred.
\end{proposition}
\begin{proof}
Solve the complementary block equation of $(z-T)\psi=f$, substitute it into
the retained equation, and expand $\langle f,\psi\rangle$. The two cross
terms assemble into $b_f(z)$ and give \eqref{f16v43:source-resolvent}; keeping
only $W_0^*Pf$ would discard a source return.
At $z=\Lambda$ the exact relation $Ju_0=\Omega$ gives
$\langle u_0,b_f(\Lambda)\rangle=\langle\Omega,f\rangle=0$.
Thus $Su=b_f(\Lambda)$ is solvable, with ambiguity only in $\mathbb C u_0$.
Substitution into the completed equations yields
$\langle Qf,D^{-1}Qf\rangle+\langle b_f,S^\dagger b_f\rangle$;
the kernel ambiguity contributes zero. Changing variable by $M^{1/2}$ in
this variational inverse gives \eqref{f16v43:source-top}. The transformed
source is orthogonal to $e_0$. This argument proves the quadratic inverse
identity; it does not use an invalid general congruence formula for
Moore--Penrose matrices. No physical time or new stationary probability is
introduced.
\end{proof}

\subsection{The next spectral estimate is now confined to one complete region}

The exponential estimate is no longer limited to seven character tops. It
matches every regional root separated from the lower cutoff to one native
eigenvalue in each sector, with multiplicity and a uniform error. An isolated
regional cluster of $k$ roots with a gap larger than that error would therefore
produce $8k$ physical eigenvalues. Whether $k=1$, or whether a second neutral
level approaches the top, is \emph{not} decided by this statement. Center
symmetry by itself gives neither answer.

The return-metric formula identifies a precise alternative target: bound or
asymptotically identify $g_\gamma$ on $\mathfrak h$, including all its
nonquadratic coordinates and its actual vacuum. The forbidden complement
already has the positive constant $d_\gamma\ge3\delta$, so if the regional
metric gap is small, its relative transfer to the native neutral gap loses
only $O(g_\gamma/d_\gamma)$. A proposed quartic or other local operator must
still be compared with this pencil and metric; it is not defined as the answer
by matching a classical Taylor polynomial.

Nothing here improves the common shift $\Lambda_\gamma-\lambda_*$. All bounds
fix $L$ and the separating regions before the coupling limit. At an independently
specified physical step $a_t$, the neutral energy is still
$-a_t^{-1}\log(\lambda_{0,2}/\Lambda)$. The supplied memoranda's exact-return
proposal and separate physical-transfer requirement remain distinct from the
unproved local soft spectrum, regulator removal, and continuum reconstruction.
No new numerical neutral-gap bound or continuum mass is claimed.

\subsection{Verification and preservation}

The diagnostic is
\begin{center}\small\path{verify_ym_f16_v43_sector_pencil.py}.\end{center}
It checks character isometries, a positive finite transfer with two retained
states per region, differentiated returns, rational inertia and root brackets,
the reconstruction metric, and source inverses. The finite reference is not
a discretization or enclosure of the native Wilson spectrum. Proofs and actual
execution records are reported separately.

Only the literal title version and this terminal insertion change V42.
Removing them recovers its bytes exactly. Available predecessor checkers are
retained unchanged; the audit states which were actually rerun. Native
conclusions use V42's proved-as-stated complement bounds as inherited premises;
the full earlier analytic lineage is not independently recertified. The next
scheduled companion checkpoint remains V45.

\begin{firewall}
V43 proves complete character-pencil, spectral-count and level comparisons,
and a same-top neutral-gap bound in the exact return metric. It retains
complementary sources. It does not evaluate the regional soft spectrum,
isolate exactly eight states, determine a neutral scaling law, remove the
spatial regulator, calibrate physical time, or construct a continuum mass gap.
The native operator and its physical spectral top remain unchanged.
\end{firewall}
% END F16 V43 APPEND

% BEGIN F16 V44 APPEND -- 2026-09-18
% Insert before the final document-ending command of cumulative V43.
% The predecessor, original transfer, endpoint laws and clocks are unchanged.

\clearpage
\section{V44: physical soft trial spaces and a confining quartic comparison}
\label{f16v44:context}

V43 leaves open whether its second regional level approaches the top, and its
relative neutral-gap comparison is conditional on $g_\gamma\to0$. We first
settle that qualitative question without assigning a quartic Hamiltonian to the
native operator. The lower trial in V41 proves only the largest eigenvalue:
an arbitrary fine-link trial is not a physical excited-state test. Here a
based-gauge tube, with its actual volume density, produces arbitrarily large
\emph{gauge-invariant} soft trial spaces in one central region. Their center
translates give such a space in every character sector. Every fixed level
consequently approaches the same native limiting top, including every fixed
neutral level. This also proves the previously unverified premise
$g_\gamma\to0$ in V43.

A separate calculation records the constant-field quartic model with the
original kinetic coefficient. Its commuting valleys do not destroy quantum
compactness: an explicit transverse-oscillator inequality proves compact
resolvent, also on the residual conjugation-invariant space. This is a theorem
about that declared comparison operator, not its identification with the native
regional pencil. The two conclusions have different logical status and are
kept separate throughout.

\subsection{Scope, source review, and the unchanged fixed-regulator problem}

The supplied V43 appendix and analytic audit were read completely. The native
inputs used below are V41's exact fine-link transfer and fixed-$L$ top limit,
together with the elementary normalized Wilson displacement estimates used in
V42. The regional consequences additionally use V43's proved-as-stated ladder
and return-metric comparison. No endpoint asymptotic or auxiliary sampling gap
is required. The f14 V54 local based-gauge chart and constant commutator formula
are precedents; their four-dimensional static integral is not substituted for
this three-dimensional slice transfer. The f15 V98 compact kinetic result is
likewise a different operator. Exact archive searches and those identified
passages were reviewed, not the entire archives. The next scheduled checkpoint
remains V45.

For literature context, L\"uscher's
\href{https://doi.org/10.1016/0550-3213(83)90436-4}{\emph{Some analytic results
concerning the mass spectrum of Yang--Mills gauge theories on a torus}}
addresses a small-volume low-energy expansion. Simon's
\href{https://doi.org/10.1016/0003-4916(83)90057-X}{\emph{Some quantum operators
with discrete spectrum but classically continuous spectrum}} is a classical
precedent for confinement by transverse quantum motion. Only their primary
abstracts and bibliographic scope were checked in this iteration. No expansion,
constant, or spectral theorem from either paper is imported. The elementary
trial and oscillator arguments needed here are written out. Neither the
finite-dimensional quartic model nor the general min--max principle is claimed
as a new general construction.

Fix $L=2m$, $m\ge4$, put $V=L^3=8B$, $E=3V$, and retain
\[
 T_\gamma=M_{b_\gamma}C_\gamma^{\otimes E}M_{b_\gamma},\qquad
 b_\gamma=e^{-\gamma S/2},\qquad
 S(U)=\tfrac12|\mathcal W(U)|^2,\quad
 \mathcal W(U)=(U_p-I)_p.
\]
The reference measure is probability Haar on $\mathcal X_L=\SU(2)^E$.
Let $\mathcal H_a$ be V42's physical center-character spaces, and enumerate
$T_\gamma|_{\mathcal H_a}$ in decreasing order, starting at one:
$\lambda_{a,1}\ge\lambda_{a,2}\ge\cdots$. In the neutral sector
$\lambda_{0,1}=\Lambda_\gamma$. Write
\[
 \varepsilon=\gamma^{-1/2},\qquad \lambda_*:=\Lambda_{{\rm G},B}>0.
\]
V41 supplies $\Lambda_\gamma\to\lambda_*$. All limits below fix $L$ first;
no bound uniform over growing lattice sizes or a physical time scale is given.

\subsection{A physical central tube with no division by a stabilizer}

Use real color-valued lattice cochains at the identity configuration. In the
orthonormal fine-link metric let $d_0$ be the vertex-to-link gradient and $d_1$
the link-to-plaquette curl. Set
\[
 \mathcal N=\ker d_0^*,\qquad
 \mathcal Z=\ker d_1\cap\mathcal N,\qquad
 \mathcal Y=\mathcal N\ominus\mathcal Z.
\]
The inherited periodic cochain calculation gives
\begin{equation}
 \dim\mathcal N=6V+3,\qquad \dim\mathcal Z=9,\qquad
 r:=\dim\mathcal Y=6(V-1),\qquad H_+=d_1^*d_1|_{\mathcal Y}\succ0.
 \label{f16v44:dimensions}
\end{equation}
The nine coordinates in $\mathcal Z$ are the constant three-direction,
three-color fields. They are not nine gauge-invariant coordinates individually.
The residual simultaneous adjoint action acts on them, and will be imposed on
functions rather than divided out as a translation.

\begin{proposition}[An exact based-gauge tube and its invariant half-density]
\label{f16v44:tube}
Let $\mathcal G_o=\{g\in\SU(2)^V:g_o=I\}$. For some fixed $r_0>0$ the map
\begin{equation}
 \Phi:\mathcal G_o\times\{s\in\mathcal N:|s|<r_0\}\longrightarrow\mathcal X_L,
 \qquad \Phi(g,s)=g\cdot\operatorname{Exp}(s)
 \label{f16v44:tube-map}
\end{equation}
is a diffeomorphism onto a tube about the trivial central-flat gauge orbit.
It can be chosen inside V42's selected trivial central region. Its exact
measure is
\begin{equation}
 dH(\Phi(g,s))=dH_{\mathcal G_o}(g)\,J_L(s)\,ds,
 \qquad J_L>0\ \hbox{smooth}.
 \label{f16v44:tube-measure}
\end{equation}
The density $J_L$ is invariant under simultaneous adjoint rotation of $s$.
Every compactly supported invariant $\phi\in L^2(\mathcal N,ds)$ in this ball
has an isometric physical extension
\begin{equation}
 (\mathcal E\phi)(\Phi(g,s))=J_L(s)^{-1/2}\phi(s),
 \qquad \mathcal E\phi=0\ \hbox{off the tube}.
 \label{f16v44:embedding}
\end{equation}
No value of $J_L$ is replaced by one in this definition.
\end{proposition}
\begin{proof}
The based group acts freely on a connected lattice: a transformation fixing
all links and one root is identity successively along a spanning tree. Its
orbit at the identity is a compact embedded submanifold; its tangent is
$\operatorname{Ran}d_0$. The action is isometric for the product round metric,
and its normal space at the identity is $\mathcal N$. The normal bundle along
this orbit is trivialized by $Dg\mathcal N$. The normal exponential of this
bundle is exactly \eqref{f16v44:tube-map}, because the product exponential at
the identity is the componentwise group exponential and the action is
isometric. The local inverse theorem gives invertibility near every zero
section point. Compactness of the orbit gives a common small radius and
injectivity: otherwise two distinct normal representatives tending to the
zero section would contradict a local inverse at their common limiting orbit
point. Shrink the radius to lie in the prescribed open region.

The pulled-back positive smooth volume has a positive density. Invariance
under the left action of the based group makes that density independent of
$g$ when its reference is probability Haar, proving
\eqref{f16v44:tube-measure}. Global conjugation fixes the identity, preserves
both cochain subspaces, and sends $(g,s)$ to $(kgk^{-1},\operatorname{Ad}_k s)$.
It preserves $ds$ and based Haar, proving the stated symmetry of $J_L$.
Any full gauge transformation is a based transformation followed by a global
one. Thus \eqref{f16v44:embedding} is invariant under the complete periodic
gauge group and is isometric by direct integration. The global stabilizer
has not been fixed by a singular orientation slice. At a central orbit the
based action is still free, which is the only freeness used.
\end{proof}

Let $\pi(U)=s$ be the tube coordinate modulo the based action. At the identity,
$D\pi$ is the orthogonal projection $\Pi_{\mathcal N}$, not the identity on
all link variables. For bounded $z$ and small $s,\varepsilon$, smoothness gives
\begin{equation}
 \pi(\operatorname{Exp}(s)\operatorname{Exp}(\varepsilon z))
 =s+\varepsilon\Pi_{\mathcal N}z+
 O_L(\varepsilon |s||z|+\varepsilon^2|z|^2).
 \label{f16v44:coordinate-step}
\end{equation}
The product and exponential are componentwise. This formula is local near the
identity representative; the associated gauge-invariant integral below makes
it sufficient on the whole orbit tube.

\subsection{The native transfer on arbitrarily many invariant soft tests}

Write $s=(y,z)\in\mathcal Y\oplus\mathcal Z$. The normalized hard Gaussian
transfer on $L^2(\mathcal Y,dy)$ is
\begin{equation}
 (\mathcal T_+h)(u)=(2\pi)^{-r/2}
 \int e^{-|u-u'|^2/2-u^{\mathsf T}H_+u/4-u'^{\mathsf T}H_+u'/4}
                         h(u')\,du'.
 \label{f16v44:hard-transfer}
\end{equation}
V41's matching of fine and rooted kinetic metrics gives
$\|\mathcal T_+\|=\lambda_*$. Its unique normalized positive Gaussian ground
function is invariant under global adjoint rotations. Smooth invariant compact
cutoffs approximate it in norm and in its bounded transfer quotient.

Fix any $\beta$ with $1/2<\beta<1$, put $t_\varepsilon=\varepsilon^\beta$, and
for smooth compactly supported invariant $h$ and $g$ define
\begin{equation}
 \phi_{\varepsilon,h,g}(y,z)
   =\varepsilon^{-r/2}t_\varepsilon^{-9/2}
                 h(y/\varepsilon)g(z/t_\varepsilon),\qquad
 \Psi_{\varepsilon,h,g}=\mathcal E\phi_{\varepsilon,h,g}.
 \label{f16v44:trials}
\end{equation}
For sufficiently small $\varepsilon$ their entire support is in the tube.
Their norms and inner products are exactly the product norms, before any limit.
There are arbitrarily large orthonormal families of such invariant soft $g$'s:
choose normalized radial functions on disjoint annuli of $\mathbb R^9$.

\begin{theorem}[A physical hard--soft trial limit for the original kernel]
\label{f16v44:trial-limit}
For every fixed pair of tests as above,
\begin{equation}
 \boxed{\quad
 \lim_{\gamma\to\infty}
 \langle\Psi_{\varepsilon,h,g},T_\gamma\Psi_{\varepsilon,k,f}\rangle
       =\langle h,\mathcal T_+k\rangle\langle g,f\rangle.
 \quad}
 \label{f16v44:trial-bilinear}
\end{equation}
The convergence is uniform over unit vectors in any fixed finite-dimensional
span of the chosen tests. This is a variational bilinear limit, not a claim of
norm-resolvent convergence on the full changing state space.
\end{theorem}
\begin{proof}
At $s=(\varepsilon u,t_\varepsilon v)$ on a fixed rescaled support, the
constraint map has the expansion
\begin{equation}
 \varepsilon^{-1}\mathcal W(\operatorname{Exp}s)
 =d_1u+O_L\bigl(\varepsilon+t_\varepsilon+t_\varepsilon^2/\varepsilon\bigr)
                     \longrightarrow d_1u.
 \label{f16v44:magnetic-limit}
\end{equation}
Here $d_1v=0$; Taylor's remainder is bounded on the original finite-dimensional
compact neighborhood. In particular the true multiplier tends to
$e^{-u^{\mathsf T}H_+u/4}$. No nonlinear magnetic word is deleted from the
native trial quotient.

For clarity, evaluate the inner kinetic integral as an expectation of the
actual product Wilson displacement:
$U'=U\operatorname{Exp}(\varepsilon Z)$, with its full principal product and
its original normalized density. This density converges on bounded sets to
the standard Gaussian on the full $3E$-dimensional fine-link tangent. Its
probability outside $|Z|>R$ is at most $A_Le^{-b_LR^2}$, uniformly for small
$\varepsilon$, by V42's elementary displacement estimate. The same bound is a
row and column Schur bound for the kernel truncated to these displacements.
Since both magnetic factors are at most one, the discarded bilinear form is
bounded by $A_Le^{-b_LR^2}$ times the exact test norms. This tail estimate is
independent of their shrinking supports.

On bounded $Z$, use \eqref{f16v44:coordinate-step}. Its hard coordinate after
rescaling tends to $u+Z_+$; its soft coordinate tends to $v$ because
$\varepsilon/t_\varepsilon\to0$. The coordinate errors in these two rescalings
also tend to zero. Both input and output stay in the common tube. When the
outer Haar integral is disintegrated by \eqref{f16v44:tube-measure}, the two
half-densities leave the ratio
$\sqrt{J_L(s)/J_L(s')}\to1$ in the displacement integral. Their constant
normalization is therefore accounted for, rather than omitted. The exact
powers $\varepsilon^r t_\varepsilon^9$ from $ds$ cancel the two trial
normalizations.

Orthogonal projection of the full standard Gaussian onto
$\mathcal Y\oplus\mathcal Z$ has identity covariance on that subspace. Its
vertical part and its soft displacement each integrate to one. The hard
part, with the two limiting magnetic multipliers, gives precisely
\eqref{f16v44:hard-transfer}. Dominated convergence on the bounded rescaled
supports and bounded displacements proves the claimed limit with cutoff $R$.
Then send $R\to\infty$ using the paid Schur tail. Both requirements on
$\beta$ were used: $t_\varepsilon^2/\varepsilon\to0$ for magnetism and
$\varepsilon/t_\varepsilon\to0$ for the soft displacement. Entrywise convergence
of a fixed finite matrix gives its operator-norm convergence and the final
assertion. No uniformity in the number of test functions is claimed.
\end{proof}

\subsection{Every fixed physical level approaches the top}

\begin{theorem}[An unbounded number of near-top levels in every center sector]
\label{f16v44:all-levels}
For each fixed $L$ and every fixed integer $j\ge1$, the original physical
sector eigenvalues satisfy
\begin{equation}
 \boxed{\qquad
 \lim_{\gamma\to\infty}\lambda_{a,j}(\gamma)=\lambda_*
              \quad\hbox{for every }a\in\{0,1\}^3.
 \qquad}
 \label{f16v44:sector-limits}
\end{equation}
They exist and are positive for all sufficiently large coupling. In particular,
for $j\ge2$ in the neutral sector,
\begin{equation}
 \Lambda_\gamma-\lambda_{0,j}\longrightarrow0,
 \qquad -\log(\lambda_{0,j}/\Lambda_\gamma)\longrightarrow0.
 \label{f16v44:neutral-limits}
\end{equation}
For every fixed $0<q<1$, the number of neutral eigenvalues strictly above
$q\Lambda_\gamma$ tends to infinity. This gives no rate or counting asymptotic.
\end{theorem}
\begin{proof}
Fix $N$ and $\eta>0$. Choose a normalized invariant hard cutoff $h$ with
$\langle h,\mathcal T_+h\rangle>\lambda_*-\eta$. Choose $N$ orthonormal
invariant soft tests $g_1,\ldots,g_N$. Their physical extensions
$\Psi_1,\ldots,\Psi_N$ are exactly orthonormal in the one-region tube. The
preceding theorem says that the matrix of $T_\gamma$ on their span tends to
$\langle h,\mathcal T_+h\rangle I_N$.

For a center character use the exact folded tests
\[
 \Psi_{a,i}=8^{-1/2}\sum_{\sigma\in\{\pm1\}^3}
                    \chi_a(\sigma)\mathsf Z_\sigma\Psi_i.
\]
Their eight translated supports lie in disjoint fixed central regions, so they
are orthonormal in $i$ and lie in $\mathcal H_a$. Each diagonal translated
matrix element is the same as in the trivial tube. Different regions have a
fixed positive distance; the actual separated-displacement Schur estimate makes
each cross-region matrix element $O_L(e^{-c_L\gamma})$. There are finitely many
of these for the chosen $N$. Thus the infimum of the Rayleigh quotient on this
$N$-dimensional physical sector space is at least $\lambda_*-2\eta$ eventually.

The compact self-adjoint min--max principle gives
$\liminf\lambda_{a,N}\ge\lambda_*-2\eta$. Its upper bound is
$\lambda_{a,N}\le\Lambda_\gamma\to\lambda_*$. Send $\eta\downarrow0$, with
$N$ fixed, to obtain \eqref{f16v44:sector-limits}. Positivity follows from the
positive lower bound eventually. For each prescribed count $N$, these $N$
levels exceed $q\Lambda_\gamma$ eventually; this is exactly the asserted
divergence of the count. No non-invariant trial is used for an excited level,
and no physical excited level is identified merely from equality of full and
physical Perron tops.
\end{proof}

The trial scale can, for example, be $t_\varepsilon=\varepsilon^{3/4}$.
It is not claimed to be a stationary fluctuation scale. The limits fix $j,L$
before coupling; they do not assert an estimate uniform in a growing index.
At each finite coupling compactness still makes every count above a positive
threshold finite.

\begin{proposition}[Consequences for V43's actual regional metric]
\label{f16v44:regional}
Every fixed ordered eigenvalue of V43's one-region compression tends to
$\lambda_*$. Consequently all its fixed roots $r_j$ eventually lie in V43's
matching interval. Its exact metric gap obeys
\begin{equation}
 \boxed{\qquad
 g_\gamma\longrightarrow0,\qquad
 \frac{\Delta_\gamma^{\rm neu}}{g_\gamma}\longrightarrow1,\qquad
 -\log(\lambda_{0,2}/\Lambda_\gamma)
       =\frac{g_\gamma}{\Lambda_\gamma}(1+o(1)).
 \qquad}
 \label{f16v44:metric-limit}
\end{equation}
No fixed-rank neutral truncation can leave a coupling-independent additive
complement gap at the native top.
\end{proposition}
\begin{proof}
The unfolded $N$ trial functions are supported in the chosen region and hence
in the space of $T^{\rm one}$. Their Rayleigh matrix is unchanged by that
compression. The same lower argument, and its upper norm bound
$\|T^{\rm one}\|\le\Lambda_\gamma$, prove the root limits. They are eventually
above the fixed V43 cutoff by a fixed positive margin.

The preceding theorem gives $\Delta_\gamma^{\rm neu}\to0$. V43's bracket,
with $d_\gamma\ge3\delta$, gives, once $\Delta_\gamma^{\rm neu}<d_\gamma$,
\[
 \Delta_\gamma^{\rm neu}\le g_\gamma
 \le\frac{d_\gamma\Delta_\gamma^{\rm neu}}
                {d_\gamma-\Delta_\gamma^{\rm neu}}.
\]
This proves all of \eqref{f16v44:metric-limit}, including the formerly unpaid
premise $g_\gamma\to0$. It retains the metric $I-F_0'(\Lambda_\gamma)$ and the
actual vacuum; no bare regional norm is substituted.

For any orthogonal neutral projection $R_\gamma$ of rank at most a fixed $N$,
min--max gives
$\|(I-R_\gamma)T_\gamma(I-R_\gamma)|_{\mathcal H_0}\|
\ge\lambda_{0,N+1}$. Its distance below $\Lambda_\gamma$ tends to zero. This is
why an isolated eight-state, or any fixed finite-state, global cluster cannot
have a fixed positive complement gap in this limit. The spatially supported
forbidden complement of V42 is different and its proved gap is not contradicted.
An eight-state cluster in a \emph{shrinking} spectral interval is not excluded;
no comparison of the neutral spacing with the charged exponential splitting
has been established here.
\end{proof}

\subsection{The quartic constant-field comparison: its coefficient and its limits}

The qualitative native theorem does not determine a scaling law. The following
separate calculation identifies a concrete nonquadratic candidate and proves
its own analytic properties, without asserting a native spectral approximation.
For constant fine links $U_{x,\mu}=\operatorname{Exp}(a_\mu)$, $a_\mu\in\mathbb R^3$,
write $v(a)=(\sin|a|/|a|)a$. The exact quaternion commutator identity gives
\begin{equation}
 S(\operatorname{Exp}a)=2V\sum_{\mu<\nu}|v(a_\mu)\times v(a_\nu)|^2
              =2V\mathcal Q(a)+O_L(|a|^6),\qquad
 \mathcal Q(a)=\sum_{\mu<\nu}|a_\mu\times a_\nu|^2.
 \label{f16v44:constant-action}
\end{equation}
The identity and its f14 V54 analogue are not new general quaternion formulas.
Here there are three spatial directions, nine coordinates and $V=L^3$ copies
of each direction pair, not the twelve-coordinate four-torus static integral.

The tangent metric on constant fields is $V\sum_\mu|da_\mu|^2$.
The normalized kinetic displacement has covariance $\gamma^{-1}I$ per fine
link at leading order. Its constant tangent component therefore has generator
$(2\gamma V)^{-1}\Delta_a$. Combining this principal term with the classical
quartic defines the following \emph{comparison}, not an exact restriction:
\begin{equation}
 H_{\rm cmp}(\gamma,V)
 =-\frac1{2\gamma V}\Delta_a+2\gamma V\mathcal Q(a).
 \label{f16v44:comparison-raw}
\end{equation}
The constant configurations are Haar null and are not preserved by the native
kinetic operator. Thus \eqref{f16v44:comparison-raw} is not obtained by conditioning
on them, freezing their complement, or equating a classical Hessian to V43's
spectral pencil.

Under the exact unitary dilation $a=(\gamma V)^{-1/3}q$ this comparison becomes
\begin{equation}
 H_{\rm cmp}(\gamma,V)\simeq(\gamma V)^{-1/3}H_{\rm mat},\qquad
 H_{\rm mat}=-\tfrac12\Delta_q+2\mathcal Q(q)
                \quad\hbox{on }L^2(\mathbb R^9,dq).
 \label{f16v44:model-scaling}
\end{equation}
The residual Gauss condition is simultaneous rotation of all three $q_\mu$.
It is imposed by restriction to invariant functions, not by deleting three
coordinates with an uncomputed Jacobian.

\begin{theorem}[Quantum confinement of the declared quartic model]
\label{f16v44:model-confinement}
The Friedrichs realization of $H_{\rm mat}$ has compact resolvent, a simple
strictly positive ground function and a strictly positive ground energy. In
quadratic-form sense,
\begin{equation}
 \boxed{\quad
 H_{\rm mat}\succeq2\sum_{\mu=1}^3|q_\mu|,
 \qquad
 H_{\rm mat}\succeq-\tfrac14\Delta_q+\sum_{\mu=1}^3|q_\mu|.
 \quad}
 \label{f16v44:oscillator-floor}
\end{equation}
Its simultaneous-rotation-invariant restriction has compact resolvent and
infinitely many eigenvalues $e_0<e_1\le e_2\le\cdots$, with $e_1-e_0>0$.
These are model eigenvalues, not evaluated native levels or numerical enclosures.
\end{theorem}
\begin{proof}
For each $i\in\{1,2,3\}$ define on smooth compact tests
\[
 A_i=-\tfrac14\sum_{j\ne i}\Delta_{q_j}
                        +\sum_{j\ne i}|q_i\times q_j|^2.
\]
Their sum is exactly $H_{\rm mat}$: each derivative occurs twice and each
unordered potential pair occurs twice. Hold $q_i$ fixed and put $r=|q_i|$.
For each $j\ne i$, an orthogonal rotation in the $q_j$ variable splits one
nonnegative free derivative and the two-dimensional oscillator
$-\tfrac14\Delta_y+r^2|y|^2$. Its lower form bound is $r$, including $r=0$.
For example expand the two nonnegative squares
$\|\tfrac12\partial_{y_k}f+r y_k f\|^2$ and integrate their cross terms.
Thus $A_i\succeq2|q_i|$. No derivative in $q_i$ is taken in this disintegration,
so the parameter-dependent rotation adds no derivative term. Sum the three
bounds to obtain the first inequality. Average it with
$H_{\rm mat}\succeq-\tfrac12\Delta_q$ to obtain the second.

The closed form therefore controls a local $H^1$ norm and
$\int |q|\,|f(q)|^2dq$. A bounded form sequence has uniformly small $L^2$ mass
outside large balls, and Rellich compactness gives convergence on each ball.
A diagonal extraction and the tail bound make the full form embedding compact;
this is equivalent to compact resolvent. Zero energy would force a constant
$L^2(\mathbb R^9)$ function and hence the zero vector, so the attained ground
energy is positive. The heat kernel is strictly positive: along any finite
Brownian bridge the continuous polynomial potential has finite integral, so
its Feynman--Kac weight is strictly positive. Positivity improvement makes the
ground eigenvalue simple with a positive eigenfunction. Rotational invariance
and uniqueness force that function to be invariant. Restriction to the closed
invariant subspace preserves compact resolvent. That subspace is infinite
dimensional (it contains all radial smooth compact tests), so its eigenvalues
form an infinite discrete sequence; simplicity of the ground gives its positive
first invariant gap. This proves properties of the specified model only.
\end{proof}

There is a coefficient check against V41's cusp, not a second potential to add.
At a commuting rank-one triple $a_\mu=c_\mu n$, $|n|=1$, the Hessian of
$2\mathcal Q$ has four positive eigenvalues $4|c|^2$. The transverse harmonic
ground contribution for $-\frac12\Delta+2\mathcal Q$ is therefore $4|c|$.
V41's flat-threshold variation has precisely this leading term when
$\varphi=2L c$:
$2|\varphi|/L=4|c|$. Keeping the full quartic coordinates and also adding that
cusp as a separate potential would count those transverse modes twice.
The model still has classical commuting valleys; the proved form bound shows
why classical flatness alone is not a test for its quantum spectrum.

\subsection{The remaining native-to-model estimate}

The new native result settles the qualitative regional alternative:
there are arbitrarily many near-top physical levels in every sector, and the
neutral return-metric gap tends to zero. It does not determine how rapidly.
The model calculation makes $(\gamma V)^{-1/3}$ a natural \emph{candidate}
scale, but this dilation is not a theorem about the native neutral gap.

For an actual scaling limit one must compare the complete regional metric
operator of V43, after its actual vacuum subtraction, with
$H_{\rm mat}-e_0$ (or a corrected operator if elimination produces leading
terms). This needs an upper trial comparison at the proposed scale and a
matching lower form estimate with compactness of low-energy states. The
nonzero-mode response, the true quotient density, the energy-dependent return,
and its metric must all be retained. In particular a lower bound on a discarded
restriction does not control the norm of a reconstructed full vector. The
forbidden outer region already has its V42 inverse; it is not the missing
within-region hard/soft comparison.

The intermediate trial width used above satisfies
$\varepsilon\ll t_\varepsilon\ll\sqrt\varepsilon$ so its nonlinear magnetic
cost simply tends to zero. At the candidate balanced scale those terms are
\emph{of the same order} as the soft kinetic deficit; the present bilinear
limit cannot be divided by that smaller scale. No big-O remainder at that
scale, tunneling prefactor, rate of the common top shift, volume-uniform bound,
or continuum time is supplied. With any independently specified $a_t(\gamma)$,
the physical neutral energy remains
$-a_t(\gamma)^{-1}\log(\lambda_{0,2}/\Lambda_\gamma)$. A vanishing uncalibrated
one-step deficit does not decide its continuum limit.

\subsection{Verification and preservation}

The new diagnostic is
\begin{center}\small\path{verify_ym_f16_v44_physical_soft_spaces.py}.\end{center}
It checks finite periodic cochains and harmonic dimensions, the exact gauge
normal/tangent decomposition, noncommutative constant plaquettes, character
isometries, hard--soft Gaussian trial limits, the oscillator allocation and
scaling, and the algebra taking V43's bracket to its now established limit.
Finite fixtures do not certify a tubular-neighborhood argument, a compact
Wilson spectral limit, or a native-to-matrix-model approximation. Actual
execution records and scope are reported separately in the audit.

Only the title version and this terminal insertion change V43. Removing those
changes recovers the predecessor byte for byte. The main native conclusion
inherits V41's fixed-regulator top limit; the invariant trial construction is
proved here. The earlier analytic lineage is not independently recertified.
The next scheduled companion checkpoint remains V45.

\begin{firewall}
V44 proves that every fixed physical level, in every center sector, approaches
the same limiting native top at fixed spatial regulator. It establishes
$g_\gamma\to0$ for V43's actual neutral return metric and rules out a fixed-rank
near-top reduction with a uniform complementary gap. Separately it proves
compact resolvent and a positive invariant excitation gap for an explicitly
defined quartic comparison operator. No native scaling law is inferred from
that model, no original transfer is replaced, and no physical time, unrestricted
regulator limit, or interacting continuum mass-gap conclusion is supplied.
\end{firewall}
% END F16 V44 APPEND

% BEGIN F16 V45 APPEND -- 2026-09-18
% Insert immediately before the final document-ending command of V44.
% No predecessor assertion, native measure, scalar, or clock is edited.

\clearpage
\section{V45: critical-scale recovery and corrected physical quasimodes}
\label{f16v45:context}

V44 proves qualitative collapse of every fixed physical level, but its trial
error is only $o(1)$. It cannot be divided by the proposed quartic scale. We
now expand the \emph{actual} transfer on that scale. A term linear in the soft
field changes the hard Gaussian operator. Its ground expectation vanishes by
color symmetry, but the term need not vanish as a vector. A fixed hard-band
inverse corrects that vector defect. This gives physical quasimodes with error
smaller than the quartic scale, not just a favorable Rayleigh quotient.

The result is spectral inclusion, with multiplicities, and a one-sided ordered
level bound. It does not prove that the included levels exhaust the native low
spectrum. In particular it does not identify the native Perron correction or
the first neutral spacing with the two lowest quartic eigenvalues. The missing
lower estimate and low-energy compactness are stated at the V45 checkpoint.

\subsection{Checkpoint, dependencies, and the balanced variables}

The supplied V44 appendix and audit were read in full. We retain its physical
tube, exact density, hard Mehler transfer, and separately proved quartic model.
The f14 V54 constant-field stationarity argument was checked; its different
four-dimensional static integral is not imported. The Grand Tensor's signed
compression and return packet were also checked. Searches of the active source,
f14 V100, f15 V100 and Grand Tensor V076 are recorded separately from proof
verification. No archive-wide novelty claim is made. The next checkpoint is V50.

The correction principle is classical. Panati--Spohn--Teufel,
\href{https://arxiv.org/html/math-ph/0201055}{\emph{Space-Adiabatic Perturbation
Theory}}, separates an isolated fast band from the effective slow dynamics.
M\'atyus--Teufel,
\href{https://arxiv.org/html/1904.00042}{\emph{Effective non-adiabatic
Hamiltonians}}, Sections III--IV, explains why a bare diagonal compression
retains an off-diagonal defect and how a reduced inverse corrects it. Their
whole-space Hamiltonian assumptions are not applied to our compact integral
operator. The needed one-step calculation is given below. L\"uscher's torus
expansion remains relevant context, but no continuum coefficient or time
normalization is substituted for the present Wilson convention.

Fix $L=2m$, $m\ge4$, and $V=L^3$. All constants in this section may depend on
$L$. Retain
\[
 T_\gamma=M_{b_\gamma}C_\gamma^{\otimes3V}M_{b_\gamma},\qquad
 b_\gamma=e^{-\gamma S/2},\qquad \lambda_* =\Lambda_{{\rm G},B}>0.
\]
The Hilbert space, Haar measure, gauge action and eight center sectors are
unchanged. Use V44's orthogonal decomposition
$\mathcal N=\mathcal Y\oplus\mathcal Z$, with $r=6(V-1)$ hard coordinates and
nine constant soft coordinates. Write $\mathsf K:\mathbb R^9\to\mathcal Z$ for
constant extension to all fine links, so
\[
 (\mathsf K a)_{x,\mu}=a_\mu,\qquad \mathsf K^*\mathsf K=V I_9.
\]
The normalized isometry $\mathsf K/\sqrt V$ identifies $\mathbb R^9$ with the
orthonormal soft coordinates. Define
\begin{equation}
 h=(\gamma V)^{-1/3},\qquad
 \varepsilon=\gamma^{-1/2}=\sqrt V\,h^{3/2},\qquad
 \tau=\sqrt V\,h,\qquad
 s=\varepsilon u+h\mathsf K q.
 \label{f16v45:scales}
\end{equation}
Thus $a=hq$ is the constant link angle and $z=\tau q$ its orthonormal coordinate.
In particular
\begin{equation}
 \varepsilon^2/\tau^2=h,\qquad \gamma Vh^4=h,
 \qquad \varepsilon=o(h),\quad h^2=o(h).
 \label{f16v45:balance}
\end{equation}
The symbol $h$ is an asymptotic parameter, not physical time.

Let $\mathcal T_+$ be exactly \eqref{f16v44:hard-transfer}, and let
$\varphi>0$ be its unit Gaussian ground function. Put $P_+=|\varphi\rangle
\langle\varphi|$. At this fixed lattice,
\begin{equation}
 \mathcal T_+\varphi=\lambda_*\varphi,\qquad
 R_+=(\lambda_* -\mathcal T_+)^{-1}(I-P_+),\qquad
 \|R_+\|\le[\lambda_*(1-e^{-\xi_{\min}})]^{-1}<\infty.
 \label{f16v45:hard-inverse}
\end{equation}
Here $\xi_{\min}>0$ is the smallest positive hard Mehler frequency. This inverse
is only the fixed Gaussian hard-band inverse; it is not an inverse on a native
neutral complement or on the nine soft coordinates.

\subsection{The two first variations that must be retained}

Use V44's exact tube coordinate $\pi(U)=s$ modulo the based gauge group and its
positive density $J_L(s)$. For small $a\in\mathbb R^9$ set
\[
 H(a)=D_y^2S(\operatorname{Exp}(\mathsf K a+y))\big|_{y=0,\ y\in\mathcal Y},
 \qquad
 R(a)Z=\left.\frac d{dt}\pi(\operatorname{Exp}(\mathsf K a)
                                      \operatorname{Exp}(tZ))\right|_{t=0}.
\]
Let $R_Y(a)=\Pi_{\mathcal Y}R(a)$ and $A(a)=R_Y(a)R_Y(a)^*$. These are smooth
finite-dimensional matrices on a fixed neighborhood. They obey
$H(0)=H_+$ and $A(0)=I_{\mathcal Y}$, and are positive for a sufficiently small
neighborhood. Define their directional derivatives
\[
 H_1(q)=DH(0)[q],\qquad A_1(q)=DA(0)[q].
\]
No quotient covariance is assigned the value one away from $a=0$.

\begin{proposition}[Exact hard stationarity and the normalized hard derivative]
\label{f16v45:jets}
For every small constant field $a$, including noncommuting triples,
\begin{equation}
 D_yS(\operatorname{Exp}(\mathsf K a+y))\big|_{y=0}=0
                 \quad\hbox{on }\mathcal Y.
 \label{f16v45:stationarity}
\end{equation}
Consequently, on fixed rescaled compact sets,
\begin{equation}
 \frac{S(\operatorname{Exp}(\varepsilon u+h\mathsf Kq))}{2\varepsilon^2}
 =\frac14u^{\mathsf T}H_+u+
 h\left\{\mathcal Q(q)+\frac14u^{\mathsf T}H_1(q)u\right\}
 +O_{L}(\varepsilon P(u,q)+h^2P(u,q)),
 \label{f16v45:magnetic-expansion}
\end{equation}
where $P$ is a fixed polynomial bound and
$\mathcal Q(q)=\sum_{\mu<\nu}|q_\mu\times q_\nu|^2$.

For the frozen, normalized hard Gaussian operator
\[
 \mathcal T(a)(u,v)=
 \frac{\exp[-(u-v)^{\mathsf T}A(a)^{-1}(u-v)/2
                 -u^{\mathsf T}H(a)u/4-v^{\mathsf T}H(a)v/4]}
 {(2\pi)^{r/2}\sqrt{\det A(a)}},
\]
its derivative $\mathcal C(q)=D\mathcal T(0)[q]$ has kernel
\begin{equation}
 \begin{split}
 \mathcal C(q)(u,v)=\mathcal T_+(u,v)\bigl\{
 &-\tfrac12\Tr A_1(q)+\tfrac12(u-v)^{\mathsf T}A_1(q)(u-v)\\
 &-\tfrac14u^{\mathsf T}H_1(q)u-\tfrac14v^{\mathsf T}H_1(q)v\bigr\}.
 \end{split}
 \label{f16v45:hard-derivative}
\end{equation}
It is a bounded self-adjoint Hilbert--Schmidt operator, linear in $q$, and
\begin{equation}
 \langle\varphi,\mathcal C(q)\varphi\rangle=0.
 \label{f16v45:singlet-cancellation}
\end{equation}
\end{proposition}
\begin{proof}
Translation invariance makes the logarithmic first derivative of $S$ at a
constant field independent of the base vertex, in each direction and color.
Every $y\in\mathcal Y$ has zero mean in each such component, so its pairing
with that derivative vanishes. This proves \eqref{f16v45:stationarity} exactly;
it is the same stationarity mechanism as f14 V54, now on the spatial slice.
Taylor expansion in $y$ therefore starts with its Hessian, not a linear hard
force. Its cubic remainder is bounded by $C_L|y|^3$ on the fixed chart. V44's
exact constant-field formula gives
$S(\operatorname{Exp}(\mathsf K a))=2V\mathcal Q(a)+O_L(|a|^6)$.
Expand $H(a)=H_++H_1(a)+O_L(|a|^2)$ and use
\eqref{f16v45:balance}. The constant sixth-order contribution is $O_L(h^3)$.
These facts give \eqref{f16v45:magnetic-expansion}.

Differentiate the complete Gaussian kernel, including its determinant. This
gives \eqref{f16v45:hard-derivative}; Gaussian domination proves the claimed
operator differentiability and boundedness. For a simultaneous color rotation
$k$, let $O_k$ denote its orthogonal action on $\mathcal Y$. Equivariance of the
tube and action gives $A_1(kq)=O_kA_1(q)O_k^*$ and the same formula for $H_1$.
The ground function $\varphi$ is invariant. Hence the scalar
$q\mapsto\langle\varphi,\mathcal C(q)\varphi\rangle$ is a rotation-invariant
linear functional on three copies of $\mathbb R^3$. It is zero, since none of
those vector representations has a fixed nonzero covector. This proves
\eqref{f16v45:singlet-cancellation}. It does not assert
$\mathcal C(q)\varphi=0$.
\end{proof}

The determinant term $-\Tr A_1/2$ is retained even when a choice of normal
coordinates makes $A_1$ vanish. Likewise the mean-zero stationarity is essential:
without it a constant-field quadratic hard force could appear on a larger scale.
Neither cancellation follows merely from smallness of the original action.

\subsection{A strong local expansion at the balanced scale}

Choose a smooth invariant cutoff $\chi(y)$, equal to one near zero, whose hard
support and a small fixed soft ball have product compactly contained in V44's
tube and in V42's central region $\mathcal O_e$. For a function $F(u,q)$ smooth
and compactly supported in $q$, with a polynomial times the hard ground Gaussian
in $u$, define
\begin{equation}
 (\mathcal U_\gamma F)(\Phi(g,y+z))=
 J_L(y+z)^{-1/2}\varepsilon^{-r/2}\tau^{-9/2}
 \chi(y)F(y/\varepsilon,z/\tau),
 \label{f16v45:dilation}
\end{equation}
with zero extension. The formula uses the orthonormal identification of
$\mathcal Z$ specified above. For small $h$ it is supported inside the original
tube. If $F(O_ku,kq)=F(u,q)$ it is a physical vector. Its exact Gram is
\begin{equation}
 \langle\mathcal U_\gamma F,\mathcal U_\gamma G\rangle
       =\int\chi(\varepsilon u)^2\overline{F(u,q)}G(u,q)\,du\,dq.
 \label{f16v45:Gram}
\end{equation}
Thus on any fixed finite family of these functions it differs from the full
product Gram by $O(e^{-c/\varepsilon^2})$. No density factor is dropped.

\begin{theorem}[The complete one-step expansion on the test core]
\label{f16v45:strong-expansion}
For every such $F$, with $\mathcal T_+$ and $\mathcal C(q)$ acting only in $u$,
\begin{equation}
 \begin{split}
 T_\gamma\mathcal U_\gamma F
 =\mathcal U_\gamma\bigl[\mathcal T_+F+h\{\tfrac12\mathcal T_+\Delta_qF
       -2\mathcal Q(q)\mathcal T_+F+\mathcal C(q)F\}\bigr]
       +o_{L,F}(h)
 \end{split}
 \label{f16v45:strong-local}
\end{equation}
in the original physical $L^2$ norm for invariant $F$, and in the based-invariant
space otherwise. Convergence is uniform on a fixed finite-dimensional span of
test functions. This is a core-wise expansion, not operator-norm convergence on
an entire changing Hilbert space.
\end{theorem}
\begin{proof}
We give the error accounting needed beyond V44's $o(1)$ limit. In the actual
kinetic integral write $U'=U\operatorname{Exp}(\varepsilon Z)$. Its normalized
scaled density is the standard full-link Gaussian density times
$1+O_L(\varepsilon^2P(Z))$ on logarithmically growing $Z$-balls. Its raw tail has
both row and column Schur bounds $C_Le^{-c_LR^2}$ outside $|Z|\le R$.
These follow directly from the original Wilson cosine and Haar factors and
V37's normalization; no heat kernel is substituted for the native kernel.

Write $s=\varepsilon u+h\mathsf Kq$. Smoothness of the tube coordinate gives,
uniformly on a fixed neighborhood,
\[
 s'=s+\varepsilon R(hq)Z+
 O_L(\varepsilon^2(1+|u|+|Z|)^2).
\]
Here the additional $\varepsilon u$ changes the differential by $O_L(\varepsilon
|u|)$, and is included in the remainder. In hard rescaled coordinates this is
\[
 u'=u+R_Y(hq)Z+O_L(\varepsilon P(u,Z)).
\]
In orthonormal soft coordinates it gives
\[
 q'=q+\sqrt h\,Z_0+O_L(h^{3/2}P(u,q,Z)),
 \qquad Z_0=(\mathsf K/\sqrt V)^*\Pi_{\mathcal Z}Z.
\]
At $a=0$ the random variables $Z_0$ and $\Pi_{\mathcal Y}Z$ are independent
standard Gaussians; all based-vertical components are integrated as well.
The change of their covariance at $a=hq$ is $O_L(h)$, which affects a first
soft derivative only at order $h^{3/2}$. The second differential of $\pi$ gives
an even smaller soft error $O_L(\varepsilon^2/\tau)=O_L(h^2)$.

After the outer Haar integral and the two physical half-densities have been
assembled, their exact ratio is $\sqrt{J_L(s)/J_L(s')}$. Since
$s'-s=O_L(\varepsilon |Z|)$ and $J_L$ is smooth positive, its difference from
one is $O_L(\varepsilon P(u,q,Z))=o(h)$ on the sets used below. This is why no
separate order-$h$ Jacobian potential occurs in this flat half-density
representation. It is not an assertion that the quotient Haar density is flat.

Taylor-expand the soft test through order two. Its first displacement has zero
Gaussian mean; its second gives $(h/2)\Delta_qF$. At this order the hard integral
is exactly $\mathcal T_+$. At both ends the constant magnetic multiplier is
$1-h\mathcal Q(q)+o(h)$, since the difference between $q'$ and $q$ changes this
term only at order $h^{3/2}$. The two factors give
$-2h\mathcal Q\mathcal T_+F$. Freeze the remaining hard displacement at $a=hq$.
Its hard marginal is the normalized Gaussian of covariance $A(hq)$, including
$\det A(hq)^{-1/2}$. Together with the two hard magnetic factors of
\eqref{f16v45:magnetic-expansion}, its first variation is exactly
$h\mathcal C(q)F$. No term proportional to $h^{1/2}$ survives. This yields the
three displayed terms, with all normalizations still included.

Here is a norm justification of the truncations. A polynomial times the hard
Gaussian has $L^2$ tails, and finitely many derivative tails, bounded by
$C_Fe^{-c_FR^2}$. Cut that input and the raw displacement to radius $R$.
Contraction of $T_\gamma$ pays the first omitted input; the raw row/column Schur
bound pays the second. A fixed physical cutoff boundary is separated from the
remaining input and contributes $O(e^{-c/\varepsilon^2})$. On the retained
pieces the output has hard coordinate of size $O_L(R)$ and soft coordinate in
a fixed compact enlargement of the test support. Taylor remainders, including
three soft derivatives, the hard cubic, smooth coordinate errors and the second
frozen hard variation, have $L^2$ norm at most
\[
 C_{L,F}(h^{3/2}+\varepsilon+h^2)(1+R)^{N_F}.
\]
The finite exponent also includes the volume of this rescaled output set.
All intermediate Gaussian integrals have uniformly positive hard quadratic
parts. Choose $R=(K_F\log(1/h))^{1/2}$ with fixed $K_F$ large enough that both
omitted tails are $o(h)$. Equation \eqref{f16v45:balance} makes the last display
$o(h)$. The native principal cut lies outside these sets and is part of the
same paid displacement or input tail. This proves a strong $L^2$ expansion,
not merely pointwise convergence. For a fixed finite family the constants and
truncation can be chosen in common, proving uniformity on its span.
\end{proof}

\subsection{Correcting the hard-band leakage rather than discarding it}

Gaussian integration in \eqref{f16v45:hard-derivative} shows that
$\mathcal C(q)\varphi$ is a hard polynomial of degree at most two times
$\varphi$, linear in $q$. Equation \eqref{f16v45:singlet-cancellation} removes
its ground component. The finite Hermite space of degrees at most two is
preserved by $\mathcal T_+$, so the following corrector has the same regularity:
\begin{equation}
 k(q)=R_+\mathcal C(q)\varphi,
 \qquad \langle\varphi,k(q)\rangle=0,
 \qquad (\lambda_* -\mathcal T_+)k(q)=\mathcal C(q)\varphi.
 \label{f16v45:corrector}
\end{equation}
It is equivariant under simultaneous color rotations. For an invariant
$f\in C_c^\infty(\mathbb R^9)$ put
\begin{equation}
 \Psi_\gamma f=\mathcal U_\gamma[(\varphi+h k(q))f(q)],
 \qquad H_{\rm mat}=-\tfrac12\Delta_q+2\mathcal Q(q).
 \label{f16v45:corrected-embedding}
\end{equation}
No native complement inverse is assumed in this definition.

\begin{theorem}[Corrected physical recovery at the quartic scale]
\label{f16v45:recovery}
For invariant compact smooth $f,g$,
\begin{align}
 \langle\Psi_\gamma f,\Psi_\gamma g\rangle
 &=\langle f,g\rangle+
 h^2\int\overline{f(q)}g(q)\|k(q)\|_2^2\,dq
                       +O_{L,f,g}(e^{-c/\varepsilon^2}),
 \label{f16v45:corrected-Gram}\\
 \big\|T_\gamma\Psi_\gamma f-
       \lambda_*\Psi_\gamma(f-hH_{\rm mat}f)\big\|_2&=o_{L,f}(h).
 \label{f16v45:intertwining}
\end{align}
Both conclusions are uniform on every fixed finite test space. In particular
\begin{equation}
 \frac{\lambda_*\langle\Psi_\gamma f,\Psi_\gamma g\rangle
          -\langle\Psi_\gamma f,T_\gamma\Psi_\gamma g\rangle}
 {\lambda_*h}
       \longrightarrow\langle f,H_{\rm mat}g\rangle.
 \label{f16v45:form-recovery}
\end{equation}
The same statements hold in each center sector after the exact character
folding of V44.
\end{theorem}
\begin{proof}
The Gram follows from \eqref{f16v45:Gram} and the pointwise orthogonality of
$k(q)$ to $\varphi$. Its order-$h$ cross term is exactly zero before a limit.
Apply Theorem~\ref{f16v45:strong-expansion} to $\varphi f$. Since
$\mathcal T_+\varphi=\lambda_*\varphi$, its order-$h$ term is
$-\lambda_*\varphi H_{\rm mat}f+\mathcal C(q)\varphi f$.
Apply the zeroth-order part of the same theorem to $k(q)f$ and multiply by $h$.
The identity in \eqref{f16v45:corrector} cancels the entire vector
$\mathcal C(q)\varphi f$, not just its expectation. The extra term from
$h\,k(q)H_{\rm mat}f$ on the right side of \eqref{f16v45:intertwining} has
coefficient $h^2$ and is $o(h)$. This proves the strong residual estimate and
then its bilinear consequence.

Every constructed vector is physical by the exact equivariant embedding.
Choose the tube closure inside a smaller neighborhood of the central orbit.
Its eight translates have a fixed positive separation. Folding by
$8^{-1/2}\sum_\sigma\chi_a(\sigma)\mathsf Z_\sigma$ is an exact isometry on
this supported space. The same residual estimate holds by commutation; cross
matrix elements between distinct translates are $O_L(e^{-c_L\gamma})$ by the
native displacement estimate. No regional trial state is substituted into an
intermediate transfer power. No conclusion about an almost invariant band for
all native low-energy states is made.
\end{proof}

\paragraph{A finite check on the role of the correction.}
Take two isotropic three-component hard oscillators with distinct Mehler ratios
$r_1,r_2\in(0,1)$, and let $p_q(u)=u^{\mathsf T}H_1(q)u$ be a rotation-covariant
cross quadratic proportional to $q\cdot(u_1\times u_2)$. With $A_1=0$,
\[
 \mathcal C(q)\varphi
   =-\frac{\lambda_*}{4}(1+r_1r_2)p_q\varphi,
 \qquad
 k(q)=-\frac{1+r_1r_2}{4(1-r_1r_2)}p_q\varphi.
\]
The expectation of the first vector is zero, but its norm is generally nonzero.
This example checks the sign and denominator in \eqref{f16v45:corrector}; it is
not a numerical evaluation of the native derivative matrices. A bare Rayleigh
calculation would miss the uncancelled order-$h$ residual.

\subsection{Native spectral inclusion and its precise one-sided ordering}

Write $e_0<e_1\le e_2\le\cdots$ for the invariant eigenvalues of V44's same
quartic model, repeated with multiplicity. They are not numerically evaluated
here. A compactly supported invariant smooth graph core is available: for a
model eigenfunction, multiplication by radial cutoffs converges in graph norm
because $\nabla f\in L^2$ and
\[
 [H_{\rm mat},\chi_R]f=-\nabla\chi_R\cdot\nabla f
                              -\tfrac12(\Delta\chi_R)f\longrightarrow0.
\]
Local elliptic regularity and mollification on the compact support finish the
approximation. Averaging over simultaneous rotations preserves it and the
operator. This proof also treats a finite eigenspace in one common graph core.

\begin{theorem}[Every fixed quartic level is realized to smaller than its scale]
\label{f16v45:inclusion}
Fix $L$, a model eigenvalue $e$ of invariant multiplicity $d_e$, and a center
character $a$. There is $\eta_\gamma\to0$ such that the original sector
$T_\gamma|_{\mathcal H_a}$ has at least $d_e$ eigenvalues, counted with
multiplicity, in
\begin{equation}
 [\,\lambda_*(1-he)-h\eta_\gamma,
                  \lambda_*(1-he)+h\eta_\gamma\,].
 \label{f16v45:spectral-window}
\end{equation}
One may choose a common sequence for any fixed finite list of model levels and
all eight sectors. The corresponding eigenvalues lie in V43's near-top interval
for sufficiently large coupling. They need not be the first $d_e$ or any
specified ordered eigenvalues of their native sector.
\end{theorem}
\begin{proof}
Approximate an orthonormal model eigenbasis at $e$ by invariant compact smooth
tests in graph norm. For each fixed approximation Theorem~\ref{f16v45:recovery}
gives its native residual divided by $h$, bounded by the model graph error plus
a term tending to zero. Choose the approximation index to increase sufficiently
slowly with coupling. This diagonal choice is also required to keep the physical
supports inside the common tube. The resulting native Gram tends to identity;
orthonormalizing on this fixed-dimensional space preserves a residual bound
$o(h)$ uniformly on its unit sphere. Thus there is $\zeta_\gamma\to0$ with
\[
 \|(T_\gamma-\lambda_*(1-he))u\|\le h\zeta_\gamma\|u\|
\]
on a $d_e$-dimensional physical sector space. The spectral theorem bounds the
norm outside a window of radius $h\sqrt{\zeta_\gamma}$ by
$\sqrt{\zeta_\gamma}\|u\|$. Its spectral projection is injective on that space
for large coupling, proving the multiplicity assertion. A zero residual needs
no enlargement; increasing $\eta_\gamma$ slightly handles both cases. The
windows are eventually positive and above the fixed cutoff $c<\lambda_*$, so
the compact operator has only discrete eigenvalues there. A common finite
maximum of errors and thresholds handles finitely many levels and characters.
No assertion of spectral exhaustion enters this argument.
\end{proof}

\begin{theorem}[One-sided ordered level estimate at the balanced scale]
\label{f16v45:minmax}
With indexing starting at one in each native center sector, for every fixed $j$,
\begin{equation}
 \boxed{\quad
 \limsup_{\gamma\to\infty}
     \frac{1-\lambda_{a,j}(\gamma)/\lambda_*}{h}
                  \le e_{j-1}.
 \quad}
 \label{f16v45:ordered-upper}
\end{equation}
The same statement holds for V43's ordered one-region eigenvalues $r_j$.
In particular
\begin{equation}
 \Lambda_\gamma\ge\lambda_*[1-h(e_0+o(1))],\qquad
 \limsup_{\gamma\to\infty}
     \frac{\log(\lambda_*/\Lambda_\gamma)}{h}\le e_0.
 \label{f16v45:top-one-side}
\end{equation}
These expressions are not assumed nonnegative. No upper estimate for
$\Lambda_\gamma-\lambda_*$ is proved.
\end{theorem}
\begin{proof}
Take a $j$-dimensional model graph-core space whose maximum Rayleigh quotient
is at most $e_{j-1}+\eta$. Equation \eqref{f16v45:form-recovery} and the paid
Gram give a physical sector space on which the native transfer quotient is at
least $\lambda_*[1-h(e_{j-1}+\eta)+o(h)]$. Min--max proves the displayed
limsup after $\eta\downarrow0$. The unfurled space is supported in the chosen
region and hence has the same quadratic form in the one-region compression.
This proves its assertion too. Take $j=1$ and the neutral character, and use
$-\log(1-hc)=hc+O(h^2)$ on the lower bound, for the final inequality. This
argument uses a lower bound on the top, not convergence of its first correction.
\end{proof}

\begin{proposition}[Explicit variational constants without a model diagonalization]
\label{f16v45:explicit-bounds}
For the invariant model eigenvalues and the native neutral ordered levels,
\begin{equation}
 e_0\le\frac{27}{4},\qquad e_1\le\frac{47}{4},\qquad
 \limsup_{\gamma\to\infty}\frac{1-\Lambda_\gamma/\lambda_*}{h}
       \le\frac{27}{4},\qquad
 \limsup_{\gamma\to\infty}\frac{1-\lambda_{0,2}/\lambda_*}{h}
       \le\frac{47}{4}.
 \label{f16v45:explicit-values}
\end{equation}
These are variational upper bounds on deficits relative to $\lambda_*$, not an
upper bound $5h$ on the actual neutral gap and not evaluated model eigenvalues.
\end{proposition}
\begin{proof}
Let $d=9$, $r^2=|q|^2$, and
\[
 \phi_0=(2/\pi)^{9/4}e^{-r^2},\qquad
 \phi_1=\frac{2\sqrt2}{3}(r^2-9/4)\phi_0.
\]
They are orthonormal, invariant, and in the model form domain. Under
$\phi_0^2dq$, the coordinates are independent with variance $1/4$, and $r^2$
has the gamma moments with shape $9/2$ and rate $2$. Direct integration gives
\[
 \big(\langle\phi_i,-\tfrac12\Delta\phi_j\rangle\big)_{i,j=0}^1
 =\begin{pmatrix}9/2&-3/\sqrt2\\-3/\sqrt2&13/2\end{pmatrix},\qquad
 \big(\langle\phi_i,2\mathcal Q\phi_j\rangle\big)_{i,j=0}^1
 =\begin{pmatrix}9/4&3/\sqrt2\\3/\sqrt2&21/4\end{pmatrix}.
\]
For the potential moments, homogeneity of $\mathcal Q$ separates the angular
factor, and the remaining radial moments are those just specified; its base
expectation is $9/4$. The summed matrix is
$\operatorname{diag}(27/4,47/4)$. Model min--max proves the first two bounds.
Theorem~\ref{f16v45:minmax} gives the native inequalities. Subtracting two
one-sided ordered bounds would not bound their difference, which explains the
qualification in the statement.
\end{proof}

\subsection{The existing outer return is negligible on these packets only}

Let $Q$ be V42's forbidden complement. The supports of the newly constructed
packets lie in a fixed compact subset of the union of retained regions. The
same is true for the slowly varying graph-core approximants above, by their
diagonal choice. A fixed distance $d_0>0$ therefore separates their support
from $Q$. The original displacement estimate gives
\begin{equation}
 \|QT_\gamma\Psi_{\gamma,a}f\|
                \le C_Le^{-c_L\gamma}\|\Psi_{\gamma,a}f\|.
 \label{f16v45:outer-leakage}
\end{equation}
For $u=W_a^*\Psi_{\gamma,a}f$ and $z\ge c$, V43's exact reconstruction then
satisfies
\begin{equation}
 \|J_a(z)u-W_a u\|+\|(I-F_a'(z))u-u\|
                   \le C_Le^{-c_L\gamma}\|u\|.
 \label{f16v45:outer-metric}
\end{equation}
Indeed the first term contains one bounded complement inverse and the leakage
in \eqref{f16v45:outer-leakage}; the second contains its actual inverse-square
return. This proves that the \emph{outer} metric does not alter the leading
coefficient on these recovery vectors. It is not operator-norm convergence of
that metric to identity on the whole regional space, and does not remove
within-region hard-band correction \eqref{f16v45:corrector}.

\subsection{V45 checkpoint: spectral inclusion is not a complete scaling theorem}

The new estimates establish the native upper-recovery half at its proposed
scale, including a corrected vector residual. They turn the comparison-model
levels into actual native spectral points with $o(h)$ errors and supply ordered
one-sided bounds. They do not prove the converse inclusion or exclude additional
near-top states with a different concentration scale. The all-level $o(1)$
limit of V44 alone cannot provide that exclusion.

A simple compact-operator example shows the logical distinction. On
$\ell^2\oplus\mathbb C$ let a positive contraction have eigenvalues
$\lambda_*e^{-h e_j}$ on the model coordinates, but eigenvalue
$\lambda_*(1+\sqrt h)$ on the final coordinate, for small $h$ and fixed
$\lambda_*<1$. Its model coordinate vectors have all the recovery and
spectral-inclusion properties above, and every fixed ordered eigenvalue tends
to $\lambda_*$. Nevertheless its top correction is of order $\sqrt h$, not
$h$, and the lowest ordered gap is not given by $e_1-e_0$. This is not a native
counterexample; it rules out inferring exhaustion or vacuum subtraction from
the proved one-sided data. Center copies can be added without changing the
point of the example.

A sufficient next native target is a lower bound on the rescaled deficit and
compactness of every sequence with bounded deficit after centering at
$\lambda_*$. In tube variables this must pay: tightness of the soft coordinate
$q$ along the commuting valleys, exclusion of hard excited mass on the same
scale, and all regions or coordinate tails not represented by the recovery
map. V44's coercive model inequality is not a substitute for those estimates
for the changing native operators. The exact regional Feshbach map and its
return metric remain the correct objects for such a proof.

Only after that completeness step could one identify
$\Lambda_\gamma=\lambda_*[1-he_0+o(h)]$ and then the neutral deficit with
$\lambda_*h(e_1-e_0)$. Neither formula is a theorem of V45. Nor is an independent
physical time step supplied. In particular the scale $h$ in
\eqref{f16v45:scales} is not defined to be physical Euclidean time. The supplied
memoranda's separation between residual power and locality cost, and their
requirement to keep the full Feshbach return, remain operative. The checkpoint
file records the current branch without recertifying the accumulated archive.

\subsection{Verification and preservation}

The new exact diagnostic is
\begin{center}\small\path{verify_ym_f16_v45_critical_recovery.py}.\end{center}
It checks the balanced powers and normalization, finite translation
stationarity, normalized Gaussian covariance derivatives, color-singlet
cancellation with nonzero hard leakage, the reduced-inverse sign, corrected
Grams, spectral projection counting, and the non-exhaustion example. Its
finite reference matrices are not discretizations or enclosures of the native
Wilson spectrum. Actual executions, source hashes and review scope are in the
audit and checkpoint; rendering checks concern presentation, not proof validity.

The supplied V44 source is unchanged except for its literal title version and
this terminal insertion. Removing those edits recovers it byte for byte. The
hard-band construction uses the exact V44 tube and the original one-step
Wilson kernel; it does not use the endpoint pressure branch. V43 enters only
the outer-return interpretation. The complete earlier analytic lineage has not
been independently recertified. The next scheduled checkpoint is V50.

\begin{firewall}
V45 proves critical-scale physical recovery, an explicit hard-band correction,
and native spectral inclusion at the quartic scale with multiplicities. It
also proves one-sided ordered-level bounds. It does not prove complete
rescaled spectral convergence, identify the native Perron correction or first
neutral coefficient, exclude additional soft states, calibrate physical time,
remove the fixed spatial regulator, or establish an interacting continuum
mass gap. All original probabilities, transfer normalizations, gauge and center
constraints, and the exact returning-state metric are retained.
\end{firewall}
% END F16 V45 APPEND

% BEGIN F16 V46 APPEND -- 2026-09-18
% Insert before the final document-ending command of the supplied V45.
% No predecessor proof, normalization, probability, or clock is changed.

\clearpage
\section{V46: discrete quartic confinement and a native profile alternative}
\label{f16v46:context}

V45 constructs native recovery vectors, but its error estimate is on a test
core. It does not control arbitrary states near the spectral top. We go
upstream to the form that a lower comparison must control. A finite transfer
step has a bounded, frequency-saturating deficit; it cannot simply be replaced
by an unbounded differential energy on every test. We prove an all-vector
confinement estimate for the exact Gaussian-step quartic sandwich, including
its commuting valleys. It gives asymptotic compactness, a lower form limit,
and complete spectral convergence for that explicitly declared comparison.

A separate native deduction uses V45's strong recovery and V42's outer
localization: every nonzero weak critical profile of a native eigenvector is
an eigenfunction of the same quartic model. Missing native states must
therefore lose compactness, or have deficits unbounded below on this scale.
We finish with one precise nonlocal domination inequality sufficient to
exclude both possibilities. That inequality is a remaining native target,
not an estimate proved by identifying the comparison with Wilson transfer.

\subsection{Sources, nonduplication, and the two different operators}

The supplied V45 appendix, analytic audit, and checkpoint were read completely.
The current source and supplied f14/f15/Grand Tensor archives were searched.
The Grand Tensor already states abstract Mosco and compact-screen principles;
its identified toron passage concerns an $\SU(3)$ family with a separate
localization premise. Neither is used as a missing $\SU(2)$ Wilson bound.
V44 already proves differential quartic confinement. The new estimate below
is for the \emph{finite nonlocal step}, before its differential limit, and
pays both position and frequency tails. General compactness and min--max
principles are credited rather than claimed new. The next checkpoint remains V50.

For primary context, Ichinose--Tamura,
\href{https://arxiv.org/abs/0905.3458}{\emph{Results on Convergence in Norm of
Exponential Product Formulas}}, discusses norm convergence of symmetric
products. Its introductory setting and Theorems 2.1--2.2 were inspected in
parsed primary text. We do not import a rate from that theorem; the lower
and compactness argument required here is proved explicitly. The normalized
oscillator calculation can also be checked against
\href{https://dlmf.nist.gov/18.18.E28}{DLMF 18.18.28}. Simon's
\href{https://doi.org/10.1016/0003-4916(83)90057-X}{discrete-spectrum work}
is the classical precedent already cited in V44. No new Navier--Stokes
statement is a premise. The supplied memoranda's separate compactness and
residual-scale requirements remain in force.

On $\mathcal K=L^2(\mathbb R^9,dq)$ write
\[
 \mathcal Q(q)=\sum_{i<j}|q_i\times q_j|^2,\qquad
 H_{\rm mat}=-\tfrac12\Delta+2\mathcal Q,\qquad q_i\in\mathbb R^3.
\]
For an independent numerical step $h>0$ define
\begin{equation}
 m_h=e^{-h\mathcal Q},\qquad G_h=e^{h\Delta/2},\qquad
 S_h=m_hG_hm_h,\qquad D_h=h^{-1}(I-S_h),\qquad
 d_h[f]=\langle f,D_hf\rangle.
 \label{f16v46:comparison}
\end{equation}
Multiplication is meant in $m_h$. The heat kernel is exactly
$(2\pi h)^{-9/2}e^{-|q-q'|^2/(2h)}$. The sandwich is positive self-adjoint,
$0\preceq S_h\preceq I$. It is not $e^{-hH_{\rm mat}}$, is not the original
Wilson transfer, and is not a conditional law on constant fine links. The
Gaussian step is explicit here; it is never substituted for a native Wilson
step without an additional comparison. Every assertion also restricts to the
closed simultaneous-rotation-invariant subspace $\mathcal K_{\rm inv}$.

\subsection{The exact deficit retains its saturated kinetic part}

Use the unitary Euclidean Fourier transform, with variable $p$.
\begin{proposition}[A positive two-term identity on all of $L^2$]
\label{f16v46:identity}
For every $f\in L^2$ and every $h>0$,
\begin{equation}
 \boxed{\quad
 d_h[f]=\int\frac{1-e^{-2h\mathcal Q(q)}}h|f(q)|^2dq
       +\int\frac{1-e^{-h|p|^2/2}}h
                         |\widehat{m_hf}(p)|^2dp.
 \quad}
 \label{f16v46:exact-energy}
\end{equation}
In particular $\|f-m_hf\|_2^2\le h d_h[f]$. If $hR^2\le1$, then
\begin{equation}
 \int_{|p|>R}|\widehat f(p)|^2dp
             \le 2h d_h[f]+8R^{-2}d_h[f].
 \label{f16v46:Fourier-tail}
\end{equation}
Neither estimate requires an $H^1$ test or a bounded quartic potential.
\end{proposition}
\begin{proof}
Add and subtract $\|m_hf\|^2$ in
$\|f\|^2-\langle m_hf,G_hm_hf\rangle$ and apply Plancherel. For
$0\le m\le1$, $(1-m)^2\le1-m^2$, proving the norm estimate. On $|p|>R$,
the second integrand multiplier is at least
$(1-e^{-hR^2/2})/h$. For $hR^2\le1$ this is at least $R^2/4$.
Apply the squared triangle inequality to
$\widehat f=\widehat{m_hf}+\widehat{f-m_hf}$. This gives
\eqref{f16v46:Fourier-tail}. All terms are finite since the multipliers are
bounded at fixed $h$.
\end{proof}

The kinetic multiplier tends to $|p|^2/2$ only at fixed frequency. It is at
most $1/h$. Thus a bound $d_h[f]\ge c\|\nabla f\|^2$ on every test, with
$c>0$, would be false even at one fixed $h$. The following confinement
estimate has the corresponding necessary saturation in position.

\subsection{A transverse-oscillator estimate before the differential limit}

Put
\[
 b(t)=\exp[-4\operatorname{arsinh}(t)]
       =\bigl(\sqrt{1+t^2}-t\bigr)^4,\qquad t\ge0.
\]
\begin{theorem}[Nonlocal confinement through all commuting valleys]
\label{f16v46:confinement}
For each $i=1,2,3$, as bounded quadratic forms,
\begin{equation}
 \boxed{\qquad S_h\preceq M_{b(h|q_i|)}.\qquad}
 \label{f16v46:fibre-bound}
\end{equation}
Consequently, for every $f\in L^2$,
\begin{align}
 d_h[f]&\ge\frac1{3h}\sum_{i=1}^3
       \int[1-b(h|q_i|)]|f(q)|^2dq,\label{f16v46:radial-exact}\\
 d_h[f]&\ge\frac16\int\sum_{i=1}^3
                  \min\{|q_i|,h^{-1}\}|f(q)|^2dq.\label{f16v46:radial-floor}
\end{align}
In particular, if $hR\le1$,
\begin{equation}
 \int_{|q|>R}|f(q)|^2dq\le\frac{6\sqrt3}{R}\,d_h[f].
 \label{f16v46:position-tail}
\end{equation}
No small-field window or restriction away from rank-one commuting triples is
used in these bounds.
\end{theorem}
\begin{proof}
Fix $i$. The heat factors in different three-vector variables commute and
are positive contractions. Hence
\[
 G_h\preceq I_i\otimes\exp\bigl[(h/2)\sum_{j\ne i}\Delta_{q_j}\bigr].
\]
Congruence by $m_h$ preserves this operator inequality. The right side now
disintegrates over $q_i=x$. At fixed $x$, put $r=|x|$ and remove only the
remaining pair potential $|q_j\times q_k|^2$, where $j,k\ne i$. The exact
fibre operator has the form $M_a B_x M_a$, with
\[
 0<a=e^{-h|q_j\times q_k|^2}\le1,\qquad
 B_x=e^{-h\sum_{j\ne i}|x\times q_j|^2}
       e^{(h/2)\sum_{j\ne i}\Delta_{q_j}}
       e^{-h\sum_{j\ne i}|x\times q_j|^2}.
\]
Thus its norm is at most $\|B_x\|$. This step is a contraction in operator
norm; it does not claim that increasing a multiplication factor orders two
sandwiches in Loewner order.

For $r>0$, rotate the two other vectors so that $x$ is axial. Each contributes
a free longitudinal heat factor of norm one and two transverse factors
\[
 e^{-hr^2y^2}e^{h\partial_y^2/2}e^{-hr^2y^2}.
\]
After $y=\sqrt h\,u$, one such factor has normalized kernel
\[
 (2\pi)^{-1/2}
 e^{-(u-v)^2/2-h^2r^2(u^2+v^2)}.
\]
Its ground eigenvalue is $e^{-\operatorname{arsinh}(hr)}$ by the Mehler
formula: its parameter is $\nu=4h^2r^2$ and
$\operatorname{arcosh}(1+\nu/2)=2\operatorname{arsinh}(hr)$.
This normalization can be checked by inserting the positive Gaussian of
frequency $\sqrt{\nu+\nu^2/4}$ and integrating. Its higher eigenvalues have
strictly smaller modulus. The four transverse factors therefore give
$\|B_x\|=b(hr)$. At $r=0$ all factors are free and the norm is one, the
same formula by continuity. The rotation is used only inside this fixed
fibre; no derivative in $q_i$ is taken or lost.

For any complex fibre test its Rayleigh quotient is bounded by that norm
times its squared norm. Integrating over $x$ proves
\eqref{f16v46:fibre-bound}. Average its three consequences for $I-S_h$.
For $0\le t\le1$, $\operatorname{arsinh}t\ge t/2$ and
$1-e^{-2t}\ge(1-e^{-2})t\ge t/2$; for $t\ge1$ monotonicity gives the same
lower bound $1/2$. Thus $1-b(t)\ge\frac12\min(t,1)$, proving
\eqref{f16v46:radial-floor}. If $|q|>R$, at least one $|q_i|>R/\sqrt3$.
When $hR\le1$, the sum of the saturated weights is at least $R/\sqrt3$ on
that event. This proves the last estimate.
\end{proof}

\begin{proposition}[Compact transfer, but not a compact resolvent at fixed step]
\label{f16v46:compact-step}
For every fixed $h>0$, $S_h$ is compact, injective, and positivity improving,
with $0<\|S_h\|<1$. The bounded operator $D_h$ does not have compact
resolvent on the infinite-dimensional space: its essential spectrum is
$\{h^{-1}\}$.
\end{proposition}
\begin{proof}
Partition $\{|q|>R\}$ into three disjoint measurable sets $E_i$ on which
$|q_i|>R/\sqrt3$. For $f=\sum_i1_{E_i}f$, positivity and the Hilbert-space
Cauchy--Schwarz inequality give
\[
 \|S_h^{1/2}f\|^2\le3\sum_i\|S_h^{1/2}1_{E_i}f\|^2
               \le3b(hR/\sqrt3)\|f\|^2.
\]
This tends to zero at fixed $h$ as $R\to\infty$. For $P_R=1_{\{|q|\le R\}}$,
it follows that $\|S_h-P_RS_hP_R\|\le2\sqrt{3b(hR/\sqrt3)}\to0$.
The bounded continuous kernel restricted to two bounded balls is
Hilbert--Schmidt. Hence $S_h$ is compact. Strict positivity of its kernel
gives positivity improvement. Its factors are injective, and
$\langle f,S_hf\rangle=\|G_h^{1/2}m_hf\|^2$ gives injectivity directly.
If its attained norm were one, \eqref{f16v46:exact-energy} would force
$\widehat{m_hf}$ to be supported at $p=0$, impossible for a nonzero $L^2$
function. Finally $D_h=h^{-1}I-h^{-1}S_h$ is a compact perturbation of
$h^{-1}I$. Its resolvent contains a nonzero scalar multiple of identity at
fixed $h$. The same conclusions hold in the invariant subspace.
\end{proof}

\subsection{All-vector compactness and the lower form limit}

\begin{theorem}[Asymptotic compactness, including the nonlocal ultraviolet tail]
\label{f16v46:asymptotic-compactness}
If $h_k\downarrow0$ and
\[
 \sup_k\bigl(\|f_k\|_2^2+d_{h_k}[f_k]\bigr)<\infty,
\]
then a subsequence converges strongly in $L^2(\mathbb R^9)$. If
$f_k\rightharpoonup f$, then
\begin{equation}
 \boxed{\quad
 \langle f,H_{\rm mat}f\rangle_{\rm form}
                 \le\liminf_k d_{h_k}[f_k].
 \quad}
 \label{f16v46:liminf}
\end{equation}
There are strongly convergent recovery sequences for every vector in the
model form domain. These are statements on all bounded-deficit sequences,
not only Gaussian or compactly supported trial functions.
\end{theorem}
\begin{proof}
For fixed large $R$, eventually $h_kR\le1$ and $h_kR^2\le1$.
Equations \eqref{f16v46:position-tail} and \eqref{f16v46:Fourier-tail} give
uniformly small position and Fourier tails in the successive limits
$k\to\infty$, then $R\to\infty$. For a direct compactness argument, the
operator
\[
 K_R=1_{\{|q|\le R\}}\mathcal F^{-1}1_{\{|p|\le R\}}\mathcal F
\]
is Hilbert--Schmidt. The norm of $f_k-K_Rf_k$ is bounded by the sum of the
two tail norms. At a fixed $R$ the sequence $K_Rf_k$ is relatively compact;
a diagonal extraction proves strong compactness of $f_k$. This argument
requires $h_k\to0$. Bounded form balls at one fixed $h$ are not compact.

For the lower bound select a subsequence attaining the finite liminf. It has
a further strongly convergent subsequence, whose limit must be the stated
weak limit $f$. Put $g_k=m_{h_k}f_k$. The first proposition gives $g_k\to f$
strongly. On a bounded position ball the potential multiplier in
\eqref{f16v46:exact-energy} converges uniformly to $2\mathcal Q$; on a bounded
Fourier ball its kinetic multiplier converges uniformly to $|p|^2/2$.
Retain both nonnegative truncated integrals, pass to the strong limit, and
then increase the two radii. This proves \eqref{f16v46:liminf}, including
membership in the model form domain.

If $f\in C_c^\infty$, $m_hf\to f$ in $H^1$, and
$(1-e^{-h|p|^2/2})/h\le|p|^2/2$. The exact energy identity therefore gives
$d_h[f]\to\frac12\|\nabla f\|^2+2\int\mathcal Q|f|^2$.
Compact smooth functions are a form core: spatial cutoff followed by local
mollification suffices for this nonnegative continuous polynomial potential.
A slow diagonal approximation supplies recovery for every form-domain
vector. Rotation averaging preserves that construction in the invariant
subspace. No derivative of an arbitrary $f_k$ was assumed.
\end{proof}

\begin{theorem}[Complete spectral convergence for the declared discrete comparison]
\label{f16v46:norm-resolvent}
On the full space, and separately on $\mathcal K_{\rm inv}$,
\begin{equation}
 \boxed{\quad
 \|(D_h+1)^{-1}-(H_{\rm mat}+1)^{-1}\|\longrightarrow0.
 \quad}
 \label{f16v46:resolvent}
\end{equation}
Let $\mu_0(h)\ge\mu_1(h)\ge\cdots>0$ be the invariant eigenvalues of $S_h$.
Then for each fixed $j$,
\begin{equation}
 \frac{1-\mu_j(h)}h\longrightarrow e_j,
 \qquad -h^{-1}\log\mu_j(h)\longrightarrow e_j.
 \label{f16v46:complete-eigenvalues}
\end{equation}
Isolated finite spectral clusters have convergent projections in norm and
exactly the limiting multiplicity eventually. Also, for $0<t_0<T<\infty$,
\begin{equation}
 \sup_{t\in[t_0,T]}
 \|S_h^{\lfloor t/h\rfloor}-e^{-tH_{\rm mat}}\|\longrightarrow0.
 \label{f16v46:step-semigroup}
\end{equation}
No rate is asserted, and these are not Wilson eigenvalue or physical-time
statements.
\end{theorem}
\begin{proof}
The lower bound and recovery just proved give convergence of minimizers of
$d_h[u]+\|u\|^2-2\Re\langle g,u\rangle$. This is the usual form-to-resolvent
argument already recorded in the Grand Tensor, here with its hypotheses
verified for the present operators. To prove the stronger norm statement,
let $g_h$ be any unit-bounded sequence. The resolvent vectors
$u_h=(D_h+1)^{-1}g_h$ have uniformly bounded norm and energy. Select
$g_h\rightharpoonup g$ and $u_h\to u$ strongly by the preceding theorem.
The minimizer argument works with these varying right sides because every
recovery test converges strongly; it identifies
$u=(H_{\rm mat}+1)^{-1}g$. V44's compact-resolvent theorem for $H_{\rm mat}$
also gives $(H_{\rm mat}+1)^{-1}g_h\to u$ strongly. Failure of
\eqref{f16v46:resolvent} would contradict this conclusion on a sequence of
almost norm-maximizing right sides. This avoids asserting that each
$(D_h+1)^{-1}$ is compact; its essential scalar part is $h/(1+h)$ and vanishes.

Apply norm-continuous spectral projection calculus to the bounded positive
resolvents, at intervals separated from the limiting spectrum. Min--max for
their decreasing eigenvalues gives
$[1+(1-\mu_j(h))/h]^{-1}\to(1+e_j)^{-1}$. This proves the first limit in
\eqref{f16v46:complete-eigenvalues}; the scalar logarithm proves the second.
It also proves the counting and projection assertions.

Norm resolvent convergence of these nonnegative operators gives convergence
of $e^{-tD_h}$ for $t>0$, uniformly on compact positive time intervals. This
can be seen by polynomial approximation of $e^{-t(x^{-1}-1)}$, extended by
zero at $x=0$, on the bounded resolvent spectral interval. The approximation
is uniform for $t\in[t_0,T]$. Finally, uniformly for $0\le x\le h^{-1}$,
\[
 |(1-hx)^{\lfloor t/h\rfloor}-e^{-tx}|\le C_{t_0,T}h.
\]
For example split $hx\le1/2$ from its complement and use
$0\le-\log(1-y)-y\le y^2$ on the first part; the other part is exponentially
small in $\lfloor t/h\rfloor$. The rounding error in $t$ is bounded using
$\sup_{x\ge0}xe^{-t_0x/2}<\infty$. Functional calculus proves
\eqref{f16v46:step-semigroup}.
\end{proof}

\subsection{A reference that also detects hard-band loss}

Retain V45's exact hard Gaussian operator and let
\[
 A_+=\mathcal T_+/\lambda_*,\qquad
 P_+=|\varphi\rangle\langle\varphi|,\qquad
 r_+=\|A_+|_{\varphi^\perp}\|<1.
\]
On $\mathcal X=L^2(\mathcal Y\times\mathbb R^9)$, or its joint invariant
subspace, define the bounded nonlocal comparison form
\begin{equation}
 \mathbb D_h=h^{-1}(I-A_+\otimes S_h),\qquad
 \mathcal I f=\varphi\otimes f.
 \label{f16v46:hard-reference}
\end{equation}
\begin{proposition}[Exact hard separation for the comparison]
\label{f16v46:hard-separation}
Writing $F=\mathcal I f+F_\perp$ gives
\begin{equation}
 \boxed{\quad
 \langle F,\mathbb D_hF\rangle
 \ge d_h[f]+\frac{1-r_+}{h}\|F_\perp\|^2.
 \quad}
 \label{f16v46:hard-floor}
\end{equation}
Every sequence with bounded norm and bounded $\mathbb D_h$ energy as
$h\downarrow0$ has a strongly convergent subsequence of the form
$F_h\to\varphi\otimes f$. Moreover
\begin{equation}
 \left\|(\mathbb D_h+1)^{-1}
       -\mathcal I(H_{\rm mat}+1)^{-1}\mathcal I^*\right\|
 \le\|(D_h+1)^{-1}-(H_{\rm mat}+1)^{-1}\|
       +\frac{h}{1-r_++h}\longrightarrow0.
 \label{f16v46:hard-resolvent}
\end{equation}
\end{proposition}
\begin{proof}
The ground projection reduces the exact tensor comparison. On its range the
form is $D_h$, and on the complement $A_+\otimes S_h\preceq r_+I$.
This proves \eqref{f16v46:hard-floor} and the resolvent bound. The hard
complement norm is $O(\sqrt h)$ for a bounded-energy sequence, and its ground
profile has bounded $d_h$ energy. Apply Theorem~\ref{f16v46:asymptotic-compactness}.
The joint symmetry commutes with both projections; its ground component is
exactly the invariant soft subspace. No native hard-band floor is asserted.
\end{proof}

\subsection{A native profile theorem without an assumed lower bound}

We now return to the original $T_\gamma$ and fix $L$ throughout. Use
$h=(\gamma V)^{-1/3}$, $\varepsilon=\gamma^{-1/2}$ and
$\tau=\sqrt Vh$, exactly as in V45. A new symbol $\mathbb D_h$ above denotes a
comparison only. It is not identified with
\[
 \mathfrak A_\gamma=h^{-1}(I-T_\gamma/\lambda_*).
\]
The latter need not be nonnegative. Neither V41 nor V45 proves a uniform
lower bound for it.

Choose fixed nested, center-translated invariant cutoffs in V44's tubes,
$\chi_\sigma$, equal to one near each central orbit and with disjoint supports.
Put $\chi=\sum_\sigma\chi_\sigma$. Use the same factor $8^{-1/2}$ as in V43
to unfold a vector of character $a$ to one regional vector. Remove the exact
tube half-density and dilate by $s=\varepsilon u+h\mathsf Kq$. More explicitly,
if $\widehat\psi=W_a^*(\chi\psi)$, define
\begin{equation}
 (\mathfrak E_{\gamma,a}\psi)(u,q)
 =\varepsilon^{r/2}\tau^{9/2}J_L(s)^{1/2}
                 \widehat\psi(\Phi(1,s)),
 \qquad s=\varepsilon u+h\mathsf Kq,
 \label{f16v46:extraction}
\end{equation}
with zero extension outside the dilated tube. The reference is
$du\,dq$ in orthonormal hard and soft coordinates; all based orbit variables
have already been integrated against their probability Haar measure. This
is an exact contraction satisfying
$\|\mathfrak E_{\gamma,a}\psi\|=\|\chi\psi\|$. Its outputs are invariant under
joint color rotations.

V42 can be applied to fixed still smaller neighborhoods inside the plateaux.
It supplies a fixed cutoff $c<\lambda_*$ such that on the native band
$E_{\gamma,a}=1_{[c,1]}(T_\gamma|_{\mathcal H_a})$,
\begin{equation}
 \|(1-\chi)E_{\gamma,a}\|\le C_Le^{-c_L\sqrt\gamma}.
 \label{f16v46:extraction-mass}
\end{equation}
Thus extraction is asymptotically isometric on this entire band, with no
rank factor. The neighborhoods and cutoff are fixed before coupling; no
shrinking-neighborhood bound is imported from V42.

\begin{theorem}[Every nonzero native critical profile is a model eigenfunction]
\label{f16v46:native-profile}
Let $\psi_\gamma\in\mathcal H_a$ be normalized native eigenvectors with
\[
 T_\gamma\psi_\gamma=\lambda_*(1-hE_\gamma)\psi_\gamma,
 \qquad E_\gamma\to E\in\mathbb R.
\]
Along every subsequence on which
$\mathfrak E_{\gamma,a}\psi_\gamma\rightharpoonup F$ in $\mathcal X$,
there is $f\in\mathcal K_{\rm inv}$, $\|f\|\le1$, such that
\begin{equation}
 \boxed{\qquad F=\varphi\otimes f,\qquad H_{\rm mat}f=Ef.\qquad}
 \label{f16v46:profile-equation}
\end{equation}
The value $f=0$ is allowed. If $f\ne0$, then $E$ is an invariant model
eigenvalue and $E\ge e_0$. In particular, any native branch with bounded
scaled deficit converging to $E<e_0$ has only zero weak critical profiles.
This theorem does not assert that any such extra native branch exists.
\end{theorem}
\begin{proof}
The eigenvalues tend to $\lambda_*$, so \eqref{f16v46:extraction-mass} applies
and the extracted norms tend to one. Weak limits nevertheless can lose norm.
For every invariant test consisting of a hard Hermite polynomial times
$\varphi$ and a compact smooth soft test, the zeroth-order part of V45's
strong expansion gives
\[
 (T_\gamma-\lambda_*)\mathcal U_{\gamma,a}G
   =\mathcal U_{\gamma,a}[(\mathcal T_+-\lambda_*)G]+O_G(h).
\]
The tests' non-plateau hard tails are exponentially small, and the soft
support lies in the plateau for small $h$. Thus pairing with $\psi_\gamma$
and passing to its weak extracted limit gives
$\langle F,(\mathcal T_+-\lambda_*)G\rangle=0$.
These tests are dense in the jointly invariant space. The hard inverse on
$\varphi^\perp$ is bounded, so $F=\varphi\otimes f$; joint invariance forces
$f$ to be invariant. This uses no estimate on the hard mass of
$\psi_\gamma$ before passage to a weak limit.

For invariant $g\in C_c^\infty(\mathbb R^9)$ use V45's \emph{corrected}
recovery vector. Its strong residual gives exactly
\[
 \langle\psi_\gamma,\Psi_{\gamma,a}H_{\rm mat}g\rangle
 -E_\gamma\langle\psi_\gamma,\Psi_{\gamma,a}g\rangle=o(1).
\]
The $hk(q)$ correction disappears in each fixed pairing; the exact extracted
half-density normalization has already been retained. Passing to the weak
limit yields $\langle f,H_{\rm mat}g\rangle=E\langle f,g\rangle$.
Averaging any compact smooth test over rotations extends this identity to
all such tests, since $f$ is invariant.

For completeness this distributional equation is an operator eigenvalue
equation, not just a formal identity. Local elliptic regularity gives
$f\in H^2_{\rm loc}$. Testing with $\zeta_R^2f$, for standard cutoffs, yields
\[
 \tfrac12\|\nabla(\zeta_Rf)\|^2
 +2\int\mathcal Q\,\zeta_R^2|f|^2
 =E\|\zeta_Rf\|^2+\tfrac12\int|\nabla\zeta_R|^2|f|^2.
\]
The right side is bounded above uniformly in $R$. Letting $R$ increase proves
membership in the Friedrichs form domain. The distributional identity
extends by form density, giving $f\in D(H_{\rm mat})$ and
$H_{\rm mat}f=Ef$. V44's compact resolvent identifies $E$ as a model eigenvalue
when $f\ne0$. Neither strong convergence of the native profiles nor preservation
of their full norm has been used or established.
\end{proof}

\subsection{A specific sufficient native comparison, still to be proved}

The all-vector nonlocal reference is now available to formulate one missing
inequality without an impossible unsaturated gradient bound. The following is
an implication theorem with an additional hypothesis, not a native estimate.
\begin{proposition}[A coarse lower comparison would complete the fixed-volume scaling]
\label{f16v46:conditional-completeness}
Suppose there are fixed constants $C>0$, $b\ge0$, independent of $\gamma$,
such that on the neutral band $\operatorname{Ran}E_{\gamma,0}$,
\begin{equation}
 \boxed{\quad
 \langle\mathfrak E_{\gamma,0}\psi,
            \mathbb D_h\mathfrak E_{\gamma,0}\psi\rangle
 \le C\left[\langle\psi,\mathfrak A_\gamma\psi\rangle
                                    +b\|\psi\|^2\right]
 \quad\hbox{for every }\psi.
 }
 \label{f16v46:missing-domination}
\end{equation}
Then every fixed ordered native level in every center sector satisfies
\begin{equation}
 \frac{1-\lambda_{a,j}(\gamma)/\lambda_*}{h}\longrightarrow e_{j-1}.
 \label{f16v46:conditional-levels}
\end{equation}
In particular this additional hypothesis would imply
\begin{equation}
 \begin{split}
 \Lambda_\gamma&=\lambda_*[1-he_0+o(h)],\\
 -h^{-1}\log(\lambda_{0,2}/\Lambda_\gamma)&\longrightarrow e_1-e_0,\\
 \frac{g_\gamma}{\lambda_*h}&\longrightarrow e_1-e_0,
 \end{split}
 \label{f16v46:conditional-gap}
\end{equation}
with V43's actual returning-state metric retained in $g_\gamma$.
\end{proposition}
\begin{proof}
The left side of \eqref{f16v46:missing-domination} is nonnegative, so it
first gives $\mathfrak A_\gamma\succeq-bI$ on the band. This excludes the
unbounded negative-deficit possibility left by V45. For any fixed number $j$
of ordered neutral eigenvectors, V45 supplies an upper bound on their scaled
deficits, and the new hypothesis supplies the lower bound and uniformly
bounded $\mathbb D_h$ energies of their extracted profiles. Proposition
\ref{f16v46:hard-separation} gives strong subsequential compactness. The exact
band localization \eqref{f16v46:extraction-mass} preserves their entire Gram,
so their limiting profiles are $j$ orthonormal invariant functions.
Theorem~\ref{f16v46:native-profile} makes these model eigenfunctions with
the limiting ordered deficits. Model min--max gives the missing liminf
$e_{j-1}$, while V45 supplies the matching limsup. This proves the neutral
part of \eqref{f16v46:conditional-levels} on every subsequence, hence in full.
V43's sector ladders differ by $O(e^{-c_L\sqrt\gamma})=o(h)$ for every fixed
level, so the same conclusion holds in all sectors. Finally subtract the
first two neutral expansions, use the scalar logarithm, and apply V43--V44's
proved $\Delta_\gamma^{\rm neu}/g_\gamma\to1$.
\end{proof}

The constants in \eqref{f16v46:missing-domination} need not approach one.
The sharp coefficient would come from V45's local equation; the coarse bound
would only prevent loss of norm and energies escaping below the model.
Its left side has the correct bounded-transfer saturation and detects both
hard excited mass and soft escape. This is weaker in purpose than a global
$o(h)$ operator approximation, but it is not a consequence of the recovery
map. The factors arising from the actual quotient kinetic covariance,
nonlinear magnetic action, half-density, and within-region return remain to
be controlled when proving it. The Gaussian reference has not been installed
as the native conditional transition.

\subsection{Current scope and verification}

The new unconditional model result supplies complete nonlocal compactness and
spectral convergence without adding a boundary or deleting commuting valleys.
The new native profile result classifies every nonzero weak critical profile,
but permits zero profiles. Neither statement proves
\eqref{f16v46:missing-domination}. Thus the native ordered asymptotics and
neutral coefficient in \eqref{f16v46:conditional-gap} remain conditional,
not accomplishments of this iteration. The next calculation should establish
that inequality, or its separate position, frequency, and hard-band bounds,
on arbitrary native near-top states rather than on another selected packet.

The appended diagnostic checks the exact deficit algebra, normalized Mehler
factor, finite-step saturation, hard-sector decomposition, and the ordering
logic on finite fixtures. The audit distinguishes these tests from the
analytic compactness proofs. Source preservation and PDF rendering are
reported separately. Only this terminal insertion and the literal title
version change the supplied V45; no historical proof has been rewritten.
The current native profile deduction explicitly inherits the V45 strong core
expansion and V42 band localization as stated. The complete lineage has not
been independently recertified, and the next checkpoint remains V50.

\begin{firewall}
V46 proves all-vector confinement and spectral completeness for the declared
finite-step quartic comparison, and a nonzero-profile rigidity statement for
actual native critical eigenvectors. It does not prove native low-energy
tightness or the stated sufficient domination inequality, identify a native
Perron correction or neutral scaling coefficient unconditionally, remove the
fixed spatial regulator, or calibrate physical time. The comparison step $h$
is not a replacement Wilson kernel or a supplied physical clock.
\end{firewall}
% END F16 V46 APPEND

% BEGIN F16 V47 APPEND -- 2026-09-18
% Insert before the final document-ending command of the supplied cumulative V46.
% The predecessor, original Wilson operator, probabilities and clocks are preserved.

\clearpage
\section{V47: native hard-band capture and local soft compactness}
\label{f16v47:context}

V46 leaves possible norm loss in native critical profiles. Here a local
operator-norm comparison captures the hard ground band for every native
near-top vector. On bounded critical soft regions its error is $O_{L,R}(h)$;
the exact nonlocal localization identity then gives bounded comparison energy
for arbitrary localized native eigenvectors. Frequency escape at bounded soft
position is excluded. Only escape to large critical soft positions remains.
The coarser $O_L(r+\varepsilon)$ error on a physical angular strip cannot be
divided by $h$ globally. No radial tightness or native scaling coefficient is
inferred without that remaining estimate.

\subsection{Inputs and the unchanged physical extraction}

V46's appendix and audit were read in full. The native inputs are V44's
measured tube, V45's exact stationarity, and V42's band localization. V46
supplies the nonlocal comparison estimates and the finite-profile equation.
Targeted archive searches, ownership and review depth are recorded in the
audit; they do not independently certify the lineage. The next checkpoint
remains V50.

For primary context, Panati--Spohn--Teufel,
\href{https://arxiv.org/html/math-ph/0201055}{\emph{Space-Adiabatic Perturbation
Theory}}, Section 2, distinguishes an isolated fast band from the dynamics
remaining within it. Its global symbol and gap assumptions are not applied to
the Wilson kernel below. Hanche-Olsen--Holden,
\href{https://arxiv.org/html/0906.4883}{\emph{The Kolmogorov--Riesz compactness
theorem}}, describes the classical position/translation compactness criterion.
Here its Fourier-screen version is already implemented in V46; the new work
is to verify local bounds for native eigenvectors before using that criterion.
No external Navier--Stokes theorem is a premise.

Fix $L=2m$, $m\ge4$, $V=L^3$, and keep exactly
\[
 T_\gamma=M_{b_\gamma}C_\gamma^{\otimes3V}M_{b_\gamma},\qquad
 b_\gamma=e^{-S/(2\varepsilon^2)},\qquad
 \varepsilon=\gamma^{-1/2},\quad h=(\gamma V)^{-1/3},\quad
 \tau=\sqrt Vh.
\]
Thus $\varepsilon^2=Vh^3$. All constants below may depend on the fixed $L$.
Let $s=y+\mathsf K a$ be V44's tube coordinate, with
$y\in\mathcal Y$, $a\in\mathbb R^9$ and $\mathsf K^*\mathsf K=VI$.
The critical dilation is $y=\varepsilon u$, $a=hq$.
The tube measure is the \emph{same} exact measure
\[
 dH(\Phi(g,s))=dg\,J_L(s)\,ds.
\]
Only the based gauge variables are removed. Joint simultaneous color invariance
of functions of $(u,q)$ imposes the residual Gauss condition.

Write $\mathcal X=L^2(\mathcal Y\times\mathbb R^9)$, restricted to that joint
invariant subspace when required. Retain V46's operators
\[
 A_+=\mathcal T_+/\lambda_*,\quad A_+\varphi=\varphi,\quad
 r_+=\|A_+|_{\varphi^\perp}\|<1,\quad
 \mathcal I f=\varphi\otimes f,
\]
\[
 S_h=e^{-h\mathcal Q}e^{h\Delta_q/2}e^{-h\mathcal Q},\qquad
 \mathbb D_h=h^{-1}(I-A_+\otimes S_h),\qquad
 d_h[f]=h^{-1}\langle f,(I-S_h)f\rangle.
\]
Here $\lambda_*$ remains the original Gaussian threshold, not an assigned
native top. No statement that $T_\gamma/\lambda_*\preceq I$ is assumed.

Choose a fixed product tube $|y|<d_0$, $|a|<r_0$ with compact closure inside
V44's tube and the selected central region. Let $\mathcal V_\gamma$ be the
partial isometry from $\mathcal X$ to this one-region physical space:
\[
 (\mathcal V_\gamma F)(\Phi(g,s))
 =J_L(s)^{-1/2}\varepsilon^{-r_Y/2}\tau^{-9/2}F(u,q),
 \qquad r_Y=\dim\mathcal Y,
\]
with zero extension and with the input restricted to that product tube.
Its Gram is exactly the corresponding multiplication projection. Let
$\mathcal V_{\gamma,a}$ be its normalized eight-character folding in sector
$a$. In the rest of this section a sector label on $\mathcal V$ or an
extraction is distinguished from the constant angle $a$ by its subscript.
Fix V46's extraction in these tubes:
\[
 \mathfrak E_{\gamma,a}\psi=\mathcal V_{\gamma,a}^*(\chi\psi),
 \qquad 0\le\chi\le1,
\]
where $\chi$ is fixed, invariant, has disjoint center-translated supports,
and equals one on smaller tubes. Shrinking the auxiliary radii below does not
change this final extraction.

\subsection{Uniform hard tails from the actual magnetic multiplier}

\begin{proposition}[A hard pin in the central tube and a whole-band moment bound]
\label{f16v47:hard-tail}
After decreasing $d_0,r_0$, there is $c_L>0$ such that
\begin{equation}
 S(\operatorname{Exp}(y+\mathsf K a))
       \ge S(\operatorname{Exp}(\mathsf K a))+c_L|y|^2
       \ge c_L|y|^2.
 \label{f16v47:hard-pin}
\end{equation}
For any fixed $c>0$, every spectral band
$\mathsf P_\gamma\le1_{[c,1]}(T_\gamma|_{\mathcal H_a})$ satisfies
\begin{equation}
 \big\|e^{\beta_L|u|^2}\mathfrak E_{\gamma,a}\mathsf P_\gamma\big\|
       \le c^{-1}
 \label{f16v47:hard-moment}
\end{equation}
for some $\beta_L>0$, independent of $\gamma$ and of the rank of the band.
The same bound holds for smaller supported cutoffs. It concerns rescaled hard
coordinates only; no critical soft moment is asserted.
\end{proposition}
\begin{proof}
V45's exact mean-zero stationarity gives $D_yS(\mathsf K a)=0$. At the origin,
$D_y^2S=H_+\succ0$. Smoothness on a fixed small product neighborhood makes this
Hessian uniformly positive along every segment from $0$ to $y$, also for
noncommuting constant triples. Taylor's integral formula proves
\eqref{f16v47:hard-pin}. Center translation preserves it.

On the band, write $\mathsf P_\gamma=T_\gamma\mathsf P_\gamma
(T_\gamma|_{\operatorname{Ran}\mathsf P_\gamma})^{-1}$.
The inverse has norm at most $c^{-1}$. In the physical tube, multiplication by
$\chi e^{\beta_L|y|^2/\varepsilon^2}b_\gamma$ has norm at most one when
$\beta_L\le c_L/2$. The remaining two factors $C_\gamma^{\otimes3V}$ and
$b_\gamma$ are contractions. The exact extraction isometry and center folding
therefore give \eqref{f16v47:hard-moment}. This is an operator estimate on the
whole band, not a sum of individual eigenvector bounds.
\end{proof}

\subsection{The reduced kinetic kernel, with its orbit normalization}

Let $d=\dim\mathcal N=6V+3$ and $D=9V$. On the fixed tube, the flat
half-density kernel of the \emph{raw} kinetic compression is exactly
\begin{equation}
 k_\varepsilon^{\rm red}(s,t)
  =J_L(s)^{1/2}J_L(t)^{1/2}
    \int_{\mathcal G_o}
       c_\gamma(\operatorname{Exp}s,g\cdot\operatorname{Exp}t)\,dg.
 \label{f16v47:quotient-kernel}
\end{equation}
This follows by disintegrating the original input Haar integral; $c_\gamma$
contains all original kinetic normalizers. The based integral is not a new
conditional transition. Formula \eqref{f16v47:quotient-kernel} also acts on
non-invariant quotient tests before restricting to joint color invariants.

\begin{proposition}[A normalized local kernel comparison with an integrable error]
\label{f16v47:kinetic}
In a sufficiently small fixed tube there are constants $C_L,c_L>0$ and an
integer $M_L$ such that, with
$g_\varepsilon(s-t)=(2\pi\varepsilon^2)^{-d/2}
 e^{-|s-t|^2/(2\varepsilon^2)}$,
\begin{equation}
 \begin{split}
 |k_\varepsilon^{\rm red}(s,t)-g_\varepsilon(s-t)|
 &\le C_L(|s|+|t|+\varepsilon)\varepsilon^{-d}
       (1+|s-t|/\varepsilon)^{M_L}e^{-c_L|s-t|^2/\varepsilon^2}\\
 &\quad+C_L\varepsilon^{-D}e^{-c_L/\varepsilon^2}.
 \end{split}
 \label{f16v47:kernel-error}
\end{equation}
The constants are uniform as a smaller soft angular strip shrinks inside this
tube. Both row and column integrals of the Gaussian-polynomial factor are
bounded independently of $\varepsilon$.
\end{proposition}
\begin{proof}
Choose local based-group coordinates $v\in\mathbb R^{D-d}$ whose orbit
differential at the identity is an isometry into $\operatorname{Ran}d_0$.
The remaining differential is the orthogonal inclusion of $\mathcal N$.
Let $c_X$ be the full Haar density in an orthonormal fine-link chart and let
$c_G$ be the based Haar density in these orbit-orthonormal coordinates.
The exact tube Jacobian gives $J_L(0)c_G=c_X$.
The normalized raw kernel has leading constant
$c_X^{-1}(2\pi\varepsilon^2)^{-D/2}$. Therefore integration of its vertical
Gaussian gives exactly
\[
 J_L(0)c_Gc_X^{-1}(2\pi\varepsilon^2)^{-D/2}
       \int e^{-|v|^2/(2\varepsilon^2)}dv
       =(2\pi\varepsilon^2)^{-d/2}.
\]
All orbit metric determinants are included in $c_G$ and $J_L(0)$; none is
replaced by one separately.

Put $t=s+\xi$ and let $A(s,\xi,v)$ be the sum of the actual displacement
cosines. Smoothness, its zero and zero first derivative at $(\xi,v)=0$, and
orthogonality at $s=0$ give
\[
 A(s,\xi,v)=\tfrac12(|\xi|^2+|v|^2)+\mathcal R,
 \quad
 |\mathcal R|\le C_L\{|s|(|\xi|^2+|v|^2)+(|\xi|+|v|)^3\}.
\]
After one fixed shrink, $A\ge c_L(|\xi|^2+|v|^2)$ there. The exact smooth
volume factors differ from their value at zero by
$O_L(|s|+|t|+|v|)$, and the original normalization has relative error
$O_L(\varepsilon^2)$. The mean-value estimate for two exponentials, using the
positive quadratic lower bounds for both of them, bounds their difference by
$|\mathcal R|/\varepsilon^2$ times a Gaussian majorant. Integrate in $v$;
each cubic vertical moment contributes one factor $\varepsilon$ and a fixed
polynomial in $|\xi|/\varepsilon$. This proves the first term in
\eqref{f16v47:kernel-error}.

Outside the selected based-group coordinate neighborhood, freeness of the
based action and compactness give a fixed positive displacement from the
identity orbit representative. This stays positive for all small $s,t$.
The original cosine bound then gives the last term. A fixed finite number of
charts bounds its integral. The row and column assertions follow by the
change $\xi=\varepsilon z$; all integrations are over a fixed bounded tube.
This proves an absolute kernel error and its Schur bounds, not a Loewner
ordering from pointwise kernel comparisons.
\end{proof}

\subsection{An all-vector strip comparison, including shrinking critical strips}

Let $\Pi_{\gamma,r}$ be multiplication in profile space by the strip
$|\varepsilon u|<d_0$, $|hq|<r$, with $0<r\le r_0$. Let
$K_{\gamma,a}^{\rm red}=\mathcal V_{\gamma,a}^*T_\gamma\mathcal V_{\gamma,a}$.

\begin{theorem}[Native strip norm comparison before dividing by the scale]
\label{f16v47:strip}
There are $C_L<\infty$ and a fixed coupling threshold such that, uniformly for
$0<r\le r_0$ and all center sectors,
\begin{equation}
 \boxed{\quad
 \big\|\Pi_{\gamma,r}
       (K_{\gamma,a}^{\rm red}-\mathcal T_+\otimes S_h)
          \Pi_{\gamma,r}\big\|
       \le C_L(r+\varepsilon).
 \quad}
 \label{f16v47:strip-bound}
\end{equation}
The norm is on the complete supported $L^2$ space, not a chosen test core.
In particular, on $|q|\le2R$, with fixed $R$ and $2hR\le r_0$, the error is
at most $C_L(2hR+\varepsilon)$.
\end{theorem}
\begin{proof}
Write $S_c(a)=S(\operatorname{Exp}(\mathsf K a))$. Exact stationarity and
bounded third derivatives give
\[
 S(\operatorname{Exp}(y+\mathsf K a))
 =S_c(a)+\tfrac12 y^{\mathsf T}H_+y+\mathcal R_m,
 \qquad |\mathcal R_m|\le C_L(|a||y|^2+|y|^3).
\]
Both the actual exponent and the split exponent are bounded below by
$S_c(a)+c_L|y|^2$ on the chosen tube. Their half-multiplier difference hence
has supremum at most $C_L(r+\varepsilon)$ on the strip: set $y=\varepsilon u$
and use the boundedness of each polynomial times $e^{-c_L|u|^2}$.
The constant-field identity gives, with $Q(a)=\mathcal Q(a)$,
\[
 (1-C_Lr^2)Q(a)
 \le\sum_{i<j}|v(a_i)\times v(a_j)|^2\le Q(a),
 \qquad v(a)=\frac{\sin|a|}{|a|}a.
\]
Since $\gamma Vh^4=h$, the supremum difference between
$e^{-S_c(hq)/(2\varepsilon^2)}$ and $e^{-h\mathcal Q(q)}$ is at most $C_Lr^2$.
This bound is uniform even when $h\mathcal Q(q)$ is large: use
$xe^{-c x}\le C$ instead of bounding the exponent difference alone.

Now sandwich \eqref{f16v47:kernel-error} by the actual magnetic factors.
On the strip, $|s|+|t|\le |y|+|y'|+C_Lr$. The hard pin gives
$|y|b_\gamma(s)\le C_L\varepsilon$ and the same estimate at $t$.
The row and column Gaussian-polynomial integrals therefore bound the weighted
kinetic error by $C_L(r+\varepsilon)$. The remote based-group term is
exponentially small. Comparing next the two magnetic factors with their split
versions uses only their supremum bounds and contraction of the flat Gaussian
convolution. Thus the complete difference has the claimed norm.

Finally dilate by $y=\varepsilon u$, $a=hq$. The hard convolution and its two
quadratic multipliers give exactly $\mathcal T_+$. The orthonormal soft step
has covariance $\varepsilon^2/\tau^2=h$; its two multipliers give exactly
$S_h$. These are comparisons after assembling the original measure and
kernel, not replacements of the original transfer. Center folding adds only
cross terms between fixed separated tubes, of norm $O_L(e^{-c_L\gamma})$.
They are absorbed in the displayed error. No derivative of an arbitrary
supported test is used.
\end{proof}

The estimate is deliberately coarse. The first hard variation is charged in
operator norm, so it costs $O(r)$, even though its ground expectation vanishes.
On a critical strip $r=2hR$ this is $O(hR)$, which suffices for local energy
bounds. It does not yield V46's global $O(h)$ domination as $R\to\infty$.

\subsection{Capture of the hard ground band for every near-top vector}

Fix $K\ge0$ and put
\[
 \mathsf B_{\gamma,a}(K)
       =1_{[\lambda_*(1-Kh),\,1]}(T_\gamma|_{\mathcal H_a}),\qquad
 F_\gamma(\psi)=\mathfrak E_{\gamma,a}\psi,\quad
 f_\gamma(\psi)=\mathcal I^*F_\gamma(\psi).
\]
A zero band gives zero operator norms. The native top is not assumed to be
below $\lambda_*$; negative scaled deficits are included in this band.

\begin{theorem}[Whole-band hard capture and unscaled soft-step regularity]
\label{f16v47:capture}
At fixed $L,K$, uniformly over unit vectors in $\operatorname{Ran}\mathsf B_{\gamma,a}(K)$,
\begin{equation}
 \boxed{\quad
 \|F_\gamma-\mathcal I f_\gamma\|\longrightarrow0,\qquad
 \|f_\gamma\|\longrightarrow1,\qquad
 h\,d_h[f_\gamma]\longrightarrow0.
 \quad}
 \label{f16v47:capture-result}
\end{equation}
Equivalently the first assertion is convergence to zero of
$\|(I-\mathcal I\mathcal I^*)\mathfrak E_{\gamma,a}\mathsf B_{\gamma,a}(K)\|$.
There is no spectral-rank factor. No $O(h)$ bound for the squared hard loss,
or boundedness of $d_h[f_\gamma]$, is asserted globally.
\end{theorem}
\begin{proof}
Choose, for each fixed small $r>0$, an invariant center-translated cutoff
$\chi_r$ supported in $|y|<d_0$, $|a|<r$, equal to one near the central orbits,
and with support inside the plateau of $\chi$. V42 applied to still smaller
fixed neighborhoods gives, eventually,
\[
 \|(1-\chi_r)\mathsf B_{\gamma,a}(K)\|
       \le C_{L,r}e^{-c_{L,r}\sqrt\gamma}=:t_{\gamma,r}.
\]
This uses $\lambda_*(1-Kh)\to\lambda_*$, so the shrinking spectral band lies
above the fixed neighborhood's cutoff. It does not use a shrinking spatial
neighborhood estimate with uniform constants.
For a unit band vector let $F_r=\mathcal V_{\gamma,a}^*(\chi_r\psi)$.
The extraction norm differs from one by at most $O(t_{\gamma,r})$, and
$\|F_r-F_\gamma\|\le t_{\gamma,r}$. Changing its transfer quotient to the
uncut one costs at most $2t_{\gamma,r}$ because $\|T_\gamma\|\le1$.
By \eqref{f16v47:strip-bound} and the lower spectral endpoint,
\[
 0\le\langle F_r,(I-A_+\otimes S_h)F_r\rangle
 \le Kh+C_L(r+\varepsilon+t_{\gamma,r}).
\]
The bounded reference form changes by $O(t_{\gamma,r})$ when $F_r$ is replaced
by $F_\gamma$. Its exact hard separation is
\[
 \langle F_\gamma,(I-A_+\otimes S_h)F_\gamma\rangle
 \ge(1-r_+)\|F_\gamma-\mathcal I f_\gamma\|^2
       +h d_h[f_\gamma].
\]
First let $\gamma\to\infty$ at each fixed $r$, and then $r\downarrow0$.
This proves both vanishing terms, uniformly on the band. V46's extraction-mass
bound, followed by Pythagoras, proves $\|f_\gamma\|\to1$.
\end{proof}

For normalized band sequences the extracted joint probability consequently
factorizes in the following precise sense:
\[
 \big\||F_\gamma(u,q)|^2-|\varphi(u)|^2|f_\gamma(q)|^2\big\|_{L^1(du\,dq)}
       \longrightarrow0.
\]
This is convergence of a joint density after the specified extraction. It is
not a claim that the compact native measure factors exactly, nor a control of
the remaining soft density. The exponential hard moment bound additionally
prevents hard-position mass from escaping behind this norm limit.

\subsection{A localized native energy bound at order \texorpdfstring{$h$}{h}}

Let $\vartheta\in C_c^\infty(\mathbb R^9)$ be real, radial, equal to one on
the unit ball and zero outside the ball of radius two, with $0\le\vartheta\le1$.
For $R\ge1$ put $\vartheta_R(q)=\vartheta(q/R)$. On each tube define the
physical multiplier
\[
 \chi_{\gamma,R}(U)=\chi(U)\vartheta\bigl(a(U)/(hR)\bigr),
\]
and extend by zero and center translation. It is gauge and center invariant.
Its Lipschitz constant in the original compact metric is at most
$C_L[1+(hR)^{-1}]$. The corresponding extracted vector is exactly
$\vartheta_R F_\gamma$.

\begin{theorem}[Local comparison energies for arbitrary native eigenvectors]
\label{f16v47:local-energy}
Let $\|\psi_\gamma\|=1$ and
\[
 T_\gamma\psi_\gamma=\lambda_*(1-hE_\gamma)\psi_\gamma,
 \qquad E_\gamma\le K.
\]
There is no lower bound assumed for $E_\gamma$. For every fixed $R\ge1$, at
sufficiently large coupling,
\begin{equation}
 \boxed{\quad
 \langle\vartheta_RF_\gamma,\mathbb D_h\vartheta_RF_\gamma\rangle
 \le [E_\gamma+C_L(R+1)]\|\vartheta_RF_\gamma\|^2
                       +C_L(R^{-2}+h^2).
 \quad}
 \label{f16v47:local-energy-sharp}
\end{equation}
In particular the left side is bounded above by a finite $C_{L,K,R}$.
It follows that
\begin{align}
 \|(I-\mathcal I\mathcal I^*)\vartheta_RF_\gamma\|^2
      &\le \frac{h C_{L,K,R}}{1-r_+},\\
 d_h[\vartheta_Rf_\gamma]&\le C_{L,K,R},
 \label{f16v47:local-consequences}\\
 \int_{|p|>P}|\widehat{\vartheta_R f_\gamma}(p)|^2dp
      &\le C_{L,K,R}(2h+8P^{-2}),\qquad hP^2\le1.
 \label{f16v47:local-frequency}
\end{align}
Thus every such sequence has a subsequence with strongly convergent soft
profiles on every bounded $q$-region. High-frequency loss at bounded soft
position is excluded. None of the constants is asserted bounded as $R\to\infty$.
\end{theorem}
\begin{proof}
For a real Lipschitz multiplier $b$, symmetry of the actual kernel gives
\[
 \langle b\psi,(\lambda_*-T_\gamma)b\psi\rangle
 =\lambda_*hE_\gamma\|b\psi\|^2+
 \frac12\iint(b(U)-b(V))^2\overline{\psi(U)}
             t_\gamma(U,V)\psi(V)\,dH(U)dH(V).
\]
The last integral is real but need not be nonnegative for a complex test.
Its absolute value is at most $C_L\varepsilon^2\operatorname{Lip}(b)^2$
by the actual displacement second moment and a row/column Schur estimate.
For $b=\chi_{\gamma,R}$ this is at most
$C_L(h/R^2+h^3)$ because $\varepsilon^2=Vh^3$.
This is a localization estimate for the bounded native kernel, not a diffusion
Dirichlet identity.

The localized vector is supported in $|a|\le2hR$, so
Theorem~\ref{f16v47:strip} bounds its quotient difference from the tensor
comparison by $C_L(2hR+\varepsilon)\|b\psi\|^2$.
The exact extraction preserves $\|b\psi\|$ and its quadratic form in the
compressed tube. Divide the resulting inequality by $\lambda_*h$, and use
$\varepsilon/h=\sqrt V\sqrt h\le C_L$ eventually. This proves
\eqref{f16v47:local-energy-sharp}, including its dependence on $E_\gamma$.
Its uniform upper bound uses only $E_\gamma\le K$ and $\|b\psi\|\le1$.

V46's exact hard separation and Fourier estimate give the next displays.
The position support of $\vartheta_R f_\gamma$ is bounded; the Fourier bound
makes it precompact by V46's Hilbert--Schmidt phase-space screen. Alternatively
apply its nonlocal compactness theorem. A diagonal extraction over integer
$R$ gives strong local convergence, with consistent limits on overlapping
balls. No derivative of $\psi_\gamma$ was required.
\end{proof}

If $E_\gamma\to-\infty$, nonnegativity of the left side of
\eqref{f16v47:local-energy-sharp} gives the additional quantitative fact
\begin{equation}
 \|\vartheta_RF_\gamma\|^2
 \le\frac{C_L(R^{-2}+h^2)}{|E_\gamma|-C_L(R+1)}
 \longrightarrow0
 \label{f16v47:negative-escape}
\end{equation}
for each fixed $R$, once the denominator is positive. Such a branch, if present,
must therefore escape in soft position; hard or purely local frequency escape
cannot account for it.

\subsection{Only soft position remains in the compactness alternative}

\begin{proposition}[Native profile decomposition in local norm]
\label{f16v47:profiles}
For the eigenvector sequences of Theorem~\ref{f16v47:local-energy}, a subsequence
has
\[
 F_\gamma-\mathcal I f_\gamma\to0\quad\hbox{in }L^2,
 \qquad f_\gamma\to f\quad\hbox{locally in }L^2,
 \qquad \|f_\gamma\|\to1,\quad \|f\|\le1.
\]
If $E_\gamma\to E\in\mathbb R$, V46's corrected profile equation gives
$H_{\rm mat}f=Ef$. If $E_\gamma\to-\infty$, then $f=0$.
The convergence is global and norm preserving precisely when this sequence
is tight in soft position, namely
\begin{equation}
 \lim_{R\to\infty}\limsup_{\gamma\to\infty}
             \int_{|q|>R}|f_\gamma(q)|^2dq=0.
 \label{f16v47:position-tightness}
\end{equation}
No separate Fourier-tightness hypothesis is needed.
\end{proposition}
\begin{proof}
The whole-band theorem applies since $E_\gamma\le K$; it supplies hard capture
and the total soft norm. The local-energy theorem supplies local compactness.
For finite $E$, its limit is also the weak extracted limit, so V46's theorem
applies with its original domain justification. The negative case is
\eqref{f16v47:negative-escape}. Local strong convergence and
\eqref{f16v47:position-tightness} give global strong convergence by a cutoff
and triangle inequality. Conversely, a strongly convergent $L^2$ sequence has
uniformly small tails, and its norm here tends to one. This proves the claim.
\end{proof}

\begin{proposition}[A position-only completeness criterion and an escape witness]
\label{f16v47:completion}
Suppose \eqref{f16v47:position-tightness} holds uniformly, for each fixed $K$,
over the normalized neutral native eigenvectors with $E_\gamma\le K$.
Then the ordered scaling and same-top conclusions in
Proposition~\ref{f16v46:conditional-completeness} follow. This position-tightness
premise is \emph{not proved here}.
In particular, failure of those ordered asymptotics yields, after subsequence
selection, normalized native eigenvectors with $E_\gamma\le K$, radii
$R_j\to\infty$ and $\eta>0$ such that
\begin{equation}
 \int_{|q|>R_j}|f_{\gamma_j}(q)|^2dq\ge\eta.
 \label{f16v47:escape-witness}
\end{equation}
Their hard-ground loss tends to zero, and their soft profiles are compact on
every bounded region. The witness is a statement about actual eigenvectors,
not evidence that such an escaping branch exists.
\end{proposition}
\begin{proof}
V45 bounds above the scaled deficits of every fixed number of ordered neutral
eigenvectors. Position tightness and Proposition~\ref{f16v47:profiles} make
those profiles strongly precompact. A sequence of deficits tending to
$-\infty$ would force every local profile to vanish and contradict norm
preservation. Thus the deficits are also bounded below along those sequences.
Take a common convergent subsequence of any finite orthonormal eigenbasis.
Extraction and hard capture preserve its entire Gram, so the limiting soft
profiles are orthonormal model eigenfunctions. V46's profile equation and
model min--max give the missing ordered liminf; V45 supplies the limsup.
V43 transfers the neutral conclusion to every center sector with its $o(h)$
ladder error, and its exact return metric supplies the neutral $g_\gamma$
conclusion. This is the same spectral implication as V46, now with a single
position-tightness premise instead of its stronger unproved global form
comparison. The contrapositive and the failure of the displayed tail condition
give \eqref{f16v47:escape-witness} by a diagonal choice of radii.
\end{proof}

\paragraph{Why local compactness is not enough.}
Even in the declared comparison a normalized invariant annular profile can be
$f_h(q)=R_h^{-9/2}f(q/R_h)$, where $f$ is smooth radial, supported in
$1<|q|<2$, and $R_h=h^{-1/8}$. Then $f_h\rightharpoonup0$ and all mass escapes
in position, although $\langle f_h,(I-S_h)f_h\rangle\to0$. Indeed the positive
energy identity and $1-e^{-x}\le x$ bound this last quantity by
\[
 C_f\{hR_h^4+hR_h^{-2}+h^3R_h^6\}\longrightarrow0.
\]
This is not a native eigenvector or a counterexample to the remaining estimate.
It shows why the proved unscaled step regularity cannot itself be divided by
$h$ or used as position tightness.

\subsection{Remaining task and verification}

The new native results remove hard-band norm loss and frequency loss at bounded
critical positions from the completeness problem. They do not bound the radial
tail in \eqref{f16v47:position-tightness}. The strip estimate pays a first hard
variation by its absolute norm, giving $C_LR$ in the localized energy; its
constant cannot be sent to infinity. A useful next bound must exploit the
singlet cancellation or a controlled hard-band return on arbitrary states,
while retaining transverse soft confinement along the commuting valleys.
Another fixed finite recovery space would not supply this radial estimate.
This keeps the supplied memoranda's no-vanishing requirement distinct from
residual improvement and from physical transfer calibration.

The diagnostic, audit and reconstruction report accompany the source. Only
this appendix and the literal title version differ from V46; removing them
recovers its exact bytes. Finite checks do not certify the general analytic
proofs or the inherited lineage. The next checkpoint remains V50.

\begin{firewall}
V47 proves a native strip operator-norm comparison, whole-band hard-ground
capture, and local compactness of critical native eigenprofiles. It reduces
the remaining completeness problem to escape in soft position. It does not
prove that radial tightness, a global lower bound at scale $h$, the native
Perron or neutral-gap coefficient, a growing-volume theorem, a physical-time
calibration, or an interacting continuum mass gap. The original transfer,
Gauss constraint, center sectors and returning-state norm remain unchanged.
\end{firewall}
% END F16 V47 APPEND

% BEGIN F16 V48 APPEND -- 2026-09-18
% Insert before the final document-ending command of cumulative V47.
% This insertion leaves every predecessor body and the native operator unchanged.

\clearpage
\section{V48: quadratic hard-return costs and fixed-volume spectral completeness}
\label{f16v48:context}

V47 reduces possible loss of critical spectral mass to escape in soft position.
Its strip estimate pays the linear hard variation in absolute operator norm,
which gives an error proportional to the angular radius. That error competes
with, rather than improves on, transverse confinement. We return to the exact
quotient metric and keep the first magnetic variation as an operator. The
quotient metric has zero first jet in the already chosen exponential normal
slice. The remaining linear operator has zero hard-ground compression. Its
coupling to the separated hard complement costs quadratically in the radius.
On a soft annulus the finite-step transverse confinement is linear in that
radius, so it dominates this quadratic return cost.

A dyadic localization pays the original kinetic cutoff errors without a factor
counting annuli. It supplies the position tightness missing in V47, and hence
the fixed-spatial-regulator ordered spectral limit. The proof does not replace
the Wilson transfer by its comparison, identify the asymptotic parameter with
physical time, or pass to growing spatial volume. Its new analytic inputs are
the second-order quotient-kernel estimate and the complete strip first jet;
the concluding spectral identification also uses V45's corrected core equation
and V46--V47's proved-as-stated profile results.

\subsection{Dependencies and unchanged notation}

The supplied V47 appendix and audit were read in full. V44's exact measured
based-gauge tube, V45's stationarity and hard derivative, V46's finite-step
valley bound, and V42's fixed-neighborhood forbidden norm are retained.
Targeted searches of the current cumulative source and the supplied f14/f15
and Grand Tensor archives were recorded; they are not an exhaustive novelty
or proof audit. The two import memoranda remain methodological proposals, not
premises supplying any missing estimate. The next scheduled checkpoint is V50.

Normal-slice geometry, quadratic elimination of an off-diagonal perturbation,
and compactness plus min--max are classical mechanisms. For context,
M\'atyus--Teufel, \href{https://arxiv.org/html/1904.00042}{\emph{Effective
non-adiabatic Hamiltonians}}, Section III, distinguishes diagonal compression
from the effect of the omitted coupling. Hanche-Olsen--Holden,
\href{https://arxiv.org/html/0906.4883}{\emph{The Kolmogorov--Riesz compactness
theorem}}, Corollary 7, gives the position/Fourier criterion already used in
V46--V47. No molecular gap assumption or quantitative compactness rate is
imported. The metric and kernel calculations needed here are given directly
for the native based action.

Fix $L=2m$, $m\ge4$, and $V=L^3$. Keep
\[
 T=T_\gamma=M_{b_\gamma}C_\gamma^{\otimes3V}M_{b_\gamma},\qquad
 b_\gamma=e^{-S/(2\varepsilon^2)},\qquad
 \varepsilon=\gamma^{-1/2},\quad h=(\gamma V)^{-1/3},\quad
 \varepsilon^2=Vh^3.
\]
All constants may depend on this fixed $L$. Write $\lambda_*=\Lambda_{{\rm G},B}$,
not an assigned native top. In the unchanged tube put
\[
 s=y+\mathsf K a,\quad y\in\mathcal Y,\quad a\in\mathbb R^9,\quad
 \mathsf K^*\mathsf K=VI,\quad y=\varepsilon u,\quad a=hq,
 \qquad dH=dg\,J_L(s)\,ds.
\]
Only the based action is removed. Residual simultaneous conjugation is imposed
on functions, not eliminated by a singular orientation slice. Let
\[
 \mathfrak P=|\varphi\rangle\langle\varphi|\otimes I,\quad
 \mathcal T_+\varphi=\lambda_*\varphi,\quad
 \kappa_+=\lambda_*\bigl(1-r_+\bigr)>0,
\]
where $r_+=\|\mathcal T_+|_{\varphi^\perp}\|/\lambda_*$. Retain
\[
 S_h=e^{-h\mathcal Q}e^{h\Delta_q/2}e^{-h\mathcal Q},\qquad
 \mathcal Q(q)=\sum_{i<j}|q_i\times q_j|^2.
\]
This is V46's explicit comparison. It is never declared to be a native
conditional transition. Write $K_\gamma^{\rm red}$ for the one-tube compression
in the exact flat half-density and critical coordinates, extended by zero in
the hard variable outside its fixed physical tube. Sector-folded compressions
have the same estimates below, with exponentially small cross-tube terms.

\subsection{The quotient metric has no linear error}

The following calculation strengthens the first-order bound in V47 without
changing its coordinates. Let $X_s$ be the differential of
$s\mapsto\operatorname{Exp}(s)$ on $\mathcal N$, and let $Y_s$ be the differential
of the based action at that point, using a fixed basis of its Lie algebra.
Inner products in the target are the original product round metric. Set
\begin{equation}
 V_s=Y_s^*Y_s,\qquad C_s=X_s^*Y_s,\qquad
 \mathsf g_s=X_s^*X_s-C_sV_s^{-1}C_s^*.
 \label{f16v48:quotient-metric}
\end{equation}
The letter $V_s$ here is a vertical Gram matrix, not the number of vertices.

\begin{proposition}[Second-order metric and the normalized quotient kernel]
\label{f16v48:metric}
On one sufficiently small fixed tube,
\begin{equation}
 \mathsf g_0=I,\qquad D\mathsf g_0=0,\qquad
 \|\mathsf g_s-I\|+\|\mathsf g_s^{-1}-I\|\le C_L|s|^2.
 \label{f16v48:metric-jet}
\end{equation}
For the actual quotient kernel of V47 there are fixed $C_L,c_L>0$ and $M_L$
such that
\begin{align}
 |k_\varepsilon^{\rm red}(s,t)-g_\varepsilon(s-t)|
 &\le C_L(|s|^2+|t|^2+\varepsilon)\varepsilon^{-d}
 (1+|s-t|/\varepsilon)^{M_L}e^{-c_L|s-t|^2/\varepsilon^2}
 \nonumber\\
 &\quad+C_L\varepsilon^{-D}e^{-c_L/\varepsilon^2},
 \label{f16v48:kinetic-error}
\end{align}
where $d=6V+3$, $D=9V$ and
$g_\varepsilon(\xi)=(2\pi\varepsilon^2)^{-d/2}e^{-|\xi|^2/(2\varepsilon^2)}$.
Both Gaussian-polynomial Schur integrals are uniformly bounded. In V45's
specified normal slice the complete horizontal covariance is
$\mathsf g_s^{-1}$; consequently its hard block satisfies $DA(0)=0$.
This last identity refines, rather than assumes away, the covariance derivative
that V45 retained.
\end{proposition}
\begin{proof}
Based freeness gives $V_0\succ0$ and a uniform inverse on a small fixed tube.
At the identity the slice is orthogonal to the orbit, so $C_0=0$. The product
exponential is a Riemannian normal chart at the identity. More explicitly, on
one round $\SU(2)$ factor its pulled-back metric is one in the radial direction
and $(\sin|s|/|s|)^2$ in the two transverse directions. Thus
$X_s^*X_s=I+O_L(|s|^2)$, including after restriction to $\mathcal N$.
Smoothness gives $C_s=O_L(|s|)$. Equation \eqref{f16v48:quotient-metric}
proves \eqref{f16v48:metric-jet}. The covariance assertion follows because the
differential of the quotient is an isometry on the horizontal tangent and
annihilates the vertical tangent. Equivalently its product with its adjoint is
the inverse Schur metric. No freeness of the residual global conjugation is
used.

Here are the measure details needed to improve a kernel estimate as well.
Take the same fixed local based-group coordinates $v$ as in V47, with Haar
density $c_G$ at the identity, and let $c_X$ be the full probability-Haar density
in an orthonormal tangent frame. The exact Jacobian identity at $(1,s)$ is
\begin{equation}
 J_L(s)c_G=c_X\sqrt{\det V_s}\sqrt{\det\mathsf g_s}.
 \label{f16v48:coarea}
\end{equation}
It is the determinant of the Gram matrix of $[Y_s\ X_s]$, followed by its
Schur determinant. This retains the varying orbit volume; $J_L(s)$ is not
assumed flat or constant.

Put $t=s+\xi$. Taylor expansion of the actual Wilson displacement cost in
$(\xi,v)$ gives its exact quadratic part
\[
 A_s(\xi,v)=\tfrac12\|X_s\xi+Y_sv\|^2+
ho_s(\xi,v),\qquad
 |\rho_s(\xi,v)|\le C_L(|\xi|+|v|)^3.
\]
A sign change of $v$ would give the alternative convention. On the selected
small chart both this quadratic form and the exact cost are bounded below by
$c_L(|\xi|^2+|v|^2)$. All constants are uniform in $s$. Complete the vertical
square by $v\mapsto v+V_s^{-1}C_s^*\xi$. Its integral has determinant
$(\det V_s)^{-1/2}$ and horizontal quadratic form $\xi^*\mathsf g_s\xi$.
Use \eqref{f16v48:coarea} and the two exact half-densities. The resulting
principal kernel is
\[
 \sqrt{J_L(t)/J_L(s)}\sqrt{\det\mathsf g_s}\,
 (2\pi\varepsilon^2)^{-d/2}
 e^{-\xi^*\mathsf g_s\xi/(2\varepsilon^2)}.
\]
Thus neither a vertical determinant nor an endpoint density remains as an
uncomputed order-$|s|$ scalar. The density ratio differs from one by
$O_L(|\xi|)$. Smooth based Haar factors differ from $c_G$ by $O_L(|v|)$.
The true kinetic normalizer has relative error $O_L(\varepsilon^2)$.
The cubic cost remainder divided by $\varepsilon^2$, integrated against a
positive Gaussian majorant, gives $\varepsilon$ times a fixed polynomial in
$|\xi|/\varepsilon$. The same conclusion follows directly from the mean-value
identity for two exponentials; no small supremum of their extensive exponents
is required.

Finally \eqref{f16v48:metric-jet} bounds the difference of the principal
Gaussian from $g_\varepsilon$ by $C_L|s|^2$ times such a polynomial majorant.
For $s,t$ in the fixed tube and $v$ outside the chosen based chart, compactness
and freeness give the same positive displacement floor as V47; its entire
contribution is the last term of \eqref{f16v48:kinetic-error}. Integrating the
majorants in either endpoint proves their Schur bounds. This proves an
absolute kernel estimate, not a Loewner inference from entrywise signs.
\end{proof}

\subsection{Keep the linear hard operator until its return is bounded}

Write $H(a)=H_++\sum_{\ell=1}^9a_\ell H_\ell+O_L(|a|^2)$ for the actual hard
Hessian in V45. Define hard multiplication operators
\[
 B_\ell(u)=\tfrac14u^*H_\ell u.
\]
The products $B_\ell\mathcal T_+$ and $\mathcal T_+B_\ell$ are bounded
Hilbert--Schmidt operators. Let $M_\ell$ mean multiplication by $hq_\ell$ in
the soft variable. On a supported angular strip define the bounded first jet
\begin{equation}
 \mathcal L_\gamma=-\sum_{\ell=1}^9
 \left[B_\ell\mathcal T_+\otimes(M_\ell S_h)
       +\mathcal T_+B_\ell\otimes(S_hM_\ell)\right].
 \label{f16v48:linear-operator}
\end{equation}
Products with the strip projection on both sides are understood; this avoids
asserting boundedness of $M_\ell$ on the whole soft space.

\begin{theorem}[Complete strip first jet with a quadratic remainder]
\label{f16v48:strip-jet}
Let $P_r$ multiply by $|hq|\le r$ on profile space, with no hard-coordinate
cutoff. There is a fixed $r_*>0$ such that, for $0<r\le4r_*$ and small
$\varepsilon$,
\begin{equation}
 \left\|P_r(K_\gamma^{\rm red}-\mathcal T_+\otimes S_h
                    -\mathcal L_\gamma)P_r\right\|
                 \le C_L(r^2+\varepsilon).
 \label{f16v48:quadratic-remainder}
\end{equation}
Moreover
\begin{equation}
 \|P_r\mathcal L_\gamma P_r\|\le C_Lr,\qquad
 \mathfrak P\mathcal L_\gamma\mathfrak P=0
 \label{f16v48:linear-cancellation}
\end{equation}
on every such strip. These are all-vector operator estimates, not core-wise
expansions or assertions that the off-diagonal first jet vanishes.
\end{theorem}
\begin{proof}
Exact hard stationarity and the uniform hard pin give, for $y=\varepsilon u$,
\[
 b_\gamma(y+\mathsf Ka)
 =m_{\gamma,c}(a)e^{-u^*H_+u/4}
       \left(1-\sum_\ell a_\ell B_\ell(u)\right)
       +\mathcal E_\gamma(a,u),
\]
where $m_{\gamma,c}(a)=e^{-S(\operatorname{Exp}(\mathsf Ka))/(2\varepsilon^2)}$
and
\[
 |\mathcal E_\gamma(a,u)|
 \le C_L(r^2+\varepsilon)m_{\gamma,c}(a)
                    (1+|u|)^M e^{-c_L|u|^2},\qquad |a|\le r.
\]
To obtain this bound on the whole hard tube, first compare the exact multiplier
with the one using $H(a)$, charging the cubic hard remainder by
$C_L\varepsilon|u|^3e^{-c_L|u|^2}$. Expand the latter Gaussian twice in $a$;
uniformly positive intermediate Hessians bound its second derivatives by a
polynomial times a Gaussian. This yields the $r^2$ term without requiring
$|u|$ bounded. The omitted hard exterior has exponentially small Gaussian
operator tails at the fixed physical hard radius.

The exact constant commutator formula gives the supremum multiplier bound
\[
 \sup_{|a|\le r}
 |m_{\gamma,c}(a)-e^{-h\mathcal Q(a/h)}|\le C_Lr^2.
\]
Indeed its action lies between $(1-C_Lr^2)\mathcal Q(a)$ and
$\mathcal Q(a)$, with the common factor $\gamma V$, and
$xe^{-cx}\le C$ pays arbitrarily large quartic exponents. This replacement
in a linear term costs at most $C_Lr^3$.

Sandwich \eqref{f16v48:kinetic-error} by the native magnetic factors. The hard
pin turns each $|y|^2$ factor into $O_L(\varepsilon^2)$ in its row/column bound;
the soft contribution is $O_L(r^2)$. The resulting norm error is
$C_L(r^2+\varepsilon)$. Expanding the two magnetic factors around the split
Gaussian now gives exactly \eqref{f16v48:linear-operator}; their product of
linear corrections is $O_L(r^2)$ in Schur norm. The hard quadratic polynomial
factors are controlled by their Gaussian kernels, while the soft Gaussian
step and both quartic multipliers are contractions. This proves
\eqref{f16v48:quadratic-remainder}, including zero extension in the hard
coordinate. No factor is proportional to the soft strip volume. Center-folded
cross terms are $O_L(e^{-c_L/\varepsilon^2})$ and are absorbed as well.

The norm bound in \eqref{f16v48:linear-cancellation} follows from
$\|P_rM_\ell\|\le r$. For the ground compression, the scalar
\[
 a\longmapsto\int\tfrac14u^*\Bigl(\sum a_\ell H_\ell\Bigr)u
                                      |\varphi(u)|^2du
\]
is a simultaneous-color-invariant linear functional on three vectors. It is
zero. Hence $\langle\varphi,B_\ell\varphi\rangle=0$ for each $\ell$ and
$P_+B_\ell\mathcal T_+P_+=P_+\mathcal T_+B_\ell P_+=0$.
This proves the exact compression cancellation for every pair of soft
endpoints. It does not require their equality or slow variation of a test.
\end{proof}

\begin{proposition}[Quadratic hard-return cost and annular coercivity]
\label{f16v48:annulus}
After a fixed decrease of $r_*$, there are $c_L,C_L>0$ such that every physical
vector supported in one tube with $r\le|a|\le4r$, $0<r\le r_*$, satisfies
\begin{equation}
 \boxed{\quad
 \langle\psi,(\lambda_*-T)\psi\rangle
            \ge(c_Lr-C_L\varepsilon)\|\psi\|^2.
 \quad}
 \label{f16v48:annular-floor}
\end{equation}
Every physical vector supported in $|a|\le r\le4r_*$ satisfies
\begin{equation}
 \langle\psi,(\lambda_*-T)\psi\rangle
                   \ge-C_L(r^2+\varepsilon)\|\psi\|^2.
 \label{f16v48:inner-floor}
\end{equation}
The supports in the hard variable lie in the same fixed small tube. Neither
estimate requires hard-ground support or a bound on derivatives of $\psi$.
\end{proposition}
\begin{proof}
In the exact profile coordinates write $F=\mathcal I f+F_\perp$. Soft strip
and annulus projections commute with $\mathfrak P$. The reference has hard
deficit at least $\kappa_+\|F_\perp\|^2$, and its ground deficit is
$\lambda_*\langle f,(I-S_h)f\rangle$. By
\eqref{f16v48:linear-cancellation}, the entire first-jet contribution is bounded
in absolute value by
\[
 C_Lr\|F_\perp\|^2+2C_Lr\|f\|\|F_\perp\|.
\]
Young's inequality charges the mixed term to
$\kappa_+\|F_\perp\|^2/4+C_Lr^2\|f\|^2$.
Choose the fixed radius small enough to absorb the first term as well.
Together with \eqref{f16v48:quadratic-remainder}, this proves
\begin{equation}
 \langle F,(\lambda_*-K_\gamma^{\rm red})F\rangle
 \ge\lambda_*\langle f,(I-S_h)f\rangle
     +\tfrac12\kappa_+\|F_\perp\|^2
     -C_L(r^2+\varepsilon)\|F\|^2
 \label{f16v48:block-payment}
\end{equation}
for a strip of radius bounded by a fixed multiple of $r$.

On the annulus V46's exact finite-step bound gives
\[
 \langle f,(I-S_h)f\rangle
 \ge\tfrac16\int\sum_i\min\{|a_i|,1\}|f(q)|^2dq
 \ge\tfrac r6\|f\|^2,
\]
after choosing $4r_*<1$. The hard complementary constant is also at least a
fixed constant times $r$ on this interval. A further fixed decrease of $r_*$
absorbs $C_Lr^2$ into half that radial lower bound, giving
\eqref{f16v48:annular-floor}. Without the annulus, discard the nonnegative soft
term in \eqref{f16v48:block-payment} to get \eqref{f16v48:inner-floor}.
Exact half-density isometry returns both statements to the native Hilbert
space. This is a paid quadratic coupling estimate; no native hard inverse or
unproved global soft gap has been used.
\end{proof}

\subsection{Dyadic localization pays every annulus once}

Choose fixed, center-translated, gauge-invariant tube cutoffs $\zeta_\sigma$,
with disjoint supports, equal to one on smaller neighborhoods of the central
orbits. Their supports have $|a|<r_*$ and lie in the plateau of V47's final
extraction cutoff. They may be chosen with a smooth exterior member
$\zeta_{\rm out}$ so that
\[
 \zeta_{\rm out}^2+\sum_\sigma\zeta_\sigma^2=1.
\]
For example use a sine/cosine transition on each of the disjoint tubes.
V42's supported forbidden-norm result, applied to their fixed plateaux, gives
\begin{equation}
 \langle\zeta_{\rm out}\psi,(\lambda_*-T)\zeta_{\rm out}\psi\rangle
                 \ge\delta_L\|\zeta_{\rm out}\psi\|^2
 \label{f16v48:outer-floor}
\end{equation}
eventually, with $\delta_L>0$ independent of coupling. No shrinking-neighborhood
version of that result is required. Define, for $0<\rho<r_*/8$,
\[
 \mathcal W_\rho[\psi]
 =\sum_\sigma\int \zeta_\sigma(U)^2|a_\sigma(U)|
       1_{\{|a_\sigma(U)|>4\rho\}}|\psi(U)|^2\,dH(U).
\]
The constant angles $a_\sigma$ are the exact translated tube coordinates.

\begin{theorem}[A global native radial inequality]
\label{f16v48:radial-form}
For all physical $\psi\in L^2$, all $0<\rho<r_*/8$, and sufficiently large
coupling, with the threshold independent of $\rho$,
\begin{equation}
 \boxed{\quad
 c_L\mathcal W_\rho[\psi]+\delta_L\|\zeta_{\rm out}\psi\|^2
 \le\langle\psi,(\lambda_*-T)\psi\rangle
       +C_L\left(\rho^2+\varepsilon+
                    \frac{\varepsilon^2}{\rho^2}\right)\|\psi\|^2.
 \quad}
 \label{f16v48:radial-form-bound}
\end{equation}
There is no factor counting the dyadic annuli. In particular
\begin{equation}
 h^{-1}(I-T/\lambda_*)\succeq-C_LI
 \label{f16v48:semibounded}
\end{equation}
eventually on the complete physical space.
\end{theorem}
\begin{proof}
Inside each tube use a smooth radial quadratic partition at scales
$\rho,2\rho,4\rho,\ldots$, restricted to the support of $\zeta_\sigma$.
The first member is supported on $|a|<2\rho$; each other member is supported
on $r_j\le|a|\le4r_j$, with $r_j$ a dyadic multiple of $\rho$ and
$r_j\le r_*$. Start from fixed radial bump functions with finite overlap and
divide by the square root of the sum of their squares. On every point where
a cutoff is needed that denominator has a fixed positive lower bound. Their
squares sum to one, their gradients obey the pointwise bound
$\sum_j|\nabla\theta_j|^2\le C_L\rho^{-2}$, and
\[
 \sum_{j\ge1}r_j\theta_j(a)^2
       \ge c|a|1_{\{|a|>4\rho\}}.
\]
Only finitely many members meet the fixed tube; their number is not used in
these estimates. Multiply them by the fixed $\zeta_\sigma$ and add
$\zeta_{\rm out}$ to get a global quadratic partition $(\chi_j)$.
The smooth vector-valued map $U\mapsto(\chi_j(U))_j$ has Lipschitz constant
at most $C_L/\rho$, independent of the number of its components. This follows
by integrating its pointwise squared-gradient bound along a minimizing path.

The exact native localization error is the kernel
\[
 T-\sum_j\chi_jT\chi_j
   =\tfrac12\sum_j(\chi_j(U)-\chi_j(U'))^2t_\gamma(U,U').
\]
Its operator norm is at most $C_L\varepsilon^2/\rho^2$ by the original
normalized displacement second moment and the row/column Schur bound.
Its Loewner sign is not asserted. Apply \eqref{f16v48:inner-floor} to the
inner members, \eqref{f16v48:annular-floor} to every annular member, and
\eqref{f16v48:outer-floor} to the exterior member. The $C_L\varepsilon$
errors sum against squares of a partition, hence to at most
$C_L\varepsilon\|\psi\|^2$, not their number times this norm.
The positive annular terms dominate $c_L\mathcal W_\rho$, and the inner loss
is at most $C_L\rho^2\|\psi\|^2$. This proves the displayed inequality.
Take $\rho=h$ and use $\varepsilon^2=Vh^3$, $\varepsilon=o(h)$, to obtain
\eqref{f16v48:semibounded}. No eigenvector equation was used in this theorem.
\end{proof}

\begin{theorem}[Native position tightness on the whole critical band]
\label{f16v48:tightness}
For fixed $K\ge0$ let
$\mathsf B_{\gamma,a}(K)=1_{[\lambda_*(1-Kh),1]}(T|_{\mathcal H_a})$.
Retain V47's unchanged extraction $F_\gamma=\mathfrak E_{\gamma,a}\psi$ and
soft profile $f_\gamma=\mathcal I^*F_\gamma$. Then, uniformly over unit vectors
in this band and every sector,
\begin{equation}
 \boxed{\quad
 \limsup_{\gamma\to\infty}
 \sup_{\psi\in\operatorname{Ran}\mathsf B_{\gamma,a}(K),\ \|\psi\|=1}
 \int_{|q|>4R}|F_\gamma(u,q)|^2\,du\,dq
       \le C_L\left(\frac K R+\frac1{R^3}\right),\qquad R\ge1.
 \quad}
 \label{f16v48:tail}
\end{equation}
The same bound holds for $|f_\gamma(q)|^2$. In fact, at one common sufficiently
large coupling threshold for fixed $L,K$,
\begin{equation}
 \sup_{\psi\in\operatorname{Ran}\mathsf B_{\gamma,a}(K),\ \|\psi\|=1}
 \int |q|\,|F_\gamma(u,q)|^2\,du\,dq\le C_L(K+1).
 \label{f16v48:first-moment}
\end{equation}
An empty band contributes zero. In particular V47's uniform
position-tightness condition is satisfied.
\end{theorem}
\begin{proof}
The band gives
$\langle\psi,(\lambda_*-T)\psi\rangle\le\lambda_*Kh$.
Set $\rho=hR$ in \eqref{f16v48:radial-form-bound}; at each fixed $R$ it is
admissible eventually. The right side is at most
\[
 h\left[\lambda_*K+C_L(hR^2+\sqrt h+R^{-2})\right],
\]
with the fixed $\sqrt V$ absorbed into $C_L$. On the part of an extraction
tube where $|q|>4R$, its $\zeta_\sigma^2$-weighted mass is bounded by
$\mathcal W_{hR}/(4hR)$. Its remaining mass is bounded by
$\|\zeta_{\rm out}\psi\|^2$, because the fixed cutoff squares partition one.
This argument also covers the portions of the unchanged extraction tube
outside the smaller cutoffs used in this proof. Both terms are controlled by
\eqref{f16v48:radial-form-bound}. Divide the first by $hR$, let
$\gamma\to\infty$ at fixed $R$, and note that the second tends to zero.
Exact extraction and center normalization give \eqref{f16v48:tail}.
Cauchy--Schwarz in the hard variable gives the same bound for $f_\gamma$.
Finally send $R\to\infty$. For \eqref{f16v48:first-moment}, use
$\rho=h$ instead. The same inequality bounds $\mathcal W_h/h$ and
$h^{-1}\|\zeta_{\rm out}\psi\|^2$ by $C_L(K+1)$. In the inner part
$|q|\le4$, while on the fixed extraction support $|q|\le C_L/h$.
Split by the same quadratic cutoff pair; these two bounds pay its entire
first moment. No support restriction is imposed on the original band vector
and no rank is counted.
\end{proof}

\subsection{The native fixed-volume spectrum now has the model coefficients}

Let $e_0<e_1\le e_2\le\cdots$ be the invariant eigenvalues of the unchanged
$H_{\rm mat}$, with multiplicity. Their existence and $e_1-e_0>0$ are V44's
model theorems; they are not numerical evaluations. Index native eigenvalues
from one in each center sector, so $\lambda_{0,1}=\Lambda_\gamma$.

\begin{theorem}[Ordered critical spectrum and the same-top neutral gap]
\label{f16v48:spectrum}
At every fixed spatial regulator and every fixed $j\ge1$,
\begin{equation}
 \boxed{\quad
 \frac{1-\lambda_{a,j}(\gamma)/\lambda_*}{h}\longrightarrow e_{j-1}
 \quad\hbox{in every center sector}.\quad}
 \label{f16v48:levels}
\end{equation}
The ordered one-region eigenvalues of V43 have the same limits. In particular
\begin{equation}
 \begin{split}
 \Lambda_\gamma&=\lambda_*[1-he_0+o_L(h)],\\
 \Lambda_\gamma-\lambda_{0,2}
        &=\lambda_*h(e_1-e_0)+o_L(h),\\
 -\log(\lambda_{0,2}/\Lambda_\gamma)
        &=h(e_1-e_0)+o_L(h),\\
 g_\gamma/(\lambda_*h)&\longrightarrow e_1-e_0.
 \end{split}
 \label{f16v48:gap}
\end{equation}
The last quantity is V43's actual reconstruction-metric gap, not a bare
regional compression. These statements fix $L$ and $j$ before the limit.
\end{theorem}
\begin{proof}
V45 bounds above the scaled deficits of any fixed number of ordered neutral
eigenvectors. Equation \eqref{f16v48:semibounded} bounds them below. On any
subsequence, choose a further subsequence on which those finitely many
deficits converge. V47 gives hard capture and strong local compactness of
their profiles; Theorem~\ref{f16v48:tightness} upgrades this to global strong
compactness. Extraction is asymptotically isometric on the native band, so
all inner products of this finite eigenbasis are preserved in the limit.
V46's corrected profile equation therefore produces orthonormal invariant
model eigenfunctions with these limiting deficits. Model min--max gives the
ordered liminf $e_{j-1}$; V45 gives the matching limsup. This proves the
neutral limits on every subsequence and hence in full. It is exactly the
previously conditional position-only completion argument of V47, with its
hypothesis now supplied by a native inequality.

Every fixed level lies in V43's matching interval by V44. Its exponential
sector and one-region differences are $o(h)$, so all the other limits follow.
Subtract the first two neutral expansions and use $-\log(1-x)=x+O(x^2)$.
V43--V44's exact $\Delta_\gamma^{\rm neu}/g_\gamma\to1$ gives the final
line of \eqref{f16v48:gap}. No identity metric is installed on the entire
regional space.
\end{proof}

\begin{proposition}[Exact low-cluster counts and critical vacuum profiles]
\label{f16v48:clusters}
For each fixed $E\notin\operatorname{spec}(H_{\rm mat}|_{\rm inv})$, the number
of eigenvalues in one sector satisfying
$(1-\lambda/\lambda_*)/h<E$ is eventually exactly the number of invariant
model eigenvalues below $E$. On isolated finite model eigenspaces, subsequential
extracted orthonormal native bases converge to orthonormal model bases. For
the actual neutral vacuum the whole extracted sequence satisfies
\begin{equation}
 \mathfrak E_{\gamma,0}\Omega_\gamma
       \longrightarrow\varphi\otimes\phi_0
       \quad\hbox{in }L^2(du\,dq),
 \label{f16v48:vacuum-profile}
\end{equation}
where $\phi_0$ is the unit positive invariant ground function of $H_{\rm mat}$.
Consequently the corresponding extracted density converges in $L^1$ to
$|\varphi|^2|\phi_0|^2$. This is a statement after the specified center
unfolding and critical coordinate scaling, not an unscaled $L^2$ limit of
native vacuum vectors.
\end{proposition}
\begin{proof}
Choose model indices immediately below and above the fixed threshold and use
\eqref{f16v48:levels}. Since ordered native eigenvalues exhaust the spectrum
above any positive threshold, these two indices give the exact count; possible
negative scaled deficits have already been bounded below. For a finite cluster,
extract a basis along subsequences as in the preceding proof. Compactness,
Gram preservation and the profile equation identify its limiting eigenspace
and multiplicity. For the ground, simplicity and positivity remove phase and
basis ambiguity: extraction, neutral unfolding and $\varphi$ are nonnegative.
Every subsequential limit is the same positive unit $\phi_0$. This proves
full-sequence convergence. The $L^1$ claim follows from
$\||F|^2-|G|^2\|_1\le(\|F\|+\|G\|)\|F-G\|$.
\end{proof}

\begin{proposition}[Two spectral scales and recentering the entrance law]
\label{f16v48:consequences}
Let $g_*=e_1-e_0>0$. For sufficiently large coupling the full physical transfer
has exactly eight eigenvalues above
\begin{equation}
 z_\gamma=\Lambda_\gamma-\tfrac12\lambda_*h g_*.
 \label{f16v48:eight-cut}
\end{equation}
There is one in each center sector. Their charged top splittings retain V42's
upper bound $C_Le^{-\kappa_L\sqrt\gamma}$, while their next sector levels are
separated from the top by $\lambda_*h g_*+o_L(h)$.

On every fixed-$B$ duration sequence admitted by V40, its projected spectral
law may now be centered at the actual Perron value:
\begin{equation}
 -n_\gamma\log(\lambda/\Lambda_\gamma)
       \ \Rightarrow\ \operatorname{Gamma}(9/2,1),\qquad
 n_\gamma(1-\lambda/\Lambda_\gamma)
       \ \Rightarrow\ \operatorname{Gamma}(9/2,1).
 \label{f16v48:recentered}
\end{equation}
The probability is V40's unchanged projected endpoint-weighted spectral
probability, not the stationary density of states.
\end{proposition}
\begin{proof}
The first two ordered sector limits and V42's same-top exponentially small
charged bounds place the first level of every sector above
\eqref{f16v48:eight-cut} and its second below. Ordered monotonicity excludes
any further level above that threshold. No tunneling prefactor or matching
charged lower bound is deduced.

V40's regime gives $n_\gamma h\to0$ at fixed $B$. The new top expansion gives
$n_\gamma\log(\Lambda_\gamma/\lambda_*)\to0$ and
$\Lambda_\gamma/\lambda_*\to1$. The logarithmic deficit differs from V40's
by the first deterministic vanishing quantity. The linear deficit is its
old value divided by $\Lambda_\gamma/\lambda_*$ plus a deterministic
$O(n_\gamma h)$ term. The probability limits follow. Unlike the deficits
centered at $\lambda_*$, these actual-top deficits are nonnegative at every
finite coupling. The probability measure and all endpoint normalizations
are unchanged.
\end{proof}

\clearpage
\subsection{Scope, next estimate, and validation}

The new conclusion is fixed-volume spectral completeness for the \emph{existing}
critical branch. The proof rests on three inputs of different kinds: the new
quotient-metric/first-jet and annular estimates, V42's supported forbidden-region
bound, and V45--V47's corrected recovery and profile results. It does not
independently recertify the entire analytic lineage. In particular the sharp
coefficients use V45's strong core equation as an inherited premise; the
new coarse bounds supply compactness rather than recomputing that equation.

There is no uniformity here as $L\to\infty$, no numerical enclosure of
$e_0,e_1$ or of the native thresholds, and no rate in the remaining $o_L(h)$.
The model and native kinetic normalizations are exactly those already fixed.
For a separately supplied physical time spacing $a_t(\gamma,L)$ the neutral
energy has the expression
\[
 E_{\rm neu}^{\rm phys}
   =\frac{(\gamma L^3)^{-1/3}}{a_t(\gamma,L)}
        [e_1-e_0+o_L(1)].
\]
No value of $a_t$ is chosen by this formula, and a fixed number of lattice
sites is not a fixed-physical-volume continuum limit. The next substantive
problem is uniform control of this comparison along a specified growing
regulator and physical-time trajectory, or a quantitative higher-order and
tunneling analysis at fixed regulator. Another fixed-$L$ recovery space is no
longer the missing completeness step. The next scheduled checkpoint is V50.

The diagnostic records exact finite tests of the quotient Schur metric,
coarea cancellation, color-singlet compression, quadratic hard-return cost,
annular inequalities, dyadic overlap, critical powers, and ordered spectral
consequences. These are not native spectral enclosures or formal verification
of the analytic kernel remainder. The audit records which earlier scripts
were actually executed. Removing this exact insertion and reversing the
literal title version recovers V47 byte for byte.

\begin{firewall}
V48 proves a quadratic-radius hard-return estimate, a native annular barrier,
and radial tightness at the critical scale. Together with the explicitly
identified inherited recovery and profile theorems this yields the ordered
fixed-spatial-regulator spectrum and its center-neutral gap coefficient.
It also separates the eight exponentially split sector tops from the next
algebraic sector levels. It does not establish a volume-uniform estimate,
calibrated physical mass, interacting continuum construction, or a matching
center tunneling asymptotic. No comparison operator or auxiliary clock
replaces the native transfer.
\end{firewall}
% END F16 V48 APPEND

% BEGIN F16 V49 APPEND -- 2026-09-18
% Insert before the final document-ending command of cumulative V48.
% This insertion leaves every predecessor body and the native operator unchanged.

\clearpage
\section{V49: symmetry sectors and a certified bracket for the critical gap coefficient}
\label{f16v49:context}

V48 identifies the fixed-regulator critical spectrum with the invariant
eigenvalues $e_0<e_1\le e_2\le\cdots$ of the quartic model
$H_{\rm mat}=-\frac12\Delta_q+2\mathcal Q(q)$, and the same-top neutral gap
coefficient with $g_*=e_1-e_0$. It leaves those numbers unenclosed: the only
recorded constants are V45's Gaussian upper bounds $27/4$ and $47/4$, and the
positivity $g_*>0$ of V44. This note supplies the missing quantitative
control at the model level. It records the exact symmetry of $H_{\rm mat}$
and the resulting sector structure of the invariant space, proves a
pairwise transverse lower bound that improves V44's oscillator floor,
certifies the first Airy zero from below by an exact power-series shooting
argument, and certifies Rayleigh--Ritz upper bounds in an exact rational
Gaussian-polynomial basis. The outcome is
\[
 5.547\le e_0\le 6.5489,\qquad e_1\le e_5\le 9.6678,\qquad
 e_6\le 10.5318,\qquad 0<g_*\le 4.1206,
\]
with the first excitation appearing, in the Rayleigh--Ritz spectra, as a
five-fold multiplet of the residual spatial rotation symmetry.
No native quantity is recomputed; the native statements of V48 are
sharpened only through these model constants.

\subsection{Dependencies and what is not claimed}

The supplied V48 appendix was read in full. The model inputs are V44's
confinement theorem (compact resolvent, simple positive ground function,
$e_1-e_0>0$), V45's Gaussian constants, and V48's identification of the
ordered native levels with $e_{j-1}$. No native transfer estimate is used or
modified. The lower bound below is not a two-sided enclosure of $g_*$: no
certified lower bound on $e_1$ better than $e_1\ge e_0$ is obtained, so the
bracket for the gap coefficient is one-sided from above and uses V44's
qualitative positivity from below. Temple-type lower bounds would require an
a priori lower bound on $e_1$ exceeding the Rayleigh quotient of the ground
trial; the separable comparison operators available here have second
invariant level about $5.84$, below the certified ground bound, so they do
not supply it. This is stated as the next model estimate, not solved.

For context, L\"uscher's small-volume analysis
(\href{https://doi.org/10.1016/0550-3213(83)90436-4}{Nucl.\ Phys.\ B219
(1983)}) and L\"uscher--M\"unster's weak-coupling expansion
(\href{https://doi.org/10.1016/0550-3213(84)90038-5}{Nucl.\ Phys.\ B232
(1984)}) study the constant-mode quartic Hamiltonian of $\SU(2)$ gauge theory
on a spatial torus, whose low levels scale as $g^{2/3}/L$. The operator
$H_{\rm mat}$ is of that type after V44's rescaling. No coefficient, level, or
time normalization from those papers is imported; the present enclosures are
computed and certified here, in the normalization fixed by V44--V48. The
comparison of orderings below is an observation about the present numbers.

\subsection{Bi-invariance and the sector structure of the invariant space}

Write $q=(q_1,q_2,q_3)$ as the $3\times3$ matrix with columns $q_i\in\mathbb
R^3$, and $G=q^{\mathsf T}q$ for its Gram matrix, $G_{ij}=q_i\cdot q_j$.
The exact quaternion identity of V44 gives
\begin{equation}
 \mathcal Q(q)=\sum_{i<j}|q_i\times q_j|^2
 =\sum_{i<j}\bigl(G_{ii}G_{jj}-G_{ij}^2\bigr)
 =\tfrac12\bigl[(\Tr G)^2-\Tr G^2\bigr]=\sigma_2(G),
 \label{f16v49:Q-gram}
\end{equation}
the second elementary symmetric function of the eigenvalues of $G$.

\begin{proposition}[Bi-invariance and sectors]
\label{f16v49:symmetry}
$H_{\rm mat}$ commutes with the left action $q\mapsto Rq$ of $O(3)$ on the
color index, with the right action $q\mapsto qO$ of $O(3)$ on the direction
index, and in particular with every sign change $q_i\mapsto-q_i$ and every
permutation of the three directions. The residual Gauss condition of V44 is
invariance under the left $\operatorname{SO}(3)$ action. Every function of the form
$F(G,\det q)$ is left-$\operatorname{SO}(3)$-invariant. The invariant space decomposes
into the isotypic components of the right $O(3)$ action, and the ground
function $\phi_0$ of V44 lies in the trivial component: it is a function of
the three singular values of $q$ alone, even under $\det q\mapsto-\det q$.
\end{proposition}
\begin{proof}
The Laplacian on $\mathbb R^9=\mathbb R^{3\times3}$ is invariant under
$q\mapsto RqO$ for orthogonal $R,O$, since this is an orthogonal map of
$\mathbb R^9$. Under $q\mapsto Rq$ the Gram matrix is unchanged; under
$q\mapsto qO$ it becomes $O^{\mathsf T}GO$, so both traces in
\eqref{f16v49:Q-gram} are unchanged. Sign changes and permutations are
right actions by signed permutation matrices. The Gauss condition of V44 is
simultaneous rotation of the three color vectors, the left $\operatorname{SO}(3)$ action.
Functions of $G$ and $\det q$ are invariant because $G$ and $\det q$ are.
The right action commutes with the left one, hence preserves the invariant
space, and $L^2$ decomposes into isotypic components of the compact group
$O(3)$. V44 proves that the invariant restriction of $H_{\rm mat}$ has a
simple ground eigenvalue with a strictly positive eigenfunction; each right
orthogonal map and each sign change sends this eigenfunction to a positive
eigenfunction of the same eigenvalue, hence to itself. A function invariant
under the full two-sided $O(3)\times O(3)$ action is a function of the
singular values, and $\det q\mapsto-\det q$ is realized by a right
reflection.
\end{proof}

The right $O(3)$ action is not a symmetry of the original lattice transfer,
whose direction group is cubic; it is an exact symmetry of the critical
model only. The five-fold multiplets found below therefore split, at higher
order in $h$ or at finite coupling, into cubic representations; nothing
about that splitting is asserted here.

\subsection{A pairwise transverse lower bound for the ground level}

Let $a_1<0$ be the first zero of the Airy function $\operatorname{Ai}$, so
that $|a_1|$ is the ground energy of $-\Delta+|x|$ on $\mathbb R^3$: the
ground function is radial and the reduced radial problem
$-u''+ru=Eu$ on $(0,\infty)$ with $u(0)=0$ has the solutions
$u=\operatorname{Ai}(r-E)$, square integrable exactly when
$\operatorname{Ai}(-E)=0$.

\begin{theorem}[Pairwise transverse bound]
\label{f16v49:pairwise}
As quadratic forms on $L^2(\mathbb R^9)$,
\begin{equation}
 \boxed{\quad
 H_{\rm mat}\succeq\sum_{i<j}h_{ij},\qquad
 h_{ij}=-\tfrac14\Delta_{q_i}-\tfrac14\Delta_{q_j}+2|q_i\times q_j|^2,
 \qquad h_{ij}\succeq2^{-1/3}|a_1|\,I,
 \quad}
 \label{f16v49:pairwise-bound}
\end{equation}
with equality in the first relation. Consequently
\begin{equation}
 e_0\ \ge\ 3\cdot2^{-1/3}|a_1|\ >\ 5.547.
 \label{f16v49:e0-lower}
\end{equation}
This improves the value $3\cdot4^{-1/3}|a_1|\approx4.42$ implied by V44's
oscillator floor $H_{\rm mat}\succeq-\frac14\Delta_q+\sum_\mu|q_\mu|$, and
the value $2|a_1|\approx4.676$ obtained from the best weighted average of
V44's two inequalities $H_{\rm mat}\succeq-\frac12\Delta_q$ and
$H_{\rm mat}\succeq2\sum_\mu|q_\mu|$.
\end{theorem}
\begin{proof}
Each Laplacian $\Delta_{q_i}$ occurs in exactly two of the three pairs, and
each pair potential once, so $\sum_{i<j}h_{ij}=H_{\rm mat}$ exactly. Fix
$(i,j)=(1,2)$ and a smooth compactly supported $f$ on $\mathbb R^9$. For
fixed $q_1=x\ne0$ and $q_3$, rotate the variable $q_2$ so that $x$ is along
the third axis; this is a unitary change of variables at fixed $x$ and no
derivative in $x$ acts on the rotated function. In the rotated coordinates
$y$, $2|x\times y|^2=2|x|^2(y_1^2+y_2^2)$ and
\[
 -\tfrac14\Delta_y+2|x|^2(y_1^2+y_2^2)
 =-\tfrac14\partial_{y_3}^2+\sum_{k=1}^2
 \bigl(-\tfrac14\partial_{y_k}^2+2|x|^2y_k^2\bigr)
 \succeq0+2\sqrt{\tfrac14\cdot2|x|^2}=\sqrt2\,|x|,
\]
because $-a\,d^2/dy^2+by^2$ has ground energy $\sqrt{ab}$. Integrating this
fibrewise inequality against $dx\,dq_3$ gives
$h_{12}\succeq-\frac14\Delta_{q_1}+\sqrt2|q_1|$, an operator acting in $q_1$
only. The unitary dilation $q_1=2^{-5/6}x'$ turns it into
$2^{-1/3}(-\Delta_{x'}+|x'|)$, whose ground energy is $2^{-1/3}|a_1|$. The
same bound holds for the other two pairs by relabeling. Summing the three
operator inequalities proves \eqref{f16v49:pairwise-bound}, and restricting
to the invariant subspace, which contains the ground function, proves
\eqref{f16v49:e0-lower} once $|a_1|\ge2.33$ is available; that certificate
is Proposition~\ref{f16v49:airy} below, and $3\cdot2^{-1/3}\cdot2.33\ge
3\cdot0.7936\cdot2.33=5.547264$ since $0.7936^3<\frac12$.
For the comparison with V44, $-\frac14\Delta+\sum_\mu|q_\mu|$ is a sum of
three commuting copies of $-\frac14\Delta_{x}+|x|$, each with ground energy
$4^{-1/3}|a_1|$, so its ground energy is $3\cdot4^{-1/3}|a_1|$; the weighted
average $(1-t)(-\frac12\Delta_q)+2t\sum_\mu|q_\mu|$ has ground energy
$3\cdot2^{1/3}(1-t)^{1/3}t^{2/3}|a_1|$, maximal at $t=2/3$ with value
$2|a_1|$. The pairwise bound keeps for every pair the full transverse
oscillator frequency of both partners instead of splitting the kinetic
energy evenly, which is where the factor $3\cdot2^{-1/3}=2.381$ against
$2$ and $3\cdot4^{-1/3}=1.890$ comes from.
\end{proof}

\begin{proposition}[Exact shooting certificate for the first Airy zero]
\label{f16v49:airy}
The reduced radial operator $-u''+ru$ on $(0,\infty)$ with Dirichlet
condition at $0$ has no eigenvalue below $E=233/100$. Hence
$|a_1|\ge2.33$.
\end{proposition}
\begin{proof}
By the Sturm oscillation theorem for this regular-at-zero, limit-point,
discrete-spectrum problem, the number of eigenvalues strictly below $E$
equals the number of zeros in $(0,\infty)$ of the solution $u_E$ of
$u''=(r-E)u$ with $u(0)=0$, $u'(0)=1$. That solution is the entire power
series $u_E(r)=\sum_nc_nr^n$ with $c_0=0$, $c_1=1$, $c_2=0$ and
$(n+3)(n+2)c_{n+3}=c_n-Ec_{n+1}$, all coefficients rational for rational $E$.
Let $R=6$ and $M_n=|c_n|R^n$. The recursion gives
$M_{n+3}\le T_n(R^3+ER^2)/((n+2)(n+3))$ with $T_n=\max(M_n,M_{n+1},M_{n+2})$,
hence $T_{n+3}\le T_n/2$ once $(n+2)(n+3)\ge2(R^3+ER^2)$, so
$\sum_{n\ge N}M_n\le6T_N$ and $\sum_{n\ge N}n|c_n|R^{n-1}\le(3T_N/R)(2N+10)$
for such $N$; with $N=150$ both tails are below $10^{-54}$. Three exact
rational checks follow. On $(0,1]$, $u_E(r)=r[1-\frac E6r^2+\sum_{n\ge4}
c_nr^{n-1}]$ and the bracket is at least $1-\frac E6-\sum_{n\ge4}|c_n|\ge
0.46>0$. On $[1,R]$, with the grid step $\delta=1/100$, the deviation of $u_E$
from its linear interpolation on a cell is at most $(\delta^2/8)\sup|u_E''|$,
and $|u_E''|=|r-E||u_E|\le R\sup|u_E|$, so the cell minimum is at least
$\min(u_E(a),u_E(b))-(\delta^2/8)R\max(|u_E(a)|,|u_E(b)|)/(1-R\delta^2/8)$;
this is positive on every cell, with margin at least $0.108$. Finally
$u_E(R)>0.80$, $u_E'(R)>1.47$ and $R>E$, so $u_E''=(r-E)u_E>0$ as long as
$u_E>0$ beyond $R$, which keeps $u_E'$ and then $u_E$ increasing forever.
Thus $u_E$ has no zero in $(0,\infty)$ and there is no eigenvalue below $E$.
All evaluations are exact rational computations of the truncated series with
the stated tail bounds, recorded by the V49 verifier.
\end{proof}

\subsection{Certified Rayleigh--Ritz upper bounds in an invariant basis}

For a polynomial $p$ in the six Gram entries $G_{ij}$, $i\le j$, of total
degree at most $d$, put
\begin{equation}
 f_p(q)=p(G)e^{-|q|^2},\qquad f_p^{\det}(q)=\det(q)\,p(G)e^{-|q|^2}.
 \label{f16v49:basis}
\end{equation}
These are left-invariant by Proposition~\ref{f16v49:symmetry}, lie in the
operator domain, and for $d=2$ span a $56$-dimensional invariant subspace
$\mathcal R_2$. All Gram, kinetic and potential matrix elements are
Gaussian moments of polynomials against $e^{-2|q|^2}$, hence rational
multiples of $(\pi/2)^{9/2}$; the common factor cancels in the generalized
eigenvalue problem. Write $S$ and $H$ for the exact rational Gram and
Hamiltonian matrices on $\mathcal R_2$.

\begin{theorem}[Certified upper bounds]
\label{f16v49:rr}
On the degree-two space $\mathcal R_2$,
\begin{equation}
 \boxed{\quad
 e_0\le\frac{65571}{10000},\qquad
 e_1\le e_2\le e_3\le e_4\le e_5\le\frac{100232}{10000},\qquad
 e_6\le\frac{10688}{1000},
 \quad}
 \label{f16v49:rr-bounds}
\end{equation}
and on the $168$-dimensional degree-three space $\mathcal R_3$,
\begin{equation}
 \boxed{\quad
 e_0\le 6.5489,\qquad
 e_5\le 9.6678,\qquad
 e_6\le 10.5318.
 \quad}
 \label{f16v49:rr-bounds-3}
\end{equation}
The Rayleigh quotient of the pure Gaussian is exactly $27/4$, V45's
constant, and the degree-one subspace reproduces V45's $47/4$ for the second
radial level; the degree-two space lowers the ground bound from $6.75$ to
$6.5571$ and the first-excitation bound from $10.25$ (degree one) to
$10.0232$, and the degree-three space to 6.5489 and 9.6678.
\end{theorem}
\begin{proof}
For a rational $\lambda$ the number of negative eigenvalues of the exact
symmetric matrix $H-\lambda S$ equals, by Sylvester's law of inertia, the
number of negative pivots of its exact $LDL^{\mathsf T}$ factorization with
diagonal pivoting, provided no zero pivot occurs; the verifier records that
none does. If that number is $k$, the negative subspace $N$ of $H-\lambda S$
has dimension $k$, and for $c\in N\setminus\{0\}$ the function
$f_c=\sum_ic_if_i$ satisfies $\langle f_c,H_{\rm mat}f_c\rangle<\lambda\|f_c\|^2$;
in particular $f_c\ne0$, so $c\mapsto f_c$ is injective on $N$ and its image
is a $k$-dimensional subspace of the invariant form domain with Rayleigh
quotients below $\lambda$. Min--max gives $e_{k-1}<\lambda$. The verifier
finds $k=1$ at $\lambda=65571/10000$, $k=6$ at $\lambda=100232/10000$ and
$k=7$ at $\lambda=10688/1000$ on $\mathcal R_2$, proving
\eqref{f16v49:rr-bounds}, and the same three counts on $\mathcal R_3$ at the
three shifts of \eqref{f16v49:rr-bounds-3}, which are the floating
eigenvalues of the exact degree-three matrices rounded up to four decimals
plus $2\cdot10^{-4}$. The Gaussian quotient $27/4$ is the exact entry
$H_{00}/S_{00}$; the degree-one statements are V45's calculation, contained
in the same matrices.
\end{proof}

In floating evaluation the exact degree-two matrices have lowest
eigenvalues
\[
 6.55701,\quad 10.02312\ (\times5),\quad 10.68789,\quad 13.75,
\]
with vanishing $\det$-odd content in all of them. The five-fold level is
the $\ell=2$ isotypic component of the right $O(3)$ action of
Proposition~\ref{f16v49:symmetry}, an exact degeneracy of the model, and the
level $10.688$ lies in the trivial component (a radial excitation of the
ground); the $\det$-odd sector starts at $13.75$ in this basis. Optimizing
the Gaussian width changes the degree-two values only in the
third digit. The exact degree-three matrices have lowest floating
eigenvalues $6.5487$, $9.6675$ ($\times5$), $10.5315$, $13.189$, with the same sector content: the ground moves by
$0.008$, the five-fold level by $0.36$ and the radial excitation by $0.16$,
so the first excitation is not yet converged at degree three and the
certified upper bound on $g_*$ from \eqref{f16v49:rr-bounds-3} is
$4.1206$. These floating values are diagnostics; only the rational
certificates are hashed.

\subsection{Consequences for the fixed-regulator spectrum of V48}

\begin{proposition}[Certified brackets for V48's constants]
\label{f16v49:consequences}
In V48's ordered spectral limits, at every fixed $L=2m$, $m\ge4$,
\begin{equation}
 \boxed{\quad
 5.547\le\lim_{\gamma\to\infty}\frac{1-\Lambda_\gamma/\lambda_*}{h}
                       =e_0\le 6.5489,\qquad
 0<g_*=\lim_{\gamma\to\infty}\frac{g_\gamma}{\lambda_*h}\le 4.1206,
 \quad}
 \label{f16v49:brackets}
\end{equation}
and the seventh ordered level of every center sector has limiting scaled
deficit at most $10.5318$. The neutral deficit
$-\log(\lambda_{0,2}/\Lambda_\gamma)=h(g_*+o_L(1))$ therefore satisfies
$\limsup h^{-1}(-\log(\lambda_{0,2}/\Lambda_\gamma))\le 4.1206$.
\end{proposition}
\begin{proof}
Insert \eqref{f16v49:e0-lower} and \eqref{f16v49:rr-bounds-3} into V48's
limits; $g_*=e_1-e_0\le 9.6678-5.5472$, and $g_*>0$ is V44. The last
statement is the third line of V48's gap display.
\end{proof}

Two remarks keep the scope honest. First, $h=(\gamma V)^{-1/3}$ is the
asymptotic parameter of V45; with the magnetic and kinetic couplings of the
transfer as normalized in V44 ($b_\gamma=e^{-\gamma S/2}$ with
$S=2\sum_p(1-\frac12\Tr U_p)$ over spatial plaquettes, applied twice per
step, and $e^{-\gamma s}$ for the kinetic kernel), the transfer is an
anisotropic Wilson weight, and no physical time step is fixed by these
numbers; the deficits above are per transfer step. Second, in the
Rayleigh--Ritz spectra the lowest excitation is the five-fold multiplet,
below the first radial excitation of the ground; in the small-volume
literature the corresponding statement is that the tensor glueball channel
lies below the excited scalar channel in sufficiently small volumes. Here it
is an observation about certified upper bounds and floating values, not a
certified ordering of the true levels $e_1$ and $e_6$.

\subsection{The next model estimate}

The missing side is a certified lower bound on $e_1$ above the ground
Rayleigh quotient, or directly on $g_*$. Two routes are open. A Temple or
Lehmann--Goerisch bound needs an a priori lower bound on $e_1$ larger than the
ground Rayleigh quotient, about $6.55$; the separable comparison $H_{\rm mat}\succeq\sum_\mu
(-\frac16\Delta_{q_\mu}+\frac43|q_\mu|)$ gives only $\frac23(|a_2|+2|a_1|)
\approx5.84$ for the second invariant level, and the pairwise bound of
Theorem~\ref{f16v49:pairwise} does not separate a second level because each
$h_{ij}$ has continuous spectrum in the free longitudinal directions. A
Barta-type bound on the bi-invariant sector, $e_0\ge\inf(H_{\rm mat}\phi)/\phi$
for a positive bi-invariant $\phi$, reduces to three singular-value
variables and could sharpen the ground enclosure, but it does not bound
$e_1$. A genuinely two-sided enclosure of $g_*$ therefore needs a sector-wise
lower bound for the $\ell=2$ component, for instance through a covariant
comparison operator; this is recorded as the next model task. It is
independent of, and does not replace, the native problems named in V48:
volume-uniform control and physical-time calibration.

\subsection{Verification and preservation}

The new exact diagnostic is
\begin{center}\small\path{scripts/verify_ym_f16_v49_gap_coefficient.py}.\end{center}
It performs the exact rational shooting certificate of
Proposition~\ref{f16v49:airy} ($E=233/100$, $R=6$, $150$ series terms, grid
step $1/100$), the cube-root check $0.7936^3<\frac12$ and the product
$3\cdot\frac{496}{625}\cdot\frac{233}{100}=\frac{86676}{15625}$, the exact
degree-two Gram, kinetic and quartic matrices on the $56$ invariant functions
\eqref{f16v49:basis}, the exact identity $H_{00}/S_{00}=27/4$, and the three
$LDL^{\mathsf T}$ inertia certificates of Theorem~\ref{f16v49:rr}; invoked
with the argument $3$ it builds the exact $168$-dimensional degree-three
matrices and certifies the three shifts of \eqref{f16v49:rr-bounds-3} in a
second payload. Floating spectra are recorded separately and are not
hashed. It has 15096 bytes and SHA--256
\begin{center}\scriptsize\ttfamily E55A8BAA7B85CC8B9D1FE1CA3BADA839A69D8A46A20DD0C37F0F27797FF1F624.\end{center}
Its canonical hashed payloads have SHA--256
\begin{center}\scriptsize\ttfamily 09C308CEF5C4B72C72770F7D4C1F3B63235320BFE458F54EC7981A470994215E\ (degree two),\\ 635822A2A99B752B16B7ABBC94D03BBE549393742A4A08BA5B998EE497877338\ (degree three).\end{center}
Only this terminal insertion and the literal title version differ from V48;
removing them recovers its bytes exactly. The predecessor checkers were not
rerun; the model inputs of V44, V45 and V48 are used as stated. The next
scheduled checkpoint remains V50.

\begin{firewall}
V49 proves an exact symmetry decomposition of the quartic model, a pairwise
transverse lower bound $e_0\ge3\cdot2^{-1/3}|a_1|>5.547$ with an exact
certificate of the Airy zero, and certified Rayleigh--Ritz upper bounds
$e_0\le 6.5489$, $e_5\le 9.6678$, $e_6\le 10.5318$, hence
$0<g_*\le 4.1206$ for V48's gap coefficient. It does not certify a lower bound
on $g_*$, the sector of the true first excitation, a native estimate uniform
in the spatial regulator, a physical time step, or an interacting continuum
mass gap. The original transfer, its constraints, and V48's spectral
identification are unchanged.
\end{firewall}
% END F16 V49 APPEND

% BEGIN F16 V50 APPEND -- 2026-09-18
% Insert before the final document-ending command of cumulative V49.
% This insertion leaves every predecessor body and the native operator unchanged.

\clearpage
\section{V50: companion checkpoint, ten-version review, and the explicit hard spectrum}
\label{f16v50:context}

V50 is the scheduled companion checkpoint. It records the rereading of the
companion programmes, reviews V40--V49, corrects one normalization remark of
V49, and adds one structural result toward the volume-uniform problem named
in V48: the hard Gaussian transfer of V44 has an explicit spectrum, its
smallest Mehler frequency is $\operatorname{arcosh}(1+2\sin^2(\pi/L))$, and
consequently the ratio between the critical soft scale $h$ and the hard-band
gap is bounded by a fixed multiple of $\gamma^{-1/3}$ uniformly in the
spatial regulator. The hard/soft separation that underlies V45--V48 is
therefore not one of the places where uniformity in $L$ is lost. Which
constants do lose it is listed at the end.

\subsection{V50 companion checkpoint}

Fresh inventories of \path{latex_notes} found the same latest companion
versions as at the f15 V100 checkpoint; their sizes and hashes are recorded
here and were unchanged since that checkpoint. The latest conclusions of each
were reread.

\begin{center}\small
\begin{tabular}{@{}p{0.27\textwidth}p{0.68\textwidth}@{}}
\toprule
Input (bytes; SHA--256 prefix) & Relevant scope retained at this checkpoint\\
\midrule
SM continuum V72 (607372; f44f9745) &
Fixed-mode, fixed-time many-copy limit; neither the physical finite-species
continuum nor a YM source-law estimate.\\[1mm]
NS stability V26 (278283; 82c887f8) &
Pointwise rank-one kernel norms; no global-regularity theorem usable here.\\[1mm]
Shell mixing V16 (623261; bf138e63) &
Measured flat-branch one-loop assembly with a local stationary coefficient and
its refinement limit; no noncommuting integral remainder, compact entry
probability, or physical-time contraction. Its measured convention is kept
distinct from the present transfer normalization.\\[1mm]
Parallel YM V5 (54174; 306f1bb8) &
Extensive total-charge tightness fails on its tensorized target; its terminal
decision does not stop this programme.\\[1mm]
GRH cofinal visibility V3 (23851; 4dd31580) &
A separate arithmetic programme; no input to the transfer analysis.\\
\bottomrule
\end{tabular}
\end{center}

No companion supplies a volume-uniform native spectral comparison or a
physical-time calibration. The cross-programme lesson retained from f15 V100,
to control densities and local quantities rather than extensive counts, is
applied below: the new statements concern the hard band gap and a vacuum
energy density, both local in the volume. The next companion checkpoint is V55.

\subsection{Ten-version review, V40--V49}

\begin{center}\small
\begin{tabular}{@{}p{0.07\textwidth}p{0.88\textwidth}@{}}
\toprule
Version & Established at fixed spatial regulator (scope limits in each firewall)\\
\midrule
V40 & Positive root-averaged endpoint representation; explicit Gaussian root
determinant; projected endpoint Gamma edge of shape $9/2$.\\
V41 & The native Perron value converges to V38's Gaussian threshold; the
stationary probability concentrates on the eight central-flat orbits.\\
V42 & Exponential localization outside central-flat neighborhoods; same-top
Schur reduction; one exponentially close top in each center sector.\\
V43 & Exact character pencils and inertia counting; every regional root
matches a level in each sector; the neutral gap in the returning-state metric.\\
V44 & Invariant soft trial spaces: every fixed level tends to the top, and
$g_\gamma\to0$; the quartic model $H_{\rm mat}$ has compact resolvent and
$e_1>e_0$.\\
V45 & Critical-scale recovery with the hard-band corrector; spectral inclusion
with multiplicities; one-sided ordered bounds; the constants $27/4$, $47/4$.\\
V46 & All-vector confinement and spectral completeness of the finite-step
comparison $S_h$; every nonzero native critical profile is a model
eigenfunction.\\
V47 & Strip operator-norm comparison; whole-band hard-ground capture; local
soft compactness; the completeness problem reduced to soft-position escape.\\
V48 & Quadratic hard-return cost, annular coercivity, dyadic localization,
radial tightness; ordered spectrum $(1-\lambda_{a,j}/\lambda_*)/h\to e_{j-1}$
and gap coefficient $g_*=e_1-e_0$.\\
V49 & Symmetry sectors of $H_{\rm mat}$; $e_0\ge3\cdot2^{-1/3}|a_1|>5.547$
with an exact Airy certificate; certified $e_0\le6.5489$, $e_5\le9.6678$,
$e_6\le10.5318$; $0<g_*\le4.1206$.\\
\bottomrule
\end{tabular}
\end{center}

The ten versions complete the fixed-regulator critical branch: the native
low spectrum at fixed $L$ and $\gamma\to\infty$ is the invariant spectrum of
an explicit nine-dimensional quartic model, with certified one-sided
constants. Nothing in them is uniform in $L$, and the model's first excitation
is not yet converged or bounded below. The direction for V51--V59 is fixed
accordingly: (i) a certified lower bound on $e_1$, sector-wise, and a
symmetry-reduced basis to converge the $\ell=2$ level; (ii) the tunneling
exponent between center sectors along the commuting valleys, where the
transverse zero-point energy supplies an effective linear potential and the
exponent scales as $h^{-3/2}=\sqrt{\gamma V}$, matching the form of V42's
upper bound; (iii) explicit tracking of the $L$-dependence of the constants of
V44--V48, beginning with the hard band below.

\subsection{Correction of V49's normalization remark}

V49's remark after Proposition~\ref{f16v49:consequences} described the
transfer as an anisotropic Wilson weight, with the parenthetical
normalization $S=2\sum_p(1-\frac12\Tr U_p)$. Both statements are withdrawn.
In quaternion norm, $|U_p-I|^2=2(1-\Re U_p)$ for a unit quaternion, so V44's
$S=\frac12|\mathcal W|^2$ equals V41's $S=\sum_p(1-\Re U_p)
=\sum_p(1-\frac12\Tr U_p)$. Two magnetic factors $b_\gamma=e^{-\gamma S/2}$
per step give each spatial plaquette the weight $e^{-\gamma(1-\Re U_p)}$, and
the kinetic kernel gives each temporal plaquette, in temporal gauge, the same
weight $e^{-\gamma(1-\Re UU'^{-1})}$. The transfer is therefore the isotropic
Wilson transfer at coupling $\gamma$ per plaquette, with one transfer step
equal to one lattice spacing. The remark's conclusion that no physical scale
is fixed by V48's formulas stands in the following precise form: the lattice
spacing in physical units, hence the physical value of one transfer step, is a
separate calibration, and V48's deficits are per lattice spacing. The
remaining content of V49 is unaffected.

\subsection{The explicit hard spectrum}

Let $\mathsf d_0$ and $\mathsf d_1$ be V44's vertex-to-link gradient and
link-to-plaquette curl on real color-valued cochains of the torus
$\mathbb Z_L^3$, in the orthonormal fine-link metric, and let
$H_+=\mathsf d_1^*\mathsf d_1$ on $\mathcal Y=\ker\mathsf d_0^*\ominus\ker
\mathsf d_1$. For $k\in\frac{2\pi}L\mathbb Z_L^3$ write
$\hat k^2=\sum_{\mu=1}^34\sin^2(k_\mu/2)$.

\begin{proposition}[Hard spectrum and Mehler frequencies]
\label{f16v50:spectrum}
The spectrum of $H_+$ is $\{\hat k^2:k\ne0\}$, each nonzero momentum carrying
multiplicity six (two transverse polarizations and three colors), for a total
of $6(V-1)$ hard modes. The hard transfer $\mathcal T_+$ of V44 is the tensor
product of the one-mode Mehler kernels $t_{\hat k^2}$ of V38, so its
eigenvalues are $\lambda_*\prod_{k\ne0}e^{-n_k\xi_k}$ over occupation
numbers $n_k\in\mathbb Z_{\ge0}^6$, with
\begin{equation}
 \boxed{\quad
 \xi_k=\operatorname{arcosh}\bigl(1+\tfrac12\hat k^2\bigr),\qquad
 \lambda_*=\Lambda_{{\rm G},B}=\exp\Bigl(-3\sum_{k\ne0}\xi_k\Bigr),\qquad
 r_+=e^{-\xi_{\min}},\quad
 \xi_{\min}=\operatorname{arcosh}\bigl(1+2\sin^2(\pi/L)\bigr).
 \quad}
 \label{f16v50:mehler}
\end{equation}
In particular the hard band gap of V46--V48 satisfies, with $x=2\sin^2(\pi/L)$,
\begin{equation}
 1-r_+=\sqrt{x^2+2x}-x,\qquad
 2\sin\tfrac\pi L\Bigl(1-\sin\tfrac\pi L\Bigr)\le1-r_+\le2\sin\tfrac\pi L,
 \label{f16v50:gap-bounds}
\end{equation}
and V45's hard inverse obeys
$\|R_+\|\le[\lambda_*(1-r_+)]^{-1}\le\lambda_*^{-1}L/(4(1-\sin(\pi/L)))$.
\end{proposition}
\begin{proof}
On the flat torus the Hodge Laplacian on one-cochains,
$\mathsf d_1^*\mathsf d_1+\mathsf d_0\mathsf d_0^*$, acts componentwise as the
scalar lattice Laplacian $\Delta$, $(\Delta a)(s,\mu)=6a(s,\mu)-\sum_{\nu,\pm}
a(s\pm\hat\nu,\mu)$; this is the standard identity, checked exactly on every
link basis vector by the V50 verifier at $L=4$, and it holds for every $L$
by the same local computation, since both sides are translation invariant and
involve only nearest neighbors. On $\mathcal N=\ker\mathsf d_0^*$ the second
term vanishes, so $H_+$ is the restriction of $\Delta$ to
$\mathcal N\ominus\ker\mathsf d_1$. The plane waves $a(s,\mu)=e_\mu e^{ik\cdot s}$
diagonalize $\Delta$ with eigenvalue $\hat k^2$; $a$ is divergence free
exactly when $\sum_\mu e_\mu(e^{-ik_\mu}-1)=0$, a single linear condition that
is nontrivial for $k\ne0$, leaving two polarizations per nonzero $k$ and per
color; for $k=0$ every $e_\mu$ is allowed and $\mathsf d_1a=0$, which is
$\ker\mathsf d_1\cap\mathcal N=\mathcal Z$, V44's nine constants. Counting
gives $2(V-1)$ per color and $6(V-1)$ in total, which is V44's $\dim\mathcal Y$,
so the transverse waves exhaust $\mathcal Y$. For $L=4$ all phases are
Gaussian integers and the verifier confirms the eigenvector, divergence and
counting statements exactly, with the multiplicities
$12,30,40,30,12,2$ of the values $2,4,6,8,10,12$ per color.

The hard kernel of V44 is $(2\pi)^{-r/2}\exp[-|u-u'|^2/2-u^{\mathsf T}H_+u/4
-u'^{\mathsf T}H_+u'/4]$; in an orthonormal eigenbasis of $H_+$ it factorizes
into V38's one-mode kernels $t_\nu$ with $\nu=\hat k^2$, whose eigenvalues are
$e^{-(n+1/2)\xi_\nu}$, $\xi_\nu=\operatorname{arcosh}(1+\nu/2)$, by V38's Mehler
calculation; the ground factor is $(1+\nu/2+\sqrt{\nu+\nu^2/4})^{-1/2}
=e^{-\xi_\nu/2}$. The verifier checks the underlying fixed-point identity of
the Gaussian ansatz exactly: with $s=e^{\xi}$ the condition that
$e^{-au^2}$ be an eigenfunction of $t_\nu$ is $s^2-(2+\nu)s+1=0$. Taking
the product over the $6(V-1)$ modes gives $\lambda_*=\|\mathcal T_+\|
=\prod e^{-\xi_k/2}$, which is V41's harmonic threshold at the identity and,
by V44's matching of the fine and rooted kinetic metrics, $\Lambda_{{\rm G},B}$.
The largest eigenvalue on $\varphi^\perp$ is attained by one excitation of
the softest hard mode, giving $r_+=e^{-\xi_{\min}}$ with $\hat k^2$ minimal at
$k=(2\pi/L)(1,0,0)$, i.e.\ $\hat k^2=4\sin^2(\pi/L)$.
Finally $e^{-\xi}=1+x-\sqrt{x^2+2x}$ for $\cosh\xi=1+x$, which gives the
formula for $1-r_+$; the bounds follow from $\sqrt{2x}\le\sqrt{x^2+2x}\le
\sqrt{2x}+x$ and $\sqrt{2x}=2\sin(\pi/L)$, and the inverse bound from V45's
formula and $\sin(\pi/L)\ge2/L$.
\end{proof}

\begin{proposition}[Volume-uniform hard/soft separation]
\label{f16v50:adiabatic}
For every $L=2m\ge8$ and every $\gamma$,
\begin{equation}
 \boxed{\quad
 \frac{h}{1-r_+}\le\frac{\gamma^{-1/3}}{4(1-\sin(\pi/8))}<0.41\,\gamma^{-1/3},
 \qquad h=(\gamma L^3)^{-1/3}.
 \quad}
 \label{f16v50:ratio}
\end{equation}
Consequently, in V46's hard separation
$\langle F,\mathbb D_hF\rangle\ge d_h[f]+h^{-1}(1-r_+)\|F_\perp\|^2$, the
hard complementary coefficient $h^{-1}(1-r_+)$ is at least
$2.4\,\gamma^{1/3}$ for all $L\ge8$, and in V48's annular estimate the
absorption of the mixed first-jet term costs $C_L^2r^2/\kappa_+$ with
$\kappa_+/\lambda_*=1-r_+\ge4(1-\sin(\pi/8))/L$, so that this particular
step loses at most one power of $L$, through $C_L^2L$.
\end{proposition}
\begin{proof}
By \eqref{f16v50:gap-bounds} and Jordan's inequality $\sin(\pi/L)\ge2/L$,
$1-r_+\ge(4/L)(1-\sin(\pi/L))\ge(4/L)(1-\sin(\pi/8))$ for $L\ge8$, while
$h=\gamma^{-1/3}/L$. Divide. The two consequences are V46's inequality and
V48's Young-inequality step with the explicit value of $\kappa_+$.
\end{proof}

Two remarks. First, the vacuum value itself is extensive: $-\log\lambda_*
=3\sum_{k\ne0}\xi_k=V\varepsilon_L$ with $\varepsilon_L\to
3\int_{[-\pi,\pi]^3}\operatorname{arcosh}(1+\tfrac12\hat k^2)\,\frac{d^3k}{(2\pi)^3}$,
numerically $\varepsilon_L=5.9979$ at $L=8$ and $5.9991$ at $L=64$. This is
the Gaussian vacuum energy per site; it confirms that every spectral
statement of the lineage must be a ratio to $\lambda_*$, as it is. Second,
$(1-r_+)L$ increases from $4.21$ at $L=8$ to $5.98$ at $L=64$ toward $2\pi$;
the hard gap in lattice units is the smallest lattice momentum, as expected,
and the softest hard mode is the $k=2\pi/L$ transverse gluon.

\subsection{Where uniformity in \texorpdfstring{$L$}{L} is still lost}

The explicit hard spectrum removes $L$-dependence from three inputs: the hard
band gap $1-r_+$, the hard inverse $R_+$ up to one power of $L$, and the
Gaussian threshold $\lambda_*$ as a normalization. The following constants of
V44--V48 remain $L$-dependent without a tracked rate: the tube radius $r_0$
and the density $J_L$ of V44's based-gauge tube; the constants $C_L,c_L,M_L$
of the kernel comparisons in V47--V48, which arise from third derivatives of
the action on the tube and from the based-group volume; V42's forbidden-region
constant $\delta_L$ and the exponential separation rates $c_L,\kappa_L$; and
the hard Gaussian moments, which are products over $6(V-1)$ modes and are
therefore extensive unless taken relative to $\varphi$. A volume-uniform
version of the critical-branch theorem would need each of these tracked as a
power of $L$ at fixed $\gamma$, together with a statement about the regime
$\gamma\sim\log L$, where the small-volume expansion parameter
$\gamma^{-1/3}$ no longer tends to zero at a rate. Neither is attempted here.

\subsection{Verification and preservation}

The new exact diagnostic is
\begin{center}\small\path{scripts/verify_ym_f16_v50_hard_spectrum.py}.\end{center}
On the $L=4$ torus it checks, in Gaussian-integer arithmetic, $\mathsf d_1
\mathsf d_0=0$, the Hodge identity on every link basis vector, the
divergence-free transverse plane waves as exact eigenvectors of
$\mathsf d_1^{\mathsf T}\mathsf d_1$ with eigenvalue $\hat k^2$, the harmonic
constants, and the count $2(V-1)=126$ per color with its multiplicities; and it
checks the Mehler fixed-point identity as a rational-function identity at
three rational $\nu$. Its floating diagnostics tabulate $\xi_{\min}$,
$(1-r_+)L$ and $\varepsilon_L$ for $L=8,16,32,64$ and are not hashed. It has
12141 bytes and SHA--256
\begin{center}\scriptsize\ttfamily 51FCE6358A9E02EC3E41692870CD9F86A88BD5D19F674FD62DFB1A3CB33755D7.\end{center}
Its canonical hashed payload has SHA--256
\begin{center}\scriptsize\ttfamily A30314BB090A954C13D7BD75168A7939F4E887C80F0EEA1D8F0481FB89FF93B6.\end{center}
Only this terminal insertion and the literal title version differ from V49;
removing them recovers its bytes exactly. The predecessor checkers were not
rerun. The next companion checkpoint is V55; the next ten-version review is
V60.

\begin{firewall}
V50 records an unchanged companion checkpoint and a ten-version review,
corrects V49's normalization remark (the transfer is the isotropic Wilson
transfer at coupling $\gamma$ with one step per lattice spacing), and proves
the explicit hard spectrum, Mehler frequencies, Gaussian threshold and hard
band gap, with the volume-uniform bound $h/(1-r_+)<0.41\gamma^{-1/3}$. It does
not make any estimate of V44--V48 uniform in $L$, does not calibrate the
lattice spacing in physical units, does not bound the model gap coefficient
from below, and does not establish an interacting continuum mass gap.
\end{firewall}
% END F16 V50 APPEND

% BEGIN F16 V51 APPEND -- 2026-09-18
% Insert before the final document-ending command of cumulative V50.
% This insertion leaves every predecessor body and the native operator unchanged.

\clearpage
\section{V51: polar coordinates, centrifugal sector gaps, and sector-resolved certificates}
\label{f16v51:context}

V49 certified one-sided constants for the quartic model and left the first
excitation without a lower bound. Here the singular-value decomposition of the
$3\times3$ matrix $q$ turns the model into a radial operator on three singular
values plus an angular part carried by the right $O(3)$ action. For a
left-invariant function the angular kinetic energy is a rigid-rotor form with
singular-value-dependent inertia, bounded below by the Casimir over $|q|^2$.
Two consequences follow. Every non-scalar sector of the right action lies
above the ground level by at least $2\ell(\ell+1)/E_\ell^2$, where $E_\ell$ is
any upper bound on that sector's ground; for the five-fold tensor multiplet
this gives a certified positive gap. And the exact Rayleigh--Ritz computation
can be organized sector by sector in the six Gram variables, with Gaussian
moments obtained by an exact Stein recursion instead of monomial expansion;
this is faster by orders of magnitude and converges the sector spectra. One
floating remark of V49 is corrected on the way.

\subsection{Dependencies and scope}

The inputs are V44's confinement theorem (compact resolvent, simple positive
ground function), V49's symmetry proposition and Airy certificate, and the
elementary lower bound $H_{\rm mat}\succeq2\sum_\mu|q_\mu|$ of V44. No native
transfer statement is used. The certified sector gap is far from the floating
value of the gap and is not an enclosure of $g_*$; the scalar sectors
$(0,\pm)$ receive no lower bound here. The next companion checkpoint is V55.

\subsection{Singular-value coordinates and the angular kinetic form}

Write $q=U\Sigma O^{\mathsf T}$ with $U,O\in O(3)$ and
$\Sigma=\operatorname{diag}(\sigma_1,\sigma_2,\sigma_3)$, $\sigma_i\ge0$.
The left action is $U\mapsto RU$, the right action $O\mapsto O R'$, and
$G=q^{\mathsf T}q=O\Sigma^2O^{\mathsf T}$. Let $E_{ij}$, $i<j$, be the
antisymmetric matrix with entries $\pm1$ at $(i,j),(j,i)$, and for a function
$F$ on $\mathbb R^9$ let $X_{ij}F(q)=\frac d{dt}F(qe^{tE_{ij}})|_{t=0}$ be the
generators of the right action.

\begin{proposition}[Polar metric and the left-invariant Dirichlet form]
\label{f16v51:polar}
Off the null set where two singular values coincide or one vanishes, the
Euclidean metric of $\mathbb R^9$ in the coordinates $(\sigma,A,B)$,
$A=U^{\mathsf T}dU=\sum_{i<j}a_{ij}E_{ij}$, $B=O^{\mathsf T}dO=\sum_{i<j}b_{ij}E_{ij}$,
is
\begin{equation}
 |dq|^2=\sum_i d\sigma_i^2+\sum_{i<j}\Bigl[(\sigma_i^2+\sigma_j^2)(a_{ij}^2+b_{ij}^2)
 -4\sigma_i\sigma_j\,a_{ij}b_{ij}\Bigr],
 \label{f16v51:metric}
\end{equation}
and $dq=c\prod_{i<j}|\sigma_i^2-\sigma_j^2|\,d\sigma\,dU\,dO$ with Haar
measures on $O(3)$. For a smooth left-invariant $F$,
\begin{equation}
 \boxed{\quad
 |\nabla F|^2=\sum_i|\partial_{\sigma_i}F|^2
 +\sum_{i<j}\frac{\sigma_i^2+\sigma_j^2}{(\sigma_i^2-\sigma_j^2)^2}\,|X_{ij}F|^2,
 \qquad
 \frac{\sigma_i^2+\sigma_j^2}{(\sigma_i^2-\sigma_j^2)^2}\ge\frac1{|q|^2}.
 \quad}
 \label{f16v51:dirichlet}
\end{equation}
\end{proposition}
\begin{proof}
From $dq=dU\,\Sigma O^{\mathsf T}+U\,d\Sigma\,O^{\mathsf T}+U\Sigma\,dO^{\mathsf T}$
and orthogonal invariance of the Frobenius norm,
$|dq|^2=|A\Sigma+d\Sigma-\Sigma B|^2$, using $dO^{\mathsf T}O=-B$. The
diagonal of $A\Sigma+d\Sigma-\Sigma B$ is $d\Sigma$, since $A,B$ have zero
diagonal. Its $(i,j)$ and $(j,i)$ entries are $a_{ij}\sigma_j-\sigma_ib_{ij}$
and $-a_{ij}\sigma_i+\sigma_jb_{ij}$; the sum of their squares is the bracket
in \eqref{f16v51:metric}. The volume form is the square root of the
determinant of this block-diagonal metric: each pair block
$\begin{psmallmatrix}s&-c\\-c&s\end{psmallmatrix}$, $s=\sigma_i^2+\sigma_j^2$,
$c=2\sigma_i\sigma_j$, has determinant $(\sigma_i^2-\sigma_j^2)^2$. This is
the classical singular-value Jacobian for square real matrices. For a
left-invariant $F$ the derivatives along $a_{ij}$ vanish, so $|\nabla F|^2$
retains the $\sigma$ part and the $b_{ij}b_{ij}$ entries of the inverse
metric, which are $s/(s^2-c^2)$ by inversion of the pair block; the
$\sigma$--angle cross terms of the inverse vanish because the metric is block
diagonal. The derivative along $b_{ij}$ at $O$ is $\frac d{dt}F(U\Sigma
(Oe^{tE_{ij}})^{\mathsf T})$, which is $X_{ij}$ applied to $F$. Finally
$(\sigma_i^2-\sigma_j^2)^2\le(\sigma_i^2+\sigma_j^2)^2$ and
$\sigma_i^2+\sigma_j^2\le\Tr G=|q|^2$.
\end{proof}

\subsection{Centrifugal lower bounds for every non-scalar sector}

Recall from V49 that the invariant space decomposes into the isotypic
components of the right $O(3)$ action, labeled $(\ell,\pm)$. For a function of
$G$ the element $-I$ acts trivially, and polynomial covariants of a symmetric
tensor are true tensors, so functions of $G$ carry only $(\ell\ {\rm even},+)$
and $\det(q)$ times functions of $G$ only $(\ell\ {\rm even},-)$. The two scalar
sectors are $(0,+)$, the bi-invariant functions, and $(0,-)$, the
$\det$-odd functions.

\begin{theorem}[Centrifugal separation of the non-scalar sectors]
\label{f16v51:centrifugal}
For every $F$ in the form domain of $H_{\rm mat}$ lying in the sector
$(\ell,\pm)$,
\begin{equation}
 \boxed{\quad
 \langle F,H_{\rm mat}F\rangle\ \ge\ e_0\|F\|^2+\frac{\ell(\ell+1)}2
 \Bigl\langle F,\frac1{|q|^2}F\Bigr\rangle.
 \quad}
 \label{f16v51:sector-form}
\end{equation}
Consequently, if $E_\ell$ is any upper bound on the ground level
$e_0^{(\ell,\pm)}$ of that sector,
\begin{equation}
 e_0^{(\ell,\pm)}-e_0\ \ge\ \frac{2\ell(\ell+1)}{E_\ell^2}.
 \label{f16v51:sector-gap}
\end{equation}
\end{theorem}
\begin{proof}
It suffices to prove \eqref{f16v51:sector-form} on the form core of smooth
compactly supported functions in the sector, obtained by averaging smooth
compactly supported functions with the character of $(\ell,\pm)$ over the
compact group; the averaging is continuous in form norm. By
\eqref{f16v51:dirichlet} and the polar measure,
\[
 \langle F,H_{\rm mat}F\rangle
 =\int\Bigl[\tfrac12|\partial_\sigma F|^2+2\mathcal Q\,|F|^2\Bigr]dq
 +\tfrac12\int\sum_{i<j}w_{ij}(\sigma)\,|X_{ij}F|^2\,dq,
 \qquad w_{ij}\ge|q|^{-2}.
\]
For fixed $\sigma$, the function $O\mapsto F(U\Sigma O^{\mathsf T})$ lies in
the $\ell$-isotypic component of $L^2(O(3))$ for the right regular action, on
which the Casimir $-\sum_{i<j}X_{ij}^2$ acts as $\ell(\ell+1)$; integrating
by parts on the compact group at fixed $(U,\sigma)$,
$\int\sum_{i<j}|X_{ij}F|^2dO=\ell(\ell+1)\int|F|^2dO$. Since $w_{ij}$ depends
only on $\sigma$ and $\min_{ij}w_{ij}\ge|q|^{-2}$, the angular integral is at
least $\frac{\ell(\ell+1)}2\langle F,|q|^{-2}F\rangle$. For the radial part
fix $O$ and define the bi-invariant function
$G_O(q)=F(\Sigma(q),O)$, where $\Sigma(q)$ is the ordered singular-value
vector; it is Lipschitz, its almost-everywhere gradient in $\mathbb R^9$ is the
$\sigma$-gradient of $F$ transported by the isometric differential of the
singular values, and $\mathcal Q$ depends on $\sigma$ only. Hence
$\int[\frac12|\nabla G_O|^2+2\mathcal Q|G_O|^2]dq\ge e_0\|G_O\|^2$ by V44,
because the positive ground function of $H_{\rm mat}$ is its global ground on
$L^2(\mathbb R^9)$ and $e_0$ is that ground energy. Integrating over $O$ and
$U$ proves \eqref{f16v51:sector-form}.

For \eqref{f16v51:sector-gap}, let $F$ be a normalized ground function of the
sector, which exists by compact resolvent. V44's inequality
$H_{\rm mat}\succeq2\sum_\mu|q_\mu|$ gives
$2\langle|q|\rangle_F\le2\langle\sum_\mu|q_\mu|\rangle_F\le e_0^{(\ell,\pm)}\le E_\ell$,
so $\langle|q|\rangle_F\le E_\ell/2$. Jensen's inequality twice,
$\langle|q|^{-2}\rangle\ge\langle|q|^{-1}\rangle^2\ge\langle|q|\rangle^{-2}$,
gives $\langle|q|^{-2}\rangle_F\ge4/E_\ell^2$. Insert in
\eqref{f16v51:sector-form}.
\end{proof}

The bound is coarse: in the floating tensor ground of the sector computation
below, $3\langle|q|^{-2}\rangle\approx0.95$, while Jensen's step replaces
$\langle|q|^{-2}\rangle\approx0.32$ by $4/E_2^2\approx0.044$. The
$\det$-odd sector has no angular term, since $\det O$ is locally constant on
$O(3)$; its profile vanishes where a singular value vanishes, and a lower bound
for it, as for the scalar excitation, remains open.

\subsection{An exact Gram-polynomial calculus}

Write $g=(g_{11},g_{22},g_{33},g_{12},g_{13},g_{23})$ for the entries of $G$
and $\sigma_1=\Tr G$, $\sigma_2=\mathcal Q$, $\sigma_3=\det G$ for its
invariants. Every sector function used below is $P(g)e^{-\sigma_1}$ or
$\det(q)P(g)e^{-\sigma_1}$ with $P$ a polynomial in $g$; the Gaussian weight
$e^{-2|q|^2}$ makes the entries $q_{ia}$ independent centered normals of
variance $1/4$.

\begin{proposition}[Stein recursion and gradient identities]
\label{f16v51:stein}
With $E$ the expectation against $e^{-2|q|^2}$ normalized to $E[1]=1$, for
every polynomial $M(g)$,
\begin{equation}
 E[g_{ij}M]=\tfrac14\Bigl[3\delta_{ij}E[M]+\sum_{k\le l}E\Bigl[\partial_{g_{kl}}M\,
 (\delta_{ki}g_{jl}+\delta_{li}g_{jk})\Bigr]\Bigr],
 \label{f16v51:stein-formula}
\end{equation}
which determines every moment $E[\prod g_{kl}^{n_{kl}}]$ recursively in the
total degree. Moreover, as polynomials in $g$,
\[
 \nabla g_{kl}\cdot\nabla g_{mn}=\delta_{km}g_{ln}+\delta_{kn}g_{lm}+\delta_{lm}g_{kn}+\delta_{ln}g_{km},
 \quad q\cdot\nabla g_{kl}=2g_{kl},\quad
 |\nabla\det q|^2=\sigma_2,\quad
 \nabla\det q\cdot\nabla g_{kl}=2\det(q)\delta_{kl},\quad
 q\cdot\nabla\det q=3\det q,
\]
and for a monomial $M$ of degree $d$ in $q$, $E[|q|^{-2}M]=\frac4{d+7}E[M]$.
Consequently the Gram, kinetic and potential matrix elements of the sector
functions are exact rationals computable without expanding in the nine
coordinates; explicitly, for $f=Pe^{-\sigma_1}$, $f'=P'e^{-\sigma_1}$,
\[
 \nabla f\cdot\nabla f'=e^{-2\sigma_1}\Bigl[\nabla P\cdot\nabla P'
 -4P'\mathcal EP-4P\mathcal EP'+4\sigma_1PP'\Bigr],\qquad
 \mathcal EP=\sum_{k\le l}g_{kl}\partial_{g_{kl}}P,
\]
and for $\det(q)Pe^{-\sigma_1}$ the same bracket is multiplied by $\sigma_3$
and completed by $\sigma_3[2P'\Delta_{\rm d}P+2P\Delta_{\rm d}P'-12PP']+\sigma_2PP'$,
$\Delta_{\rm d}P=\sum_k\partial_{g_{kk}}P$.
\end{proposition}
\begin{proof}
Gaussian integration by parts, $E[q_{ia}\Phi]=\frac14E[\partial_{q_{ia}}\Phi]$,
applied to $\Phi=q_{ja}M$ and summed over $a$, with
$\partial_{q_{ia}}g_{kl}=\delta_{ki}q_{la}+\delta_{li}q_{ka}$ for the
independent variables $k\le l$, gives \eqref{f16v51:stein-formula}: the term
$\delta_{ij}$ comes from differentiating $q_{ja}$, and
$\sum_aq_{ja}(\delta_{ki}q_{la}+\delta_{li}q_{ka})=\delta_{ki}g_{jl}+\delta_{li}g_{jk}$.
The gradient identities are the same contraction; those for $\det q$ use the
cofactor matrix, $q\operatorname{adj}(q)=\det(q)I$ and
$\operatorname{adj}(q)^{\mathsf T}\operatorname{adj}(q)=\operatorname{adj}(qq^{\mathsf T})$,
whose trace is $\sigma_2$. The radial rule is the ratio of the Gaussian radial
integrals $\int r^{d+6}e^{-2r^2}dr/\int r^{d+8}e^{-2r^2}dr=4/(d+7)$ after
separating the angular integral. The expansions of $\nabla f\cdot\nabla f'$
are direct.
\end{proof}

\subsection{Sector-resolved certified bounds}

For Gram degree $D$ let the sector bases be: $(0,+)$, the monomials
$\sigma_1^a\sigma_2^b\sigma_3^c$ with $a+2b+3c\le D$; $(2,+)$, the
$(1,2)$-components $g_{12}\,\sigma^\alpha$ with $|\alpha|\le D-1$ and
$(G^2)_{12}\,\sigma^\alpha$ with $|\alpha|\le D-2$, which span all $(1,2)$
components of $\ell=2$ covariants of $G$ up to that degree by the
Cayley--Hamilton reduction of matrix polynomials; $(0,-)$ and $(2,-)$, the
same multiplied by $\det q$ with the degree lowered by one. All are
multiplied by $e^{-\sigma_1}$.

\begin{theorem}[Certified sector bounds at Gram degree 8]
\label{f16v51:rr}
With exact rational matrices and $LDL^{\mathsf T}$ inertia certificates as in
V49, the ground of $(0,+)$, the second level of $(0,+)$, the ground of
$(2,+)$, and the grounds of $(0,-)$ and $(2,-)$ satisfy
\begin{equation}
 \boxed{\quad
 e_0\le 6.5352,\qquad e_0^{(0,+),2}\le 10.1409,\qquad
 e_0^{(2,+)}\le 9.5487,\qquad e_0^{(0,-)}\le 13.9484,\qquad e_0^{(2,-)}\le 17.9898.
 \quad}
 \label{f16v51:sector-bounds}
\end{equation}
For the ordered invariant levels, $e_0\le 6.5352$, $e_5\le 9.5487$ and
$e_6\le 10.1409$. Combined with Theorem~\ref{f16v51:centrifugal},
\begin{equation}
 \boxed{\quad
 e_0^{(2,+)}-e_0\ \ge\ \frac{12}{(9.5487)^2}\ >\ 0.1316.
 \quad}
 \label{f16v51:tensor-gap}
\end{equation}
The degree-two sector values reproduce V49's exact degree-two numbers
($27/4$ for the Gaussian, $6.55701$ and $10.02312$ for the scalar and tensor
grounds), and the single $\ell=4$ component of the degree-two space has the
exact quotient $55/4=13.75$.
\end{theorem}
\begin{proof}
The sector functions are left-invariant, lie in the operator domain, and the
components of one covariant multiplet are mutually orthogonal with equal
Rayleigh quotients, since $H_{\rm mat}$ commutes with the right action and the
multiplet is irreducible. If $k$ negative pivots occur at shift $\lambda$ in a
sector block, there are $k$ orthogonal functions in that sector with quotients
below $\lambda$, hence $k$ sector levels below $\lambda$ by min--max within the
sector, which is invariant under $H_{\rm mat}$. For the ordered invariant
levels: the five rotated copies of the tensor function and the scalar ground
function give six orthogonal invariant functions with quotients at most
$\max(6.5352,9.5487)=9.5487$, so $e_5\le 9.5487$; adding the second scalar
function gives $e_6\le 10.1409$ since $10.1409\ge 9.5487$. The gap statement is
\eqref{f16v51:sector-gap} with $\ell=2$ and $E_2=9.5487$. The $\ell=4$
identification projects $g_{12}g_{13}$, in the exact Gaussian inner product,
off the twelve lower-$\ell$ functions of the degree-two space
($\sigma_1^2,\sigma_2$, the five components of $\sigma_1T$ and of
$(G^2)_{\rm tl}$) and evaluates the quotient of the remainder exactly.
\end{proof}

\begin{record}[Correction of a floating remark in V49]
\label{f16v51:correction}
V49 described the eighth level $13.75$ of its degree-two basis as the start of
the $\det$-odd sector. It is the $\ell=4$ level $55/4$ of the even sector,
nine-fold degenerate; the $\det$-odd ground of the degree-two space is
$14.15232$. V49's certified statements are unaffected, since they concern the
first seven ordered levels only.
\end{record}

\subsection{Convergence of the sector spectra}

\begin{center}\small
\begin{tabular}{@{}rrrrrr@{}}
\toprule
Gram degree & $(0,+)$ ground & $(2,+)$ ground & $(0,+)$ second & $(0,-)$ ground & $(2,-)$ ground\\
\midrule
2&6.55701&10.02312&10.68789&14.15232&20.00000\\
4&6.53734&9.60993&10.23392&13.96839&18.18664\\
6&6.53510&9.55557&10.15532&13.94944&18.00832\\
8&6.53491&9.54846&10.14063&13.94813&17.98959\\
\bottomrule
\end{tabular}
\end{center}
These are floating eigenvalues of the exact sector matrices (not hashed);
the last row's certified shifts are those of Theorem~\ref{f16v51:rr}. The
ground and the tensor level move by less than $10^{-2}$ over the last two
degrees, so the floating picture is $e_0\approx 6.5349$,
$e_1=\cdots=e_5\approx 9.5485$, $g_*\approx 3.014$, and the scalar
excitation at $\approx 10.1409F$; these are estimates, not enclosures. In the
floating tensor ground $\langle|q|^{-2}\rangle\approx 0.314$, so the
angular energy actually carried by that state is about $0.94$, well
below the floating gap $3.014$: most of the tensor gap is radial, not
centrifugal, which is why the certified bound is coarse.

\subsection{What the sector picture leaves open}

Certified now: $e_0\in[5.547,6.5352]$; the tensor multiplet lies at least
$0.1316$ above the ground and at most $9.5487$; every sector with $\ell\ge2$
lies at least $2\ell(\ell+1)/E_\ell^2$ above the ground. Not certified: any
lower bound on the scalar excitation $e_0^{(0,+),2}$ or on the $\det$-odd
ground, hence on $g_*$ itself if a scalar state were the first excitation.
The floating spectra put both scalar candidates well above the tensor level,
so a certified two-sided $g_*$ would follow from a lower bound on the two
scalar sectors of the form $e^{(0,\pm)}\ge 9.5487$; the natural route is a
Temple inequality inside each scalar sector, which needs an a priori bound on
the sector's next level, or a Dirichlet-wall argument for the $\det$-odd
profile. Both are model tasks for V52.

\subsection{Verification and preservation}

The new exact diagnostic is
\begin{center}\small\path{scripts/verify_ym_f16_v51_sector_gaps.py}.\end{center}
It implements the Stein recursion and gradient identities of
Proposition~\ref{f16v51:stein}, checks the degree-two consistency with V49
and the exact $55/4$, builds the four sector matrices at the requested Gram
degree, certifies the five shifts of Theorem~\ref{f16v51:rr} by exact
inertia, and records the exact rational gap $12/\lambda_2^2$. Floating
spectra and the $\langle|q|^{-2}\rangle$ diagnostics are not hashed. It has
15818 bytes and SHA--256
\begin{center}\scriptsize\ttfamily 5BCBFE65A8266C8801E20AB05633A687DAD33A6E754E92FBAD8C65AFB40536F5.\end{center}
Its canonical hashed payload at Gram degree 8 has SHA--256
\begin{center}\scriptsize\ttfamily 6D25DBE80CFD89436CBF92E79D8DD211C54425269C6C43E51339B6EAAE54F8E9.\end{center}
Only this terminal insertion and the literal title version differ from V50;
removing them recovers its bytes exactly. The next companion checkpoint is
V55.

\begin{firewall}
V51 proves the polar-coordinate Dirichlet form of the quartic model, the
centrifugal separation of every non-scalar right-$O(3)$ sector from the
ground, an exact Gram-polynomial moment calculus, and sector-resolved
certified upper bounds with a certified positive gap for the tensor
multiplet. It does not certify a lower bound for the scalar sectors, an
enclosure of $g_*$, any native volume-uniform estimate, a physical lattice
spacing, or an interacting continuum mass gap. The original transfer and the
V48 spectral identification are unchanged.
\end{firewall}
% END F16 V51 APPEND

% BEGIN F16 V52 APPEND -- 2026-09-18
% Insert before the final document-ending command of cumulative V51.
% This insertion leaves every predecessor body and the native operator unchanged.

\clearpage
\section{V52: residual inclusion intervals and the small-volume scope of the tube}
\label{f16v52:context}

Two loose ends of V49--V51 are closed and one scope fact is made explicit.
First, an ownership correction: V50's explicit hard spectrum restates, at the
trivial flat point, formulas that V41 already proved; V50's new content is
narrower than its text suggests and is delimited below. Second, the exact
Gram-polynomial calculus of V51 extends to $H_{\rm mat}$ applied to a basis
function, so that residuals $\|(H_{\rm mat}-\rho)f\|$ are exact rationals;
Weinstein's theorem then places an eigenvalue of the model in an explicit
interval around each sector Rayleigh--Ritz value, without any a priori
spectral information. This is an inclusion, not an exclusion: it locates
eigenvalues near $9.55$ and near $10.14$ but does not prove them to be the
first and second excitations. Third, the exact hard Gaussian ground of V44
has a second moment of order $V/\gamma$, so its mass inside the fixed tube of
V44 is exponentially small unless $\gamma$ exceeds a fixed multiple of the
volume; the fixed-regulator critical branch is a small-volume branch in this
precise sense.

\subsection{Ownership correction for V50}

Proposition~\ref{f16v50:spectrum} states the spectrum $\{\hat k^2\}$ of the
hard Hessian, the Mehler frequencies $\xi_k=\operatorname{arcosh}(1+\hat k^2/2)$,
and $\lambda_*=\exp(-3\sum_{k\ne0}\xi_k)$. All three were proved in V41:
Proposition~\ref{f16v41:flat-spectrum} gives the curl symbol
$d_1^*d_1=|d|^2I-dd^*$ with eigenvalues $\lambda_L(k,\varphi),\lambda_L,0$ for
every flat twist $\varphi$, and Theorem~\ref{f16v41:main} displays
$\Lambda_{{\rm G},B}=\exp[-3\sum_k\operatorname{arcosh}(1+\lambda_L(k,0)/2)]$.
V50's proposition is the special case $\varphi=0$ written in the transverse
plane-wave basis. The content that V50 added is the closed form
$1-r_+=\sqrt{x^2+2x}-x$, $x=2\sin^2(\pi/L)$, its two-sided elementary bounds,
the inverse bound for $R_+$, and Proposition~\ref{f16v50:adiabatic}, the
volume-uniform ratio $h/(1-r_+)<0.41\gamma^{-1/3}$. V50's firewall sentence
``proves the explicit hard spectrum, Mehler frequencies, Gaussian threshold''
should be read as ``restates V41's explicit hard spectrum, Mehler frequencies
and Gaussian threshold at the trivial flat point''. V41's winding theorem also
already supplies the toron-valley effective potential
$\mathcal E_L(\varphi)-\mathcal E_L(0)$ with positive coefficients and the
cusp $\frac2L|\varphi|$; no separate computation of that potential is needed.

\subsection{The Hamiltonian on Gram polynomials and exact residuals}

\begin{proposition}[Exact action of $H_{\rm mat}$ on the sector functions]
\label{f16v52:H-action}
For a polynomial $P$ in the Gram entries,
\begin{equation}
 H_{\rm mat}\bigl(Pe^{-\sigma_1}\bigr)
 =\Bigl[-\tfrac12\Delta_gP+4\mathcal EP-2\sigma_1P+9P+2\sigma_2P\Bigr]e^{-\sigma_1},
 \qquad
 \Delta_gP=\sum_{v,w}\partial_v\partial_wP\;\Gamma_{vw}+6\sum_k\partial_{g_{kk}}P,
 \label{f16v52:H-even}
\end{equation}
with $\Gamma_{vw}=\nabla g_v\cdot\nabla g_w$ the polynomial of
Proposition~\ref{f16v51:stein} and $\mathcal E$ its Euler operator, and
\begin{equation}
 H_{\rm mat}\bigl(\det(q)Pe^{-\sigma_1}\bigr)
 =\det(q)\Bigl[-\tfrac12\Delta_gP+4\mathcal EP-2\sigma_1P+15P
 -2\textstyle\sum_k\partial_{g_{kk}}P+2\sigma_2P\Bigr]e^{-\sigma_1}.
 \label{f16v52:H-odd}
\end{equation}
Consequently, for every rational vector $c$ in a sector basis and
$f=\sum_ic_if_i$, the numbers $\rho=\langle f,H_{\rm mat}f\rangle/\|f\|^2$ and
$s^2=\|H_{\rm mat}f\|^2/\|f\|^2-\rho^2$ are exact rationals.
\end{proposition}
\begin{proof}
$\Delta(Pe^{-\sigma_1})=(\Delta P)e^{-\sigma_1}+2\nabla P\cdot\nabla e^{-\sigma_1}
+P\Delta e^{-\sigma_1}$ with $\nabla e^{-\sigma_1}=-2qe^{-\sigma_1}$,
$q\cdot\nabla P=2\mathcal EP$ and $\Delta e^{-\sigma_1}=(4\sigma_1-18)e^{-\sigma_1}$
on $\mathbb R^9$; $\Delta P$ is the chain rule with $\Delta g_{kk}=6$ and
$\Delta g_{kl}=0$ for $k\ne l$. For the odd sector add
$2\nabla\det q\cdot\nabla(Pe^{-\sigma_1})=2\det(q)e^{-\sigma_1}
[\sum_k\partial_{g_{kk}}P-3P]$ and $\Delta\det q=0$. The moments of the
resulting Gram polynomials, times $1$ or $\sigma_3$, are the exact rationals
of Proposition~\ref{f16v51:stein}. As a check, the pure Gaussian has
$H_{\rm mat}e^{-\sigma_1}=(9-2\sigma_1+2\sigma_2)e^{-\sigma_1}$, $\rho=27/4$ and
$s^2=27/16$.
\end{proof}

\begin{theorem}[Residual inclusion intervals at Gram degree 8]
\label{f16v52:weinstein}
Let $f$ be the rounded Rayleigh--Ritz vector of a sector at Gram degree
8, with $\rho$ and $s$ computed exactly by
Proposition~\ref{f16v52:H-action}. Then $H_{\rm mat}$ has at least one
eigenvalue in $[\rho-s,\rho+s]$, in the sector of $f$. The certified intervals
are
\begin{center}\small
\begin{tabular}{@{}lrrl@{}}
\toprule
sector, level & $\rho$ & $s$ & eigenvalue in\\
\midrule
$(0,+)$, ground&6.5349&0.0325&$[6.5024,\,6.5674]$\\
$(0,+)$, second&10.1406&0.3578&$[9.7828,\,10.4984]$\\
$(2,+)$, ground&9.5485&0.1993&$[9.3491,\,9.7478]$\\
$(0,-)$, ground&13.9481&0.0751&$[13.8730,\,14.0233]$\\
$(2,-)$, ground&17.9896&0.3295&$[17.6601,\,18.3191]$\\
\bottomrule
\end{tabular}
\end{center}
\end{theorem}
\begin{proof}
For a self-adjoint operator and a normalized $f$ in its domain,
$s^2=\|(H-\rho)f\|^2\ge\operatorname{dist}(\rho,\operatorname{spec}H)^2$ by
the spectral theorem; the spectrum of $H_{\rm mat}$ is discrete by V44, so
the nearest spectral point is an eigenvalue. Since $f$ lies in one sector,
which reduces $H_{\rm mat}$, the eigenvalue lies in that sector. The rounding
of the floating eigenvector to a rational vector only changes $f$; the
certificate is for the rounded vector.
\end{proof}

These intervals, together with V49's $e_0\ge5.547$ and V51's certified
upper bounds, are the complete certified information: the ground lies in
$[5.547,6.5352]$; the tensor sector has an eigenvalue within $0.1993$ of
$9.5485$ and its ground is at most $9.5487$; the scalar sector has an
eigenvalue within $0.3578$ of $10.1406$. What is not certified is that the
interval around $9.5485$ contains the lowest tensor level, or that no scalar
eigenvalue lies between the ground interval and the interval around
$10.1406$. The residual $s$ decreases with the Gram degree
($0.273$ and $0.033$ for the scalar ground at degrees $4$, 8), but Weinstein's
theorem cannot turn inclusion into exclusion; that step needs a lower bound
on the next level, as recorded in V51.

\subsection{The hard Gaussian ground does not fit in one tube unless \texorpdfstring{$\gamma\gtrsim V$}{gamma >~ V}}

Let $\varphi$ be the unit ground function of V44's hard transfer
$\mathcal T_+$ in the hard coordinate $u\in\mathcal Y$, $y=\varepsilon u$, and
let $\omega_k=\sinh\xi_k=\sqrt{\hat k^2+\hat k^4/4}$ be V38's oscillator
frequencies.

\begin{proposition}[Second moment and tube mass of the hard ground]
\label{f16v52:tube-mass}
$|\varphi|^2$ is the product of centered Gaussians of variance $1/(2\omega_k)$,
six for each nonzero momentum, so
\begin{equation}
 \boxed{\quad
 \int|y|^2\varphi(u)^2du=\varepsilon^2\,c_LV,\qquad
 c_L=\frac3V\sum_{k\ne0}\frac1{\sinh\xi_k},
 \qquad c_8=0.9087,\ c_{16}=0.9243,\ c_{64}=0.9293.
 \quad}
 \label{f16v52:second-moment}
\end{equation}
For every $r>0$ with $\gamma r^2<\frac12c_LV$,
\begin{equation}
 \boxed{\quad
 \int_{|y|\le r}\varphi(u)^2du\ \le\
 \exp\Bigl[-\omega_{\min}\Bigl(\tfrac12c_LV-\gamma r^2\Bigr)\Bigr],
 \qquad \omega_{\min}=\sinh\xi_{\min}=2\sin\tfrac\pi L\sqrt{1+\sin^2\tfrac\pi L}.
 \quad}
 \label{f16v52:tube-mass-bound}
\end{equation}
In particular, at fixed tube radius $r_0$ and $\gamma\le c_LV/(4r_0^2)$, the
mass of the hard Gaussian ground inside V44's tube is at most
$e^{-\omega_{\min}c_LV/4}\le e^{-c\,L^2}$.
\end{proposition}
\begin{proof}
V38's Mehler calculation gives the ground $e^{-\omega u^2/2}$ for the one-mode
kernel $t_\nu$ with $\omega=\sqrt{\nu+\nu^2/4}$; on $\mathcal Y$ the modes are
the transverse waves of V41/V50 with $\nu=\hat k^2$, six per nonzero $k$, so
$\varphi^2$ is the stated product and $\int|u|^2\varphi^2=\sum_{\rm modes}
1/(2\omega)=3\sum_{k\ne0}\omega_k^{-1}$; the values of $c_L$ are the
verifier's sums, and $\omega_k=\sinh\xi_k$ follows from
$\cosh\xi_k=1+\hat k^2/2$. For the tail, Markov's inequality on
$e^{-\lambda|y|^2}$ gives $\int_{|y|\le r}\varphi^2\le e^{\lambda r^2}
\prod_{\rm modes}(1+\lambda\varepsilon^2/\omega_k)^{-1/2}$. Choose
$\lambda=\omega_{\min}/\varepsilon^2$, so that every ratio
$\lambda\varepsilon^2/\omega_k\le1$ and $\log(1+x)\ge x/2$ applies:
the product is at most $\exp[-\frac12\lambda\varepsilon^2\cdot3\sum_{k\ne0}
\omega_k^{-1}]=\exp[-\frac12\lambda\varepsilon^2c_LV]$. With
$\varepsilon^2=\gamma^{-1}$ this is \eqref{f16v52:tube-mass-bound}. The last
statement uses $\gamma r_0^2\le c_LV/4$ and $\omega_{\min}\ge2\sin(\pi/L)
\ge4/L$ with $V=L^3$.
\end{proof}

The proposition does not contradict V44--V48, whose limits fix $L$ and send
$\gamma\to\infty$: for $\gamma\ge c_LV/r_0^2$ the Gaussian bulk is inside the
tube and the cutoff $\chi(y)$ of V45 costs only an exponentially small
amount. It does fix the regime in which those theorems are informative:
the hard fluctuation $|y|\approx\sqrt{c_LV/\gamma}$ must be small compared
with the fixed tube radius, i.e.\ the coupling must exceed a fixed multiple of
the number of lattice sites. In the physical continuum regime, where the
Wilson coupling grows only logarithmically in the number of sites, the
Gaussian bulk leaves every fixed tube, and a construction that keeps the
fine-link fluctuation in a single based-gauge chart around a central-flat
orbit cannot be uniform in $L$. Any volume-uniform version of the critical
branch must therefore localize gauge fixing and the harmonic comparison
site by site, controlling densities rather than the total displacement
$|y|$; this is the same lesson recorded at the f15 V100 checkpoint, now
with an explicit exponent.

\subsection{Verification and preservation}

The new exact diagnostic is
\begin{center}\small\path{scripts/verify_ym_f16_v52_residual_intervals.py}.\end{center}
It imports V51's engine, implements \eqref{f16v52:H-even}--\eqref{f16v52:H-odd},
checks the Gaussian identity $\rho=27/4$, $s^2=27/16$, and computes the exact
$\rho$, $s^2$ and a rational upper bound on $s$ for the rounded
Rayleigh--Ritz vectors of the five sector levels at the requested Gram
degree; the second-moment sums $c_L$ are floating diagnostics. It has
7994 bytes and SHA--256
\begin{center}\scriptsize\ttfamily 6FF698614A22B04DF0549989ECE765638263C0B6A8D97BBAA3C51F619E912F08.\end{center}
Its canonical hashed payload at Gram degree 8 has SHA--256
\begin{center}\scriptsize\ttfamily 951DDB6418950424019C6AAC5685BB8CF1B6D4D09BB566508EBF358FBF83274B.\end{center}
Only this terminal insertion and the literal title version differ from V51;
removing them recovers its bytes exactly. The next companion checkpoint is
V55.

\begin{firewall}
V52 corrects the ownership of V50's spectral formulas (V41), proves exact
residual inclusion intervals for the sector Rayleigh--Ritz states of the
quartic model, and proves that the hard Gaussian ground of V44 has mass at
most $e^{-\omega_{\min}(c_LV/2-\gamma r^2)}$ in a tube of radius $r$. It does
not prove exclusion intervals, a lower bound on the model gap, a
volume-uniform native estimate, or an interacting continuum mass gap; it
records that the fixed-regulator critical branch requires $\gamma\gtrsim c_LV$.
\end{firewall}
% END F16 V52 APPEND

% BEGIN F16 V53 APPEND -- 2026-09-18
% Insert before the final document-ending command of cumulative V52.
% This insertion leaves every predecessor body and the native operator unchanged.

\clearpage
\section{V53: the exact hard-band corrector and the true small-volume condition}
\label{f16v53:context}

V45 corrects the recovery vectors by $hk(q)=hR_+\mathcal C(q)\varphi$ and
records the Gram $\langle\Psi_\gamma f,\Psi_\gamma g\rangle=\langle f,g\rangle
+h^2\int\bar fg\,\|k(q)\|^2dq+O(e^{-c/\varepsilon^2})$ without evaluating
$\|k(q)\|^2$. With the explicit hard spectrum of V41/V50 and the Mehler
structure of the hard ground, $\mathcal C(q)\varphi$ and $k(q)$ are explicit:
the first variation of the hard Hessian under a constant field is a colour
rotation tensored with a lattice central difference, and its Gaussian second
moments reduce to one momentum sum. The result is
$\|k(q)\|^2=K_L|q|^2$ with $K_L=\kappa_LV$ and $\kappa_L\to0.148$. Hence the
second-order Gram term of V45 equals $\kappa_L\gamma^{-2/3}L\,\langle f,|q|^2g\rangle$:
it is small exactly when $L\ll\gamma^{2/3}$, the genuine small-volume
condition of the weak-coupling expansion, and unlike V52's tube condition
$\gamma\gtrsim c_LV$ it is not a bookkeeping artifact. The extensive parts of
the second-order energy cancel, as they must by the periodicity of V41's flat
threshold in the twist; that cancellation is stated, not yet used.

\subsection{Dependencies}

The inputs are V41's flat-twist spectrum $\lambda_L(k,\varphi)$ with the
adjoint twist $\varphi_\mu=2La_\mu$ for a constant link angle $a_\mu$, V44's
hard transfer and its Mehler ground, V45's definitions of $H_1(q)$,
$\mathcal C(q)$, $k(q)$ and its corrected Gram, and V48's identity $DA(0)=0$
in the normal slice, which removes the covariance-derivative terms from
$\mathcal C(q)$. The next companion checkpoint is V55.

\subsection{The first variation of the hard Hessian is a rotation times a difference}

Let $J_a$, $a=1,2,3$, be the antisymmetric $3\times3$ colour generators,
$(J_a)_{bc}=-\epsilon_{abc}$, and $D_\mu$ the central difference on
functions of the site, $(D_\mu y)(x)=\frac12[y(x+\hat\mu)-y(x-\hat\mu)]$,
acting on each link component; on the plane wave $e^{ik\cdot x}$ it has the
eigenvalue $i\sin k_\mu$.

\begin{proposition}[First variation on the transverse space]
\label{f16v53:H1}
On $\mathcal Y$, with $H(a)=D_y^2S(\operatorname{Exp}(\mathsf Ka+y))|_{y=0}$
compressed to $\mathcal Y$ as in V45,
\begin{equation}
 \boxed{\quad
 H_1(q)=DH(0)[q]=-4\sum_{\mu=1}^3(q_\mu\cdot J)\otimes D_\mu\Big|_{\mathcal Y},
 \qquad (q_\mu\cdot J)=\sum_aq_\mu^aJ_a,
 \quad}
 \label{f16v53:H1-formula}
\end{equation}
where the overall sign depends only on the orientation conventions of $J$ and
$D$. Each summand commutes with $H_+$ and with the divergence, preserves
$\mathcal Y$, and is symmetric; for a single colour axis and direction it has
eigenvalues $\pm4\sin k_\mu$ on the two charged real colours at momentum $k$
and $0$ on the neutral colour.
\end{proposition}
\begin{proof}
$H(a)$ is linear in the derivative, so it suffices to differentiate along an
abelian direction $a_\mu=\alpha e_a$, $\alpha\to0$, for a fixed $(\mu,a)$.
For such a background the configuration is flat, and V41's
Proposition~\ref{f16v41:flat-spectrum} gives the Hessian on the complex
charged colour line as the twisted covariant curl, with symbol
$\lambda_L(k,\varphi)=4\sum_\nu\sin^2((2\pi k_\nu+\varphi_\nu)/(2L))$ on the
transverse polarizations and the adjoint twist $\varphi_\mu=2L\alpha$. Hence
$\partial_\alpha\lambda_L=2L\cdot\frac2L\sin(2\pi k_\mu/L)=4\sin k_\mu$ for
the charge-$(+1)$ line and $-4\sin k_\mu$ for its conjugate, and $0$ for the
neutral colour. The real operator with these eigenvalues on the charged
pair is $\mp4J_a\otimes D_\mu$: $J_a$ has eigenvalues $\pm i$ on the charged
pair, $D_\mu$ has eigenvalue $i\sin k_\mu$, and their product is real
symmetric. The variation of the transverse projector contributes
$d_\varphi(\partial d_\varphi^*)$ on $\mathcal Y$, which maps into the gauge
directions and is removed by the compression; $D_\mu$ and $J_a$ commute with
the divergence and with the Laplacian, so the compression is the restriction.
The sign and the value of the coefficient were also checked by central finite
differences of the literal Wilson action on the $L=4$ torus: for transverse
charged test vectors the second directional derivative changes at first order
in $\alpha$ by $-4\alpha\,y^{\mathsf T}(J_3\otimes D_1)y$ to six digits, and
its value at $\alpha=0$ is $\hat k^2|y|^2$.
\end{proof}

\subsection{The corrector and its exact norm}

Let $\Omega=\sinh\xi(H_+)$ on $\mathcal Y$, so that the hard ground is
$\varphi\propto e^{-u^{\mathsf T}\Omega u/2}$ with covariance
$C=(2\Omega)^{-1}$, and put $\hat B(q)=H_1(q)(1+e^{-2\xi(H_+)})$.

\begin{theorem}[Explicit corrector]
\label{f16v53:corrector}
With V48's $DA(0)=0$,
\begin{equation}
 \mathcal C(q)\varphi=-\frac{\lambda_*}4\bigl[u^{\mathsf T}\hat B(q)u-\Tr(\hat B(q)C)\bigr]\varphi,
 \qquad
 k(q)=-\frac14\sum_{k\ne0}\frac{1}{1-e^{-2\xi_k}}
 \bigl[u^{\mathsf T}\hat B_k(q)u-\Tr(\hat B_k(q)C)\bigr]\varphi,
 \label{f16v53:C-and-k}
\end{equation}
where $\hat B_k$ is the block of $\hat B$ at momentum $k$, and
\begin{equation}
 \boxed{\quad
 \|\mathcal C(q)\varphi\|^2=2\lambda_*^2|q|^2\sum_{k\ne0}
 \frac{\sin^2k_1\,(1+e^{-2\xi_k})^2}{\sinh^2\xi_k},
 \qquad
 \|k(q)\|^2=K_L|q|^2,\quad
 K_L=2\sum_{k\ne0}\frac{\sin^2k_1\coth^2\xi_k}{\sinh^2\xi_k}.
 \quad}
 \label{f16v53:norms}
\end{equation}
Both are isotropic in the nine components of $q$. Numerically
$K_L/V=0.1303,\ 0.1405,\ 0.1453,\ 0.1477$ and
$\|\mathcal C\varphi\|^2/(\lambda_*^2V|q|^2)=0.1140,\ 0.1169,\ 0.1174,\ 0.1174$
at $L=8,16,32,64$.
\end{theorem}
\begin{proof}
By V45, $\mathcal C(q)$ has kernel $\mathcal T_+(u,v)\{-\frac14u^{\mathsf T}H_1u
-\frac14v^{\mathsf T}H_1v\}$ once $A_1=0$. In an orthonormal eigenbasis of
$H_+$ in which the fixed abelian summand of $H_1$ is diagonal, with entries
$\eta_m$, the one-mode Mehler identities
$t_\nu\varphi_0=e^{-\xi/2}\varphi_0$ and
$t_\nu[u^2\varphi_0]=e^{-\xi/2}[\frac{1-e^{-2\xi}}{2\omega}+e^{-2\xi}u^2]\varphi_0$
(checked exactly by the Gaussian integral at rational $e^{\xi}$ in the
verifier) give
\[
 \mathcal C\varphi=-\frac{\lambda_*}4\sum_m\eta_m
 \Bigl[u_m^2+\frac{1-e^{-2\xi_m}}{2\omega_m}+e^{-2\xi_m}u_m^2\Bigr]\varphi
 =-\frac{\lambda_*}4\sum_m\eta_m(1+e^{-2\xi_m})\Bigl[u_m^2-\frac1{2\omega_m}\Bigr]\varphi,
\]
since $\sum_m\eta_m/\omega_m=0$ by the pairing $\pm\eta$ at each momentum.
This is the first formula of \eqref{f16v53:C-and-k}, and it extends to a
general $q$ by linearity in $H_1$. Each bracket $u_m^2-\frac1{2\omega_m}$, and
more generally $u_iu_j-C_{ij}$ for two modes at the same momentum $k$, is an
eigenfunction of $\mathcal T_+$ with eigenvalue $\lambda_*e^{-2\xi_k}$, so
$R_+$ acts on the momentum-$k$ block by $[\lambda_*(1-e^{-2\xi_k})]^{-1}$;
this gives $k(q)$. For the norms, Wick's identity
$\operatorname{Var}(u^{\mathsf T}Bu)=2\Tr(BCBC)$ for the centred Gaussian
$\varphi^2$ (checked exactly on a rational example) gives
$\|\mathcal C\varphi\|^2=\frac{\lambda_*^2}{8}\Tr(\hat BC\hat BC)$, and by
\eqref{f16v53:H1-formula}
\[
 \Tr(\hat BC\hat BC)=16\sum_{\mu\nu ab}q_\mu^aq_\nu^b\,\Tr(J_aJ_b)\,
 \Tr_{\mathcal Y'}\bigl(D_\mu GD_\nu G\bigr),\qquad
 G=\frac{1+e^{-2\xi(H_+)}}{2\Omega},
\]
with $\Tr(J_aJ_b)=-2\delta_{ab}$ and, over the two transverse spatial
polarizations at each $k$, $\Tr(D_\mu GD_\nu G)=-2\sum_k\sin k_\mu\sin k_\nu G_k^2
=-2\delta_{\mu\nu}\sum_k\sin^2k_1G_k^2$ by the reflection $k_\nu\to-k_\nu$
and cubic symmetry. Collecting factors gives the first norm; inserting the
resolvent factors $(1-e^{-2\xi_k})^{-2}$ per momentum block gives $K_L$, with
$(1+e^{-2\xi})/(1-e^{-2\xi})=\coth\xi$. The numerical values are the
verifier's sums.
\end{proof}

\subsection{The second-order Gram term and the true small-volume condition}

\begin{proposition}[Exact size of V45's Gram correction]
\label{f16v53:gram}
For invariant compact smooth $f,g$,
\begin{equation}
 \boxed{\quad
 \langle\Psi_\gamma f,\Psi_\gamma g\rangle-\langle f,g\rangle
 =h^2K_L\langle f,|q|^2g\rangle+O(e^{-c/\varepsilon^2})
 =\kappa_L\,\gamma^{-2/3}L\,\langle f,|q|^2g\rangle+O(e^{-c/\varepsilon^2}),
 \qquad \kappa_L=K_L/V\in[0.13,0.15].
 \quad}
 \label{f16v53:gram-size}
\end{equation}
At fixed $L$ this is $O(h^2)$ as V45 states; uniformly in $L$ it is small
exactly when $L\ll\gamma^{2/3}$.
\end{proposition}
\begin{proof}
Insert \eqref{f16v53:norms} into V45's corrected Gram and use
$h^2V=\gamma^{-2/3}L^{-2}L^3$.
\end{proof}

Three remarks fix the interpretation. First, the condition $L\ll\gamma^{2/3}$
is the natural regime of the weak-coupling small-volume expansion: the
leading soft scale $he_0\sim\gamma^{-1/3}/L$ exceeds the coupling scale
$1/\gamma$, at which the hard band feeds back on the soft dynamics, exactly
when $L\ll\gamma^{2/3}$. The same power appears in the diagonal
Born--Oppenheimer term $\frac h2\|\partial_q\varphi_q\|^2=\frac{h^3}2K'V
=O(\gamma^{-1})$, an $L$-uniform constant. Second, the extensive size of
$\|\mathcal C\varphi\|^2\propto V$ is the orthogonality catastrophe of the
hard ground under a constant background, $\langle\varphi_a,\varphi_0\rangle
=e^{-O(V|a|^2)}$; the first-order corrector $\varphi+hk$ is therefore an
accurate representative of $\varphi_{hq}$ only for $L\ll\gamma^{2/3}$, while
the adiabatic ansatz with the exact $\varphi_{hq}$ has Gram $\langle f,g\rangle$
at every $L$. Third, the second-order energy is not extensive: the term
$-h^2\langle\mathcal C\varphi,R_+\mathcal C\varphi\rangle/\lambda_*$, which is
$O(V|q|^2)$, is cancelled up to $O(L|q|^2)$ by the expectation of the second
variation $\frac12D^2H(0)[q,q]$, because the exact flat threshold
$\mathcal E_L(\varphi)$ of V41 is a periodic function of the twist whose
quadratic coefficient is $\sum_za_L(z)z_\mu^2=O(1/L)$, not $O(V)$; in the
constant-field variable this is $O(L)|a|^2=O(\gamma^{-2/3}L^{-1})|q|^2$,
relative order $\gamma^{-1/3}|q|^2$ to the leading quartic term, uniform in
$L$. Making this cancellation explicit mode by mode, and thereby extracting
the $L$-uniform $O(h^2)$ correction to the soft Hamiltonian, is the next
model task; it is the lattice form of the first correction of the
small-volume expansion.

\subsection{Where this leaves the scope map}

Three conditions now bound the informative regime of the critical branch at
lattice size $L$: the tube containment of V52, $\gamma\gtrsim c_LV/(2r_0^2)$,
which is an artifact of the global based-gauge tube; the corrector condition
$\gamma\gg L^{3/2}$ of this note, which is intrinsic; and the unconverted
fixed-$L$ remainders of V45--V48. Removing the first requires a local gauge
fixing; the second is the physical boundary of the small-volume expansion and
cannot be removed, only crossed by a different method. Volume-uniform work
should therefore aim at the regime $L\le\gamma^{2/3-\delta}$ with local charts,
and at the crossover beyond it with a construction that does not expand
around constant modes.

\subsection{Verification and preservation}

The new diagnostic is
\begin{center}\small\path{scripts/verify_ym_f16_v53_corrector_norm.py}.\end{center}
Its hashed exact content is the one-mode Mehler identities at four rational
values of $e^{\xi}$ through the exact Gaussian integral, Wick's variance
identity on a rational example, and the colour and momentum trace identities
at $L=4$. Its floating diagnostics, not hashed, are the finite-difference
check of \eqref{f16v53:H1-formula} on the literal $L=4$ Wilson action and the
mode sums of \eqref{f16v53:norms}. It has 12055 bytes and SHA--256
\begin{center}\scriptsize\ttfamily DAA2EBC78BACB7BA8052803C8824F45E88DFA5BD76DAA3387FC7D4F3EEDC1180.\end{center}
Its canonical hashed payload has SHA--256
\begin{center}\scriptsize\ttfamily F88E665F909C1F63F64C1278AEE20B6064D90A7A285A58A283D5F8367B716ABC.\end{center}
Only this terminal insertion and the literal title version differ from V52;
removing them recovers its bytes exactly. The next companion checkpoint is
V55.

\begin{firewall}
V53 proves the explicit first variation of the hard Hessian, the explicit
corrector of V45 and its exact norm $K_L|q|^2$ with $K_L\asymp0.15V$, and
identifies V45's second-order Gram term as $\kappa_L\gamma^{-2/3}L\langle f,|q|^2g\rangle$,
small exactly for $L\ll\gamma^{2/3}$. It does not compute the $L$-uniform
second-order soft Hamiltonian, does not make any estimate of V44--V48
uniform in $L$, and does not establish an interacting continuum mass gap.
\end{firewall}
% END F16 V53 APPEND

% BEGIN F16 V54 APPEND -- 2026-09-18
% Insert before the final document-ending command of cumulative V53.
% This insertion leaves every predecessor body and the native operator unchanged.

\clearpage
\section{V54: the second-order hard-band potential and the first correction to the critical model}
\label{f16v54:context}

V53 left one calculation open: the second-order energy of the hard band
under a constant field, whose extensive pieces must cancel. Here it is done
exactly. For an abelian constant background the hard Gaussian transfer of V45
is the $k\ne0$ part of V41's flat-twist threshold, so its norm is an explicit
periodic function of the twist; its quadratic coefficient $m_L$ is an
explicit momentum sum, negative, of size $1.20L$, while its two
Rayleigh--Schr\"odinger pieces are each of size $0.13V$. Colour rotation,
cubic symmetry and parity make the second-order energy isotropic in all nine
constant-mode coordinates. The consequence for the critical branch is the
first correction to the quartic model: the soft effective potential acquires
the harmonic term $\gamma^{-1/3}(m_L/L)|q|^2$, of relative order
$\gamma^{-1/3}$ uniformly in $L$, with $m_L/L\to-1.203$. This is the lattice
form of the leading correction of the small-volume expansion; it is a theorem
about the frozen hard Gaussian operator, and its insertion into the native
spectral asymptotics is stated as the next-order formula, not proved.

\subsection{Dependencies}

V41's Proposition~\ref{f16v41:flat-spectrum} (twisted spectrum
$\lambda_L(k,\varphi)$ for flat backgrounds with adjoint twist
$\varphi_\mu=2La_\mu$), V38/V41's Mehler product formula for the norm of a
Gaussian transfer, V45's frozen operator $\mathcal T(a)$ with its data
$(A(a),H(a))$ on $\mathcal Y$, V48's $DA(0)=0$, and V53's explicit
$H_1(q)$ and corrector. The next companion checkpoint is V55.

\subsection{The hard threshold along constant abelian backgrounds}

\begin{proposition}[Exact hard norm on the flat abelian valley]
\label{f16v54:abelian}
Let $a=\alpha\,e_{\mu}\otimes n$ be a constant abelian field, direction
$\mu$, unit colour $n$, and let $\mathcal T(a)$ be V45's frozen hard Gaussian
transfer. Then
\begin{equation}
 \boxed{\quad
 -\log\frac{\|\mathcal T(a)\|}{\lambda_*}=\Delta_L(\alpha)
 :=2\sum_{k\ne0}\Bigl[\xi\bigl(\lambda_L(k,2L\alpha e_\mu)\bigr)-\xi\bigl(\lambda_L(k,0)\bigr)\Bigr],
 \quad}
 \label{f16v54:Delta}
\end{equation}
an even smooth periodic function of $\alpha$ with period $\pi/L$, and
\begin{equation}
 \Delta_L(\alpha)=m_L\alpha^2+O_L(\alpha^4),\qquad
 \boxed{\ m_L=4\sum_{k\ne0}\frac{\cos k_\mu\sinh^2\xi_k-\sin^2k_\mu\cosh\xi_k}{\sinh^3\xi_k}\ }
 \label{f16v54:mL}
\end{equation}
with $\cosh\xi_k=1+\hat k^2/2$. Numerically $m_L=-9.081,\ -19.013,\ -38.409,\
-77.006$ and $m_L/L=-1.135,\ -1.188,\ -1.200,\ -1.203$ at $L=8,16,32,64$; the
finite-difference quadratic coefficient of \eqref{f16v54:Delta} agrees with
$m_L$ to six digits at $L=8$ and $L=16$.
\end{proposition}
\begin{proof}
The background is flat, and by V41 the Hessian of the action at it is the
twisted covariant curl, block diagonal in momentum, with the two charged real
colour lines seeing the twists $\pm\varphi$, $\varphi=2L\alpha e_\mu$, and the
neutral line untwisted. The frozen operator $\mathcal T(a)$ is the Gaussian
transfer with kinetic covariance $A(a)$ and magnetic Hessian $H(a)$ on
$\mathcal Y$; by V38's calculation its norm is the product of the Mehler ground
factors $e^{-\xi(\nu)/2}$ over the eigenvalues $\nu$ of
$A(a)^{1/2}H(a)A(a)^{1/2}$, which is the spectrum of the quotient Hessian in
the metric of the quotient at the flat point. Momentum conservation of the
constant background makes the $\mathcal Y$ (nonzero momentum) and
$\mathcal Z$ (zero momentum) blocks of both $A(a)$ and $H(a)$ decouple, so
the $\mathcal Y$ part of that spectrum is the $k\ne0$ part of V41's spectrum:
$\lambda_L(k,\varphi)$ and $\lambda_L(k,-\varphi)$ with two polarizations each
for the charged lines, $\lambda_L(k,0)$ twice for the neutral line, over
$k\ne0$. Since $\lambda_L(k,-\varphi)$ runs over the same values as
$\lambda_L(k,\varphi)$ under $k\to-k$, the product gives
$\lambda_*\exp(-\Delta_L(\alpha))$ with \eqref{f16v54:Delta}. Periodicity
in $\alpha$ with period $\pi/L$ is the shift $k_\mu\to k_\mu+2\pi/L$ of the
sum over $k\ne0$ together with the removal of the $k=0$ term, whose own
contribution is the cusp of V41; evenness is $k\to-k$. The removed $k=0$
term is the only nonsmooth term, so $\Delta_L$ is smooth and its Taylor
expansion has even powers. Differentiating $\xi(\lambda_L)$ twice,
$\partial_\alpha\lambda_L=4\sin k_\mu$ and $\partial_\alpha^2\lambda_L=8\cos k_\mu$
at $\alpha=0$ (exact half-angle identities, checked on the $L=4$ grid), and
$\xi'(t)=1/\sqrt{t(t+4)}=1/(2\sinh\xi)$, $\xi''(t)=-(t+2)/(t(t+4))^{3/2}
=-\cosh\xi/(4\sinh^3\xi)$ give
$m_L=\sum_{k\ne0}[8\cos k_\mu\,\xi'+16\sin^2k_\mu\,\xi'']$, which is
\eqref{f16v54:mL}.
\end{proof}

\subsection{Isotropy of the second-order energy}

\begin{theorem}[Second-order hard-band potential]
\label{f16v54:isotropy}
For every constant field $a\in\mathbb R^9$, $\|\mathcal T(a)\|$ is smooth near
$a=0$ and
\begin{equation}
 \boxed{\quad
 -\log\frac{\|\mathcal T(a)\|}{\lambda_*}=m_L|a|^2+O_L(|a|^4).
 \quad}
 \label{f16v54:isotropic}
\end{equation}
\end{theorem}
\begin{proof}
The norm of a Gaussian transfer with smoothly varying positive data
$(A(a),H(a))$ is $\prod e^{-\xi(\nu_j(a))/2}$, and $\sum_j\xi(\nu_j(a))$ is a
smooth function of $a$ near $0$ because $\xi$ is smooth on $(0,\infty)$ and
the spectrum of $A^{1/2}HA^{1/2}$ stays in a fixed compact subset of
$(0,\infty)$ for small $a$ ($H(0)=H_+\succ0$ on $\mathcal Y$). Its Taylor
expansion at $0$ has no linear term by V45's singlet cancellation (or by the
evenness proved next), a quadratic form $M(a,a)$, and no cubic term:
the map $a\mapsto-a$ is induced by the lattice parity $x\mapsto-x$ composed
with link inversion, a symmetry of the transfer, so the function is even.
The quadratic form is invariant under global colour rotations $a_\mu\mapsto
Ra_\mu$, which are gauge transformations, and under the cubic group acting on
the direction index, which are lattice symmetries; the only symmetric
bilinear forms on $\mathbb R^3_{\rm colour}\otimes\mathbb R^3_{\rm direction}$
invariant under $\operatorname{SO}(3)\times O_h$ are multiples of the identity, since both
factors are irreducible. Hence $M(a,a)=m|a|^2$, and $m=m_L$ by evaluation on
the abelian direction of Proposition~\ref{f16v54:abelian}. The remainder is
$O_L(|a|^4)$ by evenness.
\end{proof}

\begin{proposition}[Cancellation of the extensive pieces]
\label{f16v54:cancellation}
In the second-order expansion of $E(a)=-\log\|\mathcal T(a)\|$, the corrector
piece $-\langle\mathcal C(a)\varphi,R_+\mathcal C(a)\varphi\rangle/\lambda_*$
contributes $-P_L|a|^2$ with
\[
 P_L=2\sum_{k\ne0}\frac{\sin^2k_\mu\,(1+e^{-2\xi_k})^2}{\sinh^2\xi_k\,(1-e^{-2\xi_k})},
 \qquad P_L/V=0.1216,\ 0.1271,\ 0.1286,\ 0.1291\ (L=8,16,32,64),
\]
and the remaining second-order pieces, the expectation of the second
variation of the kernel including the squared first variation, contribute
$(m_L+P_L)|a|^2$ to $E$. Per mode, with $x=e^{-2\xi}$ and $H_1$-eigenvalue $\eta$,
the two pieces are $\eta^2(2+x)/(32\omega^2)$ and
$\eta^2(1+x)^2/(32\omega^2(1-x))$, and their sum is
$\eta^2(2\coth\xi+1)/(32\omega^2)=-\frac14\xi''\eta^2+\frac18\xi'^2\eta^2$,
the exact second-order coefficient of $e^{-\xi(\nu(\alpha))/2}$ without the
$\nu''$ term; the identity $(2+x)+(1+x)^2/(1-x)=2\coth\xi+1$ is checked
exactly at rational $e^{\xi}$ by the verifier.
\end{proposition}
\begin{proof}
The corrector piece is V53's $\|k\|$ computation with one resolvent factor
instead of two. The per-mode expansion of the Mehler ground factor
$e^{-\xi(\nu+\eta\alpha+\frac12\nu''\alpha^2)/2}$ to second order is
$-\frac14\xi'\nu''-\frac14\xi''\eta^2+\frac18\xi'^2\eta^2$, and the product
over modes reassembles the cross terms into $\exp(-\Delta_L)$, whose
quadratic coefficient is $-m_L$ since the linear terms cancel in pairs.
Everything else is bookkeeping of \eqref{f16v54:mL}.
\end{proof}

The two pieces are extensive because the first-order response of the hard
ground to a constant field is extensive (V53); their difference is not,
because $\Delta_L$ is a periodic function of the twist whose smooth part is
controlled by V41's winding coefficients $a_L(z)=O(1/L)$. The negative sign of
$m_L$ means that the hard band lowers the energy of a slowly varying
constant field at second order: the two charged colour lines see twists of
opposite sign, and the concavity of $\xi$ makes the average of
$\xi(\lambda(k+\varphi/L))$ and $\xi(\lambda(k-\varphi/L))$ smaller than
$\xi(\lambda(k))$ for most momenta, which the cusp of the removed $k=0$
term does not compensate.

\subsection{The first correction to the critical model}

In V45's critical variables $a=hq$, the hard-band energy
$m_L|a|^2=h^2m_L|q|^2$ enters the scaled deficits, which are measured in
units of $h$, as the potential $hm_L|q|^2$. With $h=\gamma^{-1/3}/L$,
\begin{equation}
 \boxed{\quad
 H_{\rm mat}\ \longrightarrow\ H_{\rm mat}+\gamma^{-1/3}\frac{m_L}L|q|^2,
 \qquad \frac{m_L}L\longrightarrow-1.203,
 \quad}
 \label{f16v54:correction}
\end{equation}
so the first correction to the quartic model is a negative harmonic term of
relative order $\gamma^{-1/3}$, uniform in $L$. Its first-order effect on the
model levels is $\delta e_j=\gamma^{-1/3}(m_L/L)\langle j|\,|q|^2\,|j\rangle$;
in the floating sector ground of V51, $\langle|q|^2\rangle$ can be read from
the same Gram-polynomial calculus. Two statements are separated here. The
formula for $m_L$ and the isotropy are theorems about the frozen hard Gaussian
operator. The insertion \eqref{f16v54:correction} into the ordered native
asymptotics of V48 would require the next order of V45--V48: the corrected
recovery vectors at order $h^2$, the soft kinetic corrections from the
quotient metric on $\mathcal Z$ (order $h^2$ relative), the sixth-order
constant-field action (order $h^2$ relative), and a coarse lower bound at that
order. Those are not proved; \eqref{f16v54:correction} is the natural
next-order formula, consistent with the structure of the small-volume
expansion, in which the first correction to the zero-mode Hamiltonian is a
one-loop harmonic term of relative order $g^{2/3}$. The regime of validity
is V53's: $L\ll\gamma^{2/3}$, where the correction is small; at
$L\sim\gamma^{2/3}$ the term $\gamma^{-1/3}(m_L/L)|q|^2$ is of order one
relative to $H_{\rm mat}$ on the soft scale, and the expansion around
constant modes stops being informative.

\subsection{Verification and preservation}

The new diagnostic is
\begin{center}\small\path{scripts/verify_ym_f16_v54_second_order_potential.py}.\end{center}
Its hashed exact content is the per-mode second-order identity and the
$\sinh/\cosh$ forms of $\xi',\xi''$ at four rational values of $e^{\xi}$, and
the derivative identities $\partial_\alpha\lambda_L=4\sin k_\mu$,
$\partial_\alpha^2\lambda_L=8\cos k_\mu$ on the $L=4$ grid. Its floating
diagnostics, not hashed, are $m_L$, $m_L/L$, $P_L$ for $L=8,16,32,64$ and the
finite-difference quadratic coefficients of $\Delta_L$ at $L=8,16$. It has
6254 bytes and SHA--256
\begin{center}\scriptsize\ttfamily 5626F4E495C3371559927D52DC637B854A2F37CF932F2911AF5505C92B002A5F.\end{center}
Its canonical hashed payload has SHA--256
\begin{center}\scriptsize\ttfamily 80A1ECEB707A89780FD10EB96700C1AE42CCC1D6153800BB721298A3129D8937.\end{center}
Only this terminal insertion and the literal title version differ from V53;
removing them recovers its bytes exactly. The next companion checkpoint is
V55.

\begin{firewall}
V54 proves the exact hard Gaussian norm along constant abelian backgrounds,
the isotropic second-order hard-band potential $m_L|a|^2$ with the explicit
negative coefficient $m_L\asymp-1.2L$, and the cancellation of its extensive
Rayleigh--Schr\"odinger pieces. It identifies the resulting first correction
$\gamma^{-1/3}(m_L/L)|q|^2$ to the quartic model but does not prove its
insertion into the native spectral asymptotics, does not establish any
$L$-uniform native estimate, and does not establish an interacting continuum
mass gap.
\end{firewall}
% END F16 V54 APPEND

% BEGIN F16 V55 APPEND -- 2026-09-18
% Insert before the final document-ending command of cumulative V54.
% This insertion leaves every predecessor body and the native operator unchanged.

\clearpage
\section{V55: companion checkpoint, five-version review, and the first-order corrected levels}
\label{f16v55:context}

V55 is the scheduled companion checkpoint. It records the rereading of the
companion programmes, reviews V50--V54, and adds one model-level result: the
first-order effect of V54's harmonic correction $\gamma^{-1/3}(m_L/L)|q|^2$
on the levels of the quartic model, computed from the exact expectation of
$|q|^2$ in the sector Rayleigh--Ritz states with the Gram-polynomial calculus.
The correction is negative and larger in the excited states, which are more
extended, so it lowers every gap coefficient at first order by an amount of
relative order $\gamma^{-1/3}$; the numbers are recorded, and the domain of
validity of the first-order formula is proved by an Agmon decay estimate for
the model eigenfunctions along the commuting valleys.

\subsection{V55 companion checkpoint}

Fresh inventories of \path{latex_notes} found no companion version newer than
those recorded at the V50 checkpoint: SM continuum V72 (f44f9745), NS
stability V26 (82c887f8), shell mixing V16 (bf138e63), parallel YM V5
(306f1bb8) and GRH cofinal visibility V3 (4dd31580) are unchanged in size and
hash; the only file changed since V50 in that directory is this lineage. Their
latest conclusions were reread against the V50 table and no new input to the
transfer analysis appears. In particular no companion supplies a
volume-uniform native spectral comparison, a physical-time calibration, or a
lower bound for the scalar sectors of the quartic model. The next companion
checkpoint is V60, which coincides with the next ten-version review.

\subsection{Five-version review, V50--V54}

\begin{center}\small
\begin{tabular}{@{}p{0.07\textwidth}p{0.88\textwidth}@{}}
\toprule
Version & Established (scope limits in each firewall)\\
\midrule
V50 & Checkpoint and V40--V49 review; correction of V49's normalization remark
(isotropic Wilson transfer at $\beta=\gamma$); closed form of the hard band
gap $1-r_+$ and the volume-uniform ratio $h/(1-r_+)<0.41\gamma^{-1/3}$.\\
V51 & Polar coordinates of $q$; centrifugal separation of every non-scalar
right-$O(3)$ sector, $e^{(\ell,\pm)}-e_0\ge2\ell(\ell+1)/E_\ell^2$; exact
Gram-polynomial calculus; sector-resolved certified upper bounds; the $\ell=4$
identification $55/4$.\\
V52 & Ownership correction for V50 (V41 already had the spectral formulas);
exact residual inclusion intervals (ground in $[6.5024,6.5674]$, tensor level
in $[9.349,9.748]$); the hard Gaussian ground has mass
$\le e^{-\omega_{\min}(c_LV/2-\gamma r^2)}$ in a tube of radius $r$.\\
V53 & Explicit first variation $H_1(q)=-4\sum_\mu(q_\mu\cdot J)\otimes D_\mu$;
explicit corrector of V45 with $\|k(q)\|^2=K_L|q|^2$, $K_L\asymp0.15V$; V45's
Gram correction $\kappa_L\gamma^{-2/3}L\langle f,|q|^2g\rangle$, small exactly
for $L\ll\gamma^{2/3}$.\\
V54 & Exact hard norm along constant abelian backgrounds; isotropic
second-order potential $m_L|a|^2$, $m_L\asymp-1.2L$; cancellation of the
extensive Rayleigh--Schr\"odinger pieces; the first correction
$\gamma^{-1/3}(m_L/L)|q|^2$ to the quartic model.\\
\bottomrule
\end{tabular}
\end{center}

The five versions fix the quantitative content of the fixed-regulator critical
branch at the model level, identify the intrinsic small-volume condition
$L\ll\gamma^{2/3}$, and separate it from the tube artifact
$\gamma\gtrsim c_LV$. The two remaining model tasks are unchanged: a certified
lower bound for the scalar sectors, and the next order of V45--V48 needed to
insert V54's correction into the native asymptotics. The native task is
unchanged: a local gauge fixing that removes the tube condition.

\subsection{Agmon decay of the model eigenfunctions along the valleys}

\begin{proposition}[Exponential decay of model eigenfunctions]
\label{f16v55:agmon}
Let $\psi$ be a normalized eigenfunction of $H_{\rm mat}$ with eigenvalue $e$.
For every $0<\theta<1$ there is $C_{e,\theta}$ such that
\begin{equation}
 \int e^{2\theta\rho(q)}|\psi(q)|^2dq\le C_{e,\theta},\qquad
 \rho(q)=\int_0^{|q|}\sqrt{2\bigl(2t-e\bigr)_+}\,dt,
 \label{f16v55:agmon-bound}
\end{equation}
so that $|\psi|^2$ has total mass at most $C_{e,\theta}e^{-2\theta\rho(R)}$
outside $\{|q|\le R\}$, with $\rho(R)\ge\frac{2\sqrt2}{3}\theta'(R-e)^{3/2}$
for $R>e$. In particular every polynomial moment $\langle\psi,|q|^{2n}\psi\rangle$
is finite.
\end{proposition}
\begin{proof}
V44's form bound $H_{\rm mat}\succeq2\sum_\mu|q_\mu|\succeq2|q|$ gives the
effective confining potential $W(q)=2|q|$, and the standard Agmon estimate for
$-\frac12\Delta+V$ with $V\ge W$ applies verbatim with the Agmon metric of
$2(W-e)_+$, since the exponential weight $e^{\theta\rho}$ is Lipschitz with
$|\nabla\rho|^2=2(W-e)_+$ and the commutator identity
$\Re\langle e^{\theta\rho}\psi,(H-e)e^{\theta\rho}\psi\rangle
=\langle\psi,(V-e-\tfrac12\theta^2|\nabla\rho|^2)e^{2\theta\rho}\psi\rangle$
holds for the true potential $V=2\mathcal Q$ only after replacing $V$ by the
form lower bound; to keep the argument within form bounds, apply the identity
to $H_{\rm mat}$ and use $\langle\phi,H_{\rm mat}\phi\rangle\ge\langle\phi,2|q|\phi\rangle$
on $\phi=e^{\theta\rho}\psi$ cut off at large $|q|$; the region where
$2|q|-e-\theta^2(2|q|-e)_+\le\frac12(1-\theta^2)e$ is a fixed ball, on which
$\psi$ has finite weighted mass. Letting the cutoff go to infinity proves
\eqref{f16v55:agmon-bound}; the lower bound on $\rho(R)$ is the elementary
integral.
\end{proof}

\subsection{First-order shifts under the harmonic correction}

\begin{theorem}[First-order corrected levels]
\label{f16v55:shifts}
Let $\varepsilon=\gamma^{-1/3}m_L/L<0$ and $R=\pi\gamma^{1/3}$, and let
$\chi_R$ be a smooth cutoff equal to one on $|q|\le R/2$ and zero outside
$|q|\le R$. For every isolated eigenvalue $e_j$ of $H_{\rm mat}$ (in its right
$O(3)$ multiplet) the operator $H_\varepsilon=H_{\rm mat}+\varepsilon\chi_R|q|^2$
has, for $|\varepsilon|R^2$ small compared with the isolation distance, an
eigenvalue group near $e_j$ of the same multiplicity, with
\begin{equation}
 \boxed{\quad
 e_j(\varepsilon)=e_j+\varepsilon\,\langle\psi_j,|q|^2\psi_j\rangle
 +O\bigl(\varepsilon^2R^4\bigr)+O\bigl(|\varepsilon|R^2e^{-c\,R^{3/2}}\bigr),
 \quad}
 \label{f16v55:first-order}
\end{equation}
where the multiplet remains degenerate because $|q|^2$ commutes with the right
action. In the refined sector Rayleigh--Ritz states at Gram degree 8 the
exact expectations are
\begin{center}\small
\begin{tabular}{@{}lrr@{}}
\toprule
sector, level & Rayleigh quotient & $\langle|q|^2\rangle$\\
\midrule
$(0,+)$, ground&6.5349&2.4577\\
$(0,+)$, second&10.1406&4.1119\\
$(2,+)$, ground&9.5485&3.7678\\
$(0,-)$, ground&13.9481&3.1074\\
$(2,-)$, ground&17.9896&3.9918\\
\bottomrule
\end{tabular}
\end{center}
and the pure Gaussian has $\langle|q|^2\rangle=9/4$ exactly.
\end{theorem}
\begin{proof}
The perturbation $\varepsilon\chi_R|q|^2$ is bounded with norm at most
$|\varepsilon|R^2$, so Kato's analytic perturbation theory applies to each
isolated eigenvalue group, giving the first two terms and the $O(\varepsilon^2R^4)$
remainder with the isolation distance in the constant. The difference between
$\langle\psi_j,\chi_R|q|^2\psi_j\rangle$ and $\langle\psi_j,|q|^2\psi_j\rangle$
is at most $\int_{|q|>R/2}|q|^2|\psi_j|^2$, which is $O(R^2e^{-cR^{3/2}})$ by
Proposition~\ref{f16v55:agmon}. The multiplet statement is Schur's lemma for
the right action, which commutes with $|q|^2$ and $\chi_R$. The expectations
are exact rationals of the refined vectors by the Stein recursion applied to
$\sigma_1PP'$.
\end{proof}

The cutoff at $R=\pi\gamma^{1/3}$ is where the constant link angle $a=hq$
reaches $\pi/L$, the period of V54's exact hard potential $\Delta_L$; beyond
it the quadratic truncation $m_L|a|^2$ is not meaningful and the full periodic
function must be used, which is bounded, so the corrected model is well defined
even though $H_{\rm mat}+\varepsilon|q|^2$ without the cutoff is not bounded
below. The first-order formula is therefore an asymptotic statement for
$\gamma\to\infty$ at fixed $L$, in which $R\to\infty$ and
$\varepsilon R^2=\pi^2(m_L/L)\gamma^{1/3}$ is not small; the useful form of
\eqref{f16v55:first-order} is with a fixed cutoff radius $R_0$ large compared
with the localization length of $\psi_j$ but independent of $\gamma$, where all
error terms are $O(\varepsilon^2)+O(e^{-cR_0^{3/2}})$.

\begin{record}[First-order corrected gap coefficients]
\label{f16v55:numbers}
With the degree-8 states, $\langle|q|^2\rangle$ increases from the ground
to the tensor multiplet by 1.310 and to the scalar excitation by 1.654.
Hence, to first order in $\varepsilon=-1.203\gamma^{-1/3}$,
\[
 g_*(\gamma)=3.014-1.576\,\gamma^{-1/3}+\cdots,\qquad
 g_*^{(0,+)}(\gamma)=3.606-1.990\,\gamma^{-1/3}+\cdots,
\]
numerically $\gamma=10$: tensor gap $2.282$, scalar gap $2.682$; $\gamma=100$: tensor gap $2.674$, scalar gap $3.177$; $\gamma=1000$: tensor gap $2.856$, scalar gap $3.407$. The correction lowers both gaps and preserves the
tensor multiplet as the first excitation at these couplings. These are
floating consequences of the exact expectations; the insertion into the native
asymptotics is the unproved next order of V45--V48 recorded in V54.
\end{record}

\subsection{Verification and preservation}

The new diagnostic is
\begin{center}\small\path{scripts/verify_ym_f16_v55_first_order_shifts.py}.\end{center}
Its hashed exact content is the Gaussian check $\langle\sigma_1\rangle=9/4$
and the exact Rayleigh quotients and $\sigma_1$-expectations of the five
refined sector states at the requested Gram degree; the corrected gap
coefficients at $\gamma=10,100,1000$ are floating diagnostics. It has
4676 bytes and SHA--256
\begin{center}\scriptsize\ttfamily 94818AEC03BF82634630E792FA063C9C96FE86CB5383540490CCF166C0E2C92B.\end{center}
Its canonical hashed payload at Gram degree 8 has SHA--256
\begin{center}\scriptsize\ttfamily B1F4DAB314BAF88076B2BEE1A1FDF9EFE113BA8C3B30E6E63E01CCA69E2FF1D8.\end{center}
Only this terminal insertion and the literal title version differ from V54;
removing them recovers its bytes exactly. The next companion checkpoint and
ten-version review is V60.

\begin{firewall}
V55 records an unchanged companion checkpoint and a five-version review,
proves Agmon decay of the quartic-model eigenfunctions along the commuting
valleys, and proves the first-order shift formula for the harmonic correction
with a cutoff, with exact expectations of $|q|^2$ in the sector states. It
does not certify the shifted levels beyond first order, does not insert the
correction into the native spectral asymptotics, and does not establish a
volume-uniform estimate or an interacting continuum mass gap.
\end{firewall}
% END F16 V55 APPEND

% BEGIN F16 V56 APPEND -- 2026-09-18
% Insert before the final document-ending command of cumulative V55.
% This insertion leaves every predecessor body and the native operator unchanged.

\clearpage
\section{V56: the explicit center-splitting exponent}
\label{f16v56:context}

V42 proves that the seven charged center sectors have tops within
$C_Le^{-\kappa_L\sqrt\gamma}$ of the neutral top, with an unspecified rate
$\kappa_L$, and lists the matching tunneling asymptotic as open. This note
supplies the sharp exponent of that upper bound. The flat set of V41 carries an
explicit loss function, the twist dependence of the harmonic threshold, and
the exponential weights of V42 can be adapted to it: a weight whose gradient
is bounded by the Agmon metric of the loss function keeps the weighted
transfer strictly below the band cutoff on the forbidden region. The Agmon
distance between the trivial central orbit and a center-twisted one is
attained along the abelian toron valley, by a product inequality for twisted
heat traces, and equals
\[
 S_L=\sqrt L\int_0^{2\pi}\sqrt{F_L(\varphi e_1)-F_L(0)}\,d\varphi,
 \qquad S_8=6.267,\ S_{16}=6.236,\ S_{32}=6.230,\ S_{64}=6.220,
\]
an $L$-independent quantity in the limit. The result is
$1-\Lambda_{\gamma,a}/\Lambda_\gamma\le C_{L,\delta}e^{-(1-\delta)S_L\sqrt\gamma}$
for every nontrivial character $a$ and every $\delta>0$. No lower bound on the
splitting is proved; the exponent is the natural one of a tunneling
asymptotic and the value of $S_L$ is the point of the note.

\subsection{Dependencies}

V41's harmonic principle, its flat-twist threshold $\mathfrak h(\varphi)
=e^{-\mathcal E_L(\varphi)}$, the winding theorem with its positive coefficients
and the strict loss bound; V42's displacement estimates, IMS localization,
forbidden-region barrier, weighted Sylvester argument and Feshbach map; V43's
sector ladders. The proof below modifies V41's chart argument by letting the
chart radius shrink with the coupling, and V42's weight by adapting it to the
loss function. The next companion checkpoint is V60.

\subsection{The loss function and its Agmon metric on the flat moduli}

Write $\ell_L(\varphi)=\mathcal E_L(\varphi)-\mathcal E_L(0)=2[F_L(\varphi)-F_L(0)]\ge0$
for the threshold loss of a flat connection with adjoint twists
$\varphi\in[0,2\pi]^3$; by V41, $\ell_L$ vanishes exactly at the eight corners,
which are the eight central orbits $\mathcal C_\sigma$, and near the origin
$\ell_L(\varphi)=\frac2L|\varphi|+O_L(|\varphi|^2)$.

\begin{proposition}[Axis minimality of the loss]
\label{f16v56:axis}
For all $\varphi\in[0,2\pi]^3$,
$\ell_L(\varphi_1,\varphi_2,\varphi_3)\ge\ell_L(\varphi_1,0,0)$, and the same
holds for each coordinate axis.
\end{proposition}
\begin{proof}
V41's subordination formula gives
$F_L(\varphi)=\int_0^\infty[V-H_L(t,\varphi)]e^{-2t}I_0(2t)\,dt/t$ with the twisted
heat trace $H_L(t,\varphi)=\sum_ke^{-t\lambda_L(k,\varphi)}$, which factorizes
over directions, $H_L(t,\varphi)=\prod_\mu h_L(t,\varphi_\mu)$, with the
one-dimensional traces $h_L(t,\varphi)=Le^{-2t}\sum_{n\in\mathbb Z}I_{Ln}(2t)
e^{in\varphi}$, a cosine series with nonnegative coefficients (V41's Bessel
Fourier series). For such a series $h(0)\ge|h(\varphi)|$, hence
$h_L(t,0)^2\ge h_L(t,\varphi_2)h_L(t,\varphi_3)$, and $h_L(t,\varphi_1)\ge0$; so
$H_L(t,\varphi_1,0,0)\ge H_L(t,\varphi)$ for every $t$, and integrating against
the positive kernel gives $F_L(\varphi)\ge F_L(\varphi_1,0,0)$. The verifier
checks the underlying inequality $h(0)^2-h(\varphi)h(\varphi')\ge0$ exactly
for a rational nonnegative cosine series.
\end{proof}

\begin{proposition}[Agmon distance between central orbits on the flat moduli]
\label{f16v56:distance}
Equip $[0,2\pi]^3$ with the degenerate Riemannian metric
$ds^2=\frac L2\ell_L(\varphi)\,|d\varphi|^2$. The distance from the corner $0$
to the corner $2\pi e_1$ is
\begin{equation}
 \boxed{\quad
 S_L=\int_0^{2\pi}\sqrt{\tfrac L2\,\ell_L(\varphi e_1)}\,d\varphi
 =\sqrt L\int_0^{2\pi}\sqrt{F_L(\varphi e_1)-F_L(0)}\,d\varphi,
 \quad}
 \label{f16v56:S}
\end{equation}
attained on the axis, and every corner with first coordinate $2\pi$ is at
distance at least $S_L$ from every corner with first coordinate $0$. Moreover
$S_L\ge4\sqrt{LA_L}$ with V41's $A_L=a_L(1,0,0)>0$, and numerically
$S_L=6.267,6.236,6.230,6.220$ and $4\sqrt{LA_L}=1.317,1.283,1.276,1.274$ at
$L=8,16,32,64$.
\end{proposition}
\begin{proof}
For any Lipschitz path from a corner with $\varphi_1=0$ to one with
$\varphi_1=2\pi$, Proposition~\ref{f16v56:axis} and $|\dot\varphi|\ge|\dot\varphi_1|$
give length $\ge\int\sqrt{\frac L2\ell_L(\varphi_1(t),0,0)}|\dot\varphi_1|dt
\ge\int_0^{2\pi}\sqrt{\frac L2\ell_L(s,0,0)}ds$, with equality on the axis. The
lower bound uses V41's $\ell_L(\varphi)\ge4A_L\sum_\mu(1-\cos\varphi_\mu)$ and
$\int_0^{2\pi}\sqrt{1-\cos s}\,ds=4\sqrt2$:
$S_L\ge\sqrt{2LA_L}\cdot4\sqrt2=4\sqrt{LA_L}$. The numbers are the verifier's
quadratures of the exact sums and Bessel integrals.
\end{proof}

The factor $\frac L2$ is the mass of the twist coordinate: a constant abelian
field $a_\mu=\alpha n$ has twist $\varphi_\mu=2L\alpha$ and moves all $V$
links, so its round-metric length is $\sqrt V|d\alpha|=\sqrt V|d\varphi|/(2L)$,
and the Agmon metric $2\ell_L$ per unit round length becomes
$2\ell_L\cdot V/(4L^2)=\frac L2\ell_L$ per unit twist. Since
$\ell_L=O(1/L)$ by the winding coefficients, $S_L$ has a finite limit; the
barrier height $\max_\varphi L\ell_L(\varphi e_1)\to3.21$ at the twist $\pi$.

\subsection{An Agmon weight on the configuration space}

Let $\vartheta_\mu^{(x)}(U)\in[0,\pi]$ be the angle of the holonomy of $U$
along the $\mu$-line through $x$, $\vartheta^{(x)}_\mu=\arccos\frac12\Tr
H_\mu^{(x)}$, and let $\bar\vartheta_\mu(U)=L^{-2}\sum_{x\perp\mu}\vartheta^{(x)}_\mu(U)$
be its average over the $L^2$ parallel lines. Fix $\delta\in(0,1)$, let
$\hat\psi_\delta$ be the Agmon distance to the set of eight corners for the
metric $(1-\delta)^2\frac L2\ell_L|d\varphi|^2$, mollified at a scale $\eta$
chosen below and multiplied by a smooth cutoff vanishing on fixed small
neighborhoods of the corners, and put
\begin{equation}
 \psi(U)=\hat\psi_\delta\bigl(2\bar\vartheta(U)\bigr).
 \label{f16v56:weight}
\end{equation}

\begin{proposition}[Properties of the weight]
\label{f16v56:weight-properties}
$\psi$ is gauge invariant and Lipschitz on $\mathcal X_L$ with a constant
depending on $L$ and $\delta$ only. At every flat point $p$ with twist
$\varphi(p)$ outside the corner neighborhoods, $\psi$ is smooth, its gradient
lies in the constant-mode directions, and
\begin{equation}
 |\nabla\psi(p)|^2=\frac4L|\nabla_\varphi\hat\psi_\delta(\varphi(p))|^2
 \le2(1-\delta)^2\ell_L(\varphi(p))+o_\eta(1),
 \label{f16v56:agmon-condition}
\end{equation}
where $o_\eta(1)\to0$ as the mollification scale $\eta\to0$, uniformly on the
forbidden region. Between the neighborhoods $\mathcal N_{1,0}$ and
$\mathcal N_{1,\sigma}$ of the trivial and a first-coordinate-twisted central
orbit, $\psi$ changes by at least $(1-\delta)S_L-o(1)$ as the neighborhood
radius $r_0$ and $\eta$ tend to zero.
\end{proposition}
\begin{proof}
Holonomy traces are invariant under periodic gauge transformations, so
$\psi$ is gauge invariant. The angle of an $\SU(2)$ element is its round
distance to the identity, hence $1$-Lipschitz in the element, and the
holonomy along a line of $L$ links is $\sqrt L$-Lipschitz in the product
metric by the Cauchy--Schwarz inequality; averaging over $L^2$ lines and a
second Cauchy--Schwarz give Lipschitz constant $O(L^{-1/2})$ for $\bar\vartheta$,
and $\hat\psi_\delta$ is Lipschitz on the torus, so $\psi$ is Lipschitz. At a
flat point all line holonomies in direction $\mu$ are conjugate with the same
angle $\vartheta_\mu=\varphi_\mu/2$, so $\psi(p)=\hat\psi_\delta(\varphi(p))$. The
derivative of the angle of a product of commuting elements with respect to
each factor is $1$ in the direction of the common axis and $0$ transversally;
thus $\nabla\vartheta^{(x)}_\mu$ has $L$ unit entries along the line, and
$\nabla\bar\vartheta_\mu$ is the vector with the entry $L^{-2}n$ on every
$\mu$-link, $n$ the holonomy axis: it is the constant-mode direction and has
norm $L^{3/2}/L^2=L^{-1/2}$. The chain rule gives
\eqref{f16v56:agmon-condition}, where the mollification changes the gradient
bound by $o_\eta(1)$ because $\ell_L$ is smooth away from the corners.
Smoothness at flat points with some $\varphi_\mu=0$ follows because
$\hat\psi_\delta$ is even in each $\varphi_\mu$ and a smooth even function of
$\vartheta$ is a smooth function of $\cos\vartheta$. The last claim is
Proposition~\ref{f16v56:distance} with the cutoff and mollification errors.
\end{proof}

\subsection{The weighted transfer stays below the band on the forbidden region}

Let $\mathcal O$ be the union of fixed neighborhoods $\mathcal N_{1,\sigma}$
of the central orbits, $Q=1-1_{\mathcal O}$, $W=e^{\sqrt\gamma\psi}$, and
$\lambda_c=\lambda_*-\delta_0$ with $\delta_0$ as in V42.

\begin{theorem}[Sharp weighted barrier]
\label{f16v56:barrier}
There are $\delta_1>0$ and $\gamma_0<\infty$, depending on $L$, $\delta$,
$\mathcal O$ and $\eta$, such that for $\gamma\ge\gamma_0$
\begin{equation}
 \boxed{\quad\|QWT_\gamma W^{-1}Q\|\le\lambda_c-\delta_1.\quad}
 \label{f16v56:weighted-barrier}
\end{equation}
\end{theorem}
\begin{proof}
Let $r=r_\gamma=\gamma^{-1/4-\delta'}$ with a small fixed $\delta'>0$, and cover
$\mathcal X_L$ by a smooth quadratic partition of unity $\{\chi_j\}$ of bounded
overlap with supports of diameter at most $r$; those meeting the
$r/2$-neighborhood of the flat set are centered at flat points and use V41's
normal charts, the others are supported where $S\ge c_Lr^2$. As in V41's
Proposition~\ref{f16v41:IMS}, applied to the weighted kernel, the localization
error is bounded by $\sup_j\|\nabla\chi_j\|_\infty^2$ times the second
displacement moment of the weighted kinetic kernel; the moment is
$O(\gamma^{-1})$ uniformly by V42's exponential-moment bound, since $\psi$ is
Lipschitz with a fixed constant, so the error is $O(\gamma^{-1}r^{-2})
=O(\gamma^{-1/2+2\delta'})\to0$.

On a chart of the second kind, both magnetic factors are at most
$e^{-c_L\gamma r^2/2}=e^{-c_L\gamma^{1/2-2\delta'}/2}$ and the weighted kinetic
kernel has row and column sums bounded by a fixed $K_{L,\psi}$, so the local
weighted norm tends to zero.

On a chart of the first kind, centered at a flat point $p$, use V41's
normal coordinates $x=(x_+,x_0)$ and the two-sided kernel domination
\eqref{f16v41:local-kernel}, whose relative error $\delta_r=O(r)$ now tends
to zero, together with the exact magnetic lower bound
\eqref{f16v41:normal-lower}. Since $\psi$ is smooth near $p$ with Lipschitz
gradient, $\psi(U)-\psi(V)=g\cdot(x-y)+\rho(x,y)$ with $g=\nabla\psi(p)$ and
$|\rho|\le C_Lr|x-y|$. The factor $e^{\sqrt\gamma\rho}$ is bounded by
$e^{C_Lr\sqrt\gamma|x-y|}$; against the Gaussian $e^{-(1-\delta_r)|x-y|^2/(2\varepsilon^2)}$
it costs at most a factor $e^{O(C_L^2r^2)}\to1$ after completing the square, and
enlarges the variance by $O(r^2)$. It remains to bound the norm of the
dominating Gaussian operator conjugated by the linear weight $e^{\sqrt\gamma g\cdot x}$.
That operator is a tensor product of one-dimensional kernels: in each hard
direction a Mehler kernel $m(x)G(x-y)m(y)$ with $m$ the magnetic Gaussian
factor and $G$ the heat kernel of variance $\varepsilon^2/(1-\delta_r)$, and in
each remaining direction the heat kernel alone. Conjugation multiplies the
kernel by $e^{\sqrt\gamma g_i(x_i-y_i)}$, and completing the square gives
$e^{\sqrt\gamma g_i(x_i-y_i)}G(x_i-y_i)=e^{g_i^2/(2(1-\delta_r))}G(x_i-y_i-c_i)$
with the shift $c_i=g_i\varepsilon^2\sqrt\gamma/(1-\delta_r)$. A shifted
convolution factors as $G_c=G^{1/2}\tau_cG^{1/2}$ with $\tau_c$ the unitary
translation, so $\|mG_cm\|=\|(mG^{1/2})\tau_c(G^{1/2}m)\|\le\|mG^{1/2}\|^2=\|mGm\|$:
the shift does not increase the norm. Hence the local weighted norm is at most
\[
 (1+\delta_r)(1-\delta_r)^{-D_L/2}e^{|g|^2/(2(1-\delta_r))}\,\mathfrak h_{\delta_r}(p)
 \longrightarrow\mathfrak h(p)\,e^{|\nabla\psi(p)|^2/2}
 \le\lambda_*e^{-\ell_L(\varphi(p))}e^{(1-\delta)^2\ell_L(\varphi(p))+o_\eta(1)}
 =\lambda_*e^{-(2\delta-\delta^2)\ell_L(\varphi(p))+o_\eta(1)}
\]
by \eqref{f16v56:agmon-condition} and V41's $\mathfrak h=e^{-\mathcal E_L}$,
where $\mathfrak h_{\delta_r}$ is V41's chart-perturbed threshold, which
converges to $\mathfrak h(p)$ as $\delta_r\to0$ uniformly on the compact flat
set. On the forbidden region $\ell_L\ge\ell_{\min}(\mathcal O)>0$, so after
fixing $\eta$ small the right side is at most $\lambda_*(1-2\delta_1)$ for some
$\delta_1>0$, which may be taken smaller than $\delta_0$. Combining the three
kinds of charts with the vanishing localization error proves
\eqref{f16v56:weighted-barrier} for large $\gamma$. This is the sharp form of
V42's conjugation bound: the crude estimate $K_L\alpha$ is replaced by the
exact Gaussian moment generating function of the kinetic kernel, which the
verifier confirms numerically to be $e^{|g|^2/2}(1+O(\gamma^{-1}))$ for the
literal Wilson displacement.
\end{proof}

\subsection{The splitting bound}

\begin{theorem}[Explicit exponent for the charged sector tops]
\label{f16v56:splitting}
For every $\delta>0$ there are $C_{L,\delta}<\infty$ and $\gamma_0$ such that,
for $\gamma\ge\gamma_0$ and every nontrivial character $a$,
\begin{equation}
 \boxed{\quad
 0<1-\frac{\Lambda_{\gamma,a}}{\Lambda_\gamma}\le C_{L,\delta}\,e^{-(1-\delta)S_L\sqrt\gamma},
 \qquad S_L=\sqrt L\int_0^{2\pi}\sqrt{F_L(\varphi e_1)-F_L(0)}\,d\varphi .
 \quad}
 \label{f16v56:splitting-bound}
\end{equation}
The same bound holds for all levels matched by V43's ladders: for every fixed
$j$, $|\lambda_{a,j}-\lambda_{0,j}|\le C_{L,\delta}e^{-(1-\delta)S_L\sqrt\gamma}$.
\end{theorem}
\begin{proof}
With the barrier \eqref{f16v56:weighted-barrier} in place of V42's
$\|A_\alpha\|\le\lambda_*-3\delta$, V42's Sylvester argument applies verbatim
to $X=WQ\mathsf E_\gamma$ and gives $\|e^{\sqrt\gamma\psi}\mathsf E_\gamma\|\le C$;
here $\psi$ vanishes on $\mathcal O$, which is what the argument uses. For
the off-diagonal Feshbach blocks, replace $\psi$ by the analogous weight
$\psi_\tau$ built from the Agmon distance to the single corner $\tau$; the
barrier holds for it by the same proof, and V42's conjugation of the
Feshbach map gives $\|P_\sigma\mathfrak F_\gamma(z)P_\tau\|\le C
e^{-\sqrt\gamma\inf_{\mathcal N_{1,\sigma}}\psi_\tau}$, with the infimum at
least $(1-\delta)S_L-o(1)$ for every pair of distinct corners by
Proposition~\ref{f16v56:distance}, since any two distinct corners differ in
some coordinate. V43's Proposition~\ref{f16v43:Fourier} then gives
$\epsilon_\gamma\le7Ce^{-(1-\delta)S_L\sqrt\gamma+o(\sqrt\gamma)}$, and its
ladder theorem transfers this to every matched level; the top statement is
V42's argument with this $\epsilon_\gamma$. Absorb the $o(\sqrt\gamma)$ and the
neighborhood errors into a smaller $\delta$.
\end{proof}

Three remarks. First, the exponent is the Agmon length of the toron valley in
the transfer's own kinetic metric; in the continuum language of the
small-volume expansion it is the tunneling action between electric flux
sectors through the twisted flat connection, with the $\sqrt\gamma=2/g$
scaling of that literature and the value $S_\infty\approx6.22$, i.e.\ a
suppression $e^{-12.4/g}$ up to prefactors. Second, the exponent is
independent of $L$ in the limit, although V42's constant $\kappa_L$ was not:
the loss per unit twist is $O(1/L)$ and the mass of the twist is $O(L)$, and
these compensate. Third, no lower bound is proved: a matching asymptotic
would need the leading behavior of the vacuum near the sphaleron point
$\varphi=\pi$ and the prefactor of the Feshbach block, which are not
estimated here.

\subsection{Verification and preservation}

The new diagnostic is
\begin{center}\small\path{scripts/verify_ym_f16_v56_agmon_exponent.py}.\end{center}
Its hashed exact content is the product inequality for nonnegative cosine
series at the rational cosine table, the gradient norm $1/L$ of the
line-averaged holonomy angle on the $L=4$ torus, and the Agmon exponent
algebra. Its floating diagnostics, not hashed, are $S_L$, $4\sqrt{LA_L}$, the
barrier height $L\ell_L(\pi)$ and the off-axis monotonicity minimum for
$L=8,16,32,64$, and the moment generating function of the exact Wilson
displacement at $\gamma=50,200,800$. It has 6559 bytes and SHA--256
\begin{center}\scriptsize\ttfamily 408053D768BC84895E395A3B68F9FBE20EECF92D37CDF2737AFAD35C1872764A.\end{center}
Its canonical hashed payload has SHA--256
\begin{center}\scriptsize\ttfamily 956F41F08EAD09C358986FE8333E95491958BDB8DB9E1C2AD196B21943191DDE.\end{center}
Only this terminal insertion and the literal title version differ from V55;
removing them recovers its bytes exactly. The next companion checkpoint and
ten-version review is V60.

\begin{firewall}
V56 proves the axis minimality of V41's loss function, the Agmon distance
$S_L$ between central orbits on the flat moduli, a sharp weighted barrier for
the transfer on the forbidden region, and the explicit upper bound
$e^{-(1-\delta)S_L\sqrt\gamma}$ for all charged sector splittings at fixed
spatial regulator. It does not prove a lower bound or a prefactor for the
splitting, a volume-uniform estimate, a physical-time calibration, or an
interacting continuum mass gap.
\end{firewall}
% END F16 V56 APPEND

% BEGIN F16 V57 APPEND -- 2026-09-18
% Insert before the final document-ending command of cumulative V56.
% This insertion leaves every predecessor body and the native operator unchanged.

\clearpage
\section{V57: Temple enclosure of the model ground level and a certified bracket for the gap coefficient}
\label{f16v57:context}

V49 named as its next model estimate an a priori lower bound on the second
left-invariant level $e_1$ of $H_{\rm mat}$ exceeding the ground Rayleigh
quotient $\rho\approx6.535$, since with it Temple's inequality turns the
one-sided Rayleigh--Ritz bound into a two-sided enclosure of $e_0$, and it
recorded that the available separable comparison gives only $5.84$. That
record was wrong. The intermediate operator in V49's proof of the pairwise
bound, the sum of the three fibrewise inequalities, is the separable
operator
\[
 H_{\rm half}=\sum_\mu\bigl(-\tfrac14\Delta_{q_\mu}+\sqrt2\,|q_\mu|\bigr)
 =2^{-1/3}\sum_\mu(-\Delta_{x_\mu}+|x_\mu|)\quad(q_\mu=2^{-5/6}x_\mu),
\]
whose spectrum is a sum of three Airy-type spectra, and whose second
left-invariant level is $2^{-1/3}(2|a_1|+|a_2|)\approx6.956$, not $5.84$:
the value $5.84$ is the second level of the weighted average of V44's two
inequalities at $t=2/3$, a weaker comparison that V49 already superseded for
the ground. This note certifies the levels of $H_{\rm half}$ by exact
Sturm shooting with zero counting, enumerates its left-invariant and
sector-resolved spectrum, and obtains
\[
 e_0\in[\,\text{6.5322},\ \text{6.5350}\,],\qquad
 \text{0.4011}\le g_*\le\text{3.0162},
\]
the first certified lower bound on the gap coefficient of the quartic model
that is quantitative rather than the qualitative positivity of V44. The
lower bound $g_*\ge$ 0.4011 is far from the floating value $3.01$ and the
reason is identified: the separable comparison discards the transverse
oscillator excitations, and every excited level of $H_{\rm mat}$ lies above
the third left-invariant level of $H_{\rm half}$, so no Temple argument
reaches them. The two-sided identification of the tensor multiplet as the
first excitation remains open.

\subsection{Dependencies}

V44's oscillator floor, V49's pairwise bound
\eqref{f16v49:pairwise-bound} and its shooting certificate
(Proposition~\ref{f16v49:airy}), V51's sector bases and Stein calculus,
V52's exact residual certificates and rational inverse iteration
(Proposition~\ref{f16v52:weinstein}). The model $H_{\rm mat}$ and its
relation to the native transfer are unchanged from V44--V48.

\subsection{Ownership record}

V49 proved $e_0\ge3\cdot2^{-1/3}|a_1|$ by bounding each pair operator
$h_{ij}$ by its ground energy; the same proof bounds $H_{\rm mat}$ from
below by the separable operator $H_{\rm half}$, which is the content of
Proposition~\ref{f16v57:separable} and is not new. V49's remark that the
separable comparisons available have second invariant level $5.84$
referred to the $t=2/3$ average and overlooked $H_{\rm half}$; V51's
statement that a Temple inequality inside a sector needs an a priori bound
on the sector's next level is correct and is what is used here, in the
left-invariant space rather than in the bi-invariant sector. The Sturm certificates with zero counting extend V49's
first-zero certificate.

\subsection{The separable comparison operator and its spectrum}

\begin{proposition}[Separable comparison]
\label{f16v57:separable}
As quadratic forms on $L^2(\mathbb R^9)$,
$H_{\rm mat}\succeq H_{\rm half}=\sum_\mu(-\frac14\Delta_{q_\mu}+\sqrt2|q_\mu|)$.
Both operators commute with the left $\operatorname{SO}(3)$ action and with
the right action of the hyperoctahedral group $B_3$ of signed permutations
of the index $\mu$, but $H_{\rm half}$ does not commute with the right
$O(3)$ action.
\end{proposition}
\begin{proof}
V49's proof of Theorem~\ref{f16v49:pairwise} shows
$h_{12}\succeq-\frac14\Delta_{q_1}+\sqrt2|q_1|$ using the kinetic energy of
$q_2$ only, and by relabeling $h_{13}\succeq-\frac14\Delta_{q_3}+\sqrt2|q_3|$
using the kinetic energy of $q_1$, $h_{23}\succeq-\frac14\Delta_{q_2}+\sqrt2|q_2|$
using that of $q_3$; summing and using $\sum_{i<j}h_{ij}=H_{\rm mat}$ gives
the inequality. The symmetry statements are immediate from the forms of
the two operators.
\end{proof}

Let $a_1<a_2<a_3<\dots$ be the eigenvalues of the reduced radial operator
$-u''+ru$ on $(0,\infty)$ with $u(0)=0$, which are the $s$-wave levels of
$-\Delta+|x|$ on $\mathbb R^3$ (the negatives of the Airy zeros), and let
$p_1$, $d_1$ be the lowest eigenvalues of $-u''+\ell(\ell+1)r^{-2}u+ru$ for
$\ell=1,2$; the levels increase with $\ell$ at fixed radial number because
the centrifugal term is monotone in $\ell$.

\begin{proposition}[Sturm certificates]
\label{f16v57:sturm}
$a_1\ge2.33$, $a_2\ge4.08$, $a_3\ge5.51$, $p_1\ge3.3$, $d_1\ge4.2$.
\end{proposition}
\begin{proof}
For a confining potential the regular solution $u_E$ of
$-u''+\ell(\ell+1)r^{-2}u+ru=Eu$, $u_E\sim r^{\ell+1}$ at $0$, has exactly as
many zeros on $(0,\infty)$ as there are eigenvalues below $E$ (Sturm's
oscillation theorem). The verifier computes $u_E$ as an exact power series
with rational coefficients, $c_n[n(n-1)-\ell(\ell+1)]=c_{n-3}-Ec_{n-2}$,
with a rational tail bound on $[0,R]$, and certifies the sign pattern of
$u_E$: positivity on $(0,r_0]$ from the leading bracket, on each cell of a
grid of step $1/100$ on $[r_0,R]$ either a definite sign from the endpoint
values and the interpolation bound $|u''|\le(\max(R,E)+\ell(\ell+1)r_0^{-2})|u|$,
or a certified sign of $u'$ from the same bound, so that the cell contains
at most one zero; consecutive monotone cells form a transition block, and
the certified signs on either side of each block differ, so each block
contains exactly one zero. Beyond $R>E$ the sign of $u_E''$ equals that of
$u_E$, so if $u_E(R)$ and $u_E'(R)$ have the same sign there is no further
zero. The zero counts $0,1,2,0,0$ for
$(E,\ell)=(2.33,0),(4.08,0),(5.51,0),(3.3,1),(4.2,2)$ give the five bounds.
\end{proof}

\begin{proposition}[Left-invariant and sector-resolved levels of $H_{\rm half}$]
\label{f16v57:levels}
The eigenfunctions of $H_{\rm half}$ are products of three radial functions
with spherical harmonics of degrees $(\ell_1,\ell_2,\ell_3)$ in the three
colour vectors, and a joint eigenspace lies in the left-invariant subspace
exactly when the three angular momenta admit a coupling to total $J=0$,
which requires either $\ell_1=\ell_2=\ell_3=0$ or at least two $\ell_\mu\ge1$.
Consequently, with $c=2^{-1/3}>0.7936$,
\begin{align}
 e_0(H_{\rm half}|_{\rm inv})&=3c\,a_1,\qquad
 e_1(H_{\rm half}|_{\rm inv})=c\min(2a_1+a_2,\ 2p_1+a_1)\ \ge\ 0.7936\cdot8.74=\text{6.936064},
 \label{f16v57:alpha}\\
 e^{(2,\pm)}(H_{\rm mat})&\ \ge\ c\,(2p_1+a_1)\ \ge\ \text{7.086},\qquad
 e^{(0,-)}(H_{\rm mat})\ \ge\ 3c\,p_1\ \ge\ \text{7.856},
 \label{f16v57:sectors}\\
 e_2(H_{\rm half}|_{\rm bi\text{-}inv})&=c\min(2a_1+a_3,\ a_1+2a_2,\ 2d_1+a_1)\ \ge\ \text{8.070}.
 \label{f16v57:third}
\end{align}
\end{proposition}
\begin{proof}
Separability gives the product eigenfunctions and the sum of levels; the
left action rotates all three vectors simultaneously, so the invariant
subspace is the $J=0$ part of the tensor product of the three angular
representations, and $\ell_1\otimes\ell_2\otimes\ell_3$ contains $J=0$
exactly when the triangle inequalities $\ell_\mu\le\ell_\nu+\ell_\kappa$ hold,
which excludes a single nonzero $\ell_\mu$. The ground has all $\ell_\mu=0$
in the lowest radial state. For the second level, either all $\ell_\mu=0$
with one radial excitation, energy $c(2a_1+a_2)$, or at least two factors
with $\ell_\mu\ge1$, each at least $p_1$ by monotonicity in $\ell$ and in
the radial number, energy at least $c(2p_1+a_1)$. Since the two arguments
of the minimum are at least $8.74$ and $8.93$ by
Proposition~\ref{f16v57:sturm}, \eqref{f16v57:alpha} follows with the cube
root bound $0.7936^3<\frac12$ of V49.

For the sector bounds use the right $B_3$ symmetry. A sign flip
$q_\mu\to-q_\mu$ acts on a product eigenfunction by $(-1)^{\ell_\mu}$, and a
permutation permutes the factors. Every $(2,\pm)$ multiplet of $H_{\rm mat}$
contains, under the restriction of the right $O(3)$ to the octahedral group,
a component of type $E$ (functions like $q_1^2-q_2^2$, even under every
flip) and one of type $T_2$ (functions like $q_1\cdot q_2$, odd under two
flips), each multiplied by $\det q$ in the $(2,-)$ case, which is odd under
every flip. Min-max in the subspace of the left-invariant functions with a
prescribed flip parity, which reduces both operators, gives for the $T_2$
component of $(2,+)$ the energy at least $c(2p_1+a_1)$ (two factors of odd
$\ell$), for the $E\otimes\det$ component of $(2,-)$ and for the $(0,-)$
sector, whose functions $\det(q)P(G)$ are odd under every flip, at least
$c\cdot3p_1$. The bi-invariant sector is even under every flip and
symmetric under permutations, so all $\ell_\mu$ are even; its third level
is the smallest of the all-$s$ states with two radial quanta,
$c(2a_1+a_3)$ and $c(a_1+2a_2)$, and the states with two $d$-wave factors,
at least $c(2d_1+a_1)$; the numerical bounds are Proposition~\ref{f16v57:sturm}.
\end{proof}

\subsection{Temple's inequality and the enclosure of the ground level}

\begin{theorem}[Two-sided enclosure of $e_0$]
\label{f16v57:temple}
Let $f=Pe^{-\sigma_1}$ be the refined bi-invariant Rayleigh--Ritz ground
vector of Gram degree 8 (V52's construction), with exact Rayleigh
quotient $\rho=\langle f,H_{\rm mat}f\rangle/\langle f,f\rangle$ and exact
residual variance $s^2=\langle H_{\rm mat}f,H_{\rm mat}f\rangle/\langle f,f\rangle-\rho^2$,
and let $\alpha=0.7936\cdot8.74=$ 6.936064. Then $\rho<\alpha$ and
\begin{equation}
 \boxed{\quad
 \rho-\frac{s^2}{\alpha-\rho}\ \le\ e_0\ \le\ \rho,\qquad
 \text{6.5322}\le e_0\le\text{6.5350},
 \quad}
 \label{f16v57:enclosure}
\end{equation}
an enclosure of width 0.0027. In particular V52's Weinstein interval for
the ground vector contains no eigenvalue other than $e_0$.
\end{theorem}
\begin{proof}
$H_{\rm mat}$ is reduced by the left-invariant subspace $\mathcal H_{\rm inv}$,
which contains $f$ and the ground function, and by min-max applied to
Proposition~\ref{f16v57:separable} on $\mathcal H_{\rm inv}$ the second
eigenvalue of $H_{\rm mat}|_{\rm inv}$ is at least
$e_1(H_{\rm half}|_{\rm inv})\ge\alpha$ by \eqref{f16v57:alpha}. Temple's
inequality for the self-adjoint operator $H_{\rm mat}|_{\rm inv}$ with the
trial vector $f\in D(H_{\rm mat})$ and any $\alpha\le e_1$ with $\rho<\alpha$
gives $e_0\ge\rho-s^2/(\alpha-\rho)$; the upper bound is the Rayleigh
principle. The exact values $\rho=$ 6.53491 and $s=$ 0.0325 of the verifier,
computed with V52's Stein-recursion moments and the rational inverse
iteration, give the numbers. The last claim holds because
$\rho+s<\alpha\le e_1$.
\end{proof}

\begin{theorem}[Certified bracket for the gap coefficient]
\label{f16v57:bracket}
Let $\rho_T=$ 9.5485 be the exact Rayleigh quotient of the refined $(2,+)$
Rayleigh--Ritz ground vector of Gram degree 8. Then
\begin{equation}
 \boxed{\quad
 \alpha-\rho\ \le\ g_*=e_1-e_0\ \le\ \rho_T-\Bigl(\rho-\frac{s^2}{\alpha-\rho}\Bigr),
 \qquad \text{0.4011}\le g_*\le\text{3.0162}.
 \quad}
 \label{f16v57:gap-bracket}
\end{equation}
\end{theorem}
\begin{proof}
$e_1\ge\alpha$ was shown in the previous proof and $e_0\le\rho$; the
$(2,+)$ Rayleigh quotient is an upper bound on the lowest tensor level, which
is an eigenvalue of $H_{\rm mat}|_{\rm inv}$ above $e_0$, so $e_1\le\rho_T$,
and \eqref{f16v57:enclosure} bounds $e_0$ from below.
\end{proof}

\subsection{What the comparison cannot reach}

The certified sector lower bounds are now
\begin{center}\small
\begin{tabular}{@{}lccc@{}}
\toprule
sector, level & certified lower bound & Rayleigh--Ritz upper bound & source of the lower bound\\
\midrule
$(0,+)$, ground & 6.5322 & 6.5350 & Temple, Theorem~\ref{f16v57:temple}\\
$(0,+)$, second & 8.070 & $10.141$ & \eqref{f16v57:third}\\
$(2,+)$, ground & 7.086 & 9.5485 & \eqref{f16v57:sectors}\\
$(0,-)$, ground & 7.856 & $13.949$ & \eqref{f16v57:sectors}\\
$(2,-)$, ground & 7.856 & $17.990$ & \eqref{f16v57:sectors}\\
\bottomrule
\end{tabular}
\end{center}
A Temple bound for the tensor ground would need an a priori bound above
$\rho_T\approx9.55$ on the second $(2,+)$ level, and one for the second
scalar level an a priori bound above $10.14$ on the third; the comparison
operator supplies $7.09$ and $8.07$. The loss is structural: $H_{\rm half}$
replaces each transverse oscillator $-\frac14\Delta_\perp+|q_\mu|^2\rho^2$ by
its ground energy $|q_\mu|$, so its excited levels are radial or angular
excitations of loosely bound Airy factors at spacing about $1.4$, whereas
the excitations of $H_{\rm mat}$ cost about $3$. Retaining the transverse
oscillator quanta in a comparison operator that still has computable
spectrum and commutes with the sector symmetries is the model task that
would close the two-sided bracket; a covariant candidate is the operator
$-\frac12\Delta+2\sigma_2(G)$ with $\sigma_2$ replaced by a lower bound
depending on the two largest eigenvalues of $G$ only, which has separable
structure in singular-value coordinates. This is recorded, not done.

\subsection{Verification and preservation}

The new diagnostic is
\begin{center}\small\path{scripts/verify_ym_f16_v57_temple_enclosure.py}.\end{center}
Its hashed exact content is the five Sturm certificates with zero counts,
the cube-root bound, the value of $\alpha$ and the sector levels of
$H_{\rm half}$, the exact $\rho$, $s^2$ and Temple lower bound for the
refined bi-invariant ground vector, and the exact tensor Rayleigh quotient
with the resulting bracket for $g_*$, all at the requested Gram degree. It
has 11993 bytes and SHA--256
\begin{center}\scriptsize\ttfamily AB0A69D6E2235C7E4E055657515AA64FF576AE1EA420C64BB23251C48096C03B.\end{center}
Its canonical hashed payload at Gram degree 8 has SHA--256
\begin{center}\scriptsize\ttfamily 18AEAD57FA469C8676CF011A754F8E693938778FE84029AB2F061AC515C5FFB9.\end{center}
Only this terminal insertion and the literal title version differ from V56;
removing them recovers its bytes exactly. The next companion checkpoint and
ten-version review is V60.

\begin{firewall}
V57 corrects V49's record on separable comparisons, certifies the
Airy-type levels $a_1,a_2,a_3,p_1,d_1$ from below by exact Sturm shooting,
proves the separable comparison $H_{\rm mat}\succeq H_{\rm half}$ with its
left-invariant and sector-resolved spectrum, and obtains a two-sided
enclosure of the model ground level and a certified bracket for the gap
coefficient $g_*$. It does not certify the tensor multiplet as the first
excitation, does not bound any excited level from below beyond the
comparison values in the table, and does not touch the native transfer,
volume uniformity, or an interacting continuum mass gap.
\end{firewall}
% END F16 V57 APPEND

% BEGIN F16 V58 APPEND -- 2026-09-18
% Insert before the final document-ending command of cumulative V57.
% This insertion leaves every predecessor body and the native operator unchanged.

\clearpage
\section{V58: the continuum limit of the toron loss function and of the splitting exponent}
\label{f16v58:context}

V41's winding theorem writes the twist dependence of the harmonic threshold
as $\ell_L(\varphi)=2\sum_{z\ne0}a_L(z)\{1-\cos(z\cdot\varphi)\}$ with
explicit positive Bessel integrals $a_L(z)$, and states that these
coefficients are fixed-$L$ objects, not uniform constants. This note shows
that they have a clean large-$L$ limit,
\[
 L\,a_L(z)\longrightarrow\frac1{\pi^2|z|^4},
\]
with an $L$-uniform majorant, so that $L\ell_L(\varphi)$ converges
uniformly to the explicit lattice sum
$\ell_\infty(\varphi)=\frac2{\pi^2}\sum_{z\ne0}(1-\cos z\cdot\varphi)/|z|^4$.
This is the continuum one-loop energy of a flat connection with holonomy
twists $\varphi$ relative to the trivial one, in units of the inverse box
size, and it carries the whole small-volume content of the lineage's flat
moduli: the cusp $2|\varphi|$ of V41, the harmonic coefficient
$m_L/L\to-1.204$ of V54, the barrier $L\ell_L(\pi)\to3.210$ and the action
$S_L\to S_\infty$ of V56 are all values or Taylor coefficients of
$\ell_\infty$. The harmonic coefficient is identified with the Epstein zeta
value of the cubic lattice, $4c_2=4Z_3(1)/(3\pi^2)=-1.2042$, and the
center-splitting exponent of V56 acquires a continuum form,
$\exp[-(1-\delta)S_\infty\sqrt\gamma]$ with $S_\infty=$ 6.2319, i.e.\
$\exp[-2(1-\delta)S_\infty/g]$ in terms of the bare coupling $\gamma=4/g^2$.
Nothing here is volume-uniform in the native transfer: the limit is a
statement about the explicit one-loop function, taken after the
fixed-$L$ theorems.

\subsection{Dependencies}

V41's subordination formula \eqref{f16v41:subordination} and winding
theorem \eqref{f16v41:winding-formula}, V54's second-order potential and
V56's Agmon action \eqref{f16v56:S}. Nothing in the native chain is used
or modified.

\subsection{A Chernoff-shifted bound for the lattice heat kernel}

Write $p_x(n)=e^{-x}I_n(x)$ for $x>0$ and $n\in\mathbb Z$; this is the
transition probability of the continuous-time simple random walk on
$\mathbb Z$ with total jump rate one, and
$p_x(n)=\frac1\pi\int_0^\pi e^{-x(1-\cos\theta)}\cos(n\theta)\,d\theta$. Let
$I(v)=v\operatorname{arsinh}v-\sqrt{1+v^2}+1$ be its rate function.

\begin{proposition}[Chernoff-shifted heat kernel bound]
\label{f16v58:chernoff}
For $x>0$ and $n\ge0$,
\begin{equation}
 p_x(n)\le\sqrt{\frac{\pi}{8x}}\;e^{-xI(n/x)},\qquad
 I(v)\ge\min\Bigl(\frac{11}{24}v^2,\ I(1)\,v\Bigr),\quad I(1)=0.4672\ldots
 \label{f16v58:chernoff-bound}
\end{equation}
Consequently, for $L\ge1$, $s>0$ and $z\in\mathbb Z$,
$p_{2L^2s}(Lz)\le\sqrt{\pi/(16L^2s)}\,e^{-(11/48)z^2/s}$.
\end{proposition}
\begin{proof}
$p_x(n)=\frac1{2\pi}\int_{-\pi}^{\pi}e^{x(\cos\theta-1)}e^{-in\theta}d\theta$,
and the integrand is entire and $2\pi$-periodic in $\theta$, so the contour
may be shifted by $-i\lambda$ for any real $\lambda$:
$p_x(n)=e^{-n\lambda}\frac1{2\pi}\int_{-\pi}^{\pi}e^{x(\cos\theta\cosh\lambda+i\sin\theta\sinh\lambda-1)}e^{-in\theta}d\theta$,
whose modulus is at most
$e^{-n\lambda+x(\cosh\lambda-1)}\frac1{2\pi}\int_{-\pi}^{\pi}e^{-x\cosh\lambda(1-\cos\theta)}d\theta$.
With $\lambda=\operatorname{arsinh}(n/x)$ the exponent is $-xI(n/x)$, and
since $1-\cos\theta\ge2\theta^2/\pi^2$ on $[-\pi,\pi]$ the integral is at
most $\frac1{2\pi}\int_{\mathbb R}e^{-2y\theta^2/\pi^2}d\theta=\sqrt{\pi/(8y)}$
with $y=x\cosh\lambda\ge x$. For the rate function, $I(0)=I'(0)=0$ and
$I''(v)=(1+v^2)^{-1/2}\ge1-v^2/2$, so $I(v)\ge v^2/2-v^4/24\ge\frac{11}{24}v^2$
on $[0,1]$; $I$ is convex, so $I(v)\ge I(1)+I'(1)(v-1)\ge I(1)v$ for $v\ge1$
because $I'(1)=\operatorname{arsinh}1>I(1)$. For the consequence put
$x=2L^2s$, $n=L|z|$, $v=|z|/(2Ls)$: if $v\le1$ then
$xI(v)\ge2L^2s\cdot\frac{11}{24}\cdot\frac{z^2}{4L^2s^2}=\frac{11}{48}\frac{z^2}s$;
if $v>1$ then $L>|z|/(2s)$ and
$xI(v)\ge I(1)\,n=I(1)L|z|>I(1)\frac{z^2}{2s}\ge\frac{11}{48}\frac{z^2}s$.
\end{proof}

\subsection{The limit of the winding coefficients}

\begin{theorem}[Continuum limit of the winding coefficients]
\label{f16v58:coefficients}
For every $z\in\mathbb Z^3\setminus\{0\}$ and every $L\ge1$,
\begin{equation}
 L\,a_L(z)\ \le\ \frac{C_0}{|z|^4},\qquad C_0=\frac{\pi^2}{256}\Bigl(\frac{48}{11}\Bigr)^2<0.735,
 \qquad\text{and}\qquad
 \boxed{\ \lim_{L\to\infty}L\,a_L(z)=\frac1{\pi^2|z|^4}.\ }
 \label{f16v58:limit}
\end{equation}
\end{theorem}
\begin{proof}
By \eqref{f16v41:winding-formula} and $e^{-8t}=(e^{-2t})^4$,
$a_L(z)=L^3\int_0^\infty p_{2t}(0)\prod_\mu p_{2t}(Lz_\mu)\,dt/t$; the
substitution $t=L^2s$ gives
$La_L(z)=L^4\int_0^\infty p_{2L^2s}(0)\prod_\mu p_{2L^2s}(Lz_\mu)\,ds/s$.
Proposition~\ref{f16v58:chernoff} bounds the integrand by
$(\pi/(16L^2s))^2e^{-(11/48)|z|^2/s}$, so
$La_L(z)\le\frac{\pi^2}{256}\int_0^\infty s^{-3}e^{-(11/48)|z|^2/s}ds
=\frac{\pi^2}{256}\bigl(\frac{48}{11|z|^2}\bigr)^2$, using
$\int_0^\infty s^{-3}e^{-c/s}ds=c^{-2}$. For the limit, the angular
representation gives
\[
 L\,p_{2L^2s}(Lz)=\frac1\pi\int_0^{\pi L}e^{-2L^2s(1-\cos(\vartheta/L))}\cos(z\vartheta)\,d\vartheta,
\]
whose integrand converges pointwise to $e^{-s\vartheta^2}\cos(z\vartheta)$
and is dominated by $e^{-4s\vartheta^2/\pi^2}$; hence
$Lp_{2L^2s}(Lz)\to\frac1\pi\int_0^\infty e^{-s\vartheta^2}\cos(z\vartheta)d\vartheta
=\frac1{2\sqrt{\pi s}}e^{-z^2/(4s)}$ for every $s>0$ and $z\in\mathbb Z$,
including $z=0$. The product of the four factors converges pointwise in
$s$ to $(16\pi^2s^2)^{-1}e^{-|z|^2/(4s)}$, and the majorant
$\frac{\pi^2}{256}s^{-3}e^{-(11/48)|z|^2/s}$ is integrable, so dominated
convergence gives
$La_L(z)\to(16\pi^2)^{-1}\int_0^\infty s^{-3}e^{-|z|^2/(4s)}ds=(16\pi^2)^{-1}\cdot16/|z|^4$.
\end{proof}

The verifier's quadratures give $La_L(z)\pi^2|z|^4=$ $1.07030, 1.01602, 1.00393, 1.00098$ for $z=e_1$ at $L=8,16,32,64$, with the
expected $O(L^{-2})$ approach to one.

\subsection{Erratum record for V56's numerical values}

V56 quoted $S_L=6.267,6.236,6.230,6.220$ at $L=8,16,32,64$ from a plain
trapezoidal rule on $\sqrt{\Delta F_L}$, which has a square-root cusp at
both endpoints; with the substitution $\varphi=\pi u^2$ and $801$ points
the values are $S_L=$ $6.2710, 6.2408, 6.2340, 6.2324$, decreasing to the limit
$S_\infty=$ 6.2319 with $O(L^{-2})$ differences. V56's theorems define
$S_L$ exactly and are unaffected; only its quoted digits, and the
corresponding values in its verifier's floating output, were low by up to
$0.013$. The barrier values $L\ell_L(\pi e_1)=$ $3.2668, 3.2232, 3.2136, 3.2113$ are as quoted.

\subsection{The continuum toron potential}

\begin{theorem}[Continuum limit of the loss function]
\label{f16v58:loss}
Let $\ell_\infty(\varphi)=\frac2{\pi^2}\sum_{z\in\mathbb Z^3\setminus\{0\}}
\frac{1-\cos(z\cdot\varphi)}{|z|^4}$, an even, $2\pi\mathbb Z^3$-periodic
continuous function vanishing exactly on $(2\pi\mathbb Z)^3$. Then
\begin{equation}
 \boxed{\quad\sup_{\varphi\in\mathbb R^3}\bigl|L\,\ell_L(\varphi)-\ell_\infty(\varphi)\bigr|\longrightarrow0
 \qquad(L\to\infty).\quad}
 \label{f16v58:uniform}
\end{equation}
Consequently V56's Agmon action converges,
$S_L\to S_\infty=\int_0^{2\pi}\sqrt{\ell_\infty(\varphi e_1)/2}\,d\varphi=$ 6.2319,
the barrier converges, $L\ell_L(\pi e_1)\to\ell_\infty(\pi e_1)
=\frac4{\pi^2}\sum_{z_1\ {\rm odd}}|z|^{-4}=$ 3.21048, and V41's strict-loss
constant satisfies $4\sqrt{LA_L}\to4/\pi=1.2732$. The axis minimality of
V56 passes to the limit: $\ell_\infty(\varphi)\ge\ell_\infty(\varphi_1e_1)$.
\end{theorem}
\begin{proof}
$L\ell_L(\varphi)=2\sum_{z\ne0}La_L(z)\{1-\cos z\cdot\varphi\}$ and each term
is bounded by $4C_0/|z|^4$ uniformly in $L$ and $\varphi$, which is summable
over $\mathbb Z^3\setminus\{0\}$; termwise convergence by
Theorem~\ref{f16v58:coefficients} and dominated convergence for the sum give
\eqref{f16v58:uniform}. The remaining claims are the uniform convergence
of $\sqrt{L\ell_L/2}$ and the two special values, with $\cos(z_1\pi)=-1$
exactly for odd $z_1$; the limit of $4\sqrt{LA_L}$ is
$4\sqrt{1/\pi^2}$. Axis minimality is inherited from V56's
Proposition~\ref{f16v56:axis} by passing to the limit at fixed $\varphi$.
\end{proof}

\begin{proposition}[Ewald form, cusp and harmonic coefficient]
\label{f16v58:ewald}
With $\theta(s)=\sum_{n\in\mathbb Z}e^{-sn^2}$ and
$\theta_\varphi(s)=\sum_ne^{-sn^2}\cos n\varphi$,
\begin{equation}
 \ell_\infty(\varphi e_1)=\frac2{\pi^2}\int_0^\infty s\,\theta(s)^2\bigl[\theta(s)-\theta_\varphi(s)\bigr]ds,
 \label{f16v58:ewald-formula}
\end{equation}
and as $\varphi\to0$
\begin{equation}
 \ell_\infty(\varphi e_1)=2|\varphi|+c_2\varphi^2+O(|\varphi|^3),\qquad
 c_2=-\frac1{\sqrt\pi}+\frac2{\pi^2}\Bigl[\int_0^1\sqrt{\pi s}\,\frac{R(s)}{4s}ds
 +\int_0^1\sqrt{\pi s}\,\theta(s)^2\sum_{m\ne0}\Bigl(\frac1{4s}-\frac{\pi^2m^2}{2s^2}\Bigr)e^{-\pi^2m^2/s}ds
 +\int_1^\infty s\,\theta(s)^2\sum_n\frac{n^2}2e^{-sn^2}ds\Bigr],
 \label{f16v58:c2}
\end{equation}
where $R(s)=\theta(s)^2-\pi/s$. Moreover
\begin{equation}
 \boxed{\quad c_2=\frac{Z_3(1)}{3\pi^2}=\text{-0.301047},\quad}
 \label{f16v58:c2-zeta}
\end{equation}
where $Z_3(1)=$ -8.913633 is the analytically continued Epstein zeta value
$\sum'_{z\in\mathbb Z^3}|z|^{-2}$, given for all $\sigma\ne\frac32,0$ by
$\Gamma(\sigma)Z_3(\sigma)=\int_1^\infty t^{\sigma-1}(\theta^3-1)dt
+\int_0^1t^{\sigma-1}(\theta^3-(\pi/t)^{3/2})dt+\frac{\pi^{3/2}}{\sigma-3/2}-\frac1\sigma$.
In particular $4c_2=$ -1.20419, which is the limit of V54's $m_L/L$
($-1.1351,-1.1883,-1.2003,-1.2032$ at $L=8,16,32,64$, with $O(L^{-2})$
differences extrapolating to $-1.2042$).
\end{proposition}
\begin{proof}
$|z|^{-4}=\int_0^\infty se^{-s|z|^2}ds$ and Tonelli give
\eqref{f16v58:ewald-formula}, the sum over $z$ factorizing into
$\theta(s)^2$ for the transverse components and $\theta(s)-\theta_\varphi(s)$
for the first; the integrand is $O(s^{-1/2})$ at $0$ by the Poisson forms
$\theta(s)=\sqrt{\pi/s}\sum_me^{-\pi^2m^2/s}$,
$\theta_\varphi(s)=\sqrt{\pi/s}\sum_me^{-(\varphi+2\pi m)^2/(4s)}$, and
exponentially small at infinity. Split the integral at $s=1$. On $[1,\infty)$
the integrand is an even analytic function of $\varphi$ whose
$\varphi^2$-coefficient is the last integral in \eqref{f16v58:c2}. On
$(0,1]$ insert the Poisson forms: the $m\ne0$ terms
$\sqrt{\pi/s}\,[e^{-\pi^2m^2/s}-e^{-(\varphi+2\pi m)^2/(4s)}]$ are analytic
in $|\varphi|<\pi$ uniformly in $s\in(0,1]$ and vanish to all orders at
$s=0$; their $\varphi^2$-coefficient is the middle integral, from
$e^{-(2\pi m+\varphi)^2/(4s)}=e^{-\pi^2m^2/s}e^{-(\pi m\varphi+\varphi^2/4)/s}$.
The $m=0$ term is $\sqrt{\pi/s}(1-e^{-\varphi^2/(4s)})$, and with
$\theta(s)^2=\pi/s+R(s)$, where $R$ vanishes to all orders at $0$, the
$R$-part contributes $\frac2{\pi^2}\int_0^1\sqrt{\pi s}R(s)(1-e^{-\varphi^2/4s})ds
=\varphi^2\frac2{\pi^2}\int_0^1\sqrt{\pi s}\frac{R(s)}{4s}ds+O(\varphi^4)$. The
explicit part is
$\frac2{\sqrt\pi}\int_0^1s^{-1/2}(1-e^{-\varphi^2/(4s)})ds
=\frac{|\varphi|}{\sqrt\pi}\int_{\varphi^2/4}^\infty u^{-3/2}(1-e^{-u})du$
after $s=\varphi^2/(4u)$, and $\int_0^\infty u^{-3/2}(1-e^{-u})du=2\sqrt\pi$
while $\int_0^{\varphi^2/4}u^{-3/2}(1-e^{-u})du=|\varphi|+O(|\varphi|^3)$;
this gives $2|\varphi|-\varphi^2/\sqrt\pi+O(|\varphi|^3)$ and
\eqref{f16v58:c2}.

For \eqref{f16v58:c2-zeta} write $\theta(s)=\sqrt{\pi/s}\,T(s)$ with
$T(s)=\sum_me^{-\pi^2m^2/s}$, so that $R=(\pi/s)(T^2-1)$,
$\theta^3=(\pi/s)^{3/2}T^3$ and
$\sum_{m\ne0}(\frac1{4s}-\frac{\pi^2m^2}{2s^2})e^{-\pi^2m^2/s}=\frac{T-1}{4s}-\frac{T'}2$.
Call the three integrals in \eqref{f16v58:c2} $A$, $B$, $C$. Then
$6A=\frac32\pi^{3/2}\int_0^1s^{-3/2}(T^2-1)ds$,
$6B=\pi^{3/2}\int_0^1[\frac32s^{-3/2}(T^3-T^2)-s^{-1/2}(T^3)']ds$, and, since
$\sum_n\frac{n^2}2e^{-sn^2}=-\theta'/2$,
$6C=-\int_1^\infty s(\theta^3)'ds=\theta(1)^3-1+\int_1^\infty(\theta^3-1)ds$.
Integrating by parts on $(0,1]$, where $T-1$ vanishes to all orders at
$0$, $\int_0^1s^{-1/2}(T^3)'ds=T(1)^3-1+\frac12\int_0^1s^{-3/2}(T^3-1)ds$,
hence $6(A+B)=\pi^{3/2}\int_0^1s^{-3/2}(T^3-1)ds-\pi^{3/2}(T(1)^3-1)
=\int_0^1(\theta^3-(\pi/s)^{3/2})ds-\theta(1)^3+\pi^{3/2}$. Adding $6C$ and
$-3\pi^{3/2}$,
\[
 3\pi^2c_2=\int_0^1\bigl(\theta^3-(\pi/s)^{3/2}\bigr)ds+\int_1^\infty(\theta^3-1)ds-2\pi^{3/2}-1,
\]
which is $\Gamma(1)Z_3(1)$ by the continuation formula: the Mellin
transform $\Gamma(\sigma)Z_3(\sigma)=\int_0^\infty t^{\sigma-1}(\theta(t)^3-1)dt$
holds for $\Re\sigma>\frac32$, and splitting at $t=1$ with the
subtractions of $(\pi/t)^{3/2}$ and of $1$ continues it meromorphically.
The verifier evaluates the explicit formula \eqref{f16v58:c2}, the
Richardson limit of $(\ell_\infty-2|\varphi|)/\varphi^2$ and $Z_3(1)$
independently.
\end{proof}

\subsection{Consequences and reading}

The three flat-moduli constants of the lineage are now values of one
explicit function: the cusp slope $2$ (V41), the harmonic coefficient
$4c_2=-1.2042$ (V54) and the barrier $\ell_\infty(\pi)=$ 3.21048 and
action $S_\infty=$ 6.2319 (V56). With $\gamma=4/g^2$ the exponent of
Theorem~\ref{f16v56:splitting} reads, for large $L$,
\begin{equation}
 1-\frac{\Lambda_{\gamma,a}}{\Lambda_\gamma}\ \lesssim\ \exp\Bigl[-(1-\delta)\frac{2S_\infty}{g}\Bigr],
 \qquad 2S_\infty=\text{12.464},
 \label{f16v58:physical}
\end{equation}
where the reading is at fixed $L$ with the fixed-$L$ constants of V56 and
the limit $L\to\infty$ is taken only in the exponent's coefficient. In the
continuum language of the small-volume expansion, $\ell_\infty(\varphi)/L_{\rm phys}$
is the one-loop energy of the constant abelian toron with holonomy angles
$\varphi/2$, $\ell_\infty(\pi e_1)/L_{\rm phys}$ the sphaleron-type barrier
between the trivial and the center-twisted flat connections, and
$S_\infty$ the tunneling action between them; the electric flux
energies are exponentially small in $1/g$ with this action, up to the
prefactor that neither V56 nor this note controls. The Epstein zeta value
$Z_3(1)$ is the same lattice constant that governs the finite-volume
zero-point energies of the small-volume expansion, which is the expected
place for it.

\subsection{Verification and preservation}

The new diagnostic is
\begin{center}\small\path{scripts/verify_ym_f16_v58_continuum_toron_potential.py}.\end{center}
Unlike the previous verifiers its hashed content is not exact rational
arithmetic but high-precision floating evaluation (mpmath, thirty digits,
rounded to twelve): $\ell_\infty(\pi)$, $S_\infty$, the Richardson and
explicit-formula values of $c_2$, $Z_3(1)$ and their quotients. Its floating diagnostics, not hashed,
are the table of $La_L(z)\pi^2|z|^4$, grid checks of
\eqref{f16v58:chernoff-bound} and of the rate inequality, the majorant
constant, the cusp ratio $\ell_\infty/(2|\varphi|)$, and a box-sum check of
the barrier. It has 9336 bytes and SHA--256
\begin{center}\scriptsize\ttfamily 7AEBB1BB992F21C75247AE385579A7276FBE47663540CF4BF3F96CFFC7CFF8E6.\end{center}
Its canonical hashed payload has SHA--256
\begin{center}\scriptsize\ttfamily 80BCADACA475872905480FB281A78B152019AA78A73AFB07195E5E0EEA6B434E.\end{center}
Only this terminal insertion and the literal title version differ from V57;
removing them recovers its bytes exactly. The next companion checkpoint and
ten-version review is V60.

\begin{firewall}
V58 proves the Chernoff-shifted bound for the lattice heat kernel, the
limit $La_L(z)\to1/(\pi^2|z|^4)$ with an $L$-uniform majorant, the uniform
convergence $L\ell_L\to\ell_\infty$ with the explicit continuum toron
potential, the convergence of V56's action and barrier, the Ewald form and
the cusp expansion of $\ell_\infty$ with $c_2=Z_3(1)/(3\pi^2)$. Nothing here makes any
native estimate volume-uniform, controls a prefactor, calibrates physical
time, or bears on an interacting continuum mass gap.
\end{firewall}
% END F16 V58 APPEND

% BEGIN F16 V59 APPEND -- 2026-09-18
% Insert before the final document-ending command of cumulative V58.
% This insertion leaves every predecessor body and the native operator unchanged.

\clearpage
\section{V59: gauge-invariant soft coordinates, their fluctuations, and where the volume enters}
\label{f16v59:context}

The fixed-regulator chain V44--V48 reads the soft field off a based-gauge
tube of fixed $\ell^2$ radius, and V52 showed that the hard Gaussian ground
has almost no mass in that tube once $\gamma\lesssim c_LV$: the tube
condition $\gamma\gtrsim L^3$ is a chart artifact. V53 identified the
intrinsic small-volume parameter $z=L\gamma^{-2/3}$ through the Gram
correction of the recovery. This note asks the upstream question directly:
is there a gauge-invariant function of the configuration that reads off the
soft field at the critical scale $h=(\gamma V)^{-1/3}$ with errors
controlled uniformly in $L$, and where does the volume enter? The answer is
quantitative. The line-averaged holonomy trace
$T_\mu(U)=L^{-2}\sum_{\rm lines\parallel\mu}\frac12\Tr H_{\rm line}$ has, on
$U=\operatorname{Exp}(a+u)$ with $a$ constant and $u$ hard, the exact
second-order structure
\[
 T_\mu=1-\tfrac12L^2|a_\mu|^2-\tfrac12L^{-1}\sum_{k_\perp\ne0}|\hat u^{(\mu)}_{(0,k_\perp)}|^2+R,
\]
with no cross term at all: the first variation in $u$ vanishes identically
because $u$ has zero mean, and the hard term is a Parseval sum over the
line-constant modes only. Under the hard Gaussian the hard term has mean
$\beta=c_\beta L/\gamma$ with $c_\beta\to$ 0.3791 and standard deviation
$\gamma^{-1}\sqrt{\frac3{16\pi}\log L+O(1)}$, and the full nonlinear $T_\mu$
has variance at most $1/(12\gamma)$ for every $L\ge8$ by the Gaussian
Poincar\'e inequality. Against the soft signal $\frac12L^2h^2|q_\mu|^2
=\frac12\gamma^{-2/3}|q_\mu|^2$ this gives: the second-order bias is
computable and subtractable, the second-order fluctuation is relatively
$O(\gamma^{-1/3}\sqrt{\log L})$, uniformly in $L$, and the first uncontrolled
correction is the fourth-order bias of relative size $z^2=(L\gamma^{-2/3})^2$.
The intrinsic parameter of V53 therefore reappears, from a purely
gauge-invariant construction, as the order at which the toron coordinate
itself becomes ambiguous. What is proved is the exact algebra, the exact
Gaussian moments of the quadratic part, the Poincar\'e bound and the
$L$-scaling of the constants; the sizes of the cubic and quartic remainders
under the Gaussian are computed heuristically and are stated as such. The
structural lesson is recorded: deterministic charts, in $\ell^2$ or in sup
norm, cannot be volume-uniform for this problem, and every $L$-uniform
statement about the native transfer must be a statement about Gaussian
moments of chaos expansions along lines.

\subsection{Dependencies}

V44's tube and quartic model, V45's hard Gaussian ground $\varphi$ with its
mode variances (Proposition~\ref{f16v45:jets}), V52's second-moment sum
\eqref{f16v52:second-moment}, V53's transverse waves and quaternion tools,
V56's line-averaged holonomy angles. Nothing in the native chain is
modified.

\subsection{The line-averaged trace and its exact second-order structure}

Fix a direction $\mu$. For a configuration $U$ and a transverse position
$x_\perp\in(\mathbb Z/L)^2$, let $H_{x_\perp}(U)$ be the ordered product of
the $L$ links of the line through $x_\perp$ in direction $\mu$ (a
noncontractible loop), and put
\begin{equation}
 T_\mu(U)=\frac1{L^2}\sum_{x_\perp}\tfrac12\Tr H_{x_\perp}(U)\in[-1,1].
 \label{f16v59:T}
\end{equation}
$T_\mu$ is gauge invariant, since a periodic gauge transformation conjugates
each loop holonomy. On a flat constant configuration
$\operatorname{Exp}(\mathsf Ka)$ it equals $\cos(L|a_\mu|)$, and on the
central orbits $\mathcal C_\sigma$ it equals $\sigma_\mu$; V56's angle
coordinate is $\bar\vartheta_\mu=L^{-2}\sum\arccos\frac12\Tr H$, the average
of the angles rather than of the cosines, and the two differ at second order
in the hard field, which is exactly where the difference matters below.

\begin{theorem}[Exact second-order structure]
\label{f16v59:structure}
Let $U=\operatorname{Exp}(x)$ with $x_{\ell}=a_\mu+u_\ell$ on the links in
direction $\mu$, where $a_\mu\in\mathbb R^3$ is constant and $u$ has zero
mean over the lattice in each colour component. Write
$\hat u_k=V^{-1/2}\sum_xu_xe^{-ik\cdot x}$ for the $\mu$-links. Then
\begin{equation}
 \boxed{\quad
 T_\mu(U)=1-\frac12L^2|a_\mu|^2-\frac1{2L}\sum_{k_\perp\ne0}\bigl|\hat u^{(\mu)}_{(0,k_\perp)}\bigr|^2+\bar R,
 \qquad|\bar R|\le\frac C{L^2}\sum_{x_\perp}\Bigl(\sum_{j\in{\rm line}}|a_\mu+u_j|\Bigr)^3,
 \quad}
 \label{f16v59:second-order}
\end{equation}
whenever each line sum $\sum_j|a_\mu+u_j|$ is at most a fixed $c_0>0$;
$C$ and $c_0$ are absolute. In particular the first variation of $T_\mu$ in
$u$ vanishes identically at every constant $a$, not only in expectation.
\end{theorem}
\begin{proof}
For one line write $x_j=a_\mu+u_j$ and $w=\sum_jx_j=La_\mu+\sum_ju_j$. For
$\sum_j|x_j|\le c_0$ the Baker--Campbell--Hausdorff series of the ordered
product converges and $\log\prod_j\operatorname{Exp}(x_j)=w+\frac12\sum_{i<j}[x_i,x_j]+r_3$
with $|r_3|\le C(\sum|x_j|)^3$; since $\frac12\Tr\operatorname{Exp}(y)=\cos|y|
=1-\frac12|y|^2+O(|y|^4)$ and the commutator term is orthogonal to nothing in
particular but is of second order, $\frac12\Tr H=1-\frac12|w|^2+R_{\rm line}$
with $|R_{\rm line}|\le C(\sum_j|x_j|)^3$, the cross term
$w\cdot\sum[x_i,x_j]$ being of third order. Averaging over lines,
$L^{-2}\sum_{x_\perp}|w|^2=L^2|a_\mu|^2+2La_\mu\cdot L^{-2}\sum_{x_\perp}\sum_ju_j
+L^{-2}\sum_{x_\perp}|\sum_ju_j|^2$. The middle sum is the total of $u$ over
all $\mu$-links, which is zero. For the last, $\sum_{x_1}u(x_1,x_\perp)
=V^{-1/2}L\sum_{k_\perp}\hat u_{(0,k_\perp)}e^{ik_\perp\cdot x_\perp}$, and
Parseval on the transverse plane gives $\sum_{x_\perp}|\sum_ju_j|^2
=V^{-1}L^2\cdot L^2\sum_{k_\perp}|\hat u_{(0,k_\perp)}|^2=L\sum_{k_\perp}|\hat u_{(0,k_\perp)}|^2$,
with $\hat u_0=0$. The verifier checks the Parseval identity and the
cancellation exactly on the $L=4$ lattice in Gaussian-rational arithmetic,
and the identity by quaternion finite differences: the remainder scales as
the fourth power for a pure soft field, where $T_\mu=\cos(L|a_\mu|)$ exactly,
and as the third power for a pure hard field.
The first-variation statement is the vanishing of the cross term, which
holds for every $u$ of zero mean, before any expectation is taken; more
directly, $\frac d{d\varepsilon}\frac12\Tr\prod_j\operatorname{Exp}(a_\mu+\varepsilon u_j)|_0
=\frac12\Tr\bigl(\operatorname{Exp}(a_\mu)^{L-1}d\operatorname{Exp}_{a_\mu}(\sum_ju_j)\bigr)$
by cyclicity, a linear function of the line sum, whose average over lines is
zero.
\end{proof}

The hard term involves only the modes with $k_1=0$, constant along the
lines and varying transversally: the modes that a Polyakov line cannot
distinguish from the soft field. Everything else cancels at this order.

\subsection{Hard-Gaussian moments of the quadratic part}

Let $\varphi^2$ be V45's hard Gaussian ground at the trivial background: on
the normalized hard modes $u=\varepsilon^{-1}y$, six real Gaussians of
variance $1/(2\omega_k)$ for each $k\ne0$, $\omega_k=\sinh\xi_k$, so that in
link units each mode has variance $\varepsilon^2/(2\omega_k)=1/(2\gamma\omega_k)$
(V52). For $k=(0,k_\perp)$ the two lattice-transverse polarizations are the
line direction $e_\mu$ and one transverse vector, so
$\hat u^{(\mu)}_{(0,k_\perp)}$ carries three colour amplitudes of one
polarization.

\begin{proposition}[Bias and second-order variance]
\label{f16v59:moments}
With $\beta_{L,\gamma}=E_\varphi\bigl[\frac1{2L}\sum_{k_\perp\ne0}|\hat u^{(\mu)}_{(0,k_\perp)}|^2\bigr]$,
\begin{equation}
 \beta_{L,\gamma}=\frac3{4\gamma L}\sum_{k_\perp\ne0}\frac1{\sinh\xi_{(0,k_\perp)}}
 =c_\beta(L)\frac L\gamma,\qquad c_\beta(L)\to c_\beta=\frac34I_2,\quad
 I_2=\frac1{(2\pi)^2}\int_{[-\pi,\pi]^2}\frac{d^2k}{\sinh\xi(k)}=\text{0.50546},
 \label{f16v59:bias}
\end{equation}
numerically $c_\beta(L)=$ $0.3212, 0.3500, 0.3646, 0.3718, 0.3755, 0.3773$ at $L=8,16,32,64,128,256$, and, with
the mode-variance convention above,
\begin{equation}
 \operatorname{Var}_\varphi\Bigl[\frac1{2L}\sum_{k_\perp\ne0}|\hat u^{(\mu)}_{(0,k_\perp)}|^2\Bigr]
 =\frac3{8\gamma^2L^2}\sum_{k_\perp\ne0}\frac1{\sinh^2\xi_{(0,k_\perp)}}
 =\frac1{\gamma^2}\Bigl(\frac3{16\pi}\log L+O(1)\Bigr),
 \label{f16v59:variance}
\end{equation}
with fitted slope 0.05968 against $3/(16\pi)=0.05968$.
\end{proposition}
\begin{proof}
The mean is the sum of the three colour variances of the $e_\mu$
polarization, $3/(2\gamma\omega_k)$, over the $L^2-1$ transverse momenta,
times $\frac1{2L}$. The Riemann sum $L^{-2}\sum_{k_\perp\ne0}1/\omega_{k_\perp}$
converges to $I_2$ because $1/\omega(k)\sim1/|k|$ is integrable in two
dimensions and the excluded origin contributes $O(L^{-1})$; hence
$\beta_{L,\gamma}=\frac{3}{4\gamma L}(L^2I_2+o(L^2))$. For the variance, the
Gaussian fourth-moment identity $\operatorname{Var}(g^2)=2\sigma^4$ for the
independent real coordinates of the conjugate pairs $\{k,-k\}$ gives
$\frac32\varepsilon^4\sum_{k_\perp\ne0}\omega_k^{-2}$ for the variance of the
sum, and $1/\omega^2\sim1/|k|^2$ makes the Riemann sum logarithmic:
$L^{-2}\sum_{k_\perp\ne0}\omega_{k_\perp}^{-2}=\frac1{4\pi^2}\cdot2\pi\log L+O(1)$,
whence the slope $\frac38\cdot\frac1{2\pi}=\frac3{16\pi}$.
\end{proof}

\subsection{An $L$-uniform fluctuation bound for the full coordinate}

\begin{theorem}[Poincar\'e bound]
\label{f16v59:poincare}
For every $L\ge8$, every constant $a$ and every $\gamma$,
\begin{equation}
 \boxed{\quad\operatorname{Var}_\varphi\bigl[T_\mu(\operatorname{Exp}(\mathsf Ka+\varepsilon u))\bigr]
 \ \le\ \frac{\varepsilon^2}{2\omega_{\min}L}\ \le\ \frac1{12\gamma},\quad}
 \label{f16v59:poincare-bound}
\end{equation}
where $\omega_{\min}=2\sin\frac\pi L\sqrt{1+\sin^2\frac\pi L}$ is V52's
smallest hard frequency.
\end{theorem}
\begin{proof}
The Gaussian Poincar\'e inequality for the product measure $\varphi^2du$
gives $\operatorname{Var}(F)\le E\langle\nabla F,\Sigma\nabla F\rangle\le\|\Sigma\|E|\nabla F|^2$
with $\Sigma$ the mode covariance, $\|\Sigma\|=1/(2\omega_{\min})$. For
$F(u)=T_\mu(\operatorname{Exp}(\mathsf Ka+\varepsilon u))$ the derivative with
respect to the three components of $u_\ell$ on a $\mu$-link $\ell$ is
$\varepsilon L^{-2}$ times the derivative of $\frac12\Tr(A\operatorname{Exp}(x_\ell)B)$
for unit quaternions $A,B$, whose norm is at most $\|d\operatorname{Exp}_{x_\ell}\|\le1$;
links in other directions do not enter. Hence $|\nabla F|^2\le V\varepsilon^2L^{-4}
=\varepsilon^2/L$ pointwise, and restricting the gradient to the hard
subspace only decreases it. Finally $2\omega_{\min}L\ge4L\sin\frac\pi L\ge12$
for $L\ge8$, since $L\sin\frac\pi L$ increases with $L$ and
$8\sin\frac\pi8=3.061$; the verifier checks this through the rational bound
$\sin\frac\pi L\ge\frac\pi L(1-\frac{\pi^2}{6L^2})$.
\end{proof}

\subsection{Signal, bias, fluctuation, and where the volume enters}

At the critical scale $a=h\mathsf Kq$, $L^2|a_\mu|^2=\gamma^{-2/3}|q_\mu|^2$,
and the bias-corrected coordinate
\begin{equation}
 \rho_\mu(U)=2\bigl(1-T_\mu(U)\bigr)-2\beta_{L,\gamma}
 \label{f16v59:rho}
\end{equation}
has, by Theorem~\ref{f16v59:structure} and Proposition~\ref{f16v59:moments},
expectation $\gamma^{-2/3}|q_\mu|^2+E[\bar R]$ and, to second order,
standard deviation $2\gamma^{-1}\sqrt{\frac3{16\pi}\log L+O(1)}$. The
following table records the orders of the terms of $\rho_\mu$ relative to
the signal $\gamma^{-2/3}|q_\mu|^2$; the first three rows are proved, the
last two are the expected sizes of Gaussian chaos terms along a line,
obtained by counting which modes contribute coherently (the line-constant
modes $k_1=0$ give the coherent $L$-growth, they commute among themselves,
and the commutator terms come from $k_1\ne0$ modes whose line sums are
$O(1)$), and are not proved.
\begin{center}\small
\begin{tabular}{@{}llll@{}}
\toprule
term & absolute size & relative to $\gamma^{-2/3}$ & status\\
\midrule
second-order bias $2\beta$ & $2c_\beta L/\gamma$ & $2c_\beta L\gamma^{-1/3}$, subtracted exactly & proved\\
second-order fluctuation & $\gamma^{-1}\sqrt{\frac3{8\pi}\log L}$ & $\gamma^{-1/3}\sqrt{\log L}$ & proved\\
full fluctuation (any order) & $\le(3\gamma)^{-1/2}$ & $\le0.6\,\gamma^{1/6}$ & proved, $L$-uniform, crude\\
cubic chaos, line-averaged & $L^{1/2}\gamma^{-3/2}$ & $L^{1/2}\gamma^{-5/6}$ & expected\\
quartic bias $E|w|^4$ & $(L/\gamma)^2$ & $z^2=(L\gamma^{-2/3})^2$ & expected\\
\bottomrule
\end{tabular}
\end{center}
Deterministic control of $\bar R$ through \eqref{f16v59:second-order} and
$\sum_j|u_j|\le\sqrt L(\sum_j|u_j|^2)^{1/2}$ gives
$E|\bar R|\lesssim(L/\sqrt\gamma)^3$, relatively $L^3\gamma^{-5/6}$, i.e.\
validity only for $L\ll\gamma^{5/18}$: the deterministic remainder loses the
random-walk cancellation along the line, exactly as V44's $\ell^2$ tube
loses it across the lattice. Three conclusions follow.

First, a gauge-invariant soft coordinate exists whose fluctuation is
$L$-uniform at the critical scale, and whose second-order bias is an
explicit function of $L$ and $\gamma$; the $k_1=0$ modes that produce the
bias are precisely V53's resonant modes. Second, the intrinsic parameter
$z=L\gamma^{-2/3}$ of V53 is where the coordinate itself becomes ambiguous:
the quartic bias of the trace of a line of $L$ links is $(L/\gamma)^2$
against a signal $\gamma^{-2/3}$. The small-volume regime is the regime in
which a Polyakov line still measures the toron. Third, and this is the
upstream point, no deterministic chart, whether the $\ell^2$ tube of V44,
a sup-norm tube, or a chart in a fixed global gauge whose gauge parameter
has sup norm of order $(L/\gamma)^{1/2}$, can be volume-uniform; every
$L$-uniform statement about the transfer must control Gaussian chaos
expansions of products along lines, in which the cubic and quartic terms
are bounded through the decay of the hard covariance along a line, of order
$\gamma^{-1}/d^2$ at distance $d$. That is the technology the chain
V44--V48 lacks, and it is the reason its constants are not uniform. It is
recorded as the native task, with the line-chaos bound for
$\bar R$ as its first lemma.

\subsection{Verification and preservation}

The new diagnostic is
\begin{center}\small\path{scripts/verify_ym_f16_v59_gauge_invariant_coordinates.py}.\end{center}
Its hashed exact content is the transverse Parseval identity and the
first-variation cancellation on the $L=4$ lattice in Gaussian-rational
arithmetic, and the rational bound $2\omega_{\min}L\ge12$ for
$L=8,16,32,64$. Its floating diagnostics, not hashed, are quaternion
finite-difference checks of the first-order cancellation for V53's
transverse waves and of the second-order identity with its remainder scalings,
the tables of $c_\beta(L)$ and of the variance coefficient with the fitted
logarithmic slope, and the integral $I_2$. It has 9733 bytes and
SHA--256
\begin{center}\scriptsize\ttfamily 5A872EB6A11896C092F3EC2C1DA57C13F585C071D9047146550B7774AB17FBAF.\end{center}
Its canonical hashed payload has SHA--256
\begin{center}\scriptsize\ttfamily D9BC6A6596307549CE8A1E7E6C7FA978FBAD1A5D1CDBF7734659FB6F221294C3.\end{center}
Only this terminal insertion and the literal title version differ from V58;
removing them recovers its bytes exactly. V60 is the companion checkpoint
and the ten-version review of V50--V59.

\begin{firewall}
V59 proves the exact second-order structure of the line-averaged holonomy
trace with the identical vanishing of its first variation, the exact
hard-Gaussian mean and second-order variance of its quadratic part with
their $L$-scalings, and an $L$-uniform Poincar\'e bound on its full
variance. The sizes of the cubic and quartic Gaussian terms are counted, not
proved, and the conclusion that the intrinsic parameter $z=L\gamma^{-2/3}$
governs the coordinate is drawn from them. Nothing here modifies the native
chain, removes the tube condition from any theorem, or bears on an
interacting continuum mass gap.
\end{firewall}
% END F16 V59 APPEND

% BEGIN F16 V60 APPEND -- 2026-09-18
% Insert before the final document-ending command of cumulative V59.
% This insertion leaves every predecessor body and the native operator unchanged.

\clearpage
\section{V60: companion checkpoint, ten-version review, and the consolidated fixed-regulator ledger}
\label{f16v60:context}

V60 is the scheduled companion checkpoint and the ten-version review of
V50--V59. It adds no new theorem. It records the rereading of the companion
programmes, reviews the ten versions, consolidates every certified constant
of the quartic model and every small-volume constant of the flat moduli
into one ledger with an integrity check of all verifiers and artifacts,
corrects two forecasts made in V50, and fixes the order of attack for
V61--V69.

\subsection{V60 companion checkpoint}

A fresh inventory of \path{latex_notes}, hashed by the ledger verifier,
finds the five companion inputs unchanged in size and SHA--256 since V50:
SM continuum V72 (607372; f44f9745), NS stability V26 (278283; 82c887f8),
shell mixing V16 (623261; bf138e63), parallel YM V5 (54174; 306f1bb8) and
GRH cofinal visibility V3 (23851; 4dd31580); the archived $V100$ and
$V157$ files of the other programmes are frozen. The only file changed in
that directory since V50 is this lineage. Their latest conclusions were
reread against the V50 table: no companion supplies a volume-uniform native
spectral comparison, a physical-time calibration, a certified lower bound
for the excited sectors of the quartic model, or a probabilistic chart
estimate of the kind V59 identifies as necessary. The next companion
checkpoint is V65, the next ten-version review V70.

\subsection{Ten-version review, V50--V59}

\begin{center}\small
\begin{tabular}{@{}p{0.07\textwidth}p{0.88\textwidth}@{}}
\toprule
Version & Established at fixed spatial regulator (scope limits in each firewall)\\
\midrule
V50 & Checkpoint and V40--V49 review; correction of V49's normalization
remark; closed-form hard band gap and the volume-uniform ratio
$h/(1-r_+)<0.41\gamma^{-1/3}$.\\
V51 & Polar coordinates and centrifugal separation of every non-scalar
sector; exact Gram-polynomial calculus; sector-resolved certificates; the
$\ell=4$ level $55/4$.\\
V52 & Ownership correction; exact Weinstein intervals; the hard Gaussian
ground has mass $\le e^{-\omega_{\min}(c_LV/2-\gamma r^2)}$ in the tube, so
the tube condition is $\gamma\gtrsim c_LV$.\\
V53 & Explicit first variation $H_1(q)=-4\sum_\mu(q_\mu\cdot J)\otimes D_\mu$;
explicit corrector with $\|k(q)\|^2=K_L|q|^2$; the Gram correction
$\kappa_L\gamma^{-2/3}L$, small exactly for $L\ll\gamma^{2/3}$.\\
V54 & Exact hard norm along constant abelian backgrounds; isotropic
second-order potential $m_L|a|^2$ with $m_L/L\to-1.204$; cancellation of the
extensive pieces; first correction $\gamma^{-1/3}(m_L/L)|q|^2$.\\
V55 & Checkpoint and V50--V54 review; Agmon decay of model eigenfunctions;
first-order shift theorem; exact $\langle|q|^2\rangle$ in the sector
states; $g_*(\gamma)=3.014-1.576\gamma^{-1/3}+\cdots$.\\
V56 & Axis minimality of the loss; Agmon distance $S_L$ on the flat moduli;
sharp weighted barrier; all charged splittings $\le C_{L,\delta}e^{-(1-\delta)S_L\sqrt\gamma}$.\\
V57 & Correction of V49's separable-comparison record; Sturm certificates
with zero counting; $e_1\ge\alpha=6.936$; Temple enclosure of $e_0$;
certified $g_*\in[0.4011,3.0162]$; sector lower bounds via flip parities.\\
V58 & $La_L(z)\to1/(\pi^2|z|^4)$ with an $L$-uniform majorant; uniform
convergence $L\ell_L\to\ell_\infty$; $S_L\to S_\infty=6.2319$; Ewald form;
$c_2=Z_3(1)/(3\pi^2)$ proved; erratum for V56's quoted $S_L$.\\
V59 & Gauge-invariant Polyakov-trace coordinate with identically vanishing
first variation; exact bias $c_\beta L/\gamma$ and second-order variance;
$L$-uniform Poincar\'e bound; $z=L\gamma^{-2/3}$ as the quartic bias;
the structural lesson on deterministic charts.\\
\bottomrule
\end{tabular}
\end{center}

\subsection{Consolidated ledger}

The following statement collects what the lineage has proved about the
fixed-regulator critical branch, with the sharpest certified constants in
force after V59; every number is read by the ledger verifier from the
artifact that certifies it, and the two brackets below combine V57's exact
rationals with V51's degree-ten certificates.

\begin{theorem}[State of the fixed-regulator theory after V59]
\label{f16v60:state}
Let $L=2m\ge8$ be fixed and $\gamma\to\infty$, $h=(\gamma V)^{-1/3}$.
\begin{enumerate}
\item (V41--V48) The eight center sectors have tops $\Lambda_{\gamma,a}$
within $C_{L,\delta}e^{-(1-\delta)S_L\sqrt\gamma}$ of each other (V56), with
$S_L\ge4\sqrt{LA_L}$, $S_8=6.2710$, $S_{64}=6.2324$, $S_L\to S_\infty=6.2319$
(V58), and in each sector the ordered levels satisfy
$(1-\lambda_{a,j}/\lambda_*)/h\to e_{j-1}$, the invariant spectrum of
$H_{\rm mat}=-\frac12\Delta_q+2\sigma_2(G)$ on $L^2(\mathbb R^9)$.
\item (V49, V51, V57) The model levels satisfy, with exact certificates,
\[
 \text{6.5322}\le e_0\le\text{6.5350},\qquad e_1\ge\alpha=\text{6.936064},\qquad
 e^{(2,+)}\le\text{9.5478},\quad e^{(0,+)}_2\le\text{10.1385},\quad
 e^{(0,-)}\le\text{13.9483},\quad e^{(2,-)}\le\text{17.9878},
\]
$e^{(2,\pm)}\ge7.086$, $e^{(0,-)}\ge7.856$, $e^{(0,+)}_2\ge8.070$, and
therefore
\begin{equation}
 \boxed{\quad\text{0.4011}\ \le\ g_*=e_1-e_0\ \le\ \text{3.0156},\quad}
 \label{f16v60:bracket}
\end{equation}
with the floating value $g_*=3.0129$ from the degree-ten tensor level
$9.5478$ and ground $6.5349$.
\item (V53--V55) The first correction to the model is
$H_{\rm mat}+\gamma^{-1/3}(m_L/L)|q|^2$ with $m_L/L=-1.1351,\dots,-1.2032$
at $L=8,\dots,64$ and the numerical limit $4c_2=4Z_3(1)/(3\pi^2)=-1.2042$;
its first-order effect lowers the tensor gap to
$3.014-1.576\gamma^{-1/3}$. The formula is not inserted into the native
asymptotics, and the recovery's Gram correction $\kappa_L\gamma^{-2/3}L$
identifies the intrinsic small-volume parameter $z=L\gamma^{-2/3}$.
\item (V58, V59) The flat moduli carry the explicit continuum toron
potential $\ell_\infty(\varphi)=\frac2{\pi^2}\sum_{z\ne0}(1-\cos z\cdot\varphi)/|z|^4
=2|\varphi|+c_2\varphi^2+\dots$, and the Polyakov-trace coordinate reads the
soft field with an $L$-uniform fluctuation bound and a bias
$c_\beta L/\gamma$, $c_\beta=0.3791$, whose first uncontrolled correction is
of relative order $z^2$.
\end{enumerate}
None of these statements is uniform in $L$ except where stated (the hard
band ratio of V50, the exponent $S_L$ of V56 and V58, the Poincar\'e bound
of V59), and none concerns the interacting continuum limit.
\end{theorem}

The small-volume constants collected by the ledger are: per-site
$-\log\lambda_*\to5.9991$ (V50); $c_L\to0.929$ (V52); $K_L/V\to0.148$ (V53);
$m_L/L\to-1.2042$ (V54, V58); $S_L\to6.2319$ and $L\ell_L(\pi)\to3.2105$
(V56, V58); $c_2=-0.301047$, $Z_3(1)=-8.913633$ (V58); $c_\beta=0.3791$,
$I_2=0.50546$ (V59).

\subsection{Two corrected forecasts}

V50 forecast (i) a sector-wise certified lower bound on $e_1$ and (ii) a
tunneling exponent scaling as $h^{-3/2}=\sqrt{\gamma V}$. Both are corrected
by the record. For (i), the lower bound $e_1\ge6.936$ was obtained in V57 in
the whole left-invariant space, not sector-wise, and it is the best
available: sector-wise bounds of the needed size are out of reach of
separable comparisons. For (ii), the exponent is $S_L\sqrt\gamma$ with $S_L$
bounded and convergent in $L$: the Agmon metric on the flat moduli has
loss $O(1/L)$ per unit twist and mass $O(L)$ per unit twist, and the two
compensate exactly (V56, V58). The factor $\sqrt V$ forecast by V50 came from
treating the toron valley with the volume-extensive round metric; it is
absent. The physical reading is $e^{-2S_\infty/g}$, $2S_\infty=12.46$.

\subsection{Order of attack for V61--V69}

Model tasks. (M1) A certified lower bound for the tensor sector above
$9.5$, which would close the two-sided identification of the first
excitation; the separable route is exhausted (V57), so this needs either a
comparison operator retaining transverse oscillator quanta with computable
spectrum, or a certified three-dimensional eigenvalue solver in the reduced
singular-value coordinates of V51. (M2) The next order of V45--V48 to insert
V54's correction into the native asymptotics.

Native tasks, in the order fixed by V59. (N1) The line-chaos lemma: bound
the cubic and quartic Gaussian remainders of the Polyakov-line trace using
the $\gamma^{-1}/d^2$ decay of the hard covariance along a line, so that
V59's coordinate is accurate for $L\ll\gamma^{2/3}$ rather than
$L\ll\gamma^{5/18}$. (N2) Replace V41's deterministic $\ell^2$ chart in the
kernel comparison by a Gaussian-concentration version, with the tube of V44
replaced by a sup-norm event of probability $1-Ve^{-c\gamma r^2}$; this is
the step that would make the constants of V44--V48 uniform for $L\le\gamma^{2/3-\delta}$
and turn the fixed-$L$ theorem into a joint-limit theorem in the small-volume
regime. (N3) Physical-time calibration of the step. (N4) Volume-uniform
control beyond the small-volume regime, which is the mass gap proper and for
which no companion or version supplies a method. V61 begins with N1.

\subsection{Verification and preservation}

The new diagnostic is
\begin{center}\small\path{scripts/verify_ym_f16_v60_consolidated_ledger.py}.\end{center}
Its hashed content is the integrity table (the SHA--256 of every verifier
V49--V59 against the source hash recorded in its artifacts, all matching),
the consolidated exact rationals of item 2 above, and the companion
inventory. Its floating diagnostics are the small-volume constants. It has
7069 bytes and SHA--256
\begin{center}\scriptsize\ttfamily 011E87F03FDE9DF2BC2C9E33CAB3D2527BB992F58389FEACB93B34BDB9FF559A.\end{center}
Its canonical hashed payload has SHA--256
\begin{center}\scriptsize\ttfamily C5A4DCE1ADC3EBD7253CCC0B0C51DBFDFAC1E2091D701F14C83C5128C322CF65.\end{center}
Only this terminal insertion and the literal title version differ from V59;
removing them recovers its bytes exactly.

\begin{firewall}
V60 records an unchanged companion checkpoint, a ten-version review, a
consolidated ledger with an integrity check, two corrected forecasts and an
order of attack. It proves nothing new, and the consolidated theorem is a
restatement of results proved in V41--V59 with their own scope limits. It
does not establish volume-uniform control, a physical-time calibration, a
two-sided tensor level, or an interacting continuum mass gap.
\end{firewall}
% END F16 V60 APPEND

% BEGIN F16 V61 APPEND -- 2026-09-18
% Insert before the final document-ending command of cumulative V60.
% This insertion leaves every predecessor body and the native operator unchanged.

\clearpage
\section{V61: the line-chaos lemma and the bias of the Polyakov coordinate in the small-volume regime}
\label{f16v61:context}

V60 placed first among the native tasks the control of the cubic and
quartic Gaussian remainders of the Polyakov-line trace, which V59 could
bound only deterministically and hence only for $L\ll\gamma^{5/18}$. This
note proves the lemma. The ordered product of the link exponentials along
a line is, exactly, a multi-index power series in the link fields, because
$\operatorname{Exp}(u)=\sum_nu^n/n!$ for a pure imaginary quaternion; its
real part is a sum of homogeneous polynomial chaoses $T_N$ of the Gaussian
hard field, and Wick's theorem with three elementary counts gives
\[
 \|T_N\|_{L^2}\le\frac{(6L\|C\|_1)^{N/2}}{\sqrt{N!}},
\]
where $\|C\|_1$ is the $\ell^1$ norm of the hard covariance along the line,
a constant of size $0.25/\gamma$ uniformly in $L$. The factorial makes the
chaos series converge for every $L$ and every $\gamma$, so the line trace
is an entire function of the field in the $L^2$ sense, and the deterministic
condition $\sum_j|u_j|\le c_0$ of V59 disappears. With a constant soft
background $a$ the bound becomes $(\sqrt3L|a|+\sqrt{6L\|C\|_1})^N/\sqrt{N!}$,
and colour isotropy makes every odd chaos have zero mean, so that the bias
of V59's coordinate, corrected at second order, is controlled at fourth
order: relative to the signal $\gamma^{-2/3}|q_\mu|^2$ the uncontrolled
part is at most $72\gamma^{-2/3}|q_\mu|^2+19z^2$, $z=L\gamma^{-2/3}$. This is
the first estimate of the lineage that holds, with explicit constants,
throughout the small-volume regime $L\ll\gamma^{2/3}$, uniformly in $L$.
The per-line fluctuation bound is of order $(A+B)^3$, valid in the same
regime but sharp only for $L\ll\gamma^{5/9}$; the gain from averaging over
lines is not proved here.

\subsection{Dependencies}

V59's definition of $T_\mu$ and its second-order structure, V45's hard
Gaussian ground and V52's mode variances, V53's transverse waves for the
covariance. Nothing in the native chain is modified.

\subsection{The chaos expansion of a line trace}

Let $u_1,\dots,u_L\in\mathbb R^3$ be the fields on the links of one line,
identified with pure imaginary quaternions $u_j=\sum_au_j^ae_a$, so that
$u_j^2=-|u_j|^2$ and $\operatorname{Exp}(u_j)=\cos|u_j|+\frac{\sin|u_j|}{|u_j|}u_j
=\sum_{n\ge0}u_j^n/n!$ exactly. Hence, for the ordered product,
\begin{equation}
 \tau(u)=\tfrac12\Tr\prod_{j=1}^L\operatorname{Exp}(u_j)
 =\sum_{N\ge0}T_N(u),\qquad
 T_N(u)=\sum_{n\in\mathbb N^L,\ |n|=N}\operatorname{Re}\prod_{j=1}^L\frac{u_j^{n_j}}{n_j!},
 \label{f16v61:chaos}
\end{equation}
where the product is ordered by $j$ and $T_N$ is a homogeneous polynomial of
degree $N$ in the $3L$ components. The series converges absolutely for
every $u$. Directly,
\begin{equation}
 T_0=1,\qquad T_1=0,\qquad T_2=-\tfrac12\Bigl|\sum_ju_j\Bigr|^2,\qquad
 T_3=-\sum_{i<j<k}\det(u_i,u_j,u_k),
 \label{f16v61:low}
\end{equation}
since $\operatorname{Re}u_j=0$, $\operatorname{Re}(u_iu_j)=-u_i\cdot u_j$,
$\operatorname{Re}(u_i^2u_j)=\operatorname{Re}(u_iu_j^2)=\operatorname{Re}u_j^3=0$ and
$\operatorname{Re}(u_iu_ju_k)=-u_i\cdot(u_j\times u_k)$. The second-order
statement is V59's Theorem~\ref{f16v59:structure} without remainder; the
verifier checks it exactly on five links.

\begin{theorem}[Line-chaos bound]
\label{f16v61:bound}
Let $u$ be a centred Gaussian field on the line with colour-diagonal
covariance $E[u_i^au_j^b]=\delta^{ab}C(i-j)$, $C$ a real function on
$\mathbb Z/L$, and $\|C\|_1=\sum_{d\in\mathbb Z/L}|C(d)|$. Then for every $N$
\begin{equation}
 \boxed{\quad E\bigl[T_N(u)^2\bigr]\ \le\ \frac{(6L\|C\|_1)^N}{N!}.\quad}
 \label{f16v61:TN}
\end{equation}
Consequently \eqref{f16v61:chaos} converges in $L^2$ for every $L$ and every
covariance, $E[T_N]=0$ for odd $N$, and with $x=6L\|C\|_1\le1$
\begin{equation}
 \Bigl\|\tau-1+\tfrac12\Bigl|\sum_ju_j\Bigr|^2\Bigr\|_{L^2}\le0.77\,x^{3/2},\qquad
 \Bigl|E[\tau]-1+\tfrac12E\Bigl|\sum_ju_j\Bigr|^2\Bigr|\le0.25\,x^2.
 \label{f16v61:remainders}
\end{equation}
\end{theorem}
\begin{proof}
Three counts. (i) For colour indices $a_1,\dots,a_N$, the product
$e_{a_1}\cdots e_{a_N}$ lies in the quaternion group $\{\pm1,\pm e_1,\pm e_2,\pm e_3\}$,
so its real part is $0$ or $\pm1$. (ii) $(2N-1)!!\le2^NN!$. (iii) For any
function $f$ of $N$ positions, $\sum_{|n|=N}f(n)/n!=\frac1{N!}\sum_{x\in[L]^N}f(\operatorname{sort}x)$,
because a multi-index $n$ is the multiplicity vector of exactly $N!/n!$
sequences. Write the ordered product for the multi-index $n$ as the product
over the sorted sequence, expand in colour components,
$\operatorname{Re}\prod_su_{x_s}=\sum_{a}\epsilon(a)\prod_su^{a_s}_{x_s}$ with
$|\epsilon(a)|\le1$ by (i), and use (iii) for $T_N$ and for a second copy.
Wick's theorem for the $2N$ Gaussian variables, coincidences included,
gives $E[\prod_su^{a_s}_{x_s}\prod_su^{a'_s}_{x'_s}]=\sum_\pi\prod_{(p,q)\in\pi}\delta^{a_pa_q}C(x_p-x_q)$
over the pairings $\pi$ of the $2N$ slots. Summing the colour indices, each
pair contributes a factor $3$ at most; summing the positions, each pair
contributes $\sum_{x_p,x_q\in[L]}|C(x_p-x_q)|=L\|C\|_1$; there are $(2N-1)!!$
pairings and a factor $(N!)^{-2}$ from (iii). Hence
$E[T_N^2]\le(N!)^{-2}(2N-1)!!\,(3L\|C\|_1)^N\le(6L\|C\|_1)^N/N!$ by (ii).
Odd chaoses have zero mean because odd moments of a centred Gaussian vanish.
For \eqref{f16v61:remainders}, $\|\sum_{N\ge3}T_N\|\le\sum_{N\ge3}x^{N/2}/\sqrt{N!}
\le x^{3/2}\sum_{N\ge3}(N!)^{-1/2}$ and the mean error is bounded by the even
tail from $N=4$; the verifier certifies $\sum_{N\ge3}(N!)^{-1/2}\le0.77$ and
$\sum_{N\ge4\ {\rm even}}(N!)^{-1/2}\le0.25$ with rational upper bounds.
\end{proof}

\begin{proposition}[Shifted field]
\label{f16v61:shifted}
Let $y_j=a+u_j$ with a constant $a\in\mathbb R^3$ and $u$ as above, and
$A=\sqrt3L|a|$, $B=\sqrt{6L\|C\|_1}$. Then
\begin{equation}
 E\bigl[T_N(y)^2\bigr]\le\frac{(A+B)^{2N}}{N!},\qquad E[T_N(y)]=0\ \text{for odd }N,
 \label{f16v61:shifted-bound}
\end{equation}
and $E[T_2(y)]=-\frac12L^2|a|^2-\frac12E|\sum_ju_j|^2$.
\end{proposition}
\begin{proof}
Expand each factor $y_{x_s}=a+u_{x_s}$; a term with the deterministic $a$
in $p$ slots and Gaussian factors in the others has, after summing the
colour index of each $a$-slot, a factor at most $\sqrt3|a|$ per $a$-slot,
and its positions are free, contributing $L$ each. Two copies with $p$ and
$p'$ $a$-slots leave $2N-p-p'=2M$ Gaussian slots, paired as in the theorem:
$E[T_N(y)^2]\le(N!)^{-2}\sum_{p,p'}\binom Np\binom N{p'}(\sqrt3L|a|)^{p+p'}(2M-1)!!\,(3L\|C\|_1)^M$,
and with $(2M-1)!!(3L\|C\|_1)^M\le M!\,B^{2M}\le N!\,B^{2N-p-p'}$ the double
sum is $(N!)^{-1}(A+B)^{2N}$. For the odd means, a term of $T_N(y)$ with $p$
copies of $a$ and $N-p$ Gaussian factors has, after Wick contraction, the
form of a quaternion word in $a$ alone in which contracted pairs act by
$X\mapsto\sum_be_bXe_b=X-4\operatorname{Re}X$, a map preserving real and
imaginary parts; since $a^2=-|a|^2$, every such word is a real multiple of
$1$ or of $a$ according to the parity of $p$, and $N-p$ even with $N$ odd
forces $p$ odd, hence zero real part. The second-order mean is
\eqref{f16v61:low} with the cross term $E[a\cdot\sum u_j]=0$.
\end{proof}

\subsection{The bias of the Polyakov coordinate in the small-volume regime}

Let $C_L(d)$ be the covariance, in units of $1/\gamma$, of the
line-direction colour components of V45's hard Gaussian at lattice
distance $d$ along a line: $C_L(d)=V^{-1}\sum_{k\ne0}\cos(k_1d)P_{11}(k)/(2\omega_k)$
with $P_{11}=1-|e^{-ik_1}-1|^2/\hat k^2$ the transverse weight of the line
polarization. Numerically $C_L(0)=0.103$, $d^2C_L(d)\approx0.05$ for
$1\le d\le L/4$, and $\|C_L\|_1=$ $0.2141, 0.2334, 0.2430, 0.2479, 0.2503, 0.2515$ at $L=8,16,32,64,128,256$, so
$x_L=6L\|C_L\|_1/\gamma\le1.52\,L/\gamma$; the decay $d^{-2}$ comes from the
Lipschitz cusp at $k_1=0$ of the line spectral density
$\sum_{k_\perp}P_{11}/(2\omega)$, and the uniform bound on $\|C_L\|_1$ is
numerical, not proved.

\begin{theorem}[Bias and fluctuation of $T_\mu$ in the small-volume regime]
\label{f16v61:bias}
At the critical scale $a=h\mathsf Kq$, with $A=\sqrt3\gamma^{-1/3}|q_\mu|$
and $B=\sqrt{x_L}$, whenever $A+B\le1$,
\begin{equation}
 \boxed{\quad
 \Bigl|E_\varphi[T_\mu]-1+\tfrac12\gamma^{-2/3}|q_\mu|^2+\beta_{L,\gamma}\Bigr|
 \le0.25(A+B)^4\le72\,\gamma^{-4/3}|q_\mu|^4+19\,\frac{L^2}{\gamma^2},
 \quad}
 \label{f16v61:bias-bound}
\end{equation}
with V59's bias $\beta_{L,\gamma}=c_\beta(L)L/\gamma$, and
\begin{equation}
 \operatorname{Var}_\varphi[T_\mu]^{1/2}\le\operatorname{Var}_\varphi[T_2(y)]^{1/2}+0.77(A+B)^3,
 \qquad\operatorname{Var}[T_2(y)]^{1/2}=\gamma^{-1}\sqrt{\tfrac3{16\pi}\log L+O(1)}.
 \label{f16v61:fluct}
\end{equation}
Relative to the signal $\gamma^{-2/3}|q_\mu|^2$, the bias error is
$72\gamma^{-2/3}|q_\mu|^2+19z^2$, $z=L\gamma^{-2/3}$.
\end{theorem}
\begin{proof}
$T_\mu$ is the average over the $L^2$ lines of $\tau(y)$ with $y=a_\mu+u$,
each line having the same law, so $E[T_\mu]=E[\tau(y)]$ and, by convexity,
$\operatorname{Var}[T_\mu]\le\operatorname{Var}[\tau(y)]$. By
Proposition~\ref{f16v61:shifted}, $E[\tau(y)]=1+E[T_2(y)]+\sum_{N\ge4\ {\rm even}}E[T_N(y)]$
with $|E[T_N(y)]|\le(A+B)^N/\sqrt{N!}$, and $E[T_2(y)]=-\frac12L^2h^2|q_\mu|^2-\beta_{L,\gamma}$
by V59's Proposition~\ref{f16v59:moments}; the even tail is at most
$0.25(A+B)^4$ and $(A+B)^4\le8(A^4+B^4)=72\gamma^{-4/3}|q_\mu|^4+8x_L^2$ with
$x_L\le1.52L/\gamma$. The fluctuation bound is the triangle inequality with
\eqref{f16v61:shifted-bound} summed from $N=3$, and the variance of $T_2(y)$
is V59's \eqref{f16v59:variance}, the cross term $a\cdot\sum u_j$ having a
variance of order $L^2|a|^2\cdot L\|C\|_1$, i.e.\ $\gamma^{-2/3}L/\gamma$,
which is absorbed in the $O(1)$.
\end{proof}

Three remarks. First, the theorem removes the deterministic condition of
V59 entirely: the only condition is $A+B\le1$, i.e.\ $L\le0.6\gamma$ and
$\gamma\ge$ a constant, and the bias is then controlled to relative order
$z^2$ throughout the small-volume regime, with the expected structure of a
double expansion in $\gamma^{-2/3}|q|^2$ and $z^2$. Second, the fourth-order
mean $E[T_4(y)]$ is itself an explicit Wick polynomial in $a$ and $C_L$,
recorded by the verifier for the test covariance; including it in the bias
correction would push the error to $(A+B)^6$. Third, the per-line
fluctuation $0.77(A+B)^3$ is relatively $L^{3/2}\gamma^{-5/6}$ and is sharp
only for $L\ll\gamma^{5/9}$; the average over $L^2$ lines is expected to
reduce it, but that requires the cross-line covariance of the cubic chaos,
which decays only like the inverse square of the transverse distance and
is not $\ell^1$-summable over the plane; this is the next lemma, N1$'$, and
it is what separates the present bias theorem from a fluctuation theorem of
the same scope.

\subsection{Verification and preservation}

The new diagnostic is
\begin{center}\small\path{scripts/verify_ym_f16_v61_line_chaos.py}.\end{center}
Its hashed exact content is: the two counting facts for $N\le24$ and
$N\le6$; the second-order identity on five links; the exact Wick values of
$E[T_3^2]$ ($L=4$), $E[T_4^2]$ ($L=3$), $E[T_3]=0$ and $E[T_4]$ for the
rational circulant covariance $C(d)=\frac15 2^{-\min(d,L-d)}$, compared with
the bounds (ratios 1.16e-03 and 0.037); the same with a symbolic shift
$a$, where $E[T_3(a+u)]$ is the zero polynomial and $E[T_4(a+u)]$ is
recorded; and the rational tail constants. Its floating diagnostics are the
hard line covariance $C_L(d)$, $d^2C_L(d)$ and $\|C_L\|_1$. It has
12223 bytes and SHA--256
\begin{center}\scriptsize\ttfamily 4A4A2546531A3480E16641B297252DE963B58A790B880E3DD4565F088A0B27C7.\end{center}
Its canonical hashed payload has SHA--256
\begin{center}\scriptsize\ttfamily 3B1EA5E51D9FA6C1BB92F422D97329EFBF9A4EC1A459DE87FD75EB9B847358FD.\end{center}
Only this terminal insertion and the literal title version differ from V60;
removing them recovers its bytes exactly. The next companion checkpoint is
V65.

\begin{firewall}
V61 proves the chaos expansion of the Polyakov-line trace with the
factorial Wick bound, its extension to a constant shift with vanishing odd
means, and the resulting bias theorem for V59's coordinate with explicit
constants throughout the small-volume regime, given the computed
$\ell^1$ norm of the hard line covariance, whose uniform boundedness in $L$
is numerical. It does not prove the line-averaged fluctuation bound, does
not modify the native chain V44--V48, and does not bear on volume-uniform
control beyond the small-volume regime or on an interacting continuum mass
gap.
\end{firewall}
% END F16 V61 APPEND

% BEGIN F16 V62 APPEND -- 2026-09-18
% Insert before the final document-ending command of cumulative V61.
% This insertion leaves every predecessor body and the native operator unchanged.

\clearpage
\section{V62: line-averaged fluctuations and the complete Polyakov-coordinate theorem in the small-volume regime}
\label{f16v62:context}

V61 controlled the bias of the gauge-invariant coordinate $T_\mu$ throughout
the small-volume regime but bounded its fluctuation only line by line,
which is sharp only for $L\ll\gamma^{5/9}$. This note proves the
line-averaged fluctuation bound. The mechanism is that the cubic chaos
$T_3$, and the mixed soft--hard quadratic chaos that the soft field adds
to it, are Wick-ordered: their partial contractions vanish by colour
isotropy. For Wick-ordered chaoses the covariance between two lines
involves only complete cross-line pairings, each weighted by the $\ell^1$
norm $c(\Delta)$ of the hard covariance between links of two lines at
transverse offset $\Delta$, which decays like $0.3\|C\|_1/|\Delta|$. The
cubic cross-line covariance is therefore of order $L^3c(\Delta)^3$,
summable over the transverse plane, and the averaged cubic fluctuation is
smaller than the line fluctuation by a factor $L^{-1}$ up to a constant:
\[
 \operatorname{Var}\Bigl[L^{-2}\sum_{x_\perp}T_3^{(x_\perp)}\Bigr]\le4.5\,s_3\,\|C\|_1^3\,L,
 \qquad s_3=\sum_{\Delta}\bigl(c(\Delta)/\|C\|_1\bigr)^3=\text{1.23}\ (L=64).
\]
Together with V59 and V61 this yields the coordinate theorem: at the
critical scale and for $L\le\gamma^{2/3-\delta}$, the Polyakov coordinate
$T_\mu$ has expectation $1-\frac12\gamma^{-2/3}|q_\mu|^2-\beta_{L,\gamma}$ up
to a relative error $19z^2+72\gamma^{-2/3}|q_\mu|^2$ and a standard deviation
whose leading terms, relative to the signal, are
$\gamma^{-1/3}\sqrt{0.06\log L}$ and $0.3\sqrt{s_3}\,L^{1/2}\gamma^{-5/6}$,
both tending to zero in the regime. The native task N1 of V60 is complete.

\subsection{Dependencies}

V59's coordinate and second-order variance, V61's chaos expansion, Wick
bound and shifted-field proposition, V53's transverse waves for the
covariance. Nothing in the native chain is modified.

\subsection{Wick ordering of the cubic chaoses}

\begin{proposition}[Wick-ordered cubic and mixed chaoses]
\label{f16v62:wick}
Let $u$ be a centred Gaussian field on a line with colour-diagonal
covariance and $a\in\mathbb R^3$ constant. Then
\begin{equation}
 T_3(a+u)=T_3(u)+Q_a(u),\qquad T_3(u)=-\sum_{i<j<k}\det(u_i,u_j,u_k),
 \label{f16v62:split}
\end{equation}
with $Q_a$ the part of $T_3(a+u)$ linear in $a$, a quadratic form in $u$;
the parts quadratic and cubic in $a$ vanish identically. Both $T_3$ and
$Q_a$ are Wick-ordered: $E[T_3(u)\,u^c_v]=0$ and $E[Q_a(u)\,u^c_v]=0$ for
every link $v$ and colour $c$. Consequently, for two lines $x,x'$ with
jointly Gaussian fields, $E[T_3^{(x)}T_3^{(x')}]$ and $E[Q_a^{(x)}Q_a^{(x')}]$
are sums over complete cross-line pairings only.
\end{proposition}
\begin{proof}
Expanding $\det(a+u_i,a+u_j,a+u_k)$, the terms with two or three copies of
$a$ vanish because a determinant with two equal columns is zero; this
gives \eqref{f16v62:split} with $Q_a=-\sum_{i<j<k}[\det(a,u_j,u_k)+\det(u_i,a,u_k)+\det(u_i,u_j,a)]$.
A partial contraction of $T_3$ replaces two of its three factors by
$\delta^{ab}C(\cdot)$, and $\epsilon_{abc}\delta^{ab}=0$; the same holds for
the two $u$-factors of each term of $Q_a$. A homogeneous polynomial of
degree $m$ all of whose partial contractions vanish equals its Wick-ordered
version, and for Wick-ordered polynomials $P,Q$ of degree $m$ the diagram
formula gives $E[PQ]$ as the sum over pairings in which every pair joins a
variable of $P$ to a variable of $Q$. The verifier checks the identity
\eqref{f16v62:split}, the two Wick-ordering statements and the
cross-pairing formula exactly on four links.
\end{proof}

\subsection{Cross-line covariances}

Write $C_3(d,\Delta)$ for the covariance, in units of $1/\gamma$, between
the line-direction colour components of the hard field on two $\mu$-links
at longitudinal distance $d$ and transverse offset $\Delta\in(\mathbb Z/L)^2$,
$C_3(d,\Delta)=V^{-1}\sum_{k\ne0}\cos(k_1d+k_\perp\cdot\Delta)P_{11}(k)/(2\omega_k)$,
and
\begin{equation}
 c(\Delta)=\sum_{d}|C_3(d,\Delta)|,\qquad c(0)=\|C\|_1,\qquad
 s_p=\sum_{\Delta\in(\mathbb Z/L)^2}\Bigl(\frac{c(\Delta)}{\|C\|_1}\Bigr)^p .
 \label{f16v62:sp}
\end{equation}
Numerically, $|\Delta|\,c(\Delta)/\|C\|_1\approx0.3$ for $1\le|\Delta|\le L/4$,
so $c$ decays like the inverse transverse distance, and at $L=8,16,32,64$:
$s_1/L=$ $0.612, 0.589, 0.575, 0.566$, $s_2=$ $1.37, 1.66, 2.00, 2.37$ (growing like $0.5\log L$), $s_3=$ $1.06, 1.12, 1.18, 1.23$,
$s_4=$ $1.01, 1.03, 1.05, 1.06$. The boundedness of $s_3$, $s_4$ and the logarithmic growth
of $s_2$ follow from the $1/|\Delta|$ decay, which is the two-dimensional
Fourier transform of the $1/|k_\perp|$ singularity of the resonant modes;
as in V61 the uniform bounds are numerical.

\begin{theorem}[Cross-line covariance bounds]
\label{f16v62:cross}
For two lines at transverse offset $\Delta$,
\begin{align}
 \bigl|E[T_3^{(x)}T_3^{(x')}]\bigr|&\le\tfrac{27}{36}\cdot6\,(Lc(\Delta))^3=4.5\,L^3c(\Delta)^3,
 \label{f16v62:cov3}\\
 \bigl|E[Q_a^{(x)}Q_a^{(x')}]\bigr|&\le13.5\,L^4|a|^2c(\Delta)^2,
 \label{f16v62:covQ}\\
 \bigl|\operatorname{Cov}(T_4^{(x)},T_4^{(x')})\bigr|&\le\tfrac{81}{576}\bigl[72(L\|C\|_1)^2(Lc(\Delta))^2+24(Lc(\Delta))^4\bigr].
 \label{f16v62:cov4}
\end{align}
\end{theorem}
\begin{proof}
By Proposition~\ref{f16v62:wick} only complete cross pairings contribute
to \eqref{f16v62:cov3}: there are $3!=6$ of them, each cross pair
contributes at most $3$ from the colour sum and $\sum_{i,j}|C_3(i-j,\Delta)|=Lc(\Delta)$
from the positions, and the sorted-sequence count of V61 gives $(3!)^{-2}$;
hence $(3^3/36)\cdot6(Lc)^3$. For \eqref{f16v62:covQ}, $Q_a$ has one
deterministic slot, contributing $\sqrt3|a|$ after its colour sum and $L$
from its position, and two Gaussian slots; with $\binom31$ choices of the
$a$-slot in each copy, $2$ complete cross pairings and the factor $(3!)^{-2}$,
the bound is $\frac1{36}\cdot9\cdot3L^2|a|^2\cdot2\cdot9\,(Lc)^2=13.5L^4|a|^2c^2$.
For \eqref{f16v62:cov4}, $T_4$ is not Wick-ordered; group the pairings of
the eight slots by the number $m$ of cross pairs: $m=0$ gives exactly
$E[T_4^{(x)}]E[T_4^{(x')}]$, which the covariance subtracts; $m=2$ has
$\binom42^2\cdot2!=72$ pairings with two within-line pairs, and $m=4$ has
$4!=24$; the colour factor is $3^4$ and the count $(4!)^{-2}$.
\end{proof}

\subsection{The averaged fluctuation and the coordinate theorem}

\begin{theorem}[Fluctuation of the Polyakov coordinate]
\label{f16v62:fluct}
At the critical scale $a_\mu=hq_\mu$, with $A=\sqrt3\gamma^{-1/3}|q_\mu|$,
$B=\sqrt{6L\|C\|_1}$, $\|C\|_1$ in units of $1/\gamma$ and $A+B\le1$,
\begin{equation}
 \boxed{\
 \begin{aligned}
 \operatorname{Var}_\varphi[T_\mu]^{1/2}\le{}&\operatorname{Var}[\bar T_2]^{1/2}
 +2.13\sqrt{s_3}\,\|C\|_1^{3/2}L^{1/2}+3.7\sqrt{s_2}\,L|a_\mu|\,\|C\|_1\\
 &+L\|C\|_1^2\sqrt{10.2\,s_2+3.4\,s_4}+\sqrt{\tfrac{(A+B)^8-B^8}{24}}+0.11\,(A+B)^5,
 \end{aligned}\ }
 \label{f16v62:fluct-bound}
\end{equation}
where $\bar T_2$ is the line average of the quadratic chaos of the shifted
field, with $\operatorname{Var}[\bar T_2]^{1/2}=\gamma^{-1}\sqrt{\frac3{16\pi}\log L+O(1)}$
by V59. Relative to the signal $\gamma^{-2/3}|q_\mu|^2$, and with
$\|C\|_1\le0.253/\gamma$, the six terms are of orders
$\gamma^{-1/3}\sqrt{\log L}$, $0.27\sqrt{s_3}L^{1/2}\gamma^{-5/6}$,
$0.94\sqrt{s_2}\gamma^{-2/3}$, $0.06\sqrt{10s_2+3s_4}\,z\gamma^{-2/3}$,
$18L^{7/4}\gamma^{-5/4}$ and $5.1L^{5/2}\gamma^{-11/6}$ (the last two for
$B\ge A$), all of which tend to zero for $L\le\gamma^{2/3-\delta}$.
\end{theorem}
\begin{proof}
Decompose $\tau(y)=\sum_NT_N(y)$ on each line and average; by the triangle
inequality in $L^2$ it suffices to bound the standard deviation of each
averaged chaos. $N=2$ is V59. For $N=3$ use \eqref{f16v62:split}:
$\operatorname{Var}[L^{-2}\sum_xT_3^{(x)}]=L^{-4}\sum_{x,x'}E[T_3^{(x)}T_3^{(x')}]\le L^{-2}\sum_\Delta4.5L^3c(\Delta)^3
=4.5s_3\|C\|_1^3L$ by \eqref{f16v62:cov3}, and likewise
$\operatorname{Var}[L^{-2}\sum_xQ_a^{(x)}]\le13.5s_2L^2|a|^2\|C\|_1^2$ by
\eqref{f16v62:covQ}. For $N=4$ split $T_4(y)=T_4(u)+[T_4(y)-T_4(u)]$; the
first part is averaged with \eqref{f16v62:cov4},
$\operatorname{Var}\le L^{-2}\sum_\Delta\frac{81}{576}L^4[72\|C\|_1^2c^2+24c^4]
=L^2\|C\|_1^4(10.125s_2+3.375s_4)$, and the second is bounded line by line
by V61's shifted bound, whose terms with at least one $a$-slot sum to at
most $((A+B)^8-B^8)/4!$. For $N\ge5$ the line bound $(A+B)^N/\sqrt{N!}$ and
$\sum_{N\ge5}(N!)^{-1/2}\le0.11$ suffice. The orders follow from
$L|a_\mu|=\gamma^{-1/3}|q_\mu|$, $B^2\le1.52L/\gamma$ and, for the two
per-line terms with $B\ge A$, $(A+B)^8-B^8\le8A(2B)^7$ and $(A+B)^5\le32B^5$.
\end{proof}

\begin{theorem}[The Polyakov coordinate at the critical scale]
\label{f16v62:coordinate}
Let $L=2m\ge8$, $\gamma\ge\gamma_0$ and $L\le\gamma^{2/3-\delta}$. On
$U=\operatorname{Exp}(h\mathsf Kq+\varepsilon u)$ with $u$ distributed by
V45's hard Gaussian ground, the gauge-invariant coordinate
$T_\mu(U)=L^{-2}\sum_{x_\perp}\frac12\Tr H_{x_\perp}(U)$ satisfies
\begin{equation}
 T_\mu=1-\tfrac12\gamma^{-2/3}|q_\mu|^2-\beta_{L,\gamma}+\gamma^{-2/3}\bigl(b+\zeta\bigr),
 \qquad|b|\le72\gamma^{-2/3}|q_\mu|^4+19z^2,\quad
 \operatorname{Var}[\zeta]^{1/2}\le\eta(L,\gamma)\to0,
 \label{f16v62:coordinate-theorem}
\end{equation}
with $\beta_{L,\gamma}=c_\beta(L)L/\gamma$ explicit (V59), $z=L\gamma^{-2/3}$,
and $\eta$ the sum of the six relative orders of Theorem~\ref{f16v62:fluct}.
In particular $2(1-T_\mu)-2\beta_{L,\gamma}$ estimates $\gamma^{-2/3}|q_\mu|^2$
with relative bias $O(z^2+\gamma^{-2/3})$ and relative noise
$O(\gamma^{-1/3}\sqrt{\log L}+L^{1/2}\gamma^{-5/6})$, uniformly in the regime.
\end{theorem}
\begin{proof}
Combine V61's Theorem~\ref{f16v61:bias} for the mean with
Theorem~\ref{f16v62:fluct} for the fluctuation; the condition $A+B\le1$
holds in the regime for $\gamma_0$ large, since $A\le\sqrt3\gamma^{-1/3}|q_\mu|$
and $B^2\le1.52L/\gamma\le1.52\gamma^{-1/3-\delta}$.
\end{proof}

The theorem completes, for the coordinate, what V59 set out: the toron
coordinate can be read gauge-invariantly from Polyakov lines throughout
the small-volume regime, with the intrinsic parameter $z$ entering only
through the quartic bias and with all constants explicit up to the
numerically bounded covariance sums $\|C\|_1,s_2,s_3,s_4$. Two things it
does not do. It does not replace the tube in the native chain: the
coordinate is a function of the configuration, and V44--V48 also need the
kernel comparison in a chart, which is task N2. And the covariance sums are
bounded numerically; their uniform bounds follow from the $|k_1|$ cusp of
the line spectral density and the $1/|k_\perp|$ singularity of the resonant
modes, which is a two-page Fourier estimate not written here.

\subsection{Verification and preservation}

The new diagnostic is
\begin{center}\small\path{scripts/verify_ym_f16_v62_line_average.py}.\end{center}
Its hashed exact content is the Wick ordering of $T_3$ and $Q_a$ and the
identity \eqref{f16v62:split} on four links, the equality of
$E[T_3^{(x)}T_3^{(x')}]$ with the complete cross-pairing sum and its bound
(ratio 9.26e-03) for a rational two-line covariance, the bound on the
$Q_a$ cross covariance at a rational $a$ (ratio 1.03e-03), and the bound on
$\operatorname{Cov}(T_4^{(x)},T_4^{(x')})$ on three links (ratio 1.20e-01).
Its floating diagnostics are the three-dimensional hard covariance by fast
Fourier transform, $c(\Delta)$, its decay and the sums $s_1,\dots,s_4$. It
has 9063 bytes and SHA--256
\begin{center}\scriptsize\ttfamily F897805DD4BA41627E7E51CF1DFD070681D52BAAFC420B1D3A896E0C5B56566D.\end{center}
Its canonical hashed payload has SHA--256
\begin{center}\scriptsize\ttfamily 3F0298831E11D3988EB888865D90009FD3A8294B4A97A89C6F7B1D9036B40112.\end{center}
Only this terminal insertion and the literal title version differ from V61;
removing them recovers its bytes exactly. The next companion checkpoint is
V65.

\begin{firewall}
V62 proves the Wick ordering of the cubic and mixed chaoses, the
cross-line covariance bounds, the line-averaged fluctuation bound and the
resulting coordinate theorem for the Polyakov coordinate at the critical
scale in the small-volume regime, with the covariance sums bounded
numerically. It does not modify the native chain V44--V48, does not remove
the tube from any native theorem, and does not bear on volume-uniform
control beyond the small-volume regime or on an interacting continuum mass
gap.
\end{firewall}
% END F16 V62 APPEND

% BEGIN F16 V63 APPEND -- 2026-09-18
% Insert before the final document-ending command of cumulative V62.
% This insertion leaves every predecessor body and the native operator unchanged.

\clearpage
\section{V63: decay of the hard covariance and the unconditional coordinate theorem}
\label{f16v63:context}

V61 and V62 left one gap: the $\ell^1$ norms $\|C_L\|_1$, $c(\Delta)$ and
the transverse sums $s_p$ of the hard covariance were computed numerically
and their boundedness in $L$ was asserted from the observed decay. This
note proves the decay. The hard covariance of the line-direction
components is the Fourier series of a weight $w(k)=P_{11}(k)/(2\omega_k)$
on the momentum lattice with a single singularity, of degree $-1$, at
$k=0$; two summations by parts in one momentum direction, an exact
identity on the periodic lattice, convert the decay of $C_3(x)$ in $x_i$
into the summability of the second differences of $w$, and the second
derivatives of $w$ are bounded by $K|k|^{-3}$ away from the origin, which
is summable up to one logarithm. The result is
\[
 |C_3(x)|\le\frac{\kappa_0+\kappa_1\log L}{\max_ix_i^2}\qquad(x\ne0,\ |x_i|\le L/2),
\]
with absolute constants, hence $\|C_L\|_1$, $c(\Delta)|\Delta|$, $s_3$ and
$s_2/\log L$ are bounded by powers of $\log L$, and the coordinate theorem
of V62 holds unconditionally in the regime $L\log^3L\ll\gamma^{2/3}$. The
logarithm is an artifact of estimating the singular sum by absolute
values: the true line decay has the continuum constant $1/(2\pi^2)$, which
the numerics confirm.

\subsection{Dependencies}

V61's covariance $C_L(d)$ and $\|C\|_1$, V62's $c(\Delta)$ and $s_p$, V52's
hard spectrum, V53's transverse projector. Nothing else.

\subsection{Summation by parts on the momentum lattice}

Let $\mathcal K=(2\pi\mathbb Z/L)^3$ be the momentum lattice, $e_i'=(2\pi/L)e_i$,
and for $f:\mathcal K\to\mathbb C$ let $(\Delta_i^2f)(k)=f(k+2e_i')-2f(k+e_i')+f(k)$.

\begin{proposition}[Exact summation by parts]
\label{f16v63:sbp}
For every $f$ on $\mathcal K$ and every $x\in\mathbb Z^3$ with $x_i\not\equiv0\pmod L$,
\begin{equation}
 \sum_{k\in\mathcal K}e^{ik\cdot x}f(k)=\bigl(e^{-2\pi ix_i/L}-1\bigr)^{-2}\sum_{k\in\mathcal K}e^{ik\cdot x}(\Delta_i^2f)(k),
 \qquad\bigl|e^{-2\pi ix_i/L}-1\bigr|\ge\frac{4|x_i|}L\ \ (|x_i|\le L/2).
 \label{f16v63:sbp-identity}
\end{equation}
\end{proposition}
\begin{proof}
Shifting the summation variable, $\sum_ke^{ik\cdot x}f(k+me_i')=e^{-im\theta}\sum_ke^{ik\cdot x}f(k)$
with $\theta=2\pi x_i/L$, so the right-hand sum equals
$(e^{-2i\theta}-2e^{-i\theta}+1)\sum e^{ik\cdot x}f=(e^{-i\theta}-1)^2\sum e^{ik\cdot x}f$.
The modulus is $2|\sin(\theta/2)|\ge4|x_i|/L$ by $\sin t\ge2t/\pi$ on
$[0,\pi/2]$. The verifier checks the identity in Gaussian-rational
arithmetic on $(\mathbb Z/4)^3$ for every $x$ and $i$, and the modulus
bound exactly for $L=4,8$.
\end{proof}

\subsection{Second differences of the hard weight}

With $s_\mu=4\sin^2(k_\mu/2)$, $\lambda=\sum_\mu s_\mu=\hat k^2$ and
$\omega=\sqrt{\lambda(1+\lambda/4)}$, the weight of V61 is, for $k\ne0$,
\begin{equation}
 w(k)=\frac{P_{11}(k)}{2\omega_k}=\frac{s_2+s_3}{2\lambda\,\omega}
 =\lambda^{-1/2}F(\sigma,\lambda),\qquad
 F(\sigma,\lambda)=\frac{\sigma_2+\sigma_3}{2\sqrt{1+\lambda/4}},\quad\sigma=s/\lambda,
 \label{f16v63:weight}
\end{equation}
and $w(0):=0$; $\sigma$ lies in the simplex and $\lambda\in(0,12]$, so $F$
is smooth and bounded with bounded derivatives on a compact set.

\begin{proposition}[Derivative bounds]
\label{f16v63:derivatives}
There are absolute constants $K_i$ such that for $k\in[-\pi,\pi]^3\setminus\{0\}$
\begin{equation}
 |\partial_{k_i}w(k)|\le K_i'|k|^{-2},\qquad|\partial_{k_i}^2w(k)|\le K_i|k|^{-3}.
 \label{f16v63:d2w}
\end{equation}
Consequently, with $\Sigma_i(L)=L^2V^{-1}\sum_{k\in\mathcal K}|\Delta_i^2w(k)|$,
\begin{equation}
 \Sigma_i(L)\le c_0+c_1\log L
 \label{f16v63:Sigma}
\end{equation}
for absolute $c_0,c_1$. Numerically $\sup_k|k|^3|\partial_{k_1}^2w|=$ 4.80 and
$\sup_k|k|^3|\partial_{k_2}^2w|=$ 2.74 (the same at every $L$), and
$\Sigma_1(L)=$ $3.92, 5.13, 6.33, 7.53$ at $L=16,32,64,128$, growing by 1.19 per
doubling.
\end{proposition}
\begin{proof}
On the Brillouin zone $\frac4{\pi^2}|k|^2\le\lambda\le|k|^2$, since
$\frac{2t}{\pi}\le\sin t\le t$ on $[0,\pi/2]$. The chain rule with
$\partial_{k_i}s_\mu=2\delta_{i\mu}\sin k_i$, $|2\sin k_i|\le2\sqrt{s_i}\le2\sqrt\lambda$
and $|\partial_{k_i}^2s_\mu|\le2$ gives $|\partial_{k_i}\lambda|\le2\sqrt\lambda$,
$|\partial^2_{k_i}\lambda|\le2$, $|\partial_{k_i}\sigma_\mu|\le4\lambda^{-1/2}$
and $|\partial_{k_i}^2\sigma_\mu|\le K\lambda^{-1}$, so that
$|\partial_{k_i}F|\le K\lambda^{-1/2}$ and $|\partial^2_{k_i}F|\le K\lambda^{-1}$;
with $|\partial_{k_i}\lambda^{-1/2}|\le\lambda^{-1}$ and
$|\partial_{k_i}^2\lambda^{-1/2}|\le K\lambda^{-3/2}$ the product rule gives
\eqref{f16v63:d2w}. For \eqref{f16v63:Sigma}, split $\mathcal K$ into the
modes with $|k|\le6\pi/L$, at most $c$ of them, each with
$|\Delta_i^2w|\le4\max_{k\ne0}w\le2/\omega_{\min}\le L/\pi$, and the rest,
where $|\Delta_i^2w(k)|\le(2\pi/L)^2\sup|\partial_{k_i}^2w|$ over a segment of
length $4\pi/L$ on which $|k'|\ge|k|/3$, so that
$|\Delta_i^2w(k)|\le(2\pi/L)^2\cdot27K_i|k|^{-3}$. Then
$\sum_{|k|>6\pi/L}|k|^{-3}\le(L/2\pi)^3\sum_{3<|m|\le\sqrt3L/2}|m|^{-3}\le(L/2\pi)^3\cdot4\pi(\log L+c')$,
and multiplying by $L^2V^{-1}$ gives \eqref{f16v63:Sigma} with
$c_1=54K_i$ and $c_0$ from the near-origin modes.
\end{proof}

\subsection{Decay of the covariance and the bounded sums}

\begin{theorem}[Decay of the hard covariance]
\label{f16v63:decay}
For every $L=2m\ge8$ and every $x\in\mathbb Z^3\setminus\{0\}$ with
$|x_i|\le L/2$,
\begin{equation}
 \boxed{\quad|C_3(x)|\le\frac{\min_{i:\,x_i\ne0}\Sigma_i(L)}{16\,x_i^2}
 \le\frac{\kappa_0+\kappa_1\log L}{\max_ix_i^2},\quad}
 \label{f16v63:decay-bound}
\end{equation}
with $\kappa_j=c_j/16$. Consequently, with absolute constants,
\begin{equation}
 \|C_L\|_1\le\kappa_2+\tfrac{\pi^2}3(\kappa_0+\kappa_1\log L),\qquad
 c(\Delta)\le\frac{6(\kappa_0+\kappa_1\log L)}{|\Delta|_\infty},\qquad
 s_2\le K\log^3L,\qquad s_3,s_4\le K\log^3L,
 \label{f16v63:sums}
\end{equation}
where $\kappa_2$ bounds $C_3(0)=V^{-1}\sum_{k\ne0}w(k)\le V^{-1}\sum_{k\ne0}(2\omega_k)^{-1}$,
a Riemann sum of $1/(2|k|)$. Numerically $\max_{x\ne0}|x|_L^2|C_3(x)|=$ $0.0834, 0.0832, 0.0832, 0.0832$
and $\max_{x\ne0}(\max_ix_i)^2|C_3(x)|=$ $0.0557, 0.0574, 0.0578, 0.0579$ at $L=16$ to $128$.
\end{theorem}
\begin{proof}
Apply \eqref{f16v63:sbp-identity} to $f=w$ with the index $i$ of a nonzero
coordinate, take absolute values, and use \eqref{f16v63:Sigma}. For the
sums: $\|C_L\|_1=|C_3(0)|+2\sum_{d=1}^{L/2}|C_3(d,0)|$ and
$\sum_dd^{-2}\le\pi^2/6$; $c(\Delta)=\sum_d|C_3(d,\Delta)|$ is bounded by
$(\kappa_0+\kappa_1\log L)\sum_d\max(d,|\Delta|_\infty)^{-2}\le(\kappa_0+\kappa_1\log L)(2|\Delta|_\infty^{-1}+2|\Delta|_\infty^{-1}+|\Delta|_\infty^{-2})$;
then $s_p=\sum_\Delta(c(\Delta)/\|C_L\|_1)^p$ is bounded, using
$c(\Delta)\le\|C_L\|_1$ for $|\Delta|_\infty\le1$ and $\|C_L\|_1\ge C_3(0)\ge\kappa_3>0$,
by $9+(6(\kappa_0+\kappa_1\log L)/\kappa_3)^p\sum_{|\Delta|_\infty\ge2}|\Delta|_\infty^{-p}$,
which is $O(\log^{p}L\cdot\log L)$ for $p=2$ and $O(\log^pL)$ for
$p\ge3$; the lower bound $C_3(0)\ge\kappa_3$ holds because $w\ge0$ and the
$k_1=0$, $|k_\perp|\le\pi/2$ modes alone contribute a Riemann sum of
$1/(2|k_\perp|)$ over a quarter of the transverse zone.
\end{proof}

\begin{theorem}[Unconditional coordinate theorem]
\label{f16v63:unconditional}
Theorems \ref{f16v61:bias}, \ref{f16v62:fluct} and \ref{f16v62:coordinate}
hold with the numerical constants $\|C_L\|_1\le0.253/\gamma$, $s_2$, $s_3$,
$s_4$ replaced by the bounds \eqref{f16v63:sums}: for
$L\log^3L\le\gamma^{2/3-\delta}$ the Polyakov coordinate has relative bias
$O(z^2\log^2L+\gamma^{-2/3})$ and relative fluctuation
$O(\gamma^{-1/3}\log^{2}L+L^{1/2}\gamma^{-5/6}\log^{3}L)$, with absolute
constants, and no numerical input remains.
\end{theorem}
\begin{proof}
Every appearance of $\|C\|_1$, $c(\Delta)$ and $s_p$ in V61 and V62 enters
through an upper bound, so the proofs go through verbatim with
\eqref{f16v63:sums}; the powers of $\log L$ are absorbed as stated.
\end{proof}

\subsection{The continuum constant of the line decay}

The logarithm in \eqref{f16v63:Sigma} comes from bounding the oscillatory
sum $\sum_ke^{ik\cdot x}\Delta_i^2w$ by absolute values, and it is absent
from the true decay: the singular part of $w$ near the origin is
$w_s(k)=k_\perp^2/(2|k|^3)$, whose continuum Fourier transform is
\[
 \int\frac{d^3k}{(2\pi)^3}e^{ik\cdot x}\frac{k_\perp^2}{2|k|^3}=\frac{x_1^2}{2\pi^2|x|^4},
\]
from the transforms $1/|k|\mapsto1/(2\pi^2|x|^2)$ and
$k_1^2/|k|^3\mapsto(1/2\pi^2)(|x|^{-2}-2x_1^2|x|^{-4})$. Along the line
this is $1/(2\pi^2d^2)=0.0507/d^2$, and the verifier's
$d^2C_L(d)=$ $0.0389, 0.0578, 0.0545, 0.0493$ at $d=1,2,4,8$ ($L=64$) approach it with lattice
corrections at small $d$; transversally the singular part contributes
nothing at leading order, and the observed $0.3\|C\|_1/|\Delta|$ decay of
$c(\Delta)$ is the finite-volume effect of the excluded zero mode, which is
why $c(\Delta)$ decays only like $|\Delta|^{-1}$. Sharpening
\eqref{f16v63:decay-bound} to $K/|x|^2$ without the logarithm would
require subtracting $w_s$ and treating its lattice sum by Poisson
summation with the zero-mode correction; it is not needed for the
coordinate theorem and is not done.

\subsection{Verification and preservation}

The new diagnostic is
\begin{center}\small\path{scripts/verify_ym_f16_v63_covariance_decay.py}.\end{center}
Its hashed exact content is the summation-by-parts identity on
$(\mathbb Z/4)^3$ for all $x$ and $i$ in Gaussian-rational arithmetic and
the modulus bound for $L=4,8$. Its floating diagnostics are the suprema of
$|k|^3|\partial_{k_i}^2w|$, the sums $\Sigma_i(L)$, the decay constants and
the line values. It has 7214 bytes and SHA--256
\begin{center}\scriptsize\ttfamily EBFDFE367AA87E29F03FAA787EF3601DBFD6041C2E77BFE86DDC41E38A3700CF.\end{center}
Its canonical hashed payload has SHA--256
\begin{center}\scriptsize\ttfamily 7734A8C3F907CAB8DBBC2C73BBAFC6FC27CB39A75793AF53BDB2F118EFA9EFF1.\end{center}
Only this terminal insertion and the literal title version differ from V62;
removing them recovers its bytes exactly. The next companion checkpoint is
V65.

\begin{firewall}
V63 proves the summation-by-parts identity, the derivative bounds of the
hard weight, the decay $|C_3(x)|\le(\kappa_0+\kappa_1\log L)/\max_ix_i^2$ and
the resulting logarithmic bounds on all covariance sums, and thereby
removes the numerical inputs from the coordinate theorem of V59--V62. It
does not remove the logarithms, does not modify the native chain
V44--V48, and does not bear on volume-uniform control beyond the
small-volume regime or on an interacting continuum mass gap.
\end{firewall}
% END F16 V63 APPEND

% BEGIN F16 V64 APPEND -- 2026-09-18
% Insert before the final document-ending command of cumulative V63.
% This insertion leaves every predecessor body and the native operator unchanged.

\clearpage
\section{V64: sup-norm charts, the local structure of the transfer, and the two lemmas of the volume-uniform reduction}
\label{f16v64:context}

V60 named as task N2 the replacement of V41's deterministic $\ell^2$ chart
by a construction whose errors do not accumulate with the volume. This
note carries out the part of N2 that is a matter of exact local algebra,
and isolates precisely what remains. The kinetic kernel of the transfer is
a product of one-link kernels, and in per-link exponential coordinates on
a sup-norm ball of radius $r$ it is sandwiched between two Gaussians whose
variances differ by a factor $1-\delta(r)$, $\delta(r)\le\frac23r^2+\frac15r^4$,
with no dependence on the number of links: the relative error lives in the
exponent, where it adds, and not in a prefactor, where it would multiply.
The Haar density and the one-link normalization are likewise exact products
of per-link factors, $J(x_e)=(\sin|x_e|/|x_e|)^2$ with
$\log J=-|x_e|^2/3-|x_e|^4/90-\cdots$ and $N_\gamma=(2\pi\varepsilon^2)^{3/2}(1-\frac3{8\gamma}+\cdots)$.
Hence, on a sup-norm region, the flat-representation transfer is exactly
an extensive constant times a Gaussian operator with an explicitly
modified hard Hessian times a multiplicative factor with per-link relative
errors. The reduction of the low spectrum to such a region is then an
abstract Feshbach statement about the Doob chain $P=T_\gamma/\Lambda_\gamma$
with two hypotheses: a norm bound for $P$ on the complement and a one-step
escape bound from the region. Both are statements about the Perron vector
$\Omega_\gamma$ along single-link moves, i.e.\ a Harnack inequality with a
quadratic exponent; the crude Harnack inequality with a linear exponent
$3\gamma\theta$ is proved here and is not enough. The volume-uniform
programme is thereby reduced to one estimate on $\Omega_\gamma$, which is
stated in the form in which it will have to be proved.

\subsection{Dependencies}

V41's local kernel comparison \eqref{f16v41:local-kernel} and IMS
localization, V42's Doob chain \eqref{f16v42:Doob}, V37's one-group
normalization, V59--V63 for the regime. Nothing in the native chain is
modified.

\subsection{Sup-norm two-sided kinetic domination}

Write each link as $U_e=\operatorname{Exp}(x_e)$, $x_e\in\mathbb R^3$,
$\operatorname{Exp}(x)=\cos|x|+\sin|x|\,\hat x$, with the round metric in
which $d(1,\operatorname{Exp}x)=|x|$, and let $\theta_e=d(U_e,U'_e)$. The
kinetic kernel is $c_\gamma(U,U')=\prod_eN_\gamma^{-1}e^{-\gamma(1-\cos\theta_e)}$,
$N_\gamma=\int_{\mathrm{SU}(2)}e^{-\gamma(1-\cos\theta)}dU$.

\begin{proposition}[Distance comparison on a ball]
\label{f16v64:distance}
For $|x|,|y|\le r<\pi/2$,
\begin{equation}
 \frac{\sin r}r\,|x-y|\ \le\ d(\operatorname{Exp}x,\operatorname{Exp}y)\ \le\ |x-y|,
 \label{f16v64:dist}
\end{equation}
and consequently
\begin{equation}
 \bigl(1-\delta(r)\bigr)\frac{|x-y|^2}2\ \le\ 1-\cos d(\operatorname{Exp}x,\operatorname{Exp}y)\ \le\ \frac{|x-y|^2}2,
 \qquad 1-\delta(r)=\Bigl(\frac{\sin r}r\Bigr)^2\Bigl(1-\frac{r^2}3\Bigr),\quad\delta(r)\le\frac23r^2+\frac15r^4\ (r\le1).
 \label{f16v64:cos}
\end{equation}
\end{proposition}
\begin{proof}
$S^3$ has curvature one, so $\operatorname{Exp}$ does not increase distances
(Rauch), which is the upper bound. The geodesic ball of radius $r<\pi/2$ is
convex, so the minimizing geodesic from $\operatorname{Exp}x$ to
$\operatorname{Exp}y$ stays in it, and on the ball the differential of
$\operatorname{Exp}^{-1}$ has norm at most $r/\sin r$, since the singular
values of $d\operatorname{Exp}_x$ are $1$ and $\sin|x|/|x|$; the preimage of
the geodesic is a curve of length at most $(r/\sin r)d$, which is the lower
bound. For \eqref{f16v64:cos}, $1-\cos\theta\le\theta^2/2$ and
$1-\cos\theta\ge\frac{\theta^2}2(1-\frac{\theta^2}{12})\ge\frac{\theta^2}2(1-\frac{r^2}3)$
for $\theta\le2r$; the explicit bound on $\delta$ uses
$(\sin r/r)^2\ge1-r^2/3-r^4/45$ for $r\le1$, which the verifier certifies
from the exact series $2\log(\sin t/t)=-\frac{t^2}3-\frac{t^4}{90}-\frac{2t^6}{2835}-\cdots$
and $e^u\ge1+u$. The verifier also samples \eqref{f16v64:dist} on $20000$
pairs at $r=\frac12$ and $r=1$, finding minimal ratios $0.9593,0.8430$ against
$\sin r/r=0.9589,0.8415$.
\end{proof}

\begin{theorem}[Sup-norm kinetic sandwich]
\label{f16v64:sandwich}
For all configurations with $\max_e|x_e|\le r$ and $\max_e|x'_e|\le r$, $r\le1$,
\begin{equation}
 \boxed{\quad
 \prod_eN_\gamma^{-1}e^{-\gamma|x_e-x'_e|^2/2}\ \le\ c_\gamma(U,U')\ \le\
 \prod_eN_\gamma^{-1}e^{-\gamma(1-\delta(r))|x_e-x'_e|^2/2},
 \quad}
 \label{f16v64:sandwich-bound}
\end{equation}
with $\delta(r)$ as in \eqref{f16v64:cos}, independent of $L$, $V$ and
$\gamma$. Moreover $N_\gamma=(2\pi\varepsilon^2)^{3/2}\bigl(1-\frac3{8\gamma}+O(\gamma^{-2})\bigr)$,
the verifier's value of $\gamma(1-N_\gamma/(2\pi\varepsilon^2)^{3/2})$ being
$0.3774, 0.3756, 0.3751$ at $\gamma=50,200,800$.
\end{theorem}
\begin{proof}
Multiply \eqref{f16v64:cos} over the links. For the normalization,
$N_\gamma/\operatorname{vol}=\frac2\pi\int_0^\pi e^{-\gamma(1-\cos\theta)}\sin^2\theta\,d\theta$,
and the Laplace expansion with $1-\cos\theta=\theta^2/2-\theta^4/24+\cdots$,
$\sin^2\theta=\theta^2-\theta^4/3+\cdots$ and the Gaussian moments
$\langle\theta^2\rangle=3/\gamma$, $\langle\theta^4\rangle=15/\gamma^2$ gives
the relative correction $\frac{15}{24}\gamma^{-1}-\gamma^{-1}=-\frac38\gamma^{-1}$.
\end{proof}

The contrast with V41 is the point. V41's \eqref{f16v41:local-kernel} has
the form $(1+\delta)\times$Gaussian with a variance factor $1-\delta$ on a
chart of fixed $\ell^2$ radius, and both errors were fixed-$L$ quantities.
Here the prefactor is exact, the variance error is a sup-norm quantity, and
the exponent's relative error $\delta(r)$ adds over links inside the
Gaussian exponent. On the sup-norm region of radius $r=\gamma^{-1/2}\sqrt{c\log V}$,
which by V61's covariance $C_3(0)=0.103/\gamma$ and a union bound carries
all but $Ve^{-c'\log V}$ of the hard Gaussian mass, the variance error is
$\delta\le c\log V/\gamma$, tending to zero in every regime of interest.

\subsection{The Haar density and the local structure of the transfer}

\begin{proposition}[Local factors]
\label{f16v64:local}
In exponential coordinates the Haar measure is $dU_e=J(x_e)dx_e$ with
$J(x)=(\sin|x|/|x|)^2$, and for $|x|\le1$
\begin{equation}
 -\frac{|x|^2}3-\frac{|x|^4}{45}\ \le\ \log J(x)\ \le\ -\frac{|x|^2}3,\qquad
 \log J(x)=-\frac{|x|^2}3-\frac{|x|^4}{90}-\frac{2|x|^6}{2835}-\cdots .
 \label{f16v64:J}
\end{equation}
Consequently, in the flat half-density representation
$T^{\rm flat}(x,x')=\sqrt{J(x)}\,T_\gamma(U,U')\sqrt{J(x')}$ with
$J(x)=\prod_eJ(x_e)$, on the sup-norm region of radius $r\le1$,
\begin{equation}
 T^{\rm flat}(x,x')=N_\gamma^{-3V}\,b_\gamma(U)b_\gamma(U')\,
 \exp\Bigl[-\frac\gamma2\sum_e(1-\delta_e)|x_e-x'_e|^2-\frac{|x|^2+|x'|^2}6\Bigr]\,\rho(x,x'),
 \label{f16v64:structure}
\end{equation}
with $\delta_e\in[0,\delta(r)]$ depending on $(x_e,x'_e)$ and
$\exp[-\frac1{90}\sum_e(|x_e|^4+|x'_e|^4)]\le\rho\le1$; in particular
$\rho\ge\exp[-\frac{r^2}{90}(|x|^2+|x'|^2)]$.
\end{proposition}
\begin{proof}
$J$ is the Jacobian of $\operatorname{Exp}$ on $S^3$ in the round metric.
The series is that of $2\log(\sin t/t)$, whose coefficients are
$-2^{2k}|B_{2k}|/(k(2k)!)$, all negative and decreasing in ratio by more
than $5$ at $t\le1$, which gives the two-sided bound; the verifier
certifies the coefficients and the tail exactly. Then
\eqref{f16v64:structure} is \eqref{f16v64:sandwich-bound} together with
$\sqrt{J(x)J(x')}$ and the identification of $N_\gamma^{-3V}$.
\end{proof}

Three consequences. First, $N_\gamma^{-3V}$ is an exact extensive constant,
cancelling in every ratio $\lambda_j/\lambda_0$. Second, the quadratic Haar
term $-(|x|^2+|x'|^2)/6$ is a shift of the hard Hessian by $\frac2{3\gamma}$
in V45's normalization, an exact modification of the frozen hard Gaussian
whose effect on the hard frequencies $\omega_k=\sinh\xi_k$ is relatively
$O(\gamma^{-1}/\omega_k^2)$, at most $O(L^2/\gamma)$ for the softest hard
mode and $O(1/\gamma)$ on average. Third, the residual factor $\rho$ is
a product of per-link factors $1-O(|x_e|^4)$: under the hard Gaussian its
logarithm has mean $-\frac{2}{90}\cdot3V\,E|x_e|^4\approx-\frac{3V}{45}\cdot\frac{5C_3(0)^2}{\gamma^2}$,
an extensive constant of order $V/\gamma^2$, and standard deviation of order
$\sqrt V/\gamma^2$; it is a bounded multiplicative perturbation whose
extensive part cancels in ratios and whose fluctuating part is small in
$L^2$ for $V\ll\gamma^4$. None of these three items accumulates with the
volume in a ratio, which is what V41's $\ell^2$ chart could not deliver.

\subsection{The Feshbach reduction for the Doob chain and its two hypotheses}

Let $P=T_\gamma/\Lambda_\gamma$ on $L^2(dU)$, a self-adjoint positive
contraction with top eigenvalue $1$, unitarily equivalent to V42's Doob
chain on $L^2(\mu_\gamma)$; multiplication by an indicator commutes with
the unitary, so norms of restrictions may be computed in either
representation.

\begin{theorem}[Abstract reduction]
\label{f16v64:feshbach}
Let $A$ be a measurable region, $Q=1_{A^c}$, and suppose
$\|QPQ\|\le1-\kappa$ for some $\kappa\in(0,1]$. Then for $z\in(1-\kappa/2,1]$,
$z$ is an eigenvalue of $P$ if and only if it is an eigenvalue of
\[
 F(z)=1_AP1_A+1_APQ\,(z-QPQ)^{-1}QP1_A\qquad\text{on }L^2(A),
\]
with equal multiplicities, and $\|F(z)-1_AP1_A\|\le2\|1_APQ\|^2/\kappa$.
If moreover $P$ is the Doob chain of a reversible positive kernel with
stationary law $\mu$, then $\|1_APQ\|^2\le\sup_{U\in A}P(U,A^c)$ and
$\|QPQ\|\le\sup_{U\in A^c}P(U,A^c)$.
\end{theorem}
\begin{proof}
The first part is the Schur complement (Feshbach) identity for the block
decomposition $L^2=L^2(A)\oplus L^2(A^c)$, valid since $z-QPQ$ is
invertible with norm at most $2/\kappa$ on $L^2(A^c)$ for $z>1-\kappa/2$.
The second part is the Schur test: for a Markov kernel with reversible
$\mu$, the row sums of $1_APQ$ are $P(U,A^c)$ for $U\in A$ and its column
sums, computed against $\mu$, are $P(U',A)$ for $U'\in A^c$, both at most
one, and $\|1_APQ\|^2\le\sup_AP(\cdot,A^c)\cdot\sup_{A^c}P(\cdot,A)$;
the bound for $QPQ$ is the same test with both sums over $A^c$.
\end{proof}

For the transfer, with $A$ a sup-norm region (in a fixed gauge, made
gauge invariant by taking the union over the gauge orbit) and the soft
gap $g_*h=g_*(\gamma V)^{-1/3}$, the theorem reduces the low spectrum to
$1_AP1_A$ up to $2\|1_APQ\|^2/\kappa$ provided (H1) $\|QPQ\|\le1-\kappa$
with $\kappa\gg g_*h$, and (H2) $\sup_{U\in A}P(U,A^c)\ll\kappa g_*h$. Both
are one-step statements about the Doob kernel
$P(U,dU')=T_\gamma(U,U')\Omega_\gamma(U')dU'/(\Lambda_\gamma\Omega_\gamma(U))$,
and both need the ratio of the Perron vector along a move.

\begin{proposition}[Crude Harnack inequality]
\label{f16v64:harnack}
If $U'$ differs from $U$ in one link $e$ by the angle $\theta=d(U_e,U'_e)$,
then $e^{-3\gamma\theta}\le\Omega_\gamma(U')/\Omega_\gamma(U)\le e^{3\gamma\theta}$.
\end{proposition}
\begin{proof}
$\Omega_\gamma(U)=\Lambda_\gamma^{-1}\int T_\gamma(U,U'')\Omega_\gamma(U'')dU''$
and $T_\gamma(U,U'')/T_\gamma(U',U'')=e^{-\gamma[\Delta_{\rm kin}+\Delta S/2]}$
with $|\Delta_{\rm kin}|=|\cos d(U'_e,U''_e)-\cos d(U_e,U''_e)|\le\theta$ by
the triangle inequality on $S^3$, and $|\Delta S|\le\sum_{p\ni e}|\operatorname{Re}U_p-\operatorname{Re}U'_p|\le4|U_e-U'_e|\le4\theta$,
since a plaquette product changes by the change of one factor in the
quaternion norm. Hence the ratio of the integrands is between
$e^{-3\gamma\theta}$ and $e^{3\gamma\theta}$.
\end{proof}

The crude inequality has a linear exponent, so it bounds the Doob jump
probability of a link by $e^{3\gamma\theta}$ times the kinetic jump
probability $\asymp e^{-\gamma\theta^2/2}$, which is useless for
$\theta\lesssim6$. What (H1) and (H2) require is the quadratic version:
\begin{equation}
 \Bigl|\log\frac{\Omega_\gamma(U')}{\Omega_\gamma(U)}\Bigr|\le c_1\sqrt\gamma\,\theta+c_2\gamma\theta^2
 \qquad\text{(single-link move, } \theta\le r),
 \label{f16v64:harnack-target}
\end{equation}
uniformly in $L$ on the sup-norm region, which is the statement that the
logarithmic gradient of the Perron vector is of order $\sqrt\gamma$ per
link, as it is for the frozen hard Gaussian where
$\log\varphi=-\frac\gamma2u^{\mathsf T}\omega u$ and $|\nabla_e\log\varphi|\asymp\gamma|u_e|\asymp\sqrt\gamma$.
With \eqref{f16v64:harnack-target}, (H2) follows from
Theorem~\ref{f16v64:sandwich} and a union bound over links, with
$P(U,A^c)\le3Ve^{-(1-\delta)\gamma r^2/4+c_1\sqrt\gamma r}$ for $U$ at
sup-distance $r/2$ inside $A$, and (H1) follows with $\kappa$ of the order
of the hard relaxation rate. The self-consistent route to
\eqref{f16v64:harnack-target} is through the eigen-equation itself: the
ratio is an average of $T(U',U'')/T(U,U'')$ against the measure
$T(U,U'')\Omega(U'')dU''$, whose displacement $U''_e-U_e$ has, under the
same hypothesis, a Gaussian moment generating function at scale
$\gamma^{-1/2}$, and the averaged ratio is then $e^{O(\sqrt\gamma\theta+\gamma\theta^2)}$;
closing this bootstrap on the sup-norm region is the single remaining
lemma of N2, recorded as N2$'$.

\subsection{Verification and preservation}

The new diagnostic is
\begin{center}\small\path{scripts/verify_ym_f16_v64_supnorm_chart.py}.\end{center}
Its hashed exact content is the series of $2\log(\sin t/t)$ to order
$t^{12}$ with the ratio and tail certificates, and the quaternion distance
identity $|p-q|^2=2(1-\cos d(p,q))$ on rational unit quaternions. Its
floating diagnostics are the distance comparison samples, the
normalization ratios and the free one-link autocorrelation
$E[\cos\theta]=1-\frac3{2\gamma}+\cdots$ (verifier: $\gamma(1-E\cos\theta)=$ $1.4923, 1.4981, 1.4995$).
It has 7718 bytes and SHA--256
\begin{center}\scriptsize\ttfamily FFB8C9B2832230A414AA800EF042D5ECEC1B247CAB781861EE879E7CFDCC2646.\end{center}
Its canonical hashed payload has SHA--256
\begin{center}\scriptsize\ttfamily E6DD9A392A5FA3BFF5DA48DB71EA1492D7F020545124CD04CB13C5E79C8B86C9.\end{center}
Only this terminal insertion and the literal title version differ from V63;
removing them recovers its bytes exactly. V65 is the companion checkpoint.

\begin{firewall}
V64 proves the sup-norm two-sided kinetic domination with dimension-free
constants, the exact local structure of the flat-representation transfer,
the abstract Feshbach reduction for the Doob chain with its Schur-test
bounds, and the crude Harnack inequality for the Perron vector. It does
not prove the quadratic Harnack inequality \eqref{f16v64:harnack-target},
hence neither hypothesis (H1), (H2), and it does not modify any theorem of
V41--V48, remove the tube from them, or bear on an interacting continuum
mass gap.
\end{firewall}
% END F16 V64 APPEND

% BEGIN F16 V65 APPEND -- 2026-09-18
% Insert before the final document-ending command of cumulative V64.
% This insertion leaves every predecessor body and the native operator unchanged.

\clearpage
\section{V65: companion checkpoint, programme redirection, five-version review, and the transport identity of the Perron vector}
\label{f16v65:context}

V65 is the scheduled companion checkpoint. It records three things. First,
a second lineage of these notes, branched from the common V48 and continued
in parallel to V56, exists as a separate file; its V49--V56 establish a
joint-regulator critical branch, the bare Wilson calibration, and the
physical fork according to which the finite-torus toron branch is not the
object to be identified with a continuum local mass. Second, in consequence,
the role of the present lineage is redefined: it is the gauge-invariant
uniformity and small-volume RG-support programme, whose products are
regulator-uniform local estimates to be exported to the local-vacuum
programme, and whose primary target is the native concentrated-kernel
theorem N2. Third, one exact tool for N2 is proved: the transport identity of
the Perron vector under a symmetry of the kinetic kernel, with its
Feynman--Kac form and the Poisson equation for the logarithmic gradient.
From this version on every iteration states separately what is proved,
inherited, conditional and open.

\subsection{Companion checkpoint and the concurrent lineage}

The five companion inputs are unchanged since V50 in size and SHA--256: SM
continuum V72 (607372; f44f9745), NS stability V26 (278283; 82c887f8), shell
mixing V16 (623261; bf138e63), parallel YM V5 (54174; 306f1bb8), GRH cofinal
visibility V3 (23851; 4dd31580); the ledger verifier hashes them. New at this
checkpoint is the concurrent file
\begin{center}\small\path{Downloads/Renewal_Yang_Mills_Mass_Gap_Research_Notes_V56.tex}
(1741312 bytes; 888f3cac),\end{center}
whose sections V1--V48 coincide with the present file and whose V49--V56 were
read completely before this checkpoint. Its content, with the labels of that
file, is as follows; these are inherited results, cited and not reproved.
\begin{itemize}
\item Its V49 proves $|A(x)-\pi^{-2}|x|^{-4}|\le C|x|^{-6}$ for the
four-dimensional walk kernel with $a_L(z)=L^3A(0,Lz)$, hence the volume-uniform
winding asymptotics and the common limiting potential. The present V58 proves
the same limit $La_L(z)\to1/(\pi^2|z|^4)$ by a different route, the
Chernoff-shifted majorant, with the uniform bound $C_0/|z|^4$ but without the
$|x|^{-6}$ error term; V58's continuum potential $\ell_\infty$, its Ewald form
and the identity $c_2=Z_3(1)/(3\pi^2)$ are not in that file. Ownership of the
limit itself is shared, with the concurrent V49 earlier.
\item Its V50 proves a size-independent based normal tube in the total link
$\ell^2$ norm (Theorem \emph{A size-independent based normal tube}), the
anchored evaluation bound $\|\eta\|_\infty\le C\|d\eta\|_2$, dimension-free
trace bounds for the quotient metric and the relative coarea density with zero
first jet. This controls the based-gauge \emph{geometry} uniformly; the present
V64 controls the \emph{kinetic kernel} in per-link sup norm. The two are
complementary inputs to N2.
\item Its V51--V53 compare the Wilson kinetic step with heat and subordinated
references uniformly in the link count, normalized at the actual Perron tops,
at polynomial and then linear volume cost.
\item Its V54 proves the joint-regulator critical branch: with the ledger
$\Theta_{\gamma,L}=L^2\gamma^{-1/3}+L^4\gamma^{-2/3}+L^3\gamma^{-1}\to0$, the
actual-top normalized levels of the original Wilson transfer satisfy
$-h^{-1}\log(\lambda^{\rm W}_{a,j}/\Lambda^{\rm W})\to e_{j-1}-e_0$ (Theorem
\emph{Growing-lattice critical spectrum of the original Wilson transfer}).
This is the existing native joint-limit theorem, on the domain $L=O(\gamma^b)$,
$b<1/6$, and it is the baseline that N2 must extend.
\item Its V55 proves the bare calibration $a_t=a_s$, $\gamma=4/g_0^2$, and
that the fixed-physical-box asymptotically free trajectory has
$\log L=\frac{3\pi^2}{11}\gamma+O(\log\gamma)$, exponential in $\gamma$; a
bounded toron coefficient then gives a physical energy $\gamma^{-1/3}\mu/\ell\to0$.
Fixed contractible loops are blind to the soft coordinate.
\item Its V56 proves uniform soft blindness of physically bounded contractible
loops, $0\le1-W\le\frac{\kappa^4}8\gamma^{-4/3}|q|^4e^{2\kappa\gamma^{-1/3}|q|}$,
scalarization on tight soft families, and the exact excision of a reducing
toron band from centered local Grams at cost $4m\epsilon_j^2$, given the
scalar-compression error $\epsilon_j$. It leaves open the complement
contraction $V_j^*Q_je^{-2t_jH_j}Q_jV_j\preceq q_j^2V_j^*Q_jV_j+E_j$ and the
noncollapse and convergence of the local source family.
\end{itemize}
Two readings in the present lineage are corrected by these inputs. The
physical form $e^{-2S_\infty/g}$ of V58 and the small-volume theorems of
V59--V63 concern the toron branch at fixed physical volume in bare units;
after the concurrent V55 they are not statements about the continuum local
mass, and the regime $L\log^3L\ll\gamma^{2/3}$ is not the physical continuum
trajectory, which has $L$ exponential in $\gamma$. Nothing in V49--V64 is
retracted; their scope is as stated in their firewalls.

\subsection{Programme redirection}

The present lineage is henceforth the \emph{gauge-invariant uniformity and
small-volume RG-support programme}, complementary to the local-vacuum
programme of the concurrent file. It does not treat the toron spectral gap
as the physical mass and does not pursue toron spectroscopy beyond what a
uniform comparison inequality requires. Its primary target is
\begin{equation}
 \boxed{\
 \begin{aligned}
 &\text{(N2) an explicit gauge-compatible event }G_{L,\gamma}\text{ with }
 \mathbb P(G_{L,\gamma}^c)\le\varepsilon_{L,\gamma},\\
 &\bigl\|1_{G_{L,\gamma}}(T^{\rm native}_{L,\gamma}-T^{\rm effective}_{L,\gamma})1_{G_{L,\gamma}}\bigr\|_{\rm relevant}\le\eta_{L,\gamma},
 \qquad\varepsilon_{L,\gamma}+\eta_{L,\gamma}\to0
 \end{aligned}\ }
 \label{f16v65:N2}
\end{equation}
along an explicit joint trajectory $L=L(\gamma)$ approaching
$L\lesssim\gamma^{2/3-\delta}$ up to justified logarithms, in the norm actually
used by the V45--V48 comparison, with the bad-event contribution absorbable in
the same operator or form norm and every $L$-dependence explicit; where a
whole-lattice event is too expensive, a local or block event of bounded
overlap with the comparison proved by local composition. The estimate must
pay the correlations of the hard Gaussian field, the gauge constraints, the
soft component, the exponential-map remainders, the Haar factors, the link
count and the norm, and it must use the covariance decay of V63 rather than a
union bound wherever that removes a power or a logarithm. Its products are
to be exported: covariance decay, Wilson-word chaos estimates, source
noncollapse and convergence, local kernel comparisons, and explicit rates for
the scalarization error $\epsilon_{L,\gamma}$ of the concurrent V56. Every
obstruction met outside the small-volume regime is to be classified as a
coordinate failure, a probability-tail failure, an operator-norm failure, a
soft--hard scale collision, a renormalization failure or a physical-source
failure, and the first genuine one is attacked upstream. If N2 closes, the
output splits into a handoff of uniform local estimates to the local-vacuum
programme and a genuine scale-to-scale recursion
$\mathcal F_{j+1}(z)=\mathcal R_j\mathcal F_j(z)\mathcal R_j^*+\mathcal E_j(z)$
with summable errors, rather than a larger simultaneous weak-coupling box.
The two files stay separate until a theorem of this lineage supplies a
missing hypothesis of the local-Gram chain.

\subsection{Five-version review, V60--V64}

\begin{center}\small
\begin{tabular}{@{}p{0.07\textwidth}p{0.88\textwidth}@{}}
\toprule
Version & Established (scope limits in each firewall)\\
\midrule
V60 & Checkpoint and V50--V59 review; consolidated ledger with integrity
check; certified $g_*\in[0.4011,3.0156]$; two corrected forecasts; order of
attack.\\
V61 & Line-chaos lemma: $\|T_N\|_2\le(6L\|C\|_1)^{N/2}/\sqrt{N!}$, shifted
field, vanishing odd means; bias of the Polyakov coordinate to relative
order $19z^2+72\gamma^{-2/3}|q|^2$.\\
V62 & Wick ordering of the cubic chaoses; cross-line covariance bounds;
line-averaged fluctuation; the coordinate theorem in the small-volume
regime with numerically bounded covariance sums.\\
V63 & Summation by parts on the momentum lattice; derivative bounds of the
hard weight; $|C_3(x)|\le(\kappa_0+\kappa_1\log L)/\max_ix_i^2$; all
covariance sums bounded up to logarithms; the coordinate theorem
unconditional for $L\log^3L\ll\gamma^{2/3}$.\\
V64 & Dimension-free sup-norm kinetic sandwich; exact local factors of the
flat transfer; abstract Feshbach reduction for the Doob chain; crude Harnack
inequality; N2 reduced to the quadratic Harnack lemma and the complement
bound.\\
\bottomrule
\end{tabular}
\end{center}

Under the new classification the small-volume regime ends where it does for
identifiable reasons. The coordinate theorem fails at $L\gamma^{-2/3}\to\infty$
because the quartic bias of a Polyakov line of $L$ links, of size $(L/\gamma)^2$,
overtakes the soft signal $\gamma^{-2/3}$: a soft--hard scale collision, which
is V53's intrinsic condition and not a chart artifact. The native chain
V44--V48 fails earlier, at $\gamma\lesssim c_LV$ (V52) and, in the concurrent
V54, at $L\gtrsim\gamma^{1/6}$, because its chart is deterministic in $\ell^2$:
an operator-norm failure of the chart, which is exactly N2. The physical
trajectory $\log L\asymp\gamma$ lies beyond both, and there the failure is
a renormalization failure: the bare soft coefficient does not acquire the
growth $\gamma^{1/3}$ that a fixed-box energy needs. The order of attack is
therefore N2 first, then the RG recursion, and the lineage stops advancing the
polynomial exponent once N2 is closed.

\subsection{The transport identity of the Perron vector}

The Harnack lemma of V64 asks how the Perron vector $\Omega_\gamma$ changes
along a single-link move. The following exact identity is the tool for that
question. Let $\tau$ be a measure-preserving bijection of the configuration
space which leaves the kinetic kernel invariant, $c_\gamma(\tau U,\tau U')=c_\gamma(U,U')$;
left multiplication of one link by a fixed group element is such a map, by
bi-invariance of the one-link kernel and of Haar measure.

\begin{theorem}[Transport identity]
\label{f16v65:transport}
Let $T=bcb$ with $c$ a symmetric positive kernel invariant under $\tau$, $b>0$,
$T\Omega=\Lambda\Omega$ with $\Omega>0$, and $P(U,dU')=T(U,U')\Omega(U')dU'/(\Lambda\Omega(U))$
the Doob chain. Then $w(U)=\Omega(\tau U)/\Omega(U)$ satisfies
\begin{equation}
 \boxed{\quad w(U)=\rho(U)\int P(U,dU')\,\rho(U')\,w(U'),\qquad\rho(U)=\frac{b(\tau U)}{b(U)},\quad}
 \label{f16v65:identity}
\end{equation}
and for every $n\ge1$ the Feynman--Kac form
$w(U)=E_U\bigl[\rho(X_0)\rho(X_1)^2\cdots\rho(X_{n-1})^2\rho(X_n)w(X_n)\bigr]$
along the chain $(X_t)$ started at $U$. For a one-parameter family $\tau_\theta$
with $\rho_\theta=e^{-\theta s+O(\theta^2)}$, the logarithmic derivative
$g=\partial_\theta\log w_\theta|_{\theta=0}$ solves the Poisson equation
\begin{equation}
 (1-P)g=-(1+P)s .
 \label{f16v65:poisson}
\end{equation}
For the single-link rotation of $U_e$ by $\operatorname{Exp}(\theta n)$,
$s=\frac\gamma2\partial_{n}S$ is the directional derivative of the Wilson
action at the link, a function of the four plaquettes containing $e$.
\end{theorem}
\begin{proof}
$\Omega(\tau U)=\Lambda^{-1}\int T(\tau U,U'')\Omega(U'')dU''$; substitute
$U''=\tau U'''$, use the invariance of the measure, and
$T(\tau U,\tau U''')=b(\tau U)c(U,U''')b(\tau U''')=T(U,U''')\rho(U)\rho(U''')$.
Dividing by $\Omega(U)$ gives \eqref{f16v65:identity}; iterating it gives the
Feynman--Kac form; differentiating at $\theta=0$, where $w_0=1$ and $\rho_0=1$,
gives \eqref{f16v65:poisson}. The verifier checks the identity, the
Feynman--Kac form for $n\le5$ and the row sums of $P$ on a cyclic-group model
with a bi-invariant kinetic kernel and a magnetic factor to $10^{-12}$.
\end{proof}

The Poisson equation places the quadratic Harnack lemma of V64 in a definite
form: $g$ is the response of the Doob chain to the local source $-(1+P)s$,
with $|s|\lesssim\gamma\max_{p\ni e}\theta_p$, of order $\sqrt\gamma$ on the
typical region. The sup-norm bound $|g|\le c\sqrt\gamma$ is the statement that
the chain's Poisson kernel applied to a local source of size $\sqrt\gamma$ does
not amplify it, which is a local relaxation statement about the hard modes
around one link and does not involve the soft gap: the soft component of the
source is the derivative of the effective soft potential and is of relative
order $h$. This is the form in which N2$'$ will be attacked, by the hard
relaxation rate rather than by the spectral gap; it is stated here, not
proved.

\subsection{Status}

\emph{Proved in V65}: Theorem~\ref{f16v65:transport}. \emph{Inherited}: V41--V48
(common), V49--V64 (this file), the concurrent V49--V56 as listed, in
particular the joint-regulator theorem with ledger $\Theta_{\gamma,L}$ and
the bare calibration $\gamma=4/g_0^2$. \emph{Conditional}: the Feshbach
reduction of V64 under (H1), (H2); the physical reading of any small-volume
statement. \emph{Open}: N2 in the form \eqref{f16v65:N2}; the quadratic
Harnack lemma N2$'$; local-source noncollapse and convergence; the
complement local-Gram contraction; scale-to-scale RG transport; the
continuum construction; the OS Hamiltonian mass gap.

\subsection{Order of attack for V66--V69}

V66: the Gaussian maximum theorem for the hard link field with the actual
covariance of V63, giving the event $G_{L,\gamma}$ and the sharp scale
$r(L,\gamma)$ without a union bound. V67: Wilson-word concentration, the
exponential-moment and Schur-norm control of the nonlinear remainders of the
plaquette words in $b_\gamma$ and of the kinetic words on $G_{L,\gamma}$,
extending the line-chaos machinery of V61--V63 to the finite list of words in
the local comparison. V68: assembly of the native concentrated-kernel theorem
\eqref{f16v65:N2} on the largest domain the two previous results allow,
compared with the concurrent V54 ledger. V69: handoff tests, in particular an
explicit uniform rate for the scalarization error $\epsilon_{L,\gamma}$ of the
concurrent V56 from the chaos bounds for $O(L)$-edge writers, and the
hard-sector propagation between separated local writers from the covariance
decay. V70 is the ten-version review and the next companion checkpoint.

\subsection{Verification and preservation}

The new diagnostic is
\begin{center}\small\path{scripts/verify_ym_f16_v65_transport_identity.py}.\end{center}
Its hashed content is the companion inventory including the concurrent file,
and the three booleans of the cyclic-group check of
Theorem~\ref{f16v65:transport}. It has 5557 bytes and SHA--256
\begin{center}\scriptsize\ttfamily 877BA5D7498FEC8DF48B2EEB48E778AE0115A86C1049D7E12E76A4408E0B7443.\end{center}
Its canonical hashed payload has SHA--256
\begin{center}\scriptsize\ttfamily B35BD073441F1366C211E68619EF0954D6045668AF40F15D29A9803DFB5909DB.\end{center}
Only this terminal insertion and the literal title version differ from V64;
removing them recovers its bytes exactly.

\begin{firewall}
V65 records the companion checkpoint with the concurrent lineage, the
redirection of this lineage to regulator-uniform gauge-invariant estimates
in support of the local-vacuum programme, the five-version review, the
classification of the obstructions at the boundary of the small-volume
regime, and proves the transport identity of the Perron vector. It does not
prove N2, the quadratic Harnack lemma, any local-Gram statement, a
renormalized soft coupling, or an interacting continuum mass gap, and it
does not identify the toron gap with a physical local mass.
\end{firewall}
% END F16 V65 APPEND

% BEGIN F16 V66 APPEND -- 2026-09-18
% Insert before the final document-ending command of cumulative V65.
% This insertion leaves every predecessor body and the native operator unchanged.

\clearpage
\section{V66: the Gaussian maximum of the hard field and the kinetic kernel as a mean-one local perturbation}
\label{f16v66:context}

This is the first of the N2 iterations. It proves the two primitive
statements that any concentrated-kernel theorem must start from. First, the
Gaussian maximum theorem: under the hard Gaussian ground the largest link
angle over all $9V$ colour components is of order
$\sigma\sqrt{2\log(18V)/\gamma}$ with $\sigma^2=C_3(0)=0.103$, with a
Gaussian tail beyond that scale, so that the sup-norm event
$G_{L,\gamma}$ of V64 has an explicit tail probability and, on it, the
kinetic sandwich has variance error $O((\log V+\lambda)/\gamma)$. Second, the
Wilson one-link kernel is exactly a Gaussian-in-distance reference kernel
times a mean-one factor $\rho_\gamma(\theta)$ whose logarithm has mean
$-5/(8\gamma^2)$ and variance $5/(4\gamma^2)$ under the reference; the
native transfer is therefore the reference transfer times
$\exp\sum_e\log\rho_\gamma(\theta_e)$, a product of $3V$ local mean-one
factors. The deterministic budget for this perturbation is $3V/\gamma$,
which is the linear volume cost of the concurrent V53; under the frozen
hard Gaussian chain the connected budget is $\sqrt V/\gamma$ times an
explicit constant, because the one-step displacement covariance is
square-summable over the lattice. This is the first instance of the rule
that extensive error budgets of the form $V\times$small are artifacts of
summing before using locality, and it fixes the norm question that V67 and
V68 must answer: the perturbation is small in $L^2$ of the hard band but not
in operator norm, and the comparison must be organized through local
connected quantities.

\subsection{Programme clarification}

The redirection of V65 is completed by a clarification of the endpoint.
This lineage is a genuine parallel perturbative Yang--Mills and ultraviolet
renormalization-group programme, not only a supplier of lemmas to the
local-vacuum programme. Its endpoint is to construct and control the
renormalized local theory from the bare Wilson transfer down to the first
genuinely nonperturbative scale, by an explicit scale-to-scale theorem of
the form $\mathcal F_{j+1}(z)=\mathcal R_j\mathcal F_j(z)\mathcal R_j^*+\mathcal E_j(z)$,
$\|\mathcal E_j\|_{\rm local}\le\varepsilon_j$, $\sum_{j\le j_{\rm pert}}\varepsilon_j<\infty$,
tracking the running coupling, the kinetic normalization, the generated
local gauge-invariant operators, the vacuum-energy subtraction separately
from spectral ratios, the remainder operators in a locality norm, the gauge
and positivity structures, the lattice-to-physical scale relation, and the
domain of validity; at the first scale where perturbative control ceases,
the controlled local data are handed to the local-vacuum programme. The
order of attack is: N2; a uniform local block or Feshbach statement; one
explicit coarse-graining step for the Wilson transfer; its locality and
remainder; the transformed coupling; iteration while the coupling is
small; the first scale where control fails; export. The toron results
remain ultraviolet diagnostics. Every version continues to state what is
proved, inherited, conditional and open.

\subsection{Dependencies}

V45's hard Gaussian ground and mode variances, V52's hard spectrum
($r_+=e^{-\xi_{\min}}$), V61's $C_3(0)$, V63's decay and the bound
$C_3(0)\le\kappa_2$, V64's sandwich and $\delta(r)$, V65's status
discipline. From the concurrent file: V51--V53 (Wilson--heat comparison at
linear volume cost) as the deterministic baseline.

\subsection{The Gaussian maximum of the hard link field}

Let $y=\varepsilon u$ be the link-angle field of V45's hard Gaussian ground,
with $u\in\mathcal Y$ the transverse modes; every colour component $y_e^a$
of every link is a centred Gaussian of variance $\sigma^2/\gamma$,
$\sigma^2=C_3(0)=V^{-1}\sum_{k\ne0}P_{\mu\mu}(k)/(2\omega_k)$, the same for
all links by translation and cubic symmetry.

\begin{theorem}[Gaussian maximum theorem]
\label{f16v66:maximum}
Let $M=\max_{e,a}|y_e^a|$ over the $3V$ links and three colours. Then
\begin{equation}
 \boxed{\quad
 E_\varphi[M]\le\sigma\sqrt{\frac{2\log(18V)}\gamma},\qquad
 \mathbb P_\varphi\bigl(M>E_\varphi[M]+t\bigr)\le e^{-\gamma t^2/(2\sigma^2)},\qquad
 \mathbb P_\varphi\Bigl(M>\sigma\sqrt{\tfrac{2(\log(18V)+\lambda)}\gamma}\Bigr)\le e^{-\lambda},
 \quad}
 \label{f16v66:max-bound}
\end{equation}
and $\sigma^2\le\kappa_2$ uniformly in $L$ (V63), numerically $\sigma^2=$ $0.1010, 0.1027, 0.1031, 0.1033$
at $L=8,16,32,64$. Conversely $E_\varphi[M]\ge c\,\sigma\sqrt{\log(V/d_0^3)/\gamma}$
with $d_0$ the distance beyond which the correlation $C_3(x)/C_3(0)$ is at
most $\frac12$, $d_0\le\sqrt{2(\kappa_0+\kappa_1\log L)/\kappa_3}$ by V63, so
the logarithm is necessary.
\end{theorem}
\begin{proof}
The first bound is the standard estimate $E\max_{i\le N}X_i\le\sigma_\gamma\sqrt{2\log N}$
for $N=18V$ centred Gaussians $\pm y_e^a$ of variance $\sigma_\gamma^2=\sigma^2/\gamma$,
which uses no independence. The second is the Borell--TIS inequality for
the supremum of a centred Gaussian process with maximal variance
$\sigma_\gamma^2$. The third is the union bound
$\mathbb P(M>t)\le18V\cdot\frac12e^{-\gamma t^2/(2\sigma^2)}$ at
$t^2=2\sigma^2(\log(18V)+\lambda)/\gamma$. For the lower bound, the links at
mutual distance at least $d_0$ form a set of $N'\ge V/d_0^3$ variables with
pairwise correlations at most $\frac12$, so $E|y_e-y_{e'}|^2\ge\sigma_\gamma^2$
for distinct pairs and Sudakov's minoration gives $E\max\ge c\sigma_\gamma\sqrt{\log N'}$.
\end{proof}

\begin{proposition}[The sup-norm event and its cost]
\label{f16v66:event}
For $\lambda\ge0$ put $r_{L,\gamma}(\lambda)=\sigma\sqrt{2(\log(18V)+\lambda)/\gamma}$
and $G_{L,\gamma}(\lambda)=\{M\le r_{L,\gamma}(\lambda)\}$. Then
$\mathbb P_\varphi(G^c)\le e^{-\lambda}$, and on pairs of configurations in
$G$ the sandwich of Theorem~\ref{f16v64:sandwich} holds with
\begin{equation}
 \delta\bigl(r_{L,\gamma}(\lambda)\bigr)\le\frac{4\sigma^2(\log(18V)+\lambda)}{3\gamma}+\frac{4\sigma^4(\log(18V)+\lambda)^2}{5\gamma^2}
 =\frac{0.138(\log(18V)+\lambda)}\gamma\bigl(1+o(1)\bigr).
 \label{f16v66:delta-event}
\end{equation}
Numerically, $r=$ $0.1781, 0.0563, 0.0178$ and $\delta(r)=$ $0.0214, 0.0021, 0.0002$ at $L=64$ and
$\gamma=10^2,10^3,10^4$ ($\lambda=0$). Thus the chart error on the event is
logarithmic in the volume and vanishes in every regime with $\gamma\gg\log V$,
in particular far beyond the small-volume regime; the event is defined in
V45's transverse coordinates and is gauge compatible in the sense that the
transverse representative is unique on V50's tube of the concurrent file;
its uniqueness on the sup-norm region is not proved here.
\end{proposition}
\begin{proof}
Insert $r$ into $\delta(r)\le\frac23r^2+\frac15r^4$ of Proposition~\ref{f16v64:distance}.
\end{proof}

\subsection{The Wilson kernel as a Gaussian-in-distance reference times a mean-one factor}

Let $k_\gamma(\theta)=N_\gamma^{-1}e^{-\gamma(1-\cos\theta)}$ be the one-link
Wilson kernel and define the reference kernel
$g_\gamma(\theta)=Z_\gamma^{-1}e^{-\gamma\theta^2/2}$, both probability
densities with respect to Haar measure, whose radial density is
$\frac2\pi\sin^2\theta\,d\theta$; $\theta=d(U_e,U'_e)$. Write
${\rm flat}=(2\pi/\gamma)^{3/2}/(2\pi^2)$.

\begin{theorem}[Mean-one factorization]
\label{f16v66:factorization}
$k_\gamma=g_\gamma\rho_\gamma$ with
$\rho_\gamma(\theta)=(Z_\gamma/N_\gamma)\exp[\gamma(\theta^2/2-1+\cos\theta)]$
and $\int g_\gamma\rho_\gamma\,dU=1$. With exact rational Laplace
coefficients,
\begin{equation}
 \frac{N_\gamma}{\rm flat}=1-\frac3{8\gamma}-\frac{15}{128\gamma^2}+O(\gamma^{-3}),\qquad
 \frac{Z_\gamma}{\rm flat}=1-\frac1\gamma+\frac2{3\gamma^2}+O(\gamma^{-3}),
 \label{f16v66:NZ}
\end{equation}
\begin{equation}
 E_{g}[\log\rho_\gamma]=-\frac5{8\gamma^2}+O(\gamma^{-3}),\qquad
 \operatorname{Var}_{g}[\log\rho_\gamma]=\frac5{4\gamma^2}+O(\gamma^{-3}),
 \label{f16v66:logrho}
\end{equation}
and on $\{\theta\le2r\}$, $|\log\rho_\gamma(\theta)-\log(Z_\gamma/N_\gamma)|\le\frac{2\gamma r^4}3$.
Numerically $\gamma^2\operatorname{Var}_g[\log\rho]=$ $1.166, 1.228, 1.245$ at $\gamma=50,200,800$.
\end{theorem}
\begin{proof}
The factorization is the definition and the normalization is
$\int k_\gamma=1$. For the expansions write $\theta^2=u/\gamma$; the Haar
density is $\theta^2(\sin^2\theta/\theta^2)$ with
$\sin^2\theta/\theta^2=\sum_k(-1)^k2^{2k+1}\theta^{2k}/(2k+2)!$, the Wilson
exponent is $\gamma(\theta^2/2-1+\cos\theta)=\sum_{k\ge2}(-1)^k\gamma^{1-k}u^k/(2k)!$,
and the three-dimensional Gaussian moments are $E[u^k]=(2k+1)!!$; the
verifier expands all four quantities in $1/\gamma$ in exact rational
arithmetic and hashes the coefficients. The first coefficient of
$\log\rho$ vanishes, as it must since $E_g[\rho]=1$ forces
$E\log\rho=-\frac12\operatorname{Var}\log\rho+O(\gamma^{-3})$, which the
coefficients confirm: $-\frac58=-\frac12\cdot\frac54$. The bound on the event
is $\gamma\sum_{k\ge2}\theta^{2k}/(2k)!\le\gamma\theta^4/24\cdot(1+\theta^2/30+\cdots)\le\frac{2\gamma r^4}3$
for $\theta\le2r\le1$.
\end{proof}

\begin{theorem}[The native transfer as a local mean-one perturbation of the reference transfer]
\label{f16v66:local-perturbation}
Let $T^{\rm ref}_\gamma=b_\gamma\,g_\gamma^{\otimes3V}\,b_\gamma$, with the
original magnetic factor. Then, exactly,
\begin{equation}
 \boxed{\quad
 T_\gamma(U,U')=T^{\rm ref}_\gamma(U,U')\exp\Bigl[\sum_e\log\rho_\gamma\bigl(d(U_e,U'_e)\bigr)\Bigr],
 \quad}
 \label{f16v66:structure}
\end{equation}
where the perturbation is a sum of $3V$ terms, each a function of one link
displacement, each of mean one under the reference one-link kernel. The
deterministic budget is $\|g^{\otimes3V}-k^{\otimes3V}\|\le3V\sup_\ell|g_{\gamma,\ell}-k_{\gamma,\ell}|=O(V/\gamma)$
by telescoping over the Fourier coefficients. Under the frozen hard
Gaussian chain of V45, whose stationary two-step law has the displacement
covariance $E[(y_e-y'_e)^a(y_{e'}-y'_{e'})^b]=\delta^{ab}D(e-e')/\gamma$,
$D(x)=V^{-1}\sum_{k\ne0}e^{ik\cdot x}P_{\mu\mu}(k)(1-e^{-\xi_k})/\omega_k$,
the connected budget is
\begin{equation}
 \operatorname{Var}\Bigl[\sum_e\log\rho_\gamma(\theta_e)\Bigr]
 \le\frac{K}{\gamma^2}\,V\,D(0)^4s_D+O(V\gamma^{-3}),\qquad
 s_D=\sum_x\Bigl(\frac{D(x)}{D(0)}\Bigr)^2,
 \label{f16v66:connected}
\end{equation}
with an absolute $K$, $D(0)=$ 0.1667 and $s_D=$ 1.329 (numerically constant in
$L$). Hence the standard deviation of the perturbation exponent is
$O(\sqrt V/\gamma)$, against a mean $3VE_g[\log\rho]=-15V/(8\gamma^2)+\dots$
which is a constant, while its deterministic size is $O(V/\gamma)$.
\end{theorem}
\begin{proof}
\eqref{f16v66:structure} is the product of the one-link factorizations
with the magnetic factors untouched. For the connected budget, on the
frozen chain the displacements are jointly Gaussian three-vectors with the
stated covariance, the per-mode one-step correlation being $e^{-\xi_k}$ by
V52's Mehler structure; $\log\rho_\gamma(\theta_e)=\log(Z/N)+\gamma|\theta_e|^4/24+O(\gamma|\theta_e|^6)$,
and Wick's theorem gives, for jointly Gaussian three-vectors $X,Y$ with
$\operatorname{Cov}(X^a,Y^b)=\delta^{ab}D$,
$\operatorname{Cov}(|X|^4,|Y|^4)\le K'D(0)^2D^2$, since the pairings with one
cross pair vanish by parity and those with two or four cross pairs carry
$D^2$; summing over the $3V$ pairs of links with $\sum_x D(x)^2=D(0)^2s_D$
gives \eqref{f16v66:connected} with $K=9K'/576$. The deterministic
telescoping bound is the concurrent V51's mechanism.
\end{proof}

Two readings. First, in the $L^2$ sense of the hard band the kinetic
nonlinearity is a perturbation of relative size $\sqrt V/\gamma$, which is
$o(1)$ exactly for $L\ll\gamma^{2/3}$: the small-volume regime is also the
regime in which the Wilson kinetics is a small perturbation of Gaussian
kinetics in the connected sense, and $V/\gamma$ is an artifact. Second, the
effect on spectral ratios is smaller still: the level shift at first order
is $\sum_e(\langle\psi_j,\log\rho_e\,\psi_j\rangle-\langle\psi_0,\log\rho_e\,\psi_0\rangle)$,
a sum of $3V$ local terms each depending on the soft state only through the
soft displacement of one link, of variance $O(1/(\gamma V))$, so each term is
$O(1/(\gamma V))$ and the sum $O(1/\gamma)$, relatively $O(L\gamma^{-2/3})=O(z)$
to the gap $h$. That first-order computation is conditional on the
perturbation theory being justified on the hard band, which is exactly what
the operator-norm question of N2 asks; it is recorded as the expected size,
consistent with V53's Gram correction of the same order.

\subsection{Status}

\emph{Proved}: Theorems \ref{f16v66:maximum}, \ref{f16v66:factorization},
\ref{f16v66:local-perturbation} and Proposition~\ref{f16v66:event}, the
last two in the frozen hard Gaussian chain where the covariance is
Gaussian. \emph{Inherited}: V45's hard Gaussian, V52's hard spectrum,
V63's decay, V64's sandwich, and the concurrent V50 tube for gauge
compatibility of the transverse representative on the $\ell^2$ tube.
\emph{Conditional}: the first-order level-shift size $O(z)$; the uniqueness
of the transverse representative on the sup-norm region.
\emph{Open}: the operator-norm or form-norm insertion of the local
perturbation on the hard band (N2 proper), the plaquette words of the
magnetic factor beyond the quadratic order (V67), and the assembly (V68).

\subsection{Verification and preservation}

The new diagnostic is
\begin{center}\small\path{scripts/verify_ym_f16_v66_gaussian_maximum.py}.\end{center}
Its hashed exact content is the rational Laplace expansions
\eqref{f16v66:NZ}--\eqref{f16v66:logrho} to second order. Its floating
diagnostics are the numerical one-link integrals, the scales $r_{L,\gamma}$
and $\delta(r)$, and $D(0)$, $s_D$ by fast Fourier transform. It has
9990 bytes and SHA--256
\begin{center}\scriptsize\ttfamily 8268E837199A9F48DA7EEAF14014AC03E0ADEA44C0A4EA49FC1F1825FE7D7230.\end{center}
Its canonical hashed payload has SHA--256
\begin{center}\scriptsize\ttfamily E4394FBEE9D39C936D6CB29893FA7E17BE99EE15AB9BA513E1978E5638141859.\end{center}
Only this terminal insertion and the literal title version differ from V65;
removing them recovers its bytes exactly. The next companion checkpoint and
ten-version review is V70.

\begin{firewall}
V66 proves the Gaussian maximum theorem for the hard link field with its
tail and the explicit sup-norm event, the exact mean-one factorization of
the Wilson one-link kernel with its Laplace expansions, and the local
mean-one structure of the native transfer with its connected budget in
the frozen hard Gaussian chain. It does not prove an operator-norm
comparison of the native and reference transfers, does not prove
uniqueness of the transverse gauge on the sup-norm region, and does not
bear on the physical continuum trajectory or an interacting continuum mass
gap.
\end{firewall}
% END F16 V66 APPEND

% BEGIN F16 V67 APPEND -- 2026-09-18
% Insert before the final document-ending command of cumulative V66.
% This insertion leaves every predecessor body and the native operator unchanged.

\clearpage
\section{V67: the cubic vertex of the magnetic factor and the norm obstruction of N2}
\label{f16v67:context}

V66 treated the kinetic side of the native transfer as a product of local
mean-one factors. This note does the same for the magnetic factor
$b_\gamma=e^{-\gamma S/2}$ and reaches a sharper conclusion, negative for one
formulation of N2 and decisive for the right one. Each plaquette word is,
exactly to third order in the four link fields around it,
$1-\frac12\Tr U_p=\frac12|F_p|^2+F_p\cdot B_p+R_4$, with $F_p$ the linear
field strength and $B_p$ the sum of six cross products, so that
$\gamma S/2=\gamma S_2/2+\sum_p\xi_p+\gamma R_4/2$ with the frozen Gaussian
quadratic part $S_2$ and the three-gluon vertex $\xi_p=\frac\gamma2F_p\cdot B_p$.
The vertex is Wick-ordered under the hard Gaussian by colour isotropy, has
mean zero, and its total has variance exactly $\frac34v_3V/\gamma$ with an
explicit lattice constant $v_3=0.047$. Hence the perturbation
$e^{-\sum_p\xi_p}$ of the frozen transfer has $L^2$ size $0.19\sqrt{V/\gamma}$
on the hard band, from below as well as from above: a core-wise $L^2$
comparison between the native and the frozen transfer, in the norm of
V45 and of the concurrent V54, is impossible beyond $V\lesssim\gamma$, i.e.\
$L\lesssim\gamma^{1/3}$, and the ledger term $L^3\gamma^{-1}$ of the concurrent
V54 is sharp for that norm. This is an operator-norm failure in the
classification of V65, and it is the first genuine obstruction: N2 must be
formulated in connected norms, where the same vertex has a per-plaquette
activity $O((\log V)^{3/2}\gamma^{-1/2})$ on the sup-norm event and extensive
cumulants with finite local densities, which is the setting of a
transfer-matrix cluster expansion and, in the language of V66's programme
clarification, of the first renormalization step.

\subsection{Dependencies}

V44's constant-field identity, V45's frozen hard Gaussian and its magnetic
expansion, V53's transverse waves and quaternion tools, V61's exact
polynomial engine, V62's Wick-ordering argument, V63's decay, V66's
sup-norm event and mean-one factorization. From the concurrent file: the
V54 ledger as the object being sharpened.

\subsection{The plaquette word to third order}

Let the four links of the plaquette $p$ in the $(\mu,\nu)$ plane at $x$ be
$y_1=y_\mu(x)$, $y_2=y_\nu(x+\hat\mu)$, $y_3=y_\mu(x+\hat\nu)$, $y_4=y_\nu(x)$,
imaginary quaternions, so that $U_p=e^{y_1}e^{y_2}e^{-y_3}e^{-y_4}$. Put
\begin{equation}
 F_p=y_1+y_2-y_3-y_4,\qquad
 B_p=y_1\times y_2-y_1\times y_3-y_1\times y_4-y_2\times y_3-y_2\times y_4+y_3\times y_4 .
 \label{f16v67:FB}
\end{equation}

\begin{proposition}[Exact third-order structure of the plaquette word]
\label{f16v67:word}
For imaginary quaternions $[a,b]=2a\times b$, and
$\log U_p=F_p+B_p+O(|y_p|^3)$, so that
\begin{equation}
 \boxed{\quad 1-\tfrac12\Tr U_p=\tfrac12|F_p|^2+F_p\cdot B_p+R_{4,p},\qquad
 |R_{4,p}|\le K|y_p|^4\ \ (|y_p|\le1),\quad}
 \label{f16v67:S3}
\end{equation}
with $|y_p|=\max_i|y_i|$ and an absolute $K$. The verifier checks the
commutator identity and the degree-one, degree-two and degree-three parts of
$U_p$ as exact polynomial identities in the twelve components, and the
fourth-order remainder by quaternion finite differences (ratio 15.93
under halving).
\end{proposition}
\begin{proof}
The Baker--Campbell--Hausdorff formula for four factors gives the first two
orders of $\log U_p$ as the sum of the exponents and half the sum of the
ordered commutators, which is \eqref{f16v67:FB} after $[a,b]=2a\times b$;
then $\frac12\Tr e^X=\cos|X|=1-\frac12|X|^2+O(|X|^4)$ with
$|X|^2=|F_p|^2+2F_p\cdot B_p+O(|y|^4)$. The remainder is a smooth function
of the twelve components vanishing to fourth order.
\end{proof}

Summing over plaquettes, $S=S_2+S_3+R_4$ with $S_2=\frac12\sum_p|F_p|^2$ the
quadratic form of V45's hard Hessian, $S_3=\sum_pF_p\cdot B_p$ the
three-gluon vertex and $R_4=\sum_pR_{4,p}$. On the sup-norm event
$G_{L,\gamma}(\lambda)$ of V66, with $r^2=2\sigma^2(\log(18V)+\lambda)/\gamma$,
every plaquette obeys $|\gamma F_p\cdot B_p/2|\le6\gamma r^3$ and
$|\gamma R_{4,p}/2|\le\frac{K}2\gamma r^4$: the \emph{local} activities are
$O((\log V+\lambda)^{3/2}\gamma^{-1/2})$ and $O((\log V+\lambda)^2\gamma^{-1})$,
small throughout every polynomial regime, while their sums over $3V$
plaquettes are not.

\subsection{The Wick-ordered vertex and its two-sided fluctuation}

Let $C_{\mu\nu}(x)=V^{-1}\sum_{k\ne0}e^{ik\cdot x}P_{\mu\nu}(k)/(2\omega_k)$ be
the covariance, in units of $1/\gamma$, between a $\mu$-link at $0$ and a
$\nu$-link at $x$ under V45's hard Gaussian, with the transverse projector
$P_{\mu\nu}=\delta_{\mu\nu}-c_\mu\bar c_\nu/\hat k^2$, $c_\mu=e^{-ik_\mu}-1$;
$C_{11}$ is V63's $C_3$.

\begin{theorem}[Wick-ordered cubic vertex and its variance]
\label{f16v67:vertex}
Under the hard Gaussian, $\xi_p=\frac\gamma2F_p\cdot B_p$ has mean zero and
vanishing partial contractions, $E[\xi_p\,y^c_v]=0$ for every link component;
hence $E[\xi_p\xi_{p'}]$ is a sum over complete cross pairings only, decays
like the cube of the covariance, and
\begin{equation}
 \boxed{\quad
 \operatorname{Var}_\varphi\Bigl[\sum_p\xi_p\Bigr]=\frac{3V}{4\gamma}\,v_3(L),\qquad
 v_3(L)=\sum_{x\in(\mathbb Z/L)^3}E\bigl[(F_0\cdot B_0)(F_x\cdot B_x)\bigr]\gamma^3,
 \quad}
 \label{f16v67:variance}
\end{equation}
the sum over the plaquettes of one orientation, absolutely convergent
uniformly in $L$ by V63; numerically $v_3(8)=$ 0.0441 and $v_3(16)=$ 0.0468,
with the on-site term $\operatorname{Var}=$ 0.0380 and box partial sums
converging at the third shell. Consequently the standard deviation of
$\gamma S_3/2$ on the hard band is $\sqrt{3v_3/4}\sqrt{V/\gamma}=$ 0.187$\sqrt{V/\gamma}$.
\end{theorem}
\begin{proof}
The vertex is $\epsilon_{abc}F_p^ay_i^by_j^c$ summed over the six ordered
pairs; contracting $F^a$ with $y_i^b$ or $y_j^c$, or $y_i^b$ with $y_j^c$,
produces $\delta$ against $\epsilon$, which vanishes. A cubic polynomial
with vanishing partial contractions is Wick-ordered, so for two copies only
the six complete cross pairings survive (V62, Proposition~\ref{f16v62:wick}),
each a product of three covariances between links of the two plaquettes,
bounded by $K(\kappa_0+\kappa_1\log L)^3/|x|_\infty^6$ at distance $x$ by
V63's decay applied to all nine components $C_{\mu\nu}$, which obey the same
summation-by-parts bound. The sum over $x$ converges, and translation
invariance gives \eqref{f16v67:variance} with the factor $\frac\gamma2$
squared against three orientations of $V$ plaquettes each. The verifier
computes $C_{\mu\nu}$ by fast Fourier transform, the on-site variance and
the shell sums by an exact Wick enumeration of the fifteen pairings for each
pair of monomials of the cubic polynomial, and reports the convergence.
\end{proof}

\subsection{The norm obstruction of N2}

\begin{theorem}[The core-wise $L^2$ comparison cannot pass $V\asymp\gamma$]
\label{f16v67:obstruction}
Let $T^{\rm frz}=b_2\,k^{\otimes3V}b_2$, $b_2=e^{-\gamma S_2/2}$, be the transfer
with the frozen quadratic magnetic factor, and let $\Psi=\varphi\otimes f$ be
a hard-ground vector with a fixed soft test $f$. Then, on the sup-norm event
and in the frozen hard Gaussian model,
\begin{equation}
 \frac{\|(T_\gamma-c\,T^{\rm frz})\Psi\|}{\|T^{\rm frz}\Psi\|}\ \ge\ c'\sqrt{\frac V\gamma}\,(1+o(1))
 \qquad\text{for every constant }c,
 \label{f16v67:lower}
\end{equation}
with $c'=\sqrt{3v_3/4}=$ 0.187, since $T_\gamma\Psi=T^{\rm frz}\bigl[e^{-\sum_p\xi_p(U)-\sum_p\xi_p(U')}(1+O(\gamma r^4V))\bigr]\Psi$
and a log-normal multiplier of variance $s^2$ has relative $L^2$ distance at
least $s(1+o(1))$ from every constant. Consequently no core-wise $L^2$
expansion of the type of V45's Theorem~\ref{f16v45:strong-expansion}, or of
the concurrent V54, can hold with errors $o(1)$ beyond $V=O(\gamma)$, and the
term $L^3\gamma^{-1}$ of the concurrent ledger $\Theta_{\gamma,L}$ is sharp for
that norm up to the constant.
\end{theorem}
\begin{proof}
The multiplier is $e^{-Z}$ with $Z=\sum_p[\xi_p(U)+\xi_p(U')]$, a centred
Gaussian in the frozen model with variance $s^2\ge\frac34v_3V/\gamma$; for any
constant $c$, $E[(e^{-Z}-c)^2]\ge\operatorname{Var}(e^{-Z})=e^{s^2}(e^{s^2}-1)\ge s^2$,
and the fourth-order remainder is $O(\gamma r^4V)=O((\log V)^2V/\gamma)$ in
$L^\infty$ on the event, smaller than $s$ when $V\lesssim\gamma$ and irrelevant
to the lower bound otherwise. The vectors $T^{\rm frz}\Psi$ carry the frozen
Gaussian law, on which the variance computation of
Theorem~\ref{f16v67:vertex} applies to both endpoints.
\end{proof}

The reading is the one anticipated by V66 and V65. In the $L^2$ norm of the
hard band the cubic vertex is not small in large volume, for the same
reason that the Boltzmann weight of any weakly coupled field theory is not
close to a Gaussian in $L^2$: its logarithm has extensive variance. What is
small is the vertex \emph{per plaquette} on the sup-norm event, and what is
finite is the \emph{connected} content: $E[\sum_p\xi_p]=0$, the second
cumulant is $\frac34v_3V/\gamma$ with a finite local density, and by the
Wick-ordered structure and V63's decay every higher cumulant of
$\sum_p\xi_p$ is extensive with a finite density, i.e.\ the vertex field has
the cluster property. Spectral ratios and local observables of the transfer
depend on these cumulant densities, not on the $L^2$ distance. N2 is
therefore reformulated: the native concentrated-kernel theorem must compare
$T_\gamma$ and its effective version through a connected norm, i.e.\ through
the generating function of local cumulants along the Doob chain, and the
extensive first cumulant is the vacuum-energy subtraction that V66's
programme clarification requires to be tracked separately from spectral
ratios. This is exactly the structure of a transfer-matrix cluster
expansion, whose first step is the one-block elimination of the
perturbative programme. The regime is then not limited by
$V\lesssim\gamma$ but by the convergence of that expansion, which requires
the per-plaquette activity $6\gamma r^3\asymp(\log V)^{3/2}\gamma^{-1/2}$ and
its connected iterates to be small: a condition of the form
$\gamma\gg\log^3V$, compatible with every polynomial trajectory and with
the small-volume treatment of the zero modes as long as $L\ll\gamma^{2/3}$.

\subsection{Status}

\emph{Proved}: Proposition~\ref{f16v67:word}, Theorem~\ref{f16v67:vertex}
in the frozen hard Gaussian, Theorem~\ref{f16v67:obstruction} in the same
model on the sup-norm event. \emph{Inherited}: V45's frozen Gaussian, V63's
decay extended to $C_{\mu\nu}$, V66's event, the concurrent V54 ledger.
\emph{Conditional}: the cluster property of all cumulants of the vertex
field (stated, its proof for cumulants of order $\ge3$ follows the same
Wick-ordered counting and is not written). \emph{Open}: N2 in connected
form, the one-block elimination, the running coupling.

\subsection{Verification and preservation}

The new diagnostic is
\begin{center}\small\path{scripts/verify_ym_f16_v67_cubic_vertex.py}.\end{center}
Its hashed exact content is the six polynomial identities of
Proposition~\ref{f16v67:word} and the Wick ordering of $F_p\cdot B_p$ under
an arbitrary colour-diagonal covariance, all in exact rational
arithmetic. Its floating diagnostics are the remainder scaling, the
transverse covariance tensor and the vertex variance and shell sums. It has
10357 bytes and SHA--256
\begin{center}\scriptsize\ttfamily EFCB44CEB09BA149F4D735F0E4629F15CAE8132AFD11B4ABFBCC7B281404BCD8.\end{center}
Its canonical hashed payload has SHA--256
\begin{center}\scriptsize\ttfamily 799EB73AFAD3675B507400FF58BE997063AFAD60826FA9612908CA9FD973662C.\end{center}
Only this terminal insertion and the literal title version differ from V66;
removing them recovers its bytes exactly. The next companion checkpoint and
ten-version review is V70.

\begin{firewall}
V67 proves the exact third-order structure of the plaquette word, the
Wick ordering and the extensive variance of the three-gluon vertex under
the hard Gaussian, and the lower bound showing that a core-wise $L^2$
comparison of the native and frozen transfers fails beyond $V\asymp\gamma$.
It does not prove a connected comparison, a cluster expansion, a
renormalization step, or anything about the physical continuum trajectory
or an interacting continuum mass gap.
\end{firewall}
% END F16 V67 APPEND

% BEGIN F16 V68 APPEND -- 2026-09-18
% Insert before the final document-ending command of cumulative V67.
% This insertion leaves every predecessor body and the native operator unchanged.

\clearpage
\section{V68: exact horizontal reduction of the based-slice frozen transfer and the cancellation of the extensive density terms}
\label{f16v68:context}

The native chain reads the transfer in based-slice coordinates, and its
uniform versions pay the slice density $J_L(s)$ as an error: the concurrent
V50 bounds $|\log J_L(s)/J_L(0)|$ by $C\|s\|_2^2$, which for the constant
soft field $s=\mathsf Ka$ is $CV|a|^2$, relatively $L^2\gamma^{-1/3}$ at the
critical scale. This note shows that the density is not an error term.
Once the gauge integral of the second argument is performed, the frozen
Gaussian transfer on the slice has the \emph{horizontally projected} kinetic
metric $M_h$ and a normalization factor $\det M_h^{1/2}$, and the density
factorizes exactly as $J_L^2=\det(G^{\mathsf T}G)\det M_h$ into the orbit volume,
which cancels against the vertical Gaussian integral, and the horizontal
Jacobian, which cancels against the slice normalization. What remains is
the covariant frozen transfer at the background, whose top is the
twisted-mode zero-point energy of V54, with no extensive $|a|^2$ term: the
$a$-dependence is $m_L|a|^2+O(a^4)$, $m_L=O(L)$, plus the cusp of the lifted
charged constant modes. The statement is exact at the Gaussian level and is
verified on the $L=4$ lattice to finite-difference precision, including the
identity between the generalized eigenvalues of the slice pair
$(H_s,M_h)$ and the horizontal spectrum of the covariant Hessian, and the
identity of the latter with V54's twisted spectrum, which is thereby
confirmed by an independent brute-force computation. Any joint-limit
bookkeeping that treats the density as an error can drop the corresponding
term; the reference operator for N2 is the horizontally reduced one.

\subsection{Dependencies}

V44's based tube and its exact invariant half-density, V45's frozen hard
Gaussian, V54's twisted spectrum $\lambda_L(k,\pm\phi)$ and $\Delta_L(\alpha)$,
V53's quaternion tools. From the concurrent file: the V50 tube and density
theorem as the object being reinterpreted, and the V54 ledger.

\subsection{Setting}

Let $a=(a_\mu)$ be a constant field, $U^{(a)}_e=\operatorname{Exp}(a_\mu)$, and
use the right chart $U_e=U^{(a)}_e\operatorname{Exp}(y_e)$, $y\in\mathbb R^{9V}$,
in which the round link metric is orthonormal at $y=0$. Let $H_c(a)$ be the
Hessian of the Wilson action at $y=0$, $G(a)$ the $9V\times3(V-1)$ frame of
vertical (gauge) directions $y_e=R(U^{(a)}_e)^{\mathsf T}\eta_x-\eta_{x+\hat\mu}$,
$\eta_o=0$, and $P_h(a)$ the orthogonal projector onto their complement. The
based slice is $s=\mathsf Ka+n$, $n\in N=\ker(d^*\otimes I_3)$, the transverse
fields at $a=0$, with orthonormal basis $B$ ($9V\times(6V+3)$); the based
parametrization $\Phi(g,s)=g\cdot\operatorname{Exp}(s)$ has $dH=dg\,J_L(s)\,ds$
with $J_L(s)=|\det[G(s)\,|\,D_sB]|$, where $D_s$ is the right-trivialized
differential of $\operatorname{Exp}$ at $s$, $\operatorname{Exp}(s+v)=\operatorname{Exp}(s)\operatorname{Exp}(D_sv+O(v^2))$.
Let $H_s(a)$ be the Hessian of $S(\operatorname{Exp}(\mathsf Ka+n))$ in $n\in N$,
and put
\begin{equation}
 M(a)=(D_aB)^{\mathsf T}(D_aB),\qquad M_h(a)=(D_aB)^{\mathsf T}P_h(a)(D_aB).
 \label{f16v68:metrics}
\end{equation}

\subsection{The reduction theorem}

\begin{theorem}[Exact horizontal reduction at the Gaussian level]
\label{f16v68:reduction}
Let $K_a$ be the frozen Gaussian transfer at the background $a$, with kinetic
kernel $\prod_eN_\gamma^{-1}e^{-\gamma|y_e-y'_e|^2/2}$ and magnetic factors
$e^{-\gamma y^{\mathsf T}H_c(a)y/4}$ at both ends, acting on gauge-invariant
functions. Then:
\begin{enumerate}
\item $J_L(s)^2=\det\bigl(G(s)^{\mathsf T}G(s)\bigr)\det M_h(s)$, the orbit volume
times the horizontal Jacobian of the slice.
\item The gauge integral of the second argument gives, on the slice,
\[
 \int_{\mathcal G_o}K_a(\operatorname{Exp}(s),g\operatorname{Exp}(s'))\,dg\;J_L(s')
 =c_\gamma\,\det M_h(s')^{1/2}\exp\Bigl[-\frac\gamma2(n-n')^{\mathsf T}M_h(n-n')-\frac\gamma4n^{\mathsf T}H_sn-\frac\gamma4n'^{\mathsf T}H_sn'\Bigr]\bigl(1+O(|n|^3+|n'|^3)\bigr),
\]
with $c_\gamma$ independent of $a$ and $s$: the orbit volume of $J_L$ is
cancelled by the vertical Gaussian integral and the horizontal Jacobian is
what remains.
\item $H_s(a)=(D_aB)^{\mathsf T}H_c(a)(D_aB)$ and $H_c=P_hH_cP_h$, so the
generalized eigenvalues of the pair $(H_s(a),M_h(a))$ are exactly the
eigenvalues of $H_c(a)$ on the horizontal space.
\item Consequently the top of the reduced slice operator is
$\exp[-\sum_j\xi(\nu_j(a))]$ over the horizontal spectrum, the same as the
covariant top, and
\begin{equation}
 \boxed{\quad-\log\frac{\Lambda^{\rm frz}(a)}{\Lambda^{\rm frz}(0)}=\Delta_L(a)=\sum_j\bigl[\xi(\nu_j(a))-\xi(\nu_j(0))\bigr],\quad}
 \label{f16v68:top}
\end{equation}
with no contribution from $\log J_L(\mathsf Ka)/J_L(0)$. For an abelian
constant background the horizontal spectrum is V54's twisted spectrum and
$\Delta_L(a)=m_L|a|^2+\text{(cusp of the lifted charged constant modes)}+O(a^4)$.
\end{enumerate}
\end{theorem}
\begin{proof}
(1) is the block determinant identity $\det[G\,|\,X]^2=\det(G^{\mathsf T}G)\det(X^{\mathsf T}P_hX)$
for the frame $[G|X]$ with $P_h$ the projector orthogonal to the columns of
$G$. (2): write $g\operatorname{Exp}(s')$ in the right chart at
$\operatorname{Exp}(s)$ as $y'=D_s(n'-n)+G(s)\eta+O(\eta^2,\eta n)$ for
$g=\operatorname{Exp}(\eta)$; the kinetic exponent
$-\frac\gamma2|D_s\Delta n+G\eta|^2$ integrated over $\eta$ with the flat
measure of $\mathcal G_o$ at the identity gives
$(2\pi/\gamma)^{3(V-1)/2}\det(G^{\mathsf T}G)^{-1/2}\exp[-\frac\gamma2|P_hD_s\Delta n|^2]$,
since the minimum over $\eta$ of $|v+G\eta|^2$ is $|P_hv|^2$; the magnetic
factor is gauge invariant and equals $e^{-\gamma n'^{\mathsf T}H_sn'/4}$ up to
cubic terms; multiplying by $J_L(s')$ and using (1) leaves $\det M_h(s')^{1/2}$.
(3) is the chain rule for $S\circ\operatorname{Exp}$ at $\mathsf Ka$, where the
first derivative of $S$ vanishes (V45's stationarity), and the gauge
invariance of $S$; if $P_hD_aB$ is invertible onto the horizontal space, which
holds for small $a$ by continuity from $a=0$ where $N$ is horizontal, then
$H_sv=\nu M_hv$ with $w=P_hD_aBv$ reads $H_cw=\nu w$. (4): after the
substitution $m=M_h^{1/2}n$ the reduced kernel is the orthonormal Mehler
kernel with Hessian $M_h^{-1/2}H_sM_h^{-1/2}$ and the factor $\det M_h^{1/2}$
cancels the Jacobian $dn=\det M_h^{-1/2}dm$; the top of a Mehler product is
$\prod_j\bigl(1+\frac{\nu_j}2+\sqrt{\nu_j+\nu_j^2/4}\bigr)^{-1/2}=e^{-\sum\xi(\nu_j)}$
(V41). The last statement is V54's Theorem on the exact hard norm along
constant abelian backgrounds, whose spectrum is the twisted one.
\end{proof}

\subsection{Numerical verification on the $L=4$ lattice}

The verifier assembles $H_c(a)$ and $H_s(a)$ from exact per-plaquette
finite-difference Hessians, the frame $G(a)$, the projector $P_h(a)$, the
differentials $D_a$, the metrics \eqref{f16v68:metrics} and the density, at
$a=\alpha ne_1$ for $\alpha=0,0.02,0.04,0.06$ and a random unit colour vector
$n$. It finds:
\begin{itemize}
\item the horizontal spectrum of $H_c(a)$ equals the twisted multiset
$\{\lambda_L(k,\pm\phi),\lambda_L(k,0)\}$, $\phi=2L\alpha$, two polarizations
each, in number ($378$ hard modes and, for $\alpha>0$, $4$ lifted charged
constant modes) and value (maximal difference 9.6e-08);
\item the generalized eigenvalues of $(H_s(a),M_h(a))$ equal that spectrum
(maximal difference 1.8e-05), while with the full metric $M(a)$ the hard
zero-point energy differs by -32.7$\,\alpha^2$;
\item $\log J_L=\frac12\log\det(G^{\mathsf T}G)+\frac12\log\det M_h$ to 2.1e-13,
with the two parts $+$45.3$\,\alpha^2$ (orbit volume) and -86.6$\,\alpha^2$
(horizontal Jacobian), both extensive at $V=64$, summing to -41.2$\,\alpha^2$;
\item the hard zero-point shift of the reduced slice equals the covariant
one (difference 1.2e-04), whose value is -6.57$\,\alpha^2$, and the four
lifted modes contribute the cusp 7.996$\,\alpha$, i.e.\ $2\alpha$ each,
which is $\xi(4\sin^2(\phi/2L))\approx\phi/L$ per charged constant mode.
\end{itemize}

\subsection{Consequences}

First, the twisted-spectrum theorem of V54 is confirmed by an independent
construction, and the coefficient $m_L$ of V54 is the second-order hard
zero-point shift in the covariant frame, whose limit $4Z_3(1)/(3\pi^2)$ was
found in V58. Second, in any uniform version of V45--V48 the frozen reference
operator should be the horizontally reduced one: its top has no extensive
$|a|^2$ term, and the $a$-dependence of the slice density, which the
concurrent V50 bounds by $C\|s\|_2^2$ and which at the critical scale would
cost $L^2\gamma^{-1/3}$ relative to the soft gap, need not be paid at all.
Which term of the concurrent V54 ledger $\Theta_{\gamma,L}$ this affects
depends on that ledger's bookkeeping, which is not re-derived here; the
statement proved is that the density's extensive parts cancel exactly
against the vertical kinetic integral and the slice normalization. Third,
for N2 the same reduction identifies the correct kinetic metric of the slice
as the horizontal $M_h$, not the full $M$; the difference is extensive and
using $M$ would introduce a spurious $|a|^2$ term of size -32.7$\,\alpha^2$
at $L=4$.

\subsection{Status}

\emph{Proved}: Theorem~\ref{f16v68:reduction} at the Gaussian level, with
all identities verified numerically at $L=4$. \emph{Inherited}: V44's tube
and density, V45's frozen Gaussian and stationarity, V54's twisted
spectrum, the concurrent V50 density theorem. \emph{Conditional}: the
insertion of the reduced reference into a uniform version of V45--V48,
which also needs the connected treatment of V67. \emph{Open}: N2 in
connected form; the one-loop kernel at nonzero momentum (the running
coupling), for which the present reduction fixes the reference.

\subsection{Verification and preservation}

The new diagnostic is
\begin{center}\small\path{scripts/verify_ym_f16_v68_horizontal_reduction.py}.\end{center}
Its hashed content is the five booleans of the $L=4$ verification (twisted
spectrum, generalized eigenvalues, density factorization, reduced
zero-point identity, and the failure of the full metric) at the stated
tolerances; its floating diagnostics are the coefficients above. It has
10737 bytes and SHA--256
\begin{center}\scriptsize\ttfamily A1941FD0DC7E40F8A2FB161C27FEE5863579BD73D5D2B72D39659B6FD5FA91EC.\end{center}
Its canonical hashed payload has SHA--256
\begin{center}\scriptsize\ttfamily C84DDE38D3C5309B3B88D1A2733F817329CED88BCCCB2699EA1637D3D61CEEA5.\end{center}
Only this terminal insertion and the literal title version differ from V67;
removing them recovers its bytes exactly. The next companion checkpoint and
ten-version review is V70.

\begin{firewall}
V68 proves, at the Gaussian level, the exact horizontal reduction of the
based-slice frozen transfer at a constant background, the factorization of
the based density and its exact cancellation, and the identity of the
reduced slice spectrum with the covariant twisted spectrum, verified
numerically at $L=4$. It does not rederive the concurrent V54 ledger, does
not prove a uniform version of V45--V48, and does not bear on the physical
continuum trajectory or an interacting continuum mass gap.
\end{firewall}
% END F16 V68 APPEND

% BEGIN F16 V69 APPEND -- 2026-09-18
% Insert before the final document-ending command of cumulative V68.
% This insertion leaves every predecessor body and the native operator unchanged.

\clearpage
\section{V69: the static one-loop kernel at nonzero momentum and the first appearance of the running coupling}
\label{f16v69:context}

This is the first renormalization-group computation of the lineage. V68
identified the covariant frozen reference at a background: the hard modes
are the horizontal directions at the background, and their zero-point
energy is a gauge-invariant function of the background with no chart
terms. Here the background is a transverse plane wave of momentum
$k=2\pi m/L\,e_1$ and small amplitude $\varepsilon$, and the second-order
shift $\tfrac12\Pi_L(k)\varepsilon^2$ of the hard zero-point energy
$E_0=\tfrac12\sum_j\xi(\nu_j)$ is the static one-loop kernel of the
transfer: the one-loop correction to the magnetic coefficient $\gamma$ at
momentum $k$,
\[
 \gamma_{\rm eff}(k)=\gamma+\delta_L(k),\qquad
 \delta_L(k)=\frac{\Pi_L(k)}{\lambda_k|l|^2}=\frac{2\Pi_L(k)}{\lambda_kV},
\]
per unit background magnetic energy $\gamma S_2=\gamma\lambda_k|l|^2/2$. Three things are established. First,
the object is defined exactly: the horizontal projection at the background
must remove, besides the gauge directions, the low sector consisting of
the nine constant modes, the background's own two modes and the two modes
of its global colour orbit, which become gauge directions at
$\varepsilon\ne0$; with this definition the mode count is continuous in
$\varepsilon$ and the shift is exactly even and quadratic,
$Z(\varepsilon)-Z(0)=\Pi\varepsilon^2(1+O(\varepsilon^2))$ with ratio
$4.000$ under halving. Second, the computation is made exact and scalable:
translation invariance transverse to $k$ block-diagonalizes the projected
Hessian into $9L\times9L$ Hermitian blocks per transverse momentum, so
that $L=64$ is reached with exact diagonalization; at $L=6$ the block
computation reproduces the brute-force $1944$-dimensional values to all
printed digits. Third, the result: $\delta_L(k)$ is negative, of order
unity, and increases with $|k|$, with a logarithmic slope at small $|k|$
of 0.307 against the one-loop asymptotic-freedom coefficient
$11/(3\pi^2)=0.3715$ in the normalization $\gamma=4/g_0^2$, $a_t=a_s$ of
the concurrent V55. The sign and the order of the slope are those of
asymptotic freedom; the fit is from a finite momentum window on finite
lattices, and the identification is recorded as numerical, not proved.

\subsection{Dependencies}

V68's horizontal reduction and its verified twisted spectrum, V45's frozen
Gaussian and stationarity (used only at $\varepsilon=0$), V53's quaternion
tools, V41's Mehler exponent $\xi$. From the concurrent file: V55's
calibration $\gamma=4/g_0^2$ for the expected coefficient. The
background-field one-loop effective action and the Hamiltonian derivation
of asymptotic freedom from the vacuum energy in a background field are
classical; nothing is imported from them except the expected number.

\subsection{Definition of the static one-loop kernel}

Let $b$ be the transverse plane wave $b_e=\cos(kx_1)\,e_{\rm pol}\otimes e_{\rm col}$
on the links in direction ${\rm pol}\in\{2,3\}$, $k=2\pi m/L\,e_1$, and
$U^{(\varepsilon)}=\operatorname{Exp}(\varepsilon b)$. In the right chart at
$U^{(\varepsilon)}$ let $H(\varepsilon)$ be the Hessian of the Wilson action,
$G(\varepsilon)$ the vertical frame of the full gauge group (constant
transformations included), and $E(\varepsilon)$ the low-sector frame: the
nine constant modes, the modes $\cos(kx_1)e_{\rm pol}\otimes e_{\rm col}$ and
$\sin(kx_1)e_{\rm pol}\otimes e_{\rm col}$, and, at $\varepsilon=0$ only, the two
modes $\cos(kx_1)e_{\rm pol}\otimes e_{c}$, $c\ne{\rm col}$. Let $P(\varepsilon)$
project orthogonally onto the complement of the span of $G(\varepsilon)$ and
$E(\varepsilon)$, and let $\nu_j(\varepsilon)$ be the eigenvalues of
$P(\varepsilon)H(\varepsilon)P(\varepsilon)$ above a threshold below the smallest
hard frequency. Define
\begin{equation}
 Z_L(\varepsilon;k)=\sum_j\xi\bigl(\nu_j(\varepsilon)\bigr),\qquad
 \Pi_L(k)=\frac{Z_L(\varepsilon)+Z_L(-\varepsilon)-2Z_L(0)}{2\varepsilon^2},\qquad
 \delta_L(k)=\frac{2\Pi_L(k)}{\lambda_k\,V},
 \label{f16v69:definition}
\end{equation}
with $\lambda_k=4\sin^2(k/2)$ and $|l|^2=\varepsilon^2V/2$ the squared norm of
the background, so that $S_2=\tfrac12\sum_p|F_p|^2=\lambda_k|l|^2/2$ and the
hard zero-point energy is $E_0=\tfrac12Z_L$: the top eigenvalue of the
single-mode transfer $e^{-\lambda q^2/4}e^{\partial^2/2}e^{-\lambda q^2/4}$ is
exactly $e^{-\xi/2}$, $\cosh\xi=1+\lambda/2$ (V70's verifier checks this to
$10^{-6}$), so each hard mode contributes $\xi/2$ to the vacuum energy and
$\xi$ to the excitation gap.

\begin{proposition}[Well-posedness]
\label{f16v69:wellposed}
At $\varepsilon\ne0$ the two colour-orbit modes lie in the span of
$G(\varepsilon)$, since the constant rotation by $\eta$ acts on the background
links as $(R(U_e)^{\mathsf T}-1)\eta=-2\varepsilon\cos(kx_1)\,e_{\rm col}\times\eta+O(\varepsilon^2)$,
while at $\varepsilon=0$ they are physical modes of frequency $\lambda_k$;
excluding them at $\varepsilon=0$ makes the number of retained modes
$6V-10$ for all $\varepsilon$, and then $Z_L(\varepsilon)$ is even and analytic in
$\varepsilon$ near $0$ and $Z_L(\varepsilon)-Z_L(0)=\Pi_L(k)\varepsilon^2+O(\varepsilon^4)$.
The background is not a stationary point of $S$; the coordinate curvature
term $\nabla S\cdot y''$ of the gauge directions is removed by the
projection, and the linear term $\nabla S\cdot y$ has no component in the
retained space at first order, since $\nabla S(\varepsilon b)$ is proportional
to the background's own mode.
\end{proposition}
\begin{proof}
The first statement is the expansion of $R(\operatorname{Exp}v)^{\mathsf T}$ to
first order in $v$. Continuity of the count is verified directly; evenness
follows from the symmetry $b\to-b$ composed with the colour reflection
$e_{\rm col}\to-e_{\rm col}$, an automorphism of the action, and analyticity
from the analytic dependence of $P$ and $H$ on $\varepsilon$ with a spectral
gap above the threshold. The exploration finds the ratios
$4.000,4.002,4.009$ of $Z(\varepsilon)-Z(0)$ under doubling of $\varepsilon$
from $0.0125$ to $0.1$ at $L=8$, $m=1$, and a vanishing odd part.
\end{proof}

\subsection{Exact transverse block reduction}

\begin{proposition}[Block structure]
\label{f16v69:blocks}
Since the background depends on $x_1$ only, $H(\varepsilon)$, $G(\varepsilon)$ and
$E(\varepsilon)$ commute with translations in $x_\perp=(x_2,x_3)$. For each
transverse momentum $q_\perp\in(2\pi\mathbb Z/L)^2$ the projected Hessian
restricts to the $9L$-dimensional space of fields
$y_{(x_1,x_\perp),\mu}=\hat y_{x_1,\mu}e^{iq_\perp\cdot x_\perp}$, with the
Hermitian block
$\hat H_{q_\perp}[(x_1,\mu),(x_1',\mu')]=\sum_{d\in\{-1,0,1\}^2}e^{iq_\perp\cdot d}H[(x_1,0,\mu),(x_1',d,\mu')]$,
assembled from the per-plaquette Hessians, which depend on $x_1$ only;
the vertical frame restricts to $\hat G_{q_\perp}$ with the phases $e^{iq_\mu}$
on the transverse links, and $E$ lives at $q_\perp=0$. Hence
$Z_L(\varepsilon;k)=\sum_{q_\perp}\sum_j\xi(\nu_j^{(q_\perp)}(\varepsilon))$, a sum
of $L^2$ exact diagonalizations of size $9L$.
\end{proposition}
\begin{proof}
Translation covariance and the locality of the plaquette couplings, which
reach transverse distance one, so that an auxiliary $L\times3\times3$
lattice carries all distinct couplings. The verifier's block computation
agrees with the full $1944$-dimensional diagonalization at $L=6$:
$\Pi_6(2\pi/6)=-122.85673$ and $\Pi_6(4\pi/6)=-325.86060$ in both.
\end{proof}

\subsection{Results}

\begin{center}\small
\begin{tabular}{@{}lrrrr@{}}
\toprule
$L$ & $m$ & $|k|$ & $\Pi_L(k)$ & $\delta_L(k)$\\
\midrule
6 & 1 & 1.0472 & -122.857 & -1.1376\\
6 & 2 & 2.0944 & -325.861 & -1.0057\\
8 & 1 & 0.7854 & -189.292 & -1.2623\\
8 & 2 & 1.5708 & -571.260 & -1.1157\\
8 & 3 & 2.3562 & -884.585 & -1.0121\\
16 & 1 & 0.3927 & -482.261 & -1.5468\\
16 & 2 & 0.7854 & -1686.760 & -1.4060\\
16 & 3 & 1.1781 & -3220.884 & -1.2738\\
16 & 4 & 1.5708 & -4803.713 & -1.1728\\
16 & 5 & 1.9635 & -6225.285 & -1.0992\\
16 & 6 & 2.3562 & -7334.879 & -1.0490\\
32 & 1 & 0.1963 & -1149.912 & -1.8263\\
32 & 2 & 0.3927 & -4229.065 & -1.6955\\
32 & 3 & 0.5890 & -8645.264 & -1.5655\\
32 & 4 & 0.7854 & -14028.430 & -1.4617\\
32 & 5 & 0.9817 & -20049.051 & -1.3767\\
32 & 6 & 1.1781 & -26411.617 & -1.3057\\
64 & 1 & 0.0982 & -2694.295 & -2.1344\\
\bottomrule
\end{tabular}
\end{center}
The fit $\delta=A+b\log|k|+c\,k^2$ over the points with $|k|\le1$ from the
largest available $L$ gives $b=$ 0.307, $A=$ -1.390, $c=$ 0.010, against the
expected one-loop value $b=11/(3\pi^2)=0.3715$ for the running of
$\gamma=4/g^2$ with the scale: $\gamma_{\rm eff}(k)-\gamma_{\rm eff}(k')=\frac{11}{3\pi^2}\log(k/k')$
for $SU(2)$, since $\mu\,dg/d\mu=-\frac{11}{24\pi^2}g^3$ gives
$d\gamma/d\log\mu=\frac{11}{3\pi^2}$.

Three remarks. First, the sign: $\delta_L(k)<0$ and decreasing toward small
$|k|$ means the effective magnetic coefficient is smaller at long
wavelengths, i.e.\ the coupling grows in the infrared; this is the
direction of asymptotic freedom, obtained here from the vacuum energy of
the hard modes in a static background, the mechanism of the Hamiltonian
derivations. Second, the static kernel renormalizes only the magnetic
term; the electric (kinetic) term's renormalization requires a
time-dependent background and is not computed, so $\gamma_{\rm eff}$ here is
the magnetic coefficient, which coincides with the running coupling only
where the one-loop effective action is Lorentz invariant, i.e.\ at
$|k|\ll1$ up to $O(k^2)$. Third, the finite-volume dependence at fixed
$|k|$ (visible between $L=8$ and $L=16$ at $|k|=\pi/4$) is the zero-mode
and discretization effect of the transverse loop momenta, of relative order
$1/(L|k|)$, and the fit window is chosen at the largest $L$ per momentum.
The slope is therefore a finite-lattice estimate of the coefficient; the
asymptotic statement requires $L\to\infty$ at fixed $|k|$ and then
$|k|\to0$, which the block reduction makes accessible but which is not
carried out to its end here.

\subsection{What this is and is not}

$\Pi_L(k)$ is the exact one-loop static kernel of the Wilson transfer at
finite $L$, in the frozen Gaussian sense: the hard modes are quadratic and
the background is classical. It is the first explicit renormalization-group
quantity of the lineage: the coefficient of the low-mode magnetic term
after the hard modes are integrated at one loop, as a function of the
momentum, with the scale dependence that the running coupling must have.
It is not yet a renormalization step in the sense of V66's programme
clarification: no scale is integrated out separately, no effective action
for the remaining modes is written with a controlled remainder, and the
electric coefficient is not renormalized. Those are the next items: the
same block computation with the hard modes restricted to a momentum shell
(a genuine one-shell elimination), the time-dependent background for the
electric term, and the cubic and quartic remainders beyond the frozen
Gaussian, whose sizes V67 estimated.

\subsection{Erratum to V67: the cubic vertex sum is a third chaos, not a Gaussian}

The concurrent lineage, in its V57 (dependency audit), correctly observes
that the proof of V67's Theorem~\ref{f16v67:obstruction} calls the Wick-ordered
cubic vertex sum $Z=\sum_p[\xi_p(U)+\xi_p(U')]$ ``a centred Gaussian in the
frozen model'' and evaluates $\operatorname{Var}(e^{-Z})=e^{s^2}(e^{s^2}-1)$.
That is wrong: $Z$ lies in the third Gaussian chaos, its law is not Gaussian
(for $Z=X_1X_2X_3$ the kurtosis is $27$, not $3$), and its unconditional
exponential moments are infinite, since conditionally on $X_1,X_2$ the
inner expectation $e^{t^2X_1^2X_2^2/2}$ has infinite $X_2$-average on the
event $t^2X_1^2\ge1$. The statement of the theorem survives with a
corrected constant and a corrected volume range, by an argument that uses
only the reflection symmetry $y\to-y$ of the frozen Gaussian law and of
the sup-norm event, under which $Z$ is odd and the quartic remainder is
even.

\begin{proposition}[Corrected lower bound]
\label{f16v69:erratum}
Let $W=Z+R'$ on a reflection-symmetric bounded event $E$, where $Z$ is odd,
$R'$ is even, $E[Z\,|\,E]=0$, $s_E^2=E[Z^2\,|\,E]$ and $|R'|\le\rho$ on $E$.
Then for every constant $c$,
\begin{equation}
 E\bigl[(e^{-W}-c)^2\,\big|\,E\bigr]\ \ge\ \operatorname{Var}(e^{-W}\,|\,E)\ \ge\
 \frac{\operatorname{Cov}(e^{-W},Z\,|\,E)^2}{s_E^2}\ \ge\ e^{-2\rho}s_E^2 .
 \label{f16v69:erratum-bound}
\end{equation}
\end{proposition}
\begin{proof}
The first two inequalities are the definition of the variance and
Cauchy--Schwarz. For the third, pair $y$ with $-y$:
$E[Ze^{-W}]=\tfrac12E[Ze^{-R'}(e^{-Z}-e^{Z})]=-E[Z\sinh(Z)e^{-R'}]$, and
$Z\sinh Z\ge Z^2\ge0$, $e^{-R'}\ge e^{-\rho}$ give
$|\operatorname{Cov}(e^{-W},Z)|\ge e^{-\rho}s_E^2$.
\end{proof}

In V67's setting $R'$ is the centred quartic-and-higher remainder, whose
mean is absorbed into the constant $c$ and whose oscillation on the
sup-norm event is $\rho_4=O(\gamma r^4V)=O(V(\log V)^2/\gamma)$ by V64's
sandwich and V66's sup-norm scale $r^2\asymp C_3(0)\log(9V)/\gamma$; and
$s_E^2=s^2(1+o(1))$ because the event has probability $1-o(1)$ and
$\|Z\|_4\le3^{3/2}\|Z\|_2$ by hypercontractivity of the third chaos. The
corrected form of \eqref{f16v67:lower} is therefore
\begin{equation}
 \frac{\|(T_\gamma-c\,T^{\rm frz})\Psi\|}{\|T^{\rm frz}\Psi\|}\ \ge\
 e^{-\rho_4}\,c'\sqrt{\frac V\gamma}\,(1+o(1)),\qquad c'=\sqrt{3v_3/4}=0.187,
 \label{f16v69:corrected-lower}
\end{equation}
which is of order one, and not $o(1)$, as soon as $V\gtrsim\gamma/(\log V)^2$.
The conclusion of Theorem~\ref{f16v67:obstruction} is thus retained with the
volume threshold $V\asymp\gamma$ replaced by $V\asymp\gamma/(\log V)^2$: no
core-wise $L^2$ expansion with $o(1)$ errors can hold beyond that volume, and
the concurrent ledger term $L^3\gamma^{-1}$ is sharp for that norm up to
logarithms. The reformulation of N2 through connected norms, which was the
point of V67, is unaffected; its cluster-property claim was and remains
conditional. The diagnostic
\path{scripts/verify_ym_f16_v69_cubic_erratum.py} checks the reflection
identity and \eqref{f16v69:erratum-bound} by exact quadrature on the
truncated third Hermite chaos $X^3-3X$ with an even perturbation, records the
kurtosis and the divergence mechanism, and confirms the inequality by Monte
Carlo on a six-variable truncated cubic chaos. It has 5065 bytes and
SHA--256
\begin{center}\scriptsize\ttfamily 282291D4B5B84BCCBFC66A61722A79880F092F1E721B0FA1B7E4371EBB1402CA,\end{center}
with canonical hashed payload SHA--256
\begin{center}\scriptsize\ttfamily 322DFA938A3F0D839EF610FBD8180CC930BC92469C63B4CE0D96A74C29C2A2AE.\end{center}

\subsection{Companion note}

The concurrent file has advanced to V58 (its V57: normalization-invariant
soft fraction, nonreducing source-return identity, finite-bank native
resolvent certificate; its V58: private-link source packet with explicit
kernel, capture margin polynomial in $\gamma$, innovation identity and the
noniteration obstruction). Its V57 audit of V67 is accepted above; nothing
else from it is imported, and the full checkpoint is V70's.

\subsection{Status}

\emph{Proved}: Propositions \ref{f16v69:wellposed}, \ref{f16v69:blocks} and \ref{f16v69:erratum} (the corrected V67 bound);
the exact values $\Pi_L(k)$ at the listed $L,k$ are finite-dimensional
computations with the stated finite-difference accuracy.
\emph{Inherited}: V68's reduction, V45's frozen Gaussian, the concurrent
V55 calibration. \emph{Conditional}: the identification of the small-$k$
logarithmic slope with the one-loop $\beta$-function coefficient, which
requires the electric renormalization and the limit $L\to\infty$.
\emph{Open}: the one-shell elimination with remainder; the electric term;
the assembly of a scale-to-scale recursion.

\subsection{Verification and preservation}

The new diagnostic is
\begin{center}\small\path{scripts/verify_ym_f16_v69_one_loop_kernel.py}.\end{center}
Its hashed content is the agreement of the block computation with the
brute-force values at $L=6$, the vanishing of the odd part, and the
continuity of the mode count. Its floating diagnostics are the values
$\Pi_L(k)$ and a column \texttt{delta} equal to $4\Pi_L(k)/(\lambda_kV)$,
which is $2\delta_L(k)$ in the normalization of \eqref{f16v69:definition}:
the script counts the zero-point energy as $\sum\xi$ instead of
$\tfrac12\sum\xi$, an error caught by V71's proper-time density while the
computation was running and corrected here in the definition and in
the table and fit, which are computed from $\Pi_L(k)$; the script is
left as run so that its recorded hash is the one of the artifact. It has 12859 bytes and SHA--256
\begin{center}\scriptsize\ttfamily 77D39AC60D0EA2D0AD7EEF56CE40F992DB4DD029EC9C975F37E1F671CEB0D3B8.\end{center}
Its canonical hashed payload has SHA--256
\begin{center}\scriptsize\ttfamily 338E72D33D483468F9954A2AB78C558B4B4B14EC660EB08FC167540D3899135A.\end{center}
Only this terminal insertion and the literal title version differ from V68;
removing them recovers its bytes exactly. V70 is the ten-version review and
companion checkpoint.

\begin{firewall}
V69 defines the static one-loop kernel of the Wilson transfer at a
transverse plane-wave background, proves its well-posedness and its exact
transverse block reduction, validates the reduction against brute force,
and computes the kernel for $L\le64$, finding the sign and order of
magnitude of asymptotic freedom in its logarithmic slope; it corrects the
Gaussian misstatement in V67's obstruction proof with a symmetry argument
that preserves the conclusion up to logarithms. It does not
prove the one-loop $\beta$-function coefficient, does not renormalize the
electric term, does not perform a scale-to-scale elimination with a
controlled remainder, and does not bear on the physical continuum
trajectory or an interacting continuum mass gap.
\end{firewall}
% END F16 V69 APPEND

% BEGIN F16 V70 APPEND -- 2026-09-18
% Insert before the final document-ending command of cumulative V69.
% This insertion leaves every predecessor body and the native operator unchanged.

\clearpage
\section{V70: companion checkpoint, ten-version review, and the frozen-Gaussian source decrement}
\label{f16v70:context}

This is the scheduled ten-version review and companion checkpoint. It
proves one interface statement, the leading-order form of the local-source
Gram and its delayed decrements that the concurrent lineage's V57--V58
target asks for, and otherwise consolidates. The concurrent file has
advanced to V58; its audit of V67 was accepted and corrected in V69. The
review finds that the ten versions V60--V69 completed the small-volume
coordinate theorem, built the local primitives of the volume-uniform
reduction, located the obstruction of N2 in the norm rather than in the
estimates, fixed the covariant reference, and produced the first
renormalization-group quantity of the lineage, the static one-loop kernel
with the sign and order of asymptotic freedom. It also records what did not
happen: no native concentrated-kernel theorem in the connected norm, no
scale-to-scale elimination, no electric renormalization.

\subsection{Companion checkpoint}

The concurrent file read at this checkpoint is
\path{Renewal_Yang_Mills_Mass_Gap_Research_Notes_V58(1).tex}, 1812139
bytes, 37792 lines, SHA--256
\begin{center}\scriptsize\ttfamily 598F786D17109EAEF7E18007413F573CDBA8EEE9FD99F9CC39E07244AFE97578,\end{center}
whose V1--V56 body is the file read at V65 and V60 (SHA--256 prefix
\texttt{888F3CAC}) with two appended sections. Its V57 introduces the
normalization-invariant soft fraction $\eta^2=\inf\{c:S\preceq cG\}$, an
exact rectangle example showing that variance normalization can retain a
toron component that raw scalarization hides, the nonreducing
source-return identity of a Feshbach elimination that does not reduce the
transfer, a native resolvent certificate at the physical clock built from
zero-, one- and two-step vacuum correlations of a trial bank, and a
Chebyshev depth bound with the cutoff cost exposed. Its V58 populates a
native source/trial packet from mutually noninteracting private links,
proves the covariance bounds $d_\gamma I\preceq G\preceq rI$ and
$G-C_j\succeq(d_\gamma/2)I$ with $d_\gamma$ polynomial in $(1+\gamma)^{-1}$,
identifies the exact source kernel, replaces the microscopic difference by
a difference over a fixed physical interval with positive correlation
weights, and proves the innovation identity
$G-V^*\mathcal T^{2n}V=\sum_{j<n}V^*\mathcal T^j(I-\mathcal T^2)\mathcal T^jV$
together with a two-dimensional counterexample to iterating a one-step
decrement. Its stated next target is a source-marked lower bound
$\sum_{j<n}J_{a,j}\succeq\vartheta_*G_a$ with $\vartheta_*$ independent of
the regulator.

Its V57 dependency audit of V67 is correct and was acted on in V69
(Proposition~\ref{f16v69:erratum}). Nothing else is imported. The four
interface obligations are answered as follows. \emph{Local-source
noncollapse}: at leading order in $\gamma^{-1}$ the entrance Gram of centred
plaquette writers is $\frac32c_0^2\gamma^{-2}$ per writer
(Proposition~\ref{f16v70:decrement}), so the power $\gamma^{-2}$ of the
concurrent lower bound $d_\gamma$ is the correct one, while the concurrent
\emph{relative} decrement $d_\gamma/(2r)\asymp\gamma^{-2}$ is far from the
leading-order value $\theta_1=$ 0.981, which is $\gamma$-free.
\emph{Convergence}: not addressed by this lineage. \emph{Hard-sector
propagation}: in the frozen Gaussian the writer correlation $c_n$ is a
positive combination of exponentials $e^{-n\xi_k}$, hence completely
monotone and log-convex, and every innovation $J_j$ is explicit; the
leading-order decrement is an ultraviolet fact, since the plaquette writer
has no soft content at that order, and says nothing about the margin at a
physical delay $n=\tau/a\to\infty$, which is precisely the content the
concurrent programme needs. \emph{Rate for $\varepsilon_{L,\gamma}$}: V69's
kernel shows that the one-loop renormalization of $\gamma$ is an $O(1)$
shift, i.e.\ relative order $\gamma^{-1}$, with a logarithmic momentum
dependence; this is the first quantitative statement of the lineage on
the running of the bare coefficient and it enters the concurrent
calibration $\gamma=4/g_0^2$ as the one-loop correction.

\subsection{The frozen-Gaussian source decrement}

Let $u=\gamma^{1/2}y$ be the hard link field of V45's frozen Gaussian
ground, with transverse Fourier modes of variance $1/(2\omega_k)$,
$\omega_k=\sinh\xi_k$, $\xi_k=\operatorname{arccosh}(1+\hat k^2/2)$,
$\hat k^2=\sum_\mu|e^{ik_\mu}-1|^2$, $k\ne0$, and let the Doob chain of the
frozen transfer be the stationary Gaussian Markov chain of these modes.
For a plaquette $p$ in the plane $(\mu\nu)$ at $x$ let
$F_p^a=\Delta_\mu u_\nu^a-\Delta_\nu u_\mu^a$ be the linearized field
strength and $w_p=1-\tfrac12\operatorname{Tr}U_p=\tfrac{\gamma^{-1}}2|F_p|^2+O(\gamma^{-3/2})$
the plaquette writer (V67, Proposition~\ref{f16v67:word}).

\begin{proposition}[Leading-order source Gram and decrement]
\label{f16v70:decrement}
In the frozen Gaussian chain, for plaquettes $p,q$ and delay $n\ge0$,
\begin{equation}
 E\bigl[F_p^a(0)F_q^b(n)\bigr]=\delta^{ab}c_n(p,q),\qquad
 c_n(p,q)=\frac1V\sum_{k\ne0}\frac{e^{-n\xi_k}}{2\omega_k}\,
 \overline{a_p^T(k)}\cdot a_q^T(k)\,e^{ik\cdot(x_q-x_p)},
 \label{f16v70:c}
\end{equation}
where $a_p^T(k)$ is the transverse projection of the vector with
components $a_\nu=w_\mu$, $a_\mu=-w_\nu$, $w_\rho(k)=e^{ik_\rho}-1$. The
centred writers $\tilde w_p=w_p-Ew_p$ have, at leading order,
\begin{equation}
 \operatorname{Cov}\bigl(\tilde w_p(0),\tilde w_q(n)\bigr)=\frac{3}{2\gamma^2}\,c_n(p,q)^2,
 \label{f16v70:Gram}
\end{equation}
so the entrance Gram $G$ and the delayed Grams $C_n$ of any finite
family are $\gamma^{-2}$ times $\gamma$-free matrices, and every relative
decrement, in particular the capture margin
$\vartheta_n=1-\lambda_{\max}(G^{-1/2}C_nG^{-1/2})$, is $\gamma$-free at
leading order. For a single writer $c_n(p,p)$ is a positive combination of
the exponentials $e^{-n\xi_k}$, hence completely monotone in $n$ and
log-convex, $c_{n+1}c_{n-1}\ge c_n^2$, and every innovation
$C_{2j}-C_{2j+2}$ is positive and explicit.
\end{proposition}
\begin{proof}
The single-mode transfer $e^{-\lambda q^2/4}e^{\partial^2/2}e^{-\lambda q^2/4}$
has stationary variance $1/(2\sinh\xi)$ and $n$-step correlation
$e^{-n\xi}$, $\cosh\xi=1+\lambda/2$ (V38, V41; checked to $2\cdot10^{-4}$
on a grid by the verifier). Modes are independent, so \eqref{f16v70:c}
follows from $F_p=a_p\cdot\hat u$ in Fourier space with the transverse
projection, and the colour structure is diagonal. For jointly Gaussian
centred vectors, $\operatorname{Cov}(|F_p|^2,|F_q|^2)=2\sum_{ab}(E F_p^aF_q^b)^2=6c_n(p,q)^2$,
which with the factor $(\gamma^{-1}/2)^2$ gives \eqref{f16v70:Gram}. The
weights $|a_p^T|^2/(2\omega_kV)$ are positive, which gives complete
monotonicity, and log-convexity is Cauchy--Schwarz on the weights. The
equal-time value $C_3(0)=V^{-1}\sum_{k\ne0}P_{\mu\mu}(k)/(2\omega_k)$ in
this normalization is $0.1032$ at $L=48$, against V61's $0.1033$.
\end{proof}

\begin{center}\small
\begin{tabular}{@{}lrrrrrr@{}}
\toprule
$n$ & 1 & 2 & 3 & 4 & 6 & 8\\
\midrule
$c_n/c_0$ & 0.1389 & 0.0231 & 0.0049 & 0.0014 & 0.0002 & 0.0001\\
$\theta_n$ & 0.9807 & 0.9995 & 1.0000 & 1.0000 & 1.0000 & 1.0000\\
$\vartheta_n$ (3) & 0.9804 & 0.9995 & 1.0000 & 1.0000 & 1.0000 & 1.0000\\
$\vartheta_n$ (6) & 0.9793 & 0.9994 & 1.0000 & 1.0000 & 1.0000 & 1.0000\\
\bottomrule
\end{tabular}
\end{center}
The first two rows are $c_n/c_0$ and $\theta_n=1-(c_n/c_0)^2$ for one
writer at $L=64$; the last two are the capture margins $\vartheta_n$ of the
three-orientation packet at one site and of the six-writer two-site family
at $L=32$. All are ultraviolet numbers: a plaquette writer in the frozen
Gaussian is a hard observable whose correlation decays at the rate of
the hard frequencies, and its residual at a physical delay is governed by
the slowest modes $\xi_{\min}\approx2\pi/L$ with weight $O(1/V)$. The
$\gamma$-dependence of the true Grams enters at relative order
$\gamma^{-1}$ through V67's cubic vertex and the quartic terms, and the
soft (toron and low-momentum) content of the writers, which is what a
physical capture margin would have to control, is not visible at this
order. The proposition is therefore a calibration of the concurrent
target, not a contribution to it.

\subsection{Ten-version review, V60--V69}

\begin{center}\small
\begin{tabular}{@{}p{0.8cm}p{7.6cm}p{4.6cm}@{}}
\toprule
& Result & Status\\
\midrule
V60 & Checkpoint; consolidated fixed-regulator ledger; integrity check & Consolidation\\
V61 & Line-chaos lemma $\|T_N\|_2\le(6L\|C\|_1)^{N/2}/\sqrt{N!}$; bias theorem of the Polyakov coordinate & Proved (given $\|C\|_1$, numerical)\\
V62 & Wick-ordered cubic and mixed chaoses; line-averaged fluctuations; coordinate theorem at the critical scale & Proved (given covariance sums)\\
V63 & Summation by parts; $|C_3(x)|\le(\kappa_0+\kappa_1\log L)/\max x_i^2$; coordinate theorem unconditional for $L\log^3L\ll\gamma^{2/3}$ & Proved\\
V64 & Dimension-free sup-norm kinetic sandwich; local structure of the transfer; abstract Feshbach for the Doob chain; crude Harnack & Proved; quadratic Harnack open\\
V65 & Checkpoint; programme redirection; transport identity $w=\rho P(\rho w)$ & Consolidation; identity proved\\
V66 & Gaussian maximum theorem; mean-one factorization $k=g\rho$; connected budget $\sqrt{V}/\gamma$; programme clarification & Proved in the frozen chain\\
V67 & Cubic vertex Wick-ordered; $\operatorname{Var}\sum\xi_p=\frac34v_3V/\gamma$; norm obstruction of N2 & Proved; bound corrected in V69\\
V68 & Exact horizontal reduction; both extensive density terms cancel; twisted spectrum recovered & Proved at the Gaussian level\\
V69 & Static one-loop kernel $\Pi_L(k)$; exact transverse block reduction; log slope 0.307 vs $11/(3\pi^2)$; V67 erratum & Proved; identification conditional\\
\bottomrule
\end{tabular}
\end{center}

Three conclusions. First, the small-volume programme (V59--V63) is
complete as far as it was meant to go: the gauge-invariant Polyakov
coordinate is controlled unconditionally for $L\log^3L\ll\gamma^{2/3}$, and
this regime is not the continuum trajectory and was never called so.
Second, the volume-uniform reduction (V64--V68) produced its primitives
but not N2: the obstruction is the norm, since the magnetic factor's
cubic vertex has extensive variance in the $L^2$ norm of the hard band
(V67, corrected in V69 to the threshold $V\asymp\gamma/\log^2V$), so the
native concentrated-kernel theorem has to be stated through connected or
cluster quantities along the Doob chain. No such statement has been
proved, and the mean-one factorization of V66, the transport identity of
V65 and the Feshbach reduction of V64 are the tools it will use. Third,
the perturbative renormalization-group programme of the V66 clarification
has started: V68 fixed the covariant reference (the horizontal Hessian at
a background, with no chart terms) and V69 computed the first
renormalization-group quantity, the static one-loop kernel, whose
logarithmic slope has the sign and order of asymptotic freedom in the
concurrent calibration. Against the order of attack of V66: step 1 (N2)
is open in connected form; step 2 (a uniform local block or Feshbach
statement) has the abstract form of V64 and the block structure of V69
but no uniform bounds; step 3 (one explicit coarse-graining step) has the
static magnetic part of its one-loop content and nothing else; steps 4--8
are untouched.

Two errata were recorded in the ten versions: V58 corrected V56's
trapezoid values of the Agmon exponent, and V69 corrected V67's Gaussian
misstatement. Both corrections preserved the conclusions with changed
constants. The concurrent lineage's audit function was useful in the
second case, and the reciprocal audit (V68's finding that the concurrent
V50/V54 density term is a chart artifact) has not been answered in its
V57--V58.

\subsection{Integrity inventory}

The verifier checks that each of the eleven artifacts of V60--V69
(including V69's erratum diagnostic) records the SHA-256 of its present
verifier source and a passing hashed payload. All eleven do. The scripts
and artifacts are those listed in the preservation subsections of the
respective versions; no other script of the lineage was modified in
V60--V69.

\subsection{Next: the one-shell elimination}

V71 begins step 3 proper. In V69's block structure the transverse momentum
$q_\perp$ is exact, so a shell can be defined as a set of transverse
momenta, and the hard modes of the shell can be eliminated by V64's
Feshbach map at the Gaussian level with the background as the soft
variable: the result is the shell-restricted kernel $\Pi^{(\Lambda)}_L(k)$,
additive over disjoint shells at one loop, and the first check is whether
the sum over shells reproduces V69's kernel and whether each shell's
contribution is local in $x_1$ with a computable remainder. The cubic
vertex enters at the next order through the mean-one factorization of
V66; the electric coefficient requires the time-dependent background,
which the Doob chain supplies. The target theorem is
$\mathcal F_{j+1}=\mathcal R_j\mathcal F_j\mathcal R_j^*+\mathcal E_j$ with
$\mathcal R_j$ the shell elimination and $\|\mathcal E_j\|_{\rm local}$ bounded by
the shell's cubic and quartic budgets, and it is to be proved first in
the frozen Gaussian with the vertex as a perturbation, where it is a
finite computation, before any uniformity in $L$ is attempted.

\subsection{Status}

\emph{Proved}: Proposition~\ref{f16v70:decrement} in the frozen Gaussian
chain; the integrity of the V60--V69 artifacts. \emph{Inherited}: the
status ledgers of V60--V69 as stated in their sections; the concurrent
V57--V58 as read-only comparison inputs. \emph{Conditional}: nothing new.
\emph{Open}: N2 in connected form; the one-shell elimination; the electric
renormalization; the uniform block statement.

\subsection{Verification and preservation}

The new diagnostic is
\begin{center}\small\path{scripts/verify_ym_f16_v70_review_checkpoint.py}.\end{center}
Its hashed content is the artifact integrity, the single-mode Mehler
check, the equal-time covariance against V61, and the log-convexity of
the writer correlation. Its floating diagnostics are the tables above. It
has 8713 bytes and SHA--256
\begin{center}\scriptsize\ttfamily 75E0C6C8D1E58F7722EFCB98ABDF608D3FE17A86F57C8456F6FA2DC2B5D8C234.\end{center}
Its canonical hashed payload has SHA--256
\begin{center}\scriptsize\ttfamily 238047D19FF025F6C13D073871034ED6C46F16F584E306CC49D5AF4861D706F1.\end{center}
Only this terminal insertion and the literal title version differ from V69;
removing them recovers its bytes exactly. The next checkpoint is V75, the
next ten-version review V80.

\begin{firewall}
V70 records the companion checkpoint at the concurrent V58, proves the
leading-order form of the local-source Gram and its delayed decrements in
the frozen Gaussian chain as a calibration of the concurrent capture
target, reviews V60--V69 against the order of attack, and verifies the
integrity of their artifacts. It proves nothing about the physical capture
margin, N2, a scale-to-scale elimination, the physical continuum
trajectory or an interacting continuum mass gap.
\end{firewall}
% END F16 V70 APPEND

% BEGIN F16 V71 APPEND -- 2026-09-18
% Insert before the final document-ending command of cumulative V70.
% This insertion leaves every predecessor body and the native operator unchanged.

\clearpage
\section{V71: the proper-time scale decomposition of the one-loop kernel and the failure of momentum shells}
\label{f16v71:context}

This is the first attempt at step 3 of the order of attack, one explicit
scale-by-scale elimination, at one loop and in the static magnetic
sector. Three things are found. First, the two obvious notions of a
shell fail: attributing the one-loop kernel to shells of transverse
momentum (the exact block label of V69) or to shells of the eigenvalues of
the horizontal Hessian gives per-shell contributions that alternate in
sign and are several times larger than the total, because a momentum
shell is not gauge compatible and an eigenvalue shell carries boundary
terms from the motion of eigenvalues across its edges; the mode-pair
formula that makes this precise is proved and validated. Second, the
proper-time decomposition of the same kernel, through the heat trace of
the horizontal Hessian at the background, is gauge invariant, exact,
additive, and has a smooth scale density; the density is universal in the
ultraviolet, vanishes below the background momentum, and in the window
between the two is $-0.365$ per e-fold of the scale, against the
one-loop coefficient $-11/(3\pi^2)=-0.3715$ of asymptotic freedom in the
concurrent calibration. This is the scale-by-scale form of V69's result
and replaces its fit. Third, a normalization error in V69's verifier is
recorded and corrected: the hard zero-point energy is $\frac12\sum\xi$, not
$\sum\xi$, which V69's definition already states.

\subsection{Dependencies}

V69's horizontal Hessian at the background, its exclusions and its block
structure, with the same finite-difference amplitude $\varepsilon=0.05$;
V68's gauge invariance of the horizontal spectrum; V38/V41's Mehler
exponent. The heat-kernel representation of the square root and the
second-order perturbation formula for spectral traces are classical.

\subsection{The mode-pair formula and the two failing shells}

Let $A(\varepsilon)=P(\varepsilon)H(\varepsilon)P(\varepsilon)$ on one transverse block, with
$A(\varepsilon)=A_0+\varepsilon A_1+\varepsilon^2A_2+O(\varepsilon^3)$, eigenpairs
$(\nu_i,e_i)$ of $A_0$, hard set $\{\nu_i>\tau\}$ and null set $\{\nu_i=0\}$
(the projected-out directions, which stay exactly null for all $\varepsilon$
while their span rotates).

\begin{proposition}[Second-order trace formula with a rotating null space]
\label{f16v71:trace-formula}
For $f$ smooth on $[\tau,\infty)$ the coefficient of $\varepsilon^2$ in
$\sum_{\nu_j(\varepsilon)>\tau}f(\nu_j(\varepsilon))$ is
\begin{equation}
 \sum_{i\,\rm hard}f'(\nu_i)(A_2)_{ii}
 +\frac12\sum_{i,j\,\rm hard}|(A_1)_{ij}|^2\,\frac{f'(\nu_i)-f'(\nu_j)}{\nu_i-\nu_j}
 +\sum_{i\,\rm hard,\ j\,\rm null}|(A_1)_{ij}|^2\,\frac{f'(\nu_i)}{\nu_i},
 \label{f16v71:pairs}
\end{equation}
with $f''(\nu_i)$ for the diagonal and degenerate terms of the middle sum.
\end{proposition}
\begin{proof}
Write the hard trace as $\frac1{2\pi i}\oint f(z)\operatorname{Tr}(z-A(\varepsilon))^{-1}dz$
on a contour enclosing $[\tau,\|A\|]$ and not $0$, expand the resolvent to
second order, and take residues: the residue at $\nu_i$ of
$f(z)/((z-\nu_i)^2(z-\nu_j))$ is $f'(\nu_i)/(\nu_i-\nu_j)-f(\nu_i)/(\nu_i-\nu_j)^2$,
and the second term cancels against the residue at $\nu_j$ when $\nu_j$ is
inside the contour and survives as stated when $\nu_j=0$ is outside;
symmetrizing over $i\leftrightarrow j$ gives the divided difference of
$f'$. The verifier checks \eqref{f16v71:pairs} against direct finite
differences on random smooth families $P(\varepsilon)H(\varepsilon)P(\varepsilon)$ with
a rotating null space, to relative accuracy 2e-06.
\end{proof}

The three terms of \eqref{f16v71:pairs} are individually large and of
both signs: at $L=16$, $m=1$, in units of $\Pi_L(k)=-482$, the hard--hard,
hard--null and diagonal terms are $-1077$, $+1704$ and $-1113$. The
hard--null term is the coupling of the hard modes to the rotating gauge
and low-sector directions, i.e.\ a chart quantity; only the total is
gauge invariant. Any attribution of the pair terms to scales through
the modes' momenta inherits this. Attributing the pair $(i,j)$ to the
harder of its two modes and grouping in dyadic shells of $|\hat k'|$
gives, at $L=16$, $m=1$, the shell contributions
$-2.66$, $+0.82$, $+0.18$, $+0.10$ to $\delta_L(k)=-1.55$, and attributing to the softer mode
changes them individually by amounts of order one. Two other definitions
fail for a different reason: grouping V69's blocks by dyadic shells of
transverse momentum gives $-3.41$, $+1.17$, $+0.56$, $+0.13$, and grouping the eigenvalues of
$A(\varepsilon)$ themselves by dyadic shells (a gauge-invariant definition)
gives $-0.56$, $+17.48$, $-18.06$, $-0.40$. The last alternation is the boundary effect: an
eigenvalue moved by $O(\varepsilon)$ across a shell edge contributes
$O(\varepsilon^2)\times$(density of states at the edge)$\times\xi'$ to the two
adjacent shells with opposite signs, so sharp or bump-shaped spectral
shells carry terms that cancel only in the sum.

\subsection{The proper-time decomposition}

For the square-root kernel,
$\sqrt\nu=(2\sqrt\pi)^{-1}\int_0^\infty t^{-3/2}(1-e^{-t\nu})\,dt$, so with the
heat trace $\Theta(t;\varepsilon)=\sum_{\rm hard}e^{-t\nu_j(\varepsilon)}$ of the
horizontal Hessian at the background and its second-order coefficient
$\Theta_2(t)$,
\begin{equation}
 \Pi^{\sqrt{}}_L(k)=\int_0^\infty h(t)\,\frac{dt}t,\qquad
 h(t)=-\frac{t^{-1/2}}{2\sqrt\pi}\,\Theta_2(t),
 \label{f16v71:proper-time}
\end{equation}
and $2h(t)$ is the contribution per e-fold of the scale $\Lambda=t^{-1/2}$.
The Mehler kernel $\xi(\nu)=\sqrt\nu(1-\nu/24+\dots)$ differs from $\sqrt\nu$
only at lattice scales, so below $\Lambda\approx1$ the two densities
coincide. With the zero-point energy $E_0=\frac12\sum\xi$ and the
background energy $\gamma\lambda_k|l|^2/2$, the normalized density
\begin{equation}
 \rho_L(\Lambda;k)=\frac{2h(t)}{\lambda_kV/2}\Big|_{t=\Lambda^{-2}},\qquad
 \delta^{\sqrt{}}_L(k)=\int_0^\infty\rho_L(\Lambda;k)\,\frac{d\Lambda}\Lambda,
 \label{f16v71:density}
\end{equation}
is the contribution to the one-loop shift of $\gamma$ per e-fold of
$\Lambda$.

\begin{proposition}[Properties of the density]
\label{f16v71:density-properties}
$\Theta(t;\varepsilon)$ is a gauge-invariant function of the background for
every $t$ (V68), so $\rho_L(\Lambda;k)$ is gauge invariant, exact at one loop
in the frozen Gaussian, and additive over disjoint intervals of $\Lambda$.
It is even in $\varepsilon$, and the integral \eqref{f16v71:proper-time}
reproduces the direct second-order coefficient of $\sum\sqrt{\nu_j}$ (the
verifier finds agreement to better than one percent on a grid of ten
points per e-fold over twenty-five e-folds of $t$).
\end{proposition}

\begin{center}\small
\begin{tabular}{@{}rrrrrr@{}}
\toprule
$\Lambda$ & $L{=}16$, $|k|{=}0.39$ & $L{=}16$, $|k|{=}0.79$ & $L{=}32$, $|k|{=}0.20$ & $L{=}32$, $|k|{=}0.39$ & $L{=}64$, $|k|{=}0.10$\\
\midrule
3.008 & -0.684 & -0.686 & -0.687 & -0.691 & -0.686\\
2.015 & -0.640 & -0.633 & -0.648 & -0.651 & -0.649\\
1.420 & -0.513 & -0.490 & -0.529 & -0.529 & -0.533\\
1.000 & -0.398 & -0.353 & -0.428 & -0.421 & -0.435\\
0.704 & -0.326 & -0.252 & -0.377 & -0.357 & -0.391\\
0.496 & -0.261 & -0.163 & -0.347 & -0.305 & -0.372\\
0.350 & -0.177 & -0.082 & -0.311 & -0.240 & -0.358\\
0.246 & -0.072 & -0.023 & -0.257 & -0.164 & -0.338\\
0.182 & -0.014 & -0.003 & -0.192 & -0.100 & -0.311\\
0.122 & -0.000 & -0.000 & -0.082 & -0.032 & -0.252\\
0.090 & -0.000 & -0.000 & -0.022 & -0.007 & -0.185\\
0.061 & -0.000 & -0.000 & -0.001 & -0.000 & -0.067\\
\bottomrule
\end{tabular}
\end{center}
The table lists $\rho_L(\Lambda;k)$. Three regimes are visible. For
$\Lambda\gtrsim2$ the density is universal, identical for all $k$ and $L$
to three digits, and lattice-dominated, with a maximum $-1.37$ at
$\Lambda\approx3$ and a decay $\propto\Lambda^{-1}$ beyond the lattice cutoff
$\sqrt{12}$. For $\Lambda\lesssim|k|/2$ it vanishes: modes softer than the
background decouple from its renormalization. In between, for
$|k|\ll\Lambda\ll1$, the density is slowly varying; at $L=32$, $|k|=0.20$,
over $\Lambda\in[0.4,0.94]$ it averages $-0.365$, and at $L=64$, $|k|=0.098$, over $\Lambda\in[0.2,0.94]$ it averages $-0.367$.
The value expected from the one-loop $\beta$-function in the
normalization $\gamma=4/g_0^2$ is $-11/(3\pi^2)=-0.3715$: with a constant
density $-b$ between $|k|$ and the cutoff, $\delta_L(k)=-b\log(\Lambda_{\rm UV}/|k|)+{\rm const}$,
i.e.\ $\gamma_{\rm eff}$ grows with $|k|$ at rate $b$, which is asymptotic
freedom. The residual $\Lambda$-dependence in the window is of the size of
the $O(\Lambda^2)$ lattice corrections and the $O(|k|^2/\Lambda^2)$
infrared corrections at its two ends, and the window is only one e-fold
wide at $L=32$; the identification of the plateau with $11/(3\pi^2)$ is
numerical and is recorded as such.

\subsection{Correction to V69's verifier}

V69's definition \eqref{f16v69:definition} takes the hard zero-point
energy as $\frac12\sum_j\xi(\nu_j)$, but its verifier's column
\texttt{delta} and fit use $\sum_j\xi(\nu_j)$, so they are twice
$\delta_L(k)$; V69's table, slope and memory were written from $\Pi_L(k)$
with the correct factor, and the discrepancy was noticed through the
density of this section, whose first computation gave a plateau of twice
the expected size. The single-mode transfer
$e^{-\lambda q^2/4}e^{\partial^2/2}e^{-\lambda q^2/4}$ has top eigenvalue exactly
$e^{-\xi/2}$ and gap ratio $e^{-\xi}$ (V70's verifier, to $10^{-6}$); the
vacuum energy per mode is $\xi/2$.

A second defect concerns V69's $L=64$ entry only. V69's hard threshold
$0.2\lambda_{\min}$ was chosen when the gauge directions were not yet
projected exactly; with the exact projection the null eigenvalues are
below $10^{-13}$, and the threshold serves no purpose, but at $L=64$ it
drops four physical low modes that the background of amplitude
$\varepsilon=0.05$ shifts below $0.2\lambda_{\min}=0.0019$: V69's artifact records
the hard counts $1572854$ at $\varepsilon=0$ and $1572850$ at $\varepsilon=\pm0.05$,
against the count $6V-10$ that Proposition~\ref{f16v69:wellposed}
requires. The dropped modes change $\Pi_{64}(2\pi/64)$ by about one
percent but make the small-$t$ heat trace meaningless, which is how the
defect was found (a positive ultraviolet density at $L=64$ against the
universal negative one). The lowest modes at $L=64$ have $\nu\approx0.0096$
and are shifted by a large fraction of their value at $\varepsilon=0.05$, so
the $L=64$ point is also outside the range where the finite-difference
amplitude is perturbative for them; the verifier of this section uses
the threshold $10^{-6}$, asserts the count $6V-10$ at every $\varepsilon$, and
recomputes $L=64$ with $\varepsilon=0.0125$. The $\varepsilon$-halving test at $L=32$,
$m=1$ gives $\Pi=-1149.9,-1142.1,-1140.5$ at $\varepsilon=0.05,0.025,0.0125$,
so V69's $L=32$ entries carry a finite-amplitude error of about one
percent, within the accuracy claimed there; the $L=64$ row of V69's
table is superseded by $\Pi_{64}(2\pi/64)=-2615.9$, $\delta_{64}=-2.0723$ (at $\varepsilon=0.0125$).

\subsection{What this is and is not}

The proper-time decomposition is the gauge-invariant scale separation
available in this framework: an interval of proper time is a shell, the
shells are additive at one loop, and the density in the window is the
$\beta$-function coefficient. It is the one-loop content of a scale-by-scale
elimination, not the elimination itself: no effective transfer for the
modes below $\Lambda$ has been written, the cubic vertex of V67 has not
been inserted, the electric coefficient is not renormalized, and the
remainder of a genuine Feshbach step is not bounded. The obstruction
found is of the coordinate class: momentum shells, the natural variables
of the transverse block reduction, are not gauge compatible, and any
scale-by-scale theorem of the lineage must be formulated with a
gauge-covariant regulator, of which proper time is the simplest; this
fixes the form of step 3. The next version inserts the time-dependent
background for the electric coefficient, for which the Doob chain of the
frozen transfer supplies the dynamics, and states the first
scale-to-scale identity in proper-time form.

\subsection{Status}

\emph{Proved}: Propositions \ref{f16v71:trace-formula} and
\ref{f16v71:density-properties}. \emph{Inherited}: V69's kernel and its
validation, V68's gauge invariance. \emph{Conditional}: the identification
of the window density with $11/(3\pi^2)$ (numerical; needs $L\to\infty$ and
a wider window). \emph{Open}: the electric renormalization; the effective
transfer below $\Lambda$; the remainder; N2 in connected form.

\subsection{Verification and preservation}

The new diagnostic is
\begin{center}\small\path{scripts/verify_ym_f16_v71_proper_time.py}.\end{center}
Its hashed content is the trace formula, the proper-time integral
identity, the ultraviolet universality, the infrared decoupling, and the
sign of the density above $|k|$. Its floating diagnostics are the tables.
It has 10439 bytes and SHA--256
\begin{center}\scriptsize\ttfamily F50F31A8590479B72C09AD8C7D0532BB7951C381F49A40620CC20468332C96A4.\end{center}
Its canonical hashed payload has SHA--256
\begin{center}\scriptsize\ttfamily 8AD6717A698973EF8FDE334D6474574E3D97C892251DC524B5EBA07D181C48FE.\end{center}
Only this terminal insertion and the literal title version differ from V70;
removing them recovers its bytes exactly.

\begin{firewall}
V71 proves the second-order trace formula with a rotating null space,
shows that momentum and eigenvalue shells give chart-dependent or
boundary-dominated attributions of the one-loop kernel, defines the
gauge-invariant proper-time scale density of the kernel, and finds it
universal in the ultraviolet, vanishing in the infrared, and close to the
asymptotic-freedom coefficient in the window between; it corrects the
zero-point normalization of V69's verifier. It does not write an
effective transfer below a scale, does not renormalize the electric term,
does not bound a remainder, and does not bear on the physical continuum
trajectory or an interacting continuum mass gap.
\end{firewall}
% END F16 V71 APPEND

% BEGIN F16 V72 APPEND -- 2026-09-18
% Insert before the final document-ending command of cumulative V71.
% This insertion leaves every predecessor body and the native operator unchanged.

\clearpage
\section{V72: the covariant kernel, the finite-time identity, the electric coefficient, and the scheme dependence of the off-shell kernel}
\label{f16v72:context}

V69 and V71 renormalize the magnetic coefficient of the transfer from the
zero-point energy of the hard modes in a static background. The
programme clarification requires the kinetic (electric) normalization to
be tracked separately. This section computes the analogous one-loop
quantity from the four-dimensional Euclidean Wilson action at a
space-time plane wave $b=\varepsilon\cos(kx_1+\omega t)\,e_{\rm pol}\otimes e_{\rm col}$,
with the full four-dimensional gauge group projected out, and
establishes four things. First, the finite-time identity: the Gaussian
free energy of $\operatorname{Tr}T^N$ over $N$ slices is, mode by mode, $N$
times V69's zero-point energy plus $\sum_j\log(1-e^{-N\xi_j})$, and the
second-order coefficient computed from V69's spectrum converges to V69's
kernel as $N$ grows. Second, the
electric coefficient: at $a_t=a_s$ the covariant kernel of a purely
electric background $\cos(\omega t)$ equals that of the magnetic background
$\cos(kx_1)$ at $\omega=k$ exactly, by the lattice symmetry $x_1\leftrightarrow t$.
Third, and unexpectedly, the covariant horizontal-slice determinant is
\emph{not} equal to the transfer kernel at the same $L$ and $N$: it
differs by 5.4 percent at $L_1=N=8$ and 3.1 percent at $L_1=N=16$, the finite-$N$
remainder being negligible, and the naive Faddeev--Popov (vertical-Gram)
factor does not account for the difference. The reading is that the
one-loop kernel at a non-stationary background is a Born--Oppenheimer
quantity that depends on the choice of soft coordinates (the slice), as
the off-shell one-loop effective action depends on the gauge parameter
through terms proportional to the equations of motion; the transfer
scheme (V69's slice) and the covariant scheme (the four-dimensional
horizontal slice) are two such choices, and only the logarithmic
coefficient is common to them. Fourth, Euclidean covariance holds in the
covariant scheme up to lattice corrections of a few percent at
$|k|\approx0.8$. The consequence for the programme is stated: the lineage's
running coupling is scheme-dependent in its finite part at one loop, the
scheme is the transfer's own (V69), and the electric coefficient in that
scheme is not fixed by the symmetry and remains to be computed there.

\subsection{Dependencies}

V69's chart, exclusions and block reduction, now on the lattice
$L_1\times N\times L\times L$ with directions $(x_1,t,x_2,x_3)$; V68's
horizontal reduction; V41's Mehler exponent and V70's single-mode check.
The Gaussian free energy of a periodic chain of oscillators is classical.

\subsection{The covariant kernel}

On the lattice $L_1\times N\times L\times L$, periodic, with the background
on the direction-$2$ links, let $H(\varepsilon)$ be the right-chart Hessian of
the four-dimensional Wilson action, $G(\varepsilon)$ the vertical frame of the
full gauge group (gauge transformations at every site, including
time-dependent ones), $E(\varepsilon)$ the low sector of V69's rule (twelve
constant modes, the background's cos and sin, the colour orbit at
$\varepsilon=0$), $P(\varepsilon)$ the projector onto their orthogonal complement,
and $\nu_j(\varepsilon)$ the eigenvalues of $P H P$ above the threshold. The
Gaussian free energy and the covariant kernel are
\begin{equation}
 F(\varepsilon)=\tfrac12\sum_j\log\nu_j(\varepsilon),\qquad
 \Pi^{(4)}(k,\omega)=\frac{F(\varepsilon)+F(-\varepsilon)-2F(0)}{2\varepsilon^2},\qquad
 \delta^{(4)}(k,\omega)=\frac{4\,\Pi^{(4)}(k,\omega)}{V_4(\lambda_k+\lambda_\omega)},
 \label{f16v72:definition}
\end{equation}
with $V_4=L_1NL^2$ and $\lambda_k+\lambda_\omega$ from the background's
magnetic and electric plaquettes, $S_2=\varepsilon^2V_4(\lambda_k+\lambda_\omega)/4$.
Transverse translation invariance gives blocks of dimension $12L_1N$ per
transverse momentum, conjugate pairs isospectral. The mode count is
$9V_4+3-14$ and is continuous in $\varepsilon$; the odd part vanishes.

\subsection{The finite-time identity}

\begin{proposition}[Finite-time free energy of the transfer]
\label{f16v72:identity}
The Gaussian free energy $F_N=-\log\operatorname{Tr}T^N$ of the frozen
transfer over $N$ periodic time slices at a static background, with V69's
hard modes, is
\begin{equation}
 F_N(\varepsilon)=N\,E_0(\varepsilon)+\sum_j\log\bigl(1-e^{-N\xi_j(\varepsilon)}\bigr)+{\rm const},\qquad
 E_0(\varepsilon)=\tfrac12\sum_j\xi\bigl(\nu^{(3)}_j(\varepsilon)\bigr),
 \label{f16v72:FN}
\end{equation}
where $\nu^{(3)}_j(\varepsilon)$ are V69's horizontal eigenvalues and
$\cosh\xi_j=1+\nu^{(3)}_j/2$. Its second-order coefficient satisfies
$\Pi_N=\tfrac N2\Pi_L(k)+O(Ne^{-N\xi_{\min}})$.
\end{proposition}
\begin{proof}
At the frozen level $T$ is a product of single-mode Mehler transfers with
eigenvalues $e^{-(n+1/2)\xi_j}$ (V70's check), so
$\operatorname{Tr}T^N=\prod_j\sum_ne^{-N(n+1/2)\xi_j}=\prod_j(2\sinh(N\xi_j/2))^{-1}$;
equivalently, the Matsubara product of the mode is
$\prod_{\omega'}(2\cosh\xi-2\cos\omega')=(2\sinh(N\xi/2))^2$, whose
half-logarithm is $N\xi/2+\log(1-e^{-N\xi})$. The verifier
evaluates \eqref{f16v72:FN} on V69's spectrum: the resulting
$\delta$ at $N=8,16,32$ is $-1.2645,-1.2623,-1.2623$ at $L=8$ and $-1.5936,-1.5489,-1.5468$ at $L=16$, against
$E_0$'s $-1.2623$ and $-1.5468$.
\end{proof}

\subsection{The electric coefficient}

\begin{proposition}[Electric equals magnetic at $a_t=a_s$]
\label{f16v72:electric}
For $L_1=N$, $\Pi^{(4)}(2\pi m/L_1,0)=\Pi^{(4)}(0,2\pi m/N)$ exactly.
\end{proposition}
\begin{proof}
The map $x_1\leftrightarrow t$ is a symmetry of the hypercubic lattice, of
the Wilson action with $a_t=a_s$, of the gauge group, of the chart, of
the exclusions, and of the transverse polarization; it carries the
magnetic background to the electric one. The verifier confirms equality
to $10^{-9}$ at $L_1=N=8$, $m=1,2$.
\end{proof}

The same symmetry argument applies to any hypercubically symmetric
scheme, but not to the transfer scheme, in which time is
distinguished; there the electric coefficient must be computed from a
time-dependent background, which is left open. One route is closed
here: the naive temporal-gauge Gaussian, i.e.\ the four-dimensional
Wilson quadratic form with the temporal links frozen at the identity,
the spatial-link fluctuations of every slice retained, and only the
time-independent gauge transformations projected out, is not the
transfer's Gaussian even at a static background: an exploration at
$L_1=N=8$ gives $\delta=-1.462$ for the static background $(1,0)$ against
the transfer's $-1.264$, and $-1.553$ against $-1.116$ for $(2,0)$. The
time-dependent vertical directions, which are gauge directions only
when the temporal links fluctuate, carry the chart curvature term
$\nabla S\cdot y''$ of V69 as if it were a potential. Two further
explorations close two more routes. Restricting each slice's spatial
fluctuations to V69's horizontal space at that slice's own background
phase, with the temporal links frozen, reproduces the transfer's static
value exactly ($\delta=-1.26446$ at $L_1=N=8$, the finite-$N$ value of
\eqref{f16v72:FN} to five digits) but gives an electric kernel
$\delta=-0.245$ for $(0,1)$, five times smaller than the magnetic one:
the relative rotation of the horizontal spaces of consecutive slices is
counted as electric energy, which the Gauss law would remove. Adding
the temporal links as free Gaussian variables to that construction is
not a gauge fixing: its zero-mode count jumps from $24$ to $0$ with
$\varepsilon$ and the second-order coefficient is meaningless. The
single-step Born--Oppenheimer overlap
$\langle\Omega_s|M_bCM_b|\Omega_{s'}\rangle$ of the hard ground states gives
$\delta=-0.149$, smaller still, since it lacks the propagation of hard
excitations between steps. The transfer scheme's electric coefficient
therefore requires the Gauss-law projection at a time-dependent
background done exactly, which is left open; in the covariant scheme
the question does not arise.

\subsection{Scheme dependence of the off-shell kernel}

\begin{center}\small
\begin{tabular}{@{}lrrrrr@{}}
\toprule
$L_1\times N$ & $(m,n)$ & $\lambda_k+\lambda_\omega$ & $\Pi^{(4)}$ & $\delta^{(4)}$ & transfer $\delta$ at the same $L,N$\\
\midrule
$8\times 8$ & $(0,1)$ & 0.5858 & -799.114 & -1.3322 & \\
$8\times 8$ & $(0,2)$ & 2.0000 & -2368.099 & -1.1563 & \\
$8\times 8$ & $(1,0)$ & 0.5858 & -799.114 & -1.3322 & -1.2645\\
$8\times 8$ & $(1,1)$ & 1.1716 & -1530.722 & -1.2759 & \\
$8\times 8$ & $(2,0)$ & 2.0000 & -2368.099 & -1.1563 & \\
$16\times 16$ & $(1,0)$ & 0.1522 & -3981.820 & -1.5964 & -1.5489\\
\bottomrule
\end{tabular}
\end{center}
At the same $L$ and $N$ the covariant $\delta^{(4)}(k,0)$ exceeds the
transfer value in magnitude by 5.4 percent ($L_1=N=8$) and 3.1 percent
($L_1=N=16$); the finite-$N$ remainder of \eqref{f16v72:FN} is below
$0.3$ percent at $N=8$ and negligible at $N=16$, so the difference is
between the two Gaussian problems themselves. The four-dimensional
horizontal slice has $9V_4+3$ modes against the $6V_4$ of the temporal
gauge; the extra modes are the third polarization, and the naive
Faddeev--Popov factor of the slice, $\det(G(\varepsilon)^{\mathsf T}G(\varepsilon))^{1/2}$,
has second-order coefficient $-2118$ at $L_1=N=8$, an order of magnitude
larger than the discrepancy and of the wrong form. The discrepancy is
instead the slice dependence of the Hessian of a gauge-invariant
function at a non-stationary point: on the quotient, the second
derivative of $S$ at a point where $\nabla S\ne0$ depends on the
transversal used to define the soft coordinate, by terms proportional
to $\nabla S$, i.e.\ to the equations of motion of the background, which
for the plane wave are of the same order $\varepsilon\lambda_k$ as the kernel
itself. In continuum language this is the gauge-parameter dependence of
the off-shell one-loop effective action, which is proportional to the
equations of motion and vanishes on shell; its coefficient of
$\log k^2$ is universal. The two schemes are therefore both legitimate
Born--Oppenheimer definitions of the one-loop coupling at scale $k$,
differing in their finite parts, and the programme must fix one: the
transfer scheme, since the object being renormalized is the transfer.
The decrease of the difference from $L=8$ to $16$ is consistent with the
universality of the logarithm, and is not proved.

\subsection{What this is and is not}

The identity \eqref{f16v72:FN} closes the loop between the Hamiltonian
and Euclidean forms of the same Gaussian problem; the symmetry settles
the electric coefficient in the covariant scheme; the scheme
dependence is a genuine finding about off-shell one-loop kernels that
V69--V71 did not state and that any scale-to-scale theorem must respect:
the effective coupling is defined relative to a slice, and the
remainder of the elimination is slice-dependent by equation-of-motion
terms. This is the same obstruction class as V71's momentum shells,
coordinate dependence, now at the level of the one-loop finite parts.
All statements are one-loop, frozen-Gaussian. The next target is the
transfer scheme's electric coefficient (time-dependent background in
the Doob chain) and the first connected correction.

\subsection{Status}

\emph{Proved}: Propositions \ref{f16v72:identity} and
\ref{f16v72:electric}. \emph{Inherited}: V69's kernel and spectrum, V68's
reduction, V71's density. \emph{Conditional}: the interpretation of the
scheme difference as equation-of-motion terms (consistent with its
size and decrease, not derived); Euclidean covariance beyond lattice
corrections (numerical). \emph{Open}: the transfer scheme's electric
coefficient; the two-loop kernel; the effective transfer below a scale;
N2 in connected form.

\subsection{Verification and preservation}

The new diagnostic is
\begin{center}\small\path{scripts/verify_ym_f16_v72_covariant_kernel.py}.\end{center}
Its hashed content is evenness and mode-count continuity, the exact
electric--magnetic equality, the finite-$N$ identity on V69's spectrum,
and the bound and decrease of the scheme difference. Its floating
diagnostics are the table, the finite-$N$ values and the vertical-Gram
term. It has 13255 bytes and SHA--256
\begin{center}\scriptsize\ttfamily F610C363CFC1A776D5C9EDB40CCA236A0A4D9D5441B97A8E634C1D53EA94C42D.\end{center}
Its canonical hashed payload has SHA--256
\begin{center}\scriptsize\ttfamily 89C5D52914950FC58563B5E0A39775818F5E739736E92AB1D59DECC497942686.\end{center}
Only this terminal insertion and the literal title version differ from V71;
removing them recovers its bytes exactly.

\begin{firewall}
V72 computes the covariant four-dimensional one-loop kernel at a
space-time plane wave, proves the finite-time identity of the transfer's Gaussian
with V69's zero-point energy, proves that the electric and magnetic
coefficients renormalize identically in the covariant scheme at
$a_t=a_s$, and finds that the covariant and transfer schemes differ in the
finite part of the off-shell kernel by a few percent, decreasing with
$L$. It does not go beyond one loop, does not compute the transfer
scheme's electric coefficient, does not bound a remainder, does not
write an effective transfer below a scale, and does not bear on the
physical continuum trajectory or an interacting continuum mass gap.
\end{firewall}
% END F16 V72 APPEND

% BEGIN F16 V73 APPEND -- 2026-09-19
% Insert before the final document-ending command of cumulative V72.
% This insertion leaves every predecessor body and the native operator unchanged.

\clearpage
\section{V73: the one-loop flow of the transfer coefficient, the two scales, and the export interface}
\label{f16v73:context}

V69--V72 produced the pieces of a one-loop renormalization of the Wilson
transfer: the static kernel (V69), its gauge-invariant scale density
(V71), the covariant form and the scheme dependence of its finite part
(V72). This section assembles them into the explicit one-loop flow of
the transfer coefficient, $\gamma_{\rm eff}(\Lambda)$, and draws the two
consequences that the programme clarification asks for at steps 5--8 of
the order of attack: the transformed coupling as a function of the
scale, and the first scale at which control fails. There are two such
scales, and they are far apart. The Landau-type scale $\Lambda_*$, where
the one-loop $\gamma_{\rm eff}$ would reach a value of order one, is
$\log(1/\Lambda_*)=(\gamma-\gamma_*+c)/b$ with $b=$ $0.367$ from the window
density of V71, against $11/(3\pi^2)=0.3715$; its exponent $1/b=$
$2.73$ against $3\pi^2/11=2.692$ is the exponent of the concurrent V55
physical trajectory $L\sim\exp(3\pi^2\gamma/11)$, which the lineage thus
recovers from its own one-loop density. The control scale $\Lambda_c$,
where the small-volume regime of V63 ends,
$(1/\Lambda_c)\log^3(1/\Lambda_c)=\gamma_{\rm eff}(\Lambda_c)^{2/3}$, is at
$\log(1/\Lambda_c)\approx1.3$--$2.0$ for $\gamma$ between $20$ and $400$, i.e.\ at
$L\approx4$--$7$ lattice units, where $\gamma_{\rm eff}$ has moved from
$\gamma$ by less than two units. Everything between $\Lambda_c$ and
$\Lambda_*$, some $\gamma/b$ e-folds, is the range over which the one-loop
flow is computable but the lineage's rigorous control is not
established; the perturbative RG programme is the attempt to extend the
control from $\Lambda_c$ toward $\Lambda_*$ scale by scale, and its
first step, the one-loop flow itself, is now explicit.

\subsection{Dependencies}

V69's kernel and normalization, V71's density and its $L=64$ point, V72's
scheme statement, V63's small-volume condition $L\log^3L\ll\gamma^{2/3}$,
V65--V66's programme statements. From the concurrent file: V55's
trajectory exponent and calibration.

\subsection{The one-loop flow}

\begin{proposition}[One-loop flow of the transfer coefficient]
\label{f16v73:flow}
In the transfer scheme, for the magnetic coefficient at momentum $k$ and
in the frozen Gaussian, the shift accumulated from the ultraviolet down
to the scale $\Lambda$ is
\begin{equation}
 \gamma_{\rm eff}(\Lambda;k)-\gamma=\delta_L(k)-\int_0^\Lambda\rho_L(\Lambda';k)\,\frac{d\Lambda'}{\Lambda'},
 \label{f16v73:flow-eq}
\end{equation}
which decreases monotonically from $0$ at $\Lambda=\infty$ to $\delta_L(k)$
at $\Lambda=0$ (V71: $\rho\le0$ everywhere at $L=64$, $m=1$), and in the
window $|k|\ll\Lambda\ll1$ has the form $c-b\log(1/\Lambda)$ with
$b=$ $0.367$ and $c=$ $-1.209$ (fixed at $\Lambda=0.94$ from the $L=64$ flow).
The identification $b=11/(3\pi^2)$ is numerical, within 1.3 percent.
\end{proposition}
\begin{proof}
\eqref{f16v73:flow-eq} is V71's decomposition \eqref{f16v71:density}
integrated, with the normalization of V69; monotonicity is the sign of
the density; the window form is the plateau of V71's table, read at
$L=64$, $|k|=0.098$.
\end{proof}

\begin{center}\small
\begin{tabular}{@{}rrr@{}}
\toprule
$\Lambda$ & $\gamma_{\rm eff}-\gamma$, $L{=}32$, $|k|{=}0.20$ & $\gamma_{\rm eff}-\gamma$, $L{=}64$, $|k|{=}0.10$\\
\midrule
3.008 & -0.5578 & -0.5572\\
2.015 & -0.8301 & -0.8294\\
1.420 & -1.0372 & -1.0373\\
1.000 & -1.2034 & -1.2053\\
0.704 & -1.3433 & -1.3487\\
0.496 & -1.4700 & -1.4822\\
0.350 & -1.5855 & -1.6102\\
0.246 & -1.6855 & -1.7323\\
0.182 & -1.7534 & -1.8301\\
0.122 & -1.8086 & -1.9438\\
0.090 & -1.8233 & -2.0098\\
0.061 & -1.8263 & -2.0605\\
0.030 & -1.8263 & -2.0723\\
\bottomrule
\end{tabular}
\end{center}

The flow is the running coupling of the lineage in the transfer scheme:
$\gamma_{\rm eff}$ is smaller at longer wavelengths, i.e.\ $g_{\rm eff}^2=4/\gamma_{\rm eff}$
grows toward the infrared at the one-loop rate, with the lattice
ultraviolet part of the density ($\Lambda\gtrsim1$) contributing a
$k$-independent constant, $c$, which is the one-loop matching between
the bare Wilson coefficient and the coefficient at the scale $\Lambda=1$.
V72 showed that $c$ is scheme-dependent by a few percent of itself and
$b$ is not.

\subsection{The two scales}

\begin{proposition}[Landau-type scale and control scale]
\label{f16v73:scales}
With the window form of the flow and a threshold $\gamma_*$ of order one,
$\gamma_{\rm eff}(\Lambda_*)=\gamma_*$ at
\begin{equation}
 \log\frac1{\Lambda_*}=\frac{\gamma-\gamma_*+c}{b},
 \label{f16v73:landau}
\end{equation}
whose exponent $1/b$ is $2.73$ against the concurrent trajectory
exponent $3\pi^2/11=2.692$. The small-volume regime of V63,
$L\log^3L\ll\gamma_{\rm eff}^{2/3}$ with $L=1/\Lambda$, ends at $\Lambda_c$ with
$(1/\Lambda_c)\log^3(1/\Lambda_c)=\gamma_{\rm eff}(\Lambda_c)^{2/3}$, and for
$\gamma\in\{20,50,100,400\}$ the values are
\begin{center}\small
\begin{tabular}{@{}rrrrr@{}}
\toprule
$\gamma$ & $\log(1/\Lambda_c)$ & $\gamma_{\rm eff}(\Lambda_c)$ & $\log(1/\Lambda_*)$ & ratio\\
\midrule
20 & 1.26 & 18.3 & 48.5 & 0.0259\\
50 & 1.46 & 48.3 & 130.3 & 0.0112\\
100 & 1.62 & 98.2 & 266.7 & 0.0061\\
400 & 1.96 & 398.1 & 1085.0 & 0.0018\\
\bottomrule
\end{tabular}
\end{center}
so that $\Lambda_c\gg\Lambda_*$ by a factor $\exp(\gamma/b)$ up to logarithms.
\end{proposition}
\begin{proof}
Solve the window form for $\Lambda_*$; solve the V63 condition with the
flow inserted, by bisection in $\log(1/\Lambda)$; the verifier records
both. The regime between the two scales is where one-loop perturbation
theory is self-consistent ($\gamma_{\rm eff}\gg1$) but V63's control is
not available.
\end{proof}

Three remarks. First, the recovery of the trajectory exponent
$3\pi^2/11$ from the lineage's own density is a consistency check
between the two programmes at one loop: the concurrent V55 defines the
physical trajectory by the one-loop $\beta$-function in the calibration
$\gamma=4/g_0^2$, $a_t=a_s$; the lineage computes the same coefficient
from the vacuum energy of the hard modes and finds it within 1.3 percent.
Second, the constant $c$ is the transfer scheme's one-loop matching;
handed to the concurrent programme it must carry V72's scheme caveat.
Third, the control scale is tiny in lattice units: the lineage's
rigorous statements (V59--V63) cover boxes of a few lattice spacings at
the couplings of interest, and the RG programme's task is exactly the
$\gamma/b$ e-folds beyond, with no step of it yet proved. What is
proved is the one-loop flow across those e-folds in the frozen
Gaussian, and what is not is the control of the remainder along it.

\subsection{Export interface}

For the concurrent programme's uniformity constant $\varepsilon_{L,\gamma}$ and
its calibration the following are the lineage's one-loop statements,
each with its status. \emph{Proved (frozen Gaussian, transfer scheme)}:
the shift of $\gamma$ at momentum $k$ is $\delta_L(k)$, of order one,
negative, with the values of V69's table, and its scale density is
V71's. \emph{Numerical}: $b=$ $0.367$, $c=$ $-1.209$. \emph{Proved}: the
covariant scheme has equal electric and magnetic shifts (V72); the
transfer scheme's electric shift is open. \emph{Consequence}: the
relative correction to the bare coefficient at scale $\Lambda$ is
$(c-b\log(1/\Lambda))/\gamma$, so the concurrent trajectory
$L\sim e^{3\pi^2\gamma/11}$ is the locus where this correction is of order
one, and any uniformity constant defined at fixed $L$ and $\gamma\to\infty$
sees only the $O(1/\gamma)$ relative shift.

\subsection{Status}

\emph{Proved}: Propositions \ref{f16v73:flow} and \ref{f16v73:scales} as
statements about the frozen-Gaussian one-loop flow with V71's density as
input. \emph{Inherited}: V69, V71, V72, V63. \emph{Conditional}: the
identification $b=11/(3\pi^2)$ and hence the exact exponent
(numerical, $L=64$ window). \emph{Open}: every step between $\Lambda_c$
and $\Lambda_*$ beyond one loop: the connected correction, the remainder
of an elimination, the electric coefficient in the transfer scheme,
N2.

\subsection{Verification and preservation}

The new diagnostic is
\begin{center}\small\path{scripts/verify_ym_f16_v73_one_loop_flow.py}.\end{center}
Its hashed content is the sign of the density, the window slope and its
exponent within fifteen percent of the one-loop values, and the
ordering $\Lambda_c>\Lambda_*$ for four couplings. Its floating diagnostics
are the flow and the scales. It has 5329 bytes and SHA--256
\begin{center}\scriptsize\ttfamily F9FEA1DE5DAE7C5C545B365BDF67418922016C082808842806D0DFA42427A842.\end{center}
Its canonical hashed payload has SHA--256
\begin{center}\scriptsize\ttfamily 0C1CA78E494B09042F615C220F49497BFB737925CB7133B419EF6642AC25B942.\end{center}
Only this terminal insertion and the literal title version differ from V72;
removing them recovers its bytes exactly. V75 is the next companion
checkpoint.

\begin{firewall}
V73 assembles the one-loop flow of the transfer coefficient from V69--V72,
recovers the concurrent trajectory exponent from the lineage's density,
and locates the control scale of V63 far above the Landau-type scale.
It proves nothing beyond one loop, does not extend the rigorous control
past the small-volume regime, and does not bear on the physical
continuum trajectory or an interacting continuum mass gap.
\end{firewall}
% END F16 V73 APPEND

% BEGIN F16 V74 APPEND -- 2026-09-19
% Insert before the final document-ending command of cumulative V73.
% This insertion leaves every predecessor body and the native operator unchanged.

\clearpage
\section{V74: the first connected correction to the static kernel}
\label{f16v74:context}

V67 located the obstruction of N2 in the norm: beyond the frozen Gaussian
the transfer is controlled by connected quantities along the Doob
chain, not by $L^2$ distances. This section computes the first such
quantity: the order-$\gamma^{-1}$ correction to the hard vacuum energy
and to the static one-loop kernel of V69, from the cubic and quartic
vertices of the plaquette action at the background and from the Haar
density, by the cumulant expansion of the stationary Mehler chain. Two
things are established. First, the expansion itself: the correction is
$\gamma^{-1}[\langle S_4\rangle+\langle|u|^2\rangle/3-\tfrac12\sum_{n\in\mathbb Z}\langle S_3(u_0)S_3(u_n)\rangle_c]$,
a finite sum of connected Gaussian correlations of the chain, whose
cubic part is the sunset sum over hard-mode triples with the factor
$\coth((\xi_a+\xi_b+\xi_c)/2)$ and whose linear (tadpole) remainder is a
static susceptibility; the formula is validated on the one-mode
anharmonic transfer to $0.2$ percent. Second, the numbers: at $L=8$,
$m=1$, the two-loop shift of the magnetic coefficient is
$\delta_2/\gamma$ with $\delta_2=$ $-0.089$, against the one-loop
$\delta_1=-1.262$, so the two-loop term is $0.070/\gamma$, i.e.\ 0.1 percent of the one-loop term
at $\gamma=$ $100$; the vacuum-energy correction per volume is
$-0.480$ per site at $\varepsilon=0$, dominated by the quartic term, with the cubic
sunset 15.1 percent of it. This is the first quantity of the lineage
beyond the frozen Gaussian, and it fixes the size of the remainder that
any scale-to-scale statement at one loop must carry: relative order
$\gamma^{-1}$ with the coefficient computed here. The orbit-space measure
term is not included and is the stated gap.

\subsection{Dependencies}

V69's horizontal modes, exclusions and block reduction; V70's Mehler chain
(variance $1/(2\omega)$, correlation $e^{-n\xi}$); V64's Haar expansion
$\log J=-|y|^2/3-\dots$; V67's third-order structure of the plaquette
word. The cumulant expansion of a product of transfer matrices is
classical.

\subsection{The cumulant expansion}

Write the per-slice weight of $\operatorname{Tr}T^N$ in the scaled hard
variables $u=\gamma^{1/2}y$ as
$e^{-u^{\mathsf T}Au/2}\,e^{-W(u)}$,
$W=\gamma^{-1/2}S_3(u)+\gamma^{-1}S_4(u)+\gamma^{-1}|u|^2/3+O(\gamma^{-3/2})$,
where $S_3,S_4$ are the cubic and quartic Taylor terms of
$\sum_p(1-\tfrac12\operatorname{Tr}U_p)$ at the background in the right chart,
restricted to the horizontal modes, and the last term is the Haar
density of V64.

\begin{proposition}[First connected correction]
\label{f16v74:cumulant}
The vacuum energy of the frozen transfer with the vertices is
$E_0+\gamma^{-1}E_2+O(\gamma^{-3/2})$ (the odd order vanishes by the
reflection $u\to-u$ at $\varepsilon=0$ and is even in $\varepsilon$), with
\begin{equation}
 E_2=\langle S_4\rangle+\tfrac13\langle|u|^2\rangle
 -\tfrac12\sum_{n\in\mathbb Z}\langle S_3(u_0)S_3(u_n)\rangle_c ,
 \label{f16v74:E2}
\end{equation}
expectations in the stationary Gaussian chain of the hard modes. In the
mode basis, with $T_{abc}$ the cubic vertex contracted with three hard
modes (transverse momentum conserved) and $t_c$ its equal-time
self-contraction,
\begin{equation}
 \sum_{n}\langle S_3S_3\rangle_c
 =6\sum_{abc}\frac{|T_{abc}|^2\coth\frac{\xi_a+\xi_b+\xi_c}2}{8\omega_a\omega_b\omega_c}
 +9\sum_c\frac{|t_c|^2\coth\frac{\xi_c}2}{2\omega_c},
 \label{f16v74:sunset}
\end{equation}
and $\langle S_4\rangle=3\sum_p c^{(p)}_{ijkl}C_{ij}C_{kl}/24$ with $C$ the
equal-time link covariance of the plaquette's four links.
\end{proposition}
\begin{proof}
Second-order cumulant expansion of $\langle e^{-\sum_tW(u_t)}\rangle$ in the
chain, divided by $N$; the odd cumulants vanish at $\varepsilon=0$ and the
cross term $\gamma^{-3/2}\langle S_3S_4\rangle$ is higher order. For cubic
polynomials in jointly Gaussian variables the connected two-point
function at time separation $n$ is $6\,c\,c'\,G(n)^3+9\,(c\cdot G(0))(c'\cdot G(0))\,G(n)$
by Wick's theorem, with $G(n)=\operatorname{diag}(e^{-n\xi}/(2\omega))$ in the mode
basis; summing $e^{-a|n|}$ over $n\in\mathbb Z$ gives $\coth(a/2)$. The
verifier checks the formula against the exact top eigenvalue of the
one-mode transfer $e^{-V/2}e^{\partial^2/2}e^{-V/2}$, $V=\lambda q^2/2+g_3q^3+g_4q^4$,
finding the ratio 1.0006 for the cubic term at $g_3=0.01$, and the
first-order quartic term within three percent (its own second order).
\end{proof}

The two-loop kernel is the second-order coefficient of $E_2$ in the
background amplitude, $\Pi_2(k)$, and the two-loop shift of the magnetic
coefficient is $\delta_2(k)/\gamma$ with $\delta_2=2\Pi_2/(\lambda_kV)$, in the
normalization of V69 (energy, not $\tfrac12\sum\xi$).

\subsection{Results}

\begin{center}\small
\begin{tabular}{@{}lrrrrr@{}}
\toprule
$L$ & $E_2(0)/V$ & $\langle S_4\rangle/V$ & Haar$/V$ & sunset$/V$ & $\delta_2(2\pi/L)$\\
\midrule
4 & -0.4532 & -0.6760 & 0.2776 & 0.1097 & -0.0722\\
8 & -0.4803 & -0.7093 & 0.3013 & 0.1446 & -0.0886\\
\bottomrule
\end{tabular}
\end{center}
At $L=8$, $m=1$ the second-order coefficients of the pieces are: quartic -14.697, Haar 29.292, sunset 55.771, tadpole 0.004 (in units of $\Pi_2$; $\Pi_2=-13.293$).

Three remarks. First, the vacuum-energy correction is negative and
dominated by the quartic term, as for the plaquette expectation in
lattice perturbation theory, where the two-loop coefficient of
$\langle1-\tfrac12\operatorname{Tr}U_p\rangle$ is of the same sign and order
relative to the one-loop one; the sunset is a small fraction because
the cubic vertex is Wick-ordered up to its tadpole (V67) and its
three-propagator contraction carries $1/(8\omega^3)$. Second, the
two-loop kernel $\delta_2$ enters the flow of V73 as
$\gamma_{\rm eff}=\gamma+\delta_1(k)+\delta_2(k)/\gamma+\dots$, and its ratio to
$\delta_1$ gives the coupling below which the one-loop flow is
quantitatively reliable: for $|\delta_2/(\gamma\delta_1)|<0.1$ one needs
$\gamma>$ $1$. Third, what is missing at this order is the
orbit-space measure, the vertical Gram of V68 expanded to second order
in the fluctuation; it is of the same order $\gamma^{-1}$ and is recorded
as the gap of this section, together with the finite-difference
accuracy of the vertex tensors (relative $10^{-4}$).

\subsection{Status}

\emph{Proved}: Proposition~\ref{f16v74:cumulant} for the frozen transfer
with the vertices, on the hard modes of V69. \emph{Inherited}: V69, V70,
V64, V67. \emph{Conditional}: the numerical values (finite differences,
$L=4,8$). \emph{Open}: the orbit-space measure term; the two-loop
density in proper time; the remainder bound; N2.

\subsection{Verification and preservation}

The new diagnostic is
\begin{center}\small\path{scripts/verify_ym_f16_v74_two_loop.py}.\end{center}
Its hashed content is the one-mode validation of the cumulant formula,
evenness at $L=4$, and the signs of the three pieces. Its floating
diagnostics are the table. It has 12066 bytes and SHA--256
\begin{center}\scriptsize\ttfamily 462E4B26A5D1A13DB8910C3EBCA43B97718ABF25B81167E029CA9E72A59C9EA4.\end{center}
Its canonical hashed payload has SHA--256
\begin{center}\scriptsize\ttfamily 970960B9B6C7D3589E39C1AE4D944C682A2C7423C2706F8FE1352FC6C7E695D1.\end{center}
Only this terminal insertion and the literal title version differ from V73;
removing them recovers its bytes exactly. V75 is the companion
checkpoint; the concurrent file has advanced to V62.

\begin{firewall}
V74 proves the first connected correction to the hard vacuum energy and
to the static one-loop kernel as a cumulant expansion of the Mehler
chain, validates it on a one-mode transfer, and computes it at $L=4$ and
$8$ without the orbit-space measure term. It does not bound the
remainder of the expansion, does not include the measure term, and does
not bear on the physical continuum trajectory or an interacting
continuum mass gap.
\end{firewall}
% END F16 V74 APPEND

% BEGIN F16 V75 APPEND -- 2026-09-19
% Insert before the final document-ending command of cumulative V74.
% This insertion leaves every predecessor body and the native operator unchanged.

\clearpage
\section{V75: companion checkpoint at the concurrent V62 and five-version review}
\label{f16v75:context}

This is the scheduled companion checkpoint. The concurrent file has
advanced from V58 to V62 since V70; its V1--V58 body is preserved
verbatim in the V62 file (checked byte-wise up to the title). Its four
new sections stay on the local-vacuum branch: V59 frees spatial links
along a time ribbon and proves a boundary-uniform conditional
contraction with factor $\tanh\gamma$ together with the exterior-variance
payment that a frozen first slice forces to vanish; V60 rules out V59's
proposed uniform independent-ribbon capture fraction on a continuous
physical source core under its fast clock, and proves the exact mixed
predictor-Gram identity and the retained-marginal preservation of gauge
invariance and reflection positivity; V61 proves alignment domination
for compact vector-spin integrals, the all-exterior ribbon frustration
bound and the transverse Ward susceptibility identity; V62 proves the raw
second-moment graph Ward lower bound, the coherent Hessian with open-path
mismatches and closed holonomies, and the literal two-endpoint Wilson
holonomy identity. None of these imports anything from V68--V74; the
parallel V67 cubic-exponential step remains explicitly excluded there,
and V69's correction of it stands. The five versions V70--V74 of this
lineage are reviewed against the order of attack, and the export
interface of V73 is updated with V74's two-loop term.

\subsection{Companion checkpoint}

The concurrent file read is
\path{Renewal_Yang_Mills_Mass_Gap_Research_Notes_V62.tex}, 1951050
bytes, 40718 lines, SHA--256
\begin{center}\scriptsize\ttfamily 09A8B55FDDE1A8E4D47936CAFECF25A478BCADC5700AB961A7331D892F207832,\end{center}
with sections ``V59: buffered Wilson ribbons, transported innovations, and exterior predictability; V60: the microscopic-ribbon obstruction and a source-preserving exterior reduction; V61: exact retained-ribbon curvature and transverse Ward cancellation; V62: joint Wilson ladders and paid transport holonomy''. The four interface obligations. \emph{Local-source
noncollapse}: unchanged from V70, the leading-order Gram
$\frac32c_0^2\gamma^{-2}$ per writer with $\gamma$-free relative
decrements; the concurrent V59--V60 now work with ribbon-conditional
contractions whose factor $\tanh\gamma$ is the one-link Wilson law's own
and is $\gamma$-independent in the limit, consistent with V70's
observation that the entrance decrement is not where the
$\gamma$-dependence lies. \emph{Convergence}: the concurrent V60's
obstruction (vanishing conditional variance on a continuous source core
under a fast clock) is a statement about physical source families and
does not involve this lineage. \emph{Hard-sector propagation}: V71's
proper-time density is the lineage's statement about which scales
renormalize the hard sector; nothing in V59--V62 uses it yet.
\emph{Rate for $\varepsilon_{L,\gamma}$}: the export interface now reads
\begin{center}\small
\begin{tabular}{@{}lrl@{}}
\toprule
quantity & value & status\\
\midrule
$\delta_1(2\pi/8)$ at $L=8$ & $-1.2623$ & proved (frozen Gaussian, transfer scheme)\\
$\delta_2(2\pi/8)$ at $L=8$ & $-0.0886$ & proved (cumulant expansion), measure term omitted\\
$\delta_2/\delta_1$ & $0.070$ & numerical\\
$b$ (window) & $0.367$ & numerical, vs $11/(3\pi^2)=0.3715$\\
$c$ (matching at $\Lambda=1$) & $-1.209$ & numerical, scheme dependent\\
exponent $1/b$ & $2.73$ & numerical, vs $3\pi^2/11=2.692$\\
\bottomrule
\end{tabular}
\end{center}
so that the relative correction to the bare coefficient at scale
$\Lambda$ is $(c-b\log(1/\Lambda))/\gamma+O(\gamma^{-2})$ with the
$O(\gamma^{-2})$ coefficient of relative size $0.070$ of the one-loop one.

\subsection{Five-version review, V70--V74}

V70 calibrated the concurrent capture target at leading order and
verified the V60--V69 artifacts. V71 found that momentum and eigenvalue
shells fail as scale decompositions and that proper time is the
gauge-invariant one, with the density's window value at the
asymptotic-freedom coefficient; it also corrected V69's zero-point
normalization and its $L=64$ entry. V72 computed the covariant kernel,
proved the finite-time identity of the transfer's Gaussian and the
electric--magnetic equality in the covariant scheme, and found the
off-shell kernel scheme-dependent by a few percent, closing four routes
to the transfer scheme's electric coefficient. V73 assembled the one-loop
flow, recovered the concurrent trajectory exponent from the lineage's
density, and located the control scale of V63 at a few lattice spacings
against the Landau-type scale at $e^{\gamma/b}$. V74 computed the first
connected correction by the cumulant expansion of the Mehler chain,
validated on a one-mode transfer, and found the two-loop kernel at
$0.070$ of the one-loop one.

Against the order of attack of V66: step 1 (N2) remains open, now with
the connected quantities it needs identified and the first of them
computed; step 2 (uniform block statement) unchanged; step 3 (one
elimination step) has its one-loop content in proper time and its
first two-loop coefficient, but no elimination with a remainder; step 4
(locality and remainder) untouched; step 5 (transformed coupling) is
V73's flow; step 6 (iteration) is trivial at one loop; step 7 (first
failing scale) is V73's two scales, with the honest reading that the
rigorous control ends at $\Lambda_c$; step 8 (export) is the table above.
The obstructions met were of the coordinate class three times (momentum
shells, the temporal-gauge Gaussian, the slice dependence of the
off-shell kernel) and of the numerical-control class twice (the
zero-point factor, the $L=64$ threshold), and each was recorded in the
version that met it.

Two errata were recorded in the five versions: V69's factor of two and
its $L=64$ row (V71). The V69 note itself already carried the correct
definition; the verifier column was mislabeled and is documented as such.

\subsection{Integrity inventory}

The verifier checks that the five artifacts of V70--V74 record the SHA-256
of their present verifier sources and passing hashed payloads, and that
the concurrent V62 file preserves the V58 body. All checks pass.

\subsection{Next}

V76 begins the connected programme proper: the orbit-space measure term
at order $\gamma^{-1}$, which completes V74, and the proper-time density
of the two-loop kernel, which tells whether the two-loop term is also
scale-local. V80 is the next ten-version review.

\subsection{Status}

\emph{Proved}: nothing new; the integrity and preservation checks.
\emph{Inherited}: V70--V74 as stated. \emph{Conditional}: the interface
numbers as marked. \emph{Open}: as in V74.

\subsection{Verification and preservation}

The new diagnostic is
\begin{center}\small\path{scripts/verify_ym_f16_v75_checkpoint.py}.\end{center}
Its hashed content is the artifact integrity, the concurrent body
preservation, and the smallness of the two-loop term. It has 5173
bytes and SHA--256
\begin{center}\scriptsize\ttfamily F1B4EA08838833A24123302CB24CA08D9E3318F732A493ADC8DFA98D538710DF.\end{center}
Its canonical hashed payload has SHA--256
\begin{center}\scriptsize\ttfamily 8DDA922C314A27600978A5F29BC4C8A25A7BE56225B52F4F4BDED8CDC0BA14EC.\end{center}
Only this terminal insertion and the literal title version differ from V74;
removing them recovers its bytes exactly.

\begin{firewall}
V75 records the companion checkpoint at the concurrent V62, reviews
V70--V74, updates the export interface with the two-loop term, and
verifies the artifacts. It proves nothing new and does not bear on the
physical continuum trajectory or an interacting continuum mass gap.
\end{firewall}
% END F16 V75 APPEND

% BEGIN F16 V76 APPEND -- 2026-09-19
% Insert before the final document-ending command of cumulative V75.
% This insertion leaves every predecessor body and the native operator unchanged.

\clearpage
\section{V76: the two-dimensional test of the one-loop recipe, the slice dependence of the constrained determinant, and the kinetic quartic term}
\label{f16v76:context}

V72 found that two evaluations of the one-loop kernel at the same
background differ by a few percent and attributed the difference to the
choice of slice through a non-stationary point. Before going further,
that attribution is tested where everything is exactly solvable: the
two-dimensional $SU(2)$ Wilson theory, whose plaquette variables carry the
product Haar measure, so that the constrained one-loop energy of a
background can be computed exactly in plaquette space and compared with
the link-space recipes of V69--V72. Four things are found. First, the
orbit-volume factor (the vertical Gram, or Faddeev--Popov determinant of
the horizontal slice) is not a one-loop energy term: including it
changes the kernel by a factor of order $10^3$ and the exact answer does
not contain it; the link-to-plaquette Jacobian equals it to 0.29 percent,
which is the Hodge identity behind the cancellation. Second, the
link-slice determinant of V69's recipe and the exact plaquette-slice
determinant differ by a factor $31$: the finite part of the
constrained one-loop determinant at a non-stationary background depends
on the slice by more than an order of magnitude, not by a few percent,
and V72's closeness was fortuitous. Third, the exact two-dimensional
answer is the local compactness term $-\tfrac1{12}\sum_p|F_p|^2$, i.e.\ a
shift $\delta\gamma=-1/6$ with no scale dependence, as it must be in two
dimensions; the plaquette slice is the gauge-invariant soft coordinate,
and the link slice is not. Fourth, V74's order-$\gamma^{-1}$ correction is
completed by the quartic term of the native cosine kinetic kernel, which
V74 omitted; it is $-41.11$ at $L=8$, against V74's quartic action term
$-363.16$, and the corrected two-loop kernel is $\delta_2=$ $-0.1503$ (V74:
$-0.0886$). The consequences for the programme are drawn: the universal
content of V69--V74 is the logarithmic coefficient $b$ (V71, V73), which
is slice independent; the matching constant $c$, the finite parts of
$\delta_1(k)$, and $\delta_2$ are quantities of the link slice, and a
transformed coupling fit for export must be defined through a
gauge-invariant soft coordinate, of which the plaquette field is the
two-dimensional model and V59's Polyakov coordinate the
three-dimensional candidate.

\subsection{Dependencies}

V69's chart, exclusions and recipe; V72's slice statement; V74's cumulant
expansion and pieces; V64's Haar expansion; the exact solubility of
two-dimensional lattice gauge theory (plaquette variables independent
with product Haar measure on a torus up to the holonomy constraint) is
classical.

\subsection{The two-dimensional test}

On the $L\times L$ torus, $L=8$, with the background
$b=\varepsilon\cos(kx_0)e_{\rm col}$ on the direction-$1$ links, $k=2\pi/L$, the
plaquette field is $F_p=\varepsilon[\cos(k(x_0+1))-\cos(kx_0)]e_{\rm col}$. In
plaquette space the constrained one-loop energy is
\begin{equation}
 E_F(\varepsilon)=\tfrac12\log\det{}'\bigl(PH_FP\bigr)-\sum_p\log J(F_p),\qquad
 H_p=\cos|F|\,\hat F\hat F^{\mathsf T}+\frac{\sin|F|}{|F|}(1-\hat F\hat F^{\mathsf T}),\quad
 J=\Bigl(\frac{\sin|F|}{|F|}\Bigr)^2,
 \label{f16v76:EF}
\end{equation}
with $P$ removing the background's cos and sin profiles, its two
colour-orbit profiles, and the three torus constraints. Its second-order
coefficient is $-1.221$, against the local estimate
$-\tfrac1{12}\sum_p|F_p|^2/\varepsilon^2=$ $-1.562$, the difference being the
projected directions. The link-space quantities are V69's recipe
$A=\tfrac12\sum\log\nu$ over the horizontal hard modes of the chart
Hessian, the vertical Gram $g=\tfrac12\log\det'(G^{\mathsf T}G)$, the link
Haar density $h_L$, and the Jacobian $j_F=\tfrac12\log\det'(DF\,DF^{\mathsf T})$
of the link-to-plaquette map at the background; their second-order
coefficients are
\begin{center}\small
\begin{tabular}{@{}lrrrrr@{}}
\toprule
& $[E_F]_2$ (exact) & $[A]_2$ & $[g]_2$ & $[j_F]_2$ & $[h_L]_2$\\
\midrule
$L=8$, $m=1$ & $-1.221$ & $-37.78$ & $-2125.7$ & $-2119.5$ & $-10.67$\\
\bottomrule
\end{tabular}
\end{center}

\begin{proposition}[What the two-dimensional test establishes]
\label{f16v76:twod}
(i) $[j_F]_2=[g]_2$ to 0.29 percent: the link-to-plaquette Jacobian and the
orbit volume cancel in the change of variables, so the constrained
one-loop energy in link space carries neither. (ii) $[A-g]_2$ differs
from $[E_F]_2$ by a factor of order $10^3$: the orbit-volume factor is not
a one-loop energy term, as the Hamiltonian argument of V72 asserted.
(iii) $[A]_2$ differs from $[E_F]_2$ by the factor $31$: the
constrained one-loop determinant depends on the slice by more than an
order of magnitude at this background. (iv) The exact coefficient is the
compactness term of the plaquette variables, $\delta\gamma=-1/6$ up to the
projected directions, with no dependence on $k$.
\end{proposition}
\begin{proof}
All four are finite-dimensional computations of the verifier; (i) is
the Hodge identity between the plaquette and site Laplacians at
$\varepsilon=0$, deformed continuously; (iv) is the expansion of
\eqref{f16v76:EF}: $\tfrac12\log\det H_p=-\tfrac14|F|^2-\tfrac16|F|^2$ and
$-\log J=\tfrac13|F|^2$ to second order.
\end{proof}

The mechanism of (iii) is the one V72 named: at a non-stationary point
the Hessian of a function restricted to a slice depends on the slice
through $\nabla S\cdot({\rm second\ fundamental\ form})$, and here
$\nabla S$ is of the same order as the kernel; in addition, the
link-slice soft coordinate, the amplitude of a link mode, is not a
gauge-invariant function, whereas the plaquette-slice coordinate is. An
attempt to impose the plaquette-slice constraint in link space with the
Lagrange-corrected Hessian $P(H-\lambda\nabla^2g)P$ gives yet another number
(an exploration finds $-304$), because the plaquette profile $\langle w,F\rangle$
rotates under gauge transformations at $\varepsilon\ne0$ and the constraint
surfaces differ again. The lesson is not that one of these is wrong; each
is the exact one-loop energy of a different constrained integral. The
lesson is that ``the one-loop shift of $\gamma$ at momentum $k$'' is not a
slice-invariant number, and only its universal part is: in four
dimensions the coefficient of the logarithm, in two dimensions nothing,
consistent with the wild variation found here.

\subsection{Consequences for V69--V75}

The logarithmic slope $b$ of V71 and V73 is slice independent and is the
lineage's proved-numerical result at one loop. The constant $c$, the
values $\delta_1(k)$ at fixed $k$, the finite-$N$ identity's numbers, and
V74's $\delta_2$ are link-slice quantities and are so marked from here on;
V73's export interface is amended accordingly: what is exportable is
$b$, the sign and monotonicity of the flow, and the statement that the
finite parts require a gauge-invariant soft coordinate. V72's finding
that the covariant and transfer slices differ by a few percent is
consistent with this and no longer needs an explanation of its size.
The programme's step 5 (the transformed coupling) must therefore be
posed for a gauge-invariant coordinate; the two-dimensional model shows
what such a coordinate looks like (the plaquette field, exactly
integrable), and in three dimensions the candidates are the plaquette
field with its Bianchi constraint and V59's Polyakov trace, for which the
small-volume coordinate theorem of V59--V63 already exists.

\subsection{The kinetic quartic term}

V74's order-$\gamma^{-1}$ correction took the kinetic factor of the
transfer as the Gaussian $e^{-|y-y'|^2/2}$ with the Haar density of each
slice; the native kinetic factor is the Wilson cosine
$e^{-\gamma(1-\cos\theta)}$, $\theta=|\log(e^{y}e^{-y'})|$, whose quartic Taylor
term contributes at the same order:
\begin{equation}
 W_{\rm kin}=\gamma\,f_4(y,y')=\gamma^{-1}f_4(u,u'),\qquad
 \langle W_{\rm kin}\rangle=\gamma^{-1}\sum_{\rm links}\frac{3}{24}\sum c^{(4)}_{ijkl}K_{ij}K_{kl},
 \label{f16v76:kin}
\end{equation}
with $c^{(4)}$ the quartic tensor of $1-\cos\theta$ in $(y,y')$ (link
independent, by finite differences) and $K$ the two-time covariance of
the link's fluctuations in the Mehler chain. At $L=8$ it is $-41.11$ at
$\varepsilon=0$ against V74's $\langle S_4\rangle=$ $-363.16$, Haar $154.28$ and sunset
$74.02$, so the corrected two-loop vacuum energy is $E_2=$ $-287.00$ per
volume $512$ and the corrected two-loop kernel $\delta_2=$ $-0.1503$ (V74:
$-0.0886$). This completes the order-$\gamma^{-1}$ terms of the action and
the measure; the orbit-volume term, which V74 listed as its gap, is shown
above not to belong there, and no further term is known to be missing.

\subsection{Status}

\emph{Proved}: Proposition~\ref{f16v76:twod} as a finite-dimensional
statement at $L=8$; the kinetic term \eqref{f16v76:kin}. \emph{Inherited}:
V69, V72, V74. \emph{Conditional}: none. \emph{Open}: the one-loop kernel
with a gauge-invariant soft coordinate in three dimensions; N2.

\subsection{Verification and preservation}

The new diagnostic is
\begin{center}\small\path{scripts/verify_ym_f16_v76_slice_test.py}.\end{center}
Its hashed content is the Jacobian identity, the exclusion of the
orbit-volume term, the slice dependence, the compactness estimate, and
the sign and size of the kinetic term. It has 12583 bytes and
SHA--256
\begin{center}\scriptsize\ttfamily B27EB7341B46809D1B22F0C829D944D2DBE112658FA9BAA56100E7AD246063D4.\end{center}
Its canonical hashed payload has SHA--256
\begin{center}\scriptsize\ttfamily B6397132A3A854800609AD982076983C79D37EB5F1A6D13773C126BED4193359.\end{center}
Only this terminal insertion and the literal title version differ from V75;
removing them recovers its bytes exactly.

\begin{firewall}
V76 tests the one-loop recipe of V69--V72 against the exactly solvable
two-dimensional theory, finds that the orbit-volume factor is not a
one-loop term and that the finite part of the constrained determinant
is slice dependent by more than an order of magnitude, marks the
finite parts of V69--V74 as link-slice quantities with the logarithmic
coefficient as the invariant content, and completes V74 with the kinetic
quartic term. It does not compute the gauge-invariant-coordinate kernel
in three dimensions, does not bound a remainder, and does not bear on
the physical continuum trajectory or an interacting continuum mass gap.
\end{firewall}
% END F16 V76 APPEND

% BEGIN F16 V77 APPEND -- 2026-09-19
% Insert before the final document-ending command of cumulative V76.
% This insertion leaves every predecessor body and the native operator unchanged.

\clearpage
\section{V77: the structure of the semiclassical expansion of the hard energy and the volume dependence of its remainder}
\label{f16v77:context}

V74 and V76 computed the order-$\gamma^{-1}$ correction to the hard
Born--Oppenheimer energy. This section states what the expansion is as
a mathematical object, which of its properties are proved, and where
the volume enters. The main points are three. First, the cumulant
expansion of the Perron eigenvalue of the frozen transfer with the
polynomial vertices is an identity of formal power series in
$\gamma^{-1/2}$ whose coefficients are time-integrated connected
correlations of the Mehler chain; the odd coefficients vanish at
$\varepsilon=0$ by the reflection $u\to-u$, and at $\varepsilon\ne0$ the
$\gamma^{-3/2}$ coefficient is odd in $\varepsilon$ and so absent from the even
energy. The verifier confirms the resulting $O(\gamma^{-2})$ remainder on
one- and two-mode transfers, with fitted exponents $1.99$ and
$2.00$. Second, the $\gamma^{-2}$ coefficient is a sum over connected
graphs with at most four vertices, and under V63's covariance decay
every such graph's lattice sum is bounded by a constant times the
volume times a power of $\log L$: the coefficient has a finite density,
uniformly in $L$ up to logarithms. This is proved by power counting for
the cubic-vertex graphs, the tetrahedron and the double-edge ring, and
for the mixed graphs with a quartic vertex. Third, what is not proved:
that the expansion is asymptotic for the compact model, i.e.\ that the
remainder after the $\gamma^{-1}$ term is bounded by $C\,V\gamma^{-2}(\log L)^a$
with $C$ independent of $L$; the missing pieces are the contribution of
the complement of the sup-norm event and the sixth-order chart
remainder, both of which V66 and V64 bound at fixed $L$ but not with
$L$-uniform constants in the required form. This is the precise
residual content of N2 at order $\gamma^{-1}$: the frozen Gaussian plus its
first connected correction describe the hard sector up to an error
whose density is finite per unit volume if and only if the tail and
chart contributions are, and the tools to decide this exist in
V63--V66.

\subsection{Dependencies}

V45's frozen Gaussian and Mehler chain, V61--V62's chaos and Wick
machinery, V63's decay $|C_3(x)|\le(\kappa_0+\kappa_1\log L)/\max_ix_i^2$,
V64's Haar and kinetic expansions, V66's sup-norm event, V74's cumulant
formula and V76's kinetic term.

\subsection{The cumulant identity}

Let $T_0$ be the frozen transfer on the hard modes at soft coordinate
$\varepsilon$ (V69), with Perron eigenvalue $e^{-E_0^{(0)}}$, and let
$W=\gamma^{-1/2}W_3+\gamma^{-1}W_4+\gamma^{-3/2}W_5+\gamma^{-2}W_6+\dots$ be the
polynomial vertex of the per-slice weight and the kinetic factor
(V74, V76), Wick-ordered with its tadpoles carried as lower-degree
vertices.

\begin{proposition}[Cumulant identity]
\label{f16v77:cumulant}
As formal power series in $\gamma^{-1/2}$,
\begin{equation}
 E_0(\gamma)=E_0^{(0)}+\sum_{m\ge1}\frac{(-1)^{m+1}}{m!}\,\kappa_m,\qquad
 \kappa_m=\lim_{N\to\infty}\frac1N\,\kappa_m\Bigl(\sum_{t=1}^NW(u_t)\Bigr),
 \label{f16v77:series}
\end{equation}
where $\kappa_m(\cdot)$ is the $m$-th cumulant in the stationary chain, and
each $\kappa_m$ is a sum over connected graphs on $m$ vertices, of degrees
given by the vertices' orders, of products of chain propagators
$C(x-x';n)$ summed over positions and time differences. At $\varepsilon=0$
the coefficients of odd powers of $\gamma^{-1/2}$ vanish; at $\varepsilon\ne0$
the coefficient of $\gamma^{-3/2}$ is odd in $\varepsilon$, hence
\begin{equation}
 E_0(\gamma;\varepsilon)=E_0^{(0)}(\varepsilon)+\gamma^{-1}E_2(\varepsilon)+\gamma^{-2}E_4(\varepsilon)+\dots,
 \label{f16v77:even}
\end{equation}
with $E_2$ the expression of V74 and V76.
\end{proposition}
\begin{proof}
$\operatorname{Tr}(T_0e^{-W})^N=\operatorname{Tr}T_0^N\,\langle e^{-\sum_tW(u_t)}\rangle_N$
with the periodic chain expectation, and $\log\langle e^{-X}\rangle=\sum_m(-1)^m\kappa_m(X)/m!$;
dividing by $N$ and letting $N\to\infty$, the Perron eigenvalue dominates
the trace and the cumulants become time-integrated because the chain
is stationary with a spectral gap $\xi_{\min}>0$. The graph expansion is
Wick's theorem for cumulants of Wick-ordered polynomials. Parity: the
frozen chain at $\varepsilon=0$ is invariant under $u\to-u$, $W_3,W_5$ are odd,
$W_4,W_6$ even, so every cumulant with an odd total degree vanishes; at
$\varepsilon\ne0$ the combined reflection $(u,\varepsilon)\to(-u,-\varepsilon)$ with the
colour reflection is a symmetry, so the coefficient of $\gamma^{-3/2}$,
which is odd in $u$, is odd in $\varepsilon$, while $E_0$ is even in $\varepsilon$
(V69). The verifier realizes the identity on the one-mode transfer with
$V=\lambda q^2/2+\gamma^{-1/2}aq^3+\gamma^{-1}bq^4$ and on a two-mode transfer with
the coupling $\gamma^{-1/2}cq_1^2q_2$: the remainders after the $\gamma^{-1}$ term
scale as $\gamma^{-2}$ with exponents $1.99$ and $2.00$ over
$\gamma\in[50,400]$.
\end{proof}

\subsection{The volume dependence of the $\gamma^{-2}$ coefficient}

\begin{proposition}[Finite density of the second connected order]
\label{f16v77:density}
Assume the chain propagator obeys
$|C(x;n)|\le(\kappa_0+\kappa_1\log L)\,e^{-|n|\xi_{\min}}/\max(1,|x|_\infty^2)$
(V63's decay, extended in time by the Mehler factor). Then each connected
graph contributing to $E_4$ has lattice sum bounded by
$C\,V\,(\kappa_0+\kappa_1\log L)^{e}$, $e$ the number of edges, with $C$
independent of $L$; hence $|E_4|\le C'V(\log L)^6$.
\end{proposition}
\begin{proof}
The graphs are: four cubic vertices with six edges and no loops, i.e.\
the tetrahedron $K_4$ and the ring with two double edges; two cubic and
one quartic vertex with five edges; two quartic vertices with four
edges; and the single sixth-order vertex. Fix vertex $1$ at the origin
(translation invariance gives the factor $V$) and sum the others. For a
double edge the factor $(1+|x|)^{-4}$ is summable in three dimensions; for
the ring, sum the two double-edge partners first, leaving
$\sum_x(1+|x|)^{-4}$; for the tetrahedron, sum the fourth vertex against its
three edges, which gives at most $C(1+D)^{-3}$ with $D$ the diameter of the
other three, and the remaining triangle with the extra factor $(1+D)^{-3}$
has superficial degree $6-9<0$ and no divergent subsum; the mixed graphs
have a multiple edge at every vertex and are summable in the same way;
the time sums converge geometrically. Each edge carries one factor
$(\kappa_0+\kappa_1\log L)$.
\end{proof}

The proposition says that the second connected order has a finite
density per unit volume up to $(\log L)^6$, as V67 anticipated for all
cumulants of the vertex field; the growth of the power of the logarithm
with the order is the price of V63's logarithm, and the connected
structure is what keeps every order extensive rather than
super-extensive, in contrast with the $L^2$ norm of V67.

\subsection{What remains for N2 at this order}

The expansion \eqref{f16v77:even} is proved as a formal identity and
its first two coefficients have finite densities. To turn it into the
statement ``$E_0=E_0^{(0)}+\gamma^{-1}E_2+R$ with $|R|\le CV\gamma^{-2}(\log L)^a$,
$C$ independent of $L$'', which is N2 for the hard energy at order
$\gamma^{-1}$, three inputs are missing, all of the kind V63--V66 provide
at fixed $L$: (i) the complement of the sup-norm event, whose
probability V66 bounds by a power of $V$ times a Gaussian tail and
whose contribution must be shown to be $O(V\gamma^{-2})$ rather than
$O(1)$; (ii) the chart remainder beyond sixth order on the event,
$O(\gamma^{-5/2}(\log V)^{7/2}V)$ per V64's sandwich, which is smaller than
$\gamma^{-2}V$ only for $(\log V)^{7/2}\ll\gamma^{1/2}$, i.e.\ in the
small-volume regime again; (iii) the summability of the full cumulant
series, or an asymptotic-expansion argument for the compact model in
place of it. Item (ii) shows the same wall as V63: the frozen Gaussian
with its connected corrections is an expansion whose error control is
volume-uniform only up to the scale where the sup-norm of the hard
field ceases to be small, $\log V\ll\gamma^{1/7}$ here. The programme's
scale-to-scale structure is the way past this wall, since a
proper-time step (V71) never needs the sup-norm of the whole field, only
of its part above the scale; that reformulation of the event is the
next task of the connected programme.

\subsection{Status}

\emph{Proved}: Propositions \ref{f16v77:cumulant} (formal identity and
parity) and \ref{f16v77:density} (power counting under V63's decay).
\emph{Inherited}: V63, V64, V66, V74, V76. \emph{Conditional}: the
asymptotic character of the expansion for the compact model (numerical
on models). \emph{Open}: the three inputs listed; N2 proper.

\subsection{Verification and preservation}

The new diagnostic is
\begin{center}\small\path{scripts/verify_ym_f16_v77_remainder.py}.\end{center}
Its hashed content is the $\gamma^{-2}$ scaling of the remainder on the
one-mode and two-mode transfers and the smallness of the remainder
against the $\gamma^{-1}$ term. It has 4865 bytes and SHA--256
\begin{center}\scriptsize\ttfamily 6E652EDDFBB412D2F0526815D4A0DAFAF7657F6327347A4155B1F3FF137B3559.\end{center}
Its canonical hashed payload has SHA--256
\begin{center}\scriptsize\ttfamily 3BDA5753E98BA2CADBD0DAB1E1FAFFC062BA3DA97847429E4ADB54C797B53007.\end{center}
Only this terminal insertion and the literal title version differ from V76;
removing them recovers its bytes exactly.

\begin{firewall}
V77 proves the cumulant identity for the semiclassical expansion of the
hard Born--Oppenheimer energy with its parity structure, proves the
finite density of the second connected order under V63's decay, and
states the three missing inputs for a volume-uniform remainder. It does
not prove the remainder bound, does not go past the small-volume wall,
and does not bear on the physical continuum trajectory or an
interacting continuum mass gap.
\end{firewall}
% END F16 V77 APPEND

% BEGIN F16 V78 APPEND -- 2026-09-19
% Insert before the final document-ending command of cumulative V77.
% This insertion leaves every predecessor body and the native operator unchanged.

\clearpage
\section{V78: the infrared structure of the connected expansion and the role of transversality}
\label{f16v78:context}

V77 proved that the second connected order of the hard energy has a
finite density by absolute power counting under V63's decay. This
section shows that absolute power counting stops working at the third
connected order and identifies what replaces it. The obstruction is
the necklace: a chain of Wick-ordered pairs, each pair joined by a
double propagator and consecutive pairs joined by a single propagator.
With $|C(x)|\lesssim|x|^{-2}$ in three dimensions the double edge is
summable but the single bridge is not, $\sum_x|C(x)|\asymp L$, so the
absolute bound of an $n$-pair necklace grows like $L^{n-1}$: the
position-space cumulant bounds of V61--V63, which were enough for the
small-volume coordinate theorem and for V77's second order, are not
volume-uniform from the third order on. The same graphs are finite in
momentum space: a pair joined by a double edge is a self-energy
insertion, the necklace is a term of the geometric series of the
dressed propagator, and its value is $V^{-1}\sum_k\Sigma(k)^n/\lambda_k^{n+1}$,
which is bounded uniformly in $L$ provided $\Sigma(k)=O(\lambda_k\log\lambda_k)$
at small $k$, i.e.\ provided the hard self-energy is transverse and
vanishes with the momentum. That is the Ward identity of the horizontal
chain: the soft modes stay massless. V69's kernel is this self-energy at
one loop, and its data confirm the transversality: $\Pi_L(k)/\lambda_k$ is
the slowly varying $\delta_L(k)$ of V69, growing only logarithmically as
$k\to0$. The consequence for the programme is stated: a volume-uniform
control of the connected expansion beyond the second order cannot be
obtained by absolute position-space bounds and must be organized in
momentum space around the Ward identity, as in the infrared analysis of
four-dimensional perturbation theory; the lineage's proper-time density
(V71) is the natural organizing variable, and the transversality of the
one-loop self-energy is the first Ward identity the lineage has
verified.

\subsection{Dependencies}

V61--V63's chaos and decay bounds, V67's Wick-ordered vertex, V69's kernel,
V77's cumulant identity and second-order power counting.

\subsection{The necklace obstruction}

\begin{proposition}[Absolute bounds fail at the third connected order]
\label{f16v78:necklace}
Let $C$ be a propagator on the periodic cube with $|C(x)|\le\kappa/\max(1,|x|_\infty^2)$
and $|C(x)|\ge\kappa'/\max(1,|x|_\infty^2)$ on a positive fraction of
directions. The absolute-value sum of the $n$-pair necklace, a connected
graph with $2n$ cubic vertices in which vertices $2i-1,2i$ are joined by a
double edge and vertices $2i,2i+1$ by a single edge, is bounded above and
below by constants times $V\,L^{n-1}$. In particular the third connected
order of \eqref{f16v77:series} contains graphs whose absolute
position-space bounds grow with the volume.
\end{proposition}
\begin{proof}
Each double edge contributes $\sum_x C(x)^2$, which converges in three
dimensions (the verifier: $36.5,39.6,41.1$ at $L=8,16,32$ for the model
propagator), and each single bridge contributes $\sum_x|C(x)|$, which grows
linearly (fitted exponent $1.06$). The necklace has $n$ double edges
and $n-1$ bridges; translation invariance gives the factor $V$. The
necklace with $n=2$ has four cubic vertices and is not among V77's graphs
because two of its edges are single and adjacent: V77's list omitted it,
and Proposition~\ref{f16v77:density} is correct only for the graphs it
lists (the tetrahedron and the two-double-edge ring, both with every
vertex on a multiple edge); the two-pair necklace has one vertex per pair
on a single edge only, and its absolute bound is $V\cdot L$. Its actual
value is finite, by the next proposition, so the finite density of $E_4$
stands, but not by absolute power counting.
\end{proof}

\subsection{Momentum space and the Ward identity}

\begin{proposition}[Necklaces are self-energy insertions]
\label{f16v78:resummation}
Let $\Sigma(k)$ be the self-energy of the hard propagator, the amputated
Wick-ordered pair. The sum of all necklaces with $n$ pairs on a
propagator line is $V^{-1}\sum_{k\ne0}\Sigma(k)^n/\lambda_k^{n+1}$, and if
$|\Sigma(k)|\le\lambda_k(A+b|\log\lambda_k|)$ for small $k$, every such sum is
bounded uniformly in $L$, with the limit
$\int_{\rm BZ}(A+b\log\lambda)^n\lambda^{-1}\,d^3k/(2\pi)^3$.
\end{proposition}
\begin{proof}
Translation invariance diagonalizes the chain of insertions in momentum
space; $\lambda_k^{-1}$ is the frozen propagator. In three dimensions
$V^{-1}\sum_{k\ne0}\lambda_k^{-1}$ converges to the Watson value $0.2527$
(the verifier: $0.2246,0.2386,0.2457,0.2492$ at $L=8,16,32,64$), and the
powers of the logarithm are integrable at $k=0$; the lattice sums with
logarithms converge slowly, as $(\log L)^n/L$, and are recorded as
diagnostics.
\end{proof}

The bound on $\Sigma(k)$ is the transversality of the one-loop
self-energy: the horizontal chain's gauge symmetry protects the soft
modes from a mass, so $\Sigma(0)=0$ and $\Sigma(k)=\lambda_k\,\delta(k)$ with
$\delta$ at most logarithmic. V69 computed exactly this: $\Pi_L(k)$ is
$\lambda_k(V/2)\delta_L(k)$, and $\delta_L(k)$ at $L=32$ takes the values
$-1.826, -1.695, -1.565, -1.462, -1.377, -1.306$ at $m=1,\dots,6$, with consecutive differences below $0.5$ and
the logarithmic slope of V73. This is the first Ward identity of the
lineage verified on its own kernel, and it is what makes the
second-order density finite despite the two-pair necklace.

\subsection{Consequences}

The programme's remainder control beyond the second connected order
must be organized in momentum (or proper-time) space with the
self-energy resummed, i.e.\ with the running coupling of V73 inserted
into the propagators, rather than by position-space absolute bounds.
This is the standard structure of the infrared analysis of
four-dimensional gauge perturbation theory, and it fixes the form of
the next connected statements: (i) the transversality of the hard
self-energy at every order, as an identity of the horizontal chain; (ii)
the resummed propagator $(\lambda_k(1+\delta(k)/\gamma))^{-1}$ as the object
in which the higher orders are expanded; (iii) the remainder bound in
terms of the resummed propagator's decay, which is still $1/x^2$ up to
logarithms and therefore still needs the momentum-space organization at
every order. V63's decay bound and V61's chaos bounds remain the tools
for the small-volume regime and for the second order; they are not the
tools for the volume-uniform remainder, and this section records why.

\subsection{Status}

\emph{Proved}: Propositions \ref{f16v78:necklace} and
\ref{f16v78:resummation}; the correction to the graph list of
Proposition~\ref{f16v77:density}. \emph{Inherited}: V69's kernel, V77.
\emph{Conditional}: the transversality of the self-energy beyond one
loop (an identity of the chain, not yet written). \emph{Open}: the
momentum-space organization of the remainder; N2.

\subsection{Verification and preservation}

The new diagnostic is
\begin{center}\small\path{scripts/verify_ym_f16_v78_infrared.py}.\end{center}
Its hashed content is the linear growth of the bridge sum, the
boundedness of the double-edge sum, the convergence of the
momentum-space necklace to the Watson value, and the transversality of
V69's kernel. It has 4936 bytes and SHA--256
\begin{center}\scriptsize\ttfamily 170A30F25B3FEA0C5642C7CB39907E7779B2A3E91DF9DA67C96B56A54B1F1714.\end{center}
Its canonical hashed payload has SHA--256
\begin{center}\scriptsize\ttfamily 70AEE7058F3F6332BF87FD563D0D3F36956F63271E6774A5B96B952C505D6339.\end{center}
Only this terminal insertion and the literal title version differ from V77;
removing them recovers its bytes exactly.

\begin{firewall}
V78 shows that absolute position-space bounds on the connected
expansion fail from the third order by the necklace graphs, corrects
the graph list of V77, shows that the same graphs are finite in
momentum space when the hard self-energy is transverse, and verifies
that transversality on V69's one-loop kernel. It does not prove the
transversality beyond one loop, does not bound the remainder, and does
not bear on the physical continuum trajectory or an interacting
continuum mass gap.
\end{firewall}
% END F16 V78 APPEND

% BEGIN F16 V79 APPEND -- 2026-09-19
% Insert before the final document-ending command of cumulative V78.
% This insertion leaves every predecessor body and the native operator unchanged.

\clearpage
\section{V79: the exact Ward identity of the horizontal chain and the two limits of the one-loop kernel}
\label{f16v79:context}

V78 identified transversality of the hard self-energy as the property
that makes the connected expansion finite in the infrared. This section
states and verifies the exact form of that property on the lattice, and
places the one-loop kernel of V69 between its two limits. Three things.
First, the exact Ward identity: for a pure-gauge background
$U_e=g_xg_{x+\mu}^{-1}$ the horizontal spectrum of the chart Hessian equals
that at $U=1$, once the flat directions are transported by the gauge
transformation; the one-loop kernel of a longitudinal background
vanishes identically, and the verifier finds it zero to $2e-10$ at $L=8$
for two amplitudes. This is V68's gauge invariance of the horizontal
spectrum at a stationary background, made explicit, and it is the
lattice form of the statement $\Pi_{\mu\nu}k_\nu=0$. Second, the constant
(toron) limit: a spatially constant background is not the $k\to0$ limit
of the transverse plane wave but a different object, the holonomy, whose
hard zero-point energy is quadratic and negative in the amplitude
below the lowest hard frequency (exponents $2.003$ and $2.036$ at $L=8$)
and joins the toron potential of V51--V58 beyond it. Third, the
transverse kernel interpolates: $\Pi_L(k)=\lambda_k(V/2)\delta_L(k)$ with
$\delta_L$ logarithmic (V69, V73, V78), which is transversality with a
logarithm, and $\Pi_L(k)\to0$ as $k\to0$ through nonzero momenta while the
toron energy at $k=0$ is of order $L$: the kernel is discontinuous at
$k=0$ at fixed $L$, by an amount that vanishes per unit volume. The
three statements together are the lineage's Ward-identity package at
one loop, which V78 requires at every order.

\subsection{Dependencies}

V68's horizontal reduction and gauge invariance, V69's chart and
exclusions, V51--V58's toron results, V78.

\subsection{The exact Ward identity}

\begin{proposition}[Longitudinal backgrounds have no one-loop kernel]
\label{f16v79:ward}
Let $g_x\in SU(2)$ be any gauge transformation and $U_e=g_xg_{x+\mu}^{-1}$
the corresponding pure-gauge configuration. In the right chart at $U$,
the Hessian of the Wilson action projected onto the orthogonal
complement of the vertical frame at $U$ and of the transported flat
directions $y_e=\operatorname{Ad}(g_{x+\mu})c$, $c$ constant, has the same spectrum
as the corresponding projected Hessian at $U=1$. Hence $Z_L(U)=Z_L(1)$
and the one-loop kernel of a longitudinal background is zero.
\end{proposition}
\begin{proof}
The map $y\mapsto\operatorname{Ad}(g_{x+\mu})y$ on chart coordinates conjugates the
chart at $1$ to the chart at $U$: $g_xe^{y_e}g_{x+\mu}^{-1}=U_ee^{\operatorname{Ad}(g_{x+\mu})y_e}$,
it is an isometry of the flat chart metric, it carries the action to
itself by gauge invariance, the vertical frame at $1$ to the vertical
frame at $U$, and the constant directions to the transported ones.
The verifier realizes $g_x=\exp(\varepsilon\cos(kx_1)e_{\rm col})$ at $L=8$,
$k=2\pi/8$, and finds relative deviations $2e-10$ of $Z_L$ at
$\varepsilon=0.05,0.1$, with the same mode count; with untransported
constants the deviation is $10^{-6}$ and quadratic in $\varepsilon$, which
is the size of the error a wrong soft sector makes.
\end{proof}

\subsection{The toron limit}

For the constant background $b=\varepsilon e_{\rm pol}\otimes e_{\rm col}$ the
excluded sector is the nine constant modes only; the hard zero-point
energy shift at $L=8$ is
\begin{center}\small
\begin{tabular}{@{}lrrrr@{}}
\toprule
$\varepsilon$ & 0.02 & 0.04 & 0.08 & 0.16\\
\midrule
$E_0(\varepsilon)-E_0(0)$ & $-0.0036$ & $-0.0146$ & $-0.0587$ & $-0.2410$\\
\bottomrule
\end{tabular}
\end{center}
quadratic with exponents $2.003$ (first half) and $2.036$ (second half),
and negative. All amplitudes are below $\lambda_{\min}^{1/2}=0.77$, so
this is the quadratic regime of the holonomy potential; the toron
potential of V51--V58, with its cusp, is its continuation to
$L\varepsilon\gtrsim1$, where the constant background's own modes (the
torons) are no longer separable from the lowest hard modes. The constant
background is not obtained from the transverse plane wave by $k\to0$: the
plane wave's kernel vanishes with $\lambda_k$, the constant background's
does not, and the difference is the soft sector, which for the plane
wave is the pair (cos, sin) and the colour orbit and for the constant
background is the toron.

\subsection{The three statements as a package}

Longitudinal: zero, exactly. Constant: a holonomy potential, of order
$L\varepsilon^2$, not part of the local kernel. Transverse: $\lambda_k$ times a
logarithm, which is the local, universal, running part. V78's
requirement, $\Sigma(k)=O(\lambda_k\log\lambda_k)$ for $k\ne0$, is met at one
loop by the first and third statements together; the second shows that
the requirement is only about $k\ne0$ and that $k=0$ belongs to the soft
sector. Beyond one loop the first statement continues to hold exactly
(it is a property of the horizontal chain at every order, by the same
isometry), and the third becomes the transversality of the resummed
self-energy, which is where the connected programme must next prove
something.

\subsection{Status}

\emph{Proved}: Proposition~\ref{f16v79:ward}. \emph{Inherited}: V68,
V69, V51--V58. \emph{Conditional}: the logarithmic bound on the transverse
kernel (numerical, V69/V73). \emph{Open}: transversality of the
self-energy beyond one loop as a proved identity; the momentum-space
remainder.

\subsection{Verification and preservation}

The new diagnostic is
\begin{center}\small\path{scripts/verify_ym_f16_v79_ward.py}.\end{center}
Its hashed content is the vanishing of the pure-gauge kernel and the
quadratic law of the constant background. It has 6641 bytes and
SHA--256
\begin{center}\scriptsize\ttfamily 07EECBE046116997B4DDA9F7F3A2BDA6EC6810E11EE8978CC9C0E0767F2B346E.\end{center}
Its canonical hashed payload has SHA--256
\begin{center}\scriptsize\ttfamily 41044EED5AACFE4BBC974ED5E52E598A5FC155E2A61FF0CADD264A6AD8E0F461.\end{center}
Only this terminal insertion and the literal title version differ from V78;
removing them recovers its bytes exactly. V80 is the ten-version review
and companion checkpoint.

\begin{firewall}
V79 proves the exact lattice Ward identity of the horizontal chain
(pure-gauge backgrounds have no one-loop kernel), computes the
quadratic holonomy regime of the constant background, and places the
transverse kernel between the two. It does not prove transversality
beyond one loop, does not bound the remainder, and does not bear on the
physical continuum trajectory or an interacting continuum mass gap.
\end{firewall}
% END F16 V79 APPEND

% BEGIN F16 V80 APPEND -- 2026-09-19
% Insert before the final document-ending command of cumulative V79.
% This insertion leaves every predecessor body and the native operator unchanged.

\clearpage
\section{V80: ten-version review V70--V79, companion checkpoint, and the consolidated one-loop ledger}
\label{f16v80:context}

This is the scheduled ten-version review. The ten versions V70--V79
form the lineage's perturbative renormalization-group cycle: they took
the static one-loop kernel of V69 and asked what it is, what it is not,
and what its remainder needs. The answers, in order of importance.
(1) The one-loop running of the transfer coefficient is established as
the lineage's own result: the gauge-invariant proper-time density of
the kernel has the window value $0.367$ against $11/(3\pi^2)=0.3715$ (V71,
V73), and its integrated flow reproduces the concurrent trajectory
exponent $3\pi^2/11$ as $2.73$. (2) The finite part of the kernel is not
a physical number: two-dimensional lattice gauge theory, solved exactly,
shows the constrained one-loop determinant to depend on the slice
through the non-stationary background by a factor $31$ (V76), so
$c$, $\delta_1(k)$ at fixed $k$ and $\delta_2$ are link-slice quantities;
the same test shows the orbit-volume (Faddeev--Popov) factor to be a
soft-measure term and not a one-loop energy. (3) The exact Ward identity
of the horizontal chain holds on the lattice with the flat directions
transported (V79, deviation $2e-10$), and the one-loop self-energy is
transverse with a logarithm; from the third connected order on, the
finiteness of the expansion rests on this transversality and cannot be
obtained from position-space absolute bounds, because of the necklace
graphs (V78). (4) The first connected correction is computed and
validated (V74, completed in V76 by the kinetic quartic term):
$\delta_2=$ $-0.150$ at $L=8$ against $\delta_1=$ $-1.262$, so the two-loop term
is $0.12$ of the one-loop term divided by $\gamma$; the expansion's
structure is proved (V77) and its second order has a finite density
under V63's decay, while the volume-uniform remainder needs three
inputs that V63--V66 supply only at fixed $L$. (5) The electric
coefficient is fixed by symmetry in the covariant scheme (V72) and open
in the transfer scheme, four routes to it having been closed. The
control scale of the lineage's rigorous statements remains a few
lattice spacings (V73), and the mass gap is not approached by any of
this; what has changed is that the lineage now knows exactly which
quantities are invariant, which are slice-bound, and which identities
the connected programme must prove next.

\subsection{Companion checkpoint}

The newest concurrent file is
\path{Renewal\_Yang\_Mills\_Mass\_Gap\_Research\_Notes\_V62.tex}, 1951050 bytes, SHA--256
\begin{center}\scriptsize\ttfamily 09A8B55FDDE1A8E4D47936CAFECF25A478BCADC5700AB961A7331D892F207832,\end{center}
unchanged since V75 (V59--V62: ``V59: buffered Wilson ribbons, transported innovations, and exterior predictability; V60: the microscopic-ribbon obstruction and a source-preserving exterior reduction; V61: exact retained-ribbon curvature and transverse Ward cancellation; V62: joint Wilson ladders and paid transport holonomy''); its V58 body is preserved
byte-wise. Nothing in it uses V68--V79, and nothing in V70--V79 uses it
beyond V55's calibration and the V57 audit accepted in V69. The
interface table of V75 stands with V76's amendment: only $b$, the sign
and monotonicity of the flow, the Ward-identity package, and the
statement that the finite parts require an observable-based definition
are exportable.

\subsection{Consolidated ledger}

\begin{center}\small
\begin{tabular}{@{}lrl@{}}
\toprule
quantity & value & status\\
\midrule
$\delta_1(2\pi/8)$, $L=8$, transfer slice & $-1.262$ & proved, slice-bound\\
$\delta_1(2\pi/64)$, $L=64$, $\varepsilon=0.0125$ & $-2.072$ & proved, slice-bound\\
$b$ (window, $L=64$) & $0.367$ & numerical, invariant, vs $0.3715$\\
$c$ (matching at $\Lambda=1$) & $-1.209$ & numerical, slice-bound\\
$1/b$ & $2.73$ & numerical, vs $2.692$\\
$\delta_2(2\pi/8)$, $L=8$, corrected & $-0.150$ & proved (cumulant), slice-bound\\
$E_2/V$ at $\varepsilon=0$, $L=8$ & $-0.561$ & proved (cumulant)\\
2D exact vs link slice & $-1.22$ vs $-37.8$ & proved\\
pure-gauge kernel & $2e-10$ & proved (exact identity)\\
toron exponents & $2.003$, $2.036$ & numerical\\
\bottomrule
\end{tabular}
\end{center}

\subsection{Review against the order of attack}

Step 1 (N2): open; reformulated as the volume-uniform remainder of the
connected expansion, with its three missing inputs named (V77) and the
momentum-space organization identified as necessary (V78). Step 2
(uniform block statement): unchanged. Step 3 (one elimination step):
proper time is the gauge-invariant scale variable (V71); the one-loop
content is explicit; no elimination with remainder. Step 4 (locality
and remainder): the necklace analysis shows what a remainder bound
must control. Step 5 (transformed coupling): the flow is explicit (V73)
and its finite part is slice-bound (V76); an observable-based
definition is required. Step 6 (iteration): trivial at one loop. Step 7
(first failing scale): two scales (V73), rigorous control at $L\approx5$.
Step 8 (export): the table above. Obstructions met and recorded: the
coordinate class five times (momentum shells, temporal-gauge Gaussian,
slice dependence at $5$ percent, slice dependence at a factor $31$,
untransported flat directions), the numerical-control class twice
(zero-point factor, $L=64$ threshold), the probability-tail class once
(the small-volume wall of the sup-norm event, V77), the
operator-norm class once (necklaces, V78). Errata: V69's factor and
$L=64$ row (V71), V74's kinetic term (V76), V77's graph list (V78).

\subsection{What the next ten versions should do}

The lineage's invariant one-loop content is complete; the finite parts
are not physical and should not be pursued further as such. The
connected programme has three concrete targets: the transversality of
the resummed self-energy as an identity of the horizontal chain at
every order (V79's isometry argument extended to the vertices); the
momentum-space remainder bound for the cumulant expansion at fixed
order with volume-uniform constants, using the resummed propagator; and
an observable-based coupling (the plaquette expectation or the
field-strength two-point function at one loop) whose finite part is
physical and can be exported. The small-volume wall of V77 is to be
attacked by the scale-by-scale reformulation of the event. None of
these is the mass gap; all of them are what a perturbative RG programme
must do before it can hand controlled data to the local-vacuum
programme at the first nonperturbative scale.

\subsection{Status}

\emph{Proved}: the integrity and preservation checks. \emph{Inherited}:
V70--V79 as stated. \emph{Conditional}: the interface numbers as marked.
\emph{Open}: as listed.

\subsection{Verification and preservation}

The new diagnostic is
\begin{center}\small\path{scripts/verify_ym_f16_v80_review.py}.\end{center}
Its hashed content is the integrity of the V70--V79 artifacts, the
preservation of the concurrent V58 body in the newest concurrent file,
and the consistency of the review numbers. It has 5924 bytes and
SHA--256
\begin{center}\scriptsize\ttfamily C8D97874E9E22D794B75BDA692911BDEED0F1126AB47AA06F8C0179C0447A5D2.\end{center}
Its canonical hashed payload has SHA--256
\begin{center}\scriptsize\ttfamily 7F16A38B82C8D49FA340E58AE0B76B1115A8F2C32BCA51A33497C8E0CDAEE446.\end{center}
Only this terminal insertion and the literal title version differ from V79;
removing them recovers its bytes exactly. The next checkpoint is V85,
the next review V90.

\begin{firewall}
V80 reviews V70--V79, records the companion checkpoint at the concurrent
V62, consolidates the one-loop ledger with the invariant and
slice-bound quantities marked, and verifies the artifacts. It proves
nothing new and does not bear on the physical continuum trajectory or
an interacting continuum mass gap.
\end{firewall}
% END F16 V80 APPEND

% BEGIN F16 V81 APPEND -- 2026-09-19
% Insert before the final document-ending command of cumulative V80.
% This insertion leaves every predecessor body and the native operator unchanged.

\clearpage
\section{V81: the all-order longitudinal Ward identity and the two-loop kernel at three momenta}
\label{f16v81:context}

V79 proved the exact Ward identity at one loop and V78 showed that the
connected expansion's finiteness beyond the second order rests on the
transversality of the resummed self-energy and on the absence of a mass
term in it. This section does two things toward that requirement.
First, it extends V79's isometry argument to the full hard-mode problem,
with all vertices, the Haar measure and the native cosine kinetic
factor: for a pure-gauge background the whole Born--Oppenheimer problem
is unitarily equivalent to the one at $U=1$, so the longitudinal kernel
vanishes at every order. Second, it computes the two-loop kernel of
V74--V76 at the three lowest momenta of the $L=8$ lattice and tests it
for a mass term: a $k$-independent part of $\Pi_2$ would appear in
$\delta_2(k)=2\Pi_2/(\lambda_kV)$ as $2M/(\lambda_kV)$, growing fourfold from
$m=2$ to $m=1$; the values found are $-0.1503, -0.1282, -0.1161$, whose ratio between
$m=1$ and $m=2$ is $1.17$ against $\lambda_2/\lambda_1=3.41$ for a pure
mass term, and the exact three-point decomposition into a mass term, a
constant and a logarithm gives $M=$ $1.967$ against $\Pi_2(m=1)=$ $-22.539$,
i.e.\ 8.7 percent of it, within the finite-difference accuracy of the
vertices that the decomposition absorbs into $M$. No mass term dominates
at order $\gamma^{-1}$ in the transfer slice, and the two-loop kernel is,
like the one-loop one, $\lambda_k$ times a slowly varying function. This is the
first two-loop transversality statement of the lineage; it is numerical
and slice-bound in its finite part, and exact in its longitudinal half.

\subsection{Dependencies}

V79's isometry, V74's cumulant pieces, V76's kinetic term, V78's
requirement.

\subsection{The all-order longitudinal identity}

\begin{proposition}[Pure-gauge backgrounds are invisible to every order]
\label{f16v81:longitudinal}
Let $U_e=g_xg_{x+\mu}^{-1}$. The hard-mode Born--Oppenheimer problem at $U$
(the transfer restricted to the horizontal chart at $U$ with the
transported flat directions and the background's soft sector excluded,
with the full Wilson action, the Haar density, and the cosine kinetic
factor) is unitarily equivalent to the problem at $U=1$; in particular
$E_0(U;\gamma)=E_0(1;\gamma)$ exactly, and every coefficient of the expansion
\eqref{f16v77:even} vanishes for a longitudinal background.
\end{proposition}
\begin{proof}
The map $\Phi:y\mapsto(\operatorname{Ad}(g_{x+\mu})y_e)_e$ on the chart satisfies
$g_x e^{y_e}g_{x+\mu}^{-1}=U_ee^{\Phi(y)_e}$, so it intertwines the
configuration $e^{y}$ at $1$ with the configuration $Ue^{\Phi(y)}$ at $U$,
which are gauge transforms of each other by $g$. The Wilson action, the
product Haar measure, the temporal cosine factor between two slices
(gauge transformed by the same $g$ on both), the vertical frames, and
the flat directions are all carried into each other by $\Phi$ and by
the gauge transformation, and $\Phi$ is an isometry of the flat chart
metric. Hence the transfer's kernel in the chart at $U$ is the
composition of the kernel at $1$ with $\Phi$, the soft sectors
correspond, and the Perron eigenvalues coincide. No expansion is
involved, so the statement holds to every order and beyond the
expansion. V79's verifier is the one-loop instance.
\end{proof}

\subsection{The two-loop kernel at three momenta}

\begin{center}\small
\begin{tabular}{@{}lrrrrrrr@{}}
\toprule
$m$ & $\lambda_k$ & $\Pi_2$ & $\delta_2$ & quartic & Haar & sunset & kinetic\\
\midrule
1 & 0.5858 & -22.539 & -0.1503 & -14.70 & 29.29 & 55.77 & -9.25\\
2 & 2.0000 & -65.643 & -0.1282 & -37.42 & 61.89 & 135.54 & -22.34\\
3 & 3.4142 & -101.511 & -0.1161 & -58.40 & 74.21 & 175.63 & -29.51\\
\bottomrule
\end{tabular}
\end{center}
The exact three-point decomposition $\delta_2=2M/(\lambda_kV)+A'+b'\log\lambda_k$
gives $M=$ $1.967$, $A'=$ $-0.1498$, $b'=$ $0.0255$; with three points and three
parameters it is a decomposition, not a fit, and its $M$ carries the
finite-difference error of the vertices, of the order of one percent of
$E_2$. The mass term is 8.7 percent of $\Pi_2(m=1)$; the slope $b'$ is the two-loop
contribution to the running, of the size expected for the
$\gamma^{-2}$ coefficient of the $\beta$-function times the slice-bound
matching, and is not interpreted further. The pieces show that the
kinetic quartic term of V76 is as large as the action's quartic term at
this order, which is why V74 alone was incomplete.

\subsection{Consequences}

V78's requirement, $\Sigma(k)=O(\lambda_k\log\lambda_k)$, now holds at one loop
exactly in its longitudinal part and numerically in its transverse part,
and at two loops numerically in both, at $L=8$, in the sense that no mass
term dominates. The all-order
longitudinal statement is the first exact all-order property of the
connected expansion; the transverse one remains to be proved order by
order, and its natural proof is the same isometry applied to the
$\varepsilon$-derivatives of the background, i.e.\ to the gauge variation of
the plane wave, which is the next connected target.

\subsection{Status}

\emph{Proved}: Proposition~\ref{f16v81:longitudinal}. \emph{Inherited}:
V74, V76, V79. \emph{Conditional}: the two-loop transversality
(numerical, $L=8$, three momenta). \emph{Open}: the transverse identity
at every order; the momentum-space remainder.

\subsection{Verification and preservation}

The new diagnostic is
\begin{center}\small\path{scripts/verify_ym_f16_v81_two_loop_momenta.py}.\end{center}
Its hashed content is the common sign and small spread of $\delta_2$ at
the three momenta and the smallness of the fitted mass term. It has
3750 bytes and SHA--256
\begin{center}\scriptsize\ttfamily 1611F8BDC619FF0B18205E434CAAA765DFACE95CE9B2191339B2967D8E0473E5.\end{center}
Its canonical hashed payload has SHA--256
\begin{center}\scriptsize\ttfamily 152DDF60050AA9D047688229DB7A94FE90C49A4FA4B64B3C3103829FFC621E41.\end{center}
Only this terminal insertion and the literal title version differ from V80;
removing them recovers its bytes exactly.

\begin{firewall}
V81 proves that pure-gauge backgrounds have no Born--Oppenheimer kernel
at any order, and finds numerically that the two-loop kernel of the
transfer slice has no mass term at $L=8$. It does not prove the
transverse identity beyond one loop, does not bound a remainder, and
does not bear on the physical continuum trajectory or an interacting
continuum mass gap.
\end{firewall}
% END F16 V81 APPEND

% BEGIN F16 V82 APPEND -- 2026-09-19
% Insert before the final document-ending command of cumulative V81.
% This insertion leaves every predecessor body and the native operator unchanged.

\clearpage
\section{V82: the all-order transverse Ward identity}
\label{f16v82:context}

V81 proved that pure-gauge backgrounds are invisible to the hard
Born--Oppenheimer problem at every order. This section completes the
tensor form of the Ward identity: the mixed second derivative of the
hard energy along a transverse plane wave and a longitudinal
(pure-gauge) plane wave vanishes at every order, so the quadratic form
of the energy at $U=1$ on plane waves is transverse. The proof is the
isometry of V81 applied along the gauge orbit of the transverse
background, together with the absence of a linear term at $U=1$, which
follows from the inversion symmetry $U\to U^{-1}$ of the Wilson action
and measure. The verifier realizes the identity at one loop on the
two-parameter family $g\,e^{\varepsilon T}g^{-1}$ at $L=8$, with the soft
sector transported, and finds the mixed derivative $3e-08$ relative to
the kernel $\Pi(k)$, and the exact $\eta$-independence of the energy at
fixed $\varepsilon$ to $3e-10$. With V78 and V79 this closes the Ward-identity
package: the hard self-energy is transverse as a tensor at every order;
what remains open, and is dynamical rather than symmetry-based, is the
absence of a mass term in its transverse part beyond the numerical
evidence of V69 and V81.

\subsection{Dependencies}

V81's isometry, V79's transported flat directions, V69's exclusions.

\subsection{The identity}

\begin{proposition}[Transverse--longitudinal mixed kernel vanishes]
\label{f16v82:mixed}
Let $T$ be a transverse plane wave, $g^{(\eta)}_x=\exp(\eta\,\theta_x)$ a
one-parameter gauge transformation, and
$U(\varepsilon,\eta)_e=g^{(\eta)}_xe^{\varepsilon T_e}g^{(\eta)-1}_{x+\mu}$, with the soft
sector at $(\varepsilon,\eta)$ the transport by $\operatorname{Ad}(g^{(\eta)}_{x+\mu})$ of the
soft sector at $(\varepsilon,0)$. Then $E_0(\varepsilon,\eta)=E_0(\varepsilon,0)$ for all
$\eta$, exactly and at every order, and in particular
$\partial_\varepsilon\partial_\eta E_0(0,0)=0$. Consequently, for the family
$e^{\varepsilon T+\eta\nabla\theta}$ of superposed plane waves, the mixed second
derivative also vanishes, since the two families differ at order
$\varepsilon\eta$ in the background and $E_0$ has no linear term at $U=1$.
\end{proposition}
\begin{proof}
The first statement is Proposition~\ref{f16v81:longitudinal} applied
to the background $e^{\varepsilon T}$ in place of $1$: the isometry
$\Phi_\eta$ conjugates the Born--Oppenheimer problem at $U(\varepsilon,\eta)$ to
the one at $U(\varepsilon,0)$, soft sectors included, so the energies agree.
For the second, write $g\,e^{\varepsilon T}g^{-1}=e^{\varepsilon T+\eta\nabla\theta+\varepsilon\eta R+O(\varepsilon^2\eta,\varepsilon\eta^2)}$
with $R$ a fixed field; the mixed derivative along the two families
differs by $\partial_yE_0(0)[R]$, the first derivative of the energy at
$U=1$ in the direction $R$, which vanishes because $E_0$ is even under
$y\to-y$: the inversion $U_e\to U_e^{-1}$ of all links preserves the
Wilson action ($\operatorname{Tr}U_p^{-1}=\operatorname{Tr}U_p$ for $SU(2)$), the Haar measure,
the kinetic factor, and the soft sector, and acts on the chart as
$y\to-y$. The verifier evaluates the one-loop instance at $L=8$,
$k=2\pi/8$, $\varepsilon=\eta=0.05$, with the gauge transformation in the same
colour as the wave and in an orthogonal colour: the mixed differences
are $0e+00$ and $3e-08$ relative to $\Pi(k)$, and $E_0(\varepsilon,\pm\eta)$ equals
$E_0(\varepsilon,0)$ to $3e-10$ with equal mode counts.
\end{proof}

\subsection{The Ward package}

Longitudinal--longitudinal (V79, V81): zero at every order.
Transverse--longitudinal (this section): zero at every order.
Transverse--transverse: $\lambda_k$ times a slowly varying function at one
loop (V69) and at two loops (V81), numerically; its vanishing at $k\to0$
through nonzero momenta is the dynamical statement that the connected
expansion generates no mass at fixed order, which no symmetry gives and
which the verifier data support at $L=8,32,64$ (V69, V71) and at two loops
at $L=8$ (V81). V78's requirement is therefore reduced to this one
dynamical property, and the momentum-space organization of the
remainder can take the tensor transversality as given.

\subsection{Status}

\emph{Proved}: Proposition~\ref{f16v82:mixed}. \emph{Inherited}: V81,
V79. \emph{Conditional}: the absence of a transverse mass term beyond
the orders computed. \emph{Open}: the momentum-space remainder; N2.

\subsection{Verification and preservation}

The new diagnostic is
\begin{center}\small\path{scripts/verify_ym_f16_v82_transverse_ward.py}.\end{center}
Its hashed content is the vanishing of the mixed derivative and the
$\eta$-independence at fixed $\varepsilon$, for two colour choices. It has
5904 bytes and SHA--256
\begin{center}\scriptsize\ttfamily 35CA0BAA5851F0F8621379C5222484AD4223F4141C5FF95AEA6A199E189962F0.\end{center}
Its canonical hashed payload has SHA--256
\begin{center}\scriptsize\ttfamily 09BA93B5D346D0D15F3672EE0E6FBE7EBA1FC2F4C33FE05587DFB397502A0BA2.\end{center}
Only this terminal insertion and the literal title version differ from V81;
removing them recovers its bytes exactly.

\begin{firewall}
V82 proves the all-order transverse--longitudinal Ward identity of the
hard Born--Oppenheimer energy and verifies its one-loop instance. It
does not prove the absence of a transverse mass term, does not bound
the remainder, and does not bear on the physical continuum trajectory
or an interacting continuum mass gap.
\end{firewall}
% END F16 V82 APPEND

% BEGIN F16 V83 APPEND -- 2026-09-19
% Insert before the final document-ending command of cumulative V82.
% This insertion leaves every predecessor body and the native operator unchanged.

\clearpage
\section{V83: the one-loop plaquette as an external check of the frozen chain, and an erratum to V70}
\label{f16v83:context}

The lineage's frozen Mehler chain has so far been checked against
itself (brute force, block reductions, one-mode models) and against the
one-loop $\beta$-function coefficient through a fitted slope. This
section checks it against an exactly known number: the one-loop
plaquette of the four-dimensional Wilson theory,
$\langle1-\tfrac12\operatorname{Tr}U_p\rangle=3/(4\gamma)+O(\gamma^{-2})$, which is
equipartition of the quadratic action over the nine physical modes per
site and six plaquettes, exact for every lattice size up to $O(1/V)$.
In the chain the spatial plaquette at one loop is $\tfrac32c_0/\gamma$ with
$c_0=\langle F_{12}^2\rangle\gamma$ per colour, so the check is $c_0\to\tfrac12$.
V70's verifier gave $c_0=0.3884$ at every $L$, and the discrepancy is an
error in V70: the plaquette amplitude $a(k)$ on a transverse mode $v$ is
the bilinear pairing $a\cdot v$, and the sum of $|a\cdot v|^2$ over the
transverse modes is $a^{\mathsf T}P\bar a=|a|^2-|a^{\mathsf T}w|^2/|w|^2=|a|^2$, since
$a^{\mathsf T}w=w_\mu w_\nu-w_\nu w_\mu=0$; V70 projected $a$ Hermitian-wise,
$a_T=a-w(w^\dagger a)/|w|^2$, and used $a_T^\dagger a_T=a^\dagger Pa=|a|^2-|w^\dagger a|^2/|w|^2$,
which differs on the lattice because $w^\dagger a=\bar w_\nu w_\mu-\bar w_\mu w_\nu\ne0$.
With the correct contraction the verifier finds $c_0=$ $0.499859, 0.499991, 0.499999, 0.500000, 0.500000$ at
$L=8,16,32,64,96$, converging to $\tfrac12$ like $1/V$ (the excluded zero
mode), and reproduces V70's $0.3884$ with the old contraction. The chain's
equal-time covariance, its Mehler variances $1/(2\omega_k)$, its
transverse projector and its plaquette linearization are thereby
validated against lattice perturbation theory at one loop. V70's
qualitative statements (Gram $\propto\gamma^{-2}$, $\gamma$-free relative
decrements, complete monotonicity) stand; its numbers are replaced
below, and change in the fourth digit only, since the error is a
$k$-dependent factor common to all $n$.

\subsection{Dependencies}

V45's Mehler chain, V70's proposition and verifier, the one-loop
plaquette of the Wilson theory (equipartition).

\subsection{The check}

\begin{proposition}[One-loop plaquette from the chain]
\label{f16v83:plaquette}
In the frozen chain the single-plaquette equal-time coefficient is
\begin{equation}
 c_0(L)=\frac1V\sum_{k\ne0}\frac{|w_\mu(k)|^2+|w_\nu(k)|^2}{2\omega_k},\qquad
 \omega_k=\sinh\xi_k,\ \cosh\xi_k=1+\hat k^2/2,
 \label{f16v83:c0}
\end{equation}
and $c_0(L)\to\tfrac12$ as $L\to\infty$, so that the spatial plaquette of
the chain at one loop is $\tfrac32c_0/\gamma\to3/(4\gamma)$, the
four-dimensional value.
\end{proposition}
\begin{proof}
The plaquette linearization $F=w_\mu u_\nu-w_\nu u_\mu$ gives
$\langle|\hat F(k)|^2\rangle=\sum_va^{\mathsf T}\langle vv^\dagger\rangle\bar a$ over the
transverse modes with $\langle vv^\dagger\rangle=(I-ww^\dagger/|w|^2)/(2\omega_k)$, and
$a^{\mathsf T}w=0$ removes the projector, leaving $|a|^2=|w_\mu|^2+|w_\nu|^2$.
The equal-time Mehler variance $1/(2\omega_k)$ is the Matsubara sum
$\sum_{k_4}(\hat k^2+\hat k_4^2)^{-1}$ (V70), so
$c_0=\int(\hat k_\mu^2+\hat k_\nu^2)/\hat k_4^2$ over the four-dimensional
Brillouin zone up to the excluded zero mode, and by cubic symmetry each
$\hat k_\rho^2/\hat k_4^2$ averages to $\tfrac14$. The verifier finds
$\tfrac12-c_0(L)=$ $1.4e-04, 8.6e-06, 5.3e-07, 3.3e-08, 6.6e-09$ at $L=8,\dots,96$.
\end{proof}

\subsection{Corrected numbers of V70}

\begin{center}\small
\begin{tabular}{@{}lrrrrrr@{}}
\toprule
$n$ & 1 & 2 & 3 & 4 & 6 & 8\\
\midrule
$c_n/c_0$ & 0.1377 & 0.0223 & 0.0045 & 0.0012 & 0.0002 & 0.0001\\
$\theta_n$ & 0.9810 & 0.9995 & 1.0000 & 1.0000 & 1.0000 & 1.0000\\
$\vartheta_n$ (3) & 0.9807 & 0.9995 & 1.0000 & 1.0000 & 1.0000 & 1.0000\\
$\vartheta_n$ (6) & 0.9795 & 0.9994 & 1.0000 & 1.0000 & 1.0000 & 1.0000\\
\bottomrule
\end{tabular}
\end{center}
The rows are $c_n/c_0$ and $\theta_n$ for one writer at $L=64$ and the
capture margins $\vartheta_n$ of the three-orientation and six-writer
families at $L=32$, replacing V70's table. The leading-order Gram per
writer is $\tfrac32c_0^2\gamma^{-2}\to\tfrac38\gamma^{-2}$, and the interface
statement of V70 and V75 is unchanged in substance.

\subsection{What the check establishes}

The one-loop plaquette is the first quantity of the lineage compared
with an exactly known result of lattice perturbation theory, and it
agrees to $10^{-8}$ at $L=96$. It validates the frozen chain's
covariance, which enters V74's and V76's two-loop pieces, and it makes
the two-loop plaquette, $\langle1-\tfrac12\operatorname{Tr}U_p\rangle$ at order
$\gamma^{-2}$, the natural next external check: its coefficient is known
in the literature for $SU(2)$, and the chain can produce it from the
same tensors as V74 by the connected correlations of the plaquette
operator with the vertices. That computation, and the observable-based
coupling it defines, is the next target; it is the gauge-invariant
finite quantity that V76 asked for.

\subsection{Status}

\emph{Proved}: Proposition~\ref{f16v83:plaquette}; the erratum.
\emph{Inherited}: V45, V70. \emph{Conditional}: none. \emph{Open}: the
two-loop plaquette from the chain.

\subsection{Verification and preservation}

The new diagnostic is
\begin{center}\small\path{scripts/verify_ym_f16_v83_plaquette_check.py}.\end{center}
Its hashed content is the convergence $c_0\to\tfrac12$, the reproduction
of V70's erroneous value by the old contraction, and the range of the
corrected $\theta_1$. It has 5235 bytes and SHA--256
\begin{center}\scriptsize\ttfamily 7FAA1CC4D343F4747AF53097B5BD941009348C216BBA5A4E8B262AF1278981B6.\end{center}
Its canonical hashed payload has SHA--256
\begin{center}\scriptsize\ttfamily 4707E4DAF85AEE3834D262E0F7CF546720D2C852C8B0CF1DB87612369B55BE0C.\end{center}
Only this terminal insertion and the literal title version differ from V82;
removing them recovers its bytes exactly.

\begin{firewall}
V83 validates the frozen chain against the exact one-loop lattice
plaquette, corrects the contraction error of V70's writer correlations
and replaces its numbers. It proves nothing about the mass gap and does
not bear on the physical continuum trajectory or an interacting
continuum mass gap.
\end{firewall}
% END F16 V83 APPEND

% BEGIN F16 V84 APPEND -- 2026-09-19
% Insert before the final document-ending command of cumulative V83.
% This insertion leaves every predecessor body and the native operator unchanged.

\clearpage
\section{V84: the two-loop plaquette from the chain and the observable-based coupling}
\label{f16v84:context}

V83 validated the frozen chain at one loop against the exact plaquette.
This section carries the check to the next order and thereby obtains
the first gauge-invariant, slice-independent finite quantity of the
lineage beyond one loop: the two-loop coefficient of the plaquette
expectation, from which the observable-based coupling that V76 asked
for is defined. The route is thermodynamic rather than diagrammatic:
since $\operatorname{Tr}(PT_\gamma^N)$ is the four-dimensional partition function
exactly, the plaquette is the $\gamma$-derivative of the vacuum energy,
$6VP=\partial_\gamma E_0$, and the vacuum energy is the chain's
$E_2/\gamma$ of V74 and V76 plus a normalization whose $\gamma^{-1}$ terms
cancel exactly. Hence
\begin{equation}
 \bigl\langle1-\tfrac12\operatorname{Tr}U_p\bigr\rangle=\frac3{4\gamma}+\frac{p_2}{\gamma^2}+O(\gamma^{-3}),\qquad
 p_2(L)=-\frac{E_2(L)}{6V},
 \label{f16v84:p2}
\end{equation}
with $E_2$ at zero background, the sum of the action's quartic term, the
cubic sunset, the Haar term and the kinetic quartic term. The values
are $0.0873, 0.0922, 0.0934$ at $L=4,6,8$, extrapolating in $1/L^2$ to $0.0951$. Two
controls accompany the number. In two dimensions the same identity
gives $p_2=-3/8$ exactly from the Haar and kinetic terms alone, and the
verifier confirms the single-plaquette integral's expansion to
$2\cdot10^{-6}$. In four dimensions the coefficient is known in the
literature for $SU(2)$; the lineage does not import it, but records that
a crude $N$-scaling of the $SU(3)$ coefficient $1.2208/\beta^2$ gives
$0.090$, of the size found. The comparison with the exact literature
value is the first item of the next checkpoint.

\subsection{Dependencies}

V74's cumulant pieces, V76's kinetic term, V64's kernel normalization
$N_\gamma$, V83's one-loop check. The two-dimensional expansion of the
Wilson single-plaquette integral is classical.

\subsection{The thermodynamic identity}

\begin{proposition}[Plaquette from the vacuum energy]
\label{f16v84:identity}
Let $T_\gamma=M_bC_\gamma M_b$ with the unnormalized cosine kernel, so that
$\operatorname{Tr}(PT_\gamma^N)=Z_{4D}(\gamma)$ on the $N\times L^3$ lattice. Then
$6V\langle1-\tfrac12\operatorname{Tr}U_p\rangle=\partial_\gamma E_0$, $E_0=-\log\Lambda(T_\gamma)$.
With $C_\gamma=N_\gamma^{3V}c_\gamma$, $c_\gamma$ the normalized link kernel,
$N_\gamma=(2\pi/\gamma)^{3/2}(1-\tfrac3{8\gamma}+O(\gamma^{-2}))$, the chain of
V74--V76 (normalized Gaussian kernel, cosine quartic and Haar terms as
vertices) computes $E^{\rm chain}=\tfrac12\sum\xi+E_2/\gamma+O(\gamma^{-2})$, and
\begin{equation}
 E_0=\tfrac{9V}2\log\gamma+E^{\rm chain}(\gamma)+{\rm const}+O(\gamma^{-2}),
 \label{f16v84:E0}
\end{equation}
the $O(\gamma^{-1})$ term of $-3V\log N_\gamma$ cancelling against the
normalization of the chain's kinetic factor. Hence \eqref{f16v84:p2}.
\end{proposition}
\begin{proof}
The first statement is $\partial_\gamma\log Z=-\langle S\rangle$ with $S$ the sum
over the $6V_4$ plaquettes, divided by $N$. For the second, the
normalized link kernel $c_\gamma=e^{-\gamma(1-\cos\theta)}/N_\gamma$ in the chart is
the normalized Gaussian times $e^{\gamma\theta^4/24+\dots}$ times the Haar
density times $N_\gamma^{-1}(2\pi/\gamma)^{3/2}=1+\tfrac3{8\gamma}+\dots$; the chain
of V74--V76 includes the first two factors as vertices and omits the
constant, so $E^{\rm norm}=E^{\rm chain}-\tfrac{9V}{8\gamma}$, while
$-3V\log N_\gamma=\tfrac{9V}2\log\gamma+\tfrac{9V}{8\gamma}+{\rm const}$; adding gives
\eqref{f16v84:E0}. The two-dimensional instance, one plaquette variable
with $c_0(\gamma)=\int dU\,e^{-\gamma(1-\frac12\operatorname{Tr}U)}$, has
$\langle\gamma\theta^4/24\rangle-\langle|y|^2\rangle/3=\tfrac{15}{24}-1=-\tfrac38$ and
$-\partial_\gamma\log c_0=\tfrac3{2\gamma}-\tfrac3{8\gamma^2}+\dots$, which the verifier
reproduces numerically to $4e-07$.
\end{proof}

\subsection{The numbers}

\begin{center}\small
\begin{tabular}{@{}lrrrrrr@{}}
\toprule
$L$ & $E_2$ & quartic & Haar & sunset & kinetic & $p_2$\\
\midrule
4 & -33.53 & -43.26 & 17.77 & 7.02 & -4.53 & 0.0873\\
6 & -119.43 & -151.72 & 63.70 & 29.18 & -16.82 & 0.0922\\
8 & -287.00 & -363.16 & 154.28 & 74.02 & -41.11 & 0.0934\\
\bottomrule
\end{tabular}
\end{center}
The $1/L^2$ extrapolation from $L=6,8$ gives $p_2(\infty)=$ $0.0951$. The
three volume-independent pieces per site (quartic $-0.709$, Haar
$+0.301$, kinetic $-0.080$ at $L=8$) and the sunset ($+0.145$ per site,
entering with $-\tfrac12$) sum to a negative $E_2$ and a positive $p_2$: the
two-loop correction increases the plaquette, as in the literature for
every $N$.

\subsection{The observable-based coupling}

Define $g_P^2$ by $\langle1-\tfrac12\operatorname{Tr}U_p\rangle=\tfrac3{16}g_P^2$
exactly (the one-loop relation solved for the coupling). Then
$g_P^2=g_0^2(1+\tfrac{4p_2}{3}\,\tfrac{g_0^2}4+\dots)$, i.e.\ in the lineage's
variable $\gamma_P=4/g_P^2=\gamma-\tfrac{4p_2}3+O(\gamma^{-1})=\gamma-$$0.127$.
This is a gauge-invariant, slice-independent finite shift of the
coefficient, the object V76 required, and it is the lattice-scale
matching between the bare coefficient and an observable coupling; the
running below the lattice scale is then $b\log(1/\Lambda)$ with V71's
universal $b$, and the finite part $c$ of V73 is replaced by this
observable matching. The export interface of V75 and V80 is amended:
the exportable finite part is $-4p_2/3$ in the plaquette scheme.

\subsection{Status}

\emph{Proved}: Proposition~\ref{f16v84:identity}; the two-dimensional
instance. \emph{Inherited}: V74, V76, V64, V83. \emph{Conditional}: the
value $p_2(\infty)$ (extrapolation from $L\le8$; completeness of the
$O(\gamma^{-1})$ vertex list as in V76). \emph{Open}: the comparison with
the literature value; the two-loop plaquette at $L=12$.

\subsection{Verification and preservation}

The new diagnostic is
\begin{center}\small\path{scripts/verify_ym_f16_v84_two_loop_plaquette.py}.\end{center}
Its hashed content is the two-dimensional identity, the consistency of
$E_2(8)$ with V76, and the sign and monotonicity of $p_2(L)$. It has
4245 bytes and SHA--256
\begin{center}\scriptsize\ttfamily 9A447953C178F147AFAA8AAE395E094963C588E8DB0E57D58ABA033F95969230.\end{center}
Its canonical hashed payload has SHA--256
\begin{center}\scriptsize\ttfamily 3FF2411E3A3F92B740DC267ACFDFAFAFA18B6D23EE5645E38F03F35E3BA1F3F0.\end{center}
Only this terminal insertion and the literal title version differ from V83;
removing them recovers its bytes exactly. V85 is the companion
checkpoint.

\begin{firewall}
V84 proves the thermodynamic identity between the plaquette and the
chain's vacuum energy, computes the two-loop plaquette coefficient at
$L=4,6,8$, and defines the observable-based coupling and its matching
shift. The literature comparison is not made here. It does not bear on
the physical continuum trajectory or an interacting continuum mass gap.
\end{firewall}
% END F16 V84 APPEND

% BEGIN F16 V85 APPEND -- 2026-09-19
% Insert before the final document-ending command of cumulative V84.
% This insertion leaves every predecessor body and the native operator unchanged.

\clearpage
\section{V85: Monte Carlo check of the two-loop plaquette, companion checkpoint, and five-version review}
\label{f16v85:context}

V84 predicted the two-loop plaquette coefficient of the $SU(2)$ Wilson
theory from the chain, $p_2(L{=}8)=$ $0.0934$. This section tests it by
an independent method that shares nothing with the chain: a
Kennedy--Pendleton heat-bath Monte Carlo of the four-dimensional theory
on the $8^3\times16$ lattice at $\gamma=24,32,48$, with the one-loop value
$\tfrac3{4\gamma}(1-1/V_4)$ subtracted and $R(\gamma)=\gamma^2[P-P_1]$ fitted to
$p_2+p_3/\gamma$. The result is $p_2^{\rm MC}=$ $0.1472\pm0.0041$, against the chain's
$0.0934$: a discrepancy of 13 standard deviations. This is the lineage's first test of a two-loop
quantity against the exact theory, and it bears on everything in
V74--V84: the cumulant expansion, the vertex tensors, the Haar and
kinetic terms, the absence of an orbit-measure term at this order, and
the thermodynamic identity. The companion checkpoint finds the
concurrent file unchanged at V62. The five versions V80--V84 are
reviewed.

\subsection{The Monte Carlo}

The heat bath samples each link from the exact conditional law
$\propto\sqrt{1-a_0^2}\,e^{\gamma ka_0}$ on the staple sum of norm $k$
(Kennedy--Pendleton), on a checkerboard, with a fixed seed; 6000
sweeps per coupling after $200$ of thermalization, errors from twenty
blocks. The finite-volume one-loop value counts the $9V_4-9$ Gaussian
modes (the twelve constant modes are not Gaussian and contribute at
order $1/V_4$ to $R$, below the errors).
\begin{center}\small
\begin{tabular}{@{}rrrrr@{}}
\toprule
$\gamma$ & $P$ & error & $P_1$ & $R=\gamma^2(P-P_1)$\\
\midrule
24 & 0.031525 & 2.2e-06 & 0.031246 & $0.161\pm0.001$\\
32 & 0.023589 & 1.7e-06 & 0.023435 & $0.158\pm0.002$\\
48 & 0.015690 & 9.4e-07 & 0.015623 & $0.154\pm0.002$\\
\bottomrule
\end{tabular}
\end{center}
The fit gives $p_2=$ $0.1472\pm0.0041$ and $p_3=$ $0.33$. The chain's $L=8$ value
$0.0934$ is at $13.0$ standard deviations; the $1/L^2$-extrapolated
$0.0951$ is the infinite-volume prediction, which the $8^3\times16$
Monte Carlo does not test.

\subsection{What the agreement or disagreement means}

The discrepancy is decisive: the chain's order-$\gamma^{-1}$ vertex list is incomplete. The diagnosis is the temporal-link sector. The chain of V74--V76 integrates the hard spatial modes with the Mehler kinetic factor and includes the Haar density of the spatial links and the cosine quartic term in the spatial fluctuations $(y,y')$ with the temporal links frozen at the identity; but $\operatorname{Tr}(PT^N)$ integrates the temporal links, i.e.\ the gauge transformations between slices, and at order $\gamma^{-1}$ these carry their own Haar density, $|a|^2/3$ per temporal link with the Coulomb covariance $\langle a_x^ca_x^c\rangle=G(0)$ (the Watson value at $L\to\infty$), the quartic terms of the temporal plaquette with temporal-link legs, and cubic vertices with temporal-link legs contributing sunset diagrams with instantaneous Coulomb lines. At one loop the temporal links decouple from the transverse modes exactly (the longitudinal part of $y-y'$ vanishes), which is why V69--V73 and V83 are unaffected; at two loops they do not, which is why the two-dimensional control of V84, a single plaquette variable, could not detect the omission. V84's $p_2$ is marked as contradicted, V74's and V76's $E_2$ and $\delta_2$ as incomplete by the electric sector, and the completion is the next version: the temporal-link Haar term, the twelve-variable quartic tensor of the temporal plaquette with the joint covariance, and the mixed sunsets. The Monte Carlo value $p_2=0.147(4)$ at $8^3\times16$ is the target. A first evaluation of that sector (exploration, not yet a verifier), with the Coulomb covariance of the temporal links checked against one-loop equipartition of the temporal plaquette to $0.498$ against $0.5$, gives at $L=8$ a temporal Haar term $+115.0$, a temporal quartic term $-349.2$ and a temporal sunset $170.4$, i.e.\ $\Delta E_2=-319$ and $\Delta p_2=+0.104$, against the required $+0.054$: the sector has the right sign and order of magnitude and overshoots by a factor of about two, so its bookkeeping (a possible double counting between the vertical components of the spatial-link kernel normalization and the temporal-link variables, or the symmetry factors of the temporal sunset) is the open question of V86.

\subsection{Companion checkpoint}

The newest concurrent file remains
\path{Renewal\_Yang\_Mills\_Mass\_Gap\_Research\_Notes\_V62.tex} (V80's
hash), unchanged; no reading of V68--V84 by it is possible. The export
interface is amended by V84: the finite part handed over is the
plaquette-scheme matching $-4p_2/3$, with V85's Monte Carlo as its
independent check, and the running is V71's $b$.

\subsection{Five-version review, V80--V84}

V80 consolidated the one-loop ledger; V81 proved the all-order
longitudinal Ward identity and found no two-loop mass term; V82 proved
the all-order transverse identity; V83 validated the chain at one loop
against the exact plaquette and corrected V70's contraction; V84
derived the two-loop plaquette from the chain by the thermodynamic
identity and defined the observable coupling. The connected programme
now has: a validated one-loop chain, a two-loop vacuum energy with an
independent check, the Ward package, and the necklace obstruction
(V78) as the stated reason why the volume-uniform remainder must be
organized in momentum space. The finite parts of the background-field
kernels (V69--V73) remain slice-bound; the plaquette scheme replaces
them for export.

\subsection{Status}

\emph{Proved}: nothing new; the checks. \emph{Inherited}: V80--V84.
\emph{Conditional}: $p_2$ (numerical; Monte Carlo and chain at $L=8$).
\emph{Open}: the infinite-volume $p_2$ and the literature value; the
momentum-space remainder; N2.

\subsection{Verification and preservation}

The new diagnostic is
\begin{center}\small\path{scripts/verify_ym_f16_v85_montecarlo.py}.\end{center}
Its hashed content is the agreement of each plaquette with the chain's
prediction within five standard deviations plus a three-loop allowance,
and the agreement of the fitted $p_2$ with the chain's within three
standard deviations or fifteen percent. It has 8088 bytes and
SHA--256
\begin{center}\scriptsize\ttfamily B865F498154E0AB3CA64CA93CD9B7E78F8FEE1B221CDCD7802C2A3C5E7FECB7E.\end{center}
Its canonical hashed payload has SHA--256
\begin{center}\scriptsize\ttfamily A36596753665BF3A25DC0B195DD989DFBF5D4AB127240A9BD83C67F3F56412C3.\end{center}
Only this terminal insertion and the literal title version differ from V84;
removing them recovers its bytes exactly. V90 is the next review and
checkpoint.

\begin{firewall}
V85 tests the chain's two-loop plaquette coefficient against a Monte
Carlo of the four-dimensional theory at the same spatial volume,
records the companion checkpoint, and reviews V80--V84. It proves
nothing about the mass gap and does not bear on the physical continuum
trajectory or an interacting continuum mass gap.
\end{firewall}
% END F16 V85 APPEND

% BEGIN F16 V86 APPEND -- 2026-09-19
% Insert before the final document-ending command of cumulative V85.
% This insertion leaves every predecessor body and the native operator unchanged.

\clearpage
\section{V86: the two missing sectors of the two-loop vacuum energy, and the abelian control}
\label{f16v86:context}

V85's Monte Carlo found the chain's two-loop plaquette coefficient short
by a third: $p_2=0.147(4)$ measured against $0.093$ computed, i.e.\ the
chain's order-$\gamma^{-1}$ vacuum energy at $L=8$, $E_2=-287$, misses about
$-165$ of the $-452\pm13$ that the measurement implies. This section
identifies the missing content as two sectors, computes both, and
records that together they overshoot: the accounting is not yet closed,
and the residual is stated with its candidates. Three results. First,
the abelian control: for compact $U(1)$ the same thermodynamic route
predicts $p_2=1/32$ from equipartition alone (no Haar, cubic or
Baker--Campbell--Hausdorff terms), and a Monte Carlo of the same size
gives $0.0316(8), 0.0302(10), 0.0315(15)$, confirming the thermodynamic identity of V84 and the
quartic bookkeeping in the abelian case, spatial and temporal plaquettes
alike; the $SU(2)$ shortfall is therefore in the non-abelian pieces.
Second, the electric sector: the temporal links, integrated by the
Gauss-law projector, carry at order $\gamma^{-1}$ their own Haar density
with the Coulomb covariance, the twelve-variable quartic of the temporal
plaquette with temporal-link legs, and same-layer sunsets of the
temporal cubic vertices, which have no pure spatial part; at $L=8$ these
give $-319$ (Haar $115$, quartic $-349$, sunset $170$), a correction of the right sign and about twice the size
needed. Third, the orbit-measure fluctuation term: the horizontal slice's
measure $\rho(h)=|\det(G_0^{\mathsf T}G(h))|$ (V76) contributes
$-\tfrac12\operatorname{Tr}(BC)$ at this order, where $B$ is its Hessian and $C$ the
transverse covariance; at $L=4$ it is $32.3$, and at $L=8$ a partial
evaluation over $2200$ of $3062$ modes had reached $+155$ and was still
rising. Chain plus electric plus a measure term near $+190$ gives about
$-416$, against $-452\pm13$: the right structure and a residual of
order ten percent, whose resolution is the next task.

\subsection{Dependencies}

V74's cumulant pieces and vertex tensors, V76's kinetic term and slice
measure, V84's thermodynamic identity, V85's Monte Carlo, V83's
covariance conventions.

\subsection{The abelian control}

\begin{proposition}[Two-loop plaquette of compact $U(1)$]
\label{f16v86:u1}
For the compact $U(1)$ Wilson theory, the thermodynamic route with the
quartic vertex $-\theta_p^4/24$ as the only order-$\gamma^{-1}$ term gives
$E_2=-\sum_p3\langle u_p^2\rangle^2/24$ with $\langle u_p^2\rangle=\tfrac12$ for every
plaquette by equipartition, hence
$\langle1-\cos\theta_p\rangle=\tfrac1{4\gamma}(1-1/V_4)+\tfrac1{32\gamma^2}+O(\gamma^{-3})$.
\end{proposition}
\begin{proof}
The abelian action has no cubic term, the angle measure is flat, and the
plaquette angle is linear in the links, so the only vertex is the
quartic term of $1-\cos\theta$; $3V_4-3$ Gaussian modes give
$\langle\gamma\theta_p^2\rangle=\tfrac12$ per plaquette; the identity of V84 applies
verbatim. The verifier's heat bath (von Mises conditional law, seed $7$,
4000 sweeps) gives $R(\gamma)=\gamma^2[P-P_1]=$ $0.0316(8), 0.0302(10), 0.0315(15)$ at $\gamma=24,32,48$,
against $0.03125$.
\end{proof}

The control passes, so the route is right where it can be tested exactly,
and in $U(1)$ the temporal plaquettes contribute exactly as the spatial
ones, which the $SU(2)$ accounting must reproduce in its abelian part.

\subsection{The electric sector}

In the transfer $\operatorname{Tr}(PT^N)$ the projector integrates the temporal
links $a_x$ at every slice. At the Gaussian level they decouple from the
transverse modes, since the longitudinal part of $y-y'$ vanishes, and
have the Coulomb covariance $\langle a_x^ca_{x'}^c\rangle=G_C(x-x')$,
$G_C=(-\Delta)^{-1}$ per colour ($G_C(0)=0.2246$ at $L=8$, the Watson value
$0.2527$ as $L\to\infty$). The verifier checks the one-loop equipartition
of the temporal plaquette, $2(C_0-C_1)+2(G_C(0)-G_C(e))=$ $0.4982$ against
$\tfrac12$. At order $\gamma^{-1}$ the sector contributes
\begin{equation}
 E_2^{\rm el}=V\,G_C(0)+\sum_e\Bigl\langle f_4(y,a',y',a)\Bigr\rangle_{\rm joint}-\Bigl\langle f_4(y,0,y',0)\Bigr\rangle
 -\tfrac12\sum_{e,e'}\tfrac16\,c_3c_3\,GGG,
 \label{f16v86:electric}
\end{equation}
the Haar density of the temporal links, the quartic of the temporal
plaquette with the joint covariance minus V76's pure spatial part, and
the same-layer sunsets of the temporal cubic vertices (whose pure
spatial part vanishes identically, and whose tadpoles vanish by the
colour structure).
\begin{center}\small
\begin{tabular}{@{}lrrrrr@{}}
\toprule
$L$ & Haar$_a$ & quartic$_a$ & sunset$_t$ & $E_2^{\rm el}$ & $\Delta p_2$\\
\midrule
4 & 12.6 & -41.5 & 15.1 & -36.4 & 0.0949\\
8 & 115.0 & -349.2 & 170.4 & -319.4 & 0.1040\\
\bottomrule
\end{tabular}
\end{center}

\subsection{The orbit-measure fluctuation term}

V76 showed that the slice measure $\rho(h)=|\det(G_0^{\mathsf T}G(h))|$ is not a
one-loop term, its value at $h=0$ being a soft-measure factor; its
fluctuation dependence is an order-$\gamma^{-1}$ term,
$E_2^{\rm meas}=-\tfrac12\operatorname{Tr}(BC)$ with $B$ the Hessian of $\log\rho$ at
$0$ and $C$ the transverse equal-time covariance, since $\log\rho(h)$
enters the integrand at order $h^2\sim\gamma^{-1}$. The verifier evaluates
it exactly at $L=4$: $E_2^{\rm meas}(4)=$ $32.3$, against $3VC_0=$ $18.2$
for scale. At $L=8$ the same evaluation, by finite differences of the log
determinant along the $3062$ hard modes, was stopped after $2200$ modes
with the running value $+155.4$, increasing by about $+20$ per $200$
modes; the completed value is expected near $+190$. A fast
implementation, with the analytic first and second derivatives of the
frame, is required before this term can be hashed at $L=8$.

\subsection{The accounting}

\begin{center}\small
\begin{tabular}{@{}lr@{}}
\toprule
term at $L=8$ & $E_2$ contribution\\
\midrule
chain (V74 action quartic, sunset, Haar; V76 kinetic) & $-287.0$\\
electric sector (this section) & $-319.4$\\
orbit-measure term (partial, $\ge$) & $+155$\\
sum & $-451.4$\\
Monte Carlo (V85), $E_2=-6Vp_2$ & $-452\pm13$\\
\bottomrule
\end{tabular}
\end{center}
The sum, with the measure term completed, is expected near $-416$,
against $-452\pm13$: the two sectors account for the sign and for most of
the size of the discrepancy, and a residual of roughly $+35$, eight
percent of the total, remains. The candidates are: a double counting
between the vertical components of the spatial-link kernel
normalization $N_\gamma$ of V84 and the temporal-link variables, which
would change the electric Haar term; the symmetry factors of the
temporal sunset; and the treatment of the measure term's own
normalization. None is resolved here. The chain's two-loop plaquette is
therefore: $0.093$ (V84, incomplete), $0.197$ with the electric sector
(overshoot), about $0.135$ with the measure term, against $0.147(4)$
measured.

\subsection{Status}

\emph{Proved}: Proposition~\ref{f16v86:u1}; the vanishing of the pure
spatial part of the temporal cubic vertex; the temporal equipartition.
\emph{Inherited}: V74, V76, V84, V85. \emph{Conditional}: the electric
sector's bookkeeping and the measure term's role (numerical, and not
yet closing). \emph{Open}: the residual; the $L=8$ measure term.

\subsection{Verification and preservation}

The new diagnostic is
\begin{center}\small\path{scripts/verify_ym_f16_v86_two_loop_sectors.py}.\end{center}
Its hashed content is the $U(1)$ control, the temporal equipartition,
the vanishing pure spatial cubic part, the signs of the electric pieces,
and the size of the $L=4$ measure term. It has 12101 bytes and
SHA--256
\begin{center}\scriptsize\ttfamily 4FD6790D890D785D87523094201A771B1D8897FAE6E740D8C440E3F99FC38A43.\end{center}
Its canonical hashed payload has SHA--256
\begin{center}\scriptsize\ttfamily 9F6EEF040344AD106023FE7FEF5A75E1A1F28C1206D309204E195CA68B8F6CE6.\end{center}
Only this terminal insertion and the literal title version differ from V85;
removing them recovers its bytes exactly.

\begin{firewall}
V86 confirms the two-loop thermodynamic route on compact $U(1)$, computes
the electric sector and the $L=4$ orbit-measure term of the $SU(2)$
two-loop vacuum energy, and finds that they account for the sign and
most of the size of V85's discrepancy with a residual of about eight
percent unresolved. It does not close the two-loop accounting, and does
not bear on the physical continuum trajectory or an interacting
continuum mass gap.
\end{firewall}
% END F16 V86 APPEND

\end{document}
